\input{preamble}

% Bon, commençons ici.
%
\begin{document}

\title{Compléments d'algèbre}


\maketitle

\phantomsection
\label{section-phantom}

\tableofcontents



\section{Introduction}
\label{section-introduction}

\noindent
Dans ce chapitre, nous démontrons quelques résultats d'algèbre commutative
qui sont moins élémentaires que ceux du premier chapitre consacré à
l'algèbre commutative, voir
Algèbre, section \ref{algebra-section-introduction}.
On pourra consulter \cite{MatCA}.






\section{Conseils au lecteur}
\label{section-advice}

\noindent
Plus encore que dans le chapitre sur l'algèbre commutative, chacune des
sections de ce chapitre peut être lue indépendamment. À partir de la
section \ref{section-derived-modules},
nous utilisons librement la catégorie dérivée (non bornée) des modules sur
les anneaux, ainsi que tout l'appareil qui l'accompagne.




\section{Modules stablement libres}
\label{section-stably-free}

\noindent
Voici ce qui semble être la définition généralement admise.

\begin{definition}
\label{definition-stably-free}
Soit $R$ un anneau.
\begin{enumerate}
\item Deux modules $M$, $N$ sur $R$ sont dits
{\it stablement isomorphes} s'il existe $n, m \geq 0$ tels que
$M \oplus R^{\oplus m} \cong N \oplus R^{\oplus n}$
comme $R$-modules.
\item Un module $M$ est {\it stablement libre} s'il est stablement isomorphe
à un module libre.
\end{enumerate}
\end{definition}

\noindent
Remarquons qu'un module stablement libre est projectif.

\begin{lemma}
\label{lemma-exact-category-stably-free}
Soit $R$ un anneau. Soit $0 \to P' \to P \to P'' \to 0$ une suite
exacte courte de $R$-modules projectifs de type fini. Si $2$ de ces $3$
modules sont stablement libres, alors le troisième l'est aussi.
\end{lemma}

\begin{proof}
Puisque les modules sont projectifs, la suite est scindée. Nous pouvons donc
choisir un isomorphisme $P = P' \oplus P''$. Si $P' \oplus R^{\oplus n}$
et $P'' \oplus R^{\oplus m}$ sont libres, nous voyons alors que
$P \oplus R^{\oplus n + m}$ est libre. Supposons que $P'$ et $P$ soient
stablement libres, disons que $P \oplus R^{\oplus n}$ est libre et que
$P' \oplus R^{\oplus m}$ est libre. Alors
$$
P'' \oplus (P' \oplus R^{\oplus m}) \oplus R^{\oplus n} =
(P'' \oplus P') \oplus R^{\oplus m} \oplus R^{\oplus n} =
(P \oplus R^{\oplus n}) \oplus R^{\oplus m}
$$
est libre. Ainsi, $P''$ est stablement libre. Par symétrie, nous obtenons le
dernier des trois cas.
\end{proof}

\begin{lemma}
\label{lemma-lift-stably-free}
Soit $R$ un anneau. Soit $I \subset R$ un idéal. Supposons que
tout élément de $1 + I$ soit inversible (autrement dit, $I$ est contenu
dans le radical de Jacobson de $R$). Pour tout $R/I$-module $E$ de type fini
et stablement libre, il existe un $R$-module $M$ de type fini et stablement
libre tel que $M/IM \cong E$.
\end{lemma}

\begin{proof}
Choisissons $n$ et $m$, ainsi qu'un isomorphisme
$E \oplus (R/I)^{\oplus n} \cong (R/I)^{\oplus m}$.
Choisissons des applications $R$-linéaires
$\varphi : R^{\oplus m} \to R^{\oplus n}$ et
$\psi : R^{\oplus n} \to R^{\oplus m}$ relevant la
projection $(R/I)^{\oplus m} \to (R/I)^{\oplus n}$
et l'injection $(R/I)^{\oplus n} \to (R/I)^{\oplus m}$.
Alors $\varphi \circ \psi : R^{\oplus n} \to R^{\oplus n}$
se réduit à l'identité modulo $I$. Par notre hypothèse sur $I$, le déterminant
de cette application est donc inversible. Par conséquent,
$P = \Ker(\varphi)$ est stablement libre et relève $E$.
\end{proof}

\begin{lemma}
\label{lemma-lift-projective}
Soit $R$ un anneau. Soit $I \subset R$ un idéal.
Supposons que tout élément de $1 + I$ soit inversible
(autrement dit, $I$ est contenu dans le radical de Jacobson de $R$).
Soit $M$ un $R$-module plat de type fini tel que
$M/IM$ soit un $R/I$-module projectif.
Alors $M$ est un $R$-module projectif de type fini.
\end{lemma}

\begin{proof}
Par Algèbre, lemme \ref{algebra-lemma-finite-flat-local},
nous voyons que $M_\mathfrak p$ est libre de type fini pour tout idéal premier
$\mathfrak p \subset R$.
D'après
Algèbre, lemme \ref{algebra-lemma-finite-projective},
il suffit de montrer que la fonction $\rho_M : \mathfrak p \mapsto
\dim_{\kappa(\mathfrak p)} M \otimes_R \kappa(\mathfrak p)$
est localement constante sur $\Spec(R)$. Comme $M/IM$ est projectif de type
fini, c'est vrai sur $V(I) \subset \Spec(R)$. Puisque tout point fermé
de $\Spec(R)$ appartient à $V(I)$ et que
$\rho_M(\mathfrak p) = \rho_M(\mathfrak q)$
lorsque $\mathfrak p \subset \mathfrak q \subset R$
sont des idéaux premiers, nous concluons par un argument élémentaire
sur les espaces topologiques, que nous omettons.
\end{proof}

\noindent
Le relèvement du lemme \ref{lemma-lift-stably-free}
est unique à isomorphisme près d'après le lemme suivant.

\begin{lemma}
\label{lemma-isomorphic-finite-projective-lifts}
Soit $R$ un anneau. Soit $I \subset R$ un idéal. Supposons que
tout élément de $1 + I$ soit inversible (autrement dit, $I$ est contenu
dans le radical de Jacobson de $R$). Si $P$ et $P'$ sont des $R$-modules
projectifs de type fini, alors
\begin{enumerate}
\item si $\varphi : P \to P'$ est un homomorphisme de $R$-modules induisant
un isomorphisme $\overline{\varphi} : P/IP \to P'/IP'$, alors $\varphi$
est un isomorphisme ;
\item si $P/IP \cong P'/IP'$, alors $P \cong P'$.
\end{enumerate}
\end{lemma}

\begin{proof}
Preuve de (1). Comme $P'$ est projectif en tant que $R$-module, nous pouvons
choisir un relèvement $\psi : P' \to P$ de l'application
$P' \to P'/IP' \xrightarrow{\overline{\varphi}^{-1}} P/IP$.
D'après le lemme de Nakayama (Algèbre, lemme \ref{algebra-lemma-NAK}),
$\psi \circ \varphi$ et $\varphi \circ \psi$ sont surjectifs.
Ces applications sont donc des isomorphismes (Algèbre, lemme
\ref{algebra-lemma-fun}). Ainsi, $\varphi$ est un isomorphisme.

\medskip\noindent
Preuve de (2). Choisissons un isomorphisme $P/IP \cong P'/IP'$.
Puisque $P$ est projectif, nous pouvons choisir un relèvement
$\varphi : P \to P'$ de l'application $P \to P/IP \to P'/IP'$.
Alors $\varphi$ est un isomorphisme d'après (1).
\end{proof}




\section{Un commentaire sur la propriété d'Artin-Rees}
\label{section-artin-rees}

\noindent
Une partie de ce matériel est tirée de \cite{conrad-dejong}. On trouvera
dans \cite[Section 1]{Eis} une discussion générale accompagnée de références
supplémentaires.

\medskip\noindent
Soit $A$ un anneau noethérien et soit $I \subset A$ un idéal. Étant donné un
homomorphisme $f : M \to N$ de $A$-modules de type fini, il existe $c \geq 0$
tel que
$$
f(M) \cap I^nN \subset f(I^{n - c}M)
$$
pour tout $n \geq c$, voir Algèbre, lemme \ref{algebra-lemma-map-AR}. Dans
cette situation, nous dirons que {\it $c$ convient pour $f$ dans le lemme
d'Artin-Rees}.

\begin{lemma}
\label{lemma-approximate-complex}
Soit $A$ un anneau noethérien. Soit $I \subset A$ un idéal contenu dans
le radical de Jacobson de $A$. Soient
$$
S : L \xrightarrow{f} M \xrightarrow{g} N
\quad\text{et}\quad
S' : L \xrightarrow{f'} M \xrightarrow{g'} N
$$
les deux complexes de $A$-modules de type fini représentés ci-dessus.
Supposons que
\begin{enumerate}
\item $c$ convienne dans le lemme d'Artin-Rees pour $f$ et $g$ ;
\item le complexe $S$ soit exact ; et
\item $f' = f \bmod I^{c + 1}M$ et $g' = g \bmod I^{c + 1}N$.
\end{enumerate}
Alors $c$ convient dans le lemme d'Artin-Rees pour $g'$ et le
complexe $S'$ est exact.
\end{lemma}

\begin{proof}
Montrons d'abord que $g'(M) \cap I^nN \subset g'(I^{n - c}M)$ pour
$n \geq c$. Soit $a$ un élément de $M$ tel que $g'(a) \in I^nN$.
Nous voulons modifier $a$ par un élément de $f'(L)$, c'est-à-dire sans
changer $g'(a)$, afin que $a \in I^{n-c}M$. Supposons que $a \in I^rM$,
où $r < n - c$. Alors
$$
g(a) = g'(a) + (g - g')(a) \in
I^n N + I^{r + c + 1}N = I^{r + c + 1}N.
$$
Par Artin-Rees pour $g$, nous avons $g(a) \in g(I^{r + 1}M)$. Écrivons
$g(a) = g(a_1)$ avec $a_1 \in I^{r + 1}M$. Comme la suite $S$ est exacte,
$a - a_1 \in f(L)$. Nous écrivons donc $a = f(b) + a_1$ pour un certain
$b \in L$. Alors $f(b) = a - a_1 \in I^rM$. Artin-Rees pour $f$ montre que,
si $r \geq c$, nous pouvons remplacer $b$ par un élément de $I^{r - c}L$.
Dans tous les cas, $a = f'(b) + a_2$, où
$a_2 = (f - f')(b) + a_1 \in I^{r + 1}M$. (En effet, soit $c \geq r$
et $(f - f')(b) \in I^{r + 1}M$ par hypothèse, soit $c < r$ et
$b \in I^{r - c}$, d'où à nouveau $(f - f')(b) \in I^{c + 1} I^{r - c} M =
I^{r + 1}M$.) Nous pouvons donc modifier $a$ par l'élément
$f'(b) \in f'(L)$ de façon à augmenter $r$ de $1$.

\medskip\noindent
En fait, l'argument ci-dessus montre que
$(g')^{-1}(I^nN) \subset f'(L) + I^{n - c}M$ pour tout $n \geq c$.
Par conséquent, $S'$ est exact car
$$
(g')^{-1}(0) = (g')^{-1}(\bigcap I^nN) \subset
\bigcap f'(L) + I^{n - c}M = f'(L)
$$
puisque $I$ est contenu dans le radical de Jacobson de $A$, voir Algèbre,
lemme \ref{algebra-lemma-intersection-powers-ideal-module}.
\end{proof}

\noindent
Étant donnés un idéal $I \subset A$ d'un anneau $A$ et un $A$-module $M$,
nous posons
$$
\text{Gr}_I(M) = \bigoplus I^nM/I^{n + 1}M.
$$
Nous le considérons comme un $\text{Gr}_I(A)$-module gradué.

\begin{lemma}
\label{lemma-approximate-complex-graded}
Faisons les hypothèses du lemme \ref{lemma-approximate-complex}.
Soient $Q = \Coker(g)$ et $Q' = \Coker(g')$. Alors
$\text{Gr}_I(Q) \cong \text{Gr}_I(Q')$
comme $\text{Gr}_I(A)$-modules gradués.
\end{lemma}

\begin{proof}
En degré $n$, nous avons
$\text{Gr}_I(Q)_n = I^nN/(I^{n + 1}N + g(M) \cap I^nN)$
et de même pour $Q'$. Nous affirmons que
$$
g(M) \cap I^nN \subset I^{n + 1}N + g'(M) \cap I^nN.
$$
Par symétrie (la preuve de l'affirmation utilisera seulement que $c$ convient
pour $g$, ce qui vaut aussi pour $g'$ d'après le lemme), cela impliquera que
$$
I^{n + 1}N + g(M) \cap I^nN = I^{n + 1}N + g'(M) \cap I^nN
$$
d'où $\text{Gr}_I(Q)_n$ et $\text{Gr}_I(Q')_n$ coïncident comme sous-quotients
de $N$, ce qui démontre le lemme. Remarquons que l'affirmation est claire
pour $n \leq c$ puisque $g = g' \bmod I^{c + 1}N$. Si $n > c$, supposons
$b \in g(M) \cap I^nN$. Écrivons $b = g(a)$ avec $a \in I^{n - c}M$.
Posons $b' = g'(a)$. Nous avons $b - b' = (g - g')(a) \in I^{n + 1}N$,
comme voulu.
\end{proof}

\begin{lemma}
\label{lemma-works-flat-extension}
Soit $A \to B$ un homomorphisme plat d'anneaux noethériens. Soit
$I \subset A$ un idéal. Soit $f : M \to N$ un homomorphisme de $A$-modules
de type fini. Supposons que $c$ convienne pour $f$ dans le lemme
d'Artin-Rees. Alors $c$ convient pour
$f \otimes 1 : M \otimes_A B \to N \otimes_A B$ dans le lemme d'Artin-Rees
relatif à l'idéal $IB$.
\end{lemma}

\begin{proof}
Notons que
$$
(f \otimes 1)(M) \cap I^n N \otimes_A B
= (f \otimes 1)\left((f \otimes 1)^{-1}(I^n N \otimes_A B)\right)
$$
D'autre part,
\begin{align*}
(f \otimes 1)^{-1}(I^n N \otimes_A B) &
= \Ker(M \otimes_A B \to N \otimes_A B/(I^n N \otimes_A B)) \\
& =
\Ker(M \otimes_A B \to (N/I^nN) \otimes_A B)
\end{align*}
Comme $A \to B$ est plat, la formation des noyaux et des conoyaux commute
au produit tensoriel par $B$, si bien que ceci est égal à
$f^{-1}(I^nN) \otimes_A B$. Par hypothèse, $f^{-1}(I^nN)$ est contenu dans
$\Ker(f) + I^{n - c}M$. Le lemme en résulte.
\end{proof}








\section{Produits fibrés d'anneaux, I}
\label{section-fibre-products-rings}

\noindent
Les produits fibrés d'anneaux sont liés aux sommes amalgamées de schémas.
Certains cas de sommes amalgamées de schémas sont étudiés dans
Compléments sur les morphismes, section \ref{more-morphisms-section-pushouts}.

\begin{lemma}
\label{lemma-fibre-product-finite-type}
Soit $R$ un anneau. Soient $A \to B$ et $C \to B$ des homomorphismes de
$R$-algèbres. Supposons que
\begin{enumerate}
\item $R$ soit noethérien ;
\item $A$, $B$, $C$ soient de type fini sur $R$ ;
\item $A \to B$ soit surjectif ; et
\item $B$ soit fini sur $C$.
\end{enumerate}
Alors $A \times_B C$ est de type fini sur $R$.
\end{lemma}

\begin{proof}
Posons $D = A \times_B C$. Il existe un diagramme commutatif
$$
\xymatrix{
0 &
B \ar[l] &
A \ar[l] &
I \ar[l] &
0 \ar[l] \\
0 &
C \ar[l] \ar[u] &
D \ar[l] \ar[u] &
I \ar[l] \ar[u] &
0 \ar[l]
}
$$
à lignes exactes. Choisissons $y_1, \ldots, y_n \in B$ qui engendrent $B$
comme $C$-module. Choisissons $x_i \in A$ d'image $y_i$.
Alors $1, x_1, \ldots, x_n$ engendrent $A$ comme $D$-module.
L'application $D \to A \times C$ est injective et l'anneau $A \times C$ est
fini comme $D$-module (car c'est la somme directe des $D$-modules finis
$A$ et $C$). Le lemme résulte donc du lemme d'Artin-Tate
(Algèbre, lemme \ref{algebra-lemma-Artin-Tate}).
\end{proof}

\begin{lemma}
\label{lemma-formal-consequence}
Soit $R$ un anneau noethérien. Soit $I$ un ensemble fini. Considérons un
diagramme cartésien
$$
\xymatrix{
\prod B_i &
\prod A_i \ar[l]^{\prod \varphi_i} \\
Q \ar[u]^{\prod \psi_i} &
P \ar[u] \ar[l]
}
$$
où les $\psi_i$ et les $\varphi_i$ sont surjectifs, et où $Q$, les $A_i$ et
les $B_i$ sont de type fini sur $R$. Alors $P$ est de type fini sur $R$.
\end{lemma}

\begin{proof}
Cela résulte du lemme \ref{lemma-fibre-product-finite-type}
et d'une récurrence sur le cardinal de $I$.
En effet, écrivons $I = I' \amalg \{i_0\}$. Soit $P'$ l'anneau défini
par le diagramme du lemme pour $I'$. Alors $P'$ est de type fini par
l'hypothèse de récurrence. Enfin, $P$ s'insère dans un diagramme de produit
fibré
$$
\xymatrix{
B_{i_0} &
A_{i_0} \ar[l] \\
P' \ar[u] &
P \ar[u] \ar[l] &
}
$$
auquel le lemme s'applique.
\end{proof}

\begin{lemma}
\label{lemma-diagram-localize}
Considérons un diagramme cartésien d'anneaux
$$
\xymatrix{
R &
R' \ar[l]^t \\
B \ar[u]_s &
B'\ar[u] \ar[l]
}
$$
c'est-à-dire $B' = B \times_R R'$. Si $h \in B'$ correspond à $g \in B$
et à $f \in R'$ tels que $s(g) = t(f)$, alors le diagramme
$$
\xymatrix{
R_{s(g)} = R_{t(f)} &
(R')_f \ar[l]^-t \\
B_g \ar[u]_s &
(B')_h \ar[u] \ar[l]
}
$$
est lui aussi cartésien.
\end{lemma}

\begin{proof}
L'égalité $B' = B \times_R R'$ nous dit que
$$
0 \to B' \to B \oplus R' \xrightarrow{s, -t} R
$$
est une suite exacte de $B'$-modules. Nous avons $B_g = B_h$,
$R'_f = R'_h$ et $R_{s(g)} = R_{t(f)} = R_h$ comme $B'$-modules.
Par exactitude de la localisation
(Algèbre, proposition \ref{algebra-proposition-localization-exact}),
nous obtenons que
$$
0 \to B'_h \to B_g \oplus R'_f \xrightarrow{s, -t} R_{s(g)} = R_{t(f)}
$$
est une suite exacte. Cela démontre le lemme.
\end{proof}

\noindent
Considérons un diagramme commutatif d'anneaux
$$
\xymatrix{
R &
R' \ar[l] \\
B \ar[u] &
B' \ar[u] \ar[l]
}
$$
Considérons le foncteur (où le produit fibré de catégories est
celui construit dans
Catégories, exemple \ref{categories-example-2-fibre-product-categories})
\begin{equation}
\label{equation-modules}
\text{Mod}_{B'} \longrightarrow
\text{Mod}_B \times_{\text{Mod}_R} \text{Mod}_{R'},\quad
L' \longmapsto (L' \otimes_{B'} B, L' \otimes_{B'} R', can)
\end{equation}
où $can$ est l'identification canonique
$L' \otimes_{B'} B \otimes_B R = L' \otimes_{B'} R' \otimes_{R'} R$.
Dans la suite, nous noterons $(N, M', \varphi)$ un objet du membre de droite,
c'est-à-dire que $N$ est un $B$-module, $M'$ un $R'$-module et que
$\varphi : N \otimes_B R \to M' \otimes_{R'} R$ est un isomorphisme.

\begin{lemma}
\label{lemma-modules}
Étant donné un diagramme commutatif d'anneaux
$$
\xymatrix{
R &
R' \ar[l] \\
B \ar[u] &
B' \ar[u] \ar[l]
}
$$
le foncteur (\ref{equation-modules}) admet un adjoint à droite, à savoir le
foncteur
$$
F : (N, M', \varphi) \longmapsto N \times_\varphi M'
$$
(voir la preuve pour plus de précisions).
\end{lemma}

\begin{proof}
Étant donné un objet $(N, M', \varphi)$ de la catégorie
$\text{Mod}_B \times_{\text{Mod}_R} \text{Mod}_{R'}$
nous posons
$$
N \times_\varphi M' = \{(n, m') \in N \times M' \mid
\varphi(n \otimes 1) = m' \otimes 1\text{ dans }M' \otimes_{R'} R\}
$$
considéré comme un $B'$-module.
L'assertion d'adjonction signifie que, pour un $B'$-module $L'$ et
un triplet $(N, M', \varphi)$, nous avons
$$
\Hom_{B'}(L', N \times_\varphi M') =
\Hom_B(L' \otimes_{B'} B, N)
\times_{\Hom_R(L' \otimes_{B'} R, M' \otimes_{R'} R)}
\Hom_{R'}(L' \otimes_{B'} R', M')
$$
D'après Algèbre, lemme \ref{algebra-lemma-adjoint-tensor-restrict},
le membre de droite est égal à
$$
\Hom_{B'}(L', N) \times_{\Hom_{B'}(L', M' \otimes_{R'} R)} \Hom_{B'}(L', M')
$$
Il est donc clair qu'un couple $(g, f')$ d'éléments de ce produit fibré
nous donne une application $B'$-linéaire
$L' \to N \times_\varphi M'$, $l' \mapsto (g(l'), f'(l'))$.
Réciproquement, étant donnée une application $B'$-linéaire
$g' : L' \to N \times_\varphi M'$, nous pouvons prendre pour $g$ la composée
$L' \to N \times_\varphi M' \to N$ et pour $f'$ la composée
$L' \to N \times_\varphi M' \to M'$. Ces constructions sont inverses l'une
de l'autre et définissent l'isomorphisme voulu.
\end{proof}






\section{Produits fibrés d'anneaux, II}
\label{section-fibre-products-rings-II}

\noindent
Dans cette section, nous étudions les produits fibrés dans la situation
suivante.

\begin{situation}
\label{situation-module-over-fibre-product}
Dans la suite, nous considérerons des homomorphismes d'anneaux
$$
\xymatrix{
B \ar[r] & A & A' \ar[l]
}
$$
où nous supposons que $A' \to A$ est surjectif de noyau $I$.
Dans cette situation, nous posons $B' = B \times_A A'$ afin
d'obtenir un carré cartésien
$$
\xymatrix{
A & A' \ar[l] \\
B \ar[u] & B' \ar[l] \ar[u]
}
$$
\end{situation}

\begin{lemma}
\label{lemma-points-of-fibre-product}
Dans la situation \ref{situation-module-over-fibre-product},
nous avons
$$
\Spec(B') = \Spec(B) \amalg_{\Spec(A)} \Spec(A')
$$
comme espaces topologiques.
\end{lemma}

\begin{proof}
Puisque $B' = B \times_A A'$, nous obtenons un carré commutatif de
spectres, qui induit une application continue
$$
can : \Spec(B) \amalg_{\Spec(A)} \Spec(A') \longrightarrow \Spec(B')
$$
puisque la source est une somme amalgamée dans la catégorie des espaces
topologiques (laquelle existe d'après
Topologie, section \ref{topology-section-colimits}).

\medskip\noindent
Pour montrer que l'application $can$ est surjective, soit
$\mathfrak q' \subset B'$ un idéal premier. Si $I \subset \mathfrak q'$
(ici et dans la suite, nous nous permettons de considérer $I$ aussi bien
comme un idéal de $B'$ que comme un idéal de $A'$), alors $\mathfrak q'$
correspond à un idéal premier de $B$ et appartient à l'image. Sinon,
choisissons $h \in I$, $h \not \in \mathfrak q'$. Dans ce cas,
$B_h = A_h = 0$ et l'homomorphisme d'anneaux $B'_h \to A'_h$ est un
isomorphisme, voir lemme \ref{lemma-diagram-localize}. Nous voyons donc que
$\mathfrak q'$ correspond à un unique idéal premier $\mathfrak p' \subset A'$
qui ne contient pas $I$.

\medskip\noindent
Puisque $B' \to B$ est surjectif, nous voyons que $can$ est injective sur
la composante $\Spec(B)$. Nous avons vu ci-dessus que
$\Spec(A') \to \Spec(B')$ est injective sur le complémentaire de
$V(I) \subset \Spec(A')$. Comme $V(I) \subset \Spec(A')$ est exactement
l'image de $\Spec(A) \to \Spec(A')$, un argument ensembliste immédiat montre
que $can$ est injective.

\medskip\noindent
Pour achever la preuve, nous devons montrer que $can$ est ouverte. Pour cela,
remarquons qu'un ouvert de la somme amalgamée est de la forme $V \amalg U'$,
où $V \subset \Spec(B)$ et $U' \subset \Spec(A')$ sont des ouverts dont les
images réciproques dans $\Spec(A)$ coïncident. Soit $v \in V$. Nous pouvons
trouver $g \in B$ tel que $v \in D(g) \subset V$. Soit $f \in A$ son image.
Choisissons $f' \in A'$ d'image $f$. Alors
$D(f') \cap U' \cap V(I) = D(f') \cap V(I)$.
Ainsi, $V(I) \cap D(f')$ et $D(f') \cap (U')^c$ sont des parties fermées
disjointes de $D(f') = \Spec(A'_{f'})$. Écrivons $(U')^c = V(J)$
pour un certain idéal $J \subset A'$. Puisque
$A'_{f'} \to A'_{f'}/IA'_{f'} \times A'_{f'}/JA'_{f'}$
est surjectif par la disjonction qui vient d'être établie, nous pouvons trouver
$a'' \in A'_{f'}$ dont l'image vaut $1$ dans $A'_{f'}/IA'_{f'}$ et zéro dans
$A'_{f'}/JA'_{f'}$. En chassant les dénominateurs, nous trouvons un élément
$a' \in J$ d'image $f^n$ dans $A$. Alors $D(a'f') \subset U'$.
Soit $h' = (g^{n + 1}, a'f') \in B'$.
Puisque $B'_{h'} = B_{g^{n + 1}} \times_{A_{f^{n + 1}}} A'_{a'f'}$ d'après un
lemme cité précédemment, nous voyons que l'image réciproque de $D(h')$ est un
voisinage ouvert de $v$ dans la somme amalgamée, c'est-à-dire que l'image de
$V \amalg U'$ contient un voisinage ouvert de l'image de $v$. Nous omettons
la preuve (plus facile) du même fait pour $u' \in U'$ avec
$u' \not \in V(I)$.
\end{proof}

\begin{lemma}
\label{lemma-fibre-product-integral}
Dans la situation \ref{situation-module-over-fibre-product},
si $B \to A$ est entier, alors $B' \to A'$ est entier.
\end{lemma}

\begin{proof}
Soit $a' \in A'$ d'image $a \in A$. Soit
$x^d + b_1 x^{d - 1} + \ldots + b_d$ un polynôme unitaire à coefficients
dans $B$ qui annule $a$. Choisissons $b'_i \in B'$ d'image $b_i \in B$
(c'est possible). Alors $(a')^d + b'_1 (a')^{d - 1} + \ldots + b'_d$
appartient au noyau de $A' \to A$. Puisque
$\Ker(B' \to B) = \Ker(A' \to A)$, nous pouvons modifier notre choix de
$b'_d$ afin d'obtenir $(a')^d + b'_1 (a')^{d - 1} + \ldots + b'_d = 0$,
comme voulu.
\end{proof}

\noindent
Dans la situation \ref{situation-module-over-fibre-product}, nous voudrions
comprendre les $B'$-modules en termes de modules sur $A'$, $A$ et $B$.
À cette fin, nous considérons le foncteur (où le produit fibré de catégories
est celui construit dans Catégories, exemple
\ref{categories-example-2-fibre-product-categories})
\begin{equation}
\label{equation-functor}
\text{Mod}_{B'} \longrightarrow
\text{Mod}_B \times_{\text{Mod}_A} \text{Mod}_{A'},\quad
L' \longmapsto (L' \otimes_{B'} B, L' \otimes_{B'} A', can)
\end{equation}
où $can$ est l'identification canonique
$L' \otimes_{B'} B \otimes_B A = L' \otimes_{B'} A' \otimes_{A'} A$.
Dans la suite, nous noterons $(N, M', \varphi)$ un objet du membre de droite,
c'est-à-dire que $N$ est un $B$-module, $M'$ un $A'$-module et que
$\varphi : N \otimes_B A \to M' \otimes_{A'} A$ est un isomorphisme.
Il est toutefois souvent plus commode de considérer $\varphi$ comme une
application $B$-linéaire $\varphi : N \to M'/IM'$ qui induit un isomorphisme
$N \otimes_B A \to M' \otimes_{A'} A = M'/IM'$.

\begin{lemma}
\label{lemma-module-over-fibre-product}
Dans la situation \ref{situation-module-over-fibre-product},
le foncteur (\ref{equation-functor}) admet un adjoint à droite, à savoir le
foncteur
$$
F : (N, M', \varphi) \longmapsto N \times_{\varphi, M} M'
$$
où $M = M'/IM'$. De plus, la composée de $F$ avec
(\ref{equation-functor}) est le foncteur identité sur
$\text{Mod}_B \times_{\text{Mod}_A} \text{Mod}_{A'}$. Autrement dit,
en posant $N' = N \times_{\varphi, M} M'$, nous avons
$N' \otimes_{B'} B = N$ et $N' \otimes_{B'} A' = M'$.
\end{lemma}

\begin{proof}
L'assertion d'adjonction résulte du lemme plus général
\ref{lemma-modules}. Pour démontrer la dernière assertion, rappelons que
$B' = B \times_A A'$ et $N' = N \times_{\varphi, M} M'$, puis prolongeons
ces égalités en
$$
\vcenter{
\xymatrix{
A & A' \ar[l] & I \ar[l] \\
B \ar[u] & B' \ar[l] \ar[u] & J \ar[l] \ar[u]
}
}
\quad\text{et}\quad
\vcenter{
\xymatrix{
M & M' \ar[l] & K \ar[l] \\
N \ar[u]_\varphi & N' \ar[l] \ar[u] & L \ar[l] \ar[u]
}
}
$$
où $I, J, K, L$ sont les noyaux des applications horizontales des diagrammes
initiaux. Présentons la preuve comme une suite d'observations :
\begin{enumerate}
\item $K = IM'$ (voir l'énoncé du lemme) ;
\item $B' \to B$ est surjectif de noyau $J$ et $J \to I$ est bijectif ;
\item $N' \to N$ est surjectif de noyau $L$ et $L \to K$ est bijectif ;
\item $JN' \subset L$,
\item $\Im(N \to M)$ engendre $M$ comme $A$-module
(car $N \otimes_B A = M$) ;
\item $\Im(N' \to M')$ engendre $M'$ comme $A'$-module
(car cela vaut modulo $K$ et $L$ s'envoie isomorphiquement sur $K$) ;
\item $JN' = L$ (car, d'après l'observation précédente, $L \cong K = I M'$
est engendré par les images des éléments $x n'$ avec $x \in I$ et
$n' \in N'$) ;
\item $N' \otimes_{B'} B = N$ (car $N = N'/L$, $B = B'/J$, et par
l'observation précédente) ;
\item il existe une application $\gamma : N' \otimes_{B'} A' \to M'$ ;
\item $\gamma$ est surjective (voir ci-dessus) ;
\item le noyau de la composée $N' \otimes_{B'} A' \to M' \to M$
est engendré par les éléments $l \otimes 1$ et $n' \otimes x$ avec
$l \in K$, $n' \in N'$, $x \in I$ (car $M = N \otimes_B A$ par hypothèse
et parce que $N' \to N$ et $A' \to A$ sont surjectifs de noyaux
$L$ et $I$) ;
\item tout élément de $N' \otimes_{B'} A'$ appartenant au sous-module
engendré par les éléments $l \otimes 1$ et $n' \otimes x$ avec
$l \in L$, $n' \in N'$, $x \in I$ peut s'écrire $l \otimes 1$
pour un certain $l \in L$ (puisque $J$ s'envoie isomorphiquement sur $I$,
nous voyons que $n' \otimes x = n'x \otimes 1$ dans $N' \otimes_{B'} A'$ ;
de même, $x n' \otimes a' = n' \otimes xa' = n'(xa') \otimes 1$
dans $N' \otimes_{B'} A'$ lorsque $n' \in N'$, $x \in J$ et $a' \in A'$ ;
comme nous avons vu que $JN' = L$, cela démontre l'assertion) ;
\item le noyau de $\gamma$ est nul (car, d'après (10) et (11), tout élément
du noyau est de la forme $l \otimes 1$ avec $l \in L$, lequel est envoyé sur
$l \in K \subset M'$ par $\gamma$).
\end{enumerate}
Cela achève la preuve.
\end{proof}

\begin{lemma}
\label{lemma-module-over-fibre-product-bis}
Dans la situation du lemme \ref{lemma-module-over-fibre-product},
pour un $B'$-module $L'$, l'application d'adjonction
$$
L' \longrightarrow 
(L' \otimes_{B'} B) \times_{(L' \otimes_{B'} A)} (L' \otimes_{B'} A')
$$
est surjective, mais n'est pas injective en général.
\end{lemma}

\begin{proof}
Comme dans la preuve du lemme \ref{lemma-module-over-fibre-product},
soit $J \subset B'$ le noyau de l'application $B' \to B$.
Alors $L' \otimes_{B'} B = L'/JL'$. Pour démontrer la surjectivité, il suffit
donc de montrer que les éléments de la forme $(0, z)$ du produit fibré
appartiennent à l'image de l'application du lemme. Le noyau de l'application
$L' \otimes_{B'} A' \to L' \otimes_{B'} A$ est l'image de
$L' \otimes_{B'} I \to L' \otimes_{B'} A'$. Puisque l'application $J \to I$
induite par $B' \to A'$ est un isomorphisme, la composée
$$
L' \otimes_{B'} J \to L' \to
(L' \otimes_{B'} B) \times_{(L' \otimes_{B'} A)} (L' \otimes_{B'} A')
$$
induit une surjection de $L' \otimes_{B'} J$ sur l'ensemble des éléments
de la forme $(0, z)$. Pour voir que l'application n'est pas injective en
général, donnons un exemple simple. Prenons un corps $k$, posons
$B' = k[x, y]/(xy)$, $A' = B'/(x)$, $B = B'/(y)$, $A = B'/(x, y)$ et
$L' = B'/(x - y)$. Dans ce cas, la classe de $x$ dans $L'$ est non nulle,
mais elle est envoyée sur zéro par la flèche affichée.
\end{proof}

\begin{lemma}
\label{lemma-surjection-module-over-fibre-product}
Dans la situation \ref{situation-module-over-fibre-product}, soit
$(N_1, M'_1, \varphi_1) \to (N_2, M'_2, \varphi_2)$ un morphisme de
$\text{Mod}_B \times_{\text{Mod}_A} \text{Mod}_{A'}$ tel que
$N_1 \to N_2$ et $M'_1 \to M'_2$ soient surjectifs. Alors
$$
N_1 \times_{\varphi_1, M_1} M'_1 \to N_2 \times_{\varphi_2, M_2} M'_2
$$
où $M_1 = M'_1/IM'_1$ et $M_2 = M'_2/IM'_2$, est surjectif.
\end{lemma}

\begin{proof}
Prenons $(x_2, y_2) \in N_2 \times_{\varphi_2, M_2} M'_2$. Choisissons
$x_1 \in N_1$ d'image $x_2$. Puisque $M'_1 \to M_1$ est surjectif, nous
pouvons trouver $y_1 \in M'_1$ d'image $\varphi_1(x_1)$. Alors $(x_1, y_1)$
s'envoie sur $(x_2, y'_2)$ dans $N_2 \times_{\varphi_2, M_2} M'_2$. Il suffit
donc de montrer que les éléments de la forme $(0, y_2)$ appartiennent à
l'image de l'application. Nous voyons ici que $y_2 \in IM'_2$. Écrivons
$y_2 = \sum t_i y_{2, i}$ avec $t_i \in I$. Choisissons
$y_{1, i} \in M'_1$ d'image $y_{2, i}$. Alors
$y_1 = \sum t_iy_{1, i} \in IM'_1$ et l'élément $(0, y_1)$ convient.
\end{proof}

\begin{lemma}
\label{lemma-finite-module-over-fibre-product}
Soient $A, A', B, B', I, M, M', N, \varphi$ comme dans le
lemme \ref{lemma-module-over-fibre-product}.
Si $N$ est de type fini sur $B$ et $M'$ de type fini sur $A'$, alors
$N' = N \times_{\varphi, M} M'$ est de type fini sur $B'$.
\end{lemma}

\begin{proof}
Nous utiliserons sans autre mention les résultats du
lemme \ref{lemma-module-over-fibre-product}. Choisissons des générateurs
$y_1, \ldots, y_r$ de $N$ sur $B$ et des générateurs $x_1, \ldots, x_s$
de $M'$ sur $A'$. Comme $N = N' \otimes_{B'} B$ et que $B' \to B$ est
surjectif, nous pouvons trouver $u_1, \ldots, u_r \in N'$ d'images
$y_1, \ldots, y_r$ dans $N$. Comme $M' = N' \otimes_{B'} A'$, nous pouvons
trouver $v_1, \ldots, v_t \in N'$ tels que
$x_i = \sum v_j \otimes a'_{ij}$ pour certains $a'_{ij} \in A'$.
En particulier, nous voyons que les images $\overline{v}_j \in M'$ des $v_j$
engendrent $M'$ sur $A'$. Nous affirmons que
$u_1, \ldots, u_r, v_1, \ldots, v_t$ engendrent $N'$ comme $B'$-module.
En effet, prenons $\xi \in N'$. Choisissons d'abord
$b'_1, \ldots, b'_r \in B'$ tels que $\xi$ et $\sum b'_i u_i$ aient la même
image dans $N$. C'est possible car $B' \to B$ est surjectif et
$y_1, \ldots, y_r$ engendrent $N$ sur $B$. La différence
$\xi - \sum b'_i u_i$ est de la forme $(0, \theta)$ pour un certain
$\theta$ dans $IM'$. Écrivons $\theta$ égal à
$\sum t_j\overline{v}_j$ avec
$t_j \in I$. Comme $J = \Ker(B' \to B)$ s'envoie isomorphiquement sur $I$,
nous pouvons choisir $s_j \in J \subset B'$ d'image $t_j$.
Puisque $N' = N \times_{\varphi, M} M'$, il s'ensuit que
$\xi = \sum b'_i u_i + \sum s_j v_j$, comme voulu.
\end{proof}

\begin{lemma}
\label{lemma-flat-module-over-fibre-product}
Soient $A, A', B, B', I$ comme dans la
situation \ref{situation-module-over-fibre-product}.
\begin{enumerate}
\item Soit $(N, M', \varphi)$ un objet de
$\text{Mod}_B \times_{\text{Mod}_A} \text{Mod}_{A'}$.
Si $M'$ est plat sur $A'$ et $N$ plat sur $B$, alors
$N' = N \times_{\varphi, M} M'$ est plat sur $B'$.
\item Si $L'$ est un $B'$-module plat, alors
$L' = (L \otimes_{B'} B) \times_{(L \otimes_{B'} A)} (L \otimes_{B'} A')$.
\item La catégorie des $B'$-modules plats est équivalente à la sous-catégorie
pleine de $\text{Mod}_B \times_{\text{Mod}_A} \text{Mod}_{A'}$ formée des
triplets $(N, M', \varphi)$ tels que $N$ soit plat sur $B$ et $M'$ plat
sur $A'$.
\end{enumerate}
\end{lemma}

\begin{proof}
Dans la preuve, nous utiliserons sans autre mention le
lemme \ref{lemma-module-over-fibre-product}.

\medskip\noindent
Preuve de (1). Posons $J = \Ker(B' \to B)$. C'est un idéal de $B'$ qui
s'envoie isomorphiquement sur $I = \Ker(A' \to A)$.
Soit $\mathfrak b' \subset B'$ un idéal. Nous devons montrer que
$\mathfrak b' \otimes_{B'} N' \to N'$ est injectif, voir Algèbre,
lemme \ref{algebra-lemma-flat}. Nous savons que
$$
\mathfrak b'/(\mathfrak b' \cap J) \otimes_{B'} N' =
\mathfrak b'/(\mathfrak b' \cap J) \otimes_B N \to N
$$
est injectif puisque $N$ est plat sur $B$. Comme
$\mathfrak b' \cap J \to \mathfrak b' \to
\mathfrak b'/(\mathfrak b' \cap J) \to 0$ est exacte, nous en concluons qu'il
suffit de montrer que
$(\mathfrak b' \cap J) \otimes_{B'} N' \to N'$
est injectif. Nous pouvons donc supposer que $\mathfrak b' \subset J$.
Ensuite, puisque $J \to I$ est un isomorphisme, nous avons
$$
J \otimes_{B'} N' =
I \otimes_{A'} A' \otimes_{B'} N' =
I \otimes_{A'} M'
$$
qui s'injecte dans $M'$ puisque $M'$ est un $A'$-module plat.
Ainsi, $J \otimes_{B'} N' \to N'$ est injectif et nous en concluons que
$\text{Tor}_1^{B'}(B'/J, N') = 0$, voir
Algèbre, remarque \ref{algebra-remark-Tor-ring-mod-ideal}.
Nous pouvons donc appliquer Algèbre, lemme
\ref{algebra-lemma-what-does-it-mean}, à $N'$ sur $B'$ et à l'idéal $J$.
Revenons à notre idéal $\mathfrak b' \subset J$ et soit
$\mathfrak b' \subset \mathfrak b'' \subset J$ le plus petit idéal dont
l'image dans $I$ soit un sous-$A'$-module de $I$. Autrement dit, nous avons
$\mathfrak b'' = A' \mathfrak b'$ si nous considérons $J = I$ comme un
$A'$-module. Alors $\mathfrak b''/\mathfrak b'$ est annulé par $J$ et nous
obtenons une suite exacte courte
$$
0 \to \mathfrak b' \otimes_{B'} N' \to
\mathfrak b'' \otimes_{B'} N' \to
\mathfrak b''/\mathfrak b' \otimes_{B'} N' \to 0
$$
grâce à l'annulation de $\text{Tor}_1^{B'}(\mathfrak b''/\mathfrak b', N')$
que nous donne l'application du lemme. Nous pouvons donc remplacer
$\mathfrak b'$ par $\mathfrak b''$. En particulier, nous pouvons supposer que
$\mathfrak b'$ est un $A'$-module et s'envoie sur un idéal de $A'$. Alors
$$
\mathfrak b' \otimes_{B'} N' =
\mathfrak b' \otimes_{A'} A' \otimes_{B'} N' =
\mathfrak b' \otimes_{A'} M'
$$
Ce produit tensoriel s'injecte dans $M'$ par notre hypothèse que $M'$ est
plat sur $A'$. Nous en concluons que
$\mathfrak b' \otimes_{B'} N' \to N' \to M'$ est injectif, et donc que la
première application est injective, comme voulu.

\medskip\noindent
Preuve de (2). Cela résulte du produit tensoriel de la suite exacte courte
$0 \to B' \to B \oplus A' \to A \to 0$ par $L'$ sur $B'$.

\medskip\noindent
Preuve de (3). C'est une conséquence immédiate de (1) et (2).
\end{proof}

\begin{lemma}
\label{lemma-finitely-presented-module-over-fibre-product}
Soient $A, A', B, B', I$ comme dans la
situation \ref{situation-module-over-fibre-product}.
La catégorie des $B'$-modules projectifs de type fini est équivalente à la
sous-catégorie pleine de $\text{Mod}_B \times_{\text{Mod}_A} \text{Mod}_{A'}$
formée des triplets $(N, M', \varphi)$ tels que $N$ soit projectif de type
fini sur $B$ et $M'$ projectif de type fini sur $A'$.
\end{lemma}

\begin{proof}
Rappelons qu'un module est projectif de type fini si et seulement s'il est
de présentation finie et plat, voir
Algèbre, lemme \ref{algebra-lemma-finite-projective}.
En utilisant les lemmes \ref{lemma-flat-module-over-fibre-product} et
\ref{lemma-finite-module-over-fibre-product}, nous nous ramenons à montrer
que $N' = N \times_{\varphi, M} M'$ est un $B'$-module de présentation finie
si $N$ est projectif de type fini sur $B$ et $M'$ projectif de type fini
sur $A'$.

\medskip\noindent
D'après le lemme \ref{lemma-finite-module-over-fibre-product}, le module $N'$
est de type fini sur $B'$. Choisissons une surjection
$(B')^{\oplus n} \to N'$ de noyau $K'$. Par changement de base, nous obtenons
des applications $B^{\oplus n} \to N$, $(A')^{\oplus n} \to M'$ et
$A^{\oplus n} \to M$ de noyaux respectifs $K_B$, $K_{A'}$ et $K_A$.
Il existe une application canonique
$$
K' \longrightarrow K_B \times_{K_A} K_{A'}
$$
D'autre part, puisque $N' = N \times_{\varphi, M} M'$ et
$B' = B \times_A A'$, il existe aussi une application canonique
$K_B \times_{K_A} K_{A'} \to K'$ inverse de la flèche affichée.
L'application affichée est donc un isomorphisme. D'après
Algèbre, lemme \ref{algebra-lemma-extension}, les modules $K_B$ et $K_{A'}$
sont de type fini. Nous concluons du
lemme \ref{lemma-finite-module-over-fibre-product} que $K'$ est un
$B'$-module de type fini pourvu que $K_B \to K_A$ et $K_{A'} \to K_A$
induisent des isomorphismes
$K_B \otimes_B A = K_A = K_{A'} \otimes_{A'} A$.
C'est vrai car les hypothèses de platitude impliquent que les suites
$$
0 \to K_B \to B^{\oplus n} \to N \to 0
\quad\text{et}\quad
0 \to K_{A'} \to (A')^{\oplus n} \to M' \to 0
$$
restent exactes après produit tensoriel, voir
Algèbre, lemme \ref{algebra-lemma-flat-tor-zero}.
\end{proof}





\section{Produits fibrés d'anneaux, III}
\label{section-fibre-products-rings-III}

\noindent
Dans cette section, nous étudions les produits fibrés dans la situation
suivante.

\begin{situation}
\label{situation-relative-module-over-fibre-product}
Soient $A, A', B, B', I$ comme dans la
situation \ref{situation-module-over-fibre-product}.
Soit $B' \to D'$ un homomorphisme d'anneaux. Posons
$D = D' \otimes_{B'} B$, $C' = D' \otimes_{B'} A'$ et
$C = D' \otimes_{B'} A$. Cela conduit à un grand diagramme commutatif
$$
\xymatrix{
C & & & C' \ar[lll] \\
& A \ar[ul] & A' \ar[l] \ar[ru] \\
& B \ar[u] \ar[ld] & B' \ar[l] \ar[u] \ar[rd] \\
D \ar[uuu] & & & D' \ar[lll] \ar[uuu]
}
$$
d'anneaux.
Remarquons que nous ne supposons {\bf pas} que l'application
$D' \to D \times_C C'$ soit un isomorphisme\footnote{Mais
$D' \to D \times_C C'$ est surjectif d'après le
lemme \ref{lemma-module-over-fibre-product-bis}.}.
Dans cette situation, nous avons le foncteur
\begin{equation}
\label{equation-relative-functor}
\text{Mod}_{D'} \longrightarrow
\text{Mod}_D \times_{\text{Mod}_C} \text{Mod}_{C'},\quad
L' \longmapsto (L' \otimes_{D'} D, L' \otimes_{D'} C', can)
\end{equation}
analogue à (\ref{equation-functor}). Remarquons que
$L' \otimes_{D'} D = L \otimes_{D'} (D' \otimes_{B'} B) = L \otimes_{B'} B$
et, de même,
$L' \otimes_{D'} C' = L \otimes_{D'} (D' \otimes_{B'} A') = L \otimes_{B'} A'$
de sorte que le diagramme
$$
\xymatrix{
\text{Mod}_{D'} \ar[r] \ar[d] &
\text{Mod}_D \times_{\text{Mod}_C} \text{Mod}_{C'} \ar[d] \\
\text{Mod}_{B'} \ar[r] &
\text{Mod}_B \times_{\text{Mod}_A} \text{Mod}_{A'}
}
$$
est commutatif. Dans la suite, nous noterons $(N, M', \varphi)$ un objet de
$\text{Mod}_D \times_{\text{Mod}_C} \text{Mod}_{C'}$, c'est-à-dire que $N$
est un $D$-module, $M'$ un $C'$-module et que
$\varphi : N \otimes_B A \to M' \otimes_{A'} A$ est un isomorphisme de
$C$-modules. Il est toutefois souvent plus commode de considérer $\varphi$
comme une application $D$-linéaire $\varphi : N \to M'/IM'$ qui induit un
isomorphisme $N \otimes_B A \to M' \otimes_{A'} A = M'/IM'$.
\end{situation}

\begin{lemma}
\label{lemma-relative-module-over-fibre-product}
Dans la situation \ref{situation-relative-module-over-fibre-product},
le foncteur (\ref{equation-relative-functor}) admet un adjoint à droite, à
savoir le foncteur
$$
F : (N, M', \varphi) \longmapsto N \times_{\varphi, M} M'
$$
où $M = M'/IM'$. De plus, la composée de $F$ avec
(\ref{equation-relative-functor}) est le foncteur identité sur
$\text{Mod}_D \times_{\text{Mod}_C} \text{Mod}_{C'}$. Autrement dit,
en posant $N' = N \times_{\varphi, M} M'$, nous avons
$N' \otimes_{D'} D = N$ et $N' \otimes_{D'} C' = M'$.
\end{lemma}

\begin{proof}
L'assertion d'adjonction résulte du lemme plus général
\ref{lemma-modules}. La dernière assertion résulte de l'assertion
correspondante du lemme \ref{lemma-module-over-fibre-product}, car
$N' \otimes_{D'} D = N' \otimes_{D'} D' \otimes_{B'} B = N' \otimes_{B'} B$
et
$N' \otimes_{D'} C' = N' \otimes_{D'} D' \otimes_{B'} A' = N' \otimes_{B'} A'$.
\end{proof}

\begin{lemma}
\label{lemma-relative-surjection-ideals}
Dans la situation \ref{situation-relative-module-over-fibre-product},
l'application $JD' \to IC'$ est surjective, où $J = \Ker(B' \to B)$.
\end{lemma}

\begin{proof}
Puisque $C' = D' \otimes_{B'} A'$, nous avons que $IC'$ est l'image de
$D' \otimes_{B'} I = C' \otimes_{A'} I \to C'$. Comme l'homomorphisme
d'anneaux $B' \to A'$ induit un isomorphisme $J \to I$, le lemme en résulte.
\end{proof}

\begin{lemma}
\label{lemma-relative-finite-module-over-fibre-product}
Soient $A, A', B, B', C, C', D, D', I, M', M, N, \varphi$ comme dans le
lemme \ref{lemma-relative-module-over-fibre-product}.
Si $N$ est de type fini sur $D$ et $M'$ de type fini sur $C'$, alors
$N' = N \times_{\varphi, M} M'$ est de type fini sur $D'$.
\end{lemma}

\begin{proof}
Rappelons que $D' \to D \times_C C'$ est surjectif d'après le
lemme \ref{lemma-module-over-fibre-product-bis}.
Remarquons que $N' = N \times_{\varphi, M} M'$ est un module sur
$D \times_C C'$. Nous pouvons appliquer le
lemme \ref{lemma-finite-module-over-fibre-product} aux données
$C, C', D, D', IC', M', M, N, \varphi$ pour voir que
$N' = N \times_{\varphi, M} M'$ est de type fini sur $D \times_C C'$.
Il est donc de type fini sur $D'$.
\end{proof}

\begin{lemma}
\label{lemma-relative-flat-module-over-fibre-product}
Soient $A, A', B, B', C, C', D, D', I$ comme dans la
situation \ref{situation-relative-module-over-fibre-product}.
\begin{enumerate}
\item Soit $(N, M', \varphi)$ un objet de
$\text{Mod}_D \times_{\text{Mod}_C} \text{Mod}_{C'}$.
Si $M'$ est plat sur $A'$ et $N$ plat sur $B$, alors
$N' = N \times_{\varphi, M} M'$ est plat sur $B'$.
\item Si $L'$ est un $D'$-module plat sur $B'$, alors
$L' = (L \otimes_{D'} D) \times_{(L \otimes_{D'} C)} (L \otimes_{D'} C')$.
\item La catégorie des $D'$-modules plats sur $B'$ est équivalente à la
catégorie des objets $(N, M', \varphi)$ de
$\text{Mod}_D \times_{\text{Mod}_C} \text{Mod}_{C'}$ tels que $N$ soit plat
sur $B$ et $M'$ plat sur $A'$.
\end{enumerate}
\end{lemma}

\begin{proof}
La partie (1) résulte de la partie (1) du
lemme \ref{lemma-flat-module-over-fibre-product}.

\medskip\noindent
La partie (2) résulte de la partie (2) du
lemme \ref{lemma-flat-module-over-fibre-product}, compte tenu de
$L' \otimes_{D'} D = L' \otimes_{B'} B$,
$L' \otimes_{D'} C' = L' \otimes_{B'} A'$ et
$L' \otimes_{D'} C = L' \otimes_{B'} A$ ; voir la discussion dans la
situation \ref{situation-relative-module-over-fibre-product}.

\medskip\noindent
La partie (3) est une conséquence immédiate de (1) et (2).
\end{proof}

\noindent
Le lemme suivant est nettement plus intéressant que son analogue dans le cas
absolu (lemme \ref{lemma-finitely-presented-module-over-fibre-product}),
bien que la preuve soit essentiellement la même.

\begin{lemma}
\label{lemma-relative-finitely-presented-module-over-fibre-product}
Soient $A, A', B, B', C, C', D, D', I, M', M, N, \varphi$ comme dans le
lemme \ref{lemma-relative-module-over-fibre-product}. Si
\begin{enumerate}
\item $N$ est de présentation finie sur $D$ et plat sur $B$ ;
\item $M'$ est de présentation finie sur $C'$ et plat sur $A'$ ; et
\item l'homomorphisme d'anneaux $B' \to D'$ se factorise en
$B' \to D'' \to D'$, avec $B' \to D''$ plat et $D'' \to D'$ de présentation
finie,
\end{enumerate}
alors $N' = N \times_M M'$ est de présentation finie sur $D'$.
\end{lemma}

\begin{proof}
Choisissons une surjection $D''' = D''[x_1, \ldots, x_n] \to D'$ de noyau
$J$ de type fini. D'après Algèbre, lemme
\ref{algebra-lemma-finite-finitely-presented-extension}, il suffit de montrer
que $N'$ est de présentation finie comme $D'''$-module. De plus,
$D''' \otimes_{B'} B \to D' \otimes_{B'} B = D$ et
$D''' \otimes_{B'} A' \to D' \otimes_{B'} A' = C'$ sont des surjections dont
les noyaux sont engendrés par l'image de $J$ ; ainsi, $N$ est un
$D''' \otimes_{B'} B$-module de présentation finie et $M'$ un
$D''' \otimes_{B'} A'$-module de présentation finie, à nouveau d'après
Algèbre, lemme \ref{algebra-lemma-finite-finitely-presented-extension}.
Nous pouvons donc remplacer $D'$ par $D'''$ et $D$ par
$D''' \otimes_{B'} B$, etc. Puisque $D'''$ est plat sur $B'$, il s'ensuit que
nous pouvons supposer que $B' \to D'$ est plat.

\medskip\noindent
Supposons que $B' \to D'$ soit plat.
D'après le lemme \ref{lemma-relative-finite-module-over-fibre-product},
le module $N'$ est de type fini sur $D'$. Choisissons une surjection
$(D')^{\oplus n} \to N'$ de noyau $K'$. Par changement de base, nous obtenons
des applications $D^{\oplus n} \to N$, $(C')^{\oplus n} \to M'$ et
$C^{\oplus n} \to M$ de noyaux respectifs $K_D$, $K_{C'}$ et $K_C$.
Il existe une application canonique
$$
K' \longrightarrow K_D \times_{K_C} K_{C'}
$$
D'autre part, puisque $N' = N \times_M M'$ et
$D' = D \times_C C'$ (d'après le
lemme \ref{lemma-flat-module-over-fibre-product}, appliqué au $B'$-module plat
$D'$), il existe aussi une application canonique
$K_D \times_{K_C} K_{C'} \to K'$ inverse de la flèche affichée.
L'application affichée est donc un isomorphisme. D'après
Algèbre, lemme \ref{algebra-lemma-extension}, les modules $K_D$ et $K_{C'}$
sont de type fini. Nous concluons du
lemme \ref{lemma-relative-finite-module-over-fibre-product} que $K'$ est un
$D'$-module de type fini pourvu que $K_D \to K_C$ et $K_{C'} \to K_C$
induisent des isomorphismes
$K_D \otimes_B A = K_C = K_{C'} \otimes_{A'} A$.
C'est vrai car les hypothèses de platitude impliquent que les suites
$$
0 \to K_D \to D^{\oplus n} \to N \to 0
\quad\text{et}\quad
0 \to K_{C'} \to (C')^{\oplus n} \to M' \to 0
$$
restent exactes après produit tensoriel, voir
Algèbre, lemme \ref{algebra-lemma-flat-tor-zero}.
\end{proof}

\begin{lemma}
\label{lemma-properties-algebras-over-fibre-product}
Soient $A, A', B, B', I$ comme dans la
situation \ref{situation-module-over-fibre-product}.
Soit $(D, C', \varphi)$ un système constitué d'une $B$-algèbre $D$,
d'une $A'$-algèbre $C'$ et d'un isomorphisme $D \otimes_B A \to C'/IC' = C$.
Posons $D' = D \times_C C'$ (comme dans le
lemme \ref{lemma-module-over-fibre-product}). Alors
\begin{enumerate}
\item $B' \to D'$ est de type fini si et seulement si $B \to D$ et
$A' \to C'$ sont de type fini ;
\item $B' \to D'$ est plat si et seulement si $B \to D$ et $A' \to C'$ sont plats ;
\item $B' \to D'$ est plat et de présentation finie si et seulement si
$B \to D$ et $A' \to C'$ sont plats et de présentation finie ;
\item $B' \to D'$ est lisse si et seulement si $B \to D$ et $A' \to C'$ sont lisses ;
\item $B' \to D'$ est \'etale si et seulement si $B \to D$ et $A' \to C'$ sont
\'etales.
\end{enumerate}
De plus, si $D'$ est une $B'$-algèbre plate, alors
$D' \to (D' \otimes_{B'} B) \times_{(D' \otimes_{B'} A)} (D' \otimes_{B'} A')$
est un isomorphisme. De cette manière, la catégorie des $B'$-algèbres plates
est équivalente à la catégorie des systèmes $(D, C', \varphi)$ ci-dessus
avec $D$ plat sur $B$ et $C'$ plat sur $A'$.
\end{lemma}

\begin{proof}
L'implication « $\Rightarrow$ » résulte des lemmes d'Algèbre
\ref{algebra-lemma-base-change-finiteness},
\ref{algebra-lemma-flat-base-change},
\ref{algebra-lemma-base-change-smooth} et
\ref{algebra-lemma-etale}, car nous avons
$D' \otimes_{B'} B = D$ et $D' \otimes_{B'} A' = C'$
par le lemme \ref{lemma-module-over-fibre-product}.
Il suffit donc de démontrer les implications dans l'autre sens.

\medskip\noindent
Ad (1). Supposons que $D$ soit de type fini sur $B$ et $C'$ de type fini sur
$A'$. Nous utiliserons sans autre mention les résultats du
lemme \ref{lemma-module-over-fibre-product}. Choisissons des générateurs
$x_1, \ldots, x_r$ de $D$ sur $B$ et des générateurs $y_1, \ldots, y_s$
de $C'$ sur $A'$. Comme
$D = D' \otimes_{B'} B$ et que $B' \to B$ est surjectif, nous pouvons trouver
$u_1, \ldots, u_r \in D'$ s'envoyant sur $x_1, \ldots, x_r$ dans $D$.
Comme $C' = D' \otimes_{B'} A'$, nous pouvons trouver
$v_1, \ldots, v_t \in D'$ tels que $y_i = \sum v_j \otimes a'_{ij}$ pour
certains $a'_{ij} \in A'$. En particulier, les images des $v_j$ dans $C'$ engendrent
$C'$ comme $A'$-algèbre. Posons $N = r + t$ et considérons le cube d'anneaux
$$
\xymatrix{
A[x_1, \ldots, x_N] & & A'[x_1, \ldots, x_N] \ar[ll] \\
& A \ar[lu] & & A' \ar[ll] \ar[lu] \\
B[x_1, \ldots, x_N] \ar[uu] & & B'[x_1, \ldots, x_N] \ar[uu] \ar[ll] \\
& B \ar[uu] \ar[lu] & & B' \ar[ll] \ar[uu] \ar[lu]
}
$$
Remarquons que le carré arrière est lui aussi cartésien.
Considérons l'homomorphisme d'anneaux
$$
B'[x_1, \ldots, x_N] \to D',\quad
x_i \mapsto u_i \quad\text{et}\quad x_{r + j} \mapsto v_j.
$$
Nous voyons alors que les applications induites
$B[x_1, \ldots, x_N] \to D$ et $A'[x_1, \ldots, x_N] \to C'$ sont
surjectives, en particulier finies. Nous concluons du
lemme \ref{lemma-relative-finite-module-over-fibre-product} que
$B'[x_1, \ldots, x_N] \to D'$ est fini, ce qui implique que $D'$ est de
type fini sur $B'$, par exemple d'après Algèbre, lemme
\ref{algebra-lemma-compose-finite-type}.

\medskip\noindent
Ad (2). L'implication « $\Leftarrow$ » résulte du
lemme \ref{lemma-relative-flat-module-over-fibre-product}.
De plus, l'assertion finale résulte de l'assertion finale du
lemme \ref{lemma-relative-flat-module-over-fibre-product}.

\medskip\noindent
Ad (3). Supposons que $B \to D$ et $A' \to C'$ soient plats et de
présentation finie. La platitude de $B' \to D'$ a été démontrée en (2).
Nous savons par (1) que $B' \to D'$ est de type fini. Choisissons une
surjection $B'[x_1, \ldots, x_N] \to D'$. D'après Algèbre, lemme
\ref{algebra-lemma-finite-presentation-independent}, l'anneau $D$ est de
présentation finie comme $B[x_1, \ldots, x_N]$-module et l'anneau $C'$ est
de présentation finie comme $A'[x_1, \ldots, x_N]$-module. D'après le
lemme \ref{lemma-relative-finitely-presented-module-over-fibre-product},
$D'$ est de présentation finie comme $B'[x_1, \ldots, x_N]$-module,
c'est-à-dire que $B' \to D'$ est de présentation finie.

\medskip\noindent
Ad (4). Supposons que $B \to D$ et $A' \to C'$ soient lisses.
D'après (3), $B' \to D'$ est plat et de présentation finie.
D'après Algèbre, lemme \ref{algebra-lemma-flat-fibre-smooth}, il suffit de
vérifier que $D' \otimes_{B'} k$ est lisse pour tout corps $k$ sur $B'$.
Si la composée $J \to B' \to k$ est nulle, alors $B' \to k$ se factorise
par $B' \to B \to k$ et nous voyons que
$$
D' \otimes_{B'} k = D' \otimes_{B'} B \otimes_B k
= D \otimes_B k
$$
est lisse puisque $B \to D$ est lisse. Si la composée $J \to B' \to k$
n'est pas nulle, il existe alors $h \in J$ dont l'image dans $k$ n'est pas nulle.
Ainsi, $B' \to k$ se factorise par $B' \to B'_h \to k$.
Remarquons que $h$ s'envoie sur zéro dans $B$, d'où $B_h = 0$.
D'après le lemme \ref{lemma-diagram-localize}, nous avons donc $B'_h = A'_h$
et nous obtenons
$$
D' \otimes_{B'} k = D' \otimes_{B'} B'_h \otimes_{B'_h} k
= C'_h \otimes_{A'_h} k
$$
est lisse puisque $A' \to C'$ est lisse.

\medskip\noindent
Ad (5). Supposons que $B \to D$ et $A' \to C'$ soient \'etales. D'après (4),
$B' \to D'$ est lisse. Comme on peut déterminer si une application lisse est
\'etale ou non à partir de la dimension de ses fibres, (5) est vérifié
(raisonner comme dans la preuve de (4) pour identifier les fibres ; certains
détails sont omis).
\end{proof}

\begin{remark}
\label{remark-relative-modules-over-fibre-product}
Dans la situation \ref{situation-relative-module-over-fibre-product},
supposons que $B' \to D'$ soit de présentation finie et
qu'on se donne un $D'$-module $L'$. Nous affirmons qu'il existe une
correspondance bijective entre
\begin{enumerate}
\item les surjections de $D'$-modules $L' \to Q'$ avec $Q'$ de présentation
finie sur $D'$ et plat sur $B'$ ; et
\item les paires de surjections de modules
$(L' \otimes_{D'} D \to Q_1, L' \otimes_{D'} C' \to Q_2)$
avec
\begin{enumerate}
\item $Q_1$ de présentation finie sur $D$ et plat sur $B$ ;
\item $Q_2$ de présentation finie sur $C'$ et plat sur $A'$ ;
\item $Q_1 \otimes_D C = Q_2 \otimes_{C'} C$ comme quotients de
$L' \otimes_{D'} C$.
\end{enumerate}
\end{enumerate}
La correspondance est donnée par $Q \mapsto (Q_1, Q_2)$ avec
$Q_1 = Q \otimes_{D'} D$ et $Q_2 = Q \otimes_{D'} C'$. Réciproquement,
nous posons $Q = Q_1 \times_{Q_{12}} Q_2$, où $Q_{12}$ est le quotient
commun $Q_1 \otimes_D C = Q_2 \otimes_{C'} C$ de $L' \otimes_{D'} C$.
Comme application quotient, nous utilisons
$$
L' \longrightarrow
(L' \otimes_{D'} D) \times_{(L' \otimes_{D'} C)} (L' \otimes_{D'} C')
\longrightarrow Q_1 \times_{Q_{12}} Q_2 = Q
$$
où la première flèche est surjective d'après le
lemme \ref{lemma-module-over-fibre-product-bis} et la seconde d'après le
lemme \ref{lemma-surjection-module-over-fibre-product}.
L'affirmation résulte des lemmes
\ref{lemma-relative-flat-module-over-fibre-product} et
\ref{lemma-relative-finitely-presented-module-over-fibre-product}.
\end{remark}








\section{Idéaux de Fitting}
\label{section-fitting-ideals}

\noindent
Les idéaux de Fitting d'un module de type fini sont les idéaux déterminés
par la construction du lemme \ref{lemma-fitting-ideal}.

\begin{lemma}
\label{lemma-ideals-generated-by-minors}
Soit $R$ un anneau. Soit $A$ une matrice $n \times m$ à coefficients
dans $R$. Soit $I_r(A)$ l'idéal engendré par les mineurs $r \times r$
de $A$, avec la convention $I_0(A) = R$ et $I_r(A) = 0$ si
$r > \min(n, m)$. Alors
\begin{enumerate}
\item $I_0(A) \supset I_1(A) \supset I_2(A) \supset \ldots$ ;
\item si $B$ est une matrice $(n + n') \times m$ et si $A$ est constituée
des $n$ premières lignes de $B$, alors $I_{r + n'}(B) \subset I_r(A)$ ;
\item si $C$ est une matrice $n \times n$, alors $I_r(CA) \subset I_r(A)$ ;
\item si $A$ est une matrice par blocs
$$
\left(
\begin{matrix}
A_1 & 0 \\
0 & A_2 
\end{matrix}
\right)
$$
alors $I_r(A) = \sum_{r_1 + r_2 = r} I_{r_1}(A_1) I_{r_2}(A_2)$ ;
\item Ajouter d'autres propriétés ici.
\end{enumerate}
\end{lemma}

\begin{proof}
Omis. (Indication : utiliser le fait qu'un déterminant se calcule en
développant suivant une colonne ou une ligne.)
\end{proof}

\begin{lemma}
\label{lemma-fitting-ideal}
Soit $R$ un anneau. Soit $M$ un $R$-module de type fini. Choisissons une
présentation
$$
\bigoplus\nolimits_{j \in J} R \longrightarrow R^{\oplus n}
\longrightarrow M \longrightarrow 0.
$$
de $M$. Soit $A = (a_{ij})_{i = 1, \ldots, n, j \in J}$ la matrice
de l'application $\bigoplus_{j \in J} R \to R^{\oplus n}$.
L'idéal $\text{Fit}_k(M)$ engendré par les mineurs
$(n - k) \times (n - k)$ de $A$ est indépendant du choix de la présentation.
\end{lemma}

\begin{proof}
Soit $K \subset R^{\oplus n}$ le noyau de la surjection
$R^{\oplus n} \to M$. Choisissons $z_1, \ldots, z_{n - k} \in K$
et écrivons $z_j = (z_{1j}, \ldots, z_{nj})$.
Une autre description de l'idéal $\text{Fit}_k(M)$ est l'idéal engendré par
les mineurs $(n - k) \times (n - k)$ de toutes les matrices $(z_{ij})$
obtenues de cette façon.

\medskip\noindent
Supposons que nous remplacions la surjection par la surjection
$R^{\oplus n + n'} \to M$ de noyau $K'$, où nous utilisons l'application
originale sur les $n$ premiers vecteurs de base standard de
$R^{\oplus n + n'}$ et $0$ sur les $n'$ derniers. Les idéaux correspondants
sont alors les mêmes. En effet, si $z_1, \ldots, z_{n - k} \in K$ comme ci-dessus,
posons $z'_j = (z_{1j}, \ldots, z_{nj}, 0, \ldots, 0) \in K'$ pour
$j = 1, \ldots, n - k$ et
$z'_{n + j'} = (0, \ldots, 0, 1, 0, \ldots, 0) \in K'$. Nous voyons alors que
l'idéal des mineurs $(n - k) \times (n - k)$ de $(z_{ij})$ coïncide
avec l'idéal des mineurs $(n + n' - k) \times (n + n' - k)$ de
$(z'_{ij})$. Cela donne l'une des inclusions.
Réciproquement, étant donnés $z'_1, \ldots, z'_{n + n' - k}$
dans $K'$, projetons-les sur $R^{\oplus n}$ pour obtenir
$z_1, \ldots, z_{n + n' - k}$ dans $K$. D'après le
lemme \ref{lemma-ideals-generated-by-minors}, nous voyons que l'idéal engendré par les
mineurs $(n + n' - k) \times (n + n' - k)$ de
$(z'_{ij})$ est contenu dans l'idéal engendré par les
mineurs $(n - k) \times (n - k)$ de $(z_{ij})$. Cela donne l'autre inclusion.

\medskip\noindent
Soit $R^{\oplus m} \to M$ une autre surjection de noyau $L$.
D'après le lemme de Schanuel (Algèbre, lemme \ref{algebra-lemma-Schanuel})
et les résultats du paragraphe précédent, nous pouvons supposer $m = n$ et
qu'il existe un isomorphisme $R^{\oplus n} \to R^{\oplus m}$ commutant aux
surjections vers $M$. Soit $C = (c_{li})$ la matrice (inversible) de cette
application (c'est une matrice carrée puisque $n = m$). Alors, pour des
$z'_1, \ldots, z'_{n - k} \in L$ comme ci-dessus, nous pouvons trouver
$z_1, \ldots, z_{n - k} \in K$ tels que
$z_1' = Cz_1, \ldots, z'_{n - k} = Cz_{n - k}$. D'après le
lemme \ref{lemma-ideals-generated-by-minors}, nous obtenons l'une des
inclusions. Par symétrie, nous obtenons l'autre.
\end{proof}

\begin{definition}
\label{definition-fitting-ideal}
Soit $R$ un anneau. Soit $M$ un $R$-module de type fini. Soit $k \geq 0$.
L'{\it idéal de Fitting d'indice $k$} de $M$ est l'idéal $\text{Fit}_k(M)$
construit dans le lemme \ref{lemma-fitting-ideal}. Posons
$\text{Fit}_{-1}(M) = 0$.
\end{definition}

\noindent
Comme les idéaux de Fitting sont les idéaux de mineurs d'une grande matrice
(numérotés dans l'ordre inverse de celui du
lemme \ref{lemma-ideals-generated-by-minors}), nous voyons que
$$
0 = \text{Fit}_{-1}(M) \subset \text{Fit}_0(M) \subset \text{Fit}_1(M)
\subset \ldots \subset \text{Fit}_t(M) = R
$$
pour un certain $t \gg 0$. Voici quelques propriétés élémentaires des idéaux
de Fitting.

\begin{lemma}
\label{lemma-fitting-ideal-basics}
Soit $R$ un anneau. Soit $M$ un $R$-module de type fini.
\begin{enumerate}
\item Si $M$ peut être engendré par $n$ éléments, alors
$\text{Fit}_n(M) = R$.
\item Étant donné un second $R$-module $M'$ de type fini, nous avons
$\text{Fit}_0(M \oplus M') = \text{Fit}_0(M)\text{Fit}_0(M')$
et plus généralement
$$
\text{Fit}_l(M \oplus M') =
\sum\nolimits_{k + k' = l} \text{Fit}_k(M)\text{Fit}_{k'}(M')
$$
\item Si $R \to R'$ est un homomorphisme d'anneaux, alors
$\text{Fit}_k(M \otimes_R R')$ est l'idéal de $R'$ engendré par l'image de
$\text{Fit}_k(M)$ ;
\item Si $M$ est de présentation finie, alors $\text{Fit}_k(M)$ est un idéal
de type fini ;
\item Si $M \to M'$ est une surjection, alors
$\text{Fit}_k(M) \subset \text{Fit}_k(M')$.
\item Nous avons $\text{Fit}_0(M) \subset \text{Ann}_R(M)$ ;
\item Nous avons $V(\text{Fit}_0(M)) = \text{Supp}(M)$ ;
\item Si $I$ est un idéal de $R$, alors $\text{Fit}_0(R/I) = I$ ;
\item Ajouter d'autres propriétés ici.
\end{enumerate}
\end{lemma}

\begin{proof}
La partie (1) résulte du fait que $I_0(A) = R$ dans le
lemme \ref{lemma-ideals-generated-by-minors}.

\medskip\noindent
La partie (2) résulte de l'énoncé correspondant du
lemme \ref{lemma-ideals-generated-by-minors}.

\medskip\noindent
La partie (3) résulte du fait que $\otimes_R R'$ est exact à droite,
de sorte que le changement de base d'une présentation de $M$ est une
présentation de $M \otimes_R R'$.

\medskip\noindent
Preuve de (4). Soit $R^{\oplus m} \xrightarrow{A} R^{\oplus n} \to M \to 0$
une présentation. Alors $\text{Fit}_k(M)$ est l'idéal engendré par les
mineurs $(n - k) \times (n - k)$ de la matrice $A$.

\medskip\noindent
La partie (5) est immédiate d'après la définition.

\medskip\noindent
Preuve de (6). Choisissons une présentation de $M$ de matrice $A$ comme dans
le lemme \ref{lemma-fitting-ideal}. Soit $J' \subset J$ une partie de
cardinal $n$. Il suffit de montrer que
$f = \det(a_{ij})_{i = 1, \ldots, n, j \in J'}$
annule $M$. Cela est clair car le conoyau de
$$
R^{\oplus n} \xrightarrow{A' = (a_{ij})_{i = 1, \ldots, n, j \in J'}}
R^{\oplus n} \to M \to 0
$$
est annulé par $f$, puisqu'il existe une matrice $B$ telle que
$A' B = f1_{n \times n}$.

\medskip\noindent
Preuve de (7). Choisissons une présentation de $M$ de matrice $A$ comme dans
le lemme \ref{lemma-fitting-ideal}. D'après le lemme de Nakayama
(Algèbre, lemme \ref{algebra-lemma-NAK}), nous avons
$$
M_\mathfrak p \not = 0
\Leftrightarrow
M \otimes_R \kappa(\mathfrak p) \not = 0
\Leftrightarrow
\text{rang}(\text{image de }A\text{ dans }\kappa(\mathfrak p)) < n
$$
Clairement, $\text{Fit}_0(M)$ définit exactement l'ensemble des idéaux
premiers possédant cette propriété.
\end{proof}

\begin{example}
\label{example-fitting-free}
Soit $R$ un anneau.
Les idéaux de Fitting du module libre de type fini $M = R^{\oplus n}$
sont $\text{Fit}_k(M) = 0$ pour $k < n$ et $\text{Fit}_k(M) = R$
pour $k \geq n$.
\end{example}

\begin{example}
\label{example-laurent}
Soit $R$ un anneau et soient $I, J\subset R$ des idéaux.
Alors $\text{Fit}_0(R/I) = I$, $\text{Fit}_0(R/I \oplus R/J) = IJ$,
$\text{Fit}_1(R/I\oplus R/J) = I+J$.
\end{example}

\begin{lemma}
\label{lemma-fitting-ideal-generate-locally}
Soit $R$ un anneau. Soit $M$ un $R$-module de type fini. Soit $k \geq 0$.
Soit $\mathfrak p \subset R$ un idéal premier. Les conditions suivantes
sont équivalentes :
\begin{enumerate}
\item $\text{Fit}_k(M) \not \subset \mathfrak p$ ;
\item $\dim_{\kappa(\mathfrak p)} M \otimes_R \kappa(\mathfrak p) \leq k$ ;
\item $M_\mathfrak p$ peut être engendré par $k$ éléments sur $R_\mathfrak p$ ;
\item $M_f$ peut être engendré par $k$ éléments sur $R_f$ pour un certain
$f \in R$, $f \not \in \mathfrak p$.
\end{enumerate}
\end{lemma}

\begin{proof}
D'après le lemme de Nakayama (Algèbre, lemme \ref{algebra-lemma-NAK}),
$M_f$ peut être engendré par $k$ éléments sur $R_f$ pour un certain
$f \in R$, $f \not \in \mathfrak p$, si $M \otimes_R \kappa(\mathfrak p)$
peut être engendré par $k$ éléments. Ainsi, (2), (3) et (4) sont équivalentes.
En utilisant la partie (3) du lemme \ref{lemma-fitting-ideal-basics}, nous
nous ramenons au cas où $R$ est un corps et $\mathfrak p = (0)$. Dans ce cas,
le résultat découle de l'exemple \ref{example-fitting-free}.
\end{proof}

\begin{lemma}
\label{lemma-fitting-ideal-finite-locally-free}
Soit $R$ un anneau. Soit $M$ un $R$-module de type fini. Soit $r \geq 0$.
Les conditions suivantes sont équivalentes :
\begin{enumerate}
\item $M$ est localement libre de type fini de rang $r$
(Algèbre, définition \ref{algebra-definition-locally-free}) ;
\item $\text{Fit}_{r - 1}(M) = 0$ et $\text{Fit}_r(M) = R$ ;
\item $\text{Fit}_k(M) = 0$ pour $k < r$ et $\text{Fit}_k(M) = R$
pour $k \geq r$.
\end{enumerate}
\end{lemma}

\begin{proof}
Il est immédiat que (2) est équivalente à (3), car les idéaux de Fitting
forment une suite croissante d'idéaux.
Comme la formation de $\text{Fit}_k(M)$ commute au changement de base
(lemme \ref{lemma-fitting-ideal-basics}), (1) implique (2) d'après
l'exemple \ref{example-fitting-free} et les résultats de recollement
(Algèbre, section \ref{algebra-section-more-glueing}).
Réciproquement, supposons (2). D'après le
lemme \ref{lemma-fitting-ideal-generate-locally}, nous pouvons supposer que
$M$ est engendré par $r$ éléments. Nous avons donc une présentation
$\bigoplus_{j \in J} R \to R^{\oplus r} \to M \to 0$.
Mais l'hypothèse $\text{Fit}_{r - 1}(M) = 0$ implique alors que toutes les
entrées de la matrice de l'application
$\bigoplus_{j \in J} R \to R^{\oplus r}$ sont nulles.
Ainsi, $M$ est libre.
\end{proof}

\begin{lemma}
\label{lemma-principal-fitting-ideal}
Soit $R$ un anneau local. Soit $M$ un $R$-module de type fini. Soit $k \geq 0$.
Supposons que $\text{Fit}_k(M) = (f)$ pour un certain $f \in R$.
Soit $M'$ le quotient de $M$ par $\{x \in M \mid fx = 0\}$. Alors
$M'$ peut être engendré par $k$ éléments.
\end{lemma}

\begin{proof}
Choisissons des générateurs $x_1, \ldots, x_n \in M$ correspondant à la
surjection $R^{\oplus n} \to M$. Comme $R$ est local, si une famille
d'éléments $E \subset (f)$ engendre $(f)$, alors un élément $e \in E$ engendre
$(f)$, voir Algèbre, lemme \ref{algebra-lemma-NAK}. Nous pouvons donc choisir
$z_1, \ldots, z_{n - k}$ dans le noyau de $R^{\oplus n} \to M$ de sorte
qu'un mineur $(n - k) \times (n - k)$ de la matrice $n \times (n - k)$
$A = (z_{ij})$ engendre $(f)$. Après renumérotation des $x_i$, nous pouvons
supposer que le premier mineur $\det(z_{ij})_{1 \leq i, j \leq n - k}$
engendre $(f)$, c'est-à-dire que $\det(z_{ij})_{1 \leq i, j \leq n - k} = uf$
pour une unité $u \in R$. Tout autre mineur est un multiple de $f$.
D'après Algèbre, lemme \ref{algebra-lemma-matrix-right-inverse}, il existe une
matrice $n - k \times n - k$ $B$ telle que
$$
AB = f
\left(
\begin{matrix}
u 1_{n - k \times n - k} \\
C
\end{matrix}
\right)
$$
pour une matrice $C$ à coefficients dans $R$. Cela implique que, pour tout
$i \leq n - k$, l'élément $y_i = ux_i + \sum_j c_{ji}x_j$ est annulé par $f$.
Comme $M/\sum Ry_i$ est engendré par les images de
$x_{n - k + 1}, \ldots, x_n$, le résultat est démontré.
\end{proof}

\begin{lemma}
\label{lemma-fitting-ideals-and-pd1}
Soit $R$ un anneau. Soit $M$ un $R$-module de présentation finie. Soit
$k \geq 0$. Supposons que $\text{Fit}_k(M) = (f)$ pour un élément $f \in R$
qui n'est pas diviseur de zéro et que $\text{Fit}_{k - 1}(M) = 0$. Alors
\begin{enumerate}
\item $M$ est de dimension projective $\leq 1$ ;
\item $M' = \Ker(f : M \to M)$ est le sous-module de torsion $f$-primaire
de $M$ ;
\item $M'$ est de dimension projective $\leq 1$ ;
\item $M/M'$ est localement libre de type fini de rang $k$ ; et
\item $M \cong M/M' \oplus M'$.
\end{enumerate}
\end{lemma}

\begin{proof}
Choisissons une présentation
$$
R^{\oplus m} \xrightarrow{A} R^{\oplus n} \to M \to 0
$$
pour une matrice $A$ à coefficients dans $R$.

\medskip\noindent
Nous démontrons d'abord le lemme lorsque $R$ est local.
Posons $M' = \{x \in M \mid fx = 0\}$ comme dans l'énoncé. D'après le
lemme \ref{lemma-principal-fitting-ideal}, nous pouvons choisir
$x_1, \ldots, x_k \in M$ qui engendrent $M/M'$. Alors
$x_1, \ldots, x_k$ engendrent $M_f = (M/M')_f$. Ainsi, s'il existe une
relation $\sum a_ix_i = 0$ dans $M$, nous voyons que $a_1, \ldots, a_k$
s'envoient sur zéro dans $R_f$, car sinon
$\text{Fit}_{k - 1}(M) R_f = \text{Fit}_{k - 1}(M_f)$ serait non nul.
Comme $f$ n'est pas diviseur de zéro, nous concluons
$a_1 = \ldots = a_k = 0$. Ainsi $M \cong R^{\oplus k} \oplus M'$.
Après un changement de base dans la présentation ci-dessus, nous pouvons
supposer que les $n - k$ premiers vecteurs de base de $R^{\oplus n}$ s'envoient
dans le facteur $M'$ de $M$ et que les $k$ derniers vecteurs de base de
$R^{\oplus n}$ s'envoient sur des éléments de base du facteur
$R^{\oplus k}$ de $M$. Les $k$ dernières lignes de la matrice $A$ sont alors
nulles. En remplaçant $M$ par $M'$, $k$ par $0$, $n$ par $n - k$ et $A$ par la
sous-matrice obtenue en supprimant les $k$ dernières lignes, nous nous
ramenons au cas traité au paragraphe suivant.

\medskip\noindent
Supposons $R$ local, $k = 0$ et $M$ annulé par $f$.
Alors le $0$-ième idéal de Fitting de $M$ est $(f)$ et est engendré par les
mineurs $n \times n$ de la matrice $A$ de taille $n \times m$.
(Cela implique en particulier $m \geq n$.) Comme $R$ est local, un mineur
$n \times n$ de $A$ vaut $uf$ pour une unité $u \in R$. Après renumérotation,
nous pouvons supposer qu'il s'agit du premier. De plus, tous les autres
mineurs $n \times n$ de $A$ sont divisibles par $f$. Écrivons $A = (A_1 A_2)$
par blocs, où $A_1$ est une matrice $n \times n$ et $A_2$ une matrice
$n \times (m - n)$. D'après Algèbre, lemme
\ref{algebra-lemma-matrix-right-inverse}, appliqué à la transposée de $A$ (!),
il existe une matrice $n \times n$ $B$ telle que
$$
BA = B(A_1 A_2) = f
\left(
\begin{matrix}
u 1_{n \times n} & C
\end{matrix}
\right)
$$
pour une matrice $n \times (m - n)$ $C$ à coefficients dans $R$.
Nous concluons d'abord $BA_1 = fu 1_{n \times n}$. Ainsi
$$
BA_2 = fC = u^{-1}fuC = u^{-1}BA_1C
$$
Comme le déterminant de $B$ n'est pas diviseur de zéro, nous concluons
$A_2 = u^{-1}A_1C$. L'image de $A$ est donc égale à celle de $A_1$,
laquelle est isomorphe à $R^{\oplus n}$ puisque le déterminant de $A_1$ n'est
pas diviseur de zéro. Ainsi, $M$ est de dimension projective $\leq 1$.

\medskip\noindent
Revenons au cas d'un anneau général $R$. Par le cas local, nous voyons que
$M/M'$ est localement libre de type fini de rang $k$, voir Algèbre, lemme
\ref{algebra-lemma-finite-projective}. L'extension
$0 \to M' \to M \to M/M' \to 0$ est donc scindée. Il s'ensuit que $M'$ est
un module de présentation finie. Choisissons une suite exacte courte
$0 \to K \to R^{\oplus a} \to M' \to 0$. Alors $K$ est un $R$-module de type
fini, voir Algèbre, lemme \ref{algebra-lemma-extension}.
Par le cas local, nous voyons que $K_\mathfrak p \cong R_\mathfrak p^{\oplus a}$
pour tout idéal premier. Ainsi, à nouveau par Algèbre, lemme
\ref{algebra-lemma-finite-projective}, $K$ est localement libre de type fini
de rang $a$. Il s'ensuit que $M'$ est de dimension projective $\leq 1$ et le
lemme est démontré.
\end{proof}










\section{Relèvements}
\label{section-lifting}

\noindent
Dans cette section, nous rassemblons quelques lemmes concernant des
relèvements d'énoncés de la forme suivante : si $A$ est un anneau et
$I \subset A$ un idéal, et si $\overline{\xi}$ est une structure d'un certain
type sur $A/I$, alors nous pouvons relever $\overline{\xi}$ en une structure
du même type $\xi$ sur $A$ ou sur une extension \'etale de $A$. Voici quelques
types de structures pour lesquels nous avons déjà démontré des résultats :
\begin{enumerate}
\item les idempotents, voir
Algèbre, lemmes \ref{algebra-lemma-lift-idempotents} et
\ref{algebra-lemma-lift-idempotents-noncommutative},
\item les modules projectifs, voir
Algèbre, lemmes \ref{algebra-lemma-lift-projective-module} et
\ref{algebra-lemma-lift-finite-projective-module},
\item les modules stablement libres de type fini, voir le lemme
\ref{lemma-lift-stably-free} ;
\item les éléments d'une base, voir
Algèbre, lemmes \ref{algebra-lemma-local-artinian-basis-when-flat} et
\ref{algebra-lemma-lift-basis},
\item les homomorphismes d'anneaux, c'est-à-dire la preuve que certaines
algèbres sont formellement lisses, voir
Algèbre, lemme \ref{algebra-lemma-polynomial-ring-formally-smooth},
proposition \ref{algebra-proposition-smooth-formally-smooth} et
lemme \ref{algebra-lemma-smooth-strong-lift},
\item les homomorphismes d'anneaux syntomiques, voir
Algèbre, lemme \ref{algebra-lemma-lift-syntomic},
\item les homomorphismes d'anneaux lisses, voir
Algèbre, lemme \ref{algebra-lemma-lift-smooth},
\item les homomorphismes d'anneaux \'etales, voir
Algèbre, lemme \ref{algebra-lemma-lift-etale},
\item la factorisation des polynômes, voir
Algèbre, lemme \ref{algebra-lemma-factor-mod-lift-etale}, et
\item Algèbre, section \ref{algebra-section-henselian}, traite des anneaux
locaux henséliens.
\end{enumerate}
Le lecteur intéressé trouvera d'autres résultats de cette nature dans
Compléments sur les homomorphismes d'anneaux, section
\ref{smoothing-section-presentations}, en particulier la proposition
\ref{smoothing-proposition-lift-smooth}.

\medskip\noindent
Soit $A$ un anneau et soit $I \subset A$ un idéal. Soit $\overline{\xi}$
une structure d'un certain type sur $A/I$. Dans les lemmes suivants, nous
cherchons des homomorphismes d'anneaux \'etales $A \to A'$ qui induisent des
isomorphismes $A/I \to A'/IA'$ et des objets $\xi'$ sur $A'$ relevant
$\overline{\xi}$. Remarquons en général que, pour des homomorphismes d'anneaux
\'etales $A \to A' \to A''$ tels que
$A/I \cong A'/IA'$ et $A'/IA' \cong A''/IA''$, alors la composée
$A \to A''$ est également \'etale (Algèbre, lemme \ref{algebra-lemma-etale})
et satisfait aussi $A/I \cong A''/IA''$. Nous utiliserons fréquemment ce fait dans les lemmes
suivants sans autre mention. Voici un exemple trivial du type de résultat
recherché.

\begin{lemma}
\label{lemma-lift-invertible-element}
Soit $A$ un anneau, soit $I \subset A$ un idéal et soit $\overline{u} \in A/I$
un élément inversible. Il existe un homomorphisme d'anneaux \'etale $A \to A'$
qui induit un isomorphisme $A/I \to A'/IA'$ et un élément inversible
$u' \in A'$ relevant $\overline{u}$.
\end{lemma}

\begin{proof}
Choisissons un relèvement quelconque $f \in A$ de $\overline{u}$ et posons
$A' = A_f$ et $u$ égal à l'image de $f$ dans $A'$.
\end{proof}

\begin{lemma}
\label{lemma-lift-idempotent}
Soit $A$ un anneau, soit $I \subset A$ un idéal et soit $\overline{e} \in A/I$
un idempotent. Il existe un homomorphisme d'anneaux \'etale $A \to A'$ qui
induit un isomorphisme $A/I \to A'/IA'$ et un idempotent $e' \in A'$ relevant
$\overline{e}$.
\end{lemma}

\begin{proof}
Choisissons un relèvement quelconque $x \in A$ de $\overline{e}$. Posons
$$
A' = A[t]/(t^2 - t)\left[\frac{1}{t - 1 + x}\right].
$$
L'homomorphisme d'anneaux $A \to A'$ est \'etale car
$(2t - 1)\text{d}t = 0$ et $(2t - 1)(2t - 1) = 1$, qui est inversible. Nous avons
$A'/IA' = A/I[t]/(t^2 - t)[\frac{1}{t - 1 + \overline{e}}] \cong A/I$
la dernière application envoyant $t$ sur $\overline{e}$, ce qui convient car
$\overline{e}$ est une racine de $t^2 - t$. Cela montre aussi que prendre
$e'$ égal à la classe de $t$ dans $A'$ convient.
\end{proof}

\begin{lemma}
\label{lemma-lift-open-covering}
Soit $A$ un anneau, soit $I \subset A$ un idéal. Soit
$\Spec(A/I) = \coprod_{j \in J} \overline{U}_j$ un recouvrement ouvert fini
et disjoint. Il existe alors un homomorphisme d'anneaux \'etale $A \to A'$ qui
induit un isomorphisme $A/I \to A'/IA'$ et un recouvrement ouvert fini et
disjoint $\Spec(A') = \coprod_{j \in J} U'_j$ relevant le recouvrement donné.
\end{lemma}

\begin{proof}
Cela résulte du lemme \ref{lemma-lift-idempotent} et du fait que les parties
ouvertes et fermées des spectres correspondent aux idempotents, voir Algèbre,
lemme \ref{algebra-lemma-disjoint-decomposition}.
\end{proof}

\begin{lemma}
\label{lemma-localize-upstairs}
Soit $A \to B$ un homomorphisme d'anneaux et soit $J \subset B$ un idéal.
Si $A \to B$ est \'etale en tout idéal premier de $V(J)$, il existe alors
$g \in B$ dont l'image dans $B/J$ est inversible, tel que $A' = B_g$ soit
\'etale sur $A$.
\end{lemma}

\begin{proof}
L'ensemble des points de $\Spec(B)$ où $A \to B$ n'est pas \'etale est une
partie fermée de $\Spec(B)$, voir Algèbre, définition
\ref{algebra-definition-etale}. Écrivons-le $V(J')$ pour un idéal
$J' \subset B$. Alors $V(J') \cap V(J) = \emptyset$, d'où $J + J' = B$
d'après Algèbre, lemme \ref{algebra-lemma-Zariski-topology}.
Écrivons $1 = f + g$ avec $f \in J$ et $g \in J'$. Alors $g$ convient.
\end{proof}

\noindent
Nous avons ensuite trois lemmes affirmant que l'on peut relever des
factorisations de polynômes.

\begin{lemma}
\label{lemma-lift-factorization-monic}
Soit $A$ un anneau, soit $I \subset A$ un idéal et soit $f \in A[x]$ un
polynôme unitaire. Soit $\overline{f} = \overline{g} \overline{h}$ une
factorisation de $f$ dans $A/I[x]$ telle que $\overline{g}$ et $\overline{h}$
soient unitaires et engendrent l'idéal unité de $A/I[x]$. Il existe alors un
homomorphisme d'anneaux \'etale $A \to A'$ qui induit un isomorphisme
$A/I \to A'/IA'$ et une factorisation $f = g' h'$ dans $A'[x]$, avec $g'$ et
$h'$ unitaires relevant la factorisation donnée sur $A/I$.
\end{lemma}

\begin{proof}
Nous déduirons ceci des résultats sur la factorisation universelle démontrés
précédemment ; toutefois, nous encourageons le lecteur à trouver sa propre
preuve sans ce procédé. Posons $\deg(\overline{g}) = n$ et
$\deg(\overline{h}) = m$, de sorte que $\deg(f) = n + m$. Écrivons
$f = x^{n + m} + \sum \alpha_i x^{n + m - i}$ pour certains
$\alpha_1, \ldots, \alpha_{n + m} \in A$. Considérons l'homomorphisme
d'anneaux
$$
R = \mathbf{Z}[a_1, \ldots, a_{n + m}]
\longrightarrow
S = \mathbf{Z}[b_1, \ldots, b_n, c_1, \ldots, c_m]
$$
de Algèbre, exemple \ref{algebra-example-factor-polynomials-etale}.
Soit $R \to A$ l'homomorphisme d'anneaux qui envoie $a_i$ sur $\alpha_i$.
Posons
$$
B = A \otimes_R S
$$
Par construction, l'image $f_B$ de $f$ dans $B[x]$ se factorise, disons
$f_B = g_B h_B$ avec $g_B = x^n + \sum (1 \otimes b_i) x^{n - i}$
et de même pour $h_B$. Écrivons
$\overline{g} = x^n + \sum \overline{\beta}_i x^{n - i}$ et
$\overline{h} = x^m + \sum \overline{\gamma}_i x^{m - i}$.
L'homomorphisme de $A$-algèbres
$$
B \longrightarrow A/I, \quad
1 \otimes b_i \mapsto \overline{\beta}_i, \quad
1 \otimes c_i \mapsto \overline{\gamma}_i
$$
envoie $g_B$ et $h_B$ sur $\overline{g}$ et $\overline{h}$ dans $A/I[x]$.
L'application affichée est surjective ; notons $J \subset B$ son noyau.
D'après la discussion de l'algèbre, exemple
\ref{algebra-example-factor-polynomials-etale}, il est clair que $A \to B$
est \'etale en tout point de $V(J) \subset \Spec(B)$. Choisissons $g \in B$
comme dans le lemme \ref{lemma-localize-upstairs} et considérons la
$A$-algèbre $B_g$. Comme $g$ s'envoie sur une unité dans $B/J = A/I$, nous
obtenons aussi une application $B_g/I B_g \to A/I$ de $A/I$-algèbres.
Puisque $A/I \to B_g/I B_g$ est \'etale, il en va de même de
$B_g/IB_g \to A/I$ (Algèbre, lemme \ref{algebra-lemma-map-between-etale}).
Il existe donc un idempotent $e \in B_g/I B_g$ tel que
$A/I = (B_g/I B_g)_e$ (Algèbre, lemme
\ref{algebra-lemma-surjective-flat-finitely-presented}). Choisissons un
relèvement $h \in B_g$ de $e$. Alors $A \to A' = (B_g)_h$, muni de la
factorisation donnée par l'image de $f_B = g_B h_B$ dans $A'$, est une
solution au problème posé par le lemme.
\end{proof}

\begin{lemma}
\label{lemma-lift-factorization-easy}
Soit $A$ un anneau, soit $I \subset A$ un idéal et soit $f \in A[x]$ un
polynôme unitaire. Soit $\overline{f} = \overline{g} \overline{h}$ une
factorisation de $f$ dans $A/I[x]$ et supposons que
\begin{enumerate}
\item le coefficient dominant de $\overline{g}$ soit un élément inversible
de $A/I$ ;
\item $\overline{g}$ et $\overline{h}$ engendrent l'idéal unité de $A/I[x]$.
\end{enumerate}
Alors il existe un homomorphisme d'anneaux \'etale $A \to A'$ qui induit un
isomorphisme $A/I \to A'/IA'$ et une factorisation $f = g' h'$ dans $A'[x]$
relevant la factorisation donnée sur $A/I$.
\end{lemma}

\begin{proof}
En appliquant le lemme \ref{lemma-lift-invertible-element}, nous pouvons
supposer que le coefficient dominant de $\overline{g}$ est la réduction d'un
élément inversible $u \in A$. Nous pouvons alors remplacer $\overline{g}$ par
$\overline{u}^{-1}\overline{g}$ et $\overline{h}$ par
$\overline{u}\overline{h}$. Nous pouvons donc supposer que $\overline{g}$
est unitaire. Puisque $f$ est unitaire, nous concluons que $\overline{h}$
l'est aussi. Dans ce cas, le résultat découle du lemme
\ref{lemma-lift-factorization-monic}.
\end{proof}

\begin{example}
\label{example-cannot-lift-factorization}
Cet exemple montre qu'on ne peut pas supprimer la condition sur le coefficient
dominant dans le lemme \ref{lemma-lift-factorization-easy}. En effet,
prenons $A = \mathbf{Z}$, $I = 4\mathbf{Z}$, $f = 1$ et
$\bar{g} = \bar{h} = 2x^2 + 2x + 1$ dans $A/I[x]$. Si $A \to A'$ est \'etale
et $A'/I A' = A/I$, alors le complété $2$-adique de $A'$ est isomorphe
à $\mathbf{Z}_2$. Or aucun polynôme $g'$ de $\mathbf{Z}_2[x]$ congru à
$2x^2 + 2x + 1$ modulo $4$ n'a d'inverse multiplicatif polynomial dans
$\mathbf{Z}_2[x]$, d'où la conclusion du lemme qui ne peut être vraie.
\end{example}

\begin{lemma}
\label{lemma-separate-image-closed-from-closed}
Soit $R \to S$ un homomorphisme d'anneaux. Soit $I \subset R$ un idéal de
$R$ et soit $J \subset S$ un idéal de $S$. Si l'adhérence de l'image de
$V(J)$ dans $\Spec(R)$ est disjointe de $V(I)$, il existe un élément
$f \in R$ dont l'image vaut $1$ dans $R/I$ et est un élément de $J$ dans $S$.
\end{lemma}

\begin{proof}
Soit $I' \subset R$ un idéal tel que $V(I')$ soit l'adhérence de l'image de
$V(J)$. Alors $V(I) \cap V(I') = \emptyset$ par hypothèse et donc
$I + I' = R$ d'après Algèbre, lemme \ref{algebra-lemma-Zariski-topology}.
Écrivons $1 = g + f$ avec $g \in I$ et $f \in I'$. Nous avons
$V(f') \supset V(J)$, où $f'$ est l'image de $f$ dans $S$. Ainsi,
$(f')^n \in J$ pour un certain $n$, voir Algèbre, lemme
\ref{algebra-lemma-Zariski-topology}. Remplaçons $f$ par $f^n$ pour conclure.
\end{proof}

\begin{lemma}
\label{lemma-helper-integral}
Soit $I$ un idéal d'un anneau $A$. Soit $A \to B$ un homomorphisme d'anneaux
entier. Soit $b \in B$ dont l'image dans $B/IB$ est un idempotent. Il existe
alors un polynôme unitaire $f \in A[x]$ tel que $f(b) = 0$ et
$f \bmod I = x^d(x - 1)^d$ pour un certain $d \geq 1$.
\end{lemma}

\begin{proof}
Remarquons que $z = b^2 - b$ est un élément de $IB$. D'après Algèbre, lemme
\ref{algebra-lemma-integral-integral-over-ideal}, il existe un polynôme unitaire
$g(x) = x^d + \sum a_j x^j$ de degré $d$, avec $a_j \in I$, tel que
$g(z) = 0$ dans $B$. Ainsi, $f(x) = g(x^2 - x) \in A[x]$ est un polynôme
unitaire tel que $f(x) \equiv x^d(x - 1)^d \bmod I$ et tel que $f(b) = 0$
dans $B$.
\end{proof}

\begin{lemma}
\label{lemma-lift-idempotent-upstairs}
Soit $A$ un anneau, soit $I \subset A$ un idéal et soit $A \to B$ un
homomorphisme d'anneaux entier. Soit $\overline{e} \in B/IB$ un idempotent.
Il existe alors un homomorphisme d'anneaux \'etale $A \to A'$ qui induit un
isomorphisme $A/I \to A'/IA'$ et un idempotent
$e' \in B \otimes_A A'$ relevant $\overline{e}$.
\end{lemma}

\begin{proof}
Choisissons un élément $y \in B$ relevant $\overline{e}$.
Choisissons $f \in A[x]$ comme dans le lemme \ref{lemma-helper-integral} pour
$y$. D'après le lemme \ref{lemma-lift-factorization-easy}, nous pouvons
trouver un homomorphisme d'anneaux \'etale $A \to A'$ qui induit un
isomorphisme $A/I \to A'/IA'$ et tel que $f = gh$ dans $A[x]$, avec
$g(x) = x^d \bmod IA'$ et $h(x) = (x - 1)^d \bmod IA'$.
Après remplacement de $A$ par $A'$, nous pouvons supposer que la
factorisation est définie sur $A$. Dans ce cas, nous voyons que
$b_1 = g(y) \in B$ est un relèvement de $\overline{e}^d = \overline{e}$ et
$b_2 = h(y) \in B$ est un relèvement de
$(\overline{e} - 1)^d = (-1)^d (1 - \overline{e})^d = (-1)^d(1 - \overline{e})$
et, de plus, $b_1b_2 = 0$. Ainsi $(b_1, b_2)B/IB = B/IB$ et
$V(b_1, b_2) \subset \Spec(B)$ est disjoint de $V(IB)$. Comme
$\Spec(B) \to \Spec(A)$ est fermé (voir les lemmes d'Algèbre
\ref{algebra-lemma-integral-going-up} et \ref{algebra-lemma-going-up-closed}),
nous pouvons trouver $a \in A$ dont l'image dans $A/I$ est inversible et dont
l'image dans $B$ appartient à $(b_1, b_2)$, voir le lemme
\ref{lemma-separate-image-closed-from-closed}.
Après remplacement de $A$ par la localisation $A_a$, nous avons
$(b_1, b_2) = B$. Alors $\Spec(B) = D(b_1) \amalg D(b_2)$ ; cette union est
disjointe car $b_1b_2 = 0$ et couvre $\Spec(B)$ puisque $(b_1, b_2) = B$.
Soit $e \in B$ l'idempotent correspondant à la partie ouverte et fermée
$D(b_1)$, voir Algèbre, lemme \ref{algebra-lemma-disjoint-decomposition}.
Puisque $b_1$ relève $\overline{e}$ et $b_2$ relève
$\pm (1 - \overline{e})$, nous concluons que $e$ relève $\overline{e}$ par
l'assertion d'unicité du lemme d'Algèbre
\ref{algebra-lemma-disjoint-decomposition}.
\end{proof}

\begin{lemma}
\label{lemma-lift-projective-module}
Soit $A$ un anneau, soit $I \subset A$ un idéal.
Soit $\overline{P}$ un $A/I$-module projectif de type fini.
Il existe alors un homomorphisme d'anneaux \'etale $A \to A'$ qui induit
un isomorphisme $A/I \to A'/IA'$ et un $A'$-module projectif de type fini
$P'$ relevant $\overline{P}$.
\end{lemma}

\begin{proof}
Nous pouvons choisir un entier $n$ et une décomposition en somme directe
$(A/I)^{\oplus n} = \overline{P} \oplus \overline{K}$
pour un certain $R/I$-module $\overline{K}$. Choisissons un relèvement
$\varphi : A^{\oplus n} \to A^{\oplus n}$ du projecteur $\overline{p}$
associé au facteur direct $\overline{P}$.
Soit $f \in A[x]$ le polynôme caractéristique de $\varphi$.
Posons $B = A[x]/(f)$. D'après Cayley-Hamilton
(Algèbre, lemme \ref{algebra-lemma-charpoly}), il existe une application
$B \to \text{End}_A(A^{\oplus n})$ qui envoie $x$ sur $\varphi$.
Pour tout idéal premier $\mathfrak p \supset I$, l'image de $f$ dans
$\kappa(\mathfrak p)$ est $(x - 1)^rx^{n - r}$, où $r$ est la dimension de
$\overline{P} \otimes_{A/I} \kappa(\mathfrak p)$.
Ainsi $(x - 1)^nx^n$ s'envoie sur zéro dans
$B \otimes_A \kappa(\mathfrak p)$ pour tout $\mathfrak p \supset I$.
Donc $x(1 - x)$ appartient à tout idéal premier de $B/IB$. Par conséquent,
$x^N(1 - x)^N$ appartient à $IB$ pour un certain $N \geq 1$.
Il s'ensuit que $x^N + (1 - x)^N$ est une unité dans $B/IB$ et que
$$
\overline{e} = \text{image de }\frac{x^N}{x^N + (1 - x)^N}\text{ dans }B/IB
$$
est un idempotent, car les deux assertions valent dans
$\mathbf{Z}[x]/(x^N(x - 1)^N)$. L'image de $\overline{e}$ dans
$\text{End}_{A/I}((A/I)^{\oplus n})$ est
$$
\frac{\overline{p}^N}{\overline{p}^N + (1 - \overline{p})^N} = \overline{p}
$$
car $\overline{p}$ est un idempotent. Après remplacement de $A$ par une
extension \'etale $A'$ comme dans le lemme, nous pouvons supposer qu'il existe
un idempotent $e \in B$ qui s'envoie sur $\overline{e}$ dans $B/IB$, voir le
lemme \ref{lemma-lift-idempotent-upstairs}. Alors l'image de $e$ par
l'application
$$
B = A[x]/(f) \longrightarrow \text{End}_A(A^{\oplus n}).
$$
est un élément idempotent $p$ qui relève $\overline{p}$. En posant
$P = \Im(p)$, le résultat est démontré.
\end{proof}

\begin{lemma}
\label{lemma-cotangent-complex-symmetric-algebra}
Soit $A$ un anneau. Soit $0 \to K \to A^{\oplus m} \to M \to 0$
une suite de $A$-modules. Considérons la $A$-algèbre
$C = \text{Sym}^*_A(M)$ munie de sa présentation
$\alpha : A[y_1, \ldots, y_m] \to C$
issue de la surjection $A^{\oplus m} \to M$. Alors
$$
\NL(\alpha) =
(K \otimes_A C \to \bigoplus\nolimits_{j = 1, \ldots, m} C \text{d}y_j)
$$
(voir Algèbre, section \ref{algebra-section-netherlander}),
en particulier $\Omega_{C/A} = M \otimes_A C$.
\end{lemma}

\begin{proof}
Soit $J = \Ker(\alpha)$. Le lemme affirme que
$J/J^2 \cong K \otimes_A C$. Remarquons que $\alpha$ est un homomorphisme
d'algèbres graduées. Nous montrerons qu'en degré $d$ nous avons
$(J/J^2)_d = K \otimes_A C_{d - 1}$. Remarquons que
$$
J_d = \Ker(\text{Sym}^d_A(A^{\oplus m}) \to \text{Sym}^d_A(M))
= \Im(K \otimes_A \text{Sym}^{d - 1}_A(A^{\oplus m})
\to \text{Sym}^d_A(A^{\oplus m})),
$$
voir Algèbre, lemme \ref{algebra-lemma-presentation-sym-exterior}.
Il s'ensuit que $(J^2)_d = \sum_{a + b = d} J_a \cdot J_b$ est l'image de
$$
K \otimes_A K \otimes_A \text{Sym}^{d - 2}_A(A^{\otimes m})
\to \text{Sym}^d_A(A^{\oplus m}).
$$
Le conoyau de l'application
$K \otimes_A \text{Sym}^{d - 2}_A(A^{\otimes m}) \to
\text{Sym}^{d - 1}_A(A^{\oplus m})$ est $\text{Sym}^{d - 1}_A(M)$ d'après
le lemme cité ci-dessus. Il est donc clair que
$(J/J^2)_d = J_d/(J^2)_d$ est égal à
\begin{align*}
\Coker(
K \otimes_A K \otimes_A \text{Sym}^{d - 2}_A(A^{\otimes m})
\to K \otimes_A \text{Sym}^{d - 1}_A(A^{\otimes m}))
& = K \otimes_A \text{Sym}^{d - 1}_A(M) \\
& = K \otimes_A C_{d -1}
\end{align*}
comme voulu.
\end{proof}

\begin{lemma}
\label{lemma-symmetric-algebra-smooth}
Soit $A$ un anneau. Soit $M$ un $A$-module. Alors $C = \text{Sym}_A^*(M)$
est lisse sur $A$ si et seulement si $M$ est un $A$-module projectif de type
fini.
\end{lemma}

\begin{proof}
Soit $\sigma : C \to A$ la projection sur la partie de degré $0$ de $C$.
Alors $J = \Ker(\sigma)$ est la partie de degré $> 0$ et nous voyons que
$J/J^2 = M$ comme $A$-module. Ainsi, si $A \to C$ est lisse, alors $M$ est
un $A$-module projectif de type fini d'après Algèbre, lemme
\ref{algebra-lemma-section-smooth}.

\medskip\noindent
Réciproquement, supposons que $M$ soit projectif de type fini et choisissons
une surjection $A^{\oplus n} \to M$ de noyau $K$. La suite
$0 \to K \to A^{\oplus n} \to M \to 0$ est scindée puisque $M$ est projectif.
En particulier, $K$ est un $A$-module de type fini et donc $C$ est de
présentation finie sur $A$, car $C$ est un quotient de
$A[x_1, \ldots, x_n]$ par l'idéal engendré par
$K \subset \bigoplus Ax_i$. Le calcul du lemme
\ref{lemma-cotangent-complex-symmetric-algebra} montre que
$\NL_{C/A}$ est homotopiquement équivalent à $(K \to M) \otimes_A C$.
Ainsi, $\NL_{C/A}$ est quasi-isomorphe à $C \otimes_A M$ placé en degré $0$,
ce qui signifie que $C$ est lisse sur $A$ d'après Algèbre, définition
\ref{algebra-definition-smooth}.
\end{proof}

\begin{lemma}
\label{lemma-lift-section-smooth-morphism}
Soit $A$ un anneau, soit $I \subset A$ un idéal. Considérons un diagramme
commutatif
$$
\xymatrix{
B \ar[rd] \\
A \ar[u] \ar[r] & A/I
}
$$
où $B$ est une $A$-algèbre lisse. Alors il existe un morphisme d'anneaux
\'etale $A \to A'$ qui induit un isomorphisme $A/I \to A'/IA'$ et un
morphisme de $A$-algèbres $B \to A'$ relevant le morphisme d'anneaux
$B \to A/I$.
\end{lemma}

\begin{proof}
Soit $J \subset B$ le noyau de $B \to A/I$, de sorte que $B/J = A/I$.
D'après le lemme d'algèbre \ref{algebra-lemma-application-NL-smooth}, la suite
$$
0 \to I/I^2 \to J/J^2 \to \Omega_{B/A} \otimes_B B/J \to 0
$$
est exacte scindée. Ainsi $\overline{P} = J/(J^2 + IB) = \Omega_{B/A} \otimes_B B/J$
est un $A/I$-module projectif de type fini. Choisissons un entier $n$ et une
décomposition en somme directe $A/I^{\oplus n} = \overline{P} \oplus \overline{K}$.
D'après le lemme \ref{lemma-lift-projective-module}, on peut trouver un
morphisme d'anneaux \'etale $A \to A'$ qui induit un isomorphisme
$A/I \to A'/IA'$ et un $A$-module projectif de type fini $K$ qui
relève $\overline{K}$. Nous remplaçons désormais $A$ par $A'$.
Posons $B' = B \otimes_A \text{Sym}_A^*(K)$. Comme $A \to \text{Sym}_A^*(K)$
est lisse d'après le lemme \ref{lemma-symmetric-algebra-smooth}, on voit que
$B \to B'$ est lisse, ce qui implique à son tour que $A \to B'$ est lisse (voir
les lemmes d'algèbre \ref{algebra-lemma-base-change-smooth} et
\ref{algebra-lemma-locally-smooth}).
De plus, la section $\text{Sym}^*_A(K) \to A$ détermine une section
$B' \to B$ et nous définissons $B' \to A/I$ comme la composée $B' \to B \to A/I$.
Soit $J' \subset B'$ le noyau de $B' \to A/I$. Nous avons
$JB' \subset J'$ et $B \otimes_A K \subset J'$. Ces morphismes se combinent
pour donner un isomorphisme
$$
(A/I)^{\oplus n} \cong J/(J^2 + IB) \oplus \overline{K}
\longrightarrow
J'/((J')^2 + IB')
$$
Ainsi, après avoir remplacé $B$ par $B'$, nous pouvons supposer que
$J/(J^2 + IB) = \Omega_{B/A} \otimes_B B/J$ est un
$A/I$-module libre de rang $n$.

\medskip\noindent
Dans ce cas, choisissons $f_1, \ldots, f_n \in J$ dont les images forment une
base de $J/(J^2 + IB)$. Considérons la $A$-algèbre de présentation finie
$C = B/(f_1, \ldots, f_n)$. Remarquons que nous avons une suite exacte
$$
0 \to H_1(L_{C/A}) \to (f_1, \ldots, f_n)/(f_1, \ldots, f_n)^2
\to \Omega_{B/A} \otimes_B C \to \Omega_{C/A} \to 0
$$
voir le lemme d'algèbre \ref{algebra-lemma-exact-sequence-NL} (remarquons que
$H_1(L_{B/A}) = 0$ et que $\Omega_{B/A}$ est projectif de type fini,
en particulier plat, de sorte que le groupe de Tor s'annule). Pour tout idéal premier
$\mathfrak q \supset J$ de $B$, le module $\Omega_{B/A, \mathfrak q}$
est libre de rang $n$, car $\Omega_{B/A}$ est projectif de type fini et
car $\Omega_{B/A} \otimes_B B/J$ est libre de rang $n$
(voir le lemme d'algèbre \ref{algebra-lemma-finite-projective}). Par notre choix
de $f_1, \ldots, f_n$, le morphisme
$$
\left((f_1, \ldots, f_n)/(f_1, \ldots, f_n)^2\right)_{\mathfrak q}
\to
\Omega_{B/A, \mathfrak q}
$$
est surjectif modulo $J$. Nous en déduisons que ce morphisme de modules sur
l'anneau local $C_{\mathfrak q}$ doit être un isomorphisme
(d'après le lemme d'algèbre \ref{algebra-lemma-NAK}, le morphisme est surjectif,
et, par exemple d'après le lemme d'algèbre \ref{algebra-lemma-fun},
comme $((f_1, \ldots, f_n)/(f_1, \ldots, f_n)^2)_{\mathfrak q}$
est engendré par $n$ éléments, le morphisme est injectif). Ainsi
$H_1(L_{C/A})_{\mathfrak q} = 0$ et $\Omega_{C/A, \mathfrak q} = 0$. D'après
le lemme d'algèbre \ref{algebra-lemma-smooth-at-point},
$A \to C$ est lisse au premier $\overline{\mathfrak q}$
de $C$ correspondant à $\mathfrak q$. Puisque
$\Omega_{C/A, \mathfrak q} = 0$, il est en fait \'etale en
$\overline{\mathfrak q}$. Ainsi $A \to C$ est \'etale en tous les idéaux premiers
de $C$ contenant $JC$. D'après le lemme \ref{lemma-localize-upstairs},
on peut trouver un $f \in C$ dont l'image est un élément inversible de $C/JC$ tel que
$A \to C_f$ soit \'etale. Par notre choix de $f$, on a encore
$C_f/JC_f = A/I$. Le morphisme $C_f/IC_f \to A/I$ est surjectif et
\'etale d'après le lemme d'algèbre \ref{algebra-lemma-map-between-etale}.
Ainsi $A/I$ est isomorphe à la localisation de $C_f/IC_f$ en
un élément $g \in C$, voir
le lemme d'algèbre \ref{algebra-lemma-surjective-flat-finitely-presented}.
Posons $A' = C_{fg}$, ce qui conclut la démonstration.
\end{proof}








\section{Paires de Zariski}
\label{section-zariski-pairs}

\noindent
Dans cette section et la suivante, une {\it paire} est une paire $(A, I)$ où $A$
est un anneau et $I \subset A$ un idéal. Un {\it morphisme de paires}
$(A, I) \to (B, J)$ est un morphisme d'anneaux $\varphi : A \to B$ tel que
$\varphi(I) \subset J$.

\begin{definition}
\label{definition-zariski-pair}
Une {\it paire de Zariski} est une paire $(A, I)$ telle que
$I$ soit contenu dans le radical de Jacobson de $A$.
\end{definition}

\begin{lemma}
\label{lemma-idempotents-determined-modulo-radical}
Soit $(A, I)$ une paire de Zariski. Alors le morphisme des
idempotents de $A$ vers les idempotents de $A/I$ est injectif.
\end{lemma}

\begin{proof}
Un idempotent d'un anneau local vaut soit $0$, soit $1$.
Ainsi, un idempotent est déterminé par l'ensemble des idéaux maximaux
où il s'annule, d'après le lemme d'algèbre
\ref{algebra-lemma-characterize-zero-local}.
\end{proof}

\begin{lemma}
\label{lemma-check-isomorphism-zariski}
Soit $(A, I)$ une paire de Zariski. Soit $A \to B$ un morphisme d'anneaux
plat, entier et de présentation finie tel que $A/I \to B/IB$
soit un isomorphisme. Alors $A \to B$ est un isomorphisme.
\end{lemma}

\begin{proof}
Le morphisme d'anneaux $A \to B$ est fini d'après le lemme d'algèbre
\ref{algebra-lemma-characterize-finite-in-terms-of-integral}.
Ainsi $B$ est de présentation finie comme $A$-module d'après le lemme d'algèbre
\ref{algebra-lemma-finite-finitely-presented-extension}.
Ainsi $B$ est un $A$-module localement libre de type fini d'après le lemme d'algèbre
\ref{algebra-lemma-finite-projective}.
Comme le module $B$ est de rang $1$ sur $V(I)$ (voir la fonction de rang décrite dans
le lemme d'algèbre \ref{algebra-lemma-finite-projective}),
et puisque $(A, I)$ est une paire de Zariski, nous concluons que le
rang vaut $1$ partout.
Il s'ensuit que $A \to B$ est un isomorphisme :
il est facile de vérifier qu'un morphisme d'anneaux $R \to S$ tel que
$S$ soit un $R$-module localement libre de rang $1$ est un isomorphisme
(indication : regarder les anneaux locaux).
\end{proof}

\begin{lemma}
\label{lemma-helper-finite}
Soit $(A, I)$ une paire de Zariski. Soit $A \to B$ un morphisme d'anneaux fini.
Supposons
\begin{enumerate}
\item $B/IB = B_1 \times B_2$ soit un produit de $A/I$-algèbres ;
\item $A/I \to B_1/IB_1$ soit surjectif ;
\item $b \in B$ s'envoie sur $(1, 0)$ dans le produit.
\end{enumerate}
Alors il existe un polynôme unitaire $f \in A[x]$ tel que $f(b) = 0$ et
$f \bmod I = (x - 1)x^d$ pour un certain $d \geq 1$.
\end{lemma}

\begin{proof}
D'après le lemme \ref{lemma-lift-idempotent-upstairs}, on peut trouver un
morphisme d'anneaux \'etale $A \to A'$ induisant un isomorphisme $A/I \to A'/IA'$
tel que $B' = B \otimes_A A'$ contienne
un idempotent $e'$ relevant l'image de $b$ dans $B'/IB'$.
Considérons la décomposition correspondante en $A'$-algèbres
$$
B' = B'_1 \times B'_2
$$
qui est compatible avec celle donnée dans le lemme après
réduction modulo $I$.
Le morphisme $A' \to B'_1$ est surjectif modulo $IA'$. D'après le lemme de Nakayama
(lemme d'algèbre \ref{algebra-lemma-NAK}),
on peut trouver $i \in IA'$ tel qu'après avoir remplacé $A'$ par $A'_{1 + i}$,
le morphisme $A' \to B'_1$ soit surjectif.
Remarquons que l'image $b'_1 \in B'_1$ de $b$
vérifie $b'_1 - 1 \in IB'_1$.
Nous pouvons donc choisir $a' \in IA'$ s'envoyant sur $b'_1 - 1$.
D'autre part, l'image $b'_2 \in B'_2$ de $b$ appartient à $IB'_2$. D'après
le lemme d'algèbre \ref{algebra-lemma-integral-integral-over-ideal},
il existe un polynôme unitaire
$g(x) = x^d + \sum a'_j x^j$ de degré $d$, avec $a'_j \in IA'$, tel que
$g(b'_2) = 0$ dans $B'_2$. Ainsi l'image $b' = (b'_1, b'_2) \in B'$
de $b$ est une racine du polynôme $(x - 1 - a')g(x)$. Nous en déduisons que
$$
(b' - 1)(b')^d \in \sum\nolimits_{j = 0, \ldots, d} IA' \cdot (b')^j
$$
Nous affirmons que cela implique
$$
(b - 1)b^d \in \sum\nolimits_{j = 0, \ldots, d} I \cdot b^j
$$
dans $B$. Pour cela, il suffit de voir que le morphisme d'anneaux
$A \to A'$ est fidèlement plat, car la condition signifie que l'image de
$(b - 1)b^d$ est nulle dans $B/\sum_{j = 0, \ldots, d} Ib^j$
(utiliser le lemme d'algèbre \ref{algebra-lemma-faithfully-flat-universally-injective}).
Le morphisme $A \to A'$ est plat car il est \'etale
(lemme d'algèbre \ref{algebra-lemma-etale}).
D'autre part, le morphisme induit sur les spectres est ouvert
(voir la proposition d'algèbre \ref{algebra-proposition-fppf-open} et utiliser
le lemme précédent cité), et l'image
contient $V(I)$. Puisque $I$ est contenu dans le radical de Jacobson de $A$,
nous concluons.
\end{proof}

\begin{lemma}
\label{lemma-noetherian-zariski-jacobson-complement}
Soit $(A, I)$ une paire de Zariski avec $A$ noethérien. Soit $f \in I$.
Alors $A_f$ est un anneau de Jacobson.
\end{lemma}

\begin{proof}
Nous utiliserons le critère du
lemme d'algèbre \ref{algebra-lemma-noetherian-dim-1-Jacobson}.
Soit $\mathfrak p \subset A$ un idéal premier tel que
$\mathfrak p_f = \mathfrak p A_f$ soit premier et non maximal. Nous devons
montrer que $A_f/\mathfrak p_f = (A/\mathfrak p)_f$
possède une infinité d'idéaux premiers. Après avoir remplacé $A$ par $A/\mathfrak p$,
nous pouvons supposer que $A$ est un domaine, que $\dim A_f > 0$, et notre but est de
montrer que $\Spec(A_f)$ est infini. Puisque $\dim A_f > 0$, nous pouvons
trouver un idéal premier non nul $\mathfrak q \subset A$ qui ne contient pas $f$.
Choisissons un idéal maximal $\mathfrak m \subset A$ contenant $\mathfrak q$.
Comme $(A, I)$ est une paire de Zariski, on a $I \subset \mathfrak m$.
Ainsi $\mathfrak m \not = \mathfrak q$ et $\dim(A_\mathfrak m) > 1$.
Par conséquent $\Spec((A_\mathfrak m)_f) \subset \Spec(A_f)$ est infini d'après
le lemme d'algèbre \ref{algebra-lemma-Noetherian-local-domain-dim-2-infinite-opens},
ce qui conclut.
\end{proof}











\section{Paires henséliennes}
\label{section-henselian-pairs}

\noindent
Certains des résultats de la section \ref{section-lifting} peuvent être vus comme
des résultats sur les paires henséliennes. Dans cette section, une {\it paire} est
une paire $(A, I)$ où $A$ est un anneau et $I \subset A$ un idéal. Un {\it morphisme
de paires} $(A, I) \to (B, J)$ est un morphisme d'anneaux $\varphi : A \to B$ tel que
$\varphi(I) \subset J$. Comme dans
la section \ref{section-lifting}, pour un objet $\xi$ au-dessus de $A$, nous notons
$\overline{\xi}$ le « changement de base » de $\xi$ vers un objet au-dessus de $A/I$
(lorsque cela a un sens).

\begin{definition}
\label{definition-henselian-pair}
Une {\it paire hensélienne} est une paire $(A, I)$ satisfaisant
\begin{enumerate}
\item $I$ est contenu dans le radical de Jacobson de $A$ ;
\item pour tout polynôme unitaire $f \in A[T]$ et toute factorisation
$\overline{f} = g_0h_0$ avec $g_0, h_0 \in A/I[T]$ unitaires
engendrant l'idéal unité de $A/I[T]$, il
existe une factorisation $f = gh$ dans $A[T]$ avec $g, h$ unitaires
et $g_0 = \overline{g}$ et $h_0 = \overline{h}$.
\end{enumerate}
\end{definition}

\noindent
Remarquons que si $A$ est un anneau local et $I = \mathfrak m$ est l'idéal
maximal, alors $(A, I)$ est une paire hensélienne si et seulement si $A$ est un
anneau local hensélien, voir le lemme d'algèbre
\ref{algebra-lemma-characterize-henselian}.
Dans le lemme \ref{lemma-characterize-henselian-pair}, nous donnons plusieurs
caractérisations équivalentes des paires henséliennes (et nous en ajouterons
d'autres au fil du temps).

\begin{lemma}
\label{lemma-locally-nilpotent-henselian}
Soit $(A, I)$ une paire avec $I$ localement nilpotent. Alors le foncteur
$B \mapsto B/IB$ induit une équivalence entre la catégorie des
algèbres \'etale sur $A$ et la catégorie des algèbres \'etale sur $A/I$.
De plus, la paire est hensélienne.
\end{lemma}

\begin{proof}
La surjectivité essentielle résulte du lemme d'algèbre
\ref{algebra-lemma-lift-etale}. Si $B$, $B'$ sont \'etales sur $A$ et si
$B/IB \to B'/IB'$ est un morphisme d'$A/I$-algèbres, nous pouvons le relever
d'après le lemme d'algèbre \ref{algebra-lemma-smooth-strong-lift}.
Enfin, supposons que $f, g : B \to B'$ soient deux morphismes de $A$-algèbres
avec $f \bmod I = g \bmod I$. Choisissons un idempotent $e \in B \otimes_A B$
engendrant le noyau du morphisme de multiplication $B \otimes_A B \to B$,
voir les lemmes d'algèbre \ref{algebra-lemma-diagonal-unramified}
et \ref{algebra-lemma-unramified} (pour voir que \'etale est non ramifié).
Alors $(f \otimes g)(e) \in IB'$. Comme $IB'$ est localement nilpotent
(lemme d'algèbre \ref{algebra-lemma-locally-nilpotent}), cela implique
$(f \otimes g)(e) = 0$ d'après le lemme d'algèbre
\ref{algebra-lemma-lift-idempotents}. Ainsi $f = g$.

\medskip\noindent
Il est clair que $I$ est contenu dans le radical de Jacobson de $A$.
Soit $f \in A[T]$ un polynôme unitaire et soit
$\overline{f} = g_0h_0$ une factorisation
de $\overline{f} = f \bmod I$ avec $g_0, h_0 \in A/I[T]$ unitaires
engendrant l'idéal unité de $A/I[T]$. D'après
le lemme \ref{lemma-lift-factorization-monic},
il existe un morphisme d'anneaux \'etale $A \to A'$ qui
induit un isomorphisme $A/I \to A'/IA'$ tel que
la factorisation se relève en une factorisation en polynômes unitaires
sur $A'$. D'après ce qui précède, nous avons $A = A'$ et la factorisation
est sur $A$.
\end{proof}

\begin{lemma}
\label{lemma-limit-henselian}
Soit $A = \lim A_n$, où $(A_n)$ est un système inverse d'anneaux
dont les morphismes de transition sont surjectifs et ont des noyaux localement
nilpotents. Alors $(A, I_n)$ est une paire hensélienne, où
$I_n = \Ker(A \to A_n)$.
\end{lemma}

\begin{proof}
Fixons $n$. Soit $a \in A$ un élément qui s'envoie sur $1$ dans $A_n$.
D'après le lemme d'algèbre \ref{algebra-lemma-locally-nilpotent-unit},
$a$ s'envoie sur une unité de $A_m$ pour tout $m \geq n$.
Ainsi $a$ est une unité de $A$. Donc, d'après le lemme d'algèbre
\ref{algebra-lemma-contained-in-radical}, l'idéal $I_n$ est contenu dans le
radical de Jacobson de $A$. Soit $f \in A[T]$ un polynôme unitaire et soit
$\overline{f} = g_nh_n$ une factorisation
de $\overline{f} = f \bmod I_n$ avec $g_n, h_n \in A_n[T]$ unitaires
engendrant l'idéal unité de $A_n[T]$. D'après
le lemme \ref{lemma-locally-nilpotent-henselian}, nous pouvons relever
successivement cette factorisation en
$f \bmod I_m = g_m h_m$ avec $g_m, h_m$ unitaires
dans $A_m[T]$ pour tout $m \geq n$. À chaque étape, nous devons
vérifier que nos relèvements $g_m, h_m$ engendrent l'idéal unité dans
$A_n[T]$ ; cela résulte du fait correspondant pour
$g_n, h_n$ et du fait que $\Spec(A_n[T]) = \Spec(A_m[T])$
car le noyau de $A_m \to A_n$ est localement nilpotent.
Comme $A = \lim A_m$, cela conclut la démonstration.
\end{proof}

\begin{lemma}
\label{lemma-complete-henselian}
Soit $(A, I)$ une paire. Si $A$ est complet pour la topologie $I$-adique,
alors la paire est hensélienne.
\end{lemma}

\begin{proof}
D'après le lemme d'algèbre \ref{algebra-lemma-radical-completion},
l'idéal $I$ est contenu dans le radical de Jacobson de $A$.
Soit $f \in A[T]$ un polynôme unitaire et soit
$\overline{f} = g_0h_0$ une factorisation
de $\overline{f} = f \bmod I$ avec $g_0, h_0 \in A/I[T]$ unitaires
engendrant l'idéal unité de $A/I[T]$. D'après
le lemme \ref{lemma-locally-nilpotent-henselian}, nous pouvons relever
successivement cette factorisation en
$f \bmod I^n = g_n h_n$ avec $g_n, h_n$ unitaires
dans $A/I^n[T]$ pour tout $n \geq 1$.
Comme $A = \lim A/I^n$, cela conclut la démonstration.
\end{proof}

\begin{lemma}
\label{lemma-helper-finite-type}
Soit $(A, I)$ une paire. Soit $A \to B$ un morphisme d'anneaux de type fini
tel que $B/IB = C_1 \times C_2$ avec $A/I \to C_1$ fini.
Soit $B'$ la clôture intégrale de $A$ dans $B$.
Alors on peut écrire $B'/IB' = C_1 \times C'_2$ de sorte que
le morphisme $B'/IB' \to B/IB$ préserve les décompositions en produits
et qu'il existe un $g \in B'$ s'envoyant sur $(1, 0)$ dans
$C_1 \times C'_2$ tel que $B'_g \to B_g$ soit un isomorphisme.
\end{lemma}

\begin{proof}
Remarquons que $A \to B$ est quasi-fini en tout idéal premier du
fermé $T = \Spec(C_1) \subset \Spec(B)$ (cela résulte de l'examen des
anneaux des fibres, voir la définition d'algèbre
\ref{algebra-definition-quasi-finite}).
Considérons le diagramme d'espaces topologiques
$$
\xymatrix{
\Spec(B) \ar[rr]_\phi \ar[rd]_\psi & & \Spec(B') \ar[ld]^{\psi'} \\
& \Spec(A)
}
$$
D'après le théorème d'algèbre \ref{algebra-theorem-main-theorem},
pour tout $\mathfrak p \in T$, il existe un $h_\mathfrak p \in B'$, 
$h_\mathfrak p \not \in \mathfrak p$ tel que $B'_h \to B_h$ soit
un isomorphisme. L'union $U = \bigcup D(h_\mathfrak p)$ donne un ouvert
$U \subset \Spec(B')$ tel que $\phi^{-1}(U) \to U$ soit un homéomorphisme
et $T \subset \phi^{-1}(U)$. Puisque $T$ est ouvert dans $\psi^{-1}(V(I))$,
nous en déduisons que $\phi(T)$ est ouvert dans $U \cap (\psi')^{-1}(V(I))$.
Ainsi $\phi(T)$ est ouvert dans $(\psi')^{-1}(V(I))$.
D'autre part, puisque $C_1$ est fini sur $A/I$, il est
fini sur $B'$. Donc $\phi(T)$ est un fermé de $\Spec(B')$
d'après les lemmes d'algèbre \ref{algebra-lemma-going-up-closed} et
\ref{algebra-lemma-integral-going-up}. Nous concluons que
$\Spec(B'/IB') \supset \phi(T)$ est à la fois ouvert et fermé. D'après
le lemme d'algèbre \ref{algebra-lemma-disjoint-implies-product},
nous obtenons une décomposition en produit correspondante
$B'/IB' = C'_1 \times C'_2$.
Le morphisme $B'/IB' \to B/IB$ envoie $C'_1$ dans $C_1$ et $C'_2$ dans $C_2$,
comme on le voit en examinant ce qui se passe sur les spectres (indication :
l'image réciproque de $\phi(T)$ est exactement $T$ ; certains détails sont omis).
Choisissons un $g \in B'$ s'envoyant sur $(1, 0)$ dans $C'_1 \times C'_2$
tel que $D(g) \subset U$ ; cela est possible car $\Spec(C'_1)$
et $\Spec(C'_2)$ sont disjoints et fermés dans $\Spec(B')$ et
$\Spec(C'_1)$ est contenu dans $U$. Alors $B'_g \to B_g$ définit un homéomorphisme
sur les spectres et un isomorphisme sur les anneaux locaux (par notre choix de $U$).
C'est donc un isomorphisme, comme le montre par exemple le lemme d'algèbre
\ref{algebra-lemma-characterize-zero-local}.
Enfin, il s'ensuit que $C'_1 = C_1$, ce qui achève la démonstration.
\end{proof}

\begin{lemma}
\label{lemma-characterize-henselian-pair}
\begin{reference}
\cite[chapitre XI]{Henselian} et \cite[proposition 1]{Gabber-henselian}
\end{reference}
Soit $(A, I)$ une paire. Les assertions suivantes sont équivalentes
\begin{enumerate}
\item $(A, I)$ est une paire hensélienne ;
\item étant donné un morphisme d'anneaux \'etale $A \to A'$ et un morphisme de
$A$-algèbres $\sigma : A' \to A/I$, il existe un morphisme de $A$-algèbres
$A' \to A$ relevant $\sigma$ ;
\item pour toute $A$-algèbre finie $B$, le morphisme $B \to B/IB$ induit
une bijection sur les idempotents ;
\item pour toute $A$-algèbre entière $B$, le morphisme $B \to B/IB$ induit
une bijection sur les idempotents ;
\item (Gabber) $I$ est contenu dans le radical de Jacobson de $A$ et
tout polynôme unitaire $f(T) \in A[T]$ de la forme
$$
f(T) = T^n(T - 1) + a_n T^n + \ldots + a_1 T + a_0
$$
avec $a_n, \ldots, a_0 \in I$ et $n \ge 1$ possède une racine $\alpha \in 1 + I$.
\end{enumerate}
De plus, dans (5), la racine est unique.
\end{lemma}

\begin{proof}
Supposons (2). Alors $I$ est contenu dans le radical de Jacobson de $A$, car
sinon il existerait un $f \in A$ non inversible congru à $1$ modulo $I$,
et le morphisme $A \to A_f$ contredirait (2). Ainsi $IB \subset B$
est contenu dans le radical de Jacobson de $B$ pour $B$ entière sur $A$,
car $\Spec(B) \to \Spec(A)$ est fermé d'après les lemmes d'algèbre
\ref{algebra-lemma-going-up-closed} et
\ref{algebra-lemma-integral-going-up}.
Ainsi le morphisme des idempotents de $B$ vers ceux de $B/IB$
est injectif d'après le lemme \ref{lemma-idempotents-determined-modulo-radical}.
D'autre part, puisque (2) est vérifiée, tout idempotent
de $B/IB$ se relève en un idempotent de $B$
d'après le lemme \ref{lemma-lift-idempotent-upstairs}.
Nous voyons ainsi que (2) implique (4).

\medskip\noindent
L'implication (4) $\Rightarrow$ (3) est triviale.

\medskip\noindent
Supposons (3). Soit $\mathfrak m$ un idéal maximal et considérons le
morphisme fini $A \to B = A/(I \cap \mathfrak m)$. La condition selon laquelle
$B \to B/IB$ induit une bijection sur les idempotents implique que
$I \subset \mathfrak m$ (sinon $B = A/I \times A/\mathfrak m$
et $B/IB = A/I$). Nous voyons donc que $I$ est contenu dans le radical
de Jacobson de $A$. Soit $f \in A[T]$ unitaire et supposons donnée une
factorisation $\overline{f} = g_0h_0$ avec $g_0, h_0 \in A/I[T]$ unitaires
engendrant l'idéal unité dans $A/I[T]$.
Posons $B = A[T]/(f)$. Soit $\overline{e}$ l'idempotent
de $B/IB$ correspondant à la décomposition
$$
B/IB = A/I[T]/(g_0) \times A/I[T]/(h_0)
$$
en $A$-algèbres. Soit $e \in B$ un idempotent relevant $\overline{e}$,
qui existe puisque nous avons supposé (3). Cela donne une décomposition en produit
$$
B = eB \times (1 - e)B
$$
Remarquons que $B$ est libre de rang $\deg(f)$ comme $A$-module.
Ainsi $eB$ et $(1 - e)B$ sont des $A$-modules localement libres de type fini.
Cependant, puisque $eB$ et $(1 - e)B$ ont des rangs constants
$\deg(g_0)$ et $\deg(h_0)$ sur $A/I$, nous trouvons qu'il en est de même
sur $\Spec(A)$. Nous concluons que
\begin{align*}
f & = \text{CharPol}_A(T : B \to B) \\
& =
\text{CharPol}_A(T : eB \to eB)
\text{CharPol}_A(T : (1 - e)B \to (1 - e)B)
\end{align*}
est une factorisation en polynômes unitaires se réduisant sur la factorisation
donnée modulo $I$. Ici $\text{CharPol}_A$ désigne le polynôme caractéristique
d'un endomorphisme d'un module localement libre de type fini sur $A$.
Si le module est libre, $\text{CharPol}_A$ est défini comme le polynôme
caractéristique de la matrice correspondante et, en général, on utilise le lemme
d'algèbre \ref{algebra-lemma-standard-covering} pour recoller. Détails omis.
Ainsi (3) implique (1).

\medskip\noindent
Supposons (1). Soit $f$ comme dans (5). La factorisation de $f \bmod I$
en $T^n$ fois $T - 1$ se relève en une factorisation $f = gh$ avec $g$
et $h$ unitaires par la définition \ref{definition-henselian-pair}.
Alors $h$ doit être de degré $1$,
et nous voyons que $f$ possède une racine qui se réduit à $1$ modulo $I$.
Enfin, $I$ est contenu dans le radical de Jacobson par
la définition d'une paire hensélienne.
Ainsi (1) implique (5).

\medskip\noindent
Avant de donner la preuve de la dernière étape, montrons que
la racine $\alpha$ de (5), si elle existe, est unique. En effet,
en raison de la forme explicite de $f(T)$, nous avons
$f'(\alpha) \in 1 + I$, où $f'$ est la dérivée de $f$ par rapport à $T$.
Un argument élémentaire montre que
$$
f(T) = f(\alpha + T - \alpha) =
f(\alpha) + f'(\alpha) \cdot (T - \alpha) \bmod
(T - \alpha)^2 A[T]
$$
Cela montre que toute autre racine $\alpha' \in 1 + I$ de $f(T)$
vérifie $0 = f(\alpha') - f(\alpha) = (\alpha' - \alpha)(1 + i)$
pour un certain $i \in I$ ; puisque $1 + i$ est une unité de $A$,
nous avons donc $\alpha = \alpha'$.

\medskip\noindent
Supposons (5). Nous allons montrer que (2) est vérifiée, autrement dit que
pour tout morphisme \'etale $A \to A'$, toute section $\sigma : A' \to A/I$
modulo $I$ se relève en une section $A' \to A$.
Puisque $A \to A'$ est \'etale, la section $\sigma$
détermine une décomposition
\begin{equation}
\label{equation-GCHP}
A'/IA' \cong A/I \times C
\end{equation}
en $A/I$-algèbres. En effet, le morphisme d'anneaux surjectif
$A'/IA' \to A/I$ est \'etale d'après le lemme d'algèbre
\ref{algebra-lemma-map-between-etale},
et nous obtenons alors l'idempotent voulu par le lemme d'algèbre
\ref{algebra-lemma-surjective-flat-finitely-presented}.
Nous montrerons que cette décomposition se relève en une décomposition
\begin{equation}
\label{equation-GCHP-want}
A' \cong A'_1 \times A'_2
\end{equation}
en $A$-algèbres avec $A'_1$ entier sur $A$. Alors $A \to A'_1$
est entier et \'etale et $A/I \to A'_1/IA'_1$ est un isomorphisme,
donc $A \to A'_1$ est un isomorphisme d'après
le lemme \ref{lemma-check-isomorphism-zariski}
(nous utilisons aussi ici qu'un morphisme d'anneaux \'etale est plat
et de présentation finie, voir le lemme d'algèbre \ref{algebra-lemma-etale}).

\medskip\noindent
Soit $B'$ la clôture intégrale de $A$ dans $A'$. D'après
le lemme \ref{lemma-helper-finite-type}, nous pouvons décomposer
\begin{equation}
\label{equation-dec-mod-I}
B'/IB' \cong A/I \times C'
\end{equation}
en $A/I$-algèbres, de manière compatible avec (\ref{equation-GCHP}),
et trouver $b \in B'$ qui relève $(1, 0)$ tel que
$B'_b \to A'_b$ soit un isomorphisme. Si la décomposition
(\ref{equation-dec-mod-I}) se relève en une décomposition
\begin{equation}
\label{equation-want-2}
B' \cong B'_1 \times B'_2
\end{equation}
en $A$-algèbres, alors la décomposition induite
$A' = A'_1 \times A'_2$ donnera la décomposition voulue
(\ref{equation-GCHP-want}) : en effet, puisque $b$ est une unité dans $B'_1$
(détails omis), nous aurons $B'_1 \cong A'_1$, de sorte que $A'_1$ sera entier sur $A$.

\medskip\noindent
Choisissons une $A$-sous-algèbre finie $B'' \subset B'$ contenant $b$
(remarquons que toute $A$-sous-algèbre de type fini de $B'$ est
finie sur $A$). Après avoir agrandi $B''$, nous pouvons supposer que
$b$ s'envoie sur un idempotent dans $B''/IB''$ donnant
\begin{equation}
\label{equation-again-dec-mod-I}
B''/IB'' \cong C''_1 \times C''_2
\end{equation}
Puisque $B'_b \cong A'_b$, nous voyons que $B'_b$ est de type fini sur $A$.
Supposons que $B'_b$ soit engendré par $b_1/b^n, \ldots, b_t/b^n$ sur $A$
et agrandissons $B''$ de sorte que $b_1, \ldots, b_t \in B''$. Alors
$B''_b \to B'_b$ est à la fois surjectif et injectif, donc est un isomorphisme.
En particulier, nous voyons que $C''_1 = A/I$ ! Donc $A/I \to C''_1$
est un isomorphisme, en particulier surjectif.
Par le lemme \ref{lemma-helper-finite}, nous pouvons trouver
un $f(T) \in A[T]$ de la forme
$$
f(T) = T^n(T - 1) + a_n T^n + \ldots + a_1 T + a_0
$$
avec $a_n, \ldots, a_0 \in I$ et $n \ge 1$ tel que $f(b) = 0$.
En particulier, nous trouvons que $B'$ est une $A[T]/(f)$-algèbre.
D'après (5), il existe une racine $a \in 1 + I$ de $f$.
Cela produit une décomposition en produit $A[T]/(f) = A[T]/(T - a) \times D$
compatible avec la décomposition (\ref{equation-dec-mod-I}) de $B'/IB'$.
La décomposition induite de $B'$ est alors la décomposition voulue
(\ref{equation-want-2}).
\end{proof}

\begin{lemma}
\label{lemma-change-ideal-henselian-pair}
Soit $A$ un anneau. Soient $I, J \subset A$ des idéaux tels que $V(I) = V(J)$.
Alors $(A, I)$ est hensélienne si et seulement si $(A, J)$ est hensélienne.
\end{lemma}

\begin{proof}
Pour tout morphisme d'anneaux entier $A \to B$, on a $V(IB) = V(JB)$.
Les idempotents de $B/IB$ et $B/JB$ sont donc en correspondance bijective
(lemme d'algèbre \ref{algebra-lemma-disjoint-decomposition}).
Il s'ensuit que $B \to B/IB$ induit une bijection sur les ensembles des
idempotents si et seulement si $B \to B/JB$ induit une bijection sur les ensembles
d'idempotents. Nous concluons donc par le lemme
\ref{lemma-characterize-henselian-pair}.
\end{proof}

\begin{lemma}
\label{lemma-integral-over-henselian-pair}
Soit $(A, I)$ une paire hensélienne et soit $A \to B$ un morphisme d'anneaux
entier. Alors $(B, IB)$ est une paire hensélienne.
\end{lemma}

\begin{proof}
Cela résulte immédiatement de la quatrième caractérisation des paires
henséliennes dans le lemme \ref{lemma-characterize-henselian-pair} et du fait
que la composée de morphismes d'anneaux entiers est entière.
\end{proof}

\begin{lemma}
\label{lemma-henselian-henselian-pair}
Soient $I \subset J \subset A$ des idéaux d'un anneau $A$.
Les assertions suivantes sont équivalentes
\begin{enumerate}
\item $(A, I)$ et $(A/I, J/I)$ sont des paires henséliennes ;
\item $(A, J)$ est une paire hensélienne.
\end{enumerate}
\end{lemma}

\begin{proof}
Supposons (1). Soit $B$ une $A$-algèbre entière. Considérons les morphismes
d'anneaux
$$
B \to B/IB \to B/JB
$$
D'après le lemme \ref{lemma-characterize-henselian-pair}, les deux flèches
induisent des bijections sur les idempotents. Il en est donc de même de la
composée. Ainsi $(A, J)$ est une paire hensélienne
d'après le lemme \ref{lemma-characterize-henselian-pair}.

\medskip\noindent
Réciproquement, supposons (2). Alors $(A/I, J/I)$ est une paire hensélienne
d'après le lemme \ref{lemma-integral-over-henselian-pair}. Soit $B$ une
$A$-algèbre entière. Considérons les morphismes d'anneaux
$$
B \to B/IB \to B/JB
$$
D'après le lemme \ref{lemma-characterize-henselian-pair}, la composée
et la seconde flèche induisent des bijections sur les idempotents.
Il en est donc de même de la première flèche. Il s'ensuit que $(A, I)$ est une paire
hensélienne (par le même lemme).
\end{proof}

\begin{lemma}
\label{lemma-sum-henselian}
Soit $A$ un anneau et soient $(A, I)$ et $(A, I')$ des paires henséliennes.
Alors $(A, I + I')$ est une paire hensélienne.
\end{lemma}

\begin{proof}
D'après le lemme \ref{lemma-integral-over-henselian-pair}, la paire
$(A/I, (I' + I)/I)$ est hensélienne. Nous obtenons donc la conclusion du
lemme \ref{lemma-henselian-henselian-pair}.
\end{proof}

\begin{lemma}
\label{lemma-product-henselian-pairs}
Soit $J$ un ensemble et soit $\{ (A_j, I_j)\}_{j \in J}$ une collection
de paires. Alors $(\prod_{j \in J} A_j, \prod_{j\in J} I_j)$ est hensélienne
si et seulement si chaque $(A_j, I_j)$ l'est.
\end{lemma}
 
\begin{proof}
Pour tout $j \in J$, la projection $\prod_{j \in J} A_j \rightarrow A_j$
est un morphisme d'anneaux entier ; le lemme
\ref{lemma-integral-over-henselian-pair} montre donc que chaque $(A_j, I_j)$
est hensélienne si $(\prod_{j \in J} A_j, \prod_{j\in J} I_j)$ l'est.

\medskip\noindent
Réciproquement, supposons que chaque $(A_j, I_j)$ soit une paire hensélienne.
Alors tout $1 + x$ avec $x \in \prod_{j \in J} I_j$ est une unité
de $\prod_{j \in J} A_j$, car il en est ainsi composante par composante d'après
le lemme d'algèbre \ref{algebra-lemma-contained-in-radical} et la
définition \ref{definition-henselian-pair}.
Ainsi, de nouveau par le lemme d'algèbre \ref{algebra-lemma-contained-in-radical},
$\prod_{j \in J} I_j$ est contenu dans le radical de Jacobson
de $\prod_{j \in J} A_j$. En travaillant encore composante par composante, il
s'ensuit de même que pour tout polynôme unitaire
$f \in (\prod_{j \in J} A_j)[T]$ et toute factorisation
$\overline{f} = g_0h_0$ avec
$g_0, h_0 \in (\prod_{j \in J} A_j / \prod_{j \in J} I_j)[T] =
(\prod_{j \in J} A_j/I_j)[T]$ unitaires et engendrant l'idéal unité dans
$(\prod_{j \in J} A_j / \prod_{j \in J} I_j)[T]$, il existe une
factorisation $f = gh$ dans $(\prod_{j \in J} A_j)[T]$ avec $g$, $h$ unitaires
et se réduisant à $g_0$, $h_0$. En conclusion, d'après la définition
\ref{definition-henselian-pair},
$(\prod_{j \in J} A_j, \prod_{j\in J} I_j)$ est une paire hensélienne.
\end{proof}

\begin{lemma}
\label{lemma-limits-henselian}
La propriété d'être hensélienne est préservée par les limites de paires.
Plus précisément, soit $J$ un ensemble préordonné et soit $(A_j, I_j)$
un système projectif de paires henséliennes indexé par $J$.
Alors $A = \lim A_j$, muni de l'idéal $I = \lim I_j$,
est une paire hensélienne $(A, I)$.
\end{lemma}

\begin{proof}
D'après Catégories, lemme \ref{categories-lemma-limits-products-equalizers},
il suffit de considérer les produits et les égalisateurs.
Pour les produits, l'assertion découle du
lemme \ref{lemma-product-henselian-pairs}.
Considérons donc un diagramme égalisateur
$$
\xymatrix{
(A, I) \ar[r] & (A', I')
\ar@<1ex>[r]^{\varphi} \ar@<-1ex>[r]_{\psi}
&
(A'', I'')  
}
$$
dans lequel les paires $(A', I')$ et $(A'', I'')$ sont henséliennes.
Pour vérifier que la paire $(A, I)$ est elle aussi hensélienne, nous utiliserons le
critère de Gabber du lemme \ref{lemma-characterize-henselian-pair}.
Tout élément de $1 + I$ est une unité de $A$ car,
par unicité de l'inverse d'une unité,
cela peut se vérifier dans $(A', I')$.
Ainsi $I$ est contenu dans le radical de Jacobson de $A$, voir
Algèbre, lemme \ref{algebra-lemma-contained-in-radical}.
Soit donc
$$
f(T) = T^{N - 1}(T - 1) + a_{N - 1} T^{N - 1} + \dotsb + a_1 T + a_0
$$
un polynôme de $A[T]$ avec $a_{N - 1}, \dotsc, a_0 \in I$ et $N \ge 1$.
L'image de $f(T)$ dans $A'[T]$ admet une unique racine $\alpha' \in 1 + I'$
et il en est de même de son image dans $A''[T]$.
Ainsi, par unicité, $\varphi(\alpha') = \psi(\alpha')$,
de sorte que $\alpha'$ définit une racine de $f(T)$ dans $1 + I$, comme souhaité.
\end{proof}

\begin{lemma}
\label{lemma-filtered-colimits-henselian}
La propriété d'être hensélienne est préservée par les colimites filtrantes de paires.
Plus précisément, soit $J$ un ensemble dirigé et soit $(A_j, I_j)$
un système de paires henséliennes indexé par $J$.
Alors $A = \colim A_j$, muni de l'idéal $I = \colim I_j$,
est une paire hensélienne $(A, I)$.
\end{lemma}

\begin{proof}
Si $u \in 1 + I$, alors pour un certain $j \in J$, $u$ est l'image
d'un $u_j \in 1 + I_j$. Alors $u_j$ est inversible dans $A_j$ d'après
le lemme d'algèbre \ref{algebra-lemma-contained-in-radical}
et l'hypothèse que $I_j$ est contenu dans le radical de Jacobson de $A_j$.
Ainsi $u$ est inversible dans $A$. Donc $I$
est contenu dans le radical de Jacobson de $A$ (d'après le lemme).

\medskip\noindent
Soit $f \in A[T]$ un polynôme unitaire et soit $\overline{f} = g_0 h_0$
une factorisation avec $g_0, h_0 \in A/I[T]$ unitaires engendrant l'idéal unité
dans $A/I[T]$. Écrivons $1 = g_0 g'_0 + h_0 h'_0$ pour certains
$g'_0, h'_0 \in A/I[T]$. Comme $A = \colim A_j$ et $A/I = \colim A_j/I_j$
sont des colimites filtrantes, on peut trouver un $j \in J$, un $f_j \in A_j$ et
une factorisation $\overline{f}_j = g_{j, 0} h_{j, 0}$
avec $g_{j, 0}, h_{j, 0} \in A_j/I_j[T]$ unitaires
et $1 = g_{j, 0} g'_{j, 0} + h_{j, 0} h'_{j, 0}$
pour certains $g'_{j, 0}, h'_{j, 0} \in A_j/I_j[T]$, tels que
$f_j, g_{j, 0}, h_{j, 0}, g'_{j, 0}, h'_{j, 0}$
s'envoient respectivement sur $f, g_0, h_0, g'_0, h'_0$.
Puisque $(A_j, I_j)$ est une paire hensélienne, nous pouvons relever
$\overline{f}_j = g_{j, 0} h_{j, 0}$ en une factorisation
sur $A_j$ et, en prenant l'image dans $A$, nous obtenons une
factorisation correspondante dans $A$. Ainsi $(A, I)$ est hensélienne.
\end{proof}

\begin{example}[Moret-Bailly]
\label{example-moret-bailly}
Le lemme \ref{lemma-filtered-colimits-henselian} est faux si la colimite
n'est pas filtrante. Par exemple, si nous prenons le coproduit des
paires henséliennes $(\mathbf{Z}_p, (p))$ et $(\mathbf{Z}_p, (p))$, nous
obtenons $(A, pA)$ avec $A = \mathbf{Z}_p \otimes_\mathbf{Z} \mathbf{Z}_p$.
Ce n'est pas une paire hensélienne : $A/pA = \mathbf{F}_p$ ; donc si $(A, pA)$
était hensélienne, $A$ devrait être local. Cependant, $\Spec(A)$
est déconnecté ; par exemple, pour les nombres premiers impairs $p$, nous avons
l'idempotent non trivial
$$
(1/2 \otimes 1)
\left(1 \otimes 1 - (1 + p)^{-1}u \otimes u\right)
$$
où $u \in \mathbf{Z}_p$ est une racine carrée de $1 + p$. Certains détails sont omis.
\end{example}

\begin{lemma}
\label{lemma-largest-ideal-henselian}
Soit $A$ un anneau. Il existe un plus grand idéal $I \subset A$ tel que
$(A, I)$ soit une paire hensélienne.
\end{lemma}

\begin{proof}
Combinons les lemmes \ref{lemma-henselian-henselian-pair},
\ref{lemma-sum-henselian} et \ref{lemma-filtered-colimits-henselian}.
\end{proof}

\begin{lemma}
\label{lemma-irreducible-henselian-pair-connected}
Soit $(A, I)$ une paire hensélienne. Soit $\mathfrak p \subset A$
un idéal premier. Alors $V(\mathfrak p + I)$ est connexe.
\end{lemma}

\begin{proof}
D'après le lemme \ref{lemma-integral-over-henselian-pair}, nous voyons que
$(A/\mathfrak p, I + \mathfrak p/\mathfrak p)$ est une paire hensélienne.
Il suffit donc de démontrer : si $(A, I)$ est une paire hensélienne et
$A$ est un domaine, alors $\Spec(A/I) = V(I)$ est connexe. Sinon,
$A/I$ possède un idempotent non trivial d'après
le lemme d'algèbre \ref{algebra-lemma-characterize-spec-connected}.
D'après le lemme \ref{lemma-characterize-henselian-pair},
cela impliquerait que $A$ possède un idempotent non trivial. C'est une contradiction.
\end{proof}









\section{Hensélisation des paires}
\label{section-henselization}

\noindent
Nous poursuivons la discussion commencée dans
la section \ref{section-henselian-pairs}.

\begin{lemma}
\label{lemma-henselization}
Le foncteur d'inclusion
$$
\text{catégorie des paires henséliennes}
\longrightarrow
\text{catégorie des paires}
$$
admet un adjoint à gauche $(A, I) \mapsto (A^h, I^h)$.
\end{lemma}

\begin{proof}
Soit $(A, I)$ une paire. Considérons la catégorie $\mathcal{C}$ constituée
des morphismes d'anneaux \'etale $A \to B$ tels que $A/I \to B/IB$ soit un isomorphisme.
Nous montrerons que la catégorie $\mathcal{C}$ est dirigée et que
$A^h = \colim_{B \in \mathcal{C}} B$, munie de l'idéal $I^h = IA^h$,
donne l'adjoint voulu.

\medskip\noindent
Nous démontrons d'abord que $\mathcal{C}$ est dirigée
(définition de catégories \ref{categories-definition-directed}).
Elle est non vide car $\text{id} : A \to A$ est un objet.
Si $B$ et $B'$ sont deux objets de $\mathcal{C}$, alors
$B'' = B \otimes_A B'$ est un objet de $\mathcal{C}$
(utiliser le lemme d'algèbre \ref{algebra-lemma-etale})
et il existe des morphismes $B \to B''$ et $B' \to B''$.
Supposons que $f, g : B \to B'$ soient deux morphismes entre
objets de $\mathcal{C}$. Alors un coégaliseur est
$$
(B' \otimes_{f, B, g} B') \otimes_{(B' \otimes_A B')} B'
$$
qui est \'etale sur $A$ d'après les lemmes d'algèbre
\ref{algebra-lemma-etale} et
\ref{algebra-lemma-map-between-etale}.
Ainsi la catégorie $\mathcal{C}$ est dirigée.

\medskip\noindent
Comme $B/IB = A/I$ pour tout objet $B$ de $\mathcal{C}$, nous
voyons que $A^h/I^h = A^h/IA^h = \colim B/IB = \colim A/I = A/I$.

\medskip\noindent
Ensuite, nous montrons que $A^h = \colim_{B \in \mathcal{C}} B$, avec
$I^h = IA^h$, est une paire hensélienne. Pour cela, nous vérifierons
la condition (2) du lemme \ref{lemma-characterize-henselian-pair}.
Supposons donné un morphisme d'anneaux \'etale $A^h \to A'$
et un morphisme de $A^h$-algèbres $\sigma : A' \to A^h/I^h$. Alors il existe un
$B \in \mathcal{C}$ et un morphisme d'anneaux \'etale $B \to B'$ tels que
$A' = B' \otimes_B A^h$. Voir le lemme d'algèbre \ref{algebra-lemma-etale}.
Comme $A^h/I^h = A/I$, le morphisme $\sigma$ induit un morphisme de
$A$-algèbres $s : B' \to A/I$. Alors $B'/IB' = A/I \times C$ comme $A/I$-algèbre,
où $C$ est le noyau du morphisme $B'/IB' \to A/I$ induit par $s$.
Soit $g \in B'$ s'envoyant sur $(1, 0) \in A/I \times C$. Alors $B \to B'_g$
est \'etale et $A/I \to B'_g/IB'_g$ est un isomorphisme, c'est-à-dire que
$B'_g$ est un objet de $\mathcal{C}$. Nous obtenons donc un morphisme canonique
$B'_g \to A^h$ tel que
$$
\vcenter{
\xymatrix{
B'_g \ar[r] & A^h \\
B \ar[u] \ar[ur]
}
}
\quad\text{et}\quad
\vcenter{
\xymatrix{
B' \ar[r] \ar[rrd]_s & B'_g \ar[r] & A^h \ar[d] \\
& & A/I
}
}
$$
commutent. Cela induit un morphisme $A' = B' \otimes_B A^h \to A^h$
compatible avec $\sigma$, comme voulu.

\medskip\noindent
Soit $(A, I) \to (A', I')$ un morphisme de paires avec $(A', I')$ hensélienne.
Nous montrerons qu'il existe une factorisation unique $A \to A^h \to A'$, ce qui
achèvera la démonstration. En effet, pour chaque $A \to B$ dans $\mathcal{C}$,
le morphisme d'anneaux $A' \to B' = A' \otimes_A B$ est \'etale et induit
un isomorphisme $A'/I' \to B'/I'B'$. Il existe donc une section
$\sigma_B : B' \to A'$ d'après le lemme \ref{lemma-characterize-henselian-pair}.
Étant donné un morphisme $B_1 \to B_2$ dans $\mathcal{C}$, nous affirmons que le diagramme
$$
\xymatrix{
B'_1 \ar[rr] \ar[rd]_{\sigma_{B_1}} & &
B'_2 \ar[ld]^{\sigma_{B_2}} \\
& A'
}
$$
commute. Cela résulte du fait que, pour tout $B$ dans $\mathcal{C}$,
la section $\sigma_B$ est l'unique morphisme de $A'$-algèbres $B' \to A'$.
Nous avons $B' \otimes_{A'} B' = B' \times R$ pour un certain anneau $R$, voir
le lemme d'algèbre \ref{algebra-lemma-diagonal-unramified}. Dans notre cas,
$R/I'R = 0$ puisque $B'/I'B' = A'/I'$. Ainsi, étant donnés deux morphismes de
$A'$-algèbres $\sigma_B, \sigma_B' : B' \to A'$, alors
$e = (\sigma_B \otimes \sigma_B')(0, 1) \in A'$
est un idempotent contenu dans $I'$. Nous concluons que $e = 0$
d'après le lemme \ref{lemma-idempotents-determined-modulo-radical}.
Ainsi $\sigma_B = \sigma_B'$, comme voulu.
Par commutativité, nous obtenons
$$
A^h = \colim_{B \in \mathcal{C}} B \to
\colim_{B \in \mathcal{C}} A' \otimes_A B \xrightarrow{\colim \sigma_B} A'
$$
comme voulu. L'unicité des morphismes $\sigma_B$ garantit également que
ce morphisme est unique. Ainsi $(A, I) \mapsto (A^h, I^h)$ est l'adjoint voulu.
\end{proof}

\begin{lemma}
\label{lemma-henselization-flat}
Soit $(A, I)$ une paire. Soit $(A^h, I^h)$ comme dans
le lemme \ref{lemma-henselization}. Alors $A \to A^h$ est plat,
$I^h = IA^h$ et $A/I^n \to A^h/I^nA^h$ est un isomorphisme
pour tout $n$.
\end{lemma}

\begin{proof}
Dans la preuve du lemme \ref{lemma-henselization}, nous avons vu que
$A^h$ est une colimite filtrante d'$A$-algèbres \'etale $B$ telles que
$A/I \to B/IB$ soit un isomorphisme, et nous avons vu que
$I^h = IA^h$. Comme un morphisme d'anneaux \'etale est plat
(lemme d'algèbre \ref{algebra-lemma-etale}), nous concluons que
$A \to A^h$ est plat d'après le lemme d'algèbre
\ref{algebra-lemma-colimit-flat}.
Comme chaque $A \to B$ est plat, nous trouvons que les morphismes
$A/I^n \to B/I^nB$ sont également des isomorphismes (par exemple d'après
le lemme d'algèbre \ref{algebra-lemma-lift-basis}).
En prenant la colimite, nous trouvons $A/I^n = A^h/I^nA^h$,
comme voulu.
\end{proof}

\begin{lemma}
\label{lemma-henselization-local-ring}
\begin{slogan}
Compatibilité de l'hensélisation des paires et des anneaux locaux.
\end{slogan}
Le foncteur du lemme \ref{lemma-henselization} associe à un anneau local
$(A, \mathfrak m)$ son hensélisé.
\end{lemma}

\begin{proof}
Soit $(A^h, \mathfrak m^h)$ le hensélisé de la paire $(A, \mathfrak m)$
construit dans le lemme \ref{lemma-henselization}. Alors
$\mathfrak m^h = \mathfrak m A^h$ est un idéal maximal
d'après le lemme \ref{lemma-henselization-flat} et, puisqu'il
est contenu dans le radical de Jacobson, nous concluons que $A^h$ est
local d'idéal maximal $\mathfrak m^h$. Cela dit,
il y a deux façons d'achever la démonstration.

\medskip\noindent
Première preuve : remarquons que la construction comme colimite dans la preuve du
lemme d'algèbre \ref{algebra-lemma-henselization} est la même que la colimite
utilisée pour construire $A^h$ dans le lemme \ref{lemma-henselization}.
Deuxième preuve : les deux hensélisés $A \to S$ et $A \to A^h$ du
lemme \ref{lemma-henselization} sont des homomorphismes d'anneaux locaux,
$S$ et $A^h$ sont des colimites filtrantes d'$A$-algèbres \'etale,
$S$ et $A^h$ sont des anneaux locaux henséliens, et $S$ et $A^h$
ont pour corps résiduels $\kappa(\mathfrak m)$ (d'après
le lemme \ref{lemma-henselization-flat} dans le second cas).
Ils sont donc canoniquement isomorphes d'après
le lemme d'algèbre \ref{algebra-lemma-uniqueness-henselian}.
\end{proof}

\begin{lemma}
\label{lemma-henselization-Noetherian-pair}
\begin{slogan}
Le hensélisé d'une paire noethérienne est une paire noethérienne
ayant le même complété
\end{slogan}
Soit $(A, I)$ une paire avec $A$ noethérien. Soit $(A^h, I^h)$ comme dans
le lemme \ref{lemma-henselization}. Alors le morphisme des complétés $I$-adiques
$$
A^\wedge \to (A^h)^\wedge
$$
est un isomorphisme. De plus, $A^h$ est noethérien, les morphismes
$A \to A^h \to A^\wedge$ sont plats et $A^h \to A^\wedge$ est
fidèlement plat.
\end{lemma}

\begin{proof}
La première assertion est une conséquence immédiate du
lemme \ref{lemma-henselization-flat} et, en fait, reste vraie sans supposer $A$
noethérien. Dans la preuve du lemme \ref{lemma-henselization}, nous avons vu que
$A^h$ est une colimite filtrante d'$A$-algèbres \'etale $B$ telles que
$A/I \to B/IB$ soit un isomorphisme. Pour chaque tel $A \to B$,
le morphisme induit $A^\wedge \to B^\wedge$ est un isomorphisme
(voir la preuve du lemme \ref{lemma-henselization-flat}).
D'après le lemme d'algèbre \ref{algebra-lemma-completion-flat}, le morphisme
d'anneaux $B \to A^\wedge = B^\wedge = (A^h)^\wedge$ est plat pour chaque $B$.
Ainsi $A^h \to A^\wedge = (A^h)^\wedge$ est plat d'après
le lemme d'algèbre \ref{algebra-lemma-colimit-rings-flat}.
Comme $I^h = IA^h$ est contenu dans le radical de Jacobson de $A^h$
et que $A^h \to A^\wedge$ induit un isomorphisme $A^h/I^h \to A/I$,
nous voyons que $A^h \to A^\wedge$ est fidèlement plat d'après
le lemme d'algèbre \ref{algebra-lemma-ff}.
D'après le lemme d'algèbre \ref{algebra-lemma-completion-Noetherian-Noetherian},
l'anneau $A^\wedge$ est noethérien.
Nous concluons donc que $A^h$ est noethérien d'après
le lemme d'algèbre \ref{algebra-lemma-descent-Noetherian}.
\end{proof}

\begin{lemma}
\label{lemma-henselization-colimit}
Soit $(A, I) = \colim (A_i, I_i)$ une colimite filtrante de paires. Le foncteur du
lemme \ref{lemma-henselization} donne
$A^h = \colim A_i^h$ et $I^h = \colim I_i^h$.
\end{lemma}

\noindent
Ce lemme est faux pour les colimites non filtrantes, voir
l'exemple \ref{example-moret-bailly}.

\begin{proof}
D'après Catégories, lemme \ref{categories-lemma-adjoint-exact},
nous voyons que $(A^h, I^h)$ est la colimite du système $(A_i^h, I_i^h)$
dans la catégorie des paires henséliennes. Ainsi, pour une paire hensélienne $(B, J)$,
nous avons
$$
\Mor((A^h, I^h), (B, J)) =
\lim \Mor((A_i^h, I_i^h), (B, J)) =
\Mor(\colim (A_i^h, I_i^h), (B, J))
$$
Ici, la colimite est prise dans la catégorie des paires. Comme la colimite est
filtrante, nous obtenons $\colim (A_i^h, I_i^h) = (\colim A_i^h, \colim I_i^h)$
dans la catégorie des paires ; les détails sont omis. En utilisant de nouveau le fait que la colimite est filtrante,
c'est une paire hensélienne (lemme \ref{lemma-filtered-colimits-henselian}).
Il s'ensuit par le lemme de Yoneda que
$(A^h, I^h) = (\colim A_i^h, \colim I_i^h)$.
\end{proof}

\begin{lemma}
\label{lemma-henselization-change-ideal}
\begin{slogan}
L'hensélisé d'une paire ne dépend que du radical de l'idéal
\end{slogan}
Soit $A$ un anneau muni d'idéaux $I$ et $J$. Si $V(I) = V(J)$, alors le foncteur
du lemme \ref{lemma-henselization} produit le même anneau pour la paire
$(A, I)$ que pour la paire $(A, J)$.
\end{lemma}

\begin{proof}
Soit $(A', IA')$ la paire produite par le lemme \ref{lemma-henselization}
à partir de la paire $(A, I)$, voir le lemme \ref{lemma-henselization-flat}.
Soit $(A'', JA'')$ la paire produite par le lemme \ref{lemma-henselization}
à partir de la paire $(A, J)$. D'après le
lemme \ref{lemma-change-ideal-henselian-pair}, nous voyons que
$(A', JA')$ et $(A'', IA'')$ sont des paires henséliennes.
Par la propriété universelle de la construction, nous obtenons
d'uniques morphismes de $A$-algèbres $A'' \to A'$ et $A' \to A''$. Leur
unicité montre qu'ils sont inverses l'un de l'autre.
\end{proof}

\begin{lemma}
\label{lemma-henselization-integral}
\begin{slogan}
L'hensélisation commute au changement de base entier
\end{slogan}
Soit $(A, I) \to (B, J)$ un morphisme de paires tel que $V(J) = V(IB)$.
Soit $(A^h , I^h) \to (B^h, J^h)$ le morphisme induit
sur les hensélisés (lemme \ref{lemma-henselization}).
Si $A \to B$ est entier, alors le morphisme induit
$A^h \otimes_A B \to B^h$ est un isomorphisme.
\end{lemma}

\begin{proof}
D'après le lemme \ref{lemma-henselization-change-ideal}, nous pouvons supposer $J = IB$.
D'après le lemme \ref{lemma-integral-over-henselian-pair},
la paire $(A^h \otimes_A B, I^h(A^h \otimes_A B))$ est hensélienne.
Par la propriété universelle de $(B^h, IB^h)$, nous obtenons un morphisme
$B^h \to A^h \otimes_A B$. Nous omettons la démonstration
du fait que ce morphisme est l'inverse de celui du lemme.
\end{proof}


\begin{lemma}
\label{lemma-CRT-henselization}
Soient $I_1, I_2, \dots, I_n \subset A$ des idéaux d'un anneau $A$.
Supposons $I_i + I_j = A$ pour $i \not = j$, $1 \leq i, j \leq n$.
Notons $A^h$ l'hensélisé de la paire $(A, I_1 \cap \ldots \cap I_n)$
et $A_i^h$ l'hensélisé de la paire $(A, I_i)$.
Alors le morphisme naturel
$$
A^h \to \prod\nolimits_{i = 1, \ldots, n} A^h_i
$$
est un isomorphisme.
\end{lemma}

\begin{proof}
Le morphisme découle de la fonctorialité de la construction.
Si $n = 1$, le résultat est immédiat. Si $n > 1$, remarquons que
$(I_1 \cap \ldots \cap I_{n - 1}) + I_n = A$ d'après notre hypothèse
sur les idéaux. Ainsi, si nous démontrons le lemme pour $n = 2$, alors
le lemme s'ensuit pour tout $n$ par récurrence.

\medskip\noindent
Le cas $n = 2$. Choisissons $f_i \in I_i$ de sorte que $f_1 + f_2 = 1$.
Rappelons que $A^h$ est construit comme la colimite des $A$-algèbres $B$
telles que $A \to B$ soit \'etale et que $A/I_1 \cap I_2 \to B/(I_1 \cap I_2)B$
soit un isomorphisme. Voir la preuve du lemme \ref{lemma-henselization}.
Notons $Hens$ la catégorie de ces $A$-algèbres $B$.
De même, notons $Hens_i$, $i = 1, 2$, la catégorie des $A$-algèbres \'etales
$B$ telles que $A/I_i \to B/I_iB$ soit un isomorphisme.

\medskip\noindent
Considérons la sous-catégorie $Hens_1(f_2) \subset Hens_1$ formée des
$B$ tels que $f_2$ soit inversible. C'est une sous-catégorie cofinale
car, si $B$ appartient à $Hens_1$, alors $B_{f_2}$ appartient aussi à $Hens_1$
(indication : $B_{f_2}/I_1B_{f_2} = (B/I_1B)_{f_2} = B/I_1B$ car
$f_2$ est une unité modulo $I_1$). De même, la sous-catégorie
$Hens_2(f_1) \subset Hens_2$ est cofinale. Ainsi, pour calculer $A_1^h$ et
$A_2^h$, nous pouvons travailler avec ces catégories plus petites, voir
Catégories, lemme \ref{categories-lemma-cofinal}.

\medskip\noindent
Considérons le foncteur
$$
Hens_1(f_2) \times Hens_2(f_1) \to Hens,\quad (B_1, B_2) \mapsto B_1 \times B_2
$$
Ce foncteur est bien défini car
\begin{align*}
(B_1 \times B_2)/(I_1 \cap I_2)(B_1 \times B_2)
& =
B_1/(I_1 \cap I_2)B_1 \times B_2/(I_1 \cap I_2)B_2 \\
& =
B_1/I_1B_1 \times B_2/I_2B_2 \\
& = A/I_1 \times A/I_2 \\
& = A/(I_1 \cap I_2)
\end{align*}
La deuxième égalité est vraie car
$(I_1 \cap I_2)B_1 = I_1B_1 \cap I_2B_1 = I_1B_1$ puisque $f_2 \in I_2$
est inversible dans $B_1$
(nous avons aussi utilisé Algèbre, lemme \ref{algebra-lemma-flat-intersect-ideals}).
La dernière égalité résulte d'Algèbre, lemme \ref{algebra-lemma-chinese-remainder}.
Pour achever la démonstration, il suffit de montrer que le foncteur affiché
est cofinal (voir le lemme cité ci-dessus).
La première condition de
Catégories, définition \ref{categories-definition-cofinal},
est satisfaite car un objet $B$ de $Hens$ s'envoie sur $B_{f_2} \times B_{f_1}$,
qui appartient à l'image du foncteur. La seconde est toujours satisfaite car,
si nous avons $B \to B_1 \times B_2$ et $B \to B'_1 \times B'_2$ avec
$B_1, B'_1 \in Hens_1(f_2)$ et $B_2, B'_2 \in Hens_2(f_1)$, alors
nous pouvons poser $B''_1 = B_1 \otimes_B B'_1$, qui appartient à $Hens_1(f_2)$,
et de même $B''_2 = B_2 \otimes_B B''_2$ dans $Hens_2(f_1)$,
pour trouver un morphisme $B \to B''_1 \times B''_2$ s'insérant dans le diagramme commutatif
$$
\xymatrix{
& B \ar[ld] \ar[d] \ar[rd] \\
B_1 \times B_2 \ar[r] & B_1'' \times B_2'' & B_1' \times B_2' \ar[l]
}
$$
Cela achève la démonstration.
\end{proof}
\section{Relèvement et couples henséliens}
\label{section-lifting-henselian-pairs}

\noindent
Dans cette section, nous combinons principalement des résultats des
sections \ref{section-lifting} et \ref{section-henselian-pairs}.

\begin{lemma}
\label{lemma-lift-finite-projective-module}
Soit $(R, I)$ un couple hensélien. L'application
$$
P \longrightarrow P/IP
$$
induit une bijection entre les ensembles de classes d'isomorphisme de
$R$-modules projectifs de type fini et de $R/I$-modules projectifs de type
fini. En particulier, tout $R/I$-module projectif de type fini est isomorphe
à $P/IP$ pour un certain $R$-module projectif de type fini $P$.
\end{lemma}

\begin{proof}
Commençons par démontrer la dernière assertion. Soit $\overline{P}$ un
$R/I$-module projectif de type fini. Il existe un module projectif de type
fini $P'$ sur une certaine $R'$ \'etale sur $R$ telle que $R/I = R'/IR'$ et tel
que $P'/IP'$ soit isomorphe à $\overline{P}$ ; voir le
lemme \ref{lemma-lift-projective-module}.
Puisque $(R, I)$ est un couple hensélien, le morphisme d'anneaux \'etale
$R \to R'$ admet une section $\tau : R' \to R$
(lemme \ref{lemma-characterize-henselian-pair}). En posant
$P = P' \otimes_{R', \tau} R$, nous en déduisons que $P/IP$ est isomorphe
à $\overline{P}$. Cela montre bien sûr que l'application de l'énoncé du
lemme est surjective.

\medskip\noindent
Injectivité. Supposons que $P_1$ et $P_2$ soient des $R$-modules projectifs
de type fini tels que $P_1/IP_1 \cong P_2/IP_2$ comme $R/I$-modules.
Puisque $P_1$ est projectif, il existe un homomorphisme de $R$-modules
$u : P_1 \to P_2$ relevant l'isomorphisme donné. Le lemme de Nakayama montre
alors que $u$ est surjectif
(Algèbre, lemme \ref{algebra-lemma-NAK}). De même, il existe une surjection
$v : P_2 \to P_1$. D'après Algèbre, lemme \ref{algebra-lemma-fun},
l'application $v \circ u$ est un isomorphisme ; nous en concluons que $u$
est un isomorphisme.
\end{proof}

\begin{lemma}
\label{lemma-finite-etale-equivalence}
Soit $(A, I)$ un couple hensélien. Le foncteur $B \to B/IB$ définit une
équivalence entre les $A$-algèbres finies \'etales et les $A/I$-algèbres
finies \'etales.
\end{lemma}

\begin{proof}
Soient $B, B'$ deux $A$-algèbres finies \'etales sur $A$. Alors
$B' \to B'' = B \otimes_A B'$ est également fini \'etale
(Algèbre, lemmes \ref{algebra-lemma-etale} et
\ref{algebra-lemma-base-change-integral}). Nous avons des correspondances
$1$-à-$1$ bijectives entre
\begin{enumerate}
\item les homomorphismes de $A$-algèbres $B \to B'$,
\item les sections de $B' \to B''$, et
\item les idempotents $e$ de $B''$ tels que $B' \to B'' \to eB''$ soit
un isomorphisme.
\end{enumerate}
La bijection entre (2) et (3) associe à $\sigma : B'' \to B'$ l'élément
$e$ tel que $(1 - e)$ soit l'idempotent engendrant le noyau de $\sigma$ ;
celui-ci existe d'après Algèbre, lemmes
\ref{algebra-lemma-map-between-etale} et
\ref{algebra-lemma-surjective-flat-finitely-presented}.
De même, les homomorphismes de $A/I$-algèbres
$B/IB \to B'/IB'$ correspondent aux idempotents $\overline{e}$ de
$B''/IB''$ tels que
$B'/IB' \to B''/IB'' \to \overline{e}(B''/IB'')$ soit un isomorphisme.
Or tout idempotent $\overline{e}$ de $B''/IB''$ se relève de manière unique
en un idempotent $e$ de $B''$
(lemme \ref{lemma-characterize-henselian-pair}). De plus, si
$B'/IB' \to \overline{e}(B''/IB'')$ est un isomorphisme, alors le lemme de
Nakayama montre que $B' \to eB''$ est aussi un isomorphisme
(Algèbre, lemme \ref{algebra-lemma-NAK}). Ainsi, le foncteur est pleinement
fidèle.

\medskip\noindent
Surjectivité essentielle. Soit $A/I \to C$ un morphisme fini \'etale.
D'après Algèbre, lemme \ref{algebra-lemma-lift-etale}, il existe un morphisme
\'etale $A \to B$ tel que $B/IB \cong C$. Soit $B'$ la clôture intégrale de
$A$ dans $B$. D'après le lemme \ref{lemma-helper-finite-type}, on a
$B'/IB' = C \times C'$ pour un certain anneau $C'$, et $B'_g \cong B_g$
pour un certain $g \in B'$ dont l'image est $(1, 0) \in C \times C'$.
Puisque les idempotents se relèvent
(lemme \ref{lemma-characterize-henselian-pair}), nous obtenons
$B' = B'_1 \times B'_2$, avec $C = B'_1/IB'_1$ et $C' = B'_2/IB'_2$.
L'image de $g$ dans $B'_1$ est inversible. Alors
$B_g = B'_g = B'_1 \times (B_2)_g$, ce qui implique que $A \to B'_1$ est
\'etale. Nous en concluons que $B'_1$ est fini \'etale sur $A$ (par exemple,
un morphisme intégral et \'etale est fini \'etale d'après Algèbre,
lemme \ref{algebra-lemma-characterize-finite-in-terms-of-integral}), ce qui
achève la démonstration.
\end{proof}

\begin{lemma}
\label{lemma-lift-smooth-henselian}
Soient $R$ un anneau et $S$ une $R$-algèbre lisse. Supposons que $A$ soit
une $R$-algèbre et que $(A,I)$ soit un couple hensélien. Alors tout
homomorphisme de $R$-algèbres $S \to A/I$ se relève en un homomorphisme de
$R$-algèbres $S \to A$.
\end{lemma}

\begin{proof}
Soit $\tau : S \to A/I$ un homomorphisme de $R$-algèbres. Remarquons que
$S \otimes_R A$ est une $A$-algèbre lisse d'après Algèbre,
lemme \ref{algebra-lemma-base-change-smooth}. Ainsi, le
lemme \ref{lemma-lift-section-smooth-morphism} permet de relever
l'homomorphisme induit $S \otimes_R A \to A/I$ en un homomorphisme de
$A$-algèbres $S \otimes_R A \to A'$, où $A \to A'$ est \'etale et induit
un isomorphisme $A/I \to A'/IA'$. Comme $(A, I)$ est hensélien, il existe
un homomorphisme de $A$-algèbres $A' \to A$ ; voir le
lemme \ref{lemma-characterize-henselian-pair}. La composée
$S \to S \otimes_R A \to A' \to A$ est le relèvement cherché.
\end{proof}

\begin{lemma}
\label{lemma-lim-finite-projective-gives-finite-projective}
Soit $A = \lim A_n$ la limite projective d'un système projectif $(A_n)$
d'anneaux. Supposons donnés des $A_n$-modules $M_n$ et des homomorphismes de
$A_{n + 1}$-modules $M_{n + 1} \to M_n$. Supposons que
\begin{enumerate}
\item les morphismes de transition $A_{n + 1} \to A_n$ sont surjectifs et
leurs noyaux sont localement nilpotents,
\item $M_1$ est un $A_1$-module projectif de type fini,
\item $M_n$ est un $A_n$-module plat de type fini, et
\item les homomorphismes induisent des isomorphismes
$M_{n + 1} \otimes_{A_{n + 1}} A_n \to M_n$.
\end{enumerate}
Alors $M = \lim M_n$ est un $A$-module projectif de type fini et
$M \otimes_A A_n \to M_n$ est un isomorphisme pour tout $n$.
\end{lemma}

\begin{proof}
D'après le lemme \ref{lemma-limit-henselian}, le couple
$(A, \Ker(A \to A_1))$ est hensélien. D'après le
lemme \ref{lemma-lift-finite-projective-module}, nous pouvons choisir un
$A$-module projectif de type fini $P$ et un isomorphisme
$P \otimes_A A_1 \to M_1$. Comme $P$ est projectif, nous pouvons relever
successivement l'homomorphisme de $A$-modules $P \to M_1$ en des
homomorphismes de $A$-modules $P \to M_2$, $P \to M_3$, et ainsi de suite.
Nous obtenons donc un homomorphisme
$$
P \longrightarrow M
$$
Comme $P$ est projectif de type fini, on peut écrire
$A^{\oplus m} = P \oplus Q$ pour un certain $m \geq 0$ et un certain
$A$-module $Q$. Puisque $A = \lim A_n$, il vient
$P = \lim P \otimes_A A_n$. Pour montrer que l'homomorphisme de $A$-modules
affiché est un isomorphisme, il suffit donc de montrer que les homomorphismes
$P \otimes_A A_n \to M_n$ sont des isomorphismes. Le
lemme \ref{lemma-lift-projective} montre que $M_n$ est un module projectif
de type fini. D'après le
lemme \ref{lemma-isomorphic-finite-projective-lifts}, les homomorphismes
$P \otimes_A A_n \to M_n$ sont des isomorphismes.
\end{proof}






\section{Clôture intégrale absolue}
\label{section-absolute-integral-closure}

\noindent
Voici notre définition.

\begin{definition}
\label{definition-absolutely-integrally-closed}
Un anneau $A$ est {\it absolument intégralement clos} si tout polynôme
unitaire $f \in A[T]$ est un produit de facteurs linéaires.
\end{definition}

\noindent
Attention : il peut être possible d'écrire $f$ comme produit de facteurs
linéaires de nombreuses manières différentes.

\begin{lemma}
\label{lemma-absolutely-integrally-closed}
Soit $A$ un anneau. Les assertions suivantes sont équivalentes :
\begin{enumerate}
\item $A$ est absolument intégralement clos, et
\item tout polynôme unitaire $f \in A[T]$ possède une racine dans $A$.
\end{enumerate}
\end{lemma}

\begin{proof}
Démonstration omise.
\end{proof}

\begin{lemma}
\label{lemma-absolutely-integrally-closed-quotient-localization}
Soit $A$ un anneau absolument intégralement clos.
\begin{enumerate}
\item Tout anneau quotient $A/I$ de $A$ est absolument intégralement clos.
\item Tout localisé $S^{-1}A$ est absolument intégralement clos.
\end{enumerate}
\end{lemma}

\begin{proof}
Démonstration omise.
\end{proof}

\begin{lemma}
\label{lemma-integrally-closed-in-absolutely-integrally-closed}
Soit $A$ un anneau. Soit $S \subset A$ une partie multiplicative formée de
non-diviseurs de zéro. Si $S^{-1}A$ est absolument intégralement clos et si
$A \subset S^{-1}A$ est intégralement clos dans $S^{-1}A$, alors $A$ est
absolument intégralement clos.
\end{lemma}

\begin{proof}
Démonstration omise.
\end{proof}

\begin{lemma}
\label{lemma-normal-domain-absolutely-integrally-closed}
Soit $A$ un anneau intègre normal. Alors $A$ est absolument intégralement
clos si et seulement si son corps des fractions est algébriquement clos.
\end{lemma}

\begin{proof}
Remarquons qu'un corps est algébriquement clos si et seulement s'il est
absolument intégralement clos en tant qu'anneau. Le lemme résulte donc des
lemmes \ref{lemma-absolutely-integrally-closed-quotient-localization} et
\ref{lemma-integrally-closed-in-absolutely-integrally-closed}.
\end{proof}

\begin{lemma}
\label{lemma-construct-absolute-integral-closure}
Pour tout anneau $A$, il existe une extension $A \subset B$ telle que
\begin{enumerate}
\item $B$ est une colimite filtrante de $A$-algèbres libres de type fini,
\item $B$ est libre comme $A$-module, et
\item $B$ est absolument intégralement clos.
\end{enumerate}
\end{lemma}

\begin{proof}
Soit $I$ l'ensemble des polynômes unitaires sur $A$. Pour $i \in I$, notons
$x_i$ une indéterminée et $P_i$ le polynôme unitaire correspondant en
l'indéterminée $x_i$. Posons
$$
F(A) = A[x_i; i \in I]/(P_i; i \in I)
$$
Comme la notation le suggère, $F$ est un foncteur de la catégorie des
anneaux dans elle-même. Notons que $A \subset F(A)$, que $F(A)$ est libre
comme $A$-module et que $F(A)$ est une colimite filtrante de $A$-algèbres
libres de type fini. Prenons ensuite
$$
B = \colim F^n(A)
$$
où les morphismes de transition sont les inclusions
$F^n(A) \subset F(F^n(A)) = F^{n + 1}(A)$. Tout polynôme unitaire à
coefficients dans $B$ a en fait ses coefficients dans $F^n(A)$ pour un
certain $n$ et possède donc, par construction, une solution dans
$F^{n + 1}(A)$. Le lemme \ref{lemma-absolutely-integrally-closed} implique
que $B$ est absolument intégralement clos. Nous omettons la démonstration
des autres propriétés.
\end{proof}

\begin{lemma}
\label{lemma-absolutely-integrally-closed-strictly-henselian}
Soit $A$ un anneau absolument intégralement clos. Soit
$\mathfrak p \subset A$ un idéal premier. Alors l'anneau local
$A_\mathfrak p$ est strictement hensélien.
\end{lemma}

\begin{proof}
D'après le lemme
\ref{lemma-absolutely-integrally-closed-quotient-localization}, nous pouvons
supposer que $A$ est un anneau local et que $\mathfrak p$ est son idéal
maximal. Le corps résiduel est algébriquement clos d'après le
lemme \ref{lemma-absolutely-integrally-closed-quotient-localization}.
Tout polynôme unitaire se décompose entièrement en facteurs linéaires ;
Algèbre, définition \ref{algebra-definition-henselian}, s'applique donc
directement.
\end{proof}

\begin{lemma}
\label{lemma-absolutely-integrally-closed-henselian-pair}
Soit $A$ un anneau absolument intégralement clos. Soit $I \subset A$ un
idéal. Alors $(A, I)$ est un couple hensélien si (et seulement si) les
conditions suivantes sont satisfaites :
\begin{enumerate}
\item $I$ est contenu dans le radical de Jacobson de $A$,
\item $A \to A/I$ induit une bijection sur les idempotents.
\end{enumerate}
\end{lemma}

\begin{proof}
Soit $f \in A[T]$ un polynôme unitaire et soit
$f \bmod I = g_0 h_0$ une factorisation sur $A/I$, avec $g_0$ et $h_0$
unitaires, telle que $g_0$ et $h_0$ engendrent l'idéal unité de $A/I[T]$.
Cela signifie que
$$
A/I[T]/(f) = A/I[T]/(g_0) \times A/I[T]/(h_0)
$$
Notons $e \in A/I[T]/(f)$ l'élément correspondant à l'idempotent $(1, 0)$
de l'anneau de droite. Écrivons
$f = (T - a_1) \ldots (T - a_d)$ avec $a_i \in A$. Pour chaque
$i \in \{1, \ldots, d\}$, nous obtenons un homomorphisme de $A$-algèbres
$\varphi_i : A[T]/(f) \to A$, $T \mapsto a_i$, qui induit un homomorphisme
analogue de $A/I$-algèbres
$\overline{\varphi}_i : A/I[T]/(f) \to A/I$. Notons
$e_i = \overline{\varphi}_i(e) \in A/I$. Ce sont des idempotents. D'après
notre hypothèse (2), nous pouvons relever $e_i$ en un idempotent de $A$.
Nous pouvons donc écrire $A = \prod A_j$ comme produit fini d'anneaux de
sorte que, dans $A_j/IA_j$, chaque $e_i$ soit égal à $0$ ou à $1$.
Nous omettons certains détails. Remarquons que $A_j$ est absolument
intégralement clos comme anneau quotient de $A$. Il suffit de relever la
factorisation de $f$ sur $A_j/IA_j$ dans $A_j$. Nous sommes ainsi ramenés à
la situation étudiée dans le paragraphe suivant.

\medskip\noindent
Supposons que $e_i = 1$ pour $i = 1, \ldots, r$ et que $e_i = 0$ pour
$i = r + 1, \ldots, d$. Comme $(g_0, h_0) = A/I[T]$, il existe
$k_0, l_0 \in A/I[T]$ tels que $g_0 k_0 + h_0 l_0 = 1$. On voit que
$e = h_0 l_0$ et que $e_i = h_0(a_i) l_0(a_i)$. Nous en déduisons que
$h_0(a_i)$ est une unité pour $i = 1, \dots ,r$. Comme $f(a_i) = 0$, nous
avons $0 = h_0(a_i)g_0(a_i)$, d'où $g_0(a_i) = 0$ pour
$i = 1, \ldots, r$. Ainsi, $(T - a_1)$ divise $g_0$ dans $A/I[T]$ ;
écrivons $g_0 = (T - a_1) g_0'$. Posons
$f' = (T - a_2) \ldots (T - a_d)$ et $h'_0 = h_0$. Par récurrence sur $d$,
nous pouvons relever la factorisation $f' \bmod I = g'_0 h'_0$ en une
factorisation $f' = g' h'$ sur $A$, qui fournit la factorisation
$f = (T - a_1) g' h'$ relevant la factorisation
$f \bmod I = g_0 h_0$, comme voulu.
\end{proof}









\section{Anneaux auto-associés}
\label{section-auto-ass}

\noindent
Une partie de ce qui suit se trouve dans \cite{Autour}.

\begin{definition}
\label{definition-auto-ass}
On dit qu'un anneau $R$ est {\it auto-associé} si $R$ est local et si son
idéal maximal $\mathfrak m$ est faiblement associé à $R$.
\end{definition}

\begin{lemma}
\label{lemma-auto-ass-implies-P}
Un anneau auto-associé $R$ possède la propriété suivante : (P)
Tout idéal de type fini propre $I \subset R$ possède un annulateur non nul.
\end{lemma}

\begin{proof}
Par hypothèse, il existe un élément non nul $x \in R$ tel que, pour tout
$f \in \mathfrak m$, on ait $f^n x = 0$. Écrivons
$I = (f_1, \ldots, f_r)$. Alors $x$ appartient au noyau de
$R \to \bigoplus R_{f_i}$. Il existe donc un élément non nul $y \in R$ tel
que $f_i y = 0$ pour tout $i$ ; voir Algèbre,
lemme \ref{algebra-lemma-when-injective-covering}. Comme
$y \in \text{Ann}_R(I)$, le résultat est démontré.
\end{proof}

\begin{lemma}
\label{lemma-P-universally-injective}
Soit $R$ un anneau possédant la propriété (P) du
lemme \ref{lemma-auto-ass-implies-P}. Soit $u : N \to M$ un homomorphisme
de $R$-modules projectifs. Alors $u$ est universellement injectif si et
seulement si $u$ est injectif.
\end{lemma}

\begin{proof}
Supposons $u$ injectif. Il s'agit de montrer que $u$ est universellement
injectif. Choisissons d'abord un module $Q$ tel que $N \oplus Q$ soit libre.
En considérant l'homomorphisme $N \oplus Q \to M \oplus Q$, on voit qu'il
suffit de démontrer le lemme lorsque $N$ est libre. Dans ce cas, $N$ est une
colimite filtrante de $R$-modules libres de type fini. Nous sommes donc
ramenés au cas où $N$ est un $R$-module libre de type fini, disons
$N = R^{\oplus n}$. Démontrons le lemme par récurrence sur $n$. Le cas
$n = 0$ est trivial.

\medskip\noindent
Soit $u : R^{\oplus n} \to M$ un homomorphisme injectif de modules, avec
$M$ projectif. Choisissons un $R$-module $Q$ tel que $M \oplus Q$ soit
libre. En remplaçant $u$ par la composée
$R^{\oplus n} \to M \to M \oplus Q$, on voit que l'on peut supposer $M$
libre. Il existe alors un facteur direct $R^{\oplus m} \subset M$ tel que
$u(R^{\oplus n}) \subset R^{\oplus m}$. Nous pouvons donc supposer que
$M = R^{\oplus m}$. Dans ce cas, $u$ est donné par une matrice
$A = (a_{ij})$ telle que
$u(x_1, \ldots, x_n) = (\sum x_i a_{i1}, \ldots, \sum x_i a_{im})$.
Comme $u$ est injectif, on a en particulier
$u(x, 0, \ldots, 0) = (xa_{11}, xa_{12}, \ldots, xa_{1m}) \not = 0$ si
$x \not = 0$, et, puisque $R$ possède la propriété (P), on voit que
$a_{11}R + a_{12}R + \ldots + a_{1m}R = R$. Par conséquent,
$R(a_{11}, \ldots, a_{1m}) \subset R^{\oplus m}$ est un facteur direct de
$R^{\oplus m}$ ; en particulier,
$R^{\oplus m}/R(a_{11}, \ldots, a_{1m})$ est un $R$-module projectif. Nous
obtenons un diagramme commutatif
$$
\xymatrix{
0 \ar[r] &
R \ar[rr] \ar[d]^1 & & R^{\oplus n} \ar[r] \ar[d]^u &
R^{\oplus n - 1} \ar[r] \ar[d] & 0 \\
0 \ar[r] & R \ar[rr]^{(a_{11}, \ldots, a_{1m})} & &
R^{\oplus m} \ar[r] & R^{\oplus m}/R(a_{11}, \ldots, a_{1m}) \ar[r] & 0
}
$$
dont les lignes sont exactes et scindées. La flèche verticale de droite est
donc injective et l'hypothèse de récurrence montre qu'elle est
universellement injective. Il en résulte que la flèche verticale du milieu
est universellement injective.
\end{proof}

\begin{lemma}
\label{lemma-P-fPD-zero}
Soit $R$ un anneau. Les assertions suivantes sont équivalentes :
\begin{enumerate}
\item $R$ possède la propriété (P) du
lemme \ref{lemma-auto-ass-implies-P},
\item tout homomorphisme injectif de $R$-modules projectifs est
universellement injectif,
\item si $u : N \to M$ est injectif et si $N$, $M$ sont des $R$-modules
projectifs de type fini, alors $\Coker(u)$ est un $R$-module projectif de
type fini,
\item si $N \subset M$ et si $N$, $M$ sont projectifs de type fini comme
$R$-modules, alors $N$ est un facteur direct de $M$, et
\item tout homomorphisme injectif $R \to R^{\oplus n}$ est une injection
scindée.
\end{enumerate}
\end{lemma}

\begin{proof}
L'implication (1) $\Rightarrow$ (2) est le
lemme \ref{lemma-P-universally-injective}. Il est clair que (3) et (4) sont
équivalentes. L'implication (2) $\Rightarrow$ (3), (4) résulte d'Algèbre,
lemme \ref{algebra-lemma-universally-exact-split}. L'assertion (5) est un
cas particulier de (4). Supposons (5). Soit $I = (a_1, \ldots, a_n)$ un
idéal propre de type fini de $R$. Comme $I \not = R$, l'homomorphisme
$R \to R^{\oplus n}$, $x \mapsto (xa_1, \ldots, xa_n)$
n'est pas une injection scindée. Son noyau est donc non nul, et nous en
concluons que $\text{Ann}_R(I) \not = 0$. Ainsi, (1) est satisfaite.
\end{proof}

\begin{example}
\label{example-auto-ass-weird-flat}
Même lorsque les conditions équivalentes du
lemme \ref{lemma-P-fPD-zero} sont satisfaites, tout homomorphisme injectif
de $R$-modules libres n'est pas nécessairement une injection scindée. Par
exemple, supposons que $R = k[x_1, x_2, x_3, \ldots]/(x_i^2)$. C'est un
anneau auto-associé. Considérons l'homomorphisme de $R$-modules libres
$$
u :
\bigoplus\nolimits_{i \geq 1} Re_i
\longrightarrow
\bigoplus\nolimits_{i \geq 1} Rf_i, \quad
e_i \longmapsto f_i - x_if_{i + 1}.
$$
Pour tout entier $n$, la restriction de $u$ à
$\bigoplus_{i = 1, \ldots, n} Re_i$ est injective, car les images
$u(e_1), \ldots, u(e_n)$ sont $R$-linéairement indépendantes. Ainsi, $u$
est injectif et donc universellement injectif d'après le lemme. Comme
$u \otimes \text{id}_k$ est bijectif, si $u$ était une injection scindée,
alors $u$ serait surjectif. Or $u$ n'est pas surjectif, car l'antécédent de
$f_1$ serait l'élément
$$
\sum\nolimits_{i \geq 0} x_1 \ldots x_ie_{i + 1} =
e_1 + x_1e_2 + x_1x_2e_3 + \ldots
$$
qui n'appartient pas à la somme directe. Remarquons au passage que
$\Coker(u)$ est plat (car $u$ est universellement injectif), est un
$R$-module engendré par une famille dénombrable, qui n'est pas projectif
(car $u$ n'est pas scindé), et n'est donc pas de Mittag-Leffler (voir Algèbre,
lemme \ref{algebra-lemma-countgen-projective}).
\end{example}

\noindent
Le lemme suivant est un cas particulier d'Algèbre,
proposition \ref{algebra-proposition-what-exact}, lorsque l'anneau local est
noethérien.

\begin{lemma}
\label{lemma-exact-length-1}
Soit $(R, \mathfrak m)$ un anneau local. Supposons que
$\varphi : R^m \to R^n$ soit un homomorphisme de modules libres de type
fini. Les assertions suivantes sont équivalentes :
\begin{enumerate}
\item $\varphi$ est injectif,
\item le rang de $\varphi$ est $m$ et l'annulateur de $I(\varphi)$ dans
$R$ est nul.
\end{enumerate}
Si $R$ est noethérien, elles sont aussi équivalentes à
\begin{enumerate}
\item[(3)] le rang de $\varphi$ est $m$ et soit $I(\varphi) = R$, soit
il contient un non-diviseur de zéro.
\end{enumerate}
Ici, le rang de $\varphi$ et $I(\varphi)$ sont définis comme dans Algèbre,
définition \ref{algebra-definition-rank}.
\end{lemma}

\begin{proof}
Si l'un des coefficients de la matrice de $\varphi$ n'appartient pas à
$\mathfrak m$, nous appliquons Algèbre,
lemme \ref{algebra-lemma-add-trivial-complex}, pour écrire $\varphi$ comme
somme de $1 : R \to R$ et d'un homomorphisme
$\varphi' : R^{m-1} \to R^{n-1}$. On vérifie aisément que le lemme pour
$\varphi'$ implique le lemme pour $\varphi$. Nous pouvons donc supposer dès
le départ que tous les coefficients de la matrice de $\varphi$ appartiennent
à $\mathfrak m$.

\medskip\noindent
Supposons $\varphi$ injectif. Nous pouvons supposer $m > 0$. Soit
$\mathfrak q \in \text{WeakAss}(R)$, de sorte que $R_\mathfrak q$ soit un
anneau auto-associé. Alors $\varphi$ induit un homomorphisme injectif
$R_\mathfrak q^m \to R_\mathfrak q^n$ qui est universellement injectif
d'après les lemmes \ref{lemma-auto-ass-implies-P} et
\ref{lemma-P-universally-injective}. Ainsi,
$\varphi : \kappa(\mathfrak q)^m \to \kappa(\mathfrak q)^n$ est injectif.
Le rang de $\varphi \bmod \mathfrak q$ vaut donc $m$ et
$I(\varphi \otimes \kappa(\mathfrak q))$ n'est pas l'idéal nul. Comme $m$
est le rang maximal que $\varphi$ puisse avoir, nous en concluons que
$\varphi$ est également de rang $m$ (le rang d'une matrice ne peut que
décroître après un changement de base). Par conséquent,
$I(\varphi) \cdot \kappa(\mathfrak q) =
I(\varphi \otimes \kappa(\mathfrak q))$ n'est pas nul. Ainsi,
$I(\varphi)$ n'est pas contenu dans $\mathfrak q$. Aucun idéal premier
faiblement associé à $R$ n'est donc faiblement associé au $R$-module
$\text{Ann}_R I(\varphi)$. Par suite, $\text{Ann}_R I(\varphi)$ n'a aucun
idéal premier faiblement associé ; voir Algèbre,
lemme \ref{algebra-lemma-weakly-ass}. Il résulte d'Algèbre,
lemme \ref{algebra-lemma-weakly-ass-zero}, que
$\text{Ann}_R I(\varphi)$ est nul.

\medskip\noindent
Réciproquement, supposons (2). Puisque le rang vaut $m$, on a $n \geq m$.
Écrivons $I(\varphi) = (f_1, \ldots, f_r)$, ce qui est possible puisque
$I(\varphi)$ est de type fini. D'après Algèbre,
lemme \ref{algebra-lemma-matrix-left-inverse}, il existe des homomorphismes
$\psi_i : R^n \to R^m$ tels que
$\psi \circ \varphi = f_i \text{id}_{R^m}$. Ainsi, $\varphi(x) = 0$
implique $f_i x = 0$ pour $i = 1, \ldots, r$. Il en résulte que $x = 0$ et
donc que $\varphi$ est injectif.

\medskip\noindent
Pour l'équivalence de (1) et (3) dans le cas local noethérien, nous renvoyons
à Algèbre, proposition \ref{algebra-proposition-what-exact}. Si l'anneau
$R$ est noethérien mais non local, le lecteur peut la déduire du cas local ;
nous omettons les détails. On peut aussi reprendre l'argument ci-dessus en
utilisant les idéaux premiers associés, leur finitude, l'évitement des
idéaux premiers et la caractérisation des non-diviseurs de zéro comme les
éléments d'un anneau noethérien qui n'appartiennent à aucun idéal premier
associé.
\end{proof}

\begin{lemma}
\label{lemma-coker-injective-free}
Soit $R$ un anneau. Supposons que $\varphi : R^n \to R^n$ soit un
homomorphisme injectif de modules libres de type fini de même rang. Alors
$\Hom_R(\Coker(\varphi), R) = 0$.
\end{lemma}

\begin{proof}
Soit $\varphi^t : R^n \to R^n$ la transposée de $\varphi$. Le lemme affirme
que $\varphi^t$ est injectif. Avec les notations du
lemme \ref{lemma-exact-length-1}, on voit que le rang de $\varphi^t$ vaut
$n$ et que $I(\varphi) = I(\varphi^t)$. On conclut donc par l'équivalence de
(1) et (2) dans ce lemme.
\end{proof}









\section{Stratification par platification}
\label{section-flattening}

\noindent
Soient $R \to S$ un morphisme d'anneaux et $M$ un $S$-module. Pour toute
$R$-algèbre $R'$, nous pouvons considérer les changements de base
$S' = S \otimes_R R'$ et $M' = M \otimes_R R'$. Nous disons que
$R \to R'$ {\it platifie} $M$ si le module $M'$ est plat sur $R'$. Nous
voudrions comprendre la structure de la collection des morphismes d'anneaux
$R \to R'$ qui platifient $M$. Nous voudrions notamment savoir s'il existe
une {\it platification universelle $R \to R_{univ}$ de $M$}, c'est-à-dire un
morphisme d'anneaux $R \to R_{univ}$ qui platifie $M$ et tel que tout
morphisme d'anneaux $R \to R'$ qui platifie $M$ se factorise par
$R \to R_{univ}$. Il s'avère qu'une telle solution universelle n'existe
généralement pas.

\medskip\noindent
Nous étudierons les {\it platifications universelles} et les
{\it stratifications par platification} dans le cadre schématique
$\mathcal{F}/X/S$ dans Compléments sur la platitude,
section \ref{flat-section-flattening}. Si la platification universelle
$R \to R_{univ}$ existe, alors le morphisme de schémas
$\Spec(R_{univ}) \to \Spec(R)$ est la platification universelle du module
quasi-cohérent $\widetilde{M}$ sur $\Spec(S)$.

\medskip\noindent
Dans cette section et les quelques suivantes, nous démontrons des faits
algébriques élémentaires qui s'y rapportent. Le résultat le plus élémentaire
est peut-être le suivant.

\begin{lemma}
\label{lemma-intersection-flat}
Soit $R$ un anneau. Soit $M$ un $R$-module. Soient $I_1$, $I_2$ des idéaux
de $R$. Si $M/I_1M$ est plat sur $R/I_1$ et $M/I_2M$ est plat sur $R/I_2$,
alors $M/(I_1 \cap I_2)M$ est plat sur $R/(I_1 \cap I_2)$.
\end{lemma}

\begin{proof}
En remplaçant $R$ par $R/(I_1 \cap I_2)$ et $M$ par
$M/(I_1 \cap I_2)M$, nous pouvons supposer $I_1 \cap I_2 = 0$. Soit
$J \subset R$ un idéal. Pour démontrer que $M$ est plat sur $R$, nous devons
montrer que $J \otimes_R M \to M$ est injectif ; voir Algèbre,
lemme \ref{algebra-lemma-flat}. La platitude de $M/I_1M$ sur $R/I_1$
montre que l'homomorphisme
$$
J/(J \cap I_1) \otimes_R M =
(J + I_1)/I_1 \otimes_{R/I_1} M/I_1M
\longrightarrow M/I_1M
$$
est injectif. Puisque $0 \to (J \cap I_1) \to J \to J/(J \cap I_1) \to 0$
est exacte, nous obtenons un diagramme
$$
\xymatrix{
(J \cap I_1) \otimes_R M \ar[r] \ar[d] &
J \otimes_R M \ar[r] \ar[d] &
J/(J \cap I_1) \otimes_R M \ar[r] \ar[d] & 0 \\
M \ar@{=}[r] &
M \ar[r] &
M/I_1M
}
$$
Il suffit donc de montrer que $(J \cap I_1) \otimes_R M \to M$ est
injectif. Comme $I_1 \cap I_2 = 0$, l'idéal $J \cap I_1$ s'envoie
isomorphiquement sur un idéal $J' \subset R/I_2$, et l'on a
$(J \cap I_1) \otimes_R M = J' \otimes_{R/I_2} M/I_2M$. Par platitude de
$M/I_2M$ sur $R/I_2$, l'homomorphisme
$J' \otimes_{R/I_2} M/I_2M \to M/I_2M$ est injectif, ce qui implique
clairement que $(J \cap I_1) \otimes_R M \to M$ est injectif.
\end{proof}





\section{Platification sur un anneau artinien}
\label{section-flattening-artinian}

\noindent
Une platification universelle existe lorsque l'anneau de base est un anneau
local artinien. Elle existe pour un module quelconque. Par conséquent, comme
nous le verrons plus loin, une stratification par platification existe
lorsque le schéma de base est le spectre d'un anneau local artinien.

\begin{lemma}
\label{lemma-flattening-artinian}
Soit $R$ un anneau artinien. Soit $M$ un $R$-module. Alors il existe un plus
petit idéal $I \subset R$ tel que $M/IM$ soit plat sur $R/I$.
\end{lemma}

\begin{proof}
Cela résulte directement du lemme \ref{lemma-intersection-flat} et de la
propriété artinienne.
\end{proof}

\noindent
Cet idéal possède la propriété universelle suivante.

\begin{lemma}
\label{lemma-flattening-artinian-universal-property}
Soit $R$ un anneau artinien. Soit $M$ un $R$-module. Soit $I \subset R$ le
plus petit idéal $I \subset R$ tel que $M/IM$ soit plat sur $R/I$. Alors
$I$ possède la propriété universelle suivante : pour tout morphisme
d'anneaux $\varphi : R \to R'$, on a
$$
R' \otimes_R M\text{ est plat sur }R'
\Leftrightarrow
\text{on a }\varphi(I) = 0.
$$
\end{lemma}

\begin{proof}
Notons que $I$ existe d'après le lemme \ref{lemma-flattening-artinian}.
L'implication $\Rightarrow$ résulte d'Algèbre,
lemme \ref{algebra-lemma-flat-base-change}. Soit $\varphi : R \to R'$ tel
que $M \otimes_R R'$ soit plat sur $R'$. Soit $J = \Ker(\varphi)$. D'après
Algèbre,
lemme \ref{algebra-lemma-descent-flatness-injective-map-artinian-rings},
comme $R' \otimes_R M = R' \otimes_{R/J} M/JM$ est plat sur $R'$, nous en
concluons que $M/JM$ est plat sur $R/J$. Ainsi, $I \subset J$, comme voulu.
\end{proof}











\section{Platification au-dessus d'une partie fermée de la base}
\label{section-flattening-local-base}

\noindent
Soit $R \to S$ un morphisme d'anneaux. Soit $I \subset R$ un idéal. Soit
$M$ un $S$-module. Dans la suite, nous considérerons la condition suivante :
\begin{equation}
\label{equation-flat-at-primes-over}
\forall \mathfrak q \in V(IS) \subset \Spec(S) :
M_{\mathfrak q}\text{ est plat sur }R.
\end{equation}
Géométriquement, cela signifie que $M$ est plat sur $R$ le long de l'image
réciproque de $V(I)$ dans $\Spec(S)$. Si $R$ et $S$ sont des anneaux
noethériens et si $M$ est un $S$-module de type fini, alors
(\ref{equation-flat-at-primes-over}) équivaut à demander que $M/I^nM$ soit
plat sur $R/I^n$ pour tout $n \geq 1$ ; voir Algèbre,
lemme \ref{algebra-lemma-flat-module-powers}.

\begin{lemma}
\label{lemma-base-change-flat-at-primes-over}
Soit $R \to S$ un morphisme d'anneaux. Soit $I \subset R$ un idéal. Soit
$M$ un $S$-module. Soient $R \to R'$ un morphisme d'anneaux et
$IR' \subset I' \subset R'$ un idéal. Si
(\ref{equation-flat-at-primes-over}) est satisfaite pour
$(R \to S, I, M)$, alors (\ref{equation-flat-at-primes-over}) est satisfaite pour
$(R' \to S \otimes_R R', I', M \otimes_R R')$.
\end{lemma}

\begin{proof}
Supposons (\ref{equation-flat-at-primes-over}) satisfaite pour
$(R \to S, I \subset R, M)$. Soit
$I'(S \otimes_R R') \subset \mathfrak q'$ un idéal premier de
$S \otimes_R R'$. Soit $\mathfrak q \subset S$ l'idéal premier correspondant
de $S$. Alors $IS \subset \mathfrak q$. Remarquons que
$(M \otimes_R R')_{\mathfrak q'}$ est un localisé du changement de base
$M_{\mathfrak q} \otimes_R R'$. Par conséquent,
$(M \otimes_R R')_{\mathfrak q'}$ est plat sur $R'$, car c'est un localisé
d'un module plat ; voir Algèbre, lemmes
\ref{algebra-lemma-flat-base-change} et
\ref{algebra-lemma-flat-localization}.
\end{proof}

\begin{lemma}
\label{lemma-flat-descent-flat-at-primes-over}
Soit $R \to S$ un morphisme d'anneaux. Soit $I \subset R$ un idéal. Soit
$M$ un $S$-module. Soient $R \to R'$ un morphisme d'anneaux et
$IR' \subset I' \subset R'$ un idéal tels que
\begin{enumerate}
\item l'application $V(I') \to V(I)$ induite par
$\Spec(R') \to \Spec(R)$ soit surjective, et
\item $R'_{\mathfrak p'}$ soit plat sur $R$ pour tout idéal premier
$\mathfrak p' \in V(I')$.
\end{enumerate}
Si (\ref{equation-flat-at-primes-over}) est satisfaite pour
$(R' \to S \otimes_R R', I', M \otimes_R R')$, alors
(\ref{equation-flat-at-primes-over}) est satisfaite pour $(R \to S, I, M)$.
\end{lemma}

\begin{proof}
Supposons (\ref{equation-flat-at-primes-over}) satisfaite pour
$(R' \to S \otimes_R R', IR', M \otimes_R R')$. Choisissons un idéal
premier $IS \subset \mathfrak q \subset S$. Soit
$I \subset \mathfrak p \subset R$ l'idéal premier correspondant de $R$.
Par hypothèse, il existe un idéal premier $\mathfrak p' \in V(I')$ de $R'$
au-dessus de $\mathfrak p$, et $R_{\mathfrak p} \to R'_{\mathfrak p'}$ est
plat. Choisissons un idéal premier
$\overline{\mathfrak q}' \subset
\kappa(\mathfrak q) \otimes_{\kappa(\mathfrak p)} \kappa(\mathfrak p')$
qui corresponde à un idéal premier
$\mathfrak q' \subset S \otimes_R R'$ au-dessus de $\mathfrak q$ et de
$\mathfrak p'$. Remarquons que $(S \otimes_R R')_{\mathfrak q'}$ est un
localisé de $S_{\mathfrak q} \otimes_{R_{\mathfrak p}} R'_{\mathfrak p'}$.
Par hypothèse, le module $(M \otimes_R R')_{\mathfrak q'}$ est plat sur
$R'_{\mathfrak p'}$. Algèbre,
lemme \ref{algebra-lemma-base-change-flat-up-down}, implique donc que
$M_{\mathfrak q}$ est plat sur $R_{\mathfrak p}$, ce qu'il fallait
démontrer.
\end{proof}

\begin{lemma}
\label{lemma-limit-preserving-flat-at-primes-over}
Soit $R \to S$ un morphisme d'anneaux de présentation finie. Soit $M$ un
$S$-module de présentation finie. Soit
$R' = \colim_{\lambda \in \Lambda} R_\lambda$ une colimite filtrante de
$R$-algèbres. Soient $I_\lambda \subset R_\lambda$ des idéaux tels que
$I_\lambda R_\mu \subset I_\mu$ pour tout $\mu \geq \lambda$, et posons
$I' = \colim_\lambda I_\lambda$. Si
(\ref{equation-flat-at-primes-over}) est satisfaite pour
$(R' \to S \otimes_R R', I', M \otimes_R R')$, alors il existe
$\lambda \in \Lambda$ tel que
(\ref{equation-flat-at-primes-over}) soit satisfaite pour
$(R_\lambda \to S \otimes_R R_\lambda, I_\lambda,
M \otimes_R R_\lambda)$.
\end{lemma}

\begin{proof}
Notons $S_\lambda = S \otimes_R R_\lambda$,
$S' = S \otimes_R R'$, $M_\lambda = M \otimes_R R_\lambda$ et
$M' = M \otimes_R R'$. Le changement de base $S'$ est de présentation
finie sur $R'$, et $M'$ est de présentation finie sur $S'$ ; il en va de
même des versions indicées par $\lambda$ ; voir Algèbre,
lemme \ref{algebra-lemma-base-change-finiteness}. D'après Algèbre,
théorème \ref{algebra-theorem-openness-flatness}, l'ensemble
$$
U' = \{\mathfrak q' \in \Spec(S') \mid
M'_{\mathfrak q'}\text{ est plat sur }R'\}
$$
est ouvert dans $\Spec(S')$. Remarquons que $V(I'S')$ est un espace
quasi-compact contenu dans $U'$ par hypothèse. Il existe donc un nombre fini
d'éléments $g'_j \in S'$, $j = 1, \ldots, m$, tels que
$D(g'_j) \subset U'$ et $V(I'S') \subset \bigcup D(g'_j)$. En particulier,
$(M')_{g'_j}$ est un module plat sur $R'$.

\medskip\noindent
Nous allons choisir des éléments $\lambda \in \Lambda$ de plus en plus
grands. Choisissons-le d'abord assez grand pour qu'il existe des
$g_{j, \lambda} \in S_{\lambda}$ s'envoyant sur $g'_j$. L'inclusion
$V(I'S') \subset \bigcup D(g'_j)$ signifie que
$I'S' + (g'_1, \ldots, g'_m) = S'$, ce qui peut s'écrire
$1 = \sum z_sh_s + \sum f_jg'_j$ pour certains $z_s \in I'$ et
$h_s, f_j \in S'$. En agrandissant $\lambda$, nous pouvons supposer qu'une
telle égalité vaut dans $S_\lambda$. Nous pouvons donc supposer que
$V(I_\lambda S_\lambda) \subset \bigcup D(g_{j, \lambda})$. D'après
Algèbre, lemme \ref{algebra-lemma-flat-finite-presentation-limit-flat}, pour
$\lambda$ assez grand, les modules $(M_\lambda)_{g_{j, \lambda}}$ sont plats
sur $R_\lambda$. En particulier, le module $M_\lambda$ est plat sur
$R_\lambda$ en tout idéal premier au-dessus de l'idéal $I_\lambda$.
\end{proof}









\section{Platification au-dessus de parties fermées de la source et de la base}
\label{section-flattening-local-source-base}

\noindent
Dans cette section, nous généralisons légèrement l'étude de la
section \ref{section-flattening-local-base}. Nous conseillons vivement au
lecteur de commencer par lire et comprendre cette section.

\begin{situation}
\label{situation-flattening-general}
Soit $R \to S$ un morphisme d'anneaux. Soit $J \subset S$ un idéal. Soit
$M$ un $S$-module.
\end{situation}

\noindent
Dans cette situation, étant donnés une $R$-algèbre $R'$ et un idéal
$I' \subset R'$, nous posons $S' = S \otimes_R R'$ et
$M' = M \otimes_R R'$. Nous considérerons la condition
\begin{equation}
\label{equation-flat-at-primes}
\forall \mathfrak q' \in V(I'S' + JS') \subset \Spec(S') :
M'_{\mathfrak q'}\text{ est plat sur }R'.
\end{equation}
Géométriquement, cela signifie que $M'$ est plat sur $R'$ le long de
l'intersection de l'image réciproque de $V(I')$ et de l'image réciproque de
$V(J)$. Puisque $(R \to S, J, M)$ sont fixés, la condition
(\ref{equation-flat-at-primes}) ne dépend que du couple $(R', I')$, où $R'$
est considérée comme une $R$-algèbre.

\begin{lemma}
\label{lemma-base-change-flat-at-primes}
Dans la situation \ref{situation-flattening-general}, soit $R' \to R''$ un
morphisme de $R$-algèbres. Soient $I' \subset R'$ et
$I'R'' \subset I'' \subset R''$ des idéaux. Si
(\ref{equation-flat-at-primes}) est satisfaite pour $(R', I')$, alors
(\ref{equation-flat-at-primes}) est satisfaite pour $(R'', I'')$.
\end{lemma}

\begin{proof}
Supposons (\ref{equation-flat-at-primes}) satisfaite pour $(R', I')$. Soit
$I''S'' + JS'' \subset \mathfrak q''$ un idéal premier de $S''$. Soit
$\mathfrak q' \subset S'$ l'idéal premier correspondant de $S'$. Alors
$I'S' \subset \mathfrak q'$ et $JS' \subset \mathfrak q'$, puisque les
conditions correspondantes sont satisfaites pour $\mathfrak q''$.
Remarquons que $(M'')_{\mathfrak q''}$ est un localisé du changement de base
$M'_{\mathfrak q'} \otimes_R R''$. Par conséquent,
$(M'')_{\mathfrak q''}$ est plat sur $R''$, car c'est un localisé d'un
module plat ; voir Algèbre, lemmes \ref{algebra-lemma-flat-base-change} et
\ref{algebra-lemma-flat-localization}.
\end{proof}

\begin{lemma}
\label{lemma-flat-descent-flat-at-primes}
Dans la situation \ref{situation-flattening-general}, soit $R' \to R''$ un
morphisme de $R$-algèbres. Soient $I' \subset R'$ et
$I'R'' \subset I'' \subset R''$ des idéaux. Supposons que
\begin{enumerate}
\item l'application $V(I'') \to V(I')$ induite par
$\Spec(R'') \to \Spec(R')$ soit surjective, et
\item $R''_{\mathfrak p''}$ soit plat sur $R'$ pour tout idéal premier
$\mathfrak p'' \in V(I'')$.
\end{enumerate}
Si (\ref{equation-flat-at-primes}) est satisfaite pour $(R'', I'')$, alors
(\ref{equation-flat-at-primes}) est satisfaite pour $(R', I')$.
\end{lemma}

\begin{proof}
Supposons (\ref{equation-flat-at-primes}) satisfaite pour $(R'', I'')$.
Choisissons un idéal premier
$I'S' + JS' \subset \mathfrak q' \subset S'$. Soit
$I' \subset \mathfrak p' \subset R'$ l'idéal premier correspondant de
$R'$. Par hypothèse, il existe un idéal premier
$\mathfrak p'' \in V(I'')$ de $R''$ au-dessus de $\mathfrak p'$, et
$R'_{\mathfrak p'} \to R''_{\mathfrak p''}$ est plat. Choisissons un idéal
premier
$\overline{\mathfrak q}'' \subset
\kappa(\mathfrak q') \otimes_{\kappa(\mathfrak p')} \kappa(\mathfrak p'')$.
Celui-ci correspond à un idéal premier
$\mathfrak q'' \subset S'' = S' \otimes_{R'} R''$ au-dessus de
$\mathfrak q'$ et de $\mathfrak p''$. En particulier, on voit que
$I''S'' \subset \mathfrak q''$ et $JS'' \subset \mathfrak q''$.
Remarquons que $(S' \otimes_{R'} R'')_{\mathfrak q''}$ est un localisé de
$S'_{\mathfrak q'} \otimes_{R'_{\mathfrak p'}} R''_{\mathfrak p''}$. Par
hypothèse, le module $(M' \otimes_{R'} R'')_{\mathfrak q''}$ est plat sur
$R''_{\mathfrak p''}$. Algèbre,
lemme \ref{algebra-lemma-base-change-flat-up-down}, implique donc que
$M'_{\mathfrak q'}$ est plat sur $R'_{\mathfrak p'}$, ce qu'il fallait
démontrer.
\end{proof}

\begin{lemma}
\label{lemma-limit-preserving-flat-at-primes}
Dans la situation \ref{situation-flattening-general}, supposons que
$R \to S$ soit essentiellement de présentation finie et que $M$ soit un
$S$-module de présentation finie. Soit
$R' = \colim_{\lambda \in \Lambda} R_\lambda$ une colimite filtrante de
$R$-algèbres. Soient $I_\lambda \subset R_\lambda$ des idéaux tels que
$I_\lambda R_\mu \subset I_\mu$ pour tout $\mu \geq \lambda$, et posons
$I' = \colim_\lambda I_\lambda$. Si (\ref{equation-flat-at-primes}) est
satisfaite pour $(R', I')$, alors il existe $\lambda \in \Lambda$ tel
que (\ref{equation-flat-at-primes}) soit satisfaite pour
$(R_\lambda, I_\lambda)$.
\end{lemma}

\begin{proof}
Démontrons d'abord le lemme lorsque $R \to S$ est de présentation finie,
puis expliquons les modifications nécessaires dans le cas général. Notons
$S_\lambda = S \otimes_R R_\lambda$, $S' = S \otimes_R R'$,
$M_\lambda = M \otimes_R R_\lambda$ et $M' = M \otimes_R R'$. Le changement
de base $S'$ est de présentation finie sur $R'$, et $M'$ est de présentation
finie sur $S'$ ; il en va de même des versions indicées par $\lambda$ ; voir
Algèbre, lemme \ref{algebra-lemma-base-change-finiteness}. D'après Algèbre,
théorème \ref{algebra-theorem-openness-flatness}, l'ensemble
$$
U' = \{\mathfrak q' \in \Spec(S') \mid
M'_{\mathfrak q'}\text{ est plat sur }R'\}
$$
est ouvert dans $\Spec(S')$. Remarquons que $V(I'S' + JS')$ est un espace
quasi-compact contenu dans $U'$ par hypothèse. Il existe donc un nombre fini
d'éléments $g'_j \in S'$, $j = 1, \ldots, m$, tels que
$D(g'_j) \subset U'$ et $V(I'S' + JS') \subset \bigcup D(g'_j)$. En
particulier, $(M')_{g'_j}$ est un module plat sur $R'$.

\medskip\noindent
Nous allons choisir des éléments $\lambda \in \Lambda$ de plus en plus
grands. Choisissons-le d'abord assez grand pour qu'il existe des
$g_{j, \lambda} \in S_{\lambda}$ s'envoyant sur $g'_j$. L'inclusion
$V(I'S' + JS') \subset \bigcup D(g'_j)$ signifie que
$I'S' + JS' + (g'_1, \ldots, g'_m) = S'$, ce qui peut s'écrire
$$
1 = \sum y_tk_t + \sum z_sh_s + \sum f_jg'_j
$$
pour certains $z_s \in I'$, $y_t \in J$ et $k_t, h_s, f_j \in S'$. En
agrandissant $\lambda$, nous pouvons supposer qu'une telle égalité vaut dans
$S_\lambda$. Nous pouvons donc supposer que
$V(I_\lambda S_\lambda + J S_\lambda) \subset \bigcup D(g_{j, \lambda})$.
D'après Algèbre,
lemme \ref{algebra-lemma-flat-finite-presentation-limit-flat}, pour
$\lambda$ assez grand, les modules $(M_\lambda)_{g_{j, \lambda}}$ sont plats
sur $R_\lambda$. En particulier, le module $M_\lambda$ est plat sur
$R_\lambda$ en tous les idéaux premiers correspondant à des points de
$V(I_\lambda S_\lambda + J S_\lambda)$.

\medskip\noindent
Lorsque $S$ est essentiellement de présentation finie, nous pouvons écrire
$S = \Sigma^{-1}C$, où $R \to C$ est de présentation finie et
$\Sigma \subset C$ est une partie multiplicative. Nous pouvons aussi écrire
$M = \Sigma^{-1}N$ pour un certain $C$-module de présentation finie $N$ ;
voir Algèbre, lemme \ref{algebra-lemma-construct-fp-module}. Introduisons
alors $C_\lambda$, $C'$, $N_\lambda$, $N'$. Dans l'argument ci-dessus, nous
obtenons un ouvert $U' \subset \Spec(C')$ sur lequel $N'$ est plat sur $R'$.
L'hypothèse que (\ref{equation-flat-at-primes}) est satisfaite signifie que
$V(I'S' + JS')$ s'envoie dans $U'$, car, pour un idéal premier
$\mathfrak q' \subset S'$ correspondant à un idéal premier
$\mathfrak r' \subset C'$, on a $M'_{\mathfrak q'} = N'_{\mathfrak r'}$.
Il existe donc des $g'_j \in C'$ tels que $\bigcup D(g'_j)$ contienne
l'image de $V(I'S' + JS')$. La suite de la démonstration est exactement la
même que précédemment.
\end{proof}

\begin{lemma}
\label{lemma-flat-module-powers-variant}
Plaçons-nous dans la situation \ref{situation-flattening-general}. Soit
$I \subset R$ un idéal. Supposons que
\begin{enumerate}
\item $R$ soit un anneau noethérien,
\item $S$ soit un anneau noethérien,
\item $M$ soit un $S$-module de type fini, et
\item pour chaque $n \geq 1$ et tout idéal premier
$\mathfrak q \in V(J + IS)$, le module $(M/I^n M)_{\mathfrak q}$ soit plat
sur $R/I^n$.
\end{enumerate}
Alors (\ref{equation-flat-at-primes}) est satisfaite pour $(R, I)$ ;
c'est-à-dire que, pour tout idéal premier $\mathfrak q \in V(J + IS)$, le
localisé $M_{\mathfrak q}$ est plat sur $R$.
\end{lemma}

\begin{proof}
Soit $\mathfrak q \in V(J + IS)$. Algèbre,
lemme \ref{algebra-lemma-flat-module-powers}, appliqué à
$R \to S_{\mathfrak q}$ et à $M_{\mathfrak q}$, implique alors que
$M_{\mathfrak q}$ est plat sur $R$.
\end{proof}










\section{Platification sur un anneau local noethérien complet}
\label{section-flattening-Noetherian-complete-local}

\noindent
Les trois lemmes suivants donnent une démonstration entièrement algébrique
de l'existence de la stratification « locale » par platification lorsque la
base est un anneau local noethérien complet $R$ et que le module donné est
de type fini sur une $R$-algèbre de type fini $S$.

\begin{lemma}
\label{lemma-flattening-complete-local-noetherian}
Soit $R \to S$ un morphisme d'anneaux. Soit $M$ un $S$-module. Supposons que
\begin{enumerate}
\item $(R, \mathfrak m)$ soit un anneau local noethérien complet,
\item $S$ soit un anneau noethérien, et
\item $M$ soit de type fini sur $S$.
\end{enumerate}
Alors il existe un idéal $I \subset \mathfrak m$ tel que
\begin{enumerate}
\item $(M/IM)_{\mathfrak q}$ soit plat sur $R/I$ pour tout idéal premier
$\mathfrak q$ de $S/IS$ au-dessus de $\mathfrak m$, et
\item si $J \subset R$ est un idéal tel que $(M/JM)_{\mathfrak q}$ soit plat
sur $R/J$ pour tout idéal premier $\mathfrak q$ au-dessus de
$\mathfrak m$, alors $I \subset J$.
\end{enumerate}
Autrement dit, $I$ est le plus petit idéal de $R$ tel que
(\ref{equation-flat-at-primes-over}) soit satisfaite pour
$(\overline{R} \to \overline{S}, \overline{\mathfrak m}, \overline{M})$
où $\overline{R} = R/I$, $\overline{S} = S/IS$,
$\overline{\mathfrak m} = \mathfrak m/I$ et $\overline{M} = M/IM$.
\end{lemma}

\begin{proof}
Soit $J \subset R$ un idéal. Appliquons Algèbre,
lemme \ref{algebra-lemma-flat-module-powers}, au module $M/JM$ sur l'anneau
$R/J$. On voit alors que $(M/JM)_{\mathfrak q}$ est plat sur $R/J$ pour tout
idéal premier $\mathfrak q$ de $S/JS$ si et seulement si
$M/(J + \mathfrak m^n)M$ est plat sur $R/(J + \mathfrak m^n)$ pour tout
$n \geq 1$. Nous utiliserons cette remarque ci-dessous.

\medskip\noindent
Pour tout $n \geq 1$, l'anneau local $R/\mathfrak m^n$ est artinien. Le
lemme \ref{lemma-flattening-artinian} assure donc l'existence d'un plus petit
idéal $I_n \supset \mathfrak m^n$ tel que $M/I_nM$ soit plat sur $R/I_n$.
Il est clair que $I_{n + 1} + \mathfrak m^n$ contient $I_n$ et, en appliquant
le lemme \ref{lemma-intersection-flat}, on voit que
$I_n = I_{n + 1} + \mathfrak m^n$. Puisque
$R = \lim_n\ R/\mathfrak m^n$, on voit que
$I = \lim_n\ I_n/\mathfrak m^n$ est un idéal de $R$ tel que
$I_n = I + \mathfrak m^n$ pour tout $n \geq 1$. Les remarques initiales de
la démonstration montrent que $I$ vérifie (1) et (2). Nous omettons certains
détails.
\end{proof}

\begin{lemma}
\label{lemma-flattening-complete-local-noetherian-property-by-finite-type}
Reprenons les notations $R \to S$, $M$ et $I$, ainsi que les hypothèses du
lemme \ref{lemma-flattening-complete-local-noetherian}. Considérons un
homomorphisme local d'anneaux locaux
$\varphi : (R, \mathfrak m) \to (R', \mathfrak m')$
tel que $R'$ soit noethérien. Les assertions suivantes sont alors
équivalentes :
\begin{enumerate}
\item la condition (\ref{equation-flat-at-primes-over}) est satisfaite pour
$(R' \to S \otimes_R R', \mathfrak m', M \otimes_R R')$, et
\item $\varphi(I) = 0$.
\end{enumerate}
\end{lemma}

\begin{proof}
L'implication (2) $\Rightarrow$ (1) résulte du
lemme \ref{lemma-base-change-flat-at-primes-over}. Soit
$\varphi : R \to R'$ comme dans le lemme, satisfaisant (1). Nous devons
montrer que $\varphi(I) = 0$. Cela équivaut à $\varphi(I)R' = 0$. Le lemme
d'Artin-Rees dans l'anneau local noethérien $R'$ (voir Algèbre,
lemme \ref{algebra-lemma-intersect-powers-ideal-module-zero}) montre que
cette condition équivaut à
$\varphi(I)R' + (\mathfrak m')^n = (\mathfrak m')^n$ pour tout $n > 0$.
Elle équivaut donc à demander que la composée
$\varphi_n : R \to R' \to R'/(\mathfrak m')^n$ annule $I$ pour chaque $n$.
Or l'hypothèse (1) pour $\varphi$ implique l'hypothèse (1) pour $\varphi_n$
d'après le lemme \ref{lemma-base-change-flat-at-primes-over}. Nous sommes
ainsi ramenés au cas où $R'$ est local artinien.

\medskip\noindent
Supposons $R'$ artinien. Soit $J = \Ker(\varphi)$. Il faut montrer que
$I \subset J$. Par la construction de $I$ dans le
lemme \ref{lemma-flattening-complete-local-noetherian}, il suffit de montrer
que $(M/JM)_{\mathfrak q}$ est plat sur $R/J$ pour tout idéal premier
$\mathfrak q$ de $S/JS$ au-dessus de $\mathfrak m$. Puisque $R'$ est
artinien, la condition (1) signifie que $M \otimes_R R'$ est plat sur $R'$.
Comme $R'$ est artinien et que $R/J \to R'$ est un morphisme local injectif
d'anneaux, $R/J$ est également artinien. Par conséquent, la platitude de
$M \otimes_R R' = M/JM \otimes_{R/J} R'$ sur $R'$ implique celle de $M/JM$
sur $R/J$ d'après Algèbre,
lemme \ref{algebra-lemma-descent-flatness-injective-map-artinian-rings}.
Cela achève la démonstration.
\end{proof}

\begin{lemma}
\label{lemma-flattening-complete-local-universal-property}
Reprenons les notations $R \to S$, $M$ et $I$, ainsi que les hypothèses du
lemme \ref{lemma-flattening-complete-local-noetherian}. Supposons de plus que
$R \to S$ soit de type fini. Alors, pour tout homomorphisme local d'anneaux
locaux
$\varphi : (R, \mathfrak m) \to (R', \mathfrak m')$
les assertions suivantes sont équivalentes :
\begin{enumerate}
\item la condition (\ref{equation-flat-at-primes-over}) est satisfaite pour
$(R' \to S \otimes_R R', \mathfrak m', M \otimes_R R')$, et
\item $\varphi(I) = 0$.
\end{enumerate}
\end{lemma}

\begin{proof}
L'implication (2) $\Rightarrow$ (1) résulte du
lemme \ref{lemma-base-change-flat-at-primes-over}. Soit
$\varphi : R \to R'$ comme dans le lemme, satisfaisant (1). Comme $R$ est
noethérien, $R \to S$ est de présentation finie et $M$ est un $S$-module de
présentation finie. Écrivons $R' = \colim_\lambda R_\lambda$ comme colimite
filtrante de sous-$R$-algèbres locales $R_\lambda \subset R'$, d'idéaux
maximaux $\mathfrak m_\lambda = R_\lambda \cap \mathfrak m'$, tels que chaque
$R_\lambda$ soit essentiellement de type fini sur $R$. D'après le
lemme \ref{lemma-limit-preserving-flat-at-primes-over}, la condition
(\ref{equation-flat-at-primes-over}) est satisfaite pour
$(R_\lambda \to S \otimes_R R_\lambda, \mathfrak m_\lambda,
M \otimes_R R_\lambda)$ pour un certain $\lambda$. Le
lemme \ref{lemma-flattening-complete-local-noetherian-property-by-finite-type}
s'applique donc au morphisme d'anneaux $R \to R_\lambda$, et l'on voit que
$I$ s'envoie sur zéro dans $R_\lambda$ ; à plus forte raison, il s'envoie
sur zéro dans $R'$.
\end{proof}




\section{Descente de la platitude le long de morphismes intégraux}
\label{section-descent-flatness-integral}

\noindent
Commençons par quelques lemmes simples.

\begin{lemma}
\label{lemma-have-one-root}
Soit $R$ un anneau. Soit $P(T)$ un polynôme unitaire à coefficients dans
$R$. Soit $\alpha \in R$ tel que $P(\alpha) = 0$. Alors
$P(T) = (T - \alpha)Q(T)$ pour un certain polynôme unitaire
$Q(T) \in R[T]$.
\end{lemma}

\begin{proof}
Par récurrence sur le degré de $P$. Si $\deg(P) = 1$, alors
$P(T) = T - \alpha$ et le résultat est vrai. Si $\deg(P) > 1$, la division
euclidienne permet d'écrire
$P(T) = (T - \alpha)Q(T) + r$ pour un polynôme
$Q \in R[T]$ de degré $< \deg(P)$ et un certain $r \in R$. Par hypothèse,
$0 = P(\alpha) = (\alpha - \alpha)Q(\alpha) + r = r$, et nous obtenons
$r = 0$, comme voulu.
\end{proof}

\begin{lemma}
\label{lemma-adjoin-one-root}
Soit $R$ un anneau. Soit $P(T)$ un polynôme unitaire à coefficients dans
$R$. Il existe un morphisme d'anneaux fini et libre $R \to R'$ tel que
$P(T) = (T - \alpha)Q(T)$ pour un certain $\alpha \in R'$ et un certain
polynôme unitaire $Q(T) \in R'[T]$.
\end{lemma}

\begin{proof}
Écrivons $P(T) = T^d + a_1T^{d - 1} + \ldots + a_0$. Posons
$R' = R[x]/(x^d + a_1x^{d - 1} + \ldots + a_0)$. Soit $\alpha$ la classe
de congruence de $x$. Il est alors clair que $P(\alpha) = 0$. Le résultat
résulte donc du lemme \ref{lemma-have-one-root}.
\end{proof}

\begin{lemma}
\label{lemma-finite-split}
Soit $R \to S$ un morphisme d'anneaux fini. Il existe une extension
d'anneaux finie et libre $R \subset R'$ telle que $S \otimes_R R'$ soit un
quotient d'un anneau de la forme
$$
R'[T_1, \ldots, T_n]/(P_1(T_1), \ldots, P_n(T_n))
$$
où $P_i(T) = \prod_{j = 1, \ldots, d_i} (T - \alpha_{ij})$ pour certains
$\alpha_{ij} \in R'$.
\end{lemma}

\begin{proof}
Soient $x_1, \ldots, x_n \in S$ des générateurs de $S$ sur $R$. Pour chaque
$i$, nous pouvons choisir un polynôme unitaire $P_i(T) \in R[T]$ tel que
$P_i(x_i) = 0$ dans $S$ ; voir Algèbre, lemme
\ref{algebra-lemma-finite-is-integral}. Écrivons $\deg(P_i) = d_i$. D'après
le lemme \ref{lemma-adjoin-one-root} (appliqué $\sum d_i$ fois), il existe
une extension d'anneaux finie et libre $R \subset R'$ telle que chaque
$P_i$ se scinde complètement :
$$
P_i(T) = \prod\nolimits_{j = 1, \ldots, d_i} (T - \alpha_{ij})
$$
pour certains $\alpha_{ik} \in R'$. Soit
$R'[T_1, \ldots, T_n] \to S \otimes_R R'$ l'homomorphisme de
$R'$-algèbres qui envoie $T_i$ sur $x_i \otimes 1$. Comme il envoie
$P_i(T_i)$ sur zéro, il induit la surjection cherchée.
\end{proof}

\begin{lemma}
\label{lemma-split-image}
Soit $R$ un anneau. Soit $S = R[T_1, \ldots, T_n]/J$. Supposons que $J$
contienne des éléments de la forme $P_i(T_i)$, avec
$P_i(T) = \prod_{j = 1, \ldots, d_i} (T - \alpha_{ij})$ pour certains
$\alpha_{ij} \in R$. Pour $\underline{k} = (k_1, \ldots, k_n)$ tel que
$1 \leq k_i \leq d_i$, considérons l'homomorphisme d'anneaux
$$
\Phi_{\underline{k}} : R[T_1, \ldots, T_n] \to R,
\quad
T_i \longmapsto \alpha_{ik_i}
$$
Posons $J_{\underline{k}} = \Phi_{\underline{k}}(J)$. Alors l'image de
$\Spec(S) \to \Spec(R)$ est égale à $V(\bigcap J_{\underline{k}})$.
\end{lemma}

\begin{proof}
Ce lemme est immédiat. Indication :
$V(\bigcap J_{\underline{k}}) = \bigcup V(J_{\underline{k}})$.
\end{proof}

\noindent
Le résultat suivant est dû à Ferrand ; voir \cite{Ferrand}.

\begin{lemma}
\label{lemma-descent-flatness-injective-finite-Noetherian-rings}
Soit $R \to S$ un homomorphisme injectif fini d'anneaux noethériens. Soit
$M$ un $R$-module. Si $M \otimes_R S$ est un $S$-module plat, alors $M$ est
un $R$-module plat.
\end{lemma}

\begin{proof}
Soit $M$ un $R$-module tel que $M \otimes_R S$ soit plat sur $S$. D'après
Algèbre, lemme \ref{algebra-lemma-flatness-descends}, pour démontrer que $M$
est plat, nous pouvons remplacer $R$ par une extension d'anneaux
fidèlement plate quelconque. D'après le lemme \ref{lemma-finite-split}, il
existe une extension d'anneaux finie et localement libre $R \subset R'$ telle
que $S' = S \otimes_R R' = R'[T_1, \ldots, T_n]/J$ pour un idéal
$J \subset R'[T_1, \ldots, T_n]$ contenant des éléments de la forme
$P_i(T_i)$, avec $P_i(T) = \prod_{j = 1, \ldots, d_i}
(T - \alpha_{ij})$ pour certains $\alpha_{ij} \in R'$. Remarquons que
$R'$ est noethérien et que $R' \subset S'$ est une extension finie
d'anneaux. Nous pouvons donc remplacer $R$ par $R'$ et supposer que $S$
admet une présentation comme dans le lemme \ref{lemma-split-image}.
Remarquons que $\Spec(S) \to \Spec(R)$ est surjectif ; voir Algèbre,
lemme \ref{algebra-lemma-integral-overring-surjective}. Ainsi, en utilisant
le lemme \ref{lemma-split-image}, nous concluons que
$I = \bigcap J_{\underline{k}}$ est un idéal tel que
$V(I) = \Spec(R)$. Cela signifie que $I \subset \sqrt{(0)}$ et, puisque
$R$ est noethérien, que $I$ est nilpotent. Les applications
$\Phi_{\underline{k}}$ induisent des
diagrammes commutatifs
$$
\xymatrix{
S \ar[rr] & & R/J_{\underline{k}} \\
& R \ar[lu] \ar[ru]
}
$$
desquels nous concluons que $M/J_{\underline{k}}M$ est plat sur
$R/J_{\underline{k}}$. D'après le lemme \ref{lemma-intersection-flat}, on
voit que $M/IM$ est plat sur $R/I$. Enfin, en appliquant Algèbre,
lemme \ref{algebra-lemma-lift-flatness}, nous concluons que $M$ est plat sur
$R$.
\end{proof}

\begin{lemma}
\label{lemma-descent-flatness-injective-integral}
Soit $R \to S$ un morphisme d'anneaux entier injectif. Soit $M$ un module
de présentation finie sur $R[x_1, \ldots, x_n]$. Si $M \otimes_R S$ est
plat sur $S$, alors $M$ est plat sur $R$.
\end{lemma}

\begin{proof}
Choisissons une présentation
$$
R[x_1, \ldots, x_n]^{\oplus t} \to R[x_1, \ldots, x_n]^{\oplus r} \to M \to 0.
$$
Disons que le premier homomorphisme est donné par la matrice
$r \times t$ $T = (f_{ij})$, avec $f_{ij} \in R[x_1, \ldots, x_n]$.
Écrivons $f_{ij} = \sum f_{ij, I} x^I$ avec $f_{ij, I} \in R$ (notation
multi-indice). Considérons les diagrammes
$$
\xymatrix{
R \ar[r] & S \\
R_\lambda \ar[u] \ar[r] & S_\lambda \ar[u]
}
$$
où $R_\lambda$ est une sous-algèbre de type fini sur $\mathbf{Z}$ de $R$
contenant tous les $f_{ij, I}$ et où $S_\lambda$ est une
$R_\lambda$-sous-algèbre finie de $S$. Soit $M_\lambda$ le
$R_\lambda[x_1, \ldots, x_n]$-module de type fini défini par une
présentation comme ci-dessus, en utilisant la même matrice $T$, maintenant
considérée comme matrice sur $R_\lambda[x_1, \ldots, x_n]$. Remarquons que
$S$ est la colimite filtrante des $S_\lambda$ (détails omis). D'après
Algèbre, lemme \ref{algebra-lemma-flat-finite-presentation-limit-flat}, il
existe un $\lambda$ tel que le module
$M_\lambda \otimes_{R_\lambda} S_\lambda$ soit plat sur $S_\lambda$. Le
lemme \ref{lemma-descent-flatness-injective-finite-Noetherian-rings} implique
alors que $M_\lambda$ est plat sur $R_\lambda$. Comme
$M = M_\lambda \otimes_{R_\lambda} R$, le résultat découle d'Algèbre,
lemme \ref{algebra-lemma-flat-base-change}.
\end{proof}

\begin{lemma}
\label{lemma-descent-projective-injective-finite-Noetherian-rings}
Soit $R \to S$ un homomorphisme injectif fini d'anneaux noethériens. Soit
$P$ un $R$-module. Si $P \otimes_R S$ est un $S$-module projectif, alors
$P$ est un $R$-module projectif.
\end{lemma}

\begin{proof}
Soit $P$ un $R$-module tel que $P \otimes_R S$ soit projectif sur $S$. D'après
Algèbre, théorème \ref{algebra-theorem-ffdescent-projectivity}, pour
démontrer que $P$ est projectif, nous pouvons remplacer $R$ par une
extension d'anneaux fidèlement plate quelconque. D'après le lemme
\ref{lemma-finite-split}, il existe une extension d'anneaux finie et
localement libre $R \subset R'$ telle que
$S' = S \otimes_R R' = R'[T_1, \ldots, T_n]/J$ pour un idéal
$J \subset R'[T_1, \ldots, T_n]$ contenant des éléments de la forme
$P_i(T_i)$, avec $P_i(T) = \prod_{j = 1, \ldots, d_i}
(T - \alpha_{ij})$ pour certains $\alpha_{ij} \in R'$. Remarquons que
$R'$ est noethérien et que $R' \subset S'$ est une extension finie
d'anneaux. Nous pouvons donc remplacer $R$ par $R'$ et supposer que $S$
admet une présentation comme dans le lemme \ref{lemma-split-image}.
Remarquons que $\Spec(S) \to \Spec(R)$ est surjectif ; voir Algèbre,
lemme \ref{algebra-lemma-integral-overring-surjective}. Ainsi, en utilisant
le lemme \ref{lemma-split-image}, nous concluons que
$I = \bigcap J_{\underline{k}}$ est un idéal tel que
$V(I) = \Spec(R)$. Cela signifie que $I \subset \sqrt{(0)}$ et, puisque
$R$ est noethérien, que $I$ est nilpotent. Les applications
$\Phi_{\underline{k}}$ induisent des
diagrammes commutatifs
$$
\xymatrix{
S \ar[rr] & & R/J_{\underline{k}} \\
& R \ar[lu] \ar[ru]
}
$$
desquels nous concluons que $P/J_{\underline{k}}P$ est projectif sur
$R/J_{\underline{k}}$. D'après Algèbre,
lemme \ref{algebra-lemma-projective-modulo-two-ideals}, on voit que $P/IP$
est projectif sur $R/I$. Comme $P$ est plat sur $R$ d'après le
lemme \ref{lemma-descent-flatness-injective-finite-Noetherian-rings}, nous
pouvons appliquer Algèbre, lemme \ref{algebra-lemma-lift-projective}, pour
conclure que $P$ est projectif sur $R$.
\end{proof}









\section{Modules sans torsion}
\label{section-torsion-flat}

\noindent
Dans cette section, nous étudions les modules sans torsion et leur relation
avec la platitude (en particulier sur les anneaux de dimension 1).

\begin{definition}
\label{definition-torsion}
Soit $R$ un anneau intègre. Soit $M$ un $R$-module.
\begin{enumerate}
\item Nous disons qu'un élément $x \in M$ est de la {\it torsion} s'il
existe un $f \in R$ non nul tel que $fx = 0$.
\item Nous disons que $M$ est {\it sans torsion} si le seul élément de $M$
qui soit de la torsion est $0$.
\item Nous disons que $M$ est un {\it module de torsion} si tout élément de
$M$ est de la torsion.
\end{enumerate}
\end{definition}

\noindent
Soit $R$ un anneau intègre et soit $S = R \setminus \{0\}$ la partie
multiplicative des éléments non nuls de $R$. Alors un $R$-module $M$ est
sans torsion si et seulement si $M \to S^{-1}M$ est injectif. Autrement dit,
si et seulement si l'homomorphisme $M \to M \otimes_R K$ est injectif, où
$K = S^{-1}R$ est le corps des fractions de $R$.

\begin{lemma}
\label{lemma-torsion}
Soit $R$ un anneau intègre. Soit $M$ un $R$-module. L'ensemble des éléments
de torsion $M_{tors}$ de $M$ est le noyau de l'homomorphisme
$M \to M \otimes_R K$. Ainsi, $M_{tors}$ est un sous-module $R$-de $M$.
Le module quotient $M/M_{tors}$ est sans torsion.
\end{lemma}

\begin{proof}
Voir la discussion ci-dessus.
\end{proof}

\begin{lemma}
\label{lemma-localize-torsion}
Soit $R$ un anneau intègre. Soit $M$ un $R$-module sans torsion. Pour toute
partie multiplicative $S \subset R$, le module $S^{-1}M$ est un
$S^{-1}R$-module sans torsion.
\end{lemma}

\begin{proof}
Démonstration omise.
\end{proof}

\begin{lemma}
\label{lemma-flat-pullback-torsion}
Soit $R \to R'$ un homomorphisme plat d'anneaux intègres. Si $M$ est un
$R$-module sans torsion, alors $M \otimes_R R'$ est un $R'$-module sans
torsion.
\end{lemma}

\begin{proof}
Si $M$ est sans torsion, alors $M \subset M \otimes_R K$ est injectif, où
$K$ est le corps des fractions de $R$. Comme $R'$ est plat sur $R$,
$M \otimes_R R' \to (M \otimes_R K) \otimes_R R'$ est injectif. Comme
$M \otimes_R K$ est isomorphe à une somme directe de copies de $K$, il
suffit de voir que $K \otimes_R R'$ est sans torsion. Cela est vrai car il
est un localisé de $R'$.
\end{proof}

\begin{lemma}
\label{lemma-extension-torsion-free}
Soit $R$ un anneau intègre. Soit
$0 \to M \to M' \to M'' \to 0$ une suite exacte courte de $R$-modules. Si
$M$ et $M''$ sont sans torsion, alors $M'$ est sans torsion.
\end{lemma}

\begin{proof}
Démonstration omise.
\end{proof}

\begin{lemma}
\label{lemma-check-torsion}
Soit $R$ un anneau intègre. Soit $M$ un $R$-module. Alors $M$ est sans
torsion si et seulement si $M_\mathfrak m$ est un $R_\mathfrak m$-module
sans torsion pour tout idéal maximal $\mathfrak m$ de $R$.
\end{lemma}

\begin{proof}
Démonstration omise. Indication : utiliser le lemme
\ref{lemma-localize-torsion} et Algèbre, lemme
\ref{algebra-lemma-characterize-zero-local}.
\end{proof}

\begin{lemma}
\label{lemma-finite-torsion-free-submodule-free}
Soit $R$ un anneau intègre. Soit $M$ un $R$-module de type fini. Alors $M$
est sans torsion si et seulement si $M$ est un sous-module d'un module libre
de type fini.
\end{lemma}

\begin{proof}
Si $M$ est un sous-module de $R^{\oplus n}$, alors $M$ est sans torsion.
Réciproquement, supposons $M$ sans torsion. Soit $K$ le corps des fractions
de $R$. Alors $M \otimes_R K$ est un espace vectoriel de dimension finie
sur $K$. Choisissons une base $e_1, \ldots, e_r$ de cet espace vectoriel.
Soient $x_1, \ldots, x_n$ des générateurs de $M$. Écrivons
$x_i = \sum (a_{ij}/b_{ij}) e_j$ pour certains $a_{ij}, b_{ij} \in R$ avec
$b_{ij} \not = 0$. Posons $b = \prod_{i, j} b_{ij}$. Comme $M$ est sans
torsion, l'homomorphisme $M \to M \otimes_R K$ est injectif et son image
est contenue dans
$R^{\oplus r} = R e_1/b \oplus \ldots \oplus Re_r/b$.
\end{proof}

\begin{lemma}
\label{lemma-torsion-free-finite-noetherian-domain}
Soit $R$ un anneau intègre noethérien. Soit $M$ un $R$-module de type fini
non nul. Les assertions suivantes sont équivalentes :
\begin{enumerate}
\item $M$ est sans torsion,
\item $M$ est un sous-module d'un module libre de type fini,
\item $(0)$ est le seul idéal premier associé à $M$,
\item $(0)$ appartient au support de $M$ et $M$ possède la propriété $(S_1)$,
et
\item $(0)$ appartient au support de $M$ et $M$ n'a aucun idéal premier
associé immergé.
\end{enumerate}
\end{lemma}

\begin{proof}
Nous avons vu l'équivalence de (1) et (2) dans le lemme
\ref{lemma-finite-torsion-free-submodule-free}. Nous avons vu l'équivalence
de (4) et (5) dans Algèbre, lemme
\ref{algebra-lemma-criterion-no-embedded-primes}. L'équivalence de (3) et
(5) est immédiate d'après la définition. Le localisé d'un module sans
torsion est sans torsion (lemme \ref{lemma-localize-torsion}); il est donc
clair que $M$ n'a aucun idéal premier associé différent de $(0)$. Ainsi, (1)
implique (5). Réciproquement, supposons (5). Si $M$ possède de la torsion,
il existe une immersion $R/I \subset M$ pour un idéal non nul $I$ de $R$.
Ainsi, $M$ possède un idéal premier associé différent de $(0)$ (voir
Algèbre, lemmes \ref{algebra-lemma-ass} et \ref{algebra-lemma-ass-zero}).
C'est un idéal premier associé immergé, ce qui contredit l'hypothèse.
\end{proof}

\begin{lemma}
\label{lemma-flat-torsion-free}
Soit $R$ un anneau intègre. Tout $R$-module plat est sans torsion.
\end{lemma}

\begin{proof}
Si $x \in R$ est non nul, alors $x : R \to R$ est injectif, et donc, si
$M$ est plat sur $R$, $x : M \to M$ est injectif. Ainsi, si $M$ est plat
sur $R$, alors $M$ est sans torsion.
\end{proof}

\begin{lemma}
\label{lemma-valuation-ring-torsion-free-flat}
Soit $A$ un anneau de valuation. Un $A$-module $M$ est plat sur $A$ si et
seulement si $M$ est sans torsion.
\end{lemma}

\begin{proof}
L'implication « plat $\Rightarrow$ sans torsion » est le lemme
\ref{lemma-flat-torsion-free}. Réciproquement, supposons $M$ sans torsion.
D'après le critère équationnel de platitude (voir Algèbre, lemme
\ref{algebra-lemma-flat-eq}), nous devons montrer que toute relation dans $M$
est triviale. Supposons donc
$\sum_{i = 1, \ldots, n} a_i x_i = 0$ avec $x_i \in M$ et $a_i \in A$.
Après renumérotation, nous pouvons supposer que $v(a_1) \leq v(a_i)$ pour
tout $i$. Nous pouvons donc écrire $a_i = a'_i a_1$ pour un certain
$a'_i \in A$. Notons que $a'_1 = 1$. Comme $M$ est sans torsion, on voit
que $x_1 = - \sum_{i \geq 2} a'_i x_i$. Ainsi, en choisissant
$y_i = x_i$, $i = 2, \ldots, n$, on obtient
$$
x_1 = \sum\nolimits_{j \geq 2} -a'_j y_j, \quad
x_i = y_i, (i \geq 2)\quad
0 = a_1 \cdot (-a'_j) + a_j \cdot 1 (j \geq 2)
$$
ce qui montre que la relation est triviale (explicitement, les éléments
$a_{ij}$ sont définis par $a_{11} = 0$, $a_{1j} = -a'_j$ pour $j > 1$ et
$a_{ij} = \delta_{ij}$ pour $i, j \geq 2$).
\end{proof}

\begin{lemma}
\label{lemma-dedekind-torsion-free-flat}
Soit $A$ un anneau de Dedekind (par exemple un anneau de valuation discrète
ou, plus généralement, un anneau principal).
\begin{enumerate}
\item Un $A$-module est plat si et seulement s'il est sans torsion.
\item Un $A$-module de type fini sans torsion est localement libre de type
fini.
\item Un $A$-module de type fini sans torsion est libre de type fini si
$A$ est un anneau principal.
\end{enumerate}
\end{lemma}

\begin{proof}
(Pour la remarque entre parenthèses dans l'énoncé du lemme, voir Algèbre,
lemme \ref{algebra-lemma-PID-dedekind}.) Démonstration de (1). D'après le
lemme \ref{lemma-check-torsion} et Algèbre, lemme
\ref{algebra-lemma-flat-localization}, il suffit de vérifier l'assertion sur
$A_\mathfrak m$ pour $\mathfrak m \subset A$ maximal. Comme $A_\mathfrak m$
est un anneau de valuation discrète (Algèbre, lemme
\ref{algebra-lemma-characterize-Dedekind}), le résultat découle du lemme
\ref{lemma-valuation-ring-torsion-free-flat}.

\medskip\noindent
Démonstration de (2). Cela résulte d'Algèbre, lemme
\ref{algebra-lemma-finite-projective}, et de (1).

\medskip\noindent
Démonstration de (3). Soit $A$ un anneau principal et soit $M$ un module de
type fini sans torsion. D'après le lemme
\ref{lemma-finite-torsion-free-submodule-free}, on a $M \subset A^{\oplus n}$
pour un certain $n$. Montrons que $M$ est libre par récurrence sur $n$. Le
cas $n = 1$ exprime précisément le fait que $A$ est principal. Si $n > 1$,
soit $M' \subset A^{\oplus n - 1}$ l'image de la projection sur les
$n - 1$ derniers facteurs de $A^{\oplus n}$. Nous obtenons alors une suite
exacte courte $0 \to I \to M \to M' \to 0$, où $I$ est l'intersection de
$M$ avec le premier facteur $A$ de $A^{\oplus n}$. Par récurrence, $M$ est
une extension de $A$-modules libres de type fini, donc est libre de type fini.
\end{proof}

\begin{lemma}
\label{lemma-hom-into-torsion-free}
Soit $R$ un anneau intègre. Soient $M$, $N$ des $R$-modules. Si $N$ est
sans torsion, alors $\Hom_R(M, N)$ est sans torsion.
\end{lemma}

\begin{proof}
Choisissons une surjection $\bigoplus_{i \in I} R \to M$. Alors
$\Hom_R(M, N) \subset \prod_{i \in I} N$.
\end{proof}




\section{Rangs des modules}
\label{section-rank}

\noindent
Voici notre définition.

\begin{definition}
\label{definition-rank}
Soit $R$ un anneau intègre de corps des fractions $K$. Le {\it rang} d'un
$R$-module $M$ est $\dim_K(M \otimes_R K)$.
\end{definition}

\noindent
Cette définition n'entre pas en conflit avec la notion de module localement
libre de rang $r$ d'Algèbre, définition
\ref{algebra-definition-locally-free}, lorsque les deux définitions
s'appliquent.

\begin{lemma}
\label{lemma-torsion-does-not-matter-rank}
Soit $R$ un anneau intègre. Si $M \to M'$ est un homomorphisme de
$R$-modules dont le noyau et le conoyau sont de torsion, alors le rang de
$M$ est égal au rang de $M'$.
\end{lemma}

\begin{proof}
Démonstration omise. Indication : l'homomorphisme induit
$M \otimes_R K \to M' \otimes_R K$ est un isomorphisme lorsque $K$ est le
corps des fractions de $R$.
\end{proof}

\begin{lemma}
\label{lemma-rank-additive}
Soit $R$ un anneau intègre. Soit
$0 \to M \to M' \to M'' \to 0$ une suite exacte courte de $R$-modules.
Alors le rang de $M'$ est la somme des rangs de $M$ et $M''$.
\end{lemma}

\begin{proof}
Démonstration omise.
\end{proof}

\begin{lemma}
\label{lemma-rank-tensor}
Soit $R$ un anneau intègre. Soient $M$ et $N$ des $R$-modules.
\begin{enumerate}
\item Le rang de $M \otimes_R N$ est le produit des rangs de $M$ et $N$.
\item Si $M$ est un $R$-module de présentation finie, alors le rang de
$\Hom_R(M, N)$ est le produit des rangs de $M$ et $N$.
\end{enumerate}
\end{lemma}

\begin{proof}
L'assertion (1) résulte du fait analogue pour les espaces vectoriels.
Supposons que $M$ soit de présentation finie comme $R$-module. Alors
$\Hom_R(M, N) \otimes_R K$ est isomorphe à
$\Hom_K(M \otimes_R K, N \otimes_R K)$ comme espace vectoriel sur $K$,
d'après Algèbre, lemme \ref{algebra-lemma-hom-from-finitely-presented}.
Comme $M \otimes_R K$ est également de dimension finie, on conclut par le
fait correspondant pour les espaces vectoriels.
\end{proof}

\begin{lemma}
\label{lemma-dominant-pullback-rank}
Soit $R \subset R'$ une extension d'anneaux intègres. Si $M$ est un
$R$-module, alors le rang de $M$ sur $R$ est égal au rang de
$M \otimes_R R'$ sur $R'$.
\end{lemma}

\begin{proof}
Cela est vrai car la dimension d'un espace vectoriel est invariante par
extension du corps de base.
\end{proof}

\begin{lemma}
\label{lemma-finite-module-rank}
Soit $R$ un anneau intègre. Soit $M$ un $R$-module de type fini. Alors
\begin{enumerate}
\item $M$ est de rang fini $r \geq 0$,
\item il existe un homomorphisme $M \to R^{\oplus r}$ dont le noyau et le
conoyau sont des modules de torsion,
\item il existe un $f \in R$, $f \not = 0$, tel que $M_f$ soit libre de rang
$r$, et
\item il existe un homomorphisme injectif $R^{\oplus r} \to M$ dont le
conoyau est un module de torsion.
\end{enumerate}
\end{lemma}

\begin{proof}
Comme $M$ est de type fini, $M \otimes_R K$ est un espace vectoriel de
dimension finie sur le corps des fractions $K$ de $R$. Le rang $r$ est donc
fini.

\medskip\noindent
Choisissons des générateurs $x_1, \ldots, x_n$ de $M$. Choisissons une
base $e_1, \ldots, e_r$ de $M \otimes_R K$. Nous pouvons écrire
$x_i \otimes 1 = \sum (a_{ij}/b_{ij}) e_j$ pour certains
$a_{ij} , b_{ij} \in R$ avec $b_{ij} \not = 0$. En posant
$b = \prod b_{ij}$ et en remplaçant $a_{ij}$ par $a_{ij}b/b_{ij}$, nous
pouvons supposer que $x_i \otimes 1 = \sum (a_{ij}/b) e_j$. Remarquons que
si nous avons une relation $\sum r_i x_i = 0$ dans $M$, alors
$\sum r_ia_{ij} = 0$ dans $R$ pour tout $j = 1, \ldots, r$.

\medskip\noindent
Démonstration de (2). La dernière remarque montre que l'homomorphisme
$M \to \bigoplus Re_j$ qui envoie $x_i$ sur $\sum a_{ij} e_j$ est bien
défini. Le fait que $e_1, \ldots, e_r$ forment une base de
$M \otimes_R K$ montre que cet homomorphisme devient un isomorphisme après
tensorisation par $K$. C'est donc un homomorphisme comme dans (2).

\medskip\noindent
Démonstration de (3). Le conoyau de l'homomorphisme $M \to R^{\oplus r}$
construit en (2) est de type fini et de torsion. Il existe donc un
$g \in R$, $g \not = 0$, qui annule ce conoyau. Autrement dit,
$M_g \to R_g^{\oplus r}$ est surjectif. Nous pouvons alors choisir une
section $M_g = N \oplus R_g^{\oplus r}$. Le module $N$ est un
$R_g$-module de type fini tel que $N \otimes_{R_g} K = 0$. Autrement dit,
$N$ est de torsion et de type fini sur $R_g$. Il existe donc un $g' \in R$
qui annule $N$. En posant $f= gg'$, on voit que (3) est satisfaite.

\medskip\noindent
Pour démontrer (4), choisissons une base $y_1, \ldots, y_r$ de $M_f$, où
$f$ est comme dans (3). Écrivons alors $y_j = m_j/f^{t_j}$ pour un certain
$m_j \in M$ et $t_j \geq 0$. L'homomorphisme $R^{\oplus r} \to M$ qui
envoie le $j$-ième vecteur de base sur $m_j$ est celui de (4). Détails omis.
\end{proof}




\section{Modules réflexifs}
\label{section-reflexive}

\noindent
Voici notre définition.

\begin{definition}
\label{definition-reflexive}
Soit $R$ un anneau intègre. Nous disons qu'un $R$-module $M$ est
{\it réflexif} si l'homomorphisme naturel
$$
j : M \longrightarrow \Hom_R(\Hom_R(M, R), R)
$$
qui envoie $m \in M$ sur l'homomorphisme qui envoie
$\varphi \in \Hom_R(M, R)$ sur $\varphi(m) \in R$ est un isomorphisme.
\end{definition}

\noindent
On peut donner cette définition pour des anneaux plus généraux, mais la
définition ci-dessus présente déjà des inconvénients. Il serait judicieux de
se restreindre aux anneaux intègres noethériens et aux modules de type fini
sans torsion et (peut-être) d'imposer à $R$ certaines conditions de
régularité (par exemple, que $R$ soit normal).

\begin{lemma}
\label{lemma-reflexive-torsion-free}
Soit $R$ un anneau intègre et soit $M$ un $R$-module.
\begin{enumerate}
\item Si $M$ est réflexif, alors $M$ est sans torsion.
\item Si $M$ est de type fini, alors le noyau et le conoyau de
$j : M \to \Hom_R(\Hom_R(M, R), R)$ sont des modules de torsion.
\item Si $M$ est de type fini, alors $j$ est injectif si et seulement si
$M$ est sans torsion.
\end{enumerate}
\end{lemma}

\begin{proof}
Pour voir (1), utiliser le fait que le bidual de $M$ est sans torsion d'après
le lemme \ref{lemma-hom-into-torsion-free}. Supposons $M$ de type fini.
Choisissons des générateurs $x_1, \ldots, x_n$ de $M$ sur $R$. Soit $K$ le
corps des fractions de $R$. Après renumérotation, nous pouvons supposer que,
pour un certain $0 \leq r \leq n$, les éléments $x_1, \ldots, x_r$ ont pour
images une base de $M \otimes_R K$ sur $K$. Pour $r < i \leq n$, écrivons
$x_i \otimes 1 = \sum_{1 \leq j \leq r} x_j \otimes k_{ij}$, avec
$k_{ij} \in K$. Choisissons $c \in R$ tel que $a_{ij} = ck_{ij} \in R$
et tel que $cx_i = \sum_{1 \leq j \leq r} a_{ij} x_j$ dans $M$
(un petit détail est omis). Nous obtenons alors les homomorphismes
$$
R^{\oplus r} \xrightarrow{\alpha} M \xrightarrow{\beta} R^{\oplus r}
\xrightarrow{\alpha} M
$$
avec $\alpha \circ \beta$ et $\beta \circ \alpha$ donnés par la
multiplication par $c$. Ici, $\alpha$ utilise $x_1, \ldots, x_r$ et $\beta$
envoie $x_i$ sur $c$ fois le $i$-ième vecteur de base pour
$1 \leq i \leq r$, et sur le vecteur $(a_{i1}, \ldots, a_{ir})$ pour
$r < i \leq n$. Considérons le diagramme commutatif
$$
\xymatrix{
R^{\oplus r} \ar[r] \ar[d]^{\text{id}} &
M \ar[r] \ar[d]^j &
R^{\oplus r} \ar[d]^{\text{id}} \ar[r] &
M \ar[d]^j \\
R^{\oplus r} \ar[r] &
\Hom_R(\Hom_R(M, R), R) \ar[r] &
R^{\oplus r} \ar[r] &
\Hom_R(\Hom_R(M, R), R)
}
$$
Bien sûr, les composées de deux flèches consécutives de la ligne inférieure
sont égales à la multiplication par $c$. Le diagramme montre que le noyau et
le conoyau de $j$ sont annulés par $c$, ce qui démontre (2). Pour (3),
remarquons que si $M$ est sans torsion, alors $\beta$ est injectif, ce qui
implique que $j$ est injectif. La réciproque est claire puisque nous avons
déjà vu que le bidual de $M$ est sans torsion.
\end{proof}

\begin{lemma}
\label{lemma-cokernel-map-double-dual-dvr}
Soit $R$ un anneau de valuation discrète et soit $M$ un $R$-module de type
fini. Alors l'homomorphisme
$j : M \to \Hom_R(\Hom_R(M, R), R)$ est surjectif.
\end{lemma}

\begin{proof}
Soit $M_{tors} \subset M$ le sous-module de torsion. Alors
$\Hom_R(M, R) = \Hom_R(M/M_{tors}, R)$ (ceci vaut sur tout anneau intègre).
Nous pouvons donc supposer que $M$ est sans torsion. Alors $M$ est libre
d'après le lemme \ref{lemma-dedekind-torsion-free-flat}, et le lemme est
clair.
\end{proof}

\begin{lemma}
\label{lemma-check-reflexive}
Soit $R$ un anneau intègre noethérien. Soit $M$ un $R$-module de type fini.
Les assertions suivantes sont équivalentes :
\begin{enumerate}
\item $M$ est réflexif,
\item $M_\mathfrak p$ est un $R_\mathfrak p$-module réflexif pour tout
idéal premier $\mathfrak p \subset R$, et
\item $M_\mathfrak m$ est un $R_\mathfrak m$-module réflexif pour tout idéal
maximal $\mathfrak m$ de $R$.
\end{enumerate}
\end{lemma}

\begin{proof}
Le localisé de $j : M \to \Hom_R(\Hom_R(M, R), R)$ en un idéal premier
$\mathfrak p$ est l'homomorphisme correspondant pour le module
$M_\mathfrak p$ sur l'anneau local intègre noethérien $R_\mathfrak p$.
Voir Algèbre, lemme \ref{algebra-lemma-hom-from-finitely-presented}. Le
lemme résulte donc d'Algèbre, lemme
\ref{algebra-lemma-characterize-zero-local}.
\end{proof}

\begin{lemma}
\label{lemma-sequence-reflexive}
Soit $R$ un anneau intègre noethérien. Soit
$0 \to M \to M' \to M''$ une suite exacte de $R$-modules de type fini. Si
$M'$ est réflexif et $M''$ est sans torsion, alors $M$ est réflexif.
\end{lemma}

\begin{proof}
Nous utiliserons sans autre mention le fait que $\Hom_R(N, N')$ est un
$R$-module de type fini pour tous $R$-modules de type fini $N$ et $N'$ ; voir
Algèbre, lemme \ref{algebra-lemma-ext-noetherian}. Prenons les duaux pour
obtenir une suite
$$
\Hom_R(M, R) \leftarrow \Hom_R(M', R) \leftarrow \Hom_R(M'', R)
$$
En prenant de nouveau les duaux, nous obtenons un diagramme commutatif
$$
\xymatrix{
\Hom_R(\Hom_R(M, R), R) \ar[r]_j &
\Hom_R(\Hom_R(M', R), R) \ar[r] &
\Hom_R(\Hom_R(M'', R), R) \\
M \ar[u] \ar[r] & M' \ar[u] \ar[r] & M'' \ar[u]
}
$$
Nous ne savons pas si la ligne supérieure est exacte. Cependant, par
hypothèse, la flèche verticale du milieu est un isomorphisme et la flèche
verticale de droite est injective (lemme \ref{lemma-reflexive-torsion-free}).
Nous affirmons que $j$ est injectif. En admettant cette assertion, une chasse
dans le diagramme montre que la flèche verticale de gauche est un
isomorphisme, c'est-à-dire que $M$ est réflexif.

\medskip\noindent
Démonstration de l'assertion. Considérons la suite exacte
$\Hom_R(M',R)\to \Hom_R(M,R)\to Q \to 0$ définissant $Q$. On applique
Algèbre, lemme \ref{algebra-lemma-hom-from-finitely-presented}, pour obtenir
$$
\Hom_K(M'\otimes_R K,K) \to
\Hom_K(M\otimes_R K,K) \to
Q\otimes_R K\to 0
$$
Mais $M \otimes_R K \to M' \otimes_R K$ est un homomorphisme injectif
d'espaces vectoriels, donc une injection scindée ; ainsi
$Q \otimes_R K = 0$, c'est-à-dire que $Q$ est de torsion. On obtient alors
la suite exacte
$$
0 \to \Hom_R(Q,R) \to
\Hom_R(\Hom_R(M,R),R) \to
\Hom_R(\Hom_R(M',R),R)
$$
et $\Hom_R(Q,R)=0$ puisque $Q$ est de torsion.
\end{proof}

\begin{lemma}
\label{lemma-characterize-reflexive}
Soit $R$ un anneau intègre noethérien. Soit $M$ un $R$-module de type fini.
Les assertions suivantes sont équivalentes :
\begin{enumerate}
\item $M$ est réflexif,
\item il existe une suite exacte courte $0 \to M \to F \to N \to 0$ avec
$F$ libre de type fini et $N$ sans torsion.
\end{enumerate}
\end{lemma}

\begin{proof}
Remarquons qu'un module libre de type fini est réflexif. D'après le lemme
\ref{lemma-sequence-reflexive}, (2) implique (1). Supposons $M$ réflexif.
Choisissons une présentation
$R^{\oplus m} \to R^{\oplus n} \to \Hom_R(M, R) \to 0$.
En prenant les duaux, nous obtenons une suite exacte
$$
0 \to \Hom_R(\Hom_R(M, R), R) \to R^{\oplus n} \to N \to 0
$$
avec $N = \Im(R^{\oplus n} \to R^{\oplus m})$ un module sans torsion.
Comme $M = \Hom_R(\Hom_R(M, R), R)$, nous obtenons une suite exacte comme
dans (2).
\end{proof}

\begin{lemma}
\label{lemma-flat-pullback-reflexive}
Soit $R \to R'$ un homomorphisme plat d'anneaux intègres noethériens. Si
$M$ est un $R$-module réflexif de type fini, alors $M \otimes_R R'$ est un
$R'$-module réflexif de type fini.
\end{lemma}

\begin{proof}
Choisissons une suite exacte courte $0 \to M \to F \to N \to 0$ avec $F$
libre de type fini et $N$ sans torsion, voir le lemme
\ref{lemma-characterize-reflexive}. Comme $R \to R'$ est plat, nous
obtenons une suite exacte courte
$0 \to M \otimes_R R' \to F \otimes_R R' \to N \otimes_R R' \to 0$
avec $F \otimes_R R'$ libre de type fini et $N \otimes_R R'$ sans torsion
(lemme \ref{lemma-flat-pullback-torsion}). Ainsi,
$M \otimes_R R'$ est réflexif d'après le lemme
\ref{lemma-characterize-reflexive}.
\end{proof}

\begin{lemma}
\label{lemma-dual-reflexive}
Soit $R$ un anneau intègre noethérien. Soit $M$ un $R$-module de type fini.
Soit $N$ un $R$-module réflexif de type fini. Alors $\Hom_R(M, N)$ est
réflexif.
\end{lemma}

\begin{proof}
Choisissons une présentation $R^{\oplus m} \to R^{\oplus n} \to M \to 0$.
Nous obtenons alors
$$
0 \to \Hom_R(M, N) \to N^{\oplus n} \to N' \to 0
$$
avec $N' = \Im(N^{\oplus n} \to N^{\oplus m})$ sans torsion. Nous
concluons par le lemme \ref{lemma-sequence-reflexive}.
\end{proof}

\begin{definition}
\label{definition-reflexive-hull}
Soit $R$ un anneau intègre noethérien. Soit $M$ un $R$-module de type fini.
Le module $M^{**} = \Hom_R(\Hom_R(M, R), R)$ est appelé l'{\it enveloppe
réflexive} de $M$.
\end{definition}

\noindent
Cela a un sens car l'enveloppe réflexive est réflexive d'après le lemme
\ref{lemma-dual-reflexive}. L'opération $M \mapsto M^{**}$ est un foncteur.
Si $\varphi : M \to N$ est un homomorphisme de $R$-modules à valeurs dans
un $R$-module réflexif $N$, alors $\varphi$ se factorise sous la forme
$M \to M^{**} \to N$ par l'enveloppe réflexive de $M$. Autrement dit,
prendre l'enveloppe
réflexive est l'adjoint à gauche du foncteur d'inclusion
$$
\text{modules réflexifs de type fini} \subset
\text{modules de type fini}
$$
sur un anneau intègre noethérien $R$.

\begin{lemma}
\label{lemma-hom-into-depth}
Soit $R$ un anneau local noethérien. Soient $M$, $N$ des $R$-modules de
type fini.
\begin{enumerate}
\item Si $N$ est de profondeur $\geq 1$, alors $\Hom_R(M, N)$ est de
profondeur $\geq 1$.
\item Si $N$ est de profondeur $\geq 2$, alors $\Hom_R(M, N)$ est de
profondeur $\geq 2$.
\end{enumerate}
\end{lemma}

\begin{proof}
Choisissons une présentation $R^{\oplus m} \to R^{\oplus n} \to M \to 0$.
En prenant les duaux, nous obtenons une suite exacte
$$
0 \to \Hom_R(M, N) \to N^{\oplus n} \to N' \to 0
$$
avec $N' = \Im(N^{\oplus n} \to N^{\oplus m})$. Un sous-module d'un module
de profondeur $\geq 1$ est de profondeur $\geq 1$ ; cela résulte
immédiatement de la définition. L'assertion (1) est donc claire. Pour (2),
remarquons que l'hypothèse et la remarque précédente impliquent que $N'$ est
de profondeur $\geq 1$. Le module $N^{\oplus n}$ est de profondeur
$\geq 2$. D'après Algèbre, lemme \ref{algebra-lemma-depth-in-ses}, nous en
concluons que $\Hom_R(M, N)$ est de profondeur $\geq 2$.
\end{proof}

\begin{lemma}
\label{lemma-hom-into-S2}
Soit $R$ un anneau noethérien. Soient $M$, $N$ des $R$-modules de type fini.
\begin{enumerate}
\item Si $N$ possède la propriété $(S_1)$, alors $\Hom_R(M, N)$ possède la
propriété $(S_1)$.
\item Si $N$ possède la propriété $(S_2)$, alors $\Hom_R(M, N)$ possède la
propriété $(S_2)$.
\item Si $R$ est intègre, si $N$ est sans torsion et vérifie $(S_2)$, alors
$\Hom_R(M, N)$ est sans torsion et possède la propriété $(S_2)$.
\end{enumerate}
\end{lemma}

\begin{proof}
Comme la localisation en des idéaux premiers commute à la formation de
$\Hom_R$ pour des $R$-modules de type fini (Algèbre, lemme
\ref{algebra-lemma-ext-noetherian}), les assertions (1) et (2) résultent
immédiatement du lemme \ref{lemma-hom-into-depth}. L'assertion (3) résulte
de (2) et du lemme \ref{lemma-hom-into-torsion-free}.
\end{proof}

\begin{lemma}
\label{lemma-check-injective-on-ass}
 Soit $R$ un anneau noethérien. Soit $\varphi : M \to N$ une application de
 $R$-modules. Supposons que pour tout idéal premier $\mathfrak p$
 de $R$, au moins l'une des propriétés suivantes soit vérifiée
\begin{enumerate}
\item $M_\mathfrak p \to N_\mathfrak p$ est injective, ou
\item $\mathfrak p \not \in \text{Ass}(M)$.
\end{enumerate}
Alors $\varphi$ est injective.
\end{lemma}

\begin{proof}
 Soit $\mathfrak p$ un idéal premier associé de $\Ker(\varphi)$.
 Alors il existe un élément $x \in M_\mathfrak p$ qui appartient au
 noyau de $M_\mathfrak p \to N_\mathfrak p$ et dont l'annulateur est
 $\mathfrak pR_\mathfrak p$
 (Algèbre, lemme \ref{algebra-lemma-associated-primes-localize}).
 Cela est impossible dans les deux cas. Donc
 $\text{Ass}(\Ker(\varphi)) = \emptyset$ et nous concluons
 $\Ker(\varphi) = 0$ par le lemme d'Algèbre
 \ref{algebra-lemma-ass-zero}.
\end{proof}

\begin{lemma}
\label{lemma-check-isomorphism-via-depth-and-ass}
 Soit $R$ un anneau noethérien. Soit $\varphi : M \to N$ une application de
 $R$-modules. Supposons que $M$ soit de type fini et que, pour tout idéal
 premier $\mathfrak p$ de $R$, l'une des propriétés suivantes soit vérifiée
\begin{enumerate}
\item $M_\mathfrak p \to N_\mathfrak p$ est un isomorphisme, ou
\item $\text{profondeur}(M_\mathfrak p) \geq 2$ et
 $\mathfrak p \not \in \text{Ass}(N)$.
\end{enumerate}
Alors $\varphi$ est un isomorphisme.
\end{lemma}

\begin{proof}
Par le lemme \ref{lemma-check-injective-on-ass}, nous voyons que
$\varphi$ est injective. Soit $N' \subset N$ un $R$-module de type fini
contenant l'image de $M$. Alors $\text{Ass}(N_\mathfrak p) = \emptyset$
implique $\text{Ass}(N'_\mathfrak p) = \emptyset$.
Les hypothèses du lemme sont donc vérifiées pour $M \to N'$.
Pour démontrer que $\varphi$ est un isomorphisme, il suffit de démontrer
la même chose pour tout sous-module $N' \subset N$ de ce type. Nous pouvons
donc supposer que $N$ est un $R$-module de type fini. Dans ce cas,
$\mathfrak p \not \in \text{Ass}(N) \Rightarrow
\text{profondeur}(N_\mathfrak p) \geq 1$, voir le lemme d'Algèbre
\ref{algebra-lemma-ideal-nonzerodivisor}.
Considérons la suite exacte courte
$$
0 \to M \to N \to Q \to 0
$$
définissant $Q$.
En examinant les conditions, nous voyons que soit $Q_\mathfrak p = 0$
dans le cas (1), soit $\text{profondeur}(Q_\mathfrak p) \geq 1$ dans le cas (2),
par le lemme d'Algèbre \ref{algebra-lemma-depth-in-ses}.
Cela implique que $Q$ n'a aucun idéal premier associé, d'où $Q = 0$ par
le lemme d'Algèbre \ref{algebra-lemma-ass-zero}.
\end{proof}

\begin{lemma}
\label{lemma-isom-depth-2-torsion-free}
 Soit $R$ un anneau intègre noethérien. Soit $\varphi : M \to N$ une
 application de $R$-modules. Supposons que $M$ soit de type fini, que $N$
 soit sans torsion et que, pour tout idéal premier $\mathfrak p$ de $R$,
 l'une des propriétés suivantes soit vérifiée
\begin{enumerate}
\item $M_\mathfrak p \to N_\mathfrak p$ est un isomorphisme, ou
\item $\text{profondeur}(M_\mathfrak p) \geq 2$.
\end{enumerate}
Alors $\varphi$ est un isomorphisme.
\end{lemma}

\begin{proof}
Il s'agit d'un cas particulier du lemme
\ref{lemma-check-isomorphism-via-depth-and-ass}.
\end{proof}

\begin{lemma}
\label{lemma-reflexive-depth-2}
 Soit $R$ un anneau intègre noethérien. Soit $M$ un $R$-module de type fini.
Les propriétés suivantes sont équivalentes
\begin{enumerate}
\item $M$ est réflexif,
\item pour tout idéal premier $\mathfrak p$ de $R$, l'une des propriétés
 suivantes est vérifiée
\begin{enumerate}
\item $M_\mathfrak p$ est un $R_\mathfrak p$-module réflexif, ou
\item $\text{profondeur}(M_\mathfrak p) \geq 2$.
\end{enumerate}
\end{enumerate}
\end{lemma}

\begin{proof}
Si (1) est vraie, alors $M_\mathfrak p$ est un module réflexif
pour tout idéal premier $\mathfrak p$, par le lemme
\ref{lemma-check-reflexive}. Ainsi (1) $\Rightarrow$ (2). Supposons (2).
Posons $N = \Hom_R(\Hom_R(M, R), R)$, de sorte que
$$
N_\mathfrak p =
\Hom_{R_\mathfrak p}(
\Hom_{R_\mathfrak p}(M_\mathfrak p, R_\mathfrak p), R_\mathfrak p)
$$
pour tout idéal premier $\mathfrak p$ de $R$.
Voir le lemme d'Algèbre \ref{algebra-lemma-hom-from-finitely-presented}.
Nous appliquons le lemme \ref{lemma-isom-depth-2-torsion-free} à
l'application $j : M \to N$. Cela est permis car $M$ est de type fini et
$N$ est sans torsion par le lemme \ref{lemma-hom-into-torsion-free}.
Dans le cas (2)(a), l'application $M_\mathfrak p \to N_\mathfrak p$ est un
isomorphisme et, dans le cas (2)(b), nous avons
$\text{profondeur}(M_\mathfrak p) \geq 2$.
\end{proof}

\begin{lemma}
\label{lemma-reflexive-S2}
 Soit $R$ un anneau intègre noethérien. Soit $M$ un $R$-module réflexif de
 type fini. Soit $\mathfrak p \subset R$ un idéal premier.
\begin{enumerate}
\item Si $\text{profondeur}(R_\mathfrak p) \geq 2$, alors
 $\text{profondeur}(M_\mathfrak p) \geq 2$.
\item Si $R$ possède la propriété $(S_2)$, alors $M$ possède la propriété
 $(S_2)$.
\end{enumerate}
\end{lemma}

\begin{proof}
 Puisque la formation de l'enveloppe réflexive
 $\Hom_R(\Hom_R(M, R), R)$ commute à la localisation
 (Algèbre, lemme \ref{algebra-lemma-hom-from-finitely-presented}),
 la partie (1) découle du lemme \ref{lemma-hom-into-depth}.
 La partie (2) est immédiate d'après le lemme \ref{lemma-hom-into-S2}.
\end{proof}

\begin{example}
\label{example-ring-not-S2}
 Les résultats précédents et suivants suggèrent que la réflexivité est liée
 à la propriété $(S_2)$ ; voici un exemple qui empêche des conjectures trop
 optimistes. Soit $k$ un corps. Soit $R$ la $k$-sous-algèbre de $k[x, y]$
 engendrée par $1, y, x^2, xy, x^3$. Alors $R$ ne possède pas la propriété
 $(S_2)$. Ainsi, $R$ considéré comme $R$-module est un exemple de module
 $R$ réflexif qui ne possède pas la propriété $(S_2)$. Soit
 $M = k[x, y]$ considéré comme un $R$-module. Alors $M$ est un $R$-module
 réflexif car
$$
\Hom_R(M, R) = \mathfrak m = (y, x^2, xy, x^3)
\quad\text{et}\quad
\Hom_R(\mathfrak m, R) = M
$$
 et $M$ possède la propriété $(S_2)$ comme $R$-module (calculs omis).
 Ainsi, $R$ est un anneau intègre noethérien possédant un module réflexif
 ayant la propriété $(S_2)$, alors que $R$ lui-même ne possède pas $(S_2)$.
\end{example}

\begin{lemma}
\label{lemma-reflexive-over-normal}
 Soit $R$ un anneau intègre normal noethérien de corps des fractions $K$.
 Soit $M$ un $R$-module de type fini. Les propriétés suivantes sont
 équivalentes
\begin{enumerate}
\item $M$ est réflexif,
\item $M$ est sans torsion et possède la propriété $(S_2)$,
\item $M$ est sans torsion et
$M = \bigcap_{\text{hauteur}(\mathfrak p) = 1} M_{\mathfrak p}$
 où l'intersection est prise dans $M_K = M \otimes_R K$.
\end{enumerate}
\end{lemma}

\begin{proof}
 Par le lemme d'Algèbre \ref{algebra-lemma-criterion-normal},
 nous voyons que $R$ possède les propriétés $(R_1)$ et $(S_2)$.

\medskip\noindent
 Supposons (1). Alors $M$ est sans torsion par le lemme
 \ref{lemma-reflexive-torsion-free} et possède $(S_2)$ par le lemme
 \ref{lemma-reflexive-S2}. Ainsi (2) est vraie.

\medskip\noindent
 Supposons (2). Par définition,
$M' = \bigcap_{\text{hauteur}(\mathfrak p) = 1} M_{\mathfrak p}$
 est le noyau de l'application
$$
M_K
\longrightarrow
\bigoplus\nolimits_{\text{hauteur}(\mathfrak p) = 1} M_K/M_\mathfrak p
\subset
\prod\nolimits_{\text{hauteur}(\mathfrak p) = 1} M_K/M_\mathfrak p
$$
 Observons que notre application se factorise bien par la somme directe
 comme indiqué, car pour $a/b \in K$ il n'existe qu'un nombre fini d'idéaux
 premiers de hauteur $1$ $\mathfrak p$ tels que $b \in \mathfrak p$.
 Soit $\mathfrak p_0$ un idéal premier de hauteur $1$. Alors
 $(M_K/M_\mathfrak p)_{\mathfrak p_0} = 0$ sauf si
 $\mathfrak p = \mathfrak p_0$, auquel cas
 $(M_K/M_\mathfrak p)_{\mathfrak p_0} = M_K/M_{\mathfrak p_0}$.
 Ainsi, par exactitude de la localisation et puisque la localisation commute
 aux sommes directes, nous voyons que
 $M'_{\mathfrak p_0} = M_{\mathfrak p_0}$.
 Puisque $M$ est de profondeur $\geq 2$ aux idéaux premiers de hauteur
 $> 1$, l'application $M \to M'$ est un isomorphisme par le lemme
 \ref{lemma-isom-depth-2-torsion-free}. Ainsi (3) est vraie.

\medskip\noindent
 Supposons (3). Soit $\mathfrak p$ un idéal premier de hauteur $1$.
 Alors $R_\mathfrak p$ est un anneau de valuation discrète par $(R_1)$.
 Par le lemme \ref{lemma-dedekind-torsion-free-flat}, nous voyons que
 $M_\mathfrak p$ est libre de type fini, en particulier réflexif. Ainsi
 l'application $M \to M^{**}$ induit un isomorphisme en chacun des idéaux
 premiers $\mathfrak p$ de hauteur $1$. Donc la condition
$M = \bigcap_{\text{hauteur}(\mathfrak p) = 1} M_{\mathfrak p}$
 implique $M = M^{**}$ et (1) est vraie.
\end{proof}

\begin{lemma}
\label{lemma-describe-reflexive-hull}
 Soit $R$ un anneau intègre normal noethérien. Soit $M$ un $R$-module de
 type fini. Alors l'enveloppe réflexive de $M$ est l'intersection
$$
M^{**} =
\bigcap\nolimits_{\text{hauteur}(\mathfrak p) = 1}
M_{\mathfrak p}/(M_\mathfrak p)_{tors} =
\bigcap\nolimits_{\text{hauteur}(\mathfrak p) = 1}
(M/M_{tors})_\mathfrak p
$$
 prise dans $M \otimes_R K$.
\end{lemma}

\begin{proof}
 Soit $\mathfrak p$ un idéal premier de hauteur $1$.
 Le noyau de $M_\mathfrak p \to M \otimes_R K$ est le
 sous-module de torsion $(M_\mathfrak p)_{tors}$ de $M_\mathfrak p$.
 De plus, nous avons
$(M/M_{tors})_\mathfrak p = M_\mathfrak p/(M_\mathfrak p)_{tors}$
 et c'est un module libre de type fini sur l'anneau de valuation discrète
 $R_\mathfrak p$ (lemme \ref{lemma-dedekind-torsion-free-flat}).
 Alors $M_\mathfrak p/(M_\mathfrak p)_{tors} \to
 (M_\mathfrak p)^{**} = (M^{**})_\mathfrak p$ est un isomorphisme ; le lemme
 est donc une conséquence du lemme \ref{lemma-reflexive-over-normal}.
\end{proof}

\begin{lemma}
\label{lemma-integral-closure-reflexive}
 Soit $A$ un anneau intègre normal noethérien de corps des fractions $K$.
 Soit $L$ une extension finie de $K$. Si la clôture intégrale $B$ de $A$
 dans $L$ est de type fini sur $A$, alors $B$ est réflexif comme $A$-module.
\end{lemma}

\begin{proof}
 Il suffit de montrer que $B = \bigcap B_\mathfrak p$, où l'intersection
 porte sur les idéaux premiers de hauteur $1$, avec
 $\mathfrak p \subset A$ ; voir
 le lemme \ref{lemma-reflexive-over-normal}.
 Soit $b \in \bigcap B_\mathfrak p$. Soit
 $x^d + a_1x^{d - 1} + \ldots + a_d$ le polynôme minimal de $b$ sur $K$.
 Nous voulons montrer que $a_i \in A$.
 Par le lemme d'Algèbre \ref{algebra-lemma-minimal-polynomial-normal-domain},
 nous voyons que $a_i \in A_\mathfrak p$ pour tout $i$ et tout idéal premier
 $\mathfrak p$ de hauteur un. Nous obtenons donc le résultat voulu par le
 lemme d'Algèbre
 \ref{algebra-lemma-normal-domain-intersection-localizations-height-1}
 (ou par le lemme déjà cité, puisque $A$ est réflexif comme module sur
 lui-même).
\end{proof}








\section{Idéaux de contenu}
\label{section-content}

\noindent
La définition n'est peut-être pas celle à laquelle vous vous attendez.

\begin{definition}
\label{definition-content-ideal}
Soit $A$ un anneau. Soit $M$ un $A$-module plat. Soit $x \in M$.
Si l'ensemble des idéaux $I$ de $A$ tels que $x \in IM$ possède un plus
petit élément, nous l'appelons {\it idéal de contenu de $x$}.
\end{definition}

\noindent
Notons que, puisque $M$ est plat sur $A$, pour deux idéaux $I, I'$ de $A$
nous avons $IM \cap I'M = (I \cap I')M$, comme on le voit en tensorisant
la suite exacte
$0 \to I \cap I' \to I \oplus I' \to I + I' \to 0$ par $M$.

\begin{lemma}
\label{lemma-content-finitely-generated}
Soit $A$ un anneau. Soit $M$ un $A$-module plat. Soit $x \in M$.
L'idéal de contenu de $x$, s'il existe, est de type fini.
\end{lemma}

\begin{proof}
Supposons $x \in IM$. Alors nous pouvons écrire
$x = \sum_{i = 1, \ldots, n} f_i x_i$ avec $f_i \in I$ et $x_i \in M$.
Ainsi $x \in I'M$ avec $I' = (f_1, \ldots, f_n)$.
\end{proof}

\begin{lemma}
\label{lemma-equal-content}
Soit $(A, \mathfrak m)$ un anneau local. Soit $u : M \to N$ une application
de $A$-modules plats telle que $\overline{u} : M/\mathfrak m M \to
N/\mathfrak m N$ soit injective. Si $x \in M$ a pour idéal de contenu $I$,
alors $u(x)$ a également pour idéal de contenu $I$.
\end{lemma}

\begin{proof}
Il est clair que $u(x) \in IN$. Si $u(x) \in I'N$, alors
$u(x) \in (I' \cap I)N$, voir la discussion qui suit la définition
\ref{definition-content-ideal}. Il suffit donc de montrer que si
$x \in I'N$ et $I' \subset I$, $I' \not = I$, alors $u(x) \not \in I'N$.
Puisque $I/I'$ est un $A$-module non nul de type fini (lemme
\ref{lemma-content-finitely-generated}), il existe une application non nulle
$\chi : I/I' \to A/\mathfrak m$ de $A$-modules, par le lemme de Nakayama
(lemme d'Algèbre \ref{algebra-lemma-NAK}). Puisque $I$ est l'idéal de
contenu de $x$, nous voyons que $x \not \in I''M$, où
$I'' = \Ker(\chi)$. Ainsi $x$ n'appartient pas au noyau de l'application
$$
IM = I \otimes_A M \xrightarrow{\chi \otimes 1}
A/\mathfrak m \otimes M \cong M/\mathfrak m M
$$
En appliquant notre hypothèse sur $\overline{u}$, nous concluons que
$u(x)$ ne s'envoie pas sur zéro par l'application
$$
IN = I \otimes_A N \xrightarrow{\chi \otimes 1}
A/\mathfrak m \otimes N \cong N/\mathfrak m N
$$
et nous concluons.
\end{proof}

\begin{lemma}
\label{lemma-content-exists-flat-Mittag-Leffler}
Soit $A$ un anneau. Soit $M$ un module de Mittag-Leffler plat.
Alors tout élément de $M$ possède un idéal de contenu.
\end{lemma}

\begin{proof}
C'est un cas particulier du lemme d'Algèbre
\ref{algebra-lemma-flat-ML-criterion}.
\end{proof}








\section{Conditions de platitude et de finitude}
\label{section-flat-finite}

\noindent
Dans cette section, nous examinons quelques conséquences de l'implication
« plat $+$ de type fini $\Rightarrow$ de présentation finie ».
Nous reviendrons sur ce résultat dans le chapitre consacré à la platitude,
voir Plus de platitude, section \ref{flat-section-introduction}.
Un premier résultat de ce type a été démontré dans le lemme d'Algèbre
\ref{algebra-lemma-finite-flat-module-finitely-presented}.

\begin{lemma}
\label{lemma-flat-finite-type-finite-presentation-local-module}
Soit $R$ un anneau. Soit $S = R[x_1, \ldots, x_n]$ un anneau de
polynômes sur $R$. Soit $M$ un $S$-module.
Supposons que
\begin{enumerate}
\item il existe un nombre fini d'idéaux premiers $\mathfrak p_1, \ldots,
\mathfrak p_m$ de $R$ tels que l'application $R \to \prod R_{\mathfrak p_j}$
soit injective,
\item $M$ est un $S$-module de type fini,
\item $M$ est plat sur $R$, et
\item pour tout idéal premier $\mathfrak p$ de $R$, le module $M_{\mathfrak p}$
est de présentation finie sur $S_{\mathfrak p}$.
\end{enumerate}
Alors $M$ est de présentation finie sur $S$.
\end{lemma}

\begin{proof}
Choisissons une présentation
$$
0 \to K \to S^{\oplus r} \to M \to 0
$$
de $M$ comme $S$-module. Soit $\mathfrak q$ un idéal premier de $S$
au-dessus d'un idéal premier $\mathfrak p$ de $R$. Par hypothèse, il existe
un nombre fini d'éléments $k_1, \ldots, k_t \in K$ tels que, si nous posons
$K' = \sum Sk_j \subset K$, alors
$K'_{\mathfrak p} = K_{\mathfrak p}$ et
$K'_{\mathfrak p_j} = K_{\mathfrak p_j}$ pour $j = 1, \ldots, m$.
En posant $M' = S^{\oplus r}/K'$, nous déduisons en particulier que
$M'_{\mathfrak q} = M_{\mathfrak q}$. Par ouverture de la platitude, voir
le théorème d'Algèbre \ref{algebra-theorem-openness-flatness}, nous concluons
qu'il existe un $g \in S$, $g \not \in \mathfrak q$, tel que $M'_g$ soit plat
sur $R$. Ainsi $M'_g \to M_g$ est une application surjective de $R$-modules
plats. Considérons le diagramme commutatif
$$
\xymatrix{
M'_g \ar[r] \ar[d] & M_g \ar[d] \\
\prod (M'_g)_{\mathfrak p_j} \ar[r] & \prod (M_g)_{\mathfrak p_j}
}
$$
La flèche inférieure est un isomorphisme par le choix de
$k_1, \ldots, k_t$. La flèche verticale de gauche est injective car
$R \to \prod R_{\mathfrak p_j}$ est injective et $M'_g$ est plat sur $R$.
La flèche horizontale supérieure est donc injective, et par conséquent un
isomorphisme. Cela montre que $M_g$ est de présentation finie sur $S_g$.
Nous concluons en appliquant le lemme d'Algèbre
\ref{algebra-lemma-cover}.
\end{proof}

\begin{lemma}
\label{lemma-flat-finite-type-finite-presentation-local}
Soit $R \to S$ un homomorphisme d'anneaux.
Supposons que
\begin{enumerate}
\item il existe un nombre fini d'idéaux premiers
$\mathfrak p_1, \ldots, \mathfrak p_m$ de $R$ tels que
l'application $R \to \prod R_{\mathfrak p_j}$ soit injective,
\item $R \to S$ est de type fini,
\item $S$ est plat sur $R$, et
\item pour tout idéal premier $\mathfrak p$ de $R$, l'anneau $S_{\mathfrak p}$
est de présentation finie sur $R_{\mathfrak p}$.
\end{enumerate}
Alors $S$ est de présentation finie sur $R$.
\end{lemma}

\begin{proof}
Par hypothèse, $S$ est un quotient d'un anneau de polynômes sur $R$.
Le résultat découle donc directement du lemme
\ref{lemma-flat-finite-type-finite-presentation-local-module}.
\end{proof}

\begin{lemma}
\label{lemma-flat-graded-finite-type-finite-presentation-module}
Soit $R$ un anneau.
Soit $S = R[x_1, \ldots, x_n]$ une algèbre de polynômes graduée sur $R$,
c'est-à-dire $\deg(x_i) > 0$, sans que ces degrés soient nécessairement égaux
à $1$. Soit $M$ un $S$-module gradué.
Supposons que
\begin{enumerate}
\item $R$ est un anneau local,
\item $M$ est un $S$-module de type fini, et
\item $M$ est plat sur $R$.
\end{enumerate}
Alors $M$ est de présentation finie comme $S$-module.
\end{lemma}

\begin{proof}
Écrivons $M = \bigoplus M_d$ pour la graduation de $M$.
Choisissons des générateurs homogènes $m_1, \ldots, m_r \in M$ de $M$.
Posons $\deg(m_i) = d_i \in \mathbf{Z}$. Cela fournit une présentation
$$
0 \to K \to \bigoplus\nolimits_{i = 1, \ldots, r} S(-d_i) \to M \to 0
$$
qui, en chaque degré $d$, donne la suite exacte courte
$$
0 \to K_d \to \bigoplus\nolimits_{i = 1, \ldots, r} S_{d - d_i} \to
M_d \to 0.
$$
Par hypothèse, chaque $M_d$ est un $R$-module plat de type fini. Par le lemme
d'Algèbre \ref{algebra-lemma-finite-flat-local}, cela implique que chaque
$M_d$ est un $R$-module libre de type fini. Nous voyons donc que chaque
$K_d$ est un $R$-module de type fini. De plus, chaque $K_d$ est plat sur $R$
par le lemme d'Algèbre \ref{algebra-lemma-flat-ses}. Nous concluons donc que
chaque $K_d$ est libre de type fini, à nouveau par le lemme d'Algèbre
\ref{algebra-lemma-finite-flat-local}.

\medskip\noindent
Soit $\mathfrak m$ l'idéal maximal de $R$. Par platitude de $M$ sur $R$,
les suites exactes courtes ci-dessus restent exactes après tensorisation
par $\kappa = \kappa(\mathfrak m)$. Comme l'anneau $S \otimes_R \kappa$ est
noethérien, il existe des éléments homogènes $k_1, \ldots, k_t \in K$ tels
que leurs images $\overline{k}_j$ engendrent $K \otimes_R \kappa$ sur
$S \otimes_R \kappa$. Posons $\deg(k_j) = e_j$. Ainsi, pour tout $d$,
l'application
$$
\bigoplus\nolimits_{j = 1, \ldots, t} S_{d - e_j}
\longrightarrow
K_d
$$
devient surjective après tensorisation par $\kappa$. Par le lemme de Nakayama
(lemme d'Algèbre \ref{algebra-lemma-NAK}), cela implique que l'application est
surjective sur $R$. Ainsi $K$ est engendré par $k_1, \ldots, k_t$ sur $S$,
ce qui conclut la preuve.
\end{proof}

\begin{lemma}
\label{lemma-flat-graded-finite-type-finite-presentation}
Soit $R$ un anneau. Soit $S = \bigoplus_{n \geq 0} S_n$ une $R$-algèbre
graduée. Soit $M = \bigoplus_{d \in \mathbf{Z}} M_d$ un $S$-module gradué.
Supposons que $S$ soit de type fini comme $R$-algèbre, que $S_0$ soit une
$R$-algèbre de type fini, et qu'il existe un nombre fini d'idéaux premiers
$\mathfrak p_j$, $i = 1, \ldots, m$, tels que
$R \to \prod R_{\mathfrak p_j}$ soit injective.
\begin{enumerate}
\item Si $S$ est plat sur $R$, alors $S$ est une $R$-algèbre de présentation
finie.
\item Si $M$ est plat comme $R$-module et de type fini comme $S$-module,
alors $M$ est de présentation finie comme $S$-module.
\end{enumerate}
\end{lemma}

\begin{proof}
Comme $S$ est de type fini comme $R$-algèbre, il est de type fini comme
$S_0$-algèbre, disons par des éléments homogènes $t_1, \ldots, t_n \in S$
de degrés $d_1, \ldots, d_n > 0$. Posons $P = R[x_1, \ldots, x_n]$ avec
$\deg(x_i) = d_i$. L'application d'anneaux $P \to S$, $x_i \to t_i$, est
finie puisque $S_0$ est un $R$-module de type fini. Pour démontrer (1), il
suffit de démontrer que $S$ est un $P$-module de présentation finie. Pour
démontrer (2), il suffit de démontrer que $M$ est un $P$-module de
présentation finie. Il suffit donc de prouver que si $S = P$ est un anneau
de polynômes gradué et si $M$ est un $S$-module de type fini plat sur $R$,
alors $M$ est de présentation finie comme $S$-module. Par le lemme
\ref{lemma-flat-graded-finite-type-finite-presentation-module}, nous voyons
que $M_{\mathfrak p}$ est un $S_{\mathfrak p}$-module de présentation finie
pour tout idéal premier $\mathfrak p$ de $R$. Le résultat découle donc du
lemme \ref{lemma-flat-finite-type-finite-presentation-local-module}.
\end{proof}

\begin{remark}
\label{remark-when-does-condition-hold}
Soit $R$ un anneau. Quand $R$ satisfait-il la condition mentionnée dans les
lemmes \ref{lemma-flat-finite-type-finite-presentation-local-module},
\ref{lemma-flat-finite-type-finite-presentation-local} et
\ref{lemma-flat-graded-finite-type-finite-presentation} ?
C'est le cas si
\begin{enumerate}
\item $R$ est local,
\item $R$ est noethérien,
\item $R$ est intègre,
\item $R$ est un anneau réduit ayant un nombre fini d'idéaux premiers
minimaux, ou
\item $R$ possède un nombre fini d'idéaux premiers faiblement associés, voir
le lemme d'Algèbre \ref{algebra-lemma-zero-at-weakly-ass-zero}.
\end{enumerate}
Ainsi, ces lemmes s'appliquent dans tous les cas énumérés ci-dessus.
\end{remark}

\noindent
Le lemme suivant sera amélioré dans Plus de platitude, proposition
\ref{flat-proposition-flat-finite-type-finite-presentation-domain}.

\begin{lemma}
\label{lemma-flat-finite-type-valuation-ring-finite-presentation}
\begin{reference}
\cite[Théorème 3]{Nagata-Finitely}
\end{reference}
Soit $A$ un anneau de valuation. Soit $A \to B$ une application d'anneaux
de type fini. Soit $M$ un $B$-module de type fini.
\begin{enumerate}
\item Si $B$ est plat sur $A$, alors $B$ est une $A$-algèbre de présentation
finie.
\item Si $M$ est plat comme $A$-module, alors $M$ est de présentation finie
comme $B$-module.
\end{enumerate}
\end{lemma}

\begin{proof}
Nous utiliserons qu'un $A$-module est plat si et seulement s'il est sans
torsion, voir le lemme \ref{lemma-valuation-ring-torsion-free-flat}.
Par le lemme d'Algèbre \ref{algebra-lemma-homogenize}, nous pouvons trouver
une $A$-algèbre graduée $S$ avec $S_0 = A$, engendrée par un nombre fini
d'éléments de degré $1$, un élément $f \in S_1$ et un $S$-module gradué de
type fini $N$ tels que $B \cong S_{(f)}$ et $M \cong N_{(f)}$. Si $M$ est
sans torsion, nous pouvons choisir $N$ sans torsion en le remplaçant par
$N/N_{tors}$, voir le lemme \ref{lemma-torsion}. De même, si $B$ est sans
torsion, nous pouvons choisir $S$ sans torsion en le remplaçant par
$S/S_{tors}$. Ainsi, dans le cas (1), nous pouvons appliquer le lemme
\ref{lemma-flat-graded-finite-type-finite-presentation} pour voir que $S$
est une $A$-algèbre de présentation finie, ce qui implique que
$B = S_{(f)}$ est une $A$-algèbre de présentation finie. Pour (2), nous
pouvons d'abord remplacer $S$ par un anneau de polynômes gradué, puis
appliquer le lemme \ref{lemma-flat-graded-finite-type-finite-presentation-module}
pour conclure.
\end{proof}

\begin{lemma}
\label{lemma-valuation-ring-flat-essentially-finite-type}
Soit $A$ un anneau de valuation. Soit $A \to B$ un homomorphisme local
essentiellement de type fini. Soit $M$ un $B$-module de type fini.
\begin{enumerate}
\item Si $B$ est plat sur $A$, alors $B$ est essentiellement de présentation
finie sur $A$.
\item Si $M$ est plat comme $A$-module, alors $M$ est de présentation finie
comme $B$-module.
\end{enumerate}
\end{lemma}

\begin{proof}
Par hypothèse, nous pouvons écrire $B$ comme quotient de la localisation
d'une algèbre de polynômes $P = A[x_1, \ldots, x_n]$ en un idéal premier
$\mathfrak q$. Dans le cas (1), nous considérons $M = B$ comme module de type
fini sur $P_\mathfrak q$ et, dans le cas (2), nous considérons $M$ comme
module de type fini sur $P_\mathfrak q$. Dans les deux cas, nous devons
montrer qu'il s'agit d'un $P_\mathfrak q$-module de présentation finie ; voir
le lemme d'Algèbre \ref{algebra-lemma-finitely-presented-over-subring} dans
le cas (2).

\medskip\noindent
Choisissons une présentation $0 \to K \to P_\mathfrak q^{\oplus r} \to M \to 0$,
possible puisque $M$ est de type fini sur $P_\mathfrak q$. Soit
$L = P^{\oplus r} \cap K$. Alors $K = L_\mathfrak q$, voir le lemme
d'Algèbre \ref{algebra-lemma-submodule-localization}. Alors
$N = P^{\oplus r}/L$ est un sous-module de $M$ et donc plat par le lemme
\ref{lemma-valuation-ring-torsion-free-flat}. Comme $N$ est également un
$P$-module de type fini, nous voyons que $N$ est de présentation finie comme
$P$-module par le lemme
\ref{lemma-flat-finite-type-valuation-ring-finite-presentation}.
Puisque la localisation est exacte (proposition d'Algèbre
\ref{algebra-proposition-localization-exact}), nous voyons que
$N_\mathfrak q = M$ et nous concluons.
\end{proof}











\section{Éclatement et platitude}
\label{section-blowup-flat}

\noindent
Dans cette section, nous commençons l'étude de résultats de la forme :
« Après un éclatement, la transformée stricte devient plate ». D'autres
résultats de ce type se trouvent dans Diviseurs, section
\ref{divisors-section-blowup-flat}, et Plus de platitude, section
\ref{flat-section-blowup-flat}.

\begin{definition}
\label{definition-strict-transform}
Soit $R$ un anneau. Soit $I \subset R$ un idéal et $a \in I$.
Soit $R[\frac{I}{a}]$ l'algèbre d'éclatement affine, voir la définition
d'Algèbre \ref{algebra-definition-blow-up}. Soit $M$ un $R$-module.
La {\it transformée stricte de $M$ le long de $R \to R[\frac{I}{a}]$} est
le $R[\frac{I}{a}]$-module
$$
M' = \left(M \otimes_R R[\textstyle{\frac{I}{a}}]\right)/a\text{-torsion}
$$
\end{definition}

\noindent
Le résultat suivant est une version très faible de la platification par
éclatement, mais il est déjà parfois utile.

\begin{lemma}
\label{lemma-flatten-on-affine-blowup}
Soit $(R, \mathfrak m)$ un anneau local intègre de corps des fractions $K$.
Soit $S$ une $R$-algèbre de type fini. Soit $M$ un $S$-module de type fini.
Pour tout anneau de valuation $A \subset K$ dominant $R$, il existe un idéal
$I \subset \mathfrak m$ et un élément non nul $a \in I$ tels que
\begin{enumerate}
\item $I$ est de type fini,
\item $A$ a un centre sur $R[\frac{I}{a}]$,
\item l'anneau fibre de $R \to R[\frac{I}{a}]$ en $\mathfrak m$
 n'est pas nul, et
\item la transformée stricte $S_{I, a}$ de $S$ le long de
 $R \to R[\frac{I}{a}]$ est plate et de présentation finie sur $R$, et la
 transformée stricte $M_{I, a}$ de $M$ le long de
 $R \to R[\frac{I}{a}]$ est plate sur $R$ et de présentation finie sur
 $S_{I, a}$.
\end{enumerate}
\end{lemma}

\begin{proof}
Écrivons $S = R[x_1, \ldots, x_n]/J$ et notons $N = S \oplus M$,
considéré comme module sur $P = R[x_1, \ldots, x_n]$. Si nous pouvons
démontrer le lemme lorsque $S$ est une algèbre de polynômes sur $R$, alors
nous pouvons trouver $I, a$ satisfaisant (1), (2), (3) tels que la transformée
stricte $N_{I, a}$ de $N$ le long de $R \to R[\frac{I}{a}]$ soit plate sur
$R$ et de présentation finie comme module sur la transformée stricte
$P_{I, a}$ de $P$. Puisque $P_{I, a} = R[\frac{I}{a}][x_1, \ldots, x_n]$
(un détail est omis), nous voyons que le facteur direct
$S_{I, a} \subset N_{I, a}$ est plat sur $R$ et de présentation finie comme
module sur $R[\frac{I}{a}][x_1, \ldots, x_n]$. Ainsi $S_{I, a}$ est de
présentation finie comme $R[\frac{I}{a}]$-algèbre. De plus, le facteur direct
$M_{I, a} \subset N_{I, a}$ est plat sur $R$ et de présentation finie comme
module sur $P_{I, a}$, donc aussi de présentation finie comme module sur
$S_{I, a}$, voir le lemme d'Algèbre
\ref{algebra-lemma-finitely-presented-over-subring}.
Cela nous ramène au cas étudié au paragraphe suivant.

\medskip\noindent
Supposons $S = R[x_1, \ldots, x_n]$. Choisissons une présentation
$$
0 \to K \to S^{\oplus r} \to M \to 0.
$$
Soit $M_A$ le quotient de $M \otimes_R A$ par son sous-module de torsion,
voir le lemme \ref{lemma-torsion}. Alors $M_A$ est un module de type fini
sur $S_A = A[x_1, \ldots, x_n]$. Par le lemme
\ref{lemma-valuation-ring-torsion-free-flat}, nous voyons que $M_A$ est plat
sur $A$. Par le lemme \ref{lemma-flat-finite-type-valuation-ring-finite-presentation},
nous voyons que $M_A$ est de présentation finie. Il existe donc un nombre
fini d'éléments $k_1, \ldots, k_t \in S_A^{\oplus r}$ qui engendrent le
noyau de la présentation $S_A^{\oplus r} \to M_A$ comme $S_A$-module.
Pour tout choix de $a \in I \subset \mathfrak m$ satisfaisant (1), (2) et
(3), notons $M_{I, a}$ la transformée stricte de $M$ le long de
$R \to R[\frac{I}{a}]$. C'est un module de type fini sur
$S_{I, a} = R[\frac{I}{a}][x_1, \ldots, x_n]$. Par le lemme d'Algèbre
\ref{algebra-lemma-valuation-ring-colimit-affine-blowups}, nous avons
$A = \colim_{I, a} R[\frac{I}{a}]$. Cela implique que $S_A = \colim S_{I, a}$
et
$$
\colim M \otimes_R R[\textstyle{\frac{I}{a}}] = M \otimes_R A
$$
Choisissons $I, a$ et des relèvements $k_1, \ldots, k_t \in S_{I, a}^{\oplus r}$.
Puisque $M_A$ est le quotient de $M \otimes_R A$ par la torsion, nous voyons
que les images de $k_1, \ldots, k_t$ dans $M \otimes_R A$ sont annihilées par
un élément non nul $\alpha \in A$. Après avoir remplacé $I, a$ par un autre
couple (rappelons que la colimite est filtrante), nous pouvons supposer
$\alpha = x/a^n$ pour un certain $x \in I^n$ non nul. Alors
$x k_1, \ldots, x k_t$ s'envoient sur zéro dans $M \otimes_R A$.
Après un nouveau remplacement de $I, a$, nous pouvons donc supposer que
$x k_1, \ldots, x k_t$ s'envoient sur zéro dans
$M \otimes_R R[\frac{I}{a}]$ pour un certain $x \in R$ non nul. Enfin, en
remplaçant $I, a$ par $xI, xa$, nous pouvons supposer que
$k_1, \ldots, k_t$ s'envoient sur des éléments de torsion par une puissance
de $a$ de $M \otimes_R R[\frac{I}{a}]$. Pour tout couple $(I, a)$ de ce type,
posons
$$
M'_{I, a} = S_{I, a}^{\oplus r}/ \sum S_{I, a}k_j.
$$
Puisque $M_A = S_A^{\oplus r}/ \sum S_Ak_j$, nous voyons que
$M_A = \colim_{I, a} M'_{I, a}$. À ce stade, nous appliquons enfin le lemme
d'Algèbre \ref{algebra-lemma-flat-finite-presentation-limit-flat} (3) pour
conclure que $M'_{I, a}$ est plat pour un certain couple $(I, a)$ comme
ci-dessus. Ce lemme ne s'applique a priori pas au système des transformées
strictes
$$
M_{I, a} = (M \otimes_R R[\textstyle{\frac{I}{a}}])/a\text{-torsion}
$$
car les applications de transition peuvent ne pas satisfaire les hypothèses
du lemme. Mais maintenant, puisque la platitude implique l'absence de
torsion (lemme \ref{lemma-flat-torsion-free}) et puisque $M_{I, a}$ est le
quotient de $M'_{I, a}$ (car nous avons fait en sorte que les éléments
$k_1, \ldots, k_t$ s'envoient sur zéro dans $M_{I, a}$) par le sous-module de
torsion par une puissance de $a$, nous concluons également que
$M'_{I, a} = M_{I, a}$ pour un tel couple, ce qui conclut la preuve.
\end{proof}

\begin{lemma}
\label{lemma-blowup-fitting-ideal}
Soit $R$ un anneau. Soit $M$ un $R$-module de type fini.
Soit $k \geq 0$ et $I = \text{Fit}_k(M)$. Pour tout $a \in I$,
avec $R' = R[\frac{I}{a}]$, la transformée stricte
$$
M' = (M \otimes_R R')/a\text{-torsion}
$$
vérifie $\text{Fit}_k(M') = R'$.
\end{lemma}

\begin{proof}
Observons d'abord que $\text{Fit}_k(M \otimes_R R') = IR' = aR'$.
La première égalité vient du lemme \ref{lemma-fitting-ideal-basics}, partie (3),
et la seconde du lemme d'Algèbre \ref{algebra-lemma-affine-blowup}.
Par le lemme \ref{lemma-principal-fitting-ideal} et l'exactitude de la
localisation, nous voyons que $M'_{\mathfrak p'}$ peut être engendré par
$\leq k$ éléments pour tout idéal premier $\mathfrak p'$ de $R'$. Alors
$\text{Fit}_k(M') = R'$, par exemple par le lemme
\ref{lemma-fitting-ideal-generate-locally}.
\end{proof}

\begin{lemma}
\label{lemma-blowup-fitting-ideal-locally-free}
Soit $R$ un anneau. Soit $M$ un $R$-module de type fini.
Soit $k \geq 0$ et $I = \text{Fit}_k(M)$. Supposons que
$M_\mathfrak p$ soit libre de rang $k$ pour tout
$\mathfrak p \not \in V(I)$. Alors, pour tout $a \in I$,
avec $R' = R[\frac{I}{a}]$, la transformée stricte
$$
M' = (M \otimes_R R')/a\text{-torsion}
$$
est localement libre de rang $k$.
\end{lemma}

\begin{proof}
Par le lemme \ref{lemma-blowup-fitting-ideal}, nous avons
$\text{Fit}_k(M') = R'$. Par le lemme
\ref{lemma-fitting-ideal-finite-locally-free}, il suffit de montrer que
$\text{Fit}_{k - 1}(M') = 0$. Rappelons que $R' \subset R'_a = R_a$, voir
le lemme d'Algèbre \ref{algebra-lemma-affine-blowup}. Il suffit donc de
prouver que $\text{Fit}_{k - 1}(M')$ s'envoie sur zéro dans $R'_a = R_a$.
Puisque $(M')_a = M_a$, cela se réduit à montrer que
$\text{Fit}_{k - 1}(M_a) = 0$, car la formation des idéaux de Fitting commute
au changement de base d'après le lemme \ref{lemma-fitting-ideal-basics},
partie (3). Cela résulte de notre hypothèse que $M_a$ est localement libre
de type fini de rang $k$ (voir le lemme d'Algèbre
\ref{algebra-lemma-finite-projective}) et du lemme déjà cité
\ref{lemma-fitting-ideal-finite-locally-free}.
\end{proof}

\begin{lemma}
\label{lemma-blowup-module}
Soit $R$ un anneau. Soit $M$ un $R$-module de type fini. Soit $f \in R$
tel que $M_f$ soit localement libre de type fini de rang $r$.
Alors il existe un idéal $I \subset R$ de type fini avec
$V(f) = V(I)$ tel que, pour tout $a \in I$ avec $R' = R[\frac{I}{a}]$,
la transformée stricte
$$
M' = (M \otimes_R R')/a\text{-torsion}
$$
soit localement libre de rang $r$.
\end{lemma}

\begin{proof}
Choisissons une surjection $R^{\oplus n} \to M$. Choisissons un sous-module
$K \subset \Ker(R^{\oplus n} \to M)$ de type fini tel que
$R^{\oplus n}/K \to M$ devienne un isomorphisme après inversion de $f$.
C'est possible car $M_f$ est de présentation finie, par exemple par le lemme
d'Algèbre \ref{algebra-lemma-finite-projective}. Posons
$M_1 = R^{\oplus n}/K$ et supposons le lemme démontré pour $M_1$.
Soit $I \subset R$ l'idéal correspondant. Alors, pour $a \in I$,
l'application
$$
M_1' = (M_1 \otimes_R R')/a\text{-torsion}
\longrightarrow
M' = (M \otimes_R R')/a\text{-torsion}
$$
est surjective. C'est également un isomorphisme après inversion de $a$ dans
$R'$, car $R'_a = R_f$, voir le lemme d'Algèbre
\ref{algebra-lemma-blowup-in-principal}. Mais $a$ est un non-diviseur de zéro
sur $M'_1$, d'où le fait que l'application affichée est un isomorphisme.
Il suffit donc de démontrer le lemme lorsque $M$ est un $R$-module de
présentation finie.

\medskip\noindent
Supposons que $M$ soit un $R$-module de présentation finie.
Alors $J = \text{Fit}_r(M) \subset S$ est un idéal de type fini.
Nous affirmons que $I = fJ$ convient.

\medskip\noindent
Vérifions d'abord que $V(f) = V(I)$. L'inclusion $V(f) \subset V(I)$ est
claire. Réciproquement, si $f \not \in \mathfrak p$, alors
$\mathfrak p$ n'appartient pas à $V(J)$ par le lemme
\ref{lemma-fitting-ideal-generate-locally}. Ainsi
$\mathfrak p \not \in V(fJ) = V(I)$.

\medskip\noindent
Soit $a \in I$ et posons $R' = R[\frac{I}{a}]$. Nous pouvons écrire
$a = fb$ pour un certain $b \in J$. Par les lemmes d'Algèbre
\ref{algebra-lemma-affine-blowup} et \ref{algebra-lemma-blowup-add-principal},
nous voyons que $J R' = b R'$ et que $b$ est un non-diviseur de zéro dans
$R'$. Soit $\mathfrak p' \subset R' = R[\frac{I}{a}]$ un idéal premier.
Alors $JR'_{\mathfrak p'}$ est engendré par $b$. Il résulte du lemme
\ref{lemma-principal-fitting-ideal} que $M'_{\mathfrak p'}$ peut être
engendré par $r$ éléments. Puisque $M'$ est de type fini, il existe
$m_1, \ldots, m_r \in M'$ et $g \in R'$, $g \not \in \mathfrak p'$, tels
que l'application correspondante $(R')^{\oplus r} \to M'$ devienne
surjective après inversion de $g$.

\medskip\noindent
Enfin, considérons l'idéal $J' = \text{Fit}_{k - 1}(M')$.
Notons que $J' R'_g$ est engendré par les coefficients des relations entre
$m_1, \ldots, m_r$ (compatibilité de l'idéal de Fitting au changement de
base). Il suffit donc de montrer que $J' = 0$, voir le lemme
\ref{lemma-fitting-ideal-finite-locally-free}. Puisque $R'_a = R_f$
(lemme d'Algèbre \ref{algebra-lemma-blowup-in-principal}) et que
$M'_a = M_f$ est libre de rang $r$, nous voyons que $J'_a = 0$.
Comme $a$ est un non-diviseur de zéro dans $R'$, nous concluons que $J' = 0$,
ce qui achève la preuve.
\end{proof}












\section{Complétion et platitude}
\label{section-completion-flat}

\noindent
Dans cette section, nous examinons quand la complétion d'un module plat
« grand » est plate.

\begin{lemma}
\label{lemma-ui-completion-direct-sum-into-product}
Soit $R$ un anneau.
Soit $I \subset R$ un idéal.
Soit $A$ un ensemble.
Supposons que $R$ soit noethérien et complet pour $I$. Il existe une
application canonique
$$
\left(\bigoplus\nolimits_{\alpha \in A} R\right)^\wedge
\longrightarrow
\prod\nolimits_{\alpha \in A} R
$$
de la complétion $I$-adique de la somme directe vers le produit,
qui est universellement injective.
\end{lemma}

\begin{proof}
Par définition, un élément $x$ du membre de gauche est $x = (x_n)$, où
$x_n = (x_{n, \alpha}) \in \bigoplus\nolimits_{\alpha \in A} R/I^n$
et $x_{n, \alpha} = x_{n + 1, \alpha} \bmod I^n$.
Comme $R = R^\wedge$, nous voyons que, pour tout $\alpha$, il existe un
$y_\alpha \in R$ tel que $x_{n, \alpha} = y_\alpha \bmod I^n$.
Notons que, pour chaque $n$, il n'y a qu'un nombre fini de $\alpha$ tels que
les éléments $x_{n, \alpha}$ soient non nuls. Réciproquement, étant donné
$(y_\alpha) \in \prod_\alpha R$ tel que, pour chaque $n$, il n'y ait qu'un
nombre fini de $\alpha$ pour lesquels $y_{\alpha} \bmod I^n$ soit non nul,
nous obtenons un élément du membre de gauche. Nous pouvons donc voir un
élément du membre de gauche comme une « somme convergente » infinie
$\sum_\alpha y_\alpha$, avec $y_\alpha \in R$, telle que, pour chaque $n$,
il n'y ait qu'un nombre fini de $y_\alpha$ non nuls modulo $I^n$.
L'application affichée envoie cet élément sur l'élément $(y_\alpha)$ du
produit. En particulier, elle est injective.

\medskip\noindent
Soit $Q$ un $R$-module de type fini. Nous devons montrer que l'application
$$
Q \otimes_R \left(\bigoplus\nolimits_{\alpha \in A} R\right)^\wedge
\longrightarrow
Q \otimes_R \left(\prod\nolimits_{\alpha \in A} R\right)
$$
est injective, voir le théorème d'Algèbre
\ref{algebra-theorem-universally-exact-criteria}.
Choisissons une présentation $R^{\oplus k} \to R^{\oplus m} \to Q \to 0$
et notons $q_1, \ldots, q_m \in Q$ les générateurs correspondants de $Q$.
Par Artin-Rees (lemme d'Algèbre \ref{algebra-lemma-Artin-Rees}), il existe
une constante $c$ telle que
$\Im(R^{\oplus k} \to R^{\oplus m}) \cap (I^N)^{\oplus m}
\subset \Im((I^{N - c})^{\oplus k} \to R^{\oplus m})$.
Considérons le diagramme
$$
\xymatrix{
\bigoplus_{l = 1}^k \left(\bigoplus\nolimits_{\alpha \in A} R\right)^\wedge
\ar[r] \ar[d] &
\bigoplus_{j = 1}^m \left(\bigoplus\nolimits_{\alpha \in A} R\right)^\wedge
\ar[r] \ar[d] &
Q \otimes_R \left(\bigoplus\nolimits_{\alpha \in A} R\right)^\wedge
\ar[r] \ar[d] &
0 \\
\bigoplus_{l = 1}^k \left(\prod\nolimits_{\alpha \in A} R\right)
\ar[r] &
\bigoplus_{j = 1}^m \left(\prod\nolimits_{\alpha \in A} R\right)
\ar[r] &
Q \otimes_R \left(\prod\nolimits_{\alpha \in A} R\right)
\ar[r] &
0
}
$$
à lignes exactes. Prenons un élément $\sum_j \sum_\alpha y_{j, \alpha}$ de
$\bigoplus_{j = 1, \ldots, m}
\left(\bigoplus\nolimits_{\alpha \in A} R\right)^\wedge$.
Si cet élément s'envoie sur zéro dans le module
$Q \otimes_R \left(\prod\nolimits_{\alpha \in A} R\right)$,
nous voyons en particulier que
$\sum_j q_j \otimes y_{j, \alpha} = 0$ dans $Q$ pour tout $\alpha$.
Il existe donc un élément
$(z_{1, \alpha}, \ldots, z_{k, \alpha}) \in \bigoplus_{l = 1, \ldots, k} R$
qui s'envoie sur
$(y_{1, \alpha}, \ldots, y_{m, \alpha}) \in \bigoplus_{j = 1, \ldots, m} R$.
De plus, si $y_{j, \alpha} \in I^{N_\alpha}$ pour $j = 1, \ldots, m$, alors
nous pouvons supposer que $z_{l, \alpha} \in I^{N_\alpha - c}$ pour
$l = 1, \ldots, k$. Ainsi la somme $\sum_l \sum_\alpha z_{l, \alpha}$ est
« convergente » et définit un élément de
$\bigoplus_{l = 1, \ldots, k}
\left(\bigoplus\nolimits_{\alpha \in A} R\right)^\wedge$
qui s'envoie sur l'élément $\sum_j \sum_\alpha y_{j, \alpha}$ de départ.
La flèche verticale de droite est donc injective, ce qui conclut la preuve.
\end{proof}

\noindent
Le lemme suivant peut également se déduire du lemme
\ref{lemma-limit-flat} ci-dessous.

\begin{lemma}
\label{lemma-completed-direct-sum-flat}
Soit $R$ un anneau. Soit $I \subset R$ un idéal. Soit $A$ un ensemble.
Supposons que $R$ soit noethérien. La complétion
$(\bigoplus\nolimits_{\alpha \in A} R)^\wedge$
est un $R$-module plat.
\end{lemma}

\begin{proof}
Notons $R^\wedge$ la complétion de $R$ pour $I$. Comme
$R \to R^\wedge$ est plat par le lemme d'Algèbre
\ref{algebra-lemma-completion-flat}, il suffit de démontrer que
$(\bigoplus\nolimits_{\alpha \in A} R)^\wedge$ est un
$R^\wedge$-module plat (utiliser le lemme d'Algèbre
\ref{algebra-lemma-composition-flat}).
Puisque
$$
(\bigoplus\nolimits_{\alpha \in A} R)^\wedge
=
(\bigoplus\nolimits_{\alpha \in A} R^\wedge)^\wedge
$$
Nous pouvons remplacer $R$ par $R^\wedge$ et supposer que $R$ est complet
pour $I$ (voir le lemme d'Algèbre
\ref{algebra-lemma-completion-complete}). Dans ce cas, le lemme
\ref{lemma-ui-completion-direct-sum-into-product} nous dit que l'application
$(\bigoplus\nolimits_{\alpha \in A} R)^\wedge \to \prod_{\alpha \in A} R$
est universellement injective. Ainsi, par le lemme d'Algèbre
\ref{algebra-lemma-ui-flat-domain}, il suffit de montrer que
$\prod_{\alpha \in A} R$ est plat. Par la proposition d'Algèbre
\ref{algebra-proposition-characterize-coherent} (et le lemme d'Algèbre
\ref{algebra-lemma-Noetherian-coherent}), nous voyons que
$\prod_{\alpha \in A} R$ est plat.
\end{proof}

\begin{lemma}
\label{lemma-tor-strictly-pro-zero}
\begin{reference}
Ceci est \cite[Lemme 9.9]{quillenhomology} ; notons que l'auteur a oublié
le mot « strict » dans l'énoncé, bien qu'il fût clairement sous-entendu.
\end{reference}
Soit $A$ un anneau noethérien. Soit $I$ un idéal de $A$.
Soit $M$ un $A$-module de type fini. Pour tout $p > 0$, il existe un $c > 0$
tel que $\text{Tor}_p^A(M, A/I^n) \to \text{Tor}_p^A(M, A/I^{n - c})$
qui est nulle pour tout $n \geq c$.
\end{lemma}

\begin{proof}
Preuve pour $p = 1$. Choisissons une suite exacte courte
$0 \to K \to A^{\oplus t} \to M \to 0$. Alors
$\text{Tor}_1^A(M, A/I^n) = K \cap (I^n)^{\oplus t}/I^nK$.
Par Artin-Rees (lemme d'Algèbre \ref{algebra-lemma-Artin-Rees}), il existe
une constante $c \geq 0$ telle que
$K \cap (I^n)^{\oplus t} \subset I^{n - c}K$ pour $n \geq c$. Cela donne le
résultat pour $p = 1$. Pour $p > 1$, nous avons
$\text{Tor}_p^A(M, A/I^n) = \text{Tor}^A_{p - 1}(K, A/I^n)$.
Le lemme se déduit donc par récurrence.
\end{proof}

\begin{lemma}
\label{lemma-limit-flat}
Soit $A$ un anneau noethérien. Soit $I$ un idéal de $A$. Soit
$(M_n)$ un système inverse de $A$-modules tel que
\begin{enumerate}
\item $M_n$ est un $A/I^n$-module plat,
\item $M_{n + 1} \to M_n$ est surjective.
\end{enumerate}
Alors $M = \lim M_n$ est un $A$-module plat et
$Q \otimes_A M = \lim Q \otimes_A M_n$ pour tout $A$-module $Q$ de type fini.
\end{lemma}

\begin{proof}
Montrons d'abord que $Q \otimes_A M = \lim Q \otimes_A M_n$ pour tout
$A$-module $Q$ de type fini. Choisissons une résolution
$F_2 \to F_1 \to F_0 \to Q \to 0$ par des $A$-modules $F_i$ libres de type
fini. Alors
$$
F_2 \otimes_A M_n \to F_1 \otimes_A M_n \to F_0 \otimes_A M_n
$$
est un complexe de chaînes dont l'homologie en degré $0$ est
$Q \otimes_A M_n$ et dont l'homologie en degré $1$ est
$$
\text{Tor}_1^A(Q, M_n) = \text{Tor}_1^A(Q, A/I^n) \otimes_{A/I^n} M_n
$$
car $M_n$ est plat sur $A/I^n$. Par le lemme
\ref{lemma-tor-strictly-pro-zero}, nous voyons que ce système est
essentiellement constant (de valeur $0$). Il résulte du lemme d'Homologie
\ref{homology-lemma-apply-Mittag-Leffler-again} que
$\lim Q \otimes_A M_n =
\Coker(\lim F_1 \otimes_A M_n \to \lim F_0 \otimes_A M_n)$.
Puisque $F_i$ est libre de type fini, ceci vaut
$\Coker(F_1 \otimes_A M \to F_0 \otimes_A M) = Q \otimes_A M$.

\medskip\noindent
Ensuite, soit $Q \to Q'$ une application injective de $A$-modules de type
fini. Nous devons montrer que $Q \otimes_A M \to Q' \otimes_A M$ est
injective (lemme d'Algèbre \ref{algebra-lemma-flat}). D'après ce qui précède,
nous voyons
$$
\Ker(Q \otimes_A M \to Q' \otimes_A M) =
\Ker(\lim Q \otimes_A M_n \to \lim Q' \otimes_A M_n).
$$
Pour chaque $n$, nous avons une suite exacte
$$
\text{Tor}_1^A(Q', M_n) \to \text{Tor}_1^A(Q'', M_n) \to
Q \otimes_A M_n \to Q' \otimes_A M_n
$$
où $Q'' = \Coker(Q \to Q')$. Nous avons vu ci-dessus que les systèmes
inverses de Tor sont essentiellement constants de valeur $0$. Il résulte du
lemme d'Homologie \ref{homology-lemma-apply-Mittag-Leffler-again} que la
limite inverse des applications les plus à droite est injective.
\end{proof}

\begin{lemma}
\label{lemma-flat-after-completion}
Soit $R$ un anneau. Soit $I \subset R$ un idéal. Soit $M$ un
$R$-module. Supposons que
\begin{enumerate}
\item $I$ soit de type fini,
\item $R/I$ soit noethérien,
\item $M/IM$ soit plat sur $R/I$,
\item $\text{Tor}_1^R(M, R/I) = 0$.
\end{enumerate}
Alors la complétion $I$-adique $R^\wedge$ est un anneau noethérien et
$M^\wedge$ est plat sur $R^\wedge$.
\end{lemma}

\begin{proof}
Par le lemme d'Algèbre \ref{algebra-lemma-what-does-it-mean}, les modules
$M/I^nM$ sont plats sur $R/I^n$ pour tout $n$. Par le lemme d'Algèbre
\ref{algebra-lemma-hathat-finitely-generated}, nous avons (a) $R^\wedge$ et
$M^\wedge$ complets pour la topologie $I$-adique et (b)
$R/I^n = R^\wedge/I^nR^\wedge$ pour tout $n$. Par le lemme d'Algèbre
\ref{algebra-lemma-completion-Noetherian}, l'anneau $R^\wedge$ est
noethérien. En appliquant le lemme \ref{lemma-limit-flat}, nous concluons
que $M^\wedge = \lim M/I^nM$ est plat comme $R^\wedge$-module.
\end{proof}
\section{Le complexe de Koszul}
\label{section-koszul}

\noindent
Nous définissons le complexe de Koszul comme suit.

\begin{definition}
\label{definition-koszul}
Soit $R$ un anneau. Soit $\varphi : E \to R$ une application de
$R$-modules. Le {\it complexe de Koszul} $K_\bullet(\varphi)$ associé à
$\varphi$ est l'algèbre différentielle graduée commutative définie comme suit :
\begin{enumerate}
\item l'algèbre graduée sous-jacente est l'algèbre extérieure
$K_\bullet(\varphi) = \wedge(E)$,
\item la différentielle $d : K_\bullet(\varphi) \to K_\bullet(\varphi)$
est l'unique dérivation telle que $d(e) = \varphi(e)$ pour tout
$e \in E = K_1(\varphi)$.
\end{enumerate}
\end{definition}

\noindent
Explicitement, si $e_1 \wedge \ldots \wedge e_n$ est l'un des générateurs de
degré $n$ de $K_\bullet(\varphi)$, alors
$$
d(e_1 \wedge \ldots \wedge e_n) =
\sum\nolimits_{i = 1, \ldots, n} (-1)^{i + 1}
\varphi(e_i)e_1 \wedge \ldots \wedge \widehat{e_i} \wedge \ldots \wedge e_n.
$$
On voit immédiatement que cela donne une dérivation bien définie
sur l'algèbre tensorielle, qui annule $e \otimes e$ et se factorise donc
par l'algèbre extérieure.

\medskip\noindent
Nous supposons souvent que $E$ est un module libre de type fini, disons
$E = R^{\oplus n}$. Dans ce cas, l'application $\varphi$ est donnée par une
suite d'éléments
$f_1, \ldots, f_n \in R$.

\begin{definition}
\label{definition-koszul-complex}
Soit $R$ un anneau et soient $f_1, \ldots, f_r \in R$. Le
{\it complexe de Koszul sur $f_1, \ldots, f_r$} est le complexe de Koszul
associé à l'application $(f_1, \ldots, f_r) : R^{\oplus r} \to R$.
On le note $K_\bullet(f_\bullet)$, $K_\bullet(f_1, \ldots, f_r)$,
$K_\bullet(R, f_1, \ldots, f_r)$, ou $K_\bullet(R, f_\bullet)$.
\end{definition}

\noindent
Bien entendu, si $E$ est localement libre de type fini, alors
$K_\bullet(\varphi)$ est localement sur $\Spec(R)$ isomorphe à un complexe de Koszul
$K_\bullet(f_1, \ldots, f_r)$.
Ce complexe possède de nombreuses propriétés formelles intéressantes.

\begin{lemma}
\label{lemma-functorial}
Soient $\varphi : E \to R$ et $\varphi' : E' \to R$ des applications de
$R$-modules. Soit $\psi : E \to E'$ une application de $R$-modules telle que
$\varphi' \circ \psi = \varphi$. Alors $\psi$ induit un
homomorphisme d'algèbres différentielles graduées
$K_\bullet(\varphi) \to K_\bullet(\varphi')$.
\end{lemma}

\begin{proof}
Cela résulte immédiatement des définitions.
\end{proof}

\begin{lemma}
\label{lemma-change-basis}
Soit $f_1, \ldots, f_r \in R$ une suite.
Soit $(x_{ij})$ une matrice inversible $r \times r$ à
coefficients dans $R$. Alors les complexes
$K_\bullet(f_\bullet)$ et
$$
K_\bullet(\sum x_{1j}f_j, \sum x_{2j}f_j, \ldots, \sum x_{rj}f_j)
$$
sont isomorphes.
\end{lemma}

\begin{proof}
Posons $g_i = \sum x_{ij}f_j$.
La matrice $(x_{ji})$ donne un isomorphisme $x : R^{\oplus r} \to R^{\oplus r}$
tel que $(g_1, \ldots, g_r) = (f_1, \ldots, f_r) \circ x$.
L'assertion résulte donc de la fonctorialité du complexe de Koszul
décrite au
Lemme \ref{lemma-functorial}.
\end{proof}

\begin{lemma}
\label{lemma-homotopy-koszul-abstract}
Soit $R$ un anneau. Soit $\varphi : E \to R$ une application de $R$-modules.
Soit $e \in E$ d'image $f = \varphi(e)$ dans $R$. Alors
$$
f = de + ed
$$
comme endomorphismes de $K_\bullet(\varphi)$.
\end{lemma}

\begin{proof}
Cela est vrai car $d(ea) = d(e)a - ed(a) = fa - ed(a)$.
\end{proof}

\begin{lemma}
\label{lemma-homotopy-koszul}
Soit $R$ un anneau. Soit $f_1, \ldots, f_r \in R$ une suite.
La multiplication par $f_i$ sur $K_\bullet(f_\bullet)$ est homotope à
zéro et, en particulier, les modules d'homologie $H_i(K_\bullet(f_\bullet))$
sont annulés par l'idéal $(f_1, \ldots, f_r)$.
\end{lemma}

\begin{proof}
C'est un cas particulier du
Lemme \ref{lemma-homotopy-koszul-abstract}.
\end{proof}

\noindent
Dans
Catégories dérivées, Section \ref{derived-section-cones},
nous avons défini le cône d'un morphisme de complexes de cochaînes.
Le cône $C(f)_\bullet$ d'un morphisme de complexes de chaînes
$f : A_\bullet \to B_\bullet$ est le complexe $C(f)_\bullet$ donné par
$C(f)_n = B_n \oplus A_{n - 1}$ et de différentielle
\begin{equation}
\label{equation-differential-cone}
d_{C(f), n} =
\left(
\begin{matrix}
d_{B, n} & f_{n - 1} \\
0 & -d_{A, n - 1}
\end{matrix}
\right)
\end{equation}
Il est muni de morphismes canoniques de complexes
$i : B_\bullet \to C(f)_\bullet$ et
$p : C(f)_\bullet \to A_\bullet[-1]$
induits par les applications évidentes $B_n \to C(f)_n \to A_{n - 1}$.

\begin{lemma}
\label{lemma-cone-koszul-abstract}
Soit $R$ un anneau. Soit $\varphi : E \to R$ une application de $R$-modules.
Soit $f \in R$. Posons $E' = E \oplus R$ et définissons $\varphi' : E' \to R$
par $\varphi$ sur $E$ et par la multiplication par $f$ sur $R$.
Le complexe $K_\bullet(\varphi')$ est isomorphe au
cône de l'application de complexes
$$
f :
K_\bullet(\varphi)
\longrightarrow
K_\bullet(\varphi).
$$
\end{lemma}

\begin{proof}
Notons $e_0 \in E'$ l'élément $1 \in R \subset R \oplus E$.
D'après notre définition du cône ci-dessus, on voit que
$$
C(f)_n = K_n(\varphi) \oplus K_{n - 1}(\varphi) =
\wedge^n(E) \oplus \wedge^{n - 1}(E) = \wedge^n(E')
$$
où, pour le dernier $=$, on envoie $(0, e_1 \wedge \ldots \wedge e_{n - 1})$
sur $e_0 \wedge e_1 \wedge \ldots \wedge e_{n - 1}$ dans $\wedge^n(E')$.
Un calcul montre que cet isomorphisme est compatible avec les
différentielles. En effet, c'est clair pour les éléments du premier
facteur direct puisque $\varphi'|_E = \varphi$ et que la restriction de
$d_{C(f)}$ au premier facteur direct n'est autre que $d_{K_\bullet(\varphi)}$.
D'autre part, si $e_1 \wedge \ldots \wedge e_{n - 1}$
appartient au second facteur direct, alors
$$
d_{C(f)}(0, e_1 \wedge \ldots \wedge e_{n - 1}) =
fe_1 \wedge \ldots \wedge e_{n - 1}
- d_{K_\bullet(\varphi)}(e_1 \wedge \ldots \wedge e_{n - 1})
$$
et, d'autre part,
\begin{align*}
& d_{K_\bullet(\varphi')}(0, e_0 \wedge e_1 \wedge \ldots \wedge e_{n - 1}) \\
& =
\sum\nolimits_{i = 0, \ldots, n - 1}
(-1)^i \varphi'(e_i)e_0 \wedge \ldots \wedge \widehat{e_i}
\wedge \ldots \wedge e_{n - 1} \\
& =
fe_1 \wedge \ldots \wedge e_{n - 1} +
\sum\nolimits_{i = 1, \ldots, n - 1}
(-1)^i \varphi(e_i)e_0 \wedge \ldots \wedge \widehat{e_i}
\wedge \ldots \wedge e_{n - 1} \\
& =
fe_1 \wedge \ldots \wedge e_{n - 1} -
e_0 \left(\sum\nolimits_{i = 1, \ldots, n - 1}
(-1)^{i + 1} \varphi(e_i)e_1 \wedge \ldots \wedge \widehat{e_i}
\wedge \ldots \wedge e_{n - 1}\right)
\end{align*}
ce qui est l'image du résultat du calcul précédent.
\end{proof}

\begin{lemma}
\label{lemma-cone-koszul}
Soit $R$ un anneau. Soit $f_1, \ldots, f_r$ une suite d'éléments
de $R$. Le complexe $K_\bullet(f_1, \ldots, f_r)$ est isomorphe au
cône de l'application de complexes
$$
f_r :
K_\bullet(f_1, \ldots, f_{r - 1})
\longrightarrow
K_\bullet(f_1, \ldots, f_{r - 1}).
$$
\end{lemma}

\begin{proof}
C'est un cas particulier du
Lemme \ref{lemma-cone-koszul-abstract}.
\end{proof}

\begin{lemma}
\label{lemma-cone-squared}
Soit $R$ un anneau. Soit $A_\bullet$ un complexe de $R$-modules.
Soient $f, g \in R$. Soit $C(f)_\bullet$ le cône de
$f : A_\bullet \to A_\bullet$. Définissons de même $C(g)_\bullet$ et
$C(fg)_\bullet$. Alors $C(fg)_\bullet$ est homotopiquement équivalent au
cône d'une application
$$
C(f)_\bullet[1] \longrightarrow C(g)_\bullet
$$
\end{lemma}

\begin{proof}
Nous le démontrons d'abord lorsque $A_\bullet$ est le complexe constitué de
$R$ placé en degré $0$. Dans ce cas, le complexe $C(f)_\bullet$ est le
complexe
$$
\ldots \to 0 \to R \xrightarrow{f} R \to 0 \to \ldots
$$
où $R$ est placé aux degrés (homologiques) $1$ et $0$. L'application de
complexes que nous utilisons est
$$
\xymatrix{
0 \ar[r] \ar[d] &
0 \ar[r] \ar[d] &
R \ar[r]^f \ar[d]^1 &
R \ar[r] \ar[d] & 0 \ar[d] \\
0 \ar[r] & R \ar[r]^g & R \ar[r] & 0 \ar[r] & 0
}
$$
Son cône est le complexe de chaînes constitué de $R^{\oplus 2}$ placé aux
degrés $1$ et $0$, et dont la différentielle (\ref{equation-differential-cone}) est
$$
\left(
\begin{matrix}
g & 1 \\
0 & -f
\end{matrix}
\right) :
R^{\oplus 2} \longrightarrow R^{\oplus 2}
$$
Pour voir que ce complexe de chaînes est homotope à
$C(fg)_\bullet$, c'est-à-dire à $R \xrightarrow{fg} R$,
considérons les applications de complexes
$$
\xymatrix{
R \ar[d]_{(1, -g)} \ar[r]_{fg} & R \ar[d]^{(0, 1)} \\
R^{\oplus 2} \ar[r] & R^{\oplus 2}
}
\quad\quad
\xymatrix{
R^{\oplus 2} \ar[d]_{(1, 0)} \ar[r] & R^{\oplus 2} \ar[d]^{(f, 1)} \\
R \ar[r]^{fg} & R
}
$$
avec les notations évidentes.
La composée de ces deux applications dans un sens est
l'identité de $C(fg)_\bullet$, mais dans l'autre sens
ce n'est pas l'identité. Nous omettons d'écrire l'homotopie nécessaire.

\medskip\noindent
Pour voir que le résultat vaut en général, nous utilisons le foncteur
$K_\bullet \mapsto \text{Tot}(A_\bullet \otimes_R K_\bullet)$
sur la catégorie des complexes, qui est compatible aux homotopies
et aux cônes. Nous écrivons alors $C(f)_\bullet$ et $C(g)_\bullet$
comme les complexes totaux des complexes doubles
$$
(R \xrightarrow{f} R) \otimes_R A_\bullet
\quad\text{et}\quad
(R \xrightarrow{g} R) \otimes_R A_\bullet
$$
et nous déduisons ainsi le résultat du cas particulier traité ci-dessus.
Certains détails sont omis.
\end{proof}

\begin{lemma}
\label{lemma-koszul-mult-abstract}
Soit $R$ un anneau. Soit $\varphi : E \to R$ une application de $R$-modules.
Soient $f, g \in R$. Posons $E' = E \oplus R$ et définissons
$\varphi'_f, \varphi'_g, \varphi'_{fg} : E' \to R$
par $\varphi$ sur $E$ et par la multiplication par $f, g, fg$ sur $R$.
Le complexe $K_\bullet(\varphi'_{fg})$ est homotopiquement équivalent au
cône d'une application de complexes
$$
K_\bullet(\varphi'_f)[1]
\longrightarrow
K_\bullet(\varphi'_g).
$$
\end{lemma}

\begin{proof}
D'après le
Lemme \ref{lemma-cone-koszul-abstract},
le complexe $K_\bullet(\varphi'_f)$ est isomorphe au cône de la
multiplication par $f$ sur $K_\bullet(\varphi)$, et de même
dans les deux autres cas. Le lemme résulte donc du
Lemme \ref{lemma-cone-squared}.
\end{proof}

\begin{lemma}
\label{lemma-koszul-mult}
Soit $R$ un anneau. Soit $f_1, \ldots, f_{r - 1}$ une suite d'éléments
de $R$. Soient $f, g \in R$. Le complexe
$K_\bullet(f_1, \ldots, f_{r - 1}, fg)$
est homotopiquement équivalent au cône d'une application de complexes
$$
K_\bullet(f_1, \ldots, f_{r - 1}, f)[1]
\longrightarrow
K_\bullet(f_1, \ldots, f_{r - 1}, g)
$$
\end{lemma}

\begin{proof}
C'est un cas particulier du
Lemme \ref{lemma-koszul-mult-abstract}.
\end{proof}

\begin{lemma}
\label{lemma-join-sequences-koszul-complex}
Soit $R$ un anneau.
Soient $f_1, \ldots, f_r$, $g_1, \ldots, g_s$ des éléments de $R$.
Alors il existe un isomorphisme de complexes de Koszul
$$
K_\bullet(R, f_1, \ldots, f_r, g_1, \ldots, g_s) =
\text{Tot}(K_\bullet(R, f_1, \ldots, f_r) \otimes_R
K_\bullet(R, g_1, \ldots, g_s)).
$$
\end{lemma}

\begin{proof}
Omis. Indication : si $K_\bullet(R, f_1, \ldots, f_r)$ est engendré comme
algèbre différentielle graduée par $x_1, \ldots, x_r$ avec $\text{d}(x_i) = f_i$
et si $K_\bullet(R, g_1, \ldots, g_s)$ est engendré comme
algèbre différentielle graduée par $y_1, \ldots, y_s$ avec $\text{d}(y_j) = g_j$,
alors on peut considérer $K_\bullet(R, f_1, \ldots, f_r, g_1, \ldots, g_s)$
comme l'algèbre différentielle graduée engendrée par la suite d'éléments
$x_1, \ldots, x_r, y_1, \ldots, y_s$ avec $\text{d}(x_i) = f_i$
et $\text{d}(y_j) = g_j$.
\end{proof}





\section{Le complexe de {\v C}ech alterné augmenté}
\label{section-alternating-cech}

\noindent
Soit $R$ un anneau. Soient $f_1, \ldots, f_r \in R$. Le complexe de
{\v C}ech alterné augmenté de $R$ est le complexe de cochaînes
$$
R \to \bigoplus\nolimits_{i_0} R_{f_{i_0}} \to
\bigoplus\nolimits_{i_0 < i_1} R_{f_{i_0}f_{i_1}} \to
\ldots \to R_{f_1\ldots f_r}
$$
où $R$ est en degré $0$, le terme $\bigoplus_{i_0} R_{f_{i_0}}$
est en degré $1$, et ainsi de suite. Les applications sont définies comme suit :
\begin{enumerate}
\item L'application $R \to \bigoplus\nolimits_{i_0} R_{f_{i_0}}$
est donnée par les applications canoniques $R \to R_{f_{i_0}}$.
\item Étant donnés $1 \leq i_0 < \ldots < i_{p + 1} \leq r$ et
$0 \leq j \leq p + 1$, nous avons l'application canonique de localisation
$$
R_{f_{i_0} \ldots \hat f_{i_j} \ldots f_{i_{p + 1}}} \to
R_{f_{i_0} \ldots f_{i_{p + 1}}}
$$
\item Les différentielles utilisent les applications canoniques de (2) avec
le signe $(-1)^j$.
\end{enumerate}
Si $M$ est un $R$-module quelconque, le complexe de {\v C}ech alterné
augmenté de $M$ est le complexe de cochaînes construit de façon analogue
$$
M \to \bigoplus\nolimits_{i_0} M_{f_{i_0}} \to
\bigoplus\nolimits_{i_0 < i_1} M_{f_{i_0}f_{i_1}} \to
\ldots \to M_{f_1\ldots f_r}
$$
où $M$ est en degré $0$ comme précédemment.

\begin{lemma}
\label{lemma-extended-alternating-is-complex}
Les complexes de {\v C}ech alternés augmentés définis ci-dessus
sont des complexes de $R$-modules.
\end{lemma}

\begin{proof}
Omis.
\end{proof}

\begin{lemma}
\label{lemma-extended-alternating-form-module}
Soit $R$ un anneau. Soient $f_1, \ldots, f_r \in R$. Soit $M$ un $R$-module.
Le complexe de {\v C}ech alterné augmenté de $M$ est le produit tensoriel
sur $R$ de $M$ avec le complexe de {\v C}ech alterné augmenté de $R$.
\end{lemma}

\begin{proof}
Omis.
\end{proof}

\begin{lemma}
\label{lemma-extended-alternating-base-change}
Soit $R$ un anneau. Soient $f_1, \ldots, f_r \in R$. Soit $M$ un $R$-module.
Soit $R \to S$ un homomorphisme d'anneaux, notons $g_1, \ldots, g_r \in S$
les images de $f_1, \ldots, f_r$, et posons $N = M \otimes_R S$.
Le complexe de {\v C}ech alterné augmenté construit à partir de
$S$, $g_1, \ldots, g_r$ et $N$ est le produit tensoriel du
complexe de {\v C}ech alterné augmenté de $M$ avec $S$ sur $R$.
\end{lemma}

\begin{proof}
Omis.
\end{proof}

\begin{lemma}
\label{lemma-extended-alternating-homotopy-zero}
Soit $R$ un anneau. Soient $f_1, \ldots, f_r \in R$. Soit $M$ un $R$-module.
S'il existe un $i \in \{1, \ldots, r\}$ tel que $f_i$ soit inversible, alors
le complexe de {\v C}ech alterné augmenté
de $M$ est homotopiquement équivalent à $0$.
\end{lemma}

\begin{proof}
Nous utiliserons la notation suivante : une cochaîne $x$ de degré $p + 1$
du complexe de {\v C}ech alterné augmenté de $M$ est
$x = (x_{i_0 \ldots i_p})$, où $x_{i_0 \ldots i_p}$ appartient à
$M_{f_{i_0} \ldots f_{i_p}}$. Avec cette notation, nous avons
$$
d(x)_{i_0 \ldots i_{p + 1}} =
\sum\nolimits_j (-1)^j x_{i_0 \ldots \hat i_j \ldots i_{p + 1}}
$$
Comme homotopie, nous utilisons les applications
$$
h : \text{cochaînes de degré }p + 2 \to \text{cochaînes de degré }p + 1
$$
données par la règle
$$
h(x)_{i_0 \ldots i_p} = 0 \text{ si } i \in \{i_0, \ldots, i_p\}
\text{ et }
h(x)_{i_0 \ldots i_p} = 
(-1)^j x_{i_0 \ldots i_j i i_{j + 1} \ldots i_p} \text{ sinon}
$$
Ici, $j$ est l'unique indice tel que $i_j < i < i_{j + 1}$ dans le
second cas ; en outre, puisque $f_i$ est inversible, nous avons l'égalité
$$
M_{f_{i_0} \ldots f_{i_p}} =
M_{f_{i_0} \ldots f_{i_j} f_i f_{i_{j + 1}} \ldots f_{i_p}}
$$
qui permet de considérer
$(-1)^j x_{i_0 \ldots i_j i i_{j + 1} \ldots i_p}$
comme un élément de $M_{f_{i_0} \ldots f_{i_p}}$.
Nous allons montrer par un calcul que $d h + h d$ est égal
à l'opposé de l'application identité, ce qui achèvera la démonstration.
Pour ce faire, fixons une cochaîne $x$ de degré $p + 1$ et soient
$1 \leq i_0 < \ldots < i_p \leq r$.

\medskip\noindent
Cas I : $i \in \{i_0, \ldots, i_p\}$. Disons que $i = i_t$. Nous avons alors
$h(d(x))_{i_0 \ldots i_p} = 0$. D'autre part, nous avons
$$
d(h(x))_{i_0 \ldots i_p} =
\sum (-1)^j h(x)_{i_0 \ldots \hat i_j \ldots i_p} =
(-1)^t h(x)_{i_0 \ldots \hat i \ldots i_p} =
(-1)^t (-1)^{t - 1} x_{i_0 \ldots i_p}
$$
Ainsi $(dh + hd)(x)_{i_0 \ldots i_p} = -x_{i_0 \ldots i_p}$, comme voulu.

\medskip\noindent
Cas II : $i \not \in \{i_0, \ldots, i_p\}$. Soit $j$ tel que
$i_j < i < i_{j + 1}$. On voit alors que
\begin{align*}
h(d(x))_{i_0 \ldots i_p}
& =
(-1)^j d(x)_{i_0 \ldots i_j i i_{j + 1} \ldots i_p} \\
& =
\sum\nolimits_{j' \leq j} (-1)^{j + j'}
x_{i_0 \ldots \hat i_{j'} \ldots i_j i i_{j + 1} \ldots i_p} -
x_{i_0 \ldots i_p} \\
&
+ \sum\nolimits_{j' > j} (-1)^{j + j' + 1}
x_{i_0 \ldots i_j i i_{j + 1} \ldots \hat i_{j'} \ldots i_p}
\end{align*}
D'autre part, nous avons
\begin{align*}
d(h(x))_{i_0 \ldots i_p}
& =
\sum\nolimits_{j'} (-1)^{j'} h(x)_{i_0 \ldots \hat i_{j'} \ldots i_p} \\
& =
\sum\nolimits_{j' \leq j} (-1)^{j' + j - 1}
x_{i_0 \ldots \hat i_{j'} \ldots i_j i i_{j + 1} \ldots i_p} \\
& +
\sum\nolimits_{j' > j} (-1)^{j' + j}
x_{i_0 \ldots i_j i i_{j + 1} \ldots \hat i_{j'} \ldots i_p}
\end{align*}
En additionnant ces expressions, nous obtenons
$(dh + hd)(x)_{i_0 \ldots i_p} = - x_{i_0 \ldots i_p}$
comme voulu.
\end{proof}

\begin{lemma}
\label{lemma-extended-alternating-torsion}
Soit $R$ un anneau. Soient $f_1, \ldots, f_r \in R$. Soit $M$ un $R$-module.
Soit $H^q$ le $q$-ième module de cohomologie du complexe de
{\v C}ech alterné augmenté de $M$. Alors
\begin{enumerate}
\item $H^q = 0$ si $q \not \in [0, r]$,
\item pour $x \in H^i$, il existe un $n \geq 1$ tel que $f_i^n x = 0$
pour $i = 1, \ldots, r$,
\item le support de $H^q$ est contenu dans $V(f_1, \ldots, f_r)$,
\item s'il existe un $f \in (f_1, \ldots, f_r)$ qui agit inversiblement
sur $M$, alors $H^q = 0$.
\end{enumerate}
\end{lemma}

\begin{proof}
L'assertion (1) résulte du fait que le complexe de {\v C}ech alterné augmenté
est nul en degrés $< 0$ et $> r$. Pour démontrer (2), il suffit de montrer
que, pour chaque $i$, il existe un $n \geq 1$ tel que $f_i^n x = 0$. Pour
cela, il suffit de montrer que $(H^q)_{f_i} = 0$. La localisation étant
exacte, $(H^q)_{f_i}$ est le $q$-ième module de cohomologie de la localisation
du complexe alterné augmenté de $M$ en $f_i$. D'après le
Lemme \ref{lemma-extended-alternating-base-change},
cette localisation est le complexe de {\v C}ech alterné augmenté
de $M_{f_i}$ sur $R_{f_i}$ relativement aux images de
$f_1, \ldots, f_r$ dans $R_{f_i}$. On se ramène donc à montrer que
$H^q$ est nul si $f_i$ est inversible, ce qui résulte du
Lemme \ref{lemma-extended-alternating-homotopy-zero}.
L'assertion (3) résulte de l'égalité $(H^q)_{f_i} = 0$
pour tout $i$, que nous venons de démontrer. Pour voir (4), remarquons que,
dans ce cas, $f$ agit inversiblement sur $H^q$ et que $H^q$ est supporté
par $V(f)$ d'après (3). Cela force $H^q$ à être nul (un petit détail est omis).
\end{proof}

\begin{lemma}
\label{lemma-extended-alternating-Cech-is-colimit-koszul}
Soit $R$ un anneau. Soient $f_1, \ldots, f_r \in R$. Le complexe de
{\v C}ech alterné augmenté
$$
R \to \bigoplus\nolimits_{i_0} R_{f_{i_0}} \to
\bigoplus\nolimits_{i_0 < i_1} R_{f_{i_0}f_{i_1}} \to
\ldots \to R_{f_1\ldots f_r}
$$
est une colimite des complexes de Koszul $K(R, f_1^n, \ldots, f_r^n)$ ;
voir la démonstration pour un énoncé précis.
\end{lemma}

\begin{proof}
Nous invitons vivement le lecteur à le démontrer lui-même.
Notons $K(R, f_1^n, \ldots, f_r^n)$ le complexe de Koszul de la
Définition \ref{definition-koszul-complex}, considéré comme un complexe de
cochaînes placé en degrés $0, \ldots, r$. Nous avons donc
$$
K(R, f_1^n, \ldots, f_r^n) :
0 \to \wedge^r(R^{\oplus r}) \to
\wedge^{r - 1}(R^{\oplus r}) \to \ldots \to
R^{\oplus r} \to R \to 0
$$
où le terme $\wedge^r(R^{\oplus r})$ est placé en degré $0$.
Soit $e^n_1, \ldots, e^n_r$ la base canonique de $R^{\oplus r}$.
Alors les éléments $e^n_{j_1} \wedge \ldots \wedge e^n_{j_{r - p}}$, pour
$1 \leq j_1 < \ldots < j_{r - p} \leq r$, forment une base du terme de
degré $p$ du complexe de Koszul. En outre, remarquons que
$$
d(e^n_{j_1} \wedge \ldots \wedge e^n_{j_{r - p}}) =
\sum (-1)^{a + 1} f_{j_a}^n e^n_{j_1} \wedge \ldots \wedge \hat e^n_{j_a} \wedge
\ldots \wedge e^n_{j_{r - p}}
$$
d'après notre construction du complexe de Koszul à la
Section \ref{section-koszul}. Les applications de transition de notre système
$$
K(R, f_1^n, \ldots, f_r^n) \to K(R, f_1^{n + 1}, \ldots, f_r^{n + 1})
$$
sont données par la règle
$$
e^n_{j_1} \wedge \ldots \wedge e^n_{j_{r - p}}
\longmapsto
f_{i_0} \ldots f_{i_{p - 1}}
e^{n + 1}_{j_1} \wedge \ldots \wedge e^{n + 1}_{j_{r - p}}
$$
où les indices $1 \leq i_0 < \ldots < i_{p - 1} \leq r$ sont tels que
$\{1, \ldots r\} =
\{i_0, \ldots, i_{p - 1}\} \amalg \{j_1, \ldots, j_{r - p}\}$.
Nous omettons le bref calcul qui montre la compatibilité aux différentielles.
Remarquons que les applications de transition valent toujours $1$ en degré $0$
et sont égales à $f_1 \ldots f_r$ en degré $r$.

\medskip\noindent
Notons $K^p(R, f_1^n, \ldots, f_r^n)$ le terme de degré $p$ du complexe de
Koszul. Remarquons que, pour tout $f \in R$, nous avons
$$
R_f = \colim (R \xrightarrow{f} R \xrightarrow{f} R \to \ldots )
$$
Nous voyons donc qu'en degré $p$ nous obtenons
$$
\colim K^p(R, f_1^n, \ldots f_r^n) =
\bigoplus\nolimits_{1 \leq i_0 < \ldots < i_{p - 1} \leq r}
R_{f_{i_0} \ldots f_{i_{p - 1}}}
$$
Ici, l'élément $e^n_{j_1} \wedge \ldots \wedge e^n_{j_{r - p}}$
du complexe de Koszul ci-dessus s'envoie dans la colimite
sur l'élément $(f_{i_0} \ldots f_{i_{p - 1}})^{-n}$ du
facteur direct $R_{f_{i_0} \ldots f_{i_{p - 1}}}$,
où les indices sont choisis de telle sorte que
$\{1, \ldots r\} = \{i_0, \ldots, i_{p - 1}\}
\amalg \{j_1, \ldots, j_{r - p}\}$.
Ainsi, la différentielle de ce complexe est donnée par
$$
d(1\text{ dans }R_{f_{i_0} \ldots f_{i_{p - 1}}}) =
\sum\nolimits_{i \not \in \{i_0, \ldots, i_{p - 1}\}}
(-1)^{i - t}\text{ dans }
R_{f_{i_0} \ldots f_{i_t} f_i f_{i_{t + 1}} \ldots f_{i_{p - 1}}}
$$
Ainsi, si nous considérons l'application de complexes donnée en degré $p$
par l'application
$$
\bigoplus\nolimits_{1 \leq i_0 < \ldots < i_{p - 1} \leq r}
R_{f_{i_0} \ldots f_{i_{p - 1}}}
\longrightarrow
\bigoplus\nolimits_{1 \leq i_0 < \ldots < i_{p - 1} \leq r}
R_{f_{i_0} \ldots f_{i_{p - 1}}}
$$
déterminée par la règle
$$
1\text{ dans }R_{f_{i_0} \ldots f_{i_{p - 1}}}
\longmapsto
(-1)^{i_0 + \ldots + i_{p - 1} + p}\text{ dans }R_{f_{i_0} \ldots f_{i_{p - 1}}}
$$
nous obtenons alors un isomorphisme de complexes de
$\colim K(R, f_1^n, \ldots, f_r^n)$ vers le
complexe de {\v C}ech alterné augmenté défini dans cette section.
Nous omettons la vérification des signes.
\end{proof}








\section{Suites régulières au sens de Koszul}
\label{section-koszul-regular}

\noindent
Nous invitons le lecteur à consulter
Algèbre, Sections \ref{algebra-section-regular-sequences},
\ref{algebra-section-quasi-regular} et
\ref{algebra-section-depth}
avant d'aborder celle-ci.

\begin{definition}
\label{definition-koszul-regular-sequence}
Soit $R$ un anneau. Soit $r \geq 0$ et soit $f_1, \ldots, f_r \in R$
une suite d'éléments. Soit $M$ un $R$-module.
La suite $f_1, \ldots, f_r$ est dite
\begin{enumerate}
\item {\it $M$-régulière au sens de Koszul} si
$H_i(K_\bullet(f_1, \ldots, f_r) \otimes_R M) = 0$
pour tout $i \not = 0$,
\item {\it $M$-$H_1$-régulière} si
$H_1(K_\bullet(f_1, \ldots, f_r) \otimes_R M) = 0$,
\item {\it régulière au sens de Koszul} si
$H_i(K_\bullet(f_1, \ldots, f_r)) = 0$ pour
tout $i \not = 0$, et
\item {\it $H_1$-régulière} si $H_1(K_\bullet(f_1, \ldots, f_r)) = 0$.
\end{enumerate}
\end{definition}

\noindent
Nous verrons aux Lemmes \ref{lemma-regular-koszul-regular},
\ref{lemma-koszul-regular-H1-regular} et
\ref{lemma-H1-regular-quasi-regular} que, pour des éléments
$f_1, \ldots, f_r$ d'un anneau $R$, nous avons les implications suivantes :
\begin{align*}
f_1, \ldots, f_r\text{ est une suite régulière}
& \Rightarrow f_1, \ldots, f_r\text{ est une suite régulière au sens de Koszul} \\
& \Rightarrow f_1, \ldots, f_r\text{ est une suite }H_1\text{-régulière} \\
& \Rightarrow f_1, \ldots, f_r\text{ est une suite quasi-régulière.}
\end{align*}
En général, aucune de ces implications ne peut être renversée, mais si $R$
est un anneau local noethérien et si $f_1, \ldots, f_r \in \mathfrak m_R$,
alors les quatre conditions sont toutes équivalentes
(Lemme \ref{lemma-noetherian-finite-all-equivalent}).
Si $f = f_1 \in R$ est une suite de longueur $1$ et si $f$ n'est pas
inversible dans $R$, il est clair que les assertions suivantes sont équivalentes :
\begin{enumerate}
\item $f$ est une suite régulière de longueur un,
\item $f$ est une suite régulière au sens de Koszul de longueur un, et
\item $f$ est une suite $H_1$-régulière de longueur un.
\end{enumerate}
Il est également clair que ces conditions impliquent que $f$ est une suite
quasi-régulière de longueur un. Mais il existe des suites quasi-régulières de
longueur $1$ qui ne sont pas régulières. En effet, soit
$$
R = k[x, y_0, y_1, \ldots]/(xy_0, xy_1 - y_0, xy_2 - y_1, \ldots)
$$
et soit $f$ l'image de $x$ dans $R$. Alors $f$ est un diviseur de zéro, mais
$\bigoplus_{n \geq 0} (f^n)/(f^{n + 1}) \cong k[x]$ est un anneau de polynômes.

\begin{lemma}
\label{lemma-regular-koszul-regular}
Soient $R$ un anneau, $M$ un $R$-module et $f_1, \ldots, f_r \in R$
tels que, pour $i = 1, \ldots, r$, la multiplication par $f_i$
soit injective sur $M/(f_1, \ldots, f_{i - 1})M$. Alors
$f_1, \ldots, f_r$ est $M$-régulière au sens de Koszul. En particulier,
une suite $M$-régulière est $M$-régulière au sens de Koszul, et toute suite
régulière est régulière au sens de Koszul.
\end{lemma}

\begin{proof}
Soient $R$, $M$, $f_1, \ldots, f_r$ comme dans la première phrase du lemme.
Si $r = 1$, il est immédiat que $f_1$ est $M$-régulière au sens de Koszul.
Supposons $r > 1$. Puisque $f_1$ est un non-diviseur de zéro dans $M$,
nous obtenons une suite exacte courte de complexes :
$$
0 \to K_\bullet(f_2, \ldots, f_r) \otimes M
\xrightarrow{f_1}
K_\bullet(f_2, \ldots, f_r) \otimes M \to
K_\bullet(\overline{f}_2, \ldots, \overline{f}_r) \otimes M/f_1M \to 0
$$
Ici, $\overline{f}_i$ est l'image de $f_i$ dans $R/(f_1)$. D'après le
Lemme \ref{lemma-cone-koszul}, le complexe $K_\bullet(f_1, \ldots, f_r)$
est isomorphe au cône de la multiplication par $f_1$
sur $K_\bullet(f_2, \ldots, f_r)$. Ainsi
$K_\bullet(R, f_1, \ldots, f_r) \otimes M$ est isomorphe
au cône de la première application. Par conséquent,
$K_\bullet(\overline{f}_2, \ldots, \overline{f}_r) \otimes M/f_1M$
est quasi-isomorphe à $K_\bullet(f_1, \ldots, f_r) \otimes M$.
Comme $R/(f_1)$, $M/f_1M$, $\overline{f}_2, \ldots, \overline{f}_r$
satisfont aux conditions du lemme, nous concluons par récurrence que
ce complexe est acyclique en degrés strictement positifs. Cela achève la
démonstration de la première assertion. La seconde résulte immédiatement
de la première.
\end{proof}

\begin{lemma}
\label{lemma-koszul-regular-H1-regular}
Une suite $M$-régulière au sens de Koszul est $M$-$H_1$-régulière.
Une suite régulière au sens de Koszul est $H_1$-régulière.
\end{lemma}

\begin{proof}
Cela résulte immédiatement de la définition.
\end{proof}

\begin{lemma}
\label{lemma-mult-koszul-regular}
Soit $f_1, \ldots, f_{r - 1} \in R$ une suite et soient $f, g \in R$.
Soit $M$ un $R$-module.
\begin{enumerate}
\item Si $f_1, \ldots, f_{r - 1}, f$ et $f_1, \ldots, f_{r - 1}, g$
sont $M$-$H_1$-régulières, alors $f_1, \ldots, f_{r - 1}, fg$ est
également $M$-$H_1$-régulière.
\item Si $f_1, \ldots, f_{r - 1}, f$ et $f_1, \ldots, f_{r - 1}, g$
sont $M$-régulières au sens de Koszul, alors $f_1, \ldots, f_{r - 1}, fg$
est également $M$-régulière au sens de Koszul.
\end{enumerate}
\end{lemma}

\begin{proof}
D'après le
Lemme \ref{lemma-koszul-mult}, nous avons des suites exactes
$$
H_i(K_\bullet(f_1, \ldots, f_{r - 1}, f) \otimes M) \to
H_i(K_\bullet(f_1, \ldots, f_{r - 1}, fg) \otimes M) \to
H_i(K_\bullet(f_1, \ldots, f_{r - 1}, g) \otimes M)
$$
pour tout $i$.
\end{proof}

\begin{lemma}
\label{lemma-koszul-regular-flat-base-change}
Soit $\varphi : R \to S$ un homomorphisme plat d'anneaux.
Soient $f_1, \ldots, f_r \in R$.
Soit $M$ un $R$-module et posons $N = M \otimes_R S$.
\begin{enumerate}
\item Si $f_1, \ldots, f_r$ dans $R$ est une suite $M$-$H_1$-régulière,
alors $\varphi(f_1), \ldots, \varphi(f_r)$ est une suite
$N$-$H_1$-régulière dans $S$.
\item Si $f_1, \ldots, f_r$ est une suite $M$-régulière au sens de Koszul
dans $R$, alors $\varphi(f_1), \ldots, \varphi(f_r)$ est une suite
$N$-régulière au sens de Koszul dans $S$.
\end{enumerate}
\end{lemma}

\begin{proof}
Cela est vrai car
$K_\bullet(f_1, \ldots, f_r) \otimes_R S =
K_\bullet(\varphi(f_1), \ldots, \varphi(f_r))$
et donc
$(K_\bullet(f_1, \ldots, f_r) \otimes_R M) \otimes_R S =
K_\bullet(\varphi(f_1), \ldots, \varphi(f_r)) \otimes_S N$.
\end{proof}

\begin{lemma}
\label{lemma-H1-regular-quasi-regular}
Une suite $M$-$H_1$-régulière est $M$-quasi-régulière.
\end{lemma}

\begin{proof}
Soit $R$ un anneau et soit $M$ un $R$-module.
Soit $f_1, \ldots, f_r$ une suite $M$-$H_1$-régulière.
Notons $J = (f_1, \ldots, f_r)$. L'hypothèse signifie que nous avons
une suite exacte
$$
\wedge^2(R^r) \otimes M \to R^{\oplus r} \otimes M \to JM \to 0
$$
où la première flèche est donnée par
$e_i \wedge e_j \otimes m \mapsto (f_ie_j - f_je_i) \otimes m$.
En tensorisant la suite par $R/J$, nous voyons que
$$
JM/J^2M = (R/J)^{\oplus r} \otimes_R M = (M/JM)^{\oplus r}
$$
est un module libre de type fini. Pour achever la démonstration, nous devons
montrer, pour tout $n \geq 2$, l'assertion suivante : si
$$
\xi = \sum\nolimits_{|I| = n, I = (i_1, \ldots, i_r)}
m_I f_1^{i_1} \ldots f_r^{i_r} \in J^{n + 1}M
$$
alors $m_I \in JM$ pour tout $I$. Dans le paragraphe suivant, nous montrons
que $m_I \in JM$ pour $I = (0, \ldots, 0, n)$ et, dans le dernier paragraphe,
nous déduisons le cas général de ce cas particulier.

\medskip\noindent
Soit $I = (0, \ldots, 0, n)$. Soit $\xi$ comme ci-dessus. Nous pouvons écrire
$\xi = m_1 f_1 + \ldots + m_{r - 1}f_{r - 1} + m_I f_r^n$.
Puisque nous avons supposé $\xi \in J^{n + 1}M$, nous pouvons aussi écrire
$\xi = \sum_{1 \leq i \leq j \leq r - 1} m_{ij}f_if_j +
\sum_{1 \leq i \leq r - 1}m'_i f_if_r^n + m'' f_r^{n + 1}$.
On voit alors que
$$
\begin{matrix}
(m_1 - m_{11}f_1 - m'_1f_r^n)f_1 + \\
(m_2 - m_{12}f_1 - m_{22}f_2 - m'_2f_r^n)f_2 + \\
\ldots + \\
(m_{r - 1} - m_{1 r - 1}f_1 - \ldots - m_{r - 1 r - 1}f_{r - 1}
- m'_{r - 1}f_r^n)f_{r - 1} + \\
(m_I - m'' f_r)f_r^n = 0
\end{matrix}
$$
Puisque $f_1, \ldots, f_{r - 1}, f_r^n$ est $M$-$H_1$-régulière d'après le
Lemme \ref{lemma-mult-koszul-regular}, nous voyons que $m_I - m'' f_r$
appartient au sous-module $f_1M + \ldots + f_{r - 1}M + f_r^nM$.
Ainsi $m_I \in f_1M + \ldots + f_rM$.

\medskip\noindent
Soit $S = R[x_1, x_2, \ldots, x_r, 1/x_r]$. L'homomorphisme d'anneaux
$R \to S$ est fidèlement plat ; ainsi $f_1, \ldots, f_r$ est une suite
$M$-$H_1$-régulière dans $S$, voir le
Lemme \ref{lemma-koszul-regular-flat-base-change}.
D'après le
Lemme \ref{lemma-change-basis}, nous voyons que
$$
g_1 = f_1 - \frac{x_1}{x_r} f_r,
\ \ldots,
\ g_{r - 1} = f_{r - 1} - \frac{x_{r - 1}}{x_r} f_r,
\ g_r = \frac{1}{x_r}f_r
$$
est une suite $M$-$H_1$-régulière dans $S$. Enfin, remarquons que notre
élément $\xi$ peut se réécrire
$$
\xi = \sum\nolimits_{|I| = n, I = (i_1, \ldots, i_r)}
m_I (g_1 + x_1 g_r)^{i_1} \ldots (g_{r - 1} + x_{r - 1} g_r)^{i_{r - 1}}
(x_rg_r)^{i_r}
$$
et que le coefficient de $g_r^n$ dans cette expression est
$$
\sum m_I x_1^{i_1} \ldots x_r^{i_r}
$$
D'après le cas traité dans le paragraphe précédent, cette somme appartient à
$J(M \otimes_R S)$. Comme les monômes $x_1^{i_1} \ldots x_r^{i_r}$
font partie d'une $R$-base de $S$ sur $R$, nous concluons que $m_I \in J$
pour tout $I$, comme voulu.
\end{proof}

\noindent
Pour les modules non nuls de type fini sur des anneaux locaux noethériens,
toutes les notions de suite régulière introduites jusqu'ici sont équivalentes.

\begin{lemma}
\label{lemma-noetherian-finite-all-equivalent}
Soit $(R, \mathfrak m)$ un anneau local noethérien. Soit $M$ un
$R$-module non nul de type fini. Soient
$f_1, \ldots, f_r \in \mathfrak m$. Les assertions suivantes sont équivalentes :
\begin{enumerate}
\item $f_1, \ldots, f_r$ est une suite $M$-régulière,
\item $f_1, \ldots, f_r$ est une suite $M$-régulière au sens de Koszul,
\item $f_1, \ldots, f_r$ est une suite $M$-$H_1$-régulière,
\item $f_1, \ldots, f_r$ est une suite $M$-quasi-régulière.
\end{enumerate}
En particulier, la suite $f_1, \ldots, f_r$ est une suite régulière
dans $R$ si et seulement si elle est régulière au sens de Koszul, si et
seulement si elle est $H_1$-régulière, si et seulement si elle est
quasi-régulière.
\end{lemma}

\begin{proof}
L'implication (1) $\Rightarrow$ (2) est le
Lemme \ref{lemma-regular-koszul-regular}.
L'implication (2) $\Rightarrow$ (3) est le
Lemme \ref{lemma-koszul-regular-H1-regular}.
L'implication (3) $\Rightarrow$ (4) est le
Lemme \ref{lemma-H1-regular-quasi-regular}.
L'implication (4) $\Rightarrow$ (1) est
Algèbre, Lemme \ref{algebra-lemma-quasi-regular-regular}.
\end{proof}

\begin{lemma}
\label{lemma-H1-regular-in-quotient}
Soit $A$ un anneau. Soit $I \subset A$ un idéal.
Soit $g_1, \ldots, g_m$ une suite dans $A$ dont l'image dans
$A/I$ est $H_1$-régulière. Alors $I \cap (g_1, \ldots, g_m) =
I(g_1, \ldots, g_m)$.
\end{lemma}

\begin{proof}
Considérons la suite exacte de complexes
$$
0 \to I \otimes_A K_\bullet(A, g_1, \ldots, g_m)
\to K_\bullet(A, g_1, \ldots, g_m) \to
K_\bullet(A/I, g_1, \ldots, g_m) \to 0
$$
Comme le complexe de droite vérifie $H_1 = 0$ par hypothèse, nous
voyons que
$$
\Coker(I^{\oplus m} \to I)
\longrightarrow
\Coker(A^{\oplus m} \to A)
$$
est injectif. Cela équivaut à l'assertion du lemme.
\end{proof}

\begin{lemma}
\label{lemma-conormal-sequence-H1-regular}
Soit $A$ un anneau. Soient $I \subset J \subset A$ des idéaux.
Supposons que $J/I \subset A/I$ soit engendré par une suite
$H_1$-régulière. Alors $I \cap J^2 = IJ$.
\end{lemma}

\begin{proof}
Pour le démontrer, choisissons $g_1, \ldots, g_m \in J$
dont les images dans $A/I$ forment une suite $H_1$-régulière qui engendre
$J/I$. En particulier, $J = I + (g_1, \ldots, g_m)$.
Supposons que $x \in I \cap J^2$. Puisque $x \in J^2$, nous pouvons écrire
$$
x =
\sum a_{ij} g_ig_j +
\sum a_j g_j +
a
$$
avec $a_{ij} \in A$, $a_j \in I$ et $a \in I^2$.
Alors $\sum a_{ij}g_ig_j \in I \cap (g_1, \ldots, g_m)$ ;
d'après le
Lemme \ref{lemma-H1-regular-in-quotient}, nous voyons donc que
$\sum a_{ij}g_ig_j \in I(g_1, \ldots, g_m)$.
Ainsi $x \in IJ$, comme voulu.
\end{proof}

\begin{lemma}
\label{lemma-join-quasi-regular-H1-regular}
Soit $A$ un anneau. Soit $I$ un idéal engendré par une suite
quasi-régulière $f_1, \ldots, f_n$ dans $A$. Soient $g_1, \ldots, g_m \in A$
des éléments dont les images $\overline{g}_1, \ldots, \overline{g}_m$
forment une suite $H_1$-régulière dans $A/I$. Alors
$f_1, \ldots, f_n, g_1, \ldots, g_m$ est une suite quasi-régulière dans $A$.
\end{lemma}

\begin{proof}
Nous affirmons que $g_1, \ldots, g_m$ forme une suite $H_1$-régulière dans
$A/I^d$ pour tout $d$. Supposons par récurrence que cela vaille dans
$A/I^{d - 1}$. Nous avons une suite exacte courte de complexes
$$
0 \to K_\bullet(A, g_\bullet) \otimes_A I^{d - 1}/I^d
\to K_\bullet(A/I^d, g_\bullet) \to
K_\bullet(A/I^{d - 1}, g_\bullet) \to 0
$$
Comme $f_1, \ldots, f_n$ est quasi-régulière, nous voyons que le premier
complexe est une somme directe de copies de
$K_\bullet(A/I, g_1, \ldots, g_m)$, donc qu'il est acyclique en degré $1$.
Par hypothèse de récurrence, le dernier complexe est acyclique en degré $1$.
Il en est donc de même du complexe du milieu.
En particulier, la suite $g_1, \ldots, g_m$ est quasi-régulière
dans $A/I^d$ pour tout $d \geq 1$, voir le
Lemme \ref{lemma-H1-regular-quasi-regular}.
Nous sommes maintenant prêts à démontrer que
$f_1, \ldots, f_n, g_1, \ldots, g_m$ est une suite quasi-régulière dans $A$.
En effet, posons $J = (f_1, \ldots, f_n, g_1, \ldots, g_m)$ et supposons
que (avec la notation multinomiale)
$$
\sum\nolimits_{|N| + |M| = d} a_{N, M} f^N g^M \in J^{d + 1}
$$
pour certains $a_{N, M} \in A$. Nous devons montrer que $a_{N, M} \in J$
pour tous $N, M$. Soit $e \in \{0, 1, \ldots, d\}$. Alors
$$
\sum\nolimits_{|N| = d - e, \ |M| = e} a_{N, M} f^N g^M \in
(g_1, \ldots, g_m)^{e + 1} + I^{d - e + 1}
$$
Comme $g_1, \ldots, g_m$ est une suite quasi-régulière dans
$A/I^{d - e + 1}$, nous en déduisons
$$
\sum\nolimits_{|N| = d - e} a_{N, M} f^N \in
(g_1, \ldots, g_m) + I^{d - e + 1}
$$
pour tout $M$ tel que $|M| = e$. Le
Lemme \ref{lemma-H1-regular-in-quotient},
appliqué à $I^{d - e}/I^{d - e + 1}$ dans l'anneau $A/I^{d - e + 1}$,
implique que
$\sum_{|N| = d - e} a_{N, M} f^N \in I^{d - e}(g_1, \ldots, g_m)$.
Comme $f_1, \ldots, f_n$ est quasi-régulière dans $A$, cela implique que
$a_{N, M} \in J$ pour tous $N, M$ tels que $|N| = d - e$ et $|M| = e$.
Cela démontre le lemme.
\end{proof}

\begin{lemma}
\label{lemma-join-H1-regular-sequences}
Soit $A$ un anneau. Soit $I$ un idéal engendré par une suite
$H_1$-régulière $f_1, \ldots, f_n$ dans $A$.
Soient $g_1, \ldots, g_m \in A$ des éléments dont les images
$\overline{g}_1, \ldots, \overline{g}_m$ forment une suite
$H_1$-régulière dans $A/I$. Alors $f_1, \ldots, f_n, g_1, \ldots, g_m$
est une suite $H_1$-régulière dans $A$.
\end{lemma}

\begin{proof}
Nous devons montrer que $H_1(A, f_1, \ldots, f_n, g_1, \ldots, g_m) = 0$.
Pour cela, considérons le diagramme commutatif
$$
\xymatrix{
\wedge^2(A^{\oplus n + m}) \ar[r] \ar[d] &
A^{\oplus n + m} \ar[r] \ar[d] &
A \ar[r] \ar[d] & 0 \\
\wedge^2(A/I^{\oplus m}) \ar[r] &
A/I^{\oplus m} \ar[r] &
A/I \ar[r] & 0
}
$$
Considérons un élément $(a_1, \ldots, a_{n + m}) \in A^{\oplus n + m}$
qui s'envoie sur zéro dans $A$. Comme
$\overline{g}_1, \ldots, \overline{g}_m$ forme une suite $H_1$-régulière
dans $A/I$, nous voyons que
$(\overline{a}_{n + 1}, \ldots, \overline{a}_{n + m})$ est l'image
d'un élément $\overline{\alpha}$ de $\wedge^2(A/I^{\oplus m})$.
Nous pouvons relever $\overline{\alpha}$ en un élément
$\alpha \in \wedge^2(A^{\oplus n + m})$ et soustraire son image
dans $A^{\oplus n + m}$ de notre élément $(a_1, \ldots, a_{n + m})$.
Ainsi, nous pouvons supposer que $a_{n + 1}, \ldots, a_{n + m} \in I$.
Puisque $I = (f_1, \ldots, f_n)$, nous pouvons modifier notre élément
$(a_1, \ldots, a_{n + m})$ par des combinaisons linéaires des éléments
$$
(0, \ldots, g_j, 0, \ldots, 0, f_i, 0, \ldots, 0)
$$
appartenant à l'image de la flèche horizontale supérieure gauche, et nous
ramener au cas où $a_{n + 1}, \ldots, a_{n + m}$ sont nuls. Dans ce cas,
$(a_1, \ldots, a_n, 0, \ldots, 0)$ définit un élément de
$H_1(A, f_1, \ldots, f_n)$, que nous avons supposé nul.
\end{proof}

\begin{lemma}
\label{lemma-truncate-H1-regular}
Soit $A$ un anneau. Soit $f_1, \ldots, f_n, g_1, \ldots, g_m \in A$
une suite $H_1$-régulière. Alors les images
$\overline{g}_1, \ldots, \overline{g}_m$ dans $A/(f_1, \ldots, f_n)$
forment une suite $H_1$-régulière.
\end{lemma}

\begin{proof}
Posons $I = (f_1, \ldots, f_n)$. Nous devons montrer que toute relation
$\sum_{j = 1, \ldots, m} \overline{a}_j \overline{g}_j$ dans $A/I$
est une combinaison linéaire de relations triviales. Puisque
$I = (f_1, \ldots, f_n)$, nous pouvons relever cette relation en une relation
$$
\sum\nolimits_{j = 1, \ldots, m} a_j g_j +
\sum\nolimits_{i = 1, \ldots, n} b_if_i = 0
$$
dans $A$. Par hypothèse, cette relation dans $A$ est une combinaison
linéaire de relations triviales. En prenant l'image dans $A/I$, nous
obtenons l'assertion voulue.
\end{proof}

\begin{lemma}
\label{lemma-join-koszul-regular-sequences}
Soit $A$ un anneau. Soit $I$ un idéal engendré par une suite
régulière au sens de Koszul $f_1, \ldots, f_n$ dans $A$.
Soient $g_1, \ldots, g_m \in A$ des éléments dont les images
$\overline{g}_1, \ldots, \overline{g}_m$ forment une suite régulière au
sens de Koszul dans $A/I$. Alors $f_1, \ldots, f_n, g_1, \ldots, g_m$
est une suite régulière au sens de Koszul dans $A$.
\end{lemma}

\begin{proof}
Nos hypothèses disent que $K_\bullet(A, f_1, \ldots, f_n)$ est une résolution
libre de type fini de $A/I$ et que
$K_\bullet(A/I, \overline{g}_1, \ldots, \overline{g}_m)$ est une
résolution libre de type fini de $A/(f_i, g_j)$ sur $A/I$. Alors
\begin{align*}
K_\bullet(A, f_1, \ldots, f_n, g_1, \ldots, g_m)
& = \text{Tot}(K_\bullet(A, f_1, \ldots, f_n) \otimes_A
K_\bullet(A, g_1, \ldots, g_m)) \\
& \cong A/I \otimes_A K_\bullet(A, g_1, \ldots, g_m) \\
& = K_\bullet(A/I, \overline{g}_1, \ldots, \overline{g}_m) \\
& \cong A/(f_i, g_j)
\end{align*}
La première égalité résulte du
Lemme \ref{lemma-join-sequences-koszul-complex}.
Le premier quasi-isomorphisme $\cong$ résulte (du dual) de
Homologie, Lemme \ref{homology-lemma-double-complex-gives-resolution},
car la $q$-ième ligne du complexe double
$K_\bullet(A, f_1, \ldots, f_n) \otimes_A K_\bullet(A, g_1, \ldots, g_m)$
est une résolution de $A/I \otimes_A K_q(A, g_1, \ldots, g_m)$.
La seconde égalité est claire. Le dernier quasi-isomorphisme résulte de
l'hypothèse. D'où le résultat.
\end{proof}

\noindent
Pour conclure dans le lemme suivant, il est nécessaire de supposer que
$f_1, \ldots, f_n$ et $f_1, \ldots, f_n, g_1, \ldots, g_m$
soient toutes deux régulières au sens de Koszul. Un contre-exemple lorsqu'on
supprime l'hypothèse que $f_1, \ldots, f_n$ est régulière au sens de Koszul
est donné dans
Exemples, Lemme \ref{examples-lemma-strange-regular-sequence}.

\begin{lemma}
\label{lemma-truncate-koszul-regular}
Soit $A$ un anneau. Soient $f_1, \ldots, f_n, g_1, \ldots, g_m \in A$.
Si $f_1, \ldots, f_n$ et $f_1, \ldots, f_n, g_1, \ldots, g_m$
sont toutes deux des suites régulières au sens de Koszul dans $A$, alors
les images $\overline{g}_1, \ldots, \overline{g}_m$ dans
$A/(f_1, \ldots, f_n)$ forment une suite régulière au sens de Koszul.
\end{lemma}

\begin{proof}
Posons $I = (f_1, \ldots, f_n)$.
Nos hypothèses disent que $K_\bullet(A, f_1, \ldots, f_n)$ est une résolution
libre de type fini de $A/I$ et que
$K_\bullet(A, f_1, \ldots, f_n, g_1, \ldots, g_m)$ est une
résolution libre de type fini de $A/(f_i, g_j)$ sur $A$. Alors
\begin{align*}
A/(f_i, g_j) & \cong K_\bullet(A, f_1, \ldots, f_n, g_1, \ldots, g_m) \\
& = \text{Tot}(K_\bullet(A, f_1, \ldots, f_n) \otimes_A
K_\bullet(A, g_1, \ldots, g_m)) \\
& \cong A/I \otimes_A K_\bullet(A, g_1, \ldots, g_m) \\
& = K_\bullet(A/I, \overline{g}_1, \ldots, \overline{g}_m)
\end{align*}
Le premier quasi-isomorphisme $\cong$ résulte de l'hypothèse. La première
égalité résulte du
Lemme \ref{lemma-join-sequences-koszul-complex}.
Le second quasi-isomorphisme résulte (du dual) de
Homologie, Lemme \ref{homology-lemma-double-complex-gives-resolution},
car la $q$-ième ligne du complexe double
$K_\bullet(A, f_1, \ldots, f_n) \otimes_A K_\bullet(A, g_1, \ldots, g_m)$
est une résolution de $A/I \otimes_A K_q(A, g_1, \ldots, g_m)$.
La seconde égalité est claire. D'où le résultat.
\end{proof}

\begin{lemma}
\label{lemma-independence-of-generators}
Soit $R$ un anneau. Soit $I$ un idéal engendré par
$f_1, \ldots, f_r \in R$.
\begin{enumerate}
\item Si $I$ peut être engendré par une suite quasi-régulière de longueur $r$,
alors $f_1, \ldots, f_r$ est une suite quasi-régulière.
\item Si $I$ peut être engendré par une suite $H_1$-régulière de longueur $r$,
alors $f_1, \ldots, f_r$ est une suite $H_1$-régulière.
\item Si $I$ peut être engendré par une suite régulière au sens de Koszul de
longueur $r$, alors $f_1, \ldots, f_r$ est une suite régulière au sens de
Koszul.
\end{enumerate}
\end{lemma}

\begin{proof}
Si $I$ peut être engendré par une suite quasi-régulière de longueur $r$,
alors $I/I^2$ est libre de rang $r$ sur $R/I$. Puisque
$f_1, \ldots, f_r$ constituent par hypothèse un système de générateurs,
nous voyons que les images
$\overline{f}_i$ forment une base de $I/I^2$ sur $R/I$. Il s'ensuit que
$f_1, \ldots, f_r$ est une suite quasi-régulière, car, outre la liberté de
$I/I^2$, cela signifie précisément que les applications
$\text{Sym}^n_{R/I}(I/I^2) \to I^n/I^{n + 1}$ sont des isomorphismes.

\medskip\noindent
Nous continuons de supposer que $I$ peut être engendré par une suite
quasi-régulière, disons
$g_1, \ldots, g_r$. Écrivons $g_j = \sum a_{ij}f_i$. Puisque
$f_1, \ldots, f_r$ est quasi-régulière d'après le paragraphe précédent,
nous voyons que $\det(a_{ij})$ est inversible modulo $I$. La matrice
$a_{ij}$ donne une application $R^{\oplus r} \to R^{\oplus r}$ qui induit
une application de complexes de Koszul
$\alpha : K_\bullet(R, f_1, \ldots, f_r) \to K_\bullet(R, g_1, \ldots, g_r)$,
voir le
Lemme \ref{lemma-functorial}.
Cette application devient un isomorphisme après inversion de
$\det(a_{ij})$. Comme les modules d'homologie de
$K_\bullet(R, f_1, \ldots, f_r)$ et de
$K_\bullet(R, g_1, \ldots, g_r)$ sont annulés par $I$, voir le
Lemme \ref{lemma-homotopy-koszul},
nous voyons que $\alpha$ est un quasi-isomorphisme.

\medskip\noindent
Supposons maintenant que $g_1, \ldots, g_r$ soit une suite $H_1$-régulière
engendrant $I$. Alors $g_1, \ldots, g_r$ est une suite quasi-régulière
d'après le Lemme \ref{lemma-H1-regular-quasi-regular}. Le paragraphe précédent
montre que $f_1, \ldots, f_r$ est une suite $H_1$-régulière.
Le raisonnement est analogue pour les suites régulières au sens de Koszul.
\end{proof}

\begin{lemma}
\label{lemma-make-nonzero-divisor}
\begin{reference}
C'est un cas particulier de \cite[Corollary]{McCoy}.
\end{reference}
Soit $R$ un anneau. Soient $a_1, \ldots, a_n \in R$ des éléments tels que
$R \to R^{\oplus n}$, $x \mapsto (xa_1, \ldots, xa_n)$, soit injective.
Alors l'élément $\sum a_i t_i$ de l'anneau de polynômes
$R[t_1, \ldots, t_n]$ est un non-diviseur de zéro.
\end{lemma}

\begin{proof}
Si l'un des $a_i$ est inversible, il s'agit simplement du fait que tout
élément de la forme $t_1 + a_2 t_2 + \ldots + a_n t_n$ est un
non-diviseur de zéro dans l'anneau de polynômes sur $R$.

\medskip\noindent
Cas I : $R$ est noethérien. Soient $\mathfrak q_j$, $j = 1, \ldots, m$,
les idéaux premiers associés de $R$. Nous devons montrer que chacune des
applications
$$
\sum a_i t_i :
\text{Sym}^d(R^{\oplus n})
\longrightarrow
\text{Sym}^{d + 1}(R^{\oplus n})
$$
est injective. Comme $\text{Sym}^d(R^{\oplus n})$ est un $R$-module libre,
ses idéaux premiers associés sont les $\mathfrak q_j$, $j = 1, \ldots, m$.
Pour tout $j$, il existe un $i = i(j)$ tel que
$a_i \not \in \mathfrak q_j$, car il existe par hypothèse un $x \in R$
tel que $\mathfrak q_jx = 0$, mais $a_i x \not = 0$ pour un certain $i$.
Ainsi $a_i$ est inversible dans $R_{\mathfrak q_j}$ et l'application est
injective après localisation en $\mathfrak q_j$. L'application est donc
injective, voir
Algèbre, Lemme \ref{algebra-lemma-zero-at-ass-zero}.

\medskip\noindent
Cas II : $R$ est quelconque. Nous pouvons écrire $R$ comme la réunion
d'anneaux noethériens $R_\lambda$ avec
$a_1, \ldots, a_n \in R_\lambda$. Le résultat vaut pour chaque
$R_\lambda$, donc il vaut pour $R$.
\end{proof}

\begin{lemma}
\label{lemma-Koszul-regular-flat-locally-regular}
Soit $R$ un anneau. Soit $f_1, \ldots, f_n$ une suite régulière au sens de
Koszul dans $R$ telle que $(f_1, \ldots, f_n) \not = R$.
Considérons l'homomorphisme d'anneaux lisse et fidèlement plat
$$
R \longrightarrow
S = R[\{t_{ij}\}_{i \leq j}, t_{11}^{-1}, t_{22}^{-1}, \ldots, t_{nn}^{-1}]
$$
Pour $1 \leq i \leq n$, posons
$$
g_i = \sum\nolimits_{i \leq j} t_{ij} f_j \in S.
$$
Alors $g_1, \ldots, g_n$ est une suite régulière dans $S$ et
$(f_1, \ldots, f_n)S = (g_1, \ldots, g_n)$.
\end{lemma}

\begin{proof}
L'égalité des idéaux est évidente, car la matrice
$$
\left(
\begin{matrix}
t_{11} & t_{12} & t_{13} & \ldots \\
0 & t_{22} & t_{23} & \ldots \\
0 & 0 & t_{33} & \ldots \\
\ldots & \ldots & \ldots & \ldots
\end{matrix}
\right)
$$
est inversible dans $S$.
Comme $f_1, \ldots, f_n$ est une suite régulière au sens de Koszul, nous
voyons que le noyau de
$R \to R^{\oplus n}$, $x \mapsto (xf_1, \ldots, xf_n)$ est nul (il calcule
la $n$-ième homologie de Koszul de $R$ relativement à\ $f_1, \ldots, f_n$).
Par conséquent, d'après le
Lemme \ref{lemma-make-nonzero-divisor}, nous voyons que
$g_1 = f_1 t_{11} + \ldots + f_n t_{1n}$ est un non-diviseur de zéro
dans $S' = R[t_{11}, t_{12}, \ldots, t_{1n}, t_{11}^{-1}]$. Nous voyons que
$g_1, f_2, \ldots, f_n$ est une suite régulière au sens de Koszul dans
$S'$ d'après les Lemmes \ref{lemma-koszul-regular-flat-base-change} et
\ref{lemma-independence-of-generators}.
Nous en concluons que
$\overline{f}_2, \ldots, \overline{f}_n$ est une suite régulière au sens
de Koszul dans $S'/(g_1)$ d'après le
Lemme \ref{lemma-truncate-koszul-regular}.
Par récurrence sur $n$, nous voyons donc que les images
$\overline{g}_2, \ldots, \overline{g}_n$ de $g_2, \ldots, g_n$ dans
$S'/(g_1)[\{t_{ij}\}_{2 \leq i \leq j}, t_{22}^{-1}, \ldots, t_{nn}^{-1}]$
forment une suite régulière. Cela signifie à son tour que
$g_1, \ldots, g_n$ forme une suite régulière dans $S$.
\end{proof}







\section{Compléments sur les suites régulières au sens de Koszul}
\label{section-more-koszul-regular}

\noindent
Nous poursuivons l'étude commencée à la Section \ref{section-koszul-regular}.

\begin{lemma}
\label{lemma-vanishing-extended-alternating-koszul}
Soit $R$ un anneau. Soit $f_1, \ldots, f_r \in R$ une suite
régulière au sens de Koszul. Alors le complexe de {\v C}ech alterné augmenté
$R \to \bigoplus\nolimits_{i_0} R_{f_{i_0}} \to
\bigoplus\nolimits_{i_0 < i_1} R_{f_{i_0}f_{i_1}} \to
\ldots \to R_{f_1\ldots f_r}$
de la Section \ref{section-alternating-cech}
n'a de cohomologie qu'en degré $r$.
\end{lemma}

\begin{proof}
D'après le Lemme \ref{lemma-mult-koszul-regular} et par récurrence, la suite
$f_1, \ldots, f_{r - 1}, f_r^n$ est régulière au sens de Koszul pour tout
$n \geq 1$. D'après le Lemme \ref{lemma-change-basis}, toute permutation
d'une suite régulière au sens de Koszul est encore une telle suite. Nous
voyons donc que nous pouvons remplacer un (ou tous les) $f_i$ par sa
$n$-ième puissance tout en conservant une suite régulière au sens de Koszul.
Ainsi $K_\bullet(R, f_1^n, \ldots, f_r^n)$ n'a d'homologie non nulle qu'en
degré homologique $0$. Cela donne l'assertion voulue d'après le
Lemme \ref{lemma-extended-alternating-Cech-is-colimit-koszul}.
\end{proof}

\begin{lemma}
\label{lemma-blowup-regular-sequence}
Soit $a, a_2, \ldots, a_r$ une suite $H_1$-régulière dans un anneau $R$
(par exemple une suite régulière au sens de Koszul ou une suite régulière,
voir les Lemmes \ref{lemma-regular-koszul-regular} et
\ref{lemma-koszul-regular-H1-regular}).
Si $I = (a, a_2, \ldots, a_r)$, l'algèbre d'éclatement
$R' = R[\frac{I}{a}]$ est isomorphe à
$R'' = R[y_2, \ldots, y_r]/(a y_i - a_i)$.
\end{lemma}

\begin{proof}
D'après Algèbre, Lemme
\ref{algebra-lemma-affine-blowup-quotient-description},
il suffit de montrer que $R''$ est sans $a$-torsion.

\medskip\noindent
Nous affirmons que $a, ay_2 - a_2, \ldots, ay_n - a_r$ est une suite
$H_1$-régulière dans $R[y_2, \ldots, y_r]$. En effet, l'application
$$
(a, ay_2 - a_2, \ldots, ay_n - a_r) :
R[y_2, \ldots, y_r]^{\oplus r}
\longrightarrow
R[y_2, \ldots, y_r]
$$
utilisée pour définir le complexe de Koszul sur
$a, ay_2 - a_2, \ldots, ay_n - a_r$ est isomorphe à l'application
$$
(a, a_2, \ldots, a_r) :
R[y_2, \ldots, y_r]^{\oplus r} \longrightarrow
R[y_2, \ldots, y_r]
$$
utilisée pour définir le complexe de Koszul sur $a, a_2, \ldots, a_r$,
au moyen de l'isomorphisme
$$
R[y_2, \ldots, y_r]^{\oplus r}
\longrightarrow
R[y_2, \ldots, y_r]^{\oplus r}
$$
qui envoie $(b_1, \ldots, b_r)$ sur
$(b_1 - b_2y_2 \ldots - b_ry_r, -b_2, \ldots, - b_r)$.
D'après le Lemme \ref{lemma-functorial}, ces complexes de Koszul sont
isomorphes. En appliquant le
Lemme \ref{lemma-koszul-regular-flat-base-change} à l'homomorphisme plat
d'anneaux $R \to R[y_2, \ldots, y_r]$, nous obtenons notre assertion.
D'après le Lemme \ref{lemma-cone-koszul}, nous voyons que le complexe de
Koszul $K$ sur $a, ay_2 - a_2, \ldots, ay_n - a_r$ est le cône de
$a : L \to L$, où $L$ est le complexe de Koszul sur
$ay_2 - a_2, \ldots, ay_n - a_r$. Comme $H_1(K) = 0$ d'après l'assertion,
nous concluons que $a : H_0(L) \to H_0(L)$ est injective, autrement dit que
$R'' = R[y_2, \ldots, y_r]/(a y_i - a_i)$
ne possède aucun élément non nul de $a$-torsion, comme voulu.
\end{proof}

\begin{lemma}
\label{lemma-base-change-H1-regular}
Soit $A \to B$ un homomorphisme d'anneaux.
Soit $f_1, \ldots, f_r$ une suite dans $B$ telle que
$B/(f_1, \ldots, f_r)$ soit plat sur $A$. Soit $A \to A'$ un homomorphisme
d'anneaux. Alors l'application canonique
$$
H_1(K_\bullet(B, f_1, \ldots, f_r)) \otimes_A A'
\longrightarrow
H_1(K_\bullet(B', f'_1, \ldots, f'_r))
$$
est surjective. Ici $B' = B \otimes_A A'$ et $f_i' \in B'$ est l'image
de $f_i$.
\end{lemma}

\begin{proof}
La suite
$$
\wedge^2(B^{\oplus r}) \to B^{\oplus r} \to B \to B/J \to 0
$$
est un complexe de $A$-modules dans lequel $B/J$ est plat sur $A$ et dont
le groupe d'homologie $H_1 = H_1(K_\bullet(B, f_1, \ldots, f_r))$ au niveau
de $B^{\oplus r}$. Si nous le tensorisons par
$A'$, nous obtenons un complexe
$$
\wedge^2((B')^{\oplus r}) \to (B')^{\oplus r} \to B' \to B'/J' \to 0
$$
qui est exact en $B'$ et en $B'/J'$. Afin de calculer son
groupe d'homologie $H'_1 = H_1(K_\bullet(B', f'_1, \ldots, f'_r))$
en $(B')^{\oplus r}$, nous décomposons la première suite ci-dessus en
les suites exactes $0 \to J \to B \to B/J \to 0$,
$0 \to K \to B^{\oplus r} \to J \to 0$ et
$\wedge^2(B^{\oplus r}) \to K \to H_1 \to 0$.
En tensorisant sur $A$ par $A'$, nous obtenons les suites exactes
$$
\begin{matrix}
0 \to J \otimes_A A' \to B \otimes_A A' \to (B/J) \otimes_A A' \to 0 \\
K \otimes_A A' \to B^{\oplus r} \otimes_A A' \to J \otimes_A A' \to 0 \\
\wedge^2(B^{\oplus r}) \otimes_A A' \to K \otimes_A A' \to H_1 \otimes_A A'
\to 0
\end{matrix}
$$
où la première est exacte puisque $B/J$ est plat sur $A$, voir
Algèbre, Lemme \ref{algebra-lemma-flat-tor-zero}. Nous concluons que
$J' = J \otimes_A A' \subset B'$ et que
$K \otimes_A A' \to \Ker((B')^{\oplus r} \to B')$ est
surjective. Ainsi
\begin{align*}
H_1 \otimes_A A'
& =
\Coker\left(\wedge^2(B^{\oplus r}) \otimes_A A' \to K \otimes_A A'\right) \\
& \to
\Coker\left(
\wedge^2((B')^{\oplus r})  \to \Ker((B')^{\oplus r} \to B')
\right) = H'_1
\end{align*}
est également surjective.
\end{proof}

\begin{lemma}
\label{lemma-relative-regular-immersion-algebra}
Soient $A \to B$ et $A \to A'$ des homomorphismes d'anneaux.
Posons $B' = B \otimes_A A'$.
Soient $f_1, \ldots, f_r \in B$. Supposons que
$B/(f_1, \ldots, f_r)B$ soit plat sur $A$.
\begin{enumerate}
\item Si $f_1, \ldots, f_r$ est une suite quasi-régulière, alors
son image dans $B'$ est une suite quasi-régulière.
\item Si $f_1, \ldots, f_r$ est une suite $H_1$-régulière, alors
son image dans $B'$ est une suite $H_1$-régulière.
\end{enumerate}
\end{lemma}

\begin{proof}
Supposons que $f_1, \ldots, f_r$ soit quasi-régulière. Posons
$J = (f_1, \ldots, f_r)$. Par hypothèse, $J^n/J^{n + 1}$ est isomorphe à une
somme directe de copies de $B/J$, donc est plat sur $A$. Par récurrence et
d'après Algèbre, Lemme \ref{algebra-lemma-flat-ses},
nous concluons que $B/J^n$ est plat sur $A$. L'idéal $(J')^n$ est égal à
$J^n \otimes_A A'$, voir
Algèbre, Lemme \ref{algebra-lemma-flat-tor-zero}.
Ainsi $(J')^n/(J')^{n + 1} = J^n/J^{n + 1} \otimes_A A'$, ce qui implique
clairement que $f_1, \ldots, f_r$ est une suite quasi-régulière dans $B'$.

\medskip\noindent
Supposons que $f_1, \ldots, f_r$ soit $H_1$-régulière. D'après le
Lemme \ref{lemma-base-change-H1-regular}, l'annulation du groupe
d'homologie de Koszul
$H_1(K_\bullet(B, f_1, \ldots, f_r))$
implique celle de $H_1(K_\bullet(B', f'_1, \ldots, f'_r))$,
d'où le résultat.
\end{proof}

\begin{lemma}
\label{lemma-cut-by-koszul-flat}
Soit $A' \to B'$ un homomorphisme d'anneaux. Soit $I \subset A'$ un idéal.
Posons $A = A'/I$ et $B = B'/IB'$. Soient
$f'_1, \ldots, f'_r \in B'$. Supposons que
\begin{enumerate}
\item $A' \to B'$ soit plat et de présentation finie,
\item $I$ soit localement nilpotent,
\item les images $f_1, \ldots, f_r \in B$ forment une suite quasi-régulière,
\item $B/(f_1, \ldots, f_r)$ soit plat sur $A$.
\end{enumerate}
Alors $B'/(f'_1, \ldots, f'_r)$ est plat sur $A'$.
\end{lemma}

\begin{proof}
Posons $C' = B'/(f'_1, \ldots, f'_r)$. Nous devons montrer que
$A' \to C'$ est plat. Soit $\mathfrak r' \subset C'$ un idéal premier
au-dessus de $\mathfrak p' \subset A'$. Notons $\mathfrak q' \subset B'$
l'image réciproque de $\mathfrak r'$.
D'après Algèbre, Lemme \ref{algebra-lemma-flat-localization},
il suffit de montrer que $A'_{\mathfrak p'} \to C'_{\mathfrak q'}$ est plat.
Algèbre, Lemme \ref{algebra-lemma-grothendieck-regular-sequence-general},
nous dit qu'il suffit de montrer que les images de
$f'_1, \ldots, f'_r$ forment une suite régulière dans
$$
B'_{\mathfrak q'}/\mathfrak p'B'_{\mathfrak q'} =
B_\mathfrak q/\mathfrak p B_\mathfrak q =
(B \otimes_A \kappa(\mathfrak p))_\mathfrak q
$$
avec les notations évidentes. Nous savons que $f_1, \ldots, f_r$
est une suite quasi-régulière dans $B$ et que $B/(f_1, \ldots, f_r)$
est plat sur $A$. D'après le
Lemme \ref{lemma-relative-regular-immersion-algebra}, les images
$\overline{f}_1, \ldots, \overline{f}_r$ de $f'_1, \ldots, f'_r$ dans
$B \otimes_A \kappa(\mathfrak p)$ forment une suite quasi-régulière.
Comme $(B \otimes_A \kappa(\mathfrak p))_\mathfrak q$ est un anneau local
noethérien, nous concluons par le
Lemme \ref{lemma-noetherian-finite-all-equivalent}.
\end{proof}

\begin{lemma}
\label{lemma-cut-by-koszul}
Soit $A' \to B'$ un homomorphisme d'anneaux. Soit $I \subset A'$ un idéal.
Posons $A = A'/I$ et $B = B'/IB'$. Soient
$f'_1, \ldots, f'_r \in B'$. Supposons que
\begin{enumerate}
\item $A' \to B'$ soit plat et de présentation finie (par exemple lisse),
\item $I$ soit localement nilpotent,
\item les images $f_1, \ldots, f_r \in B$ forment une suite quasi-régulière,
\item $B/(f_1, \ldots, f_r)$ soit lisse sur $A$.
\end{enumerate}
Alors $B'/(f'_1, \ldots, f'_r)$ est lisse sur $A'$.
\end{lemma}

\begin{proof}
Posons $C' = B'/(f'_1, \ldots, f'_r)$ et $C = B/(f_1, \ldots, f_r)$.
Alors $A' \to C'$ est de présentation finie.
D'après le Lemme \ref{lemma-cut-by-koszul-flat}, nous voyons que
$A' \to C'$ est plat. Les anneaux des fibres de $A' \to C'$ sont égaux
à ceux de $A \to C$, et sont donc lisses d'après l'hypothèse (4).
Il s'ensuit que $A' \to C'$ est lisse d'après
Algèbre, Lemme \ref{algebra-lemma-flat-fibre-smooth}.
\end{proof}








\section{Idéaux réguliers}
\label{section-ideals}

\noindent
Nous étudierons la notion de faisceau d'idéaux réguliers dans une grande
généralité dans
Diviseurs, Section \ref{divisors-section-regular-ideal-sheaves}.
Nous définissons ici la notion correspondante dans le cas affine,
c'est-à-dire pour un idéal d'un anneau.

\begin{definition}
\label{definition-regular-ideal}
Soit $R$ un anneau et soit $I \subset R$ un idéal.
\begin{enumerate}
\item On dit que $I$ est un {\it idéal régulier} si, pour tout
$\mathfrak p \in V(I)$, il existe un $g \in R$, $g \not \in \mathfrak p$,
et une suite régulière $f_1, \ldots, f_r \in R_g$ tels que $I_g$
soit engendré par $f_1, \ldots, f_r$.
\item On dit que $I$ est un {\it idéal régulier au sens de Koszul} si,
pour tout $\mathfrak p \in V(I)$, il existe un $g \in R$,
$g \not \in \mathfrak p$, et une suite régulière au sens de Koszul
$f_1, \ldots, f_r \in R_g$ tels que $I_g$ soit engendré par
$f_1, \ldots, f_r$.
\item On dit que $I$ est un {\it idéal $H_1$-régulier} si, pour tout
$\mathfrak p \in V(I)$, il existe un $g \in R$, $g \not \in \mathfrak p$,
et une suite $H_1$-régulière $f_1, \ldots, f_r \in R_g$ tels que $I_g$
soit engendré par $f_1, \ldots, f_r$.
\item On dit que $I$ est un {\it idéal quasi-régulier} si, pour tout
$\mathfrak p \in V(I)$, il existe un $g \in R$, $g \not \in \mathfrak p$,
et une suite quasi-régulière $f_1, \ldots, f_r \in R_g$ tels que $I_g$
soit engendré par $f_1, \ldots, f_r$.
\end{enumerate}
\end{definition}

\noindent
Il est clair que, pour $I \subset R$, nous avons les implications
\begin{align*}
I\text{ est un idéal régulier}
& \Rightarrow
I\text{ est un idéal régulier au sens de Koszul} \\
& \Rightarrow
I\text{ est un idéal }H_1\text{-régulier} \\
& \Rightarrow
I\text{ est un idéal quasi-régulier}
\end{align*}
voir les Lemmes \ref{lemma-regular-koszul-regular},
\ref{lemma-koszul-regular-H1-regular} et
\ref{lemma-H1-regular-quasi-regular}. Un tel idéal est toujours de type fini.

\begin{lemma}
\label{lemma-quasi-regular-ideal-finite}
Un idéal quasi-régulier est de type fini.
\end{lemma}

\begin{proof}
Soit $I \subset R$ un idéal quasi-régulier. Puisque $V(I)$ est quasi-compact,
il existe $g_1, \ldots, g_m \in R$ tels que
$V(I) \subset D(g_1) \cup \ldots \cup D(g_m)$
et tels que $I_{g_j}$ soit engendré par une suite quasi-régulière
$g_{j1}, \ldots, g_{jr_j} \in R_{g_j}$. Écrivons
$g_{ji} = g'_{ji}/g_j^{e_{ij}}$ pour certains $g'_{ij} \in I$.
Écrivons $1 + x = \sum g_j h_j$ pour un certain $x \in I$, ce qui est
possible puisque
$V(I) \subset D(g_1) \cup \ldots \cup D(g_m)$.
Remarquons que $\Spec(R) = D(g_1) \cup \ldots \cup D(g_m) \bigcup D(x)$.
Alors $I$ est engendré par les éléments $g'_{ij}$ et $x$, car ceux-ci
engendrent sur chacun des ouverts du recouvrement, voir
Algèbre, Lemme \ref{algebra-lemma-cover}.
\end{proof}

\begin{lemma}
\label{lemma-quasi-regular-ideal-finite-projective}
Soit $I \subset R$ un idéal quasi-régulier d'un anneau.
Alors $I/I^2$ est un $R/I$-module projectif de type fini.
\end{lemma}

\begin{proof}
Cela résulte d'Algèbre, Lemme \ref{algebra-lemma-finite-projective},
et des définitions.
\end{proof}

\noindent
Nous démontrons la descente plate pour les idéaux réguliers au sens de
Koszul, $H_1$-réguliers et quasi-réguliers.

\begin{lemma}
\label{lemma-flat-descent-regular-ideal}
Soit $A \to B$ un homomorphisme fidèlement plat d'anneaux.
Soit $I \subset A$ un idéal. Si $IB$ est un idéal régulier au sens de Koszul
(resp.\ $H_1$-régulier, resp.\ quasi-régulier) dans $B$, alors
$I$ est un idéal régulier au sens de Koszul
(resp.\ $H_1$-régulier, resp.\ quasi-régulier) dans $A$.
\end{lemma}

\begin{proof}
Nous fixons l'idéal premier $\mathfrak p \supset I$ pendant toute la
démonstration. Supposons que $IB$ soit quasi-régulier. D'après le
Lemme \ref{lemma-quasi-regular-ideal-finite},
$IB$ est un module de type fini ; par conséquent, $I$ est un $A$-module de
type fini d'après
Algèbre, Lemme \ref{algebra-lemma-descend-properties-modules}.
Comme $A \to B$ est plat, nous voyons que
$$
I/I^2 \otimes_{A/I} B/IB = I/I^2 \otimes_A B = IB/(IB)^2.
$$
Comme $IB$ est quasi-régulier, le $B/IB$-module $IB/(IB)^2$ est localement
libre de type fini. Ainsi $I/I^2$ est projectif de type fini, voir
Algèbre, Proposition
\ref{algebra-proposition-ffdescent-finite-projectivity}.
En particulier, après avoir remplacé $A$ par $A_f$ pour un certain
$f \in A$, $f \not \in \mathfrak p$, nous pouvons supposer que $I/I^2$ est
libre de rang $r$. Choisissons $f_1, \ldots, f_r \in I$ qui donnent une base
de $I/I^2$. Par le lemme de Nakayama (voir
Algèbre, Lemme \ref{algebra-lemma-NAK}), nous voyons qu'après un nouveau
remplacement $A \leadsto A_f$ comme ci-dessus, $I$ est engendré par
$f_1, \ldots, f_r$.

\medskip\noindent
Démonstration du cas « quasi-régulier ». Nous avons vu ci-dessus que
$I/I^2$ est libre sur les $r$ générateurs $f_1, \ldots, f_r$.
Pour achever la démonstration dans ce cas, nous devons montrer que les
applications $\text{Sym}^d(I/I^2) \to I^d/I^{d + 1}$ sont des isomorphismes
pour tout $d \geq 2$. C'est clair, car leurs changements de base fidèlement
plats $\text{Sym}^d(IB/(IB)^2) \to (IB)^d/(IB)^{d + 1}$ sont des
isomorphismes localement sur $B$ par hypothèse.
Les détails sont omis.

\medskip\noindent
Démonstration des cas « $H_1$-régulier » et « régulier au sens de Koszul ».
Considérons la suite d'éléments $f_1, \ldots, f_r$ engendrant $I$ que nous
avons construite ci-dessus. D'après le
Lemme \ref{lemma-independence-of-generators}, nous voyons que l'image de
$f_1, \ldots, f_r$ est une suite $H_1$-régulière ou régulière au sens de
Koszul dans $B_g$, pour tout $g \in B$ tel que $IB$ soit engendré par
une suite $H_1$-régulière ou régulière au sens de Koszul. Ainsi
$K_\bullet(A, f_1, \ldots, f_r) \otimes_A B_g$ a une homologie nulle en
$H_1$ ou en $H_i$, $i > 0$. Puisque l'homologie de
$K_\bullet(B, f_1, \ldots, f_r) = K_\bullet(A, f_1, \ldots, f_r) \otimes_A B$
est annulée par $IB$ (voir le
Lemme \ref{lemma-homotopy-koszul}) et puisque
$V(IB) \subset \bigcup_{g\text{ comme ci-dessus}} D(g)$, nous concluons que
$K_\bullet(A, f_1, \ldots, f_r) \otimes_A B$ a une homologie nulle en
degré $1$, ou en tous les degrés strictement positifs. Comme $A \to B$ est
fidèlement plat, nous concluons qu'il en est de même pour
$K_\bullet(A, f_1, \ldots, f_r)$.
\end{proof}

\begin{lemma}
\label{lemma-conormal-sequence-H1-regular-ideal}
Soit $A$ un anneau. Soient $I \subset J \subset A$ des idéaux.
Supposons que $J/I \subset A/I$ soit un idéal $H_1$-régulier.
Alors $I \cap J^2 = IJ$.
\end{lemma}

\begin{proof}
Cela résulte immédiatement du
Lemme \ref{lemma-conormal-sequence-H1-regular} par localisation.
\end{proof}







\section{Morphismes d'intersection complète locale}
\label{section-lci}

\noindent
Nous pouvons utiliser les résultats précédents pour définir un morphisme
d'intersection complète locale entre anneaux au moyen de présentations par
des algèbres de polynômes (de type fini).

\begin{lemma}
\label{lemma-koszul-independence-presentation}
Soit $A \to B$ un homomorphisme d'anneaux de type fini. Si, pour une
présentation $\alpha : A[x_1, \ldots, x_n] \to B$, le noyau $I$ est un
idéal régulier au sens de Koszul, alors, pour toute présentation
$\beta : A[y_1, \ldots, y_m] \to B$, le noyau $J$ est un idéal régulier
au sens de Koszul.
\end{lemma}

\begin{proof}
Choisissons $f_j \in A[x_1, \ldots, x_n]$ tel que
$\alpha(f_j) = \beta(y_j)$ et $g_i \in A[y_1, \ldots, y_m]$ tel que
$\beta(g_i) = \alpha(x_i)$. Nous obtenons alors un diagramme commutatif
$$
\xymatrix{
A[x_1, \ldots, x_n, y_1, \ldots, y_m]
\ar[d]^{x_i \mapsto g_i} \ar[rr]_-{y_j \mapsto f_j} & &
A[x_1, \ldots, x_n] \ar[d] \\
A[y_1, \ldots, y_m] \ar[rr] & & B
}
$$
Remarquons que le noyau $K$ de $A[x_i, y_j] \to B$ est égal à
$K = (I, y_j - f_j) = (J, x_i - g_i)$. En particulier, comme
$I$ est de type fini d'après le
Lemme \ref{lemma-quasi-regular-ideal-finite}, nous voyons que
$J = K/(x_i - g_i)$ est également de type fini.

\medskip\noindent
Choisissons un idéal premier $\mathfrak q \subset B$. Puisque
$I/I^2 \oplus B^{\oplus m} = J/J^2 \oplus B^{\oplus n}$
(Algèbre, Lemme \ref{algebra-lemma-conormal-module}),
nous voyons que
$$
\dim J/J^2 \otimes_B \kappa(\mathfrak q) + n =
\dim I/I^2 \otimes_B \kappa(\mathfrak q) + m.
$$
Choisissons $p_1, \ldots, p_t \in I$ dont les images forment une base de
$I/I^2 \otimes \kappa(\mathfrak q) = I \otimes_{A[x_i]} \kappa(\mathfrak q)$.
Choisissons $q_1, \ldots, q_s \in J$ dont les images forment une base de
$J/J^2 \otimes \kappa(\mathfrak q) = J \otimes_{A[y_j]} \kappa(\mathfrak q)$.
Ainsi $s + n = t + m$. Par le lemme de Nakayama, il existe
$h \in A[x_i]$ et $h' \in A[y_j]$, dont les images sont toutes deux non
nuls dans $\kappa(\mathfrak q)$, tels que
$I_h = (p_1, \ldots, p_t)$ dans $A[x_i, 1/h]$ et
$J_{h'} = (q_1, \ldots, q_s)$ dans $A[y_j, 1/h']$.
Comme $I$ est régulier au sens de Koszul, nous pouvons aussi supposer que
$I_h$ est engendré par une suite régulière au sens de Koszul. Cette suite
doit nécessairement être de longueur
$t = \dim I/I^2 \otimes_B \kappa(\mathfrak q)$ ; nous voyons donc que
$p_1, \ldots, p_t$ est une suite régulière au sens de Koszul d'après le
Lemme \ref{lemma-independence-of-generators}.
Comme $y_1 - f_1, \ldots, y_m - f_m$ est également une suite régulière,
nous concluons que
$$
y_1 - f_1, \ldots, y_m - f_m, p_1, \ldots, p_t
$$
est une suite régulière au sens de Koszul dans $A[x_i, y_j, 1/h]$
(voir le Lemme \ref{lemma-join-koszul-regular-sequences}).
Cette suite engendre l'idéal $K_h$. Ainsi l'idéal $K_{hh'}$ est engendré
par une suite régulière au sens de Koszul de longueur $m + t = n + s$.
Mais il est aussi engendré par la suite
$$
x_1 - g_1, \ldots, x_n - g_n, q_1, \ldots, q_s
$$
de même longueur, qui est donc une suite régulière au sens de Koszul d'après le
Lemme \ref{lemma-independence-of-generators}.
Enfin, d'après le
Lemme \ref{lemma-truncate-koszul-regular}, nous concluons que les images de
$q_1, \ldots, q_s$ dans
$$
A[x_i, y_j, 1/hh']/(x_1 - g_1, \ldots, x_n - g_n)
\cong A[y_j, 1/h'']
$$
forment une suite régulière au sens de Koszul engendrant $J_{h''}$. Comme
$h''$ est l'image de $hh'$, son image dans $\kappa(\mathfrak q)$ n'est pas nulle,
et le résultat s'ensuit.
\end{proof}

\noindent
Ce lemme permet de donner la définition suivante.

\begin{definition}
\label{definition-local-complete-intersection}
Un homomorphisme d'anneaux $A \to B$ est appelé une
{\it intersection complète locale} s'il est de type fini et si, pour une
(ou, de façon équivalente, toute) présentation
$B = A[x_1, \ldots, x_n]/I$, l'idéal $I$ est régulier au sens de Koszul.
\end{definition}

\noindent
Cette notion est locale.

\begin{lemma}
\label{lemma-lci-local}
Soit $R \to S$ un homomorphisme d'anneaux. Soient $g_1, \ldots, g_m \in S$
des générateurs de l'idéal unité. Si chaque $R \to S_{g_j}$ est une
intersection complète locale, alors $R \to S$ l'est aussi.
\end{lemma}

\begin{proof}
Soit $S = R[x_1, \ldots, x_n]/I$ une présentation. Choisissons
$h_j \in R[x_1, \ldots, x_n]$ dont l'image est $g_j$ dans $S$.
Alors $R[x_1, \ldots, x_n, x_{n + 1}]/(I, x_{n + 1}h_j - 1)$
est une présentation de $S_{g_j}$. Ainsi
$I_j = (I, x_{n + 1}h_j - 1)$ est un idéal régulier au sens de Koszul dans
$R[x_1, \ldots, x_n, x_{n + 1}]$. Choisissons un idéal premier
$I \subset \mathfrak q \subset R[x_1, \ldots, x_n]$.
Alors $h_j \not \in \mathfrak q$ pour un certain $j$ et
$\mathfrak q_j = (\mathfrak q, x_{n + 1}h_j - 1)$ est un idéal premier
de $V(I_j)$ au-dessus de $\mathfrak q$.
Choisissons $f_1, \ldots, f_r \in I$ dont les images forment une base de
$I/I^2 \otimes \kappa(\mathfrak q)$. Alors
$x_{n + 1}h_j - 1, f_1, \ldots, f_r$ est une suite d'éléments de $I_j$
dont les images forment une base de $I_j \otimes \kappa(\mathfrak q_j)$, voir
Algèbre, Lemme \ref{algebra-lemma-principal-localization-NL}.
Par le lemme de Nakayama, il existe un
$h \in R[x_1, \ldots, x_n, x_{n + 1}]$ tel que $(I_j)_h$ soit engendré par
$x_{n + 1}h_j - 1, f_1, \ldots, f_r$.
Nous pouvons aussi supposer que $(I_j)_h$ est engendré par une suite
régulière au sens de Koszul de longueur $e$. En examinant la dimension de
$I_j \otimes \kappa(\mathfrak q_j)$, nous voyons que $e = r + 1$.
D'après le
Lemme \ref{lemma-independence-of-generators}, nous voyons donc que
$x_{n + 1}h_j - 1, f_1, \ldots, f_r$ est une suite régulière au sens de
Koszul engendrant $(I_j)_h$ pour un certain
$h \in R[x_1, \ldots, x_n, x_{n + 1}]$, $h \not \in \mathfrak q_j$. D'après le
Lemme \ref{lemma-truncate-koszul-regular}, nous voyons que $I_{h'}$ est
engendré par une suite régulière au sens de Koszul pour un certain
$h' \in R[x_1, \ldots, x_n]$, $h' \not \in \mathfrak q$, comme voulu.
\end{proof}

\begin{lemma}
\label{lemma-relative-global-complete-intersection-koszul}
Soit $R$ un anneau. Si $R[x_1, \ldots, x_n]/(f_1, \ldots, f_c)$
est une intersection complète globale relative, alors
$f_1, \ldots, f_c$ est une suite régulière au sens de Koszul.
\end{lemma}

\begin{proof}
Rappelons que les groupes d'homologie $H_i(K_\bullet(f_\bullet))$ sont
annulés par l'idéal $(f_1, \ldots, f_c)$. Il suffit donc de montrer que
$H_i(K_\bullet(f_\bullet))_\mathfrak q$ est nul pour tout idéal premier
$\mathfrak q \subset R[x_1, \ldots, x_n]$ contenant
$(f_1, \ldots, f_c)$. Cela résulte d'Algèbre, Lemme
\ref{algebra-lemma-relative-global-complete-intersection-conormal}
et du fait qu'une suite régulière est régulière au sens de Koszul
(Lemme \ref{lemma-regular-koszul-regular}).
\end{proof}

\begin{lemma}
\label{lemma-syntomic-lci}
Soit $R \to S$ un homomorphisme d'anneaux. Les assertions suivantes sont
équivalentes :
\begin{enumerate}
\item $R \to S$ est syntomique
(Algèbre, Définition \ref{algebra-definition-lci}), et
\item $R \to S$ est plat et est une intersection complète locale.
\end{enumerate}
\end{lemma}

\begin{proof}
Supposons (1). Alors $R \to S$ est plat par définition.
D'après Algèbre, Lemme \ref{algebra-lemma-syntomic} et le
Lemme \ref{lemma-lci-local}, il suffit de montrer qu'une intersection
complète globale relative est un homomorphisme d'intersection complète
locale, ce qui est le
Lemme \ref{lemma-relative-global-complete-intersection-koszul}.

\medskip\noindent
Supposons (2). Une intersection complète locale est de présentation finie,
car un idéal régulier au sens de Koszul est de type fini.
Soit $R \to k$ un homomorphisme vers un corps. Il suffit de montrer que
$S' = S \otimes_R k$ est une intersection complète locale sur $k$, voir
Algèbre, Définition \ref{algebra-definition-lci-field}.
Choisissons un idéal premier $\mathfrak q' \subset S'$.
Écrivons $S = R[x_1, \ldots, x_n]/I$.
Alors $S' = k[x_1, \ldots, x_n]/I'$, où
$I' \subset k[x_1, \ldots, x_n]$ est l'image de $I$.
Soient $\mathfrak p' \subset k[x_1, \ldots, x_n]$,
$\mathfrak q \subset S$ et
$\mathfrak p \subset R[x_1, \ldots, x_n]$ les idéaux premiers correspondants.
D'après la Définition \ref{definition-regular-ideal}, il existe un
$g \in R[x_1, \ldots, x_n]$, $g \not \in \mathfrak p$, et
$f_1, \ldots, f_r \in R[x_1, \ldots, x_n]_g$ qui forment une suite régulière
au sens de Koszul engendrant $I_g$.
Comme $S$, et donc $S_g$, est plat sur $R$, nous voyons que les images
$f'_1, \ldots, f'_r$ dans $k[x_1, \ldots, x_n]_g$ forment une suite
$H_1$-régulière engendrant $I'_g$, voir le
Lemme \ref{lemma-relative-regular-immersion-algebra}.
Ainsi $f'_1, \ldots, f'_r$ se transforment en une suite régulière dans
$k[x_1, \ldots, x_n]_{\mathfrak p'}$ engendrant $I'_{\mathfrak p'}$,
d'après le Lemme \ref{lemma-noetherian-finite-all-equivalent}.
En appliquant Algèbre, Lemme \ref{algebra-lemma-lci}, nous concluons que
$S'_{gg'}$, pour un certain $g' \in S$, $g' \not \in \mathfrak q'$, est une
intersection complète globale sur $k$, comme voulu.
\end{proof}

\noindent
Pour une intersection complète locale $R \to S$, nous avons
$H_n(L_{S/R}) = 0$ pour $n \geq 2$. Comme nous n'avons pas (encore) défini
le complexe cotangent complet, nous ne pouvons pas énoncer et démontrer ceci,
mais nous pouvons en déduire l'une des conséquences.

\begin{lemma}
\label{lemma-transitive-lci-at-end}
Soient $A \to B \to C$ des homomorphismes d'anneaux. Supposons que
$B \to C$ soit un homomorphisme d'intersection complète locale. Choisissons
une présentation
$\alpha : A[x_s, s \in S] \to B$ de noyau $I$. Choisissons une présentation
$\beta : B[y_1, \ldots, y_m] \to C$ de noyau $J$. Soit
$\gamma : A[x_s, y_t] \to C$ la présentation induite de $C$ de noyau
$K$. Nous obtenons alors un diagramme commutatif canonique
$$
\xymatrix{
0 \ar[r] &
\Omega_{A[x_s]/A} \otimes C \ar[r] &
\Omega_{A[x_s, y_t]/A} \otimes C \ar[r] &
\Omega_{B[y_t]/B} \otimes C \ar[r] &
0 \\
0 \ar[r] &
I/I^2 \otimes C \ar[r] \ar[u] &
K/K^2 \ar[r] \ar[u] &
J/J^2 \ar[r] \ar[u] &
0
}
$$
à lignes exactes. En particulier, la suite exacte à six termes d'Algèbre,
Lemme \ref{algebra-lemma-exact-sequence-NL}, peut être complétée par un zéro
à gauche, c'est-à-dire que la suite
$$
0 \to H_1(\NL_{B/A} \otimes_B C) \to
H_1(L_{C/A}) \to
H_1(L_{C/B}) \to
\Omega_{B/A} \otimes_B C \to
\Omega_{C/A} \to
\Omega_{C/B} \to 0
$$
est exacte.
\end{lemma}

\begin{proof}
La seule chose à démontrer est l'injectivité de l'application
$I/I^2 \otimes C \to K/K^2$.
Par hypothèse, l'idéal $J$ est régulier au sens de Koszul.
Nous avons donc $IA[x_s, y_j] \cap K^2 = IK$ d'après le
Lemme \ref{lemma-conormal-sequence-H1-regular-ideal}.
Cela signifie que le noyau de $K/K^2 \to J/J^2$ est
isomorphe à $IA[x_s, y_j]/IK$. Comme
$I/I^2 \otimes_A C = IA[x_s, y_j]/IK$ par exactitude à droite
du produit tensoriel, nous obtenons l'injectivité voulue de
$I/I^2 \otimes_A C \to K/K^2$.
\end{proof}

\begin{lemma}
\label{lemma-transitive-colimit-lci-at-end}
Soient $A \to B \to C$ des homomorphismes d'anneaux.
Si $B \to C$ est une colimite filtrante d'homomorphismes d'intersection
complète locale, alors la conclusion du
Lemme \ref{lemma-transitive-lci-at-end}
reste valable.
\end{lemma}

\begin{proof}
Cela résulte du
Lemme \ref{lemma-transitive-lci-at-end}
et d'Algèbre, Lemme \ref{algebra-lemma-colimits-NL}.
\end{proof}

\begin{lemma}
\label{lemma-henselization-NL}
Soit $A \to B$ un homomorphisme local d'anneaux locaux.
Soit $A^h \to B^h$, resp.\ $A^{sh} \to B^{sh}$, l'homomorphisme induit
sur les hensélisations, resp.\ hensélisations strictes
(Algèbre, Lemme \ref{algebra-lemma-henselian-functorial},
resp.\ Lemme \ref{algebra-lemma-strictly-henselian-functorial}).
Alors $\NL_{B/A} \otimes_B B^h \to \NL_{B^h/A^h}$ et
$\NL_{B/A} \otimes_B B^{sh} \to \NL_{B^{sh}/A^{sh}}$
induisent des isomorphismes sur les groupes de cohomologie.
\end{lemma}

\begin{proof}
Comme $A^h$ est une colimite filtrante d'algèbres \'etales sur $A$,
nous voyons que $\NL_{A^h/A}$ est un complexe acyclique d'après
Algèbre, Lemme \ref{algebra-lemma-colimits-NL} et
Algèbre, Définition \ref{algebra-definition-etale}.
Il en est de même pour $B^h/B$.
En utilisant la suite de Jacobi--Zariski
(Algèbre, Lemme \ref{algebra-lemma-exact-sequence-NL})
pour $A \to A^h \to B^h$, nous trouvons que
$\NL_{B^h/A} \to \NL_{B^h/A^h}$ induit des isomorphismes sur les groupes
de cohomologie. En outre, un homomorphisme \'etale d'anneaux est une
intersection complète locale, puisqu'il est même une intersection complète
globale, voir Algèbre, Lemme \ref{algebra-lemma-etale-standard-smooth}.
D'après le Lemme \ref{lemma-transitive-colimit-lci-at-end}, nous obtenons
une suite exacte de Jacobi--Zariski à six termes associée à
$A \to B \to B^h$, ce qui montre que
$\NL_{B/A} \otimes_B B^h \to \NL_{B^h/A}$ induit des isomorphismes
sur les groupes de cohomologie. Cela achève la démonstration pour
l'homomorphisme sur les hensélisations. Le cas des hensélisations strictes
se démontre exactement de la même manière.
\end{proof}







\section{Égalité de Cartier et régularité géométrique}
\label{section-cartier-equality}

\noindent
Une référence pour cette section et la suivante est
\cite[Section 39]{MatCA}. Pour lire cette section confortablement, le lecteur
devra être familier avec le complexe cotangent naïf et ses propriétés, voir
Algèbre, Section \ref{algebra-section-netherlander}.

\begin{lemma}[Égalité de Cartier]
\label{lemma-cartier-equality}
Soit $K/k$ une extension de corps de type fini.
Alors $\Omega_{K/k}$ et $H_1(L_{K/k})$ sont de dimension finie et
$\text{trdeg}_k(K) = \dim_K \Omega_{K/k} - \dim_K H_1(L_{K/k})$.
\end{lemma}

\begin{proof}
Nous pouvons trouver une intersection complète globale
$A = k[x_1, \ldots, x_n]/(f_1, \ldots, f_c)$
sur $k$ telle que $K$ soit isomorphe au corps des fractions de $A$, voir
Algèbre, Lemme \ref{algebra-lemma-colimit-syntomic} et sa démonstration.
Dans ce cas, nous voyons que $\NL_{K/k}$ est homotopiquement équivalent au
complexe
$$
\bigoplus\nolimits_{j = 1, \ldots, c} K \longrightarrow
\bigoplus\nolimits_{i = 1, \ldots, n} K\text{d}x_i
$$
d'après Algèbre, Lemmes \ref{algebra-lemma-NL-homotopy} et
\ref{algebra-lemma-localize-NL}.
Le degré de transcendance de $K$ sur $k$ est la dimension de $A$
(d'après Algèbre, Lemme \ref{algebra-lemma-dimension-prime-polynomial-ring}),
qui vaut $n - c$, d'où le résultat.
\end{proof}

\begin{lemma}
\label{lemma-transitivity-gamma}
Soient $M/L/K$ des extensions de corps. Alors la suite de Jacobi--Zariski
$$
0 \to H_1(L_{L/K}) \otimes_L M \to
H_1(L_{M/K}) \to
H_1(L_{M/L}) \to
\Omega_{L/K} \otimes_L M \to
\Omega_{M/K} \to
\Omega_{M/L} \to 0
$$
est exacte.
\end{lemma}

\begin{proof}
Combinons le
Lemme \ref{lemma-transitive-colimit-lci-at-end}
avec Algèbre, Lemme \ref{algebra-lemma-colimit-syntomic}.
\end{proof}

\begin{lemma}
\label{lemma-gamma-commutative-diagram}
Soit donné un diagramme commutatif de corps
$$
\xymatrix{
K \ar[r] & K' \\
k \ar[u] \ar[r] & k' \ar[u]
}
$$
tel que $k'/k$ et $K'/K$ soient des extensions de corps de type fini.
Le noyau et le conoyau des applications
$$
\alpha : \Omega_{K/k} \otimes_K K' \to \Omega_{K'/k'}
\quad\text{et}\quad
\beta : H_1(L_{K/k}) \otimes_K K' \to H_1(L_{K'/k'})
$$
 sont de dimension finie et
$$
\dim \Ker(\alpha) - \dim \Coker(\alpha)
-\dim \Ker(\beta) + \dim \Coker(\beta)
=
\text{trdeg}_k(k') - \text{trdeg}_K(K')
$$
\end{lemma}

\begin{proof}
Les suites de Jacobi--Zariski pour $k \subset k' \subset K'$ et
$k \subset K \subset K'$ sont
$$
0 \to H_1(L_{k'/k}) \otimes K'  \to
H_1(L_{K'/k}) \to
H_1(L_{K'/k'}) \to
\Omega_{k'/k} \otimes K' \to
\Omega_{K'/k} \to
\Omega_{K'/k'} \to 0
$$
et
$$
0 \to H_1(L_{K/k}) \otimes K' \to
H_1(L_{K'/k}) \to
H_1(L_{K'/K}) \to
\Omega_{K/k} \otimes K' \to
\Omega_{K'/k} \to
\Omega_{K'/K} \to 0
$$
D'après le
Lemme \ref{lemma-cartier-equality}, les espaces vectoriels
$\Omega_{k'/k}$, $\Omega_{K'/K}$, $H_1(L_{K'/K})$ et
$H_1(L_{k'/k})$ sont de dimension finie et la somme alternée de leurs
dimensions vaut $\text{trdeg}_k(k') - \text{trdeg}_K(K')$.
Le lemme s'ensuit.
\end{proof}





\section{Régularité géométrique}
\label{section-geometrically-regular}

\noindent
Soit $k$ un corps. Soit $(A, \mathfrak m, K)$ une $k$-algèbre locale
noethérienne. La suite de Jacobi--Zariski
(Algèbre, Lemme \ref{algebra-lemma-exact-sequence-NL})
est une suite exacte canonique
$$
H_1(L_{K/k}) \to \mathfrak m/\mathfrak m^2 \to \Omega_{A/k} \otimes_A K \to
\Omega_{K/k} \to 0
$$
car $H_1(L_{K/A}) = \mathfrak m/\mathfrak m^2$ d'après
Algèbre, Lemme \ref{algebra-lemma-NL-surjection}.
Nous montrerons que l'exactitude à gauche de cette suite caractérise le fait
qu'un anneau local régulier $A$ soit ou non géométriquement régulier sur $k$.
Nous relierons cela à la notion de lissité formelle à la
Section \ref{section-regular-fs}.

\begin{proposition}
\label{proposition-characterization-geometrically-regular}
Soit $k$ un corps de caractéristique $p > 0$.
Soit $(A, \mathfrak m, K)$ une $k$-algèbre locale
noethérienne. Les assertions suivantes sont équivalentes :
\begin{enumerate}
\item $A$ est géométriquement régulier sur $k$,
\item pour tout $k \subset k' \subset k^{1/p}$
fini sur $k$, l'anneau $A \otimes_k k'$ est régulier,
\item $A$ est régulier et l'application canonique
$H_1(L_{K/k}) \to \mathfrak m/\mathfrak m^2$ est injective, et
\item $A$ est régulier et l'application
$\Omega_{k/\mathbf{F}_p} \otimes_k K \to \Omega_{A/\mathbf{F}_p} \otimes_A K$
est injective.
\end{enumerate}
\end{proposition}

\begin{proof}
Démonstration de (3) $\Rightarrow$ (1).
Supposons (3). Soit $k'/k$ une extension finie purement inséparable.
Posons $A' = A \otimes_k k'$. C'est un anneau local d'idéal maximal
$\mathfrak m'$. Posons $K' = A'/\mathfrak m'$. Nous obtenons un diagramme
commutatif
$$
\xymatrix{
0 \ar[r] &
H_1(L_{K/k}) \otimes K' \ar[r] \ar[d]_\beta &
\mathfrak m/\mathfrak m^2 \otimes K' \ar[r] \ar[d] &
\Omega_{A/k} \otimes_A K' \ar[r] \ar[d]_{\cong} &
\Omega_{K/k} \otimes K' \ar[r] \ar[d]_\alpha &
0
\\
 &
H_1(L_{K'/k'}) \ar[r] &
\mathfrak m'/(\mathfrak m')^2 \ar[r] &
\Omega_{A'/k'} \otimes_{A'} K' \ar[r] &
\Omega_{K'/k'} \ar[r] &
0
}
$$
à lignes exactes. La troisième flèche verticale est un isomorphisme par
changement de base pour les modules de différentielles
(Algèbre, Lemme \ref{algebra-lemma-differentials-base-change}).
Ainsi $\alpha$ est surjective. D'après le
Lemme \ref{lemma-gamma-commutative-diagram}, nous avons
$$
\dim \Ker(\alpha) - \dim \Ker(\beta) + \dim \Coker(\beta) = 0
$$
(et toutes ces dimensions sont finies). Une chasse dans le diagramme montre
que $\dim \mathfrak m'/(\mathfrak m')^2 \leq \dim \mathfrak m/\mathfrak m^2$.
Cependant, puisque $A \to A'$ est fini plat, nous avons
$\dim(A) = \dim(A')$, voir
Algèbre, Lemme \ref{algebra-lemma-dimension-base-fibre-total}.
Ainsi $A'$ est régulier par définition.

\medskip\noindent
Démonstration de l'équivalence de (3) et (4). Considérons les suites de
Jacobi--Zariski correspondant aux lignes du diagramme commutatif
$$
\xymatrix{
\mathbf{F}_p \ar[r] & A \ar[r] & K \\
\mathbf{F}_p \ar[r] \ar[u] & k \ar[r] \ar[u] & K \ar[u]
}
$$
pour obtenir un diagramme commutatif
$$
\xymatrix{
0 \ar[r] &
\mathfrak m/\mathfrak m^2 \ar[r] &
\Omega_{A/\mathbf{F}_p} \otimes_A K \ar[r] &
\Omega_{K/\mathbf{F}_p} \ar[r] & 0 & \\
0 \ar[r] &
H_1(L_{K/k}) \ar[r] \ar[u] &
\Omega_{k/\mathbf{F}_p} \otimes_k K \ar[r] \ar[u] &
\Omega_{K/\mathbf{F}_p} \ar[r] \ar[u] &
\Omega_{K/k} \ar[r] \ar[u] &
0
}
$$
à lignes exactes. Nous avons utilisé $H_1(L_{K/A}) = \mathfrak m/\mathfrak m^2$
et $H_1(L_{K/\mathbf{F}_p}) = 0$ puisque $K/\mathbf{F}_p$ est séparable, voir
Algèbre, Proposition
\ref{algebra-proposition-characterize-separable-field-extensions}.
Il est donc clair que les noyaux de
$H_1(L_{K/k}) \to \mathfrak m/\mathfrak m^2$
et
$\Omega_{k/\mathbf{F}_p} \otimes_k K \to \Omega_{A/\mathbf{F}_p} \otimes_A K$
ont la même dimension.

\medskip\noindent
Démonstration de (2) $\Rightarrow$ (4) d'après Faltings, voir
\cite{Faltings-einfacher}. Soient $a_1, \ldots, a_n \in k$ des éléments
tels que $\text{d}a_1, \ldots, \text{d}a_n$ soient linéairement indépendants
dans $\Omega_{k/\mathbf{F}_p}$. Considérons l'extension de corps
$k' = k(a_1^{1/p}, \ldots, a_n^{1/p})$. D'après
Algèbre, Lemme \ref{algebra-lemma-size-extension-pth-roots}, nous voyons
que $k' = k[x_1, \ldots, x_n]/(x_1^p - a_1, \ldots, x_n^p - a_n)$.
En particulier, le complexe cotangent naïf de $k'/k$ est homotope au complexe
$\bigoplus_{j = 1, \ldots, n} k' \rightarrow \bigoplus_{i = 1, \ldots, n} k'$
à différentielle nulle, car
$\text{d}(x_j^p - a_j) = 0$ dans $\Omega_{k[x_1, \ldots, x_n]/k}$.
Posons $A' = A \otimes_k k'$ et $K' = A'/\mathfrak m'$ comme ci-dessus.
D'après Algèbre, Lemme \ref{algebra-lemma-change-base-NL}, nous voyons que
$\NL_{A'/A}$ est homotopiquement équivalent au complexe
$\bigoplus_{j = 1, \ldots, n} A' \rightarrow \bigoplus_{i = 1, \ldots, n} A'$
à différentielle nulle, c'est-à-dire que $H_1(L_{A'/A})$ et
$\Omega_{A'/A}$ sont libres de rang $n$. La suite de Jacobi--Zariski pour
$\mathbf{F}_p \to A \to A'$ est
$$
H_1(L_{A'/A}) \to \Omega_{A/\mathbf{F}_p} \otimes_A A'
\to \Omega_{A'/\mathbf{F}_p} \to \Omega_{A'/A} \to 0
$$
En utilisant la présentation $A[x_1, \ldots, x_n] \to A'$ de noyau
$(x_j^p - a_j)$, nous voyons, en explicitant les applications de
Algèbre, Lemme \ref{algebra-lemma-exact-sequence-NL},
que le $j$-ième vecteur de base de $H_1(L_{A'/A})$ s'envoie sur
$\text{d}a_j \otimes 1$ dans $\Omega_{A/\mathbf{F}_p} \otimes A'$.
Comme $\Omega_{A'/A}$ est libre (donc plat), en tensorisant par $K'$ nous
obtenons une suite exacte
$$
K'^{\oplus n} \to \Omega_{A/\mathbf{F}_p} \otimes_A K'
\xrightarrow{\beta} \Omega_{A'/\mathbf{F}_p} \otimes_{A'} K' \to
K'^{\oplus n} \to 0
$$
Nous concluons que les éléments $\text{d}a_j \otimes 1$ engendrent
$\Ker(\beta)$ et nous devons montrer qu'ils sont linéairement indépendants,
c'est-à-dire montrer que $\dim(\Ker(\beta)) = n$.
Considérons le grand diagramme suivant
$$
\xymatrix{
0 \ar[r] &
\mathfrak m'/(\mathfrak m')^2 \ar[r] &
\Omega_{A'/\mathbf{F}_p} \otimes K' \ar[r] &
\Omega_{K'/\mathbf{F}_p} \ar[r] & 0 \\
0 \ar[r] &
\mathfrak m/\mathfrak m^2 \otimes K' \ar[r] \ar[u]^\alpha &
\Omega_{A/\mathbf{F}_p} \otimes K' \ar[r] \ar[u]^\beta &
\Omega_{K/\mathbf{F}_p} \otimes K' \ar[r] \ar[u]^\gamma & 0
}
$$
D'après le Lemme \ref{lemma-cartier-equality} et la suite de Jacobi--Zariski
pour $\mathbf{F}_p \to K \to K'$, nous voyons que le noyau et le conoyau de
$\gamma$ ont la même dimension finie.
Par hypothèse, $A'$ est régulier (et de même dimension que $A$, voir
ci-dessus), donc le noyau et le conoyau de $\alpha$ ont la même dimension.
Il s'ensuit que le noyau et le conoyau de $\beta$ ont la même dimension,
ce qu'il fallait démontrer.

\medskip\noindent
L'implication (1) $\Rightarrow$ (2) est triviale. Cela achève la
démonstration de la proposition.
\end{proof}

\begin{lemma}
\label{lemma-geometrically-regular-over-field}
Soit $k$ un corps de caractéristique $p > 0$. Soit $(A, \mathfrak m, K)$
une $k$-algèbre locale noethérienne. Supposons que $A$ soit géométriquement
régulière sur $k$. Soit $K/F/k$ une sous-extension de type fini.
Soit $\varphi : k[y_1, \ldots, y_m] \to A$ un homomorphisme de
$k$-algèbres tel que $y_i$ s'envoie sur un élément de $F$ dans $K$ et que
$\text{d}y_1, \ldots, \text{d}y_m$ s'envoient sur une base de $\Omega_{F/k}$.
Posons $\mathfrak p = \varphi^{-1}(\mathfrak m)$. Alors
$$
k[y_1, \ldots, y_m]_\mathfrak p \to A
$$
est plat et $A/\mathfrak pA$ est régulier.
\end{lemma}

\begin{proof}
Posons $A_0 = k[y_1, \ldots, y_m]_\mathfrak p$, d'idéal maximal
$\mathfrak m_0$ et de corps résiduel $K_0$. Remarquons que
$\Omega_{A_0/k}$ est libre de rang $m$ et que
$\Omega_{A_0/k} \otimes K_0 \to \Omega_{K_0/k}$ est un isomorphisme.
Il est clair que $A_0$ est géométriquement régulier sur $k$. Ainsi
$H_1(L_{K_0/k}) \to \mathfrak m_0/\mathfrak m_0^2$ est un isomorphisme, voir
la Proposition \ref{proposition-characterization-geometrically-regular}.
Considérons maintenant
$$
\xymatrix{
H_1(L_{K_0/k}) \otimes K \ar[d] \ar[r] &
\mathfrak m_0/\mathfrak m_0^2 \otimes K \ar[d] \\
H_1(L_{K/k}) \ar[r] & \mathfrak m/\mathfrak m^2
}
$$
Comme la flèche verticale gauche est injective d'après le
Lemme \ref{lemma-transitivity-gamma} et la flèche horizontale inférieure
d'après la Proposition \ref{proposition-characterization-geometrically-regular},
nous concluons qu'il en est de même de la flèche verticale droite.
Ainsi un système régulier de paramètres de $A_0$ s'envoie sur une partie
d'un système régulier de paramètres de $A$.
Le résultat s'ensuit d'après
Algèbre, Lemmes \ref{algebra-lemma-flat-over-regular} et
\ref{algebra-lemma-regular-ring-CM}.
\end{proof}












\section{Anneaux et modules topologiques}
\label{section-topological-ring}

\noindent
Discutons rapidement de quelques propriétés des groupes abéliens topologiques.
Un groupe abélien $M$ est un {\it groupe abélien topologique} si $M$ est
muni d'une topologie telle que l'addition $M \times M \to M$,
$(x, y) \mapsto x + y$, et l'inversion $M \to M$, $x \mapsto -x$, soient
continues. Un {\it homomorphisme de groupes abéliens topologiques}
est simplement un homomorphisme de groupes abéliens qui est continu.
La catégorie des groupes topologiques commutatifs est additive et possède
des noyaux et des conoyaux, mais n'est pas abélienne (car l'axiome
$\Im = \Coim$ n'est pas vérifié). Si $N \subset M$ est un sous-groupe,
nous considérons également $N$ et $M/N$ comme des groupes topologiques,
en utilisant la topologie induite sur $N$ et la topologie quotient sur $M/N$
(c'est-à-dire telle que $M \to M/N$ soit un submersif).
Remarquons que si $N \subset M$ est un sous-groupe ouvert, alors la topologie
sur $M/N$ est discrète.

\medskip\noindent
Nous disons que la topologie sur $M$ est {\it linéaire} s'il existe un
système fondamental de voisinages de $0$ constitué de sous-groupes. Dans ce
cas, ces sous-groupes sont également ouverts. En voici un exemple. Soit $I$
un ensemble dirigé et soit $G_i$ un système inverse de groupes abéliens
(discrets) sur $I$. Alors
$$
G = \lim_{i \in I} G_i
$$
ce groupe, muni de la topologie de limite inverse, est topologisé linéairement,
avec pour système fondamental de voisinages de $0$ les
$\Ker(G \to G_i)$. Réciproquement, soit $M$ un groupe abélien topologisé
linéairement. Choisissons un système fondamental quelconque de sous-groupes
ouverts $U_i \subset M$, $i \in I$ (c'est-à-dire que les $U_i$ forment un
système fondamental de voisinages ouverts et que chaque $U_i$ est un
sous-groupe de $M$). En posant
$i \geq i' \Leftrightarrow U_i \subset U_{i'}$, nous voyons que $I$
est dirigé. Nous obtenons un homomorphisme de groupes abéliens
topologisés linéairement
$$
c : M \longrightarrow \lim_{i \in I} M/U_i.
$$
Il est clair que $M$ est {\it séparé} (comme espace topologique)
si et seulement si $c$ est injectif. Nous disons que $M$ est {\it complet}
si $c$ est un isomorphisme\footnote{Nous incluons la séparation dans la
complétude, car nous voulons disposer de limites uniques dans les groupes
complets. Il existe une définition de la complétude pour tout groupe
topologique qui coïncide, à la question de séparation près, avec celle-ci
dans notre cas particulier.}.
Nous laissons au lecteur le soin de vérifier que cette condition est
indépendante du choix du système fondamental de sous-groupes ouverts
$\{U_i\}_{i \in I}$ choisi ci-dessus. En fait, le groupe abélien topologique
$M^\wedge = \lim_{i \in I} M/U_i$ est indépendant de ce choix et est parfois
appelé le {\it complété} de $M$.
Tout $G = \lim G_i$ comme ci-dessus est complet ; en particulier, le complété
$M^\wedge$ est toujours complet.

\begin{definition}[Anneaux topologiques]
\label{definition-topological-ring}
\begin{reference}
\cite[Chapitre 0, Sections 7.1 et 7.2]{EGA1}
\end{reference}
Soit $R$ un anneau et soit $M$ un $R$-module.
\begin{enumerate}
\item Nous disons que $R$ est un {\it anneau topologique} si $R$ est muni
d'une topologie telle que l'addition et la multiplication soient continues
comme applications $R \times R \to R$, où $R \times R$ est muni de la
topologie produit. Dans ce cas, nous disons que $M$ est un {\it module
topologique} si $M$ est muni d'une topologie telle que l'addition
$M \times M \to M$ et la multiplication scalaire $R \times M \to M$ soient
continues.
\item Un {\it homomorphisme de modules topologiques} est simplement une
application de $R$-modules continue. Un {\it homomorphisme d'anneaux
topologiques} est un homomorphisme d'anneaux continu pour les topologies
données.
\item Nous disons que $M$ est {\it topologisé linéairement} si $0$ possède
un système fondamental de voisinages constitué de sous-modules. Nous disons
que $R$ est {\it topologisé linéairement} si $0$ possède un système
fondamental de voisinages constitué d'idéaux.
\item Si $R$ est topologisé linéairement, nous disons que $I \subset R$ est
un {\it idéal de définition} si $I$ est ouvert et si tout voisinage de $0$
contient $I^n$ pour un certain $n$.
\item Si $R$ est topologisé linéairement, nous disons que $R$ est
{\it pré-admissible} si $R$ possède un idéal de définition.
\item Si $R$ est topologisé linéairement, nous disons que $R$ est
{\it admissible} s'il est pré-admissible et
complet\footnote{Selon nos conventions, cela inclut séparé.}.
\item Si $R$ est topologisé linéairement, nous disons que $R$ est
{\it pré-adique} s'il existe un idéal de définition $I$ tel que
$\{I^n\}_{n \geq 0}$ forme un système fondamental de voisinages de $0$.
\item Si $R$ est topologisé linéairement, nous disons que $R$ est
{\it adique} si $R$ est pré-adique et complet.
\end{enumerate}
Remarquons qu'un anneau topologique (pré-)adique est la même chose qu'un
anneau topologique (pré-)admissible qui possède un idéal de définition $I$
tel que $I^n$ soit ouvert pour tout $n \geq 1$.
\end{definition}

\noindent
Soit $R$ un anneau et soit $M$ un $R$-module. Soit $I \subset R$ un idéal.
Nous pouvons considérer sur $R$ la topologie linéaire dont
$\{I^n\}_{n \geq 0}$ est un système fondamental de voisinages de $0$.
Cette topologie est appelée la topologie $I$-adique ; $R$ est un anneau
topologique pré-adique pour la topologie $I$-adique\footnote{Ainsi, la topologie
$I$-adique est parfois appelée topologie $I$-pré-adique.}. De plus, la
topologie linéaire sur $M$ dont $\{I^nM\}_{n \geq 0}$ est un système
fondamental de voisinages ouverts de $0$ fait de $M$ un $R$-module
topologique. C'est la {\it topologie $I$-adique} sur $M$. Nous voyons que
$M$ est complet pour la topologie $I$-adique (au sens de la
Définition \ref{algebra-definition-complete} d'Algèbre)
si et seulement si $M$ est complet dans la topologie $I$-adique\footnote{
Il peut arriver que le complété $I$-adique $M^\wedge$ ne soit pas complet
pour la topologie $I$-adique, bien que $M^\wedge$ soit toujours complet
pour la topologie limite. Si $I$ est de type fini, la topologie $I$-adique
et la topologie limite sur $M^\wedge$ coïncident, voir Algèbre,
Lemme \ref{algebra-lemma-hathat-finitely-generated} et sa démonstration.}.
En particulier, $R$ est complet pour la topologie $I$-adique si et seulement
si $R$ est un anneau topologique adique pour la topologie $I$-adique.

\medskip\noindent
Comme cas particulier, remarquons que la topologie discrète est la topologie
$0$-adique et que tout anneau muni de la topologie discrète est adique.

\begin{lemma}
\label{lemma-continuous}
Soit $\varphi : R \to S$ un homomorphisme d'anneaux.
Soient $I \subset R$ et $J \subset S$ des idéaux, et munissons $R$ de la
topologie $I$-adique et $S$ de la topologie $J$-adique. Alors $\varphi$ est
un homomorphisme d'anneaux topologiques si et seulement si
$\varphi(I^n) \subset J$ pour un certain $n \geq 1$.
\end{lemma}

\begin{proof}
Omis.
\end{proof}

\begin{lemma}[Théorème de Baire]
\label{lemma-baire-category-complete-module}
Soit $M$ un groupe abélien topologique. Supposons que $M$ soit topologisé
linéairement, complet et possède un système fondamental dénombrable de
voisinages de $0$. Si $U_n \subset M$, $n \geq 1$, sont des parties ouvertes
et denses, alors $\bigcap_{n \geq 1} U_n$ est dense.
\end{lemma}

\begin{proof}
Soient $U_n$ comme dans l'énoncé du lemme. En remplaçant $U_n$ par
$U_1 \cap \ldots \cap U_n$, nous pouvons supposer que
$U_1 \supset U_2 \supset \ldots$.
Soit $M_n$, $n \in \mathbf{N}$, un système fondamental de voisinages de $0$.
Nous pouvons supposer que $M_{n + 1} \subset M_n$. Choisissons $x \in M$.
Nous montrerons que, pour tout $k \geq 1$, il existe un
$y \in \bigcap_{n \geq 1} U_n$ tel que $x - y \in M_k$.

\medskip\noindent
Pour construire $y$, procédons comme suit.
Tout d'abord, choisissons $y_1 \in U_1$ tel que $y_1 \in x + M_k$.
C'est possible car $U_1$ est dense et $x + M_k$ est ouvert.
Choisissons ensuite $k_1 > k$ tel que $y_1 + M_{k_1} \subset U_1$.
C'est possible car $U_1$ est ouvert.
Puis choisissons $y_2 \in U_2$ tel que $y_2 \in y_1 + M_{k_1}$.
C'est possible car $U_2$ est dense et $y_2 + M_{k_1}$ est ouvert.
Choisissons ensuite $k_2 > k_1$ tel que $y_2 + M_{k_2} \subset U_2$.
C'est possible car $U_2$ est ouvert.

\medskip\noindent
En poursuivant de cette manière, nous obtenons une suite convergente
$y_i$ d'éléments de $M$, de limite $y$. Par construction,
$x - y \in M_k$. Comme
$$
y - y_i = (y_{i + 1} - y_i) + (y_{i + 2} - y_{i + 1}) + \ldots
$$
appartient à $M_{k_i}$, nous voyons que $y \in y_i + M_{k_i} \subset U_i$
pour tout $i$, comme voulu.
\end{proof}

\begin{lemma}
\label{lemma-consequence-baire-complete-module}
Avec les mêmes hypothèses que dans le Lemme
\ref{lemma-baire-category-complete-module}, si
$M = \bigcup_{n \geq 1} N_n$ pour des sous-groupes fermés $N_n$,
alors $N_n$ est ouvert pour un certain $n$.
\end{lemma}

\begin{proof}
Sinon, $U_n = M \setminus N_n$ est dense pour tout $n$, ce qui contredit le
Lemme \ref{lemma-baire-category-complete-module}.
\end{proof}

\begin{lemma}[Lemme d'application ouverte]
\label{lemma-open-mapping}
Soit $u : N \to M$ une application continue de groupes abéliens topologisés.
Supposons que $M$ soit séparé et que $N$ soit complet, topologisé
linéairement, et possède un système fondamental dénombrable de voisinages de
$0$. Alors exactement l'une des assertions suivantes est vraie :
\begin{enumerate}
\item $u$ est ouverte, ou
\item pour un sous-groupe ouvert $N' \subset N$, l'image $u(N')$ est
nulle part dense dans $M$.
\end{enumerate}
\end{lemma}

\begin{proof}
Soit $N_n$, $n \in \mathbf{N}$, un système fondamental de voisinages de $0$.
Nous pouvons supposer que $N_n$ est un sous-groupe et que
$N_{n + 1} \subset N_n$. Si (2) n'est pas satisfaite, alors
l'adhérence $M_n$ de $u(N_n)$ est un sous-groupe ouvert pour
$n = 1, 2, 3, \ldots$.
Comme $u$ est continue, nous voyons que les $M_n$, $n \in \mathbf{N}$,
forment un système fondamental de voisinages ouverts de $0$ dans $M$.
De plus, puisque $M_n$ est l'adhérence de $u(N_n)$, nous avons
$$
u(N_n) + M_{n + 1} = M_n
$$
pour tout $n \geq 1$. Choisissons $x_1 \in M_1$. Nous pouvons alors choisir
par récurrence $y_i \in N_i$ et $x_{i + 1} \in M_{i + 1}$ tels que
$$
u(y_i) + x_{i + 1} = x_i
$$
L'élément $y = y_1 + y_2 + y_3 + \ldots$ de $N$ existe car $N$ est complet.
Nous voyons alors que $x_1 = u(y)$ puisque $M$ est séparé. Ainsi
$M_1 = u(N_1)$. De la même manière, le lecteur montre que
$M_i = u(N_i)$ pour tout $i \geq 2$, et nous voyons que $u$ est ouverte.
\end{proof}









\section{Morphismes formellement lisses d'anneaux topologiques}
\label{section-formally-smooth}

\noindent
Il existe une version de la lissité formelle qui s'applique aux
homomorphismes d'anneaux topologiques.

\begin{definition}
\label{definition-formally-smooth}
Soit $R \to S$ un homomorphisme d'anneaux topologiques, $R$ et $S$ étant
topologisés linéairement. Nous disons que $S$ est {\it formellement lisse
sur $R$} si, pour tout diagramme commutatif solide
$$
\xymatrix{
S \ar[r] \ar@{-->}[rd] & A/J \\
R \ar[r] \ar[u] & A \ar[u]
}
$$
d'homomorphismes d'anneaux topologiques où $A$ est un anneau discret et
$J \subset A$ un idéal de carré nul, il existe une flèche en pointillés
rendant le diagramme commutatif.
\end{definition}

\noindent
Nous utiliserons surtout cette notion lorsqu'on dispose d'idéaux
$\mathfrak m \subset R$ et $\mathfrak n \subset S$ et que l'on munit $R$ de
la topologie $\mathfrak m$-adique et $S$ de la topologie $\mathfrak n$-adique.
La continuité de $\varphi : R \to S$ équivaut à
$\varphi(\mathfrak m^m) \subset \mathfrak n$ pour un certain $m \geq 1$, voir
le Lemme \ref{lemma-continuous}. Il s'avère que, dans ce cas, seule la
topologie sur $S$ intervient.

\begin{lemma}
\label{lemma-formally-smooth}
Soit $\varphi : R \to S$ un homomorphisme d'anneaux.
\begin{enumerate}
\item Si $R \to S$ est formellement lisse au sens de la
Définition \ref{algebra-definition-formally-smooth} d'Algèbre,
alors $R \to S$ est formellement lisse pour toute topologie linéaire sur $R$
et toute topologie pré-adique sur $S$ pour lesquelles $R \to S$ est continu.
\item Soient $\mathfrak n \subset S$ et $\mathfrak m \subset R$ des idéaux
tels que $\varphi$ soit continu pour la topologie $\mathfrak m$-adique sur
$R$ et la topologie $\mathfrak n$-adique sur $S$. Alors les assertions
suivantes sont équivalentes :
\begin{enumerate}
\item $\varphi$ est formellement lisse pour la topologie $\mathfrak m$-adique
sur $R$ et la topologie $\mathfrak n$-adique sur $S$, et
\item $\varphi$ est formellement lisse pour la topologie discrète
sur $R$ et la topologie $\mathfrak n$-adique sur $S$.
\end{enumerate}
\end{enumerate}
\end{lemma}

\begin{proof}
Supposons que $R \to S$ soit formellement lisse au sens de la
Définition \ref{algebra-definition-formally-smooth} d'Algèbre.
Si $S$ possède une topologie pré-adique, il existe un idéal
$\mathfrak n \subset S$ tel que $S$ soit muni de la topologie $\mathfrak n$-adique.
Supposons donné un diagramme commutatif solide comme dans la
Définition \ref{definition-formally-smooth}.
La continuité de $S \to A/J$ signifie que $\mathfrak n^k$ s'envoie sur zéro
dans $A/J$ pour un certain $k \geq 1$, voir le Lemme \ref{lemma-continuous}.
Nous obtenons un homomorphisme d'anneaux $\psi : S \to A$ grâce à la
lissité formelle supposée de $S$ sur $R$. Alors $\psi(\mathfrak n^k) \subset J$, d'où
$\psi(\mathfrak n^{2k}) = 0$ puisque $J^2 = 0$. Ainsi $\psi$ est continu
d'après le Lemme \ref{lemma-continuous}. Cela démontre (1).

\medskip\noindent
La démonstration de (2)(b) $\Rightarrow$ (2)(a) est la même que celle de (1).
Supposons (2)(a). Considérons un diagramme commutatif solide comme dans la
Définition \ref{definition-formally-smooth}, où nous utilisons la topologie
discrète sur $R$. Comme $\varphi$ est continu, nous avons
$\varphi(\mathfrak m^n) \subset \mathfrak n$ pour un certain $n \geq 1$.
Comme $S \to A/J$ est continu, $\mathfrak n^k$ s'envoie sur zéro dans $A/J$
pour un certain $k \geq 1$. Ainsi $\mathfrak m^{nk}$ s'envoie dans $J$ par
l'application $R \to A$. Donc $\mathfrak m^{2nk}$ s'envoie sur zéro dans $A$
et $R \to A$ est continu pour la topologie $\mathfrak m$-adique.
L'assertion (2)(a) fournit alors la flèche en pointillés voulue.
\end{proof}

\begin{definition}
\label{definition-formally-smooth-adic}
Soit $R \to S$ un homomorphisme d'anneaux. Soit $\mathfrak n \subset S$ un
idéal. Si les conditions équivalentes (2)(a) et (2)(b) du
Lemme \ref{lemma-formally-smooth} sont satisfaites, nous disons que
$R \to S$ est {\it formellement lisse pour la topologie $\mathfrak n$-adique}.
\end{definition}

\noindent
Cette propriété passe aux complétés.

\begin{lemma}
\label{lemma-formally-smooth-completion}
Soient $(R, \mathfrak m)$ et $(S, \mathfrak n)$ des anneaux munis
d'idéaux de type fini. Munissons $R$ et $S$ des topologies
$\mathfrak m$-adique et $\mathfrak n$-adique.
Soit $R \to S$ un homomorphisme d'anneaux topologiques.
Les assertions suivantes sont équivalentes :
\begin{enumerate}
\item $R \to S$ est formellement lisse pour la topologie $\mathfrak n$-adique,
\item $R \to S^\wedge$ est formellement lisse pour la topologie
$\mathfrak n^\wedge$-adique,
\item $R^\wedge \to S^\wedge$ est formellement lisse pour la topologie
$\mathfrak n^\wedge$-adique.
\end{enumerate}
Ici, $R^\wedge$ et $S^\wedge$ sont les complétés $\mathfrak m$-adique et
$\mathfrak n$-adique de $R$ et $S$.
\end{lemma}

\begin{proof}
L'hypothèse que $\mathfrak m$ est de type fini implique que
$R^\wedge$ est complet pour la topologie $\mathfrak mR^\wedge$-adique,
que $\mathfrak mR^\wedge = \mathfrak m^\wedge$ et que
$R^\wedge/\mathfrak m^nR^\wedge = R/\mathfrak m^n$,
voir Algèbre, Lemme \ref{algebra-lemma-hathat-finitely-generated}
et sa démonstration. De même pour $(S, \mathfrak n)$. Il est donc clair que
les diagrammes comme dans la Définition \ref{definition-formally-smooth}
pour les cas (1), (2) et (3) sont en correspondance biunivoque.
\end{proof}

\noindent
L'avantage de travailler avec des anneaux adiques est d'obtenir une
propriété de relèvement plus forte.

\begin{lemma}
\label{lemma-lift-continuous}
Soit $R \to S$ un homomorphisme d'anneaux. Soit $\mathfrak n$ un idéal de $S$.
Supposons que $R \to S$ soit formellement lisse pour la topologie
$\mathfrak n$-adique. Considérons un diagramme commutatif solide
$$
\xymatrix{
S \ar[r]_\psi \ar@{-->}[rd] & A/J \\
R \ar[r] \ar[u] & A \ar[u]
}
$$
d'homomorphismes d'anneaux topologiques où $A$ est adique et $A/J$ est le
quotient (comme anneau topologique) de $A$ par un idéal fermé
$J \subset A$ tel que $J^t$ soit contenu dans un idéal de définition
de $A$ pour un certain $t \geq 1$. Alors il existe une flèche en pointillés
dans la catégorie des anneaux topologiques rendant le diagramme commutatif.
\end{lemma}

\begin{proof}
Soit $I \subset A$ un idéal de définition tel que $I \supset J^t$
pour un certain $t$. Alors $A = \lim A/I^n$ et $A/J = \lim A/J + I^n$
car $J$ est supposé fermé. Considérons le diagramme suivant d'$R$-algèbres
discrètes $A_{n, m} = A/J^n + I^m$ :
$$
\xymatrix{
A/J^3 + I^3 \ar[r] \ar[d] &
A/J^2 + I^3 \ar[r] \ar[d] &
A/J + I^3 \ar[d] \\
A/J^3 + I^2 \ar[r] \ar[d] &
A/J^2 + I^2 \ar[r] \ar[d] &
A/J + I^2 \ar[d] \\
A/J^3 + I \ar[r] &
A/J^2 + I \ar[r] &
A/J + I
}
$$
Remarquons que chaque carré commutatif définit une surjection
$$
A_{n + 1, m + 1} \longrightarrow A_{n + 1, m} \times_{A_{n, m}} A_{n, m + 1}
$$
de $R$-algèbres dont le noyau est de carré nul.
Nous allons construire par récurrence des homomorphismes de $R$-algèbres
$\varphi_{n, m} : S \to A_{n, m}$.
Nous avons les applications $\varphi_{1, m} = \psi \bmod J + I^m$.
Remarquons que chacune de ces applications est continue puisque $\psi$ l'est.
Nous pouvons choisir par récurrence les applications $\varphi_{n, 1}$ en
partant de $\varphi_{1, 1}$ et en relevant, grâce à la lissité formelle de
$S$ sur $R$, le long de la ligne inférieure du diagramme ci-dessus.
Nous construisons les applications restantes $\varphi_{n, m}$ par récurrence
sur $n + m$. Nous choisissons $\varphi_{n + 1, m + 1}$ en relevant le couple
$(\varphi_{n + 1, m}, \varphi_{n, m + 1})$ le long de la surjection affichée
(toujours grâce à la lissité formelle de $S$ sur $R$). Ainsi toutes les
applications $\varphi_{n, m}$ sont compatibles avec les applications de
transition du système.
Comme $J^t \subset I$, nous voyons par exemple que
$\varphi_n = \varphi_{nt, n} \bmod I^n$ induit une application $S \to A/I^n$.
En prenant la limite $\varphi = \lim \varphi_n$, nous obtenons une application
$S \to A = \lim A/I^n$. La composée vers $A/J$ est égale à $\psi$, puisque
nous avons vu que $A/J = \lim A/J + I^n$.
Enfin, montrons que $\varphi$ est continue. Nous savons que
$\psi(\mathfrak n^r) \subset J + I/J$ pour un certain $r \geq 1$ par
notre hypothèse que $\psi$ est un morphisme d'anneaux topologiques, voir le
Lemme \ref{lemma-continuous}. Ainsi $\varphi(\mathfrak n^r) \subset J + I$,
et donc $\varphi(\mathfrak n^{rt}) \subset I$, comme voulu.
\end{proof}

\begin{lemma}
\label{lemma-increase-ideal}
Soit $R \to S$ un homomorphisme d'anneaux. Soient
$\mathfrak n \subset \mathfrak n' \subset S$ des idéaux. Si $R \to S$ est
formellement lisse pour la topologie $\mathfrak n$-adique, alors
$R \to S$ est formellement lisse pour la topologie $\mathfrak n'$-adique.
\end{lemma}

\begin{proof}
Omis.
\end{proof}

\begin{lemma}
\label{lemma-compose-formally-smooth}
La composée d'homomorphismes continus formellement lisses d'anneaux
topologisés linéairement est formellement lisse.
\end{lemma}

\begin{proof}
Omis. (Indication : c'est entièrement formel et résulte de la considération
d'un diagramme convenable.)
\end{proof}

\begin{lemma}
\label{lemma-base-change-fs}
Soient $R$, $S$ des anneaux. Soit $\mathfrak n \subset S$ un idéal.
Supposons $R \to S$ formellement lisse pour la topologie $\mathfrak n$-adique.
Soit $R \to R'$ un homomorphisme d'anneaux quelconque. Alors
$R' \to S' = S \otimes_R R'$ est formellement lisse pour la topologie
$\mathfrak n' = \mathfrak nS'$-adique.
\end{lemma}

\begin{proof}
Soit donné un diagramme solide
$$
\xymatrix{
S \ar[r] \ar@{-->}[rrd] & S' \ar[r] \ar@{-->}[rd] & A/J \\
R  \ar[u] \ar[r] & R' \ar[r] \ar[u] & A \ar[u]
}
$$
comme dans la Définition \ref{definition-formally-smooth}.
Alors la composée $S \to S' \to A/J$ est continue.
Par hypothèse, la flèche en pointillés longue existe. Par la propriété
universelle du produit tensoriel, nous obtenons la flèche en pointillés courte.
\end{proof}

\noindent
Nous avons vu la descente de la lissité formelle le long des homomorphismes
fidèlement plats dans Algèbre, Lemme
\ref{algebra-lemma-descent-formally-smooth}. Un résultat analogue vaut dans
le cadre actuel des anneaux topologiques. Ici, nous démontrons seulement la
version très simple et facile suivante, qui est déjà fort utile.

\begin{lemma}
\label{lemma-descent-fs}
Soient $R$, $S$ des anneaux. Soit $\mathfrak n \subset S$ un idéal.
Soit $R \to R'$ un homomorphisme d'anneaux. Posons
$S' = S \otimes_R R'$ et $\mathfrak n' = \mathfrak nS$. Si
\begin{enumerate}
\item l'application $R \to R'$ plonge $R$ comme facteur direct de $R'$
en tant que $R$-module, et
\item $R' \to S'$ est formellement lisse pour la topologie
$\mathfrak n'$-adique,
\end{enumerate}
alors $R \to S$ est formellement lisse pour la topologie $\mathfrak n$-adique.
\end{lemma}

\begin{proof}
Soit donné un diagramme solide
$$
\xymatrix{
S \ar[r] & A/J \\
R \ar[u] \ar[r] & A \ar[u]
}
$$
comme dans la Définition \ref{definition-formally-smooth}.
Posons $A' = A \otimes_R R'$ et $J' = \Im(J \otimes_R R' \to A')$.
Le changement de base du diagramme ci-dessus est le diagramme
$$
\xymatrix{
S' \ar[r] \ar@{-->}[rd]^{\psi'} & A'/J' \\
R' \ar[u] \ar[r] & A' \ar[u]
}
$$
à flèches continues. Par la condition (2), nous obtenons la flèche en
pointillés $\psi' : S' \to A'$. Par la condition (1), choisissons une
décomposition en facteurs directs $R' = R \oplus C$ de $R$-modules.
(Attention : $C$ n'est pas un idéal de $R'$.)
Alors $A' = A \oplus A \otimes_R C$. Posons
$$
J'' = \Im(J \otimes_R C \to A \otimes_R C) \subset J' \subset A'.
$$
Alors $J' = J \oplus J''$ comme $A$-modules. L'image de la composée
$\psi : S \to A'$ de $\psi'$ avec $S \to S'$ est contenue dans
$A + J' = A \oplus J''$. Cependant, dans l'anneau
$A + J' = A \oplus J''$, le sous-module $A$-module $J''$ est un idéal !
(Utiliser $J^2 = 0$.) Ainsi la composée
$S \to A + J' \to (A + J')/J'' = A$ est la flèche recherchée.
\end{proof}








\section{Morphismes formellement lisses d'anneaux locaux}
\label{section-formally-smooth-local}

\noindent
Dans le cas d'un homomorphisme local d'anneaux locaux, on peut
limiter les diagrammes pour lesquels il faut vérifier la propriété de
relèvement. Comparer avec Algèbre, Lemme
\ref{algebra-lemma-smooth-test-artinian}.

\begin{lemma}
\label{lemma-fs-local}
Soit $(R, \mathfrak m) \to (S, \mathfrak n)$ un homomorphisme local
d'anneaux locaux. Les conditions suivantes sont équivalentes :
\begin{enumerate}
\item $R \to S$ est formellement lisse pour la topologie
$\mathfrak n$-adique ;
\item pour tout diagramme commutatif solide
$$
\xymatrix{
S \ar[r] \ar@{-->}[rd] & A/J \\
R \ar[r] \ar[u] & A \ar[u]
}
$$
d'homomorphismes locaux d'anneaux locaux où $J \subset A$ est
un idéal de carré nul, $\mathfrak m_A^n = 0$ pour un certain $n > 0$, et
$S \to A/J$ induit un isomorphisme sur les corps résiduels, il existe
une flèche en pointillés rendant le diagramme commutatif.
\end{enumerate}
Si $S$ est noethérien, ces conditions sont également équivalentes à
\begin{enumerate}
\item[(3)] la même condition que (2), mais seulement pour les diagrammes
où, de plus, $A \to A/J$ est une petite extension
(Algèbre, Définition \ref{algebra-definition-small-extension}).
\end{enumerate}
\end{lemma}

\begin{proof}
L'implication (1) $\Rightarrow$ (2) découle des définitions.
Considérons un diagramme
$$
\xymatrix{
S \ar[r] \ar@{-->}[rd] & A/J \\
R \ar[r] \ar[u] & A \ar[u]
}
$$
comme dans la Définition \ref{definition-formally-smooth} pour la
topologie $\mathfrak m$-adique sur $R$ et la topologie
$\mathfrak n$-adique sur $S$. Choisissons $m > 0$ tel que
$\mathfrak n^m(A/J) = 0$ (ce qui est possible par continuité des
applications du diagramme). Considérons le sous-anneau $A'$ de $A$
qui est l'image réciproque de l'image de $S$ dans $A/J$. Posons
$J' = J$ vu comme idéal de $A'$. Alors $J'$ est un idéal de carré nul
de $A'$ et $A'/J'$ est un quotient de $S/\mathfrak n^m$. Ainsi $A'$ est
local et $\mathfrak m_{A'}^{2m} = 0$. Nous obtenons donc un diagramme
$$
\xymatrix{
S \ar[r] \ar@{-->}[rd] & A'/J' \\
R \ar[r] \ar[u] & A' \ar[u]
}
$$
comme dans (2). Si nous pouvons construire la flèche en pointillés dans
ce diagramme, nous obtenons alors la flèche en pointillés du diagramme
initial en la composant avec $A' \to A$. Ainsi (2) implique (1).

\medskip\noindent
Supposons $S$ noethérien. L'implication (1) $\Rightarrow$ (3) est
immédiate. Supposons (3) et donnons-nous un diagramme comme dans (2).
Alors $\mathfrak m_A^n J = 0$ pour un certain $n > 0$. En considérant les
applications
$$
A \to A/\mathfrak m_A^{n - 1}J \to \ldots \to A/\mathfrak mJ \to A/J
$$
il suffit de construire le relèvement lorsque $\mathfrak m_A J = 0$.
Supposons donc $\mathfrak m_A J = 0$ et soit $A' \subset A$ l'anneau
construit ci-dessus. Alors $A'/J'$ est artinien, comme quotient de
l'anneau local artinien $S/\mathfrak n^m$. Il suffit donc de montrer que,
grâce à la propriété (3), nous pouvons trouver la flèche en pointillés
dans les diagrammes comme dans (2) avec $A/J$ artinien et
$\mathfrak m_A J = 0$. Soit $\kappa$ le corps résiduel commun de $A$,
$A/J$ et $S$. Par (3), si $J_0 \subset J$ est un idéal tel que
$\dim_\kappa(J/J_0) = 1$, nous pouvons construire une flèche en
pointillés $S \to A/J_0$. En prenant le produit, nous obtenons
$$
S \longrightarrow \prod\nolimits_{J_0 \text{ comme ci-dessus}} A/J_0
$$
L'image de cette flèche est clairement contenue dans la sous-$R$-algèbre
$A'$ des éléments qui se projettent sur la petite diagonale
$A/J \subset \prod_{J_0} A/J$. Soit $J' \subset A'$ l'ensemble des
éléments qui s'envoient sur zéro dans $A/J$. Alors $J'$ est un idéal de
carré nul et, en tant qu'espace vectoriel sur $\kappa$, il est égal à
$$
J' = \prod\nolimits_{J_0 \text{ comme ci-dessus}} J/J_0
$$
Ainsi l'application $J \to J'$ est injective. Par la théorie des espaces
vectoriels, nous pouvons choisir une décomposition $J' = J \oplus M$. Il
en résulte que
$$
A' = A \oplus M
$$
en tant que $R$-algèbres. On peut donc composer l'application $S \to A'$
avec la projection $A' \to A$ pour obtenir la flèche en pointillés
recherchée, ce qui achève la démonstration du lemme.
\end{proof}

\noindent
Le lemme suivant sera amélioré dans la Section
\ref{section-regular-fs}.

\begin{lemma}
\label{lemma-fs-implies-regular}
Soit $k$ un corps et soit $(A, \mathfrak m, K)$ une $k$-algèbre locale
noethérienne. Si $k \to A$ est formellement lisse pour la topologie
$\mathfrak m$-adique, alors $A$ est un anneau local régulier.
\end{lemma}

\begin{proof}
Soit $k_0 \subset k$ le corps premier. Alors $k_0$ est parfait, donc
$k / k_0$ est séparable, donc formellement lisse d'après le Lemme
d'Algèbre \ref{algebra-lemma-formally-smooth-extensions-easy}. D'après les
Lemmes \ref{lemma-formally-smooth} et \ref{lemma-compose-formally-smooth},
$k_0 \to A$ est formellement lisse pour la topologie $\mathfrak m$-adique
sur $A$. On peut donc supposer $k = \mathbf{Q}$ ou $k = \mathbf{F}_p$.

\medskip\noindent
D'après les Lemmes d'Algèbre
\ref{algebra-lemma-completion-faithfully-flat} et
\ref{algebra-lemma-flat-under-regular}, il suffit de montrer que le complété
$A^\wedge$ est régulier. D'après le Lemme
\ref{lemma-formally-smooth-completion}, on peut remplacer $A$ par
$A^\wedge$. On peut donc supposer que $A$ est un anneau local noethérien
complet. Le théorème de structure de Cohen (Algèbre, Théorème
\ref{algebra-theorem-cohen-structure-theorem}) fournit une application
$K \to A$. Comme $k$ est le corps premier, nous voyons que $K \to A$ est un
homomorphisme de $k$-algèbres.

\medskip\noindent
Soient $x_1, \ldots, x_n \in \mathfrak m$ des éléments dont les images
forment une base de $\mathfrak m/\mathfrak m^2$. Posons
$T = K[[X_1, \ldots, X_n]]$. Remarquons que
$$
A/\mathfrak m^2 \cong K[x_1, \ldots, x_n]/(x_ix_j)
$$
et
$$
T/\mathfrak m_T^2 \cong K[X_1, \ldots, X_n]/(X_iX_j).
$$
Soit $A/\mathfrak m^2 \to T/m_T^2$ l'isomorphisme local de
$K$-algèbres qui envoie la classe de $x_i$ sur la classe de $X_i$.
Notons $f_1 : A \to T/\mathfrak m_T^2$ la composée de cet isomorphisme
avec la projection $A \to A/\mathfrak m^2$. L'hypothèse que
$k \to A$ est formellement lisse pour la topologie $\mathfrak m$-adique
signifie que nous pouvons relever $f_1$ en une application
$f_2 : A \to T/\mathfrak{m}_T^3$, puis en une application
$f_3 : A \to T/\mathfrak{m}_T^4$, et ainsi de suite, pour tout $n \geq 1$.
Attention : les applications $f_n$ sont des homomorphismes continus de
$k$-algèbres et peuvent ne pas être des homomorphismes de $K$-algèbres.
Nous obtenons une application induite
$f : A \to T = \lim T/\mathfrak m_T^n$ de $k$-algèbres locales.
Par notre choix de $f_1$, l'application $f$ induit un isomorphisme
$\mathfrak m/\mathfrak m^2 \to \mathfrak m_T/\mathfrak m_T^2$ ; par suite,
chaque $f_n$ est surjective et $f$ est surjective puisque $A$ est complet.
Cela implique $\dim(A) \geq \dim(T) = n$. Ainsi $A$ est régulier par
définition. (On en déduit également que $f$ est un isomorphisme.)
\end{proof}

\begin{lemma}
\label{lemma-lift-residue-field}
Soit $k$ un corps. Soit $(A, \mathfrak m, \kappa)$ une $k$-algèbre
locale complète. Si $\kappa/k$ est séparable, il existe un homomorphisme
de $k$-algèbres $\kappa \to A$ tel que $\kappa \to A \to \kappa$ soit
$\text{id}_\kappa$.
\end{lemma}

\begin{proof}
D'après la Proposition d'Algèbre
\ref{algebra-proposition-characterize-separable-field-extensions}
l'extension $\kappa/k$ est formellement lisse. D'après le Lemme
\ref{lemma-formally-smooth}, $k \to \kappa$ est formellement lisse au sens
de la Définition \ref{definition-formally-smooth}. Nous obtenons alors
$\kappa \to A$ par le Lemme \ref{lemma-lift-continuous}.
\end{proof}

\begin{lemma}
\label{lemma-power-series-over-residue-field}
Soit $k$ un corps. Soit $(A, \mathfrak m, \kappa)$ une $k$-algèbre
locale complète. Si $\kappa/k$ est séparable et si $A$ est régulier, il
existe un isomorphisme $A \cong \kappa[[t_1, \ldots, t_d]]$ de
$k$-algèbres.
\end{lemma}

\begin{proof}
Choisissons $\kappa \to A$ comme dans le Lemme
\ref{lemma-lift-residue-field} et appliquons le Lemme d'Algèbre
\ref{algebra-lemma-regular-complete-containing-coefficient-field}.
\end{proof}

\noindent
Le résultat suivant sera amélioré dans la Section
\ref{section-regular-fs}.

\begin{lemma}
\label{lemma-regular-implies-fs}
Soit $k$ un corps. Soit $(A, \mathfrak m, K)$ une $k$-algèbre locale
régulière telle que $K/k$ soit séparable. Alors $k \to A$ est
formellement lisse pour la topologie $\mathfrak m$-adique.
\end{lemma}

\begin{proof}
Il suffit de montrer que le complété de $A$ est formellement lisse sur
$k$, voir le Lemme \ref{lemma-formally-smooth-completion}. On peut donc
supposer que $A$ est une $k$-algèbre locale régulière complète dont le
corps résiduel $K$ est séparable sur $k$. D'après le Lemme
\ref{lemma-power-series-over-residue-field}, on a
$A = K[[x_1, \ldots, x_n]]$.

\medskip\noindent
L'anneau de séries formelles $K[[x_1, \ldots, x_n]]$ est formellement
lisse sur $k$. En effet, $K$ est formellement lisse sur $k$ et
$K[x_1, \ldots, x_n]$ est formellement lisse sur $K$ en tant qu'algèbre
polynomiale. Ainsi $K[x_1, \ldots, x_n]$ est formellement lisse sur $k$
d'après le Lemme d'Algèbre
\ref{algebra-lemma-compose-formally-smooth}. Il s'ensuit que
$k \to K[x_1, \ldots, x_n]$ est formellement lisse pour la topologie
$(x_1, \ldots, x_n)$-adique d'après le Lemme
\ref{lemma-formally-smooth}. Enfin, $k \to K[[x_1, \ldots, x_n]]$ est
formellement lisse pour la topologie $(x_1, \ldots, x_n)$-adique d'après
le Lemme \ref{lemma-formally-smooth-completion}.
\end{proof}

\begin{lemma}
\label{lemma-formally-smooth-finite-type}
Soit $A \to B$ un homomorphisme d'anneaux de type fini, avec $A$
noethérien. Soit $\mathfrak q \subset B$ un idéal premier au-dessus de
$\mathfrak p \subset A$. Les conditions suivantes sont équivalentes :
\begin{enumerate}
\item $A \to B$ est lisse en $\mathfrak q$ ;
\item $A_\mathfrak p \to B_\mathfrak q$ est formellement lisse pour
la topologie $\mathfrak q$-adique.
\end{enumerate}
\end{lemma}

\begin{proof}
L'implication (2) $\Rightarrow$ (1) découle du Lemme d'Algèbre
\ref{algebra-lemma-smooth-test-artinian}. Réciproquement, si $A \to B$
est lisse en $\mathfrak q$, alors $A \to B_g$ est lisse pour un certain
$g \in B$, $g \not \in \mathfrak q$. Alors $A \to B_g$ est formellement
lisse d'après la Proposition d'Algèbre
\ref{algebra-proposition-smooth-formally-smooth}. Ainsi
$A_\mathfrak p \to B_\mathfrak q$ est formellement lisse, car la
localisation conserve la lissité formelle (par exemple d'après le critère
de la Proposition d'Algèbre
\ref{algebra-proposition-characterize-formally-smooth} et le fait que le
complexe cotangent se comporte bien vis-à-vis de la localisation, voir
Algèbre, Lemmes \ref{algebra-lemma-NL-localize-bottom} et
\ref{algebra-lemma-localize-NL}).
Enfin, le Lemme \ref{lemma-formally-smooth} implique que
$A_\mathfrak p \to B_\mathfrak q$ est formellement lisse pour la topologie
$\mathfrak q$-adique.
\end{proof}









\section{Quelques résultats sur les anneaux de séries formelles}
\label{section-power-series}

\noindent
Les questions concernant les morphismes formellement lisses entre anneaux
locaux noethériens peuvent souvent se réduire à des questions sur des
morphismes entre anneaux de séries formelles. Dans cette section, nous
démontrons quelques lemmes auxiliaires pour faciliter ce type d'argument.

\begin{lemma}
\label{lemma-power-series-ring-over-Cohen-fs}
Soit $K$ un corps de caractéristique $0$ et $A = K[[x_1, \ldots, x_n]]$.
Soit $L$ un corps de caractéristique $p > 0$ et
$B = L[[x_1, \ldots, x_n]]$. Soit $\Lambda$ un anneau de Cohen. Posons
$C = \Lambda[[x_1, \ldots, x_n]]$.
\begin{enumerate}
\item $\mathbf{Q} \to A$ est formellement lisse pour la topologie
$\mathfrak m_A$-adique ;
\item $\mathbf{F}_p \to B$ est formellement lisse pour la topologie
$\mathfrak m_B$-adique ;
\item $\mathbf{Z} \to C$ est formellement lisse pour la topologie
$\mathfrak m_C$-adique.
\end{enumerate}
\end{lemma}

\begin{proof}
Par la propriété universelle des anneaux de séries formelles, il suffit de
démontrer :
\begin{enumerate}
\item $\mathbf{Q} \to K$ est formellement lisse ;
\item $\mathbf{F}_p \to L$ est formellement lisse ;
\item $\mathbf{Z} \to \Lambda$ est formellement lisse pour la topologie
$\mathfrak m_\Lambda$-adique.
\end{enumerate}
Les deux premières assertions résultent de la Proposition d'Algèbre
\ref{algebra-proposition-characterize-separable-field-extensions}.
La troisième résulte du Lemme d'Algèbre
\ref{algebra-lemma-cohen-ring-formally-smooth}, car, pour tout diagramme
test comme dans la Définition \ref{definition-formally-smooth}, une
certaine puissance de $p$ est nulle dans $A/J$ et donc une certaine
puissance de $p$ est nulle dans $A$.
\end{proof}

\begin{lemma}
\label{lemma-quotient-power-series-ring-over-Cohen}
Soit $K$ un corps et $A = K[[x_1, \ldots, x_n]]$.
Soit $\Lambda$ un anneau de Cohen et soit
$B = \Lambda[[x_1, \ldots, x_n]]$.
\begin{enumerate}
\item Si $y_1, \ldots, y_n \in A$ est un système régulier de paramètres,
alors $K[[y_1, \ldots, y_n]] \to A$ est un isomorphisme.
\item Si $z_1, \ldots, z_r \in A$ font partie d'un système régulier de
paramètres de $A$, alors $r \leq n$ et
$A/(z_1, \ldots, z_r) \cong K[[y_1, \ldots, y_{n - r}]]$.
\item Si $p, y_1, \ldots, y_n \in B$ est un système régulier de paramètres,
alors $\Lambda[[y_1, \ldots, y_n]] \to B$ est un isomorphisme.
\item Si $p, z_1, \ldots, z_r \in B$ font partie d'un système régulier de
paramètres de $B$, alors $r \leq n$ et
$B/(z_1, \ldots, z_r) \cong \Lambda[[y_1, \ldots, y_{n - r}]]$.
\end{enumerate}
\end{lemma}

\begin{proof}
Démonstration de (1). Posons $A' = K[[y_1, \ldots, y_n]]$. Il est clair
que l'application $A' \to A$ induit un isomorphisme
$A'/\mathfrak m_{A'}^n \to A/\mathfrak m_A^n$ pour tout $n \geq 1$.
Comme $A$ et $A'$ sont tous deux complets, on en déduit que $A' \to A$
est un isomorphisme.
Démonstration de (2). Complétons $z_1, \ldots, z_r$ en un système régulier
de paramètres $z_1, \ldots, z_r, y_1, \ldots, y_{n - r}$ de $A$. Considérons
l'application
$A' = K[[z_1, \ldots, z_r, y_1, \ldots, y_{n - r}]] \to A$.
Il s'agit d'un isomorphisme d'après (1). L'assertion (2) s'ensuit, car il
est clair que
$A'/(z_1, \ldots, z_r) \cong K[[y_1, \ldots, y_{n - r}]]$.
Les démonstrations de (3) et (4) sont exactement les mêmes que celles de
(1) et (2).
\end{proof}

\begin{lemma}
\label{lemma-embed-map-Noetherian-complete-local-rings}
Soit $A \to B$ un homomorphisme local d'anneaux locaux noethériens
complets. Il existe alors un diagramme commutatif
$$
\xymatrix{
S \ar[r] & B \\
R \ar[u] \ar[r] & A \ar[u]
}
$$
possédant les propriétés suivantes :
\begin{enumerate}
\item les flèches horizontales sont surjectives ;
\item si la caractéristique de $A/\mathfrak m_A$ est nulle, alors $S$ et
$R$ sont des anneaux de séries formelles sur des corps ;
\item si la caractéristique de $A/\mathfrak m_A$ est $p > 0$, alors $S$ et
$R$ sont des anneaux de séries formelles sur des anneaux de Cohen ;
\item $R \to S$ envoie un système régulier de paramètres de $R$ sur une
partie d'un système régulier de paramètres de $S$.
\end{enumerate}
En particulier, $R \to S$ est plat (voir le Lemme d'Algèbre
\ref{algebra-lemma-flat-over-regular}) et sa fibre régulière est
$S/\mathfrak m_R S$ (voir le Lemme d'Algèbre
\ref{algebra-lemma-regular-ring-CM}).
\end{lemma}

\begin{proof}
Utilisons le théorème de structure de Cohen
(Algèbre, Théorème \ref{algebra-theorem-cohen-structure-theorem}) pour
choisir une surjection $S \to B$ comme dans l'énoncé du lemme, où nous
prenons $S$ égal à un anneau de séries formelles sur un anneau de Cohen si
la caractéristique résiduelle est $p > 0$, et à un anneau de séries
formelles sur un corps sinon. Soit $J \subset S$ le noyau de $S \to B$.
Choisissons ensuite une surjection $R = \Lambda[[x_1, \ldots, x_n]] \to A$,
où $\Lambda$ est un anneau de Cohen si la caractéristique résiduelle de
$A$ est $p > 0$, et $\Lambda$ est le corps résiduel de $A$ sinon. Relevons la composée
$\Lambda[[x_1, \ldots, x_n]] \to A \to B$ en une application
$\varphi : R \to S$. Cela est possible car
$\Lambda[[x_1, \ldots, x_n]]$ est formellement lisse sur $\mathbf{Z}$ pour
la topologie $\mathfrak m$-adique (voir le Lemme
\ref{lemma-power-series-ring-over-Cohen-fs}), par application du Lemme
\ref{lemma-lift-continuous}. Enfin, remplaçons $\varphi$ par l'application
$\varphi' : R = \Lambda[[x_1, \ldots, x_n]] \to S' = S[[y_1, \ldots, y_n]]$
avec $\varphi'|_\Lambda = \varphi|_\Lambda$ et
$\varphi'(x_i) = \varphi(x_i) + y_i$. Remplaçons également $S \to B$
par l'application $S' \to B$ qui envoie $y_i$ sur zéro. Après ce
remplacement, il est clair qu'un système régulier de paramètres de $R$ est
envoyé sur une partie d'une suite régulière dans $S'$, ce qui conclut.
\end{proof}

\noindent
Il devrait exister une démonstration élémentaire du lemme suivant.

\begin{lemma}
\label{lemma-dominate-two-surjections}
Soient $S \to R$ et $S' \to R$ des applications surjectives d'anneaux
locaux noethériens complets. Alors $S \times_R S'$ est un anneau local
noethérien complet.
\end{lemma}

\begin{proof}
Soit $k$ le corps résiduel de $R$. Si la caractéristique de $k$ est
$p > 0$, notons $\Lambda$ un anneau de Cohen (Algèbre, Définition
\ref{algebra-definition-cohen-ring}) de corps résiduel $k$ (Algèbre,
Lemme \ref{algebra-lemma-cohen-rings-exist}). Si la caractéristique de
$k$ est $0$, posons $\Lambda = k$. Choisissons une surjection
$\Lambda[[x_1, \ldots, x_n]] \to R$ (comme dans le théorème de structure
de Cohen, voir Algèbre, Théorème
\ref{algebra-theorem-cohen-structure-theorem}) et relevons-la en des
applications $\Lambda[[x_1, \ldots, x_n]] \to S$,
$\varphi : \Lambda[[x_1, \ldots, x_n]] \to S$ et
$\varphi' : \Lambda[[x_1, \ldots, x_n]] \to S'$ à l'aide des Lemmes
\ref{lemma-power-series-ring-over-Cohen-fs} et
\ref{lemma-lift-continuous}. Choisissons ensuite
$f_1, \ldots, f_m \in S$ engendrant le noyau de $S \to R$ et
$f'_1, \ldots, f'_{m'} \in S'$ engendrant le noyau de $S' \to R$.
Alors l'application
$$
\Lambda[[x_1, \ldots, x_n, y_1, \ldots, y_m, z_1, \ldots, z_{m'}]]
\longrightarrow S \times_R S,
$$
qui envoie $x_i$ sur $(\varphi(x_i), \varphi'(x_i))$, $y_j$ sur
$(f_j, 0)$ et $z_{j'}$ sur $(0, f'_j)$ est surjective. Ainsi
$S \times_R S'$ est un quotient d'un anneau local complet, donc il est
complet.
\end{proof}
\section{Régularité géométrique et lissité formelle}
\label{section-regular-fs}

\noindent
Dans cette section, nous combinons les résultats des sections précédentes
pour établir la caractérisation suivante des anneaux locaux géométriquement
réguliers sur un corps. Nous reprenons ensuite certains de nos arguments afin
d'obtenir une caractérisation des homomorphismes formellement lisses pour la
topologie $\mathfrak m$-adique entre anneaux locaux noethériens.

\begin{theorem}
\label{theorem-regular-fs}
Soit $k$ un corps. Soit $(A, \mathfrak m, K)$ une $k$-algèbre locale
noethérienne. Si $k$ est de caractéristique nulle, les conditions suivantes
sont équivalentes :
\begin{enumerate}
\item $A$ est un anneau local régulier ;
\item $k \to A$ est formellement lisse pour la topologie $\mathfrak m$-adique.
\end{enumerate}
Si $k$ est de caractéristique $p > 0$, les conditions suivantes sont
équivalentes :
\begin{enumerate}
\item $A$ est géométriquement régulier sur $k$ ;
\item $k \to A$ est formellement lisse pour la topologie $\mathfrak m$-adique ;
\item pour tout $k \subset k' \subset k^{1/p}$ fini sur $k$, l'anneau
$A \otimes_k k'$ est régulier ;
\item $A$ est régulier et l'application canonique
$H_1(L_{K/k}) \to \mathfrak m/\mathfrak m^2$ est injective ;
\item $A$ est régulier et l'application
$\Omega_{k/\mathbf{F}_p} \otimes_k K \to \Omega_{A/\mathbf{F}_p} \otimes_A K$
est injective.
\end{enumerate}
\end{theorem}

\begin{proof}
Si $k$ est de caractéristique nulle, l'équivalence de (1) et (2) résulte des
lemmes \ref{lemma-fs-implies-regular} et \ref{lemma-regular-implies-fs}.

\medskip\noindent
Si $k$ est de caractéristique $p > 0$, la proposition
\ref{proposition-characterization-geometrically-regular} montre que (1),
(3), (4) et (5) sont équivalentes. Supposons (2). D'après le lemme
\ref{lemma-base-change-fs}, l'homomorphisme
$k' \to A' = A \otimes_k k'$ est formellement lisse pour la topologie
$\mathfrak m' = \mathfrak mA'$-adique. Ainsi, si $k \subset k'$ est une
extension finie purement inséparable, le lemme
\ref{lemma-fs-implies-regular} assure que $A'$ est un anneau local régulier.
On obtient donc (1).

\medskip\noindent
Enfin, montrons que (5) implique (2). Choisissons un diagramme dont les
flèches pleines sont données
$$
\xymatrix{
A \ar[r]_{\bar\psi} \ar@{-->}[rd] & B/J \\
k \ar[u]^i \ar[r]^\varphi & B \ar[u]_\pi
}
$$
comme dans la définition \ref{definition-formally-smooth}. Comme $J^2 = 0$,
$J$ possède une structure canonique de $B/J$-module et, par $\bar\psi$, une
structure de $A$-module. La continuité de $\bar\psi$ pour la topologie
$\mathfrak m$-adique entraîne $\mathfrak m^nJ = 0$ pour un certain $n$.
On peut donc filtrer $J$ par des sous-$B/J$-modules
$0 \subset J_1 \subset J_2 \subset \ldots \subset J_n = J$
de sorte que chaque quotient $J_{t + 1}/J_t$ soit annulé par $\mathfrak m$.
En considérant la suite d'homomorphismes d'anneaux
$B \to B/J_1 \to B/J_2 \to \ldots \to B/J$
on voit qu'il suffit d'établir l'existence de la flèche en pointillés lorsque
$J$ est annulé par $\mathfrak m$, c'est-à-dire lorsque $J$ est un espace
vectoriel sur $K$.

\medskip\noindent
Supposons donné un diagramme comme ci-dessus dans lequel $J$ est annulé par
$\mathfrak m$. Le lemme \ref{lemma-regular-implies-fs} montre que
$\mathbf{F}_p \to A$ est formellement lisse pour la topologie
$\mathfrak m$-adique. Il existe donc un homomorphisme d'anneaux
$\psi : A \to B$ tel que $\pi \circ \psi = \bar \psi$. Les deux
homomorphismes $\psi \circ i, \varphi : k \to B$ ont alors même composée
avec $\pi$. Par conséquent,
$D = \psi \circ i - \varphi : k \to J$ est une dérivation. D'après Algèbre,
lemme \ref{algebra-lemma-universal-omega}, on peut écrire
$D = \xi \circ \text{d}$ pour une application $k$-linéaire
$\xi : \Omega_{k/\mathbf{F}_p} \to J$. À l'aide de la structure d'espace
vectoriel sur $K$ de $J$, prolongeons $\xi$ en une application $K$-linéaire
$\xi' : \Omega_{k/\mathbf{F}_p} \otimes_k K \to J$.
En utilisant (5), on peut trouver une application $K$-linéaire
$\xi'' : \Omega_{A/\mathbf{F}_p} \otimes_A K$ dont la restriction à
$\Omega_{k/\mathbf{F}_p} \otimes_k K$ est $\xi'$. Posons
$$
D' : A \xrightarrow{\text{d}} \Omega_{A/\mathbf{F}_p}
\to \Omega_{A/\mathbf{F}_p} \otimes_A K \xrightarrow{\xi''} J.
$$
Enfin, posons $\psi' = \psi - D' : A \to B$. Le lecteur vérifiera que
$\psi'$ est un homomorphisme d'anneaux tel que
$\pi \circ \psi' = \bar \psi$ et $\psi' \circ i = \varphi$, comme souhaité.
\end{proof}

\begin{example}
\label{example-fs}
Soit $k$ un corps de caractéristique $p > 0$. Supposons que $a \in k$ ne
soit pas une puissance $p$-ième. Un exemple classique de $k$-algèbre locale
géométriquement régulière dont le corps résiduel est purement inséparable sur
$k$ est l'anneau
$$
A = k[x, y]_{(x, y^p - a)}/(y^p - a - x)
$$
En effet, $A$ est un localisé d'une algèbre lisse sur $k$ ; ainsi $k \to A$
est formellement lisse, donc $k \to A$ est formellement lisse pour la topologie
$\mathfrak m$-adique. Voici un exemple étroitement apparenté. Posons
$k = \mathbf{F}_p(s)$ et $K = \mathbf{F}_p(t)^{perf}$. Nous affirmons que
l'homomorphisme d'anneaux
$$
k \longrightarrow A = K[[x]],\quad s \longmapsto t + x
$$
est formellement lisse pour la topologie $(x)$-adique sur $A$. En effet,
$\Omega_{k/\mathbf{F}_p}$ est de dimension $1$, de base $\text{d}s$.
Il s'envoie sur l'élément $\text{d}x + \text{d}t = \text{d}x$ de
$\Omega_{A/\mathbf{F}_p}$. Nous laissons au lecteur le soin de montrer que
$\Omega_{A/\mathbf{F}_p}$ est libre de base $\text{d}x$ comme $A$-module.
La condition (5) du théorème \ref{theorem-regular-fs} est donc
satisfaite, et l'on conclut que $k \to A$ est formellement lisse pour la
topologie $(x)$-adique.
\end{example}

\begin{lemma}
\label{lemma-formally-smooth-flat}
Soit $A \to B$ un homomorphisme local d'anneaux locaux noethériens.
Supposons $A \to B$ formellement lisse pour la topologie
$\mathfrak m_B$-adique. Alors $A \to B$ est plat.
\end{lemma}

\begin{proof}
On peut supposer que $A$ et $B$ sont des anneaux locaux noethériens complets,
d'après le lemme \ref{lemma-formally-smooth-completion} et Algèbre, lemme
\ref{algebra-lemma-completion-Noetherian-Noetherian}. On utilise également
Algèbre, lemmes \ref{algebra-lemma-flatness-descends-more-general} et
\ref{algebra-lemma-completion-faithfully-flat}, pour voir que la platitude de
l'homomorphisme entre les complétés entraîne celle de $A \to B$.
Choisissons un diagramme commutatif
$$
\xymatrix{
S \ar[r] & B \\
R \ar[u] \ar[r] & A \ar[u]
}
$$
comme dans le lemme \ref{lemma-embed-map-Noetherian-complete-local-rings},
avec $R \to S$ plat. Soit $I \subset R$ le noyau de $R \to A$.
La lissité formelle de $B$ sur $A$ montre que l'homomorphisme de
$A$-algèbres
$$
S/IS \longrightarrow B
$$
admet une section, voir le lemme \ref{lemma-lift-continuous}.
Ainsi $B$ est facteur direct du $A$-module plat $S/IS$ (par changement de
base de la platitude, voir Algèbre, lemme
\ref{algebra-lemma-flat-base-change}) ; il est donc plat.
\end{proof}

\begin{lemma}
\label{lemma-formally-smooth-JZ}
Soit $A \to B$ un homomorphisme local d'anneaux locaux noethériens.
Supposons $A \to B$ formellement lisse pour la topologie
$\mathfrak m_B$-adique. Soit $K$ le corps résiduel de $B$. La suite de
Jacobi--Zariski associée à $A \to B \to K$ fournit alors une suite exacte
$$
0 \to H_1(\NL_{K/A}) \to \mathfrak m_B/\mathfrak m_B^2
\to \Omega_{B/A} \otimes_B K \to \Omega_{K/A} \to 0
$$
\end{lemma}

\begin{proof}
D'après Algèbre, lemme \ref{algebra-lemma-NL-surjection}, on a
$\mathfrak m_B/\mathfrak m_B^2 = H_1(\NL_{K/B})$. Par Algèbre, lemme
\ref{algebra-lemma-exact-sequence-NL}, il reste à montrer l'injectivité de
$H_1(\NL_{K/A}) \to \mathfrak m_B/\mathfrak m_B^2$. Si $k$ désigne le
corps résiduel de $A$, la suite de Jacobi--Zariski pour $A \to k \to K$
donne $\Omega_{K/A} = \Omega_{K/k}$ ainsi qu'une suite exacte
$$
\mathfrak m_A/\mathfrak m_A^2 \otimes_k K \to
H_1(\NL_{K/A}) \to
H_1(\NL_{K/k}) \to 0
$$
Posons $\overline{B} = B \otimes_A k$. Comme $\overline{B}$ est régulier,
l'idéal $\mathfrak m_{\overline{B}}$ est engendré par une suite régulière.
En appliquant les lemmes \ref{lemma-conormal-sequence-H1-regular} et
\ref{lemma-noetherian-finite-all-equivalent} à
$\mathfrak m_A B \subset \mathfrak m_B$, on obtient
$\mathfrak m_A B / (\mathfrak m_AB \cap \mathfrak m_B^2) =
\mathfrak m_A B / \mathfrak m_A \mathfrak m_B$, lequel est égal à
$\mathfrak m_A/\mathfrak m_A^2 \otimes_k K$, car $A \to B$ est plat par le
lemme \ref{lemma-formally-smooth-flat}. On a donc une suite exacte courte
$$
0 \to
\mathfrak m_A/\mathfrak m_A^2 \otimes_k K \to
\mathfrak m_B/\mathfrak m_B^2 \to
\mathfrak m_{\overline{B}}/\mathfrak m_{\overline{B}}^2 \to 0
$$
La fonctorialité des suites de Jacobi--Zariski donne le diagramme commutatif
$$
\xymatrix{
&
\mathfrak m_A/\mathfrak m_A^2 \otimes_k K \ar[d] \ar[r] &
H_1(\NL_{K/A}) \ar[d] \ar[r] &
H_1(\NL_{K/k}) \ar[d] \ar[r] & 0 \\
0 \ar[r] &
\mathfrak m_A/\mathfrak m_A^2 \otimes_k K \ar[r] &
\mathfrak m_B/\mathfrak m_B^2 \ar[r] &
\mathfrak m_{\overline{B}}/\mathfrak m_{\overline{B}}^2 \ar[r] & 0
}
$$
La flèche verticale de gauche est injective d'après le théorème
\ref{theorem-regular-fs}, puisque $k \to \overline{B}$ est formellement
lisse pour la topologie $\mathfrak m_{\overline{B}}$-adique par le lemme
\ref{lemma-base-change-fs}. Le lemme du serpent achève alors la démonstration.
\end{proof}

\begin{proposition}
\label{proposition-fs-flat-fibre-fs}
Soit $A \to B$ un homomorphisme local d'anneaux locaux noethériens.
Soient $k$ le corps résiduel de $A$ et
$\overline{B} = B \otimes_A k$ la fibre spéciale. Les conditions suivantes
sont équivalentes :
\begin{enumerate}
\item $A \to B$ est plat et $\overline{B}$ est géométriquement régulier
sur $k$ ;
\item $A \to B$ est plat et $k \to \overline{B}$ est formellement lisse
pour la topologie $\mathfrak m_{\overline{B}}$-adique ;
\item $A \to B$ est formellement lisse pour la topologie
$\mathfrak m_B$-adique.
\end{enumerate}
\end{proposition}

\begin{proof}
L'équivalence de (1) et (2) résulte du théorème
\ref{theorem-regular-fs}.

\medskip\noindent
Supposons (3). Le lemme \ref{lemma-formally-smooth-flat} montre que
$A \to B$ est plat. Le lemme \ref{lemma-base-change-fs} montre que
$k \to \overline{B}$ est formellement lisse pour la topologie
$\mathfrak m_{\overline{B}}$-adique. Ainsi (2) est satisfaite.

\medskip\noindent
Supposons (2). Le lemme \ref{lemma-formally-smooth-completion} assure que la
lissité formelle est préservée par complétion. Il en va de même de la
platitude, d'après Algèbre, lemme
\ref{algebra-lemma-completion-faithfully-flat}. On peut donc remplacer $A$
et $B$ par leurs complétés respectifs et supposer que $A$ et $B$ sont des
anneaux locaux noethériens complets. Choisissons alors un diagramme
$$
\xymatrix{
S \ar[r] & B \\
R \ar[u] \ar[r] & A \ar[u]
}
$$
comme dans le lemme \ref{lemma-embed-map-Noetherian-complete-local-rings}.
Nous utiliserons sans autre mention toutes les propriétés de ce diagramme.
Fixons un système régulier de paramètres $t_1, \ldots, t_d$ de $R$, avec
$t_1 = p$ lorsque $k$ est de caractéristique $p > 0$. Posons
$\overline{S} = S \otimes_R k$ et considérons la suite exacte courte
$$
0 \to J \to S \to B \to 0
$$
Comme $\overline{B}$ et $\overline{S}$ sont réguliers, le noyau de
$\overline{S} \to \overline{B}$ est engendré par des éléments
$\overline{x}_1, \ldots, \overline{x}_r$ faisant partie d'un système
régulier de paramètres de $\overline{S}$, voir Algèbre, lemme
\ref{algebra-lemma-regular-quotient-regular}. Relevons-les en des éléments
$x_1, \ldots, x_r \in J$. Alors
$t_1, \ldots, t_d, x_1, \ldots, x_r$ fait partie d'un système régulier de
paramètres de $S$. Par conséquent, $S/(x_1, \ldots, x_r)$ est un anneau de
séries formelles sur un corps (si $k$ est de caractéristique nulle) ou sur
un anneau de Cohen (si $k$ est de caractéristique $p > 0$), voir le lemme
\ref{lemma-quotient-power-series-ring-over-Cohen}. De plus,
$R \to S/(x_1, \ldots, x_r)$ envoie encore $t_1, \ldots, t_d$ sur une
partie d'un système régulier de paramètres de
$S/(x_1, \ldots, x_r)$. Autrement dit, on peut remplacer $S$ par
$S/(x_1, \ldots, x_r)$ et supposer que l'on dispose d'un diagramme
$$
\xymatrix{
S \ar[r] & B \\
R \ar[u] \ar[r] & A \ar[u]
}
$$
comme dans le lemme \ref{lemma-embed-map-Noetherian-complete-local-rings},
avec en outre $\overline{S} = \overline{B}$. Dans ce cas, l'application
$$
S \otimes_R A \longrightarrow B
$$
est un isomorphisme : elle est surjective, induit un isomorphisme sur les
fibres spéciales, et sa source comme son but sont plats sur $A$. On peut par
exemple appliquer Algèbre, lemme \ref{algebra-lemma-mod-injective}, ou
tensoriser la suite exacte courte
$0 \to I \to S \otimes_R A \to B \to 0$ sur $A$ par $k$ ; on obtient
$I \otimes_A k = 0$, donc $I = 0$ par le lemme de Nakayama. Le lemme
\ref{lemma-base-change-fs} ramène ainsi la démonstration à celle de la
lissité formelle de $R \to S$ pour la topologie $\mathfrak m_S$-adique.
Puisque $\overline{S} = \overline{B}$, l'algèbre $\overline{S}$ est bien
formellement lisse sur $k = R/\mathfrak m_R$.

\medskip\noindent
Choisissons des éléments $y_1, \ldots, y_m \in S$ tels que
$t_1, \ldots, t_d, y_1, \ldots, y_m$ soit un système régulier de paramètres
de $S$. Si $k$ est de caractéristique nulle, choisissons un corps de
coefficients $K \subset S$ ; si $k$ est de caractéristique $p > 0$,
choisissons un anneau de Cohen $\Lambda \subset S$ de corps résiduel $K$.
L'homomorphisme $K[[t_1, \ldots, t_d, y_1, \ldots, y_m]] \to S$ dans le
premier cas, ou
$\Lambda[[t_2, \ldots, t_d, y_1, \ldots, y_m]] \to S$ dans le cas de
caractéristique $p > 0$ est alors un isomorphisme, voir le lemme
\ref{lemma-quotient-power-series-ring-over-Cohen}. Désormais, nous
identifierons $S$ à cet anneau de séries formelles.

\medskip\noindent
La suite de la démonstration est analogue à l'argument employé dans la
démonstration du théorème \ref{theorem-regular-fs}. Choisissons un diagramme
dont les flèches pleines sont données
$$
\xymatrix{
S \ar[r]_{\bar\psi} \ar@{-->}[rd] & N/J \\
R \ar[u]^i \ar[r]^\varphi & N \ar[u]_\pi
}
$$
comme dans la définition \ref{definition-formally-smooth}. Comme $J^2 = 0$,
$J$ possède une structure canonique de $N/J$-module et, par $\bar\psi$, une
structure de $S$-module. La continuité de $\bar\psi$ pour la topologie
$\mathfrak m_S$-adique donne $\mathfrak m_S^nJ = 0$ pour un certain $n$.
On peut donc filtrer $J$ par des sous-$N/J$-modules
$0 \subset J_1 \subset J_2 \subset \ldots \subset J_n = J$
de sorte que chaque quotient $J_{t + 1}/J_t$ soit annulé par
$\mathfrak m_S$. En considérant la suite d'homomorphismes d'anneaux
$N \to N/J_1 \to N/J_2 \to \ldots \to N/J$
on voit qu'il suffit d'établir l'existence de la flèche en pointillés lorsque
$J$ est annulé par $\mathfrak m_S$, c'est-à-dire lorsque $J$ est un espace
vectoriel sur $K$.

\medskip\noindent
Supposons donné un tel diagramme où $J$ est annulé par $\mathfrak m_S$.
Comme $\mathbf{Q} \to S$ dans le cas de caractéristique nulle, ou
$\mathbf{Z} \to S$ dans le cas de caractéristique $p > 0$, est formellement
lisse pour la topologie $\mathfrak m_S$-adique (voir le lemme
\ref{lemma-power-series-ring-over-Cohen-fs}), il existe un homomorphisme
d'anneaux $\psi : S \to N$ tel que $\pi \circ \psi = \bar \psi$.
Puisque $S$ est un anneau de séries formelles en
$t_1, \ldots, t_d$ dans le premier cas, ou en $t_2, \ldots, t_d$ dans le
cas de caractéristique $p > 0$, sur un sous-anneau, la propriété universelle
des anneaux de séries formelles permet de modifier le choix de $\psi$ de
manière que $\psi(t_i)$ soit égal à $\varphi(t_i)$ (ce qui est automatique
pour $t_1 = p$ en
caractéristique $p$). Les deux homomorphismes
$\psi \circ i$ et $\varphi : R \to N$ ont alors même composée avec $\pi$ et
coïncident sur $t_1, \ldots, t_d$. Ainsi
$D = \psi \circ i - \varphi : R \to J$ est une dérivation qui annule
$t_1, \ldots, t_d$. D'après Algèbre, lemme
\ref{algebra-lemma-universal-omega}, on peut écrire
$D = \xi \circ \text{d}$ pour une application $R$-linéaire
$\xi : \Omega_{R/\mathbf{Z}} \to J$ qui annule
$\text{d}t_1, \ldots, \text{d}t_d$ par construction, ainsi que
$\mathfrak m_R \Omega_{R/\mathbf{Z}}$ puisque $J$ est annulé par
$\mathfrak m_R$. Par conséquent, $\xi$ se factorise en
$$
\Omega_{R/\mathbf{Z}} \to \Omega_{k/\mathbf{Z}} \xrightarrow{\xi'} J
$$
où $\xi'$ est $k$-linéaire. À l'aide de la structure d'espace vectoriel sur
$K$ de $J$, prolongeons $\xi'$ en une application $K$-linéaire
$$
\xi'' : \Omega_{k/\mathbf{Z}} \otimes_k K \longrightarrow J.
$$
La lissité formelle de $\overline{S}/k$ montre que
$$
\Omega_{k/\mathbf{Z}} \otimes_k K \to
\Omega_{\overline{S}/\mathbf{Z}} \otimes_S K
$$
est injective d'après le théorème \ref{theorem-regular-fs}. Cela vaut aussi
en caractéristique nulle, car
$\Omega_{k/\mathbf{Z}} \to \Omega_{K/\mathbf{Z}}$ y est même injective ;
voir Algèbre, proposition
\ref{algebra-proposition-characterize-separable-field-extensions}. Il existe
donc une application $K$-linéaire
$\xi''' : \Omega_{\overline{S}/\mathbf{Z}} \otimes_S K \to J$ dont la
restriction à $\Omega_{k/\mathbf{Z}} \otimes_k K$ est $\xi''$. Posons
$$
D' : S \xrightarrow{\text{d}} \Omega_{S/\mathbf{Z}}
\to \Omega_{\overline{S}/\mathbf{Z}} \to
\Omega_{\overline{S}/\mathbf{Z}} \otimes_S K \xrightarrow{\xi'''} J.
$$
Enfin, posons $\psi' = \psi - D' : S \to N$. Le lecteur vérifiera que
$\psi'$ est un homomorphisme d'anneaux tel que
$\pi \circ \psi' = \bar \psi$ et $\psi' \circ i = \varphi$, comme souhaité.
\end{proof}

\noindent
Comme application du résultat précédent, nous montrons que les déformations
d'algèbres formellement lisses sont non obstruées.

\begin{lemma}
\label{lemma-lift-fs}
Soit $A$ un anneau local noethérien complet de corps résiduel $k$.
Soit $B$ une $k$-algèbre locale noethérienne complète. Supposons
$k \to B$ formellement lisse pour la topologie $\mathfrak m_B$-adique.
Il existe alors un anneau local noethérien complet $C$ et un homomorphisme
local $A \to C$, formellement lisse pour la topologie
$\mathfrak m_C$-adique, tels que $C \otimes_A k \cong B$.
\end{lemma}

\begin{proof}
Choisissons un diagramme
$$
\xymatrix{
S \ar[r] & B \\
R \ar[u] \ar[r] & A \ar[u]
}
$$
comme dans le lemme \ref{lemma-embed-map-Noetherian-complete-local-rings}.
Soit $t_1, \ldots, t_d$ un système régulier de paramètres de $R$, avec
$t_1 = p$ lorsque $k$ est de caractéristique $p > 0$. Comme $B$ et
$\overline{S} = S \otimes_R k$ sont réguliers,
$\Ker(\overline{S} \to B)$ est engendré par des éléments
$\overline{x}_1, \ldots, \overline{x}_r$ faisant partie d'un système
régulier de paramètres de $\overline{S}$ ; voir Algèbre, lemme
\ref{algebra-lemma-regular-quotient-regular}. Relevons-les en
$x_1, \ldots, x_r \in S$. Alors
$t_1, \ldots, t_d, x_1, \ldots, x_r$ fait partie d'un système régulier de
paramètres de $S$. Par conséquent, $S/(x_1, \ldots, x_r)$ est un anneau de
séries formelles sur un corps lorsque $k$ est de caractéristique nulle, ou
sur un anneau de Cohen lorsque $k$ est de caractéristique $p > 0$ ; voir le
lemme \ref{lemma-quotient-power-series-ring-over-Cohen}. De plus,
$R \to S/(x_1, \ldots, x_r)$ envoie encore $t_1, \ldots, t_d$ sur une
partie d'un système régulier de paramètres de
$S/(x_1, \ldots, x_r)$. Autrement dit, on peut remplacer $S$ par
$S/(x_1, \ldots, x_r)$ et supposer que l'on dispose d'un diagramme
$$
\xymatrix{
S \ar[r] & B \\
R \ar[u] \ar[r] & A \ar[u]
}
$$
comme dans le lemme \ref{lemma-embed-map-Noetherian-complete-local-rings},
avec en outre $\overline{S} = B$. Dans ce cas, la proposition
\ref{proposition-fs-flat-fibre-fs} montre que $R \to S$ est formellement
lisse pour la topologie $\mathfrak m_S$-adique. Par changement de base, le
lemme \ref{lemma-base-change-fs} montre alors que
$C = S \otimes_R A$ est formellement lisse sur $A$ pour la topologie
$\mathfrak m_C$-adique.
\end{proof}

\begin{remark}
\label{remark-what-does-it-mean}
L'assertion du lemme \ref{lemma-lift-fs} est assez forte. Supposons en effet
que l'on dispose d'un diagramme
$$
\xymatrix{
& B \\
A \ar[r] & A' \ar[u]
}
$$
d'homomorphismes locaux d'anneaux locaux noethériens complets, où
$A \to A'$ induit un isomorphisme des corps résiduels
$k = A/\mathfrak m_A = A'/\mathfrak m_{A'}$ et où
$B \otimes_{A'} k$ est formellement lisse sur $k$. On peut alors le
compléter en un diagramme commutatif
$$
\xymatrix{
C \ar[r] & B \\
A \ar[r] \ar[u] & A' \ar[u]
}
$$
d'homomorphismes locaux d'anneaux locaux noethériens complets, dans lequel
$A \to C$ est formellement lisse pour la topologie $\mathfrak m_C$-adique
et $C \otimes_A k \cong B \otimes_{A'} k$. Il suffit de choisir
$A \to C$ comme dans le lemme \ref{lemma-lift-fs}, relevant
$B \otimes_{A'} k$ sur $k$. La lissité formelle fournit alors la flèche
$C \to B$ ; voir le lemme \ref{lemma-lift-continuous}.
Notons $C \otimes_A^\wedge A'$ le complété de $C \otimes_A A'$ suivant
l'idéal $C \otimes_A \mathfrak m_{A'}$. L'anneau
$C \otimes_A^\wedge A'$ est local noethérien complet (voir Algèbre, lemme
\ref{algebra-lemma-completion-Noetherian}) et plat sur $A'$ (voir Algèbre,
lemme \ref{algebra-lemma-flat-module-powers}). De plus :
\begin{enumerate}
\item $C \otimes_A^\wedge A' \to B$ est surjectif ;
\item si $A \to A'$ est surjectif, alors $C \to B$ est surjectif ;
\item si $A \to A'$ est fini, alors $C \to B$ est fini ;
\item si $A' \to B$ est plat, alors $C \otimes_A^\wedge A' \cong B$.
\end{enumerate}
En effet, le lemme de Nakayama pour les idéaux nilpotents (voir Algèbre,
lemme \ref{algebra-lemma-NAK}) montre que
$C \otimes_A k \cong B \otimes_{A'} k$ implique la surjectivité de
$C \otimes_A A'/\mathfrak m_{A'}^n \to B/\mathfrak m_{A'}^nB$ pour tout
$n$. Cela prouve (1). Les assertions (2) et (3) découlent de (1), tandis que
(4) résulte d'Algèbre, lemme \ref{algebra-lemma-mod-injective}.
\end{remark}




\section{Homomorphismes réguliers d'anneaux}
\label{section-regular}

\noindent
Soit $k$ un corps. Rappelons qu'une $k$-algèbre noethérienne $A$ est dite
{\it géométriquement régulière} sur $k$ si et seulement si
$A \otimes_k k'$ est régulier pour toute extension finie purement
inséparable $k'$ de $k$ ; voir Algèbre, définition
\ref{algebra-definition-geometrically-regular}. Lorsque c'est le cas,
$A \otimes_k k'$ est en outre régulier pour toute extension de corps de type
fini $k'/k$ ; voir Algèbre, lemme
\ref{algebra-lemma-geometrically-regular}. Nous utiliserons cette notion dans
la définition suivante.

\begin{definition}
\label{definition-regular}
Un homomorphisme d'anneaux $R \to \Lambda$ est {\it régulier} s'il est plat
et si, pour tout idéal premier $\mathfrak p \subset R$, l'anneau de la fibre
$$
\Lambda \otimes_R \kappa(\mathfrak p) =
\Lambda_\mathfrak p/\mathfrak p\Lambda_\mathfrak p
$$
est noethérien et géométriquement régulier sur $\kappa(\mathfrak p)$.
\end{definition}

\noindent
Si $R \to \Lambda$ est un homomorphisme d'anneaux avec $\Lambda$
noethérien, les anneaux des fibres sont toujours noethériens.

\begin{lemma}[La régularité est une propriété locale]
\label{lemma-regular-local}
Soit $R \to \Lambda$ un homomorphisme d'anneaux avec $\Lambda$ noethérien.
Les conditions suivantes sont équivalentes :
\begin{enumerate}
\item $R \to \Lambda$ est régulier ;
\item $R_\mathfrak p \to \Lambda_\mathfrak q$ est régulier pour tout
$\mathfrak q \subset \Lambda$ au-dessus de $\mathfrak p \subset R$ ;
\item $R_\mathfrak m \to \Lambda_{\mathfrak m'}$ est régulier pour tout
idéal maximal $\mathfrak m' \subset \Lambda$ au-dessus de $\mathfrak m$
dans $R$.
\end{enumerate}
\end{lemma}

\begin{proof}
Cela résulte du fait qu'un anneau noethérien est régulier si et seulement si
tous ses anneaux locaux sont des anneaux locaux réguliers, voir Algèbre,
définition \ref{algebra-definition-regular}, et qu'un homomorphisme d'anneaux
est plat si et seulement si tous les homomorphismes induits entre anneaux
locaux sont plats, voir Algèbre, lemme
\ref{algebra-lemma-flat-localization}.
\end{proof}

\begin{lemma}[Homomorphismes réguliers et changement de base]
\label{lemma-regular-base-change}
Soit $R \to \Lambda$ un homomorphisme régulier d'anneaux. Pour tout
homomorphisme d'anneaux de type fini $R \to R'$, le changement de base
$R' \to \Lambda \otimes_R R'$ est encore régulier.
\end{lemma}

\begin{proof}
La platitude est préservée par tout changement de base, voir Algèbre, lemme
\ref{algebra-lemma-flat-base-change}. Soit
$\mathfrak p' \subset R'$ un idéal premier au-dessus de
$\mathfrak p \subset R$. L'extension de corps résiduels
$\kappa(\mathfrak p')/\kappa(\mathfrak p)$ est de type fini puisque $R'$ est
de type fini sur $R$. Par conséquent, l'anneau de la fibre
$$
(\Lambda \otimes_R R') \otimes_{R'} \kappa(\mathfrak p') =
\Lambda \otimes_R \kappa(\mathfrak p) \otimes_{\kappa(\mathfrak p)} 
\kappa(\mathfrak p')
$$
est noethérien d'après Algèbre, lemme
\ref{algebra-lemma-Noetherian-field-extension}, et l'hypothèse sur les
anneaux des fibres de $R \to \Lambda$. La régularité géométrique des fibres
est préservée par Algèbre, lemme
\ref{algebra-lemma-geometrically-regular}.
\end{proof}

\begin{lemma}[Composition d'homomorphismes réguliers]
\label{lemma-regular-composition}
Soient $A \to B$ et $B \to C$ des homomorphismes réguliers d'anneaux.
Si les anneaux des fibres de $A \to C$ sont noethériens, alors
$A \to C$ est régulier.
\end{lemma}

\begin{proof}
Soit $\mathfrak p \subset A$ un idéal premier, et soit
$\kappa(\mathfrak p) \subset k$ une extension finie purement inséparable.
Il faut montrer que $C \otimes_A k$ est régulier. Le lemme
\ref{lemma-regular-base-change} permet de supposer $A = k$ ; il suffit alors
de prouver que $C$ est régulier. Par hypothèse, $B$ est régulier et
$B \to C$ est plat à fibres régulières. Algèbre, lemme
\ref{algebra-lemma-flat-over-regular-with-regular-fibre}, montre que $C$ est
régulier. Nous omettons quelques détails.
\end{proof}

\begin{lemma}
\label{lemma-colimit-smooth-regular}
Soit $R$ un anneau et soit $(A_i, \varphi_{ii'})$ un système inductif de
$R$-algèbres lisses. Posons $\Lambda = \colim A_i$. Si les anneaux des
fibres $\Lambda \otimes_R \kappa(\mathfrak p)$ sont noethériens pour tout
$\mathfrak p \subset R$, alors $R \to \Lambda$ est régulier.
\end{lemma}

\begin{proof}
L'algèbre $\Lambda$ est plate sur $R$ d'après Algèbre, lemmes
\ref{algebra-lemma-colimit-flat} et \ref{algebra-lemma-smooth-syntomic}.
Soit $\kappa(\mathfrak p) \subset k$ une extension finie purement
inséparable. On a
$$
\Lambda \otimes_R \kappa(\mathfrak p) \otimes_{\kappa(\mathfrak p)} k =
\Lambda \otimes_R k = \colim A_i \otimes_R k
$$
qui est une colimite de $k$-algèbres lisses, voir Algèbre, lemme
\ref{algebra-lemma-base-change-smooth}. Chaque anneau local d'une
$k$-algèbre lisse est régulier par Algèbre, lemme
\ref{algebra-lemma-characterize-smooth-over-field}. Algèbre, lemme
\ref{algebra-lemma-colimit-regular}, permet donc de conclure que tous les
anneaux locaux de $\Lambda \otimes_R k$ sont réguliers. Le lemme est démontré.
\end{proof}

\noindent
Déterminons quand une extension de corps définit un homomorphisme régulier.

\begin{lemma}
\label{lemma-regular-field-extension}
Soit $K/k$ une extension de corps. Alors $k \to K$ est un homomorphisme
régulier d'anneaux si et seulement si $K$ est une extension séparable de $k$.
\end{lemma}

\begin{proof}
Si $k \to K$ est régulier, $K$ est géométriquement réduit sur $k$, donc
$K$ est séparable sur $k$ d'après Algèbre, proposition
\ref{algebra-proposition-characterize-separable-field-extensions}.
Réciproquement, si $K/k$ est séparable, $K$ est une colimite de
$k$-algèbres lisses, voir Algèbre, lemme
\ref{algebra-lemma-colimit-syntomic}, et il est donc régulier par le lemme
\ref{lemma-colimit-smooth-regular}.
\end{proof}

\begin{lemma}
\label{lemma-regular-permanence}
Soient $A \to B \to C$ des homomorphismes d'anneaux. Si $A \to C$ est
régulier et si $B \to C$ est plat et surjectif sur les spectres, alors
$A \to B$ est régulier.
\end{lemma}

\begin{proof}
Algèbre, lemme \ref{algebra-lemma-flat-permanence}, montre que $A \to B$
est plat. Soit $\mathfrak p \subset A$ un idéal premier. L'homomorphisme
$B \otimes_A \kappa(\mathfrak p) \to C \otimes_A \kappa(\mathfrak p)$ est
plat et surjectif sur les spectres. Ainsi
$B \otimes_A \kappa(\mathfrak p)$ est géométriquement régulier d'après
Algèbre, lemme \ref{algebra-lemma-geometrically-regular-descent}.
\end{proof}




\section{Montée des propriétés par les homomorphismes réguliers}
\label{section-ascending-properties}

\noindent
Cette section est l'analogue d'Algèbre, section
\ref{algebra-section-ascending-properties}, lorsque l'homomorphisme
d'anneaux $R \to S$ est régulier.

\begin{lemma}
\label{lemma-reduced-goes-up}
Soit $\varphi : R \to S$ un homomorphisme d'anneaux. Supposons que
\begin{enumerate}
\item $\varphi$ est régulier ;
\item $S$ est noethérien ;
\item $R$ est noethérien et réduit.
\end{enumerate}
Alors $S$ est réduit.
\end{lemma}

\begin{proof}
Pour un anneau noethérien, être réduit équivaut à vérifier $(S_1)$ et
$(R_0)$ ; voir Algèbre, lemme \ref{algebra-lemma-criterion-reduced}.
On peut donc appliquer Algèbre, lemmes
\ref{algebra-lemma-Sk-goes-up} et \ref{algebra-lemma-Rk-goes-up}.
\end{proof}

\begin{lemma}
\label{lemma-normal-goes-up}
Soit $\varphi : R \to S$ un homomorphisme d'anneaux. Supposons que
\begin{enumerate}
\item $\varphi$ est régulier ;
\item $S$ est noethérien ;
\item $R$ est noethérien et normal.
\end{enumerate}
Alors $S$ est normal.
\end{lemma}

\begin{proof}
Pour un anneau noethérien, être normal équivaut à vérifier $(S_2)$ et
$(R_1)$ ; voir Algèbre, lemme \ref{algebra-lemma-criterion-normal}.
On peut donc appliquer Algèbre, lemmes
\ref{algebra-lemma-Sk-goes-up} et \ref{algebra-lemma-Rk-goes-up}.
\end{proof}

\begin{lemma}
\label{lemma-regular-goes-up}
Soit $\varphi : R \to S$ un homomorphisme d'anneaux. Supposons que
\begin{enumerate}
\item $\varphi$ est régulier ;
\item $S$ est noethérien ;
\item $R$ est noethérien et régulier.
\end{enumerate}
Alors $S$ est régulier.
\end{lemma}

\begin{proof}
Pour un anneau noethérien, être régulier équivaut à vérifier $(R_k)$ pour
tout $k$. On peut donc appliquer Algèbre, lemme
\ref{algebra-lemma-Rk-goes-up}.
\end{proof}

\begin{lemma}
\label{lemma-CM-goes-up}
Soit $\varphi : R \to S$ un homomorphisme d'anneaux. Supposons que
\begin{enumerate}
\item $\varphi$ est régulier ;
\item $S$ est noethérien ;
\item $R$ est noethérien et de Cohen--Macaulay.
\end{enumerate}
Alors $S$ est de Cohen--Macaulay.
\end{lemma}

\begin{proof}
Pour un anneau noethérien, être de Cohen--Macaulay équivaut à vérifier
$(S_k)$ pour tout $k$. On peut donc appliquer Algèbre, lemme
\ref{algebra-lemma-Sk-goes-up}.
\end{proof}




\section{Permanence des propriétés par complétion}
\label{section-permanence-completion}

\noindent
Étant donné un anneau local noethérien $(A, \mathfrak m)$, on note
$A^\wedge$ le complété de $A$ suivant $\mathfrak m$. Nous utiliserons sans
autre mention le fait que $A^\wedge$ est un anneau local noethérien complet
d'idéal maximal $\mathfrak m^\wedge = \mathfrak m A^\wedge$, et que
$A \to A^\wedge$ est fidèlement plat. Voir Algèbre, lemmes
\ref{algebra-lemma-completion-Noetherian-Noetherian},
\ref{algebra-lemma-completion-complete}, et
\ref{algebra-lemma-completion-faithfully-flat}.

\begin{lemma}
\label{lemma-completion-dimension}
Soit $A$ un anneau local noethérien. Alors
$\dim(A) = \dim(A^\wedge)$.
\end{lemma}

\begin{proof}
D'après Algèbre, lemme \ref{algebra-lemma-completion-complete},
$A \to A^\wedge$ induit des isomorphismes
$A/\mathfrak m^n = A^\wedge/(\mathfrak m^\wedge)^n$ pour $n \geq 1$.
Algèbre, lemme \ref{algebra-lemma-pushdown-module}, donne donc
$$
\text{length}_A(A/\mathfrak m^n) =
\text{length}_{A^\wedge}(A^\wedge/(\mathfrak m^\wedge)^n)
$$
pour tout $n \geq 1$. Ainsi $d(A) = d(A^\wedge)$, et la conclusion résulte
d'Algèbre, proposition \ref{algebra-proposition-dimension}. On peut aussi
utiliser Algèbre, lemme
\ref{algebra-lemma-dimension-base-fibre-equals-total}.
\end{proof}

\begin{lemma}
\label{lemma-completion-depth}
Soit $A$ un anneau local noethérien. Alors
$\text{depth}(A) = \text{depth}(A^\wedge)$.
\end{lemma}

\begin{proof}
Voir Algèbre, lemme \ref{algebra-lemma-apply-grothendieck}.
\end{proof}

\begin{lemma}
\label{lemma-completion-CM}
Soit $A$ un anneau local noethérien. Alors $A$ est de Cohen--Macaulay si et
seulement si $A^\wedge$ l'est.
\end{lemma}

\begin{proof}
Un anneau local $A$ est de Cohen--Macaulay si et seulement si
$\dim(A) = \text{depth}(A)$. Ces deux invariants étant préservés par
complétion (lemmes \ref{lemma-completion-dimension} et
\ref{lemma-completion-depth}), l'assertion en découle.
\end{proof}

\begin{lemma}
\label{lemma-completion-regular}
Soit $A$ un anneau local noethérien. Alors $A$ est régulier si et seulement
si $A^\wedge$ l'est.
\end{lemma}

\begin{proof}
Si $A^\wedge$ est régulier, Algèbre, lemme
\ref{algebra-lemma-flat-under-regular}, montre que $A$ est régulier.
Supposons $A$ régulier et notons $\mathfrak m$ l'idéal maximal de $A$. Alors
$\dim_{\kappa(\mathfrak m)} \mathfrak m/\mathfrak m^2 =
\dim(A) = \dim(A^\wedge)$, d'après le lemme
\ref{lemma-completion-dimension}. D'autre part, $\mathfrak mA^\wedge$ est
l'idéal maximal de $A^\wedge$ ; ainsi $\mathfrak m_{A^\wedge}$ est engendré
par au plus $\dim(A^\wedge)$ éléments. Donc $A^\wedge$ est régulier.
On peut aussi appliquer Algèbre, lemme
\ref{algebra-lemma-flat-over-regular-with-regular-fibre}.
\end{proof}

\begin{lemma}
\label{lemma-completion-dvr}
Soit $A$ un anneau local noethérien. Alors $A$ est un anneau de valuation
discrète si et seulement si $A^\wedge$ l'est.
\end{lemma}

\begin{proof}
Cela résulte des lemmes \ref{lemma-completion-dimension} et
\ref{lemma-completion-regular}, ainsi que d'Algèbre, lemme
\ref{algebra-lemma-characterize-dvr}.
\end{proof}

\begin{lemma}
\label{lemma-completion-reduced}
Soit $A$ un anneau local noethérien.
\begin{enumerate}
\item Si $A^\wedge$ est réduit, alors $A$ l'est.
\item En général, la réduction de $A$ n'entraîne pas celle de $A^\wedge$.
\item Si $A$ est de Nagata, alors $A$ est réduit si et seulement si
$A^\wedge$ est réduit.
\end{enumerate}
\end{lemma}

\begin{proof}
Comme $A \to A^\wedge$ est fidèlement plat, (1) résulte d'Algèbre, lemme
\ref{algebra-lemma-descent-reduced}. Pour (2), voir Algèbre, exemple
\ref{algebra-example-bad-dvr-char-p} ; il existe aussi des exemples en
caractéristique nulle, voir Algèbre, remarque
\ref{algebra-remark-resolution-dim-1}. Pour (3), voir Algèbre, lemmes
\ref{algebra-lemma-local-nagata-domain-analytically-unramified} et
\ref{algebra-lemma-analytically-unramified-easy}.
\end{proof}

\begin{lemma}
\label{lemma-completion-normal}
Soit $A$ un anneau local noethérien. Si $A^\wedge$ est normal, alors $A$
l'est.
\end{lemma}

\begin{proof}
Comme $A \to A^\wedge$ est fidèlement plat, cela résulte d'Algèbre, lemme
\ref{algebra-lemma-descent-normal}.
\end{proof}

\begin{lemma}
\label{lemma-flat-completion}
Soit $A \to B$ un homomorphisme local d'anneaux locaux noethériens.
L'homomorphisme induit entre les complétés $A^\wedge \to B^\wedge$ est plat
si et seulement si $A \to B$ est plat.
\end{lemma}

\begin{proof}
Considérons le diagramme commutatif
$$
\xymatrix{
A^\wedge \ar[r] & B^\wedge \\
A \ar[r] \ar[u] & B \ar[u]
}
$$
Les flèches verticales sont fidèlement plates. Supposons
$A^\wedge \to B^\wedge$ plat. Alors $A \to B^\wedge$ est plat, et $B$ est
donc plat sur $A$ d'après Algèbre, lemme
\ref{algebra-lemma-flatness-descends-more-general}.

\medskip\noindent
Supposons $A \to B$ plat. Alors $A \to B^\wedge$ est plat, de sorte que
$B^\wedge/\mathfrak m_A^n B^\wedge$ est plat sur
$A/\mathfrak m_A^n$ pour tout $n \geq 1$. Remarquons que
$\mathfrak m_A^n A^\wedge$ est la puissance $n$-ième de l'idéal maximal
$\mathfrak m_A^\wedge$ de $A^\wedge$, et que
$A/\mathfrak m_A^n = A^\wedge/(\mathfrak m_A^\wedge)^n$. Nous voyons donc
que $B^\wedge$ est plat sur $A^\wedge$ en appliquant Algèbre, lemme
\ref{algebra-lemma-flat-module-powers}, avec
$R = A^\wedge$, $I = \mathfrak m_A^\wedge$, $S = B^\wedge$ et $M = S$.
\end{proof}

\begin{lemma}
\label{lemma-flat-unramified}
Soit $A \to B$ un homomorphisme local plat d'anneaux locaux noethériens tel
que $\mathfrak m_A B = \mathfrak m_B$ et
$\kappa(\mathfrak m_A) = \kappa(\mathfrak m_B)$. Alors $A \to B$ induit un
isomorphisme $A^\wedge \to B^\wedge$ entre les complétés.
\end{lemma}

\begin{proof}
D'après Algèbre, lemme \ref{algebra-lemma-finite-after-completion},
$B^\wedge$ est le complété $\mathfrak m_A$-adique de $B$ et
$A^\wedge \to B^\wedge$ est fini. La platitude de $A \to B$ donne
$\text{Tor}_1^A(B, \kappa(\mathfrak m_A)) = 0$. Le lemme
\ref{lemma-flat-after-completion} montre alors que $B^\wedge$ est plat sur
$A^\wedge$. Ainsi $B^\wedge$ est libre comme $A^\wedge$-module par Algèbre,
lemme
\ref{algebra-lemma-finite-flat-local}. Puisque
$A^\wedge \to B^\wedge$ induit l'isomorphisme
$\kappa(\mathfrak m_A) = A^\wedge/\mathfrak m_A A^\wedge \to
B^\wedge/\mathfrak m_A B^\wedge = B^\wedge/\mathfrak m_B B^\wedge =
\kappa(\mathfrak m_B)$ par hypothèse et par Algèbre, lemme
\ref{algebra-lemma-hathat-finitely-generated}, on voit que $B^\wedge$ est
libre de rang $1$. Ainsi $A^\wedge \to B^\wedge$ est un isomorphisme.
\end{proof}





\section{Permanence des propriétés par les homomorphismes étales}
\label{section-permanence-etale}

\noindent
Dans cette section, nous considérons un homomorphisme étale d'anneaux
$\varphi : A \to B$. Nous étudions les propriétés de $A$ dont hérite $B$,
ainsi que celles de l'anneau local de $B$ en $\mathfrak q$ dont hérite
l'anneau local de $A$ en
$\mathfrak p = \varphi^{-1}(\mathfrak q)$. Cette section rassemble et passe
essentiellement en revue des résultats antérieurs ; elle n'ajoute aucun
résultat nouveau.

\medskip\noindent
Nous utiliserons sans autre mention qu'un homomorphisme étale d'anneaux est
plat (Algèbre, lemme \ref{algebra-lemma-etale}) et qu'un homomorphisme local
plat d'anneaux locaux est fidèlement plat (Algèbre, lemme
\ref{algebra-lemma-local-flat-ff}).

\begin{lemma}
\label{lemma-Noetherian-etale-extension}
Si $A \to B$ est un homomorphisme étale d'anneaux et si $\mathfrak q$ est un
idéal premier de $B$ au-dessus de $\mathfrak p \subset A$, alors
$A_{\mathfrak p}$ est noethérien si et seulement si $B_{\mathfrak q}$ l'est.
\end{lemma}

\begin{proof}
Comme $A_\mathfrak p \to B_\mathfrak q$ est fidèlement plat, le caractère
noethérien de $B_\mathfrak q$ entraîne celui de $A_\mathfrak p$ ; voir
Algèbre, lemme \ref{algebra-lemma-descent-Noetherian}. Réciproquement, si
$A_\mathfrak p$ est noethérien, alors $B_\mathfrak q$ l'est, car il est un
localisé d'une $A_\mathfrak p$-algèbre de type fini.
\end{proof}

\begin{lemma}
\label{lemma-dimension-etale-extension}
Si $A \to B$ est un homomorphisme étale d'anneaux et si $\mathfrak q$ est un
idéal premier de $B$ au-dessus de $\mathfrak p \subset A$, alors
$\dim(A_{\mathfrak p}) = \dim(B_{\mathfrak q})$.
\end{lemma}

\begin{proof}
Comme $A_{\mathfrak p} \to B_{\mathfrak q}$ est plat, la propriété de
descente des idéaux premiers donne
$\dim(A_{\mathfrak p}) \leq \dim(B_{\mathfrak q})$ ; voir Algèbre, lemme
\ref{algebra-lemma-dimension-going-up}. D'autre part, supposons que
$\mathfrak q_0 \subset \mathfrak q_1 \subset \ldots \subset \mathfrak q_n$
soit une chaîne d'idéaux premiers de $B_{\mathfrak q}$. La suite
correspondante d'idéaux premiers
$\mathfrak p_0 \subset \mathfrak p_1 \subset \ldots \subset \mathfrak p_n$
où $\mathfrak p_i = \mathfrak q_i \cap A_{\mathfrak p}$ est elle aussi une
chaîne stricte. En effet, un homomorphisme étale est quasi-fini (voir
Algèbre, lemme \ref{algebra-lemma-etale-quasi-finite}) et un homomorphisme
quasi-fini induit par définition une application de spectres à fibres
discrètes. On obtient donc
$\dim(A_{\mathfrak p}) \geq \dim(B_{\mathfrak q})$, comme voulu.
\end{proof}

\begin{lemma}
\label{lemma-regular-etale-extension}
Si $A \to B$ est un homomorphisme étale d'anneaux et si $\mathfrak q$ est un
idéal premier de $B$ au-dessus de $\mathfrak p \subset A$, alors
$A_{\mathfrak p}$ est régulier si et seulement si $B_{\mathfrak q}$ l'est.
\end{lemma}

\begin{proof}
Pour démontrer l'équivalence, le lemme
\ref{lemma-Noetherian-etale-extension} permet de supposer
$A_\mathfrak p$ et $B_\mathfrak q$ noethériens. Soit
$x_1, \ldots, x_t \in \mathfrak pA_\mathfrak p$ un système minimal de
générateurs. Comme $A_\mathfrak p \to B_\mathfrak q$ est fidèlement plat,
leurs images $y_1, \ldots, y_t$ dans $B_\mathfrak q$ forment un système
minimal de générateurs de
$\mathfrak pB_\mathfrak q = \mathfrak q B_\mathfrak q$ (Algèbre, lemme
\ref{algebra-lemma-etale-at-prime}). Par définition, la régularité de
$A_\mathfrak p$ signifie $t = \dim(A_\mathfrak p)$, et de même pour
$B_\mathfrak q$. Le lemme résulte donc de l'égalité
$\dim(A_\mathfrak p) = \dim(B_\mathfrak q)$ du lemme
\ref{lemma-dimension-etale-extension}.
\end{proof}

\begin{lemma}
\label{lemma-Dedekind-etale-extension}
Si $A \to B$ est un homomorphisme étale d'anneaux et si $A$ est un anneau de
Dedekind, alors $B$ est un produit fini d'anneaux de Dedekind. En particulier,
les localisés $B_\mathfrak q$, pour $\mathfrak q \subset B$ maximal, sont des
anneaux de valuation discrète.
\end{lemma}

\begin{proof}
L'assertion sur les anneaux locaux résulte des lemmes
\ref{lemma-dimension-etale-extension} et
\ref{lemma-regular-etale-extension}, ainsi que d'Algèbre, lemme
\ref{algebra-lemma-characterize-dvr}. L'anneau $B$ est donc noethérien,
normal et de dimension $1$. Algèbre, lemme
\ref{algebra-lemma-characterize-reduced-ring-normal}, montre que $B$ est un
produit fini d'anneaux intègres normaux de dimension $1$. Ce sont des anneaux
de Dedekind par Algèbre, lemme
\ref{algebra-lemma-characterize-Dedekind}.
\end{proof}



\section{Permanence des propriétés par hensélisation}
\label{section-permanence-henselization}

\noindent
Étant donné un anneau local $R$, on note $R^h$, respectivement $R^{sh}$,
son hensélisé, respectivement le hensélisé strict de $R$ ; voir Algèbre,
définition
\ref{algebra-definition-henselization}. Comme nous le montrerons dans cette
section, de nombreuses propriétés de $R$ se reflètent dans $R^h$ et $R^{sh}$.

\begin{lemma}
\label{lemma-dumb-properties-henselization}
Soit $(R, \mathfrak m, \kappa)$ un anneau local. Alors :
\begin{enumerate}
\item $R \to R^h \to R^{sh}$ sont des homomorphismes d'anneaux fidèlement
plats ;
\item $\mathfrak m R^h = \mathfrak m^h$ et
$\mathfrak m R^{sh} = \mathfrak m^h R^{sh} = \mathfrak m^{sh}$,
\item $R/\mathfrak m^n = R^h/\mathfrak m^nR^h$ pour tout $n$ ;
\item il existe des éléments $x_i \in R^{sh}$ tels que
$R^{sh}/\mathfrak m^nR^{sh}$ soit le $R/\mathfrak m^n$-module libre de base
$x_i \bmod \mathfrak m^nR^{sh}$.
\end{enumerate}
\end{lemma}

\begin{proof}
Par construction, $R^h$ est une colimite de $R$-algèbres étales, voir
Algèbre, lemme \ref{algebra-lemma-henselization}. Les homomorphismes étales
étant plats (Algèbre, lemme \ref{algebra-lemma-etale}), Algèbre, lemme
\ref{algebra-lemma-colimit-flat}, montre que $R^h$ est plat sur $R$.
Un homomorphisme local plat d'anneaux étant fidèlement plat (Algèbre, lemme
\ref{algebra-lemma-local-flat-ff}), $R \to R^h$ est fidèlement plat.
L'homomorphisme $R^h \to R^{sh}$ est une colimite d'homomorphismes étales
finis ; voir la démonstration d'Algèbre, lemme
\ref{algebra-lemma-strict-henselization}. Le même argument montre donc que
$R^h \to R^{sh}$ est fidèlement plat.

\medskip\noindent
L'assertion (2) résulte d'Algèbre, lemmes
\ref{algebra-lemma-henselization} et
\ref{algebra-lemma-strict-henselization}. L'assertion (3) découle d'Algèbre,
lemme \ref{algebra-lemma-local-artinian-basis-when-flat}, car
$R/\mathfrak m \to R^h/\mathfrak mR^h$ est un isomorphisme et
$R/\mathfrak m^n \to R^h/\mathfrak m^nR^h$ est plat, comme changement de
base de l'homomorphisme plat $R \to R^h$ (Algèbre, lemme
\ref{algebra-lemma-flat-base-change}). Soit $\kappa^{sep}$ le corps résiduel
de $R^{sh}$ ; c'est une clôture algébrique séparable de $\kappa$.
Choisissons des $x_i \in R^{sh}$ dont les images forment une base de
$\kappa^{sep}$ comme espace vectoriel sur $\kappa$. L'assertion (4) résulte
alors exactement de la même manière d'Algèbre, lemme
\ref{algebra-lemma-local-artinian-basis-when-flat}.
\end{proof}

\begin{lemma}
\label{lemma-henselization-formally-smooth}
Soit $(R, \mathfrak m, \kappa)$ un anneau local. Alors :
\begin{enumerate}
\item $R \to R^h$, $R^h \to R^{sh}$ et $R \to R^{sh}$ sont formellement
étales ;
\item $R \to R^h$, $R^h \to R^{sh}$, respectivement $R \to R^{sh}$, sont
formellement lisses pour la topologie $\mathfrak m^h$-adique,
$\mathfrak m^{sh}$-adique, respectivement $\mathfrak m^{sh}$-adique.
\end{enumerate}
\end{lemma}

\begin{proof}
L'assertion (1) résulte des faits suivants : $R^h$ et $R^{sh}$ sont, par
construction, des colimites filtrantes d'algèbres étales ; les algèbres
étales sont formellement étales (Algèbre, lemme
\ref{algebra-lemma-formally-etale-etale}) ; enfin, une colimite d'algèbres
formellement étales est formellement étale (Algèbre, lemme
\ref{algebra-lemma-colimit-formally-etale}). L'assertion (2) résulte du fait
qu'un homomorphisme formellement étale est formellement lisse, et du lemme
\ref{lemma-formally-smooth}.
\end{proof}

\begin{lemma}
\label{lemma-henselization-noetherian}
\begin{reference}
\cite[IV, Theorem 18.6.6 and Proposition 18.8.8]{EGA}
\end{reference}
Soit $R$ un anneau local. Les conditions suivantes sont équivalentes :
\begin{enumerate}
\item $R$ est noethérien ;
\item $R^h$ est noethérien ;
\item $R^{sh}$ est noethérien.
\end{enumerate}
Dans ce cas :
\begin{enumerate}
\item[(a)] $(R^h)^\wedge$ et $(R^{sh})^\wedge$ sont des anneaux locaux
noethériens complets ;
\item[(b)] $R^\wedge \to (R^h)^\wedge$ est un isomorphisme ;
\item[(c)] $R^h \to (R^h)^\wedge$ et
$R^{sh} \to (R^{sh})^\wedge$ sont plats ;
\item[(d)] $R^\wedge \to (R^{sh})^\wedge$ est formellement lisse pour la
topologie $\mathfrak m_{(R^{sh})^\wedge}$-adique ;
\item[(e)] $(R^\wedge)^{sh} = R^\wedge \otimes_{R^h} R^{sh}$ ;
\item[(f)] $((R^\wedge)^{sh})^\wedge = (R^{sh})^\wedge$.
\end{enumerate}
\end{lemma}

\begin{proof}
Comme $R \to R^h \to R^{sh}$ sont fidèlement plats (lemme
\ref{lemma-dumb-properties-henselization}), le caractère noethérien de
$R^h$ ou de $R^{sh}$ entraîne celui de $R$ ; voir Algèbre, lemme
\ref{algebra-lemma-descent-Noetherian}. Dans la suite de la démonstration,
nous supposons $R$ noethérien.

\medskip\noindent
L'idéal $\mathfrak m \subset R$ étant de type fini, il en va de même de
$\mathfrak m^h = \mathfrak m R^h$ et
$\mathfrak m^{sh} = \mathfrak mR^{sh}$ ; voir le lemme
\ref{lemma-dumb-properties-henselization}. Ainsi $(R^h)^\wedge$ et
$(R^{sh})^\wedge$ sont noethériens par Algèbre, lemme
\ref{algebra-lemma-complete-local-ring-Noetherian}. Cela prouve (a).

\medskip\noindent
L'assertion (b) résulte immédiatement du lemme
\ref{lemma-dumb-properties-henselization}. En particulier,
$(R^h)^\wedge$ est plat sur $R$ ; voir Algèbre, lemme
\ref{algebra-lemma-completion-faithfully-flat}.

\medskip\noindent
Montrons ensuite que $R^h \to (R^h)^\wedge$ est plat. Écrivons
$R^h = \colim_i R_i$ comme colimite filtrante de localisés de
$R$-algèbres étales. D'après Algèbre, lemme
\ref{algebra-lemma-colimit-rings-flat}, si $(R^h)^\wedge$ est plat sur chaque
$R_i$, alors $R^h \to (R^h)^\wedge$ est plat. Or $R^h = R_i^h$ par
construction ; l'assertion (b)
donne donc $R_i^\wedge = (R^h)^\wedge$, qui est plat sur $R_i$.
Pour achever la preuve de (c), montrons que
$R^{sh} \to (R^{sh})^\wedge$ est plat. Par un argument de limite analogue,
il suffit de montrer que $(R^{sh})^\wedge$ est plat sur $R$. Le lemme
\ref{lemma-dumb-properties-henselization} montre que $(R^{sh})^\wedge$ est
le complété d'un $R$-module libre. Il est plat sur $R$ d'après le lemme
\ref{lemma-completed-direct-sum-flat}. Cela achève la preuve de (c).

\medskip\noindent
Nous savons maintenant que (c) est vraie et que $(R^h)^\wedge$ et
$(R^{sh})^\wedge$ sont noethériens. Algèbre, lemme
\ref{algebra-lemma-descent-Noetherian}, montre alors que $R^h$ et $R^{sh}$
sont noethériens.

\medskip\noindent
L'assertion (d) résulte des lemmes
\ref{lemma-henselization-formally-smooth} et
\ref{lemma-formally-smooth-completion}.

\medskip\noindent
L'assertion (e) résulte d'Algèbre, lemme
\ref{algebra-lemma-sh-from-h-map}, et du fait que $R^\wedge$ est hensélien
par Algèbre, lemme \ref{algebra-lemma-complete-henselian}.

\medskip\noindent
Démonstration de (f). L'assertion (e) fournit un homomorphisme
$R^{sh} \to (R^\wedge)^{sh}$, donc après complétion un homomorphisme
$(R^{sh})^\wedge \to ((R^\wedge)^{sh})^\wedge$. Elle fournit aussi un
homomorphisme $R^\wedge \to (R^{sh})^\wedge$. Comme $(R^{sh})^\wedge$ est
strictement hensélien, ainsi qu'on l'a vu, cet homomorphisme induit
$(R^\wedge)^{sh} \to (R^{sh})^\wedge$ par Algèbre, lemme
\ref{algebra-lemma-strictly-henselian-functorial}. En complétant, on obtient
$((R^\wedge)^{sh})^\wedge \to (R^{sh})^\wedge$. Nous omettons la
vérification que ces deux homomorphismes sont inverses l'un de l'autre.
\end{proof}

\begin{lemma}
\label{lemma-henselization-reduced}
\begin{slogan}
La réduction passe à l'hensélisé (strict).
\end{slogan}
Soit $R$ un anneau local. Les conditions suivantes sont équivalentes :
$R$ est réduit, le hensélisé $R^h$ de $R$ est réduit, et le hensélisé strict
$R^{sh}$ de $R$ est réduit.
\end{lemma}

\begin{proof}
Les homomorphismes $R \to R^h \to R^{sh}$ sont fidèlement plats. Un sens
des implications résulte donc d'Algèbre, lemme
\ref{algebra-lemma-descent-reduced}. Réciproquement, supposons $R$ réduit.
Comme $R^h$ et $R^{sh}$ sont des colimites filtrantes de $R$-algèbres
étales, donc lisses, le résultat découle d'Algèbre, lemme
\ref{algebra-lemma-reduced-goes-up}.
\end{proof}

\begin{lemma}
\label{lemma-henselization-nil}
Soit $R$ un anneau local. Notons $nil(R)$ l'idéal des éléments nilpotents de
$R$. Alors $nil(R)R^h = nil(R^h)$ et
$nil(R)R^{sh} = nil(R^{sh})$.
\end{lemma}

\begin{proof}
L'idéal $nil(R)$ est le plus grand idéal formé d'éléments nilpotents pour
lequel $R/nil(R)$ est réduit. D'après Algèbre, lemme
\ref{algebra-lemma-locally-nilpotent}, $nil(R)R^h$ est formé d'éléments
nilpotents. De plus, Algèbre, lemme
\ref{algebra-lemma-quotient-henselization}, identifie
$R^h/nil(R) R^h$ à l'hensélisé de $R/nil(R)$. Le lemme
\ref{lemma-henselization-reduced} montre donc que $R^h/nil(R)R^h$ est
réduit. On en conclut $nil(R) R^h = nil(R^h)$. L'argument est identique pour
l'hensélisé strict, en utilisant Algèbre, lemme
\ref{algebra-lemma-quotient-strict-henselization}.
\end{proof}

\begin{lemma}
\label{lemma-henselization-normal}
Soit $R$ un anneau local. Les conditions suivantes sont équivalentes :
$R$ est un anneau intègre normal, le hensélisé $R^h$ de $R$ est un anneau
intègre normal, et le hensélisé strict $R^{sh}$ de $R$ est un anneau intègre
normal.
\end{lemma}

\begin{proof}
Remarquons d'abord qu'un anneau local est normal si et seulement s'il est
intègre normal ; voir Algèbre, définition
\ref{algebra-definition-ring-normal}. Les homomorphismes
$R \to R^h \to R^{sh}$ sont fidèlement plats ; un sens des implications
résulte donc d'Algèbre, lemme \ref{algebra-lemma-descent-normal}.
Réciproquement, supposons $R$ normal. Comme $R^h$ et $R^{sh}$ sont des
colimites filtrantes de $R$-algèbres étales, donc lisses, le résultat découle
d'Algèbre, lemmes \ref{algebra-lemma-normal-goes-up} et
\ref{algebra-lemma-colimit-normal-ring}.
\end{proof}

\begin{lemma}
\label{lemma-henselization-dimension}
Pour tout anneau local $R$, on a
$\dim(R) = \dim(R^h) = \dim(R^{sh})$.
\end{lemma}

\begin{proof}
Comme $R \to R^{sh}$ est fidèlement plat (lemme
\ref{lemma-dumb-properties-henselization}), la descente des idéaux premiers
donne $\dim(R^{sh}) \geq \dim(R)$ ; voir Algèbre, lemme
\ref{algebra-lemma-dimension-going-up}. Pour l'inégalité réciproque, écrivons
$R^{sh} = \colim R_i$ comme colimite filtrante d'anneaux locaux $R_i$,
chacun étant un localisé d'une $R$-algèbre étale. Si
$\mathfrak q_0 \subset \mathfrak q_1 \subset \ldots \subset \mathfrak q_n$
est une chaîne d'idéaux premiers de $R^{sh}$, alors, pour $i$ assez grand,
la suite
$$
R_i \cap \mathfrak q_0 \subset
R_i \cap \mathfrak q_1 \subset \ldots \subset
R_i \cap \mathfrak q_n
$$
est une chaîne d'idéaux premiers de $R_i$. Ainsi
$\dim(R^{sh}) \leq \sup_i \dim(R_i)$. Or le lemme
\ref{lemma-dimension-etale-extension} donne
$\dim(R_i) = \dim(R)$ pour chaque $i$, d'où le résultat.
\end{proof}

\begin{lemma}
\label{lemma-henselization-depth}
Pour tout anneau local noethérien $R$, on a
$\text{depth}(R) = \text{depth}(R^h) = \text{depth}(R^{sh})$.
\end{lemma}

\begin{proof}
Le lemme \ref{lemma-henselization-noetherian} montre que $R^h$ et $R^{sh}$
sont noethériens. L'assertion résulte alors d'Algèbre, lemme
\ref{algebra-lemma-apply-grothendieck}.
\end{proof}

\begin{lemma}
\label{lemma-henselization-CM}
Soit $R$ un anneau local noethérien. Les conditions suivantes sont
équivalentes : $R$ est de Cohen--Macaulay, le hensélisé $R^h$ de $R$ est de
Cohen--Macaulay, et le hensélisé strict $R^{sh}$ de $R$ est de
Cohen--Macaulay.
\end{lemma}

\begin{proof}
Le lemme \ref{lemma-henselization-noetherian} montre que $R^h$ et $R^{sh}$
sont noethériens, de sorte que l'énoncé a bien un sens. Comme
$\text{depth}(R) = \text{depth}(R^h) = \text{depth}(R^{sh})$
et
$\dim(R) = \dim(R^h) = \dim(R^{sh})$
d'après les lemmes \ref{lemma-henselization-depth} et
\ref{lemma-henselization-dimension}, la conclusion suit.
\end{proof}

\begin{lemma}
\label{lemma-henselization-regular}
Soit $R$ un anneau local noethérien. Les conditions suivantes sont
équivalentes : $R$ est un anneau local régulier, le hensélisé $R^h$ de $R$
est un anneau local régulier, et le hensélisé strict $R^{sh}$ de $R$ est un
anneau local régulier.
\end{lemma}

\begin{proof}
Le lemme \ref{lemma-henselization-noetherian} montre que $R^h$ et $R^{sh}$
sont noethériens, de sorte que l'énoncé a bien un sens. Soit $\mathfrak m$
l'idéal maximal de $R$, et soit
$x_1, \ldots, x_t \in \mathfrak m$ un système minimal de générateurs de
$\mathfrak m$, c'est-à-dire tel que leurs images dans
$\mathfrak m/\mathfrak m^2$ forment une base sur
$\kappa = R/\mathfrak m$. Comme $R \to R^h$ et $R \to R^{sh}$ sont
fidèlement plats, les images $x_1^h, \ldots, x_t^h$ dans $R^h$,
respectivement $x_1^{sh}, \ldots, x_t^{sh}$ dans $R^{sh}$, forment des
systèmes minimaux de générateurs de
$\mathfrak m^h = \mathfrak mR^h$, respectivement de
$\mathfrak m^{sh} = \mathfrak mR^{sh}$. Par définition, la régularité de
$R$ signifie $t = \dim(R)$, et de même pour $R^h$ et $R^{sh}$. Le résultat
découle donc de l'égalité
$\dim(R) = \dim(R^h) = \dim(R^{sh})$ du lemme
\ref{lemma-henselization-dimension}.
\end{proof}

\begin{lemma}
\label{lemma-henselization-dvr}
Soit $R$ un anneau local noethérien. Alors $R$ est un anneau de valuation
discrète si et seulement si $R^h$ en est un, si et seulement si $R^{sh}$ en
est un.
\end{lemma}

\begin{proof}
Cela résulte des lemmes \ref{lemma-henselization-dimension} et
\ref{lemma-henselization-regular}, ainsi que d'Algèbre, lemme
\ref{algebra-lemma-characterize-dvr}.
\end{proof}

\begin{lemma}
\label{lemma-filtered-colimit-etale-noetherian-fibres}
Soit $A$ un anneau, et soit $B$ une colimite filtrante de $A$-algèbres
étales. Soit $\mathfrak p$ un idéal premier de $A$. Si $B$ est noethérien,
il n'existe qu'un nombre fini d'idéaux premiers
$\mathfrak q_1, \ldots, \mathfrak q_r$ au-dessus de $\mathfrak p$, et l'on a
$B \otimes_A \kappa(\mathfrak p) = \prod \kappa(\mathfrak q_i)$, et
chacune des extensions de corps
$\kappa(\mathfrak q_i)/\kappa(\mathfrak p)$ étant algébrique séparable.
\end{lemma}

\begin{proof}
Écrivons $B$ comme colimite filtrante $B = \colim B_i$, avec
$A \to B_i$ étale. D'une part,
$B \otimes_A \kappa(\mathfrak p) = \colim B_i \otimes_A \kappa(\mathfrak p)$
est une colimite filtrante de $\kappa(\mathfrak p)$-algèbres étales ;
d'autre part, cet anneau est noethérien. Une
$\kappa(\mathfrak p)$-algèbre étale est un produit fini d'extensions finies
séparables (Algèbre, lemme \ref{algebra-lemma-etale-over-field}). Il n'y a
donc aucune spécialisation non triviale entre les idéaux premiers, tous
maximaux et minimaux, des algèbres
$B_i \otimes_A \kappa(\mathfrak p)$, ni par conséquent entre ceux de
$B \otimes_A \kappa(\mathfrak p)$. Ainsi
$B \otimes_A \kappa(\mathfrak p)$ est réduit et possède un nombre fini
d'idéaux premiers, tous minimaux. C'est donc un produit fini de corps ; voir
Algèbre, Lemme \ref{algebra-lemma-total-ring-fractions-no-embedded-points}
ou
Algèbre, Proposition \ref{algebra-proposition-dimension-zero-ring}).
Chacun de ces corps est une colimite d'extensions finies séparables, ce qui
donne la dernière assertion.
\end{proof}

\begin{lemma}
\label{lemma-fibres-henselization}
Soit $R$ un anneau local noethérien et soit
$\mathfrak p \subset R$ un idéal premier. Alors
$$
R^h \otimes_R \kappa(\mathfrak p) =
\prod\nolimits_{i = 1, \ldots, t} \kappa(\mathfrak q_i)
\quad\text{resp.}\quad
R^{sh} \otimes_R \kappa(\mathfrak p) =
\prod\nolimits_{i = 1, \ldots, s} \kappa(\mathfrak r_i)
$$
où $\mathfrak q_1, \ldots, \mathfrak q_t$, respectivement
$\mathfrak r_1, \ldots, \mathfrak r_s$, sont les idéaux premiers de $R^h$,
respectivement de $R^{sh}$, au-dessus de $\mathfrak p$. De plus, les
extensions de corps
$\kappa(\mathfrak q_i)/\kappa(\mathfrak p)$
respectivement $\kappa(\mathfrak r_i)/\kappa(\mathfrak p)$, sont
algébriques séparables.
\end{lemma}

\begin{proof}
Cela se déduit du lemme plus général
\ref{lemma-filtered-colimit-etale-noetherian-fibres}, puisque l'hensélisé et
l'hensélisé strict sont noethériens, ainsi qu'on l'a vu. Donnons néanmoins
une démonstration directe.

\medskip\noindent
Nous utiliserons sans autre mention les lemmes
\ref{lemma-dumb-properties-henselization} et
\ref{lemma-henselization-noetherian}. L'anneau
$R^h/\mathfrak pR^h$, respectivement $R^{sh}/\mathfrak pR^{sh}$, est
l'hensélisé, respectivement l'hensélisé strict, de $R/\mathfrak p$ ; voir
Algèbre, lemme \ref{algebra-lemma-quotient-henselization}, respectivement
\ref{algebra-lemma-quotient-strict-henselization}. On peut donc remplacer
$R$ par $R/\mathfrak p$ et supposer que $R$ est un anneau local noethérien
intègre et que $\mathfrak p = (0)$. Comme $R^h$, respectivement $R^{sh}$,
est noethérien, il possède un nombre fini d'idéaux premiers minimaux
$\mathfrak q_1, \ldots, \mathfrak q_t$, respectivement
$\mathfrak r_1, \ldots, \mathfrak r_s$. La platitude de
$R \to R^h$, respectivement de $R \to R^{sh}$, montre par descente que ce
sont exactement les idéaux premiers au-dessus de $\mathfrak p = (0)$.
Enfin, puisque $R$ est intègre, $R^h$, respectivement $R^{sh}$, est réduit
par le lemme \ref{lemma-henselization-reduced}. Ainsi
$R^h \otimes_R \kappa(\mathfrak p)$, respectivement
$R^{sh} \otimes_R \kappa(\mathfrak p)$, est un anneau noethérien réduit
ayant un nombre fini d'idéaux premiers, tous minimaux et donc maximaux. Ces
anneaux sont artiniens et produits de leurs localisés aux idéaux maximaux,
chacun de ces localisés étant nécessairement un corps ; voir Algèbre,
proposition \ref{algebra-proposition-dimension-zero-ring} et lemme
\ref{algebra-lemma-minimal-prime-reduced-ring}.

\medskip\noindent
La dernière assertion résulte du fait que $R \to R^h$, respectivement
$R \to R^{sh}$, est une colimite d'homomorphismes étales ; les extensions
induites des corps résiduels sont donc des colimites d'extensions finies
séparables, voir Algèbre, lemme \ref{algebra-lemma-etale-at-prime}.
\end{proof}









\section{Retour sur les extensions de corps}
\label{section-p-bases}

\noindent
Dans cette section, nous étudions quelques particularités des extensions de
corps en caractéristique $p > 0$.

\begin{definition}
\label{definition-p-basis}
Soit $p$ un nombre premier et soit $k \to K$ une extension de corps de
caractéristique $p$. Notons $kK^p$ le compositum de $k$ et $K^p$ dans $K$.
\begin{enumerate}
\item Une partie $\{x_i\} \subset K$ est dite {\it p-indépendante sur
$k$} si les éléments $x^E = \prod x_i^{e_i}$, où $0 \leq e_i < p$, sont
linéairement indépendants sur $kK^p$.
\item Une partie $\{x_i\}$ de $K$ est appelée une
{\it p-base de $K$ sur $k$} si les éléments $x^E$ forment une base de $K$
sur $kK^p$.
\end{enumerate}
\end{definition}

\noindent
Cette notion est liée à celle de $p$-base d'une $\mathbf{F}_p$-algèbre, que
nous étudierons plus loin (insérer ici un renvoi futur).

\begin{lemma}
\label{lemma-p-basis}
Soit $K/k$ une extension de corps, avec $k$ de caractéristique $p > 0$.
Soit $\{x_i\}$ une partie de $K$. Les conditions suivantes sont
équivalentes :
\begin{enumerate}
\item les éléments $\{x_i\}$ sont $p$-indépendants sur $k$ ;
\item les éléments $\text{d}x_i$ sont linéairement indépendants sur $K$
dans $\Omega_{K/k}$.
\end{enumerate}
Toute famille $p$-indépendante se complète en une $p$-base de $K$ sur $k$.
En particulier, $K$ possède une $p$-base sur $k$. De plus, les conditions
suivantes sont équivalentes :
\begin{enumerate}
\item[(a)] $\{x_i\}$ est une $p$-base de $K$ sur $k$ ;
\item[(b)] les $\text{d}x_i$ forment une base sur $K$ de l'espace vectoriel
$\Omega_{K/k}$.
\end{enumerate}
\end{lemma}

\begin{proof}
Supposons (2), et soit $\sum a_E x^E = 0$ une relation linéaire avec
$a_E \in k K^p$. Soit $\theta_i : K \to K$ une $k$-dérivation telle que
$\theta_i(x_j) = \delta_{ij}$, symbole de Kronecker. Toute $k$-dérivation de
$K$ annule $kK^p$. En appliquant $\theta_i$ à la relation donnée, on obtient
de nouvelles relations
$$
\sum\nolimits_{E, e_i > 0}
e_i a_E x_1^{e_1}\ldots x_i^{e_i - 1} \ldots x_n^{e_n} = 0
$$
Si l'on choisit une relation $\sum a_E x^E$ de degré total minimal
$|E| = \sum e_i$, avec un $a_E \not = 0$, on obtient donc une contradiction.
Ainsi (1) est satisfaite.

\medskip\noindent
Si $\{x_i\}$ est une $p$-base de $K$ sur $k$, alors
$K \cong kK^p[X_i]/(X_i^p - x_i^p)$. Les $\text{d}x_i$ forment donc une
base de $\Omega_{K/k}$ sur $K$. Ainsi (a) implique (b).

\medskip\noindent
Soit $\{x_i\}$ une partie $p$-indépendante de $K$ sur $k$. Le lemme de Zorn
permet de l'agrandir en une partie $p$-indépendante maximale de $K$ sur $k$.
Affirmons que toute partie $p$-indépendante maximale $\{x_i\}$ de $K$ est
une $p$-base de $K$ sur $k$. Cela montrera que (1) implique (2), ainsi que
l'existence des $p$-bases. Soit $L$ le sous-corps de $K$ engendré par
$kK^p$ et les $x_i$. Il faut montrer $L = K$. Si $x \in K$ mais
$x \not \in L$, alors $x^p \in L$ et $L(x) \cong L[z]/(z^p - x^p)$.
La famille $\{x_i\} \cup \{x\}$ serait donc $p$-indépendante sur $k$,
contradiction.

\medskip\noindent
Enfin, montrons que (b) implique (a). L'équivalence de (1) et (2) montre que
$\{x_i\}$ est une partie $p$-indépendante maximale de $K$ sur $k$. D'après
l'affirmation précédente, c'est donc une $p$-base.
\end{proof}

\begin{lemma}
\label{lemma-intersection-subfields-subspace}
Soit $K/k$ une extension de corps. Soit $\{K_\alpha\}_{\alpha \in A}$ une
famille de sous-corps de $K$ satisfaisant les propriétés suivantes :
\begin{enumerate}
\item $k \subset K_\alpha$ pour tout $\alpha \in A$ ;
\item $k = \bigcap_{\alpha \in A} K_\alpha$,
\item pour $\alpha, \alpha' \in A$, il existe $\alpha'' \in A$ tel que
$K_{\alpha''} \subset K_\alpha \cap K_{\alpha'}$.
\end{enumerate}
Alors, pour $n \geq 1$ et pour tout sous-espace vectoriel
$V \subset K^{\oplus n}$ sur $K$, on a
$V \cap k^{\oplus n} \not = 0$ si et seulement si
$V \cap K_\alpha^{\oplus n} \not = 0$ pour tout $\alpha \in A$.
\end{lemma}

\begin{proof}
Raisonnons par récurrence sur $n$. Le cas $n = 1$ résulte des hypothèses.
Supposons le résultat établi pour les sous-espaces de
$K^{\oplus n - 1}$, et supposons que
$V \subset K^{\oplus n}$ rencontre non trivialement
$K_\alpha^{\oplus n}$ pour tout $\alpha \in A$. Si
$V \cap 0 \oplus k^{\oplus n - 1}$ est non nul, la conclusion est acquise.
On peut donc supposer le contraire. Par l'hypothèse de récurrence, il existe
$\alpha$ tel que $V \cap 0 \oplus K_\alpha^{\oplus n - 1}$ soit nul.
Soit $v = (x_1, \ldots, x_n) \in V \cap K_\alpha^{\oplus n}$ non nul.
Par le choix de $\alpha$, on voit que $x_1$ est non nul. Remplaçons $v$ par
$x_1^{-1}v$, de sorte que $v = (1, x_2, \ldots, x_n)$. Si
$v' = (x_1', \ldots, x'_n) \in V \cap K_\alpha^{\oplus n}$, alors
$v' - x_1'v = 0$ par le choix de $\alpha$. Ainsi
$V \cap K_\alpha^{\oplus n} = K_\alpha v$. Si l'on choisit $\alpha'$ tel que
$K_{\alpha'} \subset K_\alpha$, on voit nécessairement que
$v \in V \cap K_{\alpha'}^{\oplus n}$ (par le même argument appliqué à
$\alpha'$). Par conséquent,
$$
x_2, \ldots, x_n \in
\bigcap\nolimits_{\alpha' \in A, K_{\alpha'} \subset K_\alpha} K_{\alpha'}
$$
qui est égal à $k$ d'après (2) et (3).
\end{proof}

\begin{lemma}
\label{lemma-intersection-subfields}
Soit $K$ un corps de caractéristique $p$. Soit
$\{K_\alpha\}_{\alpha \in A}$ une famille de sous-corps de $K$ satisfaisant
les propriétés suivantes :
\begin{enumerate}
\item $K^p \subset K_\alpha$ pour tout $\alpha \in A$ ;
\item $K^p = \bigcap_{\alpha \in A} K_\alpha$,
\item pour $\alpha, \alpha' \in A$, il existe $\alpha'' \in A$ tel que
$K_{\alpha''} \subset K_\alpha \cap K_{\alpha'}$.
\end{enumerate}
Alors
\begin{enumerate}
\item l'intersection des noyaux des applications
$\Omega_{K/\mathbf{F}_p} \to \Omega_{K/K_\alpha}$ est nulle ;
\item pour toute extension finie $L/K$, on a
$L^p = \bigcap_{\alpha \in A} L^pK_\alpha$.
\end{enumerate}
\end{lemma}

\begin{proof}
Démonstration de (1). Choisissons une $p$-base $\{x_i\}$ de $K$ sur
$\mathbf{F}_p$. Supposons que
$\eta = \sum_{i \in I'} y_i \text{d}x_i$ s'envoie sur zéro dans
$\Omega_{K/K_\alpha}$ pour tout $\alpha \in A$, où $I'$ est fini. D'après le
lemme \ref{lemma-p-basis}, cela signifie que, pour tout $\alpha$, il existe
une relation
$$
\sum\nolimits_E a_{E, \alpha} x^E = 0,\quad a_{E, \alpha} \in K_\alpha
$$
où $E$ parcourt les multi-indices $E = (e_i)_{i \in I'}$ tels que
$0 \leq e_i < p$. D'autre part, le lemme \ref{lemma-p-basis} assure qu'il
n'existe aucune relation de ce type $\sum a_E x^E = 0$ avec
$a_E \in K^p$. Le lemme \ref{lemma-intersection-subfields-subspace} donne
une contradiction.

\medskip\noindent
Démonstration de (2). Considérons une tour $L/M/K$ d'extensions finies de
corps. Posons $M_\alpha = M^p K_\alpha$ et
$L_\alpha = L^p K_\alpha = L^p M_\alpha$. On peut d'abord établir
$M^p = \bigcap_{\alpha \in A} M_\alpha$, puis
$L^p = \bigcap_{\alpha \in A} L_\alpha$. Il suffit donc de prouver (2) pour
les extensions simples sans sous-corps intermédiaire non trivial.
Supposons d'abord $L = K(\theta)$ séparable sur $K$. Alors $L$ est engendré
par $\theta^p$ sur $K$, et l'on peut donc supposer $\theta \in L^p$. Dans ce
cas,
$$
L^p = K^p \oplus K^p\theta \oplus \ldots K^p\theta^{d - 1}
\quad\text{et}\quad
L^pK_\alpha =
K_\alpha \oplus K_\alpha \theta \oplus \ldots K_\alpha\theta^{d - 1}
$$
où $d = [L : K]$. La conclusion est alors claire. Dans l'autre cas,
$L = K(\theta)$ avec $\theta^p = t \in K$ et $t \not \in K^p$. On a
$$
L^p = K^p \oplus K^pt \oplus \ldots K^pt^{p - 1}
\quad\text{et}\quad
L^pK_\alpha =
K_\alpha \oplus K_\alpha t \oplus \ldots K_\alpha t^{p - 1}
$$
Là encore, le résultat est clair.
\end{proof}

\begin{lemma}
\label{lemma-power-series-ring-subfields}
Soit $k$ un corps de caractéristique $p > 0$, et soit
$\{x_i\}_{i \in I}$ une $p$-base de $k$. Soient $n, m \geq 0$. Soit $K$ le
corps des fractions de
$A = k[[x_1, \ldots, x_n]][y_1, \ldots, y_m]$. Pour une partie finie
$J$ de $I$, considérons le sous-corps $k/k_J/k^p$ engendré par $k^p$ et les
$x_i$ avec $i \in I \setminus J$. Les corps des fractions $K_J$ des anneaux
$$
A_J = k_J[[x_1^p, \ldots, x_n^p]][y_1^p, \ldots, y_m^p]
$$
forment une famille de sous-corps de $K$ comme dans le lemme
\ref{lemma-intersection-subfields}. De plus, chacune des extensions
d'anneaux $A_J \subset A$ est finie.
\end{lemma}

\begin{proof}
Comme $k/k_J$ est fini, l'extension d'anneaux
$k_J[[x_1^p, \ldots, x_d^p]] \subset k[[x_1, \ldots, x_d]]$ est finie
d'après Algèbre, lemme \ref{algebra-lemma-finite-after-completion}. Il
s'ensuit que $A_J \to A$ est fini.

\medskip\noindent
Vérifions les propriétés (1), (2) et (3) du lemme
\ref{lemma-intersection-subfields}. Démonstration de (1). Pour $a \in A$,
on a $a^p \in A_J$, donc $K^p \subset K_J$. Démonstration de (2).
Supposons que $f/g^p \in K$, avec $f, g \in A$ et $g \not = 0$, appartienne
à $K_J$ pour tout choix de $J$. Fixons provisoirement $J$. Puisque
$f/g^p \in K_J$, on peut écrire $f/g^p = a/b^p$ avec $a \in A_J$ et
$b \in A$ non nul. Ainsi $b^p f \in A_J$. Pour toute $A_J$-dérivation
$D : A \to A$, on a $0 = D(b^pf) = b^p D(f)$, donc $D(f) = 0$ puisque $A$
est intègre. En prenant $D = \partial_{x_i}$ puis
$D = \partial_{y_j}$, on obtient
$f \in k[[x_1^p, \ldots, x_n^p]][y_1^p, \ldots, y_m^p]$. En appliquant une
$k_J$-dérivation $\theta : k \to k$, on voit de même que tous les
coefficients de $f$ appartiennent à $k_J$, c'est-à-dire $f \in A_J$.
Comme il est clair que $A^p = \bigcap\nolimits_J A_J$, où $J$ parcourt les
sous-corps considérés dans le lemme, on obtient $f \in A^p$, comme voulu.
Démonstration de (3). C'est immédiat puisque
$K_{J \cup J'} \subset K_J \cap K_{J'}$.
\end{proof}





\section{Le lieu singulier}
\label{section-singular-locus}

\noindent
Soit $R$ un anneau noethérien. Le {\it lieu régulier} $\text{Reg}(X)$ de
$X = \Spec(R)$ est l'ensemble des idéaux premiers $\mathfrak p$ tels que
$R_\mathfrak p$ soit un anneau local régulier. Le {\it lieu singulier}
$\text{Sing}(X)$ de $X = \Spec(R)$ est son complémentaire
$X \setminus \text{Reg}(X)$, c'est-à-dire l'ensemble des idéaux premiers
$\mathfrak p$ tels que $R_\mathfrak p$ ne soit pas un anneau local régulier.
La discussion précédant Algèbre, définition
\ref{algebra-definition-regular}, montre que $\text{Reg}(X)$ est stable par
générisation. Dans cette section, nous étudions des conditions garantissant
que $\text{Reg}(X)$ est ouvert.

\begin{definition}
\label{definition-J}
\begin{reference}
\cite[(32.B)]{MatCA}
\end{reference}
Soit $R$ un anneau noethérien et soit $X = \Spec(R)$.
\begin{enumerate}
\item On dit que $R$ est {\it J-0} si $\text{Reg}(X)$ contient un ouvert
non vide.
\item On dit que $R$ est {\it J-1} si $\text{Reg}(X)$ est ouvert.
\item On dit que $R$ est {\it J-2} si toute $R$-algèbre de type fini est
J-1.
\end{enumerate}
\end{definition}

\noindent
L'anneau $\mathbf{Q}[x]/(x^2)$ ne vérifie pas J-0, mais il vérifie J-1.
D'autre part, J-1 implique J-0 pour les anneaux intègres noethériens et,
plus généralement, pour les anneaux noethériens réduits non nuls, puisqu'un
tel anneau est régulier en ses idéaux premiers minimaux. Voici une
caractérisation de la propriété J-1.

\begin{lemma}
\label{lemma-J-1}
Soit $R$ un anneau noethérien et soit $X = \Spec(R)$. L'anneau $R$ est J-1
si et seulement si $V(\mathfrak p) \cap \text{Reg}(X)$ contient un ouvert
non vide de $V(\mathfrak p)$ pour tout
$\mathfrak p \in \text{Reg}(X)$.
\end{lemma}

\begin{proof}
Cela résulte de Topologie, lemme
\ref{topology-lemma-characterize-open-Noetherian}, et du fait que
$\text{Reg}(X)$ est stable par générisation d'après Algèbre, lemme
\ref{algebra-lemma-localization-of-regular-local-is-regular}.
\end{proof}

\begin{lemma}
\label{lemma-intersection-regular-with-closed}
Soit $R$ un anneau noethérien et soit $X = \Spec(R)$. Supposons que, pour
tout idéal premier $\mathfrak p \subset R$, l'anneau $R/\mathfrak p$ soit J-0.
Alors $R$ est J-1.
\end{lemma}

\begin{proof}
Montrons que le critère du lemme \ref{lemma-J-1} s'applique. Soit
$\mathfrak p \in \text{Reg}(X)$ un idéal premier de hauteur $r$.
Choisissons $f_1, \ldots, f_r \in \mathfrak p$ dont les images engendrent
$\mathfrak pR_\mathfrak p$. Puisque $\mathfrak p \in \text{Reg}(X)$, les
images de $f_1, \ldots, f_r$ forment une suite régulière dans
$R_\mathfrak p$ ; voir Algèbre, lemme
\ref{algebra-lemma-regular-ring-CM}. D'après Algèbre, lemme
\ref{algebra-lemma-regular-sequence-in-neighbourhood}, après avoir remplacé
$R$ par $R_g$ pour un $g \in R$, $g \not \in \mathfrak p$, la suite
$f_1, \ldots, f_r$ est régulière dans $R$. Après un nouveau remplacement,
on peut aussi supposer que $f_1, \ldots, f_r$ engendrent $\mathfrak p$.
Soit ensuite $\mathfrak p \subset \mathfrak q$ un idéal premier tel que
$(R/\mathfrak p)_\mathfrak q$ soit un anneau local régulier. L'hypothèse du
lemme fournit un ouvert non vide de $V(\mathfrak p)$ formé de tels idéaux
premiers ; il suffit donc de montrer que $R_\mathfrak q$ est régulier.
Or $f_1, \ldots, f_r$ est une suite régulière dans $R_\mathfrak q$ et
$R_\mathfrak q/(f_1, \ldots, f_r)R_\mathfrak q$ est régulier. Algèbre,
lemme \ref{algebra-lemma-regular-mod-x}, montre alors que
$R_\mathfrak q$ est régulier.
\end{proof}

\begin{lemma}
\label{lemma-J-0-goes-down}
Soit $R \to S$ un homomorphisme d'anneaux. Supposons que
\begin{enumerate}
\item $R$ est un anneau intègre noethérien ;
\item $R \to S$ est injectif et de type fini ;
\item $S$ est intègre et J-0.
\end{enumerate}
Alors $R$ est J-0.
\end{lemma}

\begin{proof}
Après avoir remplacé $S$ par $S_g$ pour un certain $g \in S$ non nul, on
peut supposer $S$ régulier. La platitude générique permet aussi de supposer
$R \to S$ fidèlement plat ; voir Algèbre, lemme
\ref{algebra-lemma-generic-flatness-Noetherian}. Algèbre, lemme
\ref{algebra-lemma-descent-regular}, montre alors que $R$ est régulier.
\end{proof}

\begin{lemma}
\label{lemma-J-0-goes-up}
Soit $R \to S$ un homomorphisme d'anneaux. Supposons que
\begin{enumerate}
\item $R$ est un anneau intègre noethérien et J-0 ;
\item $R \to S$ est injectif et de type fini ;
\item $S$ est intègre ;
\item l'extension induite des corps des fractions est séparable.
\end{enumerate}
Alors $S$ est J-0.
\end{lemma}

\begin{proof}
On peut remplacer $R$ par un localisé principal et supposer que $R$ est
régulier.
D'après Algèbre, lemme \ref{algebra-lemma-smooth-at-generic-point},
$R \to S$ est lisse en $(0)$. Après avoir remplacé $S$ par un localisé
principal, on peut donc supposer que $S$ est lisse sur $R$. Alors $S$ est lui
aussi
régulier ; voir Algèbre, lemme \ref{algebra-lemma-regular-goes-up}.
\end{proof}

\begin{lemma}
\label{lemma-J-2}
Soit $R$ un anneau noethérien. Les conditions suivantes sont équivalentes :
\begin{enumerate}
\item $R$ est J-2 ;
\item toute $R$-algèbre de type fini qui est intègre est J-0 ;
\item toute $R$-algèbre finie est J-1 ;
\item pour tout idéal premier $\mathfrak p$ et toute extension finie
purement inséparable $L/\kappa(\mathfrak p)$, il existe une $R$-algèbre
finie $R'$ qui est intègre et J-0, et dont le corps des fractions est $L$.
\end{enumerate}
\end{lemma}

\begin{proof}
Les implications (1) $\Rightarrow$ (2) et (2) $\Rightarrow$ (4) sont
claires. Rappelons qu'un anneau intègre J-1 est J-0. On a donc aussi
(1) $\Rightarrow$ (3) et (3) $\Rightarrow$ (4).

\medskip\noindent
Soit $R \to S$ un homomorphisme d'anneaux de type fini. Pour montrer que
$S$ est J-1, le lemme \ref{lemma-intersection-regular-with-closed} ramène à
prouver que $S/\mathfrak q$ est J-0 pour tout idéal premier $\mathfrak q$ de
$S$. Cela montre (2) $\Rightarrow$ (1).

\medskip\noindent
Supposons (4). Montrons (2), ce qui achèvera la démonstration. Soit
$R \to S$ un homomorphisme d'anneaux de type fini, avec $S$ intègre.
Posons $\mathfrak p = \Ker(R \to S)$ et notons $K$ le corps des fractions
de $S$. Il existe un diagramme de corps
$$
\xymatrix{
K \ar[r] & K' \\
\kappa(\mathfrak p) \ar[u] \ar[r] & L \ar[u]
}
$$
où les flèches horizontales sont des extensions finies purement
inséparables et où $K'/L$ est séparable ; voir Algèbre, lemme
\ref{algebra-lemma-make-separably-generated}. Choisissons
$R' \subset L$ comme dans (4), et soit $S'$ l'image de
$S \otimes_R R' \to K'$. Alors $S'$ est intègre, de corps des fractions
$K'$ ; $S'$ est donc J-0 par le lemme \ref{lemma-J-0-goes-up} et le choix de
$R'$. Le lemme \ref{lemma-J-0-goes-down} montre alors que $S$ est J-0,
comme voulu.
\end{proof}



\section{Régularité et dérivations}
\label{section-regularity-derivations}

\noindent
Soit $R \to S$ un homomorphisme d'anneaux. Soit $D : R \to R$ une
dérivation. Nous disons que $D$ {\it se prolonge à} $S$ s'il existe une
dérivation $D' : S \to S$ telle que le diagramme
$$
\xymatrix{
S \ar[r]_{D'} & S \\
R \ar[u] \ar[r]^D & R \ar[u]
}
$$
soit commutatif.

\begin{lemma}
\label{lemma-derivation-extends}
Soit $R$ un anneau. Soit $D : R \to R$ une dérivation.
\begin{enumerate}
\item Pour tout idéal $I \subset R$, la dérivation $D$ se prolonge
canoniquement en une dérivation $D^\wedge : R^\wedge \to R^\wedge$
du complété $I$-adique.
\item Pour toute partie multiplicative $S \subset R$, la dérivation
$D$ se prolonge de manière unique au localisé $S^{-1}R$ de $R$.
\end{enumerate}
Si $R \subset R'$ est une extension d'anneaux de type fini telle que
$R_g \cong R'_g$ pour un élément $g \in R$ qui est non-diviseur de zéro
dans $R'$, alors $g^ND$ se prolonge à $R'$ pour un certain $N \geq 0$.
\end{lemma}

\begin{proof}
Démonstration de (1). Pour $n \geq 2$, la règle de Leibniz donne
$D(I^n) \subset I^{n - 1}$. Ainsi $D$ induit des applications
$D_n : R/I^n \to R/I^{n - 1}$. Par passage à la limite, nous obtenons
$D^\wedge$. Nous omettons la vérification que $D^\wedge$ est une dérivation.

\medskip\noindent
Démonstration de (2). Pour prolonger $D$ à $S^{-1}R$, il suffit de poser
$D(r/s) = D(r)/s - rD(s)/s^2$ et de vérifier les axiomes.

\medskip\noindent
Démonstration de la dernière assertion. Soient
$x_1, \ldots, x_n \in R'$ des générateurs de $R'$ sur $R$. Choisissons
$N$ tel que $g^Nx_i \in R$. Considérons $g^{N + 1}D$. D'après (2), cette
dérivation se prolonge à $R_g$. De plus, la règle de Leibniz et la
construction précédente du prolongement donnent
$$
g^{N + 1}D(x_i) = g^{N + 1}D(g^{-N} g^Nx_i) = -Ng^Nx_iD(g) +
gD(g^Nx_i)
$$
et les deux termes appartiennent à $R$. Il s'ensuit que
$$
g^{N + 1}D(x_1^{e_1} \ldots x_n^{e_n}) =
\sum e_i x_1^{e_1} \ldots x_i^{e_i - 1} \ldots x_n^{e_n} g^{N + 1}D(x_i)
$$
appartient à $R'$. Par conséquent, tout élément de $R'$ (qui peut s'écrire
comme une somme de monômes en les $x_i$ à coefficients dans $R$) est envoyé
dans $R'$ par $g^{N + 1}D$, ce qui conclut la démonstration.
\end{proof}

\begin{lemma}
\label{lemma-quotient-regular}
\begin{slogan}
Le critère jacobien pour les hypersurfaces, correctement formulé.
\end{slogan}
Soit $R$ un anneau régulier. Soit $f \in R$. Supposons qu'il existe une
dérivation $D : R \to R$ telle que $D(f)$ soit inversible dans $R/(f)$.
Alors $R/(f)$ est régulier.
\end{lemma}

\begin{proof}
Il suffit de le démontrer lorsque $R$ est un anneau local d'idéal maximal
$\mathfrak m$ et de corps résiduel $\kappa$. Dans ce cas, il suffit de
démontrer que $f \not \in \mathfrak m^2$; voir
Algèbre, Lemme \ref{algebra-lemma-regular-ring-CM}.
Or, si $f \in \mathfrak m^2$, alors la règle de Leibniz donne
$D(f) \in \mathfrak m$, ce qui est contradictoire.
\end{proof}

\begin{lemma}
\label{lemma-quotient-sequence-regular}
Soit $(R, \mathfrak m, \kappa)$ un anneau local régulier. Soit $m \geq 1$.
Soient $f_1, \ldots, f_m \in \mathfrak m$. Supposons qu'il existe des
dérivations $D_1, \ldots, D_m : R \to R$ telles que
$\det_{1 \leq i, j \leq m}(D_i(f_j))$ soit inversible dans $R$.
Alors $R/(f_1, \ldots, f_m)$ est régulier et $f_1, \ldots, f_m$
est une suite régulière.
\end{lemma}

\begin{proof}
Il suffit de démontrer que $f_1, \ldots, f_m$ sont linéairement indépendants
sur $\kappa$ dans $\mathfrak m/\mathfrak m^2$; voir
Algèbre, Lemme \ref{algebra-lemma-regular-ring-CM}.
En effet, s'il existe une relation linéaire non triviale, alors
$\sum a_i f_i \in \mathfrak m^2$ pour certains $a_i \in R$, mais les
relations $a_i \in \mathfrak m$ ne sont pas toutes vérifiées. Observons que
$D_i(\mathfrak m^2) \subset \mathfrak m$ et
$D_i(a_j f_j) \equiv a_j D_i(f_j) \bmod \mathfrak m$
par la règle de Leibniz pour les dérivations. On en déduirait donc que
$$
\sum a_j D_i(f_j) \in \mathfrak m
$$
ce qui contredirait l'hypothèse sur le déterminant.
\end{proof}

\begin{lemma}
\label{lemma-degree-p-extension-regular}
Soit $R$ un anneau régulier. Soit $f \in R$.
Supposons qu'il existe une dérivation $D : R \to R$ telle que $D(f)$ soit
inversible dans $R$. Alors $R[z]/(z^n - f)$ est régulier pour tout entier
$n \geq 1$. Plus généralement, $R[z]/(p(z) - f)$ est régulier pour tout
$p \in \mathbf{Z}[z]$.
\end{lemma}

\begin{proof}
D'après Algèbre, Lemme \ref{algebra-lemma-regular-goes-up}, l'anneau
$R[z]$ est régulier. Appliquons le Lemme \ref{lemma-quotient-regular}
au prolongement de $D$ à $R[z]$ qui envoie $z$ sur zéro.
Cela convient parce que $D$ annule tout polynôme à coefficients entiers
et envoie $f$ sur un élément inversible.
\end{proof}

\begin{lemma}
\label{lemma-find-D}
Soit $p$ un nombre premier. Soit $B$ un anneau intègre tel que $p = 0$ dans $B$.
Soit $f \in B$ un élément qui n'est pas une puissance $p$-ième dans le corps
des fractions de $B$. Si $B$ est de type fini sur un anneau local noethérien
complet, alors il existe une dérivation $D : B \to B$ telle que $D(f)$
soit non nul.
\end{lemma}

\begin{proof}
Soit $R$ un anneau local noethérien complet muni d'un homomorphisme d'anneaux
de type fini $R \to B$. Nous pouvons évidemment remplacer $R$ par son image
dans $B$; nous pouvons donc supposer que $R$ est un anneau intègre de
caractéristique $p > 0$ (en plus d'être local noethérien complet). D'après
Algèbre, Lemme
\ref{algebra-lemma-complete-local-Noetherian-domain-finite-over-regular}
$R$ est une extension finie de $k[[x_1, \ldots, x_n]]$ pour un certain
corps $k$ et un certain entier $n$. Nous pouvons donc remplacer $R$ par
$k[[x_1, \ldots, x_n]]$. Ensuite, nous appliquons
Algèbre, Lemme \ref{algebra-lemma-Noether-normalization-over-a-domain}
pour factoriser $R \to B$ sous la forme
$$
R \subset R[y_1, \ldots, y_d] \subset B' \subset B
$$
où $B'$ est fini sur $R[y_1, \ldots, y_d]$ et $B'_g \cong B_g$ pour un
certain $g \in R$ non nul. Remarquons que
$f' = g^{pN} f \in B'$ pour un entier $N$ assez grand. Il est clair que
$f'$ n'est pas une puissance $p$-ième dans le corps des fractions de $B'$.
Si nous trouvons une dérivation $D' : B' \to B'$ telle que
$D'(f') \not = 0$, alors le Lemme \ref{lemma-derivation-extends} garantit
que $D = g^MD'$ se prolonge à $B$ pour un certain $M > 0$. Dans ce cas,
$D(f) = g^MD'(f) = g^MD'(g^{-pN}f') = g^{M - pN}D'(f')$ est non nul.
Il suffit donc de démontrer le lemme lorsque $B$ est une extension finie de
$A = k[[x_1, \ldots, x_n]][y_1, \ldots, y_m]$.

\medskip\noindent
Supposons que $B$ soit une extension finie de
$A = k[[x_1, \ldots, x_n]][y_1, \ldots, y_m]$.
Notons $L$ le corps des fractions de $B$.
Remarquons que $\text{d}f$ n'est pas nul dans
$\Omega_{L/\mathbf{F}_p}$; voir
Algèbre, Lemme \ref{algebra-lemma-derivative-zero-pth-power}.
Appliquons le Lemme \ref{lemma-power-series-ring-subfields} pour trouver un
sous-corps $k' \subset k$ d'indice fini tel que, si
$A' = k'[[x_1^p, \ldots, x_n^p]][y_1^p, \ldots, y_m^p]$
l'élément $\text{d}f$ ne s'envoie pas sur zéro dans $\Omega_{L/K'}$,
où $K'$ est le corps des fractions de $A'$.
Nous pouvons donc choisir une $K'$-dérivation $D' : L \to L$ telle que
$D'(f) \not = 0$. Puisque $A' \subset A$ et $A \subset B$ sont finis par
construction, l'extension $A' \subset B$ est finie. Choisissons
$b_1, \ldots, b_t \in B$ qui engendrent $B$ comme $A'$-module. Alors
$D'(b_i) = f_i/g_i$ pour certains $f_i, g_i \in B$ avec $g_i \not = 0$.
En posant $D = g_1 \ldots g_t D'$, nous obtenons la dérivation voulue.
\end{proof}

\begin{lemma}
\label{lemma-complete-Noetherian-domain-J-0}
Soit $A$ un anneau local noethérien complet et intègre. Alors $A$ est J-0.
\end{lemma}

\begin{proof}
D'après Algèbre, Lemme
\ref{algebra-lemma-complete-local-Noetherian-domain-finite-over-regular}
nous pouvons trouver un sous-anneau régulier $A_0 \subset A$ tel que $A$
soit fini sur $A_0$. L'extension induite $K/K_0$ des corps des fractions est
finie. Si $K/K_0$ est séparable, le Lemme \ref{lemma-J-0-goes-up} permet de
conclure. Sinon, $A_0$ et $A$ sont de caractéristique $p > 0$. Pour tout
corps intermédiaire $K/M/K_0$, il existe une sous-extension finie
$A_0 \subset B \subset A$ dont le corps des fractions est $M$. Par récurrence
sur $[K : K_0]$, nous pouvons donc supposer qu'il existe
$A_0 \subset B \subset A$ tel que $B$ soit J-0 et que $K/M$ n'ait aucun
corps intermédiaire non trivial. Dans ce cas, si $K/M$ est séparable, le
Lemme \ref{lemma-J-0-goes-up} montre que $A$ est J-0. Sinon,
$K = M[z]/(z^p - b_1/b_2)$ pour certains $b_1, b_2 \in B$ avec
$b_2 \not = 0$ et $b_1/b_2$ non puissance $p$-ième dans $M$.
Choisissons $a \in A$ non nul tel que $az \in A$. En remplaçant $z$ par
$b_2 a^p z$, nous obtenons $K = M[z]/(z^p - b)$, où $z \in A$ et où
$b \in B$ n'est pas une puissance $p$-ième dans $M$.
D'après le Lemme \ref{lemma-find-D}, il existe une dérivation
$D : B \to B$ telle que $D(b) \not = 0$. Le Lemme
\ref{lemma-degree-p-extension-regular} montre que $A_\mathfrak p$ est régulier
pour tout idéal premier $\mathfrak p$ de $A$ situé au-dessus d'un idéal
premier régulier de $B$ et ne contenant pas $D(b)$. Puisque $B$ est J-0,
nous concluons qu'il en est de même de $A$.
\end{proof}

\begin{proposition}
\label{proposition-ubiquity-J-2}
Les types d'anneaux suivants sont J-2 :
\begin{enumerate}
\item les corps,
\item les anneaux locaux noethériens complets,
\item $\mathbf{Z}$,
\item les anneaux locaux noethériens de dimension $1$,
\item les anneaux de Nagata de dimension $1$,
\item les anneaux de Dedekind dont le corps des fractions est de
caractéristique zéro,
\item les extensions d'anneaux de type fini de l'un quelconque des anneaux
précédents.
\end{enumerate}
\end{proposition}

\begin{proof}
Dans les cas (1), (3), (5) et (6), cela résulte de la vérification de la
condition (4) du Lemme \ref{lemma-J-2}. Nous ne la donnerons que lorsque
$R$ est un anneau de Nagata de dimension $1$. Soit $\mathfrak p \subset R$
un idéal premier et soit $L/\kappa(\mathfrak p)$ une extension finie
purement inséparable. Si $\mathfrak p \subset R$ est un idéal maximal,
alors $R \to L$ est fini et $L$ est un anneau régulier, ce qui vérifie la
condition. Si $\mathfrak p \subset R$ est un idéal premier minimal, alors
la condition de Nagata assure que la clôture intégrale $R' \subset L$ de
$R$ dans $L$ est finie sur $R$. Ainsi $R'$ est un anneau intègre normal de
dimension $1$ (Algèbre, Lemme \ref{algebra-lemma-integral-dim-up}), donc
régulier (Algèbre, Lemme \ref{algebra-lemma-criterion-normal}); la condition
est également vérifiée dans ce cas.

\medskip\noindent
Dans le cas (2), nous utilisons la condition (3) du Lemme
\ref{lemma-J-2}. Soit $R$ un anneau local noethérien complet.
Remarquons que, si $R \to R'$ est fini, alors $R'$ est un produit d'anneaux
locaux noethériens complets; voir
Algèbre, Lemme \ref{algebra-lemma-quotient-complete-local}.
Le Lemme \ref{lemma-intersection-regular-with-closed} ramène donc la
démonstration au fait qu'un anneau local noethérien complet et intègre est
J-0, qui est précisément le Lemme
\ref{lemma-complete-Noetherian-domain-J-0}.

\medskip\noindent
Dans le cas (4), nous utilisons encore la condition (3) du Lemme
\ref{lemma-J-2}. En effet, si $R$ est un anneau local noethérien de
dimension $1$ et si $R \to R'$ est fini, alors $\Spec(R')$ est fini. Comme
le lieu régulier est stable par générisation, nous voyons que $R'$ est J-1.
\end{proof}





\section{Lissité formelle et régularité}
\label{section-fs-regular}

\noindent
Le titre de cette section renvoie au Théorème \ref{theorem-fs-regular}
d'Andr\'e.

\begin{lemma}
\label{lemma-lift-derivation-through-fs}
Soit $A \to B$ un homomorphisme local d'anneaux locaux noethériens.
Soit $D : A \to A$ une dérivation. Supposons que $B$ soit complet et que
$A \to B$ soit formellement lisse pour la topologie $\mathfrak m_B$-adique.
Alors il existe un prolongement $D' : B \to B$ de $D$.
\end{lemma}

\begin{proof}
Notons $B[\epsilon] = B[x]/(x^2)$ l'anneau des nombres duaux sur $B$.
Considérons l'homomorphisme d'anneaux $\psi : A \to B[\epsilon]$,
$a \mapsto a + \epsilon D(a)$, et le diagramme commutatif
$$
\xymatrix{
B \ar[r]_1 & B \\
A \ar[u] \ar[r]^\psi & B[\epsilon] \ar[u]
}
$$
D'après le Lemme \ref{lemma-lift-continuous} et l'hypothèse de lissité
formelle de $B/A$, il existe un homomorphisme
$\varphi : B \to B[\epsilon]$ qui complète le diagramme. Écrivons
$\varphi(b) = b + \epsilon D'(b)$. Alors $D' : B \to B$ est le prolongement
cherché.
\end{proof}

\begin{proposition}
\label{proposition-fs-regular}
Soit $A \to B$ un homomorphisme local d'anneaux locaux noethériens complets.
Soit $k$ le corps résiduel de $A$ et soit
$\overline{B} = B \otimes_A k$ la fibre spéciale.
Les conditions suivantes sont équivalentes :
\begin{enumerate}
\item $A \to B$ est régulier ;
\item $A \to B$ est plat et $\overline{B}$ est géométriquement régulier
sur $k$ ;
\item $A \to B$ est plat et $k \to \overline{B}$ est formellement lisse
pour la topologie $\mathfrak m_{\overline{B}}$-adique ;
\item $A \to B$ est formellement lisse pour la topologie
$\mathfrak m_B$-adique.
\end{enumerate}
\end{proposition}

\begin{proof}
Nous avons démontré l'équivalence de (2), (3) et (4) dans la
Proposition \ref{proposition-fs-flat-fibre-fs}. Il est clair que (1)
implique (2). Supposons donc vérifiées les conditions équivalentes (2),
(3) et (4), et démontrons (1).

\medskip\noindent
Soit $\mathfrak p$ un idéal premier de $A$. Nous allons montrer que
$B \otimes_A \kappa(\mathfrak p)$ est géométriquement régulier sur
$\kappa(\mathfrak p)$. D'après le Lemme \ref{lemma-base-change-fs}, nous
pouvons remplacer $A$ par $A/\mathfrak p$ et $B$ par $B/\mathfrak pB$.
Nous pouvons donc supposer que $A$ est intègre et que $\mathfrak p = (0)$.

\medskip\noindent
Choisissons $A_0 \subset A$ comme dans Algèbre, Lemme
\ref{algebra-lemma-complete-local-Noetherian-domain-finite-over-regular}.
Nous utiliserons sans autre mention toutes les propriétés énoncées dans ce
lemme. Comme $A_0 \to A$ induit un isomorphisme sur les corps résiduels et
comme $B/\mathfrak m_A B$ est géométriquement régulier sur
$A/\mathfrak m_A$, il existe un diagramme
$$
\xymatrix{
C \ar[r] & B \\
A_0 \ar[r] \ar[u] & A \ar[u]
}
$$
où $A_0 \to C$ est formellement lisse pour la topologie
$\mathfrak m_C$-adique et tel que $B = C \otimes_{A_0} A$; voir la
Remarque \ref{remark-what-does-it-mean}. (Il n'est pas nécessaire de
compléter le produit tensoriel, puisque $A_0 \to A$ est fini; voir
Algèbre, Lemme \ref{algebra-lemma-completion-tensor}.) Il suffit donc de
montrer que $C \otimes_{A_0} K_0$ est une algèbre géométriquement régulière
sur le corps des fractions $K_0$ de $A_0$.

\medskip\noindent
Le paragraphe précédent montre que nous pouvons supposer que
$A = k[[x_1, \ldots, x_n]]$, où $k$ est un corps, ou que
$A = \Lambda[[x_1, \ldots, x_n]]$, où $\Lambda$ est un anneau de Cohen.
Dans ce cas, $B$ est un anneau régulier; voir
Algèbre, Lemme \ref{algebra-lemma-flat-over-regular-with-regular-fibre}.
Ainsi $B \otimes_A K$ est également un anneau régulier (où $K$ est le
corps des fractions de $A$), ce qui permet de conclure si $K$ est de
caractéristique zéro.

\medskip\noindent
Il reste donc le cas où $A = k[[x_1, \ldots, x_n]]$ et où $k$ est un corps
de caractéristique $p > 0$. Soit $L/K$ une extension finie purement
inséparable. Nous allons démontrer par récurrence sur $[L : K]$ que
$B \otimes_A L$ est régulier. Le cas initial $L = K$ a été traité ci-dessus.
Soit $K \subset M \subset L$ un sous-corps tel que $L$ soit une extension
de degré $p$ de $M$, obtenue par adjonction d'une racine $p$-ième d'un
élément $f \in M$. Soit $A'$ une $A$-sous-algèbre finie de $M$ dont le corps
des fractions est $M$. En chassant les dénominateurs, nous pouvons supposer,
et nous supposons, que $f \in A'$. Posons $A'' = A'[z]/(z^p -f)$ et
remarquons que $A' \subset A''$ est fini et que le corps des fractions de
$A''$ est $L$. Par récurrence, l'anneau $B \otimes_A M$ est régulier.
Nous avons
$$
B \otimes_A L = B \otimes_A M[z]/(z^p - f)
$$
Le Lemme \ref{lemma-find-D} fournit une dérivation $D : A' \to A'$ telle
que $D(f) \not = 0$. Puisque $A' \to B \otimes_A A'$ est formellement
lisse pour la topologie $\mathfrak m$-adique d'après le
Lemme \ref{lemma-descent-fs}, le
Lemme \ref{lemma-lift-derivation-through-fs} permet de prolonger $D$ en une
dérivation $D' : B \otimes_A A' \to B \otimes_A A'$. Remarquons que
$D'(f) = D(f)$ est inversible dans $B \otimes_A M$, puisque $D(f)$ est non
nul dans $A' \subset M$. Le Lemme \ref{lemma-degree-p-extension-regular}
montre donc que $B \otimes_A L$ est régulier, ce qui conclut la démonstration.
\end{proof}

\begin{theorem}[Andr\'e]
\label{theorem-fs-regular}
\begin{reference}
L'implication (4) $\Rightarrow$ (1) est le résultat principal de
\cite{Andre-smooth}.
\end{reference}
Soit $A \to B$ un homomorphisme local d'anneaux locaux noethériens.
Soit $k$ le corps résiduel de $A$ et soit
$\overline{B} = B \otimes_A k$ la fibre spéciale. Supposons que
$A \to A^\wedge$ soit régulier\footnote{Par exemple, si $A$ est un
G-anneau ou est (quasi-)excellent; voir les Sections \ref{section-G-ring} et
\ref{section-excellent}.}.
Les conditions suivantes sont équivalentes :
\begin{enumerate}
\item $A \to B$ est régulier ;
\item $A \to B$ est plat et $\overline{B}$ est géométriquement régulier
sur $k$ ;
\item $A \to B$ est plat et $k \to \overline{B}$ est formellement lisse
pour la topologie $\mathfrak m_{\overline{B}}$-adique ;
\item $A \to B$ est formellement lisse pour la topologie
$\mathfrak m_B$-adique.
\end{enumerate}
\end{theorem}

\begin{proof}
Nous avons démontré l'équivalence de (2), (3) et (4) dans la
Proposition \ref{proposition-fs-flat-fibre-fs}. Il est clair que (1)
implique (2). Supposons donc (4) vérifiée et démontrons (1).

\medskip\noindent
Le Lemme \ref{lemma-formally-smooth-completion} montre que
$A^\wedge \to B^\wedge$ est formellement lisse. La
Proposition \ref{proposition-fs-regular} montre alors que
$A^\wedge \to B^\wedge$ est régulier. Par hypothèse, $A \to A^\wedge$ est
régulier; le Lemme \ref{lemma-regular-composition} entraîne donc que
$A \to B^\wedge$ est régulier. Puisque $B \to B^\wedge$ est fidèlement
plat, le Lemme \ref{lemma-regular-permanence} permet de conclure que
$A \to B$ est régulier.
\end{proof}




\section{G-anneaux}
\label{section-G-ring}

\noindent
Soit $A$ un anneau local noethérien. Dans la
Section \ref{section-permanence-completion}, nous avons vu que certaines
propriétés de $A$, mais pas toutes, se reflètent dans le complété
$A^\wedge$ de $A$. Pour approfondir cette question, introduisons un peu de
terminologie. Pour un idéal premier $\mathfrak q$ de $A$, l'anneau fibre
$$
A^\wedge \otimes_A \kappa(\mathfrak q) =
(A^\wedge)_\mathfrak q/\mathfrak q(A^\wedge)_\mathfrak q =
(A/\mathfrak q)^\wedge \otimes_{A/\mathfrak q} \kappa(\mathfrak q)
$$
est appelé une {\it fibre formelle} de $A$. Nous considérons la fibre
formelle comme une algèbre sur $\kappa(\mathfrak q)$. Ainsi,
$A \to A^\wedge$ est un homomorphisme régulier d'anneaux si et seulement
si toutes les fibres formelles sont des algèbres géométriquement régulières.

\begin{definition}
\label{definition-G-ring}
Un anneau $R$ est appelé un {\it G-anneau} si $R$ est noethérien et si,
pour tout idéal premier $\mathfrak p$ de $R$, l'homomorphisme d'anneaux
$R_\mathfrak p \to (R_\mathfrak p)^\wedge$ est régulier.
\end{definition}

\noindent
La discussion précédente montre que $R$ est un G-anneau si et seulement si
chaque anneau local $R_\mathfrak p$ a des fibres formelles géométriquement
régulières. Remarquons que, si $\mathbf{Q} \subset R$, il suffit de vérifier
que les fibres formelles sont régulières. Le lemme suivant donne une autre
formulation de la condition de G-anneau.

\begin{lemma}
\label{lemma-check-G-ring-easy}
Soit $R$ un anneau noethérien. Alors $R$ est un G-anneau si et seulement si,
pour toute paire d'idéaux premiers
$\mathfrak q \subset \mathfrak p \subset R$, l'algèbre
$$
(R/\mathfrak q)_\mathfrak p^\wedge \otimes_{R/\mathfrak q} \kappa(\mathfrak q)
$$
est géométriquement régulière sur $\kappa(\mathfrak q)$.
\end{lemma}

\begin{proof}
Cela résulte de l'égalité
$$
R_\mathfrak p^\wedge \otimes_R \kappa(\mathfrak q) =
(R/\mathfrak q)_\mathfrak p^\wedge \otimes_{R/\mathfrak q} \kappa(\mathfrak q)
$$
entre algèbres sur $\kappa(\mathfrak q)$.
\end{proof}

\begin{lemma}
\label{lemma-G-ring-goes-up-quasi-finite}
Soit $R \to R'$ un homomorphisme de type fini d'anneaux noethériens, et
soient
$$
\xymatrix{
\mathfrak q' \ar[r] & \mathfrak p' \ar[r] & R' \\
\mathfrak q \ar[r] \ar@{-}[u] &
\mathfrak p \ar[r] \ar@{-}[u] & R \ar[u]
}
$$
des idéaux premiers. Supposons que $R \to R'$ soit quasi-fini en
$\mathfrak p'$.
\begin{enumerate}
\item Si la fibre formelle
$R_\mathfrak p^\wedge \otimes_R \kappa(\mathfrak q)$ est géométriquement
régulière sur $\kappa(\mathfrak q)$, alors la fibre formelle
$(R'_{\mathfrak p'})^\wedge \otimes_{R'} \kappa(\mathfrak q')$ est
géométriquement régulière sur $\kappa(\mathfrak q')$.
\item Si les fibres formelles de $R_\mathfrak p$ sont géométriquement
régulières, alors celles de $R'_{\mathfrak p'}$ le sont aussi.
\item Si $R \to R'$ est quasi-fini et si $R$ est un G-anneau, alors $R'$
est un G-anneau.
\end{enumerate}
\end{lemma}

\begin{proof}
Il est clair que (1) $\Rightarrow$ (2) $\Rightarrow$ (3). Supposons que
$R_\mathfrak p^\wedge \otimes_R \kappa(\mathfrak q)$ soit géométriquement
régulier sur $\kappa(\mathfrak q)$. D'après Algèbre, Lemme
\ref{algebra-lemma-completion-at-quasi-finite-prime}, nous avons
$$
R_\mathfrak p^\wedge \otimes_R R'
=
(R'_{\mathfrak p'})^\wedge \times B
$$
pour une certaine $R_\mathfrak p^\wedge$-algèbre $B$. Par conséquent,
$R'_{\mathfrak p'} \to (R'_{\mathfrak p'})^\wedge$ est un facteur d'un
changement de base de l'homomorphisme
$R_\mathfrak p \to R_\mathfrak p^\wedge$. Il s'ensuit que
$(R'_{\mathfrak p'})^\wedge \otimes_{R'} \kappa(\mathfrak q')$ est un
facteur de
$$
R_\mathfrak p^\wedge \otimes_R R' \otimes_{R'} \kappa(\mathfrak q') =
R_\mathfrak p^\wedge \otimes_R \kappa(\mathfrak q)
\otimes_{\kappa(\mathfrak q)} \kappa(\mathfrak q').
$$
Le résultat découle donc du fait que l'extension du corps de base préserve
la régularité géométrique; voir
Algèbre, Lemme \ref{algebra-lemma-geometrically-regular}.
\end{proof}

\begin{lemma}
\label{lemma-check-G-ring}
Soit $R$ un anneau noethérien. Alors $R$ est un G-anneau si et seulement si,
pour tout homomorphisme fini libre d'anneaux $R \to S$, les fibres formelles
de $S$ sont des anneaux réguliers.
\end{lemma}

\begin{proof}
Supposons que, pour tout homomorphisme fini libre d'anneaux $R \to S$,
l'anneau $S$ ait des fibres formelles régulières. Soient
$\mathfrak q \subset \mathfrak p \subset R$ des idéaux premiers et soit
$\kappa(\mathfrak q) \subset L$ une extension finie purement inséparable.
Pour montrer que $R$ est un G-anneau, il suffit de montrer que
$$
R_\mathfrak p^\wedge \otimes_R \kappa(\mathfrak q)
\otimes_{\kappa(\mathfrak q)} L
$$
soit un anneau régulier. Choisissons une extension finie libre $R \to R'$
telle que $\mathfrak q' = \mathfrak qR'$ soit premier et que
$\kappa(\mathfrak q')$ soit isomorphe à $L$ sur $\kappa(\mathfrak q)$; voir
Algèbre, Lemme
\ref{algebra-lemma-finite-free-given-residue-field-extension}. D'après
Algèbre, Lemme \ref{algebra-lemma-completion-finite-extension}, nous avons
$$
R_\mathfrak p^\wedge \otimes_R R' = \prod (R'_{\mathfrak p_i'})^\wedge
$$
où les $\mathfrak p_i'$ sont les idéaux premiers de $R'$ situés au-dessus de
$\mathfrak p$. Ainsi,
$$
R_\mathfrak p^\wedge \otimes_R \kappa(\mathfrak q)
\otimes_{\kappa(\mathfrak q)} L =
R_\mathfrak p^\wedge \otimes_R R'
\otimes_{R'} \kappa(\mathfrak q')
=
\prod (R'_{\mathfrak p_i'})^\wedge
\otimes_{R'_{\mathfrak p'_i}} \kappa(\mathfrak q')
$$
Par hypothèse, les anneaux du membre de droite sont réguliers; l'anneau du
membre de gauche est donc lui aussi régulier. Ainsi $R$ est un G-anneau.
La réciproque résulte du Lemme \ref{lemma-G-ring-goes-up-quasi-finite}.
\end{proof}

\begin{lemma}
\label{lemma-helper-G-ring}
Soit $k$ un corps de caractéristique $p$. Posons
$A = k[[x_1, \ldots, x_n]][y_1, \ldots, y_m]$ et notons $K$ le corps des
fractions de $A$. Soit $\mathfrak p \subset A$ un idéal premier. Alors
$A_\mathfrak p^\wedge \otimes_A K$ est géométriquement régulier sur $K$.
\end{lemma}

\begin{proof}
Soit $L/K$ une extension finie purement inséparable. Nous allons démontrer
par récurrence sur $[L : K]$ que $A_\mathfrak p^\wedge \otimes L$ est
régulier. Le cas initial est $L = K$ : comme $A$ est régulier,
$A_\mathfrak p^\wedge$ est régulier (Lemme \ref{lemma-completion-regular}),
donc le localisé $A_\mathfrak p^\wedge \otimes K$ est régulier.
Soit $K \subset M \subset L$ un sous-corps tel que $L$ soit une extension
de degré $p$ de $M$ obtenue par adjonction d'une racine $p$-ième d'un
élément $f \in M$. Soit $B$ une $A$-sous-algèbre finie de $M$ dont le corps
des fractions est $M$. En chassant les dénominateurs, nous pouvons supposer,
et nous supposons, que $f \in B$. Posons $C = B[z]/(z^p -f)$ et remarquons
que $B \subset C$ est fini et que le corps des fractions de $C$ est $L$.
Comme $A \subset B \subset C$ sont finis et comme $L/M/K$ sont purement
inséparables, une certaine puissance de tout élément de $B$ ou de $C$
appartient à $A$. Il existe donc un unique idéal premier
$\mathfrak r \subset B$, respectivement $\mathfrak q \subset C$, situé
au-dessus de $\mathfrak p$. Remarquons que
$$
A_\mathfrak p^\wedge \otimes_A M = B_\mathfrak r^\wedge \otimes_B M
$$
voir Algèbre, Lemme \ref{algebra-lemma-completion-finite-extension}.
Par récurrence, cet anneau est régulier. De la même manière,
$$
A_\mathfrak p^\wedge \otimes_A L =
C_\mathfrak q^\wedge \otimes_C L =
B_\mathfrak r^\wedge \otimes_B M[z]/(z^p - f)
$$
la dernière égalité provenant du fait que le complété de
$C = B[z]/(z^p - f)$ est $B_\mathfrak r^\wedge[z]/(z^p -f)$.
D'après le Lemme \ref{lemma-find-D}, il existe une dérivation
$D : B \to B$ telle que $D(f) \not = 0$. Autrement dit, $g = D(f)$ est
inversible dans $M$. D'après le Lemme \ref{lemma-derivation-extends}, $D$
se prolonge en une dérivation de $B_\mathfrak r$, de
$B_\mathfrak r^\wedge$, puis de $B_\mathfrak r^\wedge \otimes_B M$
(en passant successivement à un localisé, à un complété, puis à un autre
localisé). Nous obtenons ainsi une dérivation de
$B_\mathfrak r^\wedge \otimes_B M$ qui envoie $f$ sur $g$, et $g$ est
inversible dans l'anneau $B_\mathfrak r^\wedge \otimes_B M$.
Le Lemme \ref{lemma-degree-p-extension-regular} montre donc que
$A_\mathfrak p^\wedge \otimes_A L$ est régulier.
\end{proof}

\begin{proposition}
\label{proposition-Noetherian-complete-G-ring}
Tout anneau local noethérien complet est un G-anneau.
\end{proposition}

\begin{proof}
Soit $A$ un anneau local noethérien complet. D'après le
Lemme \ref{lemma-check-G-ring-easy}, il suffit de vérifier que
$B = A/\mathfrak q$ a des fibres formelles géométriquement régulières
au-dessus de l'idéal premier minimal $(0)$ de $B$. Nous pouvons donc
supposer que $A$ est intègre, et il suffit de vérifier la condition pour les
fibres formelles au-dessus de l'idéal premier minimal $(0)$ de $A$.
Soit $K$ le corps des fractions de $A$.

\medskip\noindent
Nous pouvons choisir un sous-anneau $A_0 \subset A$ qui soit local régulier
complet et tel que $A$ soit fini sur $A_0$; voir Algèbre, Lemme
\ref{algebra-lemma-complete-local-Noetherian-domain-finite-over-regular}.
De plus, nous pouvons supposer que $A_0$ est un anneau de séries formelles
sur un corps ou un anneau de Cohen. D'après le
Lemme \ref{lemma-G-ring-goes-up-quasi-finite}, il suffit de démontrer le
résultat pour $A_0$.

\medskip\noindent
Supposons que $A$ soit un anneau de séries formelles sur un corps ou un
anneau de Cohen. Comme $A$ est régulier, les localisés $A_\mathfrak p$ sont
réguliers (voir Algèbre, Définition \ref{algebra-definition-regular} et la
discussion qui la précède). Leurs complétés $A_\mathfrak p^\wedge$ sont
donc réguliers; voir le Lemme \ref{lemma-completion-regular}. La fibre
$A_{\mathfrak p}^\wedge \otimes_A K$, qui est un localisé de
$A_\mathfrak p^\wedge$, est donc elle aussi régulière. Cela conclut si $K$
est de caractéristique $0$. Le cas de caractéristique positive est
$A = k[[x_1, \ldots, x_d]]$, qui est un cas particulier du
Lemme \ref{lemma-helper-G-ring}.
\end{proof}

\begin{lemma}
\label{lemma-check-G-ring-maximal-ideals}
Soit $R$ un anneau noethérien. Alors $R$ est un G-anneau si et seulement si
$R_\mathfrak m$ a des fibres formelles géométriquement régulières pour tout
idéal maximal $\mathfrak m$ de $R$.
\end{lemma}

\begin{proof}
Supposons que $R_\mathfrak m \to R_\mathfrak m^\wedge$ soit régulier pour
tout idéal maximal $\mathfrak m$ de $R$. Soit $\mathfrak p$ un idéal
premier de $R$ et choisissons un idéal maximal
$\mathfrak p \subset \mathfrak m$. Comme
$R_\mathfrak m \to R_\mathfrak m^\wedge$ est fidèlement plat, nous pouvons
choisir un idéal premier $\mathfrak p'$ de $R_\mathfrak m^\wedge$ situé
au-dessus de $\mathfrak pR_\mathfrak m$. Considérons le diagramme commutatif
$$
\xymatrix{
R_\mathfrak m^\wedge \ar[r] &
(R_\mathfrak m^\wedge)_{\mathfrak p'} \ar[r] &
(R_\mathfrak m^\wedge)_{\mathfrak p'}^\wedge
\\
R_\mathfrak m \ar[u] \ar[r] & R_\mathfrak p \ar[u] \ar[r] &
R_\mathfrak p^\wedge \ar[u]
}
$$
Par hypothèse, l'homomorphisme d'anneaux
$R_\mathfrak m \to R_\mathfrak m^\wedge$ est régulier. D'après la
Proposition \ref{proposition-Noetherian-complete-G-ring},
$(R_\mathfrak m^\wedge)_{\mathfrak p'} \to
(R_\mathfrak m^\wedge)_{\mathfrak p'}^\wedge$ est régulier. La localisation
$R_\mathfrak m^\wedge \to (R_\mathfrak m^\wedge)_{\mathfrak p'}$ est
régulière. Le Lemme \ref{lemma-regular-composition} montre donc que
$R_\mathfrak m \to (R_\mathfrak m^\wedge)_{\mathfrak p'}^\wedge$ est
régulier. Comme cet homomorphisme se factorise par le localisé
$R_\mathfrak p$, l'homomorphisme d'anneaux
$R_\mathfrak p \to (R_\mathfrak m^\wedge)_{\mathfrak p'}^\wedge$ est lui
aussi régulier. Le Lemme \ref{lemma-regular-permanence} entraîne alors que
$R_\mathfrak p \to R_\mathfrak p^\wedge$ est régulier.
\end{proof}

\begin{lemma}
\label{lemma-henselization-G-ring}
Soit $R$ un anneau local noethérien qui est un G-anneau. Alors le hensélisé
$R^h$ et le hensélisé strict $R^{sh}$ sont des G-anneaux.
\end{lemma}

\begin{proof}
Nous utiliserons le critère du
Lemme \ref{lemma-check-G-ring-maximal-ideals}. Soit
$\mathfrak q \subset R^h$ un idéal premier et posons
$\mathfrak p = R \cap \mathfrak q$. Posons $\mathfrak q_1 = \mathfrak q$ et
soient $\mathfrak q_2, \ldots, \mathfrak q_t$ les autres idéaux premiers
de $R^h$ situés au-dessus de $\mathfrak p$, de sorte que
$R^h \otimes_R \kappa(\mathfrak p) =
\prod\nolimits_{i = 1, \ldots, t} \kappa(\mathfrak q_i)$; voir le
Lemme \ref{lemma-fibres-henselization}.
Comme $(R^h)^\wedge = R^\wedge$
(Lemme \ref{lemma-henselization-noetherian}), nous obtenons
$$
\prod\nolimits_{i = 1, \ldots, t}
(R^h)^\wedge \otimes_{R^h} \kappa(\mathfrak q_i) =
(R^h)^\wedge \otimes_{R^h} (R^h \otimes_R \kappa(\mathfrak p)) =
R^\wedge \otimes_R \kappa(\mathfrak p)
$$
Par hypothèse, $(R^h)^\wedge \otimes_{R^h} \kappa(\mathfrak q_i)$ est donc
géométriquement régulier sur $\kappa(\mathfrak p)$. Comme
$\kappa(\mathfrak q_i)$ est algébrique séparable sur
$\kappa(\mathfrak p)$, il résulte d'Algèbre, Lemme
\ref{algebra-lemma-geometrically-regular-over-separable-algebraic}, que
$(R^h)^\wedge \otimes_{R^h} \kappa(\mathfrak q_i)$ est géométriquement
régulier sur $\kappa(\mathfrak q_i)$.

\medskip\noindent
Soit $\mathfrak r \subset R^{sh}$ un idéal premier et posons
$\mathfrak p = R \cap \mathfrak r$. Posons $\mathfrak r_1 = \mathfrak r$ et
soient $\mathfrak r_2, \ldots, \mathfrak r_s$ les autres idéaux premiers
de $R^{sh}$ situés au-dessus de $\mathfrak p$, de sorte que
$R^{sh} \otimes_R \kappa(\mathfrak p) =
\prod\nolimits_{i = 1, \ldots, s} \kappa(\mathfrak r_i)$; voir le
Lemme \ref{lemma-fibres-henselization}.
Nous avons alors
$$
\prod\nolimits_{i = 1, \ldots, s}
(R^{sh})^\wedge \otimes_{R^{sh}} \kappa(\mathfrak r_i) =
(R^{sh})^\wedge \otimes_{R^{sh}} (R^{sh} \otimes_R \kappa(\mathfrak p)) =
(R^{sh})^\wedge \otimes_R \kappa(\mathfrak p)
$$
Remarquons que $R^\wedge \to (R^{sh})^\wedge$ est formellement lisse pour
la topologie $\mathfrak m_{(R^{sh})^\wedge}$-adique; voir le
Lemme \ref{lemma-henselization-noetherian}. La
Proposition \ref{proposition-fs-regular} montre donc que
$R^\wedge \to (R^{sh})^\wedge$ est régulier. Le Lemme
\ref{lemma-regular-composition} montre que
$(R^{sh})^\wedge \otimes_{R^{sh}} \kappa(\mathfrak r_i)$ est régulier sur
$\kappa(\mathfrak p)$ puisque, par hypothèse,
$R^\wedge \otimes_R \kappa(\mathfrak p)$ est régulier sur
$\kappa(\mathfrak p)$. Enfin, comme $\kappa(\mathfrak r_i)$ est algébrique
séparable sur $\kappa(\mathfrak p)$, il résulte d'Algèbre, Lemme
\ref{algebra-lemma-geometrically-regular-over-separable-algebraic}, que
$(R^{sh})^\wedge \otimes_{R^{sh}} \kappa(\mathfrak r_i)$ est
géométriquement régulier sur $\kappa(\mathfrak r_i)$.
\end{proof}

\begin{lemma}
\label{lemma-another-helper-G-ring}
Soit $p$ un nombre premier. Soit $A$ un anneau local noethérien complet et
intègre, dont le corps des fractions $K$ est de caractéristique $p$. Soit
$\mathfrak q \subset A[x]$ un idéal maximal situé au-dessus de l'idéal
maximal de $A$, et soit $(0) \not = \mathfrak r \subset \mathfrak q$ un
idéal premier situé au-dessus de $(0) \subset A$. Alors
$A[x]_\mathfrak q^\wedge \otimes_{A[x]} \kappa(\mathfrak r)$ est
géométriquement régulier sur $\kappa(\mathfrak r)$.
\end{lemma}

\begin{proof}
Remarquons que l'extension $K \subset \kappa(\mathfrak r)$ est finie.
Étant donnée une extension finie purement inséparable
$L/\kappa(\mathfrak r)$, il existe donc une extension finie d'anneaux locaux
noethériens complets et intègres $A \subset B$ telle que
$\kappa(\mathfrak r) \otimes_A B$ se surjecte sur $L$. Il suffit de prendre
pour $B \subset L$ une $A$-sous-algèbre finie dont le corps des fractions
est $L$. Notons $\mathfrak r' \subset B[x]$ le noyau de l'homomorphisme
$B[x] = A[x] \otimes_A B \to \kappa(\mathfrak r) \otimes_A B \to L$
de sorte que $\kappa(\mathfrak r') = L$. Alors
$$
A[x]_\mathfrak q^\wedge \otimes_{A[x]} L =
A[x]_\mathfrak q^\wedge \otimes_{A[x]} B[x]
\otimes_{B[x]} \kappa(\mathfrak r') =
\prod B[x]_{\mathfrak q_i}^\wedge \otimes_{B[x]} \kappa(\mathfrak r')
$$
où $\mathfrak q_1, \ldots, \mathfrak q_t$ sont les idéaux premiers de
$B[x]$ situés au-dessus de $\mathfrak q$; voir
Algèbre, Lemme \ref{algebra-lemma-completion-finite-extension}.
Il suffit donc de démontrer que les anneaux
$B[x]_{\mathfrak q_i}^\wedge \otimes_{B[x]} \kappa(\mathfrak r')$ sont
réguliers. Nous sommes ainsi ramenés à montrer que
$A[x]_\mathfrak q^\wedge \otimes_{A[x]} \kappa(\mathfrak r)$ est régulier
dans le cas particulier où $K = \kappa(\mathfrak r)$.

\medskip\noindent
Supposons $K = \kappa(\mathfrak r)$. Dans ce cas, $\mathfrak r K[x]$ est
engendré par $x - f$ pour un certain $f \in K$, et
$$
A[x]_\mathfrak q^\wedge \otimes_{A[x]} \kappa(\mathfrak r)
=
(A[x]_\mathfrak q^\wedge \otimes_A K)/(x - f)
$$
La dérivation $D = \text{d}/\text{d}x$ de $A[x]$ se prolonge à $K[x]$ et
envoie $x - f$ sur un élément inversible de $K[x]$. De plus, $D$ se
prolonge à $A[x]_\mathfrak q^\wedge \otimes_A K$ d'après le
Lemme \ref{lemma-derivation-extends}. Comme
$A \to A[x]_\mathfrak q^\wedge$ est formellement lisse (voir les
Lemmes \ref{lemma-formally-smooth} et
\ref{lemma-formally-smooth-completion}), la
Proposition \ref{proposition-fs-regular} montre que l'anneau
$A[x]_\mathfrak q^\wedge \otimes_A K$ est régulier (les arguments de la
démonstration de cette proposition se simplifient considérablement dans ce
cas particulier). Le Lemme \ref{lemma-quotient-regular} permet de conclure.
\end{proof}

\begin{proposition}
\label{proposition-finite-type-over-G-ring}
Soit $R$ un G-anneau. Si $R \to S$ est essentiellement de type fini, alors
$S$ est un G-anneau.
\end{proposition}

\begin{proof}
Comme la propriété d'être un G-anneau est une propriété des anneaux locaux,
il est clair que tout localisé d'un G-anneau est un G-anneau. Réciproquement,
si chaque localisation en un idéal premier est un G-anneau, alors l'anneau
est un G-anneau. Il suffit donc de montrer que $S_\mathfrak q$ est un
G-anneau pour toute $R$-algèbre de type fini $S$ et tout idéal premier
$\mathfrak q$ de $S$. En écrivant $S$ comme quotient de
$R[x_1, \ldots, x_n]$, le Lemme
\ref{lemma-G-ring-goes-up-quasi-finite} ramène la démonstration à établir que
$R[x_1, \ldots, x_n]$ est un G-anneau. Par récurrence sur $n$, il suffit de
montrer que $R[x]$ est un G-anneau. Soit $\mathfrak q \subset R[x]$ un idéal
maximal. D'après le Lemme \ref{lemma-check-G-ring-maximal-ideals}, il suffit
de montrer que
$$
R[x]_\mathfrak q \longrightarrow R[x]_\mathfrak q^\wedge
$$
soit régulier. Si $\mathfrak q$ est situé au-dessus de
$\mathfrak p \subset R$, nous pouvons remplacer $R$ par $R_\mathfrak p$.
Nous pouvons donc supposer que $R$ est un G-anneau local noethérien d'idéal
maximal $\mathfrak m$ et que $\mathfrak q \subset R[x]$ est situé au-dessus
de $\mathfrak m$. Remarquons qu'il existe un unique idéal premier
$\mathfrak q' \subset R^\wedge[x]$ situé au-dessus de $\mathfrak q$.
Considérons le diagramme
$$
\xymatrix{
R[x]_\mathfrak q^\wedge \ar[r] &
(R^\wedge[x]_{\mathfrak q'})^\wedge \\
R[x]_\mathfrak q \ar[r] \ar[u] & R^\wedge[x]_{\mathfrak q'} \ar[u]
}
$$
Comme $R$ est un G-anneau, la flèche horizontale inférieure est régulière
(elle est un localisé d'un changement de base de l'homomorphisme régulier
d'anneaux $R \to R^\wedge$). Supposons que la flèche verticale de droite soit
régulière. La composée
$R[x]_\mathfrak q \to (R^\wedge[x]_{\mathfrak q'})^\wedge$ est alors
régulière, et le Lemme \ref{lemma-regular-permanence} montre que la flèche
verticale de gauche est régulière. Nous pouvons donc supposer que $R$ est un
anneau local noethérien complet et que $\mathfrak q$ est un idéal premier
situé au-dessus de l'idéal maximal de $R$.

\medskip\noindent
Soit $R$ un anneau local noethérien complet et soit
$\mathfrak q \subset R[x]$ un idéal maximal situé au-dessus de l'idéal
maximal de $R$. Soit $\mathfrak r \subset \mathfrak q$ un idéal premier.
Nous voulons montrer que
$R[x]_\mathfrak q^\wedge \otimes_{R[x]} \kappa(\mathfrak r)$ est une
algèbre géométriquement régulière sur $\kappa(\mathfrak r)$. Posons
$\mathfrak p = R \cap \mathfrak r$. Nous pouvons remplacer $R$ par
$R/\mathfrak p$, et $\mathfrak q$ et $\mathfrak r$ par leurs images dans
$R/\mathfrak p[x]$; voir le Lemme \ref{lemma-check-G-ring-easy}.
Nous pouvons donc supposer que $R$ est intègre et que
$\mathfrak r \cap R = (0)$.

\medskip\noindent
Par Algèbre, Lemme
\ref{algebra-lemma-complete-local-Noetherian-domain-finite-over-regular}
nous pouvons trouver $R_0 \subset R$ régulier et tel que $R$ soit fini sur
$R_0$. Le Lemme \ref{lemma-G-ring-goes-up-quasi-finite} montre qu'il suffit
de démontrer que
$R[x]_\mathfrak q^\wedge \otimes_{R[x]} \kappa(\mathfrak r)$ est
géométriquement régulier sur $\kappa(\mathfrak r)$ lorsque, outre les
hypothèses précédentes, $R$ est un anneau local régulier complet.

\medskip\noindent
À présent, $R$ est un anneau local régulier complet,
$\mathfrak r \subset \mathfrak q \subset R[x]$,
$(0) = R \cap \mathfrak r$, et $\mathfrak q$ est un idéal maximal situé
au-dessus de l'idéal maximal de $R$. Comme $R$ est régulier, l'anneau $R[x]$
est régulier (Algèbre, Lemme \ref{algebra-lemma-regular-goes-up}). Son
localisé $R[x]_\mathfrak q$ est donc régulier, ainsi que son complété
$R[x]_\mathfrak q^\wedge$; voir le Lemme \ref{lemma-completion-regular}.
La fibre $R[x]_{\mathfrak q}^\wedge \otimes_{R[x]} \kappa(\mathfrak r)$,
qui est un localisé de $R[x]_\mathfrak q^\wedge$, est donc elle aussi
régulière. Cela conclut si le corps des fractions de $R$ est de
caractéristique $0$.

\medskip\noindent
Si $R$ est de caractéristique positive, alors
$R = k[[x_1, \ldots, x_n]]$. Nous distinguons alors deux cas :
\begin{enumerate}
\item Le cas $\mathfrak r = (0)$. Le résultat est une conséquence directe
du Lemme \ref{lemma-helper-G-ring}.
\item Le cas $\mathfrak r \not = (0)$. C'est le
Lemme \ref{lemma-another-helper-G-ring}.
\end{enumerate}
\end{proof}

\begin{remark}
\label{remark-G-does-not-survive-completion}
Soit $R$ un G-anneau et soit $I \subset R$ un idéal. En général, le complété
$I$-adique $R^\wedge$ n'est pas un G-anneau. Nishimura en a donné un exemple
dans \cite{Nishimura}. Une généralisation et, en un certain sens, une
clarification de cet exemple se trouvent dans la dernière section de
\cite{Dumitrescu}.
\end{remark}

\begin{proposition}
\label{proposition-ubiquity-G-ring}
Les types d'anneaux suivants sont des G-anneaux :
\begin{enumerate}
\item les corps,
\item les anneaux locaux noethériens complets,
\item $\mathbf{Z}$,
\item les anneaux de Dedekind dont le corps des fractions est de
caractéristique zéro,
\item les extensions d'anneaux de type fini de l'un quelconque des anneaux
précédents.
\end{enumerate}
\end{proposition}

\begin{proof}
Pour les corps, $\mathbf{Z}$ et les anneaux de Dedekind de caractéristique
zéro, cela résulte immédiatement de la définition et du fait que le complété
d'un anneau de valuation discrète est un anneau de valuation discrète. Tout
anneau local noethérien complet est un G-anneau d'après la
Proposition \ref{proposition-Noetherian-complete-G-ring}. L'assertion sur les
extensions d'anneaux de type fini est la
Proposition \ref{proposition-finite-type-over-G-ring}.
\end{proof}

\begin{lemma}
\label{lemma-henselian-local-limit-G-rings}
Soit $(A, \mathfrak m)$ un anneau local hensélien. Alors $A$ est la colimite
filtrante d'un système de G-anneaux locaux henséliens dont les homomorphismes
de transition sont locaux.
\end{lemma}

\begin{proof}
Écrivons $A = \colim A_i$ comme colimite filtrante de
$\mathbf{Z}$-algèbres de type fini. Soit $\mathfrak p_i$ l'idéal premier de
$A_i$ situé sous $\mathfrak m$. Nous pouvons remplacer $A_i$ par le localisé
de $A_i$ en $\mathfrak p_i$. Alors $A_i$ est un G-anneau local noethérien
(Proposition \ref{proposition-ubiquity-G-ring}). D'après le
Lemme \ref{lemma-henselization-colimit}, nous avons
$A = \colim A_i^h$. Enfin, le Lemme \ref{lemma-henselization-G-ring} montre
que les anneaux $A_i^h$ sont des G-anneaux.
\end{proof}

\begin{lemma}
\label{lemma-map-G-ring-to-completion-regular}
\begin{reference}
\cite[Theorem 79]{MatCA}
\end{reference}
Soit $A$ un G-anneau. Soit $I \subset A$ un idéal et soit $A^\wedge$ le
complété de $A$ pour la topologie $I$-adique. Alors $A \to A^\wedge$ est
régulier.
\end{lemma}

\begin{proof}
L'homomorphisme d'anneaux $A \to A^\wedge$ est plat d'après
Algèbre, Lemme \ref{algebra-lemma-completion-flat}. L'anneau $A^\wedge$ est
noethérien d'après Algèbre, Lemme
\ref{algebra-lemma-completion-Noetherian-Noetherian}. Il suffit donc de
vérifier la troisième condition du Lemme \ref{lemma-regular-local}.
Soit $\mathfrak m' \subset A^\wedge$ un idéal maximal situé au-dessus de
$\mathfrak m \subset A$. D'après Algèbre, Lemme
\ref{algebra-lemma-radical-completion}, nous avons
$IA^\wedge \subset \mathfrak m'$. Comme $A^\wedge/IA^\wedge = A/I$, nous
avons $I \subset \mathfrak m$,
$\mathfrak m/I = \mathfrak m'/IA^\wedge$ et
$A/\mathfrak m = A^\wedge/\mathfrak m'$. Puisque
$A^\wedge/\mathfrak m'$ est un corps, $\mathfrak m$ est lui aussi un idéal
maximal. Alors $A_\mathfrak m \to A^\wedge_{\mathfrak m'}$ est un
homomorphisme local plat d'anneaux locaux noethériens qui identifie les
corps résiduels et tel que
$\mathfrak m A^\wedge_{\mathfrak m'} =
\mathfrak m'A^\wedge_{\mathfrak m'}$. Il induit donc un isomorphisme sur les
anneaux locaux complétés; voir le Lemme \ref{lemma-flat-unramified}.
Soit $(A_\mathfrak m)^\wedge$ le complété de $A_\mathfrak m$ pour la
topologie définie par son idéal maximal.
L'homomorphisme d'anneaux
$$
(A^\wedge)_{\mathfrak m'} \to
((A^\wedge)_{\mathfrak m'})^\wedge = (A_\mathfrak m)^\wedge
$$
est fidèlement plat (Algèbre, Lemme
\ref{algebra-lemma-completion-faithfully-flat}). Nous pouvons donc appliquer
le Lemme \ref{lemma-regular-permanence} aux homomorphismes d'anneaux
$$
A_\mathfrak m \to (A^\wedge)_{\mathfrak m'} \to (A_\mathfrak m)^\wedge
$$
et conclure, puisque $A_\mathfrak m \to (A_\mathfrak m)^\wedge$ est
régulier du fait que $A$ est un G-anneau.
\end{proof}

\begin{lemma}
\label{lemma-henselization-pair-G-ring}
\begin{slogan}
La propriété d'être un G-anneau est stable par hensélisation le long d'idéaux
\end{slogan}
\begin{reference}
\cite[Theorem 5.3 i)]{Greco}
\end{reference}
Soit $A$ un G-anneau. Soit $I \subset A$ un idéal. Soit $(A^h, I^h)$ le
hensélisé de la paire $(A, I)$; voir le Lemme \ref{lemma-henselization}.
Alors $A^h$ est un G-anneau.
\end{lemma}

\begin{proof}
Soit $\mathfrak m^h \subset A^h$ un idéal maximal. Nous devons montrer que
l'homomorphisme de $A^h_{\mathfrak m^h}$ vers son complété a des fibres
géométriquement régulières; voir le
Lemme \ref{lemma-check-G-ring-maximal-ideals}. Soit $\mathfrak m$ l'image
réciproque de $\mathfrak m^h$ dans $A$. Remarquons que
$I^h \subset \mathfrak m^h$, donc que $I \subset \mathfrak m$, puisque
$(A^h, I^h)$ est une paire hensélienne. Rappelons que $A^h$ est noethérien,
que $I^h = IA^h$ et que $A \to A^h$ induit un isomorphisme sur les complétés
$I$-adiques; voir le Lemme \ref{lemma-henselization-Noetherian-pair}.
Alors l'homomorphisme local d'anneaux locaux noethériens
$$
A_\mathfrak m \to A^h_{\mathfrak m^h}
$$
induit un isomorphisme sur les complétés aux idéaux maximaux d'après le
Lemme \ref{lemma-flat-unramified} (nous omettons les détails). Soit
$\mathfrak q^h$ un idéal premier de $A^h_{\mathfrak m^h}$ situé au-dessus de
$\mathfrak q \subset A_\mathfrak m$. Posons
$\mathfrak q_1 = \mathfrak q^h$ et soient
$\mathfrak q_2, \ldots, \mathfrak q_t$ les autres idéaux premiers de $A^h$
situés au-dessus de $\mathfrak q$, de sorte que
$A^h \otimes_A \kappa(\mathfrak q) =
\prod\nolimits_{i = 1, \ldots, t} \kappa(\mathfrak q_i)$; voir le
Lemme \ref{lemma-filtered-colimit-etale-noetherian-fibres}.
Comme $(A^h)_{\mathfrak m^h}^\wedge = (A_\mathfrak m)^\wedge$, ainsi qu'il
a été expliqué ci-dessus, nous obtenons
$$
\prod\nolimits_{i = 1, \ldots, t}
(A^h_{\mathfrak m^h})^\wedge \otimes_{A^h_{\mathfrak m^h}}
\kappa(\mathfrak q_i) =
(A^h_{\mathfrak m^h})^\wedge \otimes_{A^h_{\mathfrak m^h}}
(A^h_{\mathfrak m^h} \otimes_{A_{\mathfrak m}} \kappa(\mathfrak q)) =
(A_{\mathfrak m})^\wedge \otimes_{A_{\mathfrak m}} \kappa(\mathfrak q)
$$
Par conséquent, en tant que l'une des composantes, l'anneau
$$
(A^h_{\mathfrak m^h})^\wedge \otimes_{A^h_{\mathfrak m^h}}
\kappa(\mathfrak q^h)
$$
est géométriquement régulier sur $\kappa(\mathfrak q)$ par l'hypothèse sur
$A$. Comme $\kappa(\mathfrak q^h)$ est algébrique séparable sur
$\kappa(\mathfrak q)$, il résulte d'Algèbre, Lemme
\ref{algebra-lemma-geometrically-regular-over-separable-algebraic}, que
$$
(A^h_{\mathfrak m^h})^\wedge \otimes_{A^h_{\mathfrak m^h}}
\kappa(\mathfrak q^h)
$$
est géométriquement régulier sur $\kappa(\mathfrak q^h)$, comme voulu.
\end{proof}







\section{Propriétés des fibres formelles}
\label{section-properties-formal-fibres}

\noindent
Dans cette section, nous reprenons certains des arguments de la
Section \ref{section-G-ring} afin de pouvoir étudier convenablement les
propriétés des fibres formelles des anneaux noethériens.

\medskip\noindent
Soit $P$ une propriété des homomorphismes d'anneaux $k \to R$, où $k$ est
un corps et $R$ est noethérien. Nous disons que $P$ est vérifiée par les
fibres d'un homomorphisme d'anneaux $A \to B$, avec $B$ noethérien, si la
propriété $P$ est vérifiée pour
$\kappa(\mathfrak q) \to B \otimes_A \kappa(\mathfrak q)$ pour tout idéal
premier $\mathfrak q$ de $A$. Nous utiliserons les
assertions suivantes :
\begin{enumerate}
\item[(A)] $P(k \to R) \Rightarrow P(k' \to R \otimes_k k')$ pour toute
extension de corps de type fini $k'/k$ ;
\item[(B)] $P(k \to R_\mathfrak p),\ \forall \mathfrak p \in \Spec(R)
\Leftrightarrow P(k \to R)$ ;
\item[(C)] étant donnés des homomorphismes plats $A \to B \to C$ d'anneaux
noethériens, si les fibres de $A \to B$ ont la propriété $P$ et si
$B \to C$ est régulier, alors les fibres de $A \to C$ ont la propriété $P$ ;
\item[(D)] étant donnés des homomorphismes plats $A \to B \to C$ d'anneaux
noethériens, si les fibres de $A \to C$ ont la propriété $P$ et si
$B \to C$ est fidèlement plat, alors les fibres de $A \to B$ ont la
propriété $P$ ;
\item[(E)] étant donnés $k \to k' \to R$, avec $R$ noethérien, si $k'/k$
est algébrique séparable et si $P(k \to R)$, alors $P(k' \to R)$ ;
\item[(F)] ajouter ici d'autres propriétés.
\end{enumerate}
Étant donné un anneau local noethérien $A$, nous disons que
``{\it les fibres formelles de $A$ ont la propriété $P$}'' si la propriété
$P$ est vérifiée pour les fibres de $A \to A^\wedge$. Nous disons que
{\it $R$ est un $P$-anneau} si $R$ est noethérien et si, pour tout idéal
premier $\mathfrak p$ de $R$, les fibres formelles de $R_\mathfrak p$ ont
la propriété $P$.

\begin{lemma}
\label{lemma-check-P-ring-easy}
Soit $R$ un anneau noethérien. Soit $P$ une propriété comme ci-dessus.
Alors $R$ est un $P$-anneau si et seulement si, pour toute paire d'idéaux
premiers $\mathfrak q \subset \mathfrak p \subset R$, la
$\kappa(\mathfrak q)$-algèbre
$$
(R/\mathfrak q)_\mathfrak p^\wedge \otimes_{R/\mathfrak q} \kappa(\mathfrak q)
$$
a la propriété $P$.
\end{lemma}

\begin{proof}
Cela résulte de l'égalité
$$
R_\mathfrak p^\wedge \otimes_R \kappa(\mathfrak q) =
(R/\mathfrak q)_\mathfrak p^\wedge \otimes_{R/\mathfrak q} \kappa(\mathfrak q)
$$
entre algèbres sur $\kappa(\mathfrak q)$.
\end{proof}

\begin{lemma}
\label{lemma-P-local}
Soit $R \to \Lambda$ un homomorphisme d'anneaux noethériens. Supposons que
$P$ vérifie (B). Les conditions suivantes sont équivalentes :
\begin{enumerate}
\item les fibres de $R \to \Lambda$ ont la propriété $P$ ;
\item les fibres de $R_\mathfrak p \to \Lambda_\mathfrak q$ ont la
propriété $P$ pour tout $\mathfrak q \subset \Lambda$ situé au-dessus de
$\mathfrak p \subset R$ ;
\item les fibres de $R_\mathfrak m \to \Lambda_{\mathfrak m'}$ ont la
propriété $P$ pour tout idéal maximal $\mathfrak m' \subset \Lambda$ situé
au-dessus d'un idéal maximal $\mathfrak m$ de $R$.
\end{enumerate}
\end{lemma}

\begin{proof}
Soit $\mathfrak p \subset R$ un idéal premier. La fibre au-dessus de
$\mathfrak p$ est l'anneau $\Lambda \otimes_R \kappa(\mathfrak p)$, dont
le spectre s'envoie bijectivement sur le sous-ensemble de
$\Spec(\Lambda)$ formé des idéaux premiers $\mathfrak q$ situés au-dessus
de $\mathfrak p$; voir Algèbre, Remarque
\ref{algebra-remark-fundamental-diagram}. Pour un tel $\mathfrak q$,
choisissons un idéal maximal $\mathfrak q \subset \mathfrak m'$ et posons
$\mathfrak m = R \cap \mathfrak m'$. Alors
$\mathfrak p \subset \mathfrak m$ et
$$
(\Lambda \otimes_R \kappa(\mathfrak p))_\mathfrak q \cong
(\Lambda_{\mathfrak m'} \otimes_{R_\mathfrak m}
\kappa(\mathfrak p))_\mathfrak q
$$
comme $\kappa(\mathfrak p)$-algèbres. Ainsi (1), (2) et (3) sont
équivalentes, puisque (B) permet de vérifier la propriété $P$ sur les
anneaux locaux.
\end{proof}

\begin{lemma}
\label{lemma-P-ring-goes-up-quasi-finite}
Soit $R \to R'$ un homomorphisme de type fini d'anneaux noethériens, et
soient
$$
\xymatrix{
\mathfrak q' \ar[r] & \mathfrak p' \ar[r] & R' \\
\mathfrak q \ar[r] \ar@{-}[u] &
\mathfrak p \ar[r] \ar@{-}[u] & R \ar[u]
}
$$
des idéaux premiers. Supposons que $R \to R'$ soit quasi-fini en
$\mathfrak p'$ et que $P$ vérifie (A) et (B).
\begin{enumerate}
\item Si $\kappa(\mathfrak q) \to
R_\mathfrak p^\wedge \otimes_R \kappa(\mathfrak q)$
a la propriété $P$, alors
$\kappa(\mathfrak q') \to
(R'_{\mathfrak p'})^\wedge \otimes_{R'} \kappa(\mathfrak q')$
a la propriété $P$.
\item Si les fibres formelles de $R_\mathfrak p$ ont la propriété $P$,
alors celles de $R'_{\mathfrak p'}$ ont la propriété $P$.
\item Si $R \to R'$ est quasi-fini et si $R$ est un $P$-anneau, alors $R'$
est un $P$-anneau.
\end{enumerate}
\end{lemma}

\begin{proof}
Il est clair que (1) $\Rightarrow$ (2) $\Rightarrow$ (3). Supposons que la
propriété $P$ soit vérifiée pour
$\kappa(\mathfrak q) \to R_\mathfrak p^\wedge \otimes_R \kappa(\mathfrak q)$.
D'après Algèbre, Lemme
\ref{algebra-lemma-completion-at-quasi-finite-prime}, nous avons
$$
R_\mathfrak p^\wedge \otimes_R R'
=
(R'_{\mathfrak p'})^\wedge \times B
$$
pour une certaine $R_\mathfrak p^\wedge$-algèbre $B$. Par conséquent,
$R'_{\mathfrak p'} \to (R'_{\mathfrak p'})^\wedge$ est un facteur d'un
changement de base de $R_\mathfrak p \to R_\mathfrak p^\wedge$. Il s'ensuit
que $(R'_{\mathfrak p'})^\wedge \otimes_{R'} \kappa(\mathfrak q')$ est un
facteur de
$$
R_\mathfrak p^\wedge \otimes_R R' \otimes_{R'} \kappa(\mathfrak q') =
R_\mathfrak p^\wedge \otimes_R \kappa(\mathfrak q)
\otimes_{\kappa(\mathfrak q)} \kappa(\mathfrak q').
$$
Le résultat découle donc des hypothèses sur $P$.
\end{proof}

\begin{lemma}
\label{lemma-check-P-ring-maximal-ideals}
Soit $R$ un anneau noethérien. Supposons que $P$ vérifie (C) et (D).
Alors $R$ est un $P$-anneau si et seulement si les fibres formelles de
$R_\mathfrak m$ ont la propriété $P$ pour tout idéal maximal
$\mathfrak m$ de $R$.
\end{lemma}

\begin{proof}
Supposons que les fibres formelles de $R_\mathfrak m$ aient la propriété
$P$ pour tout idéal maximal $\mathfrak m$ de $R$. Soit $\mathfrak p$ un
idéal premier de $R$ et choisissons un idéal maximal
$\mathfrak p \subset \mathfrak m$. Comme
$R_\mathfrak m \to R_\mathfrak m^\wedge$ est fidèlement plat, nous pouvons
choisir un idéal premier $\mathfrak p'$ de $R_\mathfrak m^\wedge$ situé
au-dessus de $\mathfrak pR_\mathfrak m$. Considérons le diagramme commutatif
$$
\xymatrix{
R_\mathfrak m^\wedge \ar[r] &
(R_\mathfrak m^\wedge)_{\mathfrak p'} \ar[r] &
(R_\mathfrak m^\wedge)_{\mathfrak p'}^\wedge
\\
R_\mathfrak m \ar[u] \ar[r] & R_\mathfrak p \ar[u] \ar[r] &
R_\mathfrak p^\wedge \ar[u]
}
$$
Par hypothèse, les fibres de l'homomorphisme d'anneaux
$R_\mathfrak m \to R_\mathfrak m^\wedge$ ont la propriété $P$.
D'après la Proposition \ref{proposition-Noetherian-complete-G-ring},
$(R_\mathfrak m^\wedge)_{\mathfrak p'} \to
(R_\mathfrak m^\wedge)_{\mathfrak p'}^\wedge$ est régulier. La localisation
$R_\mathfrak m^\wedge \to (R_\mathfrak m^\wedge)_{\mathfrak p'}$ est
régulière. Le Lemme \ref{lemma-regular-composition} montre donc que
$R_\mathfrak m^\wedge \to (R_\mathfrak m^\wedge)_{\mathfrak p'}^\wedge$
est régulier. Par (C), les fibres de
$R_\mathfrak m \to (R_\mathfrak m^\wedge)_{\mathfrak p'}^\wedge$ ont donc
la propriété $P$. Comme
$R_\mathfrak m \to (R_\mathfrak m^\wedge)_{\mathfrak p'}^\wedge$ se
factorise par le localisé $R_\mathfrak p$, les fibres de
$R_\mathfrak p \to (R_\mathfrak m^\wedge)_{\mathfrak p'}^\wedge$ ont aussi
la propriété $P$. Nous pouvons alors appliquer (D) pour voir que les fibres
de $R_\mathfrak p \to R_\mathfrak p^\wedge$ ont la propriété $P$.
\end{proof}

\begin{proposition}
\label{proposition-finite-type-over-P-ring}
Soit $R$ un $P$-anneau, où $P$ vérifie (A), (B), (C) et (D). Si
$R \to S$ est essentiellement de type fini, alors $S$ est un $P$-anneau.
\end{proposition}

\begin{proof}
Comme la propriété d'être un $P$-anneau est une propriété des anneaux
locaux, tout localisé d'un $P$-anneau est un $P$-anneau. Réciproquement, si
chaque localisation en un idéal premier est un $P$-anneau, alors l'anneau
est un $P$-anneau. Il suffit donc de montrer que $S_\mathfrak q$ est un
$P$-anneau pour toute $R$-algèbre de type fini $S$ et tout idéal premier
$\mathfrak q$ de $S$. En écrivant $S$ comme quotient de
$R[x_1, \ldots, x_n]$, le Lemme
\ref{lemma-P-ring-goes-up-quasi-finite} ramène la démonstration à établir que
$R[x_1, \ldots, x_n]$ est un $P$-anneau. Par récurrence sur $n$, il suffit
de montrer que $R[x]$ est un $P$-anneau. Soit
$\mathfrak q \subset R[x]$ un idéal maximal. D'après le
Lemme \ref{lemma-check-P-ring-maximal-ideals}, il suffit de montrer que les
fibres de
$$
R[x]_\mathfrak q \longrightarrow R[x]_\mathfrak q^\wedge
$$
aient la propriété $P$. Si $\mathfrak q$ est situé au-dessus de
$\mathfrak p \subset R$, nous pouvons remplacer $R$ par $R_\mathfrak p$.
Nous pouvons donc supposer que $R$ est un $P$-anneau local noethérien
d'idéal maximal $\mathfrak m$ et que $\mathfrak q \subset R[x]$ est situé
au-dessus de $\mathfrak m$. Remarquons qu'il existe un unique idéal premier
$\mathfrak q' \subset R^\wedge[x]$ situé au-dessus de $\mathfrak q$.
Considérons le diagramme
$$
\xymatrix{
R[x]_\mathfrak q^\wedge \ar[r] &
(R^\wedge[x]_{\mathfrak q'})^\wedge \\
R[x]_\mathfrak q \ar[r] \ar[u] & R^\wedge[x]_{\mathfrak q'} \ar[u]
}
$$
Comme $R$ est un $P$-anneau, les fibres de $R[x] \to R^\wedge[x]$ ont la
propriété $P$ : elles s'obtiennent à partir des fibres de
$R \to R^\wedge$ par changement de base le long d'une extension de corps de
type fini, de sorte que (A) s'applique. Les fibres de la flèche horizontale
inférieure ont donc la propriété $P$, par exemple d'après le
Lemme \ref{lemma-P-local}. La flèche verticale de droite est régulière parce
que $R^\wedge$ est un G-anneau (Propositions
\ref{proposition-Noetherian-complete-G-ring} et
\ref{proposition-finite-type-over-G-ring}). Il résulte de (C) que les fibres
de la composée $R[x]_\mathfrak q \to
(R^\wedge[x]_{\mathfrak q'})^\wedge$ ont la propriété $P$. Par (D), les
fibres de la flèche verticale de gauche ont donc la propriété $P$, ce qui
achève la démonstration.
\end{proof}

\begin{lemma}
\label{lemma-map-P-ring-to-completion-P}
Soit $A$ un $P$-anneau, où $P$ vérifie (B) et (D). Soit $I \subset A$ un
idéal et soit $A^\wedge$ le complété de $A$ pour la topologie $I$-adique.
Alors les fibres de $A \to A^\wedge$ ont la propriété $P$.
\end{lemma}

\begin{proof}
L'homomorphisme d'anneaux $A \to A^\wedge$ est plat d'après
Algèbre, Lemme \ref{algebra-lemma-completion-flat}. L'anneau $A^\wedge$ est
noethérien d'après Algèbre, Lemme
\ref{algebra-lemma-completion-Noetherian-Noetherian}. Il suffit donc de
vérifier la troisième condition du Lemme \ref{lemma-P-local}. Soit
$\mathfrak m' \subset A^\wedge$ un idéal maximal situé au-dessus de
$\mathfrak m \subset A$. D'après Algèbre, Lemme
\ref{algebra-lemma-radical-completion}, nous avons
$IA^\wedge \subset \mathfrak m'$. Comme $A^\wedge/IA^\wedge = A/I$, nous
avons $I \subset \mathfrak m$,
$\mathfrak m/I = \mathfrak m'/IA^\wedge$ et
$A/\mathfrak m = A^\wedge/\mathfrak m'$. Puisque
$A^\wedge/\mathfrak m'$ est un corps, $\mathfrak m$ est lui aussi un idéal
maximal. Alors $A_\mathfrak m \to A^\wedge_{\mathfrak m'}$ est un
homomorphisme local plat d'anneaux locaux noethériens qui identifie les
corps résiduels et tel que
$\mathfrak m A^\wedge_{\mathfrak m'} =
\mathfrak m'A^\wedge_{\mathfrak m'}$. Il induit donc un isomorphisme sur les
anneaux locaux complétés; voir le Lemme \ref{lemma-flat-unramified}.
Soit $(A_\mathfrak m)^\wedge$ le complété de $A_\mathfrak m$ pour la
topologie définie par son idéal maximal.
L'homomorphisme d'anneaux
$$
(A^\wedge)_{\mathfrak m'} \to
((A^\wedge)_{\mathfrak m'})^\wedge = (A_\mathfrak m)^\wedge
$$
est fidèlement plat (Algèbre, Lemme
\ref{algebra-lemma-completion-faithfully-flat}). Nous pouvons donc appliquer
(D) aux homomorphismes d'anneaux
$$
A_\mathfrak m \to (A^\wedge)_{\mathfrak m'} \to (A_\mathfrak m)^\wedge
$$
et conclure, puisque les fibres de
$A_\mathfrak m \to (A_\mathfrak m)^\wedge$ ont la propriété $P$ du fait
que $A$ est un $P$-anneau.
\end{proof}

\begin{lemma}
\label{lemma-henselization-pair-P-ring}
\begin{slogan}
L'hensélisation d'un anneau hérite des bonnes propriétés des fibres formelles
\end{slogan}
Soit $A$ un $P$-anneau, où $P$ vérifie (B), (C), (D) et (E). Soit
$I \subset A$ un idéal. Soit $(A^h, I^h)$ le hensélisé de la paire
$(A, I)$; voir le Lemme \ref{lemma-henselization}. Alors $A^h$ est un
$P$-anneau.
\end{lemma}

\begin{proof}
Soit $\mathfrak m^h \subset A^h$ un idéal maximal. Nous devons montrer que
les fibres de $A^h_{\mathfrak m^h} \to
(A^h_{\mathfrak m^h})^\wedge$ ont la propriété $P$; voir le
Lemme \ref{lemma-check-P-ring-maximal-ideals}. Soit $\mathfrak m$ l'image
réciproque de $\mathfrak m^h$ dans $A$. Remarquons que
$I^h \subset \mathfrak m^h$, donc que $I \subset \mathfrak m$, puisque
$(A^h, I^h)$ est une paire hensélienne. Rappelons que $A^h$ est noethérien,
que $I^h = IA^h$ et que $A \to A^h$ induit un isomorphisme sur les complétés
$I$-adiques; voir le Lemme \ref{lemma-henselization-Noetherian-pair}.
Alors l'homomorphisme local d'anneaux locaux noethériens
$$
A_\mathfrak m \to A^h_{\mathfrak m^h}
$$
induit un isomorphisme sur les complétés aux idéaux maximaux d'après le
Lemme \ref{lemma-flat-unramified} (nous omettons les détails). Soit
$\mathfrak q^h$ un idéal premier de $A^h_{\mathfrak m^h}$ situé au-dessus de
$\mathfrak q \subset A_\mathfrak m$. Posons
$\mathfrak q_1 = \mathfrak q^h$ et soient
$\mathfrak q_2, \ldots, \mathfrak q_t$ les autres idéaux premiers de $A^h$
situés au-dessus de $\mathfrak q$, de sorte que
$A^h \otimes_A \kappa(\mathfrak q) =
\prod\nolimits_{i = 1, \ldots, t} \kappa(\mathfrak q_i)$; voir le
Lemme \ref{lemma-filtered-colimit-etale-noetherian-fibres}.
Comme $(A^h)_{\mathfrak m^h}^\wedge = (A_\mathfrak m)^\wedge$, ainsi qu'il
a été expliqué ci-dessus, nous obtenons
$$
\prod\nolimits_{i = 1, \ldots, t}
(A^h_{\mathfrak m^h})^\wedge \otimes_{A^h_{\mathfrak m^h}}
\kappa(\mathfrak q_i) =
(A^h_{\mathfrak m^h})^\wedge \otimes_{A^h_{\mathfrak m^h}}
(A^h_{\mathfrak m^h} \otimes_{A_{\mathfrak m}} \kappa(\mathfrak q)) =
(A_{\mathfrak m})^\wedge \otimes_{A_{\mathfrak m}} \kappa(\mathfrak q)
$$
En examinant les anneaux locaux et en utilisant (B), nous en déduisons que
$$
\kappa(\mathfrak q) \longrightarrow
(A^h_{\mathfrak m^h})^\wedge \otimes_{A^h_{\mathfrak m^h}}
\kappa(\mathfrak q^h)
$$
a la propriété $P$, puisque
$\kappa(\mathfrak q) \to
(A_\mathfrak m)^\wedge \otimes_{A_\mathfrak m} \kappa(\mathfrak q)$
en dispose par l'hypothèse sur $A$. Comme
$\kappa(\mathfrak q^h)/\kappa(\mathfrak q)$ est algébrique séparable, (E)
montre que
$\kappa(\mathfrak q^h) \to
(A^h_{\mathfrak m^h})^\wedge \otimes_{A^h_{\mathfrak m^h}}
\kappa(\mathfrak q^h)$ a la propriété $P$, comme voulu.
\end{proof}

\begin{lemma}
\label{lemma-henselization-P-ring}
Soit $R$ un anneau local noethérien qui est un $P$-anneau, où $P$ vérifie
(B), (C), (D) et (E). Alors le hensélisé $R^h$ et le hensélisé strict
$R^{sh}$ sont des $P$-anneaux.
\end{lemma}

\begin{proof}
Nous l'avons démontré pour le hensélisé dans le
Lemme \ref{lemma-henselization-pair-P-ring}. Pour le démontrer pour le
hensélisé strict, il suffit de montrer que les fibres formelles de $R^{sh}$
ont la propriété $P$; voir le
Lemme \ref{lemma-check-P-ring-maximal-ideals}. Soit
$\mathfrak r \subset R^{sh}$ un idéal premier et posons
$\mathfrak p = R \cap \mathfrak r$. Posons $\mathfrak r_1 = \mathfrak r$ et
soient $\mathfrak r_2, \ldots, \mathfrak r_s$ les autres idéaux premiers
de $R^{sh}$ situés au-dessus de $\mathfrak p$, de sorte que
$R^{sh} \otimes_R \kappa(\mathfrak p) =
\prod\nolimits_{i = 1, \ldots, s} \kappa(\mathfrak r_i)$; voir le
Lemme \ref{lemma-fibres-henselization}.
Nous avons alors
$$
\prod\nolimits_{i = 1, \ldots, s}
(R^{sh})^\wedge \otimes_{R^{sh}} \kappa(\mathfrak r_i) =
(R^{sh})^\wedge \otimes_{R^{sh}} (R^{sh} \otimes_R \kappa(\mathfrak p)) =
(R^{sh})^\wedge \otimes_R \kappa(\mathfrak p)
$$
Remarquons que $R^\wedge \to (R^{sh})^\wedge$ est formellement lisse pour
la topologie $\mathfrak m_{(R^{sh})^\wedge}$-adique; voir le
Lemme \ref{lemma-henselization-noetherian}. La
Proposition \ref{proposition-fs-regular} montre donc que
$R^\wedge \to (R^{sh})^\wedge$ est régulier. Nous concluons que la
propriété $P$ est vérifiée pour
$\kappa(\mathfrak p) \to (R^{sh})^\wedge \otimes_R \kappa(\mathfrak p)$
par (C) et notre hypothèse sur $R$. En utilisant (B), la décomposition
précédente et les anneaux locaux, nous concluons que la propriété $P$ est
vérifiée pour $\kappa(\mathfrak p) \to
(R^{sh})^\wedge \otimes_{R^{sh}} \kappa(\mathfrak r)$.
Comme $\kappa(\mathfrak r)/\kappa(\mathfrak p)$ est algébrique séparable,
il résulte de (E) que la propriété $P$ est vérifiée pour
$\kappa(\mathfrak r) \to
(R^{sh})^\wedge \otimes_{R^{sh}} \kappa(\mathfrak r)$.
\end{proof}

\begin{lemma}
\label{lemma-formal-fibres-reduced}
Les propriétés (A), (B), (C), (D) et (E) sont vérifiées pour
$P(k \to R) =$ « $R$ est géométriquement réduit sur $k$ ».
\end{lemma}

\begin{proof}
La partie (A) résulte de la définition des algèbres géométriquement réduites
(Algèbre, Définition \ref{algebra-definition-geometrically-reduced}).
La partie (B) en résulte aussi : un anneau est réduit si et seulement si
tous ses anneaux locaux sont réduits. La partie (C) résulte du
Lemme \ref{lemma-reduced-goes-up}. La partie (D) résulte d'Algèbre, Lemme
\ref{algebra-lemma-descent-reduced}. La partie (E) résulte d'Algèbre, Lemme
\ref{algebra-lemma-geometrically-reduced-over-separable-algebraic}.
\end{proof}

\begin{lemma}
\label{lemma-formal-fibres-normal}
Les propriétés (A), (B), (C), (D) et (E) sont vérifiées pour
$P(k \to R) =$ « $R$ est géométriquement normal sur $k$ ».
\end{lemma}

\begin{proof}
La partie (A) résulte de la définition des algèbres géométriquement normales
(Algèbre, Définition \ref{algebra-definition-geometrically-normal}).
La partie (B) en résulte aussi : un anneau est normal si et seulement si
tous ses anneaux locaux sont normaux. La partie (C) résulte du
Lemme \ref{lemma-normal-goes-up}. La partie (D) résulte d'Algèbre, Lemme
\ref{algebra-lemma-descent-normal}. La partie (E) résulte d'Algèbre, Lemme
\ref{algebra-lemma-geometrically-normal-over-separable-algebraic}.
\end{proof}

\begin{lemma}
\label{lemma-formal-fibres-Sk}
Fixons $n \geq 1$. Les propriétés (A), (B), (C), (D) et (E) sont vérifiées
pour $P(k \to R) =$ « $R$ vérifie $(S_n)$ ».
\end{lemma}

\begin{proof}
Soit $k \to R$ un homomorphisme d'anneaux, où $k$ est un corps et $R$ un
anneau noethérien. Soit $k'/k$ une extension de corps de type fini. D'après
Algèbre, Lemme \ref{algebra-lemma-tensor-fields-CM}, les fibres de
l'homomorphisme d'anneaux $R \to R \otimes_k k'$ sont de Cohen--Macaulay.
Nous pouvons donc appliquer Algèbre, Lemme
\ref{algebra-lemma-Sk-goes-up}, à $R \to R \otimes_k k'$ pour voir que, si
$R$ vérifie $(S_n)$, alors $R \otimes_k k'$ la vérifie aussi. Cela
démontre (A). La partie (B) en résulte aussi : un anneau noethérien vérifie
$(S_n)$ si et seulement si tous ses anneaux locaux vérifient $(S_n)$.
La partie (C) résulte d'Algèbre, Lemme
\ref{algebra-lemma-Sk-goes-up}, puisque les fibres d'un homomorphisme
régulier sont régulières et, en particulier, de Cohen--Macaulay.
La partie (D) résulte d'Algèbre, Lemme
\ref{algebra-lemma-descent-Sk}. La partie (E) est immédiate, car la condition
ne fait pas intervenir le corps de base.
\end{proof}

\begin{lemma}
\label{lemma-formal-fibres-CM}
Les propriétés (A), (B), (C), (D) et (E) sont vérifiées pour
$P(k \to R) =$ « $R$ est de Cohen--Macaulay ».
\end{lemma}

\begin{proof}
Cela résulte immédiatement du Lemme \ref{lemma-formal-fibres-Sk} et du fait
qu'un anneau noethérien est de Cohen--Macaulay si et seulement s'il vérifie
les conditions $(S_n)$ pour tout $n$.
\end{proof}

\begin{lemma}
\label{lemma-formal-fibres-Rk}
Fixons $n \geq 0$. Les propriétés (A), (B), (C), (D) et (E) sont vérifiées
pour $P(k \to R) =$ « $R \otimes_k k'$ vérifie $(R_n)$ pour toute extension
finie $k'/k$ ».
\end{lemma}

\begin{proof}
Soit $k \to R$ un homomorphisme d'anneaux, où $k$ est un corps et $R$ un
anneau noethérien. Supposons $P(k \to R)$ vérifiée. Soit $K/k$ une extension
de corps de type fini. D'après Algèbre, Lemme
\ref{algebra-lemma-make-separable}, il existe un diagramme
$$
\xymatrix{
K \ar[r] & K' \\
k \ar[u] \ar[r] & k' \ar[u]
}
$$
où $k'/k$ et $K'/K$ sont des extensions de corps finies purement
inséparables et où $K'/k'$ est séparable. D'après Algèbre, Lemme
\ref{algebra-lemma-localization-smooth-separable}, il existe une
$k'$-algèbre lisse $B$ dont $K'$ est le corps des fractions de $B$. Nous
pouvons
alors raisonner comme suit.
Étape 1 : $R \otimes_k k'$ vérifie $(R_n)$ parce que nous avons supposé
$P$ vérifiée pour $k \to R$.
Étape 2 : $R \otimes_k k' \to R \otimes_k k' \otimes_{k'} B$ est un
homomorphisme lisse d'anneaux (Algèbre, Lemme
\ref{algebra-lemma-base-change-smooth}); nous concluons qu'alors
$R \otimes_k k' \otimes_{k'} B$ vérifie $(R_n)$ d'après Algèbre, Lemme
\ref{algebra-lemma-Rk-goes-up} (en utilisant également Algèbre, Lemme
\ref{algebra-lemma-characterize-smooth-over-field}, pour vérifier que les
hypothèses sont satisfaites).
Étape 3 : $R \otimes_k k' \otimes_{k'} K' = R \otimes_k K'$ vérifie
$(R_n)$, car c'est un localisé d'un anneau qui vérifie $(R_n)$.
Étape 4 : enfin, $R \otimes_k K$ vérifie $(R_n)$ par descente de $(R_n)$
le long de l'homomorphisme fidèlement plat d'anneaux
$K \otimes_k R \to K' \otimes_k R$
(Algèbre, Lemme \ref{algebra-lemma-descent-Rk}). Cela démontre (A).
La partie (B) en résulte aussi : un anneau noethérien vérifie $(R_n)$ si et
seulement si tous ses anneaux locaux vérifient $(R_n)$. La partie (C) résulte
d'Algèbre, Lemme \ref{algebra-lemma-Rk-goes-up}, puisque les fibres d'un
homomorphisme régulier sont régulières (nous omettons un petit détail).
La partie (D) résulte d'Algèbre, Lemme
\ref{algebra-lemma-descent-Rk} (nous omettons un petit détail).

\medskip\noindent
Partie (E). Soit $l/k$ une extension algébrique séparable de corps et soit
$l \to R$ un homomorphisme d'anneaux, avec $R$ noethérien. Supposons que
$k \to R$ ait la propriété $P$. Nous devons montrer que $l \to R$ a la
propriété $P$. Soit $l'/l$ une extension finie. Observons d'abord qu'il
existe une sous-extension finie $l/m/k$ et une extension finie $m'/m$
telles que $l' = l \otimes_m m'$. Alors
$R \otimes_l l' = R \otimes_m m'$. Il suffit donc de démontrer que
$m \to R$ a la propriété $P$; autrement dit, nous pouvons supposer que
$l/k$ est finie. Si $l/k$ est finie, alors $l'/k$ est finie et
$$
l' \otimes_l R = (l' \otimes_k R) \otimes_{l \otimes_k l} l
$$
est un localisé (d'après Algèbre, Lemme
\ref{algebra-lemma-separable-algebraic-diagonal}) de l'anneau noethérien
$l' \otimes_k R$, qui vérifie $(R_n)$ par l'hypothèse $P$ sur $k \to R$.
Ainsi $l' \otimes_l R$ vérifie $(R_n)$, comme voulu.
\end{proof}
\section{Anneaux excellents}
\label{section-excellent}

\noindent
Dans cette section, nous étudions la notion d'anneau excellent de Grothendieck.
Pour les définitions des G-anneaux, des anneaux J-2 et des anneaux
universellement caténaires, nous renvoyons à la
Définition \ref{definition-G-ring}, à la Définition \ref{definition-J} et à
Algèbre, Définition \ref{algebra-definition-universally-catenary}.

\begin{definition}
\label{definition-excellent}
Soit $R$ un anneau.
\begin{enumerate}
\item Nous disons que $R$ est {\it quasi excellent} si $R$ est noethérien,
est un G-anneau et est J-2.
\item Nous disons que $R$ est {\it excellent} si $R$ est quasi excellent
et universellement caténaire.
\end{enumerate}
\end{definition}

\noindent
Ainsi, un anneau noethérien est quasi excellent si ses fibres formelles sont
géométriquement régulières et si le lieu singulier de toute algèbre de type
fini sur cet anneau est fermé. Pour qu'un tel anneau soit excellent, nous
exigeons en outre l'existence locale d'une bonne fonction de dimension. Nous
verrons plus loin (Section \ref{section-formally-catenary}) que la propriété
d'être universellement caténaire peut se formuler comme une condition sur les
homomorphismes $R_\mathfrak m \to R_\mathfrak m^\wedge$, où $\mathfrak m$
parcourt les idéaux maximaux de $R$.

\begin{lemma}
\label{lemma-finite-type-over-excellent}
Tout localisé d'une algèbre de type fini sur un anneau (quasi) excellent est
(quasi) excellent.
\end{lemma}

\begin{proof}
Pour les algèbres de type fini, cela résulte des définitions des propriétés
J-2 et universellement caténaire. Pour les G-anneaux, voir la
Proposition \ref{proposition-finite-type-over-G-ring}. Nous omettons la
démonstration du fait que la localisation préserve la (quasi-)excellence.
\end{proof}

\begin{proposition}
\label{proposition-ubiquity-excellent}
Les types d'anneaux suivants sont excellents :
\begin{enumerate}
\item les corps ;
\item les anneaux locaux noethériens complets ;
\item $\mathbf{Z}$,
\item les anneaux de Dedekind dont le corps des fractions est de
caractéristique zéro ;
\item les algèbres de type fini sur l'un quelconque des anneaux précédents.
\end{enumerate}
\end{proposition}

\begin{proof}
Les Propositions \ref{proposition-ubiquity-G-ring} et
\ref{proposition-ubiquity-J-2} montrent que ces anneaux sont des G-anneaux
et sont J-2. Tout anneau de Cohen--Macaulay est universellement caténaire ;
voir Algèbre, Lemme \ref{algebra-lemma-CM-ring-catenary}. En particulier, les
corps, les anneaux de Dedekind et, plus généralement, les anneaux réguliers
sont universellement caténaires. Le théorème de structure de Cohen montre que
les anneaux locaux complets sont universellement caténaires ; voir Algèbre,
Remarque
\ref{algebra-remark-Noetherian-complete-local-ring-universally-catenary}.
\end{proof}

\noindent
Les résultats développés ci-dessus ont quelques conséquences pour les
anneaux de Nagata.

\begin{lemma}
\label{lemma-Nagata-local-ring}
Soit $(A, \mathfrak m)$ un anneau local noethérien. Les assertions suivantes
sont équivalentes :
\begin{enumerate}
\item $A$ est un anneau de Nagata ;
\item les fibres formelles de $A$ sont géométriquement réduites.
\end{enumerate}
\end{lemma}

\begin{proof}
Supposons (2). D'après Algèbre, Lemme
\ref{algebra-lemma-local-nagata-and-analytically-unramified}, il suffit de
montrer que, si $A \to B$ est fini, si $B$ est intègre et si
$\mathfrak m' \subset B$ est un idéal maximal, alors $B_{\mathfrak m'}$ est
analytiquement non ramifié. En combinant les Lemmes
\ref{lemma-formal-fibres-reduced} et
\ref{lemma-check-P-ring-maximal-ideals} avec la
Proposition \ref{proposition-finite-type-over-P-ring}, nous voyons que les
fibres formelles de $B_{\mathfrak m'}$ sont géométriquement réduites. En
particulier, $B_{\mathfrak m'}^\wedge \otimes_B L$ est réduit, où $L$ est le
corps des fractions de $B$. Il s'ensuit que $B_{\mathfrak m'}^\wedge$ est
réduit, c'est-à-dire que $B_{\mathfrak m'}$ est analytiquement non ramifié.

\medskip\noindent
Supposons (1). Soit $\mathfrak q \subset A$ un idéal premier et soit
$K/\kappa(\mathfrak q)$ une extension finie. Nous devons montrer que
$A^\wedge \otimes_A K$ est réduit. Soit
$A/\mathfrak q \subset B \subset K$ un sous-anneau local fini sur $A$ dont le
corps des fractions est $K$. Pour construire $B$, choisissons
$x_1, \ldots, x_n \in K$ qui engendrent $K$ sur $\kappa(\mathfrak q)$ et qui
satisfont à des polynômes unitaires
$P_i(T) = T^{d_i} + a_{i, 1} T^{d_i - 1} + \ldots + a_{i, d_i} = 0$
avec $a_{i, j} \in \mathfrak m$. Soit alors $B$ la sous-$A$-algèbre de $K$
engendrée par $x_1, \ldots, x_n$. (Pour plus de détails, voir la démonstration
d'Algèbre, Lemme
\ref{algebra-lemma-local-nagata-and-analytically-unramified}.) Nous avons
$$
A^\wedge \otimes_A K =
(A^\wedge \otimes_A B)_\mathfrak q =
B^\wedge_\mathfrak q
$$
Puisque $B^\wedge$ est réduit d'après Algèbre, Lemme
\ref{algebra-lemma-local-nagata-and-analytically-unramified}, la démonstration
est achevée.
\end{proof}

\begin{lemma}
\label{lemma-quasi-excellent-nagata}
Tout anneau quasi excellent est un anneau de Nagata.
\end{lemma}

\begin{proof}
Soit $R$ quasi excellent. Puisqu'une algèbre de type fini sur $R$ est quasi
excellente (Lemme \ref{lemma-finite-type-over-excellent}), il suffit de
montrer que tout anneau intègre quasi excellent est N-1 ; voir Algèbre,
Lemme \ref{algebra-lemma-check-universally-japanese}. En appliquant Algèbre,
Lemme \ref{algebra-lemma-characterize-N-1} (et en utilisant qu'un anneau quasi
excellent est J-2), nous nous ramenons à montrer qu'un anneau local intègre
quasi excellent $R$ est N-1. Comme $R \to R^\wedge$ est régulier, le
Lemme \ref{lemma-reduced-goes-up} montre que $R^\wedge$ est réduit. Autrement
dit, $R$ est analytiquement non ramifié. $R$ est donc N-1 d'après Algèbre,
Lemme \ref{algebra-lemma-analytically-unramified-easy}.
\end{proof}

\begin{lemma}
\label{lemma-completion-normal-local-ring}
Soit $(A, \mathfrak m)$ un anneau local noethérien. Si $A$ est normal et si
les fibres formelles de $A$ sont normales (par exemple si $A$ est excellent
ou quasi excellent), alors $A^\wedge$ est normal.
\end{lemma}

\begin{proof}
Cela résulte immédiatement d'Algèbre, Lemme
\ref{algebra-lemma-normal-goes-up-noetherian}.
\end{proof}





\section{Catégories abéliennes de modules}
\label{section-abelian-categories-modules}

\noindent
Soit $R$ un anneau. La catégorie $\text{Mod}_R$ des $R$-modules est une
catégorie abélienne. Voici quelques exemples de sous-catégories abéliennes de
$\text{Mod}_R$ (nous utilisons la terminologie introduite dans Homologie,
Définition \ref{homology-definition-serre-subcategory}, ainsi que dans
Homologie, Lemmes \ref{homology-lemma-characterize-serre-subcategory} et
\ref{homology-lemma-characterize-weak-serre-subcategory}) :
\begin{enumerate}
\item La catégorie des $R$-modules cohérents est une sous-catégorie de Serre
faible de $\text{Mod}_R$. Cela résulte d'Algèbre, Lemme
\ref{algebra-lemma-coherent}.
\item Soit $S \subset R$ une partie multiplicative. La sous-catégorie pleine
formée des $R$-modules $M$ tels que la multiplication par $s \in S$ soit un
isomorphisme de $M$ est une sous-catégorie de Serre de $\text{Mod}_R$. Cela
résulte d'Algèbre, Lemme \ref{algebra-lemma-localization-and-modules}.
\item Soit $I \subset R$ un idéal de type fini. La sous-catégorie pleine des
modules de torsion annulés localement par une puissance de $I$ est une
sous-catégorie de Serre de $\text{Mod}_R$. Voir le
Lemme \ref{lemma-I-power-torsion}.
\item Dans certains textes, un {\it module de torsion} est défini comme un
module $M$ tel que, pour tout $x \in M$, il existe un élément non diviseur de
zéro $f \in R$ tel que $fx = 0$. La sous-catégorie pleine des modules de
torsion est une sous-catégorie de Serre de $\text{Mod}_R$.
\item Si $R$ n'est pas noethérien, alors la catégorie $\text{Mod}^{fg}_R$ des
$R$-modules de type fini n'est {\bf pas} abélienne. En effet, si
$I \subset R$ est un idéal qui n'est pas de type fini, l'homomorphisme
$R \to R/I$ n'a pas de noyau dans $\text{Mod}^{fg}_R$.
\item Si $R$ est noethérien, alors les $R$-modules cohérents coïncident avec
les $R$-modules de type fini (c'est-à-dire finis) ; voir Algèbre, Lemmes
\ref{algebra-lemma-Noetherian-coherent}, \ref{algebra-lemma-coherent-ring} et
\ref{algebra-lemma-Noetherian-finite-type-is-finite-presentation}. Ainsi,
$\text{Mod}^{fg}_R$ est abélienne d'après (1) ci-dessus et, en fait, dans ce
cas, la catégorie $\text{Mod}_R^{fg}$ est une sous-catégorie de Serre (forte)
de $\text{Mod}_R$.
\end{enumerate}




\section{Groupes abéliens injectifs}
\label{section-abelian-groups}

\noindent
Dans cette section, nous montrons que la catégorie des groupes abéliens possède
assez d'injectifs. Rappelons qu'un groupe abélien $M$ est {\it divisible} si
et seulement si, pour tout $x \in M$ et tout $n \in \mathbf{N}$, il existe
$y \in M$ tel que $n y = x$.

\begin{lemma}
\label{lemma-injective-abelian}
Un groupe abélien $J$ est un objet injectif de la catégorie des groupes
abéliens si et seulement si $J$ est divisible.
\end{lemma}

\begin{proof}
Supposons que $J$ ne soit pas divisible. Il existe alors $x \in J$ et
$n \in \mathbf{N}$ tels qu'il n'existe aucun $y \in J$ vérifiant $n y = x$.
L'homomorphisme $\mathbf{Z} \to J$, $m \mapsto mx$, ne se prolonge donc pas
à $\frac{1}{n}\mathbf{Z} \supset \mathbf{Z}$. Ainsi, $J$ n'est pas injectif.

\medskip\noindent
Soient $A \subset B$ des groupes abéliens. Supposons que $J$ soit un groupe
abélien divisible. Soit $\varphi : A \to J$ un homomorphisme. Considérons
l'ensemble des homomorphismes $\varphi' : A' \to J$ tels que
$A \subset A' \subset B$ et $\varphi'|_A = \varphi$. Définissons
$(A', \varphi') \geq (A'', \varphi'')$ si et seulement si
$A' \supset A''$ et $\varphi'|_{A''} = \varphi''$. Si
$(A_i, \varphi_i)_{i \in I}$ est une famille totalement ordonnée de telles
paires, nous obtenons un homomorphisme $\bigcup_{i \in I} A_i \to J$ qui
envoie $a \in A_i$ sur $\varphi_i(a)$. Le lemme de Zorn s'applique donc. Pour
conclure, nous devons montrer que, si la paire $(A', \varphi')$ est maximale,
alors $A' = B$. Autrement dit, il suffit de montrer que, pour tout sous-groupe
$A \subset B$ tel que $A \not = B$ et tout $\varphi : A \to J$, on peut
trouver $\varphi' : A' \to J$ avec $A \subset A' \subset B$, de sorte que
(a) l'inclusion $A \subset A'$ soit stricte et (b) l'homomorphisme $\varphi'$
prolonge $\varphi$.

\medskip\noindent
Pour le démontrer, choisissons $x \in B$ avec $x \not \in A$. S'il n'existe
aucun $n\in \mathbf{N}$ tel que $nx \in A$, alors
$A \oplus \mathbf{Z} \cong A + \mathbf{Z}x$. Nous pouvons donc prolonger
$\varphi$ à $A' = A + \mathbf{Z}x$ en utilisant $\varphi$ sur $A$ et, par
exemple, en envoyant $x$ sur zéro. S'il existe au contraire
$n \in \mathbf{N}$ tel que $nx \in A$, choisissons $n$ minimal avec cette
propriété. Soit $z \in J$ tel que $nz = \varphi(nx)$. Définissons un
homomorphisme $\tilde\varphi : A \oplus \mathbf{Z} \to J$ par
$(a, m) \mapsto \varphi(a) + mz$. Par le choix de $z$, le noyau de
$\tilde \varphi$ contient celui de l'homomorphisme
$A \oplus \mathbf{Z} \to B$, $(a, m) \mapsto a + mx$. Ainsi,
$\tilde \varphi$ se factorise par l'image $A' = A + \mathbf{Z}x$, et le
facteur obtenu prolonge $\varphi$.
\end{proof}

\noindent
Ce lemme permet de montrer que tout groupe abélien se plonge dans un groupe
abélien injectif. Il s'agit cependant d'un cas particulier du résultat de la
section suivante.





\section{Modules injectifs}
\label{section-injectives-modules}

\noindent
Quelques lemmes sur les modules injectifs.

\begin{definition}
\label{definition-projective}
Soit $R$ un anneau. Un $R$-module $J$ est {\it injectif} si et seulement si
le foncteur $\Hom_R(-, J) : \text{Mod}_R \to \text{Mod}_R$ est exact.
\end{definition}

\noindent
Le foncteur $\Hom_R(- , M)$ est exact à gauche pour tout $R$-module $M$ ; voir
Algèbre, Lemme \ref{algebra-lemma-hom-exact}. Ainsi, la condition que $J$ soit
injectif signifie précisément que, pour toute injection de $R$-modules
$M \to M'$, l'application $\Hom_R(M', J) \to \Hom_R(M, J)$ est surjective.

\medskip\noindent
Avant de reformuler cela en termes de modules ${Ext}$, étudions la relation
entre $\Ext^1_R(M, N)$ et les extensions au sens d'Homologie,
Section \ref{homology-section-extensions}.

\begin{lemma}
\label{lemma-relation-ext-ext}
Soit $R$ un anneau. Soit $\mathcal{A}$ la catégorie abélienne des
$R$-modules. Il existe un isomorphisme canonique
$\Ext_\mathcal{A}(M, N) = \Ext^1_R(M, N)$ compatible avec les suites exactes
longues d'Algèbre, Lemmes \ref{algebra-lemma-long-exact-seq-ext} et
\ref{algebra-lemma-reverse-long-exact-seq-ext}, ainsi qu'avec les suites
exactes à $6$ termes d'Homologie,
Lemme \ref{homology-lemma-six-term-sequence-ext}.
\end{lemma}

\begin{proof}
Démonstration omise.
\end{proof}

\begin{lemma}
\label{lemma-characterize-injective}
Soit $R$ un anneau. Soit $J$ un $R$-module. Les assertions suivantes sont
équivalentes :
\begin{enumerate}
\item $J$ est injectif ;
\item $\Ext^1_R(M, J) = 0$ pour tout $R$-module $M$.
\end{enumerate}
\end{lemma}

\begin{proof}
Soit $0 \to M'' \to M' \to M \to 0$ une suite exacte courte de $R$-modules.
Considérons la suite exacte longue
$$
\begin{matrix}
0
\to \Hom_R(M, J)
\to \Hom_R(M', J)
\to \Hom_R(M'', J)
\\
\phantom{0\ }
\to \Ext^1_R(M, J)
\to \Ext^1_R(M', J)
\to \Ext^1_R(M'', J)
\to \ldots
\end{matrix}
$$
d'Algèbre, Lemme \ref{algebra-lemma-reverse-long-exact-seq-ext}. Nous voyons
donc que (2) implique (1). Réciproquement, si $J$ est injectif, alors le
groupe $\Ext$ est nul d'après Homologie,
Lemme \ref{homology-lemma-characterize-injectives}, et le
Lemme \ref{lemma-relation-ext-ext}.
\end{proof}

\begin{lemma}
\label{lemma-characterize-injective-bis}
Soit $R$ un anneau. Soit $J$ un $R$-module. Les assertions suivantes sont
équivalentes :
\begin{enumerate}
\item $J$ est injectif ;
\item $\Ext^1_R(R/I, J) = 0$ pour tout idéal $I \subset R$ ;
\item pour tout idéal $I \subset R$ et tout homomorphisme de modules
$I \to J$, il existe un prolongement $R \to J$.
\end{enumerate}
\end{lemma}

\begin{proof}
Si $I \subset R$ est un idéal, alors la suite exacte courte
$0 \to I \to R \to R/I \to 0$ donne une suite exacte
$$
\Hom_R(R, J) \to
\Hom_R(I, J) \to
\Ext^1_R(R/I, J) \to 0
$$
d'après Algèbre, Lemme \ref{algebra-lemma-reverse-long-exact-seq-ext}, et le
fait que $\Ext^1_R(R, J) = 0$ puisque $R$ est projectif (Algèbre,
Lemme \ref{algebra-lemma-characterize-projective}). Ainsi, (2) et (3) sont
équivalentes. Nous allons montrer dans cette démonstration que
(1) $\Leftrightarrow$ (3), résultat connu sous le nom de critère de Baer.

\medskip\noindent
Supposons (1). Étant donné un homomorphisme de modules $I \to J$ comme en
(3), nous obtenons le prolongement $R \to J$, car l'application
$\Hom_R(R, J) \to \Hom_R(I, J)$ est surjective par définition.

\medskip\noindent
Supposons (3). Soit $M \subset N$ une inclusion de $R$-modules. Soit
$\varphi : M \to J$ un homomorphisme. Nous allons montrer que $\varphi$ se
prolonge à $N$, ce qui achèvera la démonstration. Considérons l'ensemble des
homomorphismes $\varphi' : M' \to J$ tels que $M \subset M' \subset N$ et
$\varphi'|_M = \varphi$. Définissons
$(M', \varphi') \geq (M'', \varphi'')$ si et seulement si
$M' \supset M''$ et $\varphi'|_{M''} = \varphi''$. Si
$(M_i, \varphi_i)_{i \in I}$ est une famille totalement ordonnée de telles
paires, nous obtenons un homomorphisme $\bigcup_{i \in I} M_i \to J$ qui
envoie $a \in M_i$ sur $\varphi_i(a)$. Le lemme de Zorn s'applique donc. Pour
conclure, nous devons montrer que, si la paire $(M', \varphi')$ est maximale,
alors $M' = N$. Autrement dit, il suffit de montrer que, pour tout sous-module
$M \subset N$ tel que $M \not = N$ et tout $\varphi : M \to J$, on peut
trouver $\varphi' : M' \to J$ avec $M \subset M' \subset N$, de sorte que
(a) l'inclusion $M \subset M'$ soit stricte et (b) l'homomorphisme $\varphi'$
prolonge $\varphi$.

\medskip\noindent
Pour le démontrer, choisissons $x \in N$ avec $x \not \in M$. Posons
$I = \{f \in R \mid fx \in M\}$. C'est un idéal de $R$. Définissons un
homomorphisme $\psi : I \to J$ par $f \mapsto \varphi(fx)$. Prolongeons-le
en un homomorphisme $\tilde\psi : R \to J$, ce qui est possible par
l'hypothèse (3). Par le choix de $I$, le noyau de
$M \oplus R \to J$, $(y, f) \mapsto \varphi(y) + \tilde\psi(f)$, contient le
noyau de l'homomorphisme $M \oplus R \to N$, $(y, f) \mapsto y + fx$. Cet
homomorphisme se factorise donc par l'image $M' = M + Rx$, et le facteur ainsi
obtenu prolonge l'homomorphisme donné comme souhaité.
\end{proof}

\noindent
Dans la suite de cette section, nous démontrons qu'il existe assez de modules
injectifs sur un anneau $R$. Nous partons du fait que
$\mathbf{Q}/\mathbf{Z}$ est un groupe abélien injectif. Cela résulte du
Lemme \ref{lemma-injective-abelian}.

\begin{definition}
\label{definition-simple-functors}
Soit $R$ un anneau.
\begin{enumerate}
\item Pour tout $R$-module $M$ sur $R$, nous notons
$M^\vee = \Hom(M, \mathbf{Q}/\mathbf{Z})$, muni de sa structure naturelle de
$R$-module. Nous considérons {\it $M \mapsto M^\vee$} comme un foncteur
contravariant de la catégorie des $R$-modules vers elle-même.
\item Pour tout $R$-module $M$, nous notons
$$
F(M) = \bigoplus\nolimits_{m \in M} R[m]
$$
le {\it module libre} dont la base est formée des éléments $[m]$ avec
$m \in M$. Nous désignons par $F(M)\to M$,
$\sum f_i [m_i] \mapsto \sum f_i m_i$, la surjection naturelle de
$R$-modules. Nous considérons $M \mapsto (F(M) \to M)$ comme un foncteur de
la catégorie des $R$-modules vers la catégorie des flèches de $R$-modules.
\end{enumerate}
\end{definition}

\begin{lemma}
\label{lemma-vee-exact}
Soit $R$ un anneau. Le foncteur $M \mapsto M^\vee$ est exact.
\end{lemma}

\begin{proof}
Cela tient au fait que $\mathbf{Q}/\mathbf{Z}$ est un groupe abélien injectif
d'après le Lemme \ref{lemma-injective-abelian}.
\end{proof}

\noindent
Il existe une application canonique $ev : M \to (M^\vee)^\vee$, donnée par
l'évaluation : pour $x \in M$, nous définissons
$ev(x) \in (M^\vee)^\vee = \Hom(M^\vee, \mathbf{Q}/\mathbf{Z})$ comme
l'application $\varphi \mapsto \varphi(x)$.

\begin{lemma}
\label{lemma-ev-injective}
Pour tout $R$-module $M$, l'application d'évaluation
$ev : M \to (M^\vee)^\vee$ est injective.
\end{lemma}

\begin{proof}
Cela se vérifie en utilisant que $\mathbf{Q}/\mathbf{Z}$ est un groupe
abélien injectif. En effet, si $x \in M$ est non nul, soit $M' \subset M$ le
groupe cyclique qu'il engendre. Il existe un homomorphisme non nul
$M' \to \mathbf{Q}/\mathbf{Z}$ qui, nécessairement, n'annule pas $x$. Il se
prolonge en un homomorphisme $\varphi : M \to \mathbf{Q}/\mathbf{Z}$, et
alors $ev(x)(\varphi) = \varphi(x) \not = 0$.
\end{proof}

\noindent
La surjection canonique $F(M) \to M$ de $R$-modules donne, comme ci-dessus,
une injection canonique de $R$-modules
$$
(M^\vee)^\vee \longrightarrow (F(M^\vee))^\vee.
$$
Posons $J(M) = (F(M^\vee))^\vee$. La composée de $ev$ avec l'application
affichée donne $M \to J(M)$, fonctoriellement en $M$.

\begin{lemma}
\label{lemma-JM-injective}
Soit $R$ un anneau. Pour tout $R$-module $M$, le $R$-module $J(M)$ est
injectif.
\end{lemma}

\begin{proof}
Remarquons que $J(M) \cong \prod_{\varphi \in M^\vee} R^\vee$ comme
$R$-module. Comme un produit de modules injectifs est injectif, il suffit de
montrer que $R^\vee$ est injectif. Pour cela, nous utilisons les égalités
$$
\Hom_R(N, R^\vee) =
\Hom_R(N, \Hom_{\mathbf{Z}}(R, \mathbf{Q}/\mathbf{Z})) =
N^\vee
$$
et le fait que $(-)^\vee$ est un foncteur exact d'après le
Lemme \ref{lemma-vee-exact}.
\end{proof}

\begin{lemma}
\label{lemma-injectives-modules}
Soit $R$ un anneau. La construction ci-dessus définit un foncteur covariant
$M \mapsto (M \to J(M))$ de la catégorie des $R$-modules vers la catégorie
des flèches de $R$-modules, tel que, pour tout module $M$, l'application
obtenue $M \to J(M)$ soit une injection de $M$ dans un $R$-module injectif
$J(M)$.
\end{lemma}

\begin{proof}
Cela résulte de ce qui précède.
\end{proof}

\noindent
En particulier, à tout homomorphisme de $R$-modules $M \to N$ est associé un
homomorphisme $J(M) \to J(N)$ qui rend commutatif le diagramme suivant :
$$
\xymatrix{
M \ar[d] \ar[r] & N \ar[d] \\
J(M) \ar[r] & J(N) }
$$
C'est le type de construction que nous souhaiterions obtenir en général.
Dans Homologie, Section \ref{homology-section-injectives}, nous avons introduit
une terminologie pour l'exprimer : nous disons que la catégorie des
$R$-modules possède des plongements injectifs fonctoriels.











\section{Catégories dérivées de modules}
\label{section-derived-modules}

\noindent
Dans cette section, nous rassemblons quelques généralités sur la catégorie
dérivée des modules sur un anneau.

\medskip\noindent
Soit $A$ un anneau. La catégorie des $A$-modules est notée $\text{Mod}_A$.
Nous utiliserons le symbole $K(A)$ pour désigner la catégorie homotopique des
complexes de $A$-modules ; autrement dit, nous posons
$K(A) = K(\text{Mod}_A)$ comme catégorie ; voir Catégories dérivées,
Section \ref{derived-section-homotopy}. Les versions bornées sont $K^+(A)$,
$K^-(A)$ et $K^b(A)$. Nous considérons $K(A)$ comme une catégorie triangulée
au sens de Catégories dérivées,
Section \ref{derived-section-homotopy-triangulated}. La {\it catégorie
dérivée} de $A$, notée $D(A)$, est la catégorie obtenue à partir de $K(A)$ en
inversant les quasi-isomorphismes ; autrement dit, nous posons
$D(A) = D(\text{Mod}_A)$ ; voir Catégories dérivées,
Section \ref{derived-section-derived-categories}\footnote{Voir aussi
Injectifs, Remarque \ref{injectives-remark-existence-D}.}. Les versions
bornées sont $D^+(A)$, $D^-(A)$ et $D^b(A)$.

\medskip\noindent
Soit $A$ un anneau. La catégorie des $A$-modules possède des produits, et les
produits y sont exacts. La catégorie des $A$-modules possède assez d'injectifs d'après le
Lemme \ref{lemma-injectives-modules}. Par conséquent, tout complexe de
$A$-modules est quasi isomorphe à un complexe K-injectif (Catégories dérivées,
Lemme \ref{derived-lemma-enough-K-injectives-Ab4-star}). Il s'ensuit que
$D(A)$ possède des produits dénombrables (Catégories dérivées,
Lemme \ref{derived-lemma-products}) et, en fait, des produits arbitraires
(Injectifs, Lemme \ref{injectives-lemma-derived-products}). Il en résulte que
tout système projectif d'objets de $D(A)$ admet une limite dérivée (bien
définie à isomorphisme près) ; voir Catégories dérivées,
Section \ref{derived-section-derived-limit}.

\begin{lemma}
\label{lemma-K-injective-flat}
Soit $R \to S$ un homomorphisme plat d'anneaux. Si $I^\bullet$ est un
complexe K-injectif de $S$-modules, alors $I^\bullet$ est K-injectif comme
complexe de $R$-modules.
\end{lemma}

\begin{proof}
Cela résulte de l'égalité
$\Hom_{K(R)}(M^\bullet, I^\bullet) =
\Hom_{K(S)}(M^\bullet \otimes_R S, I^\bullet)$
d'après Algèbre, Lemme \ref{algebra-lemma-adjoint-tensor-restrict}, et du fait
que le produit tensoriel par $S$ est exact.
\end{proof}

\begin{lemma}
\label{lemma-K-injective-epimorphism}
Soit $R \to S$ un épimorphisme d'anneaux. Soit $I^\bullet$ un complexe de
$S$-modules. Si $I^\bullet$ est K-injectif comme complexe de $R$-modules,
alors $I^\bullet$ est un complexe K-injectif de $S$-modules.
\end{lemma}

\begin{proof}
Cela résulte de l'égalité
$\Hom_{K(R)}(N^\bullet, I^\bullet) =
\Hom_{K(S)}(N^\bullet, I^\bullet)$ pour tout complexe de $S$-modules
$N^\bullet$ ; voir Algèbre, Lemme
\ref{algebra-lemma-epimorphism-modules}.
\end{proof}

\begin{lemma}
\label{lemma-hom-K-injective}
Soit $A \to B$ un homomorphisme d'anneaux. Si $I^\bullet$ est un complexe
K-injectif de $A$-modules, alors $\Hom_A(B, I^\bullet)$ est un complexe
K-injectif de $B$-modules.
\end{lemma}

\begin{proof}
Cela résulte de l'égalité
$\Hom_{K(B)}(N^\bullet, \Hom_A(B, I^\bullet)) =
\Hom_{K(A)}(N^\bullet, I^\bullet)$
d'après Algèbre, Lemme \ref{algebra-lemma-adjoint-hom-restrict}.
\end{proof}










\section{Calcul de Tor}
\label{section-computing-tor}

\noindent
Soit $R$ un anneau. On note $D(R)$ la catégorie dérivée de la
catégorie abélienne $\text{Mod}_R$ des $R$-modules. Remarquons que
$\text{Mod}_R$ possède assez d'objets projectifs, puisque tout
$R$-module libre est projectif. On peut donc définir les foncteurs dérivés
à gauche de tout foncteur additif de $\text{Mod}_R$ vers une catégorie
abélienne quelconque.

\medskip\noindent
Cela s'applique en particulier au foncteur
$ - \otimes_R M : \text{Mod}_R \to \text{Mod}_R$
dont les foncteurs dérivés à gauche sont les foncteurs
$\text{Tor}_i^R(-, M)$ ; voir Algèbre, Section
\ref{algebra-section-tor}. Il existe aussi un foncteur dérivé total à gauche
\begin{equation}
\label{equation-derived-tensor-module}
-\otimes_R^{\mathbf{L}} M :
D^{-}(R)
\longrightarrow
D^{-}(R)
\end{equation}
noté $-\otimes_R^{\mathbf{L}} M$. Ses satellites sont les modules
Tor, c'est-à-dire que l'on a
$$
H^{-p}(N \otimes_R^{\mathbf{L}} M) = \text{Tor}_p^R(N, M).
$$

\medskip\noindent
Une situation particulière se présente lorsque l'on considère le produit
tensoriel par une $R$-algèbre $A$. Dans ce cas, on regarde
$- \otimes_R A$ comme un foncteur de $\text{Mod}_R$ vers
$\text{Mod}_A$. On obtient ainsi le foncteur dérivé total à gauche
\begin{equation}
\label{equation-derived-tensor-algebra}
-\otimes_R^{\mathbf{L}} A :
D^{-}(R)
\longrightarrow
D^{-}(A)
\end{equation}
noté $-\otimes_R^{\mathbf{L}} A$. Ses satellites sont les groupes
Tor, c'est-à-dire que l'on a
$$
H^{-p}(N \otimes_R^{\mathbf{L}} A) = \text{Tor}_p^R(N, A).
$$
En particulier, ces groupes Tor sont naturellement munis d'une structure
de $A$-module.

\medskip\noindent
Dans les quelques sections suivantes, nous généraliserons le contenu de
cette section aux complexes non bornés.






\section{Produits tensoriels de complexes}
\label{section-symmetric-monoidal}

\noindent
Soit $R$ un anneau. La catégorie $\text{Comp}(R)$ des complexes de
$R$-modules possède une structure monoïdale symétrique. Supposons donnés deux
complexes de $R$-modules, $L^\bullet$ et $M^\bullet$. À l'aide de
Homologie, Exemple
\ref{homology-example-double-complex-as-tensor-product-of} et de
Homologie, Définition \ref{homology-definition-associated-simple-complex},
on obtient un troisième complexe de $R$-modules, à savoir
$$
\text{Tot}(L^\bullet \otimes_R M^\bullet)
$$
Cette construction est manifestement fonctorielle en $L^\bullet$ et en
$M^\bullet$. La contrainte d'associativité est l'isomorphisme canonique de
complexes
$$
\text{Tot}(\text{Tot}(K^\bullet \otimes_R L^\bullet) \otimes_R M^\bullet)
\longrightarrow
\text{Tot}(K^\bullet \otimes_R \text{Tot}(L^\bullet \otimes_R M^\bullet))
$$
construit dans Homologie, Remarque \ref{homology-remark-triple-complex} à
partir du complexe triple
$K^\bullet \otimes_R L^\bullet \otimes_R M^\bullet$. La contrainte de
commutativité est l'isomorphisme canonique
$$
\text{Tot}(L^\bullet \otimes_R M^\bullet) \to
\text{Tot}(M^\bullet \otimes_R L^\bullet)
$$
qui applique le signe $(-1)^{pq}$ au facteur direct
$L^p \otimes_R M^q$. Pour voir qu'il s'agit d'un morphisme de complexes,
on calcule, pour $x \in L^p$ et $y \in M^q$, que
$$
\text{d}(x \otimes y) =
\text{d}_L(x) \otimes y + (-1)^px \otimes \text{d}_M(y)
$$
Notre règle envoie le membre de droite sur
$$
(-1)^{(p + 1)q}y \otimes \text{d}_L(x) +
(-1)^{p + p(q + 1)} \text{d}_M(y) \otimes x
$$
D'autre part, on a
$$
\text{d}((-1)^{pq}y \otimes x) =
(-1)^{pq} \text{d}_M(y) \otimes x +
(-1)^{pq + q} y \otimes \text{d}_L(x)
$$
Ces deux expressions coïncident immédiatement, comme voulu.

\begin{lemma}
\label{lemma-symmetric-monoidal-cat-complexes}
Soit $R$ un anneau. La catégorie $\text{Comp}(R)$ des complexes de
$R$-modules, munie du foncteur
$(L^\bullet, M^\bullet) \mapsto \text{Tot}(L^\bullet \otimes_R M^\bullet)$
et des contraintes d'associativité et de commutativité ci-dessus, est une
catégorie monoïdale symétrique.
\end{lemma}

\begin{proof}
Omis. Indications : pour unité $\mathbf{1}$, on prend le complexe égal à
$R$ en degré $0$ et nul dans les autres degrés, avec les isomorphismes évidents
$\text{Tot}(\mathbf{1} \otimes_R M^\bullet) = M^\bullet$ et
$\text{Tot}(K^\bullet \otimes_R \mathbf{1}) = K^\bullet$.
Pour démontrer le lemme, il faut vérifier la commutativité de divers
diagrammes ; voir Catégories, Définitions
\ref{categories-definition-monoidal-category} et
\ref{categories-definition-symmetric-monoidal-category}.
Dans chaque cas, les vérifications sont immédiates.
\end{proof}

\begin{lemma}
\label{lemma-derived-tor-homotopy}
Soit $R$ un anneau, soit $P^\bullet$ un complexe de $R$-modules et soient
$\alpha, \beta : L^\bullet \to M^\bullet$ des morphismes homotopes de
complexes. Alors $\alpha$ et $\beta$ induisent des morphismes homotopes
$$
\text{Tot}(\alpha \otimes \text{id}_P),
\text{Tot}(\beta \otimes \text{id}_P) :
\text{Tot}(L^\bullet \otimes_R P^\bullet)
\longrightarrow
\text{Tot}(M^\bullet \otimes_R P^\bullet).
$$
En particulier, la construction
$L^\bullet \mapsto \text{Tot}(L^\bullet \otimes_R P^\bullet)$
définit un endofoncteur de la catégorie d'homotopie des complexes.
\end{lemma}

\begin{proof}
Écrivons $\alpha = \beta + dh + hd$ pour une homotopie $h$ donnée par
$h^n : L^n \to M^{n - 1}$. Posons
$$
H^n = \bigoplus\nolimits_{a + b = n} h^a \otimes \text{id}_{P^b} :
\bigoplus\nolimits_{a + b = n} L^a \otimes_R P^b
\longrightarrow
\bigoplus\nolimits_{a + b = n} M^{a - 1} \otimes_R P^b
$$
Un calcul immédiat montre alors que
$$
\text{Tot}(\alpha \otimes \text{id}_P) =
\text{Tot}(\beta \otimes \text{id}_P) + dH + Hd
$$
comme morphismes $\text{Tot}(L^\bullet \otimes_R P^\bullet) \to
\text{Tot}(M^\bullet \otimes_R P^\bullet)$.
\end{proof}

\begin{lemma}
\label{lemma-symmetric-monoidal-homotopy-category}
Soit $R$ un anneau. La catégorie d'homotopie $K(R)$ des complexes de
$R$-modules, munie du foncteur
$(L^\bullet, M^\bullet) \mapsto \text{Tot}(L^\bullet \otimes_R M^\bullet)$
et des contraintes d'associativité et de commutativité ci-dessus, est une
catégorie monoïdale symétrique.
\end{lemma}

\begin{proof}
Cela résulte des Lemmes \ref{lemma-symmetric-monoidal-cat-complexes} et
\ref{lemma-derived-tor-homotopy}. Les détails sont omis.
\end{proof}

\begin{lemma}
\label{lemma-derived-tor-exact}
Soit $R$ un anneau et soit $P^\bullet$ un complexe de $R$-modules. Les
foncteurs
$$
K(R) \longrightarrow K(R), \quad
L^\bullet \longmapsto \text{Tot}(P^\bullet \otimes_R L^\bullet)
$$
and
$$
K(R) \longrightarrow K(R), \quad
L^\bullet \longmapsto \text{Tot}(L^\bullet \otimes_R P^\bullet)
$$
sont des foncteurs exacts de catégories triangulées.
\end{lemma}

\begin{proof}
Cela résulte de Catégories dérivées, Remarque
\ref{derived-remark-double-complex-as-tensor-product-of}.
\end{proof}




\section{Produit tensoriel dérivé}
\label{section-derived-tensor-product}

\noindent
On peut construire le produit tensoriel dérivé dans une plus grande
généralité. En fait, les hypothèses de bornitude ne sont pas nécessaires,
pourvu que l'on choisisse des résolutions K-plates.

\begin{definition}
\label{definition-K-flat}
Soit $R$ un anneau. Un complexe $K^\bullet$ est dit {\it K-plat} si, pour
tout complexe acyclique $M^\bullet$, le complexe total
$\text{Tot}(M^\bullet \otimes_R K^\bullet)$ est acyclique.
\end{definition}

\begin{lemma}
\label{lemma-K-flat-quasi-isomorphism}
Soit $R$ un anneau et soit $K^\bullet$ un complexe K-plat. Alors le
foncteur
$$
K(R) \longrightarrow K(R), \quad
L^\bullet \longmapsto \text{Tot}(L^\bullet \otimes_R K^\bullet)
$$
envoie les quasi-isomorphismes sur des quasi-isomorphismes.
\end{lemma}

\begin{proof}
Cela résulte du Lemme \ref{lemma-derived-tor-exact} et du fait que les
quasi-isomorphismes de $K(R)$ sont caractérisés par l'acyclicité de leurs
cônes.
\end{proof}

\begin{lemma}
\label{lemma-base-change-K-flat}
Soit $R \to R'$ un homomorphisme d'anneaux. Si $K^\bullet$ est un complexe
K-plat de $R$-modules, alors $K^\bullet \otimes_R R'$ est un complexe
K-plat de $R'$-modules.
\end{lemma}

\begin{proof}
Cela résulte des définitions et du fait que
$(K^\bullet \otimes_R R') \otimes_{R'} L^\bullet =
K^\bullet \otimes_R L^\bullet$ pour tout complexe $L^\bullet$ de
$R'$-modules.
\end{proof}

\begin{lemma}
\label{lemma-tensor-product-K-flat}
Soit $R$ un anneau. Si $K^\bullet$ et $L^\bullet$ sont des complexes
K-plats de $R$-modules, alors
$\text{Tot}(K^\bullet \otimes_R L^\bullet)$ est un complexe K-plat de
$R$-modules.
\end{lemma}

\begin{proof}
Cela résulte de l'isomorphisme
$$
\text{Tot}(M^\bullet \otimes_R \text{Tot}(K^\bullet \otimes_R L^\bullet))
=
\text{Tot}(\text{Tot}(M^\bullet \otimes_R K^\bullet) \otimes_R L^\bullet)
$$
et de la définition.
\end{proof}

\begin{lemma}
\label{lemma-K-flat-two-out-of-three}
Soit $R$ un anneau et soit $(K_1^\bullet, K_2^\bullet, K_3^\bullet)$ un
triangle distingué de $K(R)$. Si deux des trois complexes $K_i^\bullet$
sont K-plats, le troisième l'est aussi.
\end{lemma}

\begin{proof}
Cela résulte du Lemme \ref{lemma-derived-tor-exact} et du fait que, dans
un triangle distingué de $K(R)$, si deux objets sur trois sont acycliques,
le troisième l'est aussi.
\end{proof}

\begin{lemma}
\label{lemma-K-flat-two-out-of-three-ses}
Soit $R$ un anneau. Soit
$0 \to K_1^\bullet \to K_2^\bullet \to K_3^\bullet \to 0$ une suite
exacte courte de complexes. Si $K_3^n$ est plat pour tout
$n \in \mathbf{Z}$ et si deux des trois complexes $K_i^\bullet$ sont
K-plats, le troisième l'est aussi.
\end{lemma}

\begin{proof}
Soit $L^\bullet$ un complexe de $R$-modules. Alors
$$
0 \to
\text{Tot}(L^\bullet \otimes_R K_1^\bullet) \to
\text{Tot}(L^\bullet \otimes_R K_2^\bullet) \to
\text{Tot}(L^\bullet \otimes_R K_3^\bullet) \to 0
$$
est une suite exacte courte de complexes. En effet, pour tous $n, m$, la
suite de modules
$0 \to L^n \otimes_R K_1^m \to
L^n \otimes_R K_2^m \to
L^n \otimes_R K_3^m \to 0$
est exacte d'après Algèbre, Lemme \ref{algebra-lemma-flat-tor-zero}, et la
suite de complexes est somme directe de ces suites. Le lemme en résulte,
ainsi que du fait que, dans une suite exacte courte de complexes, si deux
complexes sur trois sont acycliques, le troisième l'est aussi.
\end{proof}

\begin{lemma}
\label{lemma-derived-tor-quasi-isomorphism}
Soit $R$ un anneau. Soit $P^\bullet$ un complexe de $R$-modules plats, borné
supérieurement. Alors $P^\bullet$ est K-plat.
\end{lemma}

\begin{proof}
Soit $L^\bullet$ un complexe acyclique de $R$-modules et soit
$\xi \in H^n(\text{Tot}(L^\bullet \otimes_R P^\bullet))$. Il faut montrer
que $\xi = 0$. Comme $\text{Tot}^n(L^\bullet \otimes_R P^\bullet)$ est une
somme directe de termes $L^a \otimes_R P^b$, on voit que $\xi$ provient
d'un élément de
$H^n(\text{Tot}(\tau_{\leq m}L^\bullet \otimes_R P^\bullet))$ pour un
$m \in \mathbf{Z}$. Puisque $\tau_{\leq m}L^\bullet$ est également
acyclique, on peut remplacer $L^\bullet$ par
$\tau_{\leq m}L^\bullet$. On peut donc supposer $L^\bullet$ borné
supérieurement. Dans ce cas, la suite spectrale de Homologie, Lemme
\ref{homology-lemma-first-quadrant-ss}, vérifie
$$
{}'E_1^{p, q} = H^p(L^\bullet \otimes_R P^q)
$$
qui est nul, puisque $P^q$ est plat et $L^\bullet$ acyclique. Ainsi
$H^*(\text{Tot}(L^\bullet \otimes_R P^\bullet)) = 0$.
\end{proof}

\noindent
Dans le lemme suivant, la colimite d'un système de complexes désigne la
colimite calculée terme à terme.

\begin{lemma}
\label{lemma-colimit-K-flat}
Soit $R$ un anneau et soit
$K_1^\bullet \to K_2^\bullet \to \ldots$ un système de complexes K-plats.
Alors $\colim_i K_i^\bullet$ est K-plat. Plus généralement, tout colimite
filtrant de complexes K-plats est K-plat.
\end{lemma}

\begin{proof}
Comme les colimites sont calculés terme à terme, on a
$$
\colim_i \text{Tot}(M^\bullet \otimes_R K_i^\bullet) =
\text{Tot}(M^\bullet \otimes_R \colim_i K_i^\bullet)
$$
d'après Algèbre, Lemme
\ref{algebra-lemma-tensor-products-commute-with-limits}. Le lemme résulte
donc de l'exactitude des colimites filtrantes ; voir Algèbre, Lemme
\ref{algebra-lemma-directed-colimit-exact}.
\end{proof}

\begin{lemma}
\label{lemma-universally-acyclic-K-flat}
Soit $R$ un anneau et soit $K^\bullet$ un complexe de $R$-modules. Si
$K^\bullet \otimes_R M$ est acyclique pour tout $R$-module $M$ de
présentation finie, alors $K^\bullet$ est K-plat.
\end{lemma}

\begin{proof}
Nous utiliserons à plusieurs reprises le fait que le produit tensoriel
commute aux colimites (Algèbre, Lemme
\ref{algebra-lemma-tensor-products-commute-with-limits}). On voit ainsi que
$K^\bullet \otimes_R M$ est acyclique pour tout $R$-module $M$, puisque
tout $R$-module est une colimite filtrante de $R$-modules $M$ de présentation
finie ; voir Algèbre, Lemme \ref{algebra-lemma-module-colimit-fp}. Soit
$M^\bullet$ un complexe acyclique de $R$-modules. Il faut montrer que
$\text{Tot}(M^\bullet \otimes_R K^\bullet)$ est acyclique. Comme
$M^\bullet = \colim \tau_{\leq n} M^\bullet$ (colimite terme à terme), on a
$$
\text{Tot}(M^\bullet \otimes_R K^\bullet) =
\colim \text{Tot}(\tau_{\leq n} M^\bullet \otimes_R K^\bullet)
$$
où les troncations sont celles de Homologie, Section
\ref{homology-section-truncations}. Comme les colimites filtrantes sont
exacts (Algèbre, Lemme \ref{algebra-lemma-directed-colimit-exact}), on peut
remplacer $M^\bullet$ par $\tau_{\leq n}M^\bullet$ et supposer que
$M^\bullet$ est borné supérieurement. Dans ce cas, on peut écrire
$M^\bullet = \colim \sigma_{\geq -n} M^\bullet$
où les complexes $\sigma_{\geq -n} M^\bullet$ sont bornés, mais ne sont
plus nécessairement acycliques. En raisonnant comme ci-dessus, on se ramène
au cas où $M^\bullet$ est borné. Enfin, pour un complexe borné
$M^a \to \ldots \to M^b$, on raisonne par récurrence sur sa longueur
$b - a$. Le cas $b - a = 1$ a été traité ci-dessus. Si $b - a > 1$, on
considère la suite exacte courte scindée de complexes
$$
0 \to \sigma_{\geq a + 1}M^\bullet \to M^\bullet \to M^a[-a] \to 0
$$
et l'on applique le Lemme \ref{lemma-derived-tor-exact} pour effectuer
l'étape de récurrence. Certains détails sont omis.
\end{proof}

\begin{lemma}
\label{lemma-K-flat-resolution}
Soit $R$ un anneau. Pour tout complexe $M^\bullet$, il existe un complexe
K-plat $K^\bullet$ dont les termes sont des $R$-modules plats et un
quasi-isomorphisme $K^\bullet \to M^\bullet$ surjectif terme à terme.
\end{lemma}

\begin{proof}
Soit $\mathcal{P} \subset \Ob(\text{Mod}_R)$ la classe des $R$-modules
plats. D'après Catégories dérivées, Lemme
\ref{derived-lemma-special-direct-system}, il existe un système
$K_1^\bullet \to K_2^\bullet \to \ldots$ et un diagramme
$$
\xymatrix{
K_1^\bullet \ar[d] \ar[r] &
K_2^\bullet \ar[d] \ar[r] & \ldots \\
\tau_{\leq 1}M^\bullet \ar[r] &
\tau_{\leq 2}M^\bullet \ar[r] & \ldots
}
$$
possédant les propriétés (1), (2), (3) énumérées dans ce lemme. Ces
propriétés impliquent que chaque complexe $K_i^\bullet$ est un complexe
borné supérieurement de modules plats. Ainsi, $K_i^\bullet$ est K-plat
d'après le Lemme \ref{lemma-derived-tor-quasi-isomorphism}. Par construction,
le morphisme induit $\colim_i K_i^\bullet \to M^\bullet$ est un
quasi-isomorphisme surjectif terme à terme. Le complexe
$\colim_i K_i^\bullet$ est K-plat d'après le Lemme
\ref{lemma-colimit-K-flat}. Les termes $\colim K_i^n$ sont plats, car les
colimites filtrantes de modules plats sont plates ; voir Algèbre, Lemme
\ref{algebra-lemma-colimit-flat}.
\end{proof}

\begin{remark}
\label{remark-P-resolution}
En fait, on peut obtenir mieux que dans le Lemme
\ref{lemma-K-flat-resolution}. On peut trouver un quasi-isomorphisme
$P^\bullet \to M^\bullet$, où $P^\bullet$ est un complexe de $R$-modules
muni d'une filtration
$$
0 = F_{-1}P^\bullet \subset F_0P^\bullet \subset
F_1P^\bullet \subset \ldots \subset P^\bullet
$$
par des sous-complexes tels que
\begin{enumerate}
\item $P^\bullet = \bigcup F_pP^\bullet$ ;
\item les inclusions $F_iP^\bullet \to F_{i + 1}P^\bullet$ sont des
injections scindées terme à terme ;
\item les quotients $F_{i + 1}P^\bullet/F_iP^\bullet$ sont isomorphes à des
sommes directes de décalés $R[k]$ (comme complexes, leurs différentielles
sont donc nulles).
\end{enumerate}
Cela sera démontré dans Algèbre différentielle graduée, Lemme
\ref{dga-lemma-resolve} (on peut aussi raisonner comme dans la démonstration
du Lemme \ref{lemma-resolution-filtered-complex}). De plus, étant donné un
tel complexe, on obtient un triangle distingué
$$
\bigoplus F_iP^\bullet \to \bigoplus F_iP^\bullet \to M^\bullet
\to \bigoplus F_iP^\bullet[1]
$$
dans $D(R)$. Cela permet parfois de ramener des énoncés sur des complexes
quelconques à des énoncés sur $R[k]$ (cela ne fonctionne, bien entendu, que
si l'énoncé est préservé par sommes directes). Plus précisément, soit $T$
une propriété des objets de $D(R)$. Supposons que
\begin{enumerate}
\item si $K_i \in D(R)$, $i \in I$, est une famille d'objets telle que
$T(K_i)$ pour tout $i \in I$, alors $T(\bigoplus K_i)$ ;
\item si $K \to L \to M \to K[1]$ est un triangle distingué et si $T$
vaut pour deux de ses objets, alors $T$ vaut pour le troisième ;
\item $T(R[k])$ vaut pour tout $k$.
\end{enumerate}
Alors $T$ vaut pour tout objet de $D(R)$.
\end{remark}

\begin{lemma}
\label{lemma-derived-tor-quasi-isomorphism-other-side}
Soit $R$ un anneau. Soit
$\alpha : P^\bullet \to Q^\bullet$ un quasi-isomorphisme de complexes
K-plats de $R$-modules. Pour tout complexe $L^\bullet$ de $R$-modules, le
morphisme induit
$$
\text{Tot}(\text{id}_L \otimes \alpha) :
\text{Tot}(L^\bullet \otimes_R P^\bullet)
\longrightarrow
\text{Tot}(L^\bullet \otimes_R Q^\bullet)
$$
est un quasi-isomorphisme.
\end{lemma}

\begin{proof}
Choisissons un quasi-isomorphisme $K^\bullet \to L^\bullet$ où
$K^\bullet$ est K-plat ; voir le Lemme \ref{lemma-K-flat-resolution}.
Considérons le diagramme commutatif
$$
\xymatrix{
\text{Tot}(K^\bullet \otimes_R P^\bullet) \ar[r] \ar[d] &
\text{Tot}(K^\bullet \otimes_R Q^\bullet) \ar[d] \\
\text{Tot}(L^\bullet \otimes_R P^\bullet) \ar[r] &
\text{Tot}(L^\bullet \otimes_R Q^\bullet)
}
$$
Le résultat en découle car, d'après le Lemme
\ref{lemma-K-flat-quasi-isomorphism}, les flèches verticales et la flèche
horizontale supérieure sont des quasi-isomorphismes.
\end{proof}

\noindent
Soit $R$ un anneau et soit $M^\bullet$ un objet de $D(R)$. Choisissons une
résolution K-plate $K^\bullet \to M^\bullet$ ; voir le Lemme
\ref{lemma-K-flat-resolution}. D'après les Lemmes
\ref{lemma-derived-tor-homotopy} et \ref{lemma-derived-tor-exact}, on obtient
un foncteur exact de catégories triangulées
$$
K(R) \longrightarrow K(R), \quad
L^\bullet \longmapsto \text{Tot}(L^\bullet \otimes_R K^\bullet)
$$
D'après le Lemme \ref{lemma-K-flat-quasi-isomorphism}, ce foncteur induit
un foncteur $D(R) \to D(R)$, puisque $D(R)$ est la localisation de $K(R)$
par les quasi-isomorphismes. Comme la catégorie des résolutions $K$-plates de
$M^\bullet$ est cofiltrante et en vertu du Lemme
\ref{lemma-derived-tor-quasi-isomorphism-other-side}, le foncteur obtenu, à
isomorphisme près, ne dépend pas du choix de $K^\bullet$.

\begin{definition}
\label{definition-derived-tor}
Soit $R$ un anneau et soit $M^\bullet$ un objet de $D(R)$. Le
{\it produit tensoriel dérivé}
$$
- \otimes_R^{\mathbf{L}} M^\bullet : D(R) \longrightarrow D(R)
$$
est le foncteur exact de catégories triangulées décrit ci-dessus.
\end{definition}

\noindent
Ce foncteur prolonge le foncteur (\ref{equation-derived-tensor-module}).
Nos constructions explicites montrent qu'il existe un isomorphisme (qui
fait intervenir un choix de signes ; voir ci-dessous)
$$
M^\bullet \otimes_R^{\mathbf{L}} L^\bullet
\cong
L^\bullet \otimes_R^{\mathbf{L}} M^\bullet
$$
dès que $L^\bullet$ et $M^\bullet$ appartiennent à $D(R)$. Ainsi, lorsque
nous écrivons $M^\bullet \otimes_R^{\mathbf{L}} L^\bullet$, nous ne
préciserons généralement pas la variable par rapport à laquelle le produit
tensoriel dérivé est défini.

\begin{lemma}
\label{lemma-flip-douoble-tensor-product}
Soit $R$ un anneau et soient $K^\bullet, L^\bullet$ des complexes de
$R$-modules. Il existe un isomorphisme canonique
$$
K^\bullet \otimes_R^\mathbf{L} L^\bullet \longrightarrow
L^\bullet \otimes_R^\mathbf{L} K^\bullet
$$
fonctoriel en les deux complexes, qui utilise le signe $(-1)^{pq}$ pour le
morphisme $K^p \otimes_R L^q \to L^q \otimes_R K^p$ (voir
l'explication dans la démonstration).
\end{lemma}

\begin{proof}
On peut, et l'on va, remplacer les complexes par des complexes K-plats
$K^\bullet$ et $L^\bullet$, puis utiliser la contrainte de commutativité
étudiée dans la Section \ref{section-symmetric-monoidal}.
\end{proof}

\begin{lemma}
\label{lemma-triple-tensor-product}
Soit $R$ un anneau et soient $K^\bullet, L^\bullet, M^\bullet$ des complexes
de $R$-modules. Il existe un isomorphisme canonique
$$
(K^\bullet \otimes_R^\mathbf{L} L^\bullet) \otimes_R^\mathbf{L} M^\bullet
=
K^\bullet \otimes_R^\mathbf{L} (L^\bullet \otimes_R^\mathbf{L} M^\bullet)
$$
fonctoriel en les trois complexes.
\end{lemma}

\begin{proof}
On remplace les complexes par des complexes K-plats et l'on utilise la
contrainte d'associativité de la Section \ref{section-symmetric-monoidal}.
\end{proof}

\begin{lemma}
\label{lemma-factor-through-K-flat}
Soit $R$ un anneau et soit $a : K^\bullet \to L^\bullet$ un morphisme de
complexes de $R$-modules. Si $K^\bullet$ est K-plat, il existe un complexe
$N^\bullet$ et des morphismes de complexes $b : K^\bullet \to N^\bullet$
et $c : N^\bullet \to L^\bullet$ tels que
\begin{enumerate}
\item $N^\bullet$ est K-plat ;
\item $c$ est un quasi-isomorphisme ;
\item $a$ est homotope à $c \circ b$.
\end{enumerate}
Si les termes de $K^\bullet$ sont plats, on peut choisir $N^\bullet$, $b$
et $c$ de sorte que les termes de $N^\bullet$ le soient également.
\end{lemma}

\begin{proof}
Nous utiliserons le fait que la catégorie d'homotopie $K(R)$ est une
catégorie triangulée ; voir Catégories dérivées, Proposition
\ref{derived-proposition-homotopy-category-triangulated}.
Choisissons un triangle distingué
$K^\bullet \to L^\bullet \to C^\bullet \to K^\bullet[1]$.
Choisissons un quasi-isomorphisme $M^\bullet \to C^\bullet$, où
$M^\bullet$ est K-plat à termes plats ; voir le Lemme
\ref{lemma-K-flat-resolution}. D'après les axiomes des catégories
triangulées, on peut insérer la composée
$M^\bullet \to C^\bullet \to K^\bullet[1]$ dans un triangle distingué
$K^\bullet \to N^\bullet \to M^\bullet \to K^\bullet[1]$.
D'après le Lemme \ref{lemma-K-flat-two-out-of-three}, $N^\bullet$ est
K-plat. En utilisant de nouveau les axiomes des catégories triangulées, on
peut choisir un morphisme $N^\bullet \to L^\bullet$ qui s'insère dans le
morphisme suivant de triangles distingués
$$
\xymatrix{
K^\bullet \ar[r] \ar[d] &
N^\bullet \ar[r] \ar[d] &
M^\bullet \ar[r] \ar[d] &
K^\bullet[1] \ar[d] \\
K^\bullet \ar[r] &
L^\bullet \ar[r] &
C^\bullet \ar[r] &
K^\bullet[1]
}
$$
Comme deux des trois flèches sont des quasi-isomorphismes, la troisième
$N^\bullet \to L^\bullet$ l'est aussi, d'après les suites exactes longues de
cohomologie associées à ces triangles distingués (on peut aussi considérer
l'image de ce diagramme dans $D(R)$ et utiliser Catégories dérivées, Lemme
\ref{derived-lemma-third-isomorphism-triangle}). Cela achève la démonstration
de (1), (2) et (3). Pour démontrer la dernière assertion, on peut choisir
$N^\bullet$ de sorte que $N^n \cong M^n \oplus K^n$ ; voir Catégories
dérivées, Lemme
\ref{derived-lemma-improve-distinguished-triangle-homotopy}. On obtient donc
la platitude voulue si les termes de $K^\bullet$ sont plats.
\end{proof}








\section{Changement d'anneaux dérivé}
\label{section-derived-base-change}

\noindent
Soit $R \to A$ un homomorphisme d'anneaux et soit $N^\bullet$ un complexe
de $A$-modules. On peut aussi employer des résolutions K-plates pour définir
un foncteur
$$
- \otimes_R^{\mathbf{L}} N^\bullet : D(R) \to D(A)
$$
comme foncteur dérivé à gauche du foncteur $K(R) \to K(A)$,
$M^\bullet \mapsto \text{Tot}(M^\bullet \otimes_R N^\bullet)$.
En particulier, en prenant $N^\bullet = A[0]$, on obtient un foncteur de
changement de base dérivé
$$
- \otimes_R^{\mathbf{L}} A : D(R) \to D(A)
$$
qui prolonge le foncteur (\ref{equation-derived-tensor-algebra}). En effet,
pour tout complexe de $R$-modules $M^\bullet$, on peut choisir une
résolution K-plate $K^\bullet \to M^\bullet$ et poser
$$
M^\bullet \otimes_R^{\mathbf{L}} N^\bullet =
\text{Tot}(K^\bullet \otimes_R N^\bullet).
$$
Les Lemmes \ref{lemma-K-flat-resolution} et
\ref{lemma-derived-tor-quasi-isomorphism-other-side} montrent que cette
définition est bien posée. Pour régler tous les détails, il est toutefois
peut-être plus commode de faire appel à une théorie générale.

\begin{lemma}
\label{lemma-derived-base-change}
La construction ci-dessus est indépendante des choix et définit un foncteur
exact de catégories triangulées
$- \otimes_R^\mathbf{L} N^\bullet : D(R) \to D(A)$.
Il existe un isomorphisme fonctoriel
$$
E^\bullet \otimes_R^\mathbf{L} N^\bullet =
(E^\bullet \otimes_R^\mathbf{L} A) \otimes_A^\mathbf{L} N^\bullet
$$
pour $E^\bullet$ dans $D(R)$.
\end{lemma}

\begin{proof}
Pour démontrer l'existence du foncteur dérivé
$- \otimes_R^\mathbf{L} N^\bullet$, on utilise la théorie générale
développée dans Catégories dérivées, Section
\ref{derived-section-derived-functors}. Posons $\mathcal{D} = K(R)$ et
$\mathcal{D}' = D(A)$. Notons $F : \mathcal{D} \to \mathcal{D}'$ le
foncteur exact de catégories triangulées défini par
$F(M^\bullet) = \text{Tot}(M^\bullet \otimes_R N^\bullet)$. Les propriétés
annoncées de $F$ résultent des Lemmes \ref{lemma-derived-tor-homotopy} et
\ref{lemma-derived-tor-exact}. Soit $S$ l'ensemble des quasi-isomorphismes
de $\mathcal{D} = K(R)$. On est alors dans la situation de Catégories
dérivées, Situation \ref{derived-situation-derived-functor}, de sorte que
Catégories dérivées, Définition
\ref{derived-definition-right-derived-functor-defined} s'applique. Nous
affirmons que $LF$ est défini partout. Cela résulte de Catégories dérivées,
Lemme \ref{derived-lemma-find-existence-computes}, en prenant pour
$\mathcal{P} \subset \Ob(\mathcal{D})$ la classe des complexes K-plats :
(1) découle du Lemme \ref{lemma-K-flat-resolution} et (2) du Lemme
\ref{lemma-derived-tor-quasi-isomorphism-other-side}. On obtient ainsi un
foncteur dérivé
$$
LF : D(R) = S^{-1}\mathcal{D} \longrightarrow \mathcal{D}' = D(A)
$$
voir Catégories dérivées, Équation (\ref{derived-equation-everywhere}).
Enfin, Catégories dérivées, Lemme
\ref{derived-lemma-find-existence-computes}, garantit que
$LF(K^\bullet) = F(K^\bullet) =
\text{Tot}(K^\bullet \otimes_R N^\bullet)$
lorsque $K^\bullet$ est K-plat ; autrement dit, $LF$ se calcule bien de la
manière décrite ci-dessus. De plus, d'après le Lemme
\ref{lemma-base-change-K-flat}, le complexe $K^\bullet \otimes_R A$ est un
complexe K-plat de $A$-modules. Ainsi
$$
(K^\bullet \otimes_R^\mathbf{L} A) \otimes_A^\mathbf{L} N^\bullet =
\text{Tot}((K^\bullet \otimes_R A) \otimes_A N^\bullet) =
\text{Tot}(K^\bullet \otimes_A N^\bullet) =
K^\bullet \otimes_A^\mathbf{L} N^\bullet
$$
ce qui démontre la dernière assertion du lemme.
\end{proof}

\begin{lemma}
\label{lemma-functoriality-derived-base-change}
Soit $R \to A$ un homomorphisme d'anneaux et soit
$f : L^\bullet \to N^\bullet$ un morphisme de complexes de $A$-modules.
Alors $f$ induit une transformation de foncteurs
$$
1 \otimes f :
- \otimes_A^\mathbf{L} L^\bullet
\longrightarrow
- \otimes_A^\mathbf{L} N^\bullet
$$
Si $f$ est un quasi-isomorphisme, alors $1 \otimes f$ est un isomorphisme
de foncteurs.
\end{lemma}

\begin{proof}
Comme les foncteurs se calculent en les évaluant sur les complexes K-plats
$K^\bullet$, il suffit d'utiliser la fonctorialité de
$$
\text{Tot}(K^\bullet \otimes_R L^\bullet) \to
\text{Tot}(K^\bullet \otimes_R N^\bullet)
$$
pour définir la transformation. La dernière assertion résulte du Lemme
\ref{lemma-K-flat-quasi-isomorphism}.
\end{proof}

\begin{lemma}
\label{lemma-tensor-hom-adjoint}
Soit $R \to A$ un homomorphisme d'anneaux. Le foncteur
$D(R) \to D(A)$, $E \mapsto E \otimes_R^\mathbf{L} A$
du Lemme \ref{lemma-derived-base-change} est adjoint à gauche du foncteur de
restriction $D(A) \to D(R)$.
\end{lemma}

\begin{proof}
Cela résulte de Catégories dérivées, Lemme
\ref{derived-lemma-pre-derived-adjoint-functors-general}, et du fait que
$- \otimes_R A$ et la restriction sont adjoints d'après Algèbre, Lemme
\ref{algebra-lemma-adjoint-tensor-restrict}.
\end{proof}

\begin{remark}[Avertissement]
\label{remark-warning-compute-base-change}
Soit $R \to A$ un homomorphisme d'anneaux, et soient $N$ et $N'$ des
$A$-modules. Notons $N_R$ et $N'_R$ les restrictions de $N$ et $N'$ aux
$R$-modules ; voir Algèbre, Section \ref{algebra-section-base-change}. Dans
cette situation, les objets $N_R \otimes_R^\mathbf{L} N'$ et
$N \otimes_R^\mathbf{L} N'_R$ de $D(A)$ ne sont en général pas
isomorphes ! Il faut donc faire très attention au membre dont provient la
structure de $A$-module.

\medskip\noindent
Pour un exemple explicite, posons $R = k[x, y]$, $A = R/(xy)$,
$N = R/(x)$ et $N' = A = R/(xy)$. La résolution
$0 \to R \xrightarrow{xy} R \to N'_R \to 0$
montre que $N \otimes_R^\mathbf{L} N'_R = N[1] \oplus N$ dans $D(A)$. La
résolution
$0 \to R \xrightarrow{x} R \to N_R \to 0$
montre que $N_R \otimes_R^\mathbf{L} N'$ est représenté par le complexe
$A \xrightarrow{x} A$. Pour voir que ces deux complexes ne sont pas
isomorphes, on peut montrer que le second n'est pas isomorphe dans $D(A)$ à
la somme directe de ses groupes de cohomologie ; on peut aussi montrer que
le premier n'est pas un objet parfait de $D(A)$, tandis que le second l'est.
Certains détails sont omis.
\end{remark}

\begin{lemma}
\label{lemma-double-base-change}
Soient $A \to B \to C$ des homomorphismes d'anneaux, $N^\bullet$ un
complexe de $B$-modules et $K^\bullet$ un complexe de $C$-modules. La
composée des foncteurs
$$
D(A) \xrightarrow{- \otimes_A^\mathbf{L} N^\bullet}
D(B) \xrightarrow{- \otimes_B^\mathbf{L} K^\bullet} D(C)
$$
est le foncteur
$- \otimes_A^\mathbf{L} (N^\bullet \otimes_B^\mathbf{L} K^\bullet) :
D(A) \to D(C)$. Si $M$, $N$, $K$ sont respectivement des modules sur $A$,
$B$, $C$, alors on a
$$
(M \otimes_A^\mathbf{L} N) \otimes_B^\mathbf{L} K =
M \otimes_A^\mathbf{L} (N \otimes_B^\mathbf{L} K) =
(M \otimes_A^\mathbf{L} C) \otimes_C^\mathbf{L} (N \otimes_B^\mathbf{L} K)
$$
dans $D(C)$. On a aussi un isomorphisme canonique
$$
(M \otimes_A^\mathbf{L} N) \otimes_B^\mathbf{L} K \longrightarrow
(M \otimes_A^\mathbf{L} K) \otimes_C^\mathbf{L} (N \otimes_B^\mathbf{L} C)
$$
avec des signes. Des résultats analogues valent pour les complexes.
\end{lemma}

\begin{proof}
Choisissons un complexe K-plat $P^\bullet$ de $B$-modules et un
quasi-isomorphisme $P^\bullet \to N^\bullet$ (Lemme
\ref{lemma-K-flat-resolution}). Soit $M^\bullet$ un complexe K-plat de
$A$-modules représentant un objet quelconque de $D(A)$. Alors
$$
(M^\bullet \otimes_A^\mathbf{L} P^\bullet) \otimes_B^\mathbf{L} K^\bullet
\longrightarrow
(M^\bullet \otimes_A^\mathbf{L} N^\bullet) \otimes_B^\mathbf{L} K^\bullet
$$
est un isomorphisme d'après le Lemme
\ref{lemma-functoriality-derived-base-change}, appliqué à ce qui se trouve
entre parenthèses. D'après les Lemmes \ref{lemma-base-change-K-flat} et
\ref{lemma-tensor-product-K-flat}, le complexe
$$
\text{Tot}(M^\bullet \otimes_A P^\bullet) =
\text{Tot}((M^\bullet \otimes_R A) \otimes_A P^\bullet
$$
est K-plat comme complexe de $B$-modules et représente, par construction,
le produit tensoriel dérivé dans $D(B)$. Par conséquent,
$(M^\bullet \otimes_A^\mathbf{L} P^\bullet) \otimes_B^\mathbf{L} K^\bullet$
est représenté par le complexe
$$
\text{Tot}(\text{Tot}(M^\bullet \otimes_A P^\bullet)\otimes_B K^\bullet) =
\text{Tot}(M^\bullet \otimes_A \text{Tot}(P^\bullet \otimes_B K^\bullet))
$$
de $C$-modules. L'égalité résulte de Homologie, Remarque
\ref{homology-remark-triple-complex}. En remontant les étapes précédentes,
on voit que ceci est égal à
$$
M^\bullet \otimes_A^\mathbf{L} (P^\bullet \otimes_B^\mathbf{L} K^\bullet)
\longleftarrow
M^\bullet \otimes_A^\mathbf{L} (N^\bullet \otimes_B^\mathbf{L} K^\bullet)
$$
La flèche est un isomorphisme par définition du foncteur
$-\otimes_B^\mathbf{L} K^\bullet$. Toutes ces constructions sont
fonctorielles en $M^\bullet$ ; on obtient donc l'isomorphisme de foncteurs
annoncé.

\medskip\noindent
D'après ce qui précède, on a la première égalité dans
$$
(M \otimes_A^\mathbf{L} N) \otimes_B^\mathbf{L} K =
M \otimes_A^\mathbf{L} (N \otimes_B^\mathbf{L} K) =
(M \otimes_A^\mathbf{L} C) \otimes_C^\mathbf{L} (N \otimes_B^\mathbf{L} K)
$$
La seconde égalité résulte de la dernière assertion du Lemme
\ref{lemma-derived-base-change}. De même, on peut écrire
$N \otimes_B^\mathbf{L} K = (N \otimes_B^\mathbf{L} C) \otimes_C^\mathbf{L} K$
et, en substituant, on obtient
\begin{align*}
(M \otimes_A^\mathbf{L} N) \otimes_B^\mathbf{L} K
& =
(M \otimes_A^\mathbf{L} C) \otimes_C^\mathbf{L}
((N \otimes_B^\mathbf{L} C) \otimes_C^\mathbf{L} K) \\
& =
(M \otimes_A^\mathbf{L} C) \otimes_C^\mathbf{L}
(K \otimes_C^\mathbf{L} (N \otimes_B^\mathbf{L} C)) \\
& =
((M \otimes_A^\mathbf{L} C) \otimes_C^\mathbf{L} K)
\otimes_C^\mathbf{L} (N \otimes_B^\mathbf{L} C)) \\
& =
(M \otimes_C^\mathbf{L} K)
\otimes_C^\mathbf{L} (N \otimes_B^\mathbf{L} C)
\end{align*}
d'après les Lemmes \ref{lemma-flip-douoble-tensor-product} et
\ref{lemma-triple-tensor-product}, ainsi que le lemme déjà cité.
\end{proof}






\section{Indépendance pour Tor}
\label{section-tor-independence}

\noindent
Considérons un diagramme commutatif
$$
\xymatrix{
A \ar[r] & A' \\
R \ar[r] \ar[u] & R' \ar[u]
}
$$
d'anneaux. Étant donné un objet $K$ de $D(A)$, on peut considérer son
changement de base dérivé $K \otimes_A^\mathbf{L} A'$, qui est un objet de
$D(A')$. On peut aussi restreindre $K$ en un objet de $D(R)$ et considérer
son changement de base dérivé en un objet de $D(R')$, noté
$K \otimes_R^\mathbf{L} R'$. Nous affirmons qu'il existe un morphisme de
comparaison fonctoriel
\begin{equation}
\label{equation-comparison-map}
K \otimes_R^{\mathbf{L}} R'
\longrightarrow
K \otimes_A^{\mathbf{L}} A'
\end{equation}
dans $D(R')$. Pour construire ce morphisme, choisissons un complexe K-plat
$K^\bullet$ de $A$-modules représentant $K$, puis un quasi-isomorphisme
$E^\bullet \to K^\bullet$ où $E^\bullet$ est un complexe K-plat de
$R$-modules. Le morphisme ci-dessus est
$$
K \otimes_R^{\mathbf{L}} R' =
E^\bullet \otimes_R R'
\longrightarrow
K^\bullet \otimes_A A' =
K \otimes_A^{\mathbf{L}} A'
$$
En général, ce morphisme n'a aucune raison d'être un isomorphisme.

\medskip\noindent
On rencontre toutefois souvent le cas où le diagramme ci-dessus est un
diagramme de « changement de base » d'anneaux, c'est-à-dire
$A' = A \otimes_R R'$. Alors, pour tout $A$-module $M$, on a
$M \otimes_A A' = M \otimes_R R'$. Ainsi, $- \otimes_R R'$ est égal à
$- \otimes_A A'$ comme foncteur
$\text{Mod}_A \to \text{Mod}_{A'}$. En général, cette égalité {\bf ne se
prolonge pas} aux produits tensoriels dérivés. Autrement dit, le morphisme
de comparaison n'est pas un isomorphisme. Un exemple simple consiste à
prendre $R = k[x]$, $A = R' = A' = k[x]/(x) = k$ et $K^\bullet = A[0]$.
Une condition nécessaire évidente est que $\text{Tor}_p^R(A, R') = 0$ pour
tout $p > 0$.

\begin{definition}
\label{definition-tor-independent}
Soit $R$ un anneau et soient $A$, $B$ des $R$-algèbres. On dit que $A$ et
$B$ sont {\it Tor-indépendantes sur $R$} si
$\text{Tor}_p^R(A, B) = 0$ pour tout $p > 0$.
\end{definition}

\begin{lemma}
\label{lemma-base-change-comparison}
Le morphisme de comparaison (\ref{equation-comparison-map}) est un
isomorphisme si $A' = A \otimes_R R'$ et si $A$ et $R'$ sont
Tor-indépendantes sur $R$.
\end{lemma}

\begin{proof}
Choisissons une résolution libre $F^\bullet \to R'$ de $R'$ comme
$R$-module. Puisque $A$ et $R'$ sont Tor-indépendantes sur $R$, on voit que
$F^\bullet \otimes_R A$ est, comme complexe de $A$-modules, une résolution
libre de $A'$ sur $A$.
D'après la construction générale du produit tensoriel dérivé donnée
ci-dessus, on a
$$
K^\bullet \otimes_A A' \cong
\text{Tot}(K^\bullet \otimes_A (F^\bullet \otimes_R A)) =
\text{Tot}(K^\bullet \otimes_R F^\bullet) \cong
\text{Tot}(E^\bullet \otimes_R F^\bullet) \cong
E^\bullet \otimes_R R'
$$
comme voulu.
\end{proof}

\begin{lemma}
\label{lemma-tor-independent-flat}
Considérons un diagramme commutatif d'anneaux
$$
\xymatrix{
A' & R' \ar[r] \ar[l] & B' \\
A \ar[u] & R \ar[l] \ar[u] \ar[r] & B \ar[u]
}
$$
Supposons que $R'$ soit plat sur $R$, que $A'$ soit plat sur
$A \otimes_R R'$ et que $B'$ soit plat sur $R' \otimes_R B$. Alors
$$
\text{Tor}_i^R(A, B) \otimes_{(A \otimes_R B)} (A' \otimes_{R'} B') =
\text{Tor}_i^{R'}(A', B')
$$
\end{lemma}

\begin{proof}
D'après Algèbre, Section \ref{algebra-section-functoriality-tor}, il existe
des morphismes canoniques
$$
\text{Tor}_i^R(A, B) \longrightarrow
\text{Tor}_i^{R'}(A \otimes_R R', B \otimes_R R') \longrightarrow
\text{Tor}_i^{R'}(A', B')
$$
Ils induisent un morphisme du membre de gauche vers le membre de droite de
la formule du lemme.

\medskip\noindent
Prenons une résolution libre $F_\bullet \to A$ de $A$ comme $R$-module.
Alors $F_\bullet \otimes_R R'$ est une résolution de
$A \otimes_R R'$. Par conséquent,
$\text{Tor}_i^{R'}(A \otimes_R R', B \otimes_R R')$ se calcule par
$F_\bullet \otimes_R B \otimes_R R'$. Puisque $R'$ est plat sur $R$, ceci
calcule $\text{Tor}_i^R(A, B) \otimes_R R'$. Ainsi,
$\text{Tor}_i^{R'}(A \otimes_R R', B \otimes_R R') =
\text{Tor}_i^R(A, B) \otimes_R R'$ (on n'utilise ici que la platitude de
$R'$ sur $R$).

\medskip\noindent
D'après le théorème de Lazard (Algèbre, Théorème
\ref{algebra-theorem-lazard}), on peut écrire $A'$, resp.\ $B'$, comme une
colimite filtrante de modules libres de rang fini sur $A \otimes_R R'$,
resp.\ $B \otimes_R R'$. Écrivons $A' = \colim M_i$ et
$B' = \colim N_j$. Le résultat précédent donne
$$
\text{Tor}_i^{R'}(M_i, N_j) =
\text{Tor}_i^R(A, B) \otimes_{A \otimes_R B} (M_i \otimes_{R'} N_j)
$$
comme on le vérifie en écrivant tout dans des bases. En passant à la
colimite, on obtient le résultat du lemme.
\end{proof}

\begin{lemma}
\label{lemma-flat-base-change-tor-independent}
Soient $R \to A$ et $R \to B$ des homomorphismes d'anneaux. Soit
$R \to R'$ un homomorphisme d'anneaux, et posons
$A' = A \otimes_R R'$ et $B' = B \otimes_R R'$. Si $A$ et $B$ sont
Tor-indépendantes sur $R$ et si $R \to R'$ est plat, alors $A'$ et $B'$
sont Tor-indépendantes sur $R'$.
\end{lemma}

\begin{proof}
Cela résulte immédiatement du Lemme \ref{lemma-tor-independent-flat} et de
la Définition \ref{definition-tor-independent}.
\end{proof}

\begin{lemma}
\label{lemma-lemma-tor-independent-flat-compare}
Sous les hypothèses du Lemme \ref{lemma-tor-independent-flat}, pour
$M \in D(A)$ il existe des isomorphismes canoniques
$$
H^i((M \otimes_A^\mathbf{L} A') \otimes_{R'}^\mathbf{L} B') =
H^i(M \otimes_R^\mathbf{L} B) \otimes_{(A \otimes_R B)} (A' \otimes_{R'} B')
$$
de $A' \otimes_{R'} B'$-modules.
\end{lemma}

\begin{proof}
Précisons les deux membres de l'égalité. À gauche, on compose les foncteurs
$D(A) \to D(A') \to D(R') \to D(B')$ avec le foncteur
$H^i : D(B') \to \text{Mod}_{B'}$. Comme il existe un morphisme de $A'$ vers
les endomorphismes de l'objet
$(M \otimes_A^\mathbf{L} A') \otimes_{R'}^\mathbf{L} B'$ dans $D(B')$, le
membre de gauche est bien un $A' \otimes_{R'} B'$-module. Le même
raisonnement montre que $H^i(M \otimes_R^\mathbf{L} B)$ possède une
structure de $A \otimes_R B$-module.

\medskip\noindent
Démontrons d'abord le résultat lorsque $B' = R' \otimes_R B$. Choisissons
alors une résolution $F^\bullet \to B$ par des $R$-modules libres, ainsi
qu'un complexe K-plat $M^\bullet$ de $A$-modules représentant $M$. Le
membre de gauche est représenté par
\begin{align*}
H^i(\text{Tot}((M^\bullet \otimes_A A') \otimes_{R'} (R' \otimes_R F^\bullet)))
& =
H^i(\text{Tot}(M^\bullet \otimes_A A' \otimes_R F^\bullet)) \\
& =
H^i(\text{Tot}(M^\bullet \otimes_R F^\bullet) \otimes_A A') \\
& =
H^i(M \otimes_R^\mathbf{L} B) \otimes_A A'
\end{align*}
La dernière égalité vient de ce que $A \to A'$ est plat. Le dernier module
est celui voulu, car $A' \otimes_{R'} B' = A' \otimes_R B$, puisque nous
avons supposé $B' = R' \otimes_R B$ dans ce paragraphe.

\medskip\noindent
Passons au cas général. Supposons que $B' \to B''$ soit un homomorphisme
plat d'anneaux. On voit immédiatement que
$$
H^i((M \otimes_A^\mathbf{L} A') \otimes_{R'}^\mathbf{L} B'') =
H^i((M \otimes_A^\mathbf{L} A') \otimes_{R'}^\mathbf{L} B')
\otimes_{B'} B''
$$
and
$$
H^i(M \otimes_R^\mathbf{L} B) \otimes_{(A \otimes_R B)} (A' \otimes_{R'} B'')
=
\left(
H^i(M \otimes_R^\mathbf{L} B) \otimes_{(A \otimes_R B)} (A' \otimes_{R'} B')
\right) \otimes_{B'} B''
$$
Ainsi, le résultat pour $B'$ implique celui pour $B''$. Comme le paragraphe
précédent a établi le résultat pour $R' \otimes_R B$, le cas général en
découle.
\end{proof}

\begin{lemma}
\label{lemma-tor-independent}
Soit $R$ un anneau et soient $A$, $B$ des $R$-algèbres. Les assertions
suivantes sont équivalentes :
\begin{enumerate}
\item $A$ et $B$ sont Tor-indépendantes sur $R$ ;
\item pour toute paire d'idéaux premiers $\mathfrak p \subset A$ et
$\mathfrak q \subset B$ situés au-dessus du même idéal premier
$\mathfrak r \subset R$, les anneaux $A_\mathfrak p$ et $B_\mathfrak q$
sont Tor-indépendants sur $R_\mathfrak r$ ;
\item pour tout idéal premier $\mathfrak s$ de $A \otimes_R B$, le module
$$
\text{Tor}_i^R(A, B)_\mathfrak s =
\text{Tor}_i^{R_\mathfrak r}(A_\mathfrak p, B_\mathfrak q)_\mathfrak s
$$
(où $\mathfrak p = A \cap \mathfrak s$, $\mathfrak q = B \cap \mathfrak s$
et $\mathfrak r = R \cap \mathfrak s$) est nul.
\end{enumerate}
\end{lemma}

\begin{proof}
Soit $\mathfrak s$ un idéal premier de $A \otimes_R B$ comme en (3).
L'égalité
$$
\text{Tor}_i^R(A, B)_\mathfrak s =
\text{Tor}_i^{R_\mathfrak r}(A_\mathfrak p, B_\mathfrak q)_\mathfrak s
$$
où $\mathfrak p = A \cap \mathfrak s$, $\mathfrak q = B \cap \mathfrak s$
et $\mathfrak r = R \cap \mathfrak s$, résulte du Lemme
\ref{lemma-tor-independent-flat}. Ainsi, (2) implique (3). Comme la nullité
d'un module se vérifie après localisation aux idéaux premiers (Algèbre,
Lemme \ref{algebra-lemma-characterize-zero-local}), on en déduit que (3)
implique (1). Pour (1) $\Rightarrow$ (2), on utilise
$$
\text{Tor}_i^{R_\mathfrak r}(A_\mathfrak p, B_\mathfrak q) =
\text{Tor}_i^R(A, B) \otimes_{(A \otimes_R B)}
(A_\mathfrak p \otimes_{R_{\mathfrak r}} B_\mathfrak q)
$$
de nouveau d'après le Lemme \ref{lemma-tor-independent-flat}.
\end{proof}








\section{Suites spectrales pour Tor}
\label{section-spectral-sequence-tor}


\noindent
Dans cette section, nous rassemblons diverses suites spectrales qui
apparaissent lorsqu'on considère les foncteurs Tor.

\begin{example}
\label{example-cohomology-complex-tensored}
Soit $R$ un anneau. Soit $K_\bullet$ un complexe de chaînes de $R$-modules
tel que $K_n = 0$ pour $n \ll 0$. Soit $M$ un $R$-module. Choisissons une
résolution $P_\bullet \to M$ de $M$ par des $R$-modules libres. On obtient
un complexe double de chaînes $K_\bullet \otimes_R P_\bullet$. En appliquant
les résultats de Homologie, Section \ref{homology-section-double-complex}
(en particulier Homologie, Lemme
\ref{homology-lemma-first-quadrant-ss}), traduits dans le langage des
complexes de chaînes, on trouve deux suites spectrales qui convergent vers
$H_*(K_\bullet \otimes_R^\mathbf{L} M)$. D'une part, une suite spectrale de
page $E_2$
$$
(E_2)_{i, j} = \text{Tor}^R_j(H_i(K_\bullet), M)
\Rightarrow
H_{i + j}(K_\bullet \otimes^{\mathbf{L}}_R M)
$$
dont la différentielle $d_2$ est donnée par les morphismes
$\text{Tor}^R_j(H_i(K_\bullet), M) \to
\text{Tor}^R_{j - 2}(H_{i + 1}(K_\bullet), M)$.
Et une autre suite spectrale de page $E_1$
$$
(E_1)_{i, j} = \text{Tor}^R_j(K_i, M)
\Rightarrow
H_{i + j}(K_\bullet \otimes^{\mathbf{L}}_R M)
$$
dont la différentielle $d_1$ est donnée par les morphismes
$\text{Tor}^R_j(K_i, M) \to \text{Tor}^R_j(K_{i - 1}, M)$
induits par $K_i \to K_{i - 1}$.
\end{example}

\begin{example}
\label{example-tor-change-rings}
Soit $R \to S$ un homomorphisme d'anneaux. Soit $M$ un $R$-module et soit
$N$ un $S$-module. Il existe alors une suite spectrale
$$
\text{Tor}^S_n(\text{Tor}^R_m(M, S), N) \Rightarrow
\text{Tor}^R_{n + m}(M, N).
$$
Pour la construire, choisissons une résolution $R$-libre $P^\bullet$ de
$M$. On a alors
$$
M \otimes_R^{\mathbf{L}} N = P^\bullet \otimes_R N =
(P^\bullet \otimes_R S) \otimes_S N
$$
puis on applique la première suite spectrale de l'Exemple
\ref{example-cohomology-complex-tensored}.
\end{example}

\begin{example}
\label{example-tor-base-change}
Considérons un diagramme commutatif
$$
\xymatrix{
B \ar[r] & B' = B \otimes_A A' \\
A \ar[r] \ar[u] & A' \ar[u]
}
$$
et des $B$-modules $M, N$. Posons
$M' = M \otimes_A A' = M \otimes_B B'$ et
$N' = N \otimes_A A' = N \otimes_B B'$.
{\it Supposons que $A \to B$ soit plat et que $M$ et $N$ soient plats sur
$A$.} Il existe alors une suite spectrale
$$
\text{Tor}^A_i(\text{Tor}_j^B(M, N), A')
\Rightarrow
\text{Tor}^{B'}_{i + j}(M', N')
$$
En voici la raison. Choisissons une résolution libre $F_\bullet \to M$ dans
la catégorie des $B$-modules. Comme $B$ et $M$ sont plats sur $A$, on voit que
$F_\bullet \otimes_A A'$ est une résolution $B'$-libre de $M'$. Les groupes
$\text{Tor}^{B'}_n(M', N')$ se calculent donc par le complexe
$$
(F_\bullet \otimes_A A') \otimes_{B'} N' =
(F_\bullet \otimes_B N) \otimes_A A' =
(F_\bullet \otimes_B N) \otimes^{\mathbf{L}}_A A'
$$
la dernière égalité venant de ce que $F_\bullet \otimes_B N$ est un complexe
de $A$-modules plats, puisque $N$ est plat sur $A$. On obtient donc la suite
spectrale en appliquant celle de l'Exemple
\ref{example-cohomology-complex-tensored}.
\end{example}

\begin{example}
\label{example-tor}
Soient $K^\bullet, L^\bullet$ des objets de $D^{-}(R)$. Il existe une suite
spectrale
$$
E_2^{p, q} = H^p(K^\bullet \otimes_R^{\mathbf{L}} H^q(L^\bullet))
\Rightarrow H^{p + q}(K^\bullet \otimes_R^{\mathbf{L}} L^\bullet)
$$
et une autre suite spectrale
$$
E_2^{p, q} =
H^p(H^q(K^\bullet) \otimes_R^{\mathbf{L}} L^\bullet)
\Rightarrow H^{p + q}(K^\bullet \otimes_R^{\mathbf{L}} L^\bullet)
$$
Les deux suites spectrales ont pour différentielle
$d_2^{p, q} : E_2^{p, q} \to E_2^{p + 2, q - 1}$. Après avoir remplacé
$K^\bullet$ et $L^\bullet$ par des complexes bornés supérieurement d'objets
projectifs, ce sont simplement les deux suites spectrales qui calculent la
cohomologie de $\text{Tot}(K^\bullet \otimes L^\bullet)$ étudiées dans
Homologie, Section \ref{homology-section-double-complex}.
\end{example}









\section{Produits et Tor}
\label{section-products-tor}

\noindent
L'exemple le plus simple de morphismes produits est la situation suivante.
Supposons que $K^\bullet, L^\bullet \in D(R)$. Il existe alors des morphismes
\begin{equation}
\label{equation-simple-tor-product}
H^i(K^\bullet) \otimes_R H^j(L^\bullet)
\longrightarrow
H^{i + j}(K^\bullet \otimes_R^{\mathbf{L}} L^\bullet)
\end{equation}
Pour les définir, on peut supposer que l'un des complexes
$K^\bullet, L^\bullet$ est un complexe K-plat de $R$-modules (par exemple un
complexe borné supérieurement de $R$-modules libres ou projectifs). Alors
$K^\bullet \otimes_R^{\mathbf{L}} L^\bullet$ est représenté par le complexe
$\text{Tot}(K^\bullet \otimes_R L^\bullet)$ ; voir la Section
\ref{section-derived-tensor-product} (ou la Section
\ref{section-computing-tor}). Supposons ensuite que
$\xi \in H^i(K^\bullet)$ et $\zeta \in H^j(L^\bullet)$. Choisissons
$k \in \Ker(K^i \to K^{i + 1})$ et
$l \in \Ker(L^j \to L^{j + 1})$ représentant $\xi$ et $\zeta$. Posons
$$
\xi \cup \zeta =
\text{classe de }k \otimes l\text{ dans }
H^{i + j}(\text{Tot}(K^\bullet \otimes_R L^\bullet)).
$$
Cette définition a un sens, car la formule de la différentielle $\text{d}$
sur le complexe total (voir Homologie, Définition
\ref{homology-definition-associated-simple-complex}) montre que
$k \otimes l$ est un cocycle. De plus, si $k' = d_K(k'')$ pour un
$k'' \in K^{i - 1}$, alors $k' \otimes l = \text{d}(k'' \otimes l)$, car
$l$ est un cocycle. De même, modifier le choix de $l$ représentant $\zeta$
ne change pas la classe de $k \otimes l$. Il est tout aussi clair que
$\cup$ est bilinéaire ; à un élément quelconque de
$H^i(K^\bullet) \otimes_R H^j(L^\bullet)$, on associe donc
$$
\sum \xi_i \otimes \zeta_i \longmapsto \sum \xi_i \cup \zeta_i
$$
dans $H^{i + j}(\text{Tot}(K^\bullet \otimes_R L^\bullet))$.

\medskip\noindent
Soit $R \to A$ un homomorphisme d'anneaux et soient
$K^\bullet, L^\bullet \in D(R)$. On a une identification canonique
\begin{equation}
\label{equation-pullback-derived-tensor-product}
(K^\bullet \otimes_R^{\mathbf{L}} A)
\otimes_A^{\mathbf{L}}
(L^\bullet \otimes_R^{\mathbf{L}} A)
=
(K^\bullet \otimes_R^{\mathbf{L}} L^\bullet) \otimes_R^{\mathbf{L}} A
\end{equation}
dans $D(A)$. Elle se construit comme suit. Choisissons d'abord des
résolutions K-plates $P^\bullet \to K^\bullet$ et
$Q^\bullet \to L^\bullet$ sur $R$. Le membre de gauche est alors représenté
par le complexe
$\text{Tot}((P^\bullet \otimes_R A) \otimes_A (Q^\bullet \otimes_R A))$
et le membre de droite par le complexe
$\text{Tot}(P^\bullet \otimes_R Q^\bullet) \otimes_R A$. Ces complexes sont
canoniquement isomorphes. La construction précédente induit donc des
produits
$$
\text{Tor}^R_n(K^\bullet, A) \otimes_A \text{Tor}^R_m(L^\bullet, A)
\longrightarrow
\text{Tor}_{n + m}^R(K^\bullet \otimes_R^\mathbf{L} L^\bullet, A)
$$
qui sont parfois utiles.

\medskip\noindent
Soient $M$, $N$ des $R$-modules. À l'aide de la construction générale
ci-dessus, du morphisme canonique
$M \otimes_R^\mathbf{L} N \to M \otimes_R N$ et de la fonctorialité de
$\text{Tor}$, on obtient des morphismes canoniques
\begin{equation}
\label{equation-tor-product}
\text{Tor}^R_n(M, A) \otimes_A \text{Tor}^R_m(N, A)
\longrightarrow \text{Tor}_{n + m}^R(M \otimes_R N, A)
\end{equation}
Voici une construction directe à l'aide de résolutions projectives.
Choisissons d'abord des résolutions projectives
$$
P_\bullet \to M, \quad Q_\bullet \to N, \quad T_\bullet \to M \otimes_R N
$$
sur $R$. On a
$H_0(\text{Tot}(P_\bullet \otimes_R Q_\bullet)) = M \otimes_R N$ by
exactitude à droite de $\otimes_R$. Par conséquent, Catégories dérivées,
Lemmes \ref{derived-lemma-morphisms-lift-projective} et
\ref{derived-lemma-morphisms-equal-up-to-homotopy-projective}, garantissent
l'existence et l'unicité d'un morphisme de complexes
$\mu : \text{Tot}(P_\bullet \otimes_R Q_\bullet) \to T_\bullet$ tel que
$H_0(\mu) = \text{id}_{M \otimes_R N}$. Il induit un morphisme canonique
\begin{align*}
(M \otimes_R^{\mathbf{L}} A) \otimes_A^{\mathbf{L}}
(N \otimes_R^{\mathbf{L}} A)
& =
\text{Tot}((P_\bullet \otimes_R A) \otimes_A (Q_\bullet \otimes_R A)) \\
& =
\text{Tot}(P_\bullet \otimes_R Q_\bullet) \otimes_R A \\
& \to
T_\bullet \otimes_R A \\
& = (M \otimes_R N) \otimes_R^{\mathbf{L}} A
\end{align*}
dans $D(A)$. Ainsi, les produits (\ref{equation-tor-product}) ci-dessus se
construisent en utilisant (\ref{equation-simple-tor-product}) sur $A$ pour
construire
$$
\text{Tor}^R_n(M, A) \otimes_A \text{Tor}^R_m(N, A) \to
H^{-n-m}((M \otimes_R^{\mathbf{L}} A) \otimes_A^{\mathbf{L}}
(N \otimes_R^{\mathbf{L}} A))
$$
puis en composant avec le morphisme affiché ci-dessus pour aboutir à
$\text{Tor}_{n + m}^R(M \otimes_R N, A)$.

\medskip\noindent
Un cas particulier intéressant se présente lorsque $M = N = B$, où $B$
est une $R$-algèbre. On obtient alors des morphismes
$$
\text{Tor}_n^R(B, A) \otimes_A \text{Tor}_m^R(B, A)
\longrightarrow
\text{Tor}_{n + m}^R(B \otimes_R B, A)
\longrightarrow
\text{Tor}_{n + m}^R(B, A)
$$
la seconde flèche étant induite par le morphisme de multiplication
$B \otimes_R B \to B$ grâce à la fonctorialité de $\text{Tor}$. Autrement
dit, on obtient une structure de $A$-algèbre sur
$\text{Tor}^R_{\star}(B, A)$. Cette structure possède de nombreuses propriétés intéressantes
(associativité, commutativité graduée, structure de $B$-algèbre, puissances
divisées dans certains cas, etc.) que nous étudierons ailleurs (insérer ici
une référence future).

\begin{lemma}
\label{lemma-functoriality-product-tor}
Soit $R$ un anneau. Soient $A, B, C$ des $R$-algèbres et soit $B \to C$ un
homomorphisme de $R$-algèbres. Alors le morphisme induit
$$
\text{Tor}^R_{\star}(B, A)
\longrightarrow
\text{Tor}^R_{\star}(C, A)
$$
est un homomorphisme de $A$-algèbres.
\end{lemma}

\begin{proof}
Omis. Indication : on peut le démontrer en déroulant les définitions et en
écrivant explicitement tous les complexes.
\end{proof}








\section{Suite spectrale de K\"unneth}
\label{section-kunneth-spectral-sequence}

\noindent
Soit $R$ un anneau et soient $K^\bullet$ et $L^\bullet$ des complexes
filtrés de $R$-modules (voir Homologie, Définition
\ref{homology-definition-filtered-complex} ; nos filtrations sont
décroissantes). Alors le complexe
$$
T^\bullet = \text{Tot}(K^\bullet \otimes_R L^\bullet)
$$
possède aussi une filtration décroissante définie par la formule
$$
F^nT^\bullet =
\Im\left(
\bigoplus\nolimits_{i + j = n}
\text{Tot}(F^iK^\bullet \otimes_R F^jL^\bullet) \to
\text{Tot}(K^\bullet \otimes_R L^\bullet)
\right)
$$
Sous certaines hypothèses sur nos complexes filtrés, nous déterminerons la
suite spectrale qui en résulte par la construction de Homologie, Section
\ref{homology-section-filtered-complex}.

\medskip\noindent
Supposons que chacun des modules $K^n$, $F^iK^n$, $\text{gr}^iK^n$,
$L^m$, $F^jL^m$ et $\text{gr}^jL^m$ soit un $R$-module plat. Alors
$F^iK^n \otimes_R L^m$, $K^n \otimes_R F^jL^m$ et
$F^iK^n \otimes_R F^jL^m$ sont des sous-modules de
$K^n \otimes_R L^m$. De même,
$\text{gr}^iK^n \otimes_R \text{gr}^jL^m$ est un sous-module de
$K^n/F^{i + 1}K^n \otimes_R L^m/F^{j + 1}L^m$. Considérons le morphisme
$$
F^nT^n \longrightarrow
\bigoplus\nolimits_{i + j = n}
\text{Tot}(K^\bullet / F^{i + 1}K^\bullet
\otimes_R L^\bullet / F^{j + 1}L^\bullet)
$$
Nous laissons au lecteur le soin de montrer que l'on aboutit bien à la
somme directe, et non au produit direct, en raison de notre définition de
$F^nT^\bullet$. Pour $a + b = n$, la restriction de la flèche affichée au
sous-complexe $\text{Tot}(F^aK^\bullet \otimes_R F^bL^\bullet)$ de
$F^nT^\bullet$ est à valeurs dans le facteur direct correspondant à
$i = a$ et $j = b$. De plus, par nos hypothèses de platitude, l'image est
isomorphe à
$\text{Tot}(\text{gr}^iK^\bullet \otimes_R \text{gr}^jL^\bullet)$
et le noyau de $\text{Tot}(F^iK^\bullet \otimes_R F^jL^\bullet) \to
\text{Tot}(\text{gr}^iK^\bullet \otimes_R \text{gr}^jL^\bullet)$ est
$\text{Tot}(F^{i + 1}K^\bullet \otimes_R F^jL^\bullet) +
\text{Tot}(F^iK^\bullet \otimes_R F^{j + 1}L^\bullet)$.
Il en résulte, sous nos hypothèses, une suite exacte courte de complexes
$$
0 \to F^{n + 1}T^\bullet \to F^nT^\bullet \to
\bigoplus\nolimits_{i + j = n}
\text{Tot}(\text{gr}^iK^\bullet \otimes_R \text{gr}^jL^\bullet) \to 0
$$
Autrement dit, $\text{gr}^nT^\bullet$ est une somme directe des complexes
$\text{Tot}(\text{gr}^iK^\bullet \otimes_R \text{gr}^jL^\bullet)$
pour $i + j = n$.

\medskip\noindent
Supposons en outre que les complexes de $R$-modules $K^\bullet$,
$F^iK^\bullet$, $\text{gr}^iK^\bullet$, $L^\bullet$, $F^jL^\bullet$ et
$\text{gr}^jL^\bullet$ soient K-plats. On en déduit que la suite spectrale
de Homologie, Section \ref{homology-section-filtered-complex}, associée au
complexe filtré $T^\bullet$, a pour termes
$$
E_1^{p, q} =
\bigoplus\nolimits_{i + j = p}
H^{p + q}(\text{gr}^iK^\bullet \otimes_R^\mathbf{L} \text{gr}^jL^\bullet)
$$
avec des différentielles induites par les suites exactes courtes
$$
0 \to \text{gr}^{n + 1}T^\bullet \to
F^nT^\bullet/F^{n + 2}T^\bullet \to
\text{gr}^nT^\bullet \to 0
$$
comme expliqué dans Homologie, Lemme
\ref{homology-lemma-spectral-sequence-filtered-complex-d1}. En particulier,
le lecteur peut vérifier que le facteur direct
$H^{p + q}(\text{gr}^iK^\bullet \otimes_R^\mathbf{L} \text{gr}^jL^\bullet)$
de $E_1^{p, q}$ s'envoie dans la somme
$$
H^{p + q + 1}(\text{gr}^{i + 1}K^\bullet \otimes_R^\mathbf{L}
\text{gr}^jL^\bullet) \oplus
H^{p + q + 1}(\text{gr}^iK^\bullet \otimes_R^\mathbf{L}
\text{gr}^{j + 1}L^\bullet)
$$
contenue dans $E_1^{p + 1, q}$.

\begin{lemma}
\label{lemma-kunneth-converges}
Sous les hypothèses ci-dessus, supposons en outre que
\begin{enumerate}
\item la filtration sur $K^\bullet$ et celle sur $L^\bullet$ soient finies ;
ou que
\item plus généralement, les conditions suivantes soient satisfaites :
\begin{enumerate}
\item $F^iK^\bullet$ est acyclique pour $i \gg 0$ ;
\item $F^iK^\bullet \to K^\bullet$ est un quasi-isomorphisme pour
$i \ll 0$ ;
\item $F^jL^\bullet$ est acyclique pour $j \gg 0$ ;
\item $F^jL^\bullet \to L^\bullet$ est un quasi-isomorphisme pour
$j \ll 0$.
\end{enumerate}
\end{enumerate}
Alors la suite spectrale est bornée, la filtration associée sur chaque
$H^n(T^\bullet) = H^n(K^\bullet \otimes_R^\mathbf{L} L^\bullet)$ est finie,
et l'on a la convergence
$$
E_1^{p, q} =
\bigoplus\nolimits_{i + j = p}
H^{p + q}(\text{gr}^iK^\bullet \otimes_R^\mathbf{L} \text{gr}^jL^\bullet)
\Rightarrow
H^n(K^\bullet \otimes_R^\mathbf{L} L^\bullet)
$$
\end{lemma}

\begin{proof}
Dans le cas (1), la filtration sur $T^\bullet$ est finie et le lemme résulte
immédiatement de Homologie, Lemme
\ref{homology-lemma-biregular-ss-converges}. Dans le cas (2), choisissons
$a < b$ tels que $F^iK^\bullet$ et $F^jL^\bullet$ soient acycliques pour
$i, j > b$, et que $F^iK^\bullet \to K^\bullet$ et
$F^jL^\bullet \to L^\bullet$ soient des quasi-isomorphismes pour
$i, j < a$. Nous affirmons alors que $F^nT^\bullet$ est acyclique pour
$n > 2b$ et que $F^nT^\bullet \to T^\bullet$ est un quasi-isomorphisme pour
$n < 2a - 1$. Cette affirmation permet d'appliquer Homologie, Lemme
\ref{homology-lemma-ss-converges-trivial}, et démontre donc le lemme.

\medskip\noindent
Démontrons l'affirmation. Étant donnés $n$ et des entiers $i_1 \leq i_2$,
considérons
$$
S^{n, i_1, i_2} =
\Im\left(
\bigoplus\nolimits_{i_1 \leq i \leq i_2}
\text{Tot}(F^iK^\bullet \otimes_R F^{n - i}L^\bullet)
\to
\text{Tot}(K^\bullet \otimes_R L^\bullet)
\right)
$$
vu comme sous-complexe de $\text{Tot}(K^\bullet \otimes_R L^\bullet)$.
Remarquons que $F^nT^\bullet = \colim S^{n, i_1, i_2}$ est une colimite
filtrante. Pour établir l'affirmation, il suffit donc de montrer que
$S^{n, i_1, i_2}$ est acyclique pour $n > 2b$ et que
$S^{n, i_1, i_2} \to \text{Tot}(K^\bullet \otimes L^\bullet)$
est un quasi-isomorphisme pour $n < 2a - 1$ et pour un ensemble cofinal de
choix des paires $i_1, i_2$.

\medskip\noindent
Si $i_1 = i_2 = i$ et $n > 2b$, alors
$$
S^{n, i, i} =
\text{Tot}(F^iK^\bullet \otimes_R F^{n - i}L^\bullet)
$$
est acyclique : soit $i > b$ et ce complexe est acyclique parce que
$F^{n - i}L^\bullet$ est K-plat, soit $i \leq b$, auquel cas
$n - i \geq n - b > b$ et il en va de même. À l'aide de nos hypothèses de
platitude, le lecteur vérifie qu'il existe une suite exacte courte
$$
0 \to
S^{n + 1, i_2 + 1, i_2 + 1}
\to
S^{n, i_1, i_2} \oplus S^{n, i_2 + 1, i_2 + 1}
\to
S^{n, i_1, i_2 + 1}
\to
0
$$
de complexes. On voit donc, par récurrence sur $i_2 - i_1$, que tous les
complexes $S^{n, i_1, i_2}$ sont acycliques pour $n > 2b$.

\medskip\noindent
Supposons $n < 2a - 1$. Remarquons que
$$
S^{n, a - 1, a - 1} =
\text{Tot}(F^{a - 1}K^\bullet \otimes_R F^{n - a + 1}L^\bullet)
\to \text{Tot}(K^\bullet \otimes_R L^\bullet)
$$
est un quasi-isomorphisme, car $a - 1 < a$ et $n - a + 1 < a$, et tous nos
complexes sont K-plats ; les deux membres calculent donc le même objet de
$D(R)$. Remarquons que, pour $i_2 \geq a - 1$ et $n < 2a - 1$, le morphisme
$$
S^{n + 1, i_2 + 1, i_2 + 1} \to S^{n, i_2 + 1, i_2 + 1}
$$
est un quasi-isomorphisme, car ses deux membres sont quasi-isomorphes à
$\text{Tot}(F^{i_2 + 1}K^\bullet \otimes L^\bullet)$. En utilisant la suite
exacte courte de complexes du paragraphe précédent, on en déduit que
$S^{n, a - 1, a - 1} \to S^{n, a - 1, t}$ est un quasi-isomorphisme pour
tout $t \geq a - 1$. D'autre part, on a de même des suites exactes courtes
$$
0 \to
S^{n + 1, i_1, i_1}
\to
S^{n, i_1 - 1, i_1 - 1} \oplus
S^{n, i_1, i_2}
\to
S^{n, i_1 - 1, i_2}
\to
0
$$
de complexes. Comme ci-dessus, on remarque que, pour $i_1 \leq a - 1$, le
morphisme
$$
S^{n + 1, i_1, i_1} \to S^{n, i_1 - 1, i_1 - 1}
$$
est un quasi-isomorphisme, car ses deux membres sont quasi-isomorphes à
$\text{Tot}(K^\bullet \otimes F^{n - i_1 + 1}L^\bullet)$. On en déduit que
$S^{n, s, t} \to S^{n, a - 1, t}$ est un quasi-isomorphisme pour tout
$s \leq a - 1$. En combinant ces résultats, on trouve que
$S^{n, s, t}$, pour $s \leq a - 1$ et $t \geq a - 1$, s'envoie par un
quasi-isomorphisme vers $\text{Tot}(K^\bullet \otimes_R L^\bullet)$. Cela
achève la démonstration du lemme.
\end{proof}

\begin{lemma}
\label{lemma-resolution-filtered-complex}
Soit $R$ un anneau et soit $K^\bullet$ un complexe filtré. Il existe un
morphisme $f : P^\bullet \to K^\bullet$ de complexes filtrés tel que
\begin{enumerate}
\item chacun des modules $P^n$, $F^iP^n$, $\text{gr}^iP^n$ est un
$R$-module libre ;
\item les complexes de $R$-modules $P^\bullet$, $F^iP^\bullet$ et
$\text{gr}^iP^\bullet$ sont K-plats ;
\item $f$ induit des quasi-isomorphismes
$P^\bullet \to K^\bullet$,
$F^iP^\bullet \to F^iK^\bullet$ et
$\text{gr}^iP^\bullet \to \text{gr}^iK^\bullet$.
\end{enumerate}
\end{lemma}

\begin{proof}
Disons qu'un complexe filtré $L^\bullet$ est {\it élémentaire} si chacun
des modules $L^n$, $F^iL^n$, $\text{gr}^iL^n$ est un $R$-module libre et si
toutes les différentielles sont nulles.

\medskip\noindent
Il existe un complexe filtré élémentaire $P_0^\bullet$ et un morphisme de
complexes filtrés $f_0 : P_0^\bullet \to K^\bullet$ tels que $f_0$ et
$F^if_0$, $i \in \mathbf{Z}$, soient surjectifs en cohomologie. Pour le
voir, posons
$$
P_0^n = \bigoplus\nolimits_{z \in \Ker(d_K^n)} R \xi_z
\oplus
\bigoplus\nolimits_{j \in \mathbf{Z}}
\bigoplus\nolimits_{z \in \Ker(d_{F^jK}^n)} R \xi_z
$$
avec différentielles nulles et le morphisme $f_0$ défini par
$f_0(\xi_z) = z$. Pour la filtration, on pose
$$
F^iP_0^n = \bigoplus\nolimits_{j \geq i \in \mathbf{Z}}
\bigoplus\nolimits_{z \in \Ker(d_{F^jK}^n)} R \xi_z
$$
Nous laissons au lecteur le soin de vérifier que l'on obtient ainsi le
morphisme $f_0 : P_0^\bullet \to K^\bullet$ annoncé.

\medskip\noindent
Par récurrence sur $m \geq 0$, nous allons construire des plongements
$$
P_0^\bullet \subset \ldots P_m^\bullet \subset P_{m + 1}^\bullet
$$
de complexes filtrés et des morphismes $f_m : P_m^\bullet \to K^\bullet$
possédant les propriétés suivantes :
\begin{enumerate}
\item le complexe filtré $P_{m + 1}^\bullet / P_m^\bullet$ est
élémentaire ;
\item les noyaux de $H^n(f_m)$ et de $H^n(F^if_m)$ s'envoient sur zéro
dans $H^n(P_{m + 1}^\bullet)$ et $H^n(F^iP_{m + 1}^\bullet)$ ;
\item le morphisme $f_{m + 1} : P_{m + 1}^\bullet \to K^\bullet$ prolonge
$f_m$.
\end{enumerate}
Pour cela, posons
$$
\Omega_{n, m} = \Ker\left(\Ker(d^n_{P_m}) \to H^n(K^\bullet)\right),
\text{ resp. }
\Omega_{n, i, m} = \Ker\left(\Ker(d^n_{F^iP_m}) \to H^n(F^iK^\bullet)\right),
$$
Remarquons que $\Omega_{n, m}$ se surjecte sur le noyau de $H^n(f_m)$ et
que $\Omega_{n, i, m}$ se surjecte sur le noyau de $H^n(F^if_m)$. Pour
chaque $z \in \Omega_{n, m}$, resp.\ $z \in \Omega_{n, i, m}$, choisissons
un $y_z \in K^{n - 1}$, resp.\ $y_z \in F^iK^{n - 1}$, tel que
$d_K(y_z) = f_m(z)$. Posons alors
$$
P^n_{m + 1} = P_m^n \oplus
\bigoplus\nolimits_{z \in \Omega_{n + 1, m}} R \eta_z
\oplus
\bigoplus\nolimits_{j \in \mathbf{Z}}
\bigoplus\nolimits_{z \in \Omega_{n + 1, j, m}} R \eta_z
$$
On pose $d_{P_{m + 1}}(\eta_z) = z$ et $f_{m + 1}(\eta_z) = y_z$. Enfin, on pose
$$
F^iP_{m + 1}^n = F^iP_m^n \oplus
\bigoplus\nolimits_{j \geq i}
\bigoplus\nolimits_{z \in \Omega_{n + 1, j, m}} R \eta_z
$$
Nous laissons au lecteur le soin de vérifier que l'on obtient ainsi
$P_m^\bullet \subset P_{m + 1}^\bullet$ et
$f_{m + 1} : P_{m + 1}^\bullet \to K^\bullet$ comme annoncé.

\medskip\noindent
Il suffit maintenant de prendre
$$
P^\bullet = \bigcup P_m^\bullet
$$
comme complexe filtré, muni du morphisme $f : P^\bullet \to K^\bullet$ donné
par $\bigcup f_m$.

\medskip\noindent
L'assertion (1) du lemme vaut parce que $P^n$, comme module filtré, est
isomorphe à la somme directe de $P_0^n$ et des $P_{m + 1}^n/P_m^n$ pour
$m \geq 0$. Un petit détail est omis.

\medskip\noindent
Passons à l'assertion (3). Remarquons que
$H^n(P^\bullet) = \colim H^n(P_m^\bullet)$.
Le morphisme $f_0$ est surjectif en cohomologie ; il en va donc de même de
chaque $f_m$. Pour chaque $m$, le plongement
$P_m^\bullet \subset P_{m + 1}^\bullet$ annule par construction le noyau de
$H^n(f_m)$. En combinant ces faits, le lecteur conclut aisément que
$H^n(f)$ est un isomorphisme. Il en va de même de $H^n(F^if)$. Alors
$H^n(\text{gr}^if)$ est aussi nécessairement un isomorphisme, en vertu de la
suite exacte courte $0 \to F^{i + 1} \to F^i \to \text{gr}^i \to 0$ (de
foncteurs sur la catégorie des complexes filtrés, par exemple). Un petit
détail est omis.

\medskip\noindent
Passons à l'assertion (2). Pour voir que $P^\bullet$ est K-plat, il suffit,
d'après le Lemme \ref{lemma-colimit-K-flat}, de montrer que $P_m^\bullet$
est K-plat. D'après le Lemme \ref{lemma-K-flat-two-out-of-three-ses} et par
récurrence, il suffit de remarquer qu'un complexe à différentielles nulles
et à termes libres est K-plat. Le même argument montre que
$F^iP^\bullet$ est K-plat pour tout $i \in \mathbf{Z}$. Enfin,
$\text{gr}^iP^\bullet$ est K-plat par une nouvelle application du Lemme
\ref{lemma-K-flat-two-out-of-three-ses}.
\end{proof}

\begin{proposition}
\label{proposition-kunneth-filtered}
Soit $R$ un anneau et soient $K^\bullet$ et $L^\bullet$ des complexes
filtrés de $R$-modules. Il existe alors un complexe filtré $T^\bullet$
représentant $K^\bullet \otimes_R^\mathbf{L} L^\bullet$ dans $D(R)$, tel que
la suite spectrale associée ait pour page $E_1$
$$
E_1^{p, q} =
\bigoplus\nolimits_{i + j = p}
H^{p + q}(\text{gr}^iK^\bullet \otimes_R^\mathbf{L} \text{gr}^jL^\bullet)
$$
Si
\begin{enumerate}
\item $F^iK^\bullet$ est acyclique pour $i \gg 0$ ;
\item $F^iK^\bullet \to K^\bullet$ est un quasi-isomorphisme pour
$i \ll 0$ ;
\item $F^jL^\bullet$ est acyclique pour $j \gg 0$ ;
\item $F^jL^\bullet \to L^\bullet$ est un quasi-isomorphisme pour
$j \ll 0$.
\end{enumerate}
alors la suite spectrale est bornée, la filtration associée sur chaque
$H^n(K^\bullet \otimes_R^\mathbf{L} L^\bullet)$ est finie et la suite
spectrale converge.
\end{proposition}

\begin{proof}
Choisissons $P^\bullet \to K^\bullet$ et $Q^\bullet \to L^\bullet$ comme
dans le Lemme \ref{lemma-resolution-filtered-complex}. On utilise alors la
suite spectrale du complexe filtré
$$
T^\bullet = \text{Tot}(P^\bullet \otimes_R Q^\bullet)
$$
décrite dans cette section et dans le Lemme
\ref{lemma-kunneth-converges}.
\end{proof}

\begin{lemma}[Suite spectrale de K\"unneth]
\label{lemma-kunneth-spectral-sequence}
Soit $R$ un anneau et soient $K$ et $L$ des objets de $D^b(R)$. Il existe
une suite spectrale bigraduée bornée $\{E_r\}_{r \geq 2}$ telle que
$$
E_2^{p, q} = \bigoplus\nolimits_{i + j = q} \text{Tor}^R_{-p}(H^i(K), H^j(L))
$$
et dont $d_r$ est de bidegré $(r, -r + 1)$, convergeant vers
$H^{p + q}(K \otimes_R^\mathbf{L} L)$.
\end{lemma}

\begin{proof}
Soit $K^\bullet$ un complexe de $R$-modules représentant $K$ et soit
$L^\bullet$ un complexe de $R$-modules représentant $L$. Posons
$F^iK^\bullet = \tau_{\leq -i}K^\bullet$, et de même pour $L$ ; voir
Homologie, Section \ref{homology-section-truncations}. Appliquons la
Proposition \ref{proposition-kunneth-filtered} en remarquant que
$\text{gr}^iK^\bullet = H^{-i}(K)[i]$. On obtient une suite spectrale bornée
$\{(E')_r\}_{r \geq 1}$ telle que
$$
(E')^{p, q}_1 =
\bigoplus\nolimits_{i + j = p} \text{Tor}^R_{-2p - q}(H^{-i}(K), H^{-j}(L))
$$
qui converge vers la cohomologie de $K \otimes_R^\mathbf{L} L$. Par
construction, cette suite spectrale a pour différentielles
$d_r^{p, q} : (E')_r^{p, q} \to (E')_r^{p + r, q - r + 1}$ pour $r \geq 1$.
Pour obtenir la suite spectrale du lemme, on pose, pour $r \geq 2$,
$$
E_r^{p, q} = (E')_{r - 1}^{-q, p + 2q}
$$
Nous laissons au lecteur le soin de vérifier que cela convient.
\end{proof}

\begin{example}
\label{example-kunneth-dedekind}
Si $R = \mathbf{Z}$ ou, plus généralement, si $R$ est un anneau de
Dedekind, alors $\text{Tor}_i^R(M, N) = 0$ pour $i \not \in \{0, 1\}$ et
pour tous $R$-modules $M$ et $N$. La suite spectrale du Lemme
\ref{lemma-kunneth-spectral-sequence} dégénère donc en page $E_2$, et l'on
obtient des suites exactes courtes
$$
0 \to \bigoplus_{i + j = n} H^i(K) \otimes_R H^j(L) \to
H^n(K \otimes_R^\mathbf{L} L) \to
\bigoplus_{i + j = n + 1} \text{Tor}^R_1(H^i(K), H^j(L)) \to 0
$$
pour tout $n \in \mathbf{Z}$.
\end{example}






\section{Modules pseudo-cohérents, I}
\label{section-pseudo-coherent}

\noindent
Supposons que $R$ soit un anneau. Rappelons qu'un $R$-module $M$ est de type
fini s'il existe une surjection $R^{\oplus a} \to M$, et de présentation
finie s'il existe une présentation
$R^{\oplus a_1} \to R^{\oplus a_0} \to M \to 0$. De même, on peut
considérer les $R$-modules qui possèdent une résolution de longueur $n$
\begin{equation}
\label{equation-pseudo-coherent}
R^{\oplus a_n} \to R^{\oplus a_{n - 1}} \to \ldots \to R^{\oplus a_0} \to
M \to 0
\end{equation}
par des $R$-modules libres de rang fini. Un module est dit
{\it pseudo-cohérent} si l'on peut trouver une telle résolution pour tout
$n$. Voici la définition formelle.

\begin{definition}
\label{definition-pseudo-coherent}
Soit $R$ un anneau. Notons $D(R)$ sa catégorie dérivée. Soit
$m \in \mathbf{Z}$.
\begin{enumerate}
\item Un objet $K^\bullet$ de $D(R)$ est {\it $m$-pseudo-cohérent} s'il
existe un complexe borné $E^\bullet$ de $R$-modules libres de rang fini et
un morphisme $\alpha : E^\bullet \to K^\bullet$ tel que $H^i(\alpha)$ soit
un isomorphisme pour $i > m$ et que $H^m(\alpha)$ soit surjectif.
\item Un objet $K^\bullet$ de $D(R)$ est {\it pseudo-cohérent} s'il est
quasi-isomorphe à un complexe borné supérieurement de $R$-modules libres de
rang fini.
\item Un $R$-module $M$ est dit {\it $m$-pseudo-cohérent} si $M[0]$ est un
objet $m$-pseudo-cohérent de $D(R)$.
\item Un $R$-module $M$ est dit
{\it pseudo-cohérent}\footnote{Cette terminologie diffère de celle de
module pseudo-cohérent dans \cite{Bourbaki-CA}.} si $M[0]$ est un objet
pseudo-cohérent de $D(R)$.
\end{enumerate}
\end{definition}

\noindent
Comme d'habitude, nous appliquons aussi cette terminologie aux complexes de
$R$-modules. Puisque tout morphisme $E^\bullet \to K^\bullet$ dans $D(R)$
est représenté par un véritable morphisme de complexes, voir Catégories
dérivées, Lemme \ref{derived-lemma-morphisms-from-projective-complex}, il
n'y a aucune ambiguïté. Il se trouve que $K^\bullet$ est pseudo-cohérent si
et seulement si $K^\bullet$ est $m$-pseudo-cohérent pour tout
$m \in \mathbf{Z}$ ; voir le Lemme \ref{lemma-pseudo-coherent}. De plus, si
l'anneau est noethérien, la condition s'interprète comme une condition de
type fini sur la cohomologie ; voir le Lemme
\ref{lemma-Noetherian-pseudo-coherent}. Commençons par relier ceci à la
discussion informelle ci-dessus.

\begin{lemma}
\label{lemma-cone-pseudo-coherent}
Soit $R$ un anneau, soit $m \in \mathbf{Z}$ et soit
$(K^\bullet, L^\bullet, M^\bullet, f, g, h)$ un triangle distingué de
$D(R)$.
\begin{enumerate}
\item Si $K^\bullet$ est $(m + 1)$-pseudo-cohérent et si $L^\bullet$ est
$m$-pseudo-cohérent, alors $M^\bullet$ est $m$-pseudo-cohérent.
\item Si $K^\bullet, M^\bullet$ sont $m$-pseudo-cohérents, alors
$L^\bullet$ est $m$-pseudo-cohérent.
\item Si $L^\bullet$ est $(m + 1)$-pseudo-cohérent et si $M^\bullet$ est
$m$-pseudo-cohérent, alors $K^\bullet$ est $(m + 1)$-pseudo-cohérent.
\end{enumerate}
\end{lemma}

\begin{proof}
Démontrons (1). Choisissons $\alpha : P^\bullet \to K^\bullet$, où
$P^\bullet$ est un complexe borné de modules libres de rang fini, tel que
$H^i(\alpha)$ soit un isomorphisme pour $i > m + 1$ et soit surjectif pour
$i = m + 1$. On peut remplacer $P^\bullet$ par
$\sigma_{\geq m + 1}P^\bullet$ et donc supposer que $P^i = 0$ pour
$i < m + 1$. Choisissons $\beta : E^\bullet \to L^\bullet$, où
$E^\bullet$ est un complexe borné de modules libres de rang fini, tel que
$H^i(\beta)$ soit un isomorphisme pour $i > m$ et soit surjectif pour
$i = m$. D'après Catégories dérivées, Lemme
\ref{derived-lemma-lift-map}, on peut trouver un morphisme
$\gamma : P^\bullet \to E^\bullet$ tel que le diagramme
$$
\xymatrix{
K^\bullet \ar[r] & L^\bullet \\
P^\bullet \ar[u] \ar[r]^\gamma & E^\bullet \ar[u]_\beta
}
$$
soit commutatif dans $D(R)$. Le cône $C(\gamma)^\bullet$ est un complexe
borné de $R$-modules libres de rang fini, et la commutativité du diagramme
implique qu'il existe un morphisme de triangles distingués
$$
(P^\bullet, E^\bullet, C(\gamma)^\bullet)
\longrightarrow
(K^\bullet, L^\bullet, M^\bullet).
$$
Le morphisme induit sur les suites exactes longues de cohomologie, ainsi que
Homologie, Lemmes \ref{homology-lemma-four-lemma} et
\ref{homology-lemma-five-lemma}, montrent que
$C(\gamma)^\bullet \to M^\bullet$ induit un isomorphisme en cohomologie aux
degrés $> m$ et une surjection au degré $m$. Ainsi, $M^\bullet$ est
$m$-pseudo-cohérent.

\medskip\noindent
Les assertions (2) et (3) résultent de (1) par rotation du triangle
distingué.
\end{proof}

\begin{lemma}
\label{lemma-finite-cohomology}
Soit $R$ un anneau, soit $K^\bullet$ un complexe de $R$-modules et soit
$m \in \mathbf{Z}$.
\begin{enumerate}
\item Si $K^\bullet$ est $m$-pseudo-cohérent et si $H^i(K^\bullet) = 0$
pour $i > m$, alors $H^m(K^\bullet)$ est un $R$-module de type fini.
\item Si $K^\bullet$ est $m$-pseudo-cohérent et si $H^i(K^\bullet) = 0$
pour $i > m + 1$, alors $H^{m + 1}(K^\bullet)$ est un $R$-module de
présentation finie.
\end{enumerate}
\end{lemma}

\begin{proof}
Démontrons (1). Choisissons un complexe borné $E^\bullet$ de $R$-modules
projectifs de type fini et un morphisme $\alpha : E^\bullet \to K^\bullet$
qui induit un isomorphisme en cohomologie aux degrés $> m$ et une surjection
au degré $m$. Il suffit manifestement de démontrer le résultat pour
$E^\bullet$. Soit $n$ le plus grand entier tel que $E^n \not = 0$. Si
$n = m$, le résultat est clair. Si $n > m$, alors
$E^{n - 1} \to E^n$ est surjectif puisque $H^n(E^\bullet) = 0$. Comme $E^n$
est projectif de type fini, on a $E^{n - 1} = E' \oplus E^n$. Il suffit donc
de démontrer le résultat pour le complexe $(E')^\bullet$, qui est identique
à $E^\bullet$ sauf qu'il vaut $E'$ au degré $n - 1$ et $0$ au degré $n$.
On conclut par récurrence sur $n$.

\medskip\noindent
Démontrons (2). Choisissons un complexe borné $E^\bullet$ de $R$-modules
projectifs de type fini et un morphisme $\alpha : E^\bullet \to K^\bullet$
qui induit un isomorphisme en cohomologie aux degrés $> m$ et une surjection
au degré $m$. Comme dans la démonstration de (1), on se ramène au cas où
$E^i = 0$ pour $i > m + 1$. Alors
$H^{m + 1}(K^\bullet) \cong
H^{m + 1}(E^\bullet) = \Coker(E^m \to E^{m + 1})$
qui est de présentation finie.
\end{proof}

\begin{lemma}
\label{lemma-n-pseudo-module}
Soit $R$ un anneau et soit $M$ un $R$-module. Alors
\begin{enumerate}
\item $M$ est $0$-pseudo-cohérent si et seulement si $M$ est un $R$-module
de type fini ;
\item $M$ est $(-1)$-pseudo-cohérent si et seulement si $M$ est un
$R$-module de présentation finie ;
\item $M$ est $(-d)$-pseudo-cohérent si et seulement s'il existe une
résolution
$$
R^{\oplus a_d} \to R^{\oplus a_{d - 1}} \to \ldots \to R^{\oplus a_0} \to
M \to 0
$$
de longueur $d$ ;
\item $M$ est pseudo-cohérent si et seulement s'il existe une résolution
infinie
$$
\ldots \to R^{\oplus a_1} \to R^{\oplus a_0} \to M \to 0
$$
par des $R$-modules libres de rang fini.
\end{enumerate}
\end{lemma}

\begin{proof}
Si $M$ est de type fini (resp.\ de présentation finie), alors $M$ est
$0$-pseudo-cohérent (resp.\ $(-1)$-pseudo-cohérent), d'après la discussion
qui précède la Définition \ref{definition-pseudo-coherent}. Réciproquement,
si $M$ est $0$-pseudo-cohérent, alors $M = H^0(M[0])$ est de type fini
d'après le Lemme \ref{lemma-finite-cohomology}. Si $M$ est
$(-1)$-pseudo-cohérent, il est $0$-pseudo-cohérent, donc de type fini.
Choisissons une surjection $R^{\oplus a} \to M$ et notons
$K = \Ker(R^{\oplus a} \to M)$. D'après le Lemme
\ref{lemma-cone-pseudo-coherent}, $K$ est $0$-pseudo-cohérent, donc de type
fini, et $M$ est ainsi de présentation finie.

\medskip\noindent
Les troisième et quatrième assertions se démontrent par récurrence, avec
un argument analogue à celui qui précède (détails omis).
\end{proof}

\begin{lemma}
\label{lemma-pseudo-coherent}
Soit $R$ un anneau et soit $K^\bullet$ un complexe de $R$-modules. Les
assertions suivantes sont équivalentes :
\begin{enumerate}
\item $K^\bullet$ est pseudo-cohérent ;
\item $K^\bullet$ est $m$-pseudo-cohérent pour tout $m \in \mathbf{Z}$ ;
\item $K^\bullet$ est quasi-isomorphe à un complexe borné supérieurement de
$R$-modules projectifs de type fini.
\end{enumerate}
Si (1), (2) et (3) sont satisfaites et si $H^i(K^\bullet) = 0$ pour
$i > b$, alors il existe un quasi-isomorphisme
$F^\bullet \to K^\bullet$ où les $F^i$ sont des $R$-modules libres de rang
fini et $F^i = 0$ pour $i > b$.
\end{lemma}

\begin{proof}
On voit que (1) $\Rightarrow$ (3), puisqu'un module libre de rang fini est
un $R$-module projectif de type fini. Réciproquement, supposons que
$P^\bullet$ soit un complexe borné supérieurement de $R$-modules projectifs
de type fini. Écrivons $P^i = 0$ pour $i > n_0$. Choisissons une
décomposition en somme directe $F^{n_0} = P^{n_0} \oplus C^{n_0}$, où
$F^{n_0}$ est un $R$-module libre de rang fini, puis, par récurrence,
$$
F^{n - 1} = P^{n - 1} \oplus C^n \oplus C^{n - 1}
$$
pour $n \leq n_0$, avec $F^{n - 1}$ un $R$-module libre de rang fini. Les
Le complexe $F^\bullet$ a des morphismes $F^{n - 1} \to F^n$ qui coïncident avec
$P^{n - 1} \to P^n$, induisent l'identité $C^n \to C^n$ et sont nuls sur
$C^{n - 1}$. Le morphisme $F^\bullet \to P^\bullet$ est un
quasi-isomorphisme (même une équivalence d'homotopie), donc (3) implique
(1).

\medskip\noindent
Supposons (1). Soit $E^\bullet$ un complexe borné supérieurement de
$R$-modules libres de rang fini et soit $E^\bullet \to K^\bullet$ un
quasi-isomorphisme. Les morphismes induits
$\sigma_{\geq m}E^\bullet \to K^\bullet$, de la troncation bête de
$E^\bullet$ vers $K^\bullet$, montrent que $K^\bullet$ est
$m$-pseudo-cohérent. Ainsi, (1) implique (2).

\medskip\noindent
Supposons (2). Puisque $K^\bullet$ est $0$-pseudo-cohérent, $K^\bullet$ est en
particulier borné supérieurement. Soit $b$ un entier tel que
$H^i(K^\bullet) = 0$ pour $i > b$. Par récurrence descendante sur
$n \in \mathbf{Z}$, nous allons construire des $R$-modules libres de rang
fini $F^i$ pour $i \geq n$, des différentielles
$d^i : F^i \to F^{i + 1}$ pour $i \geq n$, et des morphismes
$\alpha : F^i \to K^i$ compatibles aux différentielles, tels que :
(1) $H^i(\alpha)$ soit un isomorphisme pour $i > n$ et soit surjectif pour
$i = n$ ; (2) $F^i = 0$ pour $i > b$. Schématiquement,
$$
\xymatrix{
& F^n \ar[r] \ar[d]^\alpha & F^{n + 1} \ar[d]^\alpha \ar[r] & \ldots \\
K^{n - 1} \ar[r] & K^n \ar[r] & K^{n + 1} \ar[r] & \ldots
}
$$
Le cas initial est $n = b + 1$, où l'on peut prendre $F^i = 0$ pour tout
$i$. Pour l'étape de récurrence, soit $C^\bullet$ le cône de $\alpha$
(Catégories dérivées, Définition \ref{derived-definition-cone}). La suite
exacte longue de cohomologie
$$
0 \to H^{n - 1}(K^\bullet) \to H^{n - 1}(C^\bullet) \to
H^n(F^\bullet) \to H^n(K^\bullet) \to H^n(C^\bullet) \to \ldots
$$
montre que $H^i(C^\bullet) = 0$ pour $i \geq n$. D'après le Lemme
\ref{lemma-cone-pseudo-coherent}, $C^\bullet$ est
$(n - 1)$-pseudo-cohérent. D'après le Lemme
\ref{lemma-finite-cohomology}, $H^{n - 1}(C^\bullet)$ est un $R$-module de
type fini. En particulier, le noyau de
$H^n(F^\bullet) \to H^n(K^\bullet)$ est un $R$-module de type fini.
Choisissons un $R$-module libre de rang fini $F^{n - 1}$ et un morphisme
$F^{n - 1} \to \Ker(F^n \to F^{n + 1})$ tel que $F^{n - 1}$ se surjecte
sur le noyau de $H^n(F^\bullet) \to H^n(K^\bullet)$. Prolongeons notre
morphisme de complexes en
$$
\xymatrix{
& F^{n - 1} \ar[d]^{\alpha^{n - 1}} \ar[r] &
F^n \ar[r] \ar[d]^\alpha &
F^{n + 1} \ar[d]^\alpha \ar[r] & \ldots \\
\ldots \ar[r] &
K^{n - 1} \ar[r] & K^n \ar[r] & K^{n + 1} \ar[r] & \ldots
}
$$
Le morphisme $\alpha^{n - 1}$ existe, car l'image de
$F^{n - 1} \to F^n \to K^n$ est contenue, par construction, dans l'image de
$K^{n - 1} \to K^n$. Notons encore $C^\bullet$ le cône de ce morphisme de
complexes prolongé. On obtient alors une suite exacte
$$
H^{n - 1}(F^\bullet) \to H^{n - 1}(K^\bullet) \to H^{n - 1}(C^\bullet) \to
H^n(F^\bullet) \cong H^n(K^\bullet) \to H^n(C^\bullet) \to \ldots
$$
Autrement dit, le conoyau de $H^{n - 1}(F^\bullet) \to
H^{n - 1}(K^\bullet)$ est un $R$-module de type fini, engendré, disons, par
les classes de $\xi_1, \ldots, \xi_r \in \Ker(K^{n - 1} \to K^n)$. On
remplace alors $F^{n - 1}$ par $F^{n - 1} \oplus R^{\oplus r}$, en envoyant
les éléments de la base du facteur libre sur zéro dans $F^n$ et sur les
$\xi_i$ dans $K^{n - 1}$. Ceci achève l'étape de récurrence.
\end{proof}

\begin{lemma}
\label{lemma-two-out-of-three-pseudo-coherent}
Soit $R$ un anneau et soit $(K^\bullet, L^\bullet, M^\bullet, f, g, h)$ un
triangle distingué de $D(R)$. Si deux des trois objets
$K^\bullet, L^\bullet, M^\bullet$ sont pseudo-cohérents, le troisième l'est
aussi.
\end{lemma}

\begin{proof}
Il suffit de combiner les Lemmes \ref{lemma-cone-pseudo-coherent} et
\ref{lemma-pseudo-coherent}.
\end{proof}

\begin{lemma}
\label{lemma-recognize-pseudo-coherent}
Soit $R$ un anneau, soit $K^\bullet$ un complexe de $R$-modules et soit
$m \in \mathbf{Z}$.
\begin{enumerate}
\item Si $H^i(K^\bullet) = 0$ pour tout $i \geq m$, alors $K^\bullet$ est
$m$-pseudo-cohérent.
\item Si $H^i(K^\bullet) = 0$ pour $i > m$ et si $H^m(K^\bullet)$ est un
$R$-module de type fini, alors $K^\bullet$ est $m$-pseudo-cohérent.
\item Si $H^i(K^\bullet) = 0$ pour $i > m + 1$, si le module
$H^{m + 1}(K^\bullet)$ est de présentation finie et si
$H^m(K^\bullet)$ est de type fini, alors $K^\bullet$ est
$m$-pseudo-cohérent.
\end{enumerate}
\end{lemma}

\begin{proof}
Il suffit de démontrer (3). Posons $M = H^{m + 1}(K^\bullet)$. Remarquons
que $\tau_{\geq m + 1}K^\bullet$ est quasi-isomorphe à $M[- m - 1]$.
D'après le Lemme \ref{lemma-n-pseudo-module}, $M[- m - 1]$ est
$m$-pseudo-cohérent. Comme on dispose du triangle distingué
$$
(\tau_{\leq m}K^\bullet, K^\bullet, \tau_{\geq m + 1}K^\bullet)
$$
(Catégories dérivées, Remarque
\ref{derived-remark-truncation-distinguished-triangle}), le Lemme
\ref{lemma-cone-pseudo-coherent} ramène la preuve à montrer que
$\tau_{\leq m}K^\bullet$ est pseudo-cohérent. Par hypothèse,
$H^m(\tau_{\leq m}K^\bullet)$ est un $R$-module de type fini. On peut donc
trouver un $R$-module libre de rang fini $E$ et un morphisme
$E \to \Ker(d_K^m)$ tels que la composée
$E \to \Ker(d_K^m) \to H^m(\tau_{\leq m}K^\bullet)$ soit surjective. Le
morphisme $E[-m] \to \tau_{\leq m}K^\bullet$ montre alors que
$\tau_{\leq m}K^\bullet$ est $m$-pseudo-cohérent.
\end{proof}

\begin{lemma}
\label{lemma-summands-pseudo-coherent}
Soit $R$ un anneau et soit $m \in \mathbf{Z}$. Si
$K^\bullet \oplus L^\bullet$ est $m$-pseudo-cohérent (resp.\
pseudo-cohérent), alors $K^\bullet$ et $L^\bullet$ le sont aussi.
\end{lemma}

\begin{proof}
Dans cette démonstration, nous omettons l'exposant ${}^\bullet$. Supposons
que $K \oplus L$ soit $m$-pseudo-cohérent. Il est clair que
$K, L \in D^{-}(R)$. Remarquons qu'il existe un triangle distingué
$$
(K \oplus L, K \oplus L, L \oplus L[1]) =
(K, K, 0) \oplus (L, L, L \oplus L[1])
$$
voir Catégories dérivées, Lemme
\ref{derived-lemma-direct-sum-triangles}. D'après le Lemme
\ref{lemma-cone-pseudo-coherent}, $L \oplus L[1]$ est
$m$-pseudo-cohérent. Ainsi, $L[1] \oplus L[2]$ est aussi
$m$-pseudo-cohérent. Par récurrence,
$L[n] \oplus L[n + 1]$ est $m$-pseudo-cohérent. D'après le Lemme
\ref{lemma-recognize-pseudo-coherent}, $L[n]$ est $m$-pseudo-cohérent pour
$n$ assez grand. En remontant ensuite, à l'aide des triangles distingués
$$
(L[n], L[n] \oplus L[n - 1], L[n - 1])
$$
on conclut que $L[n], L[n - 1], \ldots, L$ sont $m$-pseudo-cohérents, comme
voulu. Le cas pseudo-cohérent résulte de ceci et du Lemme
\ref{lemma-pseudo-coherent}.
\end{proof}

\begin{lemma}
\label{lemma-complex-pseudo-coherent-modules}
Soit $R$ un anneau, soit $m \in \mathbf{Z}$ et soit $K^\bullet$ un complexe
borné supérieurement de $R$-modules tel que $K^i$ soit
$(m - i)$-pseudo-cohérent pour tout $i$. Alors $K^\bullet$ est
$m$-pseudo-cohérent. En particulier, si $K^\bullet$ est un complexe borné
supérieurement de $R$-modules pseudo-cohérents, alors $K^\bullet$ est
pseudo-cohérent.
\end{lemma}

\begin{proof}
On peut remplacer $K^\bullet$ par $\sigma_{\geq m - 1}K^\bullet$, par
exemple, et donc supposer $K^\bullet$ borné. Le complexe $K^\bullet$ est
alors $m$-pseudo-cohérent, car chaque $K^i[-i]$ est $m$-pseudo-cohérent ; on raisonne par
récurrence sur la longueur du complexe, à l'aide du Lemme
\ref{lemma-cone-pseudo-coherent} et des troncations bêtes. Pour la dernière
assertion, il suffit de montrer que $K^\bullet$ est $m$-pseudo-cohérent pour
tout $m \in \mathbf{Z}$ ; voir le Lemme \ref{lemma-pseudo-coherent}. Cela
résulte de la première partie.
\end{proof}

\begin{lemma}
\label{lemma-cohomology-pseudo-coherent}
Soit $R$ un anneau, soit $m \in \mathbf{Z}$ et soit
$K^\bullet \in D^{-}(R)$ tel que $H^i(K^\bullet)$ soit
$(m - i)$-pseudo-cohérent (resp.\ pseudo-cohérent) pour tout $i$. Alors
$K^\bullet$ est $m$-pseudo-cohérent (resp.\ pseudo-cohérent).
\end{lemma}

\begin{proof}
Supposons que $K^\bullet$ soit un objet de $D^{-}(R)$ tel que chaque
$H^i(K^\bullet)$ soit $(m - i)$-pseudo-cohérent. Soit $n$ le plus grand
entier tel que $H^n(K^\bullet)$ soit non nul. Démontrons le lemme par
récurrence sur $n$. Si $n < m$, alors $K^\bullet$ est $m$-pseudo-cohérent
d'après le Lemme \ref{lemma-recognize-pseudo-coherent}. Si $n \geq m$, on
dispose du triangle distingué
$$
(\tau_{\leq n - 1}K^\bullet, K^\bullet, H^n(K^\bullet)[-n])
$$
(Catégories dérivées, Remarque
\ref{derived-remark-truncation-distinguished-triangle}). Puisque
$H^n(K^\bullet)[-n]$ est $m$-pseudo-cohérent par hypothèse, le Lemme
\ref{lemma-cone-pseudo-coherent} ramène la preuve à montrer que
$\tau_{\leq n - 1}K^\bullet$ est $m$-pseudo-cohérent. On conclut par
récurrence sur $n$. (Le cas pseudo-cohérent résulte de ceci et du Lemme
\ref{lemma-pseudo-coherent}.)
\end{proof}

\begin{lemma}
\label{lemma-finite-push-pseudo-coherent}
Soit $A \to B$ un homomorphisme d'anneaux. Supposons que $B$ soit
pseudo-cohérent comme $A$-module. Soit $K^\bullet$ un complexe de
$B$-modules. Les assertions suivantes sont équivalentes :
\begin{enumerate}
\item $K^\bullet$ est $m$-pseudo-cohérent comme complexe de $B$-modules ;
\item $K^\bullet$ est $m$-pseudo-cohérent comme complexe de $A$-modules.
\end{enumerate}
La même équivalence vaut pour la pseudo-cohérence.
\end{lemma}

\begin{proof}
Supposons (1). Choisissons un complexe borné de $B$-modules libres de rang
fini $E^\bullet$ et un morphisme $\alpha : E^\bullet \to K^\bullet$ qui
induit un isomorphisme en cohomologie aux degrés $> m$ et une surjection au
degré $m$. Considérons le triangle distingué
$(E^\bullet, K^\bullet, C(\alpha)^\bullet)$. D'après le Lemme
\ref{lemma-recognize-pseudo-coherent}, $C(\alpha)^\bullet$ est
$m$-pseudo-cohérent comme complexe de $A$-modules. Il suffit donc de montrer
que $E^\bullet$ est pseudo-cohérent comme complexe de $A$-modules, ce qui
résulte du Lemme \ref{lemma-complex-pseudo-coherent-modules}. Le cas
pseudo-cohérent de (1) $\Rightarrow$ (2) résulte de ceci et du Lemme
\ref{lemma-pseudo-coherent}.

\medskip\noindent
Supposons (2). Soit $n$ le plus grand entier tel que
$H^n(K^\bullet) \not = 0$. Nous allons montrer que $K^\bullet$ est
$m$-pseudo-cohérent comme complexe de $B$-modules par récurrence sur $n - m$. Le cas
$n < m$ résulte du Lemme \ref{lemma-recognize-pseudo-coherent}. Choisissons
un complexe borné de $A$-modules libres de rang fini $E^\bullet$ et un
morphisme $\alpha : E^\bullet \to K^\bullet$ qui induit un isomorphisme en
cohomologie aux degrés $> m$ et une surjection au degré $m$. Considérons le
morphisme induit de complexes
$$
\alpha \otimes 1 : E^\bullet \otimes_A B \to K^\bullet.
$$
Remarquons que $C(\alpha \otimes 1)^\bullet$ est acyclique aux degrés
$\geq n$, car $H^n(E) \to H^n(E^\bullet \otimes_A B) \to H^n(K^\bullet)$
est surjectif par construction, et parce que
$H^i(E^\bullet \otimes_A B) = 0$ pour $i > n$ d'après la suite spectrale de
l'Exemple \ref{example-tor}. D'autre part,
$C(\alpha \otimes 1)^\bullet$ est $m$-pseudo-cohérent comme complexe de
$A$-modules, car $K^\bullet$ et $E^\bullet \otimes_A B$ le sont (voir le
Lemme \ref{lemma-complex-pseudo-coherent-modules}), d'après le Lemme
\ref{lemma-cone-pseudo-coherent}. Par récurrence,
$C(\alpha \otimes 1)^\bullet$ est donc $m$-pseudo-cohérent comme complexe de
$B$-modules. Enfin, une nouvelle application du Lemme
\ref{lemma-cone-pseudo-coherent} montre que $K^\bullet$ est
$m$-pseudo-cohérent comme complexe de $B$-modules (car
$E^\bullet \otimes_A B$ est manifestement pseudo-cohérent comme complexe de
$B$-modules). Le cas pseudo-cohérent de (2) $\Rightarrow$ (1) résulte de
ceci et du Lemme \ref{lemma-pseudo-coherent}.
\end{proof}

\begin{lemma}
\label{lemma-pull-pseudo-coherent}
Soit $A \to B$ un homomorphisme d'anneaux. Soit $K^\bullet$ un complexe
$m$-pseudo-cohérent (resp.\ pseudo-cohérent) de $A$-modules. Alors
$K^\bullet \otimes_A^{\mathbf{L}} B$ est un complexe $m$-pseudo-cohérent
(resp.\ pseudo-cohérent) de $B$-modules.
\end{lemma}

\begin{proof}
Remarquons d'abord que l'énoncé du lemme a un sens, car $K^\bullet$ est
borné supérieurement et $K^\bullet \otimes_A^{\mathbf{L}} B$ est donc
défini par l'Équation (\ref{equation-derived-tensor-algebra}). Choisissons
ensuite un complexe borné $E^\bullet$ de $A$-modules libres de rang fini et
$\alpha : E^\bullet \to K^\bullet$ tel que $H^i(\alpha)$ soit un
isomorphisme pour $i > m$ et soit surjectif pour $i = m$. Le cône
$C(\alpha)^\bullet$ est alors acyclique aux degrés $\geq m$. Comme
$-\otimes_A^{\mathbf{L}} B$ est un foncteur exact, on obtient un triangle
distingué
$$
(E^\bullet \otimes_A^{\mathbf{L}} B, K^\bullet \otimes_A^{\mathbf{L}} B,
C(\alpha)^\bullet \otimes_A^{\mathbf{L}} B)
$$
de complexes de $B$-modules. Par le dual de Catégories dérivées, Lemme
\ref{derived-lemma-negative-vanishing}, on voit que
$H^i(C(\alpha)^\bullet \otimes_A^{\mathbf{L}} B) = 0$ pour $i \geq m$.
Comme $E^\bullet$ est un complexe de $A$-modules projectifs, on a
$E^\bullet \otimes_A^{\mathbf{L}} B = E^\bullet \otimes_A B$, et donc
$$
E^\bullet \otimes_A B
\longrightarrow
K^\bullet \otimes_A^{\mathbf{L}} B
$$
est un morphisme de complexes de $B$-modules qui montre que
$K^\bullet \otimes_A^{\mathbf{L}} B$ est $m$-pseudo-cohérent. Le cas des
complexes pseudo-cohérents se déduit de celui des complexes
$m$-pseudo-cohérents à l'aide du Lemme \ref{lemma-pseudo-coherent}.
\end{proof}

\begin{lemma}
\label{lemma-flat-base-change-pseudo-coherent}
Soit $A \to B$ un homomorphisme plat d'anneaux. Soit $M$ un module
$m$-pseudo-cohérent (resp.\ pseudo-cohérent) sur $A$. Alors $M \otimes_A B$
est un module $m$-pseudo-cohérent (resp.\ pseudo-cohérent) sur $B$.
\end{lemma}

\begin{proof}
Cela résulte immédiatement du Lemme \ref{lemma-pull-pseudo-coherent} et du
fait que $M \otimes_A^{\mathbf{L}} B = M \otimes_A B$, puisque $B$ est plat
sur $A$.
\end{proof}

\noindent
Le lemme suivant résulte aussi du Lemme plus fort
\ref{lemma-flat-descent-pseudo-coherent}.

\begin{lemma}
\label{lemma-glue-pseudo-coherent}
Soit $R$ un anneau. Soient $f_1, \ldots, f_r \in R$ des éléments qui
engendrent l'idéal unité, soit $m \in \mathbf{Z}$ et soit $K^\bullet$ un
complexe de $R$-modules. Si, pour tout $i$, le complexe
$K^\bullet \otimes_R R_{f_i}$ est $m$-pseudo-cohérent (resp.\
pseudo-cohérent), alors $K^\bullet$ est $m$-pseudo-cohérent (resp.\
pseudo-cohérent).
\end{lemma}

\begin{proof}
Nous utiliserons sans le rappeler que $- \otimes_R R_{f_i}$ est un foncteur
exact et que, par conséquent,
$$
H^i(K^\bullet)_{f_i} =
H^i(K^\bullet) \otimes_R R_{f_i} = H^i(K^\bullet \otimes_R R_{f_i}).
$$
Supposons que $K^\bullet \otimes_R R_{f_i}$ soit $m$-pseudo-cohérent pour
$i = 1, \ldots, r$. Soit $n \in \mathbf{Z}$ le plus grand entier tel que
$H^n(K^\bullet \otimes_R R_{f_i})$ soit non nul pour un certain $i$. Cela
implique en particulier que $H^i(K^\bullet) = 0$ pour $i > n$ (et que
$H^n(K^\bullet) \not = 0$) ; voir Algèbre, Lemme
\ref{algebra-lemma-cover}. Démontrons le lemme par récurrence sur $n - m$.
Si $n < m$, il résulte du Lemme \ref{lemma-recognize-pseudo-coherent}. Si
$n \geq m$, alors $H^n(K^\bullet)_{f_i}$ est un $R_{f_i}$-module de type
fini pour tout $i$ ; voir le Lemme \ref{lemma-finite-cohomology}. Ainsi,
$H^n(K^\bullet)$ est un $R$-module de type fini ; voir Algèbre, Lemme
\ref{algebra-lemma-cover}. Choisissons un $R$-module libre de rang fini $E$
et une surjection $E \to H^n(K^\bullet)$. Comme $E$ est projectif, on peut
la relever en un morphisme de complexes
$\alpha : E[-n] \to K^\bullet$. Le cône $C(\alpha)^\bullet$ a alors une
cohomologie nulle aux degrés $\geq n$. D'autre part, les complexes
$C(\alpha)^\bullet \otimes_R R_{f_i}$ sont $m$-pseudo-cohérents pour tout
$i$ ; voir le Lemme \ref{lemma-cone-pseudo-coherent}. Par récurrence,
$C(\alpha)^\bullet$ est donc $m$-pseudo-cohérent comme complexe de
$R$-modules. Une nouvelle application du Lemme
\ref{lemma-cone-pseudo-coherent} permet de conclure.
\end{proof}

\begin{lemma}
\label{lemma-flat-descent-pseudo-coherent}
Soit $R$ un anneau, soit $m \in \mathbf{Z}$ et soit $K^\bullet$ un complexe
de $R$-modules. Soit $R \to R'$ un homomorphisme fidèlement plat
d'anneaux. Si le complexe $K^\bullet \otimes_R R'$ est $m$-pseudo-cohérent
(resp.\ pseudo-cohérent), alors $K^\bullet$ est $m$-pseudo-cohérent (resp.\
pseudo-cohérent).
\end{lemma}

\begin{proof}
Nous utiliserons sans le rappeler que $- \otimes_R R'$ est un foncteur
exact et que, par conséquent,
$$
H^i(K^\bullet) \otimes_R R' = H^i(K^\bullet \otimes_R R').
$$
Supposons que $K^\bullet \otimes_R R'$ soit $m$-pseudo-cohérent. Soit
$n \in \mathbf{Z}$ le plus grand entier tel que $H^n(K^\bullet)$ soit non
nul ; alors $n$ est aussi le plus grand entier tel que
$H^n(K^\bullet \otimes_R R')$ soit non nul. Démontrons le lemme par
récurrence sur $n - m$. Si $n < m$, le résultat découle du Lemme
\ref{lemma-recognize-pseudo-coherent}. Si $n \geq m$, alors
$H^n(K^\bullet) \otimes_R R'$ est un $R'$-module de type fini ; voir le
Lemme \ref{lemma-finite-cohomology}. Par conséquent, $H^n(K^\bullet)$ est
un $R$-module de type fini ; voir Algèbre, Lemme
\ref{algebra-lemma-descend-properties-modules}. Choisissons un $R$-module
libre de rang fini $E$ et une surjection $E \to H^n(K^\bullet)$. Comme $E$
est projectif, on peut la relever en un morphisme de complexes
$\alpha : E[-n] \to K^\bullet$. Le cône $C(\alpha)^\bullet$ a alors une
cohomologie nulle aux degrés $\geq n$. D'autre part, le complexe
$C(\alpha)^\bullet \otimes_R R'$ est $m$-pseudo-cohérent ; voir le Lemme
\ref{lemma-cone-pseudo-coherent}. Par récurrence, $C(\alpha)^\bullet$ est
donc $m$-pseudo-cohérent comme complexe de $R$-modules. Une nouvelle
application du Lemme \ref{lemma-cone-pseudo-coherent} permet de conclure.
\end{proof}

\begin{lemma}
\label{lemma-tensor-pseudo-coherent}
Soit $R$ un anneau et soient $K, L$ des objets de $D(R)$.
\begin{enumerate}
\item Si $K$ est $n$-pseudo-cohérent et $H^i(K) = 0$ pour $i > a$, et si
$L$ est $m$-pseudo-cohérent et $H^j(L) = 0$ pour $j > b$, alors
$K \otimes_R^\mathbf{L} L$ est $t$-pseudo-cohérent, où
$t = \max(m + a, n + b)$.
\item Si $K$ et $L$ sont pseudo-cohérents, alors
$K \otimes_R^\mathbf{L} L$ est pseudo-cohérent.
\end{enumerate}
\end{lemma}

\begin{proof}
Démontrons (1). On peut supposer qu'il existe des complexes bornés
$K^\bullet$ et $L^\bullet$ de $R$-modules libres de rang fini, ainsi que des
morphismes $\alpha : K^\bullet \to K$ et $\beta : L^\bullet \to L$, tels
que $H^i(\alpha)$ soit un isomorphisme pour $i > n$ et soit surjectif pour
$i = n$, et que $H^i(\beta)$ soit un isomorphisme pour $i > m$ et soit
surjectif pour $i = m$. Alors le morphisme
$$
\alpha \otimes^\mathbf{L} \beta :
\text{Tot}(K^\bullet \otimes_R L^\bullet)
\to K \otimes_R^\mathbf{L} L
$$
induit des isomorphismes en cohomologie au degré $i$ pour $i > t$ et une
surjection pour $i = t$. Cela résulte de la suite spectrale de Tor (détails
omis). L'assertion (2) découle de (1) et du Lemme
\ref{lemma-pseudo-coherent}.
\end{proof}

\begin{lemma}
\label{lemma-Noetherian-pseudo-coherent}
Soit $R$ un anneau noethérien. Alors
\begin{enumerate}
\item un complexe de $R$-modules $K^\bullet$ est $m$-pseudo-cohérent si et
seulement si $K^\bullet \in D^{-}(R)$ et si $H^i(K^\bullet)$ est un
$R$-module de type fini pour $i \geq m$ ;
\item un complexe de $R$-modules $K^\bullet$ est pseudo-cohérent si et
seulement si $K^\bullet \in D^{-}(R)$ et si $H^i(K^\bullet)$ est un
$R$-module de type fini pour tout $i$ ;
\item un $R$-module est pseudo-cohérent si et seulement s'il est de type
fini.
\end{enumerate}
\end{lemma}

\begin{proof}
Dans Algèbre, Lemme \ref{algebra-lemma-resolution-by-finite-free}, nous
avons vu que tout $R$-module de type fini est pseudo-cohérent. D'autre part,
un module pseudo-cohérent est de type fini ; voir le Lemme
\ref{lemma-n-pseudo-module}. L'assertion (3) en résulte. Supposons que
$K^\bullet$ soit un complexe $m$-pseudo-cohérent. Il existe alors un
complexe borné de $R$-modules libres de rang fini $E^\bullet$ tel que
$H^i(K^\bullet)$ soit isomorphe à $H^i(E^\bullet)$ pour $i > m$ et que
$H^m(K^\bullet)$ soit un quotient de $H^m(E^\bullet)$. Il est donc clair
que chaque $H^i(K^\bullet)$, pour $i \geq m$, est de type fini. La
réciproque de (1) résulte du Lemme
\ref{lemma-cohomology-pseudo-coherent} et de (3). L'assertion (2) découle de
(1) et du Lemme \ref{lemma-pseudo-coherent}.
\end{proof}

\begin{lemma}
\label{lemma-coherent-pseudo-coherent}
Soit $R$ un anneau cohérent (Algèbre, Définition
\ref{algebra-definition-coherent}). Soit $K \in D^-(R)$. Les assertions
suivantes sont équivalentes :
\begin{enumerate}
\item $K$ est $m$-pseudo-cohérent ;
\item $H^m(K)$ est un $R$-module de type fini et $H^i(K)$ est cohérent pour
$i > m$ ;
\item $H^m(K)$ est un $R$-module de type fini et $H^i(K)$ est de
présentation finie pour $i > m$.
\end{enumerate}
Ainsi, $K$ est pseudo-cohérent si et seulement si $H^i(K)$ est un module
cohérent pour tout $i$.
\end{lemma}

\begin{proof}
Rappelons qu'un $R$-module $M$ est cohérent si et seulement s'il est de
présentation finie (Algèbre, Lemme \ref{algebra-lemma-coherent-ring}). Cela
explique l'équivalence de (2) et (3). Dans ce cas, si l'on choisit une suite
exacte $0 \to N \to R^{\oplus m} \to M \to 0$, alors $N$ est cohérent
d'après Algèbre, Lemme \ref{algebra-lemma-coherent}. En répétant ce procédé
avec $N$, on obtient une résolution
$$
\ldots \to R^{\oplus n} \to R^{\oplus m} \to M \to 0
$$
par des $R$-modules libres de rang fini. Autrement dit, $M$ est
pseudo-cohérent. L'équivalence de (1) et (2) résulte de ceci et des Lemmes
\ref{lemma-cohomology-pseudo-coherent} et \ref{lemma-n-pseudo-module}. La
dernière assertion résulte de l'équivalence de (1) et (2), combinée avec le
Lemme \ref{lemma-pseudo-coherent}.
\end{proof}









\section{Modules pseudo-cohérents, II}
\label{section-pseudo-coherent-bis}

\noindent
Nous poursuivons la discussion commencée dans la Section
\ref{section-pseudo-coherent}.

\begin{lemma}
\label{lemma-pseudo-coherence-colimit-ext}
Soit $R$ un anneau. Soit $M = \colim M_i$ une colimite filtrante de
$R$-modules. Soit $K \in D(R)$ un objet $m$-pseudo-cohérent. Alors
$\colim \Ext^n_R(K, M_i) = \Ext^n_R(K, M)$ pour $n < -m$, et
$\colim \Ext^{-m}_R(K, M_i) \to \Ext^{-m}_R(K, M)$ est injectif.
\end{lemma}

\begin{proof}
Par définition, on peut trouver un triangle distingué
$$
E \to K \to L \to E[1]
$$
dans $D(R)$, où $E$ est représenté par un complexe borné de $R$-modules
libres de rang fini et où $H^i(L) = 0$ pour $i \geq m$. Alors
$\Ext^n_R(L, N) = 0$ pour tout $R$-module $N$ et tout $n \leq -m$ ; voir
Catégories dérivées, Lemme \ref{derived-lemma-negative-exts}. La suite
exacte longue de $\Ext$ associée au triangle distingué montre que
$\Ext^n_R(K, N) \to \Ext^n_R(E, N)$ est un isomorphisme pour $n < -m$ et
est injectif pour $n = -m$. Il suffit donc de montrer que
$M \mapsto \Ext_R^n(E, M)$ commute aux colimites filtrantes lorsque $E$ est
représenté par un complexe borné de $R$-modules libres de rang fini
$E^\bullet$. Les modules $\Ext^n_R(E, M)$ sont calculés par le complexe
$\Hom_R(E^\bullet, M)$ ; voir Catégories dérivées, Lemme
\ref{derived-lemma-morphisms-from-projective-complex}. Le foncteur
$M \mapsto \Hom_R(E^p, M)$ commute aux colimites filtrantes puisque $E^p$
est libre de rang fini. Ainsi,
$\Hom_R(E^\bullet, M) = \colim \Hom_R(E^\bullet, M_i)$ comme complexes.
Comme les colimites filtrantes sont exactes (Algèbre, Lemme
\ref{algebra-lemma-directed-colimit-exact}), on conclut.
\end{proof}

\begin{lemma}
\label{lemma-characterize-pseudo-coherent-colimit-ext}
Soit $R$ un anneau, soit $K \in D^-(R)$ et soit $m \in \mathbf{Z}$. Alors
$K$ est $m$-pseudo-cohérent si et seulement si, pour toute colimite
filtrante $M = \colim M_i$ de $R$-modules, on a
$\colim \Ext^n_R(K, M_i) = \Ext^n_R(K, M)$ pour $n < -m$ et si
$\colim \Ext^{-m}_R(K, M_i) \to \Ext^{-m}_R(K, M)$ est injectif.
\end{lemma}

\begin{proof}
Une implication a été démontrée dans le Lemme
\ref{lemma-pseudo-coherence-colimit-ext}. Supposons que, pour toute colimite
filtrante $M = \colim M_i$ de $R$-modules, on ait
$\colim \Ext^n_R(K, M_i) = \Ext^n_R(K, M)$ pour $n < -m$ et que
$\colim \Ext^{-m}_R(K, M_i) \to \Ext^{-m}_R(K, M)$ soit injectif. Montrons
que $K$ est $m$-pseudo-cohérent.

\medskip\noindent
Soit $t$ le plus grand entier tel que $H^t(K)$ soit non nul. Raisonnons par
récurrence sur $t$. Si $t < m$, alors $K$ est $m$-pseudo-cohérent d'après le
Lemme \ref{lemma-recognize-pseudo-coherent}. Si $t \geq m$, l'égalité
$\Hom_R(H^t(K), M) = \Ext^{-t}_R(K, M)$ montre que
$\colim \Hom_R(H^t(K), M_i) \to \Hom_R(H^t(K), M)$ est injectif pour toute
colimite filtrante $M = \colim M_i$. Il en résulte que $H^t(K)$ est un
$R$-module de type fini, d'après Algèbre, Lemme
\ref{algebra-lemma-characterize-finite-module-hom}. Choisissons un
$R$-module libre de rang fini $F$ et une surjection $F \to H^t(K)$. On peut
la relever en un morphisme $F[-t] \to K$ dans $D(R)$ et choisir un triangle
distingué
$$
F[-t] \to K \to L \to F[-t + 1]
$$
dans $D(R)$. Alors $H^i(L) = 0$ pour $i \geq t$. De plus, la suite exacte
longue de $\Ext$ associée à ce triangle distingué montre, par un petit
argument que nous omettons, que $L$ hérite de l'hypothèse faite sur $K$.
Par récurrence sur $t$, $L$ est $m$-pseudo-cohérent. Ainsi, $K$ est
$m$-pseudo-cohérent d'après le Lemme \ref{lemma-cone-pseudo-coherent}.
\end{proof}

\begin{lemma}
\label{lemma-pseudo-coherence-and-ext}
Soit $R$ un anneau et soient $L$, $M$, $N$ des $R$-modules.
\begin{enumerate}
\item Si $M$ est de présentation finie et si $L$ est plat, alors le
morphisme canonique
$\Hom_R(M, N) \otimes_R L \to \Hom_R(M, N \otimes_R L)$
est un isomorphisme.
\item Si $M$ est $(-m)$-pseudo-cohérent et si $L$ est plat, alors le
morphisme canonique
$\Ext^i_R(M, N) \otimes_R L \to \Ext^i_R(M, N \otimes_R L)$
est un isomorphisme pour $i < m$.
\end{enumerate}
\end{lemma}

\begin{proof}
Choisissons une résolution $F_\bullet \to M$ dont les termes sont des
$R$-modules libres ; voir Algèbre, Lemme
\ref{algebra-lemma-resolution-by-finite-free}. Le complexe
$\Hom_R(F_\bullet, N)$ calcule $\Ext^i_R(M, N)$ et le complexe
$\Hom_R(F_\bullet, N \otimes_R L)$ calcule
$\Ext^i_R(M, N \otimes_R L)$. Il existe toujours un morphisme de complexes
de cochaînes
$$
\Hom_R(F_\bullet, N) \otimes_R L
\longrightarrow
\Hom_R(F_\bullet, N \otimes_R L)
$$
qui induit des morphismes canoniques
$\Ext^i_R(M, N) \otimes_R L \to \Ext^i_R(M, N \otimes_R L)$
pour tout $i \geq 0$ (canoniques, par exemple, au sens où ils ne dépendent
pas du choix de la résolution $F_\bullet$). Si $L$ est plat, alors le
complexe $\Hom_R(F_\bullet, N) \otimes_R L$ calcule
$\Ext^i_R(M, N) \otimes_R L$, puisque la prise de cohomologie commute au
produit tensoriel par $L$.

\medskip\noindent
Si $M$ est $(-m)$-pseudo-cohérent, on peut choisir $F_\bullet$ de sorte que
$F_i$ soit libre de rang fini pour $i = 0, \ldots, m$. Le morphisme affiché
de complexes de cochaînes est alors un isomorphisme aux degrés $\leq m$, et
induit donc un isomorphisme en cohomologie aux degrés $< m$. Cela démontre
(2). Si $M$ est de présentation finie, alors $M$ est
$(-1)$-pseudo-cohérent d'après le Lemme \ref{lemma-n-pseudo-module}, et le
résultat suit de $\Hom = \Ext^0$.
\end{proof}

\begin{lemma}
\label{lemma-pseudo-coherence-and-base-change-ext}
Soit $R \to R'$ un homomorphisme plat d'anneaux et soient $M$, $N$ des
$R$-modules.
\begin{enumerate}
\item Si $M$ est un $R$-module de présentation finie, alors
$\Hom_R(M, N) \otimes_R R' = \Hom_{R'}(M \otimes_R R', N \otimes_R R')$.
\item Si $M$ est $(-m)$-pseudo-cohérent, alors
$\Ext^i_R(M, N) \otimes_R R' = \Ext^i_{R'}(M \otimes_R R', N \otimes_R R')$
pour $i < m$.
\end{enumerate}
En particulier, si $R$ est noethérien et si $M$ est de type fini, ceci vaut
pour tout $i$.
\end{lemma}

\begin{proof}
D'après Algèbre, Lemme \ref{algebra-lemma-flat-base-change-ext}, on a
$\Ext^i_{R'}(M \otimes_R R', N \otimes_R R') =
\Ext^i_R(M, N \otimes_R R')$.
En combinant ceci avec le Lemme \ref{lemma-pseudo-coherence-and-ext}, on
obtient (1) et (2). La dernière assertion en découle avec le Lemme
\ref{lemma-Noetherian-pseudo-coherent}.
\end{proof}

\begin{lemma}
\label{lemma-pseudo-coherent-tensor}
Soit $R$ un anneau et soit $K \in D^-(R)$. Les assertions suivantes sont
équivalentes :
\begin{enumerate}
\item $K$ est pseudo-cohérent ;
\item pour toute famille $(Q_{\alpha})_{\alpha \in A}$ de $R$-modules, le
morphisme canonique
$$
\alpha :
K \otimes_R^\mathbf{L} \left( \prod\nolimits_\alpha Q_{\alpha} \right)
\longrightarrow
\prod\nolimits_\alpha (K \otimes_R^\mathbf{L} Q_{\alpha})
$$
est un isomorphisme dans $D(R)$ ;
\item pour tout $R$-module $Q$ et tout ensemble $A$, le morphisme canonique
$$
\beta : K \otimes_R^\mathbf{L} Q^A \longrightarrow (K \otimes_R^\mathbf{L} Q)^A
$$
est un isomorphisme dans $D(R)$ ;
\item pour tout ensemble $A$, le morphisme canonique
$$
\gamma : K \otimes_R^\mathbf{L} R^A \longrightarrow K^A
$$
est un isomorphisme dans $D(R)$.
\end{enumerate}
Étant donné $m \in \mathbf{Z}$, les assertions suivantes sont équivalentes :
\begin{enumerate}
\item[(a)] $K$ est $m$-pseudo-cohérent ;
\item[(b)] pour toute famille $(Q_{\alpha})_{\alpha \in A}$ de
$R$-modules, avec $\alpha$ comme ci-dessus, $H^i(\alpha)$ est un
isomorphisme pour $i > m$ et est surjectif pour $i = m$ ;
\item[(c)] pour tout $R$-module $Q$ et tout ensemble $A$, avec $\beta$
comme ci-dessus, $H^i(\beta)$ est un isomorphisme pour $i > m$ et est
surjectif pour $i = m$ ;
\item[(d)] pour tout ensemble $A$, avec $\gamma$ comme ci-dessus,
$H^i(\gamma)$ est un isomorphisme pour $i > m$ et est surjectif pour
$i = m$.
\end{enumerate}
\end{lemma}

\begin{proof}
Si $K$ est pseudo-cohérent, alors $K$ peut être représenté par un complexe borné
supérieurement de $R$-modules libres de rang fini. Les produits tensoriels
dérivés se calculent alors en tensorisant par ce complexe. De plus, les
produits dans $D(R)$ s'obtiennent en prenant les produits de représentants
quelconques. Ainsi, (1) implique (2), (3), (4) par le résultat correspondant
pour les modules ; voir Algèbre, Proposition
\ref{algebra-proposition-fp-tensor}.

\medskip\noindent
De la même manière, en utilisant l'exactitude à droite du produit tensoriel,
le lecteur vérifie que (a) implique (b), (c) et (d).

\medskip\noindent
Supposons (4). Pour montrer que $K$ est pseudo-cohérent, il suffit de
montrer que $K$ est $m$-pseudo-cohérent pour tout $m$ (Lemme
\ref{lemma-pseudo-coherent}). Pour achever la démonstration, il suffit donc
de prouver que (d) implique (a).

\medskip\noindent
Supposons (d). Soit $i$ le plus grand entier tel que $H^i(K)$ soit non nul.
Si $i < m$, c'est terminé. Sinon, (d) et la description ci-dessus des
produits dans $D(R)$ montrent que
$H^i(K) \otimes_R R^A \to H^i(K)^A$ est surjectif. Par conséquent,
$H^i(K)$ est un $R$-module de type fini d'après Algèbre, Proposition
\ref{algebra-proposition-fg-tensor}. On peut donc choisir un complexe $L$
formé d'un seul module libre de rang fini placé au degré $i$, et un
morphisme de complexes $L \to K$ tel que $H^i(L) \to H^i(K)$ soit
surjectif. En particulier, $L$ vérifie (1), (2), (3) et (4). Choisissons un
triangle distingué
$$
L \to K \to M \to L[1]
$$
On voit alors que $H^j(M) = 0$ pour $j \geq i$. D'autre part, $M$ possède
encore la propriété (d), par un petit argument que nous omettons. Par
récurrence sur $i$, $M$ est $m$-pseudo-cohérent. Ainsi, $K$ est
$m$-pseudo-cohérent d'après le Lemme \ref{lemma-cone-pseudo-coherent}.
\end{proof}

\begin{lemma}
\label{lemma-detect-cohomology-pseudo-coherent}
Soit $R$ un anneau. Soit $K \in D(R)$ pseudo-cohérent et soit
$i \in \mathbf{Z}$. Il existe un $R$-module $M$ de présentation finie et un
morphisme $K \to M[-i]$ dans $D(R)$ qui induit une injection
$H^i(K) \to M$.
\end{lemma}

\begin{proof}
D'après la Définition \ref{definition-pseudo-coherent}, on peut représenter
$K$ par un complexe $P^\bullet$ de $R$-modules libres de rang fini. Posons
$M = \Coker(P^{i - 1} \to P^i)$.
\end{proof}

\begin{lemma}
\label{lemma-detect-cohomology}
Soit $A$ un anneau noethérien. Soit $K \in D(A)$ pseudo-cohérent,
c'est-à-dire $K \in D^-(A)$ à modules de cohomologie de type fini. Soit
$\mathfrak m$ un idéal maximal de $A$. Si
$H^i(K)/\mathfrak m H^i(K) \not = 0$, alors il existe un $A$-module $E$ de
type fini, annulé par une puissance de $\mathfrak m$, et un morphisme
$K \to E[-i]$ non nul sur $H^i(K)$.
\end{lemma}

\begin{proof}
La formulation équivalente de la pseudo-cohérence utilisée dans l'énoncé
est celle du Lemme \ref{lemma-Noetherian-pseudo-coherent}. Choisissons
$K \to M[-i]$ comme dans le Lemme
\ref{lemma-detect-cohomology-pseudo-coherent}. D'après Artin--Rees (Algèbre,
Lemme \ref{algebra-lemma-Artin-Rees}), il existe un $n$ tel que
$H^i(K) \cap \mathfrak m^n M \subset \mathfrak m H^i(K)$.
Prenons $E = M/\mathfrak m^n M$.
\end{proof}
\section{Dimension de Tor}
\label{section-tor}

\noindent
Au lieu d'utiliser des résolutions par des modules projectifs, nous pouvons
considérer des résolutions par des modules plats. Cela conduit à la notion
suivante.

\begin{definition}
\label{definition-tor-amplitude}
Soit $R$ un anneau. Notons $D(R)$ sa catégorie dérivée.
Soient $a, b \in \mathbf{Z}$.
\begin{enumerate}
\item Un objet $K^\bullet$ de $D(R)$ est d'
{\it amplitude de Tor dans $[a, b]$}
si $H^i(K^\bullet \otimes_R^\mathbf{L} M) = 0$ pour tout $R$-module
$M$ et tout $i \not \in [a, b]$.
\item Un objet $K^\bullet$ de $D(R)$ est de {\it dimension de Tor finie}
s'il est d'amplitude de Tor dans $[a, b]$ pour certains $a, b$.
\item Un $R$-module $M$ est de {\it dimension de Tor $\leq d$}
si $M[0]$, considéré comme objet de $D(R)$, est d'amplitude de Tor dans
$[-d, 0]$.
\item Un $R$-module $M$ est de {\it dimension de Tor finie}
si $M[0]$, considéré comme objet de $D(R)$, est de dimension de Tor finie.
\end{enumerate}
\end{definition}

\noindent
Observons que si $K^\bullet$ est de dimension de Tor finie, alors
$K^\bullet \in D^b(R)$.

\begin{lemma}
\label{lemma-last-one-flat}
Soit $R$ un anneau. Soit $K^\bullet$ un complexe borné supérieurement de
$R$-modules plats, d'amplitude de Tor dans $[a, b]$.
Alors $\Coker(d_K^{a - 1})$ est un $R$-module plat.
\end{lemma}

\begin{proof}
Puisque $K^\bullet$ est un complexe borné supérieurement de modules plats,
nous avons
$K^\bullet \otimes_R M = K^\bullet \otimes_R^{\mathbf{L}} M$.
Par conséquent, pour tout $R$-module $M$, la suite
$$
K^{a - 2} \otimes_R M \to K^{a - 1} \otimes_R M \to K^a \otimes_R M
$$
est exacte au milieu. Comme
$K^{a - 2} \to K^{a - 1} \to K^a \to \Coker(d_K^{a - 1}) \to 0$
est une résolution plate, cela implique que
$\text{Tor}_1^R(\Coker(d_K^{a - 1}), M) = 0$
pour tout $R$-module $M$. Cela signifie que
$\Coker(d_K^{a - 1})$ est plat ; voir Algèbre,
lemme \ref{algebra-lemma-characterize-flat}.
\end{proof}

\begin{lemma}
\label{lemma-tor-amplitude}
Soit $R$ un anneau. Soit $K^\bullet$ un objet de $D(R)$.
Soient $a, b \in \mathbf{Z}$. Les assertions suivantes sont équivalentes :
\begin{enumerate}
\item $K^\bullet$ est d'amplitude de Tor dans $[a, b]$.
\item $K^\bullet$ est quasi-isomorphe à un complexe
$E^\bullet$ de $R$-modules plats tel que $E^i = 0$ pour
$i \not \in [a, b]$.
\end{enumerate}
\end{lemma}

\begin{proof}
Si (2) est vraie, nous pouvons calculer
$K^\bullet \otimes_R^\mathbf{L} M = E^\bullet \otimes_R M$
et il est clair que (1) est vraie.
Supposons (1) vraie. Nous pouvons remplacer $K^\bullet$ par
une résolution projective telle que $K^i = 0$ pour $i > b$.
Voir Catégories dérivées, lemme
\ref{derived-lemma-projective-resolutions-exist}.
Posons $E^\bullet = \tau_{\geq a}K^\bullet$. Tout est clair, sauf
la platitude de $E^a$, qui résulte immédiatement du
lemme \ref{lemma-last-one-flat}
et des définitions.
\end{proof}

\begin{lemma}
\label{lemma-bounded-below-tor-amplitude}
Soit $R$ un anneau. Soit $a \in \mathbf{Z}$ et soit $K$ un objet de $D(R)$.
Les assertions suivantes sont équivalentes :
\begin{enumerate}
\item $K$ est d'amplitude de Tor dans $[a, \infty]$, et
\item $K$ est quasi-isomorphe à un complexe K-plat
$E^\bullet$ dont les termes sont des $R$-modules plats et tel que
$E^i = 0$ pour $i \not \in [a, \infty]$.
\end{enumerate}
\end{lemma}

\begin{proof}
L'implication (2) $\Rightarrow$ (1) est immédiate. Supposons (1) vraie.
Choisissons d'abord un complexe K-plat $K^\bullet$ à termes plats
représentant $K$ ; voir le lemme \ref{lemma-K-flat-resolution}.
Pour tout $R$-module $M$, la cohomologie de
$$
K^{n - 1} \otimes_R M \to K^n \otimes_R M \to K^{n + 1} \otimes_R M
$$
calcule $H^n(K \otimes_R^\mathbf{L} M)$. Celui-ci est toujours nul
pour $n < a$. En appliquant le lemme \ref{lemma-last-one-flat}
au complexe $\ldots \to K^{a - 1} \to K^a \to K^{a + 1}$,
nous concluons que $N = \Coker(K^{a - 1} \to K^a)$ est un $R$-module plat.
Posons
$$
E^\bullet = \tau_{\geq a}K^\bullet =
(\ldots \to 0 \to N \to K^{a + 1} \to \ldots )
$$
Le noyau $L^\bullet$ de $K^\bullet \to E^\bullet$ est le complexe
$$
L^\bullet = (\ldots \to K^{a - 1} \to I \to 0 \to \ldots)
$$
où $I \subset K^a$ est l'image de $K^{a - 1} \to K^a$.
Puisque nous avons la suite exacte courte $0 \to I \to K^a \to N \to 0$,
nous voyons que $I$ est un $R$-module plat. Ainsi, $L^\bullet$ est un
complexe borné supérieurement de modules plats, donc K-plat d'après le
lemme \ref{lemma-derived-tor-quasi-isomorphism}.
Il en résulte que $E^\bullet$ est K-plat d'après le
lemme \ref{lemma-K-flat-two-out-of-three-ses}.
\end{proof}

\begin{lemma}
\label{lemma-cone-tor-amplitude}
Soit $R$ un anneau.
Soit $(K^\bullet, L^\bullet, M^\bullet, f, g, h)$ un triangle
distingué de $D(R)$. Soient $a, b \in \mathbf{Z}$.
\begin{enumerate}
\item Si $K^\bullet$ est d'amplitude de Tor dans $[a + 1, b + 1]$ et
$L^\bullet$ d'amplitude de Tor dans $[a, b]$, alors $M^\bullet$ est
d'amplitude de Tor dans $[a, b]$.
\item Si $K^\bullet, M^\bullet$ sont d'amplitude de Tor dans $[a, b]$,
alors $L^\bullet$ est d'amplitude de Tor dans $[a, b]$.
\item Si $L^\bullet$ est d'amplitude de Tor dans $[a + 1, b + 1]$
et $M^\bullet$ d'amplitude de Tor dans $[a, b]$, alors
$K^\bullet$ est d'amplitude de Tor dans $[a + 1, b + 1]$.
\end{enumerate}
\end{lemma}

\begin{proof}
Omis. Indication : cela résulte simplement de la suite exacte longue de
cohomologie associée à un triangle distingué et du fait que
$- \otimes_R^{\mathbf{L}} M$ préserve les triangles distingués.
L'assertion (2) est la plus facile à démontrer, et les autres s'en déduisent
par translation.
\end{proof}

\begin{lemma}
\label{lemma-tor-dimension}
Soit $R$ un anneau. Soit $M$ un $R$-module.
Soit $d \geq 0$. Les assertions suivantes sont équivalentes :
\begin{enumerate}
\item $M$ est de dimension de Tor $\leq d$, et
\item il existe une résolution
$$
0 \to F_d \to \ldots \to F_1 \to F_0 \to M \to 0
$$
où chaque $F_i$ est un $R$-module plat.
\end{enumerate}
En particulier, un $R$-module est de dimension de Tor $0$ si et seulement
s'il est un $R$-module plat.
\end{lemma}

\begin{proof}
Supposons (2). Alors le complexe $E^\bullet$ défini par $E^{-i} = F_i$
est quasi-isomorphe à $M$. La dimension de Tor de $M$ est donc
au plus $d$ d'après le
lemme \ref{lemma-tor-amplitude}.
Réciproquement, supposons (1). Soit $P^\bullet \to M$ une résolution
projective de $M$. D'après le
lemme \ref{lemma-last-one-flat},
$\tau_{\geq -d}P^\bullet$ est une résolution plate de
$M$ de longueur $d$, c'est-à-dire que (2) est vraie.
\end{proof}

\begin{lemma}
\label{lemma-summands-tor-amplitude}
Soit $R$ un anneau. Soient $a, b \in \mathbf{Z}$.
Si $K^\bullet \oplus L^\bullet$ est d'amplitude de Tor dans $[a, b]$,
il en va de même de $K^\bullet$ et de $L^\bullet$.
\end{lemma}

\begin{proof}
Cela résulte du fait que les foncteurs Tor sont additifs.
\end{proof}

\begin{lemma}
\label{lemma-complex-finite-tor-dimension-modules}
Soit $R$ un anneau. Soit $K^\bullet$ un complexe borné de $R$-modules
tel que $K^i$ soit d'amplitude de Tor dans $[a - i, b - i]$ pour tout $i$.
Alors $K^\bullet$ est d'amplitude de Tor dans $[a, b]$. En particulier,
si $K^\bullet$ est un complexe fini de $R$-modules de dimension de Tor finie,
alors $K^\bullet$ est de dimension de Tor finie.
\end{lemma}

\begin{proof}
Cela se démontre par récurrence sur la longueur du complexe fini, en utilisant
le lemme \ref{lemma-cone-tor-amplitude}
et les troncatures brutales.
\end{proof}

\begin{lemma}
\label{lemma-cohomology-tor-amplitude}
Soit $R$ un anneau. Soient $a, b \in \mathbf{Z}$. Soit
$K^\bullet \in D^b(R)$ tel que $H^i(K^\bullet)$ soit d'amplitude de Tor dans
$[a - i, b - i]$ pour tout $i$. Alors $K^\bullet$ est d'amplitude de Tor dans
$[a, b]$. En particulier, si $K^\bullet \in D^b(R)$ et si tous ses groupes
de cohomologie sont de dimension de Tor finie, alors $K^\bullet$ est de
dimension de Tor finie.
\end{lemma}

\begin{proof}
Cela se démontre par récurrence sur la longueur du complexe fini, en utilisant
le lemme \ref{lemma-cone-tor-amplitude}
et les troncatures canoniques.
\end{proof}

\begin{lemma}
\label{lemma-push-tor-amplitude}
Soit $A \to B$ un homomorphisme d'anneaux. Soient $K^\bullet$ et $L^\bullet$
des complexes de $B$-modules. Soient $a, b, c, d \in \mathbf{Z}$. Si
\begin{enumerate}
\item $K^\bullet$, comme complexe de $B$-modules, est d'amplitude de Tor dans
$[a, b]$,
\item $L^\bullet$, comme complexe de $A$-modules, est d'amplitude de Tor dans
$[c, d]$,
\end{enumerate}
alors $K^\bullet \otimes^\mathbf{L}_B L^\bullet$, comme complexe de
$A$-modules, est d'amplitude de Tor dans $[a + c, b + d]$.
\end{lemma}

\begin{proof}
Nous pouvons supposer que $K^\bullet$ est un complexe de $B$-modules plats
tel que $K^i = 0$ pour $i \not \in [a, b]$ ; voir le
lemme \ref{lemma-tor-amplitude}.
Soit $M$ un $A$-module. Choisissons une résolution libre
$F^\bullet \to M$. Alors
$$
(K^\bullet \otimes_B^\mathbf{L} L^\bullet) \otimes_A^{\mathbf{L}} M =
\text{Tot}(\text{Tot}(K^\bullet \otimes_B L^\bullet) \otimes_A F^\bullet) =
\text{Tot}(K^\bullet \otimes_B \text{Tot}(L^\bullet \otimes_A F^\bullet))
$$
voir Homologie, remarque \ref{homology-remark-triple-complex}, pour la seconde
égalité. D'après l'hypothèse (2), le complexe
$\text{Tot}(L^\bullet \otimes_A F^\bullet)$
n'a de cohomologie non nulle qu'aux degrés de $[c, d]$. La suite spectrale
d'Homologie, lemme \ref{homology-lemma-ss-double-complex},
associée au bicomplexe
$K^\bullet \otimes_B \text{Tot}(L^\bullet \otimes_A F^\bullet)$
montre donc que
$(K^\bullet \otimes_B^\mathbf{L} L^\bullet) \otimes_A^{\mathbf{L}} M$
n'a de cohomologie non nulle qu'aux degrés de $[a + c, b + d]$.
\end{proof}

\begin{lemma}
\label{lemma-flat-push-tor-amplitude}
Soit $A \to B$ un homomorphisme d'anneaux. Supposons que $B$ soit plat
comme $A$-module. Soit $K^\bullet$ un complexe de $B$-modules.
Soient $a, b \in \mathbf{Z}$. Si $K^\bullet$, comme complexe de $B$-modules,
est d'amplitude de Tor dans $[a, b]$, alors $K^\bullet$, comme complexe de
$A$-modules, est d'amplitude de Tor dans $[a, b]$.
\end{lemma}

\begin{proof}
C'est un cas particulier du lemme \ref{lemma-push-tor-amplitude}, mais on peut
aussi le voir directement comme suit. Nous avons
$K^\bullet \otimes_A^{\mathbf{L}} M =
K^\bullet \otimes_B^{\mathbf{L}} (M \otimes_A B)$
car toute résolution projective de $K^\bullet$ comme complexe de $B$-modules
est une résolution plate de $K^\bullet$ comme complexe de $A$-modules et
peut servir à calculer $K^\bullet \otimes_A^{\mathbf{L}} M$.
\end{proof}

\begin{lemma}
\label{lemma-finite-tor-dimension-push-tor-amplitude}
Soit $A \to B$ un homomorphisme d'anneaux. Supposons que $B$ soit de
dimension de Tor $\leq d$ comme $A$-module. Soit $K^\bullet$ un complexe de
$B$-modules. Soient $a, b \in \mathbf{Z}$. Si $K^\bullet$, comme complexe
de $B$-modules, est d'amplitude de Tor dans $[a, b]$, alors $K^\bullet$,
comme complexe de $A$-modules, est d'amplitude de Tor dans $[a - d, b]$.
\end{lemma}

\begin{proof}
C'est un cas particulier du lemme \ref{lemma-push-tor-amplitude}, mais on peut
aussi le voir directement comme suit.
Soit $M$ un $A$-module. Choisissons une résolution libre
$F^\bullet \to M$. Alors
$$
K^\bullet \otimes_A^{\mathbf{L}} M =
\text{Tot}(K^\bullet \otimes_A F^\bullet) =
\text{Tot}(K^\bullet \otimes_B (F^\bullet \otimes_A B)) =
K^\bullet \otimes_B^{\mathbf{L}} (M \otimes_A^{\mathbf{L}} B).
$$
D'après notre hypothèse sur $B$ comme $A$-module,
$M \otimes_A^{\mathbf{L}} B$ n'a de cohomologie qu'aux degrés
$-d, -d + 1, \ldots, 0$. Puisque $K^\bullet$ est d'amplitude de Tor dans
$[a, b]$, la suite spectrale de
l'exemple \ref{example-tor}
montre que $K^\bullet \otimes_B^{\mathbf{L}} (M \otimes_A^{\mathbf{L}} B)$
n'a de cohomologie qu'aux degrés de $[-d + a, b]$, comme souhaité.
\end{proof}

\begin{lemma}
\label{lemma-pull-tor-amplitude}
Soit $A \to B$ un homomorphisme d'anneaux.
Soient $a, b \in \mathbf{Z}$.
Soit $K^\bullet$ un complexe de $A$-modules d'amplitude de Tor dans $[a, b]$.
Alors $K^\bullet \otimes_A^{\mathbf{L}} B$, comme complexe de $B$-modules,
est d'amplitude de Tor dans $[a, b]$.
\end{lemma}

\begin{proof}
D'après le
lemme \ref{lemma-tor-amplitude},
il existe un quasi-isomorphisme $E^\bullet \to K^\bullet$, où
$E^\bullet$ est un complexe de $A$-modules plats tel que $E^i = 0$ pour
$i \not \in [a, b]$. Par construction, $E^\bullet \otimes_A B$ calcule
$K^\bullet \otimes_A ^{\mathbf{L}} B$, et chaque $E^i \otimes_A B$ est un
$B$-module plat d'après Algèbre, lemme
\ref{algebra-lemma-flat-base-change}.
Nous concluons donc par le
lemme \ref{lemma-tor-amplitude}.
\end{proof}

\begin{lemma}
\label{lemma-flat-base-change-finite-tor-dimension}
Soit $A \to B$ un homomorphisme plat d'anneaux. Soit $d \geq 0$.
Soit $M$ un $A$-module de dimension de Tor $\leq d$.
Alors $M \otimes_A B$ est un $B$-module de dimension de Tor $\leq d$.
\end{lemma}

\begin{proof}
C'est une conséquence immédiate du
lemme \ref{lemma-pull-tor-amplitude}
et de l'égalité $M \otimes_A^{\mathbf{L}} B = M \otimes_A B$,
car $B$ est plat sur $A$.
\end{proof}

\begin{lemma}
\label{lemma-tor-amplitude-localization}
Soit $A  \to B$ un homomorphisme d'anneaux. Soit $K^\bullet$ un complexe de
$B$-modules. Soient $a, b \in \mathbf{Z}$. Les assertions suivantes sont
équivalentes :
\begin{enumerate}
\item $K^\bullet$ est d'amplitude de Tor dans $[a, b]$ comme complexe de
$A$-modules,
\item $K^\bullet_\mathfrak q$ est d'amplitude de Tor dans $[a, b]$ comme
complexe de $A_\mathfrak p$-modules, pour tout idéal premier
$\mathfrak q \subset B$ tel que $\mathfrak p = A \cap \mathfrak q$,
\item $K^\bullet_\mathfrak m$ est d'amplitude de Tor dans $[a, b]$ comme
complexe de $A_\mathfrak p$-modules, pour tout idéal maximal
$\mathfrak m \subset B$ tel que $\mathfrak p = A \cap \mathfrak m$.
\end{enumerate}
\end{lemma}

\begin{proof}
Supposons (3) et soit $M$ un $A$-module. Alors
$H^i = H^i(K^\bullet \otimes_A^\mathbf{L} M)$ est un $B$-module et
$(H^i)_\mathfrak m =
H^i(K^\bullet_\mathfrak m \otimes_{A_\mathfrak p}^\mathbf{L} M_\mathfrak p)$.
Ainsi $H^i = 0$ pour $i \not \in [a, b]$ d'après Algèbre,
lemme \ref{algebra-lemma-characterize-zero-local}. Donc
(3) $\Rightarrow$ (1). Nous omettons les démonstrations de
(1) $\Rightarrow$ (2) et de (2) $\Rightarrow$ (3).
\end{proof}

\begin{lemma}
\label{lemma-glue-tor-amplitude}
Soit $R$ un anneau. Soient $f_1, \ldots, f_r \in R$ des éléments qui
engendrent l'idéal unité. Soient $a, b \in \mathbf{Z}$. Soit $K^\bullet$
un complexe de $R$-modules. Si, pour chaque $i$, le complexe
$K^\bullet \otimes_R R_{f_i}$ est d'amplitude de Tor dans $[a, b]$,
alors $K^\bullet$ est d'amplitude de Tor dans $[a, b]$.
\end{lemma}

\begin{proof}
Cela résulte immédiatement du lemme
\ref{lemma-tor-amplitude-localization}, mais se voit également directement
comme suit. Remarquons que $- \otimes_R R_{f_i}$ est un foncteur exact et que
$$
H^i(K^\bullet)_{f_i} =
H^i(K^\bullet) \otimes_R R_{f_i} = H^i(K^\bullet \otimes_R R_{f_i}).
$$
De même, pour tout $R$-module $M$, nous avons
$$
H^i(K^\bullet \otimes_R^{\mathbf{L}} M)_{f_i} =
H^i(K^\bullet \otimes_R^{\mathbf{L}} M) \otimes_R R_{f_i} =
H^i(K^\bullet \otimes_R R_{f_i} \otimes_{R_{f_i}}^{\mathbf{L}} M_{f_i}).
$$
Le résultat découle donc du fait qu'un $R$-module $N$ est nul si et seulement
si $N_{f_i}$ est nul pour tout $i$ ; voir Algèbre,
lemme \ref{algebra-lemma-cover}.
\end{proof}

\begin{lemma}
\label{lemma-flat-descent-tor-amplitude}
Soit $R$ un anneau. Soient $a, b \in \mathbf{Z}$. Soit $K^\bullet$
un complexe de $R$-modules. Soit $R \to R'$ un homomorphisme d'anneaux
fidèlement plat. Si le complexe $K^\bullet \otimes_R R'$ est d'amplitude
de Tor dans $[a, b]$, alors $K^\bullet$ est d'amplitude de Tor dans $[a, b]$.
\end{lemma}

\begin{proof}
Soit $M$ un $R$-module. Puisque $R \to R'$ est plat, nous avons
$$
(M \otimes_R^{\mathbf{L}} K^\bullet) \otimes_R R'
=
((M \otimes_R R') \otimes_{R'}^{\mathbf{L}} (K^\bullet \otimes_R R')
$$
et la prise de cohomologie commute au produit tensoriel par $R'$.
Ainsi $\text{Tor}_i^R(M, K^\bullet) \otimes_R R' =
\text{Tor}_i^{R'}(M \otimes_R R', K^\bullet \otimes_R R')$.
Puisque $R \to R'$ est fidèlement plat, l'annulation de
$\text{Tor}_i^{R'}(M \otimes_R R', K^\bullet \otimes_R R')$ pour
$i \not \in [a, b]$ implique la même annulation pour
$\text{Tor}_i^R(M, K^\bullet)$.
\end{proof}

\begin{lemma}
\label{lemma-no-change-tor-amplitude}
Soient des homomorphismes d'anneaux $R \to A \to B$ avec $A \to B$
fidèlement plat, et soit $K \in D(A)$. L'amplitude de Tor de $K$ sur $R$
est la même que celle de $K \otimes_A^\mathbf{L} B$ sur $R$.
\end{lemma}

\begin{proof}
En effet, pour tout $R$-module $M$, nous avons
$H^i(K \otimes_R^\mathbf{L} M) \otimes_A B =
H^i((K \otimes_A^\mathbf{L} B) \otimes_R^\mathbf{L} M)$
pour tout $i$. En effet, représentons $K$ par un complexe
$K^\bullet$ de $A$-modules et choisissons une résolution libre
$F^\bullet \to M$. Nous avons alors l'égalité
$$
\text{Tot}(K^\bullet \otimes_A B \otimes_R F^\bullet) =
\text{Tot}(K^\bullet \otimes_R F^\bullet) \otimes_A B
$$
Les groupes de cohomologie du membre de gauche sont
$H^i((K \otimes_A^\mathbf{L} B) \otimes_R^\mathbf{L} M)$
et, dans le membre de droite, nous obtenons
$H^i(K \otimes_R^\mathbf{L} M) \otimes_A B$.
\end{proof}

\begin{lemma}
\label{lemma-finite-gl-dim-tor-dimension}
Soit $R$ un anneau de dimension globale finie $d$. Alors
\begin{enumerate}
\item tout module est de dimension de Tor $\leq d$,
\item tout complexe de $R$-modules $K^\bullet$ tel que
$H^i(K^\bullet) \not = 0$ seulement si $i \in [a, b]$ est d'amplitude de Tor
dans $[a - d, b]$, et
\item un complexe de $R$-modules $K^\bullet$ est de dimension de Tor finie
si et seulement si $K^\bullet \in D^b(R)$.
\end{enumerate}
\end{lemma}

\begin{proof}
L'hypothèse sur $R$ signifie que tout module possède une résolution projective
finie de longueur au plus $d$ ; en particulier, tout module est de dimension
de Tor $\leq d$. La deuxième assertion résulte du
lemme \ref{lemma-cohomology-tor-amplitude}
et des définitions. La troisième est une reformulation de la deuxième.
\end{proof}

\begin{lemma}
\label{lemma-check-tor-dimension-modulo-nilpotent}
Soit $R' \to R$ un homomorphisme surjectif d'anneaux dont le noyau est un
idéal nilpotent. Soit $K' \in D(R')$ et posons
$K = K' \otimes_{R'}^\mathbf{L} R$. Soient $a, b \in \mathbf{Z}$.
Alors $K$ est d'amplitude de Tor dans $[a, b]$ si et seulement si $K'$ l'est.
\end{lemma}

\begin{proof}
Une implication résulte du lemme \ref{lemma-pull-tor-amplitude}.
Pour l'autre, supposons que $K$ soit d'amplitude de Tor dans $[a, b]$ et
soit $M'$ un $R'$-module. Nous devons montrer que
$K' \otimes_{R'}^\mathbf{L} M'$ n'a de cohomologie non nulle qu'aux degrés
appartenant à l'intervalle $[a, b]$.

\medskip\noindent
Posons $I = \Ker(R' \to R)$. Alors $I^n = 0$ pour un certain $n$.
Si $IM' = 0$, nous pouvons considérer $M'$ comme un $R$-module et raisonner
comme suit :
$$
K' \otimes_{R'}^\mathbf{L} M' =
K' \otimes_{R'}^\mathbf{L} (R \otimes_R^\mathbf{L} M') =
(K' \otimes_{R'}^\mathbf{L} R) \otimes_R^\mathbf{L} M' =
K \otimes_R^\mathbf{L} M'
$$
qui n'a de cohomologie non nulle que dans l'intervalle $[a, b]$
par l'hypothèse sur $K$. Si $I^{t + 1}M' = 0$, considérons la suite exacte
courte
$$
0 \to IM' \to M' \to M'/IM' \to 0
$$
Par récurrence sur $t$, les deux complexes
$K' \otimes_{R'}^\mathbf{L} IM'$ et
$K' \otimes_{R'}^\mathbf{L} M'/IM'$ n'ont de cohomologie non nulle qu'aux
degrés de l'intervalle $[a, b]$. Le triangle distingué
$$
K' \otimes_{R'}^\mathbf{L} IM' \to
K' \otimes_{R'}^\mathbf{L} M' \to
K' \otimes_{R'}^\mathbf{L} M'/IM' \to
(K' \otimes_{R'}^\mathbf{L} IM')[1]
$$
montre alors qu'il en va de même de $K' \otimes_{R'}^\mathbf{L} M'$,
comme souhaité.
\end{proof}







\section{Suites spectrales pour Ext}
\label{section-spectral-sequence-ext}

\noindent
Dans cette section, nous rassemblons diverses suites spectrales qui
apparaissent lorsque l'on considère les foncteurs Ext. Pour toute paire
d'objets $L$, $K$ de la catégorie dérivée $D(R)$ d'un anneau $R$,
nous posons
$$
\Ext^n_R(L, K) = \Hom_{D(R)}(L, K[n])
$$
conformément à nos conventions générales de
Catégories dérivées, section \ref{derived-section-ext}.

\medskip\noindent
Pour un module $M$ sur $R$ et $K \in D^+(R)$, il existe une suite spectrale
\begin{equation}
\label{equation-first-ss-ext}
E_2^{i, j} = \Ext_R^i(M, H^j(K)) \Rightarrow \Ext_R^{i + j}(M, K)
\end{equation}
et, si $K$ est représenté par un complexe $K^\bullet$ de $R$-modules borné
inférieurement, il existe une suite spectrale
\begin{equation}
\label{equation-second-ss-ext}
E_1^{i, j} = \Ext_R^j(M, K^i) \Rightarrow \Ext_R^{i + j}(M, K)
\end{equation}
Ces suites spectrales proviennent de l'application du lemme de
Catégories dérivées \ref{derived-lemma-two-ss-complex-functor}
au foncteur $\Hom_R(M, -)$.




\section{Dimension projective}
\label{section-projective-dimension}

\noindent
Nous avons défini la dimension projective d'un module dans
Algèbre, définition \ref{algebra-definition-finite-proj-dim}.

\begin{definition}
\label{definition-projective-dimension}
Soit $R$ un anneau. Soit $K$ un objet de $D(R)$. Nous disons que $K$ est
de {\it dimension projective finie} si $K$ peut être représenté par un
complexe borné de modules projectifs. Nous disons que $K$ est
d'{\it amplitude projective dans $[a, b]$} si $K$ est quasi-isomorphe
à un complexe
$$
\ldots \to 0 \to P^a \to P^{a + 1} \to \ldots \to
P^{b - 1} \to P^b \to 0 \to \ldots
$$
où $P^i$ est un $R$-module projectif pour tout $i \in \mathbf{Z}$.
\end{definition}

\noindent
Il est clair que $K$ est de dimension projective finie si et seulement si
$K$ est d'amplitude projective dans $[a, b]$ pour certains
$a, b \in \mathbf{Z}$. De plus, si $K$ est de dimension projective finie,
alors $K$ est borné. Voici un lemme permettant de reconnaître de tels objets
de $D(R)$.

\begin{lemma}
\label{lemma-projective-amplitude}
Soit $R$ un anneau. Soit $K$ un objet de $D(R)$. Soient
$a, b \in \mathbf{Z}$. Les assertions suivantes sont équivalentes :
\begin{enumerate}
\item $K$ est d'amplitude projective dans $[a, b]$,
\item $\Ext^i_R(K, N) = 0$ pour tout $R$-module $N$ et tout
$i \not \in [-b, -a]$,
\item $H^n(K) = 0$ pour $n > b$ et
$\Ext^i_R(K, N) = 0$ pour tout $R$-module $N$ et tout $i > -a$, et
\item $H^n(K) = 0$ pour $n \not \in [a - 1, b]$ et
$\Ext^{-a + 1}_R(K, N) = 0$ pour tout $R$-module $N$.
\end{enumerate}
\end{lemma}

\begin{proof}
Supposons (1). Nous pouvons supposer que $K$ est le complexe
$$
\ldots \to 0 \to P^a \to P^{a + 1} \to \ldots \to
P^{b - 1} \to P^b \to 0 \to \ldots
$$
où $P^i$ est un $R$-module projectif pour tout $i \in \mathbf{Z}$.
Dans ce cas, nous pouvons calculer les groupes Ext à l'aide du complexe
$$
\ldots \to 0 \to \Hom_R(P^b, N) \to \ldots \to
\Hom_R(P^a, N) \to 0 \to \ldots
$$
et nous obtenons (2).

\medskip\noindent
Supposons (2) vraie. Choisissons une injection $H^n(K) \to I$, où $I$
est un $R$-module injectif. Puisque $\Hom_R(-, I)$ est un foncteur exact,
nous avons $\Ext^{-n}(K, I) = \Hom_R(H^n(K), I)$.
Nous concluons en particulier que $H^n(K)$ est nul pour $n > b$.
Ainsi (2) implique (3).

\medskip\noindent
Le même argument que pour le passage de (2) à (3) montre que (3) implique (4).

\medskip\noindent
Supposons (4). Le même argument que pour le passage de (2) à (3) montre que
$H^{a - 1}(K) = 0$, c'est-à-dire que $H^i(K) = 0$ sauf si
$i \in [a, b]$. En particulier, $K$ est borné supérieurement et nous pouvons
choisir un complexe $P^\bullet$ représentant $K$, avec $P^i$ projectif
(par exemple libre) pour tout $i \in \mathbf{Z}$ et $P^i = 0$ pour $i > b$.
Voir Catégories dérivées, lemme
\ref{derived-lemma-subcategory-left-resolution}.
Posons $Q = \Coker(P^{a - 1} \to P^a)$. Alors $K$ est quasi-isomorphe
au complexe
$$
\ldots \to 0 \to Q \to P^{a + 1} \to \ldots \to P^b \to 0 \to \ldots
$$
car $H^i(K) = 0$ pour $i < a$.
Notons $K' = (P^{a + 1} \to \ldots \to P^b)$ l'objet correspondant de
$D(R)$. Nous obtenons un triangle distingué
$$
K' \to K \to Q[-a] \to K'[1]
$$
dans $D(R)$. Ainsi, pour tout $R$-module $N$, nous avons une suite exacte
$$
\Ext^{-a}(K', N) \to \Ext^1(Q, N) \to \Ext^{1 - a}(K, N)
$$
Par hypothèse, le terme de droite est nul. D'après l'implication
(1) $\Rightarrow$ (2), le terme de gauche est nul. Ainsi $Q$
est un $R$-module projectif d'après Algèbre,
lemme \ref{algebra-lemma-characterize-projective}.
Donc (1) est vraie, ce qui achève la démonstration.
\end{proof}

\begin{example}
\label{example-ext-not-bounded}
Soit $k$ un corps et soit $R$ l'anneau des nombres duaux sur $k$,
c'est-à-dire $R = k[x]/(x^2)$. Notons $\epsilon \in R$ la classe de $x$.
Soit $M = R/(\epsilon)$. Alors $M$ est quasi-isomorphe au complexe
$$
R \xrightarrow{\epsilon} R \xrightarrow{\epsilon} R \to \ldots
$$
mais $M$ n'est pas de dimension projective finie au sens d'
Algèbre, définition \ref{algebra-definition-finite-proj-dim}.
Cela explique pourquoi nous considérons des complexes de modules projectifs
bornés dans les deux directions dans notre définition de la dimension
projective finie des objets de $D(R)$.
\end{example}






\section{Dimension injective}
\label{section-injective-dimension}

\noindent
Cette section est duale de celle consacrée à la dimension projective.

\begin{definition}
\label{definition-injective-dimension}
Soit $R$ un anneau. Soit $K$ un objet de $D(R)$.
Nous disons que $K$ est de {\it dimension injective finie} si $K$ peut être
représenté par un complexe fini de $R$-modules injectifs.
Nous disons que $K$ est d'{\it amplitude injective dans $[a, b]$}
si $K$ est isomorphe à un complexe
$$
\ldots \to 0 \to I^a \to I^{a + 1} \to \ldots \to
I^{b - 1} \to I^b \to 0 \to \ldots
$$
où $I^i$ est un $R$-module injectif pour tout $i \in \mathbf{Z}$.
\end{definition}

\noindent
Il est clair que $K$ est de dimension injective bornée si et seulement si
$K$ est d'amplitude injective dans $[a, b]$ pour certains
$a, b \in \mathbf{Z}$. De plus, si $K$ est de dimension injective bornée,
alors $K$ est borné. Voici le lemme obligé.

\begin{lemma}
\label{lemma-injective-amplitude}
Soit $R$ un anneau. Soit $K$ un objet de $D(R)$. Soient
$a, b \in \mathbf{Z}$. Les assertions suivantes sont équivalentes :
\begin{enumerate}
\item $K$ est d'amplitude injective dans $[a, b]$,
\item $\Ext^i_R(N, K) = 0$ pour tout $R$-module $N$ et tout
$i \not \in [a, b]$,
\item $\Ext^i(R/I, K) = 0$ pour tout idéal $I \subset R$ et
tout $i \not \in [a, b]$.
\end{enumerate}
\end{lemma}

\begin{proof}
Supposons (1). Nous pouvons supposer que $K$ est le complexe
$$
\ldots \to 0 \to I^a \to I^{a + 1} \to \ldots \to
I^{b - 1} \to I^b \to 0 \to \ldots
$$
où $I^i$ est un $R$-module injectif pour tout $i \in \mathbf{Z}$.
Dans ce cas, nous pouvons calculer les groupes Ext à l'aide du complexe
$$
\ldots \to 0 \to \Hom_R(N, I^a) \to \ldots \to
\Hom_R(N, I^b) \to 0 \to \ldots
$$
et nous obtenons (2). Il est clair que (2) implique (3).

\medskip\noindent
Supposons (3) vraie. Choisissons une application non nulle
$R \to H^n(K)$. Puisque $\Hom_R(R, -)$ est un foncteur exact, nous avons
$\Ext^n_R(R, K) = \Hom_R(R, H^n(K)) = H^n(K)$.
Nous concluons que $H^n(K)$ est nul pour $n \not \in [a, b]$.
En particulier, $K$ est borné inférieurement et nous pouvons choisir un
quasi-isomorphisme
$$
K \to I^\bullet
$$
où $I^i$ est injectif pour tout $i \in \mathbf{Z}$ et
$I^i = 0$ pour $i < a$. Voir Catégories dérivées, lemme
\ref{derived-lemma-subcategory-right-resolution}.
Posons $J = \Ker(I^b \to I^{b + 1})$. Alors $K$ est quasi-isomorphe
au complexe
$$
\ldots \to 0 \to I^a \to \ldots \to I^{b - 1} \to J \to 0 \to \ldots
$$
Notons $K' = (I^a \to \ldots \to I^{b - 1})$ l'objet correspondant de
$D(R)$. Nous obtenons un triangle distingué
$$
J[-b] \to K \to K' \to J[1 - b]
$$
dans $D(R)$. Ainsi, pour tout idéal $I \subset R$, nous avons une suite exacte
$$
\Ext^b(R/I, K') \to \Ext^1(R/I, J) \to \Ext^{1 + b}(R/I, K)
$$
Par hypothèse, le terme de droite est nul. D'après l'implication
(1) $\Rightarrow$ (2), le terme de gauche est nul. Ainsi $J$
est un $R$-module injectif d'après le
lemme \ref{lemma-characterize-injective-bis}.
\end{proof}

\begin{example}
\label{example-finite-injective-finite-global-dimension}
Soit $R$ un anneau de Dedekind. Alors tout idéal non nul $I$ est un module
projectif de type fini ; voir le
lemme \ref{lemma-dedekind-torsion-free-flat}.
Ainsi $R/I$ est de dimension projective $1$. Par conséquent, tout
$R$-module $M$ est de dimension injective $\leq 1$ d'après le
lemme \ref{lemma-injective-amplitude}. Ainsi
$\Ext^i_R(M, N) = 0$ pour $i \geq 2$ et toute paire de $R$-modules $M, N$.
Il s'ensuit que tout objet $K$ de $D^b(R)$ est isomorphe à la somme directe
de ses cohomologies : $K \cong \bigoplus H^i(K)[-i]$ ; voir
Catégories dérivées, lemme \ref{derived-lemma-ext-2-zero}.
\end{example}

\begin{example}
\label{example-ext-not-bounded-reversed}
Soit $k$ un corps et soit $R$ l'anneau des nombres duaux sur $k$,
c'est-à-dire $R = k[x]/(x^2)$. Notons $\epsilon \in R$ la classe de $x$.
Soit $M = R/(\epsilon)$. Alors $M$ est quasi-isomorphe au complexe
$$
\ldots \to R \xrightarrow{\epsilon} R \xrightarrow{\epsilon} R
$$
et $R$ est un $R$-module injectif. Cependant, dans cette situation, on ne
considère généralement pas que $M$ est de dimension injective finie.
Cela explique pourquoi nous considérons des complexes de modules injectifs
bornés dans les deux directions dans notre définition de la dimension
injective bornée des objets de $D(R)$.
\end{example}

\begin{lemma}
\label{lemma-finite-injective-dimension}
Soit $R$ un anneau. Soit $K \in D(R)$.
\begin{enumerate}
\item Si $K$ appartient à $D^b(R)$ et si $H^i(K)$ est de dimension injective
finie pour tout $i$, alors $K$ est de dimension injective finie.
\item Si $K^\bullet$ représente $K$, est un complexe borné de $R$-modules,
et si $K^i$ est de dimension injective finie pour tout $i$, alors $K$ est de
dimension injective finie.
\end{enumerate}
\end{lemma}

\begin{proof}
Omis. Indication : appliquer les suites spectrales de
Catégories dérivées, lemme \ref{derived-lemma-two-ss-complex-functor},
au foncteur $F = \Hom_R(N, -)$ afin de calculer
$\Ext^i_A(N, K)$, puis utiliser le critère du
lemme \ref{lemma-injective-amplitude}.
\end{proof}

\begin{lemma}
\label{lemma-finite-injective-dimension-Noetherian-radical}
Soit $R$ un anneau noethérien. Soit $I \subset R$ un idéal contenu dans
le radical de Jacobson de $R$. Soit $K \in D^+(R)$ à modules de cohomologie
de type fini. Les assertions suivantes sont équivalentes :
\begin{enumerate}
\item $K$ est de dimension injective finie, et
\item il existe un $b$ tel que $\Ext^i_R(R/J, K) = 0$ pour $i > b$
et tout idéal $J \supset I$.
\end{enumerate}
\end{lemma}

\begin{proof}
L'implication (1) $\Rightarrow$ (2) est immédiate. Supposons (2).
Disons que $H^i(K) = 0$ pour $i < a$. Alors $\Ext^i(M, K) = 0$ pour
$i < a$ et tout $R$-module $M$. Il suffit donc de montrer que
$\Ext^i(M, K) = 0$ pour $i > b$ et tout $R$-module de type fini $M$ ; voir
le lemme \ref{lemma-injective-amplitude}.
D'après Algèbre, lemme \ref{algebra-lemma-filter-Noetherian-module},
le module $M$ possède une filtration finie dont les quotients successifs
sont de la forme $R/\mathfrak p$, où $\mathfrak p$ est un idéal premier.
Si $0 \to M_1 \to M \to M_2 \to 0$ est une suite exacte courte et si
$\Ext^i(M_j, K) = 0$ pour $i > b$ et $j = 1, 2$, alors
$\Ext^i(M, K) = 0$ pour $i > b$. Nous pouvons donc supposer
$M = R/\mathfrak p$. Si $I \subset \mathfrak p$, l'annulation résulte de
l'hypothèse. Sinon, choisissons $f \in I$, $f \not \in \mathfrak p$.
Considérons la suite exacte courte
$$
0 \to R/\mathfrak p \xrightarrow{f} R/\mathfrak p \to R/(\mathfrak p, f) \to 0
$$
Le $R$-module $R/(\mathfrak p, f)$ possède une filtration dont les quotients
successifs sont des $R/\mathfrak q$ avec
$(\mathfrak p, f) \subset \mathfrak q$. Ainsi, par récurrence noethérienne
et par l'argument précédent, nous pouvons supposer que l'annulation vaut
pour $R/(\mathfrak p, f)$. D'autre part, les modules
$E^i = \Ext^i(R/\mathfrak p, K)$ sont de type fini d'après notre hypothèse
sur $K$ (borné inférieurement à modules de cohomologie de type fini), la
suite spectrale (\ref{equation-first-ss-ext}) et Algèbre,
lemme \ref{algebra-lemma-ext-noetherian}.
Ainsi, $E^i$, pour $i > b$, est un $R$-module de type fini tel que
$E^i/fE^i = 0$. Le lemme de Nakayama
(Algèbre, lemme \ref{algebra-lemma-NAK})
montre que $E^i$ est nul.
\end{proof}

\begin{lemma}
\label{lemma-finite-injective-dimension-Noetherian-local}
Soit $(R, \mathfrak m, \kappa)$ un anneau local noethérien.
Soit $K \in D^+(R)$ à modules de cohomologie de type fini.
Les assertions suivantes sont équivalentes :
\begin{enumerate}
\item $K$ est de dimension injective finie, et
\item $\Ext^i_R(\kappa, K) = 0$ for $i \gg 0$.
\end{enumerate}
\end{lemma}

\begin{proof}
C'est un cas particulier du
lemme \ref{lemma-finite-injective-dimension-Noetherian-radical}.
\end{proof}






\section{Modules presque projectifs}
\label{section-near-projective}

\noindent
Il semble exister dans la littérature de nombreuses définitions différentes
des « modules presque projectifs ». Dans cette section, nous n'examinons
qu'une possibilité parmi beaucoup d'autres.

\begin{lemma}
\label{lemma-ideal-factor-through-projective}
Soit $R$ un anneau. Soient $M$, $N$ des $R$-modules.
\begin{enumerate}
\item Pour un homomorphisme de $R$-modules $\varphi : M \to N$, les assertions
suivantes sont équivalentes : (a) $\varphi$ se factorise par un $R$-module
projectif, et (b) $\varphi$ se factorise par un $R$-module libre.
\item L'ensemble des $\varphi : M \to N$ qui satisfont aux conditions
équivalentes de (1) est un sous-$R$-module de $\Hom_R(M, N)$.
\item Étant donnés des homomorphismes $\psi : M' \to M$ et
$\xi : N \to N'$, si $\varphi : M \to N$ satisfait aux conditions
équivalentes de (1), alors il en va de même de
$\xi \circ \varphi \circ \psi : M' \to N'$.
\end{enumerate}
\end{lemma}

\begin{proof}
L'équivalence de (1)(a) et (1)(b) résulte d'
Algèbre, lemme \ref{algebra-lemma-characterize-projective}.
Si $\varphi : M \to N$ et $\varphi' : M \to N$ se factorisent respectivement
par les modules $P$ et $P'$, alors $\varphi + \varphi'$ se factorise par
$P \oplus P'$ et $\lambda \varphi$ se factorise par $P$ pour tout
$\lambda \in R$. Cela démontre (2). Si $\varphi : M \to N$ se factorise
par le module $P$ et si $\psi$ et $\xi$ sont comme dans (3), alors
$\xi \circ \varphi \circ \psi$ se factorise par $P$. Cela démontre (3).
\end{proof}

\begin{lemma}
\label{lemma-factor-through-finite-projective}
Soit $R$ un anneau. Soit $\varphi : M \to N$ un homomorphisme de
$R$-modules. Si $\varphi$ se factorise par un module projectif et si $M$
est un $R$-module de type fini, alors $\varphi$ se factorise par un module
projectif de type fini.
\end{lemma}

\begin{proof}
D'après le lemme \ref{lemma-ideal-factor-through-projective}, nous pouvons
écrire $\varphi = \tau \circ \sigma$, où le but de $\sigma$ est
$\bigoplus_{i \in I} R$ pour un certain ensemble $I$. Choisissons des
générateurs $x_1, \ldots, x_n$ de $M$. Écrivons
$\sigma(x_j) = (a_{ji})_{i \in I}$. Pour chaque $j$, seul un nombre fini
des $a_{ij}$ sont non nuls. Ainsi, l'image de $\sigma$ est contenue dans un
$R$-module libre de type fini, ce qui permet de conclure.
\end{proof}

\noindent
Soit $R$ un anneau. Observons qu'un $R$-module est projectif si et seulement
si l'identité de $R$ se factorise par un module projectif.

\begin{lemma}
\label{lemma-near-projective}
Soit $R$ un anneau. Soit $I \subset R$ un idéal. Soit $M$ un $R$-module.
Les conditions suivantes sont équivalentes :
\begin{enumerate}
\item pour tout $a \in I$, l'application $a : M \to M$ se factorise par un
$R$-module projectif,
\item pour tout $a \in I$, l'application $a : M \to M$ se factorise par un
$R$-module libre, et
\item $\Ext^1_R(M, N)$ est annulé par $I$ pour tout $R$-module $N$.
\end{enumerate}
\end{lemma}

\begin{proof}
L'équivalence de (1) et (2) résulte du
lemme \ref{lemma-ideal-factor-through-projective}.
Si (1) est vraie, alors (3) l'est, car $\Ext^1_R(P, N)$ est nul
pour tout $N$ et tout module projectif $P$.
Réciproquement, supposons (3) vraie. Choisissons une suite exacte courte
$0 \to N \to P \to M \to 0$ avec $P$ projectif (voire libre).
Par hypothèse, l'élément correspondant de $\Ext^1_R(M, N)$ est annulé par
$I$. Par conséquent, pour tout $a \in I$, l'application $a : M \to M$
se factorise par la surjection $P \to M$, et nous concluons que (1) est vraie.
\end{proof}

\noindent
Afin de parler commodément des modules qui satisfont aux conditions
équivalentes du lemme \ref{lemma-near-projective}, donnons un nom à cette
propriété.

\begin{definition}
\label{definition-near-projective}
Soit $R$ un anneau. Soit $I \subset R$ un idéal. Soit $M$ un $R$-module.
Nous disons que $M$ est {\it $I$-projectif}\footnote{Cette notation n'est
pas standard.} si les conditions équivalentes du
lemme \ref{lemma-near-projective} sont satisfaites.
\end{definition}

\noindent
Les modules annulés par $I$ sont $I$-projectifs.

\begin{lemma}
\label{lemma-torsion-near-projective}
Soit $R$ un anneau. Soit $I \subset R$ un idéal. Soit $M$ un $R$-module.
Si $M$ est annulé par $I$, alors $M$ est $I$-projectif.
\end{lemma}

\begin{proof}
Cela résulte immédiatement de la définition et du fait que le module nul est
projectif.
\end{proof}

\begin{lemma}
\label{lemma-ses-near-projective}
Soit $R$ un anneau. Soit $I \subset R$ un idéal. Soit
$$
0 \to K \to P \to M \to 0
$$
une suite exacte courte de $R$-modules.
Si $M$ est $I$-projectif et si $P$ est projectif, alors $K$ est
$I$-projectif.
\end{lemma}

\begin{proof}
L'élément $\text{id}_K \in \Hom_R(K, K)$ s'envoie sur la classe de
l'extension donnée dans $\Ext^1_R(M, K)$. Puisque, par hypothèse, cette
classe est annulée par tout $a \in I$, nous voyons que $a : K \to K$
se factorise par $K \to P$, ce qui permet de conclure.
\end{proof}

\begin{lemma}
\label{lemma-dual-near-projective}
Soit $R$ un anneau. Soit $I \subset R$ un idéal.
Si $M$ est de type fini et $I$-projectif comme $R$-module, alors
$M^\vee = \Hom_R(M, R)$ est $I$-projectif.
\end{lemma}

\begin{proof}
Supposons $M$ de type fini et $I$-projectif.
Choisissons une suite exacte courte
$0 \to K \to R^{\oplus r} \to M \to 0$.
Elle fournit une injection
$M^\vee \to R^{\oplus r} = (R^{\oplus r})^\vee$.
Puisque la classe d'extension dans $\Ext^1_R(M, K)$ correspondant à la suite
exacte courte est annulée par $I$, nous voyons que, pour tout $a \in I$,
il existe une application $M \to R^{\oplus r}$ dont la composée avec
l'application donnée $R^{\oplus r} \to M$ est égale à $a : M \to M$.
En passant aux duaux, nous voyons que $a : M^\vee \to M^\vee$ se factorise
par l'application $M^\vee \to R^{\oplus r}$ donnée ci-dessus, ce qui permet
de conclure.
\end{proof}





\section{Complexes Hom}
\label{section-hom-complexes}

\noindent
Soit $R$ un anneau. Soient $L^\bullet$ et $M^\bullet$ deux complexes de
$R$-modules. Construisons un complexe $\Hom^\bullet(L^\bullet, M^\bullet)$.
Pour chaque $n$, posons
$$
\Hom^n(L^\bullet, M^\bullet) =
\prod\nolimits_{n = p + q} \Hom_R(L^{-q}, M^p)
$$
Il est utile de considérer $\Hom^n$ comme le $R$-module de toutes les
applications $R$-linéaires de $L^\bullet$ vers $M^\bullet$ (considérés comme
modules gradués) qui sont homogènes de degré $n$. Avec cette terminologie,
nous définissons la différentielle par la règle
$$
\text{d}(f) = \text{d}_M \circ f - (-1)^n f \circ \text{d}_L
$$
pour $f \in \Hom^n(L^\bullet, M^\bullet)$.
Nous omettons la vérification de $\text{d}^2 = 0$.
Voir la section \ref{section-sign-rules} pour les règles de signes.
Cette construction est un cas particulier d'
Algèbre différentielle graduée, exemple
\ref{dga-example-category-complexes}.
Il résulte immédiatement de la construction que
\begin{equation}
\label{equation-cohomology-hom-complex}
H^n(\Hom^\bullet(L^\bullet, M^\bullet)) = \Hom_{K(R)}(L^\bullet, M^\bullet[n])
\end{equation}
pour tout $n \in \mathbf{Z}$.

\begin{lemma}
\label{lemma-compose}
Soit $R$ un anneau. Étant donnés des complexes
$K^\bullet, L^\bullet, M^\bullet$ de $R$-modules, il existe un isomorphisme
canonique
$$
\Hom^\bullet(K^\bullet, \Hom^\bullet(L^\bullet, M^\bullet))
=
\Hom^\bullet(\text{Tot}(K^\bullet \otimes_R L^\bullet), M^\bullet)
$$
de complexes de $R$-modules.
\end{lemma}

\begin{proof}
Soit $\alpha$ un élément de degré $n$ du membre de gauche. Alors
$$
\alpha = (\alpha^{p, q}) \in
\prod\nolimits_{p + q = n} \Hom_R(K^{-q}, \Hom^p(L^\bullet, M^\bullet))
$$
Chaque $\alpha^{p, q}$ est un élément
$$
\alpha^{p, q} = (\alpha^{r, s, q}) \in
\prod\nolimits_{r + s + q = n} \Hom_R(K^{-q}, \Hom_R(L^{-s}, M^r))
$$
Si nous effectuons les identifications
\begin{equation}
\label{equation-identification}
\Hom_R(K^{-q}, \Hom_R(L^{-s}, M^r)) = \Hom_R(K^{-q} \otimes_R L^{-s}, M^r)
\end{equation}
alors nos règles de signes donnent
\begin{align*}
\text{d}(\alpha^{r, s, q})
& =
\text{d}_{\Hom^\bullet(L^\bullet, M^\bullet)} \circ \alpha^{r, s, q}
- (-1)^n \alpha^{r, s, q} \circ \text{d}_K \\
& =
\text{d}_M \circ \alpha^{r, s, q}
- (-1)^{r + s} \alpha^{r, s, q} \circ \text{d}_L
- (-1)^{r + s + q} \alpha^{r, s, q} \circ \text{d}_K
\end{align*}
D'autre part, si $\beta$ est un élément de degré $n$ du membre de droite,
alors
$$
\beta = (\beta^{r, s, q}) \in \prod\nolimits_{r + s + q = n}
\Hom_R(K^{-q} \otimes_R L^{-s}, M^r)
$$
et notre règle de signes (Homologie, définition
\ref{homology-definition-associated-simple-complex}) donne
\begin{align*}
\text{d}(\beta^{r, s, q})
& =
\text{d}_M \circ \beta^{r, s, q}
- (-1)^n \beta^{r, s, q} \circ
\text{d}_{\text{Tot}(K^\bullet \otimes L^\bullet)} \\
& =
\text{d}_M \circ \beta^{r, s, q}
- (-1)^{r + s + q} \left(
\beta^{r, s, q} \circ \text{d}_K + (-1)^{-q} \beta^{r, s, q} \circ \text{d}_L
\right)
\end{align*}
Nous voyons ainsi que l'application induite par les identifications
(\ref{equation-identification}) est bien un morphisme de complexes.
\end{proof}

\begin{remark}
\label{remark-internal-hom}
Soit $R$ un anneau. La catégorie $\text{Comp}(R)$ des complexes de
$R$-modules est une catégorie monoïdale symétrique dont le produit tensoriel
est donné par $\text{Tot}(- \otimes_R -)$ ; voir le
lemme \ref{lemma-symmetric-monoidal-cat-complexes}.
Étant donnés $L^\bullet$ et $M^\bullet$ dans $\text{Comp}(R)$, un élément
$f \in \Hom^0(L^\bullet, M^\bullet)$ définit un morphisme de complexes
$f : L^\bullet \to M^\bullet$ si et seulement si $\text{d}(f) = 0$.
Le lemme \ref{lemma-compose} nous dit donc également que
$$
\Mor_{\text{Comp}(R)}(K^\bullet, \Hom^\bullet(L^\bullet, M^\bullet))
=
\Mor_{\text{Comp}(R)}(\text{Tot}(K^\bullet \otimes_R L^\bullet), M^\bullet)
$$
fonctoriellement en $K^\bullet, L^\bullet, M^\bullet$ dans
$\text{Comp}(R)$. Cela signifie que $\Hom^\bullet( - , -)$ est un Hom interne
pour la catégorie monoïdale symétrique $\text{Comp}(R)$, comme expliqué dans
Catégories, remarque \ref{categories-remark-internal-hom-monoidal}.
\end{remark}

\begin{lemma}
\label{lemma-composition}
Soit $R$ un anneau. Étant donnés des complexes
$K^\bullet, L^\bullet, M^\bullet$ de $R$-modules, il existe un morphisme
canonique
$$
\text{Tot}\left(
\Hom^\bullet(L^\bullet, M^\bullet) \otimes_R \Hom^\bullet(K^\bullet, L^\bullet)
\right)
\longrightarrow
\Hom^\bullet(K^\bullet, M^\bullet)
$$
de complexes de $R$-modules.
\end{lemma}

\begin{proof}
D'après la discussion de la remarque \ref{remark-internal-hom}, l'existence
d'une telle application canonique résulte de
Catégories, remarque \ref{categories-remark-internal-hom-monoidal}.
Nous en donnons également une construction directe.

\medskip\noindent
Un élément $\alpha$ de degré $n$ du membre de gauche est
$$
\alpha = (\alpha^{p, q}) \in \bigoplus\nolimits_{p + q = n}
\Hom^p(L^\bullet, M^\bullet) \otimes_R \Hom^q(K^\bullet, L^\bullet)
$$
L'élément $\alpha^{p, q}$ est une somme finie
$\alpha^{p, q} = \sum \beta^p_i \otimes \gamma^q_i$
avec
$$
\beta^p_i = (\beta^{r, s}_i)
\in \prod\nolimits_{r + s = p} \Hom_R(L^{-s}, M^r)
$$
et
$$
\gamma^q_i = (\gamma^{u, v}_i)
\in \prod\nolimits_{u + v = q} \Hom_R(K^{-v}, L^u)
$$
L'application envoie $\alpha$ sur $\delta = (\delta^{r, v})$, où
$$
\delta^{r, v} =
\sum\nolimits_{i, s}  \beta^{r, s}_i
\circ \gamma^{-s, v}_i
\in \Hom_R(K^{-v}, M^r)
$$
Pour $r + v = n$ fixé, cette somme est finie, car seul un nombre fini des
$\alpha^{p, q}$ sont non nuls, et donc seuls un nombre fini des
$\beta^p_i$ et $\gamma^q_i$ sont non nuls.
Nos règles de signes donnent
\begin{align*}
\text{d}(\alpha^{p, q})
& =
\text{d}_{\Hom^\bullet(L^\bullet, M^\bullet)}(\alpha^{p, q})
+ (-1)^p \text{d}_{\Hom^\bullet(K^\bullet, L^\bullet)}(\alpha^{p, q}) \\
& =
\sum \Big( \text{d}_M \circ \beta^p_i \circ \gamma^q_i
- (-1)^p \beta^p_i \circ \text{d}_L \circ \gamma^q_i \Big) \\
&
\quad + (-1)^p \sum \Big( \beta^p_i \circ \text{d}_L \circ \gamma^q_i
- (-1)^q \beta^p_i \circ \gamma^q_i \circ \text{d}_K \Big) \\
& =
\sum \Big( \text{d}_M \circ \beta^p_i \circ \gamma^q_i 
-(-1)^n \beta^p_i \circ \gamma^q_i \circ \text{d}_K \Big)
\end{align*}
Il en résulte que la règle $\alpha \mapsto \delta$ est compatible aux
différentielles, ce qui démontre le lemme.
\end{proof}

\begin{lemma}
\label{lemma-diagonal-better}
Soit $R$ un anneau. Étant donnés des complexes
$K^\bullet, L^\bullet, M^\bullet$ de $R$-modules, il existe un morphisme
canonique
$$
\text{Tot}(K^\bullet \otimes_R \Hom^\bullet(M^\bullet, L^\bullet))
\longrightarrow
\Hom^\bullet(M^\bullet, \text{Tot}(K^\bullet \otimes_R L^\bullet))
$$
de complexes de $R$-modules, fonctoriel en les trois complexes.
\end{lemma}

\begin{proof}
D'après la discussion de la remarque \ref{remark-internal-hom}, l'existence
d'une telle application canonique résulte de
Catégories, remarque \ref{categories-remark-internal-hom-monoidal}.
Nous en donnons également une construction directe.

\medskip\noindent
Soit $\alpha$ un élément de degré $n$ du membre de droite. Alors
$$
\alpha = (\alpha^{p, q}) \in \prod\nolimits_{p + q = n}
\Hom_R(M^{-q}, \text{Tot}^p(K^\bullet \otimes_R L^\bullet))
$$
Chaque $\alpha^{p, q}$ est un élément
$$
\alpha^{p, q} = (\alpha^{r, s, q}) \in
\Hom_R(M^{-q}, \bigoplus\nolimits_{r + s + q = n} K^r \otimes_R L^s)
$$
où nous considérons $\alpha^{r, s, q}$ comme une famille d'applications telle
que, pour tout $x \in M^{-q}$, seul un nombre fini des
$\alpha^{r, s, q}(x)$ sont non nuls. Nos règles de signes donnent
\begin{align*}
\text{d}(\alpha^{r, s, q})
& =
\text{d}_{\text{Tot}(K^\bullet \otimes_R L^\bullet)} \circ \alpha^{r, s, q}
- (-1)^n \alpha^{r, s, q} \circ \text{d}_M \\
& =
\text{d}_K \circ \alpha^{r, s, q} + (-1)^r \text{d}_L \circ \alpha^{r, s, q}
- (-1)^n \alpha^{r, s, q} \circ \text{d}_M
\end{align*}
D'autre part, si $\beta$ est un élément de degré $n$ du membre de gauche,
alors
$$
\beta = (\beta^{p, q}) \in
\bigoplus\nolimits_{p + q = n} K^p \otimes_R \Hom^q(M^\bullet, L^\bullet)
$$
et nous pouvons écrire $\beta^{p, q} = \sum \gamma_i^p \otimes \delta_i^q$,
avec $\gamma_i^p \in K^p$ et
$$
\delta_i^q = (\delta_i^{r, s}) \in
\prod\nolimits_{r + s = q} \Hom_R(M^{-s}, L^r)
$$
Nos règles de signes donnent
\begin{align*}
\text{d}(\beta^{p, q})
& =
\text{d}_K(\beta^{p, q}) +
(-1)^p \text{d}_{\Hom^\bullet(M^\bullet, L^\bullet)}(\beta^{p, q}) \\
& =
\sum \text{d}_K(\gamma_i^p) \otimes \delta_i^q +
(-1)^p \sum \gamma_i^p \otimes
(\text{d}_L \circ \delta_i^q - (-1)^q \delta_i^q \circ \text{d}_M)
\end{align*}
Nous envoyons l'élément $\beta$ sur $\alpha$, où
$$
\alpha^{r, s, q} = c^{r, s, q}(\sum \gamma_i^r \otimes \delta_i^{s, q})
$$
où $c^{r, s, q} : K^r \otimes_R \Hom_R(M^{-q}, L^s) \to
\Hom_R(M^{-q}, K^r \otimes_R L^s)$ est l'application canonique.
Pour $\beta$ et $r$ fixés, seul un nombre fini des $\gamma_i^r$ sont non nuls ;
par conséquent, seul un nombre fini des $\alpha^{r, s, q}$ sont non nuls
(pour $r$ fixé). Cette famille d'applications satisfait donc aux conditions
ci-dessus, et l'application est bien définie.
En comparant les signes, nous voyons qu'elle est compatible aux
différentielles.
\end{proof}

\begin{lemma}
\label{lemma-diagonal}
Soit $R$ un anneau. Étant donnés des complexes $K^\bullet, L^\bullet$
de $R$-modules, il existe un morphisme canonique
$$
K^\bullet
\longrightarrow
\Hom^\bullet(L^\bullet, \text{Tot}(K^\bullet \otimes_R L^\bullet))
$$
de complexes de $R$-modules, fonctoriel en les deux complexes.
\end{lemma}

\begin{proof}
D'après la discussion de la remarque \ref{remark-internal-hom}, l'existence
d'une telle application canonique résulte de
Catégories, remarque \ref{categories-remark-internal-hom-monoidal}.
Nous en donnons également une construction directe.

\medskip\noindent
Soit $\alpha$ un élément de degré $n$ du membre de droite. Alors
$$
\alpha = (\alpha^{p, q}) \in \prod\nolimits_{p + q = n}
\Hom_R(L^{-q}, \text{Tot}^p(K^\bullet \otimes_R L^\bullet))
$$
Chaque $\alpha^{p, q}$ est un élément
$$
\alpha^{p, q} = (\alpha^{r, s, q}) \in
\Hom_R(L^{-q}, \bigoplus\nolimits_{r + s + q = n} K^r \otimes_R L^s)
$$
où nous considérons $\alpha^{r, s, q}$ comme une famille d'applications telle
que, pour tout $x \in L^{-q}$, seul un nombre fini des
$\alpha^{r, s, q}(x)$ sont non nuls. Nos règles de signes donnent
\begin{align*}
\text{d}(\alpha^{r, s, q})
& =
\text{d}_{\text{Tot}(K^\bullet \otimes_R L^\bullet)} \circ \alpha^{r, s, q}
- (-1)^n \alpha^{r, s, q} \circ \text{d}_L \\
& =
\text{d}_K \circ \alpha^{r, s, q} + (-1)^r \text{d}_L \circ \alpha^{r, s, q}
- (-1)^n \alpha^{r, s, q} \circ \text{d}_L
\end{align*}
Nous envoyons maintenant un élément $\beta \in K^n$ sur $\alpha$, avec
$\alpha^{n, -q, q} = \beta \otimes \text{id}_{L^{-q}}$
et $\alpha^{r, s, q} = 0$ si $r \not = n$. Il s'agit bien d'un élément
comme ci-dessus, puisque, pour $q$ fixé, un seul
$\alpha^{r, s, q}$ est non nul. La description de la différentielle montre
que cette construction est compatible aux différentielles.
\end{proof}

\begin{lemma}
\label{lemma-evaluate-and-more}
Soit $R$ un anneau. Étant donnés des complexes
$K^\bullet, L^\bullet, M^\bullet$ de $R$-modules, il existe un morphisme
canonique
$$
\text{Tot}(\Hom^\bullet(L^\bullet, M^\bullet) \otimes_R K^\bullet)
\longrightarrow
\Hom^\bullet(\Hom^\bullet(K^\bullet, L^\bullet), M^\bullet)
$$
de complexes de $R$-modules, fonctoriel en les trois complexes.
\end{lemma}

\begin{proof}
D'après la discussion de la remarque \ref{remark-internal-hom}, l'existence
d'une telle application canonique résulte de
Catégories, remarque \ref{categories-remark-internal-hom-monoidal}.
Nous en donnons également une construction directe.

\medskip\noindent
Considérons un élément $\beta$ de degré $n$ du membre de droite. Alors
$$
\beta = (\beta^{p, s}) \in
\prod\nolimits_{p + s = n} \Hom_R(\Hom^{-s}(K^\bullet, L^\bullet), M^p)
$$
Nos règles de signes donnent
\begin{align*}
\text{d}(\beta^{p, s})
& =
\text{d}_M \circ \beta^{p, s}
- (-1)^n \beta^{p, s} \circ
\text{d}_{\Hom^\bullet(K^\bullet, L^\bullet)}
\end{align*}
Nous pouvons décrire le dernier terme comme suit :
$$
(\beta^{p, s} \circ \text{d}_{\Hom^\bullet(K^\bullet, L^\bullet)})(f) =
\beta^{p, s}(\text{d}_L \circ f - (-1)^{s + 1} f \circ \text{d}_K)
$$
si $f \in \Hom^{-s - 1}(K^\bullet, L^\bullet)$. Nous concluons que, dans un
sens que nous ne précisons pas, $\text{d}(\beta^{p, s})$ est une somme de
trois termes dont les signes sont les suivants :
\begin{equation}
\label{equation-beta}
\text{d}(\beta^{p, s}) =
\text{d}_M(\beta^{p, s})
-(-1)^n\text{d}_L(\beta^{p, s}) +
(-1)^{p + 1}\text{d}_K(\beta^{p, s})
\end{equation}

\noindent
Considérons ensuite un élément $\alpha$ de degré $n$ du membre de gauche.
Nous pouvons l'écrire sous la forme
$$
\alpha = (\alpha^{t, r}) \in
\bigoplus\nolimits_{t + r = n} \Hom^t(L^\bullet, M^\bullet) \otimes K^r
$$
Chaque $\alpha^{t, r}$ s'envoie sur un élément
$$
\alpha^{t, r} \mapsto (\alpha^{p, q, r}) \in
\prod\nolimits_{p + q = t} \Hom_R(L^{-q}, M^p) \otimes_R K^r
$$
Nos règles de signes donnent
\begin{align*}
\text{d}(\alpha^{p, q, r})
& =
\text{d}_{\Hom^\bullet(L^\bullet, M^\bullet)}(\alpha^{p, q, r})
+ (-1)^{p + q} \text{d}_K(\alpha^{p, q, r})
\end{align*}
où, si nous écrivons de plus
$\alpha^{p, q, r} = \sum g_i^{p, q} \otimes k_i^r$, nous avons
$$
\text{d}_{\Hom^\bullet(L^\bullet, M^\bullet)}(\alpha^{p, q, r}) =
\sum (\text{d}_M \circ g_i^{p, q}) \otimes k_i^r
- (-1)^{p + q} \sum (g_i^{p, q} \circ \text{d}_L) \otimes k_i^r
$$
Nous concluons que, dans un sens que nous ne précisons pas,
$\text{d}(\alpha^{p, q, r})$ est une somme de trois termes dont les signes
sont les suivants :
\begin{equation}
\label{equation-alpha}
\text{d}(\alpha^{p, q, r}) =
\text{d}_M(\alpha^{p, q, r})
-(-1)^{p + q}\text{d}_L(\alpha^{p, q, r}) +
(-1)^{p + q}\text{d}_K(\alpha^{p, q, r})
\end{equation}

\noindent
Pour définir notre application, nous utiliserons les applications canoniques
$$
c_{p, q, r} : 
\Hom_R(L^{-q}, M^p) \otimes_R K^r
\longrightarrow
\Hom_R(\Hom_R(K^r, L^{-q}), M^p)
$$
qui envoient $\varphi \otimes k$ sur l'application
$\psi \mapsto \varphi(\psi(k))$. Cette construction est fonctorielle en les
trois variables. Avec $s = q + r$, il existe une inclusion
$$
\Hom_R(\Hom_R(K^r, L^{-q}), M^p) \subset
\Hom_R(\Hom^{-s}(K^\bullet, L^\bullet), M^p)
$$
provenant de la projection
$\Hom^{-s}(K^\bullet, L^\bullet) \to \Hom_R(K^r, L^{-q})$.
Puisque $\alpha^{p, q, r}$ n'est non nul que pour un nombre fini de $r$,
nous voyons que, pour $s$ fixé, il n'y a qu'un nombre fini de $q, r$ tels
que $q + r = s$. Nous pouvons donc envoyer $\alpha$ sur l'élément $\beta$
défini par
$$
\beta^{p, s} =
\sum\nolimits_{q + r = s} \epsilon_{p, q, r} c_{p, q, r}(\alpha^{p, q, r})
$$
où la somme utilise les inclusions données ci-dessus et où
$\epsilon_{p, q, r} \in \{\pm 1\}$. En comparant les signes dans les
équations (\ref{equation-beta}) et (\ref{equation-alpha}), nous voyons que
\begin{enumerate}
\item $\epsilon_{p, q, r} = \epsilon_{p + 1, q, r}$,
\item $-(-1)^n\epsilon_{p, q, r} = -(-1)^{p + q}\epsilon_{p, q - 1, r}$
ou, de manière équivalente,
$\epsilon_{p, q, r} = (-1)^r\epsilon_{p, q - 1, r}$,
\item $(-1)^{p + 1}\epsilon_{p, q, r} = (-1)^{p + q}\epsilon_{p, q, r + 1}$
ou, de manière équivalente,
$(-1)^{q + 1}\epsilon_{p, q, r} = \epsilon_{p, q, r + 1}$.
\end{enumerate}
Un bon choix consiste à poser
$$
\epsilon_{p, r, s} = (-1)^{r + qr}
$$
Le choix de ce signe est expliqué dans la remarque qui suit la démonstration.
\end{proof}

\begin{remark}
\label{remark-sign-explanation}
Expliquons pourquoi le signe utilisé dans la construction directe de la
démonstration du lemme \ref{lemma-evaluate-and-more} coïncide avec celui que
donne la construction fondée sur la discussion de la
remarque \ref{remark-internal-hom} et de Catégories,
remarque \ref{categories-remark-internal-hom-monoidal}.
Notons $- \otimes - = \text{Tot}(- \otimes_R -)$ et
$hom(-, -) = \Hom^\bullet(-, -)$. Le langage des catégories monoïdales nous
dit d'utiliser la flèche
$$
hom(L^\bullet, M^\bullet) \otimes K^\bullet
\longrightarrow
hom(hom(K^\bullet, L^\bullet), M^\bullet)
$$
dans $\text{Comp}(R)$ correspondant à la flèche
$$
hom(L^\bullet, M^\bullet) \otimes K^\bullet \otimes hom(K^\bullet, L^\bullet)
\longrightarrow
M^\bullet
$$
obtenue en échangeant l'ordre des deux derniers produits tensoriels, puis en
utilisant les applications d'évaluation
$hom(K^\bullet, L^\bullet) \otimes K^\bullet \to L^\bullet$ et
$hom(L^\bullet, K^\bullet) \otimes L^\bullet \to M^\bullet$.
Un signe n'intervient que lors de l'échange. Plus précisément, dans
l'isomorphisme
$$
K^\bullet \otimes hom(K^\bullet, L^\bullet)
\to
hom(K^\bullet, L^\bullet) \otimes K^\bullet
$$
le signe $(-1)^{r(q + r')}$ intervient sur
$K^r \otimes_R \Hom_R(K^{-r'}, L^q)$ ; voir la
section \ref{section-sign-rules}, point (\ref{item-flip-tensor-product}).
Le lecteur vérifiera qu'en raison de la correspondance que nous utilisons
pour décrire les applications à valeurs dans un Hom interne, ce signe
n'importe que si $r = r'$ ; dans ce cas, nous obtenons
$(-1)^{r(q + r)} = (-1)^{r + qr}$, comme dans la démonstration directe.
\end{remark}







\section{Règles de signes}
\label{section-sign-rules}

\noindent
Dans cette section, nous passons en revue les règles de signes utilisées
jusqu'ici et examinons certaines de leurs conséquences. Il semble également
approprié de traiter ces questions dans la catégorie des complexes de modules
sur un anneau, puisque la plupart des phénomènes intéressants apparaissent
déjà dans ce cas. Nous espérons sincèrement que le lecteur n'aura pas besoin
d'utiliser les aspects les plus ésotériques de cette section.

\medskip\noindent
Dans toute la suite de cette section, fixons un anneau $R$ et notons
$M^\bullet$ un complexe de $R$-modules de différentielles
$d^n_M : M^n \to M^{n + 1}$.
\begin{enumerate}
\item Le $k$-ième complexe décalé $M^\bullet[k]$ a pour termes
$(M^\bullet[k])^n = M^{n + k}$ et pour différentielles
$d_{M[k]}^n = (-1)^kd^{n + k}_M$ ; voir Homologie,
définition \ref{homology-definition-shift-cochain}.
\item Étant donné un morphisme de complexes $f : M^\bullet \to N^\bullet$,
nous définissons $f[k] : M^\bullet[k] \to N^\bullet[k]$ sans faire intervenir
de signe ; voir Homologie, définition
\ref{homology-definition-shift-cochain}.
\item Nous identifions $H^n(M^\bullet[k])$ à $H^{n + k}(M^\bullet)$ sans
faire intervenir de signe ; voir Homologie,
définition \ref{homology-definition-cohomology-shift}.
\item L'application de bord d'une suite exacte courte de complexes est
définie comme dans le lemme du serpent, sans faire intervenir de signe ;
voir Homologie, lemme \ref{homology-lemma-long-exact-sequence-cochain}.
\item Le triangle distingué associé à une suite exacte courte de complexes,
scindée terme à terme,
$0 \to K^\bullet \to L^\bullet \to M^\bullet \to 0$, est donné par
$$
K^\bullet \to L^\bullet \to M^\bullet \to K^\bullet[1]
$$
où $M^n \to K^{n + 1}$ est l'application
$\pi^{n + 1} \circ d^n_L \circ s^n$ si $s$ et $\pi$ sont des scindages
compatibles terme à terme. Autrement dit, aucun signe n'intervient. Voir
Catégories dérivées, définitions
\ref{derived-definition-distinguished-triangle}
et \ref{derived-definition-split-ses}.
\item Le complexe total $\text{Tot}(M^\bullet \otimes_R N^\bullet)$ a une
différentielle $d$ qui satisfait à la règle de Leibniz
$d(x \otimes y) = d(x) \otimes y + (-1)^{\deg(x)}x \otimes d(y)$.
Voir Homologie, exemple
\ref{homology-example-double-complex-as-tensor-product-of} et
Homologie, définition \ref{homology-definition-associated-simple-complex}.
\item
\label{item-shift-tensor}
Il existe un isomorphisme canonique
$$
\text{Tot}(M^\bullet \otimes_R N^\bullet)[a + b] \to
\text{Tot}(M^\bullet[a] \otimes_R N^\bullet[b])
$$
qui utilise le signe $(-1)^{pb}$ sur le facteur $M^p \otimes_R N^q$ ;
voir Homologie, remarque \ref{homology-remark-shift-double-complex}.
Il est souvent plus commode de considérer l'application décalée correspondante
$\text{Tot}(M^\bullet \otimes_R N^\bullet) \to
\text{Tot}(M^\bullet[a] \otimes_R N^\bullet[b])[-a - b]$.
\item Il existe un isomorphisme canonique de complexes
$$
\text{Tot}(
\text{Tot}(K^\bullet \otimes_R L^\bullet) \otimes_R M^\bullet) \to
\text{Tot}(K^\bullet \otimes_R \text{Tot}(L^\bullet \otimes_R M^\bullet))
$$
défini sans faire intervenir de signe. Voir la
section \ref{section-symmetric-monoidal}.
\item
\label{item-flip-tensor-product}
Il existe un isomorphisme canonique
$$
\text{Tot}(L^\bullet \otimes_R M^\bullet) \to
\text{Tot}(M^\bullet \otimes_R L^\bullet)
$$
qui utilise le signe $(-1)^{pq}$ sur le facteur $L^p \otimes_R M^q$.
Voir la section \ref{section-symmetric-monoidal}.
\end{enumerate}
Avant d'aborder les conventions de signes concernant les complexes Hom,
construisons le dual d'un complexe relativement aux conventions précédentes.

\begin{lemma}
\label{lemma-left-dual-module}
Soit $R$ un anneau. Soit $M$ un $R$-module. Soit $N, \eta, \epsilon$
un dual à gauche de $M$ dans la catégorie monoïdale des $R$-modules ; voir
Catégories, définition \ref{categories-definition-dual}. Alors
\begin{enumerate}
\item $M$ et $N$ sont des $R$-modules projectifs de type fini,
\item l'application
$e : \Hom_R(M, R) \to N$, $\lambda \mapsto (\lambda \otimes 1)(\eta)$
est un isomorphisme,
\item on a $\epsilon(n, m) = e^{-1}(n)(m)$ pour $n \in N$ et $m \in M$.
\end{enumerate}
\end{lemma}

\begin{proof}
Les hypothèses signifient que
$$
M \xrightarrow{\eta \otimes 1} M \otimes_R N \otimes_R M
\xrightarrow{1 \otimes \epsilon} M
\quad\text{et}\quad
N \xrightarrow{1 \otimes \eta} N \otimes_R M \otimes_R N
\xrightarrow{\epsilon \otimes 1} N
$$
sont les applications identiques. Nous pouvons choisir un module libre de
type fini $F$, un homomorphisme de $R$-modules $F \to M$ et un relèvement
$\tilde \eta : R \to F \otimes_R N$ de $\eta$. Nous obtenons un diagramme
commutatif
$$
\xymatrix{
M \ar[rr]_-{\eta \otimes 1} \ar[rrd]_-{\tilde \eta \otimes 1} & &
M \otimes_R N \otimes_R M \ar[r]_-{1 \otimes \epsilon} &
M \\
& &
F \otimes_R N \otimes_R M \ar[u] \ar[r]^-{1 \otimes \epsilon} &
F \ar[u]
}
$$
Cela montre que l'identité de $M$ se factorise par un module libre de type
fini ; ainsi $M$ est projectif de type fini. Par symétrie, $N$ est projectif
de type fini. Cela démontre (1). L'assertion (2) résulte de
Catégories, lemme \ref{categories-lemma-left-dual}, et de sa démonstration.
L'assertion (3) résulte de la première égalité de la démonstration.
\end{proof}

\begin{lemma}
\label{lemma-left-dual-complex}
Soit $R$ un anneau. Soit $M^\bullet$ un complexe de $R$-modules.
Soit $N^\bullet, \eta, \epsilon$ un dual à gauche de $M^\bullet$
dans la catégorie monoïdale des complexes de $R$-modules. Alors
\begin{enumerate}
\item $M^\bullet$ et $N^\bullet$ sont bornés,
\item $M^n$ et $N^n$ sont des $R$-modules projectifs de type fini,
\item si l'on écrit $\epsilon = \sum \epsilon_n$ avec
$\epsilon_n : N^{-n} \otimes_R M^n \to R$ et $\eta = \sum \eta_n$ avec
$\eta_n : R \to  M^n \otimes_R N^{-n}$, alors
$(N^{-n}, \eta_n, \epsilon_n)$ est le dual à gauche de $M^n$ comme dans le
lemme \ref{lemma-left-dual-module},
\item la différentielle $d_N^n : N^n \to N^{n + 1}$ est égale à
$-(-1)^n$ fois l'application
$$
N^n = \Hom_R(M^{-n}, R)
\xrightarrow{d_M^{-n - 1}}
\Hom_R(M^{-n - 1}, R) = N^{n + 1}
$$
où les signes d'égalité sont les identifications de la partie (2) du
lemme \ref{lemma-left-dual-module}.
\end{enumerate}
Réciproquement, étant donné un complexe borné $M^\bullet$ de $R$-modules
projectifs de type fini, posons $N^n = \Hom_R(M^{-n}, R)$ avec les
différentielles ci-dessus, $\epsilon = \sum \epsilon_n$, où
$\epsilon_n : N^{-n} \otimes_R M^n \to R$ est donnée par l'évaluation, et
$\eta = \sum \eta_n$, où $\eta_n : R \to M^n \otimes_R N^{-n}$ envoie
$1$ sur $\text{id}_{M_n}$. Nous obtenons alors un dual à gauche de
$M^\bullet$ dans la catégorie monoïdale des complexes de $R$-modules.
\end{lemma}

\begin{proof}
Puisque $(1 \otimes \epsilon) \circ (\eta \otimes 1) =
\text{id}_{M^\bullet}$ et
$(\epsilon \otimes 1) \circ (1 \otimes \eta) =
\text{id}_{N^\bullet}$ d'après Catégories,
définition \ref{categories-definition-dual}, nous voyons immédiatement que
$(1 \otimes \epsilon_n) \circ (\eta_n \otimes 1) = \text{id}_{M^n}$
et $(\epsilon_n \otimes 1) \circ (1 \otimes \eta_n) = \text{id}_{N^{-n}}$,
ce qui démontre (3). Le lemme \ref{lemma-left-dual-module} donne (2).
Comme la somme $\eta = \sum \eta_n$ est finie, nous obtenons (1).
Puisque $\eta = \sum \eta_n$ est un morphisme de complexes
$R \to \text{Tot}(M^\bullet \otimes_R N^\bullet)$, nous avons
$$
(d_M^{-n - 1} \otimes 1) \circ \eta_{-n - 1}  +
(-1)^n (1 \otimes d_N^{-n}) \circ \eta_{-n} = 0
$$
d'après notre choix de signes pour la différentielle de
$\text{Tot}(M^\bullet \otimes_R N^\bullet)$.
En développant les définitions, cela démontre (4). Pour obtenir la dernière
assertion du lemme, il suffit de lire ce qui précède à l'envers.
\end{proof}

\noindent
Nous utiliserons la description du dual à gauche d'un complexe donnée dans le
lemme \ref{lemma-left-dual-complex} pour motiver notre règle de signes sur le
complexe $\Hom$. Plus précisément, nous choisissons les signes de sorte que
(\ref{item-compatible}) soit vraie.
Poursuivons la discussion des diverses règles de signes ci-dessus.
\begin{enumerate}
\setcounter{enumi}{9}
\item Étant donnés des complexes $K^\bullet$, $M^\bullet$, nous définissons
$\Hom^\bullet(M^\bullet, K^\bullet)$ comme le complexe de termes
$$
\Hom^n(M^\bullet, K^\bullet) = \prod\nolimits_{n = p + q} \Hom_R(M^{-q}, K^p)
$$
et de différentielle donnée par la règle
$$
d(f) = d_K \circ f - (-1)^n f \circ d_M
$$
\item
\label{item-compatible}
Le choix précédent est tel que, si $M^\bullet$ possède un dual à gauche
$N^\bullet$ comme dans le lemme \ref{lemma-left-dual-complex}, alors nous
avons un isomorphisme canonique
$$
\text{Tot}(K^\bullet \otimes_R N^\bullet)
\longrightarrow
\Hom^\bullet(M^\bullet, K^\bullet)
$$
défini sans faire intervenir de signe, qui envoie le facteur
$K^p \otimes_R N^q$ sur le facteur $\Hom_R(M^{-q}, K^p)$ au moyen de
$N^q = \Hom_R(M^{-q}, R)$ et de l'application canonique
$K^p \otimes_R \Hom_R(M^{-q}, R) \to \Hom_R(M^{-q}, K^p)$.
\item Il existe une composition
$$
\text{Tot}(
\Hom^\bullet(L^\bullet, K^\bullet) \otimes_R
\Hom^\bullet(M^\bullet, L^\bullet))
\longrightarrow
\Hom^\bullet(M^\bullet, K^\bullet)
$$
définie sans faire intervenir de signe ; voir le
lemme \ref{lemma-composition}.
\item Il existe un isomorphisme canonique
$$
\Hom^\bullet(K^\bullet, \Hom^\bullet(L^\bullet, M^\bullet))
=
\Hom^\bullet(\text{Tot}(K^\bullet \otimes_R L^\bullet), M^\bullet)
$$
défini sans faire intervenir de signe ; voir le
lemme \ref{lemma-compose}.
\item Il existe une application canonique
$$
\text{Tot}(K^\bullet \otimes_R \Hom^\bullet(M^\bullet, L^\bullet))
\longrightarrow
\Hom^\bullet(M^\bullet, \text{Tot}(K^\bullet \otimes_R L^\bullet))
$$
définie sans faire intervenir de signe ; voir le
lemme \ref{lemma-diagonal-better}.
\item Il existe une application canonique
$$
K^\bullet
\longrightarrow
\Hom^\bullet(L^\bullet, \text{Tot}(K^\bullet \otimes_R L^\bullet))
$$
définie sans faire intervenir de signe ; voir le
lemme \ref{lemma-diagonal}.
\item D'après le lemme \ref{lemma-evaluate-and-more}, il existe une
application canonique
$$
\text{Tot}(\Hom^\bullet(L^\bullet, M^\bullet) \otimes_R K^\bullet)
\longrightarrow
\Hom^\bullet(\Hom^\bullet(K^\bullet, L^\bullet), M^\bullet)
$$
qui utilise le signe $(-1)^{r + qr}$ sur le module
$\Hom_R(L^{-q}, M^p) \otimes_R K^r$ ; la raison en est expliquée dans la
remarque \ref{remark-sign-explanation}.
\item
\label{item-evaluation}
En prenant $L^\bullet = M^\bullet$ et en utilisant
$R \to \Hom^\bullet(M^\bullet, M^\bullet)$, l'application du point précédent
devient l'application d'évaluation
$$
ev : K^\bullet \longrightarrow
\Hom^\bullet(\Hom^\bullet(K^\bullet, M^\bullet), M^\bullet)
$$
Elle envoie $x \in K^n$ sur l'application qui envoie
$f \in \Hom^m(K^\bullet, M^\bullet)$ sur $(-1)^{nm}f(x)$.
\item
\label{item-shift-hom}
Il existe une identification canonique
$$
\Hom^\bullet(M^\bullet, K^\bullet)[a - b] \to
\Hom^\bullet(M^\bullet[b], K^\bullet[a])
$$
qui utilise des signes. Elle est définie comme l'application dont
l'application décalée correspondante
$$
\Hom^\bullet(M^\bullet, K^\bullet) \to
\Hom^\bullet(M^\bullet[b], K^\bullet[a])[b - a]
$$
utilise le signe $(-1)^{nb}$ sur le module $\Hom_R(M^{-q}, K^p)$, avec
$p + q = n$. En effet, si $f \in \Hom^n(M^\bullet, K^\bullet)$, alors
$$
d(f) = d_K \circ f - (-1)^n f \circ d_M
$$
à la source, tandis qu'au but $f$ appartient à
$\left(\Hom^\bullet(M^\bullet[b], K^\bullet[a])[b - a]\right)^n =
\Hom^{n + b -a}(M^\bullet[b], K^\bullet[a])$ ; nous obtenons donc
\begin{align*}
d(f) & =
(-1)^{b - a}
\left(d_{K[a]} \circ f - (-1)^{n + b - a} f \circ d_{M[b]}\right) \\
& =
(-1)^{b - a}
\left((-1)^a d_K \circ f - (-1)^{n + b - a} f \circ (-1)^b d_M \right) \\
& =
(-1)^b d_K \circ f - (-1)^{n + b} f \circ d_M
\end{align*}
et l'on voit que le signe choisi $(-1)^{nb}$ au degré $n$ fournit un
morphisme de complexes pour ces différentielles.
\end{enumerate}








\section{Hom dérivé}
\label{section-RHom}

\noindent
Soit $R$ un anneau. Le Hom dérivé que nous allons définir dans cette section
est un foncteur
$$
D(R)^{opp} \times D(R) \longrightarrow D(R),\quad
(K, L) \longmapsto R\Hom_R(K, L)
$$
Il s'agit d'un Hom interne dans la catégorie dérivée des
$R$-modules, au sens où il est caractérisé par la formule
\begin{equation}
\label{equation-internal-hom}
\Hom_{D(R)}(K, R\Hom_R(L, M))
=
\Hom_{D(R)}(K \otimes_R^\mathbf{L} L, M)
\end{equation}
pour des objets $K, L, M$ de $D(R)$. Remarquons que, par le lemme de Yoneda,
cette formule caractérise les objets à isomorphisme unique près. On peut
donner une construction comme suit. Choisissons un complexe K-injectif
$I^\bullet$ de $R$-modules représentant $M$, choisissons un complexe
$L^\bullet$ représentant $L$, et posons
$$
R\Hom_R(L, M) = \Hom^\bullet(L^\bullet, I^\bullet)
$$
avec les notations de la section \ref{section-hom-complexes}.
Une généralisation de cette construction est étudiée dans
Algèbre différentielle graduée, section \ref{dga-section-restriction}.
Il résulte de (\ref{equation-cohomology-hom-complex}) et de
Catégories dérivées, lemme \ref{derived-lemma-K-injective}, que
\begin{equation}
\label{equation-h0-RHom}
H^n(R\Hom_R(L, M)) = \Hom_{D(R)}(L, M[n])
\end{equation}
pour tout $n \in \mathbf{Z}$. En particulier, l'objet
$R\Hom_R(L, M)$ de $D(R)$ est bien défini, c'est-à-dire indépendant du
choix du complexe K-injectif $I^\bullet$.

\begin{lemma}
\label{lemma-internal-hom}
Soit $R$ un anneau. Soient $K, L, M$ des objets de $D(R)$. Il existe un
isomorphisme canonique
$$
R\Hom_R(K, R\Hom_R(L, M)) = R\Hom_R(K \otimes_R^\mathbf{L} L, M)
$$
dans $D(R)$, fonctoriel en $K, L, M$, qui redonne
(\ref{equation-internal-hom}) en prenant $H^0$.
\end{lemma}

\begin{proof}
Choisissons un complexe K-injectif $I^\bullet$ représentant
$M$ et un complexe K-plat de $R$-modules $L^\bullet$
représentant $L$. Pour tout complexe de $R$-modules $K^\bullet$,
nous avons
$$
\Hom^\bullet(K^\bullet, \Hom^\bullet(L^\bullet, I^\bullet)) =
\Hom^\bullet(\text{Tot}(K^\bullet \otimes_R L^\bullet), I^\bullet)
$$
par le lemme \ref{lemma-compose}. Le lemme résulte de la définition
de $R\Hom$ et du fait que
$\text{Tot}(K^\bullet \otimes_R L^\bullet)$
représente le produit tensoriel dérivé.
\end{proof}

\begin{lemma}
\label{lemma-RHom-out-of-projective}
Soit $R$ un anneau. Soit $P^\bullet$ un complexe borné supérieurement
de $R$-modules projectifs. Soit $L^\bullet$ un complexe de $R$-modules.
Alors $R\Hom_R(P^\bullet, L^\bullet)$ est représenté par le complexe
$\Hom^\bullet(P^\bullet, L^\bullet)$.
\end{lemma}

\begin{proof}
D'après (\ref{equation-cohomology-hom-complex}) et
Catégories dérivées, lemme
\ref{derived-lemma-morphisms-from-projective-complex},
les groupes de cohomologie du complexe sont « corrects ».
Ainsi, si nous choisissons un quasi-isomorphisme
$L^\bullet \to I^\bullet$, où $I^\bullet$ est un complexe K-injectif de
$R$-modules, alors l'application induite
$$
\Hom^\bullet(P^\bullet, L^\bullet)
\longrightarrow
\Hom^\bullet(P^\bullet, I^\bullet)
$$
est un quasi-isomorphisme. Comme le membre de droite est, par définition,
$R\Hom_R(P^\bullet, L^\bullet)$, nous avons terminé.
\end{proof}

\begin{lemma}
\label{lemma-internal-hom-evaluate}
Soit $R$ un anneau. Soient $K, L, M$ des objets de $D(R)$.
Il existe un morphisme canonique
$$
R\Hom_R(L, M) \otimes_R^\mathbf{L} K
\longrightarrow
R\Hom_R(R\Hom_R(K, L), M)
$$
dans $D(R)$, fonctoriel en $K, L, M$.
\end{lemma}

\begin{proof}
Choisissons
un complexe K-injectif $I^\bullet$ représentant $M$,
un complexe K-injectif $J^\bullet$ représentant $L$, et
un complexe K-plat $K^\bullet$ représentant $K$.
Le morphisme est défini à l'aide du morphisme
$$
\text{Tot}(\Hom^\bullet(J^\bullet, I^\bullet) \otimes_R K^\bullet)
\longrightarrow
\Hom^\bullet(\Hom^\bullet(K^\bullet, J^\bullet), I^\bullet)
$$
du lemme \ref{lemma-evaluate-and-more}. Nous omettons la démonstration
de sa fonctorialité en les trois objets de $D(R)$.
\end{proof}

\begin{lemma}
\label{lemma-internal-hom-composition}
Soit $R$ un anneau. Étant donnés $K, L, M$ dans $D(R)$, il existe un
morphisme canonique
$$
R\Hom_R(L, M) \otimes_R^\mathbf{L} R\Hom_R(K, L) \longrightarrow R\Hom_R(K, M)
$$
dans $D(R)$, fonctoriel en $K, L, M$.
\end{lemma}

\begin{proof}
Choisissons un complexe K-injectif $I^\bullet$ représentant $M$,
un complexe K-injectif $J^\bullet$ représentant $L$, et
un complexe quelconque de $R$-modules $K^\bullet$ représentant $K$.
Par le lemme \ref{lemma-composition}, il existe un morphisme de complexes
$$
\text{Tot}\left(
\Hom^\bullet(J^\bullet, I^\bullet) \otimes_R \Hom^\bullet(K^\bullet, J^\bullet)
\right)
\longrightarrow
\Hom^\bullet(K^\bullet, I^\bullet)
$$
Les complexes de $R$-modules $\Hom^\bullet(J^\bullet, I^\bullet)$,
$\Hom^\bullet(K^\bullet, J^\bullet)$ et $\Hom^\bullet(K^\bullet, I^\bullet)$
représentent $R\Hom_R(L, M)$, $R\Hom_R(K, L)$ et $R\Hom_R(K, M)$.
Si nous choisissons un complexe K-plat $H^\bullet$ et un quasi-isomorphisme
$H^\bullet \to \Hom^\bullet(K^\bullet, J^\bullet)$, alors il existe un morphisme
$$
\text{Tot}\left(
\Hom^\bullet(J^\bullet, I^\bullet) \otimes_R H^\bullet
\right)
\longrightarrow
\text{Tot}\left(
\Hom^\bullet(J^\bullet, I^\bullet) \otimes_R \Hom^\bullet(K^\bullet, J^\bullet)
\right)
$$
dont la source représente
$R\Hom_R(L, M) \otimes_R^\mathbf{L} R\Hom_R(K, L)$.
La composée des deux flèches affichées donne le morphisme souhaité.
Nous omettons la démonstration de la fonctorialité de la construction.
\end{proof}

\begin{lemma}
\label{lemma-internal-hom-diagonal-better}
Soit $R$ un anneau. Étant donnés des complexes $K, L, M$ dans $D(R)$,
il existe un morphisme canonique
$$
K \otimes_R^\mathbf{L} R\Hom_R(M, L)
\longrightarrow
R\Hom_R(M, K \otimes_R^\mathbf{L} L)
$$
dans $D(R)$, fonctoriel en $K$, $L$, $M$.
\end{lemma}

\begin{proof}
Choisissons un complexe K-plat $K^\bullet$ représentant $K$,
un complexe K-injectif $I^\bullet$ représentant $L$, et
un complexe quelconque $M^\bullet$ représentant $M$.
Choisissons un quasi-isomorphisme
$\text{Tot}(K^\bullet \otimes_R I^\bullet) \to J^\bullet$
où $J^\bullet$ est K-injectif. Nous utilisons alors le morphisme
$$
\text{Tot}\left(
K^\bullet \otimes_R \Hom^\bullet(M^\bullet, I^\bullet)
\right)
\to
\Hom^\bullet(M^\bullet, \text{Tot}(K^\bullet \otimes_R I^\bullet))
\to
\Hom^\bullet(M^\bullet, J^\bullet)
$$
où le premier morphisme est celui du lemme \ref{lemma-diagonal-better}.
\end{proof}

\begin{lemma}
\label{lemma-internal-hom-diagonal}
Soit $R$ un anneau. Étant donnés des complexes $K, L$ dans $D(R)$,
il existe un morphisme canonique
$$
K \longrightarrow R\Hom_R(L, K \otimes_R^\mathbf{L} L)
$$
dans $D(R)$, fonctoriel en $K$ et en $L$.
\end{lemma}

\begin{proof}
Il s'agit d'un cas particulier du lemme
\ref{lemma-internal-hom-diagonal-better}, mais nous allons aussi le démontrer
directement. Choisissons un complexe K-plat $K^\bullet$ représentant $K$ et
un complexe quelconque $L^\bullet$ représentant $L$. Choisissons un
quasi-isomorphisme
$\text{Tot}(K^\bullet \otimes_R L^\bullet) \to J^\bullet$
où $J^\bullet$ est K-injectif. Nous utilisons alors le morphisme
$$
K^\bullet \to
\Hom^\bullet(L^\bullet, \text{Tot}(K^\bullet \otimes_R L^\bullet))
\to \Hom^\bullet(L^\bullet, J^\bullet)
$$
où le premier morphisme est celui du lemme \ref{lemma-diagonal}.
\end{proof}









\section{Complexes parfaits}
\label{section-perfect}

\noindent
Un complexe parfait est un complexe pseudo-cohérent de dimension de Tor finie.
Nous n'utiliserons pas ceci comme définition, mais définirons directement
les complexes parfaits sur un anneau comme suit.

\begin{definition}
\label{definition-perfect}
Soit $R$ un anneau. Notons $D(R)$ la catégorie dérivée de la catégorie
abélienne des $R$-modules.
\begin{enumerate}
\item Un objet $K$ de $D(R)$ est {\it parfait} s'il est quasi-isomorphe
à un complexe borné de $R$-modules projectifs finis.
\item Un $R$-module $M$ est {\it parfait} si $M[0]$ est un objet parfait
de $D(R)$.
\end{enumerate}
\end{definition}

\noindent
Par exemple, sur un anneau noethérien, un module fini est parfait si et
seulement s'il est de dimension projective finie, voir le
lemme \ref{lemma-perfect-module} et Algèbre, définition
\ref{algebra-definition-finite-proj-dim}.

\begin{lemma}
\label{lemma-perfect}
Soit $K^\bullet$ un objet de $D(R)$. Les assertions suivantes sont équivalentes :
\begin{enumerate}
\item $K^\bullet$ est parfait, et
\item $K^\bullet$ est pseudo-cohérent et de dimension de Tor finie.
\end{enumerate}
Si (1) et (2) sont vérifiées et si $K^\bullet$ est d'amplitude de Tor
dans $[a, b]$, alors $K^\bullet$ est quasi-isomorphe à un complexe
$E^\bullet$ de $R$-modules projectifs finis tel que $E^i = 0$
pour $i \not \in [a, b]$.
\end{lemma}

\begin{proof}
Il est clair que (1) implique (2), voir les
lemmes \ref{lemma-pseudo-coherent} et \ref{lemma-tor-amplitude}.
Supposons (2) vérifiée et $K^\bullet$ d'amplitude de Tor dans $[a, b]$.
En particulier, $H^i(K^\bullet) = 0$ pour $i > b$.
Choisissons un complexe $F^\bullet$ de $R$-modules libres finis tel que
$F^i = 0$ pour $i > b$, ainsi qu'un quasi-isomorphisme
$F^\bullet \to K^\bullet$ (lemme \ref{lemma-pseudo-coherent}).
Posons $E^\bullet = \tau_{\geq a}F^\bullet$. Remarquons que $E^i$ est
libre fini sauf $E^a$, qui est un $R$-module de présentation finie. Par le
lemme \ref{lemma-last-one-flat}, $E^a$ est plat. Ainsi, par
Algèbre, lemme \ref{algebra-lemma-finite-projective},
$E^a$ est projectif fini.
\end{proof}

\begin{lemma}
\label{lemma-perfect-module}
Soit $M$ un module sur un anneau $R$. Les assertions suivantes sont
équivalentes :
\begin{enumerate}
\item $M$ est un module parfait, et
\item il existe une résolution
$$
0 \to F_d \to \ldots \to F_1 \to F_0 \to M \to 0
$$
où chaque $F_i$ est un $R$-module projectif fini.
\end{enumerate}
\end{lemma}

\begin{proof}
Supposons (2). Alors le complexe $E^\bullet$ défini par $E^{-i} = F_i$
est quasi-isomorphe à $M[0]$. Ainsi $M$ est parfait.
Réciproquement, supposons (1). D'après les
lemmes \ref{lemma-perfect} et \ref{lemma-n-pseudo-module},
nous pouvons trouver une résolution $E^\bullet \to M$ où $E^{-i}$ est un
$R$-module libre fini. Par le lemme \ref{lemma-last-one-flat},
$F_d = \Coker(E^{d - 1} \to E^d)$ est plat pour un certain
$d$ suffisamment grand. Par Algèbre, lemme
\ref{algebra-lemma-finite-projective}, $F_d$ est projectif fini.
Par conséquent,
$$
0 \to F_d \to E^{-d+1} \to \ldots \to E^0 \to M \to 0
$$
est la résolution souhaitée.
\end{proof}

\begin{lemma}
\label{lemma-two-out-of-three-perfect}
Soit $R$ un anneau. Soit $(K^\bullet, L^\bullet, M^\bullet, f, g, h)$
un triangle distingué dans $D(R)$. Si deux des trois objets
$K^\bullet, L^\bullet, M^\bullet$ sont parfaits, alors le troisième
est également parfait.
\end{lemma}

\begin{proof}
Combiner les lemmes \ref{lemma-perfect},
\ref{lemma-two-out-of-three-pseudo-coherent} et
\ref{lemma-cone-tor-amplitude}.
\end{proof}

\begin{lemma}
\label{lemma-summands-perfect}
Soit $R$ un anneau. Si $K^\bullet \oplus L^\bullet$ est parfait, alors
$K^\bullet$ et $L^\bullet$ le sont aussi.
\end{lemma}

\begin{proof}
Cela résulte des lemmes \ref{lemma-perfect},
\ref{lemma-summands-pseudo-coherent} et
\ref{lemma-summands-tor-amplitude}.
\end{proof}

\begin{lemma}
\label{lemma-complex-perfect-modules}
Soit $R$ un anneau. Soit $K^\bullet$ un complexe borné de
$R$-modules parfaits. Alors $K^\bullet$ est un complexe parfait.
\end{lemma}

\begin{proof}
Cela résulte par récurrence sur la longueur du complexe fini : utiliser le
lemme \ref{lemma-two-out-of-three-perfect} et les troncations brutales.
\end{proof}

\begin{lemma}
\label{lemma-cohomology-perfect}
Soit $R$ un anneau. Si $K^\bullet \in D^b(R)$ et si tous ses modules de
cohomologie sont parfaits, alors $K^\bullet$ est parfait.
\end{lemma}

\begin{proof}
Cela résulte par récurrence sur la longueur du complexe fini : utiliser le
lemme \ref{lemma-two-out-of-three-perfect} et les troncations canoniques.
\end{proof}

\begin{lemma}
\label{lemma-perfect-push-perfect}
Soit $A \to B$ un homomorphisme d'anneaux. Supposons que $B$ soit parfait
comme $A$-module. Soit $K^\bullet$ un complexe parfait de $B$-modules.
Alors $K^\bullet$ est parfait comme complexe de $A$-modules.
\end{lemma}

\begin{proof}
À l'aide du lemme \ref{lemma-perfect}, cela se ramène aux résultats
correspondants pour les modules pseudo-cohérents et les modules de dimension
de Tor finie. Voir les lemmes
\ref{lemma-finite-tor-dimension-push-tor-amplitude} et
\ref{lemma-finite-push-pseudo-coherent}.
\end{proof}

\begin{lemma}
\label{lemma-pull-perfect}
Soit $A \to B$ un homomorphisme d'anneaux.
Soit $K^\bullet$ un complexe parfait de $A$-modules. Alors
$K^\bullet \otimes_A^{\mathbf{L}} B$ est un complexe parfait de $B$-modules.
\end{lemma}

\begin{proof}
À l'aide du lemme \ref{lemma-perfect}, cela se ramène aux résultats
correspondants pour les modules pseudo-cohérents et les modules de dimension
de Tor finie. Voir les lemmes \ref{lemma-pull-tor-amplitude} et
\ref{lemma-pull-pseudo-coherent}.
\end{proof}

\begin{lemma}
\label{lemma-flat-base-change-perfect}
Soit $A \to B$ un homomorphisme plat d'anneaux. Soit $M$ un
$A$-module parfait. Alors $M \otimes_A B$ est un $B$-module parfait.
\end{lemma}

\begin{proof}
Par le lemme \ref{lemma-perfect-module}, l'hypothèse implique que $M$
possède une résolution finie $F_\bullet$ par des $R$-modules projectifs finis.
Comme $A \to B$ est plat, le complexe $F_\bullet \otimes_A B$ est une
résolution de longueur finie de $M \otimes_A B$ par des modules projectifs
finis sur $B$. Ainsi $M \otimes_A B$ est parfait.
\end{proof}

\begin{lemma}
\label{lemma-tensor-perfect}
Soit $R$ un anneau. Si $K$ et $L$ sont des objets parfaits de $D(R)$,
alors $K \otimes_R^\mathbf{L} L$ est également un objet parfait.
\end{lemma}

\begin{proof}
Nous pouvons le démontrer à partir de la définition comme suit. Nous pouvons
représenter $K$, resp.\ $L$, par un complexe borné $K^\bullet$, resp.\
$L^\bullet$, de $R$-modules projectifs finis. Alors
$K \otimes_R^\mathbf{L} L$ est représenté par le complexe borné
$\text{Tot}(K^\bullet \otimes_R L^\bullet)$.
Les termes de ce complexe sont des sommes directes des modules
$M^a \otimes_R L^b$. Comme $M^a$ et $L^b$ sont des facteurs directs
de $R$-modules libres finis, il en va de même de $M^a \otimes_R L^b$.
Nous en concluons donc que les termes du complexe
$\text{Tot}(K^\bullet \otimes_R L^\bullet)$ sont projectifs finis.

\medskip\noindent
Une autre démonstration s'obtient en utilisant la caractérisation des
complexes parfaits du lemme \ref{lemma-perfect}, ainsi que les lemmes
correspondants pour les complexes pseudo-cohérents (lemme
\ref{lemma-tensor-pseudo-coherent}) et pour l'amplitude de Tor (lemme
\ref{lemma-push-tor-amplitude}, appliqué avec $A = B = R$).
\end{proof}

\begin{lemma}
\label{lemma-glue-perfect}
Soit $R$ un anneau. Soient $f_1, \ldots, f_r \in R$ des éléments qui
engendrent l'idéal unité. Soit $K^\bullet$ un complexe de $R$-modules.
Si, pour chaque $i$, le complexe $K^\bullet \otimes_R R_{f_i}$ est parfait,
alors $K^\bullet$ est parfait.
\end{lemma}

\begin{proof}
À l'aide du lemme \ref{lemma-perfect}, cela se ramène aux résultats
correspondants pour les modules pseudo-cohérents et les modules de dimension
de Tor finie. Voir les lemmes \ref{lemma-glue-tor-amplitude} et
\ref{lemma-glue-pseudo-coherent}.
\end{proof}

\begin{lemma}
\label{lemma-flat-descent-perfect}
Soit $R$ un anneau. Soit $K^\bullet$ un complexe de $R$-modules.
Soit $R \to R'$ un homomorphisme fidèlement plat d'anneaux. Si le complexe
$K^\bullet \otimes_R R'$ est parfait, alors $K^\bullet$ est parfait.
\end{lemma}

\begin{proof}
À l'aide du lemme \ref{lemma-perfect}, cela se ramène aux résultats
correspondants pour les modules pseudo-cohérents et les modules de dimension
de Tor finie. Voir les lemmes \ref{lemma-flat-descent-tor-amplitude} et
\ref{lemma-flat-descent-pseudo-coherent}.
\end{proof}

\begin{lemma}
\label{lemma-regular-perfect}
Soit $R$ un anneau régulier. Alors
\begin{enumerate}
\item un $R$-module est parfait si et seulement s'il est un $R$-module fini, et
\item un complexe de $R$-modules $K^\bullet$ est parfait si et seulement si
$K^\bullet \in D^b(R)$ et chaque $H^i(K^\bullet)$ est un $R$-module fini.
\end{enumerate}
\end{lemma}

\begin{proof}
Tout $R$-module parfait est fini par définition. Réciproquement, soit $M$
un $R$-module fini. Choisissons une résolution
$$
\ldots \to F_2 \xrightarrow{d_2} F_1 \xrightarrow{d_1} F_0 \to M \to 0
$$
où les $F_i$ sont des $R$-modules libres finis
(Algèbre, lemme \ref{algebra-lemma-resolution-by-finite-free}).
Posons $M_i = \Ker(d_i)$. Notons $U_i \subset \Spec(R)$ l'ensemble des
idéaux premiers $\mathfrak p$ tels que $M_{i, \mathfrak p}$ soit libre ;
$U_i$ est ouvert d'après Algèbre, lemme
\ref{algebra-lemma-finitely-presented-localization-free}.
Nous avons une suite exacte
$0 \to M_{i + 1} \to F_{i + 1} \to M_i \to 0$.
Si $\mathfrak p \in U_i$, alors
$0 \to M_{i + 1, \mathfrak p} \to F_{i + 1, \mathfrak p} \to
M_{i, \mathfrak p} \to 0$ est scindée. Ainsi $M_{i + 1, \mathfrak p}$
est projectif fini, donc libre 
(Algèbre, lemme \ref{algebra-lemma-finite-projective}).
Ceci montre que $U_i \subset U_{i + 1}$.
Affirmons que $\Spec(R) = \bigcup U_i$. En effet, pour tout idéal premier
$\mathfrak p$, l'anneau local régulier $R_\mathfrak p$ est de dimension
globale finie d'après Algèbre, proposition
\ref{algebra-proposition-regular-finite-gl-dim}.
Il s'ensuit que $M_{i, \mathfrak p}$ est projectif fini (donc libre) pour
$i \gg 0$, par exemple d'après Algèbre, lemme
\ref{algebra-lemma-independent-resolution}.
Comme le spectre de $R$ est noethérien
(Algèbre, lemme \ref{algebra-lemma-Noetherian-topology}),
nous concluons que $U_n = \Spec(R)$ pour un certain $n$. Alors $M_n$ est
un $R$-module projectif d'après Algèbre, lemme
\ref{algebra-lemma-finite-projective}. Ainsi
$$
0 \to M_n \to F_n \to \ldots \to F_1 \to M \to 0
$$
est une résolution bornée par des modules projectifs finis, et donc $M$
est parfait. Ceci démontre la partie (1).

\medskip\noindent
Soit $K^\bullet$ un complexe de $R$-modules.
Si $K^\bullet$ est parfait, alors il appartient à $D^b(R)$ et il est
quasi-isomorphe à un complexe fini de $R$-modules projectifs finis ;
chaque $H^i(K^\bullet)$ est donc certainement un $R$-module fini (car $R$
est noethérien). Réciproquement, supposons que $K^\bullet$ appartienne à
$D^b(R)$ et que chaque $H^i(K^\bullet)$ soit un $R$-module fini. Alors,
d'après (1), chaque $H^i(K^\bullet)$ est un $R$-module parfait, d'où
$K^\bullet$ est parfait par le lemme \ref{lemma-cohomology-perfect}.
\end{proof}

\begin{lemma}
\label{lemma-dual-perfect-complex}
Soit $A$ un anneau. Soit $K \in D(A)$ parfait. Alors
$K^\vee = R\Hom_A(K, A)$ est un complexe parfait et
$K \cong (K^\vee)^\vee$. Il existe des isomorphismes fonctoriels
$$
L \otimes_A^\mathbf{L} K^\vee = R\Hom_A(K, L)
\quad\text{et}\quad
H^0(L \otimes_A^\mathbf{L} K^\vee) = \Ext_A^0(K, L)
$$
pour $L \in D(A)$.
\end{lemma}

\begin{proof}
Nous pouvons représenter $K$ par un complexe borné $K^\bullet$ de
$A$-modules projectifs finis. Par le lemme
\ref{lemma-RHom-out-of-projective}, l'objet $K^\vee$ est représenté par le
complexe $E^\bullet = \Hom^\bullet(K^\bullet, A)$. Remarquons que
$E^n = \Hom_A(K^{-n}, A)$ est lui aussi projectif fini. Ainsi $E^\bullet$
est un complexe borné de modules projectifs finis et nous concluons que
$K^\vee$ est parfait. Il existe un morphisme canonique
$$
K^\bullet = \text{Tot}(\Hom^\bullet(A, A) \otimes_A K^\bullet)
\longrightarrow
\Hom^\bullet(\Hom^\bullet(K^\bullet, A), A)
$$
qui utilise, au signe près, le morphisme d'évaluation en chaque degré,
voir le lemme \ref{lemma-evaluate-and-more}. (Pour les règles de signes,
voir la section \ref{section-sign-rules}.)
Ainsi ce morphisme définit un isomorphisme canonique
$(K^\vee)^\vee \cong K$, puisque le bidual d'un module projectif fini
s'identifie à ce module.

\medskip\noindent
La seconde égalité résulte de la première par le
lemme \ref{lemma-internal-hom} et Catégories dérivées, lemme
\ref{derived-lemma-morphisms-from-projective-complex}, ainsi que par la
définition des groupes Ext, voir Catégories dérivées, section
\ref{derived-section-ext}. Soit $L^\bullet$ un complexe de $A$-modules
représentant $L$. D'après la section \ref{section-sign-rules},
point (\ref{item-compatible}), il existe un isomorphisme canonique
$$
\text{Tot}(L^\bullet \otimes_A E^\bullet)
\longrightarrow
\Hom^\bullet(K^\bullet, L^\bullet)
$$
de complexes de $A$-modules (on utilise ici le lemme
\ref{lemma-left-dual-complex} pour voir que $E^\bullet$ est le dual à gauche
de $K^\bullet$ dans la catégorie des complexes). Ceci démontre la première
égalité affichée et achève la démonstration.
\end{proof}

\begin{remark}
\label{remark-dual-perfect-complex}
Soit $A$ un anneau. Soit $K \in D(A)$ parfait. Si $K$ est d'amplitude de Tor
dans $[a, b]$, alors son dual $K^\vee$ est d'amplitude de Tor dans
$[-b, -a]$. Ceci résulte de la dernière assertion du
lemme \ref{lemma-perfect} et de la description de $K^\vee$ dans la
démonstration du lemme \ref{lemma-dual-perfect-complex}.
\end{remark}

\begin{lemma}
\label{lemma-colim-and-lim-of-duals}
\begin{slogan}
Dualité triviale pour les systèmes d'objets parfaits.
\end{slogan}
Soit $A$ un anneau. Soit $(K_n)_{n \in \mathbf{N}}$ un système d'objets
parfaits de $D(A)$. Soit $K = \text{hocolim} K_n$ la colimite dérivée
(Catégories dérivées, définition \ref{derived-definition-derived-colimit}).
Alors, pour tout objet $E$ de $D(A)$, nous avons
$$
R\Hom_A(K, E) = R\lim E \otimes^\mathbf{L}_A K_n^\vee
$$
où $(K_n^\vee)$ est le système inverse des complexes parfaits duaux.
\end{lemma}

\begin{proof}
Par le lemme \ref{lemma-dual-perfect-complex}, nous avons
$R\lim E \otimes^\mathbf{L}_A K_n^\vee =
R\lim R\Hom_A(K_n, E)$
qui s'insère dans le triangle distingué
$$
R\lim R\Hom_A(K_n, E) \to
\prod R\Hom_A(K_n, E) \to
\prod R\Hom_A(K_n, E)
$$
Comme $K$ s'insère de même dans le triangle distingué
$\bigoplus K_n \to \bigoplus K_n \to K$, il suffit de montrer que
$\prod R\Hom_A(K_n, E) = R\Hom_A(\bigoplus K_n, E)$.
C'est une conséquence formelle de (\ref{equation-internal-hom})
et du fait que le produit tensoriel dérivé commute aux sommes directes.
\end{proof}

\begin{lemma}
\label{lemma-colimit-perfect-complexes}
Soit $R = \colim_{i \in I} R_i$ une colimite filtrante d'anneaux.
\begin{enumerate}
\item Étant donné un objet parfait $K$ de $D(R)$, il existe $i \in I$
et un objet parfait $K_i$ de $D(R_i)$ tels que
$K \cong K_i \otimes_{R_i}^\mathbf{L} R$ in $D(R)$.
\item Étant donnés $0 \in I$ et $K_0, L_0 \in D(R_0)$ avec $K_0$ parfait,
nous avons
$$
\Hom_{D(R)}(K_0 \otimes_{R_0}^\mathbf{L} R, L_0 \otimes_{R_0}^\mathbf{L} R) =
\colim_{i \geq 0}
\Hom_{D(R_i)}(K_0 \otimes_{R_0}^\mathbf{L} R_i,
L_0 \otimes_{R_0}^\mathbf{L} R_i)
$$
\end{enumerate}
Autrement dit, la catégorie triangulée des complexes parfaits sur $R$
est la colimite des catégories triangulées des complexes parfaits sur $R_i$.
\end{lemma}

\begin{proof}
Nous utiliserons sans autre mention les résultats d'Algèbre,
lemmes \ref{algebra-lemma-module-map-property-in-colimit} et
\ref{algebra-lemma-colimit-category-fp-modules}. Ces lemmes affirment
notamment que la catégorie des $R$-modules de présentation finie est la
colimite des catégories des $R_i$-modules de présentation finie. Comme les
modules projectifs finis se caractérisent comme les facteurs directs des
modules libres finis (Algèbre, lemme
\ref{algebra-lemma-finite-projective}), il en va de même pour la catégorie
des modules projectifs finis. Ceci démontre (1) d'après notre définition
des objets parfaits de $D(R)$.

\medskip\noindent
Pour démontrer (2), nous pouvons représenter $K_0$ par un complexe borné
$K_0^\bullet$ de $R_0$-modules projectifs finis. Nous pouvons représenter
$L_0$ par un complexe K-plat $L_0^\bullet$
(lemme \ref{lemma-K-flat-resolution}). Alors nous avons
$$
\Hom_{D(R)}(K_0 \otimes_{R_0}^\mathbf{L} R, L_0 \otimes_{R_0}^\mathbf{L} R) =
\Hom_{K(R)}(K_0^\bullet \otimes_{R_0} R, L_0^\bullet \otimes_{R_0} R)
$$
par Catégories dérivées, lemme
\ref{derived-lemma-morphisms-from-projective-complex}.
Il en va de même pour le $\Hom$ où $R$ est remplacé par $R_i$. Comme
seul un nombre fini de termes intervient dans le membre de droite, comme
$$
\Hom_R(K_0^p \otimes_{R_0} R, L_0^q \otimes_{R_0} R) =
\colim_{i \geq 0}
\Hom_{R_i}(K_0^p \otimes_{R_0} R_i, L_0^q \otimes_{R_0} R_i)
$$
d'après les lemmes cités au début de la démonstration, et comme les
colimites filtrantes sont exactes
(Algèbre, lemme \ref{algebra-lemma-directed-colimit-exact}),
nous concluons que (2) est également vérifiée.
\end{proof}






\section{Relèvement de complexes}
\label{section-lifting-complexes}

\noindent
Soit $R$ un anneau. Soit $I \subset R$ un idéal. Le problème de relèvement
que nous allons considérer est le suivant. Supposons donnés un objet $K$ de
$D(R)$ et un complexe $E^\bullet$ de $R/I$-modules tel que $E^\bullet$
représente $K \otimes_R^\mathbf{L} R/I$ dans $D(R)$.
Question : existe-t-il un complexe de $R$-modules $P^\bullet$ relevant
$E^\bullet$ et représentant $K$ dans $D(R)$ ?
En général, la réponse à cette question est négative, mais on peut faire
quelque chose dans les bons cas. Commençons par le relèvement des complexes
acycliques.

\begin{lemma}
\label{lemma-lift-acyclic-complex}
Soit $R$ un anneau. Soit $I \subset R$ un idéal. Soit $\mathcal{P}$
une classe de $R$-modules. Supposons que
\begin{enumerate}
\item chaque $P \in \mathcal{P}$ soit un $R$-module projectif,
\item si $P_1 \in \mathcal{P}$ et $P_1 \oplus P_2 \in \mathcal{P}$, alors
$P_2 \in \mathcal{P}$, et
\item si $f : P_1 \to P_2$, avec $P_1, P_2 \in \mathcal{P}$, est surjectif
modulo $I$, alors $f$ soit surjectif.
\end{enumerate}
Alors, pour tout complexe acyclique borné supérieurement $E^\bullet$
dont les termes sont de la forme $P/IP$ avec $P \in \mathcal{P}$,
il existe un complexe acyclique borné supérieurement $P^\bullet$ dont les
termes appartiennent à $\mathcal{P}$ et qui relève $E^\bullet$.
\end{lemma}

\begin{proof}
Disons que $E^i = 0$ pour $i > b$. Supposons donnés $n$ et un morphisme
de complexes
$$
\xymatrix{
& & P^n \ar[r] \ar[d] & P^{n + 1} \ar[r] \ar[d] & \ldots \ar[r]
& P^b \ar[r] \ar[d] & 0 \ar[r] \ar[d] & \ldots \\
\ldots \ar[r] & E^{n - 1} \ar[r] &
E^n \ar[r] & E^{n + 1} \ar[r] & \ldots \ar[r] &
E^b \ar[r] & 0 \ar[r] & \ldots
}
$$
où $P^i \in \mathcal{P}$, où
$P^n \to P^{n + 1} \to \ldots \to P^b$ est acyclique en degrés
$\geq n + 1$, et où les morphismes verticaux induisent des isomorphismes
$P^i/IP^i \to E^i$. Dans cette situation, on peut choisir par récurrence des
isomorphismes $P^i = Z^i \oplus Z^{i + 1}$ tels que les morphismes
$P^i \to P^{i + 1}$ soient donnés par
$Z^i \oplus Z^{i + 1} \to Z^{i + 1} \to Z^{i + 1} \oplus Z^{i + 2}$.
Par la propriété (2) et par récurrence, nous voyons que
$Z^i \in \mathcal{P}$. Choisissons $P^{n - 1} \in \mathcal{P}$ et un
isomorphisme $P^{n - 1}/IP^{n - 1} \to E^{n - 1}$. Comme $P^{n - 1}$ est
projectif et comme $Z^n/IZ^n = \Im(E^{n - 1} \to E^n)$, nous pouvons relever
le morphisme $P^{n - 1} \to E^{n - 1} \to E^n$ en un morphisme
$P^{n - 1} \to Z^n$. Par la propriété (3), le morphisme
$P^{n - 1} \to Z^n$ est surjectif.
Nous obtenons ainsi une extension du diagramme en ajoutant $P^{n - 1}$ et
les morphismes que nous venons de construire à gauche de $P^n$. Puisqu'un
diagramme de la forme souhaitée existe pour $n > b$, nous concluons par
récurrence sur $n$.
\end{proof}

\begin{lemma}
\label{lemma-lift-complex}
Soit $R$ un anneau. Soit $I \subset R$ un idéal. Soit $\mathcal{P}$
une classe de $R$-modules.
Soient $K \in D(R)$ et $E^\bullet$ un complexe de $R/I$-modules représentant
$K \otimes_R^\mathbf{L} R/I$. Supposons que
\begin{enumerate}
\item chaque $P \in \mathcal{P}$ soit un $R$-module projectif,
\item $P_1 \in \mathcal{P}$ et $P_1 \oplus P_2 \in \mathcal{P}$ si et
seulement si $P_1, P_2 \in \mathcal{P}$,
\item si $f : P_1 \to P_2$, avec $P_1, P_2 \in \mathcal{P}$, est surjectif
modulo $I$, alors $f$ soit surjectif,
\item $E^\bullet$ soit borné supérieurement et $E^i$ soit de la forme
$P/IP$ avec $P \in \mathcal{P}$, et
\item $K$ puisse être représenté par un complexe borné supérieurement
dont les termes appartiennent à $\mathcal{P}$.
\end{enumerate}
Alors il existe un complexe borné supérieurement $P^\bullet$ dont les termes
appartiennent à $\mathcal{P}$, tel que $P^\bullet/IP^\bullet$ soit isomorphe
à $E^\bullet$ et qui représente $K$ dans $D(R)$.
\end{lemma}

\begin{proof}
Par l'hypothèse (5), nous pouvons représenter $K$ par un complexe borné
supérieurement $K^\bullet$ dont les termes appartiennent à $\mathcal{P}$.
Alors $K \otimes_R^\mathbf{L} R/I$ est représenté par
$K^\bullet/IK^\bullet$. Comme, par (4), $E^\bullet$ est un complexe borné
supérieurement de $R/I$-modules projectifs, nous pouvons choisir un
quasi-isomorphisme $\delta : E^\bullet \to K^\bullet/IK^\bullet$
(Catégories dérivées, lemme
\ref{derived-lemma-morphisms-from-projective-complex}).
Soit $C^\bullet$ le cône de $\delta$
(Catégories dérivées, définition \ref{derived-definition-cone}).
Le module $C^i$ est la somme directe $K^i/IK^i \oplus E^{i + 1}$ ; il est
donc de la forme $P/IP$ pour un certain $P \in \mathcal{P}$, puisque (2)
affirme notamment que $\mathcal{P}$ est stable par sommes. Comme $C^\bullet$
est acyclique, nous pouvons appliquer le lemme
\ref{lemma-lift-acyclic-complex} et trouver un relèvement acyclique
$A^\bullet$ de $C^\bullet$. Le complexe $A^\bullet$ est borné supérieurement
et ses termes appartiennent à $\mathcal{P}$. Dans
$$
\xymatrix{
K^\bullet \ar@{..>}[r] \ar[d] & A^\bullet \ar[d] \\
K^\bullet/IK^\bullet \ar[r] & C^\bullet \ar[r] & E^\bullet[1]
}
$$
on peut trouver la flèche en pointillés rendant le diagramme commutatif
d'après Catégories dérivées, lemme
\ref{derived-lemma-morphisms-lift-projective}.
Nous montrerons ci-dessous qu'il résulte de (1), (2), (3) que
$K^i \to A^i$ est l'inclusion d'un facteur direct pour tout $i$. Par la
propriété (2), nous voyons que $P^i = \Coker(K^i \to A^i)$ appartient à
$\mathcal{P}$. Nous pouvons donc prendre
$P^\bullet = \Coker(K^\bullet \to A^\bullet)[-1]$ pour conclure.

\medskip\noindent
Pour achever la démonstration, nous devons établir ce qui suit : si
$f : P_1 \to P_2$, avec $P_1, P_2 \in \mathcal{P}$, et si
$P_1/IP_1 \to P_2/IP_2$ est une injection scindée dont le conoyau est de la
forme $P_3/IP_3$ pour un certain $P_3 \in \mathcal{P}$, alors $f$ est une
injection scindée. Écrivons $E_i = P_i/IP_i$. Alors
$E_2 = E_1 \oplus E_3$. Comme $P_2$ est projectif, nous pouvons choisir un
morphisme $g : P_2 \to P_3$ relevant le morphisme $E_2 \to E_3$. Par la
condition (3), $g$ est surjectif, donc scindé puisque $P_3$ est projectif.
Posons $P_1' = \Ker(g)$ et choisissons une décomposition
$P_2 = P'_1 \oplus P_3$. Alors $P'_1 \in \mathcal{P}$ par (2). Nous ne
savons pas si $g \circ f = 0$, mais nous pouvons considérer le morphisme
$$
P_1 \xrightarrow{f} P_2 \xrightarrow{projection} P'_1
$$
La composée modulo $I$ est un isomorphisme. Comme $P'_1$ est projectif,
nous pouvons scinder $P_1 = T \oplus P'_1$. Si $T = 0$, nous avons terminé,
car alors $P_2 \to P'_1$ est une rétraction de $f$. Nous voyons que
$T \in \mathcal{P}$ par (2). En réduisant modulo $I$, nous obtenons
$T/IT = 0$. Comme $0 \in \mathcal{P}$ (puisque c'est un facteur direct de
tout $P$ dans $\mathcal{P}$), le morphisme $0 \to T$ est surjectif et nous
concluons que $T = 0$, comme souhaité.
\end{proof}

\begin{lemma}
\label{lemma-lift-complex-projectives}
Soit $R$ un anneau. Soit $I \subset R$ un idéal. Soit $E^\bullet$
un complexe de $R/I$-modules. Soit $K$ un objet de $D(R)$. Supposons que
\begin{enumerate}
\item $E^\bullet$ soit un complexe borné supérieurement de
$R/I$-modules projectifs,
\item $K \otimes_R^\mathbf{L} R/I$ soit représenté par $E^\bullet$ dans
$D(R/I)$, et
\item $I$ soit un idéal nilpotent.
\end{enumerate}
Alors il existe un complexe borné supérieurement $P^\bullet$ de
$R$-modules projectifs représentant $K$ dans $D(R)$ tel que
$P^\bullet \otimes_R R/I$ soit isomorphe à $E^\bullet$.
\end{lemma}

\begin{proof}
Appliquons le lemme \ref{lemma-lift-complex} à la classe $\mathcal{P}$
de tous les $R$-modules projectifs. Les propriétés (1) et (2) du lemme
sont immédiates. La propriété (3) résulte du lemme de Nakayama
(Algèbre, lemme \ref{algebra-lemma-NAK}).
La propriété (4) résulte du fait que l'on peut relever les
$R/I$-modules projectifs en des $R$-modules projectifs, voir
Algèbre, lemme \ref{algebra-lemma-lift-projective-module}.
Pour voir que (5) est vérifiée, il suffit de montrer que $K$ appartient à
$D^{-}(R)$. Puisque $K \otimes_R^\mathbf{L} R/I$ appartient à $D^{-}(R/I)$
car $E^\bullet$ est borné supérieurement, ceci résulte du
lemme \ref{lemma-check-tor-dimension-modulo-nilpotent}.
\end{proof}

\begin{lemma}
\label{lemma-check-pseudo-coherent-modulo-nilpotent}
Soit $R' \to R$ un homomorphisme surjectif d'anneaux dont le noyau est un
idéal nilpotent. Soit $K' \in D(R')$ et posons
$K = K' \otimes_{R'}^\mathbf{L} R$. Alors $K$ est pseudo-cohérent si et
seulement si $K'$ est pseudo-cohérent.
\end{lemma}

\begin{proof}
Une implication résulte du lemme \ref{lemma-pull-pseudo-coherent}.
Pour l'autre, supposons $K$ pseudo-cohérent. Alors, par le
lemme \ref{lemma-pseudo-coherent}, nous pouvons représenter $K$ par un
complexe borné supérieurement $E^\bullet$ de $R$-modules libres finis. Par le
lemme \ref{lemma-lift-complex-projectives}, nous pouvons représenter
$K'$ par un complexe borné supérieurement $P^\bullet$ de $R'$-modules
projectifs tel que $P^n \otimes_{R'} R = E^n$. Le lemme de Nakayama montre
que $P^n$ est libre fini, et nous concluons que $K'$ est également
pseudo-cohérent.
\end{proof}

\begin{lemma}
\label{lemma-lift-complex-stably-frees}
Soit $R$ un anneau. Soit $I \subset R$ un idéal. Soit $E^\bullet$
un complexe de $R/I$-modules. Soit $K$ un objet de $D(R)$. Supposons que
\begin{enumerate}
\item $E^\bullet$ soit un complexe borné supérieurement de
$R/I$-modules libres stables finis,
\item $K \otimes_R^\mathbf{L} R/I$ soit représenté par $E^\bullet$ dans
$D(R/I)$,
\item $K^\bullet$ soit pseudo-cohérent, et
\item tout élément de $1 + I$ soit inversible.
\end{enumerate}
Alors il existe un complexe borné supérieurement $P^\bullet$ de
$R$-modules libres stables finis représentant $K$ dans $D(R)$ tel que
$P^\bullet \otimes_R R/I$ soit isomorphe à $E^\bullet$. De plus, si $E^i$
est libre, alors $P^i$ est libre.
\end{lemma}

\begin{proof}
Appliquons le lemme \ref{lemma-lift-complex} à la classe $\mathcal{P}$
de tous les $R$-modules libres stables finis. La propriété (1) du lemme
est immédiate. La propriété (2) résulte du
lemme \ref{lemma-exact-category-stably-free}. La propriété (3) résulte du
lemme de Nakayama (Algèbre, lemme \ref{algebra-lemma-NAK}).
La propriété (4) résulte du fait que l'on peut relever les $R/I$-modules
libres stables finis en des $R$-modules libres stables finis, voir le
lemme \ref{lemma-lift-stably-free}. La partie (5) est vérifiée parce qu'un
complexe pseudo-cohérent peut être représenté par un complexe borné
supérieurement de $R$-modules libres finis. La dernière assertion du lemme
résulte du lemme \ref{lemma-isomorphic-finite-projective-lifts}.
\end{proof}

\begin{lemma}
\label{lemma-lift-pseudo-coherent-from-residue-field}
Soit $(R, \mathfrak m, \kappa)$ un anneau local. Soit $K \in D(R)$
pseudo-cohérent. Posons
$d_i = \dim_\kappa H^i(K \otimes_R^\mathbf{L} \kappa)$.
Alors $d_i < \infty$ et, pour un certain $b \in \mathbf{Z}$, nous avons
$d_i = 0$ pour $i > b$. Il existe alors un complexe
$$
\ldots \to
R^{\oplus d_{b - 2}} \to
R^{\oplus d_{b - 1}} \to
R^{\oplus d_b} \to 0 \to \ldots
$$
représentant $K$ dans $D(R)$. De plus, ce complexe est unique à
isomorphisme près (!).
\end{lemma}

\begin{proof}
Remarquons que $K \otimes_R^\mathbf{L} \kappa$ est pseudo-cohérent comme
objet de $D(\kappa)$, voir le lemme \ref{lemma-pull-pseudo-coherent}.
Ainsi les espaces de cohomologie sont de dimension finie et s'annulent
au-dessus d'un certain degré. Tout objet de $D(\kappa)$ est isomorphe dans
$D(\kappa)$ à un complexe $E^\bullet$ dont les différentielles sont nulles.
En particulier, $E^i \cong \kappa^{\oplus d_i}$ est libre fini.
L'existence résulte alors du lemme
\ref{lemma-lift-complex-stably-frees}.

\medskip\noindent
Supposons que nous ayons deux complexes $F^\bullet$ et $G^\bullet$
avec $F^i$ et $G^i$ libres de rang $d_i$, représentant $K$.
Nous pouvons alors choisir un morphisme de complexes
$\beta : F^\bullet \to G^\bullet$ représentant l'isomorphisme
$F^\bullet \cong K \cong G^\bullet$, voir Catégories dérivées, lemme
\ref{derived-lemma-morphisms-from-projective-complex}.
Le morphisme de complexes induit
$\beta \otimes 1 : F^\bullet \otimes_R^\mathbf{L} \kappa \to
G^\bullet \otimes_R^\mathbf{L} \kappa$
doit être un isomorphisme de complexes, puisque les différentielles de
$F^\bullet \otimes_R^\mathbf{L} \kappa$ et de
$G^\bullet \otimes_R^\mathbf{L} \kappa$ sont nulles. Ainsi
$\beta^i : F^i \to G^i$ est un morphisme de $R$-modules libres finis dont
la réduction modulo $\mathfrak m$ est un isomorphisme. Par conséquent,
$\beta^i$ est un isomorphisme, ce qui conclut.
\end{proof}

\begin{lemma}
\label{lemma-lift-perfect-from-residue-field}
Soit $R$ un anneau. Soit $\mathfrak p \subset R$ un idéal premier. Soit
$K \in D(R)$ parfait. Posons
$d_i =
\dim_{\kappa(\mathfrak p)} H^i(K \otimes_R^\mathbf{L} \kappa(\mathfrak p))$.
Alors $d_i < \infty$ et seul un nombre fini de ces entiers sont non nuls.
Il existe alors $f \in R$, $f \not \in \mathfrak p$, et un complexe
$$
\ldots \to 0 \to R_f^{\oplus d_a} \to R_f^{\oplus d_{a + 1}} \to
\ldots \to
R_f^{\oplus d_{b - 1}} \to
R_f^{\oplus d_b} \to 0
\to \ldots
$$
représentant $K \otimes_R^\mathbf{L} R_f$ dans $D(R_f)$.
\end{lemma}

\begin{proof}
Remarquons que $K \otimes_R^\mathbf{L} \kappa(\mathfrak p)$ est parfait
comme objet de $D(\kappa(\mathfrak p))$, voir le
lemme \ref{lemma-pull-perfect}. Ainsi, seul un nombre fini de $d_i$ sont
non nuls et ils sont tous finis. En appliquant le
lemme \ref{lemma-lift-pseudo-coherent-from-residue-field}, nous obtenons
un complexe représentant $K$ et ayant la forme souhaitée sur l'anneau local
$R_\mathfrak p$. Nous avons $R_\mathfrak p = \colim R_f$ pour
$f \in R$, $f \not \in \mathfrak p$
(Algèbre, lemme \ref{algebra-lemma-localization-colimit}).
Nous concluons par le lemme \ref{lemma-colimit-perfect-complexes}.
Certains détails sont omis.
\end{proof}

\begin{lemma}
\label{lemma-compare-representatives-perfect}
Soit $R$ un anneau. Soit $\mathfrak p \subset R$ un idéal premier. Soient
$M^\bullet$ et $N^\bullet$ des complexes bornés de $R$-modules projectifs
finis représentant le même objet de $D(R)$. Alors il existe $f \in R$,
$f \not \in \mathfrak p$, tel qu'il existe un isomorphisme (!) de complexes
$$
M^\bullet_f \oplus P^\bullet \cong N^\bullet_f \oplus Q^\bullet
$$
où $P^\bullet$ et $Q^\bullet$ sont des sommes directes finies de complexes
triviaux, c'est-à-dire de complexes de la forme
$\ldots \to 0 \to R_f \xrightarrow{1} R_f \to 0 \to \ldots$
(placés en des degrés arbitraires).
\end{lemma}

\begin{proof}
Si nous avons un isomorphisme du type décrit sur la localisation
$R_\mathfrak p$, alors, en utilisant $R_\mathfrak p = \colim R_f$
(Algèbre, lemme \ref{algebra-lemma-localization-colimit}), nous pouvons
descendre cet isomorphisme en un isomorphisme sur $R_f$ pour un certain $f$.
Nous pouvons donc supposer que $R$ est local et que $\mathfrak p$ est
l'idéal maximal.
Dans ce cas, le résultat découle de l'unicité d'un complexe « minimal »
représentant un objet parfait, voir le
lemme \ref{lemma-lift-pseudo-coherent-from-residue-field}, et du fait que
tout complexe est somme directe d'un complexe trivial et d'un complexe
minimal (Algèbre, lemme \ref{algebra-lemma-add-trivial-complex}).
\end{proof}

\begin{lemma}
\label{lemma-lift-complex-finite-projectives}
Soit $R$ un anneau. Soit $I \subset R$ un idéal. Soit $E^\bullet$
un complexe de $R/I$-modules. Soit $K$ un objet de $D(R)$. Supposons que
\begin{enumerate}
\item $E^\bullet$ soit un complexe borné supérieurement de
$R/I$-modules projectifs finis,
\item $K \otimes_R^\mathbf{L} R/I$ soit représenté par $E^\bullet$ dans
$D(R/I)$,
\item $K$ soit pseudo-cohérent, et
\item $(R, I)$ soit une paire hensélienne.
\end{enumerate}
Alors il existe un complexe borné supérieurement $P^\bullet$ de
$R$-modules projectifs finis représentant $K$ dans $D(R)$ tel que
$P^\bullet \otimes_R R/I$ soit isomorphe à $E^\bullet$. De plus, si $E^i$
est libre, alors $P^i$ est libre.
\end{lemma}

\begin{proof}
Appliquons le lemme \ref{lemma-lift-complex} à la classe $\mathcal{P}$
de tous les $R$-modules projectifs finis. Les propriétés (1) et (2)
du lemme sont immédiates. La propriété (3) résulte du lemme de Nakayama
(Algèbre, lemme \ref{algebra-lemma-NAK}).
La propriété (4) résulte du fait que l'on peut relever les
$R/I$-modules projectifs finis en des $R$-modules projectifs finis, voir le
lemme \ref{lemma-lift-finite-projective-module}. La propriété (5) est
vérifiée parce qu'un complexe pseudo-cohérent peut être représenté par un
complexe borné supérieurement de $R$-modules libres finis.
Ainsi le lemme \ref{lemma-lift-complex} s'applique et fournit
$P^\bullet$ comme souhaité. La dernière assertion du lemme résulte du
lemme \ref{lemma-isomorphic-finite-projective-lifts}.
\end{proof}



\section{Scindage de complexes}
\label{section-splitting}

\noindent
Dans cette section, nous étudions des conditions impliquant qu'un objet de la
catégorie dérivée d'un anneau soit une somme directe de ses troncations.
Notre méthode consiste à utiliser le lemme suivant (sous des hypothèses
convenables) pour scinder les triangles distingués canoniques
$$
\tau_{\leq i}K \to K \to \tau_{\geq i + 1}K \to (\tau_{\leq i}K)[1]
$$
dans $D(R)$, voir Catégories dérivées, remarque
\ref{derived-remark-truncation-distinguished-triangle}.

\begin{lemma}
\label{lemma-splitting-unique}
Soit $R$ un anneau. Soient $K$ et $L$ des objets de $D(R)$.
Supposons que $L$ soit d'amplitude projective dans $[a, b]$, par exemple
si $L$ est parfait d'amplitude de Tor dans $[a, b]$.
\begin{enumerate}
\item Si $H^i(K) = 0$ pour $i \geq a$, alors
$\Hom_{D(R)}(L, K) = 0$.
\item Si $H^i(K) = 0$ pour $i \geq a + 1$, alors, pour tout triangle
distingué $K \to M \to L \to K[1]$, il existe un isomorphisme
$M \cong K \oplus L$ dans $D(R)$ compatible avec les morphismes du triangle
distingué.
\item Si $H^i(K) = 0$ pour $i \geq a$, alors l'isomorphisme de (2)
existe et est unique.
\end{enumerate}
\end{lemma}

\begin{proof}
L'hypothèse que $L$ est d'amplitude projective dans $[a, b]$ signifie que
nous pouvons représenter $L$ par un complexe $L^\bullet$ de $R$-modules
projectifs tel que $L^i = 0$ pour $i \not \in [a, b]$, voir la définition
\ref{definition-projective-dimension}.
Si $L$ est parfait d'amplitude de Tor dans $[a, b]$, alors nous pouvons
représenter $L$ par un complexe $L^\bullet$ de $R$-modules projectifs finis
tel que $L^i = 0$ pour $i \not \in [a, b]$, voir le
lemme \ref{lemma-perfect}. Si $H^i(K) = 0$ pour $i \geq a$, alors $K$ est
quasi-isomorphe à $\tau_{\leq a - 1}K$. Nous pouvons donc représenter $K$
par un complexe $K^\bullet$ de $R$-modules tel que $K^i = 0$ pour
$i \geq a$. Nous obtenons alors
$$
\Hom_{D(R)}(L, K) = \Hom_{K(R)}(L^\bullet, K^\bullet) = 0
$$
par Catégories dérivées, lemme
\ref{derived-lemma-morphisms-from-projective-complex}.
Ceci démontre (1). Sous les hypothèses de (2), nous voyons par (1) que
$\Hom_{D(R)}(L, K[1]) = 0$ ; le triangle distingué est donc scindé d'après
Catégories dérivées, lemme \ref{derived-lemma-split}.
L'unicité de (3) résulte de (1).
\end{proof}

\begin{lemma}
\label{lemma-better-cut-complex-in-two}
Soit $R$ un anneau. Soit $\mathfrak p \subset R$ un idéal premier.
Soit $K^\bullet$ un complexe pseudo-cohérent de $R$-modules.
Supposons que, pour un certain $i \in \mathbf{Z}$, le morphisme
$$
H^i(K^\bullet) \otimes_R \kappa(\mathfrak p)
\longrightarrow
H^i(K^\bullet \otimes_R^{\mathbf{L}} \kappa(\mathfrak p))
$$
soit surjectif. Alors il existe $f \in R$, $f \not \in \mathfrak p$, tel que
$\tau_{\geq i + 1}(K^\bullet \otimes_R R_f)$ soit un objet parfait de
$D(R_f)$ d'amplitude de Tor dans $[i + 1, \infty]$, ainsi qu'un
isomorphisme canonique
$$
K^\bullet \otimes_R R_f \cong
\tau_{\leq i}(K^\bullet \otimes_R R_f) \oplus
\tau_{\geq i + 1}(K^\bullet \otimes_R R_f)
$$
dans $D(R_f)$.
\end{lemma}

\begin{proof}
Dans cette démonstration, tous les produits tensoriels sont pris sur $R$ et
nous écrivons $\kappa = \kappa(\mathfrak p)$. Nous pouvons supposer que
$K^\bullet$ est un complexe borné supérieurement de $R$-modules libres
finis. Examinons ce qui se passe en degré $i$ :
$$
\ldots \to K^{i - 1} \xrightarrow{d^{i - 1}} K^i \xrightarrow{d^i}
K^{i + 1} \to \ldots
$$
Soient $0 \subset V \subset W \subset K^i \otimes \kappa$ définis
par les formules
$$
V = \Im\left(
K^{i - 1} \otimes \kappa \to K^i \otimes \kappa
\right)
\quad\text{et}\quad
W = \Ker\left(
K^i \otimes \kappa \to K^{i + 1} \otimes \kappa
\right)
$$
Posons $\dim(V) = r$, $\dim(W/V) = s$ et
$\dim(K^i \otimes \kappa/W) = t$. Nous pouvons choisir
$x_1, \ldots, x_r \in K^{i - 1}$ dont les images par $d^{i - 1}$ forment
une base de $V$. Par notre hypothèse, nous pouvons choisir
$y_1, \ldots, y_s \in \Ker(d^i)$ dont les images forment une base de $W/V$.
Enfin, choisissons $z_1, \ldots, z_t \in K^i$ dont les images forment une
base de $K^i \otimes \kappa/W$. Nous voyons alors que les éléments
$d^i(z_1), \ldots, d^i(z_t) \in K^{i + 1}$ sont linéairement indépendants
dans $K^{i + 1} \otimes_R \kappa$.
D'après Algèbre, lemme \ref{algebra-lemma-cokernel-flat}, après avoir
remplacé $R$ par $R_f$ pour un certain
$f \in R$, $f \not \in \mathfrak p$, nous pouvons supposer que
\begin{enumerate}
\item $d^i(x_a), y_b, z_c$ est une base sur $R$ de $K^i$,
\item $d^i(z_1), \ldots, d^i(z_t)$ sont linéairement indépendants sur $R$
dans $K^{i + 1}$, et
\item le quotient $E^{i + 1} = K^{i + 1}/\sum Rd^i(z_c)$ est projectif fini.
\end{enumerate}
Comme $d^i$ annule $d^{i - 1}(x_a)$ et $y_b$, nous déduisons de la
condition (2) que $E^{i + 1} = \Coker(d^i : K^i \to K^{i + 1})$.
Nous voyons ainsi que
$$
\tau_{\geq i + 1}K^\bullet =
(\ldots \to 0 \to E^{i + 1} \to K^{i + 2} \to \ldots)
$$
est un complexe borné de modules projectifs finis situé dans les degrés
$[i + 1, b]$ pour un certain $b$. Ainsi
$\tau_{\geq i + 1}K^\bullet$ est parfait d'amplitude $[i + 1, b]$.
Comme $\tau_{\leq i}K^\bullet$ n'a pas de cohomologie dans les degrés $> i$,
nous pouvons appliquer le lemme \ref{lemma-splitting-unique} au triangle
distingué
$$
\tau_{\leq i}K^\bullet \to K^\bullet \to \tau_{\geq i + 1}K^\bullet \to
(\tau_{\leq i}K^\bullet)[1]
$$
(Catégories dérivées, remarque
\ref{derived-remark-truncation-distinguished-triangle}) pour conclure.
\end{proof}

\begin{lemma}
\label{lemma-isolate-a-cohomology-group}
Soit $R$ un anneau. Soit $\mathfrak p \subset R$ un idéal premier.
Soit $K^\bullet$ un complexe pseudo-cohérent de $R$-modules.
Supposons que, pour un certain $i \in \mathbf{Z}$, les morphismes
$$
H^i(K^\bullet) \otimes_R \kappa(\mathfrak p)
\longrightarrow
H^i(K^\bullet \otimes_R^{\mathbf{L}} \kappa(\mathfrak p))
\quad\text{et}\quad
H^{i - 1}(K^\bullet) \otimes_R \kappa(\mathfrak p)
\longrightarrow
H^{i - 1}(K^\bullet \otimes_R^{\mathbf{L}} \kappa(\mathfrak p))
$$
soient surjectifs. Alors il existe $f \in R$, $f \not \in \mathfrak p$,
tel que
\begin{enumerate}
\item $\tau_{\geq i + 1}(K^\bullet \otimes_R R_f)$ soit un objet parfait
de $D(R_f)$ d'amplitude de Tor dans $[i + 1, \infty]$,
\item $H^i(K^\bullet)_f$ soit un $R_f$-module libre fini, et
\item il existe une décomposition canonique en somme directe
$$
K^\bullet \otimes_R R_f \cong
\tau_{\leq i - 1}(K^\bullet \otimes_R R_f) \oplus
H^i(K^\bullet)_f[-i] \oplus
\tau_{\geq i + 1}(K^\bullet \otimes_R R_f)
$$
dans $D(R_f)$.
\end{enumerate}
\end{lemma}

\begin{proof}
Le lemme \ref{lemma-better-cut-complex-in-two} nous donne (1), ainsi qu'un
scindage
$K^\bullet \otimes_R R_f = \tau_{\leq i}K^\bullet \otimes_R R_f \oplus
\tau_{\geq i + 1}K^\bullet \otimes_R R_f$
dans $D(R_f)$. En appliquant une nouvelle fois le
lemme \ref{lemma-better-cut-complex-in-two} à
$\tau_{\leq i}K^\bullet \otimes_R R_f$, nous obtenons (après avoir choisi
$f$ convenablement) un scindage
$\tau_{\leq i}K^\bullet \otimes_R R_f =
\tau_{\leq i - 1}K^\bullet \otimes_R R_f \oplus H^i(K^\bullet)_f$ in $D(R_f)$
ainsi que la conclusion que $H^i(K)_f$ est un module plat parfait,
c'est-à-dire projectif fini.
\end{proof}

\begin{lemma}
\label{lemma-cut-complex-in-two}
Soit $R$ un anneau. Soit $\mathfrak p \subset R$ un idéal premier.
Soit $i \in \mathbf{Z}$. Soit $K^\bullet$ un complexe pseudo-cohérent
de $R$-modules tel que
$H^i(K^\bullet \otimes_R^{\mathbf{L}} \kappa(\mathfrak p)) = 0$.
Alors il existe $f \in R$, $f \not \in \mathfrak p$, et une décomposition
canonique en somme directe
$$
K^\bullet \otimes_R R_f =
\tau_{\geq i + 1}(K^\bullet \otimes_R R_f) \oplus
\tau_{\leq i - 1}(K^\bullet \otimes_R R_f)
$$
dans $D(R_f)$, où $\tau_{\geq i + 1}(K^\bullet \otimes_R R_f)$ est un
complexe parfait d'amplitude de Tor dans $[i + 1, \infty]$.
\end{lemma}

\begin{proof}
Il s'agit d'un cas particulier souvent utilisé du
lemme \ref{lemma-better-cut-complex-in-two}.
Voici une démonstration directe.
Nous pouvons supposer que $K^\bullet$ est un complexe borné supérieurement
de $R$-modules libres finis. Examinons ce qui se passe en degré $i$ :
$$
\ldots \to K^{i - 2} \to R^{\oplus l}
\to R^{\oplus m} \to R^{\oplus n} \to K^{i + 2} \to \ldots
$$
Soit $A$ la matrice $m \times l$ correspondant à
$K^{i - 1} \to K^i$, et soit $B$ la matrice $n \times m$ correspondant à
$K^i \to K^{i + 1}$. L'hypothèse affirme que $A \bmod \mathfrak p$ est de
rang $r$ et que $B \bmod \mathfrak p$ est de rang $m - r$. Autrement dit,
il existe un mineur $r \times r$, noté $a$, de $A$, qui n'appartient pas à
$\mathfrak p$, et un mineur $(m - r) \times (m - r)$, noté $b$, de $B$,
qui n'appartient pas à $\mathfrak p$. Posons $f = ab$. Après inversion de
$f$, nous pouvons trouver des décompositions en somme directe
$K^{i - 1} = R^{\oplus l - r} \oplus R^{\oplus r}$,
$K^i = R^{\oplus r} \oplus R^{\oplus m - r}$,
$K^{i + 1} = R^{\oplus m - r} \oplus R^{\oplus n - m + r}$
telles que le morphisme de modules $K^{i - 1} \to K^i$ annule
$R^{\oplus l - r}$ et induise un isomorphisme de $R^{\oplus r}$ sur le
facteur correspondant de $K^i$, et telles que le morphisme de modules
$K^i \to K^{i + 1}$ annule $R^{\oplus r}$ et induise un isomorphisme de
$R^{\oplus m - r}$ sur le facteur correspondant de $K^{i + 1}$.
Ainsi $K^\bullet$ devient quasi-isomorphe à
$$
\ldots \to K^{i - 2} \to R^{\oplus l - r}
\to 0 \to R^{\oplus n - m + r} \to K^{i + 2} \to \ldots
$$
et tout est clair.
\end{proof}

\begin{lemma}
\label{lemma-split-using-ext-injective}
Soit $R$ un anneau. Soit $K \in D^-(R)$. Soit $a \in \mathbf{Z}$.
Supposons que, pour toute injection de $R$-modules $M \to M'$, le morphisme
$\Ext^{-a}_R(K, M) \to \Ext^{-a}_R(K, M')$ soit injectif.
Alors il existe une unique décomposition en somme directe
$K \cong \tau_{\leq a}K \oplus \tau_{\geq a + 1}K$
et $\tau_{\geq a + 1}K$ est d'amplitude projective dans $[a + 1, b]$
pour un certain $b$.
\end{lemma}

\begin{proof}
Considérons le triangle distingué
$$
\tau_{\leq a}K \to K \to \tau_{\geq a + 1}K \to (\tau_{\leq a}K)[1]
$$
dans $D(R)$, voir Catégories dérivées, remarque
\ref{derived-remark-truncation-distinguished-triangle}.
Remarquons que
$\Ext^{-a}_R(\tau_{\leq a}K, M) = \Hom_R(H^a(K), M)$
et
$\Ext^{-a - 1}_R(\tau_{\leq a}K, M) = 0$, see
voir Catégories dérivées, lemme \ref{derived-lemma-negative-exts}. Ainsi,
la suite exacte longue de $\Ext$ donne une suite exacte
$$
0 \to
\Ext^{-a}_R(\tau_{\geq a + 1}K, M) \to
\Ext^{-a}_R(K, M) \to
\Hom_R(H^a(K), M)
$$
fonctorielle en le $R$-module $M$.
Si maintenant $I$ est un $R$-module injectif, alors
$\Ext^{-a}_R(\tau_{\geq a + 1}K, I) = 0$, par exemple d'après
Catégories dérivées, lemme \ref{derived-lemma-compute-ext-resolutions}.
Comme tout module s'injecte dans un module injectif, nous concluons que
$\Ext^{-a}_R(\tau_{\geq a + 1}K, M) = 0$ pour tout $R$-module $M$.
Par le lemme \ref{lemma-projective-amplitude}, nous concluons que
$\tau_{\geq a + 1}K$ est d'amplitude projective dans $[a + 1, b]$ pour un
certain $b$ (c'est ici que nous utilisons que $K$ est borné supérieurement).
Le scindage résulte alors du lemme \ref{lemma-splitting-unique}.
\end{proof}

\begin{lemma}
\label{lemma-split-using-ext-zero}
Soit $R$ un anneau. Soit $K \in D^-(R)$. Soit $a \in \mathbf{Z}$.
Supposons que $\Ext^{-a}_R(K, M) = 0$ pour tout $R$-module $M$.
Alors il existe une unique décomposition en somme directe
$K \cong \tau_{\leq a - 1}K \oplus \tau_{\geq a + 1}K$
et $\tau_{\geq a + 1}K$ est d'amplitude projective dans $[a + 1, b]$
pour un certain $b$.
\end{lemma}

\begin{proof}
Par le lemme \ref{lemma-split-using-ext-injective}, nous avons une
décomposition en somme directe
$K \cong \tau_{\leq a}K \oplus \tau_{\geq a + 1}K$
et $\tau_{\geq a + 1}K$ est d'amplitude projective dans $[a + 1, b]$
pour un certain $b$. Il est clair que $H^a(K) = 0$, et nous concluons que
$\tau_{\leq a}K = \tau_{\leq a - 1}K$ dans $D(R)$.
\end{proof}














\section{Reconnaissance des complexes parfaits}
\label{section-recognizing-perfect}

\noindent
Quelques lemmes permettant de démontrer que certains complexes sont parfaits.

\begin{lemma}
\label{lemma-lift-bounded-pseudo-coherent-to-perfect}
Soit $R$ un anneau et soit $\mathfrak p \subset R$ un idéal premier.
Soit $K$ pseudo-cohérent et borné inférieurement. Posons
$d_i = \dim_{\kappa(\mathfrak p)}
H^i(K \otimes_R^\mathbf{L} \kappa(\mathfrak p))$.
S'il existe $a \in \mathbf{Z}$ tel que $d_i = 0$ pour $i < a$, alors il
existe $f \in R$, $f \not \in \mathfrak p$, et un complexe
$$
\ldots \to 0 \to R_f^{\oplus d_a} \to R_f^{\oplus d_{a + 1}} \to
\ldots \to
R_f^{\oplus d_{b - 1}} \to
R_f^{\oplus d_b} \to 0
\to \ldots
$$
représentant $K \otimes_R^\mathbf{L} R_f$ dans $D(R_f)$.
En particulier, $K \otimes_R^\mathbf{L} R_f$ est parfait.
\end{lemma}

\begin{proof}
Après avoir diminué $a$, nous pouvons supposer aussi que
$H^i(K^\bullet) = 0$ pour $i < a$. Par le
lemme \ref{lemma-cut-complex-in-two}, après avoir remplacé $R$ par $R_f$
pour un certain $f \in R$, $f \not \in \mathfrak p$, nous pouvons écrire
$K^\bullet = \tau_{\leq a - 1}K^\bullet \oplus
\tau_{\geq a}K^\bullet$ dans $D(R)$ avec $\tau_{\geq a}K^\bullet$ parfait.
Comme $H^i(K^\bullet) = 0$ pour $i < a$, nous voyons que
$\tau_{\leq a - 1}K^\bullet = 0$ dans $D(R)$. Ainsi $K^\bullet$ est parfait.
Nous pouvons alors conclure à l'aide du
lemme \ref{lemma-lift-perfect-from-residue-field}.
\end{proof}

\begin{lemma}
\label{lemma-check-perfect-pointwise}
Soit $R$ un anneau. Soient $a, b \in \mathbf{Z}$.
Soit $K^\bullet$ un complexe pseudo-cohérent de $R$-modules.
Les assertions suivantes sont équivalentes :
\begin{enumerate}
\item $K^\bullet$ est parfait d'amplitude de Tor dans $[a, b]$,
\item pour tout idéal premier $\mathfrak p$, nous avons
$H^i(K^\bullet \otimes_R^{\mathbf{L}} \kappa(\mathfrak p)) = 0$ pour tout
$i \not \in [a, b]$, et
\item pour tout idéal maximal $\mathfrak m$, nous avons
$H^i(K^\bullet \otimes_R^{\mathbf{L}} \kappa(\mathfrak m)) = 0$ pour tout
$i \not \in [a, b]$.
\end{enumerate}
\end{lemma}

\begin{proof}
Nous omettons la démonstration des implications
(1) $\Rightarrow$ (2) $\Rightarrow$ (3).
Supposons (3). Soit $i \in \mathbf{Z}$ tel que $i \not \in [a, b]$.
Par le lemme \ref{lemma-cut-complex-in-two}, l'hypothèse implique que
$H^i(K^\bullet)_{\mathfrak m} = 0$ pour tout idéal maximal de $R$.
Ainsi $H^i(K^\bullet) = 0$, voir Algèbre, lemme
\ref{algebra-lemma-characterize-zero-local}.
De plus, le lemme \ref{lemma-cut-complex-in-two} implique maintenant aussi
que, pour tout idéal maximal $\mathfrak m$, il existe un élément
$f \in R$, $f \not \in \mathfrak m$, tel que
$K^\bullet \otimes_R R_f$ soit parfait d'amplitude de Tor dans $[a, b]$.
Nous concluons donc à l'aide des lemmes \ref{lemma-glue-perfect} et
\ref{lemma-glue-tor-amplitude}.
\end{proof}

\begin{lemma}
\label{lemma-check-perfect-stalks}
Soit $R$ un anneau. Soit $K^\bullet$ un complexe pseudo-cohérent de
$R$-modules. Considérons les conditions suivantes :
\begin{enumerate}
\item $K^\bullet$ est parfait,
\item pour tout idéal premier $\mathfrak p$, le complexe
$K^\bullet \otimes_R R_{\mathfrak p}$ est parfait,
\item pour tout idéal maximal $\mathfrak m$, le complexe
$K^\bullet \otimes_R R_{\mathfrak m}$ est parfait,
\item pour tout idéal premier $\mathfrak p$, nous avons
$H^i(K^\bullet \otimes_R^{\mathbf{L}} \kappa(\mathfrak p)) = 0$ pour tout
$i \ll 0$,
\item pour tout idéal maximal $\mathfrak m$, nous avons
$H^i(K^\bullet \otimes_R^{\mathbf{L}} \kappa(\mathfrak m)) = 0$ pour tout
$i \ll 0$.
\end{enumerate}
Nous avons toujours les implications
$$
(1) \Rightarrow (2) \Leftrightarrow (3) \Leftrightarrow (3)
\Leftrightarrow (4) \Leftrightarrow (5)
$$
Si $K^\bullet$ est borné inférieurement, alors toutes les conditions sont
équivalentes.
\end{lemma}

\begin{proof}
Par le lemme \ref{lemma-pull-perfect}, nous voyons que (1) implique (2).
Il est immédiat que (2) $\Rightarrow$ (3). Comme tout idéal premier
$\mathfrak p$ est contenu dans un idéal maximal $\mathfrak m$, nous pouvons
appliquer le lemme \ref{lemma-pull-perfect} au morphisme
$R_\mathfrak m \to R_\mathfrak p$ pour voir que (3) implique (2).
En appliquant le lemme \ref{lemma-pull-perfect} aux morphismes résiduels
$R_\mathfrak p \to \kappa(\mathfrak p)$ et
$R_\mathfrak m \to \kappa(\mathfrak m)$, nous voyons que (2) implique
(4) et que (3) implique (5).

\medskip\noindent
Supposons $R$ local, d'idéal maximal $\mathfrak m$ et de corps résiduel
$\kappa$. Nous allons montrer que, si
$H^i(K^\bullet \otimes^\mathbf{L} \kappa) = 0$ pour $i < a$ avec un certain
$a$, alors $K$ est parfait. Ceci montrera que (4) implique (2) et que (5)
implique (3), d'où la première partie du lemme. Appliquons d'abord le
lemme \ref{lemma-cut-complex-in-two} avec $i = a - 1$ pour voir que
$K^\bullet = \tau_{\leq a - 1}K^\bullet \oplus \tau_{\geq a}K^\bullet$
dans $D(R)$, avec $\tau_{\geq a}K^\bullet$ parfait d'amplitude de Tor
contenue dans $[a, \infty]$. Pour terminer, nous devons montrer que
$\tau_{\leq a - 1}K$ est nul, c'est-à-dire que ses groupes de cohomologie
sont nuls. Sinon, soit $i$ le plus grand indice tel que
$M = H^i(\tau_{\leq a - 1}K)$ soit non nul. Alors $M$ est un $R$-module
fini, puisque $\tau_{\leq a - 1}K^\bullet$ est pseudo-cohérent
(lemmes \ref{lemma-finite-cohomology} et
\ref{lemma-summands-pseudo-coherent}). Ainsi, par le lemme de Nakayama
(Algèbre, lemme \ref{algebra-lemma-NAK}), nous trouvons que
$M \otimes_R \kappa$ est non nul. Ceci implique que
$$
H^i((\tau_{\leq a - 1}K^\bullet) \otimes_R^\mathbf{L} \kappa) =
H^i(K^\bullet \otimes_R^\mathbf{L} \kappa)
$$
est non nul, ce qui est une contradiction.

\medskip\noindent
Supposons que les conditions équivalentes (2) -- (5) soient vérifiées et
que $K^\bullet$ soit borné inférieurement. Disons que
$H^i(K^\bullet) = 0$ pour $i < a$. Choisissons un idéal maximal
$\mathfrak m$ de $R$. Il suffit de montrer qu'il existe
$f \in R$, $f \not \in \mathfrak m$, tel que
$K^\bullet \otimes_R^\mathbf{L} R_f$ soit parfait
(lemme \ref{lemma-glue-perfect} et Algèbre, lemme
\ref{algebra-lemma-quasi-compact}). Ceci résulte du
lemme \ref{lemma-lift-bounded-pseudo-coherent-to-perfect}.
\end{proof}

\begin{lemma}
\label{lemma-projective-amplitude-pseudo-coherent}
Soit $R$ un anneau. Soit $K$ un objet pseudo-cohérent de $D(R)$.
Soient $a, b \in \mathbf{Z}$. Les assertions suivantes sont équivalentes :
\begin{enumerate}
\item $K$ est d'amplitude projective dans $[a, b]$,
\item $K$ est parfait d'amplitude de Tor dans $[a, b]$,
\item $\Ext^i_R(K, N) = 0$ pour tout $R$-module de présentation finie $N$
et tout
$i \not \in [-b, -a]$,
\item $H^n(K) = 0$ pour $n > b$ et
$\Ext^i_R(K, N) = 0$ pour tout $R$-module de présentation finie $N$
et tout $i > -a$, et
\item $H^n(K) = 0$ pour $n \not \in [a - 1, b]$ et
$\Ext^{-a + 1}_R(K, N) = 0$ pour tout $R$-module de présentation finie $N$.
\end{enumerate}
\end{lemma}

\begin{proof}
La dernière assertion du lemme \ref{lemma-perfect} montre que (2) implique
(1). Si (1) est vérifiée, alors $K$ peut être représenté par un complexe
de modules projectifs $P^i$ tel que $P^i = 0$ pour $i \not \in [a, b]$.
Comme les modules projectifs sont plats (en tant que facteurs directs de
modules libres), nous voyons que $K$ est d'amplitude de Tor dans $[a, b]$,
voir le lemme \ref{lemma-tor-amplitude}. Ainsi, le
lemme \ref{lemma-perfect} montre que (2) est vérifiée.

\medskip\noindent
Dans les conditions (3), (4), (5), l'annulation supposée des groupes ext
$\Ext^i_R(K, M)$ pour $M$ de présentation finie équivaut à leur annulation
pour tout $R$-module $M$, d'après le
lemme \ref{lemma-pseudo-coherence-colimit-ext} et Algèbre, lemme
\ref{algebra-lemma-module-colimit-fp}. Ainsi, l'équivalence de (1), (3),
(4) et (5) résulte du lemme \ref{lemma-projective-amplitude}.
\end{proof}

\noindent
Le lemme suivant est utile pour trouver des complexes parfaits sur un anneau
de polynômes $B = A[x_1, \ldots, x_d]$.

\begin{lemma}
\label{lemma-perfect-over-polynomial-ring}
Soit $A \to B$ un homomorphisme d'anneaux. Soient $a, b \in \mathbf{Z}$.
Soit $d \geq 0$. Soit $K^\bullet$ un complexe de $B$-modules. Supposons que
\begin{enumerate}
\item l'homomorphisme d'anneaux $A \to B$ soit plat,
\item pour tout idéal premier $\mathfrak p \subset A$, l'anneau
$B \otimes_A \kappa(\mathfrak p)$ soit de dimension globale finie $\leq d$,
\item $K^\bullet$ soit pseudo-cohérent comme complexe de $B$-modules, et
\item $K^\bullet$ soit d'amplitude de Tor dans $[a, b]$ comme complexe
de $A$-modules.
\end{enumerate}
Alors $K^\bullet$ est parfait comme complexe de $B$-modules,
d'amplitude de Tor dans $[a - d, b]$.
\end{lemma}

\begin{proof}
Nous pouvons supposer que $K^\bullet$ est un complexe borné supérieurement
de $B$-modules libres finis. En particulier, $K^\bullet$ est plat comme
complexe de $A$-modules et
$K^\bullet \otimes_A M = K^\bullet \otimes_A^{\mathbf{L}} M$ pour tout
$A$-module $M$. Pour tout idéal premier $\mathfrak p$ de $A$, le complexe
$$
K^\bullet \otimes_A \kappa(\mathfrak p)
$$
est un complexe borné supérieurement de modules libres finis sur
$B \otimes_A \kappa(\mathfrak p)$, dont les $H^i$ s'annulent sauf pour
$i \in [a, b]$. Comme $B \otimes_A \kappa(\mathfrak p)$ est de dimension
globale $d$, le lemme \ref{lemma-finite-gl-dim-tor-dimension} montre que
$K^\bullet \otimes_A \kappa(\mathfrak p)$ est d'amplitude de Tor dans
$[a - d, b]$. Soit $\mathfrak q$ un idéal premier de $B$ au-dessus de
$\mathfrak p$. Comme $K^\bullet \otimes_A \kappa(\mathfrak p)$ est un
complexe borné supérieurement de
$B \otimes_A \kappa(\mathfrak p)$-modules libres, nous voyons que
\begin{align*}
K^\bullet \otimes_B^{\mathbf{L}} \kappa(\mathfrak q)
& = K^\bullet \otimes_B \kappa(\mathfrak q) \\
& = (K^\bullet \otimes_A \kappa(\mathfrak p))
\otimes_{B \otimes_A \kappa(\mathfrak p)} \kappa(\mathfrak q) \\
& = (K^\bullet \otimes_A \kappa(\mathfrak p))
\otimes^{\mathbf{L}}_{B \otimes_A \kappa(\mathfrak p)} \kappa(\mathfrak q)
\end{align*}
Les arguments ci-dessus impliquent donc que
$H^i(K^\bullet \otimes_B^{\mathbf{L}} \kappa(\mathfrak q)) = 0$
pour $i \not \in [a - d, b]$. Nous concluons par le
lemme \ref{lemma-check-perfect-pointwise}.
\end{proof}

\noindent
Le lemme suivant est une version locale du
lemme \ref{lemma-perfect-over-polynomial-ring}.
Il peut être utilisé pour trouver des complexes parfaits sur des anneaux
locaux réguliers.

\begin{lemma}
\label{lemma-perfect-over-regular-local-ring}
Soit $A \to B$ un homomorphisme local d'anneaux.
Soient $a, b \in \mathbf{Z}$. Soit $d \geq 0$.
Soit $K^\bullet$ un complexe de $B$-modules. Supposons que
\begin{enumerate}
\item l'homomorphisme d'anneaux $A \to B$ soit plat,
\item l'anneau $B/\mathfrak m_AB$ soit régulier de dimension $d$,
\item $K^\bullet$ soit pseudo-cohérent comme complexe de $B$-modules, et
\item $K^\bullet$ soit d'amplitude de Tor dans $[a, b]$ comme complexe
de $A$-modules ; en fait, il suffit que
$H^i(K^\bullet \otimes_A^\mathbf{L} \kappa(\mathfrak m_A))$
soit non nul seulement pour $i \in [a, b]$.
\end{enumerate}
Alors $K^\bullet$ est parfait comme complexe de $B$-modules,
d'amplitude de Tor dans $[a - d, b]$.
\end{lemma}

\begin{proof}
Par (3), nous pouvons supposer que $K^\bullet$ est un complexe borné
supérieurement de $B$-modules libres finis. Nous calculons
\begin{align*}
K^\bullet \otimes_B^{\mathbf{L}} \kappa(\mathfrak m_B)
& = K^\bullet \otimes_B \kappa(\mathfrak m_B) \\
& = (K^\bullet \otimes_A \kappa(\mathfrak m_A))
\otimes_{B/\mathfrak m_A B} \kappa(\mathfrak m_B) \\
& = (K^\bullet \otimes_A \kappa(\mathfrak m_A))
\otimes^{\mathbf{L}}_{B/\mathfrak m_A B} \kappa(\mathfrak m_B)
\end{align*}
La première égalité vient de ce que $K^\bullet$ est un complexe borné
supérieurement de $B$-modules plats. La deuxième résulte des propriétés
élémentaires du produit tensoriel. La troisième est valable parce que
$K^\bullet \otimes_A \kappa(\mathfrak m_A) =
K^\bullet/ \mathfrak m_A K^\bullet$ est un complexe borné supérieurement
de $B/\mathfrak m_A B$-modules plats. Comme, par (1), $K^\bullet$ est un
complexe borné supérieurement de $A$-modules plats, les modules de
cohomologie $H^i$ du complexe
$K^\bullet \otimes_A \kappa(\mathfrak m_A)$ ne sont non nuls que pour
$i \in [a, b]$, par l'hypothèse (4). Ainsi, la suite spectrale de
l'exemple \ref{example-cohomology-complex-tensored} et le fait que
$B/\mathfrak m_AB$ est de dimension globale finie $d$
(par (2) et Algèbre, proposition
\ref{algebra-proposition-regular-finite-gl-dim}) montrent que
$H^j(K^\bullet \otimes_B^{\mathbf{L}} \kappa(\mathfrak m_B))$ est nul
pour $j \not \in [a - d, b]$.
Ceci achève la démonstration par le
lemme \ref{lemma-check-perfect-pointwise}.
\end{proof}
\section{Caractérisation des complexes parfaits}
\label{section-perfect-compact}

\noindent
Dans cette section, nous démontrons que les complexes parfaits sont exactement
les objets compacts de la catégorie dérivée d'un anneau. Nous montrons d'abord
ce qui suit.

\begin{lemma}
\label{lemma-perfect-ring-classical-generator}
Soit $R$ un anneau. La sous-catégorie pleine $D_{perf}(R) \subset D(R)$ des objets
parfaits est la plus petite sous-catégorie triangulée, strictement pleine et saturée
contenant $R = R[0]$. Autrement dit, $D_{perf}(R) = \langle R \rangle$.
En particulier, $R$ est un générateur classique de $D_{perf}(R)$.
\end{lemma}

\begin{proof}
Pour comprendre le sens de l'énoncé, voir
Catégories dérivées, définitions \ref{derived-definition-saturated} et
\ref{derived-definition-generators}.
Il a été démontré dans les lemmes \ref{lemma-two-out-of-three-perfect} et
\ref{lemma-summands-perfect} que $D_{perf}(R) \subset D(R)$
est une sous-catégorie triangulée, strictement pleine et saturée de $D(R)$.
Bien entendu, $R \in D_{perf}(R)$.

\medskip\noindent
Rappelons que $\langle R \rangle = \bigcup \langle R \rangle_n$.
Pour achever la démonstration, nous allons montrer que si $M \in D_{perf}(R)$
est représenté par
$$
\ldots \to 0 \to M^a \to M^{a + 1} \to \ldots \to M^b \to 0 \to \ldots
$$
où les $M^i$ sont projectifs de type fini, alors $M \in \langle R \rangle_{b - a + 1}$.
La démonstration se fait par récurrence sur $b - a$.
Par définition, $\langle R \rangle_1$ contient tout $R$-module projectif
de type fini placé en un degré quelconque ; cela traite le cas initial
$b - a = 0$ de la récurrence. En général, nous considérons le triangle
distingué
$$
M_b[-b] \to M^\bullet \to \sigma_{\leq b - 1}M^\bullet \to M_b[-b + 1]
$$
Par récurrence, le complexe tronqué $\sigma_{\leq b - 1}M^\bullet$ appartient
à $\langle R \rangle_{b - a}$ et $M_b[-b]$ appartient à $\langle R \rangle_1$.
Ainsi $M^\bullet \in \langle R \rangle_{b - a + 1}$ par définition.
\end{proof}

\noindent
Soit $R$ un anneau. Rappelons que $D(R)$ admet des sommes directes, obtenues
simplement en prenant les sommes directes de complexes, voir
Catégories dérivées, lemme \ref{derived-lemma-direct-sums}.
Nous l'utiliserons sans autre mention dans les lemmes de cette section.

\begin{lemma}
\label{lemma-commutes-with-countable-sums}
Soit $R$ un anneau. Soit $K \in D(R)$ un objet tel que, pour toute
famille dénombrable d'objets $E_n \in D(R)$, l'application canonique
$$
\bigoplus \Hom_{D(R)}(K, E_n) \longrightarrow \Hom_{D(R)}(K, \bigoplus E_n)
$$
soit une bijection. Alors, pour tout système $L_n^\bullet$ de complexes indexé
par $\mathbf{N}$, l'application
$$
\colim \Hom_{D(R)}(K, L^\bullet_n) \longrightarrow \Hom_{D(R)}(K, L^\bullet)
$$
est une bijection, où $L^\bullet$ est la colimite terme à terme, c'est-à-dire
$L^m = \colim L_n^m$ pour tout $m \in \mathbf{Z}$.
\end{lemma}

\begin{proof}
Considérons la suite exacte courte de complexes
$$
0 \to \bigoplus L_n^\bullet \to \bigoplus L_n^\bullet \to L^\bullet \to 0
$$
où la première application est donnée par $1 - t_n$ en degré $n$, avec
$t_n : L_n^\bullet \to L_{n + 1}^\bullet$ l'application de transition.
D'après
Catégories dérivées, lemme \ref{derived-lemma-derived-canonical-delta-functor},
c'est un triangle distingué dans $D(R)$.
Appliquons le foncteur homologique $\Hom_{D(R)}(K, -)$, voir
Catégories dérivées, lemme \ref{derived-lemma-representable-homological}.
Nous obtenons ainsi une longue suite exacte de cohomologie
$$
\xymatrix{
& \ldots \ar[r] & \Hom_{D(R)}(K, \colim L^\bullet_n[-1]) \ar[lld] \\
\Hom_{D(R)}(K, \bigoplus L^\bullet_n) \ar[r] &
\Hom_{D(R)}(K, \bigoplus L^\bullet_n) \ar[r] &
\Hom_{D(R)}(K, \colim L^\bullet_n) \ar[lld] \\
\Hom_{D(R)}(K, \bigoplus L^\bullet_n[1]) \ar[r] & \ldots
}
$$
Puisque nous avons supposé que $\Hom_{D(R)}(K, \bigoplus L^\bullet_n)$
est égal à $\bigoplus \Hom_{D(R)}(K, L^\bullet_n)$, nous voyons que la première
application de chaque ligne du diagramme est injective (d'après sa description
explicite comme somme des applications induites par $1 - t_n$). Nous
concluons donc que $\Hom_{D(R)}(K, \colim L^\bullet_n)$ est le conoyau
de la première application de la ligne médiane du diagramme ci-dessus, ce qu'il
fallait démontrer.
\end{proof}

\noindent
La proposition suivante, qui caractérise les complexes parfaits comme les objets
compacts
(Catégories dérivées, définition \ref{derived-definition-compact-object})
de la catégorie dérivée, apparaît à plusieurs endroits. Voir par exemple
\cite[preuve de la proposition 6.3]{Rickard} (qui traite le cas borné),
\cite[théorème 2.4.3]{TT} (dont l'énoncé ne coïncide pas exactement) et
\cite[proposition 6.4]{Bokstedt-Neeman} (attention aux
épouvantables conventions de notation).

\begin{proposition}
\label{proposition-perfect-is-compact}
Soit $R$ un anneau. Pour un objet $K$ de $D(R)$, les assertions suivantes sont équivalentes
\begin{enumerate}
\item $K$ est parfait ;
\item $K$ est un objet compact de $D(R)$.
\end{enumerate}
\end{proposition}

\begin{proof}
Supposons $K$ parfait, c'est-à-dire que $K$ est quasi-isomorphe à un complexe borné
$P^\bullet$ de modules projectifs de type fini, voir la
définition \ref{definition-perfect}. Si $E_i$ est représenté par le complexe
$E_i^\bullet$, alors $\bigoplus E_i$ est représenté par le complexe
dont le terme de degré $n$ est $\bigoplus E_i^n$. D'autre part,
comme $P^n$ est projectif pour tout $n$, nous avons
$\Hom_{D(R)}(P^\bullet, K^\bullet) = \Hom_{K(R)}(P^\bullet, K^\bullet)$
pour tout complexe de $R$-modules $K^\bullet$, voir
Catégories dérivées,
lemme \ref{derived-lemma-morphisms-from-projective-complex}.
Ainsi $\Hom_{D(R)}(P^\bullet, E^\bullet)$ est la cohomologie du complexe
$$
\prod \Hom_R(P^n, E^{n - 1}) \to
\prod \Hom_R(P^n, E^n) \to
\prod \Hom_R(P^n, E^{n + 1}).
$$
Puisque $P^\bullet$ est borné, nous pouvons remplacer les signes $\prod$
par des signes $\bigoplus$ dans le complexe ci-dessus. Comme chaque $P^n$ est un
$R$-module de type fini, nous voyons que
$\Hom_R(P^n, \bigoplus_i E_i^m) = \bigoplus_i \Hom_R(P^n, E_i^m)$
pour tous $n, m$.
En combinant ces remarques, nous voyons que l'application de
Catégories dérivées, définition \ref{derived-definition-compact-object},
est une bijection.

\medskip\noindent
Réciproquement, supposons $K$ compact.
Représentons $K$ par un complexe $K^\bullet$ et considérons l'application
$$
K^\bullet
\longrightarrow
\bigoplus\nolimits_{n \geq 0} \tau_{\geq n} K^\bullet
$$
où nous avons utilisé les troncations canoniques, voir
Homologie, section \ref{homology-section-truncations}.
Cela a un sens car, en chaque degré, la somme directe à droite est finie.
Par hypothèse, cette application se factorise par une somme directe finie.
Nous concluons que $K \to \tau_{\geq n} K$ est nul pour au moins un $n$,
c'est-à-dire que $K$ appartient à $D^{-}(R)$.

\medskip\noindent
Comme $K \in D^{-}(R)$ et que tout $R$-module est quotient d'un module
libre, nous pouvons représenter $K$ par un complexe $K^\bullet$ borné supérieurement
de $R$-modules libres, voir
Catégories dérivées, lemme \ref{derived-lemma-subcategory-left-resolution}.
Remarquons que
$$
K^\bullet = \bigcup\nolimits_{n \leq 0} \sigma_{\geq n}K^\bullet
$$
où nous avons utilisé les troncations brutales, voir
Homologie, section \ref{homology-section-truncations}.
Le lemme \ref{lemma-commutes-with-countable-sums} montre donc que
$1 : K^\bullet \to K^\bullet$ se factorise par
$\sigma_{\geq n}K^\bullet \to K^\bullet$ dans $D(R)$.
Ainsi $1 : K^\bullet \to K^\bullet$ se factorise sous la forme
$$
K^\bullet \xrightarrow{\varphi} L^\bullet \xrightarrow{\psi} K^\bullet
$$
dans $D(R)$ pour un complexe $L^\bullet$ borné dont les termes sont des
$R$-modules libres. Disons que $L^i = 0$ pour $i \not \in [a, b]$.
Fixons désormais $a, b$. Soit $c$ le plus grand entier $\leq b + 1$
tel que l'on puisse trouver une factorisation de $1_{K^\bullet}$ comme ci-dessus
où $L^i$ soit libre de type fini pour $i < c$. Nous allons démontrer par récurrence que
$c = b + 1$. Écrivons $L^c = \bigoplus_{\lambda \in \Lambda} R$.
Comme $L^{c - 1}$ est libre de type fini, nous pouvons trouver une partie finie
$\Lambda' \subset \Lambda$ telle que $L^{c - 1} \to L^c$ se factorise
par $\bigoplus_{\lambda \in \Lambda'} R \subset L^c$. Considérons
l'application de complexes
$$
\pi :
L^\bullet
\longrightarrow
(\bigoplus\nolimits_{\lambda \in \Lambda \setminus \Lambda'} R)[-c]
$$
donnée par la projection sur les facteurs correspondant à
$\Lambda \setminus \Lambda'$ en degré $c$.
Notre hypothèse sur $K$ montre qu'après avoir éventuellement remplacé $\Lambda'$ par
une partie finie plus grande, nous pouvons supposer que $\pi \circ \varphi = 0$
dans $D(R)$. Soit $(L')^\bullet \subset L^\bullet$ le noyau de $\pi$.
Comme $\pi$ est surjective, nous obtenons une suite exacte courte de complexes,
qui donne un triangle distingué dans $D(R)$ (voir
Catégories dérivées, lemme \ref{derived-lemma-derived-canonical-delta-functor}).
Comme $\Hom_{D(R)}(K, -)$ est homologique (voir
Catégories dérivées, lemme \ref{derived-lemma-representable-homological})
et que $\pi \circ \varphi = 0$, il existe un morphisme
$\varphi' : K^\bullet \to (L')^\bullet$ dans $D(R)$ dont la
composée avec $(L')^\bullet \to L^\bullet$ est $\varphi$.
En prenant pour $\psi'$ la composée de $\psi$ avec
$(L')^\bullet \to L^\bullet$, nous obtenons une nouvelle factorisation.
Comme $(L')^\bullet$ coïncide avec $L^\bullet$ sauf en degré $c$
et que $(L')^c = \bigoplus_{\lambda \in \Lambda'} R$, l'étape de
récurrence est démontrée.

\medskip\noindent
La conclusion de la discussion du paragraphe précédent est que
$1_K : K \to K$ se factorise sous la forme
$$
K \xrightarrow{\varphi} L \xrightarrow{\psi} K
$$
dans $D(R)$, où $L$ peut être représenté par un complexe fini de
$R$-modules libres. En particulier, $L$ est parfait.
Remarquons que $e = \varphi \circ \psi \in \text{End}_{D(R)}(L)$
est un idempotent. D'après Catégories dérivées,
lemme \ref{derived-lemma-projectors-have-images-triangulated},
nous avons $L = \Ker(e) \oplus \Ker(1 - e)$.
L'application $\varphi : K \to L$ induit un isomorphisme avec
$\Ker(1 - e)$ dans $D(R)$. Nous concluons finalement que
$K$ est parfait par le lemme \ref{lemma-summands-perfect}.
\end{proof}

\begin{lemma}
\label{lemma-perfect-modulo-nilpotent-ideal}
Soit $R$ un anneau. Soit $I \subset R$ un idéal.
Soit $K$ un objet de $D(R)$. Supposons que
\begin{enumerate}
\item $K \otimes_R^\mathbf{L} R/I$ soit parfait dans $D(R/I)$ ;
\item $I$ soit un idéal nilpotent.
\end{enumerate}
Alors $K$ est parfait dans $D(R)$.
\end{lemma}

\begin{proof}
Choisissons un complexe fini $\overline{P}^\bullet$ de $R/I$-modules
projectifs de type fini représentant $K \otimes_R^\mathbf{L} R/I$, voir la
définition \ref{definition-perfect}. D'après le
lemme \ref{lemma-lift-complex-projectives},
il existe un complexe $P^\bullet$ de $R$-modules projectifs
représentant $K$ tel que $\overline{P}^\bullet = P^\bullet/IP^\bullet$.
Il résulte du lemme de Nakayama (Algèbre, lemme \ref{algebra-lemma-NAK})
que $P^\bullet$ est un complexe fini de $R$-modules
projectifs de type fini.
\end{proof}

\begin{lemma}
\label{lemma-perfect-modulo-two-ideals}
Soit $R$ un anneau. Soient $I, J \subset R$ des idéaux.
Soit $K$ un objet de $D(R)$. Supposons que
\begin{enumerate}
\item $K \otimes_R^\mathbf{L} R/I$ soit parfait dans $D(R/I)$ ;
\item $K \otimes_R^\mathbf{L} R/J$ soit parfait dans $D(R/J)$.
\end{enumerate}
Alors $K \otimes_R^\mathbf{L} R/IJ$ est parfait dans $D(R/IJ)$.
\end{lemma}

\begin{proof}
Il est clair que nous pouvons remplacer $R$ par $R/IJ$ et $K$ par
$K \otimes_R^\mathbf{L} R/IJ$. Alors $R \to R/(I \cap J)$ est
une surjection dont le noyau est de carré nul. Ainsi, d'après le
lemme \ref{lemma-perfect-modulo-nilpotent-ideal},
il suffit de démontrer que $K \otimes_R^\mathbf{L} R/(I \cap J)$ est
parfait. Nous pouvons donc supposer que $I \cap J = 0$.

\medskip\noindent
Nous démontrons le lemme dans le cas où $I \cap J = 0$. Tout d'abord, nous pouvons représenter $K$
par un complexe K-plat $K^\bullet$ dont tous les $K^n$ sont plats, voir le
lemme \ref{lemma-K-flat-resolution}. Nous obtenons alors une suite
exacte courte de complexes
$$
0 \to
K^\bullet \to K^\bullet/IK^\bullet \oplus K^\bullet/JK^\bullet \to
K^\bullet/(I + J)K^\bullet \to 0
$$
Remarquons que $K^\bullet/IK^\bullet$ représente $K \otimes^\mathbf{L}_R R/I$
par construction du produit tensoriel dérivé. Il en va de même de
$K^\bullet/JK^\bullet$ et $K^\bullet/(I + J)K^\bullet$.
Remarquons que $K^\bullet/(I + J)K^\bullet$ est un complexe parfait
de $R/(I + J)$-modules, voir le
lemme \ref{lemma-pull-perfect}.
Ainsi les complexes $K^\bullet/IK^\bullet$,
$K^\bullet/JK^\bullet$ et $K^\bullet/(I + J)K^\bullet$
n'ont qu'un nombre fini de groupes de cohomologie non nuls
(puisqu'un complexe parfait a une amplitude de Tor finie, voir le
lemme \ref{lemma-perfect}). La longue suite exacte de cohomologie associée à la suite
exacte courte de complexes affichée ci-dessus montre que $K \in D^b(R)$.
En particulier, nous pouvons supposer que $K^\bullet$ est un complexe de $R$-modules
libres borné supérieurement (voir
Catégories dérivées, lemme \ref{derived-lemma-subcategory-left-resolution}).

\medskip\noindent
Nous allons maintenant montrer que $K$ est parfait en utilisant le critère de la
proposition \ref{proposition-perfect-is-compact}. Soit donc
$E_j \in D(R)$ une famille d'objets paramétrée par un ensemble $J$.
Choisissons des complexes $E_j^\bullet$ à termes plats
représentant les $E_j$, voir par exemple le lemme \ref{lemma-K-flat-resolution}.
Il est clair que
$$
0 \to
E_j^\bullet \to
E_j^\bullet/IE_j^\bullet \oplus E_j^\bullet/JE_j^\bullet \to
E_j^\bullet/(I + J)E_j^\bullet \to 0
$$
est une suite exacte courte de complexes. En prenant les sommes directes, nous obtenons une
suite exacte courte analogue
$$
0 \to
\bigoplus E_j^\bullet \to
\bigoplus E_j^\bullet/IE_j^\bullet \oplus E_j^\bullet/JE_j^\bullet \to
\bigoplus E_j^\bullet/(I + J)E_j^\bullet \to 0
$$
(Remarquons que $- \otimes_R R/I$ commute aux sommes directes.)
Cette suite exacte courte détermine un triangle distingué dans $D(R)$, voir
Catégories dérivées, lemme \ref{derived-lemma-derived-canonical-delta-functor}.
Appliquons le foncteur homologique $\Hom_{D(R)}(K, -)$ (voir
Catégories dérivées, lemme \ref{derived-lemma-representable-homological})
pour obtenir un diagramme commutatif
$$
\xymatrix{
\bigoplus \Hom_{D(R)}(K^\bullet, E_j^\bullet/(I + J))[-1] \ar[r] \ar[d] &
\Hom_{D(R)}(K^\bullet, \bigoplus E_j^\bullet/(I + J))[-1] \ar[d] \\
\bigoplus \Hom_{D(R)}(K^\bullet, E_j^\bullet/I \oplus E_j^\bullet/J)[-1]
\ar[r] \ar[d] &
\Hom_{D(R)}(K^\bullet, \bigoplus E_j^\bullet/I \oplus E_j^\bullet/J)[-1]
\ar[d] \\
\bigoplus \Hom_{D(R)}(K^\bullet, E_j^\bullet) \ar[r] \ar[d] &
\Hom_{D(R)}(K^\bullet, \bigoplus E_j^\bullet) \ar[d] \\
\bigoplus \Hom_{D(R)}(K^\bullet, E_j^\bullet/I \oplus E_j^\bullet/J)
\ar[r] \ar[d] &
\Hom_{D(R)}(K^\bullet, \bigoplus E_j^\bullet/I \oplus E_j^\bullet/J)
\ar[d] \\
\bigoplus \Hom_{D(R)}(K^\bullet, E_j^\bullet/(I + J)) \ar[r] &
\Hom_{D(R)}(K^\bullet, \bigoplus E_j^\bullet/(I + J))
}
$$
à colonnes exactes. Il est clair que, pour tout complexe $E^\bullet$
de $R$-modules, nous avons
\begin{align*}
\Hom_{D(R)}(K^\bullet, E^\bullet/I) & =
\Hom_{K(R)}(K^\bullet, E^\bullet/I) \\
& =
\Hom_{K(R/I)}(K^\bullet/IK^\bullet, E^\bullet/I) \\
& =
\Hom_{D(R/I)}(K^\bullet/IK^\bullet, E^\bullet/I)
\end{align*}
et de même après quotient par $J$ ou $I + J$, voir
Catégories dérivées,
lemme \ref{derived-lemma-morphisms-from-projective-complex}.
Ainsi toutes les flèches horizontales, sauf peut-être celle du milieu,
sont des isomorphismes car les complexes
$K^\bullet/IK^\bullet$, $K^\bullet/JK^\bullet$, $K^\bullet/(I + J)K^\bullet$
sont des complexes parfaits de $R/I$-, $R/J$- et $R/(I + J)$-modules, voir la
proposition \ref{proposition-perfect-is-compact}.
Il résulte du $5$-lemme (Homologie, lemme \ref{homology-lemma-five-lemma})
que l'application médiane est un isomorphisme, et le lemme découle alors de la
proposition \ref{proposition-perfect-is-compact}.
\end{proof}






\section{Générateurs forts et anneaux réguliers}
\label{section-strong-generators-regular}

\noindent
Soit $R$ un anneau. Notons $D(R)_c$ la sous-catégorie triangulée pleine et saturée
de $D(R)$. Nous savons déjà que
$$
\langle R \rangle = D_{perf}(R) = D(R)_c
$$
Voir le lemme \ref{lemma-perfect-ring-classical-generator}
et la proposition \ref{proposition-perfect-is-compact}.
Il se trouve que si $R$ est régulier, alors $R$ est un générateur fort
(Catégories dérivées, définition \ref{derived-definition-generators}).

\begin{lemma}
\label{lemma-ghost-lemma}
\begin{reference}
\cite{Kelly}
\end{reference}
Soit $R$ un anneau. Soit $n \geq 1$. Soit $K \in \langle R \rangle_n$, avec
les notations de Catégories dérivées, section \ref{derived-section-generators}.
Considérons des applications
$$
K \xrightarrow{f_1} K_1 \xrightarrow{f_2} K_2
\xrightarrow{f_3} \ldots \xrightarrow{f_n} K_n
$$
dans $D(R)$. Si $H^i(f_j) = 0$ pour tous $i, j$, alors
$f_n \circ \ldots \circ f_1 = 0$.
\end{lemma}

\begin{proof}
Si $n = 1$, alors $K$ est un facteur direct dans $D(R)$ d'un complexe borné
$P^\bullet$ dont les termes sont des $R$-modules libres de type fini et dont les différentielles
sont nulles. Il suffit donc de montrer que tout morphisme $f : P^\bullet \to K_1$
dans $D(R)$ tel que $H^i(f) = 0$ pour tout $i$ est nul. Comme $P^\bullet$ est
une somme directe finie $P^\bullet = \bigoplus R[m_j]$, il suffit de montrer que
tout morphisme $g : R[m] \to K_1$ tel que $H^{-m}(g) = 0$ dans $D(R)$ est nul.
Cela résulte de l'égalité $\Hom_{D(R)}(R[-m], K) = H^m(K)$.

\medskip\noindent
Pour $n > 1$, nous procédons par récurrence sur $n$. En effet, nous savons que $K$
est un facteur direct dans $D(R)$ d'un objet $P$ qui s'insère dans un triangle distingué
$$
P' \xrightarrow{i} P \xrightarrow{p} P'' \to P'[1]
$$
avec $P' \in \langle R \rangle_1$ et $P'' \in \langle R \rangle_{n - 1}$.
Comme ci-dessus, nous pouvons remplacer $K$ par $P$ et supposer que nous avons
$$
P \xrightarrow{f_1} K_1 \xrightarrow{f_2} K_2
\xrightarrow{f_3} \ldots \xrightarrow{f_n} K_n
$$
dans $D(R)$, les $f_j$ induisant l'application nulle en cohomologie.
Le cas $n = 1$ montre que la composée $f_1 \circ i$ est nulle.
Ainsi, d'après Catégories dérivées, lemme \ref{derived-lemma-representable-homological},
il existe un morphisme $h : P'' \to K_1$ tel que $f_1 = h \circ p$.
Remarquons que $f_2 \circ h$ est nul en cohomologie.
Par récurrence, nous obtenons donc $f_n \circ \ldots \circ f_2 \circ h = 0$,
ce qui implique $f_n \circ \ldots \circ f_1 =
f_n \circ \ldots \circ f_2 \circ h \circ p = 0$, comme souhaité.
\end{proof}

\begin{lemma}
\label{lemma-not-regular-not-strong}
Soit $R$ un anneau noethérien. Si $R$ est
un générateur fort de $D_{perf}(R)$, alors $R$ est régulier
de dimension finie.
\end{lemma}

\begin{proof}
Supposons $D_{perf}(R) = \langle R \rangle_n$ pour un certain $n \geq 1$.
Pour tout $R$-module de type fini $M$, nous pouvons choisir un complexe
$$
P = (
P^{-n - 1} \xrightarrow{d^{-n - 1}}
P^{-n} \xrightarrow{d^{-n}}
P^{-n + 1} \xrightarrow{d^1}
\ldots \xrightarrow{d^{-1}} P^0)
$$
de $R$-modules libres de type fini avec $H^i(P) = 0$ pour $i = -n, \ldots, - 1$
et $M \cong \Coker(d^{-1})$. Remarquons que $P$ appartient à $D_{perf}(R)$.
Pour tout $R$-module $N$, nous pouvons calculer $\Ext^n_R(M, N)$ à l'aide de
la résolution libre finie $P$ de $M$, voir
Algèbre, section \ref{algebra-section-ext}, et comparer avec
Catégories dérivées, section \ref{derived-section-ext}.
En particulier, la suite ci-dessus définit un élément
$$
\xi \in \Ext^n_R(\Coker(d^{-1}), \Coker(d^{-n - 1})) =
\Ext^n_R(M, \Coker(d^{-n - 1}))
$$
et, pour tout élément $\overline{\xi}$ de $\Ext^n_R(M, N)$,
il existe une application de $R$-modules $\varphi : \Coker(d^{-n - 1}) \to N$ telle
que $\varphi$ envoie $\xi$ sur $\overline{\xi}$.
Pour $j = 1, \ldots, n - 1$, considérons les complexes
$$
K_j = (\Coker(d^{-n - 1}) \to P^{-n + 1} \to \ldots \to P^{-j})
$$
avec $\Coker(d^{-n - 1})$ en degré $-n$ et $P^t$ en degré $t$.
Posons aussi $K_n = \Coker(d^{-n - 1})[n]$. Nous avons alors des applications
$$
P \to K_1 \to K_2 \to \ldots \to K_n
$$
qui induisent des applications nulles en cohomologie. Comme
$P \in D_{perf}(R) = \langle R \rangle_n$, le lemme \ref{lemma-ghost-lemma}
montre que la composée de ces applications est nulle dans $D(R)$. Puisque
$\Hom_{D(R)}(P, K_n) = \Hom_{K(R)}(P, K_n)$ d'après
Catégories dérivées, lemme \ref{derived-lemma-morphisms-from-projective-complex},
nous concluons que $\xi = 0$. Ainsi $\Ext^n_R(M, N) = 0$ pour tout $R$-module
$N$, d'après la discussion ci-dessus.
Il s'ensuit que $M$ est de dimension projective $\leq n - 1$ d'après
Algèbre, lemme \ref{algebra-lemma-projective-dimension-ext}.
Comme cela vaut pour tout $R$-module de type fini $M$, nous concluons que
$R$ est de dimension globale finie, voir
Algèbre, lemme \ref{algebra-lemma-finite-gl-dim}.
Nous concluons finalement par Algèbre, lemme
\ref{algebra-lemma-finite-gl-dim-finite-dim-regular}.
\end{proof}

\begin{lemma}
\label{lemma-ext-regular}
Soit $R$ un anneau noethérien régulier de dimension $d < \infty$.
Soient $K, L \in D^-(R)$. Supposons qu'il existe un $k$ tel que
$H^i(K) = 0$ pour $i \leq k$ et $H^i(L) = 0$ pour $i \geq k - d + 1$.
Alors $\Hom_{D(R)}(K, L) = 0$.
\end{lemma}

\begin{proof}
Soit $K^\bullet$ un complexe borné supérieurement représentant $K$, disons
$K^i = 0$ pour $i \geq n + 1$. Après avoir remplacé $K^\bullet$ par
$\tau_{\geq k + 1}K^\bullet$, nous pouvons supposer $K^i = 0$ pour $i \leq k$.
Nous pouvons alors utiliser le triangle distingué
$$
K^n[-n] \to K^\bullet \to \sigma_{\leq n - 1}K^\bullet
$$
pour voir qu'il suffit de démontrer le lemme pour $K^n[-n]$ et
$\sigma_{\leq n - 1}K^\bullet$. Par récurrence sur $n$, nous concluons
qu'il suffit de démontrer le lemme lorsque $K$ est représenté par
le complexe $M[-m]$ pour un $R$-module $M$ et un $m \geq k + 1$.
Comme $R$ est de dimension globale $d$ d'après Algèbre, lemme
\ref{algebra-lemma-finite-gl-dim-finite-dim-regular},
nous voyons que $M$ admet une résolution projective
$0 \to P_d \to \ldots \to P_0 \to M \to 0$.
Le complexe $P^\bullet$ dont $P_i$ est placé en degré $m - i$
est alors un complexe borné de projectifs représentant $M[-m]$.
D'autre part, nous pouvons choisir un complexe $L^\bullet$ représentant
$L$ avec $L^i = 0$ pour $i \geq k - d  + 1$.
Toute application de complexes $P^\bullet \to L^\bullet$ est donc nulle.
Le lemme en résulte par Catégories dérivées, lemme
\ref{derived-lemma-morphisms-from-projective-complex}.
\end{proof}

\begin{lemma}
\label{lemma-split-complex-regular}
Soit $R$ un anneau noethérien régulier de dimension $1 \leq d < \infty$.
Soit $K \in D(R)$ parfait et soit $k \in \mathbf{Z}$ tel que
$H^i(K) = 0$ pour $i = k - d + 2, \ldots, k$ (condition vide si $d = 1$).
Alors $K = \tau_{\leq k - d + 1}K \oplus \tau_{\geq k + 1}K$.
\end{lemma}

\begin{proof}
L'annulation de la cohomologie montre que nous avons un triangle distingué
$$
\tau_{\leq k - d + 1}K \to K \to \tau_{\geq k + 1}K \to
(\tau_{\leq k - d + 1}K)[1]
$$
Par Catégories dérivées, lemme \ref{derived-lemma-split}, il suffit de montrer que
la troisième flèche est nulle. Il suffit donc de montrer que
$\Hom_{D(R)}(\tau_{\geq k + 1}K, (\tau_{\leq k - d + 1}K)[1]) = 0$
ce qui résulte du lemme \ref{lemma-ext-regular}.
\end{proof}

\begin{lemma}
\label{lemma-regular-strong-generator}
Soit $R$ un anneau noethérien régulier de dimension finie.
Alors $R$ est un générateur fort de la sous-catégorie pleine
$D_{perf}(R) \subset D(R)$ des objets parfaits.
\end{lemma}

\begin{proof}
Nous utiliserons qu'un objet $K$ de $D(R)$ est parfait si et seulement si
$K$ est borné et possède des modules de cohomologie de type fini, voir le
lemme \ref{lemma-regular-perfect}.
Les générateurs forts des catégories triangulées sont définis dans
Catégories dérivées, définition \ref{derived-definition-generators}.
Soit $d = \dim(R)$.

\medskip\noindent
Soit $K \in D_{perf}(R)$. Nous allons montrer que $K \in \langle R \rangle_{d + 1}$.
D'après Algèbre, lemme \ref{algebra-lemma-finite-gl-dim-finite-dim-regular},
tout $R$-module de type fini est de dimension projective $\leq d$.
Nous allons montrer par récurrence sur $0 \leq i \leq d$ que,
si $H^n(K)$ est de dimension projective $\leq i$ pour tout
$n \in \mathbf{Z}$, alors $K$ appartient à $\langle R \rangle_{i + 1}$.

\medskip\noindent
Cas initial $i = 0$. Dans ce cas, $H^n(K)$ est un $R$-module de type fini
de dimension projective $0$. Autrement dit, chaque module de cohomologie
est un $R$-module projectif. Ainsi $\Ext^i_R(H^n(K), H^m(K)) = 0$
pour tout $i > 0$ et tous $m, n \in \mathbf{Z}$. D'après
Catégories dérivées, lemme \ref{derived-lemma-ext-2-zero-pre},
nous obtenons que $K$ est isomorphe à la somme directe des décalages de ses
modules de cohomologie. Comme chacun de ceux-ci est un $R$-module
projectif de type fini, il est facteur direct d'une somme directe
de copies de $R$. Par définition, nous voyons donc que $K$
appartient à $\langle R \rangle_1$.

\medskip\noindent
Étape de récurrence. Supposons l'assertion vraie pour $i < d$ et soit
$K \in D_{perf}(R)$ tel que $H^n(K)$ soit de dimension
projective $\leq i + 1$ pour tout $n \in \mathbf{Z}$.
Choisissons $a \leq b$ tels que $H^n(K)$ soit nul pour $n \not \in [a, b]$.
Pour chaque $n \in [a, b]$, choisissons une surjection $F^n \to H^n(K)$,
où $F^n$ est un $R$-module libre de type fini. Comme $F^n$ est projectif,
nous pouvons relever $F^n \to H^n(K)$ en une application $F^n[-n] \to K$ dans $D(R)$
(un petit détail est omis). Nous obtenons ainsi un morphisme
$\bigoplus_{a \leq n \leq b} F^n[-n] \to K$
qui est surjectif sur les modules de cohomologie.
Choisissons un triangle distingué
$$
K' \to \bigoplus\nolimits_{a \leq n \leq b} F^n[-n] \to K \to K'[1]
$$
dans $D(R)$. Bien entendu, l'objet $K'$ est borné et possède des
modules de cohomologie de type fini. La longue suite exacte de cohomologie
se décompose en suites exactes courtes
$$
0 \to H^n(K') \to F^n \to H^n(K) \to 0
$$
grâce à nos choix. D'après
Algèbre, lemme \ref{algebra-lemma-exact-sequence-projective-dimension},
nous voyons que la dimension projective de $H^n(K')$ est $\leq \max(0, i)$.
Ainsi $K' \in \langle R \rangle_{i + 1}$. Par définition, cela signifie
que $K$ appartient à $\langle R \rangle_{i + 1 + 1}$, comme souhaité.
\end{proof}

\begin{proposition}
\label{proposition-regular-strong-generator}
Soit $R$ un anneau noethérien. Les assertions suivantes sont équivalentes
\begin{enumerate}
\item $R$ est régulier de dimension finie ;
\item $D_{perf}(R)$ admet un générateur fort ;
\item $R$ est un générateur fort de $D_{perf}(R)$.
\end{enumerate}
\end{proposition}

\begin{proof}
C'est une conséquence formelle des
lemmes \ref{lemma-perfect-ring-classical-generator},
\ref{lemma-not-regular-not-strong} et
\ref{lemma-regular-strong-generator},
ainsi que de Catégories dérivées, lemme
\ref{derived-lemma-classical-generator-strong-generator}.
\end{proof}







\section{Modules de présentation finie relative}
\label{section-relative-finite-presentation}

\noindent
Soit $R$ un anneau. Soit $A \to B$ un morphisme fini de $R$-algèbres de type fini.
Soit $M$ un $B$-module de type fini. Dans ce cas, il n'est {\bf pas vrai} que
$$
M\text{ est de présentation finie sur }B
\Leftrightarrow
M\text{ est de présentation finie sur }A
$$
Un contre-exemple est donné par $R = k[x_1, x_2, x_3, \ldots]$, $A = R$, $B = R/(x_i)$
et $M = B$. Pour « corriger » cela, nous introduisons une notion relative de
présentation finie.

\begin{lemma}
\label{lemma-relatively-finitely-presented}
Soit $R \to A$ un morphisme d'anneaux de type fini.
Soit $M$ un $A$-module.
Les assertions suivantes sont équivalentes
\begin{enumerate}
\item pour une présentation $\alpha : R[x_1, \ldots, x_n] \to A$, le module
$M$ est un $R[x_1, \ldots, x_n]$-module de présentation finie ;
\item pour toute présentation $\alpha : R[x_1, \ldots, x_n] \to A$, le module
$M$ est un $R[x_1, \ldots, x_n]$-module de présentation finie ;
\item pour toute surjection $A' \to A$ où $A'$ est une $R$-algèbre de
présentation finie, le module $M$ est de présentation finie comme $A'$-module.
\end{enumerate}
Dans ce cas, $M$ est un $A$-module de présentation finie.
\end{lemma}

\begin{proof}
Si $\alpha : R[x_1, \ldots, x_n] \to A$ et
$\beta : R[y_1, \ldots, y_m] \to A$ sont des présentations,
choisissons $f_j \in R[x_1, \ldots, x_n]$ avec $\alpha(f_j) = \beta(y_j)$
et $g_i \in R[y_1, \ldots, y_m]$ avec $\beta(g_i) = \alpha(x_i)$.
Nous obtenons alors un diagramme commutatif
$$
\xymatrix{
R[x_1, \ldots, x_n, y_1, \ldots, y_m]
\ar[d]^{x_i \mapsto g_i} \ar[rr]_-{y_j \mapsto f_j} & &
R[x_1, \ldots, x_n] \ar[d] \\
R[y_1, \ldots, y_m] \ar[rr] & & A
}
$$
L'équivalence de (1) et (2) résulte donc de l'application des
lemmes d'algèbre \ref{algebra-lemma-finitely-presented-over-subring} et
\ref{algebra-lemma-finite-finitely-presented-extension}.
L'équivalence de (2) et (3) résulte du choix d'une présentation
$A' = R[x_1, \ldots, x_n]/(f_1, \ldots, f_m)$ et de l'utilisation du
lemme d'algèbre \ref{algebra-lemma-finite-finitely-presented-extension}
pour montrer que $M$ est de présentation finie comme $A'$-module si et seulement si
$M$ est de présentation finie comme $R[x_1, \ldots, x_n]$-module.
\end{proof}

\begin{definition}
\label{definition-relatively-finitely-presented}
Soit $R \to A$ un morphisme d'anneaux de type fini. Soit $M$ un $A$-module.
Nous disons que $M$ est un $A$-module {\it de présentation finie relativement à $R$}
si les conditions équivalentes du
lemme \ref{lemma-relatively-finitely-presented}
sont satisfaites.
\end{definition}

\noindent
Remarquons que si $R \to A$ est de présentation finie, alors $M$ est un
$A$-module de présentation finie relativement à $R$ si et seulement si $M$
est un $A$-module de présentation finie. Il est tout aussi clair que $A$, considéré
comme $A$-module, est de présentation finie relativement à $R$ si et seulement si
$A$ est de présentation finie sur $R$. Si $R$ est noethérien, cette notion
ne présente aucun intérêt. Nous pouvons maintenant formuler le résultat recherché.

\begin{lemma}
\label{lemma-finite-extension}
Soit $R$ un anneau. Soit $A \to B$ un morphisme fini de $R$-algèbres de type fini.
Soit $M$ un $B$-module. Alors
$M$ est un $A$-module de présentation finie relativement à $R$
si et seulement si
$M$ est un $B$-module de présentation finie relativement à $R$.
\end{lemma}

\begin{proof}
Choisissons une surjection $R[x_1, \ldots, x_n] \to A$.
Choisissons $y_1, \ldots, y_m \in B$ engendrant $B$ sur $A$.
Comme $A \to B$ est fini, chaque $y_i$ satisfait une équation unitaire à
coefficients dans $A$. Nous pouvons donc trouver des polynômes unitaires
$P_j(T) \in R[x_1, \ldots, x_n][T]$ tels que $P_j(y_j) = 0$ dans $B$.
Nous obtenons alors un diagramme commutatif
$$
\xymatrix{
R[x_1, \ldots, x_n] \ar[d] \ar[r] &
R[x_1, \ldots, x_n, y_1, \ldots, y_m]/(P_j(y_j)) \ar[d] \\
A \ar[r] & B
}
$$
Comme la flèche supérieure est un morphisme d'anneaux fini et de présentation finie,
nous concluons par le
lemme d'algèbre \ref{algebra-lemma-finite-finitely-presented-extension}
et la définition.
\end{proof}

\noindent
Grâce à ce résultat, nous voyons que la notion relative a un sens
et se comporte bien vis-à-vis des morphismes finis d'anneaux de type fini
sur $R$. Elle est aussi stable par localisation et par
changement de base, et se « recolle » bien.

\begin{lemma}
\label{lemma-localize-relative-finite-presentation}
Soient $R$ un anneau, $f \in R$ un élément, $R_f \to A$ un morphisme d'anneaux de type fini,
$g \in A$ et $M$ un $A$-module. Si $M$ est de présentation finie relativement
à $R_f$, alors $M_g$ est un $A_g$-module de présentation finie relativement
à $R$.
\end{lemma}

\begin{proof}
Choisissons une présentation $R_f[x_1, \ldots, x_n] \to A$. Écrivons
$R_f = R[x_0]/(fx_0 - 1)$. Considérons la présentation
$R[x_0, x_1, \ldots, x_n, x_{n + 1}] \to A_g$ qui prolonge le morphisme
donné, envoie $x_0$ sur l'image de $1/f$ et $x_{n + 1}$ sur $1/g$.
Choisissons $g' \in R[x_0, x_1, \ldots, x_n]$ s'envoyant sur $g$ (cela est
possible). Supposons que
$$
R_f[x_1, \ldots, x_n]^{\oplus s} \to
R_f[x_1, \ldots, x_n]^{\oplus t} \to M \to 0
$$
soit une présentation de $M$ donnée par une matrice $(h_{ij})$. Choisissons
$h'_{ij} \in R[x_0, x_1, \ldots, x_n]$ s'envoyant sur $h_{ij}$.
Alors
$$
R[x_0, x_1, \ldots, x_n, x_{n + 1}]^{\oplus s + 2t} \to
R[x_0, x_1, \ldots, x_n, x_{n + 1}]^{\oplus t} \to M_g \to 0
$$
est une présentation de $M_f$.
Ici, la matrice $t \times (s + 2t)$ définissant l'application possède un premier
bloc $t \times s$ formé de la matrice $h'_{ij}$, un deuxième
bloc $t \times t$ égal à $(x_0f - )I_t$, et un troisième bloc
égal à $(x_{n + 1}g' - 1)I_t$.
\end{proof}

\begin{lemma}
\label{lemma-base-change-relative-finite-presentation}
Soit $R \to A$ un morphisme d'anneaux de type fini. Soit $M$ un $A$-module de
présentation finie relativement à $R$. Pour tout morphisme d'anneaux $R \to R'$, le
$A \otimes_R R'$-module
$$
M \otimes_A A' = M \otimes_R R'
$$
est de présentation finie relativement à $R'$.
\end{lemma}

\begin{proof}
Choisissons une surjection $R[x_1, \ldots, x_n] \to A$. Choisissons une présentation
$$
R[x_1, \ldots, x_n]^{\oplus s} \to
R[x_1, \ldots, x_n]^{\oplus t} \to M \to 0
$$
Alors
$$
R'[x_1, \ldots, x_n]^{\oplus s} \to
R'[x_1, \ldots, x_n]^{\oplus t} \to M \otimes_R R' \to 0
$$
est une présentation du changement de base, ce qui conclut.
\end{proof}

\begin{lemma}
\label{lemma-pull-relative-finite-presentation}
Soit $R \to A$ un morphisme d'anneaux de type fini.
Soit $M$ un $A$-module de présentation finie relativement à $R$.
Soit $A \to A'$ un morphisme d'anneaux de présentation finie.
Le $A'$-module $M \otimes_A A'$ est de présentation finie relativement à $R$.
\end{lemma}

\begin{proof}
Choisissons une surjection $R[x_1, \ldots, x_n] \to A$. Choisissons une présentation
$A' = A[y_1, \ldots, y_m]/(g_1, \ldots, g_l)$.
Choisissons $g'_i \in R[x_1, \ldots, x_n, y_1, \ldots, y_m]$ s'envoyant sur $g_i$.
Disons que
$$
R[x_1, \ldots, x_n]^{\oplus s} \to
R[x_1, \ldots, x_n]^{\oplus t} \to M \to 0
$$
est une présentation de $M$ donnée par une matrice $(h_{ij})$.
Alors
$$
R[x_1, \ldots, x_n, y_1, \ldots, y_m]^{\oplus s + tl} \to
R[x_0, x_1, \ldots, x_n, y_1, \ldots, y_m]^{\oplus t} \to M \otimes_A A' \to 0
$$
est une présentation de $M \otimes_A A'$.
Ici, la matrice $t \times (s + lt)$ définissant l'application possède un premier
bloc $t \times s$ formé de la matrice $h_{ij}$, suivi de
$l$ blocs de taille $t \times t$ égaux à $g'_iI_t$.
\end{proof}

\begin{lemma}
\label{lemma-composition-relative-finite-presentation}
Soient $R \to A \to B$ des morphismes d'anneaux de type fini. Soit $M$ un $B$-module.
Si $M$ est de présentation finie relativement à $A$ et si $A$ est de présentation finie
sur $R$, alors $M$ est de présentation finie relativement à $R$.
\end{lemma}

\begin{proof}
Choisissons une surjection $A[x_1, \ldots, x_n] \to B$.
Choisissons une présentation
$$
A[x_1, \ldots, x_n]^{\oplus s} \to
A[x_1, \ldots, x_n]^{\oplus t} \to M \to 0
$$
donnée par une matrice $(h_{ij})$. Choisissons une présentation
$$
A = R[y_1, \ldots, y_m]/(g_1, \ldots, g_u).
$$
Choisissons $h'_{ij} \in R[y_1, \ldots, y_m, x_1, \ldots, x_n]$
s'envoyant sur $h_{ij}$. Nous obtenons alors la présentation
$$
R[y_1, \ldots, y_m, x_1, \ldots, x_n]^{\oplus s + tu} \to
R[y_1, \ldots, y_m, x_1, \ldots, x_n]^{\oplus t} \to M \to 0
$$
où la matrice $t \times (s + tu)$ est constituée d'un premier bloc $t \times s$
formé des $h'_{ij}$, suivi de $u$ blocs de taille $t \times t$ donnés
par $g_iI_t$, $i = 1, \ldots, u$.
\end{proof}

\begin{lemma}
\label{lemma-glue-relative-finite-presentation}
Soit $R \to A$ un morphisme d'anneaux de type fini. Soit $M$ un $A$-module.
Soient $f_1, \ldots, f_r \in A$ engendrant l'idéal unité.
Les assertions suivantes sont équivalentes
\begin{enumerate}
\item chaque $M_{f_i}$ est de présentation finie relativement à $R$ ;
\item $M$ est de présentation finie relativement à $R$.
\end{enumerate}
\end{lemma}

\begin{proof}
L'implication (2) $\Rightarrow$ (1) se trouve dans le
lemme \ref{lemma-localize-relative-finite-presentation}.
Supposons (1). Écrivons $1 = \sum f_ig_i$ dans $A$.
Choisissons une surjection
$R[x_1, \ldots, x_n, y_1, \ldots, y_r, z_1, \ldots, z_r] \to A$.
telle que $y_i$ s'envoie sur $f_i$ et $z_i$ sur $g_i$. Nous
voyons alors qu'il existe une surjection
$$
P = R[x_1, \ldots, x_n, y_1, \ldots, y_r, z_1, \ldots, z_r]/(\sum y_iz_i - 1)
\longrightarrow
A.
$$
D'après le
lemme \ref{lemma-relatively-finitely-presented},
$M_{f_i}$ est un $A_{f_i}$-module de présentation finie ; par conséquent, d'après
Algèbre, lemme \ref{algebra-lemma-cover},
$M$ est un $A$-module de présentation finie.
Ainsi $M$ est un $P$-module de type fini (avec $P$ comme ci-dessus).
Choisissons une surjection $P^{\oplus t} \to M$.
Nous devons montrer que le noyau $K$ de cette application est un
$P$-module de type fini. Comme $P_{y_i}$ se surjecte sur
$A_{f_i}$, le
lemme \ref{lemma-relatively-finitely-presented}
et
Algèbre, lemme \ref{algebra-lemma-extension},
montrent que la localisation $K_{y_i}$ est un
$P_{y_i}$-module de type fini. Choisissons des éléments
$k_{i, j} \in K$, $i = 1, \ldots, r$, $j = 1, \ldots, s_i$, tels
que les images des $k_{i, j}$ engendrent $K_{y_i}$.
Soit $K' \subset K$ le $P$-module
engendré par les éléments $k_{i, j}$. Alors $K/K'$ est un module
dont la localisation en $y_i$ est nulle pour tout $i$. Comme $(y_1, \ldots, y_r) = P$,
nous voyons que $K/K' = 0$, comme souhaité.
\end{proof}

\begin{lemma}
\label{lemma-ses-relatively-finite-presentation}
Soit $R \to A$ un morphisme d'anneaux de type fini. Soit $0 \to M' \to M \to M'' \to 0$
une suite exacte courte de $A$-modules.
\begin{enumerate}
\item Si $M', M''$ sont de présentation finie relativement à $R$, alors $M$ l'est aussi.
\item Si $M'$ est un $A$-module de type fini et si $M$ est de présentation finie
relativement à $R$, alors $M''$ est de présentation finie relativement à $R$.
\end{enumerate}
\end{lemma}

\begin{proof}
Cela résulte immédiatement d'
Algèbre, lemme \ref{algebra-lemma-extension}.
\end{proof}

\begin{lemma}
\label{lemma-sum-relatively-finite-presentation}
Soit $R \to A$ un morphisme d'anneaux de type fini.
Soient $M, M'$ des $A$-modules. Si $M \oplus M'$ est de
présentation finie relativement à $R$, alors $M$ et $M'$ le sont aussi.
\end{lemma}

\begin{proof}
La démonstration est omise.
\end{proof}






\section{Modules pseudo-cohérents relativement à la base}
\label{section-relative-pseudo-coherent}

\noindent
Cette section est l'analogue de la
section \ref{section-relative-finite-presentation}
pour la pseudo-cohérence.

\begin{lemma}
\label{lemma-pull-push}
Soit $R$ un anneau. Soit $K^\bullet$ un complexe de $R$-modules.
Considérons le morphisme de $R$-algèbres $R[x] \to R$ qui envoie $x$ sur zéro.
Alors
$$
K^\bullet \otimes_{R[x]}^{\mathbf{L}} R \cong K^\bullet \oplus K^\bullet[1]
$$
dans $D(R)$.
\end{lemma}

\begin{proof}
Choisissons une résolution K-plate $P^\bullet \to K^\bullet$ sur $R$
telle que $P^n$ soit un $R$-module plat pour tout $n$, voir le
lemme \ref{lemma-K-flat-resolution}. Alors $P^\bullet \otimes_R R[x]$
est un complexe K-plat de $R[x]$-modules dont les termes sont des $R[x]$-modules plats, voir le
lemme \ref{lemma-base-change-K-flat} et
Algèbre, lemme \ref{algebra-lemma-flat-base-change}.
En particulier, $x : P^n \otimes_R R[x] \to P^n \otimes_R R[x]$
est injectif de conoyau isomorphe à $P^n$. Ainsi
$$
P^\bullet \otimes_R R[x] \xrightarrow{x} P^\bullet \otimes_R R[x]
$$
est un bicomplexe de $R[x]$-modules dont le complexe total associé
est quasi-isomorphe à $P^\bullet$, et donc à $K^\bullet$.
De plus, ce complexe total associé est un complexe K-plat
de $R[x]$-modules, par exemple d'après le
lemme \ref{lemma-tensor-product-K-flat} ou le
lemme \ref{lemma-K-flat-two-out-of-three}.
Par conséquent
\begin{align*}
K^\bullet \otimes_{R[x]}^{\mathbf{L}} R
& \cong
\text{Tot}(P^\bullet \otimes_R R[x] \xrightarrow{x} P^\bullet \otimes_R R[x])
\otimes_{R[x]} R =
\text{Tot}(P^\bullet \xrightarrow{0} P^\bullet) \\
& = P^\bullet \oplus P^\bullet[1] \cong K^\bullet \oplus K^\bullet[1]
\end{align*}
comme souhaité.
\end{proof}

\begin{lemma}
\label{lemma-add-variable-pseudo-coherent}
Soient $R$ un anneau et $K^\bullet$ un complexe de $R$-modules.
Soit $m \in \mathbf{Z}$. Considérons le morphisme de $R$-algèbres $R[x] \to R$
qui envoie $x$ sur zéro. Alors $K^\bullet$ est $m$-pseudo-cohérent comme
complexe de $R$-modules si et seulement si $K^\bullet$ est $m$-pseudo-cohérent
comme complexe de $R[x]$-modules.
\end{lemma}

\begin{proof}
C'est un cas particulier du
lemme \ref{lemma-finite-push-pseudo-coherent}.
Nous en donnons aussi la démonstration suivante.

\medskip\noindent
Remarquons que $0 \to R[x] \to R[x] \to R \to 0$ est exacte. Ainsi $R$ est
pseudo-cohérent comme $R[x]$-module. Une implication du lemme résulte donc du
lemme \ref{lemma-finite-push-pseudo-coherent}.
Pour démontrer l'autre implication, supposons que $K^\bullet$ soit
$m$-pseudo-cohérent comme complexe de $R[x]$-modules. D'après le
lemme \ref{lemma-pull-pseudo-coherent},
$K^\bullet \otimes^{\mathbf{L}}_{R[x]} R$ est
$m$-pseudo-cohérent comme complexe de $R$-modules. D'après le
lemme \ref{lemma-pull-push},
$K^\bullet \oplus K^\bullet[1]$ est $m$-pseudo-cohérent
comme complexe de $R$-modules.
Finalement, le lemme \ref{lemma-summands-pseudo-coherent}
montre que $K^\bullet$ est $m$-pseudo-cohérent
comme complexe de $R$-modules.
\end{proof}

\begin{lemma}
\label{lemma-relatively-pseudo-coherent}
Soit $R \to A$ un morphisme d'anneaux de type fini.
Soit $K^\bullet$ un complexe de $A$-modules.
Soit $m \in \mathbf{Z}$.
Les assertions suivantes sont équivalentes
\begin{enumerate}
\item pour une présentation $\alpha : R[x_1, \ldots, x_n] \to A$, le complexe
$K^\bullet$ est un complexe $m$-pseudo-cohérent de
$R[x_1, \ldots, x_n]$-modules ;
\item pour toute présentation $\alpha : R[x_1, \ldots, x_n] \to A$, le complexe
$K^\bullet$ est un complexe $m$-pseudo-cohérent de
$R[x_1, \ldots, x_n]$-modules.
\end{enumerate}
En particulier, la même équivalence vaut pour la pseudo-cohérence.
\end{lemma}

\begin{proof}
Si $\alpha : R[x_1, \ldots, x_n] \to A$ et
$\beta : R[y_1, \ldots, y_m] \to A$ sont des présentations,
choisissons $f_j \in R[x_1, \ldots, x_n]$ avec $\alpha(f_j) = \beta(y_j)$
et $g_i \in R[y_1, \ldots, y_m]$ avec $\beta(g_i) = \alpha(x_i)$.
Nous obtenons alors un diagramme commutatif
$$
\xymatrix{
R[x_1, \ldots, x_n, y_1, \ldots, y_m]
\ar[d]^{x_i \mapsto g_i} \ar[rr]_-{y_j \mapsto f_j} & &
R[x_1, \ldots, x_n] \ar[d] \\
R[y_1, \ldots, y_m] \ar[rr] & & A
}
$$
Après un changement de coordonnées, l'homomorphisme d'anneaux
$R[x_1, \ldots, x_n, y_1, \ldots, y_m] \to R[x_1, \ldots, x_n]$
est isomorphe à l'homomorphisme d'anneaux qui envoie
chaque $y_i$ sur zéro. Il en va de même pour la flèche verticale gauche du
diagramme. Ainsi, par récurrence sur le nombre de variables, le lemme résulte du
lemme \ref{lemma-add-variable-pseudo-coherent}.
Le cas pseudo-cohérent en résulte avec le
lemme \ref{lemma-pseudo-coherent}.
\end{proof}

\begin{definition}
\label{definition-relatively-pseudo-coherent}
Soit $R \to A$ un morphisme d'anneaux de type fini.
Soit $K^\bullet$ un complexe de $A$-modules.
Soit $M$ un $A$-module.
Soit $m \in \mathbf{Z}$.
\begin{enumerate}
\item Nous disons que $K^\bullet$ est {\it $m$-pseudo-cohérent relativement à $R$}
si les conditions équivalentes du
lemme \ref{lemma-relatively-pseudo-coherent}
sont satisfaites.
\item Nous disons que $K^\bullet$ est {\it pseudo-cohérent relativement à $R$}
si $K^\bullet$ est $m$-pseudo-cohérent relativement à $R$ pour tout
$m \in \mathbf{Z}$.
\item Nous disons que $M$ est {\it $m$-pseudo-cohérent relativement à $R$}
si $M[0]$ est $m$-pseudo-cohérent relativement à $R$.
\item Nous disons que $M$ est {\it pseudo-cohérent relativement à $R$}
si $M[0]$ est pseudo-cohérent relativement à $R$.
\end{enumerate}
\end{definition}

\noindent
La partie (2) signifie que $K^\bullet$ est pseudo-cohérent comme complexe
de $R[x_1, \ldots, x_n]$-modules pour toute surjection
$R[y_1, \ldots, y_m] \to A$, voir le
lemme \ref{lemma-pseudo-coherent}.
Cette définition possède la propriété agréable suivante.

\begin{lemma}
\label{lemma-finite-extension-pseudo-coherent}
Soit $R$ un anneau. Soit $A \to B$ un morphisme fini de $R$-algèbres de type fini.
Soit $m \in \mathbf{Z}$. Soit $K^\bullet$ un complexe de $B$-modules.
Alors $K^\bullet$ est $m$-pseudo-cohérent (resp.\ pseudo-cohérent)
relativement à $R$ si et seulement si $K^\bullet$, considéré comme complexe de $A$-modules,
est $m$-pseudo-cohérent (resp.\ pseudo-cohérent) relativement à $R$.
\end{lemma}

\begin{proof}
Choisissons une surjection $R[x_1, \ldots, x_n] \to A$.
Choisissons $y_1, \ldots, y_m \in B$ engendrant $B$ sur $A$.
Comme $A \to B$ est fini, chaque $y_i$ satisfait une équation unitaire à
coefficients dans $A$. Nous pouvons donc trouver des polynômes unitaires
$P_j(T) \in R[x_1, \ldots, x_n][T]$ tels que $P_j(y_j) = 0$ dans $B$.
Nous obtenons alors un diagramme commutatif
$$
\xymatrix{
& R[x_1, \ldots, x_n, y_1, \ldots, y_m] \ar[d] \\
R[x_1, \ldots, x_n] \ar[d] \ar[r] &
R[x_1, \ldots, x_n, y_1, \ldots, y_m]/(P_j(y_j)) \ar[d] \\
A \ar[r] & B
}
$$
La flèche horizontale supérieure et la flèche verticale supérieure droite
satisfont les hypothèses du
lemme \ref{lemma-finite-push-pseudo-coherent}.
Ainsi $K^\bullet$ est $m$-pseudo-cohérent (resp.\ pseudo-cohérent) comme complexe
de $R[x_1, \ldots, x_n]$-modules si et seulement si $K^\bullet$ est
$m$-pseudo-cohérent (resp.\ pseudo-cohérent) comme complexe de
$R[x_1, \ldots, x_n, y_1, \ldots, y_m]$-modules.
\end{proof}

\begin{lemma}
\label{lemma-cone-relatively-pseudo-coherent}
Soit $R$ un anneau. Soit $R \to A$ un morphisme d'anneaux de type fini.
Soit $m \in \mathbf{Z}$. Soit $(K^\bullet, L^\bullet, M^\bullet, f, g, h)$
un triangle distingué dans $D(A)$.
\begin{enumerate}
\item Si $K^\bullet$ est $(m + 1)$-pseudo-cohérent relativement à $R$ et si
$L^\bullet$ est $m$-pseudo-cohérent relativement à $R$, alors $M^\bullet$ est
$m$-pseudo-cohérent relativement à $R$.
\item Si $K^\bullet, M^\bullet$ sont $m$-pseudo-cohérents relativement à $R$,
alors $L^\bullet$ est $m$-pseudo-cohérent relativement à $R$.
\item Si $L^\bullet$ est $(m + 1)$-pseudo-cohérent relativement à $R$
et si $M^\bullet$ est $m$-pseudo-cohérent relativement à $R$, alors
$K^\bullet$ est $(m + 1)$-pseudo-cohérent relativement à $R$.
\end{enumerate}
De plus, si deux des trois complexes $K^\bullet, L^\bullet, M^\bullet$
sont pseudo-cohérents relativement à $R$, alors le troisième l'est aussi.
\end{lemma}

\begin{proof}
Cela résulte immédiatement du
lemme \ref{lemma-cone-pseudo-coherent}
et des définitions.
\end{proof}

\begin{lemma}
\label{lemma-rel-n-pseudo-module}
Soit $R \to A$ un morphisme d'anneaux de type fini. Soit $M$ un $A$-module.
Alors
\begin{enumerate}
\item $M$ est $0$-pseudo-cohérent relativement à $R$ si et seulement si
$M$ est un $A$-module de type fini ;
\item $M$ est $(-1)$-pseudo-cohérent relativement à $R$ si et seulement si
$M$ est de présentation finie relativement à $R$ ;
\item $M$ est $(-d)$-pseudo-cohérent relativement à $R$ si et seulement si,
pour toute surjection $R[x_1, \ldots, x_n] \to A$, il existe une
résolution
$$
R[x_1, \ldots, x_n]^{\oplus a_d} \to R[x_1, \ldots, x_n]^{\oplus a_{d - 1}}
\to \ldots \to R[x_1, \ldots, x_n]^{\oplus a_0} \to M \to 0
$$
de longueur $d$ ;
\item $M$ est pseudo-cohérent relativement à $R$ si et seulement si,
pour toute présentation $R[x_1, \ldots, x_n] \to A$, il existe une
résolution infinie
$$
\ldots \to R[x_1, \ldots, x_n]^{\oplus a_1} \to
R[x_1, \ldots, x_n]^{\oplus a_0} \to M \to 0
$$
par des $R[x_1, \ldots, x_n]$-modules libres de type fini.
\end{enumerate}
\end{lemma}

\begin{proof}
Cela résulte immédiatement du
lemme \ref{lemma-n-pseudo-module}
et des définitions.
\end{proof}

\begin{lemma}
\label{lemma-summands-relative-pseudo-coherent}
Soit $R \to A$ un morphisme d'anneaux de type fini.
Soit $m \in \mathbf{Z}$. Soient $K^\bullet, L^\bullet \in D(A)$.
Si $K^\bullet \oplus L^\bullet$
est $m$-pseudo-cohérent (resp.\ pseudo-cohérent) relativement à $R$,
alors $K^\bullet$ et $L^\bullet$ le sont aussi.
\end{lemma}

\begin{proof}
Cela résulte immédiatement du
lemme \ref{lemma-summands-pseudo-coherent}
et des définitions.
\end{proof}

\begin{lemma}
\label{lemma-complex-relative-pseudo-coherent-modules}
Soit $R \to A$ un morphisme d'anneaux de type fini.
Soit $m \in \mathbf{Z}$. Soit $K^\bullet$ un complexe borné
supérieurement de $A$-modules tel que $K^i$ soit $(m - i)$-pseudo-cohérent
relativement à $R$ pour tout $i$. Alors $K^\bullet$ est $m$-pseudo-cohérent
relativement à $R$. En particulier, si $K^\bullet$ est un complexe borné supérieurement de
$A$-modules pseudo-cohérents relativement à $R$, alors $K^\bullet$ est
pseudo-cohérent relativement à $R$.
\end{lemma}

\begin{proof}
Cela résulte immédiatement du
lemme \ref{lemma-complex-pseudo-coherent-modules}
et des définitions.
\end{proof}

\begin{lemma}
\label{lemma-cohomology-relative-pseudo-coherent}
Soit $R \to A$ un morphisme d'anneaux de type fini. Soit $m \in \mathbf{Z}$.
Soit $K^\bullet \in D^{-}(A)$ tel que $H^i(K^\bullet)$ soit
$(m - i)$-pseudo-cohérent (resp.\ pseudo-cohérent) relativement à $R$
pour tout $i$. Alors $K^\bullet$ est $m$-pseudo-cohérent
(resp.\ pseudo-cohérent) relativement à $R$.
\end{lemma}

\begin{proof}
Cela résulte immédiatement du
lemme \ref{lemma-cohomology-pseudo-coherent}
et des définitions.
\end{proof}

\begin{lemma}
\label{lemma-localize-relative-pseudo-coherent}
Soient $R$ un anneau, $f \in R$ un élément, $R_f \to A$ un morphisme d'anneaux de type fini,
$g \in A$ et $K^\bullet$ un complexe de $A$-modules.
Si $K^\bullet$ est $m$-pseudo-cohérent (resp.\ pseudo-cohérent)
relativement à $R_f$, alors $K^\bullet \otimes_A A_g$ est
$m$-pseudo-cohérent (resp.\ pseudo-cohérent) relativement à $R$.
\end{lemma}

\begin{proof}
Montrons d'abord que $K^\bullet$ est $m$-pseudo-cohérent relativement à $R$.
Supposons en effet $R_f[x_1, \ldots, x_n] \to A$ surjectif. Écrivons
$R_f = R[x_0]/(fx_0 - 1)$. Alors $R[x_0, x_1, \ldots, x_n] \to A$
est surjectif, et $R_f[x_1, \ldots, x_n]$ est pseudo-cohérent comme
$R[x_0, \ldots, x_n]$-module. Ainsi, d'après le
lemme \ref{lemma-finite-push-pseudo-coherent},
$K^\bullet$ est $m$-pseudo-cohérent comme complexe de
$R[x_0, x_1, \ldots, x_n]$-modules.

\medskip\noindent
Choisissons un élément $g' \in R[x_0, x_1, \ldots, x_n]$ s'envoyant sur $g \in A$. D'après le
lemme \ref{lemma-pull-pseudo-coherent},
nous voyons que
\begin{align*}
K^\bullet \otimes_{R[x_0, x_1, \ldots, x_n]}^{\mathbf{L}}
R[x_0, x_1, \ldots, x_n, \frac{1}{g'}] & =
K^\bullet \otimes_{R[x_0, x_1, \ldots, x_n]}
R[x_0, x_1, \ldots, x_n, \frac{1}{g'}] \\
& = K^\bullet \otimes_A A_f
\end{align*}
est $m$-pseudo-cohérent comme complexe de
$R[x_0, x_1, \ldots, x_n, \frac{1}{g'}]$-modules.
Écrivons
$$
R[x_0, x_1, \ldots, x_n, \frac{1}{g'}] =
R[x_0, \ldots, x_n, x_{n + 1}]/(x_{n + 1}g' - 1).
$$
Comme $R[x_0, x_1, \ldots, x_n, \frac{1}{g'}]$ est pseudo-cohérent comme
$R[x_0, \ldots, x_n, x_{n + 1}]$-module, nous concluons (voir le
lemme \ref{lemma-finite-push-pseudo-coherent})
que $K^\bullet \otimes_A A_g$ est $m$-pseudo-cohérent comme complexe de
$R[x_0, \ldots, x_n, x_{n + 1}]$-modules, comme souhaité.
\end{proof}

\begin{lemma}
\label{lemma-base-change-relative-pseudo-coherent}
Soit $R \to A$ un morphisme d'anneaux de type fini. Soit $m \in \mathbf{Z}$.
Soit $K^\bullet$ un complexe de $A$-modules qui est $m$-pseudo-cohérent
(resp.\ pseudo-cohérent) relativement à $R$. Soit $R \to R'$ un morphisme
d'anneaux tel que $A$ et $R'$ soient Tor-indépendants sur $R$. Posons
$A' = A \otimes_R R'$. Alors
$K^\bullet \otimes_A^{\mathbf{L}} A'$
est $m$-pseudo-cohérent (resp.\ pseudo-cohérent) relativement à $R'$.
\end{lemma}

\begin{proof}
Choisissons une surjection $R[x_1, \ldots, x_n] \to A$.
Remarquons que
$$
K^\bullet \otimes_A^{\mathbf{L}} A' =
K^\bullet \otimes_R^{\mathbf{L}} R' =
K^\bullet \otimes_{R[x_1, \ldots, x_n]}^{\mathbf{L}} R'[x_1, \ldots, x_n]
$$
d'après le
lemme \ref{lemma-base-change-comparison}
appliqué deux fois. Nous concluons donc par le
lemme \ref{lemma-pull-pseudo-coherent}.
\end{proof}

\begin{lemma}
\label{lemma-pull-relative-pseudo-coherent}
Soient $R \to A \to B$ des morphismes d'anneaux de type fini.
Soit $m \in \mathbf{Z}$.
Soit $K^\bullet$ un complexe de $A$-modules.
Supposons que $B$, comme $B$-module, soit pseudo-cohérent relativement à $A$.
Si $K^\bullet$ est $m$-pseudo-cohérent (resp.\ pseudo-cohérent)
relativement à $R$, alors $K^\bullet \otimes_A^{\mathbf{L}} B$ est
$m$-pseudo-cohérent (resp.\ pseudo-cohérent) relativement à $R$.
\end{lemma}

\begin{proof}
Choisissons une surjection $A[y_1, \ldots, y_m] \to B$.
Choisissons une surjection $R[x_1, \ldots, x_n] \to A$.
En les combinant, nous obtenons une surjection $R[x_1, \ldots, x_n, y_1, \ldots y_m] \to B$.
Choisissons une résolution $E^\bullet \to B$ de $B$ par un complexe de
$A[y_1, \ldots, y_n]$-modules libres de type fini (ce qui est possible
par notre hypothèse sur le morphisme d'anneaux $A \to B$). Nous pouvons supposer
que $K^\bullet$ est un complexe borné supérieurement de $A$-modules plats. Alors
\begin{align*}
K^\bullet \otimes_A^{\mathbf{L}} B & =
\text{Tot}(K^\bullet \otimes_A B[0]) \\
& = \text{Tot}(K^\bullet \otimes_A A[y_1, \ldots, y_m]
\otimes_{A[y_1, \ldots, y_m]} B[0]) \\
& \cong
\text{Tot}\left(
(K^\bullet \otimes_A A[y_1, \ldots, y_m])
\otimes_{A[y_1, \ldots, y_m]} E^\bullet
\right) \\
& =
\text{Tot}(K^\bullet \otimes_A E^\bullet)
\end{align*}
dans $D(A[y_1, \ldots, y_m])$. Le quasi-isomorphisme $\cong$ provient
d'une application du
lemme \ref{lemma-derived-tor-quasi-isomorphism}.
Nous devons donc montrer que
$\text{Tot}(K^\bullet \otimes_A E^\bullet)$ est $m$-pseudo-cohérent
comme complexe de $R[x_1, \ldots, x_n, y_1, \ldots y_m]$-modules.
Remarquons que $\text{Tot}(K^\bullet \otimes_A E^\bullet)$ possède une filtration par
sous-complexes dont les quotients successifs sont les complexes
$K^\bullet \otimes_A E^i[-i]$. Pour $i \ll 0$, les complexes
$K^\bullet \otimes_A E^i[-i]$ ont une cohomologie nulle
en degrés $\leq m$ et sont donc $m$-pseudo-cohérents (sur tout anneau).
Ainsi, en appliquant le
lemme \ref{lemma-cone-relatively-pseudo-coherent}
et en raisonnant par récurrence, il suffit de montrer que $K^\bullet \otimes_A E^i[-i]$ est
pseudo-cohérent relativement à $R$ pour tout $i$. Remarquons que $E^i = 0$ pour
$i > 0$. Comme $E^i$ est en outre libre de type fini, cela
se ramène à montrer que $K^\bullet \otimes_A A[y_1, \ldots, y_m]$ est
$m$-pseudo-cohérent relativement à $R$, ce qui résulte par exemple du
lemme \ref{lemma-base-change-relative-pseudo-coherent}.
\end{proof}

\begin{lemma}
\label{lemma-pull-relative-pseudo-coherent-module}
Soient $R \to A \to B$ des morphismes d'anneaux de type fini.
Soit $m \in \mathbf{Z}$. Soit $M$ un $A$-module.
Supposons $B$ plat sur $A$ et $B$, comme $B$-module,
pseudo-cohérent relativement à $A$.
Si $M$ est $m$-pseudo-cohérent (resp.\ pseudo-cohérent)
relativement à $R$, alors $M \otimes_A B$ est
$m$-pseudo-cohérent (resp.\ pseudo-cohérent) relativement à $R$.
\end{lemma}

\begin{proof}
Cela résulte immédiatement du
lemme \ref{lemma-pull-relative-pseudo-coherent}.
\end{proof}

\begin{lemma}
\label{lemma-composition-relative-pseudo-coherent}
Soit $R$ un anneau. Soit $A \to B$ un morphisme de $R$-algèbres de type fini.
Soit $m \in \mathbf{Z}$. Soit $K^\bullet$ un complexe de $B$-modules.
Supposons $A$ pseudo-cohérent relativement à $R$. Alors les assertions suivantes sont
équivalentes
\begin{enumerate}
\item $K^\bullet$ est $m$-pseudo-cohérent (resp.\ pseudo-cohérent)
relativement à $A$ ;
\item $K^\bullet$ est $m$-pseudo-cohérent (resp.\ pseudo-cohérent)
relativement à $R$.
\end{enumerate}
\end{lemma}

\begin{proof}
Choisissons une surjection $R[x_1, \ldots, x_n] \to A$.
Choisissons une surjection $A[y_1, \ldots, y_m] \to B$.
Nous obtenons alors une surjection
$$
R[x_1, \ldots, x_n, y_1, \ldots, y_m] \to A[y_1, \ldots, y_m]
$$
qui est un changement de base plat de $R[x_1, \ldots, x_n] \to A$.
Par hypothèse, $A$ est un module pseudo-cohérent sur $R[x_1, \ldots, x_n]$ ;
ainsi, d'après le
lemme \ref{lemma-flat-base-change-pseudo-coherent},
$A[y_1, \ldots, y_m]$ est pseudo-cohérent sur
$R[x_1, \ldots, x_n, y_1, \ldots, y_m]$. Le lemme résulte donc du
lemme \ref{lemma-finite-push-pseudo-coherent}
et des définitions.
\end{proof}

\begin{lemma}
\label{lemma-glue-relative-pseudo-coherent}
Soit $R \to A$ un morphisme d'anneaux de type fini.
Soit $K^\bullet$ un complexe de $A$-modules.
Soit $m \in \mathbf{Z}$.
Soient $f_1, \ldots, f_r \in A$ engendrant l'idéal unité.
Les assertions suivantes sont équivalentes
\begin{enumerate}
\item chaque $K^\bullet \otimes_A A_{f_i}$ est
$m$-pseudo-cohérent relativement à $R$ ;
\item $K^\bullet$ est $m$-pseudo-cohérent relativement à $R$.
\end{enumerate}
La même équivalence vaut pour la pseudo-cohérence relativement à $R$.
\end{lemma}

\begin{proof}
L'implication (2) $\Rightarrow$ (1) se trouve dans le
lemme \ref{lemma-localize-relative-pseudo-coherent}.
Supposons (1). Écrivons $1 = \sum f_ig_i$ dans $A$.
Choisissons une surjection
$R[x_1, \ldots, x_n, y_1, \ldots, y_r, z_1, \ldots, z_r] \to A$.
telle que $y_i$ s'envoie sur $f_i$ et $z_i$ sur $g_i$. Nous
voyons alors qu'il existe une surjection
$$
P = R[x_1, \ldots, x_n, y_1, \ldots, y_r, z_1, \ldots, z_r]/(\sum y_iz_i - 1)
\longrightarrow
A.
$$
Remarquons que $P$ est pseudo-cohérent comme
$R[x_1, \ldots, x_n, y_1, \ldots, y_r, z_1, \ldots, z_r]$-module
et que $P[1/y_i]$ est pseudo-cohérent comme
$R[x_1, \ldots, x_n, y_1, \ldots, y_r, z_1, \ldots, z_r, 1/y_i]$-module.
Ainsi, d'après le
lemme \ref{lemma-finite-push-pseudo-coherent},
$K^\bullet \otimes_A A_{f_i}$ est un complexe $m$-pseudo-cohérent
de $P[1/y_i]$-modules pour tout $i$.
Le lemme \ref{lemma-glue-pseudo-coherent}
montre donc que $K^\bullet$ est pseudo-cohérent comme complexe de
$P$-modules, et le
lemme \ref{lemma-finite-push-pseudo-coherent}
montre que $K^\bullet$ est pseudo-cohérent comme complexe de
$R[x_1, \ldots, x_n, y_1, \ldots, y_r, z_1, \ldots, z_r]$-modules.
\end{proof}

\begin{lemma}
\label{lemma-Noetherian-relative-pseudo-coherent}
Soit $R$ un anneau noethérien. Soit $R \to A$ un morphisme d'anneaux de type fini. Alors
\begin{enumerate}
\item Un complexe de $A$-modules $K^\bullet$ est $m$-pseudo-cohérent
relativement à $R$ si et seulement si $K^\bullet \in D^{-}(A)$ et
$H^i(K^\bullet)$ est un $A$-module de type fini pour $i \geq m$.
\item Un complexe de $A$-modules $K^\bullet$ est pseudo-cohérent relativement à $R$
si et seulement si $K^\bullet \in D^{-}(A)$ et
$H^i(K^\bullet)$ est un $A$-module de type fini pour tout $i$.
\item Un $A$-module est pseudo-cohérent relativement à $R$
si et seulement s'il est de type fini.
\end{enumerate}
\end{lemma}

\begin{proof}
C'est une conséquence immédiate du
lemme \ref{lemma-Noetherian-pseudo-coherent}
et des définitions.
\end{proof}










\section{Morphismes d'anneaux pseudo-cohérents et parfaits}
\label{section-pseudo-coherent-perfect-ring-map}

\noindent
Nous pouvons définir ces types de morphismes d'anneaux comme suit.

\begin{definition}
\label{definition-pseudo-coherent-perfect}
Soit $A \to B$ un morphisme d'anneaux.
\begin{enumerate}
\item Nous disons que $A \to B$ est un {\it morphisme d'anneaux pseudo-cohérent} s'il est de
type fini et si $B$, comme $B$-module, est pseudo-cohérent relativement à $A$.
\item Nous disons que $A \to B$ est un {\it morphisme d'anneaux parfait} s'il est un
morphisme d'anneaux pseudo-cohérent tel que $B$, comme $A$-module, soit de
dimension de Tor finie.
\end{enumerate}
\end{definition}

\noindent
Cette terminologie n'est peut-être pas standard. En utilisant le
lemme \ref{lemma-rel-n-pseudo-module},
nous voyons que $A \to B$ est pseudo-cohérent si et seulement si
$B = A[x_1, \ldots, x_n]/I$ et si $B$, comme $A[x_1, \ldots, x_n]$-module,
admet une résolution par des $A[x_1, \ldots, x_n]$-modules libres de type fini.
La motivation de la définition d'un morphisme d'anneaux parfait est le
lemme \ref{lemma-perfect}.
Le lemme suivant donne une caractérisation plus utile et intuitive
d'un morphisme d'anneaux parfait.

\begin{lemma}
\label{lemma-perfect-ring-map}
Un morphisme d'anneaux $A \to B$ est parfait si et seulement si $B = A[x_1, \ldots, x_n]/I$
et si $B$, comme $A[x_1, \ldots, x_n]$-module, admet une résolution finie par des
$A[x_1, \ldots, x_n]$-modules projectifs de type fini.
\end{lemma}

\begin{proof}
Si $A \to B$ est parfait, alors $B = A[x_1, \ldots, x_n]/I$,
$B$ est pseudo-cohérent comme $A[x_1, \ldots, x_n]$-module et est de
dimension de Tor finie comme $A$-module. Ainsi le
lemme \ref{lemma-perfect-over-polynomial-ring}
implique que $B$ est parfait comme $A[x_1, \ldots, x_n]$-module, c'est-à-dire
qu'il admet une résolution finie par des $A[x_1, \ldots, x_n]$-modules projectifs de type fini
(lemme \ref{lemma-perfect-module}).
Réciproquement, si $B = A[x_1, \ldots, x_n]/I$
et si $B$, comme $A[x_1, \ldots, x_n]$-module, admet une résolution finie par des
$A[x_1, \ldots, x_n]$-modules projectifs de type fini, alors
$B$ est pseudo-cohérent comme $A[x_1, \ldots, x_n]$-module,
donc $A \to B$ est pseudo-cohérent. De plus, la résolution donnée
sur $A[x_1, \ldots, x_n]$ est une résolution finie par des
$A$-modules plats ; ainsi $B$ est de dimension de Tor finie comme $A$-module.
\end{proof}

\noindent
De nombreux résultats des sections précédentes peuvent être reformulés
dans cette terminologie. Nous renvoyons aussi à
Compléments sur les morphismes, sections \ref{more-morphisms-section-pseudo-coherent}
et \ref{more-morphisms-section-perfect},
pour la discussion correspondante concernant les morphismes de schémas.

\begin{lemma}
\label{lemma-Noetherian-pseudo-coherent-ring-map}
Un morphisme d'anneaux de type fini entre anneaux noethériens est pseudo-cohérent.
\end{lemma}

\begin{proof}
Voir le
lemme \ref{lemma-Noetherian-relative-pseudo-coherent}.
\end{proof}

\begin{lemma}
\label{lemma-flat-finite-presentation-perfect}
Un morphisme d'anneaux plat et de présentation finie est parfait.
\end{lemma}

\begin{proof}
Soit $A \to B$ un morphisme d'anneaux plat et de présentation finie.
Il est clair que $B$ est de dimension de Tor finie. D'après
Algèbre, lemme \ref{algebra-lemma-flat-finite-presentation-limit-flat},
il existe une $\mathbf{Z}$-algèbre de type fini $A_0 \subset A$
et un morphisme d'anneaux plat de type fini $A_0 \to B_0$ tels que
$B = B_0 \otimes_{A_0} A$. D'après le
lemme \ref{lemma-Noetherian-relative-pseudo-coherent},
$A_0 \to B_0$ est pseudo-cohérent.
Comme $A_0 \to B_0$ est plat, $B_0$ et $A$ sont Tor-indépendants
sur $A_0$ ; nous pouvons donc appliquer le
lemme \ref{lemma-base-change-relative-pseudo-coherent}
pour conclure que $A \to B$ est pseudo-cohérent.
\end{proof}

\begin{lemma}
\label{lemma-regular-perfect-ring-map}
Soit $A \to B$ un morphisme d'anneaux de type fini, avec $A$ régulier
de dimension finie. Alors $A \to B$ est parfait.
\end{lemma}

\begin{proof}
D'après
Algèbre, lemme \ref{algebra-lemma-finite-gl-dim-finite-dim-regular},
l'hypothèse sur $A$ signifie que $A$ est de dimension globale finie.
Tout module est donc de dimension de Tor finie, voir le
lemme \ref{lemma-finite-gl-dim-tor-dimension},
en particulier $B$. D'après le
lemme \ref{lemma-Noetherian-pseudo-coherent-ring-map},
le morphisme est pseudo-cohérent.
\end{proof}

\begin{lemma}
\label{lemma-lci-perfect}
Un homomorphisme d'intersection complète locale est parfait.
\end{lemma}

\begin{proof}
Soit $A \to B$ un homomorphisme d'intersection complète locale.
D'après la définition \ref{definition-local-complete-intersection}, cela
signifie que $B = A[x_1, \ldots, x_n]/I$, où $I$ est un idéal de Koszul
dans $A[x_1, \ldots, x_n]$.
D'après les lemmes \ref{lemma-perfect-ring-map} et \ref{lemma-perfect-module},
il suffit de montrer que $I$ est un module parfait sur $A[x_1, \ldots, x_n]$.
D'après le lemme \ref{lemma-glue-perfect}, c'est une question locale. Nous
pouvons donc supposer que $I$ est engendré par une suite de Koszul régulière (d'après la
définition \ref{definition-regular-ideal}).
Cela signifie bien entendu que $I$ admet une résolution libre finie, ce qui conclut.
\end{proof}

\begin{lemma}
\label{lemma-relative-pseudo-coherent-is-moot}
Soit $R \to A$ un morphisme d'anneaux pseudo-cohérent. Soit $K \in D(A)$.
Les assertions suivantes sont équivalentes
\begin{enumerate}
\item $K$ est $m$-pseudo-cohérent (resp.\ pseudo-cohérent) relativement à $R$ ;
\item $K$ est $m$-pseudo-cohérent (resp.\ pseudo-cohérent) dans $D(A)$.
\end{enumerate}
\end{lemma}

\begin{proof}
C'est une reformulation d'un cas particulier du
lemme \ref{lemma-composition-relative-pseudo-coherent}.
\end{proof}

\begin{lemma}
\label{lemma-more-relative-pseudo-coherent-is-moot}
Soient $R \to B \to A$ des morphismes d'anneaux, avec $\varphi : B \to A$ surjectif et
$R \to B$ et $R \to A$ plats et de présentation finie. Pour $K \in D(A)$,
notons $\varphi_*K \in D(B)$ la restriction.
Les assertions suivantes sont équivalentes
\begin{enumerate}
\item $K$ est pseudo-cohérent ;
\item $K$ est pseudo-cohérent relativement à $R$ ;
\item $K$ est pseudo-cohérent relativement à $A$ ;
\item $\varphi_*K$ est pseudo-cohérent ;
\item $\varphi_*K$ est pseudo-cohérent relativement à $R$.
\end{enumerate}
Il en va de même pour la $m$-pseudo-cohérence.
\end{lemma}

\begin{proof}
Remarquons que $R \to A$ et $R \to B$ sont des morphismes d'anneaux parfaits
(lemme \ref{lemma-flat-finite-presentation-perfect}),
donc a fortiori des morphismes d'anneaux pseudo-cohérents.
Ainsi (1) $\Leftrightarrow$ (2) et (4) $\Leftrightarrow$ (5)
d'après le lemme \ref{lemma-relative-pseudo-coherent-is-moot}.

\medskip\noindent
Comme $A$ est pseudo-cohérent relativement à $R$, le
lemme \ref{lemma-composition-relative-pseudo-coherent}
montre que (2) $\Leftrightarrow$ (3).
Cependant, puisque $A \to B$ est surjectif, la
définition \ref{definition-relatively-pseudo-coherent}
montre directement que (3) est équivalent à (4).
\end{proof}



\section{Modules relativement parfaits}
\label{section-relatively-perfect}

\noindent
Cette section est l'analogue de la
section \ref{section-relative-pseudo-coherent}
pour les objets parfaits de la catégorie dérivée.
Nous ne définissons cette notion que dans une généralité
limitée, car nous ne savons pas quelle est la bonne
définition en général. Voir
Catégories dérivées des schémas, remarque
\ref{perfect-remark-discuss-rel-perfect},
pour une discussion.

\begin{definition}
\label{definition-relatively-perfect}
Soit $R \to A$ un morphisme d'anneaux plat et de présentation finie.
Un objet $K$ de $D(A)$ est {\it $R$-parfait} ou {\it parfait relativement à $R$}
si $K$ est pseudo-cohérent
(définition \ref{definition-pseudo-coherent})
et de dimension de Tor finie sur $R$
(définition \ref{definition-tor-amplitude}).
\end{definition}

\noindent
D'après le lemme \ref{lemma-more-relative-pseudo-coherent-is-moot},
il aurait été équivalent de demander que $K$ soit
pseudo-cohérent relativement à $R$.
Voici quelques lemmes incontournables.

\begin{lemma}
\label{lemma-cone-relatively-perfect}
Soit $R \to A$ un morphisme d'anneaux plat et de présentation finie.
Les objets $R$-parfaits de $D(A)$ forment une sous-catégorie
triangulée, strictement pleine et
saturée\footnote{Catégories dérivées, définition
\ref{derived-definition-saturated}.}.
\end{lemma}

\begin{proof}
Cela résulte des
lemmes \ref{lemma-cone-pseudo-coherent},
\ref{lemma-summands-pseudo-coherent},
\ref{lemma-cone-tor-amplitude} et
\ref{lemma-summands-tor-amplitude}.
\end{proof}

\begin{lemma}
\label{lemma-perfect-relatively-perfect}
Soit $R \to A$ un morphisme d'anneaux plat et de présentation finie.
Un objet parfait de $D(A)$ est $R$-parfait. Si $K, M \in D(A)$,
alors $K \otimes_A^\mathbf{L} M$ est $R$-parfait si $K$ est parfait
et $M$ est $R$-parfait.
\end{lemma}

\begin{proof}
La première assertion résulte de la seconde en prenant $M = A$.
La seconde résulte des lemmes \ref{lemma-perfect},
\ref{lemma-push-tor-amplitude} et \ref{lemma-tensor-pseudo-coherent}.
\end{proof}

\begin{lemma}
\label{lemma-structure-relatively-perfect}
Soit $R \to A$ un morphisme d'anneaux plat et de présentation finie.
Soit $K \in D(A)$. Les assertions suivantes sont équivalentes
\begin{enumerate}
\item $K$ est $R$-parfait ;
\item $K$ est isomorphe à un complexe fini de modules plats sur $R$
et de présentation finie sur $A$.
\end{enumerate}
\end{lemma}

\begin{proof}
Pour démontrer que (2) implique (1), il suffit, d'après le
lemme \ref{lemma-cone-relatively-perfect},
de montrer qu'un module plat sur $R$ et de présentation finie sur $A$, noté $M$, définit
un objet $R$-parfait de $D(A)$. Comme $M$ est de dimension de Tor finie
sur $R$, il suffit de montrer que $M$ est pseudo-cohérent. D'après
Algèbre, lemme \ref{algebra-lemma-flat-finite-presentation-limit-flat},
il existe une $\mathbf{Z}$-algèbre de type fini $R_0 \subset R$,
un morphisme d'anneaux plat de type fini $R_0 \to A_0$ et
un $A_0$-module de type fini $M_0$, plat sur $R_0$, tels que
$A = A_0 \otimes_{R_0} R$ et $M = M_0 \otimes_{R_0} R$. D'après le
lemme \ref{lemma-Noetherian-pseudo-coherent},
$M_0$ est un $A_0$-module pseudo-cohérent.
Choisissons une résolution $P_0^\bullet \to M_0$ par des
$A_0$-modules libres de type fini $P_0^n$. Comme $A_0$ est plat sur $R_0$,
c'est une résolution plate. Comme $M_0$ est plat sur $R_0$,
nous voyons que $P^\bullet = P_0^\bullet \otimes_{R_0} R$
résout encore $M = M_0 \otimes_{R_0} R$. (On peut utiliser le
lemme \ref{lemma-base-change-comparison} pour le voir.)
Ainsi $P^\bullet$ est une résolution libre finie de $M$
sur $A$, et nous concluons que $M$ est pseudo-cohérent.

\medskip\noindent
Supposons (1). Nous pouvons représenter $K$ par un complexe borné supérieurement
$P^\bullet$ de $A$-modules libres de type fini. Supposons que
$K$, considéré comme objet de $D(R)$, ait une amplitude de Tor dans $[a, b]$.
D'après le lemme \ref{lemma-last-one-flat},
$\tau_{\geq a}P^\bullet$ est un complexe de modules plats sur $R$ et
de présentation finie sur $A$, représentant $K$.
\end{proof}

\begin{lemma}
\label{lemma-base-change-relatively-perfect}
Soit $R \to A$ un morphisme d'anneaux plat et de présentation finie.
Soit $R \to R'$ un morphisme d'anneaux et posons $A' = A \otimes_R R'$.
Si $K \in D(A)$ est $R$-parfait, alors $K \otimes_A^\mathbf{L} A'$ est
$R'$-parfait.
\end{lemma}

\begin{proof}
D'après le lemme \ref{lemma-pull-pseudo-coherent},
$K \otimes_A^\mathbf{L} A'$ est pseudo-cohérent.
D'après le lemme \ref{lemma-base-change-comparison},
$K \otimes_A^\mathbf{L} A'$ est égal à $K \otimes_R^\mathbf{L} R'$ dans $D(R')$.
Nous pouvons alors appliquer le lemme \ref{lemma-pull-tor-amplitude}
pour voir que $K \otimes_R^\mathbf{L} R'$ dans $D(R')$ est de
dimension de Tor finie.
\end{proof}

\begin{lemma}
\label{lemma-compute-RHom-relatively-perfect}
Soit $R \to A$ un morphisme d'anneaux plat. Soient $K, L \in D(A)$, avec $K$
pseudo-cohérent et $L$ de dimension de Tor finie sur $R$. Nous pouvons choisir
\begin{enumerate}
\item un complexe $P^\bullet$ borné supérieurement
de $A$-modules libres de type fini représentant $K$ ;
\item un complexe borné de $R$-modules plats, qui sont aussi des
$A$-modules, noté $F^\bullet$, représentant $L$.
\end{enumerate}
Avec ces choix, nous avons
\begin{enumerate}
\item[(a)] $E^\bullet = \Hom^\bullet(P^\bullet, F^\bullet)$
est un complexe borné inférieurement de $R$-modules plats munis d'une structure de
$A$-module
représentant $R\Hom_A(K, L)$ ;
\item[(b)] pour tout morphisme d'anneaux $R \to R'$ avec $A' = A \otimes_R R'$,
le complexe $E^\bullet \otimes_R R'$ représente
$R\Hom_{A'}(K \otimes_A^\mathbf{L} A', L \otimes_A^\mathbf{L} A')$.
\end{enumerate}
Si, de plus, $R \to A$ est de présentation finie et $L$ est
$R$-parfait, alors nous pouvons choisir les $F^p$ parmi les
$A$-modules de présentation finie ; par conséquent, les $E^n$ seront aussi des
$A$-modules de présentation finie.
\end{lemma}

\begin{proof}
L'existence de $P^\bullet$ est la définition d'un complexe pseudo-cohérent.
Nous représentons d'abord $L$
par un complexe $F^\bullet$ borné supérieurement de $A$-modules libres
(cela est possible car une dimension de Tor bornée implique en particulier
que l'objet est borné). Disons ensuite que $L$, considéré comme objet de $D(R)$,
a une amplitude de Tor dans $[a, b]$. Après avoir remplacé
$F^\bullet$ par $\tau_{\geq a}F^\bullet$, nous obtenons un complexe
comme en (2). Cela résulte du lemme \ref{lemma-last-one-flat}.

\medskip\noindent
Preuve de (a).
Comme $F^\bullet$ est borné et $P^\bullet$ borné supérieurement,
nous voyons que $E^n = 0$ pour $n \ll 0$ et que
$E^n$ est une somme directe finie (!)
$$
E^n = \bigoplus\nolimits_{p + q = n} \Hom_A(P^{-q}, F^p)
$$
et, comme $P^{-q}$ est libre de type fini, c'est bien un module plat sur $R$ et un
$A$-module.
Le fait que $E^\bullet$ représente $R\Hom_A(K, L)$
résulte du lemme \ref{lemma-RHom-out-of-projective}.

\medskip\noindent
Preuve de (b).
Soit $R \to R'$ un morphisme d'anneaux et $A' = A \otimes_R R'$.
D'après le lemme \ref{lemma-base-change-comparison}, l'objet
$L \otimes_A^\mathbf{L} A'$ est représenté par
$F^\bullet \otimes_R R'$, considéré comme complexe de $A'$-modules
(par platitude de $F^p$ sur $R$). De même pour $P^\bullet \otimes_R R'$.
Comme ci-dessus, $R\Hom_{A'}(K \otimes_A^\mathbf{L} A', L \otimes_A^\mathbf{L} A')$
est représenté par
$$
\Hom^\bullet(P^\bullet \otimes_R R', F^\bullet \otimes_R R') =
E^\bullet \otimes_R R'
$$
L'égalité se vérifie terme à terme dans le complexe en utilisant
$\Hom_{A'}(P^{-q} \otimes_R R', F^p \otimes_R R') =
\Hom_A(P^{-q}, F^p) \otimes_R R'$.
\end{proof}

\begin{lemma}
\label{lemma-colimit-relatively-perfect}
Soit $R = \colim_{i \in I} R_i$ une colimite filtrante d'anneaux.
Soit $0 \in I$ et soit $R_0 \to A_0$ un morphisme d'anneaux plat et de
présentation finie. Pour $i \geq 0$, posons $A_i = R_i \otimes_{R_0} A_0$
et $A = R \otimes_{R_0} A_0$.
\begin{enumerate}
\item Étant donné un complexe $R$-parfait $K$ dans $D(A)$, il existe $i \in I$
et un complexe $R_i$-parfait $K_i$ dans $D(A_i)$ tels que
$K \cong K_i \otimes_{A_i}^\mathbf{L} A$ dans $D(A)$.
\item Étant donnés $K_0, L_0 \in D(A_0)$, avec $K_0$ pseudo-cohérent
et $L_0$ de dimension de Tor finie sur $R_0$, nous avons
$$
\Hom_{D(A)}(K_0 \otimes_{A_0}^\mathbf{L} A, L_0 \otimes_{A_0}^\mathbf{L} A) =
\colim_{i \geq 0}
\Hom_{D(A_i)}(K_0 \otimes_{A_0}^\mathbf{L} A_i,
L_0 \otimes_{A_0}^\mathbf{L} A_i)
$$
\end{enumerate}
En particulier, la catégorie triangulée des complexes $R$-parfaits sur $A$
est la colimite des catégories triangulées des
complexes $R_i$-parfaits sur $A_i$.
\end{lemma}

\begin{proof}
D'après Algèbre, lemme \ref{algebra-lemma-colimit-category-fp-modules},
la catégorie des $A$-modules de présentation finie est la colimite des
catégories des $A_i$-modules de présentation finie.
Cela étant, Algèbre, lemme
\ref{algebra-lemma-flat-finite-presentation-limit-flat},
énonce que la catégorie des modules plats sur $R$ et de présentation finie sur $A$
est la colimite des catégories des
modules plats sur $R_i$ et de présentation finie sur $A_i$.
Ainsi, la caractérisation du
lemme \ref{lemma-structure-relatively-perfect}
démontre (1).

\medskip\noindent
Pour démontrer (2), choisissons $P_0^\bullet$ représentant $K_0$ et
$F_0^\bullet$ représentant $L_0$ comme dans le
lemme \ref{lemma-compute-RHom-relatively-perfect}.
Alors $E_0^\bullet = \Hom^\bullet(P_0^\bullet, F_0^\bullet)$
vérifie
$$
H^0(E_0^\bullet \otimes_{R_0} R_i) =
\Hom_{D(A_i)}(K_0 \otimes_{A_0}^\mathbf{L} A_i,
L_0 \otimes_{A_0}^\mathbf{L} A_i)
$$
et
$$
H^0(E_0^\bullet \otimes_{R_0} R) =
\Hom_{D(A)}(K_0 \otimes_{A_0}^\mathbf{L} A, L_0 \otimes_{A_0}^\mathbf{L} A)
$$
d'après le lemme. Le résultat en découle, car le produit tensoriel commute
aux colimites et les colimites filtrantes sont exactes
(Algèbre, lemme \ref{algebra-lemma-directed-colimit-exact}).
\end{proof}

\begin{lemma}
\label{lemma-thickening-relatively-perfect}
Soit $R' \to A'$ un morphisme d'anneaux plat et de présentation finie.
Soit $R' \to R$ un morphisme d'anneaux surjectif dont le noyau est un idéal nilpotent.
Posons $A = A' \otimes_{R'} R$. Soit $K' \in D(A')$ et posons
$K = K' \otimes_{A'}^\mathbf{L} A$ dans $D(A)$.
Si $K$ est $R$-parfait, alors $K'$ est $R'$-parfait.
\end{lemma}

\begin{proof}
Remarquons que $A' \to A$ a un noyau nilpotent et que, par platitude de
$R' \to A'$, nous avons $K = K' \otimes_{R'}^\mathbf{L} R$ (voir la
section \ref{section-tor-independence}). Le lemme résulte donc de la combinaison des
lemmes \ref{lemma-check-pseudo-coherent-modulo-nilpotent} et
\ref{lemma-check-tor-dimension-modulo-nilpotent}.
\end{proof}

\begin{lemma}
\label{lemma-lift-from-fibre-relatively-perfect}
Soit $R$ un anneau. Soit $A = R[x_1, \ldots, x_d]/I$
plate et de présentation finie sur $R$.
Soit $\mathfrak q \subset A$ un idéal premier au-dessus de
$\mathfrak p \subset R$. Soit $K \in D(A)$ pseudo-cohérent.
Soient $a, b \in \mathbf{Z}$. Si
$H^i(K_\mathfrak q \otimes_{R_\mathfrak p}^\mathbf{L} \kappa(\mathfrak p))$
n'est non nul que pour $i \in [a, b]$, alors
$K_\mathfrak q$ a une amplitude de Tor dans $[a - d, b]$ sur $R$.
\end{lemma}

\begin{proof}
D'après le lemme \ref{lemma-more-relative-pseudo-coherent-is-moot},
$K$ est pseudo-cohérent comme complexe de $R[x_1, \ldots, x_d]$-modules.
Nous pouvons donc supposer $A = R[x_1, \ldots, x_d]$.
En appliquant le lemme \ref{lemma-perfect-over-regular-local-ring}
à $R_\mathfrak p \to A_\mathfrak q$ et au complexe $K_\mathfrak q$
avec notre hypothèse, nous trouvons que $K_\mathfrak q$ est parfait
dans $D(A_\mathfrak q)$, d'amplitude de Tor contenue dans $[a - d, b]$.
Comme $R_\mathfrak p \to A_\mathfrak q$ est plat, nous concluons
par le lemme \ref{lemma-flat-push-tor-amplitude}.
\end{proof}

\begin{lemma}
\label{lemma-bounded-on-fibres-relatively-perfect}
Soit $R \to A$ un morphisme d'anneaux plat et de présentation finie.
Soit $K \in D(A)$ pseudo-cohérent. Les assertions suivantes sont équivalentes
\begin{enumerate}
\item $K$ est $R$-parfait ;
\item $K$ est borné inférieurement et, pour tout idéal premier $\mathfrak p \subset R$,
l'objet $K \otimes_R^\mathbf{L} \kappa(\mathfrak p)$ est borné inférieurement.
\end{enumerate}
\end{lemma}

\begin{proof}
Remarquons que (1) implique (2), car un complexe $R$-parfait est par définition
de dimension de Tor bornée comme complexe de $R$-modules. Démontrons
l'autre implication.

\medskip\noindent
Écrivons $A = R[x_1, \ldots, x_d]/I$. Notons $L$ dans $D(R[x_1, \ldots, x_d])$
la restriction de $K$. D'après le
lemme \ref{lemma-more-relative-pseudo-coherent-is-moot},
$L$ est pseudo-cohérent.
Comme $L$ et $K$ ont la même image dans $D(R)$, nous voyons que
$L$ est $R$-parfait si et seulement si $K$ est $R$-parfait.
De plus, $L \otimes_R^\mathbf{L} \kappa(\mathfrak p)$
et $K \otimes_R^\mathbf{L} \kappa(\mathfrak p)$ sont les mêmes objets de
$D(\kappa(\mathfrak p))$. Nous sommes donc ramenés au cas
$A = R[x_1, \ldots, x_d]$.

\medskip\noindent
Supposons $A = R[x_1, \ldots, x_d]$ et $K$ satisfaisant (2).
Soit $\mathfrak q \subset A$ un idéal premier au-dessus d'un
idéal premier $\mathfrak p \subset R$. En appliquant le
lemme \ref{lemma-perfect-over-regular-local-ring}
à $R_\mathfrak p \to A_\mathfrak q$ et au complexe $K_\mathfrak q$
avec notre hypothèse, nous trouvons que $K_\mathfrak q$ est parfait
dans $D(A_\mathfrak q)$. Comme $K$ est borné inférieurement, le
lemme \ref{lemma-check-perfect-stalks}
montre que $K$ est parfait dans $D(A)$.
Il s'ensuit que $K$ est $R$-parfait d'après le
lemme \ref{lemma-perfect-relatively-perfect},
ce qui achève la démonstration.
\end{proof}











\section{Complexes à deux termes}
\label{section-two-term}

\noindent
Dans cette section, nous démontrons quelques résultats sur les complexes à deux termes de modules
qui nous aideront à comprendre les conditions portant sur le complexe cotangent naïf.

\begin{lemma}
\label{lemma-ext-1-zero}
Soit $R$ un anneau. Soit $K \in D(R)$ avec $H^i(K) = 0$ pour
$i \not \in \{-1, 0\}$. Les assertions suivantes sont équivalentes
\begin{enumerate}
\item $H^{-1}(K) = 0$ et $H^0(K)$ est un module projectif ;
\item $\Ext^1_R(K, M) = 0$ pour tout $R$-module $M$.
\end{enumerate}
Si $R$ est noethérien et si $H^i(K)$ est un $R$-module de type fini pour
$i = -1, 0$, ces assertions sont aussi équivalentes à
\begin{enumerate}
\item[(3)] $\Ext^1_R(K, M) = 0$ pour tout $R$-module de type fini $M$.
\end{enumerate}
\end{lemma}

\begin{proof}
L'équivalence de (1) et (2) résulte du
lemme \ref{lemma-projective-amplitude}.
Si $R$ est noethérien et si $H^i(K)$ est un $R$-module de type fini pour
$i = -1, 0$, alors $K$ est pseudo-cohérent, voir le
lemme \ref{lemma-Noetherian-pseudo-coherent}.
L'équivalence de (1) et (3) résulte donc du
lemme \ref{lemma-projective-amplitude-pseudo-coherent}.
\end{proof}

\begin{remark}
\label{remark-smoothness-ext-1-zero}
Les deux assertions suivantes résultent du lemme \ref{lemma-ext-1-zero},
d'Algèbre, définition \ref{algebra-definition-smooth}, et
d'Algèbre, proposition \ref{algebra-proposition-characterize-formally-smooth}.
\begin{enumerate}
\item Un morphisme d'anneaux $A \to B$ est lisse si et seulement si $A \to B$ est
de présentation finie et $\Ext^1_B(\NL_{B/A}, N) = 0$
pour tout $B$-module $N$.
\item Un morphisme d'anneaux $A \to B$ est formellement lisse si et seulement si
$\Ext^1_B(\NL_{B/A}, N) = 0$ pour tout $B$-module $N$.
\end{enumerate}
\end{remark}

\begin{lemma}
\label{lemma-represent-two-term-complex}
Soit $R$ un anneau. Soit $K$ un objet de $D(R)$ avec $H^i(K) = 0$
pour $i \not \in \{-1, 0\}$. Alors
\begin{enumerate}
\item $K$ peut être représenté par un complexe à deux termes
$K^{-1} \to K^0$, avec $K^0$ libre ;
\item si $R$ est noethérien et si $H^i(K)$ est un $R$-module de type fini pour
$i = -1, 0$, alors $K$ peut être représenté par un complexe à deux termes
$K^{-1} \to K^0$, avec $K^0$ libre de type fini et $K^{-1}$ de type fini.
\end{enumerate}
\end{lemma}

\begin{proof}
Preuve de (1). Supposons que $K$ soit donné par le complexe de modules $M^\bullet$.
Nous pouvons d'abord remplacer $M^\bullet$ par $\tau_{\leq 0}M^\bullet$. Nous
pouvons donc supposer $M^i = 0$ pour $i > 0$. Ensuite, choisissons une
résolution libre $P^\bullet \to M^\bullet$ avec $P^i = 0$ pour $i > 0$, voir
Catégories dérivées, lemme \ref{derived-lemma-subcategory-left-resolution}.
Enfin, nous pouvons poser $K^\bullet = \tau_{\geq -1}P^\bullet$.

\medskip\noindent
Preuve de (2). Supposons $R$ noethérien et $H^i(K)$ un $R$-module de type fini pour
$i = -1, 0$. D'après le lemme \ref{lemma-pseudo-coherent}, nous pouvons choisir un
quasi-isomorphisme $F^\bullet \to M^\bullet$ avec $F^i = 0$ pour $i > 0$
et $F^i$ libre de type fini. Nous pouvons alors poser $K^\bullet = \tau_{\geq -1}F^\bullet$.
\end{proof}

\noindent
Les morphismes, dans la catégorie dérivée, issus du complexe cotangent naïf
$\NL_{B/A}$ ou $\NL(\alpha)$ (voir
Algèbre, section \ref{algebra-section-netherlander})
sont faciles à comprendre grâce au lemme suivant.

\begin{lemma}
\label{lemma-map-out-of-almost-free}
Soit $R$ un anneau. Soit $M^\bullet$ un complexe de modules sur $R$
avec $M^i = 0$ pour $i > 0$ et $M^0$ un $R$-module projectif.
Soit $K^\bullet$ un second complexe.
\begin{enumerate}
\item Supposons $K^i = 0$ pour $i \leq -2$. Alors
$\Hom_{D(R)}(M^\bullet, K^\bullet) = \Hom_{K(R)}(M^\bullet, K^\bullet)$.
\item Supposons $K^i = 0$ pour $i \not \in [-1, 0]$ et
$K^0$ un $R$-module projectif. Alors, pour une application de complexes
$a^\bullet : M^\bullet \to K^\bullet$, les assertions suivantes sont équivalentes
\begin{enumerate}
\item $a^\bullet$ induit l'application nulle $\Ext^1_R(K^\bullet, N) \to
\Ext^1_R(M^\bullet, N)$ pour tout $R$-module $N$ ;
\item il existe une application $h^0 : M^0 \to K^{-1}$ telle que
$a^{-1} + h^0 \circ d^{-1}_K = 0$.
\end{enumerate}
\item Supposons $K^i = 0$ pour $i \leq -3$. Soit
$\alpha \in \Hom_{D(R)}(M^\bullet, K^\bullet)$. Si la
composée de $\alpha$ avec
$K^\bullet \to K^{-2}[2]$ provient d'une application de $R$-modules
$a : M^{-2} \to K^{-2}$ avec $a \circ d_M^{-3} = 0$, alors
$\alpha$ peut être représenté par une application de complexes
$a^\bullet : M^\bullet \to K^\bullet$ avec $a^{-2} = a$.
\item Dans (2), pour toute seconde application de complexes
$(a')^\bullet : M^\bullet \to K^\bullet$
représentant $\alpha$ avec $a = (a')^{-2}$,
il existe $h^i : M^i \to K^{i - 1}$ pour $i = 0, -1$ tels que
$$
h^{-1} \circ d_M^{-2} = 0, \quad
(a')^{-1} = a^{-1} + d_K^{-2} \circ h^{-1} + h^0 \circ d_M^{-1},\quad
(a')^0 = a^0 + d_K^{-1} \circ h^0
$$
\end{enumerate}
\end{lemma}

\begin{proof}
Posons $F^0 = M^0$.
Choisissons un $R$-module libre $F^{-1}$ et une surjection $F^{-1} \to M^{-1}$.
Choisissons un $R$-module libre $F^{-2}$ et une surjection
$F^{-2} \to M^{-2} \times_{M^{-1}} F^{-1}$. En poursuivant ainsi,
nous obtenons un quasi-isomorphisme $p^\bullet : F^\bullet \to M^\bullet$
surjectif terme à terme et tel que $F^i$ soit projectif pour tout $i$.

\medskip\noindent
Preuve de (1). D'après
Catégories dérivées, lemme \ref{derived-lemma-morphisms-from-projective-complex},
nous avons
$$
\Hom_{D(R)}(M^\bullet, K^\bullet) = \Hom_{K(R)}(F^\bullet, K^\bullet)
$$
Si $K^i = 0$ pour $i \leq -2$, alors tout morphisme de complexes
$F^\bullet \to K^\bullet$ se factorise par $p^\bullet$. De même, toute
homotopie $\{h^i : F^i \to K^{i - 1}\}$ se factorise par $p^\bullet$.
Ainsi (1) est vérifié.

\medskip\noindent
Preuve de (2). Si (2)(b) est vérifié, alors $a^\bullet$ est homotope à une application
de complexes $(a')^\bullet : M^\bullet \to K^\bullet$
qui est nulle en degré $-1$.
D'autre part, soit $N \to I^\bullet$ une résolution injective.
Nous avons
$$
\Ext^1_R(K^\bullet, N) = \Hom_{D(R)}(K^\bullet, I^\bullet[1]) =
\Hom_{K(R)}(K^\bullet, I^\bullet[1])
$$
d'après Catégories dérivées, lemme
\ref{derived-lemma-morphisms-into-injective-complex}.
Soit $b^\bullet : K^\bullet \to I^\bullet[1]$ une application de complexes.
Comme $K^1 = 0$, l'application $b^0 : K^0 \to I^1$ prend ses valeurs dans le noyau de
$I^1 \to I^2$, qui est l'image de $I^0 \to I^1$. Comme $K^0$ est projectif,
nous pouvons relever $b^0$ en une application $h : K^0 \to I^0$. Ainsi $b^\bullet$
est homotope à une application de complexes $(b')^\bullet$ avec $(b')^0 = 0$. Comme
$K^i = 0$ pour $i \not \in [-1, 0]$, il s'ensuit que
$(b')^\bullet \circ (a')^\bullet = 0$ comme application de complexes.
Ainsi l'application $\Ext^1_R(K^\bullet, N) \to
\Ext^1_R(M^\bullet, N)$ est nulle. Nous voyons donc que (2)(b)
implique (2)(a). Réciproquement, supposons (2)(a). L'élément
canonique de $\Ext^1_R(K^\bullet, K^{-1})$ s'envoie sur
zéro dans $\Ext^1_R(M^\bullet, K^{-1})$. En utilisant (1), nous
obtenons immédiatement une application $h^0$ comme en (2)(b).

\medskip\noindent
Preuve de (3). Choisissons $b^\bullet : F^\bullet \to K^\bullet$ représentant
$\alpha$. La composée de $\alpha$ avec $K^\bullet \to K^{-2}[2]$ est
représentée par $b^{-2} : F^{-2} \to K^{-2}$. Comme celle-ci est homotope à
$a \circ p^{-2} : F^{-2} \to M^{-2} \to K^{-2}$, il existe une application
$h : F^{-1} \to K^{-2}$ telle que $b^{-2} = a \circ p^{-2} + h \circ d_F^{-2}$.
En modifiant $b^\bullet$ par $h$, considéré comme une homotopie de $F^\bullet$
vers $K^\bullet$, nous obtenons $b^{-2} = a \circ p^{-2}$. Ainsi $b^{-2}$
se factorise par $p^{-2}$. Comme $F^0 = M^0$, le noyau de $p^{-2}$
se surjecte sur le noyau de $p^{-1}$ (par exemple parce que le noyau
de $p^\bullet$ est un complexe acyclique, ou par chasse au diagramme). Ainsi $b^{-1}$
se factorise nécessairement lui aussi par $p^{-1}$, et nous voyons que (3)
est vérifié pour ces factorisations et $a^0 = b^0$.

\medskip\noindent
La preuve de (4) est omise. Indication : il existe une homotopie entre
$a^\bullet \circ p^\bullet$ et $(a')^\bullet \circ p^\bullet$,
et nous raisonnons comme précédemment pour montrer qu'elle se factorise par $p^\bullet$.
\end{proof}

\noindent
Soit $A \to B$ un morphisme d'anneaux de présentation finie.
Étant donné un idéal $I \subset B$, nous pouvons considérer la condition
\begin{enumerate}
\item[(*)]
$\Ext^1_B(\NL_{B/A}, N)$ est annulé par $I$ pour tout $B$-module $N$.
\end{enumerate}
Cette condition est une formulation mathématique précise possible de
la notion « le lieu singulier de $A \to B$ est contenu schématiquement
dans $V(I)$ ». Comparer avec la
remarque \ref{remark-smoothness-ext-1-zero} et les lemmes suivants.

\begin{lemma}
\label{lemma-ext-1-annihilated-definite}
Soit $R$ un anneau et soit $I \subset R$ un idéal.
Soit $K \in D(R)$. Supposons $H^i(K) = 0$ pour $i \not \in \{-1, 0\}$.
Les assertions suivantes sont équivalentes
\begin{enumerate}
\item $\Ext^1_R(K, N)$ est annulé par $I$ pour tout $R$-module $N$ ;
\item $K$ peut être représenté par un complexe $K^{-1} \to K^0$
avec $K^0$ libre, tel que pour tout $a \in I$ l'application
$a : K^{-1} \to K^{-1}$ se factorise par $d_K^{-1} : K^{-1} \to K^0$ ;
\item chaque fois que $K$ est représenté par un complexe à deux termes
$K^{-1} \to K^0$ avec $K^0$ projectif, alors, pour tout $a \in I$, l'application
$a : K^{-1} \to K^{-1}$ se factorise par $d_K^{-1} : K^{-1} \to K^0$.
\end{enumerate}
Si $R$ est noethérien et si $H^i(K)$ est un $R$-module de type fini pour $i = -1, 0$,
alors ces assertions sont aussi équivalentes à
\begin{enumerate}
\item[(4)] $\Ext^1_R(K, N)$ est annulé par $I$ pour tout
$R$-module de type fini $N$ ;
\item[(5)] $K$ peut être représenté par un complexe $K^{-1} \to K^0$
avec $K^0$ libre de type fini et $K^{-1}$ de type fini, tel que pour tout $a \in I$ l'application
$a : K^{-1} \to K^{-1}$ se factorise par $d_K^{-1} : K^{-1} \to K^0$.
\end{enumerate}
\end{lemma}

\begin{proof}

\medskip\noindent
Supposons (1) et soit $K^{-1} \to K^0$
un complexe à deux termes représentant $K$,
avec $K^0$ projectif.
Nous utiliserons sans autre mention la description donnée dans le
lemme \ref{lemma-map-out-of-almost-free} des morphismes dans $D(R)$ issus de $K^\bullet$.
En prenant $N = K^{-1}$, considérons l'élément $\xi$ de $\Ext^1_R(K, N)$
donné par $\text{id}_{K^{-1}} : K^{-1} \to K^{-1}$. Comme il est annulé
par $a \in I$, nous obtenons la flèche pointillée s'insérant dans le
diagramme commutatif suivant
$$
\xymatrix{
K^{-1} \ar[d]_a \ar[r]_{d_K^{-1}} & K^0 \ar@{..>}[ld]^h \\
K^{-1}
}
$$
Cela démontre (3). La partie (3) implique (2) d'après le
lemme \ref{lemma-represent-two-term-complex}, partie (1). Supposons
$K^\bullet$ comme en (2) et soit $N$ un $R$-module quelconque.
Tout élément $\xi$ de $\Ext^1_R(K, N)$ est donné par la classe d'une
application $\varphi : K^{-1} \to N$. Alors, pour $a \in I$, nous pouvons par hypothèse
choisir une application $h$ comme dans le diagramme ci-dessus, et nous voyons que
$a\varphi = \varphi \circ a = \varphi \circ h \circ \text{d}_K^{-1}$
ce qui montre que $a \xi$ est nul dans $\Ext^1_R(K, N)$.
Ainsi (1), (2) et (3) sont équivalentes.

\medskip\noindent
Supposons $R$ noethérien et $H^i(K)$ un $R$-module de type fini pour $i = -1, 0$.
La partie (3) implique (5) d'après le
lemme \ref{lemma-represent-two-term-complex}, partie (2).
Il est clair que (5) implique (2). Trivialement, (1) implique (4).
Pour achever la démonstration, il suffit donc de montrer que (4) implique
l'une quelconque des autres conditions. Soit $K^{-1} \to K^0$ un complexe
représentant $K$, avec $K^0$ libre de type fini et $K^{-1}$ de type fini, comme dans le
lemme \ref{lemma-represent-two-term-complex}, partie (2). L'argument
donné dans la preuve de (2) $\Rightarrow$ (1)
montre que si $\Ext^1_R(K, K^{-1})$ est annulé par $I$,
alors (1) est vérifié. Ainsi (4) implique (1), ce qui
achève la démonstration.
\end{proof}

\begin{lemma}
\label{lemma-two-term-base-change}
Soit $R$ un anneau. Soit $K$ un objet de $D(R)$ avec $H^i(K) = 0$
pour $i \not \in \{-1, 0\}$. Soit $K^{-1} \to K^0$ un complexe à deux termes
de $R$-modules représentant $K$, tel que $K^0$ soit un $R$-module plat
(par exemple projectif ou libre). Soit $R \to R'$ un morphisme d'anneaux.
Alors le complexe $K^\bullet \otimes_R R'$ représente
$\tau_{\geq -1}(K \otimes_R^\mathbf{L} R')$.
\end{lemma}

\begin{proof}
Nous avons un triangle distingué
$$
K^0 \to K^\bullet \to K^{-1}[1] \to K^0[1]
$$
dans $D(R)$. Il détermine un morphisme de triangles distingués
$$
\xymatrix{
K^0 \otimes_R^\mathbf{L} R' \ar[d] \ar[r] &
K^\bullet \otimes_R^\mathbf{L} R' \ar[r] \ar[d] &
K^{-1} \otimes_R^\mathbf{L} R'[1] \ar[r] \ar[d] &
K^0 \otimes_R^\mathbf{L} R'[1] \ar[d] \\
K^0 \otimes_R R' \ar[r] &
K^\bullet \otimes_R R' \ar[r] &
K^{-1} \otimes_R R'[1] \ar[r] &
K^0 \otimes_R R'[1]
}
$$
Les flèches verticales de gauche et de droite sont des isomorphismes puisque $K^0$ est plat.
Comme $K^{-1} \otimes_R^\mathbf{L} R' \to K^{-1} \otimes_R R'$
est un isomorphisme sur la cohomologie en degré $0$, nous concluons.
\end{proof}

\begin{lemma}
\label{lemma-base-change-property-ext-1-annihilated}
Soit $I$ un idéal d'un anneau $R$. Soit $K$ un objet de $D(R)$ avec
$H^i(K) = 0$ pour $i \not \in \{-1, 0\}$. Soit $R \to R'$ un morphisme d'anneaux.
Si $K$ satisfait les conditions équivalentes (1), (2) et (3)
du lemme \ref{lemma-ext-1-annihilated-definite} relativement à $(R, I)$,
alors $\tau_{\geq -1}(K \otimes_R^\mathbf{L} R')$
satisfait les conditions équivalentes (1), (2) et (3)
du lemme \ref{lemma-ext-1-annihilated-definite} relativement à $(R', IR')$.
\end{lemma}

\begin{proof}
Nous pouvons supposer que $K$ est représenté par un complexe à deux termes $K^{-1} \to K^0$
avec $K^0$ libre, tel que pour tout $a \in I$ l'application $a : K^{-1} \to K^{-1}$
soit égale à $h_a \circ d_K^{-1}$ pour une application $h_a : K^0 \to K^{-1}$.
D'après le lemme \ref{lemma-two-term-base-change},
$\tau_{\geq -1}(K \otimes_R^\mathbf{L} R')$ est représenté par
$K^\bullet \otimes_R R'$. Alors, pour tout $a \in I$, nous voyons bien entendu que
$a \otimes 1 : K^{-1} \otimes_R R' \to K^{-1} \otimes_R R'$
est égal à $(h_a \otimes 1) \circ (d_K^{-1} \otimes 1)$.
Comme l'ensemble des applications $K^{-1} \otimes_R R' \to K^{-1} \otimes_R R'$
qui se factorisent par $\text{d}_K^{-1} \otimes 1$
forme un $R'$-module, nous concluons.
\end{proof}

\begin{lemma}
\label{lemma-two-term-surjection-map-zero}
Soit $R$ un anneau. Soit $\alpha : K \to K'$ un morphisme de $D(R)$. Supposons
\begin{enumerate}
\item $H^i(K) = H^i(K') = 0$ pour $i \not \in \{-1, 0\}$ ;
\item $H^0(\alpha)$ est un isomorphisme et $H^{-1}(\alpha)$ est surjectif.
\end{enumerate}
Pour tout $f \in R$, si $f : K \to K$ est $0$, alors $f : K' \to K'$ est $0$.
\end{lemma}

\begin{proof}
Posons $M = \Ker(H^{-1}(\alpha))$. Alors $\alpha$ s'insère dans un triangle
distingué
$$
M[1] \to K \to K' \to M[2]
$$
Comme $K \to K' \xrightarrow{f} K'$ est nul par hypothèse, nous voyons
que $f : K' \to K'$ se factorise par une application $M[2] \to K'$. Cependant,
$\Hom(M[2], K') = 0$, par exemple d'après
Catégories dérivées, lemme \ref{derived-lemma-negative-exts}.
\end{proof}

\begin{lemma}
\label{lemma-surjection-property-ext-1-annihilated}
Soit $I$ un idéal d'un anneau $R$. Soit $\alpha : K \to K'$
un morphisme de $D(R)$. Supposons
\begin{enumerate}
\item $H^i(K) = H^i(K') = 0$ pour $i \not \in \{-1, 0\}$ ;
\item $H^0(\alpha)$ est un isomorphisme et $H^{-1}(\alpha)$ est surjectif.
\end{enumerate}
Si $K$ satisfait les conditions équivalentes (1), (2) et (3)
du lemme \ref{lemma-ext-1-annihilated-definite},
alors $K'$ les satisfait aussi.
\end{lemma}

\begin{proof}
Posons $M = \Ker(H^{-1}(\alpha))$. Alors $\alpha$ s'insère dans un triangle
distingué
$$
M[1] \to K \to K' \to M[2]
$$
Pour tout $R$-module $N$, cela détermine une suite exacte
$$
\Ext^0_R(M[1], N) \to
\Ext^1_R(K', N) \to
\Ext^1_R(K, N)
$$
Comme $\Ext^0_R(M[1], N) = \Ext^{-1}_R(M, N) = 0$, nous voyons que
$\Ext^1_R(K', N)$ est un sous-module de $\Ext^1_R(K, N)$. Ainsi, si
$\Ext^1_R(K, N)$ est annulé par $I$, il en va de même de $\Ext^1_R(K', N)$.
\end{proof}

\begin{lemma}
\label{lemma-ext-1-annihilated}
Soit $R$ un anneau et soit $I \subset R$ un idéal.
Soit $K \in D(R)$ tel que $H^i(K) = 0$ pour $i \not \in \{-1, 0\}$.
Les assertions suivantes sont équivalentes :
\begin{enumerate}
\item il existe $c \geq 0$ tel que les conditions équivalentes
(1), (2), (3) du Lemme \ref{lemma-ext-1-annihilated-definite}
soient satisfaites pour $K$ et l'idéal $I^c$,
\item il existe $c \geq 0$ tel que (a) $I^c$ annule
$H^{-1}(K)$ et (b) $H^0(K)$ soit un module $I^c$-projectif (voir la
Section \ref{section-near-projective}).
\end{enumerate}
Si $R$ est noethérien et si $H^i(K)$ est un $R$-module fini
pour $i = -1, 0$, ces assertions sont aussi équivalentes aux suivantes :
\begin{enumerate}
\item[(3)] il existe $c \geq 0$ tel que les conditions équivalentes
(4), (5) du Lemme \ref{lemma-ext-1-annihilated-definite}
soient satisfaites pour $K$ et l'idéal $I^c$,
\item[(4)] $H^{-1}(K)$ est de $I$-torsion et il existe
$f_1, \ldots, f_s \in R$ tels que $V(f_1, \ldots, f_s) \subset V(I)$
et que les localisés $H^0(K)_{f_i}$ soient des
$R_{f_i}$-modules projectifs,
\item[(5)] $H^{-1}(K)$ est de $I$-torsion et il existe
$f_1, \ldots, f_s \in I$ tels que $V(f_1, \ldots, f_s) = V(I)$
et que les localisés $H^0(K)_{f_i}$ soient des
$R_{f_i}$-modules projectifs, et
\item[(6)] $H^{-1}(K)$ est de $I$-torsion et, pour tous
$f_1, \ldots, f_s \in I$ tels que $V(f_1, \ldots, f_s) = V(I)$,
les localisés $H^0(K)_{f_i}$ sont des $R_{f_i}$-modules projectifs.
\end{enumerate}
\end{lemma}

\begin{proof}
Le triangle distingué $H^{-1}(K)[1] \to K \to H^0(K)[0] \to H^{-1}(K)[2]$
détermine une suite exacte
$$
0 \to \Ext^1_R(H^0(K), N) \to \Ext^1_R(K, N) \to \Hom_R(H^{-1}(K), N) \to
\Ext^2_R(H^0(K), N)
$$
Ainsi, (2) implique que $I^{2c}$ annule $\Ext^1_R(K, N)$ pour tout
$R$-module $N$. Si l'on suppose (1), on voit immédiatement que $H^0(K)$
est $I^c$-projectif. D'autre part, on peut choisir une application injective
$H^{-1}(K) \to N$ vers un $R$-module injectif $N$. Cette application
est alors l'image d'un élément de $\Ext^1_R(K, N)$, par l'annulation
du $\Ext^2$ dans la suite, et l'on conclut que $H^{-1}(K)$
est annulé par $I^c$.

\medskip\noindent
Supposons $R$ noethérien et $H^i(K)$ un $R$-module fini
pour $i = -1, 0$. Par le Lemme \ref{lemma-ext-1-annihilated-definite},
on voit que (3) est équivalent à (1) et (2). De plus, si (3)
est satisfaite, alors, pour $f \in I$, la multiplication par $f$ sur
$H^0(K)$ se factorise par un module projectif, ce qui implique
que $H^0(K)_f$ est un facteur direct d'un $R_f$-module projectif et est donc
lui-même un $R_f$-module projectif. Les conditions équivalentes
(1), (2) et (3) impliquent donc (6).
Bien entendu, (6) implique (5), et (5) implique trivialement (4).

\medskip\noindent
Supposons (4). Puisque $H^{-1}(K)$ est un $R$-module fini et de $I$-torsion,
on voit que $I^{c_1}$ annule $H^{-1}(K)$ pour un certain $c_1 \geq 0$.
Choisissons une suite exacte courte
$$
0 \to M \to R^{\oplus r} \to H^0(K) \to 0
$$
qui détermine un élément $\xi \in \Ext^1_R(H^0(K), M)$.
Pour tout $f \in I$, on a $\Ext^1_R(H^0(K), M)_f = \Ext^1_{R_f}(H^0(K)_f, M_f)$
par le Lemme \ref{lemma-pseudo-coherence-and-base-change-ext}.
Par conséquent, si $H^0(K)_f$ est projectif, une puissance de $f$ annule $\xi$.
On en conclut que $\xi$ est annulé par $(f_1, \ldots, f_s)^{c_2}$
pour un certain $c_2 \geq 0$. Puisque $V(f_1, \ldots, f_s) \subset V(I)$, on a
$\sqrt{I} \subset (f_1, \ldots, f_s)$
(Algèbre, Lemme \ref{algebra-lemma-Zariski-topology}).
Comme $R$ est noethérien, on trouve $I^{c_3} \subset (f_1, \ldots, f_s)$
pour un certain $c_3 \geq 0$ (Algèbre, Lemme \ref{algebra-lemma-Noetherian-power}).
Ainsi, $I^{c2c3}$ annule $\xi$.
Ceci signifie à son tour que $H^0(K)$ est $I^{c_2c_3}$-projectif (car la multiplication
par $a \in I$ qui annule $\xi$ se factorise par $R^{\oplus r}$).
En prenant donc $c = \max(c_1, c_2c_3)$, on voit que (2) est satisfaite.
\end{proof}

\begin{lemma}
\label{lemma-zero-in-derived}
Soit $R$ un anneau. Soient $K_j \in D(R)$, $j = 1, 2, 3$, tels que $H^i(K_j) = 0$
pour $i \not \in \{-1, 0\}$. Soient $\varphi : K_1 \to K_2$ et
$\psi : K_2 \to K_3$ des morphismes dans $D(R)$.
Si $H^0(\varphi) = 0$ et $H^{-1}(\psi) = 0$, alors
$\varphi \circ \psi = 0$.
\end{lemma}

\begin{proof}
Appliquons Catégories dérivées, Lemme \ref{derived-lemma-trick-vanishing-composition},
pour voir que $\varphi \circ \psi$ se factorise par $\tau_{\leq -2}K_2 = 0$.
\end{proof}

\begin{lemma}
\label{lemma-silly}
Soit $R$ un anneau. Soit $K \in D(R)$ donné par un complexe à deux termes
de la forme $R^{\oplus n} \to R^{\oplus n}$. Notons
$A \in \text{Mat}(n \times n, R)$ la matrice de la différentielle.
Alors $\det(a) : K \to K$ est nul dans $D(R)$.
\end{lemma}

\begin{proof}
Omis. Bon exercice.
\end{proof}






\section{Le complexe cotangent naïf}
\label{section-netherlander}

\noindent
Dans cette section, nous poursuivons la discussion commencée dans
Algèbre, Section \ref{algebra-section-netherlander}.
Nous commençons par étudier le changement de base.
Le premier lemme montre que prendre le produit tensoriel naïf
du complexe cotangent naïf avec une extension d'anneaux
n'est pas tout à fait aussi naïf que l'on pourrait le penser.

\begin{lemma}
\label{lemma-tensor-NL}
Soient $R \to S$ et $S \to S'$ des morphismes d'anneaux. Le morphisme canonique
$\NL_{S/R} \otimes_S^\mathbf{L} S' \to \NL_{S/R} \otimes_S S'$
induit un isomorphisme
$\tau_{\geq -1}(\NL_{S/R} \otimes_S^\mathbf{L} S') \to \NL_{S/R} \otimes_S S'$
dans $D(S')$. De même, étant donnée une présentation $\alpha$ de $S$ sur $R$,
le morphisme canonique
$\NL(\alpha) \otimes_S^\mathbf{L} S' \to \NL(\alpha) \otimes_S S'$
induit un isomorphisme $\tau_{\geq -1}(\NL(\alpha) \otimes_S^\mathbf{L} S') \to
\NL(\alpha) \otimes_S S'$ dans $D(S')$.
\end{lemma}

\begin{proof}
Cas particulier du Lemme \ref{lemma-two-term-base-change}.
\end{proof}

\begin{lemma}
\label{lemma-base-change-NL}
Soient $R \to S$ et $R \to R'$ des morphismes d'anneaux.
Soit $\alpha : P \to S$ une présentation de $S$ sur $R$.
Alors $\alpha' : P \otimes_R R' \to S \otimes_R R'$ est une
présentation de $S' = S \otimes_R R'$ sur $R'$.
Le morphisme canonique
$$
NL(\alpha) \otimes_S S' \to \NL(\alpha')
$$
est un isomorphisme sur $H^0$ et est surjectif sur $H^{-1}$. En particulier,
le morphisme canonique
$$
\NL_{S/R} \otimes_S S' \to \NL_{S'/R'}
$$
est un isomorphisme sur $H^0$ et est surjectif sur $H^{-1}$.
\end{lemma}

\begin{proof}
Notons $I = \Ker(P \to S)$. Notons $P' = P \otimes_R R'$ et
$I' = \Ker(P' \to S')$. Supposons que $P$ soit l'algèbre polynomiale
en les $x_j$ pour $j \in J$. Le morphisme affiché dans le lemme devient
$$
\xymatrix{
\bigoplus_{j \in J} S' \text{d}x_j \ar[r] &
\bigoplus_{j \in J} S' \text{d}x_j \\
I/I^2 \otimes_S S' \ar[r] \ar[u] &
I'/(I')^2 \ar[u]
}
$$
où la colonne de gauche est $\NL(\alpha) \otimes_S S'$ et celle de droite
est $\NL(\alpha')$. Par l'exactitude à droite du produit tensoriel, on
voit que $I \otimes_R R' \to I'$ est surjectif. La flèche inférieure
est donc une surjection. Ceci démontre la première assertion du lemme.
L'assertion concernant $\NL_{S/R} \otimes_S S' \to \NL_{S'/R'}$
en découle, car ces complexes sont homotopes à $\NL(\alpha) \otimes_S S'$ et
$\NL(\alpha')$.
\end{proof}

\begin{lemma}
\label{lemma-base-change-NL-flat}
Considérons un diagramme cocartésien d'anneaux
$$
\xymatrix{
B \ar[r] & B' \\
A \ar[r] \ar[u] & A' \ar[u]
}
$$
Si $B$ est plat sur $A$, alors le morphisme canonique
$\NL_{B/A} \otimes_B B' \to \NL_{B'/A'}$ est un quasi-isomorphisme.
Si, de plus, $\NL_{B/A}$ est d'amplitude de Tor contenue dans $[-1, 0]$,
alors $\NL_{B/A} \otimes_B^\mathbf{L} B' \to \NL_{B'/A'}$
est aussi un quasi-isomorphisme.
\end{lemma}

\begin{proof}
Choisissons une présentation $\alpha : P \to B$ comme dans
Algèbre, Section \ref{algebra-section-netherlander}.
Posons $I = \Ker(\alpha)$. Posons $P' = P \otimes_A A'$ et notons
$\alpha' : P' \to B'$ la présentation correspondante de $B'$ sur $A'$.
Comme $B$ est plat sur $A$, on voit que $I' = \Ker(\alpha')$ est égal
à $I \otimes_A A'$. Par conséquent,
$$
I'/(I')^2 = \Coker(I^2 \otimes_A A' \to I \otimes_A A') =
I/I^2 \otimes_A A' = I/I^2 \otimes_B B'
$$
On a $\Omega_{P'/A'} = \Omega_{P/A} \otimes_A A'$
car les deux membres ont la même base. Il s'ensuit que
$\Omega_{P'/A'} \otimes_{P'} B' = \Omega_{P/A} \otimes_P B \otimes_B B'$.
Ceci démontre que $\NL(\alpha) \otimes_B B' \to \NL(\alpha')$
est un isomorphisme de complexes, d'où la première assertion.

\medskip\noindent
On a
$$
\NL(\alpha) = I/I^2 \longrightarrow \Omega_{P/A} \otimes_P B
$$
comme complexe de $B$-modules, avec $I/I^2$ placé en degré $-1$.
Puisque le terme en degré $0$ est libre, ce complexe est d'amplitude de Tor
contenue dans $[-1, 0]$ si et seulement si $I/I^2$ est un $B$-module plat, voir le
Lemme \ref{lemma-last-one-flat}.
Si tel est le cas, alors $\NL(\alpha) \otimes_B^\mathbf{L} B' =
\NL(\alpha) \otimes_B B'$, et l'on obtient la seconde assertion.
\end{proof}

\begin{lemma}
\label{lemma-lci-NL}
Soit $A \to B$ une intersection complète locale au sens de la
Définition \ref{definition-local-complete-intersection}.
Alors $\NL_{B/A}$ est un objet parfait de
$D(B)$ d'amplitude de Tor contenue dans $[-1, 0]$.
\end{lemma}

\begin{proof}
Écrivons $B = A[x_1, \ldots, x_n]/I$. Alors $\NL_{B/A}$ est représenté par
le complexe
$$
I/I^2 \longrightarrow \bigoplus B \text{d}x_i
$$
de $B$-modules, avec $I/I^2$ placé en degré $-1$. Puisque le terme en
degré $0$ est libre fini, ce complexe est d'amplitude de Tor contenue dans
$[-1, 0]$ si et seulement si $I/I^2$ est un $B$-module plat, voir le
Lemme \ref{lemma-last-one-flat}. Par définition, $I$ est un idéal régulier
au sens de Koszul et donc un idéal quasi-régulier, voir la Section \ref{section-ideals}.
Ainsi, $I/I^2$ est un $B$-module projectif fini
(Lemme \ref{lemma-quasi-regular-ideal-finite-projective}),
et l'on conclut à la fois que $\NL_{B/A}$ est parfait et qu'il est d'amplitude de Tor
contenue dans $[-1, 0]$.
\end{proof}

\begin{lemma}
\label{lemma-base-change-lci-bis}
Considérons un diagramme cocartésien d'anneaux
$$
\xymatrix{
B \ar[r] & B' \\
A \ar[r] \ar[u] & A' \ar[u]
}
$$
Si $A \to B$ et $A' \to B'$ sont des intersections complètes locales au sens de la
Définition \ref{definition-local-complete-intersection}, alors
le noyau de $H^{-1}(\NL_{B/A} \otimes_B B') \to H^{-1}(\NL_{B'/A'})$
est un $B'$-module projectif fini.
\end{lemma}

\begin{proof}
Par le Lemme \ref{lemma-lci-NL}, les complexes $\NL_{B/A}$ et $\NL_{B'/A'}$
sont parfaits d'amplitude de Tor contenue dans $[-1, 0]$.
En combinant les Lemmes \ref{lemma-tensor-NL}, \ref{lemma-pull-perfect} et
\ref{lemma-pull-tor-amplitude}, on a
$\NL_{B/A} \otimes_B B' = \NL_{B/A} \otimes_B^\mathbf{L} B'$
et ce complexe est aussi parfait d'amplitude de Tor contenue dans $[-1, 0]$.
Choisissons un triangle distingué
$$
C \to \NL_{B/A} \otimes_B B'  \to \NL_{B'/A'} \to C[1]
$$
dans $D(B')$. Par les Lemmes \ref{lemma-two-out-of-three-perfect} et
\ref{lemma-cone-tor-amplitude}, on conclut que $C$ est parfait
d'amplitude de Tor contenue dans $[-1, 1]$. Par le Lemme \ref{lemma-base-change-NL},
le complexe $C$ n'a qu'un seul module de cohomologie non nul, à savoir le module
du lemme placé en degré $-1$. Ce module est de présentation finie
(Lemme \ref{lemma-n-pseudo-module}) et plat
(Lemme \ref{lemma-tor-dimension}). Il est donc projectif fini par
Algèbre, Lemme \ref{algebra-lemma-finite-projective}.
\end{proof}












\section{Rlim de groupes abéliens}
\label{section-Rlim}

\noindent
Nous examinons brièvement $R\lim$ sur les groupes abéliens.
Dans cette section, nous noterons $\textit{Ab}(\mathbf{N})$ la
catégorie abélienne des systèmes projectifs de groupes abéliens.
Cette notation est compatible avec celle des faisceaux de groupes abéliens
sur un site, car un système projectif de groupes abéliens est
la même chose qu'un faisceau de groupes sur la catégorie $\mathbf{N}$
(avec un unique morphisme $i \to j$ si $i \leq j$), voir la
Remarque \ref{remark-Rlim-cohomology}. De nombreux arguments de cette
section reprennent ceux employés pour construire la machinerie cohomologique
des faisceaux de groupes abéliens sur les sites.

\begin{lemma}
\label{lemma-compute-Rlim}
Le foncteur $\lim : \textit{Ab}(\mathbf{N}) \to \textit{Ab}$
admet un foncteur dérivé droit
\begin{equation}
\label{equation-Rlim}
R\lim : D(\textit{Ab}(\mathbf{N})) \longrightarrow D(\textit{Ab})
\end{equation}
Comme d'habitude, nous posons $R^p\lim(K) = H^p(R\lim(K))$. En outre :
\begin{enumerate}
\item pour tout $(A_n)$ dans $\textit{Ab}(\mathbf{N})$, on a
$R^p\lim A_n = 0$ pour $p > 1$,
\item l'objet $R\lim A_n$ de $D(\textit{Ab})$ est représenté
par le complexe
$$
\prod A_n \to \prod A_n,\quad (x_n) \mapsto (x_n - f_{n + 1}(x_{n + 1}))
$$
placé en degrés $0$ et $1$,
\item si $(A_n)$ est ML, alors $R^1\lim A_n = 0$, c'est-à-dire que $(A_n)$
est acyclique à droite pour $\lim$,
\item tout $K^\bullet \in D(\textit{Ab}(\mathbf{N}))$ est quasi-isomorphe
à un complexe dont les termes sont acycliques à droite pour $\lim$, et
\item si chaque $K^p = (K^p_n)$ est acyclique à droite pour $\lim$, c'est-à-dire si
$R^1\lim_n K^p_n = 0$, alors $R\lim K$ est représenté par le
complexe dont le terme en degré $p$ est $\lim_n K_n^p$.
\end{enumerate}
\end{lemma}

\begin{proof}
Soit $(A_n)$ un système projectif quelconque. Soit $(B_n)$ le système
projectif défini par
$$
B_n = A_n \oplus A_{n - 1} \oplus \ldots \oplus A_1
$$
et dont les morphismes de transition sont donnés par les projections. Définissons
$A_n \to B_n$ par $(1, f_n, f_{n - 1} \circ f_n, \ldots, f_2 \circ \ldots \circ f_n)$,
où les $f_i : A_i \to A_{i - 1}$ sont les morphismes de transition.
On voit ainsi que tout système projectif est un sous-objet d'un
système ML (Homologie, Section \ref{homology-section-inverse-systems}).
Il résulte de
Catégories dérivées, Lemme \ref{derived-lemma-subcategory-right-acyclics},
en utilisant Homologie, Lemme \ref{homology-lemma-Mittag-Leffler},
que tout système ML est acyclique à droite pour $\lim$, c'est-à-dire que (3) est satisfaite.
Ceci implique déjà que $RF$ est défini sur $D^+(\textit{Ab}(\mathbf{N}))$,
voir Catégories dérivées, Proposition \ref{derived-proposition-enough-acyclics}.
Posons $C_n = A_{n - 1} \oplus \ldots \oplus A_1$ pour $n > 1$ et
$C_1 = 0$, avec également pour morphismes de transition les projections.
Il existe alors une suite exacte courte de systèmes projectifs
$0 \to (A_n) \to (B_n) \to (C_n) \to 0$, où $B_n \to C_n$
est donné par $(x_i) \mapsto (x_i - f_{i + 1}(x_{i + 1}))$.
Puisque $(C_n)$ est aussi ML, on conclut que (2) est satisfaite
(par la proposition citée ci-dessus), ce qui implique aussi (1).
Enfin, ceci implique, par Catégories dérivées, Lemme
\ref{derived-lemma-unbounded-right-derived}
que $R\lim$ est en fait défini sur tout $D(\textit{Ab}(\mathbf{N}))$.
En effet, la preuve de Catégories dérivées, Lemme
\ref{derived-lemma-unbounded-right-derived}
procède en démontrant les assertions (4) et (5).
\end{proof}

\begin{lemma}
\label{lemma-six-term-Rlim}
Soit
$$
0 \to (A_i) \to (B_i) \to (C_i) \to 0
$$
une suite exacte courte de systèmes projectifs de groupes abéliens.
Il lui est alors associée une suite exacte à $6$ termes
$0 \to \lim A_i \to \lim B_i \to \lim C_i \to
R^1\lim A_i \to R^1\lim B_i \to R^1\lim C_i \to 0$.
\end{lemma}

\begin{proof}
Cela résulte de l'annulation du Lemme \ref{lemma-compute-Rlim}.
\end{proof}

\noindent
Voici la formulation « correcte » de
Homologie, Lemme \ref{homology-lemma-apply-Mittag-Leffler-again}.

\begin{lemma}
\label{lemma-apply-Mittag-Leffler-again}
Soit
$$
(A^{-2}_n \to A^{-1}_n \to A^0_n \to A^1_n)
$$
un système projectif de complexes de groupes abéliens, et notons
$A^{-2} \to A^{-1} \to A^0 \to A^1$ sa limite. Notons
$(H_n^{-1})$, $(H_n^0)$ les systèmes projectifs de cohomologies, et
notons $H^{-1}$, $H^0$ les cohomologies de $A^{-2} \to A^{-1} \to A^0 \to A^1$.
Si
\begin{enumerate}
\item $(A^{-2}_n)$ et $(A^{-1}_n)$ ont un $R^1\lim$ nul,
\item $(H^{-1}_n)$ a un $R^1\lim$ nul,
\end{enumerate}
alors $H^0 = \lim H_n^0$.
\end{lemma}

\begin{proof}
Soit $K \in D(\textit{Ab}(\mathbf{N}))$ l'objet représenté
par le système de complexes dont la $n$-ième composante
est le complexe $A^{-2}_n \to A^{-1}_n \to A^0_n \to A^1_n$.
Nous allons calculer $H^0(R\lim K)$ à l'aide des deux suites
spectrales\footnote{Pour utiliser ces suites spectrales, il faut montrer
que $\textit{Ab}(\mathbf{N})$ possède assez d'injectifs.
Un système projectif $(I_n)$ de groupes abéliens est injectif si et seulement si
chaque $I_n$ est un groupe abélien injectif et si les morphismes de transition sont
des surjections scindées. Tout système se plonge dans un tel système. Détails omis.} de
Catégories dérivées, Lemme \ref{derived-lemma-two-ss-complex-functor}.
La première a pour page $E_1$
$$
\begin{matrix}
0 & 0 & R^1\lim A^0_n  & R^1\lim A^1_n \\
A^{-2} & A^{-1} & A^0  & A^1
\end{matrix}
$$
avec des différentielles horizontales, et toutes les différentielles supérieures sont nulles.
La seconde a pour page $E_2$
$$
\begin{matrix}
R^1\lim H^{-2}_n & 0 & R^1\lim H^0_n & R^1 \lim H^1_n \\
\lim H^{-2}_n &
\lim H^{-1}_n &
\lim H^0_n &
\lim H^1_n
\end{matrix}
$$
et dégénère à ce stade. Le résultat en découle.
\end{proof}

\begin{lemma}
\label{lemma-pro-isomorphism-iso-Rlim}
Soient $(A_n)$ et $(B_n)$ des systèmes projectifs de groupes abéliens.
Un morphisme de pro-systèmes $\varphi : (A_n) \to (B_n)$ détermine des
morphismes $\lim A_n \to \lim B_n$ et $R^1\lim A_n \to R^1\lim B_n$.
Ces morphismes sont des isomorphismes si $\varphi$ est un pro-isomorphisme.
\end{lemma}

\begin{proof}
Voir
Catégories, Exemple \ref{categories-example-pro-morphism-inverse-systems},
pour une discussion des morphismes de pro-systèmes.
Le morphisme $\varphi$ est donné par $1 \leq m_1 < m_2 < m_3 < \ldots$
et par des morphismes $\varphi_n : A_{m_n} \to B_n$ compatibles aux morphismes
de transition. Posons $\varphi', \varphi'' : \prod A_n \to \prod B_n$
égaux à
$\varphi'((a_n)) = (\varphi_n(a_{m_n}))$ et
$$
\varphi''((a_n)) =
(\varphi_n(a_{m_n} + f(a_{m_n + 1}) + \ldots + f(a_{m_{n + 1} - 1})))
$$
où chaque occurrence de $f$ désigne un morphisme de transition convenable du
système projectif $(A_n)$. Alors le diagramme
$$
\xymatrix{
\prod A_n \ar[r]_\delta \ar[d]^{\varphi'} &
\prod A_n \ar[d]^{\varphi''} \\
\prod B_n \ar[r]^\delta & \prod B_n
}
$$
où les flèches horizontales sont celles du Lemme \ref{lemma-compute-Rlim}
est commutatif. On obtient ainsi les morphismes voulus. La construction
est fonctorielle au sens suivant : si l'on se donne un système projectif $(C_n)$,
des entiers $1 \leq m'_1 < m'_2 < m'_3 < \ldots$ et des morphismes
$\psi_n : B_{m'_n} \to C_n$ compatibles aux morphismes de transition, alors
$(\psi \circ \varphi)' = \psi' \circ \varphi'$ and
$(\psi \circ \varphi)'' = \psi'' \circ \varphi''$
où la composée $\psi \circ \varphi$ correspond aux entiers
$1 \leq m_{m'_1} < m_{m'_2} < \ldots$ et aux morphismes
$\psi_n \circ \varphi_{m'_n} : A_{m_{m'_n}} \to C_n$. Nous allons montrer
que si $B_n = A_n$ et si $\varphi_n$ est le morphisme de transition, alors
les morphismes obtenus sont les morphismes identiques. Ceci démontrera à la fois que
la construction est indépendante du choix du représentant
$(m_n, \varphi_n)$ de $\varphi$ et la dernière assertion du lemme.

\medskip\noindent
Prenons donc $B_n = A_n$ et pour $\varphi_n$ le morphisme de transition.
Soit $(a_n) \in \lim A_n$ un élément. Alors $\varphi'((a_n))$
est l'élément dont la composante d'indice $n$ est l'image de $a_{m_n}$,
qui est égale à $a_n$. Ceci démontre l'assertion pour $\lim A_n$.
Soit $\xi \in R^1\lim A_n$ la classe de l'élément
$(a_n)$ de $\prod A_n$. Considérons l'élément $(b_n)$ défini par
$$
b_i = a_i + f(a_{i + 1}) + \ldots + f(a_{m_n - 1})
$$
si $m_{n - 1} < i < m_n$, et par $0$ si $i = m_n$ pour un certain $n$.
Alors $(a'_n) = (a_n) - \delta((b_n))$ définit la même classe
dans $R^1\lim A_n$, et un calcul montre que $a'_i = 0$
sauf si $i = m_n$ pour un certain $n$. On peut donc supposer, et l'on suppose, que
$a_i = 0$ sauf si $i = m_n$ pour un certain $n$. Remarquons que
$$
(0, \ldots, -f(a_{m_j}), 0, \ldots, 0, a_{m_j}, 0, \ldots)
$$
dont les composantes non nulles sont aux places $j$ et $m_j$, est l'image par $\delta$ de
$$
c_j = (0, \ldots, 0, f(a_{m_j}), f(a_{m_j}), \ldots, a_{m_j}, 0, \ldots)
$$
dont les composantes non nulles sont aux places $j + 1, \ldots, m_j$. La somme
$c = \sum c_j$ a un sens dans $\prod A_n$. En rappelant que
$\varphi''((a_n)) = (a_{m_1}, a_{m_2}, \ldots)$, on voit que
$$
(a_n) - \varphi''((a_n)) = \delta(c)
$$
ce qui achève la preuve.
\end{proof}

\begin{lemma}
\label{lemma-map-into-Rlim}
Soit $\mathcal{D}$ une catégorie triangulée. Soit
$(K_n)$ un système projectif d'objets de $\mathcal{D}$.
Soit $K$ une limite dérivée du système $(K_n)$.
Alors, pour tout $L$ dans $\mathcal{D}$, on a une suite exacte courte
$$
0 \to R^1\lim \Hom_\mathcal{D}(L, K_n[-1]) \to
\Hom_\mathcal{D}(L, K) \to
\lim \Hom_\mathcal{D}(L, K_n) \to 0
$$
\end{lemma}

\begin{proof}
Cela résulte de
Catégories dérivées, Définition \ref{derived-definition-derived-limit} et
Lemme \ref{derived-lemma-representable-homological},
ainsi que de la description de $\lim$ et $R^1\lim$ dans le
Lemme \ref{lemma-compute-Rlim} ci-dessus.
\end{proof}

\begin{lemma}
\label{lemma-pro-isomorphism-Rlim}
Soit $\mathcal{D}$ une catégorie triangulée. Soient
$(K_n)$ et $(M_n)$ des systèmes projectifs d'objets de $\mathcal{D}$
de limites dérivées $K$ et $M$. Soit $a : (K_n) \to (M_n)$ un
pro-isomorphisme de pro-objets. On peut alors utiliser $a$
pour produire un isomorphisme (non canonique) $K \to M$.
\end{lemma}

\begin{proof}
On obtient une flèche $K \to M$ s'insérant dans un morphisme
$$
\xymatrix{
K \ar[r] \ar[d] & \prod K_n \ar[r] \ar[d]^{a'} &
\prod K_n \ar[d]^{a''} \\
M \ar[r] & \prod M_n \ar[r] & \prod M_n
}
$$
des triangles distingués définissants, par
Catégories dérivées, Remarque \ref{derived-remark-functoriality-derived-limit}.
Ainsi, pour tout objet $L$ de $\mathcal{D}$, on obtient un morphisme
de suites exactes courtes
$$
\xymatrix{
0 \ar[r] &
R^1\lim \Hom_\mathcal{D}(L, K_n[-1]) \ar[r] \ar[d] &
\Hom_\mathcal{D}(L, K) \ar[r] \ar[d] &
\lim \Hom_\mathcal{D}(L, K_n) \ar[r] \ar[d] &
0 \\
0 \ar[r] &
R^1\lim \Hom_\mathcal{D}(L, M_n[-1]) \ar[r] &
\Hom_\mathcal{D}(L, M) \ar[r] &
\lim \Hom_\mathcal{D}(L, M_n) \ar[r] &
0
}
$$
voir le Lemme \ref{lemma-map-into-Rlim} et sa preuve. Par le
Lemme \ref{lemma-pro-isomorphism-iso-Rlim},
les flèches verticales de gauche et de droite sont des isomorphismes\footnote{Si
le pro-isomorphisme $a$ est donné par des morphismes $a_n : K_{m_n} \to M_n$,
on obtient alors un pro-isomorphisme sur $\Hom$ à l'aide des morphismes correspondants
$\Hom_\mathcal{D}(L, K_{m_n}) \to \Hom_\mathcal{D}(L, M_n)$.
Ainsi, la construction des flèches sur $\lim$ et $R^1\lim$
dans la preuve du Lemme \ref{lemma-pro-isomorphism-iso-Rlim}
coïncide avec la construction de $a'$ et $a''$ dans
Catégories dérivées, Remarque \ref{derived-remark-functoriality-derived-limit}.}.
La flèche du milieu est donc un isomorphisme. Par le lemme de Yoneda,
on voit que le morphisme $K \to M$ est un isomorphisme.
\end{proof}

\begin{lemma}
\label{lemma-map-from-hocolim}
Soit $\mathcal{D}$ une catégorie triangulée. Soit
$(K_n)$ un système d'objets de $\mathcal{D}$.
Soit $K$ une colimite dérivée du système $(K_n)$.
Alors, pour tout $L$ dans $\mathcal{D}$, on a une suite exacte courte
$$
0 \to R^1\lim \Hom_\mathcal{D}(K_n, L[-1]) \to
\Hom_\mathcal{D}(K, L) \to
\lim \Hom_\mathcal{D}(K_n, L) \to 0
$$
\end{lemma}

\begin{proof}
Cela résulte de
Catégories dérivées, Définition \ref{derived-definition-derived-colimit} et
Lemme \ref{derived-lemma-representable-homological},
ainsi que de la description de $\lim$ et $R^1\lim$ dans le
Lemme \ref{lemma-compute-Rlim} ci-dessus.
\end{proof}

\begin{remark}[Rlim comme cohomologie]
\label{remark-Rlim-cohomology}
Considérons la catégorie $\mathbf{N}$ dont les objets sont les entiers naturels et
dont les morphismes sont les uniques flèches $i \to j$ si $j \geq i$. Munissons $\mathbf{N}$
de la topologie chaotique (Sites, Exemple \ref{sites-example-indiscrete}), de sorte
qu'un faisceau $\mathcal{F}$ soit la même chose qu'un système projectif
$$
\mathcal{F}_1 \leftarrow \mathcal{F}_2 \leftarrow \mathcal{F}_3
\leftarrow \ldots
$$
d'ensembles sur $\mathbf{N}$. Remarquons que
$\Gamma(\mathbf{N}, \mathcal{F}) = \lim \mathcal{F}_n$. Pour un système
projectif de groupes abéliens $\mathcal{F}_n$, on a
$$
R^p\lim \mathcal{F}_n = H^p(\mathbf{N}, \mathcal{F})
$$
car les deux membres sont les foncteurs dérivés droits supérieurs de
$\mathcal{F} \mapsto \lim \mathcal{F}_n = H^0(\mathbf{N}, \mathcal{F})$.
Ainsi, l'existence de $R\lim$ résulte également des résultats généraux de
Cohomologie sur les sites, Sections
\ref{sites-cohomology-section-cohomology-sheaves} et
\ref{sites-cohomology-section-unbounded}.
\end{remark}

\noindent
Les produits du lemme suivant peuvent être vus comme des produits terme à terme
de complexes ou comme des produits dans la catégorie dérivée $D(\textit{Ab})$, voir
Catégories dérivées, Lemme \ref{derived-lemma-products}.

\begin{lemma}
\label{lemma-distinguished-triangle-Rlim}
Soit $K = (K_n^\bullet)$ un objet de $D(\textit{Ab}(\mathbf{N}))$.
Il existe un triangle distingué canonique
$$
R\lim K \to \prod\nolimits_n K_n^\bullet \to \prod\nolimits_n K_n^\bullet
\to R\lim K[1]
$$
dans $D(\textit{Ab})$. Autrement dit, $R\lim K$ est une limite dérivée
du système projectif $(K_n^\bullet)$ de $D(\textit{Ab})$, voir
Catégories dérivées, Définition \ref{derived-definition-derived-limit}.
\end{lemma}

\begin{proof}
Supposons que, pour chaque $p$, le système projectif $(K_n^p)$ soit acyclique
à droite pour $\lim$. Par le Lemme \ref{lemma-compute-Rlim},
ceci donne une suite exacte courte
$$
0 \to \lim_n K^p_n \to \prod\nolimits_n K^p_n \to \prod\nolimits_n K^p_n \to 0
$$
pour chaque $p$. Puisque le complexe formé des $\lim_n K^p_n$
calcule $R\lim K$ par le Lemme \ref{lemma-compute-Rlim}, on voit que le
lemme est vrai dans ce cas.

\medskip\noindent
Supposons maintenant $K = (K_n^\bullet)$ quelconque. Par le Lemme
\ref{lemma-compute-Rlim}, il existe un quasi-isomorphisme $K \to L$ dans
$D(\textit{Ab}(\mathbf{N}))$ tel que $(L_n^p)$ soit acyclique pour chaque $p$.
Alors $\prod K_n^\bullet$ est quasi-isomorphe à $\prod L_n^\bullet$, car les produits
sont exacts dans $\textit{Ab}$ ; le résultat pour $L$, démontré ci-dessus,
implique donc celui pour $K$.
\end{proof}

\begin{lemma}
\label{lemma-break-long-exact-sequence}
Avec les notations du Lemme \ref{lemma-distinguished-triangle-Rlim},
la suite exacte longue de cohomologie associée au triangle distingué
se décompose en suites exactes courtes
$$
0 \to R^1\lim_n H^{p - 1}(K_n^\bullet) \to
H^p(R\lim K) \to
\lim_n H^p(K_n^\bullet) \to 0
$$
\end{lemma}

\begin{proof}
La suite exacte longue du triangle distingué est
$$
\ldots \to H^p(R\lim K) \to \prod\nolimits_n H^p(K_n^\bullet)
\to \prod\nolimits_n H^p(K_n^\bullet) \to
H^{p + 1}(R\lim K) \to \ldots
$$
Le morphisme du milieu a pour noyau $\lim_n H^p(K_n^\bullet)$, d'après sa
description explicite donnée dans le lemme.
Le conoyau de ce morphisme est $R^1\lim_n H^p(K_n^\bullet)$
par le Lemme \ref{lemma-compute-Rlim}.
\end{proof}

\noindent
{\bf Avertissement.} Un objet de $D(\textit{Ab}(\mathbf{N}))$ est un complexe de
systèmes projectifs de groupes abéliens. On peut aussi le considérer comme un système
projectif $(K_n^\bullet)$ de complexes. Cependant, ce n'est {\bf pas} la
même chose qu'un système projectif d'objets de $D(\textit{Ab})$ ; le lemme et la
remarque suivants expliquent la différence.

\begin{lemma}
\label{lemma-lift-to-system-complexes-Ab}
Soit $(K_n)$ un système projectif d'objets de $D(\textit{Ab})$.
Il existe alors un objet $M = (M_n^\bullet)$
de $D(\textit{Ab}(\mathbf{N}))$ et des isomorphismes
$M_n^\bullet \to K_n$ dans $D(\textit{Ab})$ tels que les diagrammes
$$
\xymatrix{
M_{n + 1}^\bullet \ar[d] \ar[r] &
M_n^\bullet \ar[d] \\
K_{n + 1} \ar[r] & K_n
}
$$
soient commutatifs dans $D(\textit{Ab})$.
\end{lemma}

\begin{proof}
En effet, soit $M_1^\bullet$ un complexe de groupes abéliens représentant $K_1$.
Supposons avoir construit
$M_e^\bullet \to M_{e - 1}^\bullet \to \ldots \to M_1^\bullet$
et des morphismes $\psi_i : M_i^\bullet \to K_i$ tels que
les diagrammes de l'énoncé du lemme soient commutatifs pour tout $n < e$.
Considérons alors le diagramme
$$
\xymatrix{
& M_n^\bullet \ar[d]^{\psi_n} \\
K_{n + 1} \ar[r] & K_n
}
$$
dans $D(\textit{Ab})$. Par la définition des morphismes de $D(\textit{Ab})$,
on peut trouver un complexe $M_{n + 1}^\bullet$ de groupes abéliens, un isomorphisme
$M_{n + 1}^\bullet \to K_{n + 1}$ dans $D(\textit{Ab})$, et un
morphisme de complexes $M_{n + 1}^\bullet \to M_n^\bullet$
représentant la composée
$$
K_{n + 1} \to K_n \xrightarrow{\psi_n^{-1}} M_n^\bullet
$$
dans $D(\textit{Ab})$.
Le lemme en résulte par récurrence.
\end{proof}

\begin{remark}
\label{remark-compare-derived-limit}
Soit $(K_n)$ un système projectif d'objets de $D(\textit{Ab})$.
Soit $K = R\lim K_n$ une limite dérivée de ce système (voir
Catégories dérivées, Section \ref{derived-section-derived-limit}). Une telle
limite dérivée existe parce que $D(\textit{Ab})$ possède des produits dénombrables
(Catégories dérivées, Lemme \ref{derived-lemma-products}).
Par le Lemme \ref{lemma-lift-to-system-complexes-Ab}, on peut aussi relever
$(K_n)$ en un objet $M$ de $D(\textit{Ab}(\mathbf{N}))$.
Alors $K \cong R\lim M$, où $R\lim$ est le foncteur (\ref{equation-Rlim}),
car $R\lim M$ est aussi une limite dérivée du système $(K_n)$
par le Lemme \ref{lemma-distinguished-triangle-Rlim}.
Ainsi, bien qu'il puisse exister de nombreuses classes d'isomorphisme de relèvements
$M$ du système $(K_n)$, le type d'isomorphisme de
$R\lim M$ est indépendant du choix, car il est isomorphe
à la limite dérivée $K = R\lim K_n$ du système. On peut donc
appliquer aux limites dérivées les résultats sur $R\lim$ démontrés dans cette section.
Par exemple, pour tout $p \in \mathbf{Z}$, il existe une
suite exacte courte canonique
$$
0 \to R^1\lim H^{p - 1}(K_n) \to H^p(K) \to \lim H^p(K_n) \to 0
$$
car on peut appliquer le Lemme \ref{lemma-break-long-exact-sequence} à $M$.
On peut aussi le voir directement, sans invoquer l'existence de $M$,
en appliquant l'argument de la preuve du
Lemme \ref{lemma-break-long-exact-sequence} au triangle distingué (définissant)
$K \to \prod K_n \to \prod K_n \to K[1]$.
\end{remark}

\begin{lemma}
\label{lemma-Rlim-pro-equal}
Soit $E \to D$ un morphisme de $D(\textit{Ab}(\mathbf{N}))$.
Soient $(E_n)$, resp.\ $(D_n)$, le système d'objets de
$D(\textit{Ab})$ associé à $E$, resp.\ à $D$.
Si $(E_n) \to (D_n)$ est un isomorphisme de pro-objets, alors
$R\lim E \to R\lim D$ est un isomorphisme dans $D(\textit{Ab})$.
\end{lemma}

\begin{proof}
L'hypothèse implique en particulier que les pro-objets
$H^p(E_n)$ et $H^p(D_n)$ sont isomorphes. Le résultat découle
des suites exactes courtes du
Lemme \ref{lemma-break-long-exact-sequence}
et du Lemme \ref{lemma-pro-isomorphism-iso-Rlim}.
\end{proof}

\begin{lemma}[Emmanouil]
\label{lemma-emmanouil}
\begin{reference}
Tiré de \cite{Emmanouil}.
\end{reference}
Soit $(A_n)$ un système projectif de groupes abéliens.
Les assertions suivantes sont équivalentes :
\begin{enumerate}
\item $(A_n)$ est de Mittag-Leffler,
\item $R^1\lim A_n = 0$, et
il en va de même de $\bigoplus_{i \in \mathbf{N}} (A_n)$.
\end{enumerate}
\end{lemma}

\begin{proof}
Posons $B = \bigoplus_{i \in \mathbf{N}} (A_n)$ ; ainsi,
$B = (B_n)$ avec $B_n = \bigoplus_{i \in \mathbf{N}} A_n$.
Si $(A_n)$ est ML, alors $B$ est ML et, par conséquent,
$R^1\lim A_n = 0$ et $R^1\lim B_n = 0$ par le
Lemme \ref{lemma-compute-Rlim}.

\medskip\noindent
Réciproquement, supposons que $(A_n)$ ne soit pas ML. On peut alors choisir un entier
$m$, une suite d'entiers $m < m_1 < m_2 < \ldots$ et des éléments
$x_i \in A_{m_i}$ dont l'image $y_i \in A_m$ n'appartient pas à
l'image de $A_{m_i + 1} \to A_m$.
Nous allons utiliser les éléments $x_i$ et $y_i$ pour
montrer de deux manières que $R^1\lim B_n \not = 0$.
Ceci achèvera la preuve du lemme.

\medskip\noindent
Première preuve.
Posons $C = (C_n)$ avec $C_n = \prod_{i \in \mathbf{N}} A_n$.
Il existe un morphisme injectif canonique $B_n \to C_n$ de conoyau
$Q_n$. Posons $Q = (Q_n)$. On peut considérer, et l'on considère, les éléments
$q_n$ de $Q_n$ comme des suites d'éléments
$q_n = (q_{n, 1}, q_{n, 2}, \ldots)$, avec $q_{n, i} \in A_n$,
modulo les suites dont la queue est nulle (autrement dit, on identifie
les suites qui diffèrent en un nombre fini de places).
On a une suite exacte courte de systèmes projectifs
$$
0 \to (B_n) \to (C_n) \to (Q_n) \to 0
$$
Considérons l'élément $q_n \in Q_n$ donné par
$$
q_{n, i} =
\left\{
\begin{matrix}
\text{image de }x_i &\text{si}& m_i \geq n \\
0 & \text{sinon}
\end{matrix}
\right.
$$
Il est alors clair que $q_{n + 1}$ s'envoie sur $q_n$.
On obtient donc $q = (q_n) \in \lim Q_n$.
D'autre part, nous affirmons que $q$ n'appartient pas à l'image
de $\lim C_n \to \lim Q_n$. En effet, supposons que $c = (c_n)$
s'envoie sur $q$. On peut écrire $c_n = (c_{n, i})$ et, puisque
$c_{n', i} \mapsto c_{n, i}$ pour $n' \geq n$, on voit que
$c_{n, i} \in \Im(C_{n'} \to C_n)$ pour tous $n, i, n' \geq n$.
En particulier, l'image de $c_{m, i}$ dans $A_m$ appartient à
$\Im(A_{m_i + 1} \to A_m)$ et ne peut donc pas être égale à $y_i$.
Ainsi, $c_m$ et $q_m = (y_1, y_2, y_3, \ldots)$
diffèrent en une infinité de places, ce qui est une contradiction.
En considérant la suite exacte longue de cohomologie
$$
0 \to \lim B_n \to \lim C_n \to \lim Q_n \to R^1\lim B_n
$$
on conclut que le dernier groupe est non nul, comme voulu.

\medskip\noindent
Seconde preuve. Pour $n' \geq n$, notons $A_{n, n'} = \Im(A_{n'} \to A_n)$.
On a alors $y_i \in A_m$, $y_i \not \in A_{m, m_i + 1}$.
Soit $\xi = (\xi_n) \in \prod B_n$ l'élément tel que
$\xi_n = 0$ sauf si $n = m_i$, et $\xi_{m_i} = (0, \ldots, 0, x_i, 0, \ldots)$,
où $x_i$ est placé dans le $i$-ième facteur. Nous affirmons que $\xi$ n'appartient pas à
l'image du morphisme $\prod B_n \to \prod B_n$ du Lemme \ref{lemma-compute-Rlim}.
Ceci montre que $R^1\lim B_n$ est non nul et achève la preuve.
En effet, supposons que $\xi$ soit l'image de $\eta = (z_1, z_2, \ldots)$,
avec $z_n = \sum z_{n, i} \in \bigoplus_i A_n$.
Remarquons que $x_i = z_{m_i, i} \bmod A_{m_i, m_i + 1}$.
Alors $z_{m_i - 1, i}$ est l'image de $z_{m_i, i}$ par
$A_{m_i} \to A_{m_i - 1}$, et ainsi de suite ; on conclut que
$z_{m, i}$ est l'image de $z_{m_i, i}$ par $A_{m_i} \to A_m$.
On conclut que $z_{m, i}$ est congru à $y_i$ modulo
$A_{m, m_i + 1}$. En particulier, $z_{m, i} \not = 0$.
Ceci est impossible car $\sum z_{m, i} \in \bigoplus_i A_m$ ;
seul un nombre fini des $z_{m, i}$ peut donc être non nul.
\end{proof}

\begin{lemma}
\label{lemma-Mittag-Leffler}
Soit
$$
0 \to (A_i) \to (B_i) \to (C_i) \to 0
$$
une suite exacte courte de systèmes projectifs de groupes abéliens.
Si $(A_i)$ et $(C_i)$ sont ML, alors $(B_i)$ l'est aussi.
\end{lemma}

\begin{proof}
Cela résulte du Lemme \ref{lemma-emmanouil}, du fait que
prendre des sommes directes infinies est exact, et de la suite exacte longue
de cohomologie associée à $R\lim$.
\end{proof}

\begin{lemma}
\label{lemma-Rlim-zero-of-direct-sums}
Soit $(A_n)$ un système projectif de groupes abéliens.
Les assertions suivantes sont équivalentes :
\begin{enumerate}
\item $(A_n)$ est nul comme pro-objet,
\item $\lim A_n = 0$ et $R^1\lim A_n = 0$, et
il en va de même de $\bigoplus_{i \in \mathbf{N}} (A_n)$.
\end{enumerate}
\end{lemma}

\begin{proof}
Il résulte du Lemme \ref{lemma-pro-isomorphism-iso-Rlim}
que (1) implique (2). Supposons (2).
Alors $(A_n)$ est ML par le Lemme \ref{lemma-emmanouil}.
Pour $m \geq n$, posons $A_{n, m} = \Im(A_m \to A_n)$,
de sorte que $A_n = A_{n, n} \supset A_{n, n + 1} \supset \ldots$.
Remarquons que $(A_n)$ est nul comme pro-objet si et seulement si, pour tout
$n$, il existe $m \geq n$ tel que $A_{n, m} = 0$.
Remarquons que $(A_n)$ est ML si et seulement si, pour tout $n$, il existe
$m_n \geq n$ tel que $A_{n, m} = A_{n, m + 1} = \ldots$. Dans le cas ML, il est
clair que $\lim A_n = 0$ implique $A_{n, m_n} = 0$,
car les morphismes $A_{n + 1, m_{n + 1}} \to A_{n, m}$ sont surjectifs.
Ceci achève la preuve.
\end{proof}
\section{Rlim des modules}
\label{section-Rlim-modules}

\noindent
Nous discutons brièvement de $R\lim$ pour les modules. Une grande partie des
arguments de cette section reprend les arguments utilisés pour construire le
formalisme cohomologique des modules sur les sites annelés.

\medskip\noindent
Soit $(A_n)$ un système projectif d'anneaux. Nous noterons
$\textit{Mod}(\mathbf{N}, (A_n))$ la catégorie des systèmes projectifs
$(M_n)$ de groupes abéliens tels que chaque $M_n$ soit muni d'une structure de
$A_n$-module et que les applications de transition
$M_{n + 1} \to M_n$ soient des homomorphismes de $A_{n + 1}$-modules.
C'est une catégorie abélienne. Posons $A = \lim A_n$.
Si $(M_n)$ est un objet de $\textit{Mod}(\mathbf{N}, (A_n))$, alors la limite
$\lim M_n$ est un $A$-module.

\begin{lemma}
\label{lemma-compute-Rlim-modules}
Dans la situation ci-dessus, le foncteur
$\lim : \textit{Mod}(\mathbf{N}, (A_n)) \to \text{Mod}_A$
admet un foncteur dérivé droit
$$
R\lim :
D(\textit{Mod}(\mathbf{N}, (A_n)))
\longrightarrow
D(A)
$$
Comme d'habitude, nous posons $R^p\lim(K) = H^p(R\lim(K))$. De plus,
\begin{enumerate}
\item pour tout $(M_n)$ dans $\textit{Mod}(\mathbf{N}, (A_n))$, on a
$R^p\lim M_n = 0$ pour $p > 1$ ;
\item l'objet $R\lim M_n$ de $D(\text{Mod}_A)$ est représenté par le complexe
$$
\prod M_n \to \prod M_n,\quad
(x_n) \mapsto (x_n - f_{n + 1}(x_{n + 1}))
$$
placé en degrés $0$ et $1$ ;
\item si $(M_n)$ vérifie la condition ML, alors $R^1\lim M_n = 0$, autrement
dit $(M_n)$ est acyclique à droite pour $\lim$ ;
\item tout $K^\bullet \in D(\textit{Mod}(\mathbf{N}, (A_n)))$ est
quasi-isomorphe à un complexe dont les termes sont acycliques à droite pour
$\lim$ ;
\item si chaque $K^p = (K^p_n)$ est acyclique à droite pour $\lim$,
c'est-à-dire si $R^1\lim_n K^p_n = 0$, alors $R\lim K$ est représenté par le
complexe dont le terme de degré $p$ est $\lim_n K_n^p$.
\end{enumerate}
\end{lemma}

\begin{proof}
La preuve est mot pour mot la même que celle du
Lemme \ref{lemma-compute-Rlim}.
\end{proof}

\begin{remark}
\label{remark-Rlim-cohomology-modules}
Cette remarque prolonge la Remarque \ref{remark-Rlim-cohomology}.
Un faisceau d'anneaux sur $\mathbf{N}$ n'est autre qu'un système projectif
d'anneaux $(A_n)$. Un faisceau de modules sur $(A_n)$ est exactement un objet
de la catégorie $\textit{Mod}(\mathbf{N}, (A_n))$ définie ci-dessus. Le
foncteur dérivé $R\lim$ du Lemme \ref{lemma-compute-Rlim-modules} est simplement
$R\Gamma(\mathbf{N}, -)$, de la catégorie dérivée des modules vers la catégorie
dérivée des modules sur les sections globales du faisceau structural. En
général, les cohomologies des groupes et des modules coïncident ; voir
Cohomologie sur les sites, Lemme
\ref{sites-cohomology-lemma-cohomology-modules-abelian-agree}.
\end{remark}

\noindent
Dans le lemme suivant, les produits peuvent être vus comme des produits terme
à terme de complexes ou comme des produits dans la catégorie dérivée $D(A)$ ;
voir Catégories dérivées, Lemme \ref{derived-lemma-products}.

\begin{lemma}
\label{lemma-distinguished-triangle-Rlim-modules}
Soit $K = (K_n^\bullet)$ un objet de
$D(\textit{Mod}(\mathbf{N}, (A_n)))$.
Il existe un triangle distingué canonique
$$
R\lim K \to \prod\nolimits_n K_n^\bullet \to \prod\nolimits_n K_n^\bullet
\to R\lim K[1]
$$
dans $D(A)$. Autrement dit, $R\lim K$ est une limite dérivée du système
projectif $(K_n^\bullet)$ dans $D(A)$ ; voir Catégories dérivées, Définition
\ref{derived-definition-derived-limit}.
\end{lemma}

\begin{proof}
La preuve est exactement la même que celle du
Lemme \ref{lemma-distinguished-triangle-Rlim}, en utilisant le
Lemme \ref{lemma-compute-Rlim-modules} à la place du
Lemme \ref{lemma-compute-Rlim}.
\end{proof}

\begin{lemma}
\label{lemma-break-long-exact-sequence-modules}
Avec les notations du Lemme \ref{lemma-distinguished-triangle-Rlim-modules}, la
suite exacte longue de cohomologie associée au triangle distingué se décompose
en suites exactes courtes
$$
0 \to R^1\lim_n H^{p - 1}(K_n^\bullet) \to
H^p(R\lim K) \to
\lim_n H^p(K_n^\bullet) \to 0
$$
de $A$-modules.
\end{lemma}

\begin{proof}
La preuve est exactement la même que celle du
Lemme \ref{lemma-break-long-exact-sequence}, en utilisant le
Lemme \ref{lemma-compute-Rlim-modules} à la place du
Lemme \ref{lemma-compute-Rlim}.
\end{proof}

\noindent
{\bf Avertissement.} Comme pour les groupes abéliens, un objet
$M = (M_n^\bullet)$ de
$D(\textit{Mod}(\mathbf{N}, (A_n)))$ est un système projectif de complexes de
modules, ce qui n'est {\bf pas} la même chose qu'un système projectif d'objets
des catégories dérivées. Dans le lemme suivant, nous montrons qu'un système
projectif d'objets de catégories dérivées se relève toujours en un objet de
$D(\textit{Mod}(\mathbf{N}, (A_n)))$.

\begin{lemma}
\label{lemma-lift-to-system-complexes}
Soit $(A_n)$ un système projectif d'anneaux. Supposons donnés
\begin{enumerate}
\item pour tout $n$, un objet $K_n$ de $D(A_n)$ ;
\item pour tout $n$, un morphisme $\varphi_n : K_{n + 1} \to K_n$ dans
$D(A_{n + 1})$, où l'on considère $K_n$ comme un objet de $D(A_{n + 1})$ par
restriction des scalaires le long de $A_{n + 1} \to A_n$.
\end{enumerate}
Il existe un objet
$M = (M_n^\bullet) \in D(\textit{Mod}(\mathbf{N}, (A_n)))$
et des isomorphismes $\psi_n : M_n^\bullet \to K_n$ dans $D(A_n)$ tels que les
diagrammes
$$
\xymatrix{
M_{n + 1}^\bullet \ar[d]_{\psi_{n + 1}} \ar[r] & M_n^\bullet \ar[d]^{\psi_n} \\
K_{n + 1} \ar[r]^{\varphi_n} & K_n
}
$$
commutent dans $D(A_{n + 1})$.
\end{lemma}

\begin{proof}
Détaillons la preuve. Pour un $A_n$-module $T$, nous notons
$T_{A_{n + 1}}$ le même module considéré comme un $A_{n + 1}$-module.
Supposons que $K_n^\bullet$ soit un complexe de $A_n$-modules représentant
$K_n$. Alors $K_{n, A_{n + 1}}^\bullet$ est le même complexe, considéré comme
un complexe de $A_{n + 1}$-modules. Par construction de la catégorie dérivée,
le morphisme $\psi_n$ peut s'écrire
$$
\psi_n = \tau_n \circ \sigma_n^{-1}
$$
où $\sigma_n : L_{n + 1}^\bullet \to K_{n + 1}^\bullet$ est un
quasi-isomorphisme de complexes de $A_{n + 1}$-modules et
$\tau_n : L_{n + 1}^\bullet \to K_{n, A_{n + 1}}^\bullet$ un morphisme de
complexes de $A_{n + 1}$-modules.

\medskip\noindent
Construisons maintenant les complexes $M_n^\bullet$ par récurrence. Pour
initialiser, posons $M_1^\bullet = K_1^\bullet$. Supposons déjà construits
$M_e^\bullet \to M_{e - 1}^\bullet \to \ldots \to M_1^\bullet$
ainsi que des morphismes de complexes $\psi_i : M_i^\bullet \to K_i^\bullet$
tels que les diagrammes
$$
\xymatrix{
M_{n + 1}^\bullet \ar[d]_{\psi_{n + 1}} \ar[rr] & &
M_{n, A_{n + 1}}^\bullet \ar[d]^{\psi_{n, A_{n + 1}}} \\
K_{n + 1}^\bullet & L_{n + 1}^\bullet \ar[l]_{\sigma_n} \ar[r]^{\tau_n} &
K_{n, A_{n + 1}}^\bullet
}
$$
ci-dessus commutent dans $D(A_{n + 1})$ pour tout $n < e$. Considérons alors
le diagramme
$$
\xymatrix{
& & M_{e, A_{e + 1}}^\bullet \ar[d]^{\psi_{e, A_{e + 1}}} \\
K_{e + 1}^\bullet & L_{e + 1}^\bullet \ar[r]^{\tau_e} \ar[l]_{\sigma_e} &
K_{e, A_{e + 1}}^\bullet
}
$$
dans $D(A_{e + 1})$. Puisque $\psi_e$ est un quasi-isomorphisme,
$\psi_{e, A_{e + 1}}$ en est également un. Par définition des morphismes dans
$D(A_{e + 1})$, on peut trouver un quasi-isomorphisme
$\psi_{e + 1} : M_{e + 1}^\bullet \to K_{e + 1}^\bullet$ de complexes de
$A_{e + 1}$-modules tel qu'il existe un morphisme de complexes
$M_{e + 1}^\bullet \to M_{e, A_{e + 1}}^\bullet$ de $A_{e + 1}$-modules
représentant la composée
$\psi_{e, A_{e + 1}}^{-1} \circ \tau_e \circ \sigma_e^{-1}$ dans
$D(A_{e + 1})$. Le lemme en résulte par récurrence.
\end{proof}

\begin{remark}
\label{remark-how-unique}
Plaçons-nous sous les hypothèses du Lemme \ref{lemma-lift-to-system-complexes}.
A priori, de nombreuses classes d'isomorphisme d'objets $M$ de
$D(\textit{Mod}(\mathbf{N}, (A_n)))$ peuvent donner le système
$(K_n, \varphi_n)$ du lemme. Pour chacun de ces $M$, on peut considérer le
complexe $R\lim M \in D(A)$, où $A = \lim A_n$. Le
Lemme \ref{lemma-distinguished-triangle-Rlim-modules} montre que $R\lim M$ est
une limite dérivée du système projectif $(K_n)$ dans $D(A)$. La classe
d'isomorphisme de $R\lim M$ dans $D(A)$ ne dépend donc pas des choix effectués
pour construire $M$. En particulier, on peut appliquer les résultats sur
$R\lim$ établis dans cette section aux limites dérivées de systèmes projectifs
dans $D(A)$. Par exemple, pour tout $p \in \mathbf{Z}$, il existe une suite
exacte courte canonique
$$
0 \to R^1\lim H^{p - 1}(K_n) \to H^p(R\lim K_n) \to \lim H^p(K_n) \to 0
$$
car on peut appliquer le Lemme \ref{lemma-distinguished-triangle-Rlim-modules}
à $M$. On peut aussi le voir directement, sans invoquer l'existence de $M$, en
appliquant l'argument de la preuve du
Lemme \ref{lemma-distinguished-triangle-Rlim-modules} au triangle distingué
(définissant)
$R\lim K_n \to \prod K_n \to \prod K_n \to (R\lim K_n)[1]$
de la limite dérivée.
\end{remark}

\begin{lemma}
\label{lemma-get-ML-system}
Soit $(A_n)$ un système projectif d'anneaux. Tout
$K \in D(\textit{Mod}(\mathbf{N}, (A_n)))$ peut être représenté par un système
de complexes $(M_n^\bullet)$ dont toutes les applications de transition
$M_{n + 1}^\bullet \to M_n^\bullet$ sont surjectives.
\end{lemma}

\begin{proof}
Représentons $K$ par le système $(K_n^\bullet)$. Posons
$M_1^\bullet = K_1^\bullet$. Supposons construits des morphismes surjectifs de
complexes $M_n^\bullet \to M_{n - 1}^\bullet \to \ldots \to M_1^\bullet$
ainsi que des équivalences d'homotopie
$\psi_e : K_e^\bullet \to M_e^\bullet$ tels que les diagrammes
$$
\xymatrix{
K_{e + 1}^\bullet \ar[d] \ar[r] & K_e^\bullet \ar[d] \\
M_{e + 1}^\bullet \ar[r] & M_e^\bullet
}
$$
commutent pour tout $e < n$. Considérons alors le diagramme
$$
\xymatrix{
K_{n + 1}^\bullet \ar[r] & K_n^\bullet \ar[d] \\
& M_n^\bullet
}
$$
D'après Catégories dérivées, Lemme \ref{derived-lemma-make-surjective}, on peut
factoriser la composée $K_{n + 1}^\bullet \to M_n^\bullet$ sous la forme
$K_{n + 1}^\bullet \to M_{n + 1}^\bullet \to M_n^\bullet$, de sorte que la
première flèche soit une équivalence d'homotopie et la seconde une surjection
scindée terme à terme. Le lemme en résulte par récurrence.
\end{proof}

\begin{lemma}
\label{lemma-get-K-flat-system}
Soit $(A_n)$ un système projectif d'anneaux. Tout
$K \in D(\textit{Mod}(\mathbf{N}, (A_n)))$ peut être représenté par un système
de complexes $(K_n^\bullet)$ tels que chacun des $K_n^\bullet$ soit K-plat.
\end{lemma}

\begin{proof}
Utilisons d'abord le Lemme \ref{lemma-get-ML-system} pour représenter $K$ par un
système de complexes $(M_n^\bullet)$ dont toutes les applications de
transition $M_{n + 1}^\bullet \to M_n^\bullet$ sont surjectives. Prenons ensuite
un quasi-isomorphisme $K_1^\bullet \to M_1^\bullet$, où $K_1^\bullet$ est un
complexe K-plat de $A_1$-modules
(Lemme \ref{lemma-K-flat-resolution}). Supposons construits
$K_n^\bullet \to K_{n - 1}^\bullet \to \ldots \to K_1^\bullet$
et des morphismes de complexes $\psi_e : K_e^\bullet \to M_e^\bullet$ tels que
$$
\xymatrix{
K_{e + 1}^\bullet \ar[d] \ar[r] & K_e^\bullet \ar[d] \\
M_{e + 1}^\bullet \ar[r] & M_e^\bullet
}
$$
commute pour tout $e < n$. Considérons alors le diagramme
$$
\xymatrix{
C^\bullet \ar@{..>}[d] \ar@{..>}[r] & K_n^\bullet \ar[d]^{\psi_n} \\
M_{n + 1}^\bullet \ar[r]^{\varphi_n} & M_n^\bullet
}
$$
dans $D(A_{n + 1})$. Comme $M_{n + 1}^\bullet \to M_n^\bullet$ est surjectif
terme à terme, le complexe $C^\bullet$ placé dans le coin supérieur gauche et
dont les termes sont
$$
C^p = M_{n + 1}^p \times_{M_n^p} K_n^p
$$
est quasi-isomorphe à $M_{n + 1}^\bullet$ (détails omis). Choisissons un
quasi-isomorphisme $K_{n + 1}^\bullet \to C^\bullet$ avec
$K_{n +1}^\bullet$ K-plat. Le lemme en résulte par récurrence.
\end{proof}

\begin{lemma}
\label{lemma-derived-tensor-product-systems}
Soit $(A_n)$ un système projectif d'anneaux. Étant donnés
$K, L \in D(\textit{Mod}(\mathbf{N}, (A_n)))$, il existe un produit tensoriel
dérivé canonique $K \otimes^\mathbf{L} L$ dans $D(\mathbf{N}, (A_n))$,
compatible avec les foncteurs vers $D(A_n)$. La construction est symétrique en
$K$ et $L$ et définit un foncteur exact de catégories triangulées en chaque
variable.
\end{lemma}

\begin{proof}
Choisissons un représentant $(K_n^\bullet)$ de $K$ tel que chaque
$K_n^\bullet$ soit un complexe K-plat
(Lemme \ref{lemma-get-K-flat-system}). On peut alors définir
$K \otimes^\mathbf{L} L$ comme l'objet représenté par le système de complexes
$$
(\text{Tot}(K_n^\bullet \otimes_{A_n} L_n^\bullet))
$$
pour n'importe quel choix de représentant $(L_n^\bullet)$ de $L$. Cette
définition est indépendante des choix dans les deux variables d'après les
Lemmes \ref{lemma-K-flat-quasi-isomorphism} et
\ref{lemma-derived-tor-quasi-isomorphism-other-side}. La compatibilité avec le
foncteur vers $D(A_n)$ est claire. L'exactitude se démontre exactement comme
dans le Lemme \ref{lemma-derived-tor-exact}.
\end{proof}

\begin{remark}
\label{remark-constructing-tensor-with-limits-functorially}
Soit $A$ un anneau et soit $(E_n)$ un système projectif d'objets de $D(A)$.
Nous avons vu ci-dessus qu'une limite dérivée $R\lim E_n$ existe. Ainsi, pour
tout objet $K$ de $D(A)$, la limite dérivée
$R\lim( K \otimes_A^\mathbf{L} E_n )$ existe également. On peut en fait
construire ces limites dérivées de manière fonctorielle en $K$ et obtenir un
foncteur exact
$$
R\lim(- \otimes_A^\mathbf{L} E_n) : D(A) \longrightarrow D(A)
$$
de catégories triangulées. En effet, relevons d'abord $(E_n)$ en un objet $E$
de $D(\mathbf{N}, A)$ ; voir Lemme \ref{lemma-lift-to-system-complexes}.
(Le foncteur dépendra du choix de ce relèvement.) Observons ensuite qu'il
existe un foncteur « diagonal » ou « constant »
$$
\Delta : D(A) \longrightarrow D(\mathbf{N}, A)
$$
qui envoie le complexe $K^\bullet$ sur le système projectif constant de
complexes de valeur $K^\bullet$. On définit alors simplement
$$
R\lim(K \otimes_A^\mathbf{L} E_n) = R\lim(\Delta(K)\otimes^\mathbf{L} E)
$$
où, au membre de droite, on utilise le foncteur $R\lim$ du
Lemme \ref{lemma-compute-Rlim-modules} et le foncteur
$- \otimes^\mathbf{L} -$ du
Lemme \ref{lemma-derived-tensor-product-systems}.
\end{remark}

\begin{lemma}
\label{lemma-tensor-Rlim-exact}
Soit $A$ un anneau. Soit $E \to D \to F \to E[1]$ un triangle distingué de
$D(\mathbf{N}, A)$. Soient $(E_n)$, resp.\ $(D_n)$, resp.\ $(F_n)$ les systèmes
d'objets de $D(A)$ associés à $E$, resp.\ $D$, resp.\ $F$. Alors, pour tout
$K \in D(A)$, il existe un triangle distingué canonique
$$
R\lim (K \otimes^\mathbf{L}_A E_n) \to
R\lim (K \otimes^\mathbf{L}_A D_n) \to
R\lim (K \otimes^\mathbf{L}_A F_n) \to
R\lim (K \otimes^\mathbf{L}_A E_n)[1]
$$
dans $D(A)$, avec les notations de la
Remarque \ref{remark-constructing-tensor-with-limits-functorially}.
\end{lemma}

\begin{proof}
Cela résulte immédiatement de la construction de la
Remarque \ref{remark-constructing-tensor-with-limits-functorially} et du fait
que $\Delta : D(A) \to D(\mathbf{N}, A)$, $- \otimes^\mathbf{L} -$ et $R\lim$
sont des foncteurs exacts de catégories triangulées.
\end{proof}

\begin{lemma}
\label{lemma-tensor-Rlim-pro-equal}
Soit $A$ un anneau. Soit $E \to D$ un morphisme de $D(\mathbf{N}, A)$. Soient
$(E_n)$, resp.\ $(D_n)$ les systèmes d'objets de $D(A)$ associés à $E$,
resp.\ $D$. Si $(E_n) \to (D_n)$ est un isomorphisme de pro-objets, alors, pour
tout $K \in D(A)$, le morphisme correspondant
$$
R\lim (K \otimes^\mathbf{L}_A E_n)
\longrightarrow
R\lim (K \otimes^\mathbf{L}_A D_n)
$$
dans $D(A)$ est un isomorphisme (notations de la
Remarque \ref{remark-constructing-tensor-with-limits-functorially}).
\end{lemma}

\begin{proof}
Cela résulte des définitions et du Lemme \ref{lemma-Rlim-pro-equal}.
\end{proof}






\section{Modules de torsion}
\label{section-torsion}

\noindent
Dans cette section, l'expression « modules de torsion » désigne les modules à
support dans un sous-ensemble fermé donné $V(I)$ d'un schéma affine
$\Spec(R)$. Cette notion est différente de celle d'un module de torsion sur un
anneau intègre, mais lui est analogue (Définition \ref{definition-torsion}).

\begin{definition}
\label{definition-f-power-torsion}
Soit $R$ un anneau et soit $M$ un $R$-module.
\begin{enumerate}
\item Soit $I \subset R$ un idéal. On dit que $M$ est un
{\it module de torsion annulée par une puissance de $I$} si, pour tout
$m \in M$, il existe $n > 0$ tel que $I^n m = 0$.
\item Soit $f \in R$. On dit que $M$ est
{\it un module de torsion annulée par une puissance de $f$} si, pour tout
$m \in M$, il existe $n > 0$ tel que $f^n m = 0$.
\end{enumerate}
\end{definition}

\noindent
Ainsi, un module de torsion annulée par une puissance de $f$ est la même chose
qu'un module de torsion annulée par une puissance de $I$ pour $I = (f)$. Nous
utiliserons les notations
$$
M[I^n] = \{m \in M \mid I^nm = 0\}
$$
et
$$
M[I^\infty] = \bigcup M[I^n]
$$
pour un $R$-module $M$. Ainsi, $M$ est de torsion annulée par une puissance de
$I$ si et seulement si
$M = M[I^\infty]$, ou encore si et seulement si $M = \bigcup M[I^n]$.

\begin{lemma}
\label{lemma-I-power-torsion-presentation}
Soient $R$ un anneau, $I$ un idéal de $R$ et $M$ un module de torsion annulée
par une puissance de $I$. Alors $M$ admet une résolution
$$
\ldots \to K_2 \to K_1 \to K_0 \to M \to 0
$$
où chaque $K_i$ est une somme directe de copies de $R/I^n$, avec $n$ variable.
\end{lemma}

\begin{proof}
Il existe une surjection canonique
$$
\oplus_{m \in M} R/I^{n_m} \to M \to 0
$$
où $n_m$ est le plus petit entier strictement positif tel que
$I^{n_m} \cdot m = 0$. Le noyau de cette surjection est lui aussi un module de
torsion annulée par une puissance de $I$. En poursuivant par récurrence, on
construit la résolution voulue de $M$.
\end{proof}

\begin{lemma}
\label{lemma-torsion-free}
Soient $R$ un anneau et $I$ un idéal de $R$. Pour tout $R$-module $M$, posons
$M[I^n] = \{m \in M \mid I^nm = 0\}$. Si $I$ est de type fini, les conditions
suivantes sont équivalentes :
\begin{enumerate}
\item $M[I] = 0$ ;
\item $M[I^n] = 0$ pour tout $n \geq 1$ ;
\item si $I = (f_1, \ldots, f_t)$, alors l'application
$M \to \bigoplus M_{f_i}$ est injective.
\end{enumerate}
\end{lemma}

\begin{proof}
Cela résulte de Algèbre, Lemme
\ref{algebra-lemma-when-injective-covering}.
\end{proof}

\begin{lemma}
\label{lemma-divide-by-torsion}
Soient $R$ un anneau et $I$ un idéal de type fini de $R$.
\begin{enumerate}
\item Pour tout $R$-module $M$, on a $(M/M[I^\infty])[I] = 0$.
\item Toute extension de modules de torsion annulée par une puissance de $I$
est encore de torsion annulée par une puissance de $I$.
\end{enumerate}
\end{lemma}

\begin{proof}
Soit $m \in M$. Si l'image de $m$ appartient à $(M/M[I^\infty])[I]$, alors
$Im \subset M[I^\infty]$. Écrivons $I = (f_1, \ldots, f_t)$. On a donc
$f_i m \in M[I^\infty]$, c'est-à-dire $I^{n_i}f_i m = 0$ pour un certain
$n_i > 0$. Par conséquent, $I^Nm = 0$ pour $N = \sum n_i + 2$. L'image de $m$
dans $(M/M[I^\infty])$ est donc nulle, ce qui démontre la première assertion du
lemme.

\medskip\noindent
Pour la seconde assertion, supposons que
$0 \to M' \to M \to M'' \to 0$ soit une suite exacte courte de modules dans
laquelle $M'$ et $M''$ sont tous deux de torsion annulée par une puissance de
$I$. Alors $M[I^\infty] \supset M'$ ; ainsi, $M/M[I^\infty]$ est un quotient de
$M''$ et est donc de torsion annulée par une puissance de $I$. La première
assertion et le Lemme \ref{lemma-torsion-free} impliquent qu'il est nul.
\end{proof}

\begin{lemma}
\label{lemma-I-power-torsion}
Soit $I$ un idéal de type fini d'un anneau $R$. Les modules de torsion annulée
par une puissance de $I$ forment une sous-catégorie de Serre de la catégorie
abélienne $\text{Mod}_R$ ; voir Homologie, Définition
\ref{homology-definition-serre-subcategory}.
\end{lemma}

\begin{proof}
Il est clair qu'un sous-module ou un module quotient d'un module de torsion
annulée par une puissance de $I$ est encore de torsion annulée par une
puissance de $I$. De plus, le Lemme \ref{lemma-divide-by-torsion} montre qu'une
extension de deux modules de torsion annulée par une puissance de $I$ est
encore de torsion annulée par une puissance de $I$. L'assertion résulte donc
de Homologie, Lemme
\ref{homology-lemma-characterize-serre-subcategory}.
\end{proof}

\begin{lemma}
\label{lemma-local-cohomology-closed}
Soient $R$ un anneau et $I \subset R$ un idéal de type fini. La sous-catégorie
$I^\infty\text{-torsion} \subset \text{Mod}_R$ ne dépend que du sous-ensemble
fermé $Z = V(I) \subset \Spec(R)$. Plus précisément, un $R$-module $M$ est de
torsion annulée par une puissance de $I$ si et seulement si son support est
contenu dans $Z$.
\end{lemma}

\begin{proof}
Soient $M$ un $R$-module et $x \in M$. Si $x \in M[I^\infty]$, alors l'image
de $x$ dans $M_f$ est nulle pour tout $f \in I$. L'image de $x$ dans
$M_\mathfrak p$ est donc nulle pour tout $\mathfrak p \not \supset I$.
Réciproquement, si l'image de $x$ dans $M_\mathfrak p$ est nulle pour tout
$\mathfrak p \not \supset I$, alors l'image de $x$ dans $M_f$ est nulle pour tout
$f \in I$. Par suite, si $I = (f_1, \ldots, f_r)$, alors
$f_i^{n_i}x = 0$ pour un certain $n_i \geq 1$. Il s'ensuit que
$x \in M[I^{\sum n_i}]$. Ainsi, $M[I^\infty]$ est le noyau de
$M \to \prod_{\mathfrak p \not \in Z} M_\mathfrak p$. La seconde assertion du
lemme en résulte, et elle implique la première.
\end{proof}

\noindent
Les deux lemmes suivants devraient probablement figurer ailleurs.

\begin{lemma}
\label{lemma-bounded-above-mod-I}
Soient $R$ un anneau, $I \subset R$ un idéal et $K$ un objet de $D(R)$ tel que
$H^i(K \otimes_R^\mathbf{L} R/I) = 0$ pour $i > 0$. Alors :
\begin{enumerate}
\item $H^i(K \otimes_R^\mathbf{L} R/I^n) = 0$ pour tous $n \geq 1$ et
$i > 0$ ;
\item $H^i(K \otimes_R^\mathbf{L} N) = 0$ pour tout module de torsion annulée
par une puissance de $I$, qui soit un $R$-module $N$, et tout $i > 0$ ;
\item pour tout $M \in D^b(R)$ dont les modules de cohomologie $H^i(M)$ sont de
torsion annulée par une puissance de $I$ et égaux à $0$ pour $i > 0$, on a
$H^i(K \otimes_R^\mathbf{L} M) = 0$ pour $i > 0$.
\end{enumerate}
\end{lemma}

\begin{proof}
Démontrons (2). On peut écrire $N = \bigcup N[I^n]$. On a
$K \otimes_R^\mathbf{L} N = \text{hocolim}_n K \otimes_R^\mathbf{L} N[I^n]$,
car les produits tensoriels commutent aux colimites (détails omis ; indication :
représenter $K$ par un complexe K-plat et calculer directement). On peut donc
supposer que $N$ est annulé par $I^n$. Considérons la $R$-algèbre
$R' = R/I^n \oplus N$, où $N$ est un idéal de carré nul. Il suffit de montrer
que l'objet $K' = K \otimes_R^\mathbf{L} R'$ de $D(R')$ a une cohomologie nulle
en degrés strictement positifs. On dispose d'une surjection $R' \to R/I$ de
$R$-algèbres dont le noyau $J$ est nilpotent (tout produit de $n$ éléments du
noyau est nul). On a
$$
K \otimes_R^\mathbf{L} R/I =
(K \otimes_R^\mathbf{L} R') \otimes_{R'}^\mathbf{L} R/I =
K' \otimes_{R'}^\mathbf{L} R/I
$$
Par hypothèse, le complexe $K \otimes_R^\mathbf{L} R/I$ est d'amplitude de Tor
contenue dans $[-\infty, 0]$. La conclusion résulte donc du
Lemme \ref{lemma-check-tor-dimension-modulo-nilpotent}.

\medskip\noindent
La partie (1) résulte immédiatement de la partie (2). La partie (3) résulte de
la partie (2), d'une récurrence sur le nombre de modules de cohomologie non nuls
de $M$ et des triangles distingués de troncation de Catégories dérivées,
Remarque \ref{derived-remark-truncation-distinguished-triangle}. Les détails
sont omis.
\end{proof}

\begin{lemma}
\label{lemma-derived-vanishing-mod-I}
Soient $R$ un anneau, $I \subset R$ un idéal et $K$ un objet de $D(R)$ tel que
$K \otimes_R^\mathbf{L} R/I = 0$ dans $D(R)$. Alors :
\begin{enumerate}
\item $K \otimes_R^\mathbf{L} R/I^n = 0$ pour tout $n \geq 1$ ;
\item $K \otimes_R^\mathbf{L} N = 0$ pour tout module de torsion annulée par
une puissance de $I$, qui soit un $R$-module $N$ ;
\item $K \otimes_R^\mathbf{L} M = 0$ pour tout $M \in D^b(R)$ dont les modules
de cohomologie sont de torsion annulée par une puissance de $I$.
\end{enumerate}
\end{lemma}

\begin{proof}
C'est une conséquence du Lemme \ref{lemma-bounded-above-mod-I} (certains
détails sont omis).
\end{proof}

\begin{lemma}
\label{lemma-torsion-completion}
\begin{reference}
Légère généralisation de \cite[Lemme~1]{Beauville-Laszlo}.
\end{reference}
Soit $R \to R'$ un homomorphisme d'anneaux. Soit $I \subset R$ un idéal tel
que $R/I^n \to R'/I^nR'$ soit un isomorphisme pour $n > 0$. Pour tout module
de torsion annulée par une puissance de $I$, qui soit un $R$-module $M$,
l'application
$M \to M \otimes_R R'$ est un isomorphisme. Par exemple, si $I$ est de type
fini et si $R^\wedge$ est le complété de $R$ pour la topologie $I$-adique,
alors $M \cong M \otimes_R R^\wedge$.
\end{lemma}

\begin{proof}
Si $M$ est annulé par $I^n$, alors
$$
M \otimes_R R' \cong
M \otimes_{R/I^n} R'/I^n R' \cong
M \otimes_{R/I^n} R/I^n \cong M.
$$
Si $M$ est de torsion annulée par une puissance de $I$, alors
$M = \bigcup M[I^n]$. Comme les produits tensoriels commutent aux limites
inductives (Algèbre, Lemme
\ref{algebra-lemma-tensor-products-commute-with-limits}), on obtient
l'isomorphisme voulu. La dernière assertion est un cas particulier de la
première d'après Algèbre, Lemme
\ref{algebra-lemma-hathat-finitely-generated}.
\end{proof}






\section{Recollement formel de catégories de modules}
\label{section-formal-glueing}

\noindent
Fixons un schéma noethérien $X$ et un sous-schéma fermé $Z$ de complémentaire
$U$. Notre but est d'expliquer comment construire, de manière unique, les
faisceaux cohérents sur $X$ à partir de faisceaux cohérents sur le complété
formel de $X$ le long de $Z$ et de faisceaux cohérents sur $U$, munis d'une
compatibilité convenable sur l'intersection. Nous le faisons d'abord en
n'utilisant que l'algèbre commutative (dans cette section), puis nous
l'expliquerons dans le cadre des espaces algébriques (Sommes amalgamées
d'espaces, Section \ref{spaces-pushouts-section-formal-glueing}).

\medskip\noindent
Voici quelques références qui traitent d'une partie des résultats de cette
section : \cite[Section 2]{ArtinII}, \cite[Appendix]{Ferrand-Raynaud},
\cite{Beauville-Laszlo}, \cite{MB} et \cite[Section 4.6]{dJ-crystalline}.

\begin{lemma}
\label{lemma-characterize-flatness-on-torsion}
Soient $\varphi : R \to S$ un homomorphisme d'anneaux et $I \subset R$ un
idéal. Les conditions suivantes sont équivalentes :
\begin{enumerate}
\item $\varphi$ est plat et $R/I \to S/IS$ est fidèlement plat ;
\item $\varphi$ est plat et l'application
$\Spec(S/IS) \to \Spec(R/I)$ est surjective ;
\item $\varphi$ est plat et le foncteur de changement de base
$M \mapsto M \otimes_R S$ est fidèle sur les modules annulés par $I$ ;
\item $\varphi$ est plat et le foncteur de changement de base
$M \mapsto M \otimes_R S$ est fidèle sur les modules de torsion annulée par une
puissance de $I$.
\end{enumerate}
\end{lemma}

\begin{proof}
Si $R \to S$ est plat, alors $R/I^n \to S/I^nS$ est plat pour tout $n$ ; voir
Algèbre, Lemme \ref{algebra-lemma-flat-base-change}. Les conditions (1) et (2)
sont donc équivalentes d'après Algèbre, Lemme
\ref{algebra-lemma-ff-rings}. L'équivalence de (1) et (3) résulte de
l'identification des modules de torsion annulée par $I$ parmi les $R$-modules
aux $R/I$-modules, de
l'égalité
$$
M \otimes_R S = M \otimes_{R/I} S/IS
$$
pour tout $R$-module $M$ annulé par $I$, et de Algèbre, Lemme
\ref{algebra-lemma-easy-ff}. L'implication (4) $\Rightarrow$ (3) est immédiate.
Supposons (3). Nous avons vu que $R/I^n \to S/I^nS$ est plat ; par hypothèse,
il induit en outre une surjection sur les spectres, puisque
$\Spec(R/I^n) = \Spec(R/I)$, et de même pour $S$. Le foncteur de changement de
base est donc fidèle sur les modules annulés par $I^n$. Comme tout module de
torsion annulée par une puissance de $I$, noté $M$, est la réunion
$M = \bigcup M_n$, où $M_n$ est annulé par $I^n$, le foncteur de changement de
base est fidèle sur la catégorie de tous les modules de torsion par une
puissance de $I$ (le produit tensoriel
commutant aux colimites).
\end{proof}

\begin{lemma}
\label{lemma-neighbourhood-isomorphism}
Supposons que $(\varphi : R \to S, I)$ satisfasse aux conditions équivalentes du
Lemme \ref{lemma-characterize-flatness-on-torsion}. Les conditions suivantes
sont équivalentes :
\begin{enumerate}
\item pour tout module de torsion annulée par une puissance de $I$, noté $M$,
l'application naturelle $M \to M \otimes_R S$ est un isomorphisme ;
\item $R/I \to S/IS$ est un isomorphisme.
\end{enumerate}
\end{lemma}

\begin{proof}
L'implication (1) $\Rightarrow$ (2) est immédiate. Supposons (2). Supposons
d'abord $M$ annulé par $I$. Dans ce cas, $M$ est un $R/I$-module, et l'on a un
isomorphisme
$$
M \otimes_R S = M \otimes_{R/I} S/IS = M \otimes_{R/I} R/I = M
$$
ce qui démontre l'assertion. Montrons ensuite par récurrence que
$M \to M \otimes_R S$ est un isomorphisme pour tout module $M$ annulé par
$I^n$. Supposons le résultat établi pour $n$ et $M$ annulé par $I^{n + 1}$.
On dispose alors d'une suite exacte courte
$$
0 \to I^nM \to M \to M/I^nM \to 0
$$
Comme $R \to S$ est plat, elle donne une suite exacte courte
$$
0 \to I^nM \otimes_R S \to M \otimes_R S \to M/I^nM \otimes_R S \to 0
$$
L'application canonique étant un isomorphisme pour $M' = I^nM$ et
$M'' = M/I^nM$ par hypothèse de récurrence, elle l'est aussi pour $M$. Enfin,
si $M$ est un module quelconque de torsion annulée par une puissance de $I$,
alors $M = \bigcup M_n$, où $M_n$ est annulé par $I^n$, et on conclut en
utilisant le fait que les produits tensoriels commutent aux colimites.
\end{proof}

\begin{lemma}
\label{lemma-neighbourhood-equivalence}
Supposons que $\varphi : R \to S$ soit un homomorphisme plat, que
$I \subset R$ soit un idéal de type fini et que $R/I \to S/IS$ soit un
isomorphisme. Alors :
\begin{enumerate}
\item pour tout $R$-module $M$, l'application $M \to M \otimes_R S$ induit un
isomorphisme
$M[I^\infty] \to (M \otimes_R S)[(IS)^\infty]$ entre les sous-modules de
torsion annulée par une puissance de $I$ ;
\item l'application naturelle
$$
\Hom_R(M, N) \longrightarrow \Hom_S(M \otimes_R S, N \otimes_R S)
$$
est un isomorphisme si $M$ ou $N$ est de torsion annulée par une puissance de
$I$ ;
\item le foncteur de changement de base $M \mapsto M \otimes_R S$ définit une
équivalence de catégories entre les modules de torsion annulée par une
puissance de $I$ et ceux annulés par une puissance de $IS$.
\end{enumerate}
\end{lemma}

\begin{proof}
Remarquons que les conditions équivalentes du
Lemme \ref{lemma-characterize-flatness-on-torsion} et du
Lemme \ref{lemma-neighbourhood-isomorphism} sont satisfaites. Nous les
utiliserons sans autre mention. Démontrons d'abord (1). Soit $M$ un
$R$-module quelconque. Posons $M' = M/M[I^\infty]$ et considérons la suite
exacte
$$
0 \to M[I^\infty] \to M \to M' \to 0
$$
Comme $M[I^\infty] = M[I^\infty] \otimes_R S$, il suffit de montrer que
$(M' \otimes_R S)[(IS)^\infty] = 0$. Écrivons
$I = (f_1, \ldots, f_t)$. Le Lemme \ref{lemma-divide-by-torsion} donne
$M'[I^\infty] = 0$. Par conséquent, pour tout $n > 0$, l'application
$$
M' \longrightarrow \bigoplus\nolimits_{i = 1, \ldots t} M',
\quad
x \longmapsto (f_1^n x, \ldots, f_t^n x)
$$
est injective. Puisque $S$ est plat sur $R$, l'application correspondante
$M' \otimes_R S \to \bigoplus_{i = 1, \ldots t} M' \otimes_R S$ est elle aussi
injective. Cela signifie que $(M' \otimes_R S)[I^n] = 0$, comme voulu.

\medskip\noindent
Démontrons ensuite (2). Si $N$ est de torsion annulée par une puissance de
$I$, alors $N \otimes_R S = N$, et l'application affichée en (2) est un
isomorphisme d'après Algèbre, Lemme
\ref{algebra-lemma-adjoint-tensor-restrict}. Si $M$ est de torsion annulée par
une puissance de $I$, alors l'image de tout morphisme $M \to N$ se factorise
par $N[I^\infty]$, et l'image de tout morphisme
$M \otimes_R S \to N \otimes_R S$ se factorise par
$(N \otimes_R S)[(IS)^\infty]$. Dans ce cas, la partie (1) permet donc de
remplacer $N$ par $N[I^\infty]$, et le résultat découle du cas où $N$ est de
torsion annulée par une puissance de $I$, déjà traité.

\medskip\noindent
Démontrons enfin (3). Le foncteur est pleinement fidèle par (2). Pour
l'essentielle surjectivité, il suffit de remarquer que, pour tout module de
torsion annulée par une puissance de $IS$, qui soit un $S$-module $N$,
l'application naturelle
$N \otimes_R S \to N$ est un isomorphisme.
\end{proof}

\begin{lemma}
\label{lemma-map-identifies-koszul-and-cech-complexes}
Supposons que $\varphi : R \to S$ soit un homomorphisme plat, que
$I \subset R$ soit un idéal de type fini et que $R/I \to S/IS$ soit un
isomorphisme. Pour tous $f_1, \ldots, f_r \in R$ tels que
$V(f_1, \ldots, f_r) = V(I)$ :
\begin{enumerate}
\item le morphisme de complexes de Koszul
$K(R, f_1, \ldots, f_r) \to K(S, f_1, \ldots, f_r)$ est un
quasi-isomorphisme ;
\item le morphisme de complexes de {\v C}ech alternés augmentés
$$
\xymatrix{
R \to \prod_{i_0} R_{f_{i_0}} \to \prod_{i_0 < i_1} R_{f_{i_0}f_{i_1}}
\to \ldots \to R_{f_1\ldots f_r} \ar[d] \\
S \to \prod_{i_0} S_{f_{i_0}} \to \prod_{i_0 < i_1} S_{f_{i_0}f_{i_1}}
\to \ldots \to S_{f_1\ldots f_r}
}
$$
est un quasi-isomorphisme.
\end{enumerate}
\end{lemma}

\begin{proof}
Dans les deux cas, on dispose d'un complexe $K_\bullet$ de $R$-modules, et il
s'agit de montrer que $K_\bullet \to K_\bullet \otimes_R S$ est un
quasi-isomorphisme. D'après le Lemme
\ref{lemma-neighbourhood-isomorphism} et la platitude de $R \to S$, il suffit
que tous les groupes d'homologie de $K$ soient de torsion annulée par une
puissance de $I$. C'est le cas du complexe de Koszul d'après le
Lemme \ref{lemma-homotopy-koszul}, et du complexe de {\v C}ech alterné augmenté
d'après le Lemme \ref{lemma-extended-alternating-torsion}.
\end{proof}

\begin{lemma}
\label{lemma-naive-Koszul-complex}
Soient $R$ un anneau et $I = (f_1, \ldots, f_n)$ un idéal de type fini de $R$.
Soit $M$ le $R$-module engendré par des éléments $e_1, \ldots, e_n$ soumis aux
relations $f_i e_j - f_j e_i = 0$. Il existe une suite exacte courte
$$
0 \to K \to M \to I \to 0
$$
telle que $K$ soit annulé par $I$.
\end{lemma}

\begin{proof}
Il s'agit simplement d'une troncation du complexe de Koszul. L'application
$M \to I$ est définie par la règle $e_i \mapsto f_i$. Si
$m = \sum a_i e_i$ appartient au noyau de $M \to I$, c'est-à-dire si
$\sum a_i f_i = 0$, alors
$f_j m = \sum f_j a_i e_i = (\sum f_i a_i) e_j = 0$.
\end{proof}

\begin{lemma}
\label{lemma-explicit-ext}
Soient $R$ un anneau et $I = (f_1, \ldots, f_n)$ un idéal de type fini de $R$.
Pour tout $R$-module $N$, posons
$$
H_1(N, f_\bullet) =
\frac{\{(x_1, \ldots, x_n) \in N^{\oplus n} \mid f_i x_j = f_j x_i \}}
{\{f_1x, \ldots, f_nx) \mid x \in N\}}
$$
Il existe, pour tout $R$-module $N$, une suite exacte courte canonique
$$
0 \to \Ext_R(R/I, N) \to H_1(N, f_\bullet) \to \Hom_R(K, N)
$$
où $K$ est celui du Lemme \ref{lemma-naive-Koszul-complex}.
\end{lemma}

\begin{proof}
La notation ci-dessus désigne les groupes $\Ext$ dans $\text{Mod}_R$ tels
qu'ils sont définis dans Homologie, Section
\ref{homology-section-extensions}. Ceux-ci sont notés $\Ext_R(M, N)$. La suite
exacte longue de Homologie, Lemme
\ref{homology-lemma-six-term-sequence-ext}, associée à la suite exacte courte
$0 \to I \to R \to R/I \to 0$, ainsi que l'égalité $\Ext_R(R, N) = 0$ donnent
$$
\Ext_R(R/I, N) =
\Coker(N \longrightarrow \Hom(I, N))
$$
En utilisant la suite exacte courte du
Lemme \ref{lemma-naive-Koszul-complex}, on obtient un complexe
$$
N \to \Hom(M, N) \to \Hom_R(K, N)
$$
dont l'homologie au milieu est canoniquement isomorphe à
$\Ext_R(R/I, N)$. La preuve est achevée puisque le conoyau de la première
application est canoniquement isomorphe à $H_1(N, f_\bullet)$.
\end{proof}

\begin{lemma}
\label{lemma-koszul-homology-annihilated}
Soient $R$ un anneau et $I = (f_1, \ldots, f_n)$ un idéal de type fini de $R$.
Pour tout $R$-module $N$, le groupe d'homologie de Koszul
$H_1(N, f_\bullet)$ défini dans le Lemme \ref{lemma-explicit-ext} est annulé
par $I$.
\end{lemma}

\begin{proof}
Soit $(x_1, \ldots, x_n) \in N^{\oplus n}$ tel que
$f_i x_j = f_j x_i$. Alors
$f_i(x_1, \ldots, x_n) = (f_i x_i, \ldots, f_i x_n)$. Autrement dit, $f_i$
annule $H_1(N, f_\bullet)$.
\end{proof}

\noindent
On peut renforcer la pleine fidélité du
Lemme \ref{lemma-neighbourhood-equivalence} en montrant que les groupes $\Ext$
dont la source est de torsion annulée par une puissance de $I$ sont eux aussi
insensibles au passage à $S$. Pour une version dérivée du lemme suivant, voir
Complexes dualisants, Lemme
\ref{dualizing-lemma-neighbourhood-extensions}.

\begin{lemma}
\label{lemma-neighbourhood-extensions}
Supposons que $\varphi : R \to S$ soit un homomorphisme plat, que
$I \subset R$ soit un idéal de type fini et que $R/I \to S/IS$ soit un
isomorphisme. Soient $M$ et $N$ des $R$-modules, et supposons $M$ de torsion
annulée par une puissance de $I$. Étant donnée une suite exacte courte
$$
0 \to N \otimes_R S \to \tilde E \to M \otimes_R S \to 0
$$
il existe un diagramme commutatif
$$
\xymatrix{
0 \ar[r] &
N \ar[r] \ar[d] &
E \ar[r] \ar[d] &
M \ar[r] \ar[d] &
0 \\
0 \ar[r] &
N \otimes_R S \ar[r] &
\tilde E \ar[r] &
M \otimes_R S \ar[r] &
0
}
$$
dont les lignes sont exactes.
\end{lemma}

\begin{proof}
Comme $M$ est de torsion annulée par une puissance de $I$, on a
$M \otimes_R S = M$ d'après le
Lemme \ref{lemma-neighbourhood-isomorphism}. Nous utiliserons cette
identification sans autre mention. Puisque $R \to S$ est plat, le foncteur de
changement de base est exact et fournit un morphisme fonctoriel de groupes
$\Ext$
$$
\Ext_R(M, N)
\longrightarrow
\Ext_S(M \otimes_R S, N \otimes_R S),
$$
voir Homologie, Lemme \ref{homology-lemma-exact-functor-ext}. L'assertion du
lemme est que ce morphisme est surjectif lorsque $M$ est de torsion annulée par
une puissance de $I$. Nous allons en fait montrer qu'il s'agit d'un
isomorphisme. D'après le Lemme \ref{lemma-I-power-torsion-presentation}, il
existe une surjection $M' \to M$ où $M'$ est une somme directe de modules de la
forme $R/I^n$. En utilisant la suite exacte longue de Homologie, Lemme
\ref{homology-lemma-six-term-sequence-ext}, ainsi que le
Lemme \ref{lemma-neighbourhood-equivalence}, on voit qu'il suffit d'établir le
lemme pour $M'$. La compatibilité de $\Ext$ aux sommes directes (détails omis)
ramène au cas $M = R/I^n$ pour un certain $n$.

\medskip\noindent
Soient $f_1, \ldots, f_t$ des générateurs de $I^n$. Le
Lemme \ref{lemma-explicit-ext} donne un diagramme commutatif
$$
\xymatrix{
0 \ar[r] &
\Ext_R(R/I^n, N) \ar[r] \ar[d] &
H_1(N, f_\bullet) \ar[r] \ar[d] &
\Hom_R(K, N) \ar[d] \\
0 \ar[r] &
\Ext_S(S/I^nS, N \otimes S) \ar[r] &
H_1(N \otimes S, f_\bullet) \ar[r] &
\Hom_S(K \otimes S, N \otimes S)
}
$$
à lignes exactes, où $K$ est celui du
Lemme \ref{lemma-naive-Koszul-complex}. Il suffit donc de montrer que les deux
flèches verticales de droite sont des isomorphismes. Puisque $K$ est annulé par
$I^n$, on a
$\Hom_R(K, N) = \Hom_S(K \otimes_R S, N \otimes_R S)$ d'après le
Lemme \ref{lemma-neighbourhood-equivalence}. La platitude de $R \to S$ donne
$H_1(N, f_\bullet) \otimes_R S = H_1(N \otimes_R S, f_\bullet)$. Comme
$H_1(N, f_\bullet)$ est annulé par $I^n$ d'après le
Lemme \ref{lemma-koszul-homology-annihilated}, on a
$H_1(N, f_\bullet) \otimes_R S = H_1(N, f_\bullet)$ d'après le
Lemme \ref{lemma-neighbourhood-isomorphism}.
\end{proof}

\noindent
Soit $R \to S$ un homomorphisme d'anneaux. Soient
$f_1, \ldots, f_t \in R$ et $I = (f_1, \ldots, f_t)$. Pour tout $R$-module
$M$, on peut définir un complexe
\begin{equation}
\label{equation-glueing-complex}
0 \to M \xrightarrow{\alpha}
M \otimes_R S \times \prod M_{f_i} \xrightarrow{\beta}
\prod (M \otimes_R S)_{f_i}
\times
\prod M_{f_if_j}
\end{equation}
où $\alpha(m) = (m \otimes 1, m/1, \ldots, m/1)$ et
$$
\beta(m', m_1, \ldots, m_t) =
((m'/1 - m_1 \otimes 1, \ldots, m'/1 - m_t \otimes 1),
(m_1 - m_2, \ldots, m_{t - 1} - m_t).
$$
Nous voudrions savoir quand ce complexe est exact.

\begin{lemma}
\label{lemma-recover-module-from-glueing-data}
Supposons que $\varphi : R \to S$ soit un homomorphisme plat et que
$I = (f_1, \ldots, f_t) \subset R$ soit un idéal tel que
$R/I \to S/IS$ soit un isomorphisme. Soit $M$ un $R$-module. Alors le
complexe (\ref{equation-glueing-complex}) est exact.
\end{lemma}

\begin{proof}
Première preuve. Notons
$\check{\mathcal{C}}_R \to \check{\mathcal{C}}_S$ le quasi-isomorphisme de
complexes de {\v C}ech alternés augmentés du
Lemme \ref{lemma-map-identifies-koszul-and-cech-complexes}. Comme ces complexes
sont bornés et à termes plats, le morphisme
$M \otimes_R \check{\mathcal{C}}_R \to
M \otimes_R \check{\mathcal{C}}_S$ est lui aussi un quasi-isomorphisme
(Lemmes \ref{lemma-derived-tor-quasi-isomorphism} et
\ref{lemma-derived-tor-quasi-isomorphism-other-side}). Or le complexe
(\ref{equation-glueing-complex}) est une troncation du cône du morphisme
$M \otimes_R \check{\mathcal{C}}_R \to
M \otimes_R \check{\mathcal{C}}_S$, ce qui conclut.

\medskip\noindent
Seconde preuve, par le calcul. Soit $m \in M$. Si $\alpha(m) = 0$, alors
$m \in M[I^\infty]$ d'après le Lemme \ref{lemma-torsion-free}. Choisissons $n$
tel que $I^n m = 0$ et considérons l'application $\varphi : R/I^n \to M$.
Si $m \otimes 1 = 0$, alors $\varphi \otimes 1_S = 0$, donc $\varphi = 0$
(voir Lemme \ref{lemma-neighbourhood-equivalence}), et par conséquent $m = 0$.
Ainsi, $\alpha$ est injective.

\medskip\noindent
Soit $(m', m'_1, \ldots, m'_t) \in \Ker(\beta)$. Écrivons
$m'_i = m_i/f_i^n$ pour un certain $n > 0$ et avec $m_i \in M$. Quitte à
augmenter $n$, on peut supposer que $f_i^n m' = m_i \otimes 1$ dans
$M \otimes_R S$ et que $f_j^nm_i - f_i^nm_j = 0$ dans $M$. En particulier,
$(m_1, \ldots, m_t)$ définit un élément $\xi$ de
$H_1(M, (f_1^n, \ldots, f_t^n))$. Comme
$H_1(M, (f_1^n, \ldots, f_t^n))$ est annulé par $I^{tn + 1}$ (voir
Lemme \ref{lemma-koszul-homology-annihilated}) et que $R \to S$ est plat, on a
$$
H_1(M, (f_1^n, \ldots, f_t^n)) =
H_1(M, (f_1^n, \ldots, f_t^n)) \otimes_R S =
H_1(M \otimes_R S, (f_1^n, \ldots, f_t^n))
$$
d'après le Lemme \ref{lemma-neighbourhood-isomorphism}. L'existence de $m'$
implique que $\xi$ s'envoie sur zéro dans le dernier groupe, c'est-à-dire que
$\xi$ est nul. Il existe donc $m \in M$ tel que $m_i = f_i^n m$. Alors
$(m', m'_1, \ldots, m'_t) - \alpha(m) = (m'', 0, \ldots, 0)$ pour un certain
$m'' \in (M \otimes_R S)[(IS)^\infty]$. D'après le
Lemme \ref{lemma-neighbourhood-equivalence}, on a
$m'' \in M[I^\infty]$, et cela conclut.
\end{proof}
\begin{remark}
\label{remark-glueing-data}
Dans cette remarque, nous définissons une catégorie de données de recollement.
Soit $R \to S$ un homomorphisme d'anneaux. Soient
$f_1, \ldots, f_t \in R$ et $I = (f_1, \ldots, f_t)$. La catégorie
$\text{Glue}(R \to S, f_1, \ldots, f_t)$ est la catégorie dont :
\begin{enumerate}
\item les objets sont les systèmes $(M', M_i, \alpha_i, \alpha_{ij})$, où
$M'$ est un $S$-module, $M_i$ est un $R_{f_i}$-module,
$\alpha_i : (M')_{f_i} \to M_i \otimes_R S$ est un isomorphisme et les
$\alpha_{ij} : (M_i)_{f_j} \to (M_j)_{f_i}$ sont des isomorphismes tels que :
\begin{enumerate}
\item $\alpha_{ij} \circ \alpha_i = \alpha_j$ comme applications
$(M')_{f_if_j} \to (M_j)_{f_i}$ ;
\item $\alpha_{jk} \circ \alpha_{ij} = \alpha_{ik}$ comme applications
$(M_i)_{f_jf_k} \to (M_k)_{f_if_j}$ (condition de cocycle) ;
\end{enumerate}
\item les morphismes
$(M', M_i, \alpha_i, \alpha_{ij}) \to (N', N_i, \beta_i, \beta_{ij})$
sont donnés par des applications $\varphi' : M' \to N'$ et
$\varphi_i : M_i \to N_i$ compatibles avec les applications
$\alpha_i, \beta_i, \alpha_{ij}, \beta_{ij}$ données.
\end{enumerate}
Il existe un foncteur canonique
$$
\text{Can} : \text{Mod}_R
\longrightarrow
\text{Glue}(R \to S, f_1, \ldots, f_t),
\quad
M \longmapsto (M \otimes_R S, M_{f_i}, \text{can}_i, \text{can}_{ij})
$$
où $\text{can}_i : (M \otimes_R S)_{f_i} \to M_{f_i} \otimes_R S$ et
$\text{can}_{ij} : (M_{f_i})_{f_j} \to (M_{f_j})_{f_i}$ sont les
isomorphismes canoniques. Pour tout objet
$\mathbf{M} = (M', M_i, \alpha_i, \alpha_{ij})$ de la catégorie
$\text{Glue}(R \to S, f_1, \ldots, f_t)$, nous définissons
$$
H^0(\mathbf{M}) =
\{(m', m_i) \mid \alpha_i(m') = m_i \otimes 1, \alpha_{ij}(m_i) = m_j\}
$$
autrement dit par la suite exacte
$$
0 \to H^0(\mathbf{M}) \to
M' \times \prod M_i \to
\prod M'_{f_i}
\times
\prod (M_i)_{f_j}
$$
analogue à (\ref{equation-glueing-complex}). Nous considérons
$H^0(\mathbf{M})$ comme un $R$-module. Nous obtenons ainsi un foncteur
$$
H^0 :
\text{Glue}(R \to S, f_1, \ldots, f_t)
\longrightarrow
\text{Mod}_R
$$
Notre prochain objectif est de montrer que les foncteurs $\text{Can}$ et $H^0$
sont parfois quasi inverses l'un de l'autre.
\end{remark}

\begin{lemma}
\label{lemma-H0-Can-adjoint}
Dans la Remarque \ref{remark-glueing-data}, le foncteur
$H^0 : \text{Glue}(R \to S, f_1, \ldots, f_t) \to \text{Mod}_R$ est adjoint à
droite du foncteur
$\text{Can} : \text{Mod}_R \to \text{Glue}(R \to S, f_1, \ldots, f_t)$.
\end{lemma}

\begin{proof}
Soit $\mathbf{M} = (M', M_i, \alpha_i, \alpha_{ij})$ un objet de
$\text{Glue}(R \to S, f_1, \ldots, f_t)$. Pour tout $R$-module $N$, il existe
une application
$$
\Hom_{\text{Glue}(R \to S, f_1, \ldots, f_t)}(\text{Can}(N), \mathbf{M})
\to
\Hom_R(N, H^0(\mathbf{M}))
$$
qui envoie $\psi$ sur la composée de $H^0(\psi)$ et de l'application évidente
$N \to H^0(\text{Can}(N))$. Par construction, l'application affichée est un
isomorphisme pour $N = R$ (même si $R \to H^0(\text{Can}(R))$ n'est pas un
isomorphisme en général). La catégorie
$\text{Glue}(R \to S, f_1, \ldots, f_t)$ admet des sommes directes et des
conoyaux. Le foncteur $\text{Can}$ commute aux sommes directes et aux conoyaux.
Ces observations montrent que l'application affichée est bijective, en
écrivant $N$ comme conoyau d'un morphisme entre $R$-modules libres. Nous
omettons les détails.
\end{proof}

\begin{lemma}
\label{lemma-H0-inverse}
Supposons que $\varphi : R \to S$ soit un homomorphisme plat et que
$I = (f_1, \ldots, f_t) \subset R$ soit un idéal tel que
$R/I \to S/IS$ soit un isomorphisme. Alors le foncteur $H^0$ est un quasi
inverse à gauche du foncteur $\text{Can}$ de la
Remarque \ref{remark-glueing-data}.
\end{lemma}

\begin{proof}
C'est une reformulation du
Lemme \ref{lemma-recover-module-from-glueing-data}.
\end{proof}

\begin{lemma}
\label{lemma-exact}
Supposons que $\varphi : R \to S$ soit un homomorphisme plat et soit
$I = (f_1, \ldots, f_t) \subset R$ un idéal. Alors
$\text{Glue}(R \to S, f_1, \ldots, f_t)$ est une catégorie abélienne, et le
foncteur $\text{Can}$ est exact et commute aux colimites quelconques.
\end{lemma}

\begin{proof}
Étant donné un morphisme
$(\varphi', \varphi_i) :
(M', M_i, \alpha_i, \alpha_{ij})
\to
(N', N_i, \beta_i, \beta_{ij})$
de la catégorie $\text{Glue}(R \to S, f_1, \ldots, f_t)$, son noyau existe et
est l'objet $(\Ker(\varphi'), \Ker(\varphi_i), \alpha_i, \alpha_{ij})$, tandis
que son conoyau existe et est l'objet
$(\Coker(\varphi'), \Coker(\varphi_i), \beta_i, \beta_{ij})$. Cela fonctionne
parce que $R \to S$ est plat, de sorte que la prise de noyaux et de conoyaux
commute avec $- \otimes_R S$. Les détails sont omis. L'exactitude résulte de la
platitude sur $R$ de $R_{f_i}$ et de $S$, tandis que la commutation aux
colimites vient de ce que les produits tensoriels commutent aux colimites.
\end{proof}

\begin{lemma}
\label{lemma-equivalence-I-unit}
Soit $\varphi : R \to S$ un homomorphisme plat et supposons
$(f_1, \ldots, f_t) = R$. Alors $\text{Can}$ et $H^0$ sont des équivalences de
catégories quasi inverses
$$
\text{Mod}_R = \text{Glue}(R \to S, f_1, \ldots, f_t)
$$
\end{lemma}

\begin{proof}
Considérons un objet $\mathbf{M} = (M', M_i, \alpha_i, \alpha_{ij})$ de
$\text{Glue}(R \to S, f_1, \ldots, f_t)$. D'après Algèbre, Lemme
\ref{algebra-lemma-glue-modules}, il existe un unique module $M$ et des
isomorphismes $M_{f_i} \to M_i$ qui redonnent les données de recollement
$\alpha_{ij}$. Alors $M'$ et $M \otimes_R S$ sont deux $S$-modules qui redonnent
deux les modules $M_i \otimes_R S$ après localisation en $f_i$. Il existe donc
un isomorphisme canonique $M \otimes_R S \to M'$. Cela montre que $\mathbf{M}$
appartient à l'image essentielle de $\text{Can}$. Le lemme résulte alors du
Lemme \ref{lemma-H0-inverse}.
\end{proof}

\begin{lemma}
\label{lemma-base-change-glue}
Soit $\varphi : R \to S$ un homomorphisme plat et soit
$I = (f_1, \ldots, f_t)$ un idéal. Soit $R \to R'$ un homomorphisme plat, et
posons $S' = S \otimes_R R'$. On obtient alors un diagramme commutatif de
catégories et de foncteurs
$$
\xymatrix{
\text{Mod}_R \ar[r]_-{\text{Can}} \ar[d]_{-\otimes_R R'} &
\text{Glue}(R \to S, f_1, \ldots, f_t) \ar[r]_-{H^0} \ar[d]^{-\otimes_R R'} &
\text{Mod}_R \ar[d]^{-\otimes_R R'} \\
\text{Mod}_{R'} \ar[r]^-{\text{Can}} &
\text{Glue}(R' \to S', f_1, \ldots, f_t) \ar[r]^-{H^0} &
\text{Mod}_{R'}
}
$$
\end{lemma}

\begin{proof}
Omis.
\end{proof}

\begin{proposition}
\label{proposition-equivalence}
Supposons que $\varphi : R \to S$ soit un homomorphisme plat et que
$I = (f_1, \ldots, f_t) \subset R$ soit un idéal tel que
$R/I \to S/IS$ soit un isomorphisme. Alors $\text{Can}$ et $H^0$ sont des
équivalences de catégories quasi inverses
$$
\text{Mod}_R = \text{Glue}(R \to S, f_1, \ldots, f_t)
$$
\end{proposition}

\begin{proof}
Nous avons déjà vu que $H^0 \circ \text{Can}$ est isomorphe au foncteur
identité ; voir Lemme \ref{lemma-H0-inverse}. Considérons un objet
$\mathbf{M} = (M', M_i, \alpha_i, \alpha_{ij})$ de
$\text{Glue}(R \to S, f_1, \ldots, f_t)$. On obtient un morphisme naturel
$$
\Psi :
(H^0(\mathbf{M}) \otimes_R S, H^0(\mathbf{M})_{f_i},
\text{can}_i, \text{can}_{ij})
\longrightarrow
(M', M_i, \alpha_i, \alpha_{ij}).
$$
En effet, par définition, $H^0(\mathbf{M})$ est muni d'applications
compatibles de $R$-modules $H^0(\mathbf{M}) \to M'$ et
$H^0(\mathbf{M}) \to M_i$. Il faut montrer que ce morphisme est un
isomorphisme.

\medskip\noindent
Choisissons un indice $i$ et posons $R' = R_{f_i}$. En combinant les
Lemmes \ref{lemma-base-change-glue} et \ref{lemma-equivalence-I-unit}, on voit
que $\Psi \otimes_R R'$ est un isomorphisme. Le noyau, resp.\ le conoyau de
$\Psi$ est donc un système de la forme $(K, 0, 0, 0)$, resp.\ $(Q, 0, 0, 0)$.
Remarquons que $H^0((K, 0, 0, 0)) = K$, que $H^0$ est exact à gauche et que,
par construction, $H^0(\Psi)$ est bijectif. Ainsi $K = 0$, autrement dit le
noyau de $\Psi$ est nul.

\medskip\noindent
Il résulte de ce qui précède que l'on a une suite exacte courte
$$
0 \to H^0(\mathbf{M}) \otimes_R S \to M' \to Q \to 0
$$
et que $M_i = H^0(\mathbf{M})_{f_i}$. On peut considérer $Q$ comme un
$R$-module de torsion annulée par une puissance de $I$, de sorte que
$Q = Q \otimes_R S$. Le Lemme \ref{lemma-neighbourhood-extensions} fournit un
diagramme commutatif
$$
\xymatrix{
0 \ar[r] &
H^0(\mathbf{M}) \ar[r] \ar[d] &
E \ar[r] \ar[d] &
Q \ar[r] \ar[d] &
0 \\
0 \ar[r] &
H^0(\mathbf{M}) \otimes_R S \ar[r] &
M' \ar[r] &
Q \ar[r] &
0
}
$$
à lignes exactes. Cela détermine manifestement un isomorphisme
$\text{Can}(E) \to (M', M_i, \alpha_i, \alpha_{ij})$ dans la catégorie
$\text{Glue}(R \to S, f_1, \ldots, f_t)$, et conclut la preuve. (Bien entendu,
a posteriori, on a $Q = 0$.)
\end{proof}

\begin{lemma}
\label{lemma-application-formal-glueing}
Soient $\varphi : R \to S$ un homomorphisme plat et $I \subset R$ un idéal de
type fini tel que $R/I \to S/IS$ soit un isomorphisme.
\begin{enumerate}
\item Étant donnés un $R$-module $N$, un $S$-module $M'$ et un morphisme de
$S$-modules $\varphi : M' \to N \otimes_R S$ dont le noyau et le conoyau sont
de torsion annulée par une puissance de $I$, il existe un morphisme de
$R$-modules $\psi : M \to N$ et un isomorphisme $M \otimes_R S = M'$
compatibles avec $\varphi$ et $\psi$.
\item Étant donnés un $R$-module $M$, un $S$-module $N'$ et un morphisme de
$S$-modules $\varphi : M \otimes_R S \to N'$ dont le noyau et le conoyau sont
de torsion annulée par une puissance de $I$, il existe un morphisme de
$R$-modules $\psi : M \to N$ et un isomorphisme $N \otimes_R S = N'$
compatibles avec $\varphi$ et $\psi$.
\end{enumerate}
Dans les deux cas, on a $\Ker(\varphi) \cong \Ker(\psi)$ et
$\Coker(\varphi) \cong \Coker(\psi)$.
\end{lemma}

\begin{proof}
Démontrons (1). Écrivons $I = (f_1, \ldots, f_t)$. Il est clair que la
localisation $\varphi_{f_i}$ est un isomorphisme. Ainsi,
$(M', N_{f_i}, \varphi_{f_i}, can_{ij})$ est un objet de
$\text{Glue}(R \to S, f_1, \ldots, f_t)$ ; voir
Remarque \ref{remark-glueing-data}. La Proposition
\ref{proposition-equivalence} donne un $R$-module $M$ tel que
$M' = M \otimes_R S$ et $N_{f_i} = M_{f_i}$, de façon compatible avec les
isomorphismes $\varphi_{f_i}$ et $can_{ij}$. Il existe un morphisme
$$
(M \otimes_R S, M_{f_i}, can_i, can_{ij}) =
(M', N_{f_i}, \varphi_{f_i}, can_{ij})
\to
(N \otimes_R S, N_{f_i}, can_i, can_{ij})
$$
dans $\text{Glue}(R \to S, f_1, \ldots, f_t)$, dont la première composante est
$\varphi$. Il correspond à un morphisme de $R$-modules $\psi : M \to N$ par
l'équivalence de catégories de la Proposition
\ref{proposition-equivalence}. La composée du changement de base de
$M \to N$ avec l'isomorphisme $M' \cong M \otimes_R S$ est $\varphi$ ;
autrement dit, $M \to N$ est compatible avec $\varphi$.

\medskip\noindent
Démontrons (2). L'argument est le dual de celui qui précède. En effet, la
localisation $\varphi_{f_i}$ est un isomorphisme. Ainsi,
$(N', M_{f_i}, \varphi_{f_i}^{-1}, can_{ij})$ est un objet de
$\text{Glue}(R \to S, f_1, \ldots, f_t)$ ; voir
Remarque \ref{remark-glueing-data}. La Proposition
\ref{proposition-equivalence} donne un $R$-module $N$ tel que
$N' = N \otimes_R S$ et $N_{f_i} = M_{f_i}$, de façon compatible avec les
isomorphismes $\varphi_{f_i}^{-1}$ et $can_{ij}$. Il existe un morphisme
$$
(M \otimes_R S, M_{f_i}, can_i, can_{ij}) \to
(N', M_{f_i}, \varphi_{f_i}, can_{ij}) =
(N \otimes_R S, N_{f_i}, can_i, can_{ij})
$$
dans $\text{Glue}(R \to S, f_1, \ldots, f_t)$, dont la première composante est
$\varphi$. Il correspond à un morphisme de $R$-modules $\psi : M \to N$ par
l'équivalence de catégories de la Proposition
\ref{proposition-equivalence}. La composée du changement de base de
$M \to N$ avec l'isomorphisme $N' \cong N \otimes_R S$ est $\varphi$ ;
autrement dit, $M \to N$ est compatible avec $\varphi$.

\medskip\noindent
La dernière assertion résulte par exemple du
Lemme \ref{lemma-neighbourhood-equivalence}.
\end{proof}

\noindent
Spécialisons maintenant la Proposition \ref{proposition-equivalence} afin
d'obtenir un énoncé plus facile à utiliser. Lorsque $I = (f)$ est principal,
les objets de $\text{Glue}(R \to S, f)$ sont simplement des triplets
$(M', M_1, \alpha_1)$, et il n'y a {\it aucune} condition de cocycle à
vérifier !

\begin{theorem}
\label{theorem-formal-glueing}
Soient $R$ un anneau et $f \in R$. Soit $\varphi : R \to S$ un homomorphisme
plat induisant un isomorphisme $R/fR \to S/fS$. Alors le foncteur
$$
\text{Mod}_R
\longrightarrow
\text{Mod}_S \times_{\text{Mod}_{S_f}} \text{Mod}_{R_f},
\quad
M
\longmapsto
(M \otimes_R S, M_f, \text{can})
$$
est une équivalence.
\end{theorem}

\begin{proof}
La catégorie figurant à droite de la flèche est la catégorie des triplets
$(M', M_1, \alpha_1)$, où $M'$ est un $S$-module, $M_1$ un $R_f$-module et
$\alpha_1 : M'_f \to M_1 \otimes_R S$ un $S_f$-isomorphisme ; voir Catégories,
Exemple \ref{categories-example-2-fibre-product-categories}. Ce théorème est
donc un cas particulier de la Proposition \ref{proposition-equivalence}.
\end{proof}

\noindent
Un cas particulier utile du Théorème \ref{theorem-formal-glueing} est celui où
$R$ est noethérien et $S$ le complété de $R$ en un élément $f$. Le morphisme de
complétion $R \to S$ est plat, et le foncteur
$M \mapsto M \otimes_R S$ s'identifie au foncteur de complétion $f$-adique
lorsque $M$ est de type fini.
Plus précisément, notons $\text{Mod}^{fg}_R$ la catégorie des $R$-modules de
type fini.

\begin{proposition}
\label{proposition-formal-glueing}
Soit $R$ un anneau noethérien. Soit $f \in R$, et soit $R^\wedge$ le complété
$f$-adique de $R$. Alors le foncteur $M \mapsto (M^\wedge, M_f, \text{can})$
définit une équivalence
$$
\text{Mod}^{fg}_R
\longrightarrow
\text{Mod}^{fg}_{R^\wedge}
\times_{\text{Mod}^{fg}_{(R^\wedge)_f}}
\text{Mod}^{fg}_{R_f}
$$
\end{proposition}

\begin{proof}
L'homomorphisme $R \to R^\wedge$ est plat d'après Algèbre, Lemme
\ref{algebra-lemma-completion-flat}. Il est clair que
$R/fR = R^\wedge/fR^\wedge$. D'après Algèbre, Lemme
\ref{algebra-lemma-completion-tensor}, le complété d'un $R$-module de type fini
$M$ est égal à $M \otimes_R R^\wedge$. Le foncteur affiché dans la proposition
est donc celui du Théorème \ref{theorem-formal-glueing}. Il est en particulier
pleinement fidèle. Soit $(M_1, M_2, \psi)$ un objet du membre de droite. Le
Théorème \ref{theorem-formal-glueing} donne un $R$-module $M$ tel que
$M_1 = M \otimes_R R^\wedge$ et $M_2 = M_f$. Comme
$R \to R^\wedge \times R_f$ est fidèlement plat, Algèbre, Lemme
\ref{algebra-lemma-cover} montre que $M$ est de type fini, c'est-à-dire
$M \in \text{Mod}^{fg}_R$. Cela démontre la proposition.
\end{proof}

\begin{remark}
\label{remark-formal-glueing-algebras}
Les équivalences de la Proposition \ref{proposition-equivalence}, du
Théorème \ref{theorem-formal-glueing} et de la
Proposition \ref{proposition-formal-glueing} préservent les propriétés des
modules. Par exemple, si $M$ correspond à
$\mathbf{M} = (M', M_i, \alpha_i, \alpha_{ij})$, alors $M$ est de type fini,
de présentation finie, plat ou projectif sur $R$ si et seulement si $M'$ et
les $M_i$ possèdent la propriété correspondante sur $S$ et $R_{f_i}$. Cela
résulte de ce que $R \to S \times \prod R_{f_i}$ est fidèlement plat, ainsi que
de la descente et de l'ascension de ces propriétés le long d'homomorphismes
fidèlement plats ; voir Algèbre, Lemme
\ref{algebra-lemma-descend-properties-modules} et Théorème
\ref{algebra-theorem-ffdescent-projectivity}. Ces foncteurs préservent aussi
les structures $\otimes$ des deux côtés. Ils définissent donc des équivalences
entre diverses catégories construites à partir du couple
$(\text{Mod}_R, \otimes)$, par exemple la catégorie des algèbres.
\end{remark}

\begin{remark}
\label{remark-topological-analogue}
Soient $X$ une variété différentielle et $Z$ une sous-variété compacte et
fermée de complémentaire $U$. Se donner un faisceau sur $X$ revient à se
donner un faisceau sur $U$, un faisceau sur un voisinage tubulaire non spécifié
$T$ de $Z$ dans $X$, et un isomorphisme entre les deux faisceaux obtenus sur
$T \cap U$. Les voisinages tubulaires n'existent pas tels quels en géométrie
algébrique, mais des résultats comme la Proposition
\ref{proposition-equivalence}, le Théorème \ref{theorem-formal-glueing} et la
Proposition \ref{proposition-formal-glueing} permettent de travailler à leur
place avec des voisinages formels.
\end{remark}







\section{Le théorème de Beauville-Laszlo}
\label{section-beauville-laszlo}

\noindent
Soient $R$ un anneau et $f$ un élément de $R$. Notons
$R^\wedge = \lim R/f^n R$ le complété $f$-adique de $R$. Dans cette section,
nous étudions et généralisons légèrement un théorème de Beauville et Laszlo ;
voir \cite{Beauville-Laszlo}. Ce théorème affirme que, sous des hypothèses
convenables, un module sur $R$ peut être construit en « recollant » des modules
sur $R^\wedge$ et sur $R_f$ au moyen d'un isomorphisme entre leurs extensions
des scalaires à $(R^\wedge)_f$.

\medskip\noindent
Dans \cite{Beauville-Laszlo}, on suppose que $f$ n'est diviseur de zéro ni sur
$R$ ni sur $M$. En fait, il suffit de supposer que
$$
R[f^\infty] \longrightarrow R^\wedge[f^\infty]
$$
est bijective et que
$$
M[f^\infty] \longrightarrow M \otimes_R R^\wedge
$$
est injective. Cette amélioration a été en partie inspirée par une autre
méthode de recollement introduite dans \cite[\S 1.3]{Kedlaya-Liu-I} pour la
théorie des espaces analytiques non archimédiens.

\medskip\noindent
Nous établirons en réalité le théorème de Beauville-Laszlo dans le cadre plus
général d'un homomorphisme d'anneaux
$$
R \longrightarrow R'
$$
qui induit des isomorphismes $R/f^nR \to R'/f^nR'$ pour tout $n > 0$ et un
isomorphisme $R[f^\infty] \to R'[f^\infty]$. Cette formulation se globalise
mieux et ne découle pas formellement du cas où $R'$ est le complété de $R$ :
par exemple, la bijectivité de $R[f^\infty] \to R'[f^\infty]$ n'implique pas
celle de $R[f^\infty] \to R^\wedge[f^\infty]$.

\medskip\noindent
Le théorème de Beauville et Laszlo démontré dans cette section peut être vu
comme une version non plate du Théorème \ref{theorem-formal-glueing} et, dans
le cas $R' = R^\wedge$, comme une version non noethérienne de la
Proposition \ref{proposition-formal-glueing}. Pour une comparaison avec la
descente plate, voir la Remarque \ref{remark-not-descent}.

\medskip\noindent
On peut établir des résultats encore plus forts, par exemple sans imposer de
condition à $M$, mais il faut alors travailler au niveau des catégories
dérivées. Voir \cite[\S 5]{Bhatt-Algebraize} pour plus de détails.

\begin{lemma}
\label{lemma-same-quotients}
Soient $R$ un anneau et $f \in R$. Pour tout entier strictement positif $n$,
l'application $R/f^nR \to R^\wedge/f^n R^\wedge$ est un isomorphisme.
\end{lemma}

\begin{proof}
C'est un cas particulier de Algèbre, Lemme
\ref{algebra-lemma-hathat-finitely-generated}.
\end{proof}

\noindent
Nous utiliserons les notations introduites dans la Section
\ref{section-torsion}. Ainsi, pour un $R$-module $M$, nous notons $M[f^n]$ le
sous-module de $M$ annulé par $f^n$ et nous posons
$$
M[f^\infty] = \bigcup\nolimits_{n = 1}^\infty M[f^n] = \Ker(M \to M_f).
$$
Si $M = M[f^\infty]$, nous disons que $M$ est un module de torsion annulée par
une puissance de $f$.

\begin{lemma}
\label{lemma-BL-faithful}
Soient $R$ un anneau, $f \in R$ et $R \to R'$ un homomorphisme d'anneaux qui
induit des isomorphismes $R/f^nR \to R'/f^nR'$ pour $n > 0$. Le $R$-module
$R' \oplus R_f$ est fidèle : pour tout $R$-module non nul $M$, le module
$M \otimes_R (R' \oplus R_f)$ est lui aussi non nul. Par exemple, si $M$ est
non nul, alors $M \otimes_R (R^\wedge \oplus R_f)$ est non nul.
\end{lemma}

\noindent
Toutefois, l'application $M \to M \otimes_R (R' \oplus R_f)$ n'est pas
nécessairement injective ; voir l'Exemple \ref{example-not-glueing-pair}.

\begin{proof}
Si $M \neq 0$ mais $M \otimes_R R_f = 0$, alors $M$ est de torsion annulée par
une puissance de $f$. Le Lemme \ref{lemma-torsion-completion} donne
$M \otimes_R R' \cong M \neq 0$. La dernière assertion est un cas particulier
de la première d'après le Lemme \ref{lemma-same-quotients}.
\end{proof}

\begin{lemma}
\label{lemma-cover-spec}
Soient $R$ un anneau, $f \in R$ et $R \to R'$ un homomorphisme d'anneaux qui
induit un isomorphisme $R/fR \to R'/fR'$. L'application
$\Spec(R') \amalg \Spec(R_f) \to \Spec(R)$ est surjective. Par exemple,
l'application $\Spec(R^\wedge) \amalg \Spec(R_f) \to \Spec(R)$ est
surjective.
\end{lemma}

\begin{proof}
Rappelons que $\Spec(R) = V(f) \amalg D(f)$, où
$V(f) = \Spec(R/fR)$ et $D(f) = \Spec(R_f)$ ; voir Algèbre, Section
\ref{algebra-section-spectrum-ring}, et en particulier les Lemmes
\ref{algebra-lemma-spec-closed} et \ref{algebra-lemma-standard-open}. Le lemme
résulte donc de ce que l'application $R \to R/fR$ se factorise par $R'$. La
dernière assertion est un cas particulier de la première d'après le
Lemme \ref{lemma-same-quotients}.
\end{proof}

\begin{lemma}
\label{lemma-faithful-descent}
\begin{reference}
Légère généralisation de \cite[Lemme~2(a)]{Beauville-Laszlo}.
\end{reference}
Soient $R$ un anneau, $f \in R$ et $R \to R'$ un homomorphisme d'anneaux qui
induit des isomorphismes $R/f^nR \to R'/f^nR'$ pour $n > 0$. Un $R$-module
$M$ est de type fini si et seulement si le ($R' \oplus R_f$)-module
$M \otimes_R (R' \oplus R_f)$ est de type fini. Par exemple, si
$M \otimes_R (R^\wedge \oplus R_f)$ est de type fini comme module sur
$R^\wedge \oplus R_f$, alors $M$ est un $R$-module de type fini.
\end{lemma}

\begin{proof}
Le sens direct est clair. Supposons donc
$M \otimes_R (R' \oplus R_f)$ de type fini. En écrivant chaque générateur comme
une somme de tenseurs simples, on voit que $M \otimes_R (R' \oplus R_f)$ admet
un système fini de générateurs formé d'éléments de $M$. Il existe donc un
morphisme d'un $R$-module libre de rang fini vers $M$ dont le conoyau est
annulé après produit tensoriel par $R' \oplus R_f$. En appliquant le
Lemme \ref{lemma-BL-faithful} à ce conoyau, on en déduit que $M$ est de type
fini. La dernière assertion est un cas particulier de la première d'après le
Lemme \ref{lemma-same-quotients}.
\end{proof}

\begin{remark}
\label{remark-not-descent}
Bien que $R \to R_f$ soit toujours plat, $R \to R^\wedge$ ne l'est en général
pas, sauf si $R$ est noethérien (voir Algèbre, Lemme
\ref{algebra-lemma-completion-flat}, et la discussion dans Exemples, Section
\ref{examples-section-nonflat}). On ne peut donc pas, en général, appliquer au
morphisme $R \to R^\wedge \oplus R_f$ la descente fidèlement plate étudiée
dans Descente, Section \ref{descent-section-descent-modules}. De plus, même
dans le cas noethérien, la définition usuelle d'une donnée de descente pour ce
morphisme fait intervenir l'anneau $R^\wedge \otimes_R R^\wedge$, que nous
éviterons de considérer dans cette section.
\end{remark}

\noindent
{\bf Paires de recollement.} Soit $R \to R'$ un homomorphisme d'anneaux qui
induit des isomorphismes $R/f^nR \to R'/f^nR'$ pour $n > 0$. Considérons la
suite
\begin{equation}
\label{equation-BL-cech-re}
0 \to R \to R' \oplus R_f \to R'_f \to 0,
\end{equation}
où l'application de droite est la différence des deux homomorphismes
canoniques. Si cette suite est exacte, nous disons que $(R \to R', f)$ est une
\emph{paire de recollement}. Nous disons que $(R, f)$ est une
\emph{paire de recollement} si $(R \to R^\wedge, f)$ est une paire de
recollement ; cela a un sens d'après le Lemme \ref{lemma-same-quotients}.
Ainsi, $(R, f)$ est une paire de recollement si et seulement si la suite
\begin{equation}
\label{equation-BL-cech}
0 \to R \to R^\wedge \oplus R_f \to (R^\wedge)_f \to 0,
\end{equation}
est exacte.

\begin{lemma}
\label{lemma-same-f-torsion}
Soient $R$ un anneau, $f \in R$ et $R \to R'$ un homomorphisme d'anneaux qui
induit des isomorphismes $R/f^nR \to R'/f^nR'$ pour $n > 0$. La suite
(\ref{equation-BL-cech-re}) est :
\begin{enumerate}
\item exacte à droite ;
\item exacte à gauche si et seulement si
$R[f^\infty] \to R'[f^\infty]$ est injective ;
\item exacte au milieu si et seulement si
$R[f^\infty] \to R'[f^\infty]$ est surjective.
\end{enumerate}
En particulier, $(R \to R', f)$ est une paire de recollement si et seulement
si $R[f^\infty] \to R'[f^\infty]$ est bijective. Par exemple, $(R, f)$ est une
paire de recollement si et seulement si
$R[f^\infty] \to R^\wedge[f^\infty]$ est bijective.
\end{lemma}

\begin{proof}
Soit $x \in R'_f$. Écrivons $x = x'/f^n$ avec $x' \in R'$, puis
$x' = x'' + f^n y$ avec $x'' \in R$ et $y \in R'$. On voit alors que
$(y, -x''/f^n)$ s'envoie sur $x$. Cela démontre (1).

\medskip\noindent
La partie (2) résulte de l'égalité $\Ker(R \to R_f) = R[f^\infty]$.

\medskip\noindent
Si la suite est exacte au milieu, alors les éléments de la forme $(x, 0)$ avec
$x \in R'[f^\infty]$ appartiennent à l'image de la première flèche. Cela
implique que $R[f^\infty] \to R'[f^\infty]$ est surjective. Réciproquement,
supposons que $R[f^\infty] \to R'[f^\infty]$ soit surjective. Soit $(x, y)$
un élément du terme du
milieu dont l'image à droite est nulle. Écrivons $y = y'/f^n$ avec
$y' \in R$. Alors $f^n x - y'$ est annulé par une puissance de $f$ dans $R'$.
Par hypothèse, on peut écrire $f^nx - y' = z$ pour un certain
$z \in R[f^\infty]$. Ainsi, $y = y''/f^n$, où $y'' = y' + z$ appartient au
noyau de $R \to R/f^nR$. On peut donc représenter $y$ sous la forme $y'''/1$
avec $y''' \in R$. Alors $x - y'''$ appartient à $R'[f^\infty]$, et donc
$x - y''' = z' \in R[f^\infty]$. On a ainsi
$(x, y'''/1) = (y''' + z', (y''' + z')/1)$, comme voulu.

\medskip\noindent
La dernière assertion du lemme est un cas particulier de l'avant-dernière
d'après le Lemme \ref{lemma-same-quotients}.
\end{proof}

\begin{remark}
\label{remark-BL-special-case}
Supposons que $f$ ne soit pas diviseur de zéro. Alors Algèbre, Lemme
\ref{algebra-lemma-completion-differ-by-torsion} montre que $f$ n'est pas
diviseur de zéro dans $R^\wedge$. Par conséquent, $(R, f)$ est une paire de
recollement.
\end{remark}

\begin{remark}
\label{remark-noetherian-case}
Si $R \to R^\wedge$ est plat, alors, pour tout entier strictement positif $n$,
le produit tensoriel de la suite $0 \to R[f^n] \to R \to R$ par $R^\wedge$
donne la suite
$0 \to R[f^n] \otimes_R R^\wedge \to R^\wedge \to R^\wedge$. En combinant
ceci avec le Lemme \ref{lemma-torsion-completion}, on conclut que
$R[f^n] \to R^\wedge[f^n]$ est un isomorphisme. Ainsi, $(R, f)$ est une paire
de recollement. C'est notamment le cas si $R$ est noethérien ; voir Algèbre,
Lemme \ref{algebra-lemma-completion-flat}.
\end{remark}

\begin{example}
\label{example-not-glueing-pair}
Soit $k$ un corps et posons
$$
R = k[f, T_1, T_2, \ldots]/(fT_1, fT_2 - T_1, fT_3 - T_2, \ldots).
$$
Alors $(R, f)$ n'est pas une paire de recollement, car l'application
$R[f^\infty] \to R^\wedge[f^\infty]$ n'est pas injective : l'image de $T_1$
est divisible par $f$ dans $R^\wedge$. Pour
$$
R = k[f, T_1, T_2, \ldots]/(fT_1, f^2T_2, \ldots),
$$
l'application $R[f^\infty] \to R^\wedge[f^\infty]$ n'est pas surjective, car
l'élément $T_1 + fT_2 + f^2 T_3 + \ldots$ n'appartient pas à son image. En
particulier, la Remarque \ref{remark-noetherian-case} montre que, dans ces deux
exemples, $R \to R^\wedge$ n'est pas plat.
\end{example}

\noindent
{\bf Modules recollables.} Soit $R \to R'$ un homomorphisme d'anneaux qui
induit des isomorphismes $R/f^nR \to R'/f^nR'$ pour $n > 0$. Pour tout
$R$-module $M$, on peut tensoriser (\ref{equation-BL-cech-re}) par $M$ et
obtenir une suite
\begin{equation}
\label{equation-BL-cech-mod-re}
0 \to M \to (M \otimes_R R') \oplus (M \otimes_R R_f) \to
M \otimes_R R'_f \to 0
\end{equation}
Remarquons que $M \otimes_R R_f = M_f$ et que
$M \otimes_R R'_f = (M \otimes_R R')_f$. Si cette suite est exacte, nous
disons que $M$ est \emph{recollable pour $(R \to R', f)$}. Si $R$ est un
anneau et $f \in R$, nous disons qu'un $R$-module est
\emph{recollable} lorsque $M$ est recollable pour
$(R \to R^\wedge, f)$. Ainsi, $M$ est recollable si et seulement si la suite
\begin{equation}
\label{equation-BL-cech-mod}
0 \to M \to (M \otimes_R R^\wedge) \oplus (M \otimes_R R_f) \to
M \otimes_R (R^\wedge)_f \to 0
\end{equation}
est exacte.

\begin{lemma}
\label{lemma-same-f-torsion-module}
Soient $R$ un anneau, $f \in R$ et $R \to R'$ un homomorphisme d'anneaux qui
induit des isomorphismes $R/f^nR \to R'/f^nR'$ pour $n > 0$. La suite
(\ref{equation-BL-cech-mod-re}) est :
\begin{enumerate}
\item exacte à droite ;
\item exacte à gauche si et seulement si
$M[f^\infty] \to (M \otimes_R R')[f^\infty]$ est injective ;
\item exacte au milieu si et seulement si
$M[f^\infty] \to (M \otimes_R R')[f^\infty]$ est surjective.
\end{enumerate}
Ainsi, $M$ est recollable pour $(R \to R', f)$ si et seulement si
$M[f^\infty] \to (M \otimes_R R')[f^\infty]$ est bijective. Si
$(R \to R', f)$ est une paire de recollement, alors $M$ est recollable pour
$(R \to R', f)$ si et seulement si l'application
$M[f^\infty] \to (M \otimes_R R')[f^\infty]$ est injective. Par exemple, si
$(R, f)$ est une paire de recollement, alors $M$ est recollable si et seulement
si $M[f^\infty] \to (M \otimes_R R^\wedge)[f^\infty]$ est injective.
\end{lemma}

\begin{proof}
Nous utiliserons sans autre mention les résultats du
Lemme \ref{lemma-same-f-torsion}. Le foncteur $M \otimes_R -$ est exact à
droite (Algèbre, Lemme \ref{algebra-lemma-tensor-product-exact}), d'où (1).

\medskip\noindent
Le noyau de $M \to M \otimes_R R_f = M_f$ est $M[f^\infty]$. Cela donne (2).

\medskip\noindent
Si la suite est exacte au milieu, alors les éléments de la forme $(x, 0)$ avec
$x \in (M \otimes_R R')[f^\infty]$ appartiennent à l'image de la première
flèche. L'application $M[f^\infty] \to (M \otimes_R R')[f^\infty]$ est donc
surjective. Réciproquement, supposons que
$M[f^\infty] \to (M \otimes_R R')[f^\infty]$ soit surjective. Soit
$(x, y)$ un élément du terme du milieu dont l'image à droite est nulle.
Écrivons $y = y'/f^n$ avec $y' \in M$. Alors $f^n x - y'$ est annulé par une
puissance de $f$ dans $M \otimes_R R'$. Par hypothèse, on peut écrire
$f^nx - y' = z$ pour un certain $z \in M[f^\infty]$. Ainsi,
$y = y''/f^n$, où $y'' = y' + z$ appartient au noyau de
$M \to M/f^nM$. On peut donc représenter $y$ sous la forme $y'''/1$ avec
$y''' \in M$. Alors $x - y'''$ appartient à
$(M \otimes_R R')[f^\infty]$, donc $x - y''' = z' \in M[f^\infty]$. On a ainsi
$(x, y'''/1) = (y''' + z', (y''' + z')/1)$, comme voulu.

\medskip\noindent
Si $(R \to R', f)$ est une paire de recollement, alors
(\ref{equation-BL-cech-mod-re}) est exacte au milieu pour tout $M$ d'après
Algèbre, Lemme \ref{algebra-lemma-tensor-product-exact}. Cela donne
l'avant-dernière assertion du lemme. La dernière résulte de celle-ci et du fait
que $(R, f)$ est une paire de recollement si et seulement si
$(R \to R^\wedge, f)$ en est une.
\end{proof}

\begin{remark}
\label{remark-glueable}
Soient $(R \to R', f)$ une paire de recollement et $M$ un $R$-module. Voici
quelques observations permettant de déterminer si $M$ est recollable pour
$(R \to R', f)$.
\begin{enumerate}
\item D'après le Lemme \ref{lemma-same-f-torsion-module}, $M$ est recollable
pour $(R \to R^\wedge, f)$ si et seulement si
$M[f^\infty] \to M \otimes_R R^\wedge$ est injective. C'est notamment le cas
si $M[f] \to M^\wedge$ est injective, c'est-à-dire lorsque
$M[f] \cap \bigcap_{n = 1}^\infty f^n M = 0$.
\item Si $\text{Tor}_1^R(M, R'_f) = 0$, alors $M$ est recollable pour
$(R \to R', f)$ (utiliser Algèbre, Lemme
\ref{algebra-lemma-long-exact-sequence-tor}). Cela revient à dire que
$\text{Tor}_1^R(M, R')$ est de torsion annulée par une puissance de $f$. En
particulier, tout $R$-module plat est recollable pour $(R \to R', f)$.
\item Si $R \to R'$ est plat, alors $\text{Tor}_1^R(M, R') = 0$ pour tout
$R$-module, et tout $R$-module est donc recollable pour $(R \to R', f)$. C'est
notamment le cas lorsque $R$ est noethérien et $R' = R^\wedge$ ; voir Algèbre,
Lemme \ref{algebra-lemma-completion-flat}.
\end{enumerate}
\end{remark}

\begin{example}[Module non recollable]
\label{example-not-glueable-module}
\begin{reference}
\cite[\S 4, Remarques]{Beauville-Laszlo}
\end{reference}
Soit $R$ l'anneau des germes en $0$ de fonctions $C^\infty$ sur
$\mathbf{R}$. Soit $f \in R$ la fonction $f(x) = x$. Alors $f$ n'est pas
diviseur de zéro dans $R$ ; ainsi, $(R, f)$ est une paire de recollement et
$R^\wedge \cong \mathbf{R}[[x]]$. Soit $\varphi \in R$ la fonction
$\varphi(x) = \text{exp}(-1/x^2)$. La série de Taylor de $\varphi$ est nulle,
donc $\varphi \in \Ker(R \to R^\wedge)$. Comme $\varphi(x) \neq 0$ pour
$x \neq 0$, on voit que $\varphi$ n'est pas diviseur de zéro dans $R$. La
fonction $\varphi/f$ a elle aussi une série de Taylor nulle ; son image dans
$M = R/\varphi R$ est donc un élément non nul de $M[f]$ dont l'image dans
$M \otimes_R R^\wedge = R^\wedge/\varphi R^\wedge = R^\wedge$ est nulle. Par
conséquent, $M$ n'est pas recollable.
\end{example}

\noindent
Effectuons maintenant quelques calculs de groupes de Tor.

\begin{lemma}
\label{lemma-first-tor}
Soit $(R \to R', f)$ une paire de recollement. Alors
$\text{Tor}^R_1(R', f^n R) = 0$ pour tout $n > 0$.
\end{lemma}

\begin{proof}
La suite exacte $0 \to R[f^n] \to R \to f^n R \to 0$ montre qu'il suffit de
vérifier que $R[f^n] \otimes_R R' \to R'$ est injective. Le
Lemme \ref{lemma-torsion-completion} donne
$R[f^n] \otimes_R R' = R[f^n]$, et le Lemme
\ref{lemma-same-f-torsion} montre que $R[f^n] \to R'$ est injective puisque
$(R \to R', f)$ est une paire de recollement.
\end{proof}

\begin{lemma}
\label{lemma-first-tor-total}
Soit $(R \to R',f)$ une paire de recollement. Alors
$\text{Tor}^R_1(R', R/R[f^\infty]) = 0$.
\end{lemma}

\begin{proof}
On a $R/R[f^\infty] = \colim R/R[f^n] = \colim f^nR$. Comme la formation des
groupes de Tor commute aux colimites filtrantes (Algèbre, Lemme
\ref{algebra-lemma-tor-commutes-filtered-colimits}), on peut appliquer le
Lemme \ref{lemma-first-tor}.
\end{proof}

\begin{lemma}
\label{lemma-BL3}
\begin{reference}
Légère généralisation de \cite[Lemme 3(a)]{Beauville-Laszlo}
\end{reference}
Soit $(R \to R', f)$ une paire de recollement. Pour tout $R$-module $M$, on a
$\text{Tor}^R_1(R', \Coker(M \to M_f)) = 0$.
\end{lemma}

\begin{proof}
Posons $\overline{M} = M/M[f^\infty]$. Comme
$\Coker(M \to M_f) \cong \Coker(\overline{M} \to \overline{M}_f)$, on peut et
on va supposer que $f$ n'est pas diviseur de zéro sur $M$. Dans ce cas,
$M \subset M_f$ et $M_f/M = \colim M/f^nM$, les applications de transition
étant données par multiplication par $f$. Puisque la formation des groupes de
Tor commute aux colimites (Algèbre, Lemme
\ref{algebra-lemma-tor-commutes-filtered-colimits}), il suffit de montrer que
$\text{Tor}^R_1(R', M/f^n M) = 0$.

\medskip\noindent
Traitons d'abord le cas $M = R/R[f^\infty]$. D'après le
Lemme \ref{lemma-same-f-torsion}, on a
$M \otimes_R R' = R'/R'[f^\infty]$. La suite exacte courte
$0 \to M \to M \to M/f^nM \to 0$ donne la suite exacte
$$
\xymatrix{
\text{Tor}_1^R(R', R/R[f^\infty]) \ar[r] &
\text{Tor}_1^R(R', M/f^n M) \ar[r] &
R'/R'[f^\infty] \ar[dll]_{f^n} \\
R'/R'[f^\infty] \ar[r] &
(R'/R'[f^\infty])/(f^n
(R'/R'[f^\infty])) \ar[r] & 0
}
$$
d'après Algèbre, Lemme \ref{algebra-lemma-long-exact-sequence-tor}. La flèche
diagonale est ici injective. Comme le premier groupe
$\text{Tor}_1^R(R', R/R[f^\infty])$ est nul d'après le
Lemme \ref{lemma-first-tor-total}, on en déduit que
$\text{Tor}_1^R(R', M/f^nM) = 0$, comme voulu.

\medskip\noindent
Pour traiter le cas général, choisissons une surjection $F \to M$ avec $F$
libre comme $R/R[f^\infty]$-module, et formons une suite exacte
$$
0 \to N \to F/f^n F \to M/f^n M \to 0.
$$
D'après le Lemme \ref{lemma-torsion-completion}, cette suite reste inchangée,
et donc exacte, après tensorisation par $R'$. Puisque
$\text{Tor}^R_1(R', F/f^n F) = 0$ d'après le paragraphe précédent, on en déduit
$\text{Tor}^R_1(R', M/f^n M) = 0$, comme voulu.
\end{proof}

\noindent
Soit $(R \to R', f)$ une paire de recollement. Cela signifie que
$R/f^nR \to R'/f^nR'$ est un isomorphisme pour $n > 0$ et que la suite
$$
0 \to R \to R' \oplus R_f \to R_f' \to 0
$$
est exacte. Considérons la catégorie $\text{Glue}(R \to R', f)$ introduite
dans la Remarque \ref{remark-glueing-data}. Nous appellerons
\emph{donnée de recollement} tout objet $(M', M_1, \alpha_1)$ de
$\text{Glue}(R \to R', f)$. Il est constitué d'un $R'$-module $M'$, d'un
$R_f$-module $M_1$ et d'un isomorphisme
$\alpha_1 : (M')_f \to M_1 \otimes_R R'$. Il existe un foncteur évident
$$
\text{Can} : \text{Mod}_R \longrightarrow \text{Glue}(R \to R', f),\quad
M \longmapsto (M \otimes_R R', M_f, \text{can}),
$$
ainsi qu'un foncteur en sens inverse
$$
H^0 : \text{Glue}(R \to R', f) \longrightarrow \text{Mod}_R,\quad
(M', M_1, \alpha_1) \longmapsto \Ker(M' \oplus M_1 \to (M')_f)
$$
dont la définition précise figure dans la Remarque \ref{remark-glueing-data}.

\begin{theorem}
\label{theorem-BL-glueing}
\begin{reference}
Légère généralisation du théorème principal de \cite{Beauville-Laszlo}.
\end{reference}
Soit $(R \to R',f)$ une paire de recollement. Le foncteur
$\text{Can} : \text{Mod}_R \longrightarrow \text{Glue}(R \to R', f)$ définit
une équivalence entre la catégorie des $R$-modules recollables pour
$(R \to R', f)$ et la catégorie $\text{Glue}(R \to R', f)$ des données de
recollement.
\end{theorem}

\begin{proof}
Soit $(M', M_1, \alpha_1)$ une donnée de recollement. Montrons que
$M = H^0((M', M_1, \alpha_1))$ est recollable pour $(R \to R', f)$ et que
$(M', M_1, \alpha_1) \cong \text{Can}(M)$.

\medskip\noindent
Vérifions d'abord que l'application
$\text{d} : M' \oplus M_1 \to (M')_f$ utilisée pour définir $H^0$ est
surjective. Remarquons que $(x, y) \in M' \oplus M_1$ s'envoie sur
$\text{d}(x, y) = x/1 - \alpha_1^{-1}(y \otimes 1)$ dans $(M')_f$. Si
$z \in (M')_f$, on peut écrire $\alpha_1(z) = \sum y_i \otimes g_i$ avec
$g_i \in R'$ et $y_i \in M_1$. Écrivons
$\alpha_1^{-1}(y_i \otimes 1) = y_i'/f^n$ avec $y'_i \in M'$ et $n \geq 0$
(on peut choisir le même $n$ pour tout $i$). Écrivons encore
$g_i = a_i + f^n b_i$, avec $a_i \in R$ et $b_i \in R'$. Alors, en posant
$y = \sum a_i y_i \in M_1$ et $x = \sum b_i y'_i \in M'$, on obtient
$\text{d}(x, -y) = z$, comme voulu.

\medskip\noindent
Comme $M = H^0((M', M_1, \alpha_1)) = \Ker(\text{d})$, on obtient une suite
exacte de $R$-modules
\begin{equation}
\label{equation-define-M}
0 \to M \to M' \oplus M_1 \to (M')_f \to 0.
\end{equation}
Nous allons montrer que les applications $M \to M'$ et $M \to M_1$ induisent
des isomorphismes $M \otimes_R R' \to M'$ et
$M \otimes_R R_f \to M_1$. Il en résultera que $M$ est recollable pour
$(R \to R', f)$ et que
$\text{Can}(M) \cong (M', M_1, \alpha_1)$, comme voulu.

\medskip\noindent
Puisque $f$ n'est pas diviseur de zéro sur $M_1$, on a
$M[f^\infty] \cong M'[f^\infty]$. Il en résulte une suite exacte
\begin{equation}
\label{equation-exact-mod-torsion}
0 \to M/M[f^\infty] \to M_1 \to (M')_f/M' \to 0.
\end{equation}
Comme $R \to R_f$ est plat, on peut tensoriser cette suite exacte par $R_f$ et
en déduire que
$M \otimes_R R_f = (M/M[f^\infty]) \otimes_R R_f \to M_1$ est un
isomorphisme.

\medskip\noindent
Le Lemme \ref{lemma-BL3} donne
$\text{Tor}_1^R(R', \Coker(M' \to (M')_f)) = 0$. La suite
(\ref{equation-exact-mod-torsion}) reste donc exacte après tensorisation sur
$R$ par $R'$. En utilisant $\alpha_1$ et le
Lemme \ref{lemma-torsion-completion}, la suite exacte obtenue s'écrit
\begin{equation}
\label{equation-mod-torsion-sequence}
0 \to (M/M[f^\infty]) \otimes_R R' \to
(M')_f \to (M')_f/M' \to 0
\end{equation}
Cela donne un isomorphisme
$(M/M[f^\infty]) \otimes_R R' \cong M'/M'[f^\infty]$. On en déduit que, dans
le diagramme
$$
\xymatrix{
& M[f^\infty] \otimes_R R' \ar[r] \ar[d] &
M \otimes_R R'  \ar[r] \ar[d] &
(M/M[f^\infty]) \otimes_R R' \ar[r] \ar[d] & 0 \\
0 \ar[r] &
M'[f^\infty] \ar[r] &
M' \ar[r] &
M'/M'[f^\infty] \ar[r] & 0,
}
$$
la troisième flèche verticale est un isomorphisme. Comme les lignes sont
exactes et que la première flèche verticale est un isomorphisme d'après le
Lemme \ref{lemma-torsion-completion} et l'égalité
$M[f^\infty] = M'[f^\infty]$, le lemme des cinq montre que
$M \otimes_R R' \to M'$ est un isomorphisme.

\medskip\noindent
Ce qui précède montre que $\text{Can}$ est essentiellement surjectif et que le
foncteur $H^0$ prend ses valeurs dans la catégorie des modules recollables.
Comme (\ref{equation-BL-cech-mod-re}) est exacte pour les modules recollables,
on a $H^0 \circ \text{Can} = \text{id}$ sur cette catégorie. Il s'ensuit que
$\text{Can}$ est pleinement fidèle d'après le
Lemme \ref{lemma-H0-Can-adjoint} et Catégories, Lemme
\ref{categories-lemma-adjoint-fully-faithful}. Cela achève la preuve.
\end{proof}

\begin{remark}
\label{remark-what-you-get-for-general-modules}
Soit $(R \to R', f)$ une paire de recollement, et soit $M$ un $R$-module qui
n'est pas nécessairement recollable pour $(R \to R', f)$. En posant
$M' = M \otimes_R R'$ et $M_1 = M_f$, on obtient la donnée de recollement
$\text{Can}(M) = (M', M_1, \text{can})$. Alors
$\tilde M = H^0(M', M_1, \text{can})$ est un $R$-module recollable pour
$(R \to R', f)$, et l'application canonique $M \to \tilde M$ induit des
isomorphismes $M \otimes_R R' \to \tilde M \otimes_R R'$ et
$M_f \to \tilde M_f$ ; voir Théorème \ref{theorem-BL-glueing}. L'exactitude
des suites
$$
M \to (M \otimes_R R' )\oplus M_f \to M \otimes_R (R')_f  \to 0
$$
et
$$
0 \to \tilde M \to (\tilde M \otimes_R R') \oplus \tilde M_f \to
\tilde M \otimes_R  (R')_f \to 0
$$
montre que l'application $M \to \tilde M$ est surjective.
\end{remark}

\noindent
Rappelons que les $R$-modules plats sur une paire de recollement
$(R \to R', f)$ sont recollables (Remarque \ref{remark-glueable}). Le lemme
suivant montre donc que le Théorème \ref{theorem-BL-glueing} définit une
équivalence entre la catégorie des $R$-modules plats et la catégorie des
données de recollement $(M', M_1, \alpha_1)$ dans lesquelles $M'$ et $M_1$
sont plats sur $R'$ et $R_f$.

\begin{lemma}
\label{lemma-BL-flat}
Soit $(R \to R', f)$ une paire de recollement. Soit $M$ un $R$-module qui
n'est pas nécessairement recollable pour $(R \to R', f)$. Alors $M$ est plat
sur $R$ si et seulement si $M \otimes_R R'$ est plat sur $R'$ et $M_f$ est
plat sur $R_f$.
\end{lemma}

\begin{proof}
Un sens résulte de Algèbre, Lemme \ref{algebra-lemma-flat-base-change}. Pour
l'autre, supposons $M \otimes_R R'$ plat sur $R'$ et $M_f$ plat sur $R_f$.
Soit $\tilde M$ le module de la
Remarque \ref{remark-what-you-get-for-general-modules}. Si $\tilde M$ est plat
sur $R$, l'application de Algèbre, Lemme
\ref{algebra-lemma-flat-tor-zero}, à la suite exacte courte
$0 \to \Ker(M \to \tilde M) \to M \to \tilde M \to 0$ montre que
$\Ker(M \to \tilde M) \otimes_R (R' \oplus R_f)$ est nul. Ainsi,
$M = \tilde M$ d'après le Lemme \ref{lemma-BL-faithful}, et l'on conclut.
Autrement dit, on peut remplacer $M$ par $\tilde M$ et supposer $M$ recollable
pour $(R \to R', f)$. Soit $N$ un autre $R$-module. Il suffit de montrer que
$\text{Tor}_1^R(M, N) = 0$ ; voir Algèbre, Lemme
\ref{algebra-lemma-characterize-flat}.

\medskip\noindent
La suite exacte longue de Tor associée à la suite exacte courte
$0 \to R \to R' \oplus R_f \to (R')_f \to 0$ et à $N$ donne une suite exacte
$$
0 \to \text{Tor}_1^R(R', N) \to \text{Tor}_1^R((R')_f, N)
$$
et des isomorphismes
$\text{Tor}_i^R(R', N) = \text{Tor}_i^R((R')_f, N)$ pour $i \geq 2$. Comme
$\text{Tor}_i^R((R')_f, N) = \text{Tor}_i^R(R', N)_f$, on en déduit que $f$
n'est pas diviseur de zéro sur $\text{Tor}_1^R(R', N)$ et est inversible sur
$\text{Tor}_i^R(R', N)$ pour $i \geq 2$. La platitude de
$M \otimes_R R'$ sur $R'$ donne
$$
\text{Tor}_i^R(M \otimes_R R', N) =
(M \otimes_R R') \otimes_{R'} \text{Tor}_i^R(R', N)
$$
par la suite spectrale de l'Exemple \ref{example-tor-change-rings}. En écrivant
$M \otimes_R R'$ comme une colimite filtrante de $R'$-modules libres de rang
fini (Algèbre, Théorème \ref{algebra-theorem-lazard}), on conclut que $f$ n'est
pas diviseur de zéro sur $\text{Tor}_1^R(M \otimes_R R', N)$ et est inversible
sur $\text{Tor}_i^R(M \otimes_R R', N)$. Considérons ensuite la suite exacte
$0 \to M \to M \otimes_R R' \oplus M_f \to M \otimes_R (R')_f \to 0$
qui exprime que $M$ est recollable, ainsi que la suite exacte longue de
$\text{Tor}$ associée. La partie pertinente est
$$
\xymatrix{
\text{Tor}_1^R(M, N) \ar[r] &
\text{Tor}_1^R(M \otimes_R R', N) \ar[r] &
\text{Tor}_1^R(M \otimes_R (R')_f, N) \\
& \text{Tor}_2^R(M \otimes_R R', N) \ar[r] &
\text{Tor}_2^R(M \otimes_R (R')_f, N) \ar[llu]
}
$$
Nous concluons que $\text{Tor}_1^R(M, N) = 0$ d'après les remarques précédentes
sur l'action de $f$ sur les $\text{Tor}_i^R(M \otimes_R R', N)$.
\end{proof}

\noindent
Remarquons que nous avons déjà établi le résultat du lemme suivant pour la
propriété « de type fini » dans le Lemme \ref{lemma-faithful-descent}.

\begin{lemma}
\label{lemma-BL-properties}
Soit $(R \to R', f)$ une paire de recollement. Soit $M$ un $R$-module qui
n'est pas nécessairement recollable pour $(R \to R', f)$. Alors $M$ est un
$R$-module projectif de type fini si et seulement si $M \otimes_R R'$ est
projectif de type fini sur $R'$ et $M_f$ projectif de type fini sur $R_f$.
\end{lemma}

\begin{proof}
Supposons que $M \otimes_R R'$ soit un module projectif de type fini sur $R'$
et que $M_f$ soit un module projectif de type fini sur $R_f$. Il s'agit de
montrer que $M$ est projectif de type fini sur $R$. Nous utiliserons sans autre
mention Algèbre, Lemme \ref{algebra-lemma-finite-projective}. Le
Lemme \ref{lemma-BL-flat} montre que $M$ est plat, et le
Lemme \ref{lemma-faithful-descent} montre que $M$ est de type fini. Choisissons une
suite exacte courte $0 \to K \to R^{\oplus n} \to M \to 0$. Comme un module
projectif de type fini est de présentation finie, et comme la suite reste
exacte après tensorisation par $R'$ (d'après Algèbre, Lemme
\ref{algebra-lemma-flat-tor-zero}) et par $R_f$, on conclut que
$K \otimes_R R'$ et $K_f$ sont de type fini. Le lemme précédent montre alors
que $K$ est de type fini. Ainsi, $M$ est de présentation finie et donc
projectif de type fini.
\end{proof}

\begin{remark}
\label{remark-compare-BL}
Dans \cite{Beauville-Laszlo}, on suppose que $f$ n'est pas diviseur de zéro
dans $R$ et que $R' = R^\wedge$, ce qui donne une paire de recollement d'après
le Lemme \ref{lemma-same-f-torsion}. Même dans ce cadre, le
Théorème \ref{theorem-BL-glueing} apporte un résultat nouveau : les résultats
de \cite{Beauville-Laszlo} ne s'appliquent qu'aux modules sur lesquels $f$
n'est pas diviseur de zéro, donc aux modules recollables au sens présent (voir
Lemme \ref{lemma-same-f-torsion-module}). Le Lemme
\ref{lemma-BL-properties} donne également une légère extension des résultats
de \cite{Beauville-Laszlo} : non seulement on peut autoriser $M$ à avoir une
torsion non nulle annulée par une puissance de $f$, mais on n'exige même pas
qu'il soit recollable.
\end{remark}











\section{Complétion dérivée}
\label{section-derived-completion}

\noindent
Parmi les références concernant cette section figurent
\cite{Dwyer-Greenlees}, \cite{Greenlees-May}, \cite{PSY} et \cite{dag12}
(notamment le Chapitre 4). Notre exposition suit \cite{BS}. L'analogue, ou le
« dual », de cette section pour les modules de torsion se trouve dans
Complexes dualisants, Section \ref{dualizing-section-local-cohomology}. La
relation entre la catégorie dérivée des complexes à cohomologie de torsion et
les complexes complets au sens dérivé est étudiée dans Complexes dualisants,
Section \ref{dualizing-section-torsion-and-complete}.

\medskip\noindent
Soient $K \in D(A)$ et $f \in A$. Nous notons $T(K, f)$ une limite dérivée du
système
$$
\ldots \to K \xrightarrow{f} K \xrightarrow{f} K
$$
dans $D(A)$.

\begin{lemma}
\label{lemma-hom-from-Af}
Soient $A$ un anneau, $f \in A$ et $K \in D(A)$. Les conditions suivantes sont
équivalentes :
\begin{enumerate}
\item $\Ext^n_A(A_f, K) = 0$ pour tout $n$ ;
\item $\Hom_{D(A)}(E, K) = 0$ pour tout $E$ dans $D(A_f)$ ;
\item $T(K, f) = 0$ ;
\item pour tout $p \in \mathbf{Z}$, on a $T(H^p(K), f) = 0$ ;
\item pour tout $p \in \mathbf{Z}$, on a
$\Hom_A(A_f, H^p(K)) = 0$ et $\Ext^1_A(A_f, H^p(K)) = 0$ ;
\item $R\Hom_A(A_f, K) = 0$ ;
\item l'application $\prod_{n \geq 0} K \to \prod_{n \geq 0} K$,
$(x_0, x_1, \ldots) \mapsto (x_0 - fx_1, x_1 - fx_2, \ldots)$ est un
isomorphisme dans $D(A)$ ;
\item en ajouter d'autres ici.
\end{enumerate}
\end{lemma}

\begin{proof}
Il est clair que (2) implique (1) et que (1) équivaut à (6). Supposons (1).
Soit $I^\bullet$ un complexe K-injectif de $A$-modules représentant $K$. La
condition (1) signifie que $\Hom_A(A_f, I^\bullet)$ est acyclique. Soit
$M^\bullet$ un complexe de $A_f$-modules représentant $E$. Alors
$$
\Hom_{D(A)}(E, K) =
\Hom_{K(A)}(M^\bullet, I^\bullet) =
\Hom_{K(A_f)}(M^\bullet, \Hom_A(A_f, I^\bullet))
$$
d'après Algèbre, Lemme \ref{algebra-lemma-adjoint-hom-restrict}. Comme
$\Hom_A(A_f, I^\bullet)$ est un complexe K-injectif de $A_f$-modules d'après
le Lemme \ref{lemma-hom-K-injective}, son acyclicité implique qu'il est
homotopiquement équivalent à zéro (Catégories dérivées, Lemme
\ref{derived-lemma-K-injective}). On obtient donc (2).

\medskip\noindent
Une résolution libre du $A$-module $A_f$ est donnée par
$$
0 \to \bigoplus\nolimits_{n \in \mathbf{N}} A \to
\bigoplus\nolimits_{n \in \mathbf{N}} A
\to A_f \to 0
$$
où la première application envoie $(a_0, a_1, a_2, \ldots)$ sur
$(a_0, a_1 - fa_0, a_2 - fa_1, \ldots)$ et la seconde envoie
$(a_0, a_1, a_2, \ldots)$ sur $a_0 + a_1/f + a_2/f^2 + \ldots$. En appliquant
$\Hom_A(-, I^\bullet)$, on obtient
$$
0 \to \Hom_A(A_f, I^\bullet) \to \prod I^\bullet \to \prod I^\bullet \to 0
$$
Puisque $\prod I^\bullet$ représente $\prod_{n \geq 0} K$, ceci démontre
l'équivalence de (1) et (7). D'autre part, la construction des limites dérivées
dans Catégories dérivées, Section \ref{derived-section-derived-limit}, et la
suite exacte affichée montrent que $T(K, f)$ représente
$R\Hom_A(A_f, K)$ dans $D(A)$. On obtient ainsi l'équivalence de (1) et (3).

\medskip\noindent
Il existe une suite spectrale
$$
E_2^{p, q} = \Ext^p_A(A_f, H^q(K)) \Rightarrow
\Ext^{p + q}_A(A_f, K)
$$
Voir l'Équation (\ref{equation-first-ss-ext})\footnote{Si $K$ n'est pas borné
inférieurement, on ne peut pas utiliser la référence indiquée. Dans ce cas,
on choisit un complexe $K^\bullet$ représentant $K$ et une résolution libre à
deux termes $0 \to F^{-1} \to F^0 \to A_f \to 0$ (voir ci-dessus). On obtient
alors un complexe double de termes $\Hom_A(F^{-q}, K^p)$ dont la totalisation
calcule $R\Hom(A_f, K)$. La seconde suite spectrale associée à ce complexe
double (Homologie, Section \ref{homology-section-double-complex}) donne le
résultat voulu. Les détails sont omis.}. Cette suite spectrale dégénère en
$E_2$, car $A_f$ admet une résolution de longueur $1$ par des $A$-modules
projectifs (voir ci-dessus) ; la page $E_2$ ne possède donc que deux colonnes
non nulles. On obtient ainsi des suites exactes courtes
$$
0 \to \Ext^1_A(A_f, H^{p - 1}(K)) \to
\Ext^p_A(A_f, K) \to
\Hom_A(A_f, H^p(K)) \to 0
$$
Cela démontre que (4) et (5) équivalent à (1).
\end{proof}

\begin{lemma}
\label{lemma-ideal-of-elements-complete-wrt}
Soient $A$ un anneau et $K \in D(A)$. L'ensemble $I$ des $f \in A$ tels que
$T(K, f) = 0$ est un idéal radical de $A$.
\end{lemma}

\begin{proof}
Nous utiliserons sans autre mention les résultats du
Lemme \ref{lemma-hom-from-Af}. Si $f \in I$ et $g \in A$, alors $A_{gf}$ est un
$A_f$-module ; ainsi, $\Ext^n_A(A_{gf}, K) = 0$ pour tout $n$, donc
$gf \in I$. Supposons $f, g \in I$. Il existe une suite exacte courte
$$
0 \to A_{f + g} \to A_{f(f + g)} \oplus A_{g(f + g)} \to A_{gf(f + g)} \to 0
$$
car $f, g$ engendrent l'idéal unité dans $A_{f + g}$. Cela résulte de
Algèbre, Lemme \ref{algebra-lemma-standard-covering}, et du fait élémentaire
que la dernière flèche est surjective. La suite exacte longue de $\Ext$ et
l'annulation de $\Ext^n_A(A_{f(f + g)}, K)$,
$\Ext^n_A(A_{g(f + g)}, K)$ et $\Ext^n_A(A_{gf(f + g)}, K)$ pour tout $n$
entraînent celle de $\Ext^n_A(A_{f + g}, K)$ pour tout $n$. Enfin, si
$f^n \in I$ pour un certain $n > 0$, alors $f \in I$, car
$T(K, f) = T(K, f^n)$, ou encore parce que $A_f \cong A_{f^n}$.
\end{proof}

\begin{lemma}
\label{lemma-complete-derived-complete}
Soient $A$ un anneau, $I \subset A$ un idéal et $M$ un $A$-module.
\begin{enumerate}
\item Si $M$ est complet pour la topologie $I$-adique, alors
$T(M, f) = 0$ pour tout $f \in I$.
\item Réciproquement, si $T(M, f) = 0$ pour tout $f \in I$ et si $I$ est de
type fini, alors $M \to \lim M/I^nM$ est surjective.
\end{enumerate}
\end{lemma}

\begin{proof}
Démontrons (1). Supposons $M$ complet pour la topologie $I$-adique. D'après le
Lemme \ref{lemma-hom-from-Af}, il suffit de montrer
$\Ext^1_A(A_f, M) = 0$ et $\Hom_A(A_f, M) = 0$. Puisque
$M = \lim M/I^nM$ et $\Hom_A(A_f, M/I^nM) = 0$, il s'ensuit que
$\Hom_A(A_f, M) = 0$. Considérons une extension
$$
0 \to M \to E \to A_f \to 0
$$
Pour $n \geq 0$, choisissons $e_n \in E$ s'envoyant sur $1/f^n$. Posons
$\delta_n = fe_{n + 1} - e_n \in M$ pour $n \geq 0$. Remplaçons $e_n$ par
$$
e'_n = e_n + \delta_n + f\delta_{n + 1} + f^2 \delta_{n + 2} + \ldots
$$
La somme infinie existe puisque $M$ est complet pour la topologie $I$-adique
et que $f \in I$. Un calcul immédiat donne $fe'_{n + 1} = e'_n$. On obtient
donc une scission de l'extension en envoyant $1/f^n$ sur $e'_n$.

\medskip\noindent
Démontrons (2). Supposons $I = (f_1, \ldots, f_r)$ et $T(M, f_i) = 0$ pour
$i = 1, \ldots, r$. D'après Algèbre, Lemme
\ref{algebra-lemma-when-surjective-to-completion}, on peut supposer
$I = (f)$ et $T(M, f) = 0$. Soit $x_n \in M$ pour $n \geq 0$. Considérons
l'extension
$$
0 \to M \to E \to A_f \to 0
$$
donnée par
$$
E = M \oplus \bigoplus Ae_n\Big/\langle x_n - fe_{n + 1} + e_n\rangle
$$
où $e_n$ s'envoie sur $1/f^n$ dans $A_f$ (voir ci-dessus). Par hypothèse et
d'après le Lemme \ref{lemma-hom-from-Af}, cette extension est scindée ; il
existe donc un élément $x + e_0$ qui engendre une copie de $A_f$ dans $E$.
Alors
$$
x + e_0 = x - x_0 + fe_1 = x - x_0 - f x_1 + f^2 e_2 = \ldots
$$
Comme $M/f^nM = E/f^nE$ d'après le lemme du serpent, on voit que
$x = x_0 + fx_1 + \ldots + f^{n - 1}x_{n - 1}$ modulo $f^nM$. Autrement dit,
l'application $M \to \lim M/f^nM$ est surjective, comme voulu.
\end{proof}

\noindent
Les résultats qui précèdent motivent la définition suivante.

\begin{definition}
\label{definition-derived-complete}
Soient $A$ un anneau, $K \in D(A)$ et $I \subset A$ un idéal. On dit que $K$
est {\it complet au sens dérivé par rapport à $I$} si, pour tout $f \in I$,
on a $T(K, f) = 0$. Si $M$ est un $A$-module, on dit que $M$ est
{\it complet au sens dérivé par rapport à $I$} si $M[0] \in D(A)$ est complet
au sens dérivé par rapport à $I$.
\end{definition}

\noindent
La sous-catégorie pleine
$D_{comp}(A) = D_{comp}(A, I) \subset D(A)$ formée des objets complets au sens
dérivé est une sous-catégorie triangulée strictement pleine et saturée ; voir
Catégories dérivées, Définitions
\ref{derived-definition-triangulated-subcategory} et
\ref{derived-definition-saturated}. D'après le
Lemme \ref{lemma-ideal-of-elements-complete-wrt}, la sous-catégorie
$D_{comp}(A, I)$ ne dépend que du radical $\sqrt{I}$ de $I$, autrement dit que
du sous-ensemble fermé $Z = V(I)$ de $\Spec(A)$. La sous-catégorie
$D_{comp}(A, I)$ est stable par produits et limites homotopiques dans $D(A)$.
En revanche, elle n'est pas stable en général par sommes directes
dénombrables. Nous dirons souvent simplement que $M$ est un module complet au
sens dérivé lorsque le choix de l'idéal $I$ ressort du contexte.

\begin{proposition}
\label{proposition-derived-complete-modules}
Soient $I \subset A$ un idéal de type fini d'un anneau $A$ et $M$ un
$A$-module. Les conditions suivantes sont équivalentes :
\begin{enumerate}
\item $M$ est complet pour la topologie $I$-adique ;
\item $M$ est complet au sens dérivé par rapport à $I$ et
$\bigcap I^nM = 0$.
\end{enumerate}
\end{proposition}

\begin{proof}
Cela résulte immédiatement du
Lemme \ref{lemma-complete-derived-complete}.
\end{proof}

\noindent
Le lemme suivant montre que la catégorie $\mathcal{C}$ des modules complets au
sens dérivé est abélienne. La catégorie $\mathcal{C}$ n'est pas, en général,
une catégorie abélienne
de Grothendieck ; voir Exemples, Section
\ref{examples-section-derived-complete-modules}.

\begin{lemma}
\label{lemma-serre-subcategory}
Soit $I$ un idéal d'un anneau $A$.
\begin{enumerate}
\item Les $A$-modules complets au sens dérivé forment une sous-catégorie de
Serre faible $\mathcal{C}$ de $\text{Mod}_A$.
\item $D_\mathcal{C}(A) \subset D(A)$ est la sous-catégorie pleine des objets
complets au sens dérivé.
\end{enumerate}
\end{lemma}

\begin{proof}
La partie (2) résulte immédiatement du Lemme \ref{lemma-hom-from-Af} et des
définitions. Pour la partie (1), soit $M \to N$ un morphisme de modules
complets au sens dérivé, et notons $K = (M \to N)$ l'objet correspondant de
$D(A)$. Choisissons $f \in I$. Alors $\Ext_A^n(A_f, K)$ est nul pour tout $n$,
car $\Ext_A^n(A_f, M)$ et $\Ext_A^n(A_f, N)$ sont nuls pour tout $n$. Ainsi,
$K$ est complet au
sens dérivé. La partie (2) montre que $\Ker(M \to N)$ et
$\Coker(M \to N)$ sont des objets de $\mathcal{C}$. Enfin, supposons donnée une
suite exacte courte $0 \to M_1 \to M_2 \to M_3 \to 0$ de $A$-modules, avec
$M_1$ et $M_3$ complets au sens dérivé. La suite exacte longue de $\Ext$
montre que $M_2$ est complet au sens dérivé. Par conséquent, $\mathcal{C}$ est
une sous-catégorie de Serre faible d'après Homologie, Lemme
\ref{homology-lemma-characterize-weak-serre-subcategory}.
\end{proof}

\noindent
Nous généraliserons le lemme suivant dans le
Lemme \ref{lemma-derived-complete-zero-bis}.

\begin{lemma}
\label{lemma-derived-complete-zero}
Soit $I$ un idéal de type fini d'un anneau $A$. Soit $M$ un $A$-module complet
au sens dérivé. Si $M/IM = 0$, alors $M = 0$.
\end{lemma}

\begin{proof}
Supposons $M/IM$ nul et écrivons $I = (f_1, \ldots, f_r)$. Soit $i < r$ le
plus grand entier tel que $N = M/(f_1, \ldots, f_i)M$ soit non nul. Si un tel
$i$ n'existe pas, alors $M = 0$, ce qui est précisément l'assertion voulue.
Le module $N$ est complet au sens dérivé comme conoyau d'un morphisme entre
modules complets au sens dérivé ; voir Lemme
\ref{lemma-serre-subcategory}. Par choix de $i$, l'application
$f_{i + 1} : N \to N$ est surjective. Par suite,
$$
\lim (\ldots \to N \xrightarrow{f_{i + 1}} N \xrightarrow{f_{i + 1}} N)
$$
est non nulle, ce qui contredit la complétude dérivée de $N$.
\end{proof}

\noindent
Lorsque l'anneau est complet pour la topologie $I$-adique, on obtient une
grande réserve de complexes complets au sens dérivé.

\begin{lemma}
\label{lemma-pseudo-coherent-is-derived-complete}
Soient $A$ un anneau et $I \subset A$ un idéal. Si $A$ est complet au sens
dérivé (par exemple s'il est complet pour la topologie $I$-adique), alors tout
objet pseudo-cohérent de $D(A)$ est complet au sens dérivé.
\end{lemma}

\begin{proof}
(Le Lemme \ref{lemma-complete-derived-complete} explique l'assertion entre
parenthèses.) Soit $K$ un objet pseudo-cohérent de $D(A)$. Par définition,
$K$ est représenté par un complexe $K^\bullet$ borné supérieurement de
$A$-modules libres de rang fini. Puisque $A$ est complet au sens dérivé, la
partie (1) du Lemme \ref{lemma-serre-subcategory} montre que $H^n(K)$ est
complet au sens dérivé pour tout $n$. La partie (2) du même lemme implique à
son tour que $K$ est complet au sens dérivé.
\end{proof}

\begin{lemma}
\label{lemma-double-localize}
Soient $A$ un anneau et $f, g \in A$. Pour $K \in D(A)$, on a
$R\Hom_A(A_f, R\Hom_A(A_g, K)) = R\Hom_A(A_{fg}, K)$.
\end{lemma}

\begin{proof}
Cela résulte du Lemme \ref{lemma-internal-hom}.
\end{proof}

\begin{lemma}
\label{lemma-derived-completion}
\begin{slogan}
Les complétions dérivées le long d'idéaux de type fini existent et se
calculent par un procédé de {\v C}ech.
\end{slogan}
Soit $I$ un idéal de type fini d'un anneau $A$. Le foncteur d'inclusion
$D_{comp}(A, I) \to D(A)$ admet un adjoint à gauche. Autrement dit, pour tout
objet $K$ de $D(A)$, il existe un morphisme $K \to K^\wedge$ de $K$ vers un objet
complet au sens dérivé de $D(A)$ tel que l'application
$$
\Hom_{D(A)}(K^\wedge, E) \longrightarrow \Hom_{D(A)}(K, E)
$$
soit bijective pour tout objet $E$ de $D(A)$ complet au sens dérivé. Plus
précisément, si $I$ est engendré par $f_1, \ldots, f_r \in A$, alors
$$
K^\wedge = R\Hom\left((A \to \prod\nolimits_{i_0} A_{f_{i_0}} \to
\prod\nolimits_{i_0 < i_1} A_{f_{i_0}f_{i_1}}
\to \ldots \to A_{f_1\ldots f_r}), K\right)
$$
fonctoriellement en $K$.
\end{lemma}

\begin{proof}
Définissons $K^\wedge$ par la dernière formule affichée du lemme. Il existe un
morphisme de complexes
$$
(A \to \prod\nolimits_{i_0} A_{f_{i_0}} \to
\prod\nolimits_{i_0 < i_1} A_{f_{i_0}f_{i_1}} \to
\ldots \to A_{f_1\ldots f_r}) \longrightarrow A
$$
qui induit un morphisme $K \to K^\wedge$. Il suffit de montrer que $K^\wedge$
est complet au sens dérivé et que $K \to K^\wedge$ est un isomorphisme si $K$
est complet au sens dérivé\footnote{En effet, si $E$ est complet au sens dérivé
et si $a : K \to E$ est un morphisme, le diagramme commutatif
$$
\xymatrix{
K \ar[r]_a \ar[d] & E \ar[d]^{\cong} \\
K^\wedge \ar[r]^{a^\wedge} & E^\wedge
}
$$
montre que tout morphisme vers $E$ se factorise par $K^\wedge$. Choisissons un
triangle distingué $K \to K^\wedge \to C$ et appliquons ${}^\wedge$. On trouve
$C^\wedge = 0$. D'après ce qui précède, cela signifie que $\Hom(C, E) = 0$
pour tout $E$ complet au sens dérivé, ce qui démontre
$\Hom(K, E) = \Hom(K^\wedge, E)$, comme voulu.}.

\medskip\noindent
Soit $f \in A$. D'après le Lemme \ref{lemma-double-localize}, l'objet
$R\Hom_A(A_f, K^\wedge)$ est égal à
$$
R\Hom\left((A_f \to \prod\nolimits_{i_0} A_{ff_{i_0}} \to
\prod\nolimits_{i_0 < i_1} A_{ff_{i_0}f_{i_1}} \to
\ldots \to A_{ff_1\ldots f_r}), K\right)
$$
Si $f \in I$, alors $f_1, \ldots, f_r$ engendrent l'idéal unité dans $A_f$ ;
le complexe de {\v C}ech alterné augmenté
$$
A_f \to \prod\nolimits_{i_0} A_{ff_{i_0}} \to
\prod\nolimits_{i_0 < i_1} A_{ff_{i_0}f_{i_1}} \to
\ldots \to A_{ff_1\ldots f_r}
$$
est donc nul dans $D(A)$ d'après le
Lemme \ref{lemma-extended-alternating-torsion}. (En fait, si $f = f_i$ pour un
certain $i$, ce complexe est homotope à zéro d'après le
Lemme \ref{lemma-extended-alternating-homotopy-zero} ; c'est le seul cas dont
nous avons besoin.) Ainsi, $R\Hom_A(A_f, K^\wedge) = 0$, et le
Lemme \ref{lemma-hom-from-Af} montre que $K^\wedge$ est complet au sens dérivé.

\medskip\noindent
Réciproquement, si $K$ est complet au sens dérivé, alors
$R\Hom_A(A_f, K)$ est nul pour tout
$f = f_{i_0} \ldots f_{i_p}$, $p \geq 0$. Ainsi,
$K \to K^\wedge$ est un isomorphisme dans $D(A)$.
\end{proof}

\begin{remark}
\label{remark-derived-completion}
Soient $A$ un anneau et $I \subset A$ un idéal de type fini. L'adjoint à
gauche du foncteur d'inclusion $D_{comp}(A, I) \to D(A)$ dont l'existence est
établie par le Lemme \ref{lemma-derived-completion} est appelé la
{\it complétion dérivée}. Pour l'indiquer, nous dirons « soit $K^\wedge$ la
complétion dérivée de $K$ ». Gardons à l'esprit que l'unité de l'adjonction est
un morphisme fonctoriel $K \to K^\wedge$.
\end{remark}

\begin{lemma}
\label{lemma-derived-completion-vanishes}
Soient $A$ un anneau et $I \subset A$ un idéal de type fini. Soit $K^\bullet$
un complexe de $A$-modules tel que $f : K^\bullet \to K^\bullet$ soit un
isomorphisme pour un certain $f \in I$, c'est-à-dire que $K^\bullet$ soit un
complexe de $A_f$-modules. Alors la complétion dérivée de $K^\bullet$ est
nulle.
\end{lemma}

\begin{proof}
En effet, dans ce cas, $R\Hom_A(K, L)$ est nul pour tout complexe $L$ complet
au sens dérivé ; voir Lemme \ref{lemma-hom-from-Af}. La propriété universelle
du Lemme \ref{lemma-derived-completion} implique donc que $K^\wedge$ est nul.
\end{proof}

\begin{lemma}
\label{lemma-completion-RHom}
Soient $A$ un anneau, $I \subset A$ un idéal de type fini et
$K, L \in D(A)$. Alors
$$
R\Hom_A(K, L)^\wedge = R\Hom_A(K, L^\wedge) = R\Hom_A(K^\wedge, L^\wedge)
$$
\end{lemma}

\begin{proof}
D'après le Lemme \ref{lemma-derived-completion}, la complétion dérivée est
donnée par $R\Hom_A(C, -)$ pour un certain $C \in D(A)$. Ainsi,
\begin{align*}
R\Hom_A(C, R\Hom_A(K, L))
& =
R\Hom_A(C \otimes_A^\mathbf{L} K, L) \\
& =
R\Hom_A(K, R\Hom_A(C, L))
\end{align*}
d'après le Lemme \ref{lemma-internal-hom}. Cela démontre la première égalité.
Le morphisme $K \to K^\wedge$ induit un morphisme
$$
R\Hom_A(K^\wedge, L^\wedge) \to R\Hom_A(K, L^\wedge)
$$
qui est un isomorphisme dans $D(A)$ par définition de la complétion dérivée
comme adjointe à gauche du foncteur d'inclusion.
\end{proof}

\begin{lemma}
\label{lemma-naive-derived-completion}
Soient $A$ un anneau, $I \subset A$ un idéal et $(K_n)$ un système projectif
d'objets de $D(A)$ tel que, pour tout $f \in I$ et tout $n$, il existe
$e = e(n, f)$ tel que $f^e$ agisse par zéro sur $K_n$. Alors, pour
$K \in D(A)$, l'objet $K' = R\lim (K \otimes_A^\mathbf{L} K_n)$ est complet au
sens dérivé par rapport à $I$.
\end{lemma}

\begin{proof}
Puisque la catégorie des objets complets au sens dérivé est stable par
$R\lim$, il suffit de montrer que chaque
$K \otimes_A^\mathbf{L} K_n$ est complet au sens dérivé. Par hypothèse, pour
tout $f \in I$, il existe $e$ tel que $f^e$ agisse par zéro sur
$K \otimes_A^\mathbf{L} K_n$. Cela implique évidemment
$T(K \otimes_A^\mathbf{L} K_n, f) = 0$, ce qui conclut.
\end{proof}

\begin{situation}
\label{situation-koszul}
Soient $A$ un anneau et $I = (f_1, \ldots, f_r) \subset A$. Soit
$K_n^\bullet = K_\bullet(A, f_1^n, \ldots, f_r^n)$ le complexe de Koszul sur
$f_1^n, \ldots, f_r^n$, considéré comme complexe de cochaînes en degrés
$-r, -r + 1, \ldots, 0$. En utilisant la fonctorialité du
Lemme \ref{lemma-functorial}, on obtient un système projectif
$$
\ldots \to K_3^\bullet \to K_2^\bullet \to K_1^\bullet
$$
compatible avec le système projectif
$H^0(K_n^\bullet) = A/(f_1^n, \ldots, f_r^n)$ et avec les morphismes
$A \to K_n^\bullet$.
\end{situation}

\noindent
Un point essentiel de la discussion ci-dessous est que, pour $m > n$, le
morphisme
$$
K_m^{-p} = \wedge^p(A^{\oplus r}) \to \wedge^p(A^{\oplus r}) = K_n^{-p}
$$
est donné par multiplication par
$f_{i_1}^{m - n} \ldots f_{i_p}^{m - n}$ sur l'élément de base
$e_{i_1} \wedge \ldots \wedge e_{i_p}$.

\begin{lemma}
\label{lemma-koszul-derived-completion-complete}
Dans la Situation \ref{situation-koszul}, pour $K \in D(A)$, l'objet
$K' = R\lim (K \otimes_A^\mathbf{L} K_n^\bullet)$ est complet au sens dérivé
par rapport à $I$.
\end{lemma}

\begin{proof}
C'est un cas particulier du Lemme \ref{lemma-naive-derived-completion}, car
$f_i^n$ agit sur $K_n^\bullet$ par un endomorphisme homotope à zéro d'après le
Lemme \ref{lemma-homotopy-koszul}.
\end{proof}

\begin{lemma}
\label{lemma-characterize-derived-complete-Koszul}
Dans la Situation \ref{situation-koszul}, soit $K \in D(A)$. Les conditions
suivantes sont équivalentes :
\begin{enumerate}
\item $K$ est complet au sens dérivé par rapport à $I$ ;
\item le morphisme canonique
$K \to R\lim (K \otimes_A^\mathbf{L} K_n^\bullet)$ est un isomorphisme dans
$D(A)$.
\end{enumerate}
\end{lemma}

\begin{proof}
Si (2) est satisfaite, alors $K$ est complet au sens dérivé par rapport à $I$
d'après le Lemme \ref{lemma-koszul-derived-completion-complete}.
Réciproquement, supposons $K$ complet au sens dérivé par rapport à $I$.
Considérons les filtrations
$$
K_n^\bullet \supset
\sigma_{\geq -r + 1}K_n^\bullet \supset
\sigma_{\geq -r + 2}K_n^\bullet \supset \ldots \supset
\sigma_{\geq -1}K_n^\bullet \supset
\sigma_{\geq 0}K_n^\bullet = A
$$
par troncations bêtes (Homologie, Section
\ref{homology-section-truncations}). Comme la construction
$R\lim(K \otimes E)$ est exacte en la seconde variable
(Lemme \ref{lemma-tensor-Rlim-exact}), il suffit de montrer que
$$
R\lim \left(
K \otimes_A^\mathbf{L}
(\sigma_{\geq p}K_n^\bullet/ \sigma_{\geq p + 1}K_n^\bullet)
\right) = 0
$$
pour $p < 0$. La description explicite des complexes de Koszul ci-dessus donne
$$
R\lim \left(
K \otimes_A^\mathbf{L}
(\sigma_{\geq p}K_n^\bullet/ \sigma_{\geq p + 1}K_n^\bullet)
\right) =
\bigoplus\nolimits_{i_1, \ldots, i_{-p}}
T(K, f_{i_1}\ldots f_{i_{-p}})
$$
qui est nul pour $p < 0$ par hypothèse sur $K$.
\end{proof}

\begin{lemma}
\label{lemma-derived-completion-koszul}
Dans la Situation \ref{situation-koszul}, le foncteur qui envoie
$K \in D(A)$ sur la limite dérivée
$K' = R\lim( K \otimes_A^\mathbf{L} K_n^\bullet )$ est l'adjoint à gauche du
foncteur d'inclusion $D_{comp}(A) \to D(A)$ construit dans le
Lemme \ref{lemma-derived-completion}.
\end{lemma}

\begin{proof}[Première preuve]
La règle $K \leadsto K'$ est un foncteur, et $K'$ est complet au sens dérivé
par rapport à $I$ d'après le
Lemme \ref{lemma-koszul-derived-completion-complete}. Un argument formel
(omis) montre qu'il suffit d'établir que $K \to K'$ est un isomorphisme lorsque
$K$ est complet au sens dérivé par rapport à $I$. C'est le contenu du
Lemme \ref{lemma-characterize-derived-complete-Koszul}.
\end{proof}

\begin{proof}[Seconde preuve]
Notons $K \mapsto K^\wedge$ l'adjoint construit dans le
Lemme \ref{lemma-derived-completion}. D'après ce lemme,
$$
K^\wedge = R\Hom\left((A \to \prod\nolimits_{i_0} A_{f_{i_0}} \to
\prod\nolimits_{i_0 < i_1} A_{f_{i_0}f_{i_1}}
\to \ldots \to A_{f_1\ldots f_r}), K\right)
$$
Dans le Lemme \ref{lemma-extended-alternating-Cech-is-colimit-koszul}, nous
avons vu que le complexe de {\v C}ech alterné augmenté
$$
A \to \prod\nolimits_{i_0} A_{f_{i_0}} \to
\prod\nolimits_{i_0 < i_1} A_{f_{i_0}f_{i_1}}
\to \ldots \to A_{f_1\ldots f_r}
$$
est une colimite des complexes de Koszul
$K^n = K(A, f_1^n, \ldots, f_r^n)$ placés en degrés $0, \ldots, r$.
Remarquons que $K^n$ est un complexe de chaînes fini de $A$-modules libres de
rang fini, dont le dual (au sens du Lemme
\ref{lemma-dual-perfect-complex}) est
$R\Hom_A(K^n, A) = K_n$, où $K_n$ est le complexe de cochaînes de Koszul placé
en degrés $-r, \ldots, 0$, comme d'habitude. Il suffit donc de montrer que
$$
R\Hom_A(\text{hocolim} K^n, K) =  R\lim (K \otimes_A^\mathbf{L} K_n)
$$
Cela résulte du Lemme \ref{lemma-colim-and-lim-of-duals}.
\end{proof}

\begin{lemma}
\label{lemma-derived-complete-bounded}
Soit $I = (f_1, \ldots, f_r)$ un idéal de type fini d'un anneau $A$. Soit $K$
un objet de $D(A)$ complet au sens dérivé. Les conditions suivantes sont
équivalentes :
\begin{enumerate}
\item $H^i(K) = 0$ pour $i > 0$ ;
\item $H^i(K \otimes_A^\mathbf{L} A/I) = 0$ pour $i > 0$ ;
\item $H^i(K \otimes_A^\mathbf{L} K_1^\bullet) = 0$ pour $i > 0$, où
$K_1^\bullet$ est celui de la Situation \ref{situation-koszul}.
\end{enumerate}
\end{lemma}

\begin{proof}
L'implication (1) $\Rightarrow$ (2) est toujours vraie. L'implication
(2) $\Rightarrow$ (3) résulte du Lemme \ref{lemma-bounded-above-mod-I}.
Supposons (3). Pour $s = 0, \ldots, r$, considérons le complexe
$$
K(s) =  K \otimes_A^\mathbf{L} K_\bullet(A, f_1, \ldots, f_s)
$$
où $K_\bullet(A, f_1, \ldots, f_s)$ est le complexe de Koszul placé en degrés
cohomologiques $-s, \ldots, 0$. D'après le
Lemme \ref{lemma-cone-koszul}, on a des triangles distingués
$$
K(s - 1) \xrightarrow{f_s} K(s - 1) \to K(s) \to K(s - 1)[1]
$$
Comme $K(0) = K$ est complet au sens dérivé, une récurrence ascendante montre
que chaque $K(s)$ est un objet complet au sens dérivé de $D(A)$. Montrons par
récurrence descendante que $K(s)$ n'a pas de cohomologie non nulle en degrés
$> 0$. Cela vaut pour $s = r$ par l'hypothèse (3). Si cela vaut
pour $K(s)$ avec $s > 0$, alors
$f_s : H^i(K(s - 1)) \to H^i(K(s - 1))$ est surjective pour $i > 0$, et les
Lemmes \ref{lemma-serre-subcategory} et \ref{lemma-derived-complete-zero}
montrent que $H^i(K(s - 1))$ est nul.
\end{proof}

\begin{lemma}
\label{lemma-derived-complete-zero-bis}
\begin{slogan}
Nakayama dérivé
\end{slogan}
\begin{reference}
Un résultat apparenté est \cite[Proposition 6.5]{Dwyer-Greenlees}. Le lemme de
Nakayama dérivé se trouve par exemple dans la troisième leçon de Bhatt sur la
cohomologie prismatique à l'université Columbia, à l'automne 2018, Section 2,
propriété (2). Leonid Positselski a proposé une preuve dans
\url{https://mathoverflow.net/a/331501}. Nous suivons toutefois la preuve
suggérée par Anonymous dans les commentaires.
\end{reference}
Soit $I$ un idéal de type fini d'un anneau $A$. Soit $K$ un objet de $D(A)$
complet au sens dérivé. Si $K \otimes_A^\mathbf{L} A/I = 0$, alors $K = 0$.
\end{lemma}

\begin{proof}
C'est un cas particulier du Lemme \ref{lemma-derived-complete-bounded}, mais
nous donnons aussi une preuve directe. Choisissons des générateurs
$f_1, \ldots, f_r$ de $I$. Notons $K_n$ le complexe de Koszul sur
$f_1^n, \ldots, f_r^n$ au-dessus de $A$. Rappelons que $K_n$ est borné et que
les modules de cohomologie de $K_n$ sont annulés par
$f_1^n, \ldots, f_r^n$, donc par
$I^{nr}$. Le Lemme \ref{lemma-derived-vanishing-mod-I} montre que
$K \otimes_A^\mathbf{L} K_n = 0$. Comme $K$ est complet au sens dérivé, le
Lemme \ref{lemma-derived-completion-koszul} donne
$K = R\lim K \otimes_A^\mathbf{L} K_n = 0$, comme voulu.
\end{proof}

\noindent
La relation avec le complexe de Koszul montre notamment que la complétion
dérivée est de dimension cohomologique finie.

\begin{lemma}
\label{lemma-derived-completion-finite-cohomological-dimension}
Soient $A$ un anneau et $I \subset A$ un idéal qui peut être engendré par $r$
éléments. Alors la complétion dérivée est de dimension cohomologique finie :
\begin{enumerate}
\item Soit $K \to L$ un morphisme dans $D(A)$ tel que
$H^i(K) \to H^i(L)$ soit un isomorphisme pour $i \geq 1$ et surjectif pour
$i = 0$. Alors $H^i(K^\wedge) \to H^i(L^\wedge)$ est un isomorphisme pour
$i \geq 1$ et surjectif pour $i = 0$.
\item Soit $K \to L$ un morphisme de $D(A)$ tel que
$H^i(K) \to H^i(L)$ soit un isomorphisme pour $i \leq -1$ et injectif pour
$i = 0$. Alors $H^i(K^\wedge) \to H^i(L^\wedge)$ est un isomorphisme pour
$i \leq -r - 1$ et injectif pour $i = -r$.
\end{enumerate}
\end{lemma}

\begin{proof}
Supposons $I$ engendré par $f_1, \ldots, f_r$. Pour tout $K \in D(A)$, le
Lemme \ref{lemma-derived-completion-koszul} donne
$K^\wedge = R\lim K \otimes_A^\mathbf{L} K_n$, où $K_n$ est le complexe de
Koszul sur $f_1^n, \ldots, f_r^n$. On obtient donc une suite exacte courte
$$
0 \to R^1\lim H^{i - 1}(K \otimes_A^\mathbf{L} K_n)
\to H^i(K^\wedge) \to \lim H^i(K \otimes_A^\mathbf{L} K_n) \to 0
$$
d'après le Lemme \ref{lemma-break-long-exact-sequence-modules}.

\medskip\noindent
Démontrons (1). Choisissons un triangle distingué
$K \to L \to C \to K[1]$. Alors $H^i(C) = 0$ pour $i \geq 0$. Comme $K_n$ est
placé en degrés $\leq 0$, on a
$H^i(C \otimes_A^\mathbf{L} K_n) = 0$ pour $i \geq 0$, et
$H^{-1}(C \otimes_A^\mathbf{L} K_n) =
H^{-1}(C) \otimes_A A/(f_1^n, \ldots, f_r^n)$ est un système dont les
applications de transition sont surjectives. La suite exacte affichée ci-dessus
montre que $H^i(C^\wedge) = 0$ pour $i \geq 0$. En appliquant le triangle
distingué $K^\wedge \to L^\wedge \to C^\wedge \to K^\wedge[1]$, on obtient
(1).

\medskip\noindent
Démontrons (2). Choisissons un triangle distingué
$K \to L \to C \to K[1]$. Alors $H^i(C) = 0$ pour $i < 0$. Comme $K_n$ est
placé en degrés $-r, \ldots, 0$, on a
$H^i(C \otimes_A^\mathbf{L} K_n) = 0$ pour $i < -r$. La suite exacte affichée
ci-dessus montre que $H^i(C^\wedge) = 0$ pour $i < -r$. En appliquant le
triangle distingué $K^\wedge \to L^\wedge \to C^\wedge \to K^\wedge[1]$, on
obtient (2).
\end{proof}

\begin{lemma}
\label{lemma-derived-completion-spectral-sequence}
Soient $A$ un anneau, $I \subset A$ un idéal de type fini et $K^\bullet$ un
complexe filtré de $A$-modules. Il existe une suite spectrale canonique
$(E_r, \text{d}_r)_{r \geq 1}$ de $A$-modules bigradués complets au sens dérivé,
dont $d_r$ est de bidegré $(r, -r + 1)$ et telle que
$$
E_1^{p, q} = H^{p + q}((\text{gr}^pK^\bullet)^\wedge)
$$
Si la filtration de chaque $K^n$ est finie, alors la suite spectrale est bornée
et converge vers $H^*((K^\bullet)^\wedge)$.
\end{lemma}

\begin{proof}
D'après le Lemme \ref{lemma-derived-completion}, la complétion dérivée est
donnée par $R\Hom_A(C, -)$ pour un certain $C \in D^b(A)$. Les Lemmes
\ref{lemma-derived-completion-finite-cohomological-dimension} et
\ref{lemma-projective-amplitude} montrent que $C$ est de dimension projective
finie. On peut donc choisir un complexe borné de modules projectifs
$P^\bullet$ représentant $C$. Alors
$$
M^\bullet = \Hom^\bullet(P^\bullet, K^\bullet)
$$
est un complexe de $A$-modules représentant $(K^\bullet)^\wedge$. Il est muni
de la filtration définie par
$F^pM^\bullet = \Hom^\bullet(P^\bullet, F^pK^\bullet)$. On voit que
$F^pM^\bullet$ représente $(F^pK^\bullet)^\wedge$ et que
$\text{gr}^pM^\bullet$ représente $(\text{gr}K^\bullet)^\wedge$. La suite
spectrale cherchée est donc celle du complexe filtré $M^\bullet$ ; voir
Homologie, Section \ref{homology-section-filtered-complex}. Si la filtration de
chaque $K^n$ est finie, celle de chaque $M^n$ l'est aussi, car $P^\bullet$ est
un complexe borné. La dernière assertion résulte alors de Homologie, Lemme
\ref{homology-lemma-biregular-ss-converges}.
\end{proof}

\begin{example}
\label{example-derived-completion-spectral-sequence}
Soient $A$ un anneau, $I \subset A$ un idéal de type fini et $K^\bullet$ un
complexe de $A$-modules. On peut appliquer le
Lemme \ref{lemma-derived-completion-spectral-sequence} avec
$F^pK^\bullet = \tau_{\leq -p}K^\bullet$. On obtient alors une suite spectrale
bornée
$$
E_1^{p, q} = H^{p + q}(H^{-p}(K^\bullet)^\wedge[p]) =
H^{2p + q}(H^{-p}(K^\bullet)^\wedge)
$$
qui converge vers $H^{p + q}((K^\bullet)^\wedge)$. Après renumérotation
$p = -j$ et $q = i + 2j$, on voit que, pour tout $K \in D(A)$, il existe une
suite spectrale bornée $(E'_r, d'_r)_{r \geq 2}$ de modules bigradués complets
au sens dérivé, dont $d'_r$ est de bidegré $(r, -r + 1)$, telle que
$$
(E'_2)^{i, j} = H^i(H^j(K)^\wedge)
$$
et qui converge vers $H^{i + j}(K^\wedge)$.
\end{example}

\begin{lemma}
\label{lemma-restriction-derived-complete}
Soient $A \to B$ un homomorphisme d'anneaux et $I \subset A$ un idéal.
L'image réciproque de $D_{comp}(A, I)$ par le foncteur de restriction
$D(B) \to D(A)$ est $D_{comp}(B, IB)$.
\end{lemma}

\begin{proof}
Le Lemme \ref{lemma-ideal-of-elements-complete-wrt} montre que
$L \in D(B)$ appartient à $D_{comp}(B, IB)$ si et seulement si $T(L, f)$ est
nul pour tout $f \in I$. La cohomologie de $T(L, f)$ se calcule dans la
catégorie des groupes abéliens ; peu importe donc que l'on considère $f$ comme
un élément de $A$ ou que l'on prenne l'image de $f$ dans $B$. Le lemme résulte immédiatement de
ceci et de la définition des objets complets au sens dérivé.
\end{proof}

\begin{lemma}
\label{lemma-restriction-derived-complete-equivalence}
Soient $A \to B$ un homomorphisme d'anneaux et $I \subset A$ un idéal de type
fini. Si $A \to B$ est plat et $A/I \cong B/IB$, alors le foncteur de
restriction $D(B) \to D(A)$ induit une équivalence
$D_{comp}(B, IB) \to D_{comp}(A, I)$.
\end{lemma}

\begin{proof}
Choisissons des générateurs $f_1, \ldots, f_r$ de $I$. Notons
$\check{\mathcal{C}}^\bullet_A \to \check{\mathcal{C}}^\bullet_B$ le
quasi-isomorphisme de complexes de {\v C}ech alternés augmentés du
Lemme \ref{lemma-map-identifies-koszul-and-cech-complexes}. Soit
$K \in D_{comp}(A, I)$. Soit $I^\bullet$ un complexe K-injectif de
$A$-modules représentant $K$. Comme $\Ext^n_A(A_f, K)$ et
$\Ext^n_A(B_f, K)$ sont nuls pour tous $f \in I$ et
$n \in \mathbf{Z}$ (Lemme \ref{lemma-hom-from-Af}), on en conclut que
$\check{\mathcal{C}}^\bullet_A \to A$ et
$\check{\mathcal{C}}^\bullet_B \to B$ induisent des quasi-isomorphismes
$$
I^\bullet = \Hom_A(A, I^\bullet) \longrightarrow
\text{Tot}(\Hom_A(\check{\mathcal{C}}^\bullet_A, I^\bullet))
$$
et
$$
\Hom_A(B, I^\bullet) \longrightarrow
\text{Tot}(\Hom_A(\check{\mathcal{C}}^\bullet_B, I^\bullet))
$$
Certains détails sont omis. Puisque
$\check{\mathcal{C}}^\bullet_A \to \check{\mathcal{C}}^\bullet_B$ est un
quasi-isomorphisme et que $I^\bullet$ est K-injectif, on conclut que
$\Hom_A(B, I^\bullet) \to I^\bullet$ est un quasi-isomorphisme. Comme le
complexe $\Hom_A(B, I^\bullet)$ est un complexe de $B$-modules, $K$ appartient
à l'image du foncteur de restriction ; autrement dit, ce foncteur est
essentiellement surjectif.

\medskip\noindent
En fait, l'argument montre que
$F : D_{comp}(A, I) \to D_{comp}(B, IB)$,
$K \mapsto \Hom_A(B, I^\bullet)$ est un inverse à gauche du foncteur de
restriction. Enfin, soit $L \in D_{comp}(B, IB)$. Représentons $L$ par un
complexe K-injectif $J^\bullet$ de $B$-modules. Le complexe $J^\bullet$ est
également K-injectif comme complexe de $A$-modules
(Lemme \ref{lemma-K-injective-flat}) ; par conséquent,
$F(\text{restriction de }L) = \Hom_A(B, J^\bullet)$. Il existe un morphisme
$J^\bullet \to \Hom_A(B, J^\bullet)$ de complexes de $B$-modules dont la composée
avec $\Hom_A(B, J^\bullet) \to J^\bullet$ est l'identité. On en conclut que
$F$ est aussi un inverse à droite du foncteur de restriction, ce qui achève la
preuve.
\end{proof}






\section{La catégorie des modules complets au sens dérivé}
\label{section-derived-complete-modules}

\noindent
Soient $A$ un anneau et $I$ un idéal. Notons $\mathcal{C}$ la catégorie des
modules complets au sens dérivé ; voir Définition
\ref{definition-derived-complete}. Dans cette section, nous étudions quelques
propriétés de cette catégorie. Dans Exemples, Section
\ref{examples-section-derived-complete-modules}, nous montrons que
$\mathcal{C}$ n'est pas en général une catégorie abélienne de Grothendieck.

\medskip\noindent
D'après le Lemme \ref{lemma-serre-subcategory}, la catégorie $\mathcal{C}$ est
abélienne et le foncteur d'inclusion
$\mathcal{C} \to \text{Mod}_A$ est exact.

\medskip\noindent
Comme $D_{comp}(A) \subset D(A)$ est stable par produits (voir la discussion
qui suit la Définition \ref{definition-derived-complete}) et que les produits
dans $D(A)$ se calculent au niveau des complexes, on voit que $\mathcal{C}$
admet des produits qui coïncident avec ceux de $\text{Mod}_A$. Par conséquent,
$\mathcal{C}$ admet en fait toutes les limites, et le foncteur d'inclusion
$\mathcal{C} \to \text{Mod}_A$ commute à celles-ci ; voir Catégories, Lemme
\ref{categories-lemma-limits-products-equalizers}.

\medskip\noindent
Supposons $I$ de type fini. Notons ${}^\wedge : D(A) \to D(A)$ le foncteur de
complétion dérivée du Lemme \ref{lemma-derived-completion}. Montrons que le
foncteur
$$
\text{Mod}_A \longrightarrow \mathcal{C},\quad
M \longmapsto H^0(M^\wedge)
$$
est adjoint à gauche du foncteur d'inclusion
$\mathcal{C} \to \text{Mod}_A$. Remarquons que $H^i(M^\wedge) = 0$ pour
$i > 0$, par exemple d'après le
Lemme \ref{lemma-derived-completion-finite-cohomological-dimension}. Ainsi, si
$N$ est un $A$-module complet au sens dérivé, on a
\begin{align*}
\Hom_\mathcal{C}(H^0(M^\wedge), N)
& =
\Hom_{D_{comp}(A)}(M^\wedge, N)\\
& =
\Hom_{D(A)}(M, N) \\
& =
\Hom_A(M, N)
\end{align*}
comme voulu.

\medskip\noindent
Soit $T$ un ensemble préordonné et soit $t \mapsto M_t$ un système de
$A$-modules complets au sens dérivé, c'est-à-dire un système indexé par $T$
dans $\mathcal{C}$ ; voir Catégories, Section
\ref{categories-section-posets-limits}. Notons $\colim_{t \in T} M_t$ la
colimite de ce système dans $\text{Mod}_A$. Il résulte formellement de ce qui
précède que
$$
H^0((\colim_{t \in T} M_t)^\wedge)
$$
est la colimite du système dans $\mathcal{C}$. On voit ainsi que
$\mathcal{C}$ admet toutes les colimites. En général, le foncteur d'inclusion
$\mathcal{C} \to \text{Mod}_A$ ne commute pas aux colimites ; voir Exemples,
Section \ref{examples-section-derived-complete-modules}.

\begin{lemma}
\label{lemma-derived-complete-modules}
Soient $A$ un anneau et $I \subset A$ un idéal. La catégorie $\mathcal{C}$ des
modules complets au sens dérivé est abélienne et admet toutes les limites ; le
foncteur d'inclusion $F : \mathcal{C} \to \text{Mod}_A$ est exact et commute
aux limites. Si $I$ est de type fini, alors $\mathcal{C}$ admet toutes les
colimites et $F$ possède un adjoint à gauche.
\end{lemma}

\begin{proof}
Cela résume la discussion qui précède.
\end{proof}






\section{Complétion dérivée pour un idéal principal}
\label{section-derived-completion-principal}

\noindent
Dans cette section, nous étudions la complétion dérivée lorsque l'idéal est
engendré par un seul élément.

\begin{lemma}
\label{lemma-lift-universally}
Soit $A$ un anneau et soit $f \in A$. S'il existe un entier $c \geq 1$
tel que $A[f^c] = A[f^{c + 1}] = A[f^{c + 2}] = \ldots$ (par exemple si
$A$ est noethérien), alors, pour tout $n \geq 1$, il existe des morphismes
$$
(A \xrightarrow{f^n} A) \longrightarrow A/(f^n),
\quad\text{and}\quad
A/(f^{n + c}) \longrightarrow (A \xrightarrow{f^n} A)
$$
dans $D(A)$ qui induisent un isomorphisme entre les pro-objets
$\{A/(f^n)\}$ et $\{(f^n : A \to A)\}$ dans $D(A)$.
\end{lemma}

\begin{proof}
La première flèche affichée est évidente. On peut définir la seconde flèche
du lemme par le diagramme
$$
\xymatrix{
A/A[f^c] \ar[r]_-{f^{n + c}} \ar[d]_{f^c} & A \ar[d]^1 \\
A \ar[r]^{f^n} & A
}
$$
Puisque la flèche horizontale supérieure est injective, le complexe de la
ligne supérieure est quasi isomorphe à $A/f^{n + c}A$. Nous omettons le
calcul des composées nécessaire pour établir l'assertion sur les pro-objets.
\end{proof}

\begin{lemma}
\label{lemma-when-does-it-work}
Soit $A$ un anneau et soit $f \in A$. Posons $I = (f)$. Dans cette situation,
nous disposons de la complétion dérivée naïve
$K \mapsto K' = R\lim (K \otimes_A^\mathbf{L} A/f^nA)$ et de la complétion
dérivée
$$
K \mapsto K^\wedge = R\lim (K \otimes_A^\mathbf{L} (A \xrightarrow{f^n} A))
$$
du Lemme \ref{lemma-derived-completion-koszul}. La transformation naturelle
de foncteurs $K^\wedge \to K'$ est un isomorphisme si et seulement si la
$f$-torsion de $A$ est bornée.
\end{lemma}

\begin{proof}
Si la torsion par les puissances de $f$ est bornée, alors les pro-objets
$\{(f^n : A \to A)\}$ et $\{A/f^nA\}$ sont isomorphes d'après le
Lemme \ref{lemma-lift-universally}. Les foncteurs sont donc isomorphes d'après
le Lemme \ref{lemma-Rlim-pro-equal}. Réciproquement, le
Lemme \ref{lemma-tensor-Rlim-exact} montre que la condition équivaut exactement
à l'annulation de
$$
R\lim (K \otimes_A^\mathbf{L} A[f^n])
$$
pour tout $K \in D(A)$. Ici, les morphismes du système $(A[f^n])$ sont donnés
par la multiplication par $f$. En prenant $K = A$, puis
$K = \bigoplus_{i \in \mathbf{N}} A$, le
Lemme \ref{lemma-Rlim-zero-of-direct-sums} montre que cela entraîne que
$(A[f^n])$ est nul comme pro-objet, c'est-à-dire que
$f^{n - 1}A[f^n] = 0$ pour un certain $n$, ou encore que
$A[f^{n - 1}] = A[f^n]$ ; la torsion par les puissances de $f$ est donc
bornée.
\end{proof}

\begin{example}
\label{example-derived-complete-modules}
Soit $A$ un anneau et soit $f \in A$ un non-diviseur de zéro. L'exemple à
garder à l'esprit est $A = \mathbf{Z}_p$ et $f = p$. Soit $M$ un
$A$-module. Affirmation : $M$ est complet au sens dérivé par rapport à $f$
si et seulement s'il existe une suite exacte courte
$$
0 \to K \to L \to M \to 0
$$
où $K, L$ sont des modules complets pour la topologie $f$-adique et sans
$f$-torsion. En effet, s'il existe une telle suite exacte courte, alors
$$
M \otimes_A^\mathbf{L} (A \xrightarrow{f^n} A) = (K/f^nK \to L/f^nL)
$$
car $f$ est un non-diviseur de zéro sur $K$ et $L$. On en conclut que
$R\lim (M \otimes_A^\mathbf{L} (A \xrightarrow{f^n} A))$ est quasi isomorphe
à $K \to L$, c'est-à-dire à $M$. Le Lemme
\ref{lemma-characterize-derived-complete-Koszul} montre alors que $M$ est
complet au sens dérivé. Réciproquement, supposons $M$ complet au sens dérivé.
Choisissons une surjection $F \to M$, où $F$ est un $A$-module libre. Comme
$f$ est un non-diviseur de zéro sur $F$, la complétion dérivée de $F$ est
$L = \lim F/f^nF$. Remarquons que $L$ est sans $f$-torsion : si $(x_n)$,
avec $x_n \in F$, représente un élément $\xi$ de $L$ et si $f\xi = 0$, alors
$x_n = x_{n + 1} + f^nz_n$ et $fx_n = f^ny_n$ pour certains
$z_n, y_n \in F$. On a alors $f^n y_n = fx_n =
fx_{n + 1} + f^{n + 1}z_n = f^{n + 1}y_{n + 1} + f^{n + 1}z_n$ ; puisque
$f$ est un non-diviseur de zéro sur $F$, il vient que $y_n \in fF$, d'où
$x_n \in f^nF$, c'est-à-dire $\xi = 0$. Comme $L$ est la complétion dérivée,
la propriété universelle fournit un morphisme $L \to M$ par lequel se
factorise $F \to M$. Soit $K = \Ker(L \to M)$ son noyau. Le module $K$ est
encore sans $f$-torsion ; la complétion dérivée de $K$ est donc
$\lim K/f^nK$.
D'autre part, $M$ et $L$ sont tous deux complets au sens dérivé, donc $K$
l'est aussi d'après le Lemme \ref{lemma-serre-subcategory}. Il s'ensuit que
$K = \lim K/f^nK$, ce qui prouve l'affirmation.
\end{example}

\begin{example}
\label{example-derived-complete-not-complete}
Soit $p$ un nombre premier. Considérons le morphisme
$\mathbf{Z}_p[x] \to \mathbf{Z}_p[y]$ d'algèbres de polynômes qui envoie $x$
sur $py$. Considérons le conoyau
$M = \Coker(\mathbf{Z}_p[x]^\wedge \to \mathbf{Z}_p[y]^\wedge)$ du morphisme
induit entre les complétions $p$-adiques ordinaires. Alors $M$ est un
$\mathbf{Z}_p$-module complet au sens dérivé d'après la
Proposition \ref{proposition-derived-complete-modules} et le
Lemme \ref{lemma-serre-subcategory} ; voir aussi la discussion de
l'Exemple \ref{example-derived-complete-modules}. Cependant, $M$ n'est pas
$p$-adiquement complet, car $1 + py + p^2 y^2 + \ldots$ s'envoie sur un
élément non nul de $M$ qui appartient à $\bigcap p^nM$.
\end{example}

\begin{example}
\label{example-spectral-sequence-principal}
Soit $A$ un anneau et soit $f \in A$. Notons $K \mapsto K^\wedge$ la
complétion dérivée par rapport à $(f)$. Soit $M$ un $A$-module. En utilisant
l'égalité
$$
M^\wedge = R\lim (M \xrightarrow{f^n} M)
$$
donnée par le Lemme \ref{lemma-derived-completion-koszul}, puis le
Lemme \ref{lemma-break-long-exact-sequence-modules}, on obtient
$$
H^{-1}(M^\wedge) = \lim M[f^n] = T_f(M)
$$
le {\it module de Tate $f$-adique de $M$}. Ici, les morphismes
$M[f^n] \to M[f^{n - 1}]$ sont donnés par la multiplication par $f$. On a
ensuite une suite exacte courte
$$
0 \to R^1\lim M[f^n] \to H^0(M^\wedge) \to \lim M/f^n M \to 0
$$
qui décrit $H^0(M^\wedge)$. On a $H^1(M^\wedge) = R^1\lim M/f^nM = 0$, car
les morphismes de transition sont surjectifs (Lemme
\ref{lemma-compute-Rlim-modules}). Toutes les autres cohomologies de
$M^\wedge$ sont nulles pour des raisons triviales. Enfin, pour $K \in D(A)$
et $p \in \mathbf{Z}$, on a une suite exacte courte
$$
0 \to H^0(H^p(K)^\wedge) \to H^p(K^\wedge) \to T_f(H^{p + 1}(K)) \to 0
$$
Cela résulte de la suite spectrale de l'Exemple
\ref{example-derived-completion-spectral-sequence}, qui dégénère en $E_2$
(puisque seuls $i = -1, 0$ donnent des termes non nuls) ; le lemme suivant
fournit davantage d'informations.
\end{example}

\begin{lemma}
\label{lemma-compare-completion-derived-completion}
Soit $A$ un anneau, soit $f \in A$ et soit $K$ un objet de $D(A)$. Notons
$K_n = K \otimes_A^\mathbf{L} (A \xrightarrow{f^n} A)$. Pour tout
$p \in \mathbf{Z}$, on a un diagramme commutatif
$$
\xymatrix{
& 0 & 0 \\
0 \ar[r] &
\widehat{H^p(K)} \ar[r] \ar[u] &
\lim H^p(K_n) \ar[r] \ar[u] &
T_f(H^{p + 1}(K)) \ar[r] &
0 \\
0 \ar[r] &
H^0(H^p(K)^\wedge) \ar[r] \ar[u] &
H^p(K^\wedge) \ar[r] \ar[u] &
T_f(H^{p + 1}(K)) \ar[r] \ar@{=}[u] &
0 \\
&
R^1\lim H^p(K)[f^n] \ar[u] \ar[r]^\cong &
R^1\lim H^{p - 1}(K_n) \ar[u] \\
& 0 \ar[u] & 0 \ar[u]
}
$$
à lignes et colonnes exactes, où
$\widehat{H^p(K)} = \lim H^p(K)/f^nH^p(K)$ est la complétion $f$-adique
usuelle. La suite exacte courte verticale de gauche et la suite exacte courte
horizontale médiane proviennent de l'Exemple
\ref{example-spectral-sequence-principal}. La suite exacte courte verticale
médiane est celle du Lemme \ref{lemma-break-long-exact-sequence-modules}.
\end{lemma}

\begin{proof}
Pour construire la suite exacte courte horizontale supérieure, remarquons
que nous avons les suites exactes courtes de systèmes projectifs suivantes :
$$
0 \to H^p(K)/f^nH^p(K) \to H^p(K_n) \to H^{p + 1}(K)[f^n] \to 0
$$
elles proviennent de la construction de $K_n$ comme décalé du cône de
$f^n : K \to K$. En prenant leur limite projective, on obtient la suite
exacte courte horizontale supérieure ; voir Homologie, Lemme
\ref{homology-lemma-Mittag-Leffler}.

\medskip\noindent
Montrons que l'on dispose d'un diagramme commutatif comme dans le lemme.
Considérons le morphisme $L = \tau_{\leq p}K \to K$. En posant
$L_n = L \otimes_A^\mathbf{L} (A \xrightarrow{f^n} A)$, on obtient un
morphisme $(L_n) \to (K_n)$ de systèmes projectifs, lequel induit un
morphisme de suites exactes courtes
$$
\xymatrix{
0 & 0 \\
\lim H^p(L_n) \ar[r] \ar[u] &
\lim H^p(K_n) \ar[u] \\
H^p(L^\wedge) \ar[r] \ar[u] &
H^p(K^\wedge) \ar[u] \\
R^1\lim H^{p - 1}(L_n) \ar[r] \ar[u] &
R^1\lim H^{p - 1}(K_n) \ar[u] \\
0 \ar[u] & 0 \ar[u]
}
$$
Puisque $H^i(L) = 0$ pour $i > p$ et $H^p(L) = H^p(K)$, un calcul fondé sur
les références de l'énoncé montre que
$H^p(L^\wedge) = H^0(H^p(K)^\wedge)$ et que
$H^p(L_n) = H^p(K)/f^nH^p(K)$. D'autre part, on a
$H^{p - 1}(L_n) = H^{p - 1}(K_n)$ ; on obtient donc l'isomorphisme indiqué
dans l'énoncé, puisque l'on sait déjà que le noyau de
$H^0(H^p(K)^\wedge) \to \widehat{H^p(K)}$ est égal à
$R^1\lim H^p(K)[f^n]$. Nous omettons de vérifier que le carré situé le plus à
droite dans le diagramme commute lorsque la ligne supérieure est définie par
la construction du premier paragraphe de la démonstration.
\end{proof}

\begin{remark}
\label{remark-ML-condition}
Avec les notations du Lemme \ref{lemma-compare-completion-derived-completion},
on voit aussi que le système projectif $H^p(K_n)$ vérifie la condition ML si
et seulement si le système projectif $H^{p + 1}(K)[f^n]$ la vérifie. Cela
résulte du système projectif de suites exactes courtes
$0 \to H^p(K)/f^nH^p(K) \to H^p(K_n) \to H^{p + 1}(K)[f^n] \to 0$
(voir la démonstration du lemme), combiné avec Homologie, Lemme
\ref{homology-lemma-Mittag-Leffler}, et le
Lemme \ref{lemma-Mittag-Leffler}.
\end{remark}

\begin{lemma}[Bhatt]
\label{lemma-torsion-and-derived-complete}
Soit $I$ un idéal de type fini d'un anneau $A$. Soit $M$ un $A$-module
complet au sens dérivé. Si $M$ est un module de torsion par les puissances de
$I$, alors $I^nM = 0$ pour un certain $n$.
\end{lemma}

\begin{proof}
Écrivons $I = (f_1, \ldots, f_r)$. Il suffit de montrer que, pour tout $i$,
il existe un $n_i$ tel que $f_i^{n_i}M = 0$. On peut donc supposer que
$I = (f)$ est un idéal principal. Soit $B = \mathbf{Z}[x] \to A$ le
morphisme d'anneaux qui envoie $x$ sur $f$. D'après le
Lemme \ref{lemma-restriction-derived-complete}, $M$ est complet au sens dérivé
comme $B$-module par rapport à l'idéal $(x)$. Après avoir remplacé $A$ par
$B$, on peut supposer que $f$ est un non-diviseur de zéro dans $A$.

\medskip\noindent
Supposons $I = (f)$, où $f \in A$ est un non-diviseur de zéro. D'après
l'Exemple \ref{example-derived-complete-modules}, il existe une suite exacte
courte
$$
0 \to K \xrightarrow{u} L \to M \to 0
$$
où $K$ et $L$ sont complets pour la topologie $I$-adique comme $A$-modules et
sans $f$-torsion\footnote{Pour la démonstration, il suffit d'établir
l'existence d'une suite $K \xrightarrow{u} L \to M \to 0$ dans laquelle $K$
et $L$ sont complets pour la topologie $I$-adique comme $A$-modules. On peut le
faire en choisissant une présentation $F_1 \to F_0 \to M \to 0$ avec $F_i$
libre, puis en prenant pour $K$ et $L$ les complétions $f$-adiques de $F_1$
et $F_0$. En effet, puisque $f$ est un non-diviseur de zéro, ces complétions
seront les complétions dérivées et la suite restera exacte.}. Considérons $K$
et $L$ comme des modules topologiques munis de la topologie $I$-adique. Alors
$u$ est continu. Posons
$$
L_n = \{x \in L \mid f^n x \in u(K)\}
$$
Comme $M$ est de torsion par les puissances de $f$, on a
$L = \bigcup L_n$. Soit $N_n$ l'adhérence de $L_n$ dans $L$. D'après le
Lemme \ref{lemma-consequence-baire-complete-module}, $N_n$ est ouvert dans
$L$ pour un certain $n$. Fixons un tel $n$. Comme
$f^{n + m} : L \to L$ est un morphisme continu et ouvert, et comme
$f^{n + m} L_n \subset u(f^m K)$, on conclut que l'adhérence de $u(f^mK)$
est ouverte pour tout $m \geq 1$. Le Lemme \ref{lemma-open-mapping} montre
donc que $u$ est ouvert. Ainsi, $f^tL \subset \Im(u)$ pour un certain $t$,
et l'on conclut que $f^t$ annule $M$, comme voulu.
\end{proof}

\begin{lemma}
\label{lemma-kernel-to-completion-square-zero}
Soit $f \in A$ un élément d'un anneau. Posons $J = \bigcap f^nA$. Soit $M$
un $A$-module complet au sens dérivé par rapport à $f$. Alors $JM' = 0$, où
$M' = \Ker(M \to \lim M/f^nM)$. En particulier, si $A$ est complet au sens
dérivé, alors $J$ est un idéal de carré nul.
\end{lemma}

\begin{proof}
Prenons $x \in M'$ et $g \in J$. Pour tout $n \geq 1$, on peut écrire
$x = f^n x_n$. Comme $g$ appartient à $f^nA$, l'élément $y_n = gx_n$ de
$M'$ est indépendant du choix de $x_n$. On peut en particulier prendre
$x_n = fx_{n + 1}$, et l'on trouve $y_n = fy_{n + 1}$. On obtient ainsi un
morphisme $A_f \to M$ qui envoie $1/f^n$ sur $y_n$. Ce morphisme doit être
nul puisque $M$ est complet au sens dérivé (Lemme
\ref{lemma-hom-from-Af}) ; par conséquent, $y_n = 0$ pour tout $n$. Comme
$gx = gfx_1 = fy_1$, la démonstration est achevée.
\end{proof}

\begin{lemma}
\label{lemma-derived-complete-henselian}
Soit $A$ un anneau complet au sens dérivé par rapport à un idéal $I$. Alors
$(A, I)$ est une paire hensélienne.
\end{lemma}

\begin{proof}
Soit $f \in I$. D'après le Lemme \ref{lemma-largest-ideal-henselian}, il
suffit de montrer que $(A, fA)$ est une paire hensélienne. Remarquons que $A$
est complet au sens dérivé par rapport à $fA$ (cela résulte immédiatement de
la Définition \ref{definition-derived-complete}). D'après le
Lemme \ref{lemma-complete-derived-complete}, le morphisme de $A$ vers la
complétion $f$-adique $A'$ de $A$ est surjectif. D'après le
Lemme \ref{lemma-complete-henselian}, la paire $(A', fA')$ est hensélienne. Il
suffit donc de montrer que $(A, \bigcap f^nA)$ est une paire hensélienne ;
voir le Lemme \ref{lemma-henselian-henselian-pair}. Cela résulte des Lemmes
\ref{lemma-kernel-to-completion-square-zero} et
\ref{lemma-locally-nilpotent-henselian}.
\end{proof}

\begin{lemma}
\label{lemma-kernel-to-completion-nilpotent}
Soit $A$ un anneau complet au sens dérivé par rapport à un idéal $I$. Posons
$J = \bigcap I^n$. Si $I$ peut être engendré par $r$ éléments, alors
$J^N = 0$, où $N = 2^r$.
\end{lemma}

\begin{proof}
Pour $r = 1$, c'est le Lemme
\ref{lemma-kernel-to-completion-square-zero}. Écrivons
$I = (f_1, \ldots, f_r)$ avec $r > 1$. D'après le
Lemme \ref{lemma-serre-subcategory}, l'anneau $A_t = A/f_r^tA$ est complet au
sens dérivé par rapport à $I$, donc a fortiori complet au sens dérivé par
rapport à $I_t = (f_1, \ldots, f_{r - 1})A_t$. Remarquons que $A \to A_t$
envoie $J$ dans $J_t = \bigcap I_t^n$. Par récurrence, $J_t^{N/2} = 0$ avec
$N = 2^r$. L'idéal $\bigcap \Ker(A \to A_t) = \bigcap f_r^t A$ est de carré
nul d'après le cas $r = 1$. Cela achève la démonstration.
\end{proof}

\begin{lemma}
\label{lemma-reduced-derived-complete-complete}
Soit $A$ un anneau réduit, complet au sens dérivé par rapport à un idéal de
type fini $I$. Alors $A$ est $I$-adiquement complet.
\end{lemma}

\begin{proof}
Cela résulte du Lemme \ref{lemma-kernel-to-completion-nilpotent} et de la
Proposition \ref{proposition-derived-complete-modules}.
\end{proof}







\section{Complétion dérivée pour les anneaux noethériens}
\label{section-derived-completion-noetherian}

\noindent
Soit $A$ un anneau et soit $I \subset A$ un idéal. Pour tout $K \in D(A)$,
on peut considérer la limite dérivée
$$
K' = R\lim (K \otimes_A^\mathbf{L} A/I^n)
$$
C'est un foncteur en $K$ ; voir la Remarque
\ref{remark-constructing-tensor-with-limits-functorially}. Le système de
morphismes $A \to A/I^n$ induit un morphisme $K \to K'$, et $K'$ est complet
au sens dérivé par rapport à $I$ (Lemme
\ref{lemma-naive-derived-completion}). En général, cette construction
« naïve » de la complétion dérivée ne coïncide pas avec l'adjoint du
Lemme \ref{lemma-derived-completion}. Par exemple, si
$A = \mathbf{Z}_p \oplus \mathbf{Q}_p/\mathbf{Z}_p$, le second facteur étant
un idéal de carré nul, si $K = A[0]$ et $I = (p)$, alors la complétion dérivée
naïve donne $\mathbf{Z}_p[0]$, tandis que la construction du
Lemme \ref{lemma-derived-completion} donne
$K^\wedge \cong \mathbf{Z}_p[1] \oplus \mathbf{Z}_p[0]$ (calcul omis). Le
Lemme \ref{lemma-when-does-it-work} caractérise les cas où les deux foncteurs
coïncident lorsque $I$ est engendré par un seul élément.

\medskip\noindent
Le but principal de cette section est de montrer que la complétion dérivée
naïve est égale à la complétion dérivée lorsque $A$ est noethérien.

\begin{lemma}
\label{lemma-sequence-Koszul-complexes}
Dans la Situation \ref{situation-koszul}, si $A$ est noethérien, alors les
pro-objets $\{K_n^\bullet\}$ et $\{A/(f_1^n, \ldots, f_r^n)\}$ de $D(A)$ sont
isomorphes\footnote{En particulier, pour tout $n$, il existe un $m \geq n$
tel que $K_m^\bullet \to K_n^\bullet$ se factorise par le morphisme
$K_m^\bullet \to A/(f_1^m, \ldots, f_r^m)$.}.
\end{lemma}

\begin{proof}
Nous avons un système projectif de triangles distingués
$$
\tau_{\leq -1}K_n^\bullet \to K_n^\bullet \to A/(f_1^n, \ldots, f_r^n) \to
(\tau_{\leq -1}K_n^\bullet)[1]
$$
Voir Catégories dérivées, Remarque
\ref{derived-remark-truncation-distinguished-triangle}. D'après Catégories
dérivées, Lemme \ref{derived-lemma-pro-isomorphism}, il suffit de montrer que
le système projectif $\tau_{\leq -1}K_n^\bullet$ est pro-nul. Rappelons que
$K_n^\bullet$ n'a de termes non nuls qu'aux degrés $i$ tels que
$-r \leq i \leq 0$. D'après Catégories dérivées, Lemme
\ref{derived-lemma-essentially-constant-cohomology}, il suffit donc de montrer
que $H^p(K_n^\bullet)$ est pro-nul pour $p \leq -1$. Autrement dit, pour tout
$n \in \mathbf{N}$, il faut montrer qu'il existe un $m \geq n$ tel que
$H^p(K_m^\bullet) \to H^p(K_n^\bullet)$ soit nul. Comme $A$ est noethérien,
on a
$$
H^p(K_n^\bullet) =
\frac{\Ker(K_n^p \to K_n^{p + 1})}{\Im(K_n^{p - 1} \to K_n^p)}
$$
qui est un $A$-module de type fini. De plus, le morphisme
$K_m^p \to K_n^p$ est donné par une matrice diagonale dont les coefficients
appartiennent à l'idéal $(f_1^{m - n}, \ldots, f_r^{m - n})$, puisque
$p < 0$. Remarquons que $H^p(K_n^\bullet)$ est annulé par
$J = (f_1^n, \ldots, f_r^n)$ ; voir le Lemme
\ref{lemma-homotopy-koszul}. Or
$(f_1^{m - n}, \ldots, f_r^{m - n}) \subset J^t$ dès que
$m - n \geq tn$. Appliquons donc Algèbre, Lemme
\ref{algebra-lemma-Artin-Rees} (Artin--Rees), à l'idéal $J$, au module
$M = K_n^p$ et à son sous-module
$N = \Ker(K_n^p \to K_n^{p + 1})$. Pour $m$ assez grand, l'intersection de
l'image de $K_m^p \to K_n^p$ avec
$\Ker(K_n^p \to K_n^{p + 1})$ est contenue dans
$J \Ker(K_n^p \to K_n^{p + 1})$. Pour un tel $m$, le morphisme obtenu est
nul.
\end{proof}

\begin{proposition}
\label{proposition-noetherian-naive-completion-is-completion}
Soit $A$ un anneau noethérien et soit $I \subset A$ un idéal. Le foncteur qui
envoie $K \in D(A)$ sur la limite dérivée
$K' = R\lim( K \otimes_A^\mathbf{L} A/I^n )$ est l'adjoint à gauche du
foncteur d'inclusion $D_{comp}(A) \to D(A)$ construit dans le
Lemme \ref{lemma-derived-completion}.
\end{proposition}

\begin{proof}
Écrivons $(f_1, \ldots, f_r) = I$ et soit $K_n^\bullet$ le complexe de Koszul
associé à $f_1^n, \ldots, f_r^n$. D'après le
Lemme \ref{lemma-derived-completion-koszul}, il suffit de montrer que
$$
R\lim (K \otimes_A^\mathbf{L} K_n^\bullet) =
R\lim (K \otimes_A^\mathbf{L} A/(f_1^n, \ldots, f_r^n) ) =
R\lim (K \otimes_A^\mathbf{L} A/I^n ).
$$
D'après le Lemme \ref{lemma-sequence-Koszul-complexes}, les pro-objets
$\{K_n^\bullet\}$ et $\{A/(f_1^n, \ldots, f_r^n)\}$ de $D(A)$ sont
isomorphes. Il est clair que les pro-objets
$\{A/(f_1^n, \ldots, f_r^n)\}$ et $\{A/I^n\}$ sont isomorphes. Le morphisme
de gauche à droite est donc un isomorphisme d'après le
Lemme \ref{lemma-tensor-Rlim-pro-equal}.
\end{proof}

\begin{lemma}
\label{lemma-noetherian-calculate}
Soit $I$ un idéal d'un anneau noethérien $A$. Soit $M$ un $A$-module de
complétion dérivée $M^\wedge$. On a alors des suites exactes courtes
$$
0 \to R^1\lim \text{Tor}_{i + 1}^A(M, A/I^n) \to
H^{-i}(M^\wedge) \to \lim \text{Tor}_i^A(M, A/I^n) \to 0
$$
Un résultat analogue vaut pour $M \in D^-(A)$.
\end{lemma}

\begin{proof}
C'est une conséquence immédiate de la
Proposition \ref{proposition-noetherian-naive-completion-is-completion} et du
Lemme \ref{lemma-break-long-exact-sequence-modules}.
\end{proof}

\noindent
Comme application de la proposition précédente, nous déterminons la
complétion dérivée des complexes pseudo-cohérents dans le cas noethérien.

\begin{lemma}
\label{lemma-derived-completion-pseudo-coherent}
Soit $A$ un anneau noethérien, soit $I \subset A$ un idéal et soit $K$ un
objet de $D(A)$ tel que $H^n(K)$ soit un $A$-module de type fini pour tout
$n \in \mathbf{Z}$. Alors les modules de cohomologie $H^n(K^\wedge)$ de la
complétion dérivée sont les complétions $I$-adiques des modules de
cohomologie $H^n(K)$.
\end{lemma}

\begin{proof}
Le complexe $\tau_{\leq m}K$ est pseudo-cohérent pour tout $m$ d'après le
Lemme \ref{lemma-Noetherian-pseudo-coherent}. Ainsi, $\tau_{\leq m}K$ est
représenté par un
complexe borné supérieurement $P^\bullet$ de $A$-modules libres de rang fini.
Alors $\tau_{\leq m}K \otimes_A^\mathbf{L} A/I^n = P^\bullet/I^nP^\bullet$.
Ainsi, $(\tau_{\leq m}K)^\wedge = R\lim P^\bullet/I^nP^\bullet$
(Proposition \ref{proposition-noetherian-naive-completion-is-completion}).
Comme $R\lim$ est ici simplement donné par $\lim$ terme à terme (Lemme
\ref{lemma-compute-Rlim-modules}) et comme la complétion $I$-adique est un
foncteur exact sur les $A$-modules de type fini (Algèbre, Lemme
\ref{algebra-lemma-completion-flat}), le résultat vaut pour
$\tau_{\leq m}K$. Il vaut donc pour $K$, puisque la complétion dérivée est de
dimension cohomologique finie ; voir le
Lemme \ref{lemma-derived-completion-finite-cohomological-dimension}.
\end{proof}

\begin{lemma}
\label{lemma-derived-complete-finite}
Soit $I$ un idéal d'un anneau noethérien $A$. Soit $M$ un $A$-module complet
au sens dérivé. Si $M/IM$ est un $A/I$-module de type fini, alors
$M = \lim M/I^nM$ et $M$ est un $A^\wedge$-module de type fini.
\end{lemma}

\begin{proof}
Supposons $M/IM$ de type fini. Choisissons $x_1, \ldots, x_t \in M$ dont les
images engendrent $M/IM$. On obtient un morphisme $A^{\oplus t} \to M$ qui
envoie le $i$-ième vecteur de base sur $x_i$. D'après la
Proposition \ref{proposition-noetherian-naive-completion-is-completion}, la
complétion dérivée de $A$ est $A^\wedge = \lim A/I^n$. Comme $M$ est complet
au sens dérivé, notre morphisme se factorise par un morphisme
$q : (A^\wedge)^{\oplus t} \to M$. Le module $\Coker(q)$ est nul d'après le
Lemme \ref{lemma-derived-complete-zero}. Ainsi, $M$ est un
$A^\wedge$-module de type fini. Puisque $A^\wedge$ est noethérien et complet
par rapport à $IA^\wedge$, il s'ensuit que $M$ est $I$-adiquement complet
(utiliser Algèbre, Lemmes \ref{algebra-lemma-completion-Noetherian},
\ref{algebra-lemma-when-finite-module-complete-over-complete-ring} et
\ref{algebra-lemma-Artin-Rees}).
\end{proof}

\begin{lemma}
\label{lemma-when-derived-completion-is-completion}
Soit $I$ un idéal d'un anneau noethérien $A$.
\begin{enumerate}
\item Si $M$ est un $A$-module de type fini et $N$ un $A$-module plat, alors
la complétion $I$-adique dérivée de $M \otimes_A N$ est la complétion
$I$-adique usuelle de $M \otimes_A N$.
\item Si $M$ est un $A$-module de type fini et si $f \in A$, alors la
complétion $I$-adique dérivée de $M_f$ est la complétion $I$-adique usuelle
de $M_f$.
\end{enumerate}
\end{lemma}

\begin{proof}
Pour un $A$-module $M$, notons $M^\wedge$ sa complétion dérivée et
$\lim M/I^nM$ sa complétion usuelle. Supposons $M$ de type fini. Le système
$\text{Tor}^A_i(M, A/I^n)$ est pro-nul pour $i > 0$ ; voir le
Lemme \ref{lemma-tor-strictly-pro-zero}. Puisque
$\text{Tor}_i^A(M \otimes_A N, A/I^n) =
\text{Tor}_i^A(M, A/I^n) \otimes_A N$, car $N$ est plat, il en va de même du
système $\text{Tor}^A_i(M \otimes_A N, A/I^n)$. D'après le
Lemme \ref{lemma-noetherian-calculate}, on conclut que
$R\lim (M \otimes_A N) \otimes_A^\mathbf{L} A/I^n$ n'a de cohomologie qu'au
degré $0$, où elle est donnée par la complétion usuelle
$\lim M \otimes_A N/ I^n(M \otimes_A N)$. Cela prouve (1). L'assertion (2)
résulte de (1) et de l'égalité $M_f = M \otimes_A A_f$.
\end{proof}

\begin{lemma}
\label{lemma-derived-completion-tensor-finite}
Soit $I$ un idéal d'un anneau noethérien $A$. Notons ${}^\wedge$ la
complétion dérivée par rapport à $I$. Soit $K \in D^-(A)$.
\begin{enumerate}
\item Si $M$ est un $A$-module de type fini, alors
$(K \otimes_A^\mathbf{L} M)^\wedge = K^\wedge \otimes_A^\mathbf{L} M$.
\item Si $L \in D(A)$ est pseudo-cohérent, alors
$(K \otimes_A^\mathbf{L} L)^\wedge = K^\wedge \otimes_A^\mathbf{L} L$.
\end{enumerate}
\end{lemma}

\begin{proof}
Soit $L$ comme dans (2). On peut représenter $K$ par un complexe borné
supérieurement $P^\bullet$ de $A$-modules libres, et $L$ par un complexe borné
supérieurement $F^\bullet$ de $A$-modules libres de rang fini. Comme
$\text{Tot}(P^\bullet \otimes_A F^\bullet)$ représente
$K \otimes_A^\mathbf{L} L$, on voit que
$(K \otimes_A^\mathbf{L} L)^\wedge$ est représenté par
$$
\text{Tot}((P^\bullet)^\wedge \otimes_A F^\bullet)
$$
où $(P^\bullet)^\wedge$ est le complexe dont les termes sont les complétions
usuelles, égales aux complétions dérivées, $(P^n)^\wedge$ ; voir par exemple
la Proposition \ref{proposition-noetherian-naive-completion-is-completion} et
le Lemme \ref{lemma-when-derived-completion-is-completion}. Cela prouve (2),
et l'assertion (1) en est un cas particulier.
\end{proof}








\section{Un opérateur introduit par Berthelot et Ogus}
\label{section-eta}

\noindent
Dans cette section, nous étudions une construction introduite dans
\cite[Section 8]{Berthelot-Ogus} et généralisée dans \cite[Section 6]{BMS}.
Nous recommandons vivement au lecteur de consulter les articles originaux où
cette notion est étudiée.

\medskip\noindent
Soit $A$ un anneau, soit $f \in A$ un non-diviseur de zéro et soit $M$ un
$A$-module. D'après le Lemme \ref{lemma-torsion-free}, les conditions suivantes
sont équivalentes :
\begin{enumerate}
\item $f$ est un non-diviseur de zéro sur $M$ ;
\item $M[f] = 0$,
\item $M[f^n] = 0$ pour tout $n \geq 1$ ;
\item le morphisme $M \to M_f$ est injectif.
\end{enumerate}
Lorsque ces conditions équivalentes sont satisfaites, nous dirons dans cette
section que {\it $M$ est sans $f$-torsion}. Dans ce cas, nous notons
$f^iM \subset M_f$ le sous-module formé des éléments de la forme $f^ix$ avec
$x \in M$. Bien entendu, $f^iM$ est isomorphe à $M$ comme $A$-module. Soit
$M^\bullet$ un complexe de modules sans $f$-torsion sur $A$, de différentielles
$d^i : M^i \to M^{i + 1}$. On définit alors $\eta_fM^\bullet$ comme le
complexe de termes
$$
(\eta_fM)^i = \{x \in f^iM^i \mid d^i(x) \in f^{i + 1}M^{i + 1}\}
$$
et dont la différentielle est induite par $d^i$. Remarquons que
$\eta_fM^\bullet$ est encore un complexe de modules sans $f$-torsion sur $A$.
Si
$a^\bullet : M^\bullet \to N^\bullet$ est un morphisme de complexes de
modules sans $f$-torsion sur $A$, on obtient un morphisme de complexes
$$
\eta_fa^\bullet : \eta_fM^\bullet \longrightarrow \eta_fN^\bullet
$$
induit par les morphismes $f^iM^i \to f^iN^i$. Le lecteur vérifiera que l'on
obtient ainsi un endofoncteur de la catégorie des complexes de modules sans
$f$-torsion sur $A$. Si
$a^\bullet, b^\bullet : M^\bullet \to N^\bullet$ sont deux morphismes de
complexes formés de modules sans $f$-torsion sur $A$, et si
$h = \{h^i : M^i \to N^{i - 1}\}$ est une homotopie entre
$a^\bullet$ et $b^\bullet$, on définit $\eta_fh$ comme la famille de
morphismes
$(\eta_fh)^i : (\eta_fM)^i \to (\eta_fN)^{i - 1}$
qui envoie $x$ sur $h^i(x)$ ; cette définition a un sens, car
$x \in f^iM^i$ entraîne $h^i(x) \in f^iN^{i - 1}$, lequel est bien contenu
dans $(\eta_fN)^{i - 1}$. Le lecteur vérifiera que $\eta_fh$ est une
homotopie entre $\eta_fa^\bullet$ et $\eta_fb^\bullet$. Au total, on obtient
un foncteur
$$
\eta_f :
K(f\text{-torsion free }A\text{-modules})
\longrightarrow
K(f\text{-torsion free }A\text{-modules})
$$
sur la catégorie d'homotopie (Catégories dérivées, Section
\ref{derived-section-homotopy}) de la catégorie additive des modules sans
$f$-torsion sur $A$. Le foncteur $\eta_f$ n'est exact comme foncteur de catégories
triangulées en aucun sens ; voir l'Exemple
\ref{example-eta-not-distinguished}.

\begin{example}
\label{example-eta-not-distinguished}
Soit $A$ un anneau et soit $f \in A$ un non-diviseur de zéro. Considérons le
foncteur $\eta_f : K(f\text{-torsion free }A\text{-modules})
\to K(f\text{-torsion free }A\text{-modules})$. Soit $M^\bullet$ un complexe
de modules sans $f$-torsion sur $A$. La multiplication par $f$ définit un
isomorphisme $\eta_f(M^\bullet[1]) \to (\eta_fM^\bullet)[1]$ ; en ce sens,
$\eta_f$ est compatible aux décalages. Considérons toutefois le diagramme
$$
\xymatrix{
A \ar[r]_f &
A \ar[r]_1 &
A \ar[r] &
0 \\
0 \ar[r] \ar[u] &
0 \ar[r] \ar[u] &
A \ar[r]^{-1} \ar[u]^f &
A \ar[u]
}
$$
Considérons chaque colonne comme un complexe de modules sans $f$-torsion sur $A$,
le module supérieur étant placé au degré $1$ et le module inférieur au degré
$0$. Ce diagramme fournit alors un triangle distingué dans
$K(f\text{-torsion free }A\text{-modules})$, muni de la structure triangulée
décrite dans Catégories dérivées, Section
\ref{derived-section-homotopy-triangulated}. En effet, le troisième complexe
est le cône du morphisme entre les deux premiers. En appliquant $\eta_f$ à
chaque colonne, on obtient toutefois
$$
\xymatrix{
fA \ar[r]_f &
fA \ar[r]_1 &
fA \ar[r] &
0 \\
0 \ar[r] \ar[u] &
0 \ar[r] \ar[u] &
A \ar[r]^{-1} \ar[u]^f &
A \ar[u]
}
$$
Or le troisième complexe est acyclique, et même homotope à zéro. S'il
s'agissait d'un triangle distingué, la première flèche devrait donc être un
isomorphisme dans la catégorie d'homotopie, ce qui est faux sauf si $f$ est
inversible.
\end{example}

\begin{lemma}
\label{lemma-eta-first-property}
Soit $A$ un anneau, soit $f \in A$ un non-diviseur de zéro et soit
$M^\bullet$ un complexe de modules sans $f$-torsion sur $A$. Il existe un
isomorphisme canonique
$$
f^i : H^i(M^\bullet)/H^i(M^\bullet)[f] \longrightarrow H^i(\eta_fM^\bullet)
$$
donné par la multiplication par $f^i$.
\end{lemma}

\begin{proof}
Remarquons que $\Ker(d^i : (\eta_fM)^i \to (\eta_fM)^{i + 1})$ est égal à
$\Ker(d^i : f^iM^i \to f^iM^{i + 1}) = f^i\Ker(d^i : M^i \to M^{i + 1})$.
On obtient ainsi une surjection
$f^i : H^i(M^\bullet) \to H^i(\eta_fM^\bullet)$ qui envoie la classe de
$z \in \Ker(d^i : M^i \to M^{i + 1})$ sur celle de $f^iz$. Si l'on obtient
la classe nulle dans $H^i(\eta_fM^\bullet)$, alors
$f^i z = d^{i - 1}(f^{i - 1}y)$ pour un certain $y \in M^{i - 1}$. Puisque
$f$ est un non-diviseur de zéro sur tous les modules en jeu, cela signifie
que $f z = d^{i - 1}(y)$, ce qui dit exactement que la classe de $z$ est de
$f$-torsion, comme voulu.
\end{proof}

\begin{lemma}
\label{lemma-eta-qis}
Soit $A$ un anneau et soit $f \in A$ un non-diviseur de zéro. Si
$M^\bullet \to N^\bullet$ est un quasi-isomorphisme de complexes de
modules sans $f$-torsion sur $A$, alors le morphisme induit
$\eta_fM^\bullet \to \eta_fN^\bullet$ est lui aussi un quasi-isomorphisme.
\end{lemma}

\begin{proof}
Cela résulte de la compatibilité des isomorphismes du
Lemme \ref{lemma-eta-first-property} avec les morphismes de complexes.
\end{proof}

\begin{lemma}
\label{lemma-Leta}
Soit $A$ un anneau et soit $f \in A$ un non-diviseur de zéro. Il existe un
foncteur additif\footnote{Attention : ce foncteur n'est pas exact, c'est-à-dire
qu'il ne transforme pas les triangles distingués en triangles distingués.
Voir l'Exemple \ref{example-eta-not-distinguished}.}
$L\eta_f : D(A) \to D(A)$ tel que, si $M \in D(A)$ est représenté par un
complexe $M^\bullet$ de modules sans $f$-torsion sur $A$, alors
$L\eta_fM = \eta_fM^\bullet$ ; il en va de même pour les morphismes.
\end{lemma}

\begin{proof}
Notons $\mathcal{T} \subset \text{Mod}_A$ la sous-catégorie pleine des
modules sans $f$-torsion sur $A$. On a l'inclusion correspondante
$$
K(\mathcal{T}) \quad\subset\quad K(\text{Mod}_A) = K(A)
$$
qui fait de $K(\mathcal{T})$ une sous-catégorie triangulée pleine de $K(A)$.
Soit $S \subset \text{Arrows}(K(\mathcal{T}))$ la classe des
quasi-isomorphismes. Nous allons appliquer Catégories dérivées, Lemme
\ref{derived-lemma-localization-subcategory}, pour montrer que le foncteur
$$
S^{-1}K(\mathcal{T}) \longrightarrow D(A)
$$
est une équivalence de catégories triangulées. Le lemme montre qu'il suffit
d'établir ceci : pour tout complexe $M^\bullet$ de $A$-modules, il existe un
quasi-isomorphisme $K^\bullet \to M^\bullet$ où $K^\bullet$ est un complexe
de modules sans $f$-torsion. D'après le Lemme
\ref{lemma-K-flat-resolution}, on peut trouver un quasi-isomorphisme
$K^\bullet \to M^\bullet$ tel que le complexe $K^\bullet$ soit K-plat (ce que
nous n'utiliserons pas) et formé de $A$-modules plats $K^i$. En particulier,
$f$ est un non-diviseur de zéro sur $K^i$ pour tout $i$, comme voulu.

\medskip\noindent
Ces préliminaires étant établis, on peut définir $L\eta_f$. En effet, la
discussion du début de cette section fournit déjà un foncteur bien défini
$$
K(\mathcal{T}) \xrightarrow{\eta_f} K(\mathcal{T}) \to K(A) \to D(A)
$$
qui, d'après le Lemme \ref{lemma-eta-qis}, envoie les quasi-isomorphismes sur
des quasi-isomorphismes. Il se factorise donc par
$S^{-1}K(\mathcal{T}) = D(A)$ d'après Catégories, Lemme
\ref{categories-lemma-properties-left-localization}.
\end{proof}

\begin{remark}
\label{remark-eta-BZ}
Soit $A$ un anneau, soit $f \in A$ un non-diviseur de zéro et soit
$M^\bullet$ un complexe de modules sans $f$-torsion sur $A$. Pour tout $i$, posons
$\overline{M}^i = M^i/fM^i$. Notons
$B^i \subset Z^i \subset \overline{M}^i$ les cobords et les cocycles pour les
différentielles du complexe
$\overline{M}^\bullet = M^\bullet \otimes_A A/fA$. Nous affirmons qu'il existe
un diagramme commutatif
$$
\xymatrix{
0 \ar[r] &
B^{i + 1} \ar[r] \ar@{=}[d] &
B^{i + 1} \oplus B^i \ar[r] \ar[d]^{s, s'} &
B^i \ar[r] \ar[d] & 0 \\
0 \ar[r] &
B^{i + 1} \ar[r]^-s &
(\eta_fM)^i /f(\eta_fM)^i \ar[r]^-t &
Z^i \ar[r] & 0
}
$$
à lignes exactes. Voici la construction des morphismes :
\begin{enumerate}
\item Si $x \in (\eta_fM)^i$, alors $x = f^ix'$ avec $d^i(x') = 0$ dans
$\overline{M}^{i + 1}$. On peut donc définir $t$ en envoyant $x$ sur la classe
de $x'$.
\item Si $y \in M^{i + 1}$ a pour classe $\overline{y}$ dans
$B^{i + 1} \subset \overline{M}^{i + 1}$, alors on peut écrire
$y = fy' + d^i(x)$ avec $y' \in M^{i + 1}$ et $x \in M^i$. On peut donc
définir $s$ en envoyant $\overline{y}$ sur la classe de $f^{i + 1}x$ dans
$(\eta_fM)^i /f(\eta_fM)^i$ ; nous omettons la vérification que cette
définition est indépendante des choix.
\item Si $x \in M^i$ a pour classe $\overline{x}$ dans
$B^i \subset \overline{M}^i$, alors on peut écrire
$x = fx' + d^{i - 1}(z)$ avec $x' \in M^i$ et $z \in M^{i - 1}$. On définit
$s'$ en envoyant $\overline{x}$ sur la classe de $f^i d^{i - 1}(z)$ dans
$(\eta_fM)^i/f(\eta_fM)^i$. Cette définition est indépendante des choix, car
si $fx' + d^{i - 1}(z) = 0$, alors $f^ix'$ appartient à $(\eta_fM)^i$ et,
par conséquent, $f^id^{i - 1}(z)$ appartient à $f(\eta_fM)^i$.
\end{enumerate}
Nous omettons de vérifier que la ligne inférieure du diagramme affiché est une
suite exacte courte de modules. Ces constructions montrent immédiatement que
l'on a des diagrammes commutatifs
$$
\xymatrix{
B^{i + 1} \oplus B^i \ar[d]^{s, s'} \ar[r] &
B^{i + 2} \oplus B^{i + 1} \ar[d]^{s, s'} \\
(\eta_fM)^i /f(\eta_fM)^i \ar[r] &
(\eta_fM)^{i + 1} /f(\eta_fM)^{i + 1}
}
$$
où la flèche horizontale supérieure est donnée par l'identification des
facteurs $B^{i + 1}$ à la source et au but. Autrement dit, nous avons trouvé
un sous-complexe acyclique de
$\eta_fM^\bullet / f(\eta_fM^\bullet) = \eta_fM^\bullet \otimes_A A/fA$, et
le quotient par ce sous-complexe est un complexe dont les termes $Z^i/B^i$
sont les modules de cohomologie du complexe
$\overline{M}^\bullet = M^\bullet \otimes_A A/fA$.
\end{remark}

\noindent
Pour expliquer de manière plus canonique le phénomène observé dans la
Remarque \ref{remark-eta-BZ}, nous allons construire les opérateurs de
Bockstein. Soit $A$ un anneau, soit $f \in A$ un non-diviseur de zéro et soit
$M^\bullet$ un complexe de modules sans $f$-torsion sur $A$. Pour tout
$i \in \mathbf{Z}$, on a un diagramme commutatif (les produits tensoriels
étant pris sur $A$)
$$
\xymatrix{
0 \ar[r] &
M^\bullet \otimes f^{i + 1}A \ar[r] \ar[d] &
M^\bullet \otimes f^iA \ar[r] \ar[d] &
M^\bullet \otimes f^iA/f^{i + 1}A \ar[r] \ar@{=}[d] &
0 \\
0 \ar[r] &
M^\bullet \otimes f^{i + 1}A/f^{i + 2}A \ar[r] &
M^\bullet \otimes f^iA/f^{i + 2}A \ar[r] &
M^\bullet \otimes f^iA/f^{i + 1}A \ar[r] &
0
}
$$
dont les lignes sont des suites exactes courtes de complexes. Bien entendu,
pour les différentes valeurs de $i$, ces suites exactes courtes sont toutes
isomorphes entre elles par multiplication par des puissances convenables de
$f$. La suite exacte longue de cohomologie de la suite inférieure détermine en
particulier l'{\it opérateur de Bockstein}
$$
\beta = \beta^i : H^i(M^\bullet \otimes f^iA/f^{i + 1}A) \to
H^{i + 1}(M^\bullet \otimes f^{i + 1}A/f^{i + 2}A)
$$
pour tout $i \in \mathbf{Z}$. Pour un usage ultérieur, notons que le diagramme
commutatif précédent fournit une factorisation
\begin{equation}
\label{equation-factorization-bockstein}
\vcenter{
\xymatrix{
H^i(M^\bullet \otimes f^iA/f^{i + 1}A)
\ar[r]_\delta \ar[rd]_\beta &
H^{i + 1}(M^\bullet \otimes f^{i + 1}A) \ar[d] \\
&
H^{i + 1}(M^\bullet \otimes f^{i + 1}A/f^{i + 2}A)
}
}
\end{equation}
de l'opérateur de Bockstein, où $\delta$ est le morphisme de connexion issu de
la ligne supérieure du diagramme commutatif précédent. Montrons que l'on
obtient un complexe
\begin{equation}
\label{equation-complex-bocksteins}
H^\bullet(M^\bullet/f) =
\left[
\begin{matrix}
\ldots \\
\downarrow \\
H^{i - 1}(M^\bullet \otimes f^{i - 1}A/f^iA) \\
\downarrow \beta \\
H^i(M^\bullet \otimes f^iA/f^{i + 1}A) \\
\downarrow \beta \\
H^{i + 1}(M^\bullet \otimes f^{i + 1}A/f^{i + 2}A) \\
\downarrow \\
\ldots
\end{matrix}
\right]
\end{equation}
c'est-à-dire que $\beta \circ \beta = 0$\footnote{On peut aussi remarquer que
$\beta$ apparaît comme la différentielle de la suite spectrale du complexe
$(M^\bullet)_f$ filtré par les sous-complexes $f^iM^\bullet$. Un autre
argument encore, qui donne un résultat plus fort, consiste à considérer
d'abord le cas $M^\bullet = A$. Les suites exactes courtes
$0 \to f^{i + 1}A/f^{i + 2}A \to f^iA/f^{i + 2}A \to f^iA/f^{i + 1}A \to 0$
définissent alors des morphismes
$\beta^i : f^iA/f^{i + 1}A \to f^{i + 1}A/f^{i + 2}A[1]$ dans $D(A)$. On calcule ensuite,
par un raisonnement analogue à celui du texte, que la composée
$f^iA/f^{i + 1}A \to f^{i + 1}A/f^{i + 2}A[1] \to f^{i + 2}A/f^{i + 3}A[2]$
est nulle dans $D(A)$. Comme $M^\bullet \otimes f^iA/f^{i + 1}A =
M^\bullet \otimes^\mathbf{L} f^iA/f^{i + 1}A$ puisque les termes de
$M^\bullet$ sont sans $f$-torsion, on conclut que la composée
$$
(M^\bullet \otimes f^iA/f^{i + 1}A)
\to
(M^\bullet \otimes f^{i + 1}A/f^{i + 2}A)[1]
\to
(M^\bullet \otimes f^{i + 2}A/f^{i + 3}A)[2]
$$
est elle aussi nulle dans $D(A)$.}. En effet, en utilisant la factorisation
(\ref{equation-factorization-bockstein}), on voit qu'il suffit de montrer que
$$
H^{i + 1}(M^\bullet \otimes f^{i + 1}A)
\to
H^{i + 1}(M^\bullet \otimes f^{i + 1}A/f^{i + 2}A)
\xrightarrow{\beta^{i + 1}}
H^{i + 2}(M^\bullet \otimes f^{i + 2}A/f^{i + 3}A)
$$
est nul. Cela vient du fait que le noyau de $\beta^{i + 1}$ est formé des
classes de cohomologie qui se relèvent à
$H^{i + 1}(M^\bullet \otimes f^{i + 1}A/f^{i + 3}A)$, et que celles qui sont
dans l'image du premier morphisme se relèvent certainement.

\begin{lemma}
\label{lemma-eta-second-property}
Soit $A$ un anneau, soit $f \in A$ un non-diviseur de zéro et soit
$M^\bullet$ un complexe de modules sans $f$-torsion sur $A$. Il existe un
morphisme canonique de complexes
$$
\eta_fM^\bullet \otimes_A A/fA
\longrightarrow
H^\bullet(M^\bullet/f)
$$
qui est un quasi-isomorphisme, le membre de droite étant le complexe
(\ref{equation-complex-bocksteins}).
\end{lemma}

\begin{proof}
Soit $x \in (\eta_fM)^i$. Alors $x = f^ix' \in f^iM$ et
$d^i(x) = f^{i + 1}y \in f^{i + 1}M^{i + 1}$. Ainsi, $d^i$ envoie
$x' \otimes f^i$ sur zéro dans
$M^{i + 1} \otimes f^iA/f^{i + 1}A$. Tous les produits tensoriels de cette
démonstration sont pris sur $A$. On peut donc envoyer $x$ sur la classe de
$x' \otimes f^i$ dans $H^i(M^\bullet \otimes f^iA/f^{i + 1}A)$. Il est clair
que cette règle définit un morphisme
$$
(\eta_fM)^i \otimes A/fA
\longrightarrow
H^i(M^\bullet \otimes f^iA/f^{i + 1}A)
$$
de $A/fA$-modules. Remarquons que, dans la situation précédente, on peut voir
$x' \otimes f^i$ comme un élément de $M^i \otimes f^iA/f^{i + 2}A$, de
différentielle $d^i(x' \otimes f^i) = y \otimes f^{i + 1}$. La construction
de $\beta$ donne alors
$\beta(x' \otimes f^i) = y \otimes f^{i + 1}$ ; nos morphismes sont donc
compatibles aux différentielles, c'est-à-dire qu'ils forment un morphisme de
complexes.

\medskip\noindent
Pour achever la démonstration, remarquons que la construction du paragraphe
précédent coïncide avec les morphismes
$(\eta_fM)^i \otimes A/fA \to Z^i/B^i$ étudiés dans la
Remarque \ref{remark-eta-BZ}. Puisque nous avons vu que le noyau de ces
morphismes est un sous-complexe acyclique de
$\eta_fM^\bullet \otimes A/fA$, le lemme est démontré.
\end{proof}

\begin{lemma}
\label{lemma-vanishing-beta}
Soit $A$ un anneau, soit $f \in A$ un non-diviseur de zéro et soit
$M^\bullet$ un complexe de modules sans $f$-torsion sur $A$. Pour
$i \in \mathbf{Z}$, les conditions suivantes sont équivalentes :
\begin{enumerate}
\item $\Ker(d^i \bmod f^2)$ se surjecte sur $\Ker(d^i \bmod f)$ ;
\item $\beta : H^i(M^\bullet \otimes_A f^iA/f^{i + 1}A) \to
H^{i + 1}(M^\bullet \otimes_A f^{i + 1}A/f^{i + 2}A)$ est nul.
\end{enumerate}
Ces conditions équivalentes sont impliquées par la condition
$H^{i + 1}(M^\bullet)[f] = 0$.
\end{lemma}

\begin{proof}
L'équivalence de (1) et (2) résulte de la définition de $\beta$ comme
morphisme de connexion en cohomologie d'une suite exacte courte de complexes
isomorphe à la suite exacte courte de complexes
$0 \to fM^\bullet/f^2M^\bullet \to M^\bullet/f^2M^\bullet \to
M^\bullet/fM^\bullet \to 0$. If $\beta \not = 0$,
alors $H^{i + 1}(M^\bullet)[f] \not = 0$ en raison de la factorisation
(\ref{equation-factorization-bockstein}).
\end{proof}

\begin{lemma}
\label{lemma-eta-vanishing-beta}
Soit $A$ un anneau, soit $f \in A$ un non-diviseur de zéro et soit
$M^\bullet$ un complexe de modules sans $f$-torsion sur $A$. Si
$\Ker(d^i \bmod f^2)$ se surjecte sur $\Ker(d^i \bmod f)$, alors le morphisme
canonique
$$
(1, d^i) :
(\eta_fM)^i / f(\eta_fM)^i \longrightarrow
f^iM^i/f^{i + 1}M^i \oplus f^{i + 1}M^{i + 1}/f^{i + 2}M^{i + 1}
$$
identifie le membre de gauche à une somme directe de sous-modules du membre de
droite.
\end{lemma}

\begin{proof}
Avec les notations de la Remarque \ref{remark-eta-BZ}, définissons un
morphisme $t^{-1} : Z^i \to (\eta_fM)^i / f(\eta_fM)^i$. Pour
$x \in M^i$ tel que $d^i(x) = f^2y$, on envoie la classe de $x$ dans $Z^i$
sur la classe de $f^ix$ dans $(\eta_fM)^i / f(\eta_fM)^i$. Nous omettons la
vérification que cette définition est indépendante des choix ; l'hypothèse du
lemme signifie exactement que le domaine de cette opération est tout $Z^i$.
Alors
$t \circ t^{-1} = \text{id}_{Z^i}$. Par conséquent, $t^{-1}$
définit un scindage de la suite exacte courte de la
Remarque \ref{remark-eta-BZ}, et la décomposition en somme directe qui en
résulte
$$
(\eta_fM)^i / f(\eta_fM)^i = Z^i \oplus B^{i + 1}
$$
est compatible avec le morphisme affiché dans le lemme.
\end{proof}

\begin{lemma}
\label{lemma-eta-third-property}
Soit $A$ un anneau, soient $f, g \in A$ des non-diviseurs de zéro et soit
$M^\bullet$ un complexe de $A$-modules tel que $fg$ soit un non-diviseur de
zéro sur tout $M^i$. Alors $\eta_f\eta_gM^\bullet = \eta_{fg}M^\bullet$.
\end{lemma}

\begin{proof}
L'énoncé signifie qu'au degré $i$, on obtient le même sous-module du localisé
$M^i_{fg} = (M^i_g)_f$. Nous omettons les détails.
\end{proof}

\begin{lemma}
\label{lemma-eta-flat-base-change}
Soit $A$ un anneau et soit $f \in A$ un non-diviseur de zéro. Soit
$A \to B$ un morphisme plat d'anneaux, et soit $g \in B$ l'image de $f$. Soit
$M^\bullet$ un complexe de modules sans $f$-torsion sur $A$. Alors $g$ est un
non-diviseur de zéro, $M^\bullet \otimes_A B$ est un complexe de modules sans
$g$-torsion, et
$\eta_fM^\bullet \otimes_A B = \eta_g(M^\bullet \otimes_A B)$.
\end{lemma}

\begin{proof}
La démonstration est omise.
\end{proof}







\section{Complexes parfaits et opérateur êta}
\label{section-perfect-eta}

\noindent
Dans cette section, nous développons quelques résultats algébriques en vue de
notre version de la construction du graphe de MacPherson ; voir Compléments
sur la platitude, Section \ref{flat-section-blowup-complexes-III}. Nous
utiliserons l'opérateur $\eta_f$ introduit à la Section \ref{section-eta}.

\medskip\noindent
Soit $A$ un anneau, soit $f \in A$ un non-diviseur de zéro et soit
$M^\bullet$ un complexe borné de $A$-modules libres de rang fini. Pour tout
$i$, notons $r_i$ le rang de $M^i$ et posons
$$
I_i(M^\bullet, f) = \text{l'idéal engendré par les mineurs de taille }
r_i \times r_i\text{ de }
(f, d^i) : M^i \to M^i \oplus M^{i + 1}
$$
Remarquons que $f^{r_i} \in I_i(M^\bullet, f)$.

\begin{lemma}
\label{lemma-ideal-well-defined}
Soit $A$ un anneau et soit $f \in A$ un non-diviseur de zéro. Soient
$M^\bullet$ et $N^\bullet$ deux complexes bornés de $A$-modules libres de
rang fini représentant le même objet de $D(A)$. Alors
$$
f^m I_i(M^\bullet, f) = f^n I_i(N^\bullet, f)
$$
comme idéaux de $A$, pour tous entiers $n, m \geq 0$ tels que
$$
m + \sum\nolimits_{j \geq i} (-1)^{j - i}rk(M^j) =
n + \sum\nolimits_{j \geq i} (-1)^{j - i}rk(N^j)
$$
\end{lemma}

\begin{proof}
Il suffit de démontrer l'égalité après localisation en tout idéal premier de
$A$. Ainsi, d'après le Lemme \ref{lemma-compare-representatives-perfect} et
un raisonnement par récurrence que nous omettons, on peut supposer
$N^\bullet = M^\bullet \oplus Q^\bullet$ pour un certain complexe trivial
$Q^\bullet$, c'est-à-dire
$$
Q^\bullet = \ldots \to 0 \to A \xrightarrow{1} A \to 0 \to \ldots
$$
où $A$ est placé aux degrés $j$ et $j + 1$. Si
$j \not = i - 1, i, i + 1$, on a manifestement
$I_i(M^\bullet, f) = I_i(N^\bullet, f)$ et $m = n$, ce qui donne l'égalité
voulue. Si $j = i + 1$, les morphismes
$$
(f, d^i) : M^i \to M^i \oplus M^{i + 1}
\quad\text{and}\quad
(f, d^i, 0) : M^i \to M^i \oplus M^{i + 1} \oplus A
$$
ont les mêmes mineurs non nuls ; dans ce cas encore,
$I_i(M^\bullet, f) = I_i(N^\bullet, f)$ et $m = n$. Si $j = i$, alors
$I_i(M^\bullet, f)$ est l'idéal engendré par les mineurs de taille
$r_i \times r_i$ de
$$
(f, d^i) : M^i \to M^i \oplus M^{i + 1}
$$
et $I_i(N^\bullet, f)$ est l'idéal engendré par les mineurs de taille
$(r_i + 1) \times (r_i + 1)$ de
$$
(f \oplus f, d^i \oplus 1) :
(M^i \oplus A) \to (M^i \oplus A) \oplus (M^{i + 1} \oplus A)
$$
Avec un choix convenable de coordonnées, la matrice du second morphisme est
sous forme diagonale par blocs :
$$
T =
\left(
\begin{matrix}
T_1 & 0 \\
0 & T_2
\end{matrix}
\right),
\quad
T_1 = \text{matrice du premier morphisme},
\quad
T_2 =
\left(
\begin{matrix}
f \\
1
\end{matrix}
\right)
$$
Avec les notations du Lemme \ref{lemma-ideals-generated-by-minors}, on a
$I_0(T_2) = A$, $I_1(T_2) = A$ et $I_p(T_2) = 0$ pour $p \geq 2$. Par
conséquent, $I_{r_i + 1}(T) = I_{r_i + 1}(T_1) + I_{r_i}(T_1) =
I_{r_i}(T_1)$, ce qui signifie que
$I_i(M^\bullet, f) = I_i(N^\bullet, f)$. On a aussi $m = n$, ce qui achève le
cas $j = i$. Enfin, supposons $j = i - 1$. Alors $m = n + 1$ ; il faut donc
montrer que $fI_i(M^\bullet, f) = I_i(N^\bullet, f)$. Dans ce cas,
$I_i(M^\bullet, f)$ est l'idéal engendré par les mineurs de taille
$r_i \times r_i$ de
$$
(f, d^i) : M^i \to M^i \oplus M^{i + 1}
$$
et $I_i(N^\bullet, f)$ est l'idéal engendré par les mineurs de taille
$(r_i + 1) \times (r_i + 1)$ de
$$
(f \oplus f, d^i) : (M^i \oplus A) \to (M^i \oplus A) \oplus M^{i + 1}
$$
Avec un choix convenable de coordonnées, la matrice du second morphisme est
sous forme diagonale par blocs :
$$
T =
\left(
\begin{matrix}
T_1 & 0 \\
0 & T_2
\end{matrix}
\right),
\quad
T_1 = \text{matrice du premier morphisme},
\quad
T_2 =
\left(
\begin{matrix}
f
\end{matrix}
\right)
$$
Le même raisonnement que ci-dessus donne bien
$fI_i(M^\bullet, f) = I_i(N^\bullet, f)$.
\end{proof}

\begin{lemma}
\label{lemma-eta-change-unit}
Soit $f \in A$ un non-diviseur de zéro d'un anneau $A$, soit $u \in A$ une
unité et soit $M^\bullet$ un complexe borné de $A$-modules libres de rang
fini. Alors
$I_i(M^\bullet, f) = I_i(M^\bullet, uf)$.
\end{lemma}

\begin{proof}
La démonstration est omise.
\end{proof}

\begin{lemma}
\label{lemma-eta-base-change-pre}
Soit $A \to B$ un morphisme d'anneaux, soit $f \in A$ un non-diviseur de zéro
et soit $M^\bullet$ un complexe borné de $A$-modules libres de rang fini.
Supposons que $f$ s'envoie sur un non-diviseur de zéro $g$ dans $B$. Alors
$I_i(M^\bullet, f)B = I_i(M^\bullet \otimes_A B, g)$.
\end{lemma}

\begin{proof}
Les mineurs de $(f, d^i) : M^i \to M^i \oplus M^{i + 1}$ s'envoient sur les
mineurs correspondants de
$(g, d^i) : M^i \otimes_A B \to M^i \otimes_A B \oplus M^{i + 1} \otimes_A B$.
\end{proof}

\begin{lemma}
\label{lemma-eta-cohomology-free}
Soit $A$ un anneau, soit $\mathfrak p \subset A$ un idéal premier, soit
$f \in A$ un non-diviseur de zéro et soit $M^\bullet$ un complexe borné de
$A$-modules libres de rang fini. Si $H^i(M^\bullet)_\mathfrak p$ est libre
pour tout $i$, alors $I_i(M^\bullet, f)_\mathfrak p$ est un idéal principal,
et même engendré par une puissance de $f$, pour tout $i$.
\end{lemma}

\begin{proof}
D'après le Lemme \ref{lemma-eta-base-change-pre}, on peut supposer $A$ local
d'idéal maximal $\mathfrak p$. D'après le
Lemme \ref{lemma-ideal-well-defined}, on peut aussi remplacer $M^\bullet$ par
un complexe quasi isomorphe. L'hypothèse de liberté des modules de cohomologie
montre que $M^\bullet$ est quasi isomorphe au complexe dont le terme de degré
$i$ est $H^i(M^\bullet)$ et dont les différentielles sont nulles ; voir par
exemple Catégories dérivées, Lemme \ref{derived-lemma-ext-2-zero-pre}.
Autrement dit, on peut supposer toutes les différentielles de $M^\bullet$
nulles. Il est alors clair que $I_i(M^\bullet, f) = (f^{r_i})$ est principal.
\end{proof}

\begin{lemma}
\label{lemma-eta-locally-free}
Soit $A$ un anneau, soit $f \in A$ un non-diviseur de zéro et soit
$M^\bullet$ un complexe borné de $A$-modules libres de rang fini. Supposons
que $I_i(M^\bullet, f)$ soit un idéal principal. Alors $(\eta_fM)^i$ est
localement libre de rang $r_i$, et le morphisme
$(1, d^i) : (\eta_fM)^i \to f^iM^i \oplus f^{i + 1}M^{i + 1}$
est l'inclusion d'un facteur direct.
\end{lemma}

\begin{proof}
Choisissons un générateur $g$ de $I_i(M^\bullet, f)$. Puisque
$f^{r_i} \in I_i(M^\bullet, f)$, l'élément $g$ divise une puissance de $f$.
En particulier, $g$ est un non-diviseur de zéro dans $A$. Les mineurs de taille
$r_i \times r_i$ du morphisme
$(f, d^i) : M^i \to M^i \oplus M^{i + 1}$
engendrent l'idéal $I_i(M^\bullet, f)$, et les mineurs de taille
$(r_i + 1) \times (r_i + 1)$ de
$(f, d^i)$ sont nuls : on peut le vérifier après localisation en $f$, où le
rang du morphisme est égal à $r_i$. Considérons la surjection
$$
M^i \oplus M^{i + 1}
\longrightarrow
Q = \Coker(f, d^i)/g\text{-torsion}
$$
D'après le Lemme \ref{lemma-fitting-ideals-and-pd1}, le module $Q$ est
localement libre de type fini et de rang $r_{i + 1}$. Ainsi, $Q$ est sans
$f$-torsion, et l'on conclut que le conoyau de $(f, d^i)$ modulo sa torsion
par les puissances de $f$ est lui aussi $Q$.

\medskip\noindent
Considérons le complexe de $A$-modules libres de rang fini
$$
0 \to f^{i + 1}M^i \xrightarrow{1, d^i} f^iM^i \oplus f^{i + 1}M^{i + 1}
\xrightarrow{d^i, -1} f^iM^{i + 1} \to 0
$$
qui devient exact scindé après localisation en $f$. Le morphisme
$(1, d^i) : f^{i + 1}M^i \to f^iM^i \oplus f^{i + 1}M^{i + 1}$ est isomorphe
au morphisme $(f, d^i) : M^i \to M^i \oplus M^{i + 1}$ étudié ci-dessus.
Ainsi, l'image
$$
Q' = \Im(f^iM^i \oplus f^{i + 1}M^{i + 1} \xrightarrow{d^i, -1} f^iM^{i + 1})
$$
est isomorphe à $Q$, donc notamment projective. D'autre part, par construction
de $\eta_f$ à la Section \ref{section-eta}, l'image du morphisme injectif
$(1, d^i) : (\eta_fM)^i \to f^iM^i \oplus f^{i + 1}M^{i + 1}$ is
est le noyau de $(d^i, -1)$. On obtient donc un isomorphisme
$(\eta_fM)^i \oplus Q' = f^iM^i \oplus f^{i + 1}M^{i + 1}$, et l'on voit que
$(\eta_fM)^i$ est bien localement libre de type fini et de rang $r_i$, et que
$(1, d^i)$ est l'inclusion d'un facteur direct.
\end{proof}

\begin{lemma}
\label{lemma-eta-base-change}
Soit $A \to B$ un morphisme d'anneaux, soit $f \in A$ un non-diviseur de zéro
et soit $M^\bullet$ un complexe borné de $A$-modules libres de rang fini.
Supposons que $f$ s'envoie sur un non-diviseur de zéro $g$ dans $B$ et que
$I_i(M^\bullet, f)$ soit un idéal principal pour tout $i \in \mathbf{Z}$.
Il existe alors un isomorphisme canonique
$\eta_fM^\bullet \otimes_A B = \eta_g(M^\bullet \otimes_A B)$.
\end{lemma}

\begin{proof}
Posons $N^i = M^i \otimes_A B$. Remarquons que
$f^iM^i \otimes_A B = g^iN^i$ comme sous-modules de $(N^i)_g$. Les
morphismes
$$
(\eta_fM)^i \otimes_A B \to g^iN^i \oplus g^{i + 1}N^{i + 1}
\quad\text{and}\quad
(\eta_gN)^i \to g^iN^i \oplus g^{i + 1}N^{i + 1}
$$
sont des inclusions de facteurs directs d'après le
Lemme \ref{lemma-eta-locally-free}. Comme leurs images coïncident après
localisation en $g$, on conclut.
\end{proof}

\begin{lemma}
\label{lemma-ideal-direct-summand}
Soit $A$ un anneau. Soient $M$, $N_1$, $N_2$ des $A$-modules projectifs de
type fini, et soit $s : M \to N_1 \oplus N_2$ une injection scindée. Il existe
un idéal de type fini $J \subset A$ ayant la propriété suivante : un morphisme
d'anneaux $A \to B$ se factorise par $A/J$ si et seulement si
$s \otimes \text{id}_B$ identifie $M \otimes_A B$ à une somme directe de
sous-modules de $N_1 \otimes_A B \oplus N_2 \otimes_A B$.
\end{lemma}

\begin{proof}
Choisissons une rétraction $\pi : N_1 \oplus N_2 \to M$ de $s$. Notons
$q_i : N_1 \oplus N_2 \to N_1 \oplus N_2$ le projecteur sur $N_i$, et posons
$p_i = \pi \circ q_i \circ s$. Remarquons que
$p_1 + p_2 = \text{id}_M$. Nous affirmons que $M$ est une somme directe de
sous-modules de $N_1 \oplus N_2$ si et seulement si $p_1$ et $p_2$ sont des
projecteurs orthogonaux. Ainsi, $J$ est le plus petit idéal de $A$ tel que
$p_1 \circ p_1 - p_1$, $p_2 \circ p_2 - p_2$, $p_1 \circ p_2$, and
$p_2 \circ p_1$ appartiennent à $J \otimes_A \text{End}_A(M)$. Certains
détails sont omis.
\end{proof}

\noindent
Soit $A$ un anneau, soit $f \in A$ un non-diviseur de zéro et soit
$M^\bullet$ un complexe borné de $A$-modules libres de rang fini. Supposons
que les idéaux $I_i(M^\bullet, f)$ soient principaux pour tout
$i \in \mathbf{Z}$. Alors les morphismes
$$
(1, d^i) : (\eta_fM)^i / f(\eta_fM)^i \longrightarrow
f^iM^i/f^{i + 1}M^i \oplus f^{i + 1}M^{i + 1}/f^{i + 2}M^{i + 1}
$$
sont des injections scindées d'après le Lemme
\ref{lemma-eta-locally-free}. Notons
$J_i(M^\bullet, f) \subset A/fA$ l'idéal de type fini du
Lemme \ref{lemma-ideal-direct-summand} associé à l'injection scindée
$(1, d^i)$ affichée ci-dessus.

\begin{lemma}
\label{lemma-eta-vanishing-beta-plus-pre}
Soit $A$ un anneau et soit $f \in A$ un non-diviseur de zéro. Soient
$M^\bullet$ et $N^\bullet$ deux complexes bornés de $A$-modules libres de
rang fini représentant le même objet de $D(A)$. Supposons que
$I_i(M^\bullet, f)$ soit un idéal principal pour tout $i \in \mathbf{Z}$.
Alors $J_i(M^\bullet, f) = J_i(N^\bullet, f)$ comme idéaux de $A/fA$.
\end{lemma}

\begin{proof}
Remarquons que le fait que $I_i(M^\bullet, f)$ soit principal implique, par le
Lemme \ref{lemma-ideal-well-defined}, que $I_i(N^\bullet, f)$ est principal ;
l'énoncé a donc bien un sens. Comme dans la démonstration du
Lemme \ref{lemma-ideal-well-defined}, on peut supposer
$N^\bullet = M^\bullet \oplus Q^\bullet$ pour un certain complexe trivial
$Q^\bullet$, c'est-à-dire
$$
Q^\bullet = \ldots \to 0 \to A \xrightarrow{1} A \to 0 \to \ldots
$$
où $A$ est placé aux degrés $j$ et $j + 1$. Comme $\eta_f$ est compatible aux
sommes directes, le morphisme
$$
(1, d^i) : (\eta_fN)^i / f(\eta_fN)^i \longrightarrow
f^iN^i/f^{i + 1}N^i \oplus f^{i + 1}N^{i + 1}/f^{i + 2}N^{i + 1}
$$
est la somme directe des morphismes correspondants pour $M^\bullet$ et pour
$Q^\bullet$. Par la propriété universelle qui définit les idéaux en question,
on en conclut que
$J_i(N^\bullet, f) = J_i(M^\bullet, f) + J_i(Q^\bullet, f)$. Il suffit donc
de montrer que $J_i(Q^\bullet, f) = 0$ pour tout $i$. Nous omettons ce calcul.
\end{proof}

\begin{lemma}
\label{lemma-eta-vanishing-beta-plus}
Soit $A$ un anneau, soit $f \in A$ un non-diviseur de zéro et soit
$M^\bullet$ un complexe borné de $A$-modules libres de rang fini. Supposons
que $I_i(M^\bullet, f)$ soit un idéal principal pour tout $i \in \mathbf{Z}$.
Considérons l'idéal $J(M^\bullet, f) = \sum_i J_i(M^\bullet, f)$ de $A/fA$ et
l'ensemble d'idéaux premiers
\begin{align*}
E
& =
\{f \in \mathfrak p \subset A \mid
\Ker(d^i \bmod f^2)_\mathfrak p \text{ se surjecte sur }
\Ker(d^i \bmod f)_\mathfrak p \text{ pour tout }i \in \mathbf{Z}\} \\
& =
\{f \in \mathfrak p \subset A \mid
\text{les localisés }\beta_\mathfrak p
\text{ des opérateurs de Bockstein sont nuls}\}
\end{align*}
Alors :
\begin{enumerate}
\item $J(M^\bullet, f)$ est de type fini ;
\item $A/fA \to C = (A/fA)/J(M^\bullet, f)$
est surjectif et de présentation finie ;
\item $J(M^\bullet, f)_\mathfrak p = 0$ pour $\mathfrak p \in E$ ;
\item si $f \in \mathfrak p$ et si $H^i(M^\bullet)_\mathfrak p$ est libre
pour tout $i \in \mathbf{Z}$, alors $\mathfrak p \in E$ ;
\item les modules de cohomologie de $\eta_f M^\bullet \otimes_A C$ sont des
$C$-modules localement libres de type fini.
\end{enumerate}
\end{lemma}

\begin{proof}
L'égalité dans la définition de $E$ résulte du
Lemme \ref{lemma-vanishing-beta} ; de plus, la dernière assertion de ce lemme
implique (4).

\medskip\noindent
L'assertion (1) est vraie parce que les idéaux $J_i(M^\bullet, f)$ sont de
type fini et parce que $M^\bullet$ est borné, de sorte que
$J_i(M^\bullet, f)$ est nul pour presque tout $i$. L'assertion (2) n'est
qu'une reformulation de (1).

\medskip\noindent
Démonstration de (3). D'après le Lemme \ref{lemma-eta-locally-free},
$(\eta_fM)^i$ est localement libre de type fini et de rang $r_i$ pour tout
$i$. Considérons le morphisme
$$
(1, d^i) :
(\eta_fM)^i / f(\eta_fM)^i
\longrightarrow
f^iM^i/f^{i + 1}M^i \oplus f^{i + 1}M^{i + 1}/f^{i + 2}M^{i + 1}
$$
Choisissons $\mathfrak p \in E$. D'après le
Lemme \ref{lemma-eta-vanishing-beta} et la liberté locale des modules
$(\eta_fM)^i$, on peut écrire
$$
\left((\eta_fM)^i / f(\eta_fM)^i\right)_\mathfrak p =
(A/fA)_\mathfrak p^{\oplus m_i} \oplus (A/fA)_\mathfrak p^{\oplus n_i}
$$
de manière compatible avec la flèche $(1, d^i)$ ci-dessus. Par la propriété
universelle de l'idéal $J_i(M^\bullet, f)$, on conclut que
$J_i(M^\bullet, f)_\mathfrak p = 0$. Ainsi,
$J(M^\bullet, f)_\mathfrak p = 0$ pour $\mathfrak p \in E$.

\medskip\noindent


\medskip\noindent
Démonstration de (5). Remarquons que la différentielle de
$\eta_fM^\bullet$ s'insère dans un diagramme commutatif
$$
\xymatrix{
(\eta_fM)^i \ar[d] \ar[r] &
f^iM^i \oplus f^{i + 1}M^{i + 1}
\ar[d]^{\left(
\begin{matrix}
0 & 1 \\
0 & 0
\end{matrix}
\right)} \\
(\eta_fM)^{i + 1} \ar[r] &
f^{i + 1}M^{i + 1} \oplus f^{i + 2}M^{i + 2}
}
$$
Par construction, après tensorisation par $C$, les modules de gauche sont des
sommes directes de facteurs directs des facteurs de droite. On obtient le
diagramme
$$
\xymatrix{
(\eta_fM)^i \otimes_A C \ar[d] \ar@{=}[r] &
K^i \oplus L^i \ar[r] \ar[d] &
f^iM^i \otimes_A C \oplus f^{i + 1}M^{i + 1} \otimes_A C
\ar[d]^{\left(
\begin{matrix}
0 & 1 \\
0 & 0
\end{matrix}
\right)} \\
(\eta_fM)^{i + 1} \otimes_A C \ar@{=}[r] &
K^{i + 1} \oplus L^{i + 1} \ar[r] &
f^{i + 1}M^{i + 1} \otimes_A C \oplus f^{i + 2}M^{i + 2} \otimes_A C
}
$$
où les flèches horizontales sont compatibles à la fois aux décompositions en
sommes directes et aux inclusions de facteurs directs. Il s'ensuit que la
différentielle identifie $L^i$ à un facteur direct de $K^{i + 1}$ ; on
conclut que la cohomologie de $\eta_fM^\bullet \otimes_A C$ au degré $i$ est
le module $K^i/L^{i - 1}$, qui est projectif de type fini, comme voulu.
\end{proof}









\section{Limites de complexes}
\label{section-limits}

\noindent
Dans cette section, nous étudions ce qui se passe lorsqu'on dispose d'une
« déformation formelle » d'un complexe et qu'on en prend la limite. Nous
considérerons deux cas :
\begin{enumerate}
\item on a une limite $A = \lim A_n$ d'un système projectif d'anneaux dont
les morphismes de transition sont surjectifs à noyaux localement nilpotents,
et des objets $K_n \in D(A_n)$ qui se recollent au sens où
$K_n = K_{n + 1} \otimes_{A_{n + 1}}^\mathbf{L} A_n$ ;
\item on a un anneau $A$, un idéal $I$ et des objets $K_n \in D(A/I^n)$ qui
se recollent au sens où
$K_n = K_{n + 1} \otimes_{A/I^{n + 1}}^\mathbf{L} A/I^n$.
\end{enumerate}
Sous des hypothèses supplémentaires, on peut montrer que $K = R\lim K_n$
reproduit le système, au sens où $K_n = K \otimes_A^\mathbf{L} A_n$ ou
$K_n = K \otimes_A^\mathbf{L} A/I^n$.

\begin{lemma}
\label{lemma-Rlim-pseudo-coherent-gives-pseudo-coherent}
Soit $A = \lim A_n$ la limite d'un système projectif d'anneaux $(A_n)$.
Supposons donnés $K_n \in D(A_n)$ et des morphismes
$K_{n + 1} \to K_n$ dans $D(A_{n + 1})$. Supposons que :
\begin{enumerate}
\item les morphismes de transition $A_{n + 1} \to A_n$ soient surjectifs à
noyaux localement nilpotents ;
\item soit tous les $K_n$ soient pseudo-cohérents, soit il existe un entier
$n_0$ tel que $K_{n_0}$ soit pseudo-cohérent et que les noyaux de
$A_{n + 1} \to A_n$ soient des idéaux nilpotents pour $n \geq n_0$ ;
\item les morphismes induisent des isomorphismes
$K_{n + 1} \otimes_{A_{n + 1}}^\mathbf{L} A_n \to K_n$.
\end{enumerate}
Alors $K = R\lim K_n$ est un objet pseudo-cohérent de $D(A)$, et
$K \otimes_A^\mathbf{L} A_n \to K_n$ est un isomorphisme pour tout $n$.
\end{lemma}

\begin{proof}
Par hypothèse, on peut trouver un complexe borné supérieurement de
$A_1$-modules libres de rang fini, noté $P_1^\bullet$, représentant $K_1$ ; voir
la Définition \ref{definition-pseudo-coherent}. Le
Lemme \ref{lemma-check-pseudo-coherent-modulo-nilpotent} montre que $K_n$ est
pseudo-cohérent pour tout $n$. Puis, d'après le
Lemme \ref{lemma-lift-complex-stably-frees}, on peut construire par récurrence
sur $n > 1$ des complexes $P_n^\bullet$ de $A_n$-modules libres de rang fini
représentant $K_n$, ainsi que des morphismes
$P_n^\bullet \to P_{n - 1}^\bullet$ représentant les morphismes
$K_n \to K_{n - 1}$ et induisant des isomorphismes, au niveau même des
complexes, $P_n^\bullet \otimes_{A_n} A_{n - 1} \to P_{n - 1}^\bullet$. Ainsi,
$K = R\lim K_n$ est représenté par $P^\bullet = \lim P_n^\bullet$ ; voir le
Lemme \ref{lemma-compute-Rlim-modules} et la Remarque
\ref{remark-how-unique}. Comme $P_n^i$ est un $A_n$-module libre de rang fini
pour tout $n$ et comme $A = \lim A_n$, le module $P^i$ est libre de rang fini,
du même rang que $P_1^i$, pour tout $i$. Cela signifie que $K$ est
pseudo-cohérent. En outre, $K \otimes_A^\mathbf{L} A_n$ est représenté par
$P^\bullet \otimes_A A_n = P_n^\bullet$, ce qui prouve la dernière assertion.
\end{proof}

\begin{lemma}
\label{lemma-Rlim-pseudo-coherent-gives-complete-pseudo-coherent}
Soit $A$ un anneau et soit $I \subset A$ un idéal. Supposons donnés
$K_n \in D(A/I^n)$ et des morphismes $K_{n + 1} \to K_n$ dans
$D(A/I^{n + 1})$. Supposons que :
\begin{enumerate}
\item $A$ soit $I$-adiquement complet ;
\item $K_1$ soit pseudo-cohérent ;
\item les morphismes induisent des isomorphismes
$K_{n + 1} \otimes_{A/I^{n + 1}}^\mathbf{L} A/I^n \to K_n$.
\end{enumerate}
Alors $K = R\lim K_n$ est un objet pseudo-cohérent et complet au sens dérivé
de $D(A)$, et $K \otimes_A^\mathbf{L} A/I^n \to K_n$ est un isomorphisme pour
tout $n$.
\end{lemma}

\begin{proof}
Le Lemme \ref{lemma-Rlim-pseudo-coherent-gives-pseudo-coherent} montre déjà
que $K$ est pseudo-cohérent et que
$K \otimes_A^\mathbf{L} A/I^n \to K_n$ est un isomorphisme pour tout $n$.
Enfin, $K$ est complet au sens dérivé d'après le
Lemme \ref{lemma-naive-derived-completion}.
\end{proof}

\begin{lemma}
\label{lemma-Rlim-perfect-gives-perfect}
\begin{reference}
\cite[Lemma 4.2]{Bhatt-Algebraize}
\end{reference}
Soit $A = \lim A_n$ la limite d'un système projectif d'anneaux $(A_n)$.
Supposons donnés $K_n \in D(A_n)$ et des morphismes
$K_{n + 1} \to K_n$ dans $D(A_{n + 1})$. Supposons que :
\begin{enumerate}
\item les morphismes de transition $A_{n + 1} \to A_n$ soient surjectifs à
noyaux localement nilpotents ;
\item soit tous les $K_n$ soient parfaits, soit il existe un entier $n_0$ tel
que $K_{n_0}$ soit parfait et que les noyaux de $A_{n + 1} \to A_n$ soient
nilpotents pour $n \geq n_0$ ;
\item les morphismes induisent des isomorphismes
$K_{n + 1} \otimes_{A_{n + 1}}^\mathbf{L} A_n \to K_n$.
\end{enumerate}
Alors $K = R\lim K_n$ est un objet parfait de $D(A)$, et
$K \otimes_A^\mathbf{L} A_n \to K_n$ est un isomorphisme pour tout $n$.
\end{lemma}

\begin{proof}
Le Lemme \ref{lemma-Rlim-pseudo-coherent-gives-pseudo-coherent} montre déjà
que $K$ est pseudo-cohérent et que $K \otimes_A^\mathbf{L} A_n \to K_n$ est
un isomorphisme pour tout $n$. Considérons un morphisme surjectif
$A \to \kappa$ dont le noyau est un idéal maximal $\mathfrak m$. Tout élément
de $A$ dont l'image est une unité de $A_1$ est une unité de $A$ d'après
Algèbre, Lemme \ref{algebra-lemma-locally-nilpotent-unit}. Ainsi,
$\Ker(A \to A_1)$ est contenu dans le radical de Jacobson de $A$ d'après
Algèbre, Lemme \ref{algebra-lemma-contained-in-radical}. Par conséquent,
$A \to \kappa$ se factorise sous la forme $A \to A_1 \to \kappa$. Dès lors,
$$
K \otimes_A^\mathbf{L} \kappa =
K \otimes_A^\mathbf{L} A_1 \otimes_{A_1}^\mathbf{L} \kappa =
K_1 \otimes_{A_1}^\mathbf{L} \kappa
$$
Remarquons que $K_1$ est parfait par l'hypothèse (2). Le complexe $K_1$ est
donc de
dimension Tor finie d'après le Lemme \ref{lemma-perfect}. Il existe ainsi
$a, b \in \mathbf{Z}$ tels que
$H^i(K \otimes_A^\mathbf{L} \kappa) = 0$ pour tout $i \not \in [a, b]$. Le
Lemme \ref{lemma-check-perfect-pointwise} permet de conclure que $K$ est
parfait.
\end{proof}

\begin{lemma}
\label{lemma-Rlim-perfect-gives-complete}
Soit $A$ un anneau et soit $I \subset A$ un idéal. Supposons donnés
$K_n \in D(A/I^n)$ et des morphismes $K_{n + 1} \to K_n$ dans
$D(A/I^{n + 1})$. Supposons que :
\begin{enumerate}
\item $A$ soit $I$-adiquement complet ;
\item $K_1$ soit un objet parfait ;
\item les morphismes induisent des isomorphismes
$K_{n + 1} \otimes_{A/I^{n + 1}}^\mathbf{L} A/I^n \to K_n$.
\end{enumerate}
Alors $K = R\lim K_n$ est un objet parfait et complet au sens dérivé de
$D(A)$, et $K \otimes_A^\mathbf{L} A/I^n \to K_n$ est un isomorphisme pour
tout $n$.
\end{lemma}

\begin{proof}
Combiner les Lemmes \ref{lemma-Rlim-perfect-gives-perfect} et
\ref{lemma-Rlim-pseudo-coherent-gives-complete-pseudo-coherent} (le second
donne la complétude dérivée).
\end{proof}

\noindent
Nous ignorons si le lemme suivant vaut pour les complexes non bornés.

\begin{lemma}
\label{lemma-Rlim-gives-complete}
Soit $A$ un anneau et soit $I \subset A$ un idéal. Supposons donnés
$K_n \in D(A/I^n)$ et des morphismes $K_{n + 1} \to K_n$ dans
$D(A/I^{n + 1})$. Si :
\begin{enumerate}
\item $A$ est noethérien ;
\item $K_1$ est borné supérieurement ;
\item les morphismes induisent des isomorphismes
$K_{n + 1} \otimes_{A/I^{n + 1}}^\mathbf{L} A/I^n \to K_n$,
\end{enumerate}
alors $K = R\lim K_n$ est un objet complet au sens dérivé de $D^-(A)$, et
$K \otimes_A^\mathbf{L} A/I^n \to K_n$ est un isomorphisme pour tout $n$.
\end{lemma}

\begin{proof}
L'objet $K$ de $D(A)$ est complet au sens dérivé d'après le
Lemme \ref{lemma-naive-derived-completion}.

\medskip\noindent
Supposons $H^i(K_1) = 0$ pour $i > b$. On peut alors trouver un complexe de
$A/I$-modules libres, noté $P_1^\bullet$, représentant $K_1$, avec $P_1^i = 0$
pour $i > b$. D'après le Lemme \ref{lemma-lift-complex-projectives}, on peut
construire par récurrence sur $n > 1$ des complexes $P_n^\bullet$ de
$A/I^n$-modules libres représentant $K_n$, ainsi que des morphismes
$P_n^\bullet \to P_{n - 1}^\bullet$ représentant les morphismes
$K_n \to K_{n - 1}$ et induisant des isomorphismes, au niveau même des
complexes, $P_n^\bullet/I^{n - 1}P_n^\bullet \to P_{n - 1}^\bullet$.

\medskip\noindent
Nous sommes ainsi dans la situation où $R\lim K_n$ est représenté par
$P^\bullet = \lim P_n^\bullet$ ; voir le
Lemme \ref{lemma-compute-Rlim-modules} et la Remarque
\ref{remark-how-unique}. Les complexes $P_n^\bullet$ sont uniformément bornés
supérieurement, formés de $A/I^n$-modules plats, et leurs morphismes de
transition sont surjectifs terme à terme. Le Lemme \ref{lemma-limit-flat}
montre alors que $P^\bullet$ est un complexe borné supérieurement de
$A$-modules plats. Par conséquent, $K \otimes_A^\mathbf{L} A/I^t$ est
représenté par $P^\bullet \otimes_A A/I^t$. D'après le
Lemme \ref{lemma-limit-flat}, on a terme à terme
$P^\bullet \otimes_A A/I^t = \lim P_n^\bullet \otimes_A A/I^t$. Les morphismes
de transition
$P_{n + 1}^\bullet \otimes_A A/I^t \to P_n^\bullet \otimes_A A/I^t$ sont des
isomorphismes pour $n \geq t$ par le choix des $P_n^\bullet$. Ainsi,
$\lim P_n^\bullet \otimes_A A/I^t = P_t^\bullet \otimes_A A/I^t = P_t^\bullet$.
Comme $P_t^\bullet$ représente $K_t$, le morphisme
$K \otimes_A^\mathbf{L} A/I^t \to K_t$ est un isomorphisme.
\end{proof}

\noindent
Voici un résultat d'un autre type.

\begin{lemma}[Koll\'ar-Kov\'acs]
\label{lemma-kollar-kovacs}
\begin{reference}
Courriel de Kovács du 23/02/2018.
\end{reference}
Soit $I$ un idéal d'un anneau noethérien $A$ et soit $K \in D(A)$. Posons
$K_n = K \otimes_A^\mathbf{L} A/I^n$. Supposons que, pour tout
$i \in \mathbf{Z}$ :
\begin{enumerate}
\item $H^i(K)$ soit un $A$-module de type fini ;
\item le système $H^i(K_n)$ vérifie la condition de Mittag--Leffler.
\end{enumerate}
Alors $\lim H^i(K)/I^nH^i(K)$ est égal à $\lim H^i(K_n)$ pour tout
$i \in \mathbf{Z}$.
\end{lemma}

\begin{proof}
Rappelons que $K^\wedge = R\lim K_n$ est la complétion dérivée de $K$ ; voir
la Proposition \ref{proposition-noetherian-naive-completion-is-completion}.
D'après le Lemme \ref{lemma-derived-completion-pseudo-coherent}, on a
$H^i(K^\wedge) = \lim H^i(K)/I^nH^i(K)$. Le
Lemme \ref{lemma-break-long-exact-sequence-modules} fournit des suites exactes
courtes
$$
0 \to R^1\lim H^{i - 1}(K_n) \to H^i(K^\wedge) \to \lim H^i(K_n) \to 0
$$
La condition de Mittag--Leffler garantit que les termes de gauche sont nuls
(Lemme \ref{lemma-compute-Rlim-modules}), ce qui démontre le lemme.
\end{proof}



\section{Quelques morphismes d'évaluation}
\label{section-evaluation}

\noindent
Dans cette section, nous montrons que certains morphismes canoniques entre
$R\Hom$ sont des isomorphismes pour des classes convenables de complexes.

\begin{lemma}
\label{lemma-internal-hom-evaluate-isomorphism}
Soit $R$ un anneau et soient $K, L, M$ des objets de $D(R)$. Le morphisme
$$
R\Hom_R(L, M) \otimes_R^\mathbf{L} K \longrightarrow R\Hom_R(R\Hom_R(K, L), M)
$$
du Lemme \ref{lemma-internal-hom-evaluate} est un isomorphisme dans les deux
cas suivants :
\begin{enumerate}
\item $K$ est parfait ;
\item $K$ est pseudo-cohérent, $L \in D^+(R)$ et $M$ est de dimension
injective finie.
\end{enumerate}
\end{lemma}

\begin{proof}
Choisissons un complexe K-injectif $I^\bullet$ représentant $M$, un complexe
K-injectif $J^\bullet$ représentant $L$ et un complexe borné supérieurement
$K^\bullet$ de modules projectifs de type fini représentant $K$. Considérons
le morphisme de complexes
$$
\text{Tot}(\Hom^\bullet(J^\bullet, I^\bullet) \otimes_R K^\bullet)
\longrightarrow
\Hom^\bullet(\Hom^\bullet(K^\bullet, J^\bullet), I^\bullet)
$$
du Lemme \ref{lemma-evaluate-and-more}. Remarquons que
$$
\left(\prod\nolimits_{p + r = t} \Hom_R(J^{-r}, I^p)\right) \otimes_R K^s =
\prod\nolimits_{p + r = t} \Hom_R(J^{-r}, I^p) \otimes_R K^s
$$
parce que $K^s$ est projectif de type fini. Le morphisme est donné par les
morphismes
$$
c_{p, r, s} :
\Hom_R(J^{-r}, I^p) \otimes_R K^s
\longrightarrow
\Hom_R(\Hom_R(K^s, J^{-r}), I^p)
$$
qui sont des isomorphismes puisque $K^s$ est projectif de type fini. Pour tout
élément $\alpha = (\alpha^{p, r, s})$ de degré $n$ du membre de gauche, il n'y
a qu'un nombre fini de valeurs de $s$ pour lesquelles
$\alpha^{p, r, s}$ est non nul (pour certains $p, r$ tels que
$n = p + r + s$). Notre morphisme est donc un isomorphisme si la même
condition d'annulation est imposée aux éléments
$\beta = (\beta^{p, r, s})$ du membre de droite. Si $K^\bullet$ est un
complexe borné de modules projectifs de type fini, cela est clair. D'autre
part, si l'on peut choisir $I^\bullet$ borné et $J^\bullet$ borné
inférieurement, alors $\beta^{p, r, s}$ est nul pour $p$ hors d'un intervalle
fixe, pour $s \gg 0$ et pour $r \gg 0$. Parmi les solutions de
$n = p + r + s$ telles que $\beta^{p, r, s}$ soit non nul, il n'apparaît donc
qu'un nombre fini de valeurs de $s$.
\end{proof}

\begin{lemma}
\label{lemma-internal-hom-evaluate-isomorphism-technical}
Soit $R$ un anneau et soient $K, L, M$ des objets de $D(R)$. Le morphisme
$$
R\Hom_R(L, M) \otimes_R^\mathbf{L} K \longrightarrow R\Hom_R(R\Hom_R(K, L), M)
$$
du Lemme \ref{lemma-internal-hom-evaluate} est un isomorphisme si les trois
conditions suivantes sont satisfaites :
\begin{enumerate}
\item $L, M$ sont de dimension injective finie ;
\item $R\Hom_R(L, M)$ est de dimension Tor finie ;
\item pour tout $n \in \mathbf{Z}$, la troncation $\tau_{\leq n}K$ est
pseudo-cohérente.
\end{enumerate}
\end{lemma}

\begin{proof}
Choisissons un entier $n$ et considérons le triangle distingué
$$
\tau_{\leq n}K \to K \to \tau_{\geq n + 1}K \to \tau_{\leq n}K[1]
$$
voir Catégories dérivées, Remarque
\ref{derived-remark-truncation-distinguished-triangle}.
Par l'hypothèse (3) et le Lemme
\ref{lemma-internal-hom-evaluate-isomorphism}, le morphisme est un isomorphisme
pour $\tau_{\leq n}K$. Il suffit donc de montrer que les deux objets
$$
R\Hom_R(L, M) \otimes_R^\mathbf{L} \tau_{\geq n + 1}K
\quad\text{and}\quad
R\Hom_R(R\Hom_R(\tau_{\geq n + 1}K, L), M)
$$
ont une cohomologie nulle aux degrés $\leq n - c$ pour un certain $c$. Cela
résulte immédiatement des hypothèses (2) et (1).
\end{proof}

\begin{lemma}
\label{lemma-internal-hom-evaluate-tensor-isomorphism}
Soit $R$ un anneau et soient $K, L, M$ des objets de $D(R)$. Le morphisme
$$
K \otimes_R^\mathbf{L} R\Hom_R(M, L) \longrightarrow
R\Hom_R(M, K \otimes_R^\mathbf{L} L)
$$
du Lemme \ref{lemma-internal-hom-diagonal-better} est un isomorphisme dans les
cas suivants :
\begin{enumerate}
\item $M$ est parfait ;
\item $K$ est parfait ;
\item $M$ est pseudo-cohérent, $L \in D^+(R)$ et $K$ est d'amplitude Tor
contenue dans $[a, \infty]$.
\end{enumerate}
\end{lemma}

\begin{proof}
Démonstration lorsque $M$ est parfait. Remarquons que les deux membres de la
flèche transforment les triangles distingués en $M$ en triangles distingués
et commutent aux sommes directes. Il suffit donc de vérifier l'assertion pour
$M = R[n]$ ; voir Catégories dérivées, Remarque
\ref{derived-remark-check-on-generator}, et le
Lemme \ref{lemma-perfect-ring-classical-generator}. Dans ce cas, le résultat
est évident.

\medskip\noindent
Démonstration lorsque $K$ est parfait. On applique le même argument que dans
le cas précédent.

\medskip\noindent
Démonstration dans le cas (3). On peut représenter $K$ et $L$ par des
complexes bornés inférieurement de $R$-modules $K^\bullet$ et $L^\bullet$. On
peut supposer que $K^\bullet$ est un complexe K-plat formé de $R$-modules
plats ; voir le Lemme \ref{lemma-bounded-below-tor-amplitude}. On peut
représenter $M$ par un complexe borné supérieurement $M^\bullet$ de
$R$-modules libres de rang fini ; voir la Définition
\ref{definition-pseudo-coherent}. Le membre de gauche est alors représenté
par
$$
\text{Tot}(K^\bullet \otimes_R \Hom^\bullet(M^\bullet, L^\bullet))
$$
et le membre de droite par
$$
\Hom^\bullet(M^\bullet, \text{Tot}(K^\bullet \otimes_R L^\bullet))
$$
On utilise ici le Lemme \ref{lemma-RHom-out-of-projective}. Les deux complexes
ont au degré $n$ le module
$$
\bigoplus\nolimits_{p + q + r = n} K^p \otimes \Hom_R(M^{-r}, L^q) =
\bigoplus\nolimits_{p + q + r = n} \Hom_R(M^{-r}, K^p \otimes_R L^q)
$$
parce que $M^{-r}$ est libre de rang fini (et que ces sommes directes sont
finies). Le morphisme défini dans le
Lemme \ref{lemma-internal-hom-diagonal-better} provient du morphisme de
complexes défini dans le Lemme \ref{lemma-diagonal-better}, lequel utilise les
isomorphismes canoniques entre ces modules.
\end{proof}

\begin{lemma}
\label{lemma-hom-complex-K-flat}
Soit $R$ un anneau. Soit $P^\bullet$ un complexe borné supérieurement de
$R$-modules projectifs et soit $K^\bullet$ un complexe K-plat de $R$-modules.
Si $P^\bullet$ est un objet parfait de $D(R)$, alors
$\Hom^\bullet(P^\bullet, K^\bullet)$ est K-plat et représente
$R\Hom_R(P^\bullet, K^\bullet)$.
\end{lemma}

\begin{proof}
La dernière assertion est le Lemme \ref{lemma-RHom-out-of-projective}.
Puisque $P^\bullet$ représente un objet parfait, il existe un complexe fini de
$R$-modules projectifs de type fini, noté $F^\bullet$, tel que $P^\bullet$ et
$F^\bullet$ soient isomorphes dans $D(R)$ ; voir la Définition
\ref{definition-perfect}. Les complexes $P^\bullet$ et $F^\bullet$ sont alors
homotopiquement équivalents ; voir
Catégories dérivées, Lemme
\ref{derived-lemma-morphisms-from-projective-complex}. Par conséquent,
$\Hom^\bullet(P^\bullet, K^\bullet)$ et
$\Hom^\bullet(F^\bullet, K^\bullet)$ sont homotopiquement équivalents. Le
premier est donc K-plat si et seulement si le second l'est (cela résulte de la
Définition \ref{definition-K-flat} et du
Lemme \ref{lemma-derived-tor-homotopy}). Il est clair que
$$
\Hom^\bullet(F^\bullet, K^\bullet) =
\text{Tot}(E^\bullet \otimes_R K^\bullet)
$$
où $E^\bullet$ est le complexe dual de $F^\bullet$, de termes
$E^n = \Hom_R(F^{-n}, R)$ ; voir le Lemme
\ref{lemma-dual-perfect-complex} et sa démonstration. Comme $E^\bullet$ est un
complexe borné de modules projectifs, il est K-plat d'après le
Lemme \ref{lemma-derived-tor-quasi-isomorphism}. On conclut alors par le
Lemme \ref{lemma-tensor-product-K-flat}.
\end{proof}






\section{Changement de base pour le Hom dérivé}
\label{section-base-change-RHom}

\noindent
Nous avons déjà rencontré des résultats sur cette question dans le
Lemme \ref{lemma-pseudo-coherence-and-base-change-ext} et dans Algèbre,
Section \ref{algebra-section-functoriality-ext}.

\begin{lemma}
\label{lemma-upgrade-adjoint-tensor-RHom}
Soit $R \to R'$ un morphisme d'anneaux. Pour $K \in D(R)$ et $M \in D(R')$,
il existe un isomorphisme canonique
$$
R\Hom_R(K, M) = R\Hom_{R'}(K \otimes_R^\mathbf{L} R', M)
$$
\end{lemma}

\begin{proof}
Choisissons un complexe K-injectif de $R'$-modules, noté $J^\bullet$,
représentant $M$, un quasi-isomorphisme $J^\bullet \to I^\bullet$ où
$I^\bullet$ est un
complexe K-injectif de $R$-modules, et un complexe K-plat $K^\bullet$ de
$R$-modules représentant $K$. Considérons le morphisme
$$
\Hom^\bullet(K^\bullet \otimes_R R', J^\bullet)
\longrightarrow
\Hom^\bullet(K^\bullet, I^\bullet)
$$
Le morphisme entre les termes de degré $n$ est donné par
$$
\prod\nolimits_{n = p + q} \Hom_{R'}(K^{-q} \otimes_R R', J^p)
\longrightarrow
\prod\nolimits_{n = p + q} \Hom_R(K^{-q}, I^p)
$$
obtenu en précomposant par $K^{-q} \to K^{-q} \otimes_R R'$ et en
postcomposant par $J^p \to I^p$. Pour achever la démonstration, il suffit de
montrer que l'on obtient des isomorphismes sur les groupes de cohomologie :
$$
\Hom_{D(R)}(K, M) = \Hom_{D(R')}(K \otimes_R^\mathbf{L} R', M)
$$
ce qui est vrai parce que le changement de base
$- \otimes_R^\mathbf{L} R' : D(R) \to D(R')$ est adjoint à gauche du foncteur
de restriction $D(R') \to D(R)$ d'après le
Lemme \ref{lemma-tensor-hom-adjoint}.
\end{proof}

\noindent
Soit $R \to R'$ un morphisme d'anneaux. On a un morphisme de changement de
base
\begin{equation}
\label{equation-base-change-RHom}
R\Hom_R(K, M) \otimes_R^\mathbf{L} R'
\longrightarrow
R\Hom_{R'}(K \otimes_R^\mathbf{L} R', M \otimes_R^\mathbf{L} R')
\end{equation}
dans $D(R')$, fonctoriel en $K, M \in D(R)$. En effet, par l'adjonction entre
$- \otimes_R^\mathbf{L} R' : D(R) \to D(R')$ et le foncteur de restriction
$D(R') \to D(R)$, cela revient à donner un morphisme
$$
R\Hom_R(K, M)
\longrightarrow
R\Hom_{R'}(K \otimes_R^\mathbf{L} R', M \otimes_R^\mathbf{L} R') =
R\Hom_R(K, M \otimes_R^\mathbf{L} R')
$$
(l'égalité venant du Lemme \ref{lemma-upgrade-adjoint-tensor-RHom}), pour
lequel on peut utiliser le morphisme canonique
$M \to M \otimes_R^\mathbf{L} R'$ (l'unité de l'adjonction).

\begin{lemma}
\label{lemma-base-change-RHom}
Soit $R \to R'$ un morphisme d'anneaux et soient $K, M \in D(R)$. Le
morphisme
(\ref{equation-base-change-RHom})
$$
R\Hom_R(K, M) \otimes_R^\mathbf{L} R'
\longrightarrow
R\Hom_{R'}(K \otimes_R^\mathbf{L} R', M \otimes_R^\mathbf{L} R')
$$
est un isomorphisme dans $D(R')$ dans les cas suivants :
\begin{enumerate}
\item $K$ est parfait ;
\item $R'$ est parfait comme $R$-module ;
\item $R \to R'$ est plat, $K$ est pseudo-cohérent et $M \in D^{+}(R)$ ;
\item $R'$ est de dimension Tor finie comme $R$-module, $K$ est
pseudo-cohérent et $M \in D^{+}(R)$.
\end{enumerate}
\end{lemma}

\begin{proof}
On peut vérifier que le morphisme est un isomorphisme après application du
foncteur de restriction $D(R') \to D(R)$. Après application de ce foncteur,
notre morphisme devient
$$
R\Hom_R(K, M) \otimes_R^\mathbf{L} R'
\longrightarrow
R\Hom_R(K, M \otimes_R^\mathbf{L} R')
$$
du Lemme \ref{lemma-internal-hom-diagonal-better}. Voir la discussion qui
précède le lemme pour identifier les membres de gauche et de droite ; on y
utilise notamment le Lemme \ref{lemma-upgrade-adjoint-tensor-RHom}. On conclut
alors par le Lemme \ref{lemma-internal-hom-evaluate-tensor-isomorphism}.
\end{proof}





\section{Systèmes de modules}
\label{section-systems}

\noindent
Soit $I$ un idéal d'un anneau noethérien $A$. Dans cette section, nous
complétons notre étude de la relation entre les $A$-modules de type fini et les
systèmes de $A/I^n$-modules de type fini.

\begin{lemma}
\label{lemma-consequence-Artin-Rees}
Soit $I$ un idéal d'un anneau noethérien $A$. Soit
$
K \xrightarrow{\alpha} L \xrightarrow{\beta} M
$
un complexe de $A$-modules de type fini. Posons
$H = \Ker(\beta)/\Im(\alpha)$. Pour $n \geq 0$, soit
$$
K/I^nK \xrightarrow{\alpha_n} L/I^nL \xrightarrow{\beta_n} M/I^nM
$$
le complexe induit et posons $H_n = \Ker(\beta_n)/\Im(\alpha_n)$. Il existe
alors des morphismes canoniques de $A$-modules qui donnent un diagramme
commutatif
$$
\xymatrix{
& & & H \ar[lld] \ar[ld] \ar[d] \\
\ldots \ar[r] & H_3 \ar[r] & H_2 \ar[r] & H_1
}
$$
De plus, il existe un $c > 0$ et des morphismes canoniques de $A$-modules
$H_n \to H/I^{n - c}H$ pour $n \geq c$, tels que les composées
$$
H/I^n H \to H_n \to  H/I^{n - c}H
\quad\text{and}\quad
H_n \to H/I^{n - c}H \to H_{n - c}
$$
sont les morphismes canoniques. En outre :
\begin{enumerate}
\item $(H_n)$ et $(H/I^nH)$ sont isomorphes comme pro-objets de
$\text{Mod}_A$ ;
\item $\lim H_n = \lim H/I^n H$ ;
\item le système projectif $(H_n)$ vérifie la condition de Mittag--Leffler ;
\item l'image de $H_{n + c} \to H_n$ est égale à l'image de $H \to H_n$ ;
\item la composée
$I^cH_n \to H_n \to H/I^{n - c}H \to H_n/I^{n - c}H_n$ est l'inclusion
$I^cH_n \to H_n$ suivie du morphisme quotient
$H_n \to H_n/I^{n - c}H_n$ ;
\item le noyau et le conoyau de $H/I^nH \to H_n$ sont annulés par $I^c$.
\end{enumerate}
\end{lemma}

\begin{proof}
Remarquons que $H_n = \beta^{-1}(I^nM)/\Im(\alpha) + I^nL$. Pour $n \geq 2$,
on a $\beta^{-1}(I^nM) \subset \beta^{-1}(I^{n - 1}M)$ et
$\Im(\alpha) + I^nL \subset \Im(\alpha) + I^{n - 1}L$. On obtient ainsi le
morphisme canonique $H_n \to H_{n - 1}$. De même,
$\Ker(\beta) \subset \beta^{-1}(I^nM)$ et
$\Im(\alpha) \subset \Im(\alpha) + I^nL$, ce qui donne le morphisme canonique
$H \to H_n$. Nous omettons de vérifier la commutativité du diagramme.

\medskip\noindent
Par Artin--Rees, on peut choisir $c_1, c_2 \geq 0$ tels que
$\beta^{-1}(I^nM) \subset \Ker(\beta) + I^{n - c_1}L$ for $n \geq c_1$ and
$\Ker(\beta) \cap I^nL \subset I^{n - c_2}\Ker(\beta)$ for $n \geq c_2$, see
Algèbre, Lemmes \ref{algebra-lemma-map-AR} et
\ref{algebra-lemma-Artin-Rees}. Posons $c = c_1 + c_2$.

\medskip\noindent
Soit $n \geq c$. Définissons $\psi_n : H_n \to H/I^{n - c}H$ comme suit.
Soit $x \in H_n$. Choisissons $y \in \beta^{-1}(I^nM)$ représentant $x$ et
écrivons $y = z + w$ avec $z \in \Ker(\beta)$ et
$w \in I^{n - c_1}L$ (ce qui est possible par le choix de $c_1$). On définit
$\psi_n(x)$ comme la classe de $z$ dans $H/I^{n - c}H$. Pour vérifier que
cette définition est indépendante des choix, supposons disposer d'un second
triplet $y', z', w'$ comme ci-dessus pour $x$. Alors
$y' - y \in \Im(\alpha) + I^nL$ ; écrivons
$y' - y = \alpha(v) + u$ avec $v \in K$ et $u \in I^nL$. Ainsi,
$$
y' = z' + w' = \alpha(v) + u + z + w
\Rightarrow
z' = z + \alpha(v) + u + w - w'
$$
Comme $\beta(z' - z - \alpha(v)) = 0$, on obtient
$u + w - w' \in \Ker(\beta) \cap I^{n - c_1}L$, lequel est contenu dans
$I^{n - c_1 - c_2}\Ker(\beta) = I^{n - c}\Ker(\beta)$ par le choix de $c_2$.
Ainsi, $z'$ et $z$ ont la même image dans $H/I^{n - c}H$, comme voulu.

\medskip\noindent
La composée $H/I^n H \to H_n \to H/I^{n - c}H$ est le morphisme canonique.
En effet, si $z \in \Ker(\beta)$ représente un élément $x$ de
$H/I^nH = \Ker(\beta)/\Im(\alpha) + I^n\Ker(\beta)$, il ressort clairement de
ce qui précède que les morphismes construits envoient $x$ sur la classe de
$z$ dans $H/I^{n - c}H$.

\medskip\noindent
Considérons la composée $H_n \to H/I^{n - c}H \to H_{n - c}$. Étant donnés
$x, y, z, w$ comme dans la construction précédente de $\psi_n$, l'élément
$x$ est envoyé sur la classe de $z$ dans $H_{n - c}$. D'autre part, le
morphisme canonique $H_n \to H_{n - c}$ du premier paragraphe envoie $x$ sur
la classe de $y$. Il suffit donc de montrer que
$y - z \in \Im(\alpha) + I^{n - c}L$, ce qui résulte de
$y - z = w \in I^{n - c_1}L \subset I^{n - c}L$.

\medskip\noindent
Les assertions (1)--(4) sont des conséquences formelles de ce que nous venons
de démontrer. En effet, (1) résulte de l'existence des morphismes et de la
définition des morphismes de pro-objets dans Catégories, Remarque
\ref{categories-remark-pro-category}. L'assertion (2) vient de ce que des
pro-objets isomorphes ont des limites isomorphes. L'assertion (3) résulte
immédiatement de (4). Enfin, (4) résulte de la factorisation
$H_{n + c} \to H/I^nH \to H_n$ du morphisme canonique
$H_{n + c} \to H_n$.

\medskip\noindent
Démonstration de (5). Soit $x \in I^cH_n$. Écrivons $x = \sum f_i x_i$ avec
$x_i \in H_n$ et $f_i \in I^c$. Choisissons $y_i, z_i, w_i$ comme dans la
construction de $\psi_n$ appliquée à $x_i$. Pour le calcul de $\psi_n$ sur
$x$, on
peut alors choisir
$y = \sum f_iy_i$, $z = \sum f_i z_i$ et $w = \sum f_i w_i$
et l'on voit que $\psi_n(x)$ est donné par la classe de $z$. Son image dans
$H_n/I^{n - c}H_n$ est égale à la classe de $y$, puisque
$w = \sum f_i w_i$ appartient à $I^nL$. Cela prouve (5).

\medskip\noindent
Démonstration de (6). Soit $y \in \Ker(\beta)$ dont la classe est $x$ dans
$H$. Si l'image de $x$ dans $H_n$ est nulle, alors
$y \in I^nL + \Im(\alpha)$. Ainsi,
$y - \alpha(v) \in \Ker(\beta) \cap I^nL$ pour un certain $v \in K$. Donc
$y - \alpha(v) \in I^{n - c_2}\Ker(\beta)$, et la classe de $y$ dans
$H/I^nH$ est annulée par $I^{c_2}$. Enfin, soit $x \in H_n$ la classe d'un
$y \in \beta^{-1}(I^nM)$. Écrivons, comme ci-dessus, $y = z + w$ avec
$z \in \Ker(\beta)$ et $w \in I^{n - c_1}L$. Si $f \in I^{c_1}$, alors $fx$
est la classe de $fy + fw \equiv fy$ modulo
$\Im(\alpha) + I^nL$ ; ainsi, $fx$ est l'image de la classe de $fy$ dans $H$,
comme voulu.
\end{proof}

\begin{lemma}
\label{lemma-kollar-kovacs-pseudo-coherent}
\begin{reference}
Courriel de Kovács du 23/02/2018.
\end{reference}
Soit $I$ un idéal d'un anneau noethérien $A$ et soit $K \in D(A)$
pseudo-cohérent. Posons $K_n = K \otimes_A^\mathbf{L} A/I^n$. Alors, pour
tout $i \in \mathbf{Z}$, le système $H^i(K_n)$ vérifie la condition de
Mittag--Leffler, et $\lim H^i(K)/I^nH^i(K)$ est égal à $\lim H^i(K_n)$.
\end{lemma}

\begin{proof}
On peut représenter $K$ par un complexe borné supérieurement $P^\bullet$ de
$A$-modules libres de rang fini. Alors $K_n$ est représenté par
$P^\bullet/I^nP^\bullet$. La propriété de Mittag--Leffler résulte donc du
Lemme \ref{lemma-consequence-Artin-Rees}. La dernière assertion découle
ensuite du Lemme \ref{lemma-kollar-kovacs}.
\end{proof}

\begin{lemma}
\label{lemma-derived-completion-plain-completion}
Soit $A$ un anneau noethérien, soit $I \subset A$ un idéal et soit
$M^\bullet$ un complexe borné de $A$-modules de type fini. Le système
projectif de morphismes
$$
M^\bullet \otimes_A^\mathbf{L} A/I^n \longrightarrow M^\bullet/I^nM^\bullet
$$
définit un isomorphisme de pro-objets de $D(A)$.
\end{lemma}

\begin{proof}
Écrivons $I = (f_1, \ldots, f_r)$. Soit $K_n \in D(A)$ l'objet représenté
par le complexe de Koszul associé à $f_1^n, \ldots, f_r^n$. Rappelons que
l'on a des morphismes $K_n \to A/I^n$ qui induisent un pro-isomorphisme de
systèmes projectifs ; voir le Lemme \ref{lemma-sequence-Koszul-complexes}. Il
suffit donc de montrer que
$$
M^\bullet \otimes_A^\mathbf{L} K_n \longrightarrow M^\bullet/I^nM^\bullet
$$
définit un isomorphisme de pro-objets de $D(A)$. Comme $K_n$ est représenté
par un complexe de $A$-modules libres de rang fini placé aux degrés
$-r, \ldots, 0$, il existe $a, b \in \mathbf{Z}$ tels que la source et le but
de la flèche affichée aient une cohomologie nulle hors de $[a, b]$, pour tout
$n$. On peut donc appliquer Catégories dérivées, Lemme
\ref{derived-lemma-pro-isomorphism-bis} ; il suffit alors de montrer que les
morphismes
$$
H^i(M^\bullet \otimes_A^\mathbf{L} A/I^n) \to H^i(M^\bullet/I^nM^\bullet)
$$
définissent des isomorphismes de pro-systèmes de $A$-modules pour tout
$i \in \mathbf{Z}$. Pour le voir, choisissons un quasi-isomorphisme
$P^\bullet \to M^\bullet$ où $P^\bullet$ est un complexe borné supérieurement
de $A$-modules libres de rang fini. Les flèches précédentes sont données par
les morphismes
$$
H^i(P^\bullet/I^nP^\bullet) \to H^i(M^\bullet/I^nM^\bullet)
$$
Ceux-ci définissent un isomorphisme de pro-systèmes d'après le
Lemme \ref{lemma-consequence-Artin-Rees}. En effet, ce lemme montre qu'ils sont
tous deux isomorphes au pro-système $H^i/I^nH^i$, où
$H^i = H^i(M^\bullet) = H^i(P^\bullet)$.
\end{proof}

\begin{lemma}
\label{lemma-hom-systems-ML}
Soit $A$ un anneau noethérien, soit $I \subset A$ un idéal et soient $M$, $N$
des $A$-modules de type fini. Posons $M_n = M/I^nM$ et $N_n = N/I^nN$.
Alors :
\begin{enumerate}
\item les systèmes $(\Hom_A(M_n, N_n))$ et
$(\text{Isom}_A(M_n, N_n))$ vérifient la condition de Mittag--Leffler ;
\item il existe un $c \geq 0$ tel que les noyaux et conoyaux de
$$
\Hom_A(M, N)/I^n\Hom_A(M, N) \to \Hom_A(M_n, N_n)
$$
sont annulés par $I^c$ pour tout $n$ ;
\item on a
$\lim \Hom_A(M_n, N_n) =\Hom_A(M, N)^\wedge =
\Hom_{A^\wedge}(M^\wedge, N^\wedge)$
\item $\lim \text{Isom}_A(M_n, N_n) =
\text{Isom}_{A^\wedge}(M^\wedge, N^\wedge)$.
\end{enumerate}
Ici, ${}^\wedge$ désigne la complétion $I$-adique usuelle.
\end{lemma}

\begin{proof}
Remarquons que $\Hom_A(M_n, N_n) = \Hom_A(M, N_n)$. Choisissons une
présentation
$$
A^{\oplus t} \to A^{\oplus s} \to M \to 0
$$
En appliquant le foncteur contravariant exact à gauche $\Hom_A(-, N)$, on
obtient un complexe
$$
0 \xrightarrow{\alpha} N^{\oplus s} \xrightarrow{\beta} N^{\oplus t}
$$
dont la cohomologie au milieu est $\Hom_A(M, N)$, et tel que, pour
$n \geq 0$, la cohomologie de
$$
0 \xrightarrow{\alpha_n} N_n^{\oplus s} \xrightarrow{\beta_n} N_n^{\oplus t}
$$
soit $\Hom_A(M_n, N_n)$. Soit $c \geq 0$ comme dans le
Lemme \ref{lemma-consequence-Artin-Rees} pour $A$, $I$, $\alpha$ et $\beta$.
L'assertion (3) de ce lemme donne la propriété de Mittag--Leffler pour
$(\Hom_A(M_n, N_n))$. Le noyau et le conoyau des morphismes
$\Hom_A(M, N)/I^n\Hom_A(M, N) \to \Hom_A(M_n, N_n)$
sont annulés par $I^c$ d'après l'assertion (6) du lemme. L'assertion (2) du
lemme donne $\lim \Hom_A(M_n, N_n) = \Hom_A(M, N)^\wedge$. L'égalité
$$
\Hom_{A^\wedge}(M^\wedge, N^\wedge) = \lim \Hom_A(M_n, N_n)
$$
résulte formellement des égalités $M^\wedge = \lim M_n$ et
$M_n = M^\wedge/I^nM^\wedge$, ainsi que des égalités correspondantes pour
$N$ ; voir Algèbre, Lemme \ref{algebra-lemma-completion-complete}.

\medskip\noindent
Le résultat pour les isomorphismes résulte du cas des homomorphismes appliqué
à $(\Hom(M_n, N_n))$ et à $(\Hom(N_n, M_n))$, ainsi que du fait suivant : si
$n > m > 0$ et si l'on a des morphismes $\alpha : M_n \to N_n$ et
$\beta : N_n \to M_n$ qui induisent des isomorphismes $M_m \to N_m$ et
$N_m \to M_m$, alors $\alpha$ et $\beta$ sont des isomorphismes. En effet,
$\alpha \circ \beta$ est alors surjectif par le lemme de Nakayama (Algèbre,
Lemme \ref{algebra-lemma-NAK}), donc $\alpha \circ \beta$ est un isomorphisme
d'après Algèbre,
Lemme \ref{algebra-lemma-fun}.
\end{proof}

\begin{lemma}
\label{lemma-isomorphic-completions}
Soit $A$ un anneau noethérien, soit $I \subset A$ un idéal et soient $M$, $N$
des $A$-modules de type fini. Posons $M_n = M/I^nM$ et $N_n = N/I^nN$. Si
$M_n \cong N_n$ pour tout $n$, alors $M^\wedge \cong N^\wedge$ comme
$A^\wedge$-modules.
\end{lemma}

\begin{proof}
D'après le Lemme \ref{lemma-hom-systems-ML}, le système
$(\text{Isom}_A(M_n, N_n))$ vérifie la condition de Mittag--Leffler. Par
hypothèse, chacun des ensembles $\text{Isom}_A(M_n, N_n)$ est non vide. Ainsi,
$\lim \text{Isom}_A(M_n, N_n)$ est non vide. Puisque
$\lim \text{Isom}_A(M_n, N_n) = \text{Isom}_{A^\wedge}(M^\wedge, N^\wedge)$,
on obtient un isomorphisme.
\end{proof}

\begin{remark}
\label{remark-weird-systems}
Soit $I$ un idéal d'un anneau noethérien $A$. Posons $A_n = A/I^n$ pour
$n \geq 1$. Considérons la catégorie suivante :
\begin{enumerate}
\item un objet est une suite $\{E_n\}_{n \geq 1}$ où $E_n$ est un
$A_n$-module de type fini ;
\item un morphisme $\{E_n\} \to \{E'_n\}$ est donné par des morphismes
$$
\varphi_n : I^cE_n \longrightarrow E'_n/E'_n[I^c]
\quad\text{for }n \geq c
$$
où $E'_n[I^c]$ est le sous-module de torsion (Section
\ref{section-torsion}), modulo l'équivalence suivante : on identifie
$(c, \varphi_n)$ à $(c + 1, \overline{\varphi}_n)$, où
$\overline{\varphi}_n : I^{c + 1}E_n \longrightarrow E'_n/E'_n[I^{c + 1}]$
est le morphisme induit.
\end{enumerate}
La composée de $(c, \varphi_n) : \{E_n\} \to \{E'_n\}$ et de
$(c', \varphi'_n) : \{E'_n\} \to \{E''_n\}$ est définie par les composées
évidentes
$$
I^{c + c'}E_n \to I^{c'}E'_n/E'_n[I^{c}] \to E''_n/E''_n[I^{c + c'}]
$$
pour $n \geq c + c'$. Nous omettons de vérifier qu'il s'agit bien d'une
catégorie.
\end{remark}

\begin{lemma}
\label{lemma-iso}
Un morphisme $(c, \varphi_n)$ de la catégorie de la
Remarque \ref{remark-weird-systems} est un isomorphisme si et seulement s'il
existe un $c' \geq 0$ tel que $\Ker(\varphi_n)$ et $\Coker(\varphi_n)$ soient
de $I^{c'}$-torsion pour tout $n \gg 0$.
\end{lemma}

\begin{proof}
On peut supposer, et l'on suppose, que $c' \geq c$ et que
$\Ker(\varphi_n)$ et $\Coker(\varphi_n)$ sont de $I^{c'}$-torsion pour tout
$n$. Pour $n \geq c'$ et $x \in I^{c'}E'_n$, on peut choisir
$y \in I^cE_n$ tel que $x = \varphi_n(y) \bmod E'_n[I^c]$, puisque
$\Coker(\varphi_n)$ est annulé par $I^{c'}$. Posons $\psi_n(x)$ égal à la
classe de $y$ dans $E_n/E_n[I^{c'}]$. Pour un autre choix $y' \in I^cE_n$
tel que $x = \varphi_n(y') \bmod E'_n[I^c]$, la différence $y - y'$ s'envoie
sur zéro dans $E'_n/E'_n[I^c]$ et est donc annulée par $I^{c'}$ dans
$I^cE_n$. Les morphismes
$\psi_n : I^{c'}E'_n \to E_n/E_n[I^{c'}]$
sont donc bien définis. Nous omettons de vérifier que $(c', \psi_n)$ est
l'inverse de $(c, \varphi_n)$ dans cette catégorie.
\end{proof}

\begin{lemma}
\label{lemma-dejong-kollar-kovacs}
\begin{reference}
Correspondance électronique entre János Kollár, Sándor Kovács et Johan de
Jong, le 23/02/2018.
\end{reference}
Soit $I$ un idéal de l'anneau noethérien $A$, et soient $M$ et $N$ des
$A$-modules de type fini. Écrivons $A_n = A/I^n$, $M_n = M/I^nM$ et
$N_n = N/I^nN$. Pour tout $i \geq 0$, les objets
$$
\{\Ext^i_A(M, N)/I^n\Ext^i_A(M, N)\}_{n \geq 1}
\quad\text{and}\quad
\{\Ext^i_{A_n}(M_n, N_n)\}_{n \geq 1}
$$
sont isomorphes dans la catégorie $\mathcal{C}$ de la
Remarque \ref{remark-weird-systems}.
\end{lemma}

\begin{proof}
Choisissons une suite exacte courte
$$
0 \to K \to A^{\oplus r} \to M \to 0
$$
et posons $K_n = K/I^nK$. Pour $n \geq 1$, définissons
$K(n) = \Ker(A_n^{\oplus r} \to M_n)$, de sorte que l'on a des suites exactes
$$
0 \to K(n) \to A_n^{\oplus r} \to M_n \to 0
$$
ainsi que des surjections $K_n \to K(n)$. En fait, d'après le
Lemme \ref{lemma-consequence-Artin-Rees}, il existe un $c \geq 0$ et des
morphismes $K(n) \to K_n/I^{n - c}K_n$ qui sont « presque inverses ». Comme
$I^{n - c}K_n \subset K_n[I^c]$, ces morphismes montrent que les systèmes
$\{K(n)\}_{n \geq 1}$ et $\{K_n\}_{n \geq 1}$ sont isomorphes dans
$\mathcal{C}$.

\medskip\noindent
Nous affirmons que les systèmes
$$
\{\Ext^i_{A_n}(K(n), N_n)\}_{n \geq 1}
\quad\text{and}\quad
\{\Ext^i_{A_n}(K_n, N_n)\}_{n \geq 1}
$$
sont isomorphes dans la catégorie $\mathcal{C}$. En effet, les morphismes
surjectifs $K_n \to K(n)$ ont des noyaux annulés par $I^c$ et déterminent donc
des morphismes
$$
\Ext^i_{A_n}(K(n), N_n) \to \Ext^i_{A_n}(K_n, N_n)
$$
dont le noyau et le conoyau sont annulés par $I^c$. L'affirmation résulte donc
du Lemme \ref{lemma-iso}.

\medskip\noindent
Pour $i \geq 2$, on a des isomorphismes
$$
\Ext^{i - 1}_A(K, N) = \Ext^i_A(M, N)
\quad\text{and}\quad
\Ext^{i - 1}_{A_n}(K(n), N_n) = \Ext^i_{A_n}(M_n, N_n)
$$
On voit ainsi qu'il suffit de démontrer le lemme pour $i = 0, 1$.

\medskip\noindent
Pour $i = 0, 1$, considérons le diagramme commutatif
$$
\xymatrix{
0 \ar[r] &
\Hom(M, N) \ar[r] \ar[dd] &
N^{\oplus r} \ar[r]_-\varphi \ar[dd] &
\Hom(K, N) \ar[r] \ar[d] &
\Ext^1(M, N) \ar[r] &
0 \\
& & &
\Hom(K_n, N_n)
\\
0 \ar[r] &
\Hom(M_n, N_n) \ar[r] &
N_n^{\oplus r} \ar[r] &
\Hom(K(n), N_n) \ar[r] \ar[u] &
\Ext^1(M_n, N_n) \ar[r] &
0
}
$$
D'après le Lemme \ref{lemma-hom-systems-ML}, le noyau et le conoyau de
$\Hom(M, N)/I^n \Hom(M, N) \to \Hom(M_n, N_n)$ et de
$\Hom(K, N)/I^n \Hom(K, N) \to \Hom(K_n, N_n)$ sont de $I^c$-torsion pour un
certain $c \geq 0$ indépendant de $n$. Nous avons vu ci-dessus que les
conoyaux des morphismes injectifs
$\Hom(K(n), N_n) \to \Hom(K_n, N_n)$ sont annulés par $I^c$, quitte à
augmenter $c$. Pour un tel $c$, on obtient des morphismes
$\delta_n : I^c\Hom(K, N)/I^n\Hom(K, N) \to \Hom(K(n), N_n)$ qui s'insèrent
dans le diagramme (nous omettons la formulation précise). Le noyau et le
conoyau de $\delta_n$ sont annulés par $I^c$, quitte à augmenter encore $c$,
puisque cela vaut pour
$\Hom(K, N)/I^n \Hom(K, N) \to \Hom(K_n, N_n)$ et pour
$\Hom(K(n), N_n) \to \Hom(K_n, N_n)$. On peut alors utiliser la commutativité
du diagramme en traits pleins
$$
\xymatrix{
\varphi^{-1}(I^c\Hom(K, N)) \ar[r]_-\varphi \ar[d] &
I^c\Hom(K, N)/I^n\Hom(K, N) \ar[r] \ar[d]^{\delta_n} &
I^c\Ext^1(M, N)/I^n\Ext^1(M, N) \ar[r] \ar@{..>}[d] & 0 \\
N_n^{\oplus r} \ar[r] &
\Hom(K(n), N_n) \ar[r] &
\Ext^1(M_n, N_n) \ar[r] & 0
}
$$
pour définir la flèche pointillée. Une chasse au diagramme immédiate (omise)
montre que le noyau et le conoyau de cette flèche sont annulés par $I^c$,
quitte à augmenter $c$ une dernière fois.
\end{proof}

\begin{remark}
\label{remark-awkward}
La formulation peu commode du Lemme \ref{lemma-dejong-kollar-kovacs} vient en
partie de ce qu'il n'existe pas de morphismes évidents entre les modules
$\Ext^i_{A_n}(M_n, N_n)$ lorsque $n$ varie. Ce que l'on peut déduire du lemme
est qu'il existe un $c \geq 0$ tel que, pour $m \gg n \gg 0$, on dispose de
morphismes (canoniques)
$$
I^c\Ext^i_{A_m}(M_m, N_m)/I^n\Ext^i_{A_m}(M_m, N_m) \to
\Ext^i_{A_n}(M_n, N_n)/\Ext^i_{A_n}(M_n, N_n)[I^c]
$$
dont le noyau et le conoyau sont annulés par $I^c$. C'est en ce sens (faible)
que l'on obtient un système de modules.
\end{remark}

\begin{example}
\label{example-has-to-be-awkward}
Soit $k$ un corps et soit $A = k[[x, y]]/(xy)$. Par abus de notation, nous
notons encore $x$ et $y$ les images de $x$ et $y$ dans $A$. Posons $I = (x)$
et $M = A/(y)$. On a une résolution libre
$$
\ldots \to
A \xrightarrow{y}
A \xrightarrow{x}
A \xrightarrow{y} A \to M \to 0
$$
On en conclut que
$$
\Ext^2_A(M, N) = N[y]/xN
$$
où $N[y] = \Ker(y : N \to N)$. Notons $A_n = A/I^n$, $M_n = M/I^nM$ et
$N_n = N/I^nN$. Pour tout $n$, on a une résolution libre
$$
\ldots \to
A_n^{\oplus 2} \xrightarrow{y, x^{n - 1}}
A_n \xrightarrow{x}
A_n \xrightarrow{y} A_n \to M_n \to 0
$$
On en conclut que
$$
\Ext^2_{A_n}(M_n, N_n) = (N_n[y] \cap N_n[x^{n - 1}])/xN_n
$$
où $N_n[y] = \Ker(y : N_n \to N_n)$ et
$N_n[x^{n - 1}] = \Ker(x^{n - 1} : N_n \to N_n)$. Prenons $N = A/(y)$. On
a alors
$$
\Ext^2_A(M, N) = N[y]/xN = N/xN \cong k
$$
mais
$$
\Ext^2_{A_n}(M_n, N_n) = (N_n[y] \cap N_n[x^{n - 1}])/xN_n =
N_n[x^{n - 1}]/xN_n = 0
$$
pour tout $n$, car $N_n = k[x]/(x^n)$ et la suite
$$
N_n \xrightarrow{x} N_n \xrightarrow{x^{n - 1}} N_n
$$
est exacte. Il est donc nécessaire de négliger une certaine torsion par les
puissances de $I$ pour obtenir un résultat comme le
Lemme \ref{lemma-dejong-kollar-kovacs}.
\end{example}

\begin{lemma}
\label{lemma-not-awkward}
\begin{reference}
Correspondance électronique entre János Kollár, Sándor Kovács et Johan de
Jong, le 23/02/2018.
\end{reference}
Soit $A \to B$ un homomorphisme plat d'anneaux noethériens. Soit
$I \subset A$ un idéal, et soient $M, N$ des $B$-modules de type fini. Posons
$B_n = B/I^nB$, $M_n = M/I^nM$ et $N_n = N/I^nN$. Si $M$ est plat sur $A$,
alors on a
$$
\lim \Ext^i_B(M, N)/I^n \Ext^i_B(M, N) =
\lim \Ext^i_{B_n}(M_n, N_n)
$$
pour tout $i \in \mathbf{Z}$.
\end{lemma}

\begin{proof}
Choisissons une résolution
$$
\ldots \to P_2 \to P_1 \to P_0 \to M \to 0
$$
par des $B$-modules libres de rang fini $P_i$. Posons
$P_{i, n} = P_i/I^nP_i$. Comme $M$ et $B$ sont plats sur $A$, la suite
$$
\ldots \to P_{2, n} \to P_{1, n} \to P_{0, n} \to M_n \to 0
$$
est exacte. D'une part, le complexe
$$
\Hom_B(P_0, N) \to \Hom_B(P_1, N) \to \Hom_B(P_2, N) \to \ldots
$$
calcule les modules $\Ext^i_B(M, N)$ ; d'autre part, le complexe
$$
\Hom_{B_n}(P_{0, n}, N_n) \to \Hom_{B_n}(P_{1, n}, N_n) \to
\Hom_{B_n}(P_{2, n}, N_n) \to \ldots
$$
calcule les modules $\Ext^i_{B_n}(M_n, N_n)$. Puisque
$$
\Hom_{B_n}(P_{i, n}, N_n) = \Hom_B(P_i, N)/I^n \Hom_B(P_i, N)
$$
le résultat découle de l'assertion (2) du
Lemme \ref{lemma-consequence-Artin-Rees}.
\end{proof}








\section{Systèmes de modules, bis}
\label{section-systems-bis}

\noindent
Soit $I$ un idéal d'un anneau noethérien $A$. À la Section
\ref{section-systems}, nous avons étudié les systèmes de la forme $M/I^nM$
pour des $A$-modules de type fini $M$. Dans cette section, nous considérons
plutôt les systèmes $I^nM$.

\begin{lemma}
\label{lemma-consequence-Artin-Rees-bis}
Soit $I$ un idéal d'un anneau noethérien $A$. Soit
$
K \xrightarrow{\alpha} L \xrightarrow{\beta} M
$
un complexe de $A$-modules de type fini. Posons
$H = \Ker(\beta)/\Im(\alpha)$. Pour $n \geq 0$, soit
$$
I^nK \xrightarrow{\alpha_n} I^nL \xrightarrow{\beta_n} I^nM
$$
le complexe induit et posons $H_n = \Ker(\beta_n)/\Im(\alpha_n)$. Il existe
alors des morphismes canoniques de $A$-modules
$$
\ldots \to H_3 \to H_2 \to H_1 \to H
$$
Il existe un $c > 0$ tel que, pour $n \geq c$, l'image de $H_n \to H$ soit
contenue dans $I^{n - c}H$, et il existe un morphisme canonique de $A$-modules
$I^nH \to H_{n - c}$ tel que les composées
$$
I^n H \to H_{n - c} \to  I^{n - 2c}H
\quad\text{and}\quad
H_n \to I^{n - c}H \to H_{n - 2c}
$$
sont les morphismes canoniques. En particulier, les systèmes projectifs
$(H_n)$ et $(I^nH)$ sont isomorphes comme pro-objets de $\text{Mod}_A$.
\end{lemma}

\begin{proof}
On a $H_n = \Ker(\beta) \cap I^nL/\alpha(I^nK)$. Comme
$\Ker(\beta) \cap I^nL \subset \Ker(\beta) \cap I^{n - 1}L$ et
$\alpha(I^nK) \subset \alpha(I^{n - 1}K)$, on obtient les morphismes
$H_n \to H_{n - 1}$. Il en va de même du morphisme $H_1 \to H$.

\medskip\noindent
Par Artin--Rees, on peut choisir $c_1, c_2 \geq 0$ tels que
$\Im(\alpha) \cap I^nL \subset \alpha(I^{n - c_1}K)$ for $n \geq c_1$ and
$\Ker(\beta) \cap I^nL \subset I^{n - c_2}\Ker(\beta)$ for $n \geq c_2$, see
Algèbre, Lemmes \ref{algebra-lemma-map-AR} et
\ref{algebra-lemma-Artin-Rees}. Posons $c = c_1 + c_2$.

\medskip\noindent
Le choix de $c \geq c_2$ montre immédiatement que, pour $n \geq c$, l'image
de $H_n \to H$ est contenue dans $I^{n - c}H$.

\medskip\noindent
Soit $n \geq c$. Définissons $\psi_n : I^nH \to H_{n - c}$ comme suit. Soit
$x \in I^nH$. Choisissons $y \in I^n\Ker(\beta)$ représentant $x$ et posons
$\psi_n(x)$ égal à la classe de $y$ dans $H_{n - c}$. Pour vérifier que
cette définition est indépendante des choix, supposons qu'un autre $y'$
représente $x$ comme ci-dessus. Alors $y' - y \in \Im(\alpha)$. Comme
$c \geq c_1$, on conclut que $y' - y \in \alpha(I^{n - c}K)$, si bien que
$y$ et $y'$ représentent le même élément de $H_{n - c}$. Ainsi, $\psi_n$
est bien défini.

\medskip\noindent
Les assertions sur les composées $I^n H \to H_{n - c} \to I^{n - 2c}H$ et
$H_n \to I^{n - c}H \to H_{n - 2c}$ résultent immédiatement de nos
définitions.
\end{proof}

\begin{lemma}
\label{lemma-ext-factors}
Soit $A$ un anneau noethérien. Soit $I \subset A$ un idéal.
Soient $M$ et $N$ des $A$-modules, avec $M$ de type fini. Pour tout
$p \geq 0$, il existe un $c \geq 0$ tel que, pour $n \geq c$, l'application
$\Ext_A^p(M, N) \to \Ext_A^p(I^nM, N)$
se factorise par $\Ext^p_A(I^nM, I^{n - c}N) \to \Ext_A^p(I^nM, N)$.
\end{lemma}

\begin{proof}
Pour $p = 0$, si $\varphi : M \to N$ est une application $A$-linéaire, alors
$\varphi(\sum f_i m_i) = \sum f_i \varphi(m_i)$ pour $f_i \in A$
et $m_i \in M$. Ainsi, $\varphi$ induit une application $I^nM \to I^nN$
pour tout $n$, et le résultat est vrai avec $c = 0$.

\medskip\noindent
Choisissons une suite exacte $0 \to K \to A^{\oplus t} \to M \to 0$.
Pour chaque $n$, choisissons une suite exacte
$0 \to L_n \to A^{\oplus s_n} \to I^nM \to 0$.
On peut manifestement construire un morphisme de suites exactes
$$
\xymatrix{
0 \ar[r] &
L_n \ar[r] \ar[d] &
A^{\oplus s_n} \ar[r] \ar[d] &
I^nM \ar[r] \ar[d] & 0 \\
0 \ar[r] &
K \ar[r] &
A^{\oplus t} \ar[r] &
M \ar[r] & 0
}
$$
tel que l'image de $A^{\oplus s_n} \to A^{\oplus t}$ soit contenue dans
$(I^n)^{\oplus t}$. D'après Artin--Rees (Algèbre, Lemme
\ref{algebra-lemma-Artin-Rees}), il existe un $c \geq 0$ tel que
$L_n \to K$ se factorise par $I^{n - c}K$ si $n \geq c$.

\medskip\noindent
Pour $p = 1$, les choix ci-dessus induisent un diagramme commutatif dont les
flèches pleines sont
$$
\xymatrix{
\Hom_A(A^{\oplus s_n}, N) \ar[r] &
\Hom_A(L_n, N) \ar[r] &
\Ext_A^1(I^nM, N) \ar[r] & 0 \\
\Hom_A((I^n)^{\oplus t}, I^{n - c}N) \ar[r] \ar[u] &
\Hom_A(K \cap (I^n)^{\oplus t}, I^{n - c}N) \ar[r] \ar[u] &
\Ext_A^1(I^nM, I^{n - c}N) \ar[u] \\
\Hom_A(A^{\oplus t}, N) \ar[r] \ar[u] &
\Hom_A(K, N) \ar[r] \ar[u] &
\Ext_A^1(M, N) \ar@{..>}[u] \ar[r] & 0
}
$$
et dont les lignes sont exactes. La flèche verticale médiane inférieure
existe parce que $K \cap (I^n)^{\oplus t} \subset I^{n - c}K$ ; par conséquent,
toute application $A$-linéaire $K \to N$ induit une application $A$-linéaire
$(I^n)^{\oplus t} \to I^{n - c}N$ par l'argument du premier paragraphe.
On obtient donc la flèche pointillée souhaitée.

\medskip\noindent
Pour $p > 1$, on obtient un diagramme commutatif
$$
\xymatrix{
\Ext^{p - 1}_A(I^{n - c}K, N) \ar[r] &
\Ext^{p - 1}_A(L_n, N) \ar[r] &
\Ext_A^p(I^nM, N) \\
\Ext^{p - 1}_A(K, N) \ar[rr] \ar[u] & &
\Ext_A^p(M, N) \ar[u]
}
$$
dont la flèche horizontale inférieure est un isomorphisme. Par récurrence
sur $p$, la flèche verticale de gauche se factorise par
$\Ext^{p - 1}_A(I^{n - c}K, I^{n - c - c'}N)$ pour un certain $c' \geq 0$
et tout $n \geq c + c'$. En utilisant la composée
$\Ext^{p - 1}_A(I^{n - c}K, I^{n - c - c'}N) \to
\Ext^{p - 1}_A(L_n, I^{n - c - c'}N) \to \Ext^p_A(I^nM, I^{n - c - c'}N)$
on obtient la factorisation souhaitée (pour $n \geq c + c'$, en remplaçant
$c$ par $c + c'$).
\end{proof}

\begin{lemma}
\label{lemma-ext-annihilated}
Soit $A$ un anneau noethérien. Soit $I \subset A$ un idéal. Soient $M$ et $N$
des $A$-modules, avec $M$ de type fini et $N$ annulé par une puissance de $I$.
Pour tout $p > 0$, il existe un $n$ tel que l'application
$\Ext_A^p(M, N) \to \Ext_A^p(I^nM, N)$ soit nulle.
\end{lemma}

\begin{proof}
Cela résulte immédiatement du Lemme \ref{lemma-ext-factors} et du fait que
$I^mN = 0$ pour un certain $m > 0$.
\end{proof}

\begin{lemma}
\label{lemma-ext-induced-toplogy}
Soit $A$ un anneau noethérien. Soit $I \subset A$ un idéal.
Soit $K \in D(A)$ pseudo-cohérent et soit $M$ un $A$-module de type fini.
Pour tout $p \in \mathbf{Z}$, il existe un $c$ tel que l'image de
$\Ext_A^p(K, I^nM) \to \Ext_A^p(K, M)$ soit contenue dans
$I^{n - c}\Ext_A^p(K, M)$ pour $n \geq c$.
\end{lemma}

\begin{proof}
Choisissons un complexe borné supérieurement $P^\bullet$ de $A$-modules libres
de type fini représentant $K$. Alors $\Ext_A^p(K, M)$ est la cohomologie de
$$
\Hom_A(P^{-p + 1}, M) \xrightarrow{a}
\Hom_A(P^{-p}, M) \xrightarrow{b}
\Hom_A(P^{-p - 1}, M)
$$
et $\Ext_A^p(K, I^nM)$ se calcule en remplaçant ces $A$-modules de type fini
par leur produit par $I^n$. Le résultat découle donc du
Lemme \ref{lemma-consequence-Artin-Rees-bis}
(et un énoncé bien plus fort est vrai).
\end{proof}

\noindent
Dans la Situation \ref{situation-koszul}, définissons des complexes
$I_n^\bullet$ de sorte que l'on ait des triangles distingués
$$
I_n^\bullet \to A \to K_n^\bullet \to I_n^\bullet[1]
$$
dans la catégorie triangulée $K(A)$ des complexes de $A$-modules à homotopie
près. Plus précisément, posons
$I_n^\bullet = \sigma_{\leq -1}K_n^\bullet[-1]$. On a des suites exactes de
complexes, scindées terme à terme,
$$
0 \to A \to K_n^\bullet \to I_n^\bullet[1] \to 0
$$
qui définissent des triangles distingués par définition de la structure
triangulée de $K(A)$. Leurs rotations donnent les triangles distingués
ci-dessus. Remarquons que $I_n^0 = A^{\oplus r} \to A$ est donné par la
multiplication par $f_i^n$ sur le $i$-ième facteur. Ainsi,
$I_n^\bullet \to A$ se factorise sous la forme
$$
I_n^\bullet \to (f_1^n, \ldots, f_r^n) \to A
$$
En fait, on a une suite exacte
$$
0 \to H^{-1}(K_n^\bullet) \to H^0(I_n^\bullet) \to (f_1^n, \ldots, f_r^n) \to 0
$$
et, pour tout $i < 0$, on a
$H^i(I_n^\bullet) = H^{i - 1}(K_n^\bullet)$. Les applications
$K_{n + 1}^\bullet \to K_n^\bullet$ induisent des applications
$I_{n + 1}^\bullet \to I_n^\bullet$, et l'on obtient un diagramme commutatif
$$
\xymatrix{
\ldots \ar[r] &
I_3^\bullet \ar[d] \ar[r] &
I_2^\bullet \ar[d] \ar[r] &
I_1^\bullet \ar[d] \\
\ldots \ar[r] &
(f_1^3, \ldots, f_r^3) \ar[r] &
(f_1^2, \ldots, f_r^2) \ar[r] &
(f_1, \ldots, f_r)
}
$$
dans $K(A)$.

\begin{lemma}
\label{lemma-sequence-powers-pro-bounded}
Dans la Situation \ref{situation-koszul}, supposons $A$ noethérien. Avec les
notations ci-dessus, le système projectif $(I^n)$ est pro-isomorphe dans
$D(A)$ au système projectif $(I_n^\bullet)$.
\end{lemma}

\begin{proof}
Il est élémentaire de vérifier que le système projectif $I^n$ est pro-isomorphe
au système projectif $(f_1^n, \ldots, f_r^n)$ dans la catégorie des
$A$-modules. Considérons le système projectif de triangles distingués
$$
I_n^\bullet \to (f_1^n, \ldots, f_r^n) \to C_n^\bullet \to I_n^\bullet[1]
$$
où $C_n^\bullet$ est le cône de la première flèche. D'après Catégories
dérivées, Lemme \ref{derived-lemma-pro-isomorphism}, il suffit de montrer que
le système projectif $C_n^\bullet$ est pro-nul. Le complexe $I_n^\bullet$ n'a
de termes non nuls qu'aux degrés $i$ tels que
$-r + 1 \leq i \leq 0$ ; il en va donc de même pour les bornes de
$C_n^\bullet$. Ainsi, d'après Catégories dérivées, Lemme
\ref{derived-lemma-essentially-constant-cohomology}, il suffit de montrer que
$H^p(C_n^\bullet)$ est pro-nul. La discussion précédente donne
$H^p(C_n^\bullet) = H^p(K_n^\bullet)$ pour $p \leq -1$ et
$H^p(C_n^\bullet) = 0$ pour $p \geq 0$. La preuve du
Lemme \ref{lemma-sequence-Koszul-complexes} montre que les systèmes projectifs
$H^p(K_n^\bullet)$ sont pro-nuls (c'est ici qu'intervient l'hypothèse que $A$
est noethérien).
\end{proof}

\begin{lemma}
\label{lemma-tensoring-Deligne-system}
Soit $A$ un anneau noethérien. Soit $I \subset A$ un idéal. Soit
$M^\bullet$ un complexe borné de $A$-modules de type fini. Le système
projectif d'applications
$$
I^n \otimes_A^\mathbf{L} M^\bullet \longrightarrow I^nM^\bullet
$$
définit un isomorphisme de pro-objets de $D(A)$.
\end{lemma}

\begin{proof}
Choisissons des générateurs $f_1, \ldots, f_r \in I$ de $I$. Le système
projectif $I^n$ est pro-isomorphe au système projectif
$(f_1^n, \ldots, f_r^n)$ dans la catégorie des $A$-modules. Avec les notations
du Lemme \ref{lemma-sequence-powers-pro-bounded}, il suffit donc de montrer que
le système projectif d'applications
$$
I_n^\bullet \otimes_A^\mathbf{L} M^\bullet
\longrightarrow
(f_1^n, \ldots, f_r^n)M^\bullet
$$
définit un isomorphisme de pro-objets de $D(A)$. Soient $a \leq b$ tels que
$M^i = 0$ si $i \not \in [a, b]$. La source et le but des flèches ci-dessus
n'ont alors de cohomologie qu'aux degrés $[-r + a, b]$. Il suffit donc de
montrer que, pour tout $p \in \mathbf{Z}$, le système projectif d'applications
$$
H^p(I_n^\bullet \otimes_A^\mathbf{L} M^\bullet)
\longrightarrow
H^p((f_1^n, \ldots, f_r^n)M^\bullet)
$$
définit un isomorphisme de pro-objets de $A$-modules ; voir Catégories dérivées,
Lemme \ref{derived-lemma-pro-isomorphism-bis}. En utilisant le pro-isomorphisme
entre
$I_n^\bullet \otimes_A^\mathbf{L} M^\bullet$
et $I^n \otimes_A^\mathbf{L} M^\bullet$, ainsi que le pro-isomorphisme entre
$(f_1^n, \ldots, f_r^n)M^\bullet$ and $I^nM^\bullet$
on se ramène à montrer que le système projectif d'applications
$$
H^p(I^n \otimes_A^\mathbf{L} M^\bullet)
\longrightarrow
H^p(I^nM^\bullet)
$$
définit un isomorphisme de pro-objets de $A$-modules. Choisissons un complexe
borné supérieurement de $A$-modules libres de type fini, noté $P^\bullet$, et un
quasi-isomorphisme $P^\bullet \to M^\bullet$. Il suffit alors de montrer que le
système projectif d'applications
$$
H^p(I^nP^\bullet)
\longrightarrow
H^p(I^nM^\bullet)
$$
est un pro-isomorphisme. Cela résulte du
Lemme \ref{lemma-consequence-Artin-Rees-bis}, puisque
$H^p(P^\bullet) = H^p(M^\bullet)$.
\end{proof}

\begin{lemma}
\label{lemma-factor-through-derived-tensor-product}
Soit $A$ un anneau noethérien. Soit $I \subset A$ un idéal. Soit $M$ un
$A$-module de type fini. Il existe un entier $n > 0$ tel que
$I^nM \to M$ se factorise par l'application
$I \otimes_A^\mathbf{L} M \to M$ dans $D(A)$.
\end{lemma}

\begin{proof}
Cela résulte du Lemme \ref{lemma-tensoring-Deligne-system}. On peut également
le voir directement comme suit. Considérons le triangle distingué
$$
I \otimes_A^\mathbf{L} M \to M \to A/I \otimes_A^\mathbf{L} M \to
I \otimes_A^\mathbf{L} M[1]
$$
Les axiomes d'une catégorie triangulée ramènent la preuve à montrer que
$I^nM \to A/I \otimes_A^\mathbf{L} M$ est nul dans $D(A)$ pour un certain
$n$. Choisissons des générateurs $f_1, \ldots, f_r$ de $I$, posons
$K = K_\bullet(A, f_1, \ldots, f_r)$ pour le complexe de Koszul et considérons
la factorisation $A \to K \to A/I$ de l'application quotient. Il suffit alors
de montrer que $I^nM \to K \otimes_A M$ est nul dans $D(A)$ pour un certain
$n > 0$. Supposons avoir trouvé un $n > 0$ tel que
$I^nM \to K \otimes_A M$ se factorise par
$\tau_{\geq t}(K \otimes_A M)$ dans $D(A)$. L'obstruction à une factorisation
par $\tau_{\geq t + 1}(K \otimes_A M)$ est alors un élément de
$\Ext^t(I^nM, H_t(K \otimes_A M))$. Le $A$-module de type fini
$H_t(K \otimes_A M)$ est annulé par $I$. Le
Lemme \ref{lemma-ext-annihilated} permet donc, après avoir augmenté $n$, de
supposer cette obstruction nulle. En répétant l'opération un nombre fini de
fois, on trouve un $n$ tel que $I^nM \to K \otimes_A M$ se factorise par
$0 = \tau_{\geq r + 1}(K \otimes_A M)$ dans $D(A)$, ce qui conclut.
\end{proof}






\section{Résultats divers}
\label{section-misc}

\noindent
Quelques résultats qui ne trouvent pas place ailleurs.

\begin{lemma}
\label{lemma-ext-annihilated-into}
Soit $A$ un anneau noethérien. Soit $I \subset A$ un idéal.
Soit $K \in D(A)$ pseudo-cohérent. Soit $a \in \mathbf{Z}$.
Supposons que, pour tout $A$-module de type fini $M$, les modules
$\Ext^i_A(K, M)$ soient de $I$-torsion pour $i \geq a$.
Alors, pour $i \geq a$ et $M$ de type fini, le système
$\Ext^i_A(K, M/I^nM)$ est essentiellement constant, de valeur
$$
\Ext^i_A(K, M) = \lim \Ext^i_A(K, M/I^nM)
$$
\end{lemma}

\begin{proof}
Soit $M$ un $A$-module de type fini. Comme $K$ est pseudo-cohérent,
$\Ext^i_A(K, M)$ est un $A$-module de type fini. Ainsi, pour $i \geq a$, il
est annulé par $I^t$ pour un certain $t \geq 0$. Le
Lemme \ref{lemma-ext-induced-toplogy} montre donc que l'image de
$\Ext^i_A(K, I^nM) \to \Ext^i_A(K, M)$ est nulle pour un certain $n > 0$.
La suite exacte $0 \to I^nM \to M \to M/I^n M \to 0$ donne une suite exacte
longue
$$
\Ext^i_A(K, I^nM) \to \Ext^i_A(K, M) \to
\Ext^i_A(K, M/I^nM) \to \Ext^{i + 1}_A(K, I^nM)
$$
Les systèmes $\Ext^i_A(K, I^nM)$ et $\Ext^{i + 1}_A(K, I^nM)$ sont
essentiellement constants, de valeur $0$, d'après ce qui précède (appliqué aux
$A$-modules de type fini $I^mM$). Une chasse au diagramme montre que
$\Ext^i_A(K, M/I^nM)$ est essentiellement constant, de valeur
$\Ext^i_A(K, M)$.
\end{proof}

\begin{lemma}
\label{lemma-tor-annihilated}
Soit $A$ un anneau noethérien. Soit $I \subset A$ un idéal. Soit $M$ un
$A$-module de type fini. Soit $N$ un $A$-module annulé par $I$.
Il existe un entier $n > 0$ tel que
$\text{Tor}^A_p(I^nM, N) \to \text{Tor}^A_p(M, N)$ soit nul pour tout
$p \geq 0$.
\end{lemma}

\begin{proof}
D'après le Lemme \ref{lemma-factor-through-derived-tensor-product}, on peut
factoriser $I^nM \to M$ sous la forme
$I^nM \to M \otimes_A^\mathbf{L} I \to M$. Nous affirmons que la composée
$$
I^nM \otimes_A^\mathbf{L} N \to
(M \otimes_A^\mathbf{L} I) \otimes_A^\mathbf{L} N
\to M \otimes_A^\mathbf{L} N
$$
est nulle. En effet, le diagramme
$$
\xymatrix{
(M \otimes_A^\mathbf{L} I) \otimes_A^\mathbf{L} N \ar[rr] \ar[rd] & &
M \otimes_A^\mathbf{L} (I \otimes_A^\mathbf{L} N) \ar[ld] \\
& M \otimes_A^\mathbf{L} N
}
$$
est commutatif (les détails sont laissés au lecteur), et l'application
$I \otimes_A^\mathbf{L} N \to N$ est nulle puisque $N$ est annulé par $I$.
\end{proof}

\begin{lemma}
\label{lemma-pseudo-coherent-tensor-limit}
Soit $R$ un anneau. Soit $K \in D(R)$ pseudo-cohérent.
Soit $(M_n)$ un système projectif de $R$-modules. Alors
$R\lim K \otimes_R^\mathbf{L} M_n = K \otimes_R^\mathbf{L} R\lim M_n$.
\end{lemma}

\begin{proof}
Considérons le triangle distingué de définition
$$
R\lim M_n \to \prod M_n \to \prod M_n \to R\lim M_n[1]
$$
et appliquons le Lemme \ref{lemma-pseudo-coherent-tensor}.
\end{proof}

\begin{lemma}
\label{lemma-additivity-of-pd}
Soit $R$ un anneau local noethérien. Soit $I \subset R$ un idéal et soit $E$
un $R/I$-module de type fini non nul. Si $R/I$ est de dimension projective
finie et si $E$ est de dimension projective finie sur $R/I$, alors $E$ est de
dimension projective finie sur $R$ et
$$
\text{pd}_R(E) = \text{pd}_R(R/I) + \text{pd}_{R/I}(E)
$$
\end{lemma}

\begin{proof}
Nous utiliserons le fait que, pour un module de type fini, être de dimension
projective finie sur $R$, resp.\ sur $R/I$, équivaut à être un module parfait ;
voir la discussion qui suit la Définition \ref{definition-perfect}. Le
Lemme \ref{lemma-cohomology-perfect} montre que $E$ est de dimension projective
finie sur $R$. La formule d'Auslander--Buchsbaum (Algèbre, Proposition
\ref{algebra-proposition-Auslander-Buchsbaum}) donne alors
$$
\text{pd}_R(E) + \text{depth}(E) = \text{depth}(R),\quad
\text{pd}_{R/I}(E) + \text{depth}(E) = \text{depth}(R/I),
$$
and
$$
\text{pd}_R(R/I) + \text{depth}(R/I) = \text{depth}(R)
$$
Dans la première égalité, on prend la profondeur de $E$ comme $R$-module, et
dans la seconde comme $R/I$-module. Ces profondeurs sont toutefois égales
(c'est immédiat, mais cela résulte aussi d'Algèbre, Lemme
\ref{algebra-lemma-depth-goes-down-finite}). Ceci achève la preuve.
\end{proof}

\begin{lemma}
\label{lemma-enlarge}
Soit $A \to B$ un homomorphisme d'anneaux. Il existe un cardinal
$\kappa = \kappa(A \to B)$ possédant la propriété suivante. Soient
$M^\bullet$, resp.\ $N^\bullet$, un complexe de $A$-modules, resp.\ de
$B$-modules. Soit $a : M^\bullet \to N^\bullet$ un morphisme de complexes de
$A$-modules qui induit un isomorphisme
$M^\bullet \otimes_A^\mathbf{L} B \to N^\bullet$ dans $D(B)$. Soient
$M_1^\bullet \subset M^\bullet$, resp.\ $N_1^\bullet \subset N^\bullet$, un
sous-complexe de $A$-modules, resp.\ de $B$-modules, tels que
$a(M_1^\bullet) \subset N_1^\bullet$. Il existe alors des sous-complexes
$$
M_1^\bullet \subset M_2^\bullet \subset M^\bullet
\quad\text{and}\quad
N_1^\bullet \subset N_2^\bullet \subset N^\bullet
$$
tels que $a(M_2^\bullet) \subset N_2^\bullet$ et satisfaisant aux propriétés
suivantes :
\begin{enumerate}
\item $\Ker(H^i(M_1^\bullet \otimes_A^\mathbf{L} B) \to H^i(N_1^\bullet))$
s'envoie sur zéro dans $H^i(M_2^\bullet \otimes_A^\mathbf{L} B)$,
\item $\Im(H^i(N_1^\bullet) \to H^i(N_2^\bullet))$ est contenue dans
$\Im(H^i(M_2^\bullet \otimes_A^\mathbf{L} B) \to H^i(N_2^\bullet))$,
\item $|\bigcup M_2^i \cup \bigcup N_2^i| \leq
\max(\kappa, |\bigcup M_1^i \cup \bigcup N_1^i|)$.
\end{enumerate}
\end{lemma}

\begin{proof}
Posons $\kappa = \max(|A|, |B|, \aleph_0)$ et
$|M^\bullet| = |\bigcup M^i|$, avec une notation analogue pour les autres
complexes. On a alors
$$
\max(\kappa, |\bigcup M_1^i \cup \bigcup N_1^i|) =
\max(\kappa, |M_1^\bullet|, |N_1^\bullet|)
$$
pour la quantité qui figure dans l'énoncé du lemme. Nous utiliserons ceci,
ainsi que d'autres observations élémentaires d'arithmétique des cardinaux,
sans autre mention.

\medskip\noindent
Montrons d'abord qu'il existe suffisamment de « petits » sous-complexes.
Pour toute paire de familles $E = \{E^i\}$ et $F = \{F^i\}$ de parties
finies $E^i \subset M^i$, $i \in \mathbf{Z}$, et
$F^i \subset N^i$, $i \in \mathbf{Z}$, on peut prendre
$$
M_1^\bullet \subset M_1(E, F)^\bullet \subset M^\bullet
\quad\text{and}\quad
N_1^\bullet \subset N_1(E, F)^\bullet \subset N^\bullet
$$
pour les plus petits sous-complexes de $A$-modules et de $B$-modules tels que
$a(M_1(E, F)^\bullet) \subset N_1(E, F)^\bullet$, avec
$E^i \subset M_1(E, F)^i$ et $F^i \subset N_1(E, F)^i$. On vérifie aisément
que
$$
|M_1(E, F)^\bullet| \leq \max(\kappa, |M_1^\bullet|)
\quad\text{and}\quad
|N_1(E, F)^\bullet| \leq \max(\kappa, |N_1^\bullet|)
$$
Nous omettons les détails. Il est clair que
$$
M^\bullet = \colim_{(E, F)} M_1(E, F)^\bullet
\quad\text{and}\quad
N^\bullet = \colim_{(E, F)} N_1(E, F)^\bullet
$$
et que les colimites sont des colimites filtrantes (terme à terme).

\medskip\noindent
Il existe une résolution $\ldots \to F^{-1} \to F^0 \to B$ par des
$A$-modules libres $F^i$ tels que $|F^i| \leq \kappa$ (nous omettons les
détails). Les modules de cohomologie de
$M_1^\bullet \otimes_A^\mathbf{L} B$ se calculent au moyen de
$\text{Tot}(M_1^\bullet \otimes_A F^\bullet)$. Il s'ensuit que
$|H^i(M_1^\bullet \otimes_A^\mathbf{L} B)| \leq \max(\kappa, |M_1^\bullet|)$.

\medskip\noindent
Soient $i \in \mathbf{Z}$ et
$\xi \in H^i(M_1^\bullet \otimes_A^\mathbf{L} B)$ un élément qui s'envoie sur
zéro dans $H^i(N_1^\bullet)$. Alors $\xi$ s'envoie sur zéro dans
$H^i(N^\bullet)$, donc $\xi$ s'envoie aussi dans
$H^i(M^\bullet \otimes_A^\mathbf{L} B)$. Puisque le produit tensoriel dérivé
commute aux colimites filtrantes, on peut trouver des familles finies
$E_\xi$ et $F_\xi$ comme ci-dessus telles que $\xi$ s'envoie sur zéro dans
$H^i(M_1(E_\xi, F_\xi)^\bullet \otimes_A^\mathbf{L} B)$.

\medskip\noindent
Soient $i \in \mathbf{Z}$ et $\eta \in H^i(N_1^\bullet)$. L'image de $\eta$
dans $H^i(N^\bullet)$ appartient à l'image de
$H^i(M^\bullet \otimes_A^\mathbf{L} B) \to H^i(N^\bullet)$. Comme ci-dessus,
on peut donc trouver des familles finies $E_\eta$ et $F_\eta$ telles que
$\eta$ s'envoie sur un élément de $H^i(N_1(E_\eta, F_\eta)^\bullet)$ qui
appartient à l'image de l'application
$H^i(M_1(E_\eta, F_\eta)^\bullet \otimes_A^\mathbf{L} B) \to
H^i(N_1(E_\eta, F_\eta)^\bullet)$.

\medskip\noindent
Définissons simplement
$$
M_2^\bullet =
\sum\nolimits_\xi M_1(E_\xi, F_\xi)^\bullet +
\sum\nolimits_\eta M_1(E_\eta, F_\eta)^\bullet
$$
où les sommes portent sur les $\xi$ et les $\eta$ des deux paragraphes
précédents et sont prises dans $M^\bullet$. De même, posons
$$
N_2^\bullet =
\sum\nolimits_\xi N_1(E_\xi, F_\xi)^\bullet +
\sum\nolimits_\eta N_1(E_\eta, F_\eta)^\bullet
$$
où les sommes sont prises dans $N^\bullet$. Par construction, ces choix
satisfont aux propriétés (1) et (2). La borne (3) en résulte également,
puisque l'ensemble des $\xi$ et des $\eta$ est de cardinal au plus
$\max(\kappa, |M_1^\bullet|, |N_1^\bullet|)$.
\end{proof}

\begin{lemma}
\label{lemma-f-bounded}
Soit $R$ un anneau et soit $f \in R$. Soit $M \in D(R)$ et soit $C$ le cône
de $f : M \to M$. Si $H^i(M)_f = 0$ pour $i < 0$ et $H^i(C) = 0$ pour
$i < -1$, alors $H^i(M) = 0$ pour $i < 0$.
\end{lemma}

\begin{proof}
Notons $M_f = M \otimes_R^\mathbf{L} R_f$ et choisissons un triangle distingué
$F \to M \to M_f$. Par hypothèse, $H^i(M_f) = H^i(M)_f = 0$ pour $i < 0$.
Il suffit donc de montrer que $H^i(F) = 0$ pour $i < 0$. Remarquons que
$H^i(F)$ est de $f$-torsion pour tout $i \in \mathbf{Z}$. D'autre part, comme
$f : M_f \to M_f$ est un isomorphisme, $C$ est isomorphe au cône de
$f : F \to F$ (voir Catégories dérivées, Proposition
\ref{derived-proposition-9}). Si $H^i(F) \not = 0$, le noyau de
$f : H^i(F) \to H^i(F)$ est non nul, ce qui entraîne que
$H^{i - 1}(C)$ est non nul. Notre hypothèse exclut ce cas lorsque
$i - 1 < -1$, et la preuve est achevée.
\end{proof}

\begin{lemma}
\label{lemma-f-flat}
Soit $R$ un anneau et soit $f \in R$. Soit $M \in D(R)$. Supposons que
\begin{enumerate}
\item $H^i(M) = 0$ for $i > 0$,
\item $M \otimes_R^\mathbf{L} R_f$ soit isomorphe à un $R_f$-module plat
placé en degré $0$,
\item $M \otimes_R^\mathbf{L} R/fR$ soit isomorphe à un $R/fR$-module plat
placé en degré $0$.
\end{enumerate}
Alors $M$ est isomorphe à un $R$-module plat placé en degré $0$.
\end{lemma}

\begin{proof}
Soit $K$ un $R$-module. Il suffit de montrer que
$N = M \otimes_R^\mathbf{L} K$ est isomorphe à un $R$-module placé en degré
$0$ ; voir Section \ref{section-tor}. L'hypothèse (1) montre que la cohomologie
de $N$ ne peut être non nulle qu'aux degrés $\leq 0$. D'après le
Lemme \ref{lemma-f-bounded}, il suffit donc de montrer que
$H^i(N)_f = 0$ pour $i < 0$ et que $H^i(C) = 0$ pour $i < -1$, où $C$ est le
cône de $f : N \to N$. Pour la première assertion, remarquons que
$$
H^i(N)_f = H^i(N \otimes_R^\mathbf{L} R_f) =
H^i(M \otimes_R^\mathbf{L} K \otimes_R^\mathbf{L} R_f ) =
H^i((M \otimes_R^\mathbf{L} R_f) \otimes_{R_f}^\mathbf{L}
(K \otimes_R^\mathbf{L} R_f))
$$
D'après l'hypothèse (2), ceci est nul sauf lorsque $i = 0$, auquel cas on
obtient $H^0(M)_f \otimes_R K$. Pour la seconde assertion, soit $C'$ le cône
de $f : K \to K$. Alors
$$
C = M \otimes_R^\mathbf{L} C'
$$
La cohomologie de $C'$ n'est non nulle qu'aux degrés $0$ et $-1$, où elle vaut
respectivement $K/fK$ et $K[f]$ ; autrement dit, on a un triangle distingué
$$
(K[f])[1] \to C' \to K/fK
$$
D'autre part, pour tout $R$-module $K'$ annulé par $f$, on a
$$
M \otimes_R^\mathbf{L} K' =
M \otimes_R^\mathbf{L} R/fR \otimes_{R/fR}^\mathbf{L} K'
$$
D'après l'hypothèse (3), ceci est égal au module
$H^0(M \otimes_R^\mathbf{L} R/fR) \otimes_{R/fR} K'$ placé en degré $0$.
En combinant ce qui précède, on conclut que la cohomologie de $C$ n'est non
nulle qu'aux degrés $0$ et $-1$, ce qui achève la preuve.
\end{proof}
\section{Astuces avec les bicomplexes}
\label{section-tricks}

\noindent
Cette section poursuit la discussion d'Homologie, section
\ref{homology-section-double-complexes-abelian-groups}.

\begin{lemma}
\label{lemma-vanishing-coh-prod-totalization}
Soit $A_0^\bullet \to A_1^\bullet \to A_2^\bullet \to \ldots$
un complexe de complexes de groupes abéliens. Supposons que
$H^{-p}(A_p^\bullet) = 0$ pour tout $p \geq 0$.
Posons $A^{p, q} = A_p^q$ et considérons $A^{\bullet, \bullet}$ comme un
bicomplexe. Alors $H^0(\text{Tot}_\pi(A^{\bullet, \bullet})) = 0$.
\end{lemma}

\begin{proof}
Notons $f_p : A_p^\bullet \to A_{p + 1}^\bullet$ les morphismes de complexes
donnés. Rappelons que la différentielle sur
$\text{Tot}_\pi(A^{\bullet, \bullet})$ est donnée par
$$
\prod\nolimits_{p + q = n} (f^q_p + (-1)^p\text{d}^q_{A_p^\bullet})
$$
sur les éléments de degré $n$.
Soit $\xi \in H^0(\text{Tot}_\pi(A^{\bullet, \bullet}))$ une classe de
cohomologie. Nous allons montrer que $\xi$ est nulle. Représentons $\xi$
par la classe d'un cocycle $x = (x_p) \in \prod A^{p, -p}$.
Comme $\text{d}(x) = 0$, nous obtenons
$\text{d}_{A_0^\bullet}(x_0) = 0$.
Comme $H^0(A_0^\bullet) = 0$, il existe $y_{-1} \in A^{0, -1}$ tel que
$\text{d}_{A_0^\bullet}(y_{-1}) = x_0$.
Nous voyons alors que
$\text{d}_{A_1^\bullet}(x_1 + f_0(y_{-1})) = 0$.
Comme $H^{-1}(A_1^\bullet) = 0$, nous pouvons trouver
$y_{-2} \in A^{1, -2}$ tel que
$-\text{d}_{A_1^\bullet}(y_{-2}) = x_1 + f_0(y_{-1})$.
Par récurrence, nous pouvons trouver
$y_{-p - 1} \in A^{p, -p - 1}$ tel que
$$
(-1)^p\text{d}_{A_p^\bullet}(y_{-p - 1}) = x_p + f_{p - 1}(y_{-p})
$$
Ceci implique que $\text{d}(y) = x$, où $y = (y_{-p - 1})$.
\end{proof}

\begin{lemma}
\label{lemma-prod-qis-gives-qis}
Soit
$$
(A_0^\bullet \to A_1^\bullet \to A_2^\bullet \to \ldots)
\longrightarrow
(B_0^\bullet \to B_1^\bullet \to B_2^\bullet \to \ldots)
$$
un morphisme entre deux complexes de complexes de groupes abéliens.
Posons $A^{p, q} = A_p^q$, $B^{p, q} = B_p^q$ afin d'obtenir des
bicomplexes. Soient $\text{Tot}_\pi(A^{\bullet, \bullet})$ et
$\text{Tot}_\pi(B^{\bullet, \bullet})$ les complexes totaux par produits
associés aux bicomplexes. Si chaque
$A_p^\bullet \to B_p^\bullet$ est un quasi-isomorphisme, alors
$\text{Tot}_\pi(A^{\bullet, \bullet}) \to \text{Tot}_\pi(B^{\bullet, \bullet})$
est un quasi-isomorphisme.
\end{lemma}

\begin{proof}
Rappelons que $\text{Tot}_\pi(A^{\bullet, \bullet})$ en degré
$n$ est donné par $\prod_{p + q = n} A^{p, q} = \prod_{p + 1 = n} A^q_p$.
Soit $C_p^\bullet$ le cône du morphisme
$A_p^\bullet \to B_p^\bullet$, voir Catégories dérivées, section
\ref{derived-section-cones}. Par fonctorialité de la construction du cône,
nous obtenons un complexe de complexes
$$
C_0^\bullet \to C_1^\bullet \to C_2^\bullet \to \ldots
$$
Nous voyons alors que $\text{Tot}_\pi(C^{\bullet, \bullet})$ en degré $n$
est donné par
$$
\prod_{p + q = n} C^{p, q} = \prod_{p + q = n} C^q_p =
\prod_{p + q = n} (B^q_p \oplus A^{q + 1}_p) =
\prod_{p + q = n} B^q_p \oplus \prod_{p + q = n} A^{q + 1}_p
$$
Nous concluons que $\text{Tot}_\pi(C^{\bullet, \bullet})$ est le cône du
morphisme
$\text{Tot}_\pi(A^{\bullet, \bullet}) \to \text{Tot}_\pi(B^{\bullet, \bullet})$
(Nous omettons la vérification de la compatibilité des différentielles.)
Il suffit donc de montrer que $\text{Tot}_\pi(A^{\bullet, \bullet})$ est
acyclique si chaque $A_p^\bullet$ est acyclique. Ceci résulte du
lemme \ref{lemma-vanishing-coh-prod-totalization}.
\end{proof}






\section{Homomorphismes d'anneaux faiblement étales}
\label{section-weakly-etale}

\noindent
La plupart des résultats de cette section proviennent de l'article
\cite{Olivier-AF} d'Olivier. Voir aussi l'article connexe
\cite{Ferrand-epi}.

\begin{definition}
\label{definition-weakly-etale}
Un anneau $A$ est dit {\it absolument plat} si tout $A$-module est plat sur
$A$. Un homomorphisme d'anneaux $A \to B$ est {\it faiblement étale} ou
{\it absolument plat} si $A \to B$ et $B \otimes_A B \to B$ sont tous deux
plats.
\end{definition}

\noindent
Les anneaux absolument plats sont parfois appelés anneaux réguliers de
von Neumann (souvent dans le cadre des anneaux non commutatifs). Une
localisation est un homomorphisme d'anneaux faiblement étale. Un homomorphisme
d'anneaux étale est faiblement étale. Voici une propriété simple mais
fondamentale.

\begin{lemma}
\label{lemma-key}
Soit $A \to B$ un homomorphisme d'anneaux tel que
$B \otimes_A B \to B$ soit plat. Soit $N$ un $B$-module. Si $N$ est plat
comme $A$-module, alors $N$ est plat comme $B$-module.
\end{lemma}

\begin{proof}
Supposons $N$ plat comme $A$-module. Alors le foncteur
$$
\text{Mod}_B \longrightarrow \text{Mod}_{B \otimes_A B},\quad
N' \mapsto N \otimes_A N'
$$
est exact. Comme $B \otimes_A B \to B$ est plat, nous concluons que le
foncteur
$$
\text{Mod}_B \longrightarrow \text{Mod}_B,\quad
N' \mapsto (N \otimes_A N') \otimes_{B \otimes_A B} B = N \otimes_B N'
$$
est exact, donc $N$ est plat sur $B$.
\end{proof}

\begin{definition}
\label{definition-weak-dimension}
Soit $A$ un anneau. Soit $d \geq 0$ un entier.
Nous disons que $A$ est de {\it dimension faible $\leq d$}
si tout $A$-module est de dimension de Tor $\leq d$.
\end{definition}

\begin{lemma}
\label{lemma-weak-dimension-goes-up}
Soit $A \to B$ un homomorphisme d'anneaux faiblement étale.
Si $A$ est de dimension faible au plus $d$, alors il en va de même de $B$.
\end{lemma}

\begin{proof}
Soit $N$ un $B$-module. Si $d = 0$, alors $N$ est plat comme $A$-module,
donc plat comme $B$-module par le lemme \ref{lemma-key}.
Supposons $d > 0$. Choisissons une résolution $F_\bullet \to N$
par des $B$-modules libres. Notre hypothèse implique que
$K = \Im(F_d \to F_{d - 1})$ est plat sur $A$, voir le
lemme \ref{lemma-last-one-flat}. Il est donc plat sur $B$ par le
lemme \ref{lemma-key}. Ainsi
$0 \to K \to F_{d - 1} \to \ldots \to F_0 \to N \to 0$
est une résolution plate de longueur $d$, et nous voyons que $N$ est de
dimension de Tor au plus $d$.
\end{proof}

\begin{lemma}
\label{lemma-absolutely-flat}
Soit $A$ un anneau. Les assertions suivantes sont équivalentes :
\begin{enumerate}
\item $A$ est de dimension faible $\leq 0$,
\item $A$ est absolument plat,
\item $A$ est réduit et tout idéal premier est maximal, et
\item tout anneau local de $A$ est un corps.
\end{enumerate}
\end{lemma}

\begin{proof}
L'équivalence de (1) et (2) est immédiate.

\medskip\noindent
Démonstration de (1) $\Rightarrow$ (3).
Supposons $A$ absolument plat. Alors tout idéal de $A$ est pur, voir
Algèbre, définition \ref{algebra-definition-pure-ideal}.
Ainsi, tout idéal de type fini est engendré par un idempotent d'après
Algèbre, lemme \ref{algebra-lemma-finitely-generated-pure-ideal}.
Si $f \in A$, alors $(f) = (e)$ pour un certain idempotent $e \in A$, et
$D(f) = D(e)$ est ouvert et fermé
(Algèbre, lemme \ref{algebra-lemma-idempotent-spec}).
Ceci implique que tout idéal premier de $A$ est maximal d'après
Algèbre, lemme \ref{algebra-lemma-ring-with-only-minimal-primes}.
De plus, si $f$ est nilpotent, alors $e = 0$, donc $f = 0$.
Ainsi $A$ est réduit.

\medskip\noindent
Démonstration de (3) $\Rightarrow$ (4). Si $A$ est réduit et tout idéal
premier est maximal, alors tout anneau local est un corps d'après
Algèbre, lemme \ref{algebra-lemma-minimal-prime-reduced-ring}.

\medskip\noindent
Démonstration de (4) $\Rightarrow$ (2). Si tout anneau local de $A$ est un
corps, alors tout $A$-module est plat d'après
Algèbre, lemme \ref{algebra-lemma-flat-localization}.
\end{proof}

\begin{lemma}
\label{lemma-product-fields-absolutely-flat}
Un produit de corps est un anneau absolument plat.
\end{lemma}

\begin{proof}
Soit $K_i$ une famille de corps. Si $f = (f_i) \in \prod K_i$, alors
l'idéal engendré par $f$ est le même que celui engendré par l'idempotent
$e = (e_i)$, où $e_i = 0, 1$ suivant que $f_i$ est $0$ ou non. Ainsi
$D(f) = D(e)$ est ouvert et fermé, et nous concluons par le
lemme \ref{lemma-absolutely-flat} et Algèbre, lemme
\ref{algebra-lemma-ring-with-only-minimal-primes}.
\end{proof}

\begin{lemma}
\label{lemma-base-change-weakly-etale}
Soient $A \to B$ et $A \to A'$ des homomorphismes d'anneaux. Soit
$B' = B \otimes_A A'$ le changement de base de $B$.
\begin{enumerate}
\item Si $B \otimes_A B \to B$ est plat, alors
$B' \otimes_{A'} B' \to B'$ est plat.
\item Si $A \to B$ est faiblement étale, alors $A' \to B'$ est faiblement
étale.
\end{enumerate}
\end{lemma}

\begin{proof}
Supposons $B \otimes_A B \to B$ plat.
L'homomorphisme d'anneaux $B' \otimes_{A'} B' \to B'$ est le changement
de base de $B \otimes_A B \to B$ par $A \to A'$. Il est donc plat d'après
Algèbre, lemme \ref{algebra-lemma-flat-base-change}. Ceci démontre (1).
La partie (2) résulte de (1) et du fait (que nous venons d'utiliser) que le
changement de base d'un homomorphisme d'anneaux plat est plat.
\end{proof}

\begin{lemma}
\label{lemma-absolutely-flat-over-absolutely-flat}
Soit $A \to B$ un homomorphisme d'anneaux tel que
$B \otimes_A B \to B$ soit plat.
\begin{enumerate}
\item Si $A$ est un anneau absolument plat, alors $B$ l'est aussi.
\item Si $A$ est réduit et $A \to B$ est faiblement étale, alors $B$ est
réduit.
\end{enumerate}
\end{lemma}

\begin{proof}
La partie (1) résulte immédiatement du lemme \ref{lemma-key} et des
définitions. Si $A$ est réduit, alors il existe une injection
$A \to A' = \prod_{\mathfrak p \subset A\text{ minimal}} A_\mathfrak p$
de $A$ dans un anneau absolument plat
(Algèbre, lemme \ref{algebra-lemma-reduced-ring-sub-product-fields} et
lemme \ref{lemma-product-fields-absolutely-flat}).
Si $A \to B$ est plat, alors le morphisme induit
$B \to B' = B \otimes_A A'$ est lui aussi injectif. Par le
lemme \ref{lemma-base-change-weakly-etale}, l'homomorphisme d'anneaux
$A' \to B'$ est faiblement étale. Par la partie (1), $B'$ est absolument
plat. Par le lemme \ref{lemma-absolutely-flat}, l'anneau $B'$ est réduit.
Ainsi $B$ est réduit.
\end{proof}

\begin{lemma}
\label{lemma-composition-weakly-etale}
Soient $A \to B$ et $B \to C$ des homomorphismes d'anneaux.
\begin{enumerate}
\item Si $B \otimes_A B \to B$ et $C \otimes_B C \to C$ sont plats, alors
$C \otimes_A C \to C$ est plat.
\item Si $A \to B$ et $B \to C$ sont faiblement étales, alors $A \to C$
est faiblement étale.
\end{enumerate}
\end{lemma}

\begin{proof}
La partie (1) résulte de la factorisation
$$
C \otimes_A C \longrightarrow C \otimes_B C \longrightarrow C
$$
du morphisme de multiplication, du fait que
$$
C \otimes_B C = (C \otimes_A C) \otimes_{B \otimes_A B} B,
$$
du fait qu'un changement de base d'un morphisme plat est plat, et du fait
que la composée d'homomorphismes d'anneaux plats est plate. Voir
Algèbre, lemmes \ref{algebra-lemma-flat-base-change} et
\ref{algebra-lemma-composition-flat}. La partie (2) résulte de (1) et du fait
(que nous venons d'utiliser) que la composée d'homomorphismes d'anneaux
plats est plate.
\end{proof}

\begin{lemma}
\label{lemma-go-down}
Soient $A \to B \to C$ des homomorphismes d'anneaux.
\begin{enumerate}
\item Si $B \to C$ est fidèlement plat et $C \otimes_A C \to C$ est plat,
alors $B \otimes_A B \to B$ est plat.
\item Si $B \to C$ est fidèlement plat et $A \to C$ est faiblement étale,
alors $A \to B$ est faiblement étale.
\end{enumerate}
\end{lemma}

\begin{proof}
Supposons $B \to C$ fidèlement plat et $C \otimes_A C \to C$ plat.
Considérons le diagramme commutatif
$$
\xymatrix{
C \otimes_A C \ar[r] & C \\
B \otimes_A B \ar[r] \ar[u] & B \ar[u]
}
$$
Les flèches verticales sont plates et la flèche horizontale supérieure est
plate. Ainsi $C$ est plat comme $B \otimes_A B$-module. Le morphisme
$B \to C$ est fidèlement plat et $C = B \otimes_B C$. Ainsi $B$ est plat
comme $B \otimes_A B$-module d'après
Algèbre, lemme \ref{algebra-lemma-flatness-descends-more-general}.
Ceci démontre (1). La partie (2) résulte de (1) et du fait que
$A \to B$ est plat si $A \to C$ est plat et $B \to C$ est fidèlement plat
(Algèbre, lemme \ref{algebra-lemma-flatness-descends-more-general}).
\end{proof}

\begin{lemma}
\label{lemma-weakly-etale-permanence}
Soit $A$ un anneau. Soit $B \to C$ un homomorphisme de $A$-algèbres
faiblement étales sur $A$. Alors $B \to C$ est faiblement étale.
\end{lemma}

\begin{proof}
L'homomorphisme d'anneaux $B \to C$ est plat par le lemme \ref{lemma-key}.
L'homomorphisme d'anneaux $C \otimes_A C \to C \otimes_B C$ est surjectif,
donc est un épimorphisme. Ainsi le lemme \ref{lemma-key} implique que,
puisque $C$ est plat sur $C \otimes_A C$, $C$ est également plat sur
$C \otimes_B C$.
\end{proof}

\begin{lemma}
\label{lemma-formally-unramified}
Soit $A \to B$ un homomorphisme d'anneaux tel que
$B \otimes_A B \to B$ soit plat. Alors $\Omega_{B/A} = 0$, c'est-à-dire que
$B$ est formellement non ramifié sur $A$.
\end{lemma}

\begin{proof}
Soit $I \subset B \otimes_A B$ le noyau du morphisme plat surjectif
$B \otimes_A B \to B$. Alors $I$ est un idéal pur
(Algèbre, définition \ref{algebra-definition-pure-ideal}), donc
$I^2 = I$ (Algèbre, lemme \ref{algebra-lemma-pure}).
Comme $\Omega_{B/A} = I/I^2$
(Algèbre, lemme \ref{algebra-lemma-differentials-diagonal}), nous obtenons
l'annulation. Ceci signifie que $B$ est formellement non ramifié sur $A$
d'après Algèbre, lemme
\ref{algebra-lemma-characterize-formally-unramified}.
\end{proof}

\begin{lemma}
\label{lemma-weakly-etale-finite-type}
Soit $A \to B$ un homomorphisme d'anneaux tel que
$B \otimes_A B \to B$ soit plat.
\begin{enumerate}
\item Si $A \to B$ est de type fini, alors $A \to B$ est non ramifié.
\item Si $A \to B$ est de présentation finie et plat, alors
$A \to B$ est étale.
\end{enumerate}
En particulier, un homomorphisme d'anneaux faiblement étale de présentation
finie est étale.
\end{lemma}

\begin{proof}
La partie (1) résulte du lemme \ref{lemma-formally-unramified} et
d'Algèbre, définition \ref{algebra-definition-unramified}.
La partie (2) résulte de la partie (1) et d'Algèbre, lemme
\ref{algebra-lemma-etale-flat-unramified-finite-presentation}.
\end{proof}

\begin{lemma}
\label{lemma-when-weakly-etale}
Soit $A \to B$ un homomorphisme d'anneaux. Alors $A \to B$ est faiblement
étale dans chacun des cas suivants :
\begin{enumerate}
\item $B = S^{-1}A$ est une localisation de $A$,
\item $A \to B$ est étale,
\item $B$ est une colimite filtrante de $A$-algèbres faiblement étales.
\end{enumerate}
\end{lemma}

\begin{proof}
Un homomorphisme d'anneaux étale est plat et le morphisme
$B \otimes_A B \to B$ est lui aussi étale comme morphisme entre
$A$-algèbres étales
(Algèbre, lemme \ref{algebra-lemma-map-between-etale}).
Ceci démontre (2).

\medskip\noindent
Soit $B_i$ un système inductif de $A$-algèbres faiblement étales.
Alors $B = \colim B_i$ est plat sur $A$ d'après
Algèbre, lemme \ref{algebra-lemma-colimit-flat}.
Remarquons que les morphismes de transition $B_i \to B_{i'}$ sont plats
d'après le lemme \ref{lemma-weakly-etale-permanence}.
Ainsi $B$ est plat sur $B_i$ pour tout $i$, et nous voyons que $B$ est plat
sur $B_i \otimes_A B_i$ d'après
Algèbre, lemme \ref{algebra-lemma-composition-flat}.
Ainsi $B$ est plat sur $B \otimes_A B = \colim B_i \otimes_A B_i$
d'après Algèbre, lemme \ref{algebra-lemma-colimit-rings-flat}.

\medskip\noindent
La partie (1) peut être démontrée directement, mais résulte aussi de la
combinaison de (2) et (3).
\end{proof}

\begin{lemma}
\label{lemma-absolutely-flat-fields}
Soit $L/K$ une extension de corps. Si $L \otimes_K L \to L$ est plat,
alors $L$ est une extension algébrique séparable de $K$.
\end{lemma}

\begin{proof}
Par le lemme \ref{lemma-go-down}, nous voyons que, pour tout sous-corps
$K \subset L' \subset L$, le morphisme $L' \otimes_K L' \to L'$ est plat.
Nous pouvons donc supposer que $L$ est une extension de corps de type fini
de $K$. Dans ce cas, le fait que $L/K$ soit formellement non ramifiée
(lemme \ref{lemma-formally-unramified}) implique que $L/K$ est finie
séparable, voir Algèbre, lemme
\ref{algebra-lemma-characterize-separable-algebraic-field-extensions}.
\end{proof}

\begin{lemma}
\label{lemma-absolutely-flat-over-field}
Soit $B$ une algèbre sur un corps $K$. Les assertions suivantes sont
équivalentes :
\begin{enumerate}
\item $B \otimes_K B \to B$ est plat,
\item $K \to B$ est faiblement étale, et
\item $B$ est une colimite filtrante de $K$-algèbres étales.
\end{enumerate}
De plus, toute sous-$K$-algèbre de type fini de $B$ est étale sur $K$.
\end{lemma}

\begin{proof}
Les parties (1) et (2) sont équivalentes parce que toute $K$-algèbre est
plate sur $K$. La partie (3) implique (1) et (2) par le
lemme \ref{lemma-when-weakly-etale}.

\medskip\noindent
Supposons (1) et (2) vérifiées. Nous allons démontrer (3) et l'assertion de
finitude du lemme. Un corps est un anneau absolument plat ; ainsi $B$ est
un anneau absolument plat par le
lemme \ref{lemma-absolutely-flat-over-absolutely-flat}.
Ainsi $B$ est réduit et tout anneau local est un corps, voir le
lemme \ref{lemma-absolutely-flat}.

\medskip\noindent
Soit $\mathfrak q \subset B$ un idéal premier. L'homomorphisme d'anneaux
$B \to B_\mathfrak q$ est faiblement étale ; ainsi $B_\mathfrak q$ est
faiblement étale sur $K$ (lemme \ref{lemma-composition-weakly-etale}).
Ainsi $B_\mathfrak q$ est une extension algébrique séparable de $K$ par le
lemme \ref{lemma-absolutely-flat-fields}.

\medskip\noindent
Soit $K \subset A \subset B$ une sous-$K$-algèbre de type fini. Nous allons
montrer que $A$ est étale sur $K$, ce qui achèvera la démonstration du
lemme. Alors tout idéal premier minimal $\mathfrak p \subset A$ est l'image
d'un idéal premier $\mathfrak q$ de $B$, voir
Algèbre, lemme \ref{algebra-lemma-injective-minimal-primes-in-image}.
Ainsi $\kappa(\mathfrak p)$, comme sous-corps de
$B_\mathfrak q = \kappa(\mathfrak q)$, est algébrique séparable sur $K$.
Par conséquent, tout point générique de $\Spec(A)$ est fermé
(Algèbre, lemme \ref{algebra-lemma-finite-residue-extension-closed}).
Ainsi $\dim(A) = 0$. Alors $A$ est le produit de ses anneaux locaux, par
exemple d'après Algèbre, proposition
\ref{algebra-proposition-dimension-zero-ring}. De plus, comme $A$ est
réduit, tous les anneaux locaux sont égaux à leurs corps résiduels, qui sont
finis séparables sur $K$. Ceci signifie que $A$ est étale sur $K$ d'après
Algèbre, lemme \ref{algebra-lemma-etale-over-field}, et achève la
démonstration.
\end{proof}

\begin{lemma}
\label{lemma-weakly-etale-residue-field-extensions}
Soit $A \to B$ un homomorphisme d'anneaux. Si $A \to B$ est faiblement
étale, alors $A \to B$ induit des extensions algébriques séparables des
corps résiduels.
\end{lemma}

\begin{proof}
Soit $\mathfrak p$ un idéal premier de $A$. Alors
$\kappa(\mathfrak p) \to B \otimes_A \kappa(\mathfrak p)$ est faiblement
étale par le lemme \ref{lemma-base-change-weakly-etale}.
Ainsi $B \otimes_A \kappa(\mathfrak p)$ est une colimite filtrante de
$\kappa(\mathfrak p)$-algèbres étales par le
lemme \ref{lemma-absolutely-flat-over-field}. Ainsi, pour
$\mathfrak q \subset B$ au-dessus de $\mathfrak p$, l'extension
$\kappa(\mathfrak q)/\kappa(\mathfrak p)$ est une colimite filtrante
d'extensions finies séparables d'après
Algèbre, lemme \ref{algebra-lemma-etale-over-field}.
\end{proof}

\begin{lemma}
\label{lemma-weak-dimension-at-most-1}
Soit $A$ un anneau. Les assertions suivantes sont équivalentes :
\begin{enumerate}
\item $A$ est de dimension faible $\leq 1$,
\item tout idéal de $A$ est plat,
\item tout idéal de type fini de $A$ est plat,
\item tout sous-module d'un $A$-module plat est plat, et
\item tout anneau local de $A$ est un anneau de valuation.
\end{enumerate}
\end{lemma}

\begin{proof}
Si $A$ est de dimension faible $\leq 1$, alors la résolution
$0 \to I \to A \to A/I \to 0$ montre que tout idéal $I$ est plat par le
lemme \ref{lemma-last-one-flat}. Ainsi (1) $\Rightarrow$ (2).

\medskip\noindent
Supposons (4). Soit $M$ un $A$-module. Choisissons une surjection
$F \to M$, où $F$ est un $A$-module libre. Alors $\Ker(F \to M)$ est plat
par hypothèse, et nous voyons que $M$ est de dimension de Tor $\leq 1$ par le
lemme \ref{lemma-tor-dimension}. Ainsi (4) $\Rightarrow$ (1).

\medskip\noindent
Tout idéal est la réunion des idéaux de type fini qu'il contient. Ainsi (3)
implique (2) d'après Algèbre, lemme \ref{algebra-lemma-colimit-flat}.
Donc (3) $\Leftrightarrow$ (2).

\medskip\noindent
Supposons (2). Supposons que $N \subset M$, avec $M$ un $A$-module plat.
Nous allons démontrer que $N$ est plat. Nous pouvons écrire
$M = \colim M_i$, avec chaque $M_i$ libre fini, voir
Algèbre, théorème \ref{algebra-theorem-lazard}.
En notant $N_i \subset M_i$ l'image inverse de $N$, nous voyons que
$N = \colim N_i$. D'après Algèbre, lemme
\ref{algebra-lemma-colimit-flat}, il suffit de démontrer que $N_i$ est plat,
et nous nous ramenons au cas $M = R^{\oplus n}$. Dans ce cas, le module
$N$ possède une filtration finie par les sous-modules
$R^{\oplus j} \cap N$, dont les sous-quotients sont des idéaux.
Par (2), ces idéaux sont plats et donc $N$ est plat d'après
Algèbre, lemme \ref{algebra-lemma-flat-ses}. Ainsi (2) $\Rightarrow$ (4).

\medskip\noindent
Supposons que $A$ satisfasse (1) et soit $\mathfrak p \subset A$ un idéal
premier. D'après les lemmes \ref{lemma-when-weakly-etale} et
\ref{lemma-weak-dimension-goes-up}, nous voyons que $A_\mathfrak p$
satisfait (1). Nous allons montrer que $A$ est un anneau de valuation si
$A$ est un anneau local satisfaisant (3). Soit $f \in \mathfrak m$ un
élément non nul. Alors $(f)$ est un module plat non nul engendré par un
élément. C'est donc un $A$-module libre d'après
Algèbre, lemme \ref{algebra-lemma-finite-flat-local}.
Il s'ensuit que $f$ est un non-diviseur de zéro et que $A$ est un domaine.
Si $I \subset A$ est un idéal de type fini, nous voyons de même que $I$ est
un $A$-module libre fini, donc (en considérant le rang) libre de rang $1$,
et $I$ est un idéal principal. Ainsi $A$ est un anneau de valuation d'après
Algèbre, lemme \ref{algebra-lemma-characterize-valuation-ring}.
Donc (1) $\Rightarrow$ (5).

\medskip\noindent
Supposons (5). Soit $I \subset A$ un idéal de type fini.
Alors $I_\mathfrak p \subset A_\mathfrak p$ est un idéal de type fini dans
un anneau de valuation, donc principal
(Algèbre, lemme \ref{algebra-lemma-characterize-valuation-ring}), donc plat.
Ainsi $I$ est plat d'après
Algèbre, lemme \ref{algebra-lemma-flat-localization}.
Donc (5) $\Rightarrow$ (3). Ceci achève la démonstration du lemme.
\end{proof}

\begin{lemma}
\label{lemma-product-weak-dimension-at-most-1}
Soit $J$ un ensemble. Pour chaque $j \in J$, soit $A_j$ un anneau de
valuation de corps des fractions $K_j$. Posons $A = \prod A_j$ et
$K = \prod K_j$. Alors $A$ est de dimension faible au plus $1$ et
$A \to K$ est une localisation.
\end{lemma}

\begin{proof}
Soit $I \subset A$ un idéal de type fini. Par le
lemme \ref{lemma-weak-dimension-at-most-1}, il suffit de montrer que $I$ est
un $A$-module plat. Soit $I_j \subset A_j$ l'image de $I$.
Remarquons que $I_j = I \otimes_A A_j$, donc $I \to \prod I_j$ est
surjectif d'après Algèbre, proposition
\ref{algebra-proposition-fg-tensor}. Ainsi $I = \prod I_j$.
Comme $A_j$ est un anneau de valuation, l'idéal $I_j$ est engendré par un
seul élément (Algèbre, lemme
\ref{algebra-lemma-characterize-valuation-ring}).
Disons que $I_j = (f_j)$. Alors $I$ est engendré par l'élément $f = (f_j)$.
Soit $e \in A$ l'idempotent qui vaut $0$ ou $1$ dans $A_j$ suivant que
$f_j$ est $0$ ou non. Alors $f = g e$ pour un certain non-diviseur de zéro
$g \in A$ : prendre $g = (g_j)$ avec $g_j = 1$ si $f_j = 0$ et
$g_j = f_j$ sinon. Ainsi $I \cong (e)$ comme module. Nous concluons que
$I$ est plat, puisque $(e)$ est un facteur direct de $A$. La dernière
assertion est vraie parce que $K = S^{-1}A$, où
$S = \prod (A_j \setminus \{0\})$.
\end{proof}

\begin{lemma}
\label{lemma-product-found-valuation-rings}
Soit $A$ un domaine normal de corps des fractions $K$.
Il existe un diagramme cartésien
$$
\xymatrix{
A \ar[d] \ar[r] & K \ar[d] \\
V \ar[r] & L
}
$$
d'anneaux où $V$ est de dimension faible au plus $1$ et où $V \to L$ est
un épimorphisme d'anneaux plat et injectif.
\end{lemma}

\begin{proof}
Pour tout $x \in K$, $x \not \in A$, choisissons $V_x \subset K$ comme dans
Algèbre, lemme \ref{algebra-lemma-find-valuation-rings}.
Posons $V = \prod_{x \in K \setminus A} V_x$ et
$L = \prod_{x \in K \setminus A} K$. L'anneau $V$ est de dimension faible
au plus $1$ d'après le
lemme \ref{lemma-product-weak-dimension-at-most-1}, qui montre aussi que
$V \to L$ est une localisation. Une localisation est plate et est un
épimorphisme, voir Algèbre, lemmes \ref{algebra-lemma-flat-localization} et
\ref{algebra-lemma-epimorphism-local}.
\end{proof}

\begin{lemma}
\label{lemma-weak-dimension-at-most-1-integrally-closed}
Soit $A$ un anneau de dimension faible au plus $1$.
Si $A \to B$ est un épimorphisme d'anneaux plat et injectif, alors
$A$ est intégralement clos dans $B$.
\end{lemma}

\begin{proof}
Soit $x \in B$ entier sur $A$. Soit $A' = A[x] \subset B$.
Alors $A'$ est une extension finie d'anneaux de $A$ d'après
Algèbre, lemme
\ref{algebra-lemma-characterize-finite-in-terms-of-integral}.
Pour montrer $A = A'$, il suffit de montrer que $A \to A'$ est un
épimorphisme d'après Algèbre, lemme
\ref{algebra-lemma-finite-epimorphism-surjective}.
Remarquons que $A'$ est plat sur $A$ par l'hypothèse sur $A$ et le fait que
$B$ est plat sur $A$ (lemme \ref{lemma-weak-dimension-at-most-1}).
Par conséquent, la composée
$$
A' \otimes_A A' \to B \otimes_A A' \to B \otimes_A B \to B
$$
est injective, c'est-à-dire $A' \otimes_A A' \cong A'$, et le lemme est
démontré.
\end{proof}

\begin{lemma}
\label{lemma-normality-goes-up}
Soit $A$ un domaine normal de corps des fractions $K$.
Soit $A \to B$ faiblement étale. Alors $B$ est intégralement clos dans
$B \otimes_A K$.
\end{lemma}

\begin{proof}
Choisissons un diagramme comme dans le
lemme \ref{lemma-product-found-valuation-rings}.
Comme $A \to B$ est plat, le changement de base donne un diagramme cartésien
$$
\xymatrix{
B \ar[d] \ar[r] & B \otimes_A K \ar[d] \\
B \otimes_A V \ar[r] & B \otimes_A L
}
$$
d'anneaux. Remarquons que $V \to B \otimes_A V$ est faiblement étale
(lemme \ref{lemma-base-change-weakly-etale}), donc $B \otimes_A V$ est de
dimension faible au plus $1$ d'après le
lemme \ref{lemma-weak-dimension-goes-up}. Remarquons que
$B \otimes_A V \to B \otimes_A L$ est un épimorphisme d'anneaux plat et
injectif, comme changement de base plat d'un tel morphisme
(Algèbre, lemmes \ref{algebra-lemma-flat-base-change} et
\ref{algebra-lemma-base-change-epimorphism}).
Par le lemme \ref{lemma-weak-dimension-at-most-1-integrally-closed}, nous
voyons que $B \otimes_A V$ est intégralement clos dans $B \otimes_A L$.
Il résulte de la propriété cartésienne du diagramme que $B$ est
intégralement clos dans $B \otimes_A K$.
\end{proof}

\begin{lemma}
\label{lemma-integral-over-henselian}
Soit $A \to B$ un homomorphisme d'anneaux. Supposons que
\begin{enumerate}
\item $A$ soit un anneau local hensélien,
\item $A \to B$ soit entier,
\item $B$ soit un domaine.
\end{enumerate}
Alors $B$ est un anneau local hensélien et $A \to B$ est un homomorphisme
local. Si $A$ est strictement hensélien, alors $B$ est un anneau local
strictement hensélien et l'extension
$\kappa(\mathfrak m_B)/\kappa(\mathfrak m_A)$ des corps résiduels est
purement inséparable.
\end{lemma}

\begin{proof}
Écrivons $B$ comme une colimite filtrante $B = \colim B_i$ de
sous-$A$-algèbres finies. Si nous démontrons les résultats pour chaque
$B_i$, alors le résultat en découle pour $B$. Voir Algèbre, lemme
\ref{algebra-lemma-colimit-henselian}. Si $A \to B$ est fini, alors $B$ est
un produit d'anneaux locaux henséliens d'après Algèbre, lemme
\ref{algebra-lemma-finite-over-henselian}. Comme $B$ est un domaine, nous
voyons que $B$ est un anneau local. L'idéal maximal de $B$ est au-dessus de
l'idéal maximal de $A$ par la remontée pour $A \to B$
(Algèbre, lemme \ref{algebra-lemma-integral-going-up}).
Si $A$ est strictement hensélien, alors l'extension de corps
$\kappa(\mathfrak m_B)/\kappa(\mathfrak m_A)$, étant algébrique, doit être
purement inséparable. Bien entendu, $\kappa(\mathfrak m_B)$ est alors
séparablement algébriquement clos et $B$ est strictement hensélien.
\end{proof}

\begin{theorem}[Olivier]
\label{theorem-olivier}
Soit $A \to B$ un homomorphisme local d'anneaux locaux.
Si $A$ est strictement hensélien et $A \to B$ est faiblement étale, alors
$A = B$.
\end{theorem}

\begin{proof}
Nous allons montrer que, pour tout $\mathfrak p \subset A$, il existe un
unique idéal premier $\mathfrak q \subset B$ au-dessus de $\mathfrak p$ et
que $\kappa(\mathfrak p) = \kappa(\mathfrak q)$. Ceci implique que
$B \otimes_A B \to B$ est bijectif sur les spectres, ainsi que surjectif et
plat. C'est donc un isomorphisme, par exemple d'après la description des
idéaux purs dans Algèbre, lemme
\ref{algebra-lemma-pure-open-closed-specializations}.
Ainsi $A \to B$ est un épimorphisme d'anneaux fidèlement plat. Nous obtenons
$A = B$ par Algèbre, lemme
\ref{algebra-lemma-faithfully-flat-epimorphism}.

\medskip\noindent
Remarquons que l'anneau fibre $B \otimes_A \kappa(\mathfrak p)$ est une
colimite d'extensions étales de $\kappa(\mathfrak p)$ d'après les
lemmes \ref{lemma-base-change-weakly-etale} et
\ref{lemma-absolutely-flat-over-field}. Ainsi, s'il existe plus d'un idéal
premier au-dessus de $\mathfrak p$, ou si
$\kappa(\mathfrak p) \not = \kappa(\mathfrak q)$ pour un certain
$\mathfrak q$, alors $B \otimes_A L$ possède un idempotent non trivial pour
une certaine extension algébrique (séparable)
$L/\kappa(\mathfrak p)$.

\medskip\noindent
Soit $L/\kappa(\mathfrak p)$ une extension algébrique de corps.
Soit $A' \subset L$ la clôture intégrale de $A/\mathfrak p$ dans $L$.
Par le lemme \ref{lemma-integral-over-henselian}, nous voyons que $A'$ est
un anneau local strictement hensélien dont le corps résiduel est une
extension purement inséparable du corps résiduel de $A$. Ainsi
$B \otimes_A A'$ est un anneau local d'après
Algèbre, lemme \ref{algebra-lemma-local-tensor-with-integral}.
D'autre part, $B \otimes_A A'$ est intégralement clos dans
$B \otimes_A L$ par le lemme \ref{lemma-normality-goes-up}.
Comme $B \otimes_A A'$ est local, il s'ensuit que l'anneau
$B \otimes_A L$ ne possède pas d'idempotents non triviaux, ce que nous
voulions démontrer.
\end{proof}







\section{Algèbres faiblement étales sur les corps}
\label{section-weakly-etale-over-field}

\noindent
Si $K$ est un corps, alors une algèbre $B$ est faiblement étale sur $K$
si et seulement si elle est une colimite filtrante de $K$-algèbres étales.
C'est le lemme \ref{lemma-absolutely-flat-over-field}.

\begin{lemma}
\label{lemma-class-weakly-etale-over-field}
Soit $K$ un corps. Si $B$ est faiblement étale sur $K$, alors
\begin{enumerate}
\item $B$ est réduite,
\item $B$ est entière sur $K$,
\item toute sous-$K$-algèbre de type fini de $B$ est un produit fini
d'extensions finies séparables de $K$,
\item $B$ est un corps si et seulement si $B$ ne possède pas d'idempotent non trivial
et, dans ce cas, c'est une extension algébrique séparable de $K$,
\item toute sous-$K$-algèbre ou algèbre quotient de $B$ est faiblement étale sur $K$,
\item si $B'$ est faiblement étale sur $K$, alors $B \otimes_K B'$ est
faiblement étale sur $K$.
\end{enumerate}
\end{lemma}

\begin{proof}
L'assertion (1) découle du lemme \ref{lemma-absolutely-flat-over-absolutely-flat},
mais elle découle bien sûr aussi de l'assertion (3).
L'assertion (3) découle du lemme \ref{lemma-absolutely-flat-over-field}
et du fait que les $K$-algèbres étales sont des produits finis
d'extensions finies séparables de $K$, voir
Algèbre, lemme \ref{algebra-lemma-etale-over-field}.
L'assertion (3) implique (2).
L'assertion (4) découle de (3), car un produit de corps est
un corps si et seulement s'il ne possède pas d'idempotent non trivial.

\medskip\noindent
Si $S \subset B$ est une sous-algèbre, alors c'est la colimite
filtrante de ses sous-algèbres de type fini, qui sont
toutes étales sur $K$ d'après ce qui précède, et donc
$S$ est faiblement étale sur $K$ d'après le
lemme \ref{lemma-absolutely-flat-over-field}.
Si $B \to Q$ est une algèbre quotient, alors
$Q$ est la colimite filtrante de $K$-algèbres quotients de
produits finis $\prod_{i \in I} L_i$ d'extensions finies séparables
$L_i/K$. Un tel quotient est de la forme $\prod_{i \in J} L_i$
pour un certain sous-ensemble $J \subset I$, et donc le résultat
vaut pour les quotients par le même raisonnement.

\medskip\noindent
L'assertion sur les produits tensoriels s'obtient de manière analogue
ou en combinant les lemmes \ref{lemma-base-change-weakly-etale} et
\ref{lemma-composition-weakly-etale}.
\end{proof}

\begin{lemma}
\label{lemma-max-weakly-etale-subalgebra}
Soit $K$ un corps. Soit $A$ une $K$-algèbre. Il existe
une sous-$K$-algèbre faiblement étale maximale $B_{max} \subset A$.
\end{lemma}

\begin{proof}
Soient $B_1, B_2 \subset A$ des sous-$K$-algèbres faiblement étales.
Alors $B_1 \otimes_K B_2$ est faiblement étale sur $K$,
de même que l'image de $B_1 \otimes_K B_2 \to A$
(lemme \ref{lemma-class-weakly-etale-over-field}).
Ainsi, l'ensemble $\mathcal{B}$ des sous-$K$-algèbres faiblement étales
$B \subset A$ est filtrant et la colimite
$B_{max} = \colim_{B \in \mathcal{B}} B$ est
une $K$-algèbre faiblement étale d'après le lemme \ref{lemma-when-weakly-etale}.
L'image de $B_{max} \to A$ est donc faiblement étale sur $K$
(par le lemme précédent). Il s'ensuit que cette image appartient à $\mathcal{B}$
et donc que $\mathcal{B}$ possède un élément maximal
(et que l'image est égale à $B_{max}$).
\end{proof}

\begin{lemma}
\label{lemma-properties-of-max-weakly-etale-subalgebra}
Soit $K$ un corps. Pour une $K$-algèbre $A$, notons $B_{max}(A)$
la sous-$K$-algèbre faiblement étale maximale de $A$ comme dans le
lemme \ref{lemma-max-weakly-etale-subalgebra}. Alors
\begin{enumerate}
\item tout homomorphisme de $K$-algèbres $A' \to A$ induit un homomorphisme de $K$-algèbres
$B_{max}(A') \to B_{max}(A)$,
\item si $A' \subset A$, alors $B_{max}(A') = B_{max}(A) \cap A'$,
\item si $A = \colim A_i$ est une colimite filtrante, alors
$B_{max}(A) = \colim B_{max}(A_i)$,
\item l'homomorphisme $B_{max}(A) \to B_{max}(A_{red})$ est un isomorphisme,
\item $B_{max}(A_1 \times \ldots \times A_n) =
B_{max}(A_1) \times \ldots \times B_{max}(A_n)$,
\item si $A$ ne possède pas d'idempotent non trivial, alors $B_{max}(A)$ est un
corps et une extension algébrique séparable de $K$,
\item ajouter d'autres assertions ici.
\end{enumerate}
\end{lemma}

\begin{proof}
Démonstration de (1). Ceci est vrai car l'image de $B_{max}(A') \to A$
est faiblement étale sur $K$ d'après le lemme \ref{lemma-class-weakly-etale-over-field}.

\medskip\noindent
Démonstration de (2). D'après (1), nous avons $B_{max}(A') \subset B_{max}(A)$.
Réciproquement, $B_{max}(A) \cap A'$ est une $K$-algèbre faiblement étale
d'après le lemme \ref{lemma-class-weakly-etale-over-field}, et elle est donc
contenue dans $B_{max}(A')$.

\medskip\noindent
Preuve de (3). D'après (1), il existe un homomorphisme $\colim B_{max}(A_i) \to A$
qui est injectif car le système est filtrant et
$B_{max}(A_i) \subset A_i$. La colimite $\colim B_{max}(A_i)$
est faiblement étale sur $K$ d'après le
lemme \ref{lemma-when-weakly-etale}.
Nous obtenons donc un homomorphisme injectif $\colim B_{max}(A_i) \to B_{max}(A)$.
Supposons que $a \in B_{max}(A)$. Alors $a$ engendre une sous-$K$-algèbre
de présentation finie $B \subset B_{max}(A)$.
D'après Algèbre, lemme \ref{algebra-lemma-characterize-finite-presentation},
il existe un $i$ et un homomorphisme de $K$-algèbres $f : B \to A_i$
relevant l'homomorphisme donné $B \to A$.
Comme $B$ est faiblement étale d'après le
lemme \ref{lemma-class-weakly-etale-over-field},
nous voyons que $f(B) \subset B_{max}(A_i)$ et nous concluons que $a$
appartient à l'image de $\colim B_{max}(A_i) \to B_{max}(A)$.

\medskip\noindent
Démonstration de (4). Écrivons $B_{max}(A_{red}) = \colim B_i$
comme une colimite filtrante de $K$-algèbres étales
(lemme \ref{lemma-absolutely-flat-over-field}).
D'après Algèbre, lemme \ref{algebra-lemma-smooth-strong-lift},
pour tout $i$, il existe un homomorphisme de $K$-algèbres $f_i : B_i \to A$
relevant l'homomorphisme donné $B_i \to A_{red}$. Il s'ensuit que l'homomorphisme
canonique $B_{max}(A_{red}) \to B_{max}(A)$ est surjectif.
Le noyau est constitué d'éléments nilpotents et est donc nul,
puisque $B_{max}(A_{red})$ est réduite
(lemme \ref{lemma-class-weakly-etale-over-field}).

\medskip\noindent
Démonstration de (5). Omise.

\medskip\noindent
Démonstration de (6). Elle découle de l'assertion (4) du
lemme \ref{lemma-class-weakly-etale-over-field}.
\end{proof}

\begin{lemma}
\label{lemma-change-fields-max-weakly-etale-subalgebra}
Soit $L/K$ une extension de corps. Soit $A$ une $K$-algèbre.
Soit $B \subset A$ la sous-$K$-algèbre faiblement étale maximale de
$A$ comme dans le lemme \ref{lemma-max-weakly-etale-subalgebra}.
Alors $B \otimes_K L$ est la sous-$L$-algèbre faiblement étale maximale
de $A \otimes_K L$.
\end{lemma}

\begin{proof}
Pour une algèbre $A$ sur $K$, nous notons $B_{max}(A/K)$
la sous-$K$-algèbre faiblement étale maximale de $A$.
De même, nous notons $B_{max}(A'/L)$ la sous-$L$-algèbre faiblement étale
maximale de $A'$ si $A'$ est une $L$-algèbre.
Comme $B_{max}(A/K) \otimes_K L$ est faiblement étale sur $L$
(lemme \ref{lemma-base-change-weakly-etale})
et comme $B_{max}(A/K) \otimes_K L \subset A \otimes_K L$,
nous obtenons un homomorphisme injectif canonique
$$
B_{max}(A/K) \otimes_K L \to B_{max}((A \otimes_K L)/L)
$$
Le lemme affirme que cet homomorphisme est un isomorphisme.

\medskip\noindent
Pour démontrer le lemme pour $L$ et notre $K$-algèbre $A$, il suffit de
le démontrer pour une extension de corps quelconque $L'$ de $L$. En effet, nous avons
la factorisation
$$
B_{max}(A/K) \otimes_K L' \to
B_{max}((A \otimes_K L)/L) \otimes_L L' \to
B_{max}((A \otimes_K L')/L')
$$
et donc la composée ne peut être surjective sans que
$B_{max}(A/K) \otimes_K L \to B_{max}((A \otimes_K L)/L)$
soit surjectif. Nous pouvons donc supposer que $L$ est algébriquement clos.

\medskip\noindent
Réduction au cas d'une $K$-algèbre de type fini. Nous pouvons écrire que $A$
est la colimite filtrante de ses sous-$K$-algèbres de type fini.
En utilisant le lemme \ref{lemma-properties-of-max-weakly-etale-subalgebra},
nous voyons qu'il suffit de démontrer le lemme pour les
$K$-algèbres de type fini.

\medskip\noindent
Supposons que $A$ soit une $K$-algèbre de type fini. Puisque le noyau
de $A \to A_{red}$ est nilpotent, il en est de même de celui de
$A \otimes_K L \to A_{red} \otimes_K L$. Alors
$$
B_{max}((A \otimes_K L)/L) \to B_{max}((A_{red} \otimes_K L)/L)
$$
est injectif car le noyau est nilpotent et la
$L$-algèbre faiblement étale $B_{max}((A \otimes_K L)/L)$ est réduite
(lemme \ref{lemma-class-weakly-etale-over-field}).
Puisque $B_{max}(A/K) = B_{max}(A_{red}/K)$
d'après le lemme \ref{lemma-properties-of-max-weakly-etale-subalgebra},
nous concluons
qu'il suffit de démontrer le lemme pour $A_{red}$.

\medskip\noindent
Supposons que $A$ soit une $K$-algèbre réduite de type fini.
Soit $Q = Q(A)$ l'anneau total des fractions de $A$.
Alors $A \subset Q$ et $A \otimes_K L \subset Q \otimes_K L$,
et par conséquent
$$
B_{max}(A/K) = A \cap B_{max}(Q/K)
$$
et
$$
B_{max}((A \otimes_K L)/L) =
(A \otimes_K L) \cap B_{max}((Q \otimes_K L)/L)
$$
d'après le lemme \ref{lemma-properties-of-max-weakly-etale-subalgebra}.
Comme $-\otimes_K L$ est un foncteur exact, il s'ensuit que
si nous démontrons le résultat pour $Q$, alors il en découle pour $A$.
Comme $Q$ est un produit fini de corps (Algèbre, lemmes
\ref{algebra-lemma-total-ring-fractions-no-embedded-points},
\ref{algebra-lemma-minimal-prime-reduced-ring},
\ref{algebra-lemma-Noetherian-irreducible-components} et
\ref{algebra-lemma-Noetherian-permanence})
et comme $B_{max}$ commute aux produits
(lemme \ref{lemma-properties-of-max-weakly-etale-subalgebra}),
il suffit de démontrer le lemme lorsque $A$ est un corps.

\medskip\noindent
Supposons que $A$ soit un corps. L'argument du troisième paragraphe de la preuve
nous ramène au cas où $A$ est de type fini sur $K$.
(En fait, la façon dont nous nous sommes ramenés au cas d'un corps produit
une extension de corps de type fini de $K$.)

\medskip\noindent
Supposons que $A$ soit une extension de corps de type fini de $K$.
Alors $K' = B_{max}(A/K)$ est un corps algébrique séparable sur $K$ d'après
l'assertion (6) du lemme \ref{lemma-properties-of-max-weakly-etale-subalgebra}.
Par conséquent, $K'$ est une extension finie séparable
de $K$ et $A$ est géométriquement irréductible sur $K'$ d'après
Algèbre, lemme \ref{algebra-lemma-make-geometrically-irreducible}.
Puisque $L$ est algébriquement clos et $K'/K$ finie séparable,
nous voyons que
$$
K' \otimes_K L \to \prod\nolimits_{\sigma \in \Hom_K(K', L)} L,\quad
\alpha \otimes \beta \mapsto (\sigma(\alpha)\beta)_\sigma
$$
est un isomorphisme
(Corps, lemme \ref{fields-lemma-finite-separable-tensor-alg-closed}).
Nous concluons que
$$
A \otimes_K L = A \otimes_{K'} (K' \otimes_K L) =
\prod\nolimits_{\sigma \in \Hom_K(K', L)} A \otimes_{K', \sigma} L
$$
Comme $A$ est géométriquement irréductible sur $K'$, nous voyons que
$A \otimes_{K', \sigma} L$ possède un unique idéal premier minimal.
Puisque $L$ est algébriquement clos, il s'ensuit que
$B_{max}((A \otimes_{K', \sigma} L)/L) = L$
car cette $L$-algèbre est un corps algébrique sur $L$ d'après
l'assertion (6) du lemme \ref{lemma-properties-of-max-weakly-etale-subalgebra}.
Il s'ensuit que la sous-$K' \otimes_K L$-algèbre faiblement étale maximale
de $A \otimes_K L$ est $K' \otimes_K L$, car nous pouvons décomposer
ces sous-algèbres en produits comme ci-dessus. Ainsi, l'inclusion
$K' \otimes_K L \subset B_{max}((A \otimes_K L)/L)$ est une
égalité : l'homomorphisme d'anneaux $K' \otimes_K L \to B_{max}((A \otimes_K L)/L)$
est faiblement étale d'après le
lemme \ref{lemma-weakly-etale-permanence}.
\end{proof}












\section{Irréductibilité locale}
\label{section-unibranch}

\noindent
La définition suivante semble être celle qui est généralement admise.
Pour l'interpréter, observons que si $A \subset B$ est une extension entière d'anneaux
locaux intègres, alors $A \to B$ est un homomorphisme local d'anneaux par
montée (Algèbre, lemme \ref{algebra-lemma-integral-going-up}).

\begin{definition}
\label{definition-unibranch}
\begin{reference}
\cite[Chapter 0 (23.2.1)]{EGA4}
\end{reference}
Soit $A$ un anneau local. Nous disons que $A$ est {\it unibranche}
si la réduction $A_{red}$ est un domaine et si la clôture intégrale
$A'$ de $A_{red}$ dans son corps des fractions est locale.
Nous disons que $A$ est {\it géométriquement unibranche} si $A$ est unibranche
et si, en outre, le corps résiduel de $A'$ est purement inséparable sur
le corps résiduel de $A$.
\end{definition}

\noindent
Soit $A$ un anneau local. Voici une formulation équivalente :
\begin{enumerate}
\item $A$ est unibranche si $A$ possède un unique idéal premier minimal $\mathfrak p$ et si
la clôture intégrale de $A/\mathfrak p$ dans son corps des fractions est un
anneau local, et
\item $A$ est géométriquement unibranche si $A$ possède un unique idéal premier minimal
$\mathfrak p$ et si la clôture intégrale de $A/\mathfrak p$ dans son
corps des fractions est un anneau local dont le corps résiduel est purement inséparable
sur le corps résiduel de $A$.
\end{enumerate}
Un anneau local normal est géométriquement unibranche
(cela découle de la définition \ref{definition-unibranch} et de
Algèbre, définition \ref{algebra-definition-ring-normal}).
Les lemmes \ref{lemma-unibranch} et \ref{lemma-geometrically-unibranch}
indiquent que la propriété d'être (géométriquement) unibranche
est une propriété raisonnable à considérer.

\begin{lemma}
\label{lemma-branches}
Soit $A$ un anneau local. Supposons que $A$ possède un nombre fini d'idéaux premiers minimaux.
Soit $A'$ la clôture intégrale de $A$ dans l'anneau total des fractions
de $A_{red}$.
Soit $A^h$ l'hensélisé de $A$.
Considérons les homomorphismes
$$
\Spec(A') \leftarrow \Spec((A')^h) \rightarrow \Spec(A^h)
$$
où $(A')^h = A' \otimes_A A^h$. Alors
\begin{enumerate}
\item la flèche de gauche induit une bijection sur les idéaux maximaux,
\item la flèche de droite induit une bijection sur les idéaux premiers minimaux,
\item tout idéal premier minimal de $(A')^h$ est contenu dans un unique
idéal maximal, et tout idéal maximal contient exactement un idéal premier minimal.
\end{enumerate}
\end{lemma}

\begin{proof}
Soit $I \subset A$ l'idéal des éléments nilpotents.
Nous avons $(A/I)^h = A^h/IA^h$ d'après
(Algèbre, lemme \ref{algebra-lemma-quotient-henselization}).
Les spectres de $A$, $A^h$, $A'$ et $(A')^h$ sont les mêmes
que les spectres de $A/I$, $A^h/IA^h$, $A'$ et
$(A')^h = A' \otimes_{A/I} A^h/IA^h$. Ainsi,
nous pouvons remplacer $A$ par $A_{red} = A/I$ et supposer $A$ réduit.
Nous avons alors $A \subset A'$, ce que nous utiliserons ci-dessous sans autre mention.

\medskip\noindent
Démonstration de (1). Comme $A'$ est entière sur $A$, nous voyons que $(A')^h$
est entière sur $A^h$. Par montée
(Algèbre, lemme \ref{algebra-lemma-integral-going-up}),
tout idéal maximal de $A'$, resp.\ de $(A')^h$, est au-dessus
de l'idéal maximal $\mathfrak m$, resp.\ $\mathfrak m^h$, de
$A$, resp.\ de $A^h$. Ainsi, (1) découle de l'isomorphisme
$$
(A')^h \otimes_{A^h} \kappa^h =
A' \otimes_A A^h \otimes_{A^h} \kappa^h = A' \otimes_A \kappa
$$
car l'extension de corps résiduels $\kappa^h/\kappa$
induite par $A \to A^h$ est triviale. Nous utiliserons ci-dessous le fait que
l'anneau affiché est entier sur un corps, donc que le spectre de
cet anneau est un espace profini, voir
Algèbre, lemmes \ref{algebra-lemma-integral-over-field} et
\ref{algebra-lemma-ring-with-only-minimal-primes}.

\medskip\noindent
Démonstration de (3). L'anneau $A'$ est normal et est en fait un
produit fini de domaines normaux, voir
Algèbre, lemme \ref{algebra-lemma-characterize-reduced-ring-normal}.
Comme $A^h$ est une colimite filtrante de $A$-algèbres étales,
$(A')^h$ est une colimite filtrante de $A'$-algèbres étales ;
par conséquent, $(A')^h$ est un anneau normal d'après
Algèbre, lemmes \ref{algebra-lemma-normal-goes-up} et
\ref{algebra-lemma-colimit-normal-ring}.
Ainsi, tout anneau local de $(A')^h$ est un domaine normal,
et nous voyons que tout idéal maximal contient un unique idéal premier minimal.
En appliquant le lemme \ref{lemma-integral-over-henselian-pair}
à $A^h \to (A')^h$, nous voyons que $((A')^h, \mathfrak m(A')^h)$ est
une paire hensélienne. Si $\mathfrak q \subset (A')^h$
est un idéal premier minimal (ou un idéal premier quelconque), alors l'intersection de
$V(\mathfrak q)$ avec $V(\mathfrak m (A')^h)$ est connexe
d'après le lemme \ref{lemma-irreducible-henselian-pair-connected}.
Comme $V(\mathfrak m (A')^h) = \Spec((A')^h \otimes \kappa^h)$
est un espace profini, nous voyons qu'il existe un unique
idéal maximal contenant $\mathfrak q$.

\medskip\noindent
Démonstration de (2). Les idéaux premiers minimaux de $A'$ sont exactement les idéaux premiers
au-dessus d'un idéal premier minimal de $A$ (par construction).
Comme $A' \to (A')^h$ est plat, par descente
(Algèbre, lemme \ref{algebra-lemma-flat-going-down}),
tout idéal premier minimal de $(A')^h$ est au-dessus d'un idéal premier minimal
de $A'$. Réciproquement, tout idéal premier de $(A')^h$ au-dessus d'un
idéal premier minimal de $A'$ est minimal, car $(A')^h$ est une colimite
filtrante d'algèbres étales, donc quasi-finies, sur $A'$ (petit détail omis).
Nous concluons que les idéaux premiers minimaux de $(A')^h$
sont exactement les idéaux premiers au-dessus d'un idéal premier minimal de $A$.
De même, les idéaux premiers minimaux de $A^h$ sont exactement les
idéaux premiers au-dessus des idéaux premiers minimaux de $A$.
Par construction, nous avons $A' \otimes_A Q(A) = Q(A)$, où $Q(A)$
est l'anneau total des fractions de notre anneau local réduit $A$.
Bien sûr, $Q(A)$ est le produit fini des corps résiduels des
idéaux premiers minimaux de $A$. Il s'ensuit que
$$
(A')^h \otimes_A Q(A) = A^h \otimes_A A' \otimes_A Q(A) = A^h \otimes_A Q(A)
$$
La discussion ci-dessus montre que le spectre de l'anneau
de gauche est l'ensemble des idéaux premiers minimaux de $(A')^h$
et que le spectre de l'anneau de droite
est l'ensemble des idéaux premiers minimaux de $A^h$. Ceci achève la démonstration.
\end{proof}

\begin{lemma}
\label{lemma-unibranch}
\begin{reference}
\cite[Chapter IV Proposition 18.6.12]{EGA4}
\end{reference}
Soit $A$ un anneau local. Soit $A^h$ l'hensélisé de $A$.
Les assertions suivantes sont équivalentes :
\begin{enumerate}
\item $A$ est unibranche, et
\item $A^h$ possède un unique idéal premier minimal.
\end{enumerate}
\end{lemma}

\begin{proof}
Ceci découle du lemme \ref{lemma-branches}, mais nous allons aussi donner une
démonstration directe. Notons $\mathfrak m$ l'idéal maximal de l'anneau $A$.
Rappelons que le corps résiduel $\kappa = A/\mathfrak m$ est le même que
le corps résiduel de $A^h$.

\medskip\noindent
Supposons (2). Soit $\mathfrak p^h$ l'unique idéal premier minimal de
$A^h$. La platitude de $A \to A^h$ implique que
$\mathfrak p = A \cap \mathfrak p^h$ est l'unique idéal premier
minimal de $A$ (par descente, voir
Algèbre, lemme \ref{algebra-lemma-flat-going-down}).
De plus, puisque $A^h/\mathfrak pA^h = (A/\mathfrak p)^h$ (voir
Algèbre, lemme \ref{algebra-lemma-quotient-henselization})
est réduit d'après le lemme \ref{lemma-henselization-reduced},
nous voyons que $\mathfrak p^h = \mathfrak pA^h$.
Soit $A'$ la clôture intégrale de $A/\mathfrak p$ dans son corps des
fractions. Nous devons montrer que $A'$ est local.
Comme $A \to A'$ est entier, tout idéal maximal de $A'$ est au-dessus de
$\mathfrak m$ (par montée pour les homomorphismes entiers d'anneaux, voir
Algèbre, lemme \ref{algebra-lemma-integral-going-up}).
Si $A'$ n'est pas local, nous pouvons trouver des idéaux maximaux distincts
$\mathfrak m_1$, $\mathfrak m_2$. Choisissons des éléments $f_1, f_2 \in A'$
tels que $f_i \in \mathfrak m_i$ et $f_i \not \in \mathfrak m_{3 - i}$.
Nous obtenons une sous-algèbre finie $B = A/\mathfrak p[f_1, f_2] \subset A'$
ayant des idéaux maximaux distincts $B \cap \mathfrak m_i$, $i = 1, 2$.
Remarquons que les inclusions
$$
A/\mathfrak p \subset B \subset \kappa(\mathfrak p)
$$
donnent, après tensorisation par l'homomorphisme plat d'anneaux $A \to A^h$, les inclusions
$$
A^h/\mathfrak p^h \subset
B \otimes_A A^h \subset
\kappa(\mathfrak p) \otimes_A A^h \subset
\kappa(\mathfrak p^h)
$$
la dernière inclusion venant de ce que
$\kappa(\mathfrak p) \otimes_A A^h =
\kappa(\mathfrak p) \otimes_{A/\mathfrak p} A^h/\mathfrak p^h$
est une localisation du domaine $A^h/\mathfrak p^h$.
Remarquons que $B \otimes_A \kappa$ possède au moins deux idéaux maximaux,
car $B/\mathfrak mB$ possède deux idéaux maximaux. Par conséquent, comme
$A^h$ est hensélien, nous voyons que
$B \otimes_A A^h$ est un produit de $\geq 2$ anneaux locaux, voir
Algèbre, lemme \ref{algebra-lemma-mop-up}.
Mais nous venons de voir que $B \otimes_A A^h$ est un sous-anneau d'un domaine,
ce qui donne une contradiction.

\medskip\noindent
Supposons (1). Soit $\mathfrak p \subset A$ l'unique idéal premier
minimal et soit $A'$ la clôture intégrale de $A/\mathfrak p$
dans son corps des fractions. Soit $A \to B$ un homomorphisme local d'anneaux locaux
induisant un isomorphisme de corps résiduels et qui est une
localisation d'une $A$-algèbre étale. En particulier, $\mathfrak m_B$
est l'unique idéal premier contenant $\mathfrak m B$. Alors $B' = A' \otimes_A B$
est entière sur $B$, et l'hypothèse que $A \to A'$ est local
implique que $B'$ est local
(Algèbre, lemme \ref{algebra-lemma-local-tensor-with-integral}).
D'autre part, $A' \to B'$ est la localisation
d'un homomorphisme étale d'anneaux ; ainsi, $B'$ est normal, voir
Algèbre, lemme \ref{algebra-lemma-normal-goes-up}.
Par conséquent, $B'$ est un domaine normal (local). Enfin, nous avons
$$
B/\mathfrak pB \subset B \otimes_A \kappa(\mathfrak p)
= B' \otimes_{A'} (\text{corps des fractions de }A') \subset
\text{corps des fractions de }B'
$$
Ainsi, $B/\mathfrak pB$ est un domaine, ce qui implique que $B$ possède un unique
idéal premier minimal (car, par platitude de $A \to B$, tous ces idéaux doivent être
au-dessus de $\mathfrak p$). Comme $A^h$ est une colimite filtrante des
anneaux locaux $B$, il s'ensuit que $A^h$ possède un unique idéal premier minimal.
En effet, si $fg = 0$ dans $A^h$ pour des éléments non nilpotents
$f, g$, nous pouvons trouver un $B$ comme ci-dessus contenant à la fois $f$ et $g$,
ce qui conduit à une contradiction.
\end{proof}

\begin{lemma}
\label{lemma-geometric-branches}
Soit $(A, \mathfrak m, \kappa)$ un anneau local.
Supposons que $A$ possède un nombre fini d'idéaux premiers minimaux.
Soit $A'$ la clôture intégrale de $A$ dans l'anneau total des fractions
de $A_{red}$. Choisissons une clôture algébrique $\overline{\kappa}$ de
$\kappa$ et notons $\kappa^{sep} \subset \overline{\kappa}$
la clôture algébrique séparable de $\kappa$.
Soit $A^{sh}$ l'hensélisé strict de $A$
relatif à $\kappa^{sep}$.
Considérons les homomorphismes
$$
\Spec(A') \xleftarrow{c} \Spec((A')^{sh}) \xrightarrow{e} \Spec(A^{sh})
$$
où $(A')^{sh} = A' \otimes_A A^{sh}$. Alors
\begin{enumerate}
\item pour $\mathfrak m' \subset A'$ maximal, le corps résiduel
$\kappa'$ est algébrique sur $\kappa$ et la fibre de $c$
au-dessus de $\mathfrak m'$ s'identifie canoniquement
à $\Hom_\kappa(\kappa', \overline{\kappa})$,
\item la flèche de droite induit une bijection sur les idéaux premiers minimaux,
\item tout idéal premier minimal de $(A')^{sh}$ est contenu dans un unique
idéal maximal, et tout idéal maximal contient un unique idéal premier minimal.
\end{enumerate}
\end{lemma}

\begin{proof}
La démonstration est presque exactement la même que celle du lemme \ref{lemma-branches}.
Soit $I \subset A$ l'idéal des éléments nilpotents.
Nous avons $(A/I)^{sh} = A^{sh}/IA^{sh}$ d'après
(Algèbre, lemme \ref{algebra-lemma-quotient-henselization}).
Les spectres de $A$, $A^{sh}$, $A'$ et $(A')^{sh}$ sont les mêmes
que les spectres de $A/I$, $A^{sh}/IA^{sh}$, $A'$ et
$(A')^{sh} = A' \otimes_{A/I} A^{sh}/IA^{sh}$. Ainsi,
nous pouvons remplacer $A$ par $A_{red} = A/I$ et supposer $A$ réduit.
Nous avons alors $A \subset A'$, ce que nous utiliserons ci-dessous sans autre mention.

\medskip\noindent
Démonstration de (1). L'extension de corps $\kappa'/\kappa$ est algébrique
car $A'$ est entière sur $A$. Comme $A'$ est entière sur $A$,
nous voyons que $(A')^{sh}$ est entière sur $A^{sh}$. Par montée
(Algèbre, lemme \ref{algebra-lemma-integral-going-up}),
tout idéal maximal de $A'$, resp.\ de $(A')^{sh}$, est au-dessus
de l'idéal maximal $\mathfrak m$, resp.\ $\mathfrak m^{sh}$, de
$A$, resp.\ de $A^{sh}$. Nous avons
$$
(A')^{sh} \otimes_{A^{sh}} \kappa^{sep} =
A' \otimes_A A^{sh} \otimes_{A^{sh}} \kappa^{sep} =
(A' \otimes_A \kappa) \otimes_{\kappa} \kappa^{sep}
$$
car le corps résiduel de $A^{sh}$ est $\kappa^{sep}$.
Ainsi, la fibre de $c$ au-dessus de $\mathfrak m'$ est le spectre
de $\kappa' \otimes_\kappa \kappa^{sep}$.
Nous concluons que (1) est vraie, car il existe une bijection
$$
\Hom_\kappa(\kappa', \overline{\kappa}) \to
\Spec(\kappa' \otimes_\kappa \kappa^{sep}),\quad
\sigma \mapsto \Ker(
\sigma \otimes 1 : \kappa' \otimes_\kappa \kappa^{sep} \to \overline{\kappa}
)
$$
Nous utiliserons ci-dessous le fait que l'anneau affiché est entier sur un corps,
et donc que le spectre de cet anneau est un espace profini, voir
Algèbre, lemmes \ref{algebra-lemma-integral-over-field} et
\ref{algebra-lemma-ring-with-only-minimal-primes}.

\medskip\noindent
Démonstration de (3). L'anneau $A'$ est normal et est en fait un
produit fini de domaines normaux, voir
Algèbre, lemme \ref{algebra-lemma-characterize-reduced-ring-normal}.
Comme $A^{sh}$ est une colimite filtrante de $A$-algèbres étales,
$(A')^{sh}$ est une colimite filtrante de $A'$-algèbres étales ;
par conséquent, $(A')^{sh}$ est un anneau normal d'après
Algèbre, lemmes \ref{algebra-lemma-normal-goes-up} et
\ref{algebra-lemma-colimit-normal-ring}.
Ainsi, tout anneau local de $(A')^{sh}$ est un domaine normal,
et nous voyons que tout idéal maximal contient un unique idéal premier minimal.
En appliquant le lemme \ref{lemma-integral-over-henselian-pair}
à $A^{sh} \to (A')^{sh}$, nous voyons que $((A')^{sh}, \mathfrak m(A')^{sh})$
est une paire hensélienne. Si $\mathfrak q \subset (A')^{sh}$
est un idéal premier minimal (ou un idéal premier quelconque), alors l'intersection de
$V(\mathfrak q)$ avec $V(\mathfrak m (A')^{sh})$ est connexe
d'après le lemme \ref{lemma-irreducible-henselian-pair-connected}.
Comme $V(\mathfrak m (A')^{sh}) = \Spec((A')^{sh} \otimes \kappa^{sep})$
est un espace profini, nous voyons qu'il existe un unique
idéal maximal contenant $\mathfrak q$.

\medskip\noindent
Démonstration de (2). Les idéaux premiers minimaux de $A'$ sont exactement les idéaux premiers
au-dessus d'un idéal premier minimal de $A$ (par construction).
Comme $A' \to (A')^{sh}$ est plat, par descente
(Algèbre, lemme \ref{algebra-lemma-flat-going-down}),
tout idéal premier minimal de $(A')^{sh}$ est au-dessus d'un idéal premier minimal
de $A'$. Réciproquement, tout idéal premier de $(A')^{sh}$ au-dessus d'un
idéal premier minimal de $A'$ est minimal, car $(A')^{sh}$ est une colimite
filtrante d'algèbres étales, donc quasi-finies, sur $A'$ (petit détail omis).
Nous concluons que les idéaux premiers minimaux de $(A')^{sh}$
sont exactement les idéaux premiers au-dessus d'un idéal premier minimal de $A$.
De même, les idéaux premiers minimaux de $A^{sh}$ sont exactement les
idéaux premiers au-dessus des idéaux premiers minimaux de $A$.
Par construction, nous avons $A' \otimes_A Q(A) = Q(A)$, où $Q(A)$
est l'anneau total des fractions de notre anneau local réduit $A$.
Bien sûr, $Q(A)$ est le produit fini des corps résiduels des
idéaux premiers minimaux de $A$. Il s'ensuit que
$$
(A')^{sh} \otimes_A Q(A) =
A^{sh} \otimes_A A' \otimes_A Q(A) = A^{sh} \otimes_A Q(A)
$$
La discussion ci-dessus montre que le spectre de l'anneau
de gauche est l'ensemble des idéaux premiers minimaux de $(A')^{sh}$
et que le spectre de l'anneau de droite
est l'ensemble des idéaux premiers minimaux de $A^{sh}$. Ceci achève la démonstration.
\end{proof}

\begin{lemma}
\label{lemma-geometrically-unibranch}
\begin{reference}
\cite[Lemma 2.2]{Etale-coverings} and
\cite[Chapter IV Proposition 18.8.15]{EGA4}
\end{reference}
Soit $A$ un anneau local. Soit $A^{sh}$ un hensélisé strict de $A$.
Les assertions suivantes sont équivalentes :
\begin{enumerate}
\item $A$ est géométriquement unibranche, et
\item $A^{sh}$ possède un unique idéal premier minimal.
\end{enumerate}
\end{lemma}

\begin{proof}
Ceci découle du lemme \ref{lemma-geometric-branches},
mais nous allons aussi donner une démonstration directe ; cette démonstration
directe est presque exactement la même que la démonstration directe du
lemme \ref{lemma-unibranch}.
Notons $\mathfrak m$ l'idéal maximal de l'anneau $A$.
Notons $\kappa$, $\kappa^{sh}$ les corps résiduels de $A$, $A^{sh}$.

\medskip\noindent
Supposons (2). Soit $\mathfrak p^{sh}$ l'unique idéal premier minimal de
$A^{sh}$. La platitude de $A \to A^{sh}$ implique que
$\mathfrak p = A \cap \mathfrak p^{sh}$ est l'unique idéal premier
minimal de $A$ (par descente, voir
Algèbre, lemme \ref{algebra-lemma-flat-going-down}).
De plus, puisque $A^{sh}/\mathfrak pA^{sh} = (A/\mathfrak p)^{sh}$ (voir
Algèbre, lemme \ref{algebra-lemma-quotient-strict-henselization})
est réduit d'après le lemme \ref{lemma-henselization-reduced},
nous voyons que $\mathfrak p^{sh} = \mathfrak pA^{sh}$.
Soit $A'$ la clôture intégrale de $A/\mathfrak p$ dans son corps des
fractions. Nous devons montrer que $A'$ est local et que son corps
résiduel est purement inséparable sur $\kappa$.
Comme $A \to A'$ est entier, tout idéal maximal de $A'$ est au-dessus de
$\mathfrak m$ (par montée pour les homomorphismes entiers d'anneaux, voir
Algèbre, lemme \ref{algebra-lemma-integral-going-up}).
Si $A'$ n'est pas local, nous pouvons trouver des idéaux maximaux distincts
$\mathfrak m_1$, $\mathfrak m_2$. En choisissant des éléments $f_1, f_2 \in A'$
tels que $f_i \in \mathfrak m_i, f_i \not \in \mathfrak m_{3 - i}$, nous obtenons
une sous-algèbre finie $B = A[f_1, f_2] \subset A'$ ayant des idéaux maximaux
distincts $B \cap \mathfrak m_i$, $i = 1, 2$. Si $A'$ est local d'idéal maximal
$\mathfrak m'$, mais que $A/\mathfrak m \subset A'/\mathfrak m'$
n'est pas purement inséparable, alors nous pouvons trouver $f \in A'$ dont l'image dans
$A'/\mathfrak m'$ engendre une extension finie non purement inséparable
de $A/\mathfrak m$, et nous obtenons une sous-algèbre locale finie $B = A[f] \subset A'$
dont le corps résiduel n'est pas une extension purement inséparable de $A/\mathfrak m$.
Remarquons que les inclusions
$$
A/\mathfrak p \subset B \subset \kappa(\mathfrak p)
$$
donnent, après tensorisation par l'homomorphisme plat d'anneaux $A \to A^{sh}$, les inclusions
$$
A^{sh}/\mathfrak p^{sh} \subset
B \otimes_A A^{sh} \subset
\kappa(\mathfrak p) \otimes_A A^{sh} \subset
\kappa(\mathfrak p^{sh})
$$
la dernière inclusion venant de ce que
$\kappa(\mathfrak p) \otimes_A A^{sh} =
\kappa(\mathfrak p) \otimes_{A/\mathfrak p} A^{sh}/\mathfrak p^{sh}$
est une localisation du domaine $A^{sh}/\mathfrak p^{sh}$.
Remarquons que $B \otimes_A \kappa^{sh}$ possède au moins deux idéaux maximaux,
car $B/\mathfrak mB$ possède soit deux idéaux maximaux, soit un seul dont le
corps résiduel n'est pas purement inséparable sur $\kappa$, et parce que
$\kappa^{sh}$ est séparablement algébriquement clos. Par conséquent, comme
$A^{sh}$ est strictement hensélien, nous voyons que
$B \otimes_A A^{sh}$ est un produit de $\geq 2$ anneaux locaux, voir
Algèbre, lemme \ref{algebra-lemma-mop-up-strictly-henselian}.
Mais nous venons de voir que $B \otimes_A A^{sh}$ est un sous-anneau d'un domaine,
ce qui donne une contradiction.

\medskip\noindent
Supposons (1). Soit $\mathfrak p \subset A$ l'unique idéal premier
minimal et soit $A'$ la clôture intégrale de $A/\mathfrak p$
dans son corps des fractions. Soit $A \to B$ un homomorphisme local d'anneaux locaux qui est une
localisation d'une $A$-algèbre étale. En particulier, $\mathfrak m_B$
est l'unique idéal premier contenant $\mathfrak m B$. Alors $B' = A' \otimes_A B$
est entière sur $B$, et l'hypothèse que $A \to A'$ est local
avec une extension de corps résiduels purement inséparable implique que $B'$
est local (Algèbre, lemme \ref{algebra-lemma-local-tensor-with-integral}).
D'autre part, $A' \to B'$ est la localisation
d'un homomorphisme étale d'anneaux ; ainsi, $B'$ est normal, voir
Algèbre, lemme \ref{algebra-lemma-normal-goes-up}.
Par conséquent, $B'$ est un domaine normal (local). Enfin, nous avons
$$
B/\mathfrak pB \subset B \otimes_A \kappa(\mathfrak p)
= B' \otimes_{A'} (\text{corps des fractions de }A') \subset
\text{corps des fractions de }B'
$$
Ainsi, $B/\mathfrak pB$ est un domaine, ce qui implique que $B$ possède un unique
idéal premier minimal (car, par platitude de $A \to B$, tous ces idéaux doivent être
au-dessus de $\mathfrak p$). Comme $A^{sh}$ est une colimite filtrante des
anneaux locaux $B$, il s'ensuit que $A^{sh}$ possède un unique idéal premier minimal.
En effet, si $fg = 0$ dans $A^{sh}$ pour des éléments non nilpotents
$f, g$, nous pouvons trouver un $B$ comme ci-dessus contenant à la fois $f$ et $g$,
ce qui conduit à une contradiction.
\end{proof}

\begin{definition}
\label{definition-number-of-branches}
Soit $A$ un anneau local d'hensélisé $A^h$ et
d'hensélisé strict $A^{sh}$. Le {\it nombre de branches de $A$}
est le nombre d'idéaux premiers minimaux de $A^h$ s'il est fini, et $\infty$
sinon. Le {\it nombre de branches géométriques de $A$}
est le nombre d'idéaux premiers minimaux de $A^{sh}$ s'il est fini, et $\infty$
sinon.
\end{definition}

\noindent
Précisons la relation avec la définition \ref{definition-unibranch}.

\begin{lemma}
\label{lemma-number-of-branches-1}
Soit $(A, \mathfrak m, \kappa)$ un anneau local.
\begin{enumerate}
\item Si $A$ possède une infinité d'idéaux premiers minimaux, alors
le nombre de branches (géométriques) de $A$ est $\infty$.
\item Le nombre de branches de $A$ est $1$ si et seulement si
$A$ est unibranche.
\item Le nombre de branches géométriques de $A$ est $1$ si et seulement si
$A$ est géométriquement unibranche.
\end{enumerate}
Supposons que $A$ possède un nombre fini d'idéaux premiers minimaux et soit $A'$ la
clôture intégrale de $A$ dans l'anneau total des fractions de $A_{red}$.
Alors
\begin{enumerate}
\item[(4)] le nombre de branches de $A$ est le nombre d'idéaux maximaux
$\mathfrak m'$ de $A'$,
\item[(5)] pour obtenir le nombre de branches géométriques de $A$, il faut compter
chaque idéal maximal $\mathfrak m'$ de $A'$ avec la multiplicité donnée par le
degré séparable de $\kappa(\mathfrak m')/\kappa$.
\end{enumerate}
\end{lemma}

\begin{proof}
Ce lemme découle immédiatement des définitions,
du lemme \ref{lemma-branches},
du lemme \ref{lemma-geometric-branches} et de
Corps, lemme \ref{fields-lemma-separable-degree}.
\end{proof}

\begin{lemma}
\label{lemma-invariance-number-branches-smooth}
Soit $A \to B$ un homomorphisme local d'anneaux locaux qui
est la localisation d'un homomorphisme lisse d'anneaux.
\begin{enumerate}
\item Le nombre de branches géométriques de $A$ est égal au nombre
de branches géométriques de $B$.
\item Si $A \to B$ induit une extension purement inséparable des corps résiduels,
alors le nombre de branches de $A$ est égal au nombre de branches de $B$.
\end{enumerate}
\end{lemma}

\begin{proof}
Nous utiliserons le fait que les homomorphismes lisses d'anneaux sont plats
(Algèbre, lemme \ref{algebra-lemma-smooth-syntomic}),
que les localisations sont plates (Algèbre, lemme
\ref{algebra-lemma-flat-localization}),
que les composées d'homomorphismes plats d'anneaux sont plates
(Algèbre, lemme \ref{algebra-lemma-composition-flat}),
que le changement de base d'un homomorphisme plat d'anneaux est plat
(Algèbre, lemme \ref{algebra-lemma-flat-base-change}),
que les homomorphismes locaux plats sont fidèlement plats
(Algèbre, lemme \ref{algebra-lemma-local-flat-ff}),
que l'hensélisation (stricte) est plate
(lemme \ref{lemma-dumb-properties-henselization}) et la
descente pour les homomorphismes plats d'anneaux
(Algèbre, lemme \ref{algebra-lemma-flat-going-down}).

\medskip\noindent
Démonstration de (2). Soient $A^h$, $B^h$ les hensélisés de $A$, $B$. Alors $B^h$
est l'hensélisé de $A^h \otimes_A B$ en l'unique
idéal maximal au-dessus de $\mathfrak m_B$, voir
Algèbre, lemme \ref{algebra-lemma-henselian-functorial-improve}.
Nous pouvons donc supposer, et supposons, que $A$ est hensélien.
Comme $A \to B \to B^h$ est plat, tout idéal premier minimal de $B^h$
est au-dessus d'un idéal premier minimal de $A$ et, comme $A \to B^h$ est fidèlement plat,
tout idéal premier minimal de $A$ est au-dessous d'un idéal premier minimal de $B^h$ ;
dans les deux cas, on utilise la descente pour les homomorphismes plats d'anneaux.
Il suffit donc de montrer que, pour un idéal premier minimal donné
$\mathfrak p \subset A$, il existe au plus un idéal premier minimal
de $B^h$ au-dessus de $\mathfrak p$.
Après avoir remplacé $A$ par $A/\mathfrak p$ et $B$ par $B/\mathfrak p B$,
nous pouvons supposer que $A$ est un domaine ; $A$ reste hensélien d'après
Algèbre, lemme \ref{algebra-lemma-quotient-henselization}.
D'après le lemme \ref{lemma-unibranch}, nous voyons que la clôture intégrale
$A'$ de $A$ dans son corps des fractions est un domaine local.
Bien sûr, $A'$ est un domaine normal. D'après
Algèbre, lemme \ref{algebra-lemma-normal-goes-up},
nous voyons que $A' \otimes_A B^h$ est un anneau normal
(le lemme le donne seulement pour $A' \otimes_A B$ ; pour passer à
$A' \otimes_A B^h$, on utilise que $B^h$ est une colimite de $B$-algèbres
étales ainsi que Algèbre, lemme \ref{algebra-lemma-colimit-normal-ring}).
D'après Algèbre, lemme \ref{algebra-lemma-local-tensor-with-integral},
nous voyons que $A' \otimes_A B^h$ est local (c'est ici que nous
utilisons l'hypothèse sur les corps résiduels de $A$ et de $B$).
Ainsi, $A' \otimes_A B^h$ est un anneau local normal, donc un domaine local.
Comme $B^h \subset A' \otimes_A B^h$ par platitude de $A \to B^h$,
nous concluons que $B^h$ est un domaine, comme désiré.

\medskip\noindent
Démonstration de (1). Soient $A^{sh}$, $B^{sh}$ des hensélisés stricts
de $A$, $B$. Alors $B^{sh}$ est un hensélisé strict de
$A^{sh} \otimes_A B$ en un idéal maximal au-dessus de $\mathfrak m_B$
et de $\mathfrak m_{A^{sh}}$, voir
Algèbre, lemme \ref{algebra-lemma-strictly-henselian-functorial-improve}.
Nous pouvons donc supposer, et supposons, que $A$ est strictement hensélien.
Comme $A \to B \to B^{sh}$ est plat, tout idéal premier minimal de $B^{sh}$
est au-dessus d'un idéal premier minimal de $A$ et, comme $A \to B^{sh}$ est fidèlement plat,
tout idéal premier minimal de $A$ est au-dessous d'un idéal premier minimal de $B^{sh}$ ;
dans les deux cas, on utilise la descente pour les homomorphismes plats d'anneaux.
Il suffit donc de montrer que, pour un idéal premier minimal donné
$\mathfrak p \subset A$, il existe au plus un idéal premier minimal
de $B^{sh}$ au-dessus de $\mathfrak p$.
Après avoir remplacé $A$ par $A/\mathfrak p$ et $B$ par $B/\mathfrak p B$,
nous pouvons supposer que $A$ est un domaine ; alors $A$ reste strictement hensélien d'après
Algèbre, lemme \ref{algebra-lemma-quotient-strict-henselization}.
D'après le lemme \ref{lemma-geometrically-unibranch}, nous voyons que la clôture intégrale
$A'$ de $A$ dans son corps des fractions est un domaine local dont le corps
résiduel est une extension purement inséparable du corps résiduel de $A$.
Bien sûr, $A'$ est un domaine normal. D'après
Algèbre, lemme \ref{algebra-lemma-normal-goes-up},
nous voyons que $A' \otimes_A B^{sh}$ est un anneau normal
(le lemme le donne seulement pour $A' \otimes_A B$ ; pour passer à
$A' \otimes_A B^{sh}$, on utilise que $B^{sh}$ est une colimite de $B$-algèbres
étales ainsi que Algèbre, lemme \ref{algebra-lemma-colimit-normal-ring}).
D'après Algèbre, lemme \ref{algebra-lemma-local-tensor-with-integral},
nous voyons que $A' \otimes_A B^{sh}$ est local (car $A \subset A'$
induit une extension purement inséparable des corps résiduels).
Ainsi, $A' \otimes_A B^{sh}$ est un anneau local normal, donc un domaine local.
Comme $B^{sh} \subset A' \otimes_A B^{sh}$ par platitude de $A \to B^{sh}$,
nous concluons que $B^{sh}$ est un domaine, comme désiré.
\end{proof}







\section{Divers résultats sur les branches}
\label{section-branches}

\noindent
Quelques résultats relatifs aux branches d'anneaux locaux telles qu'elles sont définies
dans la section \ref{section-unibranch}.

\begin{lemma}
\label{lemma-ideal-inverse-nonzero}
Soient $A$ et $B$ des domaines et soit $A \to B$ un homomorphisme d'anneaux.
Supposons que $A \to B$ possède en outre au moins l'une des propriétés suivantes :
\begin{enumerate}
\item c'est la localisation d'un homomorphisme étale d'anneaux,
\item il est plat et est la localisation d'un homomorphisme non ramifié d'anneaux,
\item il est plat et est la localisation d'un homomorphisme quasi-fini d'anneaux,
\item il est plat et est la localisation d'un homomorphisme entier d'anneaux,
\item il est plat et il n'existe aucune spécialisation non triviale entre des points
des fibres de $\Spec(B) \to \Spec(A)$,
\item $\Spec(B) \to \Spec(A)$ envoie le point générique sur le point
générique et il n'existe aucune spécialisation non triviale entre des points des fibres, ou
\item exactement un point de $\Spec(B)$ est envoyé sur le point
générique de $\Spec(A)$.
\end{enumerate}
Alors $A \cap J$ est non nul pour tout idéal non nul $J$ de $B$.
\end{lemma}

\begin{proof}
Démonstration dans le cas (7).
Soient $K$, resp.\ $L$, le corps des fractions de $A$, resp.\ de $B$.
D'après Algèbre, lemme
\ref{algebra-lemma-minimal-prime-image-minimal-prime}, nous voyons que
l'unique point de $\Spec(B)$ qui s'envoie sur le point
générique $(0) \in \Spec(A)$ est $(0) \in \Spec(B)$.
Nous concluons que $B \otimes_A K$ est un anneau possédant un unique idéal
premier dont le corps résiduel est $L$ (en fait, cet anneau est égal à $L$, mais
nous n'en avons pas besoin).
Choisissons $b \in J$ non nul. Alors $b$ s'envoie sur une unité de $L$.
Par conséquent, $b$ s'envoie sur une unité de $B \otimes_A K$
(Algèbre, lemme \ref{algebra-lemma-surjective-on-spec-units}).
Comme $B \otimes_A K = \colim_{f \in A \setminus \{0\}} B_f$,
nous voyons que $b$ s'envoie sur une unité de $B_f$ pour un certain $f \in A$ non nul.
Cela signifie que $b b' = f^n$ pour un certain $b' \in B$ et un certain $n \geq 1$.
Ainsi, $f^n \in A \cap J$, comme désiré.

\medskip\noindent
Dans la suite de la démonstration, nous montrons que chacune des autres hypothèses implique (7).
Sous les hypothèses (1) -- (5), l'homomorphisme d'anneaux $A \to B$ est plat et
donc $A \to B$ est injectif (car les homomorphismes locaux plats sont
fidèlement plats d'après Algèbre, lemme \ref{algebra-lemma-local-flat-ff}).
Par conséquent, le point générique de $\Spec(B)$
s'envoie sur le point générique de $\Spec(A)$. Or, s'il n'existe aucune
spécialisation non triviale
entre des points des fibres de $\Spec(B) \to \Spec(A)$, alors ce
point générique de $\Spec(B)$ doit évidemment être l'unique point s'envoyant
sur le point générique de $\Spec(A)$. Ainsi, (6) implique (7).
Enfin, pour conclure, montrons que dans les cas (1) -- (5), il n'existe aucune
spécialisation non triviale
entre les points des fibres de $\Spec(B) \to \Spec(A)$.
En effet, voir
Algèbre, lemme \ref{algebra-lemma-integral-no-inclusion} pour
le cas entier,
Algèbre, définition \ref{algebra-definition-quasi-finite}
pour le cas quasi-fini, et utiliser le fait que les homomorphismes
non ramifiés et étales d'anneaux sont quasi-finis
(Algèbre, lemmes \ref{algebra-lemma-unramified-quasi-finite} et
\ref{algebra-lemma-etale-quasi-finite}).
\end{proof}

\begin{lemma}
\label{lemma-local-unramified-extension-unibranch-domain-is-etale}
Soit $A \to B$ un homomorphisme d'anneaux. Soit $\mathfrak q \subset B$ un idéal
premier au-dessus de l'idéal premier $\mathfrak p \subset A$. Supposons que
\begin{enumerate}
\item $A$ est un domaine,
\item $A_\mathfrak p$ est géométriquement unibranche,
\item $A \to B$ est non ramifié en $\mathfrak q$, et
\item $A_\mathfrak p \to B_\mathfrak q$ est injectif.
\end{enumerate}
Alors il existe un $g \in B$, $g \not \in \mathfrak q$, tel que
$B_g$ soit étale sur $A$.
\end{lemma}

\begin{proof}
D'après Algèbre, proposition \ref{algebra-proposition-unramified-locally-standard},
après avoir remplacé $B$ par une localisation principale, nous pouvons trouver un
homomorphisme étale standard d'anneaux $A \to B'$ et une surjection $B' \to B$.
Notons $\mathfrak q' \subset B'$ l'image réciproque de $\mathfrak q$.
Nous allons montrer que $B' \to B$ est injectif,
quitte à remplacer $B'$ par une localisation principale.

\medskip\noindent
Dans ce paragraphe, nous nous ramenons au cas où $B'$ est un domaine.
Comme $A$ est un domaine, l'anneau $B'$ est réduit, voir
le lemme \ref{lemma-reduced-goes-up}.
Soit $K$ le corps des fractions de $A$. Alors $B' \otimes_A K$
est étale sur un corps, et est donc un produit fini de corps, voir
Algèbre, lemme \ref{algebra-lemma-etale-over-field}.
Comme $A \to B'$ est étale (donc plat), les idéaux premiers minimaux de $B'$
sont au-dessus de $(0) \subset A$ (par descente pour les homomorphismes plats d'anneaux).
Nous concluons que $B'$ possède un nombre fini d'idéaux premiers minimaux, disons
$\mathfrak r_1, \ldots, \mathfrak r_r \subset B'$.
Comme $A_\mathfrak p$ est géométriquement unibranche et $A \to B'$ étale,
l'anneau $B'_{\mathfrak q'}$ est un domaine, voir
les lemmes \ref{lemma-invariance-number-branches-smooth} et
\ref{lemma-number-of-branches-1}.
Par conséquent, $\mathfrak q' \supset \mathfrak r_i$ pour exactement un $i = i_0$.
Choisissons $g' \in B'$, $g' \not \in \mathfrak r_{i_0}$, mais
$g' \in \mathfrak r_i$ pour $i \not = i_0$, voir
Algèbre, lemme \ref{algebra-lemma-silly}. Après avoir remplacé
$B'$ et $B$ par $B'_{g'}$ et $B_{g'}$, nous obtenons que $B'$
est un domaine.

\medskip\noindent
Supposons que $B'$ soit un domaine ; en particulier, $B' \subset B'_{\mathfrak q'}$.
Si $B' \to B$ n'est pas injectif,
alors $J = \Ker(B'_{\mathfrak q'} \to B_\mathfrak q)$ est non nul.
En appliquant le lemme \ref{lemma-ideal-inverse-nonzero} à
$A_\mathfrak p \to B'_{\mathfrak q'}$,
nous trouvons un élément non nul $a \in A_\mathfrak p$
qui s'envoie sur zéro dans $B_\mathfrak q$, contredisant
l'hypothèse (4). Ceci achève la démonstration.
\end{proof}

\begin{lemma}
\label{lemma-unramified-extension-unibranch-domain-is-etale}
\begin{reference}
Généralisation de \cite[Exposé I, théorème 9.5, partie (ii)]{SGA1}
\end{reference}
Soit $(A, \mathfrak m)$ un domaine local géométriquement unibranche. Soit
$A \to B$ un homomorphisme local injectif d'anneaux locaux,
essentiellement de type fini. Si $\mathfrak m B$
est l'idéal maximal de $B$ et si l'extension induite des
corps résiduels est séparable, alors $A \to B$ est la localisation
d'un homomorphisme étale d'anneaux.
\end{lemma}

\begin{proof}
Nous pouvons écrire $B = C_\mathfrak q$, où $A \to C$ est un homomorphisme d'anneaux
de type fini et $\mathfrak q \subset C$ est un idéal premier
au-dessus de $\mathfrak m$. D'après
Algèbre, lemme \ref{algebra-lemma-characterize-unramified},
l'homomorphisme d'anneaux $A \to C$ est non ramifié en $\mathfrak q$.
D'après Algèbre, proposition \ref{algebra-proposition-unramified-locally-standard},
après avoir remplacé $C$ par une localisation principale,
nous pouvons trouver un homomorphisme étale standard d'anneaux $A \to C'$
et une surjection $C' \to C$. Notons $\mathfrak q' \subset C'$
l'image réciproque de $\mathfrak q$ et posons $B' = C'_{\mathfrak q'}$.
Alors $B' \to B$ est surjectif.
Il suffit de montrer que $B' \to B$ est aussi injectif.

\medskip\noindent
Comme $A$ est un domaine, les anneaux $C'$ et $B'$ sont réduits, voir
le lemme \ref{lemma-reduced-goes-up}.
Comme $A$ est géométriquement unibranche, l'anneau $B'$ est un domaine, voir
les lemmes \ref{lemma-invariance-number-branches-smooth} et
\ref{lemma-number-of-branches-1}.
Si $B' \to B$ n'est pas injectif, alors
$A \cap \Ker(B' \to B)$ est non nul d'après le lemme \ref{lemma-ideal-inverse-nonzero},
ce qui contredit l'hypothèse que $A \to B$ est injectif.
\end{proof}

\begin{lemma}
\label{lemma-minimal-primes-tensor-strictly-henselian}
Soit $k$ un corps algébriquement clos. Soient $A$, $B$ des $k$-algèbres locales
strictement henséliennes de corps résiduel égal à $k$.
Soit $C$ l'hensélisé strict de $A \otimes_k B$ en l'idéal
maximal $\mathfrak m_A \otimes_k B + A \otimes_k \mathfrak m_B$.
Alors les idéaux premiers minimaux de $C$ sont en correspondance $1$-à-$1$ avec les couples d'idéaux
premiers minimaux de $A$ et de $B$.
\end{lemma}

\begin{proof}
Remarquons d'abord qu'un idéal premier minimal $\mathfrak r$ de $C$ s'envoie sur un idéal
premier minimal $\mathfrak p$ de $A$ et sur un idéal premier minimal $\mathfrak q$ de $B$,
car les homomorphismes d'anneaux $A \to C$ et $B \to C$ sont plats (par descente pour les
homomorphismes plats d'anneaux, Algèbre, lemme \ref{algebra-lemma-flat-going-down}).
Il suffit donc de montrer que l'hensélisé strict de
$(A/\mathfrak p \otimes_k B/\mathfrak q)_{
\mathfrak m_A \otimes_k B + A \otimes_k \mathfrak m_B}$
possède un unique idéal premier minimal. D'après
Algèbre, lemme \ref{algebra-lemma-quotient-strict-henselization},
les anneaux $A/\mathfrak p$, $B/\mathfrak q$ sont strictement henséliens.
Nous pouvons donc supposer que $A$ et $B$ sont des domaines locaux
strictement henséliens, et notre but est de montrer que $C$ possède un unique idéal premier minimal.
D'après le lemme \ref{lemma-geometrically-unibranch}, la
clôture intégrale $A'$ de $A$ dans son corps des fractions
est un domaine local normal de corps résiduel $k$. Il en est de même de la
clôture intégrale $B'$ de $B$ dans son corps des fractions. D'après
Algèbre, lemme \ref{algebra-lemma-geometrically-normal-tensor-normal},
nous voyons que $A' \otimes_k B'$ est un anneau normal. Sa localisation
$$
R = (A' \otimes_k B')_{
\mathfrak m_{A'} \otimes_k B' + A' \otimes_k \mathfrak m_{B'}}
$$
est donc un domaine local normal. Remarquons que $A \otimes_k B \to A' \otimes_k B'$
est entier (la montée vaut donc --
Algèbre, lemme \ref{algebra-lemma-integral-going-up})
et que $\mathfrak m_{A'} \otimes_k B' + A' \otimes_k \mathfrak m_{B'}$
est l'unique idéal maximal de $A' \otimes_k B'$
au-dessus de $\mathfrak m_A \otimes_k B + A \otimes_k \mathfrak m_B$.
Nous voyons donc que
$$
R = (A' \otimes_k B')_{
\mathfrak m_A \otimes_k B + A \otimes_k \mathfrak m_B}
$$
d'après
Algèbre, lemme \ref{algebra-lemma-unique-prime-over-localize-below}.
Il s'ensuit que
$$
(A \otimes_k B)_{
\mathfrak m_A \otimes_k B + A \otimes_k \mathfrak m_B}
\longrightarrow
R
$$
est entier. Nous concluons que $R$ est la clôture intégrale de
$(A \otimes_k B)_{
\mathfrak m_A \otimes_k B + A \otimes_k \mathfrak m_B}$
dans son corps des fractions, et, par le
lemme \ref{lemma-geometrically-unibranch},
nous concluons une nouvelle fois que $C$ possède un unique idéal premier minimal.
\end{proof}






\section{Branches du complété}
\label{section-branches-completion}

\noindent
Soit $(A, \mathfrak m)$ un anneau local noethérien. Considérons les homomorphismes
$A \to A^h \to A^\wedge$. En général, l'homomorphisme $A^h \to A^\wedge$
n'induit pas nécessairement une bijection sur les idéaux premiers minimaux, voir
Exemples, section \ref{examples-section-bad}. Autrement dit, le
nombre de branches de $A$ (tel qu'il est défini dans la
définition \ref{definition-number-of-branches})
peut être différent du nombre de
branches de $A^\wedge$. Toutefois, sous certaines conditions, le nombre de
branches est le même, par exemple si la dimension de $A$ est $1$.

\begin{lemma}
\label{lemma-nr-branches-completion}
Soit $(A, \mathfrak m)$ un anneau local noethérien.
\begin{enumerate}
\item L'homomorphisme $A^h \to A^\wedge$ définit une application surjective des idéaux
premiers minimaux de $A^\wedge$ vers les idéaux premiers minimaux de $A^h$.
\item Le nombre de branches de $A$ est au plus le nombre de branches
de $A^\wedge$.
\item Le nombre de branches géométriques de $A$ est au plus le nombre
de branches géométriques de $A^\wedge$.
\end{enumerate}
\end{lemma}

\begin{proof}
D'après le lemme \ref{lemma-henselization-noetherian}, l'homomorphisme $A^h \to A^\wedge$
est plat et injectif. En combinant la descente
(Algèbre, lemme \ref{algebra-lemma-flat-going-down}) et
Algèbre, lemme \ref{algebra-lemma-injective-minimal-primes-in-image},
nous voyons que l'assertion (1) est vraie. L'assertion (2) en découle, ainsi que de la
définition \ref{definition-number-of-branches} et
du fait que $A^\wedge$ est hensélien
(Algèbre, lemme \ref{algebra-lemma-complete-henselian}).
D'après le lemme \ref{lemma-henselization-noetherian},
nous avons $(A^\wedge)^{sh} = A^{sh} \otimes_{A^h} A^\wedge$.
Nous pouvons donc répéter les arguments ci-dessus en utilisant l'homomorphisme plat injectif
$A^{sh} \to (A^\wedge)^{sh}$ pour démontrer (3).
\end{proof}

\begin{lemma}
\label{lemma-equal-nr-branches-completion}
Soit $(A, \mathfrak m)$ un anneau local noethérien. Le nombre de
branches de $A$ est le même que le nombre de branches de $A^\wedge$
si et seulement si $\sqrt{\mathfrak qA^\wedge}$ est premier pour tout
idéal premier minimal $\mathfrak q \subset A^h$ de l'hensélisé.
\end{lemma}

\begin{proof}
Ceci découle du lemme \ref{lemma-nr-branches-completion}
et du fait que $A$ et $A^\wedge$ ne possèdent qu'un nombre fini de branches,
d'après
Algèbre, lemme \ref{algebra-lemma-Noetherian-irreducible-components},
et du fait que $A^h$ et $A^\wedge$ sont noethériens
(lemme \ref{lemma-henselization-noetherian}).
\end{proof}

\noindent
Un lemme élémentaire de recollement.

\begin{lemma}
\label{lemma-glueing-sum-components-open}
Soit $A$ un anneau et soit $I$ un idéal de type fini.
Soit $A \to C$ un homomorphisme d'anneaux tel que, pour tout $f \in I$,
l'homomorphisme d'anneaux $A_f \to C_f$ soit la localisation en un idempotent.
Alors il existe une surjection $A \to C'$ telle que
$A_f \to (C \times C')_f$ soit un isomorphisme pour tout $f \in I$.
\end{lemma}

\begin{proof}
Choisissons des générateurs $f_1, \ldots, f_r$ de $I$. Écrivons
$$
C_{f_i} = (A_{f_i})_{e_i}
$$
pour un certain idempotent $e_i \in A_{f_i}$. Écrivons $e_i = a_i/f_i^n$
pour un certain $a_i \in A$ et un certain $n \geq 0$ ; nous pouvons prendre le même $n$
pour tout $i = 1, \ldots, r$. Après avoir remplacé $a_i$ par $f_i^ma_i$
et $n$ par $n + m$ pour un $m \gg 0$ convenable, nous pouvons supposer que
$a_i^2 = f_i^n a_i$ pour tout $i$. Comme $e_i$ s'envoie sur $1$ dans
$C_{f_if_j} = (A_{f_if_j})_{e_j} = A_{f_if_ja_j}$, nous voyons que
$$
(f_if_ja_j)^N(f_j^n a_i  - f_i^na_j) = 0
$$
pour un certain $N$ (nous pouvons prendre le même $N$ pour tous les couples $i, j$). En utilisant
$a_j^2 = f_j^na_j$, ceci donne
$$
f_i^{N + n} f_j^{N + nN} a_j = f_i^N f_j^{N + n} a_ia_j^N
$$
Après avoir augmenté $n$ à $n + N + nN$ et remplacé $a_i$ par $f_i^{N + nN}a_i$,
nous voyons que $f_i^n a_j$ appartient à l'idéal engendré par $a_i$ pour tous les couples $i, j$.
Soit $C' = A/(a_1, \ldots, a_r)$. Alors
$$
C'_{f_i} = A_{f_i}/(a_i) = A_{f_i}/(e_i)
$$
car $a_j$ appartient à l'idéal engendré par $a_i$ après inversion de $f_i$.
Puisque, pour un idempotent $e$ d'un anneau $B$, nous avons $B = B_e \times B/(e)$,
nous voyons que la conclusion du lemme vaut lorsque $f$ est égal à l'un
des $f_1, \ldots, f_r$. En utilisant le recollement des fonctions, sous la forme de
Algèbre, lemme \ref{algebra-lemma-cover},
nous concluons que le résultat vaut pour tout $f \in I$.
En effet, pour $f \in I$, les éléments $f_1, \ldots, f_r$
engendrent l'idéal unité dans $A_f$, de sorte que $A_f \to (C \times C')_f$
est un isomorphisme si et seulement si c'est le cas après localisation
en $f_1, \ldots, f_r$.
\end{proof}

\noindent
Le lemme \ref{lemma-quotient-by-idempotent} peut être utilisé pour construire
des extensions de type fini à partir d'extensions de type fini données du
complété formel. Nous
généraliserons ce lemme dans Algébrisation des espaces formels, lemme
\ref{restricted-lemma-approximate-by-etale-over-complement}.

\begin{lemma}
\label{lemma-quotient-by-idempotent}
Soit $A$ un anneau noethérien et $I$ un idéal. Soit $B$
une $A$-algèbre de type fini. Soit $B^\wedge \to C$ un homomorphisme surjectif
d'anneaux de noyau $J$, où $B^\wedge$ est le complété $I$-adique.
Si $J/J^2$ est annulé par $I^c$ pour un certain $c \geq 0$, alors $C$ est
isomorphe au complété d'une $A$-algèbre de type fini.
\end{lemma}

\begin{proof}
Soit $f \in I$. Comme $B^\wedge$ est noethérien (Algèbre, lemme
\ref{algebra-lemma-completion-Noetherian-Noetherian}),
nous voyons que $J$ est un idéal de type fini.
Nous concluons donc de
Algèbre, lemme \ref{algebra-lemma-ideal-is-squared-union-connected},
que
$$
C_f = ((B^\wedge)_f)_e
$$
pour un certain idempotent $e \in (B^\wedge)_f$. D'après le
lemme \ref{lemma-glueing-sum-components-open},
nous pouvons trouver une surjection $B^\wedge \to C'$
telle que $B^\wedge \to C \times C'$ devienne un
isomorphisme après inversion de tout $f \in I$.
Observons que $C \times C'$ est une $B^\wedge$-algèbre finie.

\medskip\noindent
Choisissons des générateurs $f_1, \ldots, f_r \in I$. Notons
$\alpha_i : (C \times C')_{f_i} \to B_{f_i} \otimes_B B^\wedge$
l'inverse de l'isomorphisme de $(B^\wedge)_{f_i}$-algèbres
obtenu ci-dessus. Notons $\alpha_{ij} : (B_{f_i})_{f_j} \to (B_{f_j})_{f_i}$
l'isomorphisme évident de $B$-algèbres.
Considérons l'objet
$$
(C \times C', B_{f_i}, \alpha_i, \alpha_{ij})
$$
de la catégorie $\text{Glue}(B \to B^\wedge, f_1, \ldots, f_r)$ introduite
dans la remarque \ref{remark-glueing-data}. Nous omettons la vérification des
conditions (1)(a) et (1)(b).
Comme $B \to B^\wedge$ est un homomorphisme plat
(Algèbre, lemme \ref{algebra-lemma-completion-flat}) induisant un isomorphisme
$B/IB \to B^\wedge/IB^\wedge$, nous pouvons appliquer la
proposition \ref{proposition-equivalence}
et la remarque \ref{remark-formal-glueing-algebras}.
Nous concluons que $C \times C'$ est isomorphe à $D \otimes_B B^\wedge$
pour une certaine $B$-algèbre finie $D$.
Alors $D/ID \cong C/IC \times C'/IC'$. Soit $\overline{e} \in D/ID$
l'idempotent correspondant au facteur $C/IC$.
D'après le lemme \ref{lemma-lift-idempotent-upstairs}, il existe un
homomorphisme étale d'anneaux $B \to B'$ qui induit un isomorphisme
$B/IB \to B'/IB'$ tel que $D' = D \otimes_B B'$ contienne un
idempotent $e$ relevant $\overline{e}$. Comme $C \times C'$
est $I$-adiquement complet, la paire $(C \times C', IC \times IC')$
est hensélienne (lemme \ref{lemma-complete-henselian}).
Nous pouvons donc factoriser l'homomorphisme $B \to C \times C'$ par $B'$ (voir
le lemme \ref{lemma-characterize-henselian-pair}).
Ce faisant, nous pouvons remplacer $B$ par $B'$ et $D$ par $D'$ (car les
complétés $I$-adiques n'en sont pas modifiés). Alors
nous trouvons que $D = D_e \times D_{1 - e} = D/(1 - e) \times D/(e)$
est un produit de $A$-algèbres de type fini, que le complété du
premier facteur est $C$ et que le complété du second facteur est $C'$.
\end{proof}

\begin{lemma}
\label{lemma-one-dimensional-formal-branch}
Soit $(A, \mathfrak m)$ un anneau local noethérien d'hensélisé $A^h$.
Soit $\mathfrak q \subset A^\wedge$ un idéal premier minimal tel que
$\dim(A^\wedge/\mathfrak q) = 1$. Alors il existe un idéal premier
minimal $\mathfrak q^h$ de $A^h$ tel que
$\mathfrak q = \sqrt{\mathfrak q^hA^\wedge}$.
\end{lemma}

\begin{proof}
Comme les complétés de $A$ et de $A^h$ sont les mêmes, nous pouvons supposer
que $A$ est hensélien (lemme \ref{lemma-henselization-noetherian}).
Nous allons appliquer le lemme \ref{lemma-quotient-by-idempotent}
à $A^\wedge \to A^\wedge/J$, où
$J = \Ker(A^\wedge \to (A^\wedge)_{\mathfrak q})$.
Comme $\dim((A^\wedge)_\mathfrak q) = 0$, nous voyons que
$\mathfrak q^n \subset J$ pour un certain $n$. Ainsi, $J/J^2$ est
annulé par $\mathfrak q^n$. D'autre part, $(J/J^2)_\mathfrak q = 0$
car $J_\mathfrak q = 0$. Par conséquent, $\mathfrak m$ est le seul idéal premier
associé de $J/J^2$, et nous trouvons qu'une puissance
de $\mathfrak m$ annule $J/J^2$. Le lemme s'applique donc,
et nous trouvons que $A^\wedge/J = C^\wedge$ pour une certaine
$A$-algèbre de type fini $C$.

\medskip\noindent
Alors $C/\mathfrak m C = A/\mathfrak m$, car $A^\wedge/J$ possède la même
propriété. Ainsi, $\mathfrak m_C = \mathfrak m C$ est un idéal maximal
et $A \to C$ est non ramifié en $\mathfrak m_C$
(Algèbre, lemme \ref{algebra-lemma-characterize-unramified}).
Après avoir remplacé $C$ par une localisation principale, nous pouvons
supposer que $C$ est un quotient d'une $A$-algèbre étale $B$, voir
Algèbre, proposition \ref{algebra-proposition-unramified-locally-standard}.
Cependant, comme l'extension de corps résiduels de $A \to C_{\mathfrak m_C}$
est triviale et que $A$ est hensélien, nous concluons que $B = A$,
à nouveau après une localisation.
Ainsi, $C = A/I$ pour un certain idéal $I \subset A$, et il s'ensuit que
$J = IA^\wedge$ (car la complétion est exacte dans notre situation d'après
Algèbre, lemme \ref{algebra-lemma-completion-flat}) et que $I = J \cap A$
(par platitude de $A \to A^\wedge$). Comme
$\mathfrak q^n \subset J \subset \mathfrak q$, nous voyons que
$\mathfrak p = \mathfrak q \cap A$ satisfait
$\mathfrak p^n \subset I \subset \mathfrak p$.
Alors $\sqrt{\mathfrak p A^\wedge} = \mathfrak q$, ce qui achève la démonstration.
\end{proof}

\begin{lemma}
\label{lemma-completion-disconnected}
Soit $(A, \mathfrak m)$ un anneau local noethérien. Le spectre épointé
de $A^\wedge$ est non connexe si et seulement si le spectre épointé de $A^h$
est non connexe.
\end{lemma}

\begin{proof}
Comme les complétés de $A$ et de $A^h$ sont les mêmes, nous pouvons supposer
que $A$ est hensélien (lemme \ref{lemma-henselization-noetherian}).

\medskip\noindent
Comme $A \to A^\wedge$ est fidèlement plat (voir la référence qui précède),
l'application du spectre épointé de $A^\wedge$ vers le spectre
épointé de $A$ est surjective
(voir Algèbre, lemme \ref{algebra-lemma-ff-rings}).
Par conséquent, si le spectre épointé de $A$ est non connexe, alors
il en est de même de celui de $A^\wedge$.

\medskip\noindent
Supposons que le spectre épointé de $A^\wedge$ soit non connexe.
Cela signifie que
$$
\Spec(A^\wedge) \setminus \{\mathfrak m^\wedge\} = Z \amalg Z'
$$
où $Z$ et $Z'$ sont fermés. Soient
$\overline{Z}, \overline{Z}' \subset \Spec(A^\wedge)$ leurs adhérences.
Écrivons $\overline{Z} = V(J)$, $\overline{Z}' = V(J')$ pour certains idéaux
$J, J' \subset A^\wedge$. Alors $V(J + J') = \{\mathfrak m^\wedge\}$
et $V(JJ') = \Spec(A^\wedge)$. La première égalité signifie que
$\mathfrak m^\wedge = \sqrt{J + J'}$, ce qui implique que
$(\mathfrak m^\wedge)^e \subset J + J'$ pour un certain $e \geq 1$.
La seconde égalité implique que tout élément
de $JJ'$ est nilpotent, et donc que $(JJ')^n = 0$ pour un certain $n \geq 1$.
En combinant ces faits, nous obtenons que $J^n/J^{2n}$ est annulé par
$J^n$ et $(J')^n$, et donc par $(\mathfrak m^\wedge)^{2en}$.
Nous pouvons ainsi appliquer le lemme \ref{lemma-quotient-by-idempotent}
pour voir qu'il existe une $A$-algèbre de type fini $C$ et un
isomorphisme $A^\wedge/J^n = C^\wedge$.

\medskip\noindent
La suite de la démonstration est exactement la même que la seconde partie
de la démonstration du lemme \ref{lemma-one-dimensional-formal-branch} ;
bien sûr, ce lemme est un cas particulier du présent lemme !
Nous avons $C/\mathfrak m C = A/\mathfrak m$, car $A^\wedge/J^n$ possède la même
propriété. Ainsi, $\mathfrak m_C = \mathfrak m C$ est un idéal maximal
et $A \to C$ est non ramifié en $\mathfrak m_C$
(Algèbre, lemme \ref{algebra-lemma-characterize-unramified}).
Après avoir remplacé $C$ par une localisation principale, nous pouvons
supposer que $C$ est un quotient d'une $A$-algèbre étale $B$, voir
Algèbre, proposition \ref{algebra-proposition-unramified-locally-standard}.
Cependant, comme l'extension de corps résiduels de $A \to C_{\mathfrak m_C}$
est triviale et que $A$ est hensélien, nous concluons que $B = A$,
à nouveau après une localisation.
Ainsi, $C = A/I$ pour un certain idéal $I \subset A$, et il s'ensuit que
$J^n = IA^\wedge$ (car la complétion est exacte dans notre situation d'après
Algèbre, lemme \ref{algebra-lemma-completion-flat}) et que $I = J^n \cap A$
(par platitude de $A \to A^\wedge$).
Par symétrie, $I' = (J')^n \cap A$ satisfait $(J')^n = I'A^\wedge$.
Alors $\mathfrak m^e \subset I + I'$ et $II' = 0$,
et nous concluons que $V(I)$ et $V(I')$ sont des sous-schémas fermés
qui donnent la décomposition en union disjointe recherchée du
spectre épointé de $A$.
\end{proof}

\begin{lemma}
\label{lemma-one-dimensional-number-of-branches}
Soit $(A, \mathfrak m)$ un anneau local noethérien de dimension $1$.
Alors le nombre de branches (géométriques) de $A$ et de $A^\wedge$ est le même.
\end{lemma}

\begin{proof}
Pour le voir pour le nombre de branches, combinons les
lemmes \ref{lemma-nr-branches-completion},
\ref{lemma-equal-nr-branches-completion} et
\ref{lemma-one-dimensional-formal-branch}
et utilisons le fait que la dimension de $A^\wedge$ est un, voir
le lemme \ref{lemma-completion-dimension}.
Pour voir que cela est vrai pour le nombre de branches géométriques,
nous utilisons le résultat pour les branches, le fait que la dimension
ne change pas par hensélisation stricte
(lemme \ref{lemma-henselization-dimension}) et le fait que
$(A^{sh})^\wedge = ((A^\wedge)^{sh})^\wedge$
d'après le lemme \ref{lemma-henselization-noetherian}.
\end{proof}

\begin{lemma}
\label{lemma-geometrically-normal-formal-fibres-number-of-branches}
\begin{reference}
\cite[Theorem 2.3]{Beddani}
\end{reference}
Soit $(A, \mathfrak m)$ un anneau local noethérien. Si les fibres
formelles de $A$ sont géométriquement normales (par exemple si $A$ est
excellent ou quasi-excellent), alors $A$ est de Nagata
et le nombre de branches (géométriques) de $A$ et de $A^\wedge$ est le même.
\end{lemma}

\begin{proof}
Comme un anneau normal est réduit, nous voyons que $A$ est de Nagata d'après le
lemme \ref{lemma-Nagata-local-ring}. Dans la suite de la démonstration,
nous utiliserons le lemme \ref{lemma-formal-fibres-normal},
la proposition \ref{proposition-finite-type-over-P-ring} et le
lemme \ref{lemma-check-P-ring-maximal-ideals}. Ceux-ci nous disent
que $A$ est un anneau P, où $P(k \to R) = $``$R$ est géométriquement
normal sur $k$'', et qu'il en est de même de toute $A$-algèbre (essentiellement) de type fini.

\medskip\noindent
Soit $\mathfrak q \subset A$ un idéal premier minimal. Alors
$A^\wedge/\mathfrak q A^\wedge = (A/\mathfrak q)^\wedge$
et $A^h/\mathfrak qA^h = (A/\mathfrak q)^h$
(Algèbre, lemme \ref{algebra-lemma-quotient-henselization}).
Ainsi, le nombre de branches de $A$ est la somme des
nombres de branches des anneaux $A/\mathfrak q$, et
de même pour $A^\wedge$. Nous nous ramenons ainsi
au cas où $A$ est un domaine.

\medskip\noindent
Supposons que $A$ soit un domaine. Soit $A'$ la clôture intégrale de $A$
dans le corps des fractions $K$ de $A$. Comme $A$ est de Nagata, nous voyons que
$A \to A'$ est fini. Rappelons que le nombre de branches de
$A$ est le nombre d'idéaux maximaux $\mathfrak m'$ de $A'$
(lemme \ref{lemma-branches}). Rappelons aussi que
$$
(A')^\wedge = A' \otimes_A A^\wedge =
\prod\nolimits_{\mathfrak m' \subset A'} (A'_{\mathfrak m'})^\wedge
$$
d'après Algèbre, lemme \ref{algebra-lemma-completion-finite-extension}.
Comme $A'_{\mathfrak m'}$ est un anneau local dont les fibres
formelles sont géométriquement normales, nous voyons que
$(A'_{\mathfrak m'})^\wedge$ est normal
(lemme \ref{lemma-completion-normal-local-ring}).
Par conséquent, les idéaux premiers minimaux de $A' \otimes_A A^\wedge$
sont en correspondance $1$-à-$1$ avec les facteurs de la
décomposition ci-dessus. Par platitude de $A \to A^\wedge$, nous avons
$$
A^\wedge \subset A' \otimes_A A^\wedge \subset K \otimes_A A^\wedge
$$
Comme les anneaux de gauche et de droite ont le même ensemble d'idéaux
premiers minimaux, il en est de même de l'anneau du milieu (petit
détail omis), ce qui achève la démonstration.

\medskip\noindent
Pour voir que cela est vrai pour le nombre de branches géométriques,
nous utilisons le résultat pour les branches, le fait que les fibres
formelles de $A^{sh}$ sont géométriquement normales
(lemmes \ref{lemma-formal-fibres-normal} et
\ref{lemma-henselization-P-ring})
et le fait que $(A^{sh})^\wedge = ((A^\wedge)^{sh})^\wedge$
d'après le lemme \ref{lemma-henselization-noetherian}.
\end{proof}






\section{Anneaux formellement caténaires}
\label{section-formally-catenary}

\noindent
Dans cette section, nous démontrons un théorème de Ratliff
\cite{Ratliff} selon lequel un anneau local
noethérien est universellement caténaire si et seulement s'il est formellement caténaire.

\begin{definition}
\label{definition-formally-catenary}
Un anneau local noethérien $A$ est {\it formellement caténaire}
si, pour tout idéal premier minimal $\mathfrak p \subset A$, le spectre de
$A^\wedge/\mathfrak p A^\wedge$ est équidimensionnel.
\end{definition}

\noindent
Soit $A$ un anneau local noethérien formellement caténaire.
D'après le résultat de Ratliff (proposition \ref{proposition-ratliff}),
nous voyons que tout quotient de $A$ est également formellement caténaire
(car la classe des anneaux universellement caténaires est stable par quotients).
Nous concluons que le spectre de $A^\wedge/\mathfrak p A^\wedge$
est équidimensionnel pour tout idéal premier $\mathfrak p$ de $A$.

\begin{lemma}
\label{lemma-not-formally-catenary}
Soit $(A, \mathfrak m)$ un anneau local noethérien qui n'est pas
formellement caténaire. Alors $A$ n'est pas universellement caténaire.
\end{lemma}

\begin{proof}
Par hypothèse, il existe un idéal premier minimal $\mathfrak p \subset A$ tel que
le spectre de $A^\wedge /\mathfrak p A^\wedge$ ne soit pas équidimensionnel.
Après avoir remplacé $A$ par $A/\mathfrak p$, nous pouvons supposer que $A$
est un domaine et que le spectre de $A^\wedge$ n'est pas équidimensionnel.
Soit $\mathfrak q$ un idéal premier minimal de
$A^\wedge$ tel que $d = \dim(A^\wedge/\mathfrak q)$
soit minimal, et donc $0 < d < \dim(A)$. Nous démontrons le lemme par récurrence
sur $d$.

\medskip\noindent
Le cas $d = 1$. Dans ce cas, $\dim(A^\wedge_\mathfrak q) = 0$.
Ainsi, $A^\wedge_\mathfrak q$ est local artinien, et nous voyons que,
pour un certain $n > 0$, l'idéal $J = \mathfrak q^n$ s'envoie sur zéro dans
$A^\wedge_\mathfrak q$. Il s'ensuit que $\mathfrak m$ est le
seul idéal premier associé de $J/J^2$, si bien que $\mathfrak m^m$ annule
$J/J^2$ pour un certain $m > 0$. Nous pouvons donc utiliser le
lemme \ref{lemma-quotient-by-idempotent}
pour trouver $A \to B$ de type fini tel que $B^\wedge \cong A^\wedge/J$.
Il s'ensuit que $\mathfrak m_B = \sqrt{\mathfrak mB}$ est un idéal
maximal de même corps résiduel que $\mathfrak m$, et que $B^\wedge$
est le complété $\mathfrak m_B$-adique
(Algèbre, lemme \ref{algebra-lemma-finite-after-completion}).
Alors
$$
\dim(B_{\mathfrak m_B}) = \dim(B^\wedge) = 1 = d.
$$
Comme nous avons la factorisation $A \to B \to A^\wedge/J$, l'image réciproque
de $\mathfrak q/J$ est un idéal premier $\mathfrak q' \subset \mathfrak m_B$ au-dessus
de $(0)$ dans $A$. Ainsi, si $A$ était universellement caténaire, la formule
de dimension (Algèbre, lemme \ref{algebra-lemma-dimension-formula}) donnerait
\begin{align*}
\dim(B_{\mathfrak m_B})
& \geq
\dim((B/\mathfrak q')_{\mathfrak m_B}) \\
& =
\dim(A) + \text{trdeg}_A(B/\mathfrak q') -
\text{trdeg}_{\kappa(\mathfrak m)}(\kappa(\mathfrak m_B)) \\
& =
\dim(A) + \text{trdeg}_A(B/\mathfrak q')
\end{align*}
Cette contradiction achève l'argument dans le cas $d = 1$.

\medskip\noindent
Supposons $d > 1$. Soit $Z \subset \Spec(A^\wedge)$ la réunion des
composantes irréductibles distinctes de $V(\mathfrak q)$.
Soient $\mathfrak r_1, \ldots, \mathfrak r_m \subset A^\wedge$
les idéaux premiers correspondant aux composantes irréductibles de
$V(\mathfrak q) \cap Z$ de dimension $> 0$.
Choisissons $f \in \mathfrak m$, $f \not \in A \cap \mathfrak r_j$,
en utilisant l'évitement des idéaux premiers (Algèbre, lemme \ref{algebra-lemma-silly}).
Alors $\dim(A/fA) = \dim(A) - 1$ et il existe une composante irréductible
de $V(\mathfrak q, f)$ de dimension $d - 1$.
Ainsi, $A/fA$ n'est pas formellement caténaire et l'invariant $d$ a
diminué. Par récurrence, $A/fA$ n'est pas universellement caténaire, donc
$A$ n'est pas universellement caténaire.
\end{proof}

\begin{lemma}
\label{lemma-flat-under-catenary-equidimensional}
Soit $A \to B$ un homomorphisme local plat d'anneaux locaux noethériens.
Supposons que $B$ soit caténaire et que $\Spec(B)$ soit équidimensionnel. Alors
\begin{enumerate}
\item $\Spec(B/\mathfrak p B)$ est équidimensionnel pour tout
$\mathfrak p \subset A$, et
\item $A$ est caténaire et $\Spec(A)$ est équidimensionnel.
\end{enumerate}
\end{lemma}

\begin{proof}
Soit $\mathfrak p \subset A$ un idéal premier. Soit $\mathfrak q \subset B$
un idéal premier minimal au-dessus de $\mathfrak pB$. Alors $\mathfrak q \cap A = \mathfrak p$
par descente pour $A \to B$
(Algèbre, lemme \ref{algebra-lemma-flat-going-down}).
Ainsi, $A_\mathfrak p \to B_\mathfrak q$ est un homomorphisme local plat d'anneaux
dont la fibre spéciale est de dimension $0$, et donc
$$
\dim(A_\mathfrak p) = \dim(B_\mathfrak q) = \dim(B) - \dim(B/\mathfrak q)
$$
(Algèbre, lemme \ref{algebra-lemma-dimension-base-fibre-equals-total}).
La seconde égalité vaut parce que $\Spec(B)$ est équidimensionnel et $B$ est caténaire.
Ainsi, $\dim(B/\mathfrak q)$ est indépendant du choix de $\mathfrak q$,
et nous concluons que $\Spec(B/\mathfrak p B)$ est équidimensionnel de
dimension $\dim(B) - \dim(A_\mathfrak p)$. D'autre part, nous
avons
$\dim(B/\mathfrak p B) = \dim(A/\mathfrak p) + \dim(B/\mathfrak m_A B)$
et
$\dim(B) = \dim(A) + \dim(B/\mathfrak m_A B)$
par platitude (voir le lemme cité ci-dessus), et nous obtenons
$$
\dim(A_\mathfrak p) = \dim(A) - \dim(A/\mathfrak p)
$$
pour tout $\mathfrak p$ de $A$. En appliquant ceci à tous les idéaux premiers minimaux de
$A$, nous voyons que $A$ est équidimensionnel.
Si $\mathfrak p \subset \mathfrak p'$ est une inclusion stricte
sans idéal premier intermédiaire, alors nous pouvons appliquer ce qui précède à
l'idéal premier $\mathfrak p'/\mathfrak p$ de $A/\mathfrak p$,
car $A/\mathfrak p \to B/\mathfrak p B$ est plat et
$\Spec(B/\mathfrak p B)$ est équidimensionnel, pour obtenir
$$
1 = \dim((A/\mathfrak p)_{\mathfrak p'}) =
\dim(A/\mathfrak p) - \dim(A/\mathfrak p')
$$
Ainsi, $\mathfrak p \mapsto \dim(A/\mathfrak p)$ est une fonction
de dimension, et nous concluons que $A$ est caténaire.
\end{proof}

\begin{lemma}
\label{lemma-formally-catenary}
Soit $A$ un anneau local noethérien formellement caténaire.
Alors $A$ est universellement caténaire.
\end{lemma}

\begin{proof}
Nous pouvons remplacer $A$ par $A/\mathfrak p$, où $\mathfrak p$ est un idéal premier minimal
de $A$, voir Algèbre, lemme \ref{algebra-lemma-catenary-check-irreducible}.
Nous pouvons donc supposer que le spectre de $A^\wedge$ est équidimensionnel.
Il suffit de montrer que tout anneau local essentiellement de type fini
sur $A$ est caténaire (voir par exemple
Algèbre, lemme \ref{algebra-lemma-catenary-check-local}).
Il suffit donc de montrer que $A[x_1, \ldots, x_n]_\mathfrak m$ est caténaire,
où $\mathfrak m \subset A[x_1, \ldots, x_n]$ est un idéal
maximal au-dessus de $\mathfrak m_A$, voir
Algèbre, lemme \ref{algebra-lemma-localization-at-closed-point-special-fibre}
(et Algèbre, lemmes \ref{algebra-lemma-quotient-catenary} et
\ref{algebra-lemma-localization-catenary}).
Soit $\mathfrak m' \subset A^\wedge[x_1, \ldots, x_n]$ l'unique idéal
maximal au-dessus de $\mathfrak m$. Alors
$$
A[x_1, \ldots, x_n]_\mathfrak m \to A^\wedge[x_1, \ldots, x_n]_{\mathfrak m'}
$$
est local et plat (Algèbre, lemme \ref{algebra-lemma-completion-flat}).
Il suffit donc de montrer que l'anneau du membre de droite
est caténaire et de spectre équidimensionnel, voir le
lemme \ref{lemma-flat-under-catenary-equidimensional}.
Il est caténaire car les anneaux locaux complets sont universellement caténaires
(Algèbre, remarque
\ref{algebra-remark-Noetherian-complete-local-ring-universally-catenary}).
Choisissons un idéal premier minimal quelconque $\mathfrak q$ de
$A^\wedge[x_1, \ldots, x_n]_{\mathfrak m'}$. Alors
$\mathfrak q = \mathfrak p A^\wedge[x_1, \ldots, x_n]_{\mathfrak m'}$
pour un certain idéal premier minimal $\mathfrak p$ de $A^\wedge$ (petit détail omis).
Par conséquent,
$$
\dim(A^\wedge[x_1, \ldots, x_n]_{\mathfrak m'}/\mathfrak q) =
\dim(A^\wedge/\mathfrak p) + n = \dim(A^\wedge) + n
$$
la première égalité venant de
Algèbre, lemme \ref{algebra-lemma-dimension-base-fibre-equals-total},
et la seconde du fait que le spectre de $A^\wedge$ est équidimensionnel.
Ceci achève la démonstration.
\end{proof}

\begin{proposition}[Ratliff]
\label{proposition-ratliff}
\begin{reference}
\cite{Ratliff}
\end{reference}
Un anneau local noethérien est universellement caténaire si et seulement s'il
est formellement caténaire.
\end{proposition}

\begin{proof}
Combiner les lemmes \ref{lemma-not-formally-catenary} et
\ref{lemma-formally-catenary}.
\end{proof}

\begin{lemma}
\label{lemma-geometrically-normal-fibres-universally-catenary}
\begin{reference}
\cite[Corollary 2.3]{Heinzer-Rotthaus-Wiegand}
\end{reference}
Soit $(A, \mathfrak m)$ un anneau local noethérien dont les
fibres formelles sont géométriquement normales. Alors
\begin{enumerate}
\item $A^h$ est universellement caténaire, et
\item si $A$ est unibranche (par exemple normal), alors
$A$ est universellement caténaire.
\end{enumerate}
\end{lemma}

\begin{proof}
D'après le lemme \ref{lemma-geometrically-normal-formal-fibres-number-of-branches},
les nombres de branches de $A$ et de $A^\wedge$ sont les mêmes,
et le lemme \ref{lemma-equal-nr-branches-completion} s'applique donc.
Alors, pour tout idéal premier minimal $\mathfrak q \subset A^h$,
nous voyons que $A^\wedge/\mathfrak q A^\wedge$
possède un unique idéal premier minimal. Ainsi, $A^h$ est formellement caténaire
(par définition) et donc universellement caténaire d'après la
proposition \ref{proposition-ratliff}. Si $A$ est unibranche,
alors $A^h$ possède un unique idéal premier minimal, donc $A^\wedge$ possède
un unique idéal premier minimal, donc $A$ est formellement caténaire et
nous concluons de la même manière.
\end{proof}






\section{Actions de groupes et clôture intégrale}
\label{section-group-actions-integral}

\noindent
Cette section prolonge en un certain sens
Algèbre, section \ref{algebra-section-going-down-integral-over-normal}.
On trouvera davantage de résultats du même genre dans
Groupes fondamentaux, section \ref{pione-section-group-actions-integral}.

\begin{lemma}
\label{lemma-pol-lifting}
Soit $\varphi : A \to B$ une surjection d'anneaux. Soit $G$ un groupe fini
d'ordre $n$ agissant sur $\varphi : A \to B$. Si $b \in B^G$, alors
il existe un polynôme unitaire $P \in A^G[T]$ qui s'envoie sur
$(T - b)^n$ dans $B^G[T]$.
\end{lemma}

\begin{proof}
Choisissons $a \in A$ relevant $b$ et posons
$P = \prod_{\sigma \in G} (T - \sigma(a))$.
\end{proof}

\begin{lemma}
\label{lemma-invariants-modulo}
Soit $R$ un anneau. Soit $G$ un groupe fini agissant sur $R$. Soit $I \subset R$
un idéal tel que $\sigma(I) \subset I$ pour tout $\sigma \in G$.
Alors $R^G/I^G \subset (R/I)^G$ est une extension entière d'anneaux qui
induit un homéomorphisme sur les spectres et des extensions purement inséparables des
corps résiduels.
\end{lemma}

\begin{proof}
Comme $I^G = R^G \cap I$, il est clair que l'homomorphisme est injectif.
Le lemme \ref{lemma-pol-lifting} montre que
Algèbre, lemme \ref{algebra-lemma-universally-bijective},
s'applique.
\end{proof}

\begin{lemma}
\label{lemma-pol-in-ideal}
Soit $G$ un groupe fini d'ordre $n$ agissant sur un anneau $R$.
Soit $J \subset R^G$ un idéal. Pour $x \in JR$,
nous avons $\prod_{\sigma \in G} (T - \sigma(x)) =
T^n + a_1 T^{n - 1} + \ldots + a_n$, où $a_i \in J$.
\end{lemma}

\begin{proof}
Observons que le polynôme est bien unitaire et que ses coefficients appartiennent à $R^G$.
Nous pouvons écrire $x = f_1 b_1 + \ldots + f_m b_m$ avec $f_j \in J$ et $b_j \in R$.
Ainsi, en raisonnant par récurrence sur $m$, nous pouvons supposer que $x = y - fb$ avec
$f \in J$, $b \in R$ et $y \in JR$, et que le résultat vaut pour $y$.
Alors nous voyons que
$$
\prod\nolimits_{\sigma \in G} (T - \sigma(x)) =
\prod\nolimits_{\sigma \in G} (T - \sigma(y) + f\sigma(b)) =
\prod\nolimits_{\sigma \in G} (T - \sigma(y)) +
\sum_{i = 1, \ldots, n} f^i a_i
$$
où
$$
a_i =
\sum\nolimits_{S \subset G,\ |S| = i}
\prod\nolimits_{\sigma \in S} \sigma(b)
\prod\nolimits_{\sigma \not \in S} (T - \sigma(y))
$$
Un calcul que nous omettons montre que $a_i \in R^G$ (indication : l'expression donnée
est symétrique). Ainsi, le polynôme
de l'énoncé du lemme
pour $x$ est congru modulo $J$ au polynôme correspondant à $y$,
ce qui démontre l'étape de récurrence.
\end{proof}

\begin{lemma}
\label{lemma-invariants-modulo-bis}
Soit $R$ un anneau. Soit $G$ un groupe fini d'ordre $n$ agissant sur $R$.
Soit $J \subset R^G$ un idéal. Alors $R^G/J \to (R/JR)^G$ est un homomorphisme d'anneaux
tel que
\begin{enumerate}
\item pour $b \in (R/JR)^G$, il existe un polynôme unitaire
$P \in R^G/J[T]$ dont l'image dans $(R/JR)^G[T]$ est $(T - b)^n$,
\item pour $a \in \Ker(R^G/J \to (R/JR)^G)$,
nous avons $(T - a)^n = T^n$ dans $R^G/J[T]$.
\end{enumerate}
En particulier, $R^G/J \to (R/JR)^G$ est un homomorphisme entier d'anneaux qui
induit des homéomorphismes sur les spectres et des extensions purement inséparables des
corps résiduels.
\end{lemma}

\begin{proof}
L'assertion (1) découle du lemme \ref{lemma-pol-lifting} avec $I = JR$.
Si $a$ est comme dans l'assertion (2), alors $a$ est l'image de $x \in R^G \cap JR$.
Par conséquent, $(T - x)^n = \prod_{\sigma \in G} (T - \sigma(x))$
est congru à $T^n$ modulo $J$ d'après le lemme \ref{lemma-pol-in-ideal}.
Ceci démontre l'assertion (2). Pour la dernière assertion, nous pouvons appliquer
Algèbre, lemme \ref{algebra-lemma-universally-bijective}.
\end{proof}

\begin{remark}
\label{remark-when-iso}
Dans le lemme \ref{lemma-invariants-modulo-bis},
nous voyons que l'homomorphisme $R^G/J \to (R/JR)^G$
est un isomorphisme si $n$ est inversible dans $R$.
\end{remark}

\begin{lemma}
\label{lemma-functor-invariants-tensor}
Soit $R$ un anneau. Soit $G$ un groupe fini d'ordre $n$ agissant sur $R$.
Soit $A$ une $R^G$-algèbre.
\begin{enumerate}
\item pour $b \in (A \otimes_{R^G} R)^G$, il existe un polynôme unitaire
$P \in A[T]$ dont l'image dans $(A \otimes_{R^G} R)^G[T]$ est $(T - b)^n$,
\item pour $a \in \Ker(A \to (A \otimes_{R^G} R)^G)$, nous avons
$(T - a)^n = T^n$ dans $A[T]$.
\end{enumerate}
\end{lemma}

\begin{proof}
Choisissons une surjection $E \to A$, où $E$ est une algèbre polynomiale sur $R^G$.
Alors $(E \otimes_{R^G} R)^G = E$, car $E$ est libre comme $R^G$-module.
Notons $J = \Ker(E \to A)$. Comme le produit tensoriel est exact à droite, nous voyons
que $A \otimes_{R^G} R$ est le quotient de $E \otimes_{R^G} R$
par l'idéal engendré par $J$. Nous voyons ainsi que notre lemme
est un cas particulier du lemme \ref{lemma-invariants-modulo-bis}.
\end{proof}

\begin{lemma}
\label{lemma-base-change-invariants}
Soit $R$ un anneau. Soit $G$ un groupe fini agissant sur $R$.
Soit $R^G \to A$ un homomorphisme d'anneaux. L'homomorphisme
$$
A \to (A \otimes_{R^G} R)^G
$$
est un isomorphisme si $R^G \to A$ est plat. En général, cet homomorphisme
est entier, induit un homéomorphisme sur les spectres et
induit des extensions purement inséparables des corps résiduels.
\end{lemma}

\begin{proof}
Pour voir la première assertion, considérons la suite
exacte $0 \to R^G \to R \to \bigoplus_{\sigma \in G} R$,
où le second homomorphisme envoie $x$ sur $(\sigma(x) - x)_{\sigma \in G}$.
Après tensorisation par $A$, la suite reste exacte si $R^G \to A$ est plat.
Ainsi, $A$ est l'anneau des $G$-invariants de $A \otimes_{R^G} R$.

\medskip\noindent
La seconde assertion découle du lemme \ref{lemma-functor-invariants-tensor}
et de Algèbre, lemme \ref{algebra-lemma-universally-bijective}.
\end{proof}

\begin{lemma}
\label{lemma-one-orbit}
Soit $G$ un groupe fini agissant sur un anneau $R$. Pour deux idéaux premiers quelconques
$\mathfrak q, \mathfrak q' \subset R$ au-dessus du même idéal premier de $R^G$,
il existe un $\sigma \in G$ tel que $\sigma(\mathfrak q) = \mathfrak q'$.
\end{lemma}

\begin{proof}
L'extension $R^G \subset R$ est entière, car tout $x \in R$
est racine du polynôme unitaire $\prod_{\sigma \in G}(T - \sigma(x))$
de $R^G[T]$. Ainsi, il n'existe aucune relation d'inclusion entre les idéaux premiers
au-dessus d'un idéal premier donné $\mathfrak p$
(Algèbre, lemme \ref{algebra-lemma-integral-no-inclusion}).
Si le lemme était faux, alors
nous pourrions choisir $x \in \mathfrak q'$, $x \not \in \sigma(\mathfrak q)$
pour tout $\sigma \in G$. Voir Algèbre, lemme \ref{algebra-lemma-silly}.
Alors $y = \prod_{\sigma \in G} \sigma(x)$ appartient à $R^G$ et
à $\mathfrak p = R^G \cap \mathfrak q'$. D'autre part,
$x \not \in \sigma(\mathfrak q)$ pour tout $\sigma$ signifie que
$\sigma(x) \not \in \mathfrak q$ pour tout $\sigma$. Par conséquent,
$y \not \in \mathfrak q$, car $\mathfrak q$ est un idéal premier.
C'est impossible puisque $y \in \mathfrak p \subset \mathfrak q$.
\end{proof}

\begin{lemma}
\label{lemma-one-orbit-geometric}
Soit $G$ un groupe fini agissant sur un anneau $R$. Soit $\mathfrak q \subset R$
un idéal premier au-dessus de $\mathfrak p \subset R^G$. Alors
$\kappa(\mathfrak q)/\kappa(\mathfrak p)$ est une extension algébrique
normale et l'homomorphisme
$$
D = \{\sigma \in G \mid \sigma(\mathfrak q) = \mathfrak q\}
\longrightarrow
\text{Aut}(\kappa(\mathfrak q)/\kappa(\mathfrak p))
$$
est surjectif\footnote{Rappelons que nous n'utilisons la notation $\text{Gal}$
que dans le cas des extensions galoisiennes.}.
\end{lemma}

\begin{proof}
Avec $A = (R^G)_\mathfrak p$ et $B = A \otimes_{R^G} R$, nous voyons que $A = B^G$
car la localisation est plate, voir le lemme \ref{lemma-base-change-invariants}.
Observons que $\mathfrak pA$ et $\mathfrak qB$ sont des idéaux premiers,
que $D$ est le stabilisateur de $\mathfrak qB$, et que
$\kappa(\mathfrak p) = \kappa(\mathfrak pA)$ et
$\kappa(\mathfrak q) = \kappa(\mathfrak qB)$. Nous pouvons donc remplacer
$R$ par $B$ et supposer que $\mathfrak p$ est un idéal maximal.
Comme $R^G \subset R$ est une extension entière d'anneaux, nous trouvons
que les idéaux maximaux de $R$ sont exactement les idéaux premiers au-dessus de
$\mathfrak p$ (cela découle de
Algèbre, lemmes \ref{algebra-lemma-integral-no-inclusion} et
\ref{algebra-lemma-integral-going-up}).
D'après le lemme \ref{lemma-one-orbit}, ils sont en nombre fini,
$\mathfrak q = \mathfrak q_1, \mathfrak q_2, \ldots, \mathfrak q_m$
et forment une seule orbite sous $G$.
D'après le théorème chinois des restes
(Algèbre, lemme \ref{algebra-lemma-chinese-remainder}), l'homomorphisme
$R \to \prod_{j = 1, \ldots, m} R/\mathfrak q_j$ est surjectif.

\medskip\noindent
Nous démontrons d'abord que l'extension est normale. Choisissons un élément
$\alpha \in \kappa(\mathfrak q)$. Nous devons montrer que le
polynôme minimal $P$ de $\alpha$ sur $\kappa(\mathfrak p)$
est complètement scindé. D'après ce qui précède, nous pouvons choisir
$a \in \mathfrak q_2 \cap \ldots \cap \mathfrak q_m$
qui s'envoie sur $\alpha$ dans $\kappa(\mathfrak q)$.
Considérons le polynôme $Q = \prod_{\sigma \in G} (T - \sigma(a))$
de $R^G[T]$. L'image de $Q$ dans $R[T]$ est complètement scindée
en facteurs linéaires ; il en est donc de même de son
image dans $\kappa(\mathfrak q)[T]$. Comme $P$ divise
l'image de $Q$ dans $\kappa(\mathfrak p)[T]$, nous concluons
que $P$ est complètement scindé en facteurs linéaires sur
$\kappa(\mathfrak q)$, comme désiré.

\medskip\noindent
Comme $\kappa(\mathfrak q)/\kappa(\mathfrak p)$ est normale, nous pouvons écrire
$\kappa(\mathfrak q) = \kappa_1 \otimes_{\kappa(\mathfrak p)} \kappa_2$
avec $\kappa_1/\kappa(\mathfrak p)$ purement inséparable et
$\kappa_2/\kappa(\mathfrak p)$ galoisienne, voir
Corps, lemme \ref{fields-lemma-normal-case}. Choisissons
$\alpha \in \kappa_2$ qui engendre $\kappa_2$ sur $\kappa(\mathfrak p)$
si cette extension est finie, et qui engendre un sous-corps de degré $> |G|$ si elle est infinie
(afin d'obtenir une contradiction).
Ceci est possible d'après Corps, lemme \ref{fields-lemma-primitive-element}.
Choisissons $a$, $P$ et $Q$ comme dans le paragraphe précédent.
Si $\alpha' \in \kappa_2$ est un conjugué galoisien de $\alpha$ sur
$\kappa(\mathfrak p)$, alors le fait que $P$ divise l'
image de $Q$ dans $\kappa(\mathfrak p)[T]$ montre qu'il existe un
$\sigma \in G$ tel que $\sigma(a)$ s'envoie sur $\alpha'$. Par notre choix de
$a$ (qui s'annule aux autres idéaux maximaux), ceci implique que $\sigma \in D$ et
que l'image de $\sigma$ dans
$\text{Aut}(\kappa(\mathfrak q)/\kappa(\mathfrak p))$
envoie $\alpha$ sur $\alpha'$. Nous obtenons ainsi la surjectivité, ou bien
l'absurdité recherchée
dans le cas où $\alpha$ est de degré $> |G|$ sur $\kappa(\mathfrak p)$.
\end{proof}

\begin{lemma}
\label{lemma-one-orbit-geometric-galois}
Soit $A$ un domaine normal de corps des fractions $K$.
Soit $L/K$ une extension galoisienne (éventuellement infinie).
Soit $G = \text{Gal}(L/K)$ et soit
$B$ la clôture intégrale de $A$ dans $L$.
\begin{enumerate}
\item Pour deux idéaux premiers quelconques
$\mathfrak q, \mathfrak q' \subset B$ au-dessus du même idéal premier de $A$,
il existe un $\sigma \in G$ tel que $\sigma(\mathfrak q) = \mathfrak q'$.
\item Soit $\mathfrak q \subset B$ un idéal premier au-dessus de
$\mathfrak p \subset A$. Alors $\kappa(\mathfrak q)/\kappa(\mathfrak p)$
est une extension algébrique normale et l'homomorphisme
$$
D = \{\sigma \in G \mid \sigma(\mathfrak q) = \mathfrak q\}
\longrightarrow
\text{Aut}(\kappa(\mathfrak q)/\kappa(\mathfrak p))
$$
est surjectif.
\end{enumerate}
\end{lemma}

\begin{proof}
Démonstration de (1). Considérons les couples $(M, \sigma)$ où $K \subset M \subset L$
est un sous-corps tel que $M/K$ soit galoisienne, $\sigma \in \text{Gal}(M/K)$
et $\sigma(\mathfrak q \cap M) = \mathfrak q' \cap M$.
Nous disons que $(M', \sigma') \geq (M, \sigma)$ si et seulement si
$M \subset M'$ et $\sigma'|_M = \sigma$.
Observons que $(K, \text{id}_K)$ est un tel couple puisque $A = K \cap B$,
car $A$ est un domaine normal.
L'ensemble de ces couples satisfait les hypothèses du lemme de Zorn ;
il existe donc un couple maximal $(M, \sigma)$.
Si $M \not = L$, alors nous pouvons trouver
$M \subset M' \subset L$ avec $M'/M$ non triviale et finie et $M'/K$ galoisienne
(Corps, lemme \ref{fields-lemma-normal-closure-inside-normal}).
Choisissons $\sigma' \in \text{Gal}(M'/K)$ dont la restriction à $M$
est $\sigma$ (Corps, lemme \ref{fields-lemma-galois-infinite}).
Alors les idéaux premiers $\sigma'(\mathfrak q \cap M')$ et $\mathfrak q' \cap M'$
se restreignent au même idéal premier de $B \cap M$. Comme
$B \cap M = (B \cap M')^{\text{Gal}(M'/M)}$, nous pouvons
utiliser le lemme \ref{lemma-one-orbit} pour trouver $\tau \in \text{Gal}(M'/M)$
tel que $\tau(\sigma'(\mathfrak q \cap M')) = \mathfrak q' \cap M'$.
Ainsi, $(M', \tau \circ \sigma') > (M, \sigma)$,
ce qui contredit la maximalité de $(M, \sigma)$.

\medskip\noindent
L'assertion (2) se démontre exactement de la même manière que l'assertion (1). Nous
donnons les détails. Choisissons
$\overline{\sigma} \in \text{Aut}(\kappa(\mathfrak q)/\kappa(\mathfrak p))$.
Considérons les couples $(M, \sigma)$ où $K \subset M \subset L$
est un sous-corps tel que $M/K$ soit galoisienne, $\sigma \in \text{Gal}(M/K)$
et $\sigma(\mathfrak q \cap M) = \mathfrak q \cap M$, et où le diagramme
$$
\xymatrix{
\kappa(\mathfrak q \cap M) \ar[r] \ar[d]_\sigma &
\kappa(\mathfrak q) \ar[d]_{\overline{\sigma}} \\
\kappa(\mathfrak q \cap M) \ar[r] & \kappa(\mathfrak q)
}
$$
est commutatif. Nous disons que $(M', \sigma') \geq (M, \sigma)$ si et seulement si
$M \subset M'$ et $\sigma'|_M = \sigma$.
Comme ci-dessus, $(K, \text{id}_K)$ est un tel couple.
L'ensemble de ces couples satisfait les hypothèses du lemme de Zorn ;
il existe donc un couple maximal $(M, \sigma)$.
Si $M \not = L$, alors nous pouvons trouver
$M \subset M' \subset L$ avec $M'/M$ finie et $M'/K$ galoisienne
(Corps, lemme \ref{fields-lemma-normal-closure-inside-normal}).
Choisissons $\sigma' \in \text{Gal}(M'/K)$ dont la restriction à $M$
est $\sigma$ (Corps, lemme \ref{fields-lemma-galois-infinite}).
Alors les idéaux premiers $\sigma'(\mathfrak q \cap M')$ et $\mathfrak q \cap M'$
se restreignent au même idéal premier de $B \cap M$. En ajustant le choix
de $\sigma'$ comme dans le premier paragraphe, nous pouvons supposer que
$\sigma'(\mathfrak q \cap M') = \mathfrak q \cap M'$.
Alors $\sigma'$ et $\overline{\sigma}$ définissent des homomorphismes
$\kappa(\mathfrak q \cap M') \to \kappa(\mathfrak q)$
qui coïncident sur $\kappa(\mathfrak q \cap M)$. Comme
$B \cap M = (B \cap M')^{\text{Gal}(M'/M)}$, nous pouvons
utiliser le lemme \ref{lemma-one-orbit-geometric}
pour trouver $\tau \in \text{Gal}(M'/M)$ tel que
$\tau(\mathfrak q \cap M') = \mathfrak q \cap M'$
et tel que $\tau \circ \sigma'$ et $\overline{\sigma}$
induisent le même homomorphisme sur $\kappa(\mathfrak q \cap M')$.
Il y a ici un petit détail : le lemme garantit d'abord
que $\kappa(\mathfrak q \cap M')/\kappa(\mathfrak q \cap M)$
est normale, ce qui nous dit ensuite que la différence entre
les homomorphismes est un automorphisme de cette extension
(Corps, lemme \ref{fields-lemma-normal-embeddings-differ-by-aut}),
auquel nous pouvons
appliquer le lemme pour obtenir $\tau$. Ainsi, $(M', \tau \circ \sigma') > (M, \sigma)$,
ce qui contredit la maximalité de $(M, \sigma)$.
\end{proof}

\begin{lemma}
\label{lemma-one-orbit-geometric-galois-compare}
Soit $A$ un domaine normal de corps des fractions $K$.
Soit $M/L/K$ une tour d'extensions galoisiennes (éventuellement infinies) de $K$.
Soient $H = \text{Gal}(M/K)$ et $G = \text{Gal}(L/K)$, et soient
$C$ et $B$ les clôtures intégrales de $A$ dans $M$ et $L$.
Soient $\mathfrak r \subset C$ et $\mathfrak q = B \cap \mathfrak r$.
Posons
$D_\mathfrak r = \{\tau \in H \mid \tau(\mathfrak r) = \mathfrak r\}$
et
$I_\mathfrak r = \{\tau \in D_\mathfrak r \mid
\tau \bmod \mathfrak r = \text{id}_{\kappa(\mathfrak r)}\}$
et définissons de même $D_\mathfrak q$ et $I_\mathfrak q$.
Sous l'homomorphisme $H \to G$, les homomorphismes induits
$D_\mathfrak r \to D_\mathfrak q$ et
$I_\mathfrak r \to I_\mathfrak q$ sont surjectifs.
\end{lemma}

\begin{proof}
Soit $\sigma \in D_\mathfrak q$. Choisissons $\tau \in H$ s'envoyant sur $\sigma$.
Ceci est possible d'après Corps, lemme \ref{fields-lemma-galois-infinite}.
Alors $\tau(\mathfrak r)$ et $\mathfrak r$ sont tous deux au-dessus de $\mathfrak q$.
Par conséquent, d'après le lemme \ref{lemma-one-orbit-geometric-galois},
il existe un $\sigma' \in \text{Gal}(M/L)$ tel que
$\sigma'(\tau(\mathfrak r)) = \mathfrak r$. Ainsi,
$\sigma'\tau \in D_\mathfrak r$ s'envoie sur $\sigma$.
Le cas des groupes d'inertie se démontre exactement de la même
manière en utilisant la surjectivité sur les groupes d'automorphismes.
\end{proof}








\section{Extensions d'anneaux de valuation discrète}
\label{section-discrete-valuation-rings}

\noindent
Dans cette section et les quelques suivantes, nous utilisons les définitions ci-dessous.

\begin{definition}
\label{definition-extension-discrete-valuation-rings}
Nous disons que $A \to B$ ou $A \subset B$ est une
{\it extension d'anneaux de valuation discrète} si $A$ et $B$ sont des
anneaux de valuation discrète et si $A \to B$ est injectif et local.
En particulier, si $\pi_A$ et $\pi_B$ sont des uniformisants de
$A$ et de $B$, alors $\pi_A = u \pi_B^e$ pour un certain $e \geq 1$ et une unité
$u$ de $B$. L'entier $e$ ne dépend pas du choix des
uniformisants, puisqu'il est aussi l'unique entier $\geq 1$ tel que
$$
\mathfrak m_A B = \mathfrak m_B^e
$$
L'entier $e$ est appelé l'{\it indice de ramification} de $B$ sur $A$.
Nous disons que $B$ est {\it faiblement non ramifié} sur $A$ si $e = 1$.
Si l'extension des corps résiduels
$\kappa_A = A/\mathfrak m_A \subset \kappa_B = B/\mathfrak m_B$
est finie, alors nous posons $f = [\kappa_B : \kappa_A]$ et nous
l'appelons le {\it degré résiduel} ou {\it degré du corps résiduel}
de l'extension $A \subset B$.
\end{definition}

\noindent
Notons que nous n'exigeons pas que l'extension des corps des fractions soit finie.

\begin{lemma}
\label{lemma-inequality}
Soit $A \subset B$ une extension d'anneaux de valuation discrète de
corps des fractions $K \subset L$. Si l'extension $L/K$
est finie, alors l'extension des corps résiduels est finie et nous avons
$ef \leq [L : K]$.
\end{lemma}

\begin{proof}
La finitude de l'extension des corps résiduels est donnée par
Algèbre, lemme \ref{algebra-lemma-finite-extension-residue-fields-dimension-1}.
L'inégalité découle de
Algèbre, lemmes \ref{algebra-lemma-finite-length} et
\ref{algebra-lemma-pushdown-module}.
\end{proof}

\begin{lemma}
\label{lemma-multiplicative-e-f}
Soient $A \subset B \subset C$ des extensions d'anneaux de valuation discrète.
Alors le produit des indices de ramification de $B/A$ et de $C/B$ est
l'indice de ramification de $C/A$. En formule, $e_{C/A} = e_{B/A} e_{C/B}$.
Il en est de même des degrés résiduels lorsqu'ils sont finis.
\end{lemma}

\begin{proof}
Ceci découle immédiatement des définitions et de
Corps, lemme \ref{fields-lemma-multiplicativity-degrees}.
\end{proof}

\begin{lemma}
\label{lemma-ramification-index-a-power-of-p}
Soit $A \subset B$ une extension d'anneaux de valuation discrète
induisant l'extension de corps $K \subset L$. Si la caractéristique
de $K$ est $p > 0$ et si $L$ est purement inséparable sur $K$, alors
l'indice de ramification $e$ est une puissance de $p$.
\end{lemma}

\begin{proof}
Écrivons $\pi_A = u \pi_B^e$ pour un certain $u \in B^*$. D'autre part, nous avons
$\pi_B^q \in K$ pour une certaine puissance de $p$, notée $q$. Écrivons
$\pi_B^q = v \pi_A^k$ pour certains $v \in A^*$ et $k \in \mathbf{Z}$.
Alors $\pi_A^q = u^q \pi_B^{qe} = u^q v^e \pi_A^{ke}$.
En prenant les valuations dans $B$, nous concluons que $ke = q$.
\end{proof}

\noindent
Dans le lemme suivant, nous examinons ce que signifie pour une extension $A \subset B$
d'anneaux de valuation discrète d'être ``non ramifiée'', c'est-à-dire d'avoir un indice de
ramification égal à $1$ et une extension séparable (éventuellement non algébrique) des corps résiduels.
Cependant, nous ne pouvons pas employer le terme ``non ramifiée'' lui-même, car il existe déjà
une notion d'homomorphisme non ramifié d'anneaux, voir
Algèbre, section \ref{algebra-section-unramified}.

\begin{lemma}
\label{lemma-extension-dvrs-formally-smooth}
Soit $A \subset B$ une extension d'anneaux de valuation discrète.
Les assertions suivantes sont équivalentes :
\begin{enumerate}
\item $A \to B$ est formellement lisse pour la topologie $\mathfrak m_B$-adique, et
\item $A \to B$ est faiblement non ramifié et $\kappa_B/\kappa_A$
est une extension de corps séparable.
\end{enumerate}
\end{lemma}

\begin{proof}
Ceci découle de la proposition \ref{proposition-fs-flat-fibre-fs} et de
Algèbre, proposition
\ref{algebra-proposition-characterize-separable-field-extensions}.
\end{proof}

\begin{remark}
\label{remark-finite-separable-extension}
Soit $A$ un anneau de valuation discrète de corps des fractions $K$.
Soit $L/K$ une extension de corps finie séparable.
Soit $B \subset L$ la clôture intégrale de $A$ dans $L$.
Diagramme :
$$
\xymatrix{
B \ar[r] & L \\
A \ar[u] \ar[r] & K \ar[u]
}
$$
D'après Algèbre, lemme
\ref{algebra-lemma-Noetherian-normal-domain-finite-separable-extension}
l'extension d'anneaux $A \subset B$ est finie, et $B$ est donc noethérien.
D'après Algèbre, lemme \ref{algebra-lemma-integral-sub-dim-equal},
la dimension de $B$ est $1$, et $B$ est donc un domaine de Dedekind, voir
Algèbre, lemme \ref{algebra-lemma-characterize-Dedekind}.
Soient $\mathfrak m_1, \ldots, \mathfrak m_n$ les idéaux maximaux
de $B$ (c'est-à-dire les idéaux premiers au-dessus de $\mathfrak m_A$). Nous obtenons
des extensions d'anneaux de valuation discrète
$$
A \subset B_{\mathfrak m_i}
$$
et donc des indices de ramification $e_i$ et des degrés résiduels $f_i$. Nous avons
$$
[L : K] = \sum\nolimits_{i = 1, \ldots, n} e_i f_i
$$
d'après Algèbre, lemme \ref{algebra-lemma-finite-extension-dim-1},
appliqué à un uniformisant de $A$.
Nous observons que $n = 1$ si $A$ est hensélien (d'après
Algèbre, lemme \ref{algebra-lemma-finite-over-henselian}), par exemple si
$A$ est complet.
\end{remark}

\begin{definition}
\label{definition-types-of-extensions}
Soit $A$ un anneau de valuation discrète de corps des fractions $K$. Soit $L/K$
une extension finie séparable. Avec $B$ et
$\mathfrak m_i$, $i = 1, \ldots, n$
comme dans la remarque \ref{remark-finite-separable-extension}, nous disons que l'extension
$L/K$ est
\begin{enumerate}
\item {\it non ramifiée par rapport à $A$} si $e_i = 1$ et si l'extension
$\kappa(\mathfrak m_i)/\kappa_A$ est séparable pour tout $i$,
\item {\it modérément ramifiée par rapport à $A$}
si soit la caractéristique de $\kappa_A$
est $0$, soit la caractéristique de $\kappa_A$ est $p > 0$, les extensions de corps
$\kappa(\mathfrak m_i)/\kappa_A$ sont séparables
et les indices de ramification $e_i$ sont premiers à $p$, et
\item {\it totalement ramifiée par rapport à $A$}
si $n = 1$ et si l'extension des corps résiduels
$\kappa(\mathfrak m_1)/\kappa_A$ est triviale.
\end{enumerate}
Si l'anneau de valuation discrète $A$ ressort clairement du contexte, nous disons parfois,
en abrégé, que $L/K$ est non ramifiée, totalement ramifiée ou modérément ramifiée.
\end{definition}

\noindent
Pour les extensions non ramifiées, nous avons le lemme élémentaire suivant.

\begin{lemma}
\label{lemma-permanence-unramified}
Soit $A$ un anneau de valuation discrète de corps des fractions $K$.
\begin{enumerate}
\item Si $M/L/K$ sont des extensions finies séparables et si
$M$ est non ramifiée par rapport à $A$, alors $L$ est non ramifiée
par rapport à $A$.
\item Si $L/K$ est une extension finie séparable qui est
non ramifiée par rapport à $A$, alors il existe une extension galoisienne
$M/K$ contenant $L$ qui est non ramifiée par rapport à $A$.
\item Si $L_1/K$, $L_2/K$ sont des extensions finies séparables qui sont
non ramifiées par rapport à $A$, alors il existe une extension finie
séparable $L/K$ qui est non ramifiée par rapport
à $A$ et contient $L_1$ et $L_2$.
\end{enumerate}
\end{lemma}

\begin{proof}
Nous utiliserons les résultats de la discussion de la
remarque \ref{remark-finite-separable-extension}
sans autre mention.

\medskip\noindent
Démonstration de (1). Soient $C/B/A$ les clôtures intégrales de $A$ dans $M/L/K$.
Comme $C$ est une extension finie d'anneaux de $B$, nous voyons que
$\Spec(C) \to \Spec(B)$ est surjectif. Ainsi, pour tout idéal maximal
$\mathfrak m \subset B$, il existe un idéal maximal
$\mathfrak m' \subset C$ au-dessus de $\mathfrak m$.
Par multiplicativité des indices de ramification
(lemme \ref{lemma-multiplicative-e-f})
et par l'hypothèse, nous concluons que l'indice de
ramification de $B_\mathfrak m$ sur $A$ est $1$.
Comme $\kappa(\mathfrak m')/\kappa_A$ est finie séparable,
il en est de même de $\kappa(\mathfrak m)/\kappa_A$.

\medskip\noindent
Démonstration de (2). Soit $M$ la clôture normale de $L$ sur $K$, voir
Corps, définition \ref{fields-definition-normal-closure}.
Alors $M/K$ est galoisienne d'après Corps, lemme \ref{fields-lemma-normal-closure-galois}.
D'autre part, il existe une surjection
$$
L \otimes_K \ldots \otimes_K L \longrightarrow M
$$
d'algèbres sur $K$, voir Corps, lemme
\ref{fields-lemma-normal-closure-tensor-product}.
Soit $B$ la clôture intégrale de $A$ dans $L$
comme dans la remarque \ref{remark-finite-separable-extension}. La
condition que $L$ soit non ramifiée par rapport à $A$
signifie exactement que $A \to B$ est un homomorphisme étale d'anneaux, voir
Algèbre, lemme \ref{algebra-lemma-characterize-etale}.
D'après les propriétés de permanence des homomorphismes étales d'anneaux, nous voyons que
$$
B \otimes_A \ldots \otimes_A B
$$
est étale sur $A$, voir
Algèbre, lemme \ref{algebra-lemma-etale}.
L'anneau affiché est donc un produit de domaines de Dedekind, voir le
lemme \ref{lemma-Dedekind-etale-extension}.
Nous concluons que $M$ est le corps des fractions d'un
domaine de Dedekind fini étale sur $A$.
Cela signifie que $M$ est non ramifiée par rapport à $A$,
comme désiré.

\medskip\noindent
Démonstration de (3). Soit $B_i \subset L_i$ la clôture intégrale de $A$.
Raisonner de la même manière que ci-dessus pour montrer que $B_1 \otimes_A B_2$
est fini étale sur $A$. Détails omis.
\end{proof}

\begin{lemma}
\label{lemma-composition-unramified}
Soit $A$ un anneau de valuation discrète de corps des fractions $K$.
Soient $M/L/K$ des extensions finies séparables.
Soit $B$ la clôture intégrale de $A$ dans $L$.
Si $L/K$ est non ramifiée par rapport à $A$
et si $M/L$ est non ramifiée par rapport à $B_\mathfrak m$
pour tout idéal maximal $\mathfrak m$ de $B$, alors
$M/K$ est non ramifiée par rapport à $A$.
\end{lemma}

\begin{proof}
Soit $C$ la clôture intégrale de $A$ dans $M$.
Tout idéal maximal $\mathfrak m'$ de $C$ est au-dessus
d'un idéal maximal $\mathfrak m$ de $B$.
Le lemme découle alors de la multiplicativité
des indices de ramification (lemme \ref{lemma-multiplicative-e-f})
et du fait que nous avons la tour
$\kappa(\mathfrak m')/\kappa(\mathfrak m)/\kappa_A$
d'extensions finies de corps.
\end{proof}
\section{Extensions de Galois et ramification}
\label{section-ramification}

\noindent
Dans le cas des extensions de Galois, nous pouvons approfondir la discussion de la
Section \ref{section-discrete-valuation-rings}.

\begin{lemma}
\label{lemma-galois}
Soit $A$ un anneau de valuation discrète de corps des fractions $K$.
Soit $L/K$ une extension de Galois finie de groupe de Galois $G$.
Alors $G$ agit sur l'anneau $B$ de la Remarque \ref{remark-finite-separable-extension}
et agit transitivement sur l'ensemble des idéaux maximaux de $B$.
\end{lemma}

\begin{proof}
Remarquons que $A = B^G$, car $A$ est intégralement clos dans $K$ et $K = L^G$.
Ce lemme est donc un cas particulier du Lemme \ref{lemma-one-orbit}.
\end{proof}

\begin{lemma}
\label{lemma-galois-conclusion}
Soit $A$ un anneau de valuation discrète de corps des fractions $K$.
Soit $L/K$ une extension de Galois finie. Il existe alors $e \geq 1$ et
$f \geq 1$ tels que $e_i = e$ et $f_i = f$ pour tout $i$ (notations
de la Remarque \ref{remark-finite-separable-extension}). En particulier,
$[L : K] = n e f$.
\end{lemma}

\begin{proof}
Conséquence immédiate du Lemme \ref{lemma-galois} et des définitions.
\end{proof}

\begin{definition}
\label{definition-decomposition-inertia}
Soit $A$ un anneau de valuation discrète de corps des fractions $K$.
Soit $L/K$ une extension de Galois finie de groupe de Galois $G$.
Soit $B$ la clôture intégrale de $A$ dans $L$.
Soit $\mathfrak m \subset B$ un idéal maximal.
\begin{enumerate}
\item Le {\it groupe de décomposition de $\mathfrak m$}
est le sous-groupe $D = \{\sigma \in G \mid \sigma(\mathfrak m) = \mathfrak m\}$.
\item Le {\it groupe d'inertie de $\mathfrak m$} est le noyau $I$ du morphisme
$D \to \text{Aut}(\kappa(\mathfrak m)/\kappa_A)$.
\end{enumerate}
\end{definition}

\noindent
Remarquons que le corps $\kappa(\mathfrak m)$ peut être inséparable sur $\kappa_A$.
En particulier, l'extension de corps $\kappa(\mathfrak m)/\kappa_A$
n'est pas nécessairement galoisienne. Si $\kappa_A$ est parfait, elle l'est.

\begin{lemma}
\label{lemma-galois-galois}
Soit $A$ un anneau de valuation discrète de corps des fractions $K$ et de corps résiduel
$\kappa$. Soit $L/K$ une extension de Galois finie de groupe de Galois $G$.
Soit $B$ la clôture intégrale de $A$ dans $L$. Soit $\mathfrak m$ un idéal
maximal de $B$. Alors
\begin{enumerate}
\item l'extension de corps $\kappa(\mathfrak m)/\kappa$ est normale, et
\item $D \to \text{Aut}(\kappa(\mathfrak m)/\kappa)$ est surjectif.
\end{enumerate}
Si, pour un idéal maximal $\mathfrak m \subset B$ (ou, de manière équivalente, pour tout idéal maximal),
l'extension de corps $\kappa(\mathfrak m)/\kappa$ est séparable, alors
\begin{enumerate}
\item[(3)] $\kappa(\mathfrak m)/\kappa$ est galoisienne, et
\item[(4)] $D \to \text{Gal}(\kappa(\mathfrak m)/\kappa)$ est surjectif.
\end{enumerate}
Ici, $D \subset G$ est le groupe de décomposition de $\mathfrak m$.
\end{lemma}

\begin{proof}
Remarquons que $A = B^G$, car $A$ est intégralement clos dans $K$ et $K = L^G$.
Les parties (1) et (2) résultent donc du Lemme \ref{lemma-one-orbit-geometric}.
La partie « de manière équivalente, pour tout » du lemme résulte du
Lemme \ref{lemma-galois}. Supposons $\kappa(\mathfrak m)/\kappa$
séparable. Les parties (3) et (4) résultent alors immédiatement de (1) et (2).
\end{proof}

\begin{lemma}
\label{lemma-galois-inertia}
Soit $A$ un anneau de valuation discrète de corps des fractions $K$.
Soit $L/K$ une extension de Galois finie de groupe de Galois $G$.
Soit $B$ la clôture intégrale de $A$ dans $L$. Soit $\mathfrak m \subset B$
un idéal maximal. Le groupe d'inertie $I$ de $\mathfrak m$
s'insère dans une suite exacte canonique
$$
1 \to P \to I \to I_t \to 1
$$
de sorte que
\begin{enumerate}
\item si $D$ est le groupe de décomposition, on a
$$
P = \{\sigma \in I \mid \sigma|_{\mathfrak m/\mathfrak m^2} =
\text{id}_{\mathfrak m/\mathfrak m^2}\}
=
\{\sigma \in D \mid \sigma \text{ agit trivialement sur }\text{Gr}_\mathfrak m(B)\}
$$
\item $P$ est un sous-groupe normal de $D$,
\item $P$ est un $p$-groupe si la caractéristique de $\kappa_A$ est
$p > 0$, et $P = \{1\}$ si la caractéristique de $\kappa_A$ est nulle,
\item $I_t$ est cyclique d'ordre égal à la partie première à $p$ de l'entier $e$,
\item il existe un isomorphisme canonique
$\theta : I_t \to \mu_e(\kappa(\mathfrak m))$, et
\item
$P =
\{\sigma \in D \mid \sigma|_{B/\mathfrak m^2} = \text{id}_{B/\mathfrak m^2}\}$
si $\kappa(m)$ est séparable sur le corps résiduel de $A$.
\end{enumerate}
Ici, $e$ est l'entier du Lemme \ref{lemma-galois-conclusion}.
\end{lemma}

\begin{proof}
Rappelons que $|G| = [L : K] = nef$, voir le Lemme \ref{lemma-galois-conclusion}.
Puisque $G$ agit transitivement sur l'ensemble
$\{\mathfrak m_1, \ldots, \mathfrak m_n\}$ des idéaux maximaux de $B$
(Lemme \ref{lemma-galois})
et puisque $D$ est le stabilisateur d'un élément, on voit que $|D| = ef$.
Par le Lemme \ref{lemma-galois-galois}, on a
$$
ef = |D| = |I| \cdot |\text{Aut}(\kappa(\mathfrak m)/\kappa)|
$$
où $\kappa$ est le corps résiduel de $A$.
Comme $\kappa(\mathfrak m)$ est normal sur $\kappa$, l'ordre de
$\text{Aut}(\kappa(\mathfrak m)/\kappa)$ ne diffère de $f$ que par un facteur qui est
une puissance de $p$ (voir
Corps, Lemme \ref{fields-lemma-normal-and-automorphisms},
et la discussion qui suit
Corps, Définition \ref{fields-definition-insep-degree}).
Par conséquent, la partie première à $p$
de $|I|$ est égale à la partie première à $p$ de $e$.

\medskip\noindent
Posons $C = B_\mathfrak m$. Alors $I$ agit sur $C$ au-dessus de $A$ et trivialement
sur le corps résiduel de $C$. Soient $\pi_A \in A$ et $\pi_C \in C$ des
uniformisants. Écrivons $\pi_A = u \pi_C^e$ pour une unité $u$ de $C$.
Pour $\sigma \in I$, écrivons $\sigma(\pi_C) = \theta_\sigma \pi_C$ pour une
unité $\theta_\sigma$ de $C$. On a alors
$$
\pi_A = \sigma(\pi_A) = \sigma(u) (\theta_\sigma \pi_C)^e
= \sigma(u) \theta_\sigma^e \pi_C^e = \frac{\sigma(u)}{u} \theta_\sigma^e \pi_A
$$
Puisque $\sigma(u) \equiv u \bmod \mathfrak m_C$ car $\sigma \in I$,
on voit que l'image $\overline{\theta}_\sigma$
de $\theta_\sigma$ dans $\kappa_C = \kappa(\mathfrak m)$
est une racine $e$-ième de l'unité.
On obtient une application
\begin{equation}
\label{equation-inertia-character}
\theta : I \longrightarrow \mu_e(\kappa(\mathfrak m)),\quad
\sigma \mapsto \overline{\theta}_\sigma
\end{equation}
Nous affirmons que $\theta$ est un homomorphisme de groupes indépendant
du choix de l'uniformisant $\pi_C$. En effet, si $\tau$ est un second
élément de $I$, alors
$\tau(\sigma(\pi_C)) = \tau(\theta_\sigma \pi_C) =
\tau(\theta_\sigma) \theta_\tau \pi_C$, d'où
$\theta_{\tau \sigma} = \tau(\theta_\sigma) \theta_\tau$ et,
comme $\tau \in I$, on conclut que
$\overline{\theta}_{\tau \sigma} =
\overline{\theta}_\sigma \overline{\theta}_\tau$.
Si $\pi'_C$ est un second uniformisant, on voit alors
que $\pi'_C = w \pi_C$ pour une unité $w$ de $C$ et
$\sigma(\pi'_C) = w^{-1}\sigma(w)\theta_\sigma \pi'_C$,
d'où $\theta'_\sigma = w^{-1}\sigma(w)\theta_\sigma$ ;
ainsi, $\theta'_\sigma$ et $\theta_\sigma$
ont la même image dans le corps résiduel, comme précédemment.

\medskip\noindent
Puisque $\kappa(\mathfrak m)$ est de caractéristique $p$, le groupe
$\mu_e(\kappa(\mathfrak m))$ est cyclique d'ordre au plus égal à la partie
première à $p$ de $e$ (voir Corps, Section \ref{fields-section-roots-of-1}).

\medskip\noindent
Soit $P = \Ker(\theta)$. Par construction, les éléments de $P$ sont exactement
les éléments de $I$ qui agissent trivialement sur
$\mathfrak m/\mathfrak m^2 = \text{Gr}_\mathfrak m^1(B)$,
c'est-à-dire $P = \{\sigma \in I \mid \sigma|_{\mathfrak m/\mathfrak m^2} =
\text{id}_{\mathfrak m/\mathfrak m^2}\}$. De plus,
$I$ est formé des éléments de $D$ qui agissent trivialement sur
$\kappa(\mathfrak m) = B/\mathfrak m$.
Puisque l'anneau gradué $\text{Gr}_\mathfrak m(B)$ est engendré
par $\text{Gr}^1_\mathfrak m(B) = \mathfrak m/\mathfrak m^2$ sur
$\text{Gr}^0_\mathfrak m(B) = B/\mathfrak m$, on conclut que
$P = \{\sigma \in D \mid \sigma \text{ agit trivialement sur }
\text{Gr}_\mathfrak m(B)\}$.
Ainsi, (1) est vraie. Ceci implique (2), car $P$ est le noyau
de l'homomorphisme $D \to \text{Aut}(\text{Gr}_\mathfrak m(B))$.
Si nous démontrons (3), les parties (4) et (5) en découleront, car $I_t$
sera isomorphe à $\mu_e(\kappa(\mathfrak m))$, les arguments ci-dessus montrant
que $|I_t| \geq |\mu_e(\kappa(\mathfrak m))|$.

\medskip\noindent
Il suffit donc de démontrer que le
noyau $P$ de $\theta$ est un $p$-groupe.
Soit $\sigma$ un élément non trivial du noyau. Alors $\sigma - \text{id}$
envoie $\mathfrak m_C^i$ dans $\mathfrak m_C^{i + 1}$
pour tout $i$. Soit $m$ l'ordre de $\sigma$. Choisissons $c \in C$ tel
que $\sigma(c) \not = c$. Alors $\sigma(c) - c \in \mathfrak m_C^i$,
$\sigma(c) - c \not \in \mathfrak m_C^{i + 1}$ pour un certain $i$, et
on a
\begin{align*}
0
& =
\sigma^m(c) - c \\
& =
\sigma^m(c) - \sigma^{m - 1}(c) + \ldots + \sigma(c) - c \\
& =
\sum\nolimits_{j = 0, \ldots, m - 1} \sigma^j(\sigma(c) - c) \\
& \equiv
m(\sigma(c) - c) \bmod \mathfrak m_C^{i + 1}
\end{align*}
Il s'ensuit que $p | m$ (ou $m = 0$ si $p = 1$). Ainsi, tout élément du
noyau de $\theta$ a un ordre divisible par $p$, c'est-à-dire que $\Ker(\theta)$
est un $p$-groupe.

\medskip\noindent
Preuve de (6). Supposons $\kappa(\mathfrak m)/\kappa$ séparable.
Dans ce cas, $f = |D/I|$ par le Lemme \ref{lemma-galois-galois}
et $I$ est d'ordre $e$. Si $e = 1$, alors $P = I = \{\text{id}\}$
et le résultat est vrai. Si $e > 1$, alors
$B/\mathfrak m^2 = C/\pi_C^2C$ est une $\kappa$-algèbre
et une petite extension de $\kappa(\mathfrak m)$.
Comme $\kappa(\mathfrak m)$ est formellement étale sur $\kappa$
(Algèbre, Lemme
\ref{algebra-lemma-characterize-separable-algebraic-field-extensions})
il existe un unique morphisme de $\kappa$-algèbres
$\kappa(\mathfrak m) \to B/\mathfrak m^2$
donnant un inverse à droite de $B/\mathfrak m^2 \to \kappa(\mathfrak m)$.
Autrement dit, il existe un isomorphisme $D$-équivariant
$B/\mathfrak m^2 = \mathfrak m/\mathfrak m^2 \oplus \kappa(\mathfrak m)$
de $\kappa$-algèbres. Ceci montre immédiatement que
$P =
\{\sigma \in D \mid \sigma|_{B/\mathfrak m^2} = \text{id}_{B/\mathfrak m^2}\}$
dans ce cas.
\end{proof}

\begin{example}
\label{example-counter-description-P}
L'égalité de (6) dans le Lemme \ref{lemma-galois-inertia}
est fausse sans l'hypothèse sur $\kappa(\mathfrak m)$.
Un contre-exemple est l'extension de $A = Z_2[t]_{(2)}$
donnée par $B = A[x]/(x^2 - t)$. En effet, alors $\sigma(x) = -x = x - 2x$
et $2x$ n'appartient pas à $\mathfrak m^2$.
\end{example}

\begin{definition}
\label{definition-wild-inertia}
Reprenons les hypothèses et les notations du Lemme \ref{lemma-galois-inertia}.
\begin{enumerate}
\item Le {\it groupe d'inertie sauvage de $\mathfrak m$} est le sous-groupe $P$.
\item Le {\it groupe d'inertie modérée de $\mathfrak m$} est le groupe
quotient défini par $I \to I_t$.
\end{enumerate}
Nous notons $\theta : I \to \mu_e(\kappa(\mathfrak m))$ l'application surjective
(\ref{equation-inertia-character}) dont le noyau est $P$ et qui
induit l'isomorphisme $I_t \to \mu_e(\kappa(\mathfrak m))$.
\end{definition}

\begin{lemma}
\label{lemma-inertia-character}
Sous les hypothèses et avec les notations du Lemme \ref{lemma-galois-inertia},
le caractère d'inertie $\theta : I \to \mu_e(\kappa(\mathfrak m))$
satisfait la propriété suivante :
$$
\theta(\tau \sigma \tau^{-1}) = \tau(\theta(\sigma))
$$
pour $\tau \in D$ et $\sigma \in I$.
\end{lemma}

\begin{proof}
La formule a un sens car $I$ est un sous-groupe normal de $D$
et car $\tau$ agit sur $\kappa(\mathfrak m)$ par le morphisme
$D \to \text{Aut}(\kappa(\mathfrak m))$ étudié, par exemple, dans le
Lemme \ref{lemma-galois-galois}.
Rappelons la construction de $\theta$. Choisissons
un uniformisant $\pi$ de $B_\mathfrak m$ et, pour
$\sigma \in I$, écrivons $\sigma(\pi) = \theta_\sigma \pi$.
Alors $\theta(\sigma)$ est l'image $\overline{\theta}_\sigma$
de $\theta_\sigma$ dans le corps résiduel. Pour tout $\tau \in D$,
on peut écrire $\tau(\pi) = \theta_\tau \pi$ pour une unité $\theta_\tau$.
Alors $\theta_{\tau^{-1}} = \tau^{-1}(\theta_\tau^{-1})$.
On calcule
\begin{align*}
\theta_{\tau \sigma \tau^{-1}}
& =
\tau(\sigma(\tau^{-1}(\pi)))/\pi \\
& =
\tau(\sigma(\tau^{-1}(\theta_\tau^{-1}) \pi))/\pi \\
& =
\tau(\sigma(\tau^{-1}(\theta_\tau^{-1})) \theta_\sigma \pi)/\pi \\
& =
\tau(\sigma(\tau^{-1}(\theta_\tau^{-1}))) \tau(\theta_\sigma) \theta_\tau
\end{align*}
Cependant, puisque $\sigma$ agit trivialement modulo $\pi$, on voit que
le produit $\tau(\sigma(\tau^{-1}(\theta_\tau^{-1}))) \theta_\tau$
a pour image $1$ dans le corps résiduel. Ceci démontre le lemme.
\end{proof}

\noindent
Nous généraliserons le lemme suivant dans
Groupes fondamentaux, Lemme \ref{pione-lemma-inertial-invariants-unramified}.

\begin{lemma}
\label{lemma-inertial-invariants-unramified}
Soit $A$ un anneau de valuation discrète de corps des fractions $K$.
Soit $L/K$ une extension de Galois finie. Soit $\mathfrak m \subset B$
un idéal maximal de la clôture intégrale de $A$ dans $L$.
Soit $I \subset G$ le groupe d'inertie de $\mathfrak m$.
Alors $B^I$ est la clôture intégrale de $A$ dans $L^I$ et
$A \to (B^I)_{B^I \cap \mathfrak m}$ est étale.
\end{lemma}

\begin{proof}
Écrivons $B' = B^I$. Il résulte des définitions que $B' = B^I$
est la clôture intégrale de $A$ dans $L^I$. Écrivons
$\mathfrak m' = B^I \cap \mathfrak m = B' \cap \mathfrak m \subset B'$.
Par le Lemme \ref{lemma-one-orbit}, l'idéal maximal $\mathfrak m$ est l'unique
idéal premier de $B$ au-dessus de $\mathfrak m'$.
Comme $I$ agit trivialement sur $\kappa(\mathfrak m)$, le
Lemme \ref{lemma-invariants-modulo} montre que l'extension
$\kappa(\mathfrak m)/\kappa(\mathfrak m')$ est purement inséparable
(une autre possibilité, peut-être plus simple, est d'appliquer le résultat du
Lemme \ref{lemma-one-orbit-geometric}).
Puisque $D/I$ agit fidèlement sur $\kappa(\mathfrak m')$,
on conclut que $D/I$ agit fidèlement sur $\kappa(\mathfrak m)$.
Bien entendu, les éléments du corps résiduel $\kappa$ de $A$
sont fixes sous cette action.
La théorie de Galois montre que $[\kappa(\mathfrak m') : \kappa] \geq |D/I|$,
voir Corps, Lemme \ref{fields-lemma-galois-over-fixed-field}.

\medskip\noindent
Soit $\pi$ l'uniformisant de $A$. Puisque
$\text{Norm}_{L/K}(\pi) = \pi^{[L : K]}$, le
Lemme d'Algèbre \ref{algebra-lemma-finite-extension-dim-1}
montre que
$$
|G| = [L : K] = [L : K]\ \text{ord}_A(\pi) =
|G/D|\ [\kappa(\mathfrak m) : \kappa]\ \text{ord}_{B_\mathfrak m}(\pi)
$$
car il y a $n = |G/D|$ idéaux maximaux de $B$, tous
conjugués sous $G$, voir la
Remarque \ref{remark-finite-separable-extension} et le
Lemme \ref{lemma-galois}.
En appliquant le même raisonnement à
l'extension finie $L/L^I$ de degré $|I|$, on trouve
$$
|I|\ \text{ord}_{B'_{\mathfrak m'}}(\pi) =
[\kappa(\mathfrak m) : \kappa(\mathfrak m')]\ \text{ord}_{B_\mathfrak m}(\pi)
$$
On conclut que
$$
\text{ord}_{B'_{\mathfrak m'}}(\pi) =
\frac{|D/I|}{[\kappa(\mathfrak m') : \kappa]}
$$
Puisque le membre de gauche est un entier positif et que celui de droite
est $\leq 1$ d'après ce qui précède, on conclut qu'il y a égalité,
$\text{ord}_{B'_{\mathfrak m'}}(\pi) = 1$ et
$\kappa(\mathfrak m')/\kappa$ est de degré $|D/I|$.
Ainsi, $\pi B'_{\mathfrak m'} = \mathfrak m' B_\mathfrak m'$ et
$\kappa(\mathfrak m')$ est galoisien sur $\kappa$, de
groupe de Galois $D/I$, donc en particulier séparable, voir
Corps, Lemme \ref{fields-lemma-finite-Galois}.
Par Algèbre, Lemme \ref{algebra-lemma-characterize-etale},
on trouve que $A \to B'_{\mathfrak m'}$ est étale,
comme voulu.
\end{proof}

\begin{remark}
\label{remark-tower-of-rings}
Soit $A$ un anneau de valuation discrète de corps des fractions $K$.
Soit $L/K$ une extension de Galois finie. Soit $\mathfrak m \subset B$
un idéal maximal de la clôture intégrale de $A$ dans $L$.
Soit
$$
P \subset I \subset D \subset G
$$
la suite formée du groupe d'inertie sauvage, du groupe d'inertie et du groupe de décomposition de $\mathfrak m$.
Considérons le diagramme
$$
\xymatrix{
\mathfrak m \ar@{-}[d] \ar@{-}[r] &
\mathfrak m^P \ar@{-}[d] \ar@{-}[r] &
\mathfrak m^I \ar@{-}[d] \ar@{-}[r] &
\mathfrak m^D \ar@{-}[d] \ar@{-}[r] &
A \cap \mathfrak m \ar@{-}[d] \\
B & B^P \ar[l] & B^I \ar[l] & B^D \ar[l] & A \ar[l]
}
$$
Remarquons que $B^P, B^I, B^D$ sont les clôtures intégrales de
$A$ dans les corps $L^P$, $L^I$, $L^D$. On voit donc aussi que
$B^P$ est la clôture intégrale de $B^I$ dans $L^P$, et ainsi de suite.
Remarquons que $\mathfrak m^P = \mathfrak m \cap B^P$,
$\mathfrak m^I = \mathfrak m \cap B^I$, et
$\mathfrak m^D = \mathfrak m \cap B^D$. La
ligne supérieure du diagramme correspond donc aux images
de $\mathfrak m \in \Spec(B)$ par les morphismes induits entre
spectres. Cela étant dit, on a les assertions suivantes :
\begin{enumerate}
\item l'extension $L^I/L^D$ est galoisienne de groupe $D/I$,
\item l'extension $L^P/L^I$ est galoisienne de groupe $I_t = I/P$,
\item l'extension $L^P/L^D$ est galoisienne de groupe $D/P$,
\item $\mathfrak m^I$ est l'unique idéal premier de $B^I$ au-dessus de $\mathfrak m^D$,
\item $\mathfrak m^P$ est l'unique idéal premier de $B^P$ au-dessus de $\mathfrak m^I$,
\item $\mathfrak m$ est l'unique idéal premier de $B$ au-dessus de $\mathfrak m^P$,
\item $\mathfrak m^P$ est l'unique idéal premier de $B^P$ au-dessus de $\mathfrak m^D$,
\item $\mathfrak m$ est l'unique idéal premier de $B$ au-dessus de $\mathfrak m^I$,
\item $\mathfrak m$ est l'unique idéal premier de $B$ au-dessus de $\mathfrak m^D$,
\item $A \to B^D_{\mathfrak m^D}$ est étale et induit une
extension triviale des corps résiduels,
\item $B^D_{\mathfrak m^D} \to B^I_{\mathfrak m^I}$ est étale
et induit une extension galoisienne des corps résiduels de groupe de Galois
$D/I$,
\item $A \to B^I_{\mathfrak m^I}$ est étale,
\item $B^I_{\mathfrak m^I} \to B^P_{\mathfrak m^P}$
a pour indice de ramification $|I/P|$, premier à $p$, et induit une
extension triviale des corps résiduels,
\item $B^D_{\mathfrak m^D} \to B^P_{\mathfrak m^P}$
a pour indice de ramification $|I/P|$, premier à $p$, et induit une
extension séparable des corps résiduels,
\item $A \to B^P_{\mathfrak m^P}$
a pour indice de ramification $|I/P|$, premier à $p$, et induit une
extension séparable des corps résiduels.
\end{enumerate}
Les assertions (1), (2) et (3) résultent immédiatement de la théorie de Galois
(Corps, Section \ref{fields-section-galois-theory})
et du Lemme \ref{lemma-galois-inertia}.
Les assertions (4) -- (9) sont claires d'après le
Lemme \ref{lemma-galois}.
La partie (12) est le Lemme \ref{lemma-inertial-invariants-unramified}.
Puisque nous avons la factorisation
$A \to B^D_{\mathfrak m^D} \to B^I_{\mathfrak m^I}$
on en déduit l'étalité dans (10) et (11).
L'extension des corps résiduels de (10) doit être triviale, car
elle est séparable et l'homomorphisme de $D/I$ dans
$\text{Aut}(\kappa(\mathfrak m)/\kappa_A)$ est surjectif, comme le montre le
Lemme \ref{lemma-galois-galois}. Le même argument
démontre l'assertion sur les corps résiduels de (11).
Pour démontrer (13), (14) et (15), il suffit de démontrer (13).
D'après ce qui précède, l'extension $L^P/L^I$ est galoisienne,
de groupe de Galois cyclique d'ordre premier à $p$,
l'idéal premier $\mathfrak m^P$ est l'unique idéal premier au-dessus de
$\mathfrak m^I$, et l'action de $I/P$ sur le corps
résiduel est triviale. Nous pouvons donc appliquer le Lemme \ref{lemma-galois-inertia}
à cette extension et à l'anneau de valuation discrète
$B^I_{\mathfrak m^I}$ pour voir que (13) est satisfaite.
\end{remark}

\begin{lemma}
\label{lemma-compare-inertia}
Soit $A$ un anneau de valuation discrète de corps des fractions $K$.
Soit $M/L/K$ une tour telle que $M/K$ et $L/K$ soient des extensions
de Galois finies. Soient $C$ et $B$ les clôtures intégrales de $A$
dans $M$ et $L$. Soit $\mathfrak m' \subset C$ un idéal maximal et posons
$\mathfrak m = \mathfrak m' \cap B$. Soient
$$
P \subset I \subset D \subset \text{Gal}(L/K)
\quad\text{et}\quad
P' \subset I' \subset D' \subset \text{Gal}(M/K)
$$
les groupes d'inertie sauvage, d'inertie et de décomposition de
$\mathfrak m$ et $\mathfrak m'$, respectivement.
Alors la surjection canonique $\text{Gal}(M/K) \to \text{Gal}(L/K)$
induit des surjections $P' \to P$, $I' \to I$ et $D' \to D$. De plus,
elles s'insèrent dans les diagrammes commutatifs
$$
\vcenter{
\xymatrix{
D' \ar[r] \ar[d] &
\text{Aut}(\kappa(\mathfrak m')/\kappa_A) \ar[d] \\
D \ar[r] &
\text{Aut}(\kappa(\mathfrak m)/\kappa_A)
}
}
\quad\text{et}\quad
\vcenter{
\xymatrix{
I' \ar[r]_-{\theta'} \ar[d] &
\mu_{e'}(\kappa(\mathfrak m')) \ar[d]^{(-)^{e'/e}} \\
I \ar[r]^-\theta &
\mu_e(\kappa(\mathfrak m))
}
}
$$
où $e'$ et $e$ sont les indices de ramification de
$A \to C_{\mathfrak m'}$ et $A \to B_\mathfrak m$.
\end{lemma}

\begin{proof}
Le fait que l'application $\text{Gal}(M/K) \to \text{Gal}(L/K)$
envoie les groupes $P', I', D'$ dans $P, I, D$ résulte immédiatement
des définitions de ces groupes. La commutativité du premier diagramme
est claire (remarquons que, puisque $\kappa(\mathfrak m)/\kappa_A$ est normale,
tout automorphisme de $\kappa(\mathfrak m')$ au-dessus de $\kappa_A$ induit
bien un automorphisme de $\kappa(\mathfrak m)$ au-dessus de $\kappa_A$ ;
on obtient ainsi la flèche verticale droite du premier diagramme, voir le
Lemme \ref{lemma-galois-galois} et
Corps, Lemme \ref{fields-lemma-lift-maps}).

\medskip\noindent
Les applications $I' \to I$ et $D' \to D$ sont surjectives d'après le
Lemme \ref{lemma-one-orbit-geometric-galois-compare}.
La surjectivité de $P' \to P$ résulte de ce que $P'$ et $P$
sont des sous-groupes de Sylow p de $I'$ et $I$.

\medskip\noindent
Pour voir la commutativité du second diagramme, choisissons un uniformisant
$\pi'$ de $C_{\mathfrak m'}$ et un uniformisant $\pi$ de $B_\mathfrak m$.
Alors $\pi = c' (\pi')^{e'/e}$ pour une unité $c'$ de $C_{\mathfrak m'}$.
Pour $\sigma' \in I'$, l'image $\sigma \in I$ est simplement la restriction
de $\sigma'$ à $L$. Écrivons $\sigma'(\pi') = c \pi'$ pour une unité
$c \in C_{\mathfrak m'}$ et
$\sigma(\pi) = b \pi$ pour une unité $b$ de $B_\mathfrak m$.
Alors $\sigma'(\pi) = b \pi$ et nous obtenons
$$
b \pi = \sigma'(\pi) = \sigma'(c' (\pi')^{e'/e}) =
\sigma'(c') c^{e'/e} (\pi')^{e'/e} =
\frac{\sigma'(c')}{c'} c^{e'/e} \pi
$$
Comme $\sigma' \in I'$, nous voyons que $b$ et $c^{e'/e}$ ont la même
image dans le corps résiduel, ce qui démontre l'assertion voulue.
\end{proof}

\begin{remark}
\label{remark-canonical-inertia-character}
Afin d'utiliser le caractère d'inertie
$\theta : I \to \mu_e(\kappa(\mathfrak m))$
pour les extensions de Galois infinies, il convient de le normaliser.
Soient $A, K, L, B, \mathfrak m, G, P, I, D, e, \theta$
comme dans le Lemme \ref{lemma-galois-inertia} et la
Définition \ref{definition-wild-inertia}.
Alors $e = q |I_t|$, où $q$ est une puissance de la caractéristique $p$
de $\kappa(\mathfrak m)$ si celle-ci est positive, et vaut $1$ sinon.
Remarquons que $\mu_e(\kappa(\mathfrak m)) = \mu_{|I_t|}(\kappa(\mathfrak m))$
car la caractéristique de $\kappa(\mathfrak m)$ est $p$. Considérons
l'application
$$
\theta_{can} = q\theta : I \longrightarrow \mu_{|I_t|}(\kappa(\mathfrak m))
$$
Cette application induit un isomorphisme
$\theta_{can} : I_t \to \mu_{|I_t|}(\kappa(\mathfrak m))$.
On a $\theta_{can}(\tau \sigma \tau^{-1}) = \tau(\theta_{can}(\sigma))$
pour $\tau \in D$ et $\sigma \in I$
d'après le Lemme \ref{lemma-inertia-character}.
Enfin, si $M/L$ est une extension telle que $M/K$ soit galoisienne
et si $\mathfrak m'$ est un idéal premier de la clôture intégrale de $A$ dans $M$
au-dessus de $\mathfrak m$, alors on obtient le diagramme commutatif
$$
\xymatrix{
I' \ar[r]_-{\theta'_{can}} \ar[d] &
\mu_{|I'_t|}(\kappa(\mathfrak m')) \ar[d]^{(-)^{|I'_t|/|I_t|}} \\
I \ar[r]^-{\theta_{can}} &
\mu_{|I_t|}(\kappa(\mathfrak m))
}
$$
d'après le Lemme \ref{lemma-compare-inertia}.
\end{remark}













\section{Lemme de Krasner}
\label{section-krasner}

\noindent
Voici le lemme de Krasner dans le cas des corps à valuation discrète.

\begin{lemma}[Lemme de Krasner]
\label{lemma-krasner}
Soit $A$ un anneau local complet intègre de dimension $1$. Soit $P(t) \in A[t]$
un polynôme à coefficients dans $A$. Soit $\alpha \in A$ une racine
de $P$ qui n'est pas une racine de la dérivée $P' = \text{d}P/\text{d}t$.
Pour tout $c \geq 0$, il existe un entier $n$ tel que, pour tout
$Q \in A[t]$ dont les coefficients appartiennent à $\mathfrak m_A^n$, le polynôme
$P + Q$ possède une racine $\beta \in A$ telle que $\beta - \alpha \in \mathfrak m_A^c$.
\end{lemma}

\begin{proof}
Choisissons un élément non nul $\pi \in \mathfrak m$. Puisque $A$ est de dimension $1$,
on a $\mathfrak m = \sqrt{(\pi)}$. Par hypothèse, nous pouvons écrire
$P'(\alpha)^{-1} = \pi^{-m} a$ pour un entier $m \geq 0$ et un élément $a \in A$.
Nous pouvons supposer, et nous le faisons, que $c \geq m + 1$.
Choisissons $n$ tel que $\mathfrak m_A^n \subset (\pi^{c + m})$.
Choisissons un $Q$ quelconque comme dans l'énoncé. Pour la suite, remarquons que l'on peut écrire
$$
P(x + y) = P(x) + P'(x)y + R(x, y)y^2
$$
pour un certain $R(x, y) \in A[x, y]$. Nous allons montrer par récurrence que l'on peut trouver une
suite $\alpha_m, \alpha_{m + 1}, \alpha_{m + 2}, \ldots$ telle que
\begin{enumerate}
\item $\alpha_k \equiv \alpha \bmod \pi^c$,
\item $\alpha_{k + 1} - \alpha_k \in (\pi^k)$, et
\item $(P + Q)(\alpha_k) \in (\pi^{m + k})$.
\end{enumerate}
En posant $\beta = \lim \alpha_k$, on achèvera la démonstration.

\medskip\noindent
Initialisation. Puisque les coefficients de $Q$ appartiennent à $(\pi^{c + m})$, on a
$(P + Q)(\alpha) \in (\pi^{c + m})$. Ainsi, $\alpha_m = \alpha$ convient.
Ce choix garantit que $\alpha_k \equiv \alpha \bmod \pi^c$
pour tout $k \geq m$.

\medskip\noindent
Étape de récurrence. Étant donné $\alpha_k$, écrivons
$\alpha_{k + 1} = \alpha_k + \delta$ pour un certain $\delta \in (\pi^k)$.
On a alors
$$
(P + Q)(\alpha_{k + 1}) =
P(\alpha_k + \delta) + Q(\alpha_k + \delta)
$$
Puisque les coefficients de $Q$ appartiennent à $(\pi^{c + m})$, on voit que
$Q(\alpha_k + \delta) \equiv Q(\alpha_k) \bmod \pi^{c + m + k}$.
D'autre part, on a
$$
P(\alpha_k + \delta) =
P(\alpha_k) + P'(\alpha_k)\delta + R(\alpha_k, \delta)\delta^2
$$
Remarquons que $P'(\alpha_k) \equiv P'(\alpha) \bmod (\pi^{m + 1})$
car $\alpha_k \equiv \alpha \bmod \pi^{m + 1}$. On obtient donc
$$
P(\alpha_k + \delta) \equiv P(\alpha_k) + P'(\alpha) \delta
\bmod \pi^{k + m + 1}
$$
En réunissant les deux termes, on voit que
$$
(P + Q)(\alpha_{k + 1}) \equiv (P + Q)(\alpha_k) + P'(\alpha) \delta
\bmod \pi^{k + m + 1}
$$
Il suffit donc de prendre
$\delta = -P'(\alpha)^{-1} (P + Q)(\alpha_k) =  - \pi^{-m} a (P + Q)(\alpha_k)$
qui appartient à $(\pi^k)$ par l'hypothèse de récurrence.
\end{proof}

\begin{lemma}
\label{lemma-approximate-separable-extension}
Soit $A$ un anneau de valuation discrète de corps des fractions $K$.
Soit $A^\wedge$ le complété de $A$, de corps des fractions $K^\wedge$.
Si $M/K^\wedge$ est une extension séparable finie, il existe alors
une extension séparable finie $L/K$
telle que $M = K^\wedge \otimes_K L$.
\end{lemma}

\begin{proof}
Remarquons que $A^\wedge$ est lui aussi un anneau de valuation discrète (d'après les
Lemmes \ref{lemma-completion-regular} et \ref{lemma-completion-dimension}).
En particulier, $A^\wedge$ est intègre. La démonstration vaut plus généralement
pour les anneaux locaux noethériens $A$ tels que $A^\wedge$ soit un anneau local intègre
de dimension $1$.

\medskip\noindent
Soit $\theta \in M$ un élément qui engendre $M$ sur $K^\wedge$.
(Théorème de l'élément primitif.)
Soit $P(t) \in K^\wedge[t]$ le polynôme minimal de $\theta$ sur
$K^\wedge$. Soit $\pi \in \mathfrak m_A$ un élément non nul.
Après avoir remplacé $\theta$ par $\pi^n\theta$, on peut supposer que
les coefficients de $P(t)$ appartiennent à $A^\wedge$. Soit
$B = A^\wedge[\theta] = A^\wedge[t]/(P(t))$. Remarquons que $B$ est
un anneau local complet intègre de dimension $1$, car il est fini sur $A$ et
contenu dans $M$. Puisque $M$ est séparable sur $K$, l'élément $\theta$
n'est pas une racine de la dérivée de $P$. Pour tout entier $n$, on peut trouver
un polynôme unitaire $P_1 \in A[t]$ tel que les coefficients de $P - P_1$
appartiennent à $\pi^nA^\wedge[t]$. D'après le lemme de Krasner
(Lemme \ref{lemma-krasner}), $P_1$ possède une racine $\beta$ dans $B$
pour $n$ assez grand. De plus, on peut supposer (si $n$ est choisi assez grand)
que $\theta - \beta \in \pi B$. Considérons le morphisme
$\Phi : A^\wedge[t]/(P_1) \to B$ de $A^\wedge$-algèbres
qui envoie $t$ sur $\beta$. Puisque
$B = \pi B + \sum_{i < \deg(P)} A^\wedge \theta^i$, le morphisme $\Phi$
est surjectif d'après le lemme de Nakayama. Comme $\deg(P_1) = \deg(P)$,
il s'ensuit que $\Phi$ est un isomorphisme. Nous concluons que l'extension
$L = K[t]/(P_1(t))$ vérifie $K^\wedge \otimes_K L \cong M$.
Cela implique que $L$ est un corps et achève la démonstration.
\end{proof}

\begin{definition}
\label{definition-mixed}
Soit $A$ un anneau de valuation discrète. On dit que $A$ est de
{\it caractéristique mixte} si la caractéristique du corps résiduel de $A$
est $p > 0$ et celle du corps des fractions de $A$ est $0$.
Dans ce cas, on obtient une extension d'anneaux de valuation discrète
$\mathbf{Z}_{(p)} \subset A$, et l'{\it indice de ramification absolu}
de $A$ est l'indice de ramification de cette extension.
\end{definition}









\section{Lemme d'Abhyankar et ramification modérée}
\label{section-abhyankar-tame}

\noindent
Dans cette section, nous prouvons ce que nous pensons être la version la plus
générale du lemme d'Abhyankar pour les anneaux de valuation discrète. Après cela,
nous l'appliquons pour démontrer certains résultats concernant les extensions
modérément ramifiées du corps des fractions d'un anneau de valuation discrète.

\begin{remark}
\label{remark-construction}
Soit $A \to B$ une extension d'anneaux de valuation discrète avec corps des
fractions $K \subset L$. Soit $K_1/K$ une extension finie de corps. Soit $A_1 \subset K_1$ la
clôture intégrale de $A$ dans $K_1$. D'autre part, soit $L_1 = (L \otimes_K K_1)_{red}$. Alors
$L_1$ est un produit fini non vide d'extensions finies de corps de $L$. Soit
$B_1$ la clôture intégrale de $B$ dans $L_1$. Nous obtenons des
diagrammes commutatifs compatibles
$$
\vcenter{
\xymatrix{
L \ar[r] & L_1 \\
K \ar[u] \ar[r] & K_1 \ar[u]
}
}
\quad\text{et}\quad
\vcenter{
\xymatrix{
B \ar[r] & B_1 \\
A \ar[u] \ar[r] & A_1 \ar[u]
}
}
$$
Dans cette situation, nous avons ce qui suit
\begin{enumerate}
\item Par Algèbre, Lemme \ref{algebra-lemma-integral-closure-Dedekind} l'anneau $A_1$ est un anneau de Dedekind et
$B_1$ est un produit fini d'anneaux de Dedekind.
\item Notons que $L \otimes_K K_1 = (B \otimes_A A_1)_\pi$, où $\pi \in A$ est un uniformisant, et que $\pi$ est un
élément non-diviseur de zéro sur $B \otimes_A A_1$. Ainsi, l'homomorphisme d'anneaux $B \otimes_A A_1 \to B_1$
est intégral avec un noyau constitué d'éléments nilpotents. Par conséquent,
$\Spec(B_1) \to \Spec(B \otimes_A A_1)$ est surjective sur les spectres (Algèbre, Lemme \ref{algebra-lemma-integral-overring-surjective}). L'application
$\Spec(B \otimes_A A_1) \to \Spec(A_1)$ est surjective car $A_1/\mathfrak m_A A_1 \to B/\mathfrak m_AB \otimes_{\kappa_A} A_1/\mathfrak m_A A_1$ est un morphisme d'anneaux injectif avec
$A_1/\mathfrak m_A A_1$ Artinien. Nous concluons que $\Spec(B_1) \to \Spec(A_1)$ est surjective.
\item Soient $\mathfrak m_i$, $i = 1, \ldots n$ avec $n \geq 1$ les idéaux maximaux de $A_1$. Pour
chaque $i = 1, \ldots, n$, soient $\mathfrak m_{ij}$, $j = 1, \ldots, m_i$ avec $m_i \geq 1$ les idéaux maximaux de
$B_1$ au-dessus de $\mathfrak m_i$. Nous obtenons des diagrammes
$$
\xymatrix{
B \ar[r] & (B_1)_{\mathfrak m_{ij}} \\
A \ar[u] \ar[r] & (A_1)_{\mathfrak m_i} \ar[u]
}
$$
d'extensions d'anneaux à valuation discrète.
\item Si $A$ est hensélien (par exemple complet), alors $A_1$ est un
anneau à valuation discrète, c’est-à-dire $n = 1$. En effet, $A_1$ est une
union d'extensions finies de $A$ qui sont des domaines, donc locaux par
Algèbre, Lemme \ref{algebra-lemma-finite-over-henselian}.
\item Si $B$ est hensélien (par exemple complet), alors $B_1$ est un
produit d'anneaux à valuation discrète, c’est-à-dire $m_i = 1$ pour $i = 1, \ldots, n$.
\item Si $K \subset K_1$ est purement inséparable, alors $A_1$ et $B_1$ sont tous
deux des anneaux de valuation discrète, c'est-à-dire $n = 1$ et $m_1 = 1$. Ceci est
vrai parce que, pour chaque $b \in B_1$, une puissance d'exposant une puissance de $p$ de l'élément $b$ appartient à
$B$, donc $B_1$ ne peut avoir qu'un seul idéal maximal.
\item Si $K \subset K_1$ est une extension finie séparable, alors $L_1 = L \otimes_K K_1$ et est également un
produit fini d'extensions finies séparables. Par conséquent, $A \subset A_1$ et $B \subset B_1$
sont finis d'après Algèbre, Lemme \ref{algebra-lemma-Noetherian-normal-domain-finite-separable-extension}.
\item Si $A$ est Nagata, alors $A \subset A_1$ est fini.
\item Si $B$ est Nagata, alors $B \subset B_1$ est fini.
\end{enumerate}
\end{remark}

\begin{lemma}
\label{lemma-pull-root-uniformizer}
Soit $A$ un anneau de valuation discrète d'uniformisant $\pi$. Soit
$n \geq 2$. Soit $K_1 = K[\pi^{1/n}]$. Alors
\begin{enumerate}
\item $[K_1 : K] = n$.
\item la clôture intégrale $A_1$ de $A$ dans $K_1$ est l'anneau $A[\pi^{1/n}]$,
\item $A_1$ est un anneau de valuation discrète,
\item l'indice de ramification de $A_1$ sur $A$ est $n$,
\item $K_1$ est totalement ramifié par rapport à $A$, et
\item si $n$ est premier à la caractéristique du corps résiduel de $A$, alors
$K_1/K$ est modérément ramifiée et toute sous-extension de $K_1/K$ est engendrée
par $\pi^{1/d}$ pour un certain diviseur $d$ de $n$.
\end{enumerate}
\end{lemma}

\begin{proof}
Considérons l'anneau $A' = A[x]/(x^n - \pi)$ et notons $\pi'$ l'image de $x$ dans $A_1$.
Comme $A'$ est libre de rang fini $n$ en tant que $A$-module, l'élément
$\pi$ est un élément non nul-diviseur dans $A'$, et donc $\pi'$ l'est aussi.
De plus, $A'/\pi'A'$ est isomorphe à $A/\pi A$ qui est un corps. Ainsi, $A'$ est
un anneau de valuation discrète d'uniformisant $\pi'$. Il s'ensuit que le
corps des fractions $K'$ de $A'$ est une extension de degré $n$. De
toute évidence, $K'$ est obtenu de $K$ en adjoignant une racine
$n$-ième $\pi' = \pi^{1/n}$ de $\pi$, c'est-à-dire $K' \cong K_1$. Comme $A'$ est
normal, nous voyons que $A'$ est la clôture intégrale de $A$ dans
$K'$, c'est-à-dire $A' = A_1$. Cela prouve (1) -- (5).

\medskip\noindent L'affirmation selon laquelle $K_1/K$ est modérément ramifiée si
$n$ est inversible dans le corps résiduel résulte immédiatement de la
Définition \ref{definition-types-of-extensions} et des assertions (3) et (4). Pour prouver la toute dernière affirmation, il
suffit de montrer que chaque racine $n$-ième de l’unité $\zeta \in K'$ est contenue
dans $K$, voir Corps, Lemme \ref{fields-lemma-subfields-kummer}. Il est clair que $\zeta \in (A')^*$. Nous pouvons
écrire $\zeta = a + b$ avec $a \in A$ et $b \in A \pi' \oplus \ldots \oplus A (\pi')^{n - 1}$. Si $b$ est non nul, alors on peut
écrire $b = u(\pi')^t$ où $t \geq 1$ et $b$ une unité de $A'$. Alors
$$
1 = \zeta^n \equiv a^n + n a^{n - 1} u (\pi')^t \bmod (\pi')^{t + 1}
$$
Étant donné que $n$ et $a$ sont des unités de $A$, cela impliquerait
que $b = u(\pi')^t$ est congruent à $(1 - a^n)/(n a^{n - 1}) \in A$ modulo $(\pi')^{t + 1}$. La décomposition en somme
directe $A' = A \oplus \ldots \oplus A(\pi')^{n - 1}$ montre que cela est impossible, car l'image de la multiplication
par $(\pi')^{t + 1}$ est une somme directe correspondante dont l'intersection avec $A(\pi')^t$
est $(\pi A)(\pi')^t$. Ainsi, comme voulu, $b = 0$ et $\zeta \in A$.
\end{proof}

\begin{lemma}
\label{lemma-formally-smooth-goes-up}
Soit $A \to B$ une extension d'anneaux de valuation discrète avec des corps des
fractions $K \subset L$. Supposons que $A \to B$ soit formellement lisse dans la
topologie $\mathfrak m_B$-adique. Alors, pour toute extension finie $K_1/K$, nous avons
$L_1 = L \otimes_K K_1$, $B_1 = B \otimes_A A_1$, et chaque extension $(A_1)_{\mathfrak m_i} \subset (B_1)_{\mathfrak m_{ij}}$ (voir Remarque \ref{remark-construction}) est
formellement lisse dans la topologie $\mathfrak m_{ij}$-adique.
\end{lemma}

\begin{proof}
Nous utiliserons l’équivalence du Lemme \ref{lemma-extension-dvrs-formally-smooth} sans autre mention. Soient
$\pi \in A$ et $\pi_i \in (A_1)_{\mathfrak m_i}$ des uniformisants. Comme $\kappa_A \subset \kappa_B$ est séparable, l’anneau
$$
(B \otimes_A (A_1)_{\mathfrak m_i})/\pi_i (B \otimes_A (A_1)_{\mathfrak m_i}) =
B/\pi B \otimes_{A/\pi A} (A_1)_{\mathfrak m_i}/\pi_i (A_1)_{\mathfrak m_i}
$$
est un produit de corps séparables sur $\kappa_{\mathfrak m_i}$. Ainsi, l’élément $\pi_i$ dans $B \otimes_A (A_1)_{\mathfrak m_i}$
est un non-diviseur de zéro et le quotient par cet élément est un produit de corps. Il
s’ensuit que $B \otimes_A A_1$ est un anneau de Dedekind, donc en particulier réduit. Ainsi,
$B \otimes_A A_1 \subset B_1$ est une égalité.
\end{proof}

\noindent
Le lemme suivant est notre version du lemme d’Abhyankar pour les anneaux de
valuation discrète. Observons que $\kappa_B/\kappa_A$ n’est pas supposé être une extension
algébrique de corps.

\begin{lemma}[Lemme d’Abhyankar]
\label{lemma-abhyankar}
Soit $A \subset B$ une extension d'anneaux de valuation discrète. Supposons que
la caractéristique résiduelle de $A$ soit $0$, ou qu’elle soit $p$, que
l’indice de ramification $e$ soit premier à $p$ et que $\kappa_B/\kappa_A$ soit une
extension de corps séparable. Soit $K_1/K$ une extension finie. Avec les
notations de la Remarque \ref{remark-construction}, supposons que $e$ divise l’indice de ramification de
$A \subset (A_1)_{\mathfrak m_i}$ pour un certain $i$. Alors $(A_1)_{\mathfrak m_i} \subset (B_1)_{\mathfrak m_{ij}}$ est formellement lisse dans la
topologie $\mathfrak m_{ij}$-adique pour toute $j = 1, \ldots, m_i$.
\end{lemma}

\begin{proof}
Soit $\pi \in A$ un uniformisant. Soit $\pi_1$ un uniformisant de
$(A_1)_{\mathfrak m_i}$.
Écrivez $\pi = u \pi_1^{e_1}$ avec $u$ une unité de $(A_1)_{\mathfrak m_i}$ et $e_1$ l’indice de
ramification de $A \subset (A_1)_{\mathfrak m_i}$.

\medskip\noindent Affirmation : nous pouvons supposer que $u$ est une puissance
$e$-ième dans $K_1$. En effet, soit $K_2$ une extension de $K_1$ obtenue en
adjoignant une racine de $x^e = u$ ; ainsi $K_2$ est un facteur de $K_1[x]/(x^e - u)$. Alors
$K_2/K_1$ est une extension séparable finie (d'après notre hypothèse sur $e$) et
donc $A_1 \subset A_2$ est finie. Puisque $(A_1)_{\mathfrak m_i} \to (A_1)_{\mathfrak m_i}[x]/(x^e - u)$ est fini étale (car $e$ est
premier à la caractéristique résiduelle et $u$ une unité) nous concluons que
$(A_2)_{\mathfrak m_i}$ est un facteur d'une extension finie étale de $(A_1)_{\mathfrak m_i}$ et est donc elle-même finie étale
sur $(A_1)_{\mathfrak m_i}$. Le même raisonnement montre que $B_1 \subset B_2$ induit des
extensions finies étales $(B_1)_{\mathfrak m_{ij}} \subset (B_2)_{\mathfrak m_{ij}}$. Choisissez un idéal maximal $\mathfrak m'_{ij} \subset B_2$ au-dessus
de $\mathfrak m_{ij} \subset B_1$ (bien sûr il peut y en avoir plus d'un) et considérez
$$
\xymatrix{
(B_1)_{\mathfrak m_{ij}} \ar[r] & (B_2)_{\mathfrak m'_{ij}} \\
(A_1)_{\mathfrak m_i} \ar[u] \ar[r] &
(A_2)_{\mathfrak m'_i} \ar[u]
}
$$
où $\mathfrak m'_i \subset A_2$ est l’image. Maintenant, les flèches horizontales ont un indice de
ramification $1$ et induisent des extensions finies de corps résiduels
séparables. Ainsi, en utilisant l’équivalence du Lemme \ref{lemma-extension-dvrs-formally-smooth}, on voit qu’il
suffit de montrer que la flèche verticale droite est formellement lisse dans la
topologie $\mathfrak m'_{ij}$-adique. Puisque $u$ a une racine $e$-ième dans $K_2$, nous
obtenons l'affirmation.

\medskip\noindent Supposons que $u$ ait une racine $e$-ième dans $K_1$.
Puisque $e | e_1$ et puisque $u$ a une racine $e$-ième dans $K_1$, nous voyons
que $\pi = \theta^e$ pour un certain $\theta \in K_1$. Soit $K'_1 = K[\theta] \subset K_1$ le sous-corps engendré par
$\theta$. D'après le Lemme \ref{lemma-pull-root-uniformizer}, la clôture intégrale $A'_1$ de $A$
dans $K[\theta]$ est l'anneau de valuation discrète $A'_1 = A[\theta]$ qui a un indice de
ramification $e$ sur $A$. Si nous pouvons prouver le lemme pour
l'extension $K'_1/K$, alors nous concluons d'après le Lemme \ref{lemma-formally-smooth-goes-up} appliqué au
diagramme
$$
\xymatrix{
(B'_1)_{B'_1 \cap \mathfrak m_{ij}} \ar[r] & (B_1)_{\mathfrak m_{ij}} \\
A'_1 \ar[u] \ar[r] &
(A_1)_{\mathfrak m_i} \ar[u]
}
$$
pour tout $j = 1, \ldots, m_i$. Cela nous ramène au cas discuté dans le paragraphe suivant.

\medskip\noindent Supposons $K_1 = K[\pi^{1/e}]$ et posons $\theta = \pi^{1/e}$. Soit $\pi_B$ un uniformisant
pour $B$ et écrivons $\pi = w \pi_B^e$ pour une certaine unité $w$ de $B$. Alors
nous voyons que $L_1 = L \otimes_K K_1$ est obtenu en adjoignant $\theta / \pi_B$ qui est une racine
$e$-ième de l'unité $w$. Ainsi $B \subset B_1$ est fini et étale. Ainsi, pour tout
idéal maximal $\mathfrak m \subset B_1$, considérons le diagramme commutatif
$$
\xymatrix{
B \ar[r]_1 & (B_1)_{\mathfrak m} \\
A \ar[u]^e \ar[r]^e & A_1 \ar[u]_{e_\mathfrak m}
}
$$
Ici, les nombres le long des flèches sont les indices de ramification. Par la
multiplicativité des indices de ramification (Lemme \ref{lemma-multiplicative-e-f}), nous concluons
$e_\mathfrak m = 1$. En regardant les extensions des corps résiduels, nous trouvons que
$\kappa(\mathfrak m)$ est une extension finie séparable de $\kappa_B$ qui est séparable sur
$\kappa_A$. Par conséquent, $\kappa(\mathfrak m)$ est séparable sur $\kappa_A$ qui est égal au corps
résiduel de $A_1$ et le résultat découle du Lemme \ref{lemma-extension-dvrs-formally-smooth}.
\end{proof}

\begin{lemma}
\label{lemma-composition-tame}
Soit $A$ un anneau à valuation discrète avec corps des fractions $K$.
Soient $M/L/K$ des extensions finies séparables. Soit $B$ la clôture
intégrale de $A$ dans $L$. Si $L/K$ est modérément ramifiée par
rapport à $A$ et $M/L$ est modérément ramifiée par rapport à $B_\mathfrak m$ pour
chaque idéal maximal $\mathfrak m$ de $B$, alors $M/K$ est modérément ramifiée
par rapport à $A$.
\end{lemma}

\begin{proof}
Soit $C$ la clôture intégrale de $A$ dans $M$. Tout idéal maximal
$\mathfrak m'$ de $C$ est au-dessus d'un idéal maximal $\mathfrak m$ de $B$. Alors le
lemme découle de la multiplicativité des indices de ramification (Lemme \ref{lemma-multiplicative-e-f})
et du fait que nous avons la tour $\kappa(\mathfrak m')/\kappa(\mathfrak m)/\kappa_A$ d'extensions finies de corps.
\end{proof}

\begin{lemma}
\label{lemma-subextension-tame}
Soit $A$ un anneau de valuation discrète avec corps des fractions $K$. Si
$M/L/K$ sont des extensions finies séparables et que $M$ est modérément
ramifiée par rapport à $A$, alors $L$ est modérément ramifiée par
rapport à $A$.
\end{lemma}

\begin{proof}
Nous utiliserons les résultats de la discussion dans la Remarque \ref{remark-finite-separable-extension} sans
autre mention. Soient $C/B/A$ les clôtures intégrales de $A$ dans $M/L/K$.
Puisque $C$ est une extension finie d'anneaux de $B$, nous voyons que
$\Spec(C) \to \Spec(B)$ est surjective. Par conséquent, pour tout idéal maximal $\mathfrak m \subset B$, il existe
un idéal maximal $\mathfrak m' \subset C$ au-dessus de $\mathfrak m$. Par la multiplicativité des indices
de ramification (Lemme \ref{lemma-multiplicative-e-f}) et l'hypothèse, nous concluons que l'indice de
ramification de $B_\mathfrak m$ sur $A$ est premier à la caractéristique résiduelle.
Puisque $\kappa(\mathfrak m')/\kappa_A$ est finie et séparable, il en est de même pour $\kappa(\mathfrak m)/\kappa_A$.
\end{proof}

\begin{lemma}
\label{lemma-characterize-tame}
Soit $A$ un anneau de valuation discrète avec un corps des fractions $K$.
Soit $\pi \in A$ un uniformisant. Soit $L/K$ une extension finie séparable.
Les éléments suivants sont équivalents
\begin{enumerate}
\item $L$ est modérément ramifiée par rapport à $A$,
\item il existe un entier $e \geq 1$ inversible dans $\kappa_A$ et une extension $L'/K' = K[\pi^{1/e}]$ non
ramifiée par rapport à $A' = A[\pi^{1/e}]$ telle que $L$ est contenue dans $L'$, et
\item il existe un $e_0 \geq 1$ inversible dans $\kappa_A$ tel que pour tout $d \geq 1$
inversible dans $\kappa_A$ (2) est vrai avec $e = de_0$.
\end{enumerate}
\end{lemma}

\begin{proof}
Remarquez que $A'$ est un anneau de valuation discrète avec corps des fractions
$K'$, voir le Lemme \ref{lemma-pull-root-uniformizer}. Bien sûr, l'indice de ramification de $A'$ sur
$A$ est $e$. Ainsi, si (2) est vrai, alors $L'$ est modérément ramifié
par rapport à $A$ selon le Lemme \ref{lemma-composition-tame}. Par conséquent, $L$ est
modérément ramifié par rapport à $A$ selon le Lemme \ref{lemma-subextension-tame}.

\medskip\noindent L'implication (3) $\Rightarrow$ (2) est immédiate.

\medskip\noindent Supposons que (1) soit vérifié. Soit $B$ la clôture intégrale
de $A$ dans $L$ et soit $\mathfrak m_1, \ldots, \mathfrak m_n$ ses idéaux maximaux. On note $e_i$
l'indice de ramification de $A \to B_{\mathfrak m_i}$. Soit $e_0$ le plus petit commun multiple de
$e_1, \ldots, e_r$. Cela est inversible dans $\kappa_A$ selon notre hypothèse (1). Soit $e = de_0$
comme en (3). Posons $A' = A[\pi^{1/e}]$. Alors $A \to A'$ est une extension d'anneaux à
valuation discrète avec corps des fractions $K' = K[\pi^{1/e}]$, voir le Lemme \ref{lemma-pull-root-uniformizer}.
Choisissez une décomposition en produit
$$
L \otimes_K K' = \prod L'_j
$$
où les $L'_j$ sont des corps. Soit $B'_j$ la clôture intégrale de $A$ dans
$L'_j$. Soient $\mathfrak m_{ijk}$ les idéaux maximaux de $B'_j$ situés au-dessus de $\mathfrak m_i$.
On observe que $(B'_j)_{\mathfrak m_i}$ est la clôture intégrale de $B_{\mathfrak m_i}$ dans $L'_j$. D'après
le lemme d'Abhyankar (Lemme \ref{lemma-abhyankar}) appliqué à $A \subset B_{\mathfrak m_i}$ et à l'extension
$K'/K$, nous voyons que $A' \to (B'_j)_{\mathfrak m_{ijk}}$ est formellement lisse dans la topologie
$\mathfrak m_{ijk}$-adique. Cela implique que l'indice de ramification est $1$ et que
l'extension des corps résiduels est séparable (Lemme \ref{lemma-extension-dvrs-formally-smooth}). De cette manière,
nous voyons que $L'_j$ n'est pas ramifié par rapport à $A'$. Cela termine la
démonstration : nous prenons $L' = L'_j$ pour un certain $j$.
\end{proof}

\begin{lemma}
\label{lemma-permanence-tame}
Soit $A$ un anneau de valuation discrète avec le corps des fractions
$K$.
\begin{enumerate}
\item Si $L/K$ est une extension séparable finie qui est modérément ramifiée
par rapport à $A$, alors il existe une extension de Galois $M/K$ contenant
$L$ qui est modérément ramifiée par rapport à $A$.
\item Si $L_1/K$, $L_2/K$ sont des extensions séparables finies qui sont
modérément ramifiées par rapport à $A$, alors il existe une extension
séparable finie $L/K$ qui est modérément ramifiée par rapport à $A$
contenant $L_1$ et $L_2$.
\end{enumerate}
\end{lemma}

\begin{proof}
Preuve de (2). Choisissons un uniformisant $\pi \in A$. Nous pouvons choisir un entier
$e$ inversible dans $\kappa_A$ et des extensions $L_i'/K' = K[\pi^{1/e}]$ non ramifiées par
rapport à $A' = A[\pi^{1/e}]$, où $L'_i/L_i$ est une extension de $K$, voir le Lemme \ref{lemma-characterize-tame}.
D'après le Lemme \ref{lemma-permanence-unramified}, nous pouvons trouver une extension $L'/K'$ qui est non
ramifiée par rapport à $A'$ telle que $L'_i/K$ soit isomorphe à une
sous-extension de $L'/K'$ pour $i = 1, 2$. Cela termine la preuve de (3) car
$L'/K$ est modérément ramifiée (utiliser le même lemme que ci-dessus).

\medskip\noindent
Preuve de (1). Nous pouvons d'abord remplacer $L$ par une extension plus
grande et supposer que $L$ est une extension de $K' = K[\pi^{1/e}]$ non ramifiée par
rapport à $A' = A[\pi^{1/e}]$ où $e$ est inversible dans $\kappa_A$, voir le Lemme \ref{lemma-characterize-tame}.
Soit $M$ la clôture normale de $L$ sur $K$, voir Corps, Définition
\ref{fields-definition-normal-closure}. Alors $M/K$ est galoisienne d'après Corps, Lemme \ref{fields-lemma-normal-closure-galois}. D'autre
part, il existe une surjection
$$
L \otimes_K \ldots \otimes_K L \longrightarrow M
$$
d'algèbres $K$, voir Corps, Lemme \ref{fields-lemma-normal-closure-tensor-product}. Soit $B$ la clôture
intégrale de $A$ dans $L$ comme dans la Remarque \ref{remark-finite-separable-extension}. La condition
que $L$ est non ramifiée par rapport à $A' = A[\pi^{1/e}]$ signifie exactement que
$A' \to B$ est un morphisme d'anneaux étale, voir Algèbre, Lemme \ref{algebra-lemma-characterize-etale}.
Affirmation :
$$
K' \otimes_K \ldots \otimes_K K' = \prod K'_i
$$
est un produit d'extensions de corps $K'_i/K$ modérément ramifiées par rapport à
$A$. Alors si $A'_i$ est la clôture intégrale de $A$ dans $K'_i$, nous
voyons que
$$
\prod A'_i \otimes_{(A' \otimes_A \ldots \otimes_A A')}
(B \otimes_A \ldots \otimes_A B)
$$
est fini et étale sur $\prod A'_i$ et donc un produit d'anneaux de Dedekind (Lemme
\ref{lemma-Dedekind-etale-extension}). Nous en concluons que $M$ est le corps des fractions de l'un de ces
anneaux de Dedekind qui est fini et étale sur $A'_i$ pour un certain $i$. Il
s'ensuit que $M/K'_i$ est non ramifié par rapport à chaque idéal maximal de
$A'_i$ et donc $M/K$ est modérément ramifiée par le Lemme \ref{lemma-composition-tame}.

\medskip\noindent Il reste à prouver l'affirmation. Pour cela nous écrivons
$A' = A[x]/(x^e - \pi)$ et nous voyons que
$$
A' \otimes_A \ldots \otimes_A A' =
A'[x_1, \ldots, x_r]/(x_1^e - \pi, \ldots, x_r^e - \pi)
$$
La normalisation de cet anneau contient certainement les éléments $y_i = x_i/x_1$ pour
$i = 2, \ldots, r$ soumis aux relations $y_i^e - 1 = 0$ et nous obtenons
$$
A[x_1, y_2, \ldots, y_r]/(x_1^e - \pi, y_2^e - 1, \ldots, y_r^e - 1) =
A'[y_2, \ldots, y_r]/(y_2^e - 1, \ldots, y_r^e - 1)
$$
Cet anneau est fini étale sur $A'$ parce que $e$ est inversible dans
$A'$. Par conséquent, il est un produit d'anneaux de Dedekind, chacun non
ramifié sur $A'$ comme voulu (voir les références données ci-dessus pour plus
de détails).
\end{proof}

\begin{lemma}
\label{lemma-tame-goes-up}
Soit $A \subset B$ une extension d'anneaux de valuation discrète. Notons $L/K$
l'extension correspondante des corps des fractions. Soit $K'/K$ une extension
finie séparable. Alors
$$
K' \otimes_K L = \prod L'_i
$$
est un produit fini de corps et ce qui suit est vrai
\begin{enumerate}
\item Si $K'$ est non ramifié par rapport à $A$, alors chaque $L'_i$
est non ramifié par rapport à $B$.
\item Si $K'$ est modérément ramifié par rapport à $A$, alors chaque
$L'_i$ est modérément ramifié par rapport à $B$.
\end{enumerate}
\end{lemma}

\begin{proof}
L'algèbre $K' \otimes_K L$ est un produit fini de corps car c'est une algèbre finie étale
sur $L$. Soit $A'$ la clôture intégrale de $A$ dans $K'$.

\medskip\noindent Dans le cas (1) l’application d’anneaux $A \to A'$ est finie
\'etale. Ainsi, $B' = B \otimes_A A'$ est fini \'etale sur $B$ et est un produit fini de
anneaux de Dedekind (Lemme \ref{lemma-Dedekind-etale-extension}). Ainsi, $B'$ est la clôture intégrale de
$B$ dans $K' \otimes_K L$. Il s’ensuit immédiatement que chaque $L'_i$ n’est pas
ramifiée par rapport à $B$.

\medskip\noindent Choisissons un uniformisant $\pi \in A$. Pour prouver (2), on peut
remplacer $K'$ par une extension plus grande modérément ramifiée par rapport à $A$
(détails omis ; indication : utiliser le Lemme \ref{lemma-subextension-tame}). Ainsi, par le Lemme \ref{lemma-characterize-tame}, on
peut supposer qu’il existe un $e \geq 1$ inversible dans $\kappa_A$ tel que
$K'$ contient $K[\pi^{1/e}]$ et telle que $K'$ n’est pas ramifiée par rapport à
$A[\pi^{1/e}]$. Choisissons une décomposition en produit
$$
K[\pi^{1/e}] \otimes_K L = \prod L_{e, j}
$$
Pour chaque $i$, il existe un $j_i$ tel que $L'_i/L_{e, j_i}$ soit une extension finie
séparable. Soit $B_{e, j}$ la clôture intégrale de $B$ dans $L_{e, j}$. D'après (1)
appliqué à $K'/K[\pi^{1/e}]$ et $A[\pi^{1/e}] \subset (B_{e, j_i})_\mathfrak m$, nous voyons que $L'_i$ est non ramifié par rapport
à $(B_{e, j_i})_\mathfrak m$ pour tout idéal maximal $\mathfrak m \subset B_{e, j_i}$. Ainsi, la preuve sera complète si nous
pouvons montrer que $L_{e, j}$ est modérément ramifié par rapport à $B$, voir
le Lemme \ref{lemma-composition-tame}.

\medskip\noindent Choisissons un uniformisant $\theta$ dans $B$. Écrivons
$\pi = u \theta^t$ où $u$ est une unité de $B$ et $t \geq 1$. Alors nous avons
$$
A[\pi^{1/e}] \otimes_A B = B[x]/(x^e - u \theta^t)
\subset B[y, z]/(y^{e'} - \theta, z^e - u)
$$
où $e' = e/\gcd(e, t)$. L'application envoie $x$ sur $z y^{t/\gcd(e, t)}$. Comme le côté droit est
un produit d'anneaux de Dedekind, chacun modérément ramifié sur $B$, la
preuve est complète (détails omis).
\end{proof}












\section{Élimination de la ramification}
\label{section-eliminating-ramification}

\noindent
Dans cette section, nous discutons d'un résultat de Helmut Epp, voir \cite{Epp}. Nous
encourageons fortement le lecteur à lire l'original. Notre approche est légèrement
différente car nous essayons de traiter les cas de caractéristique mixte et
équicaractéristique par la même méthode. Pour des résultats liés, voir aussi
\cite{Ponomarev}, \cite{Ponomarev-Abhyankar}, \cite{Kuhlmann} et \cite{ZK}.

\medskip\noindent Soit $A \subset B$ une extension d'anneaux de valuation discrète avec
des corps des fractions $K \subset L$. Le but de cette section est de trouver une
extension finie $K_1/K$ telle qu'avec
$$
\vcenter{
\xymatrix{
L \ar[r] & L_1 \\
K \ar[u] \ar[r] & K_1 \ar[u]
}
}
\quad\text{et}\quad
\vcenter{
\xymatrix{
B \ar[r] & B_1 \ar[r] & (B_1)_{\mathfrak m_{ij}} \\
A \ar[u] \ar[r] & A_1 \ar[r] \ar[u] & (A_1)_{\mathfrak m_i} \ar[u]
}
}
$$
comme dans la Remarque \ref{remark-construction}, les extensions $(A_1)_{\mathfrak m_i} \subset (B_1)_{\mathfrak m_{ij}}$ soient toutes faiblement
non ramifiées ou même formellement lisses dans les topologies adiques pertinentes.
L'exemple le plus simple (mais non trivial) de cela est le lemme d'Abhyankar, voir
le Lemme \ref{lemma-abhyankar}.

\begin{definition}
\label{definition-solution}
Soit $A \to B$ une extension d'anneaux de valuation discrète avec des corps des
fractions $K \subset L$.
\begin{enumerate}
\item Nous disons qu'une extension de corps finie $K_1/K$ est une {\it
solution faible pour $A \subset B$} si toutes les extensions $(A_1)_{\mathfrak m_i} \subset (B_1)_{\mathfrak m_{ij}}$ de la Remarque \ref{remark-construction}
sont faiblement non ramifiées.
\item Nous disons qu'une extension de corps finie $K_1/K$ est une {\it
solution pour $A \subset B$} si chaque extension $(A_1)_{\mathfrak m_i} \subset (B_1)_{\mathfrak m_{ij}}$ de la Remarque \ref{remark-construction} est
formellement lisse dans la topologie $\mathfrak m_{ij}$-adique.
\end{enumerate}
Nous disons qu'une solution $K_1/K$ est une {\it solution séparable} si
$K_1/K$ est séparable.
\end{definition}

\noindent
En général, les solutions (faibles) n'existent pas ; il existe un exemple dans
\cite{Epp}. Sous une hypothèse faible sur l'extension du corps résiduel, nous
prouverons l'existence de solutions faibles dans le Théorème \ref{theorem-epp} en suivant
\cite{Epp}. Dans la section suivante, nous déduirons l'existence de solutions et
parfois de solutions séparables dans des cas géométriquement significatifs, voir
la Proposition \ref{proposition-epp-essentially-finite-type} et le Lemme \ref{lemma-epp-essentially-finite-type-separable}. Cependant, l'exemple suivant montre
qu'en général, il faut des extensions inséparables pour obtenir même une solution
faible.

\begin{example}
\label{example-inseparable-necessary}
Soit $k$ un corps parfait de caractéristique $p > 0$. Soient $A = k[[x]]$ et
$K = k((x))$. Soit $B = A[x^{1/p}]$. Toute solution faible $K_1/K$ pour $A \to B$ est
inséparable (et toute extension finie inséparable de $K$ est une solution).
Nous omettons la démonstration.
\end{example}

\noindent
Les solutions sont stables sous des extensions supplémentaires, voir le Lemme
\ref{lemma-solution-goes-up}. Cela peut ne pas être vrai pour les solutions faibles. Les solutions
faibles sont en quelque sorte stables sous des extensions totalement ramifiées,
voir le Lemme \ref{lemma-weakly-unramified-goes-up-along-totally-ramified}.

\begin{lemma}
\label{lemma-weakly-unramified-goes-up-along-totally-ramified}
Soit $A \to B$ une extension d'anneaux de valuation discrète avec les corps des
fractions $K \subset L$. Supposons que $A \to B$ soit faiblement non ramifiée. Alors, pour
toute extension finie séparable $K_1/K$ totalement ramifiée par rapport à $A$,
alors $L_1 = L \otimes_K K_1$ est un corps, $A_1$ et $B_1 = B \otimes_A A_1$ sont des anneaux de
valuation discrète, et l'extension $A_1 \subset B_1$ (voir la Remarque \ref{remark-construction}) est
faiblement non ramifiée.
\end{lemma}

\begin{proof}
Soient $\pi \in A$ et $\pi_1 \in A_1$ des uniformisants. Comme $K_1/K$ est
totalement ramifiée par rapport à $A$, nous avons $\pi_1^e = u_1 \pi$ pour une unité $u_1$ de
$A_1$. Ainsi, $A_1$ est engendré par $\pi_1$ sur $A$ et le polynôme
minimal $P(t)$ de $\pi_1$ sur $K$ a la forme
$$
P(t) = t^e + a_{e - 1} t^{e - 1} + \ldots + a_0
$$
avec $a_i \in (\pi)$ et $a_0 = u\pi$ pour une certaine unité $u$ de $A$. Notez que
$e = [K_1 : K]$ aussi. Puisque $A \to B$ est faiblement non ramifié, on voit que $\pi$
est un uniformisant de $B$ et donc $B_1 = B[t]/(P(t))$ est un anneau de valuation
discrète dont la classe de $t$ est un uniformisant. Ainsi, le lemme est clair.
\end{proof}

\begin{lemma}
\label{lemma-solutions-go-down}
Soient $A \to B \to C$ des extensions d’anneaux de valuation discrète de corps
des fractions $K \subset L \subset M$. Soit $K_1/K$ une extension finie.
\begin{enumerate}
\item Si $K_1$ est une solution (faible) pour $A \to C$, alors $K_1$ est
une solution (faible) pour $A \to B$.
\item Si $K_1$ est une solution (faible) pour $A \to B$ et $L_1 = (L \otimes_K K_1)_{red}$ est un
produit de corps qui sont des solutions (faibles) pour $B \to C$, alors $K_1$ est
une solution (faible) pour $A \to C$.
\end{enumerate}
\end{lemma}

\begin{proof}
Soient $L_1 = (L \otimes_K K_1)_{red}$ et $M_1 = (M \otimes_K K_1)_{red}$ et soient $B_1 \subset L_1$ et $C_1 \subset M_1$ les clôtures intégrales respectives de
$B$ et de $C$. Notez que $M_1 = (M \otimes_L L_1)_{red}$ et que $L_1$ est un produit fini (non
vide) d'extensions finies de $L$. Par conséquent, l'homomorphisme d'anneaux
$B_1 \to C_1$ est un produit fini de morphismes d'anneaux de la forme discutée dans la
Remarque \ref{remark-construction}. En particulier, l'application $\Spec(C_1) \to \Spec(B_1)$ est surjective.
Choisissez un idéal maximal $\mathfrak m \subset C_1$ et considérez les extensions d'anneaux à
valuation discrète
$$
(A_1)_{A_1 \cap \mathfrak m} \to
(B_1)_{B_1 \cap \mathfrak m} \to
(C_1)_\mathfrak m
$$
Si la composition est faiblement non ramifiée, il en est de même pour
l'application $(A_1)_{A_1 \cap \mathfrak m} \to (B_1)_{B_1 \cap \mathfrak m}$. Si l'extension du corps résiduel $\kappa_{A_1 \cap \mathfrak m} \to \kappa_\mathfrak m$ est séparable, il
en est de même pour la sous-extension $\kappa_{A_1 \cap \mathfrak m} \to \kappa_{B_1 \cap \mathfrak m}$. En tenant compte du Lemme \ref{lemma-extension-dvrs-formally-smooth},
cela prouve (1). Un argument similaire fonctionne pour (2).
\end{proof}

\begin{lemma}
\label{lemma-solution-after-strict-henselization}
Soit $A \to B$ une extension d'anneaux de valuation discrète. Il existe un diagramme
commutatif
$$
\xymatrix{
B \ar[r] & B' \\
A \ar[r] \ar[u] & A' \ar[u]
}
$$
d'extensions d'anneaux de valuation discrète tel que
\begin{enumerate}
\item les extensions $K'/K$ et $L'/L$ des corps des fractions soient
algébriques séparables,
\item les corps résiduels de $A'$ et $B'$ soient des clôtures algébriques
séparables des corps résiduels de $A$ et $B$, et
\item si une solution, une solution faible ou une solution séparable existe pour
$A' \to B'$, alors une solution, une solution faible ou une solution séparable existe
pour $A \to B$.
\end{enumerate}
\end{lemma}

\begin{proof}
Par Algèbre, Lemme \ref{algebra-lemma-colimit-finite-etale-given-residue-field}, il existe une extension $A \subset A'$ qui est une
colimite filtrée d'extensions étales finies telle que le corps résiduel
de $A'$ est une clôture algébrique séparable du corps résiduel de $A$.
Alors $A \subset A'$ est une extension d'anneaux à valuation discrète telle que
l'extension induite $K'/K$ des corps des fractions est algébrique séparable.

\medskip\noindent Soit $B \subset B'$ une hensélisation stricte de $B$. Alors $B \subset B'$
est une extension d'anneaux de valuation discrète dont l'extension des corps des
fractions est algébrique séparable. D'après Algèbre, Lemme \ref{algebra-lemma-strictly-henselian-functorial-prepare}, il existe un
diagramme commutatif comme dans l'énoncé du lemme. Les parties (1) et (2) du lemme
sont évidentes.

\medskip\noindent Soit $K'_1/K'$ une solution (faible) pour $A' \to B'$. Puisque $A'$
est une colimite, nous pouvons trouver une extension étale finie $A \subset A_1'$ et
une extension finie $K_1$ du corps des fractions $F$ de $A_1'$ telle que
$K'_1 = K' \otimes_F K_1$. Comme $A \subset A_1'$ est finie étale et $B'$ strictement
hensélien, il s'ensuit que $B' \otimes_A A_1'$ est un produit fini d'anneaux isomorphes à
$B'$. Par conséquent
$$
L' \otimes_K K_1 = L' \otimes_K F \otimes_F K_1
$$
est un produit fini d’anneaux isomorphes à $L' \otimes_{K'} K'_1$. Ainsi, on voit que $K_1/K$ est
une solution (faible) pour $A \to B'$. C’est donc aussi une solution (faible) pour
$A \to B$ par le Lemme \ref{lemma-solutions-go-down}.
\end{proof}

\begin{lemma}
\label{lemma-galois-relative}
Soit $A \to B$ une extension d'anneaux de valuation discrète avec des corps des
fractions $K \subset L$. Soit $K_1/K$ une extension normale. Posons $G = \text{Aut}(K_1/K)$. Alors $G$
agit sur les anneaux $K_1$, $L_1$, $A_1$ et $B_1$ de la Remarque \ref{remark-construction}
et agit de manière transitive sur l’ensemble des idéaux maximaux de $B_1$.
\end{lemma}

\begin{proof}
Tout est clair à part la dernière affirmation. S’il y a deux orbites ou plus de
l’action, alors on peut trouver un élément $b \in B_1$ qui appartient à tous les idéaux
maximaux d’une orbite et dont le résidu vaut $1$ modulo tous les idéaux maximaux d’une
autre orbite. Alors $b' = \prod_{\sigma \in G} \sigma(b)$ est un élément $G$-invariant de $B_1 \subset L_1 = (L \otimes_K K_1)_{red}$ qui
appartient à certains idéaux maximaux de $B_1$ mais pas à tous les idéaux
maximaux de $B_1$. En le relevant en un élément de $L \otimes_K K_1$ et en l'élevant à
une puissance suffisamment grande, on obtient un élément $G$-invariant $b''$ de $L \otimes_K K_1$
dont l'image est $(b')^N$ pour un certain $N > 0$ ;
en fait, nous n'avons besoin de faire cela que dans le cas où la caractéristique
est $p > 0$ et, dans ce cas, l'élévation à une puissance de $p$ suffisamment grande,
notée $q$, définit une application canonique $(L \otimes_K K_1)_{red} \to L \otimes_K K_1$. Comme $K = (K_1)^G$, nous concluons
que $b'' \in L$. Comme $b''$ a pour image un élément de $B_1$, nous voyons que
$b'' \in B$ (car $B$ est normal). Ensuite, d'une part il doit être vrai que
$b'' \in \mathfrak m_B$ car $b'$ est dans un certain idéal maximal de $B_1$ et d'autre
part il doit être vrai que $b'' \not \in \mathfrak m_B$ car $b'$ n'est pas dans tous les idéaux
maximaux de $B_1$. Cette contradiction termine la démonstration du lemme.
\end{proof}

\begin{lemma}
\label{lemma-make-degree-q-extension}
Soit $A$ un anneau de valuation discrète muni d'un uniformisant $\pi$. Si la
caractéristique résiduelle de $A$ est $p > 0$, alors pour tout $n > 1$ et
toute puissance de $p$, notée $q$, il existe une extension séparable de degré
$q$, notée $L/K$, totalement ramifiée par rapport à $A$, telle que la
clôture intégrale $B$ de $A$ dans $L$ ait un indice de ramification $q$ et un
uniformisant $\pi_B$ tel que $\pi_B^q = \pi + \pi^n b$ et $\pi_B^q = \pi + (\pi_B)^{nq}b'$ pour certains $b, b' \in B$.
\end{lemma}

\begin{proof}
Si la caractéristique de $K$ est nulle, alors on peut prendre l’extension
donnée par $\pi_B^q = \pi$, voir le Lemme \ref{lemma-pull-root-uniformizer}. Si la caractéristique de $K$ est
$p > 0$, alors on peut prendre l’extension de $K$ donnée par $z^q - \pi^n z = \pi^{1 - q}$. À
savoir, on voit alors que $y^q - \pi^{n + q - 1} y = \pi$ où $y = \pi z$. En prenant $\pi_B = y$, nous obtenons
le résultat souhaité.
\end{proof}

\begin{lemma}
\label{lemma-pre-purely-inseparable-case}
Soit $A$ un anneau de valuation discrète. Supposons que le corps résiduel
$\kappa_A$ ait pour caractéristique $p > 0$ et que $a \in A$ soit un élément dont la
classe résiduelle dans $\kappa_A$ n'est pas une $p$-ième puissance. Alors $a$
n'est pas une $p$-ième puissance dans $K$ et la clôture intégrale de $A$
dans $K[a^{1/p}]$ est l'anneau $A[a^{1/p}]$ qui est un anneau de valuation discrète
faiblement non ramifié sur $A$.
\end{lemma}

\begin{proof}
Ce lemme est immédiat.
\end{proof}

\begin{lemma}
\label{lemma-purely-inseparable-case}
Soient $A \subset B \subset C$ des extensions d'anneaux de valuation discrète avec des corps des
fractions $K \subset L \subset M$. Soit $\pi \in A$ un uniformisant. Supposons que
\begin{enumerate}
\item $B$ soit un anneau de Nagata,
\item $A \subset B$ soit faiblement non ramifié,
\item $M$ soit une extension purement inséparable de degré $p$ de
$L$.
\end{enumerate}
Alors l'un des cas suivants se produit :
\begin{enumerate}
\item $A \to C$ est faiblement non ramifié, soit
\item $C = B[\pi^{1/p}]$, ou
\item il existe une extension séparable de degré $p$ $K_1/K$ totalement
ramifiée par rapport à $A$ telle que $L_1 = L \otimes_K K_1$ et $M_1 = M \otimes_K K_1$ sont des corps et que
les applications des clôtures intégrales $A_1 \to B_1 \to C_1$ sont des extensions faiblement
non ramifiées d'anneaux de valuation discrète.
\end{enumerate}
\end{lemma}

\begin{proof}
Soit $e$ l'indice de ramification de $C$ sur $B$. Si $e = 1$, alors
nous avons terminé. Sinon, $e = p$ par les lemmes \ref{lemma-inequality} et \ref{lemma-ramification-index-a-power-of-p}.
Cela implique à son tour que les corps résiduels de $B$ et $C$
coïncident. Choisissons un uniformisant $\pi_C$ de $C$. Écrivons $\pi_C^p = u \pi$ pour
une certaine unité $u$ de $C$. Puisque $\pi_C^p \in L$, nous voyons que
$u \in B^*$. De plus $M = L[\pi_C]$.

\medskip\noindent Supposons qu'il existe un entier $m \geq 0$ tel que
$$
u = \sum\nolimits_{0 \leq i < m} b_i^p \pi^i + b \pi^m
$$
avec $b_i \in B$ et avec $b \in B$ un élément dont l'image dans $\kappa_B$ n'est pas une
$p$-ième puissance. Choisissez une extension $K_1/K$ comme dans le Lemme
\ref{lemma-make-degree-q-extension} avec $n = m + 2$ et notons $\pi'$ l'uniformisant de la clôture intégrale
$A_1$ de $A$ dans $K_1$ tel que $\pi = (\pi')^p + (\pi')^{np} a$ pour un certain $a \in A_1$. Soit
$B_1$ la clôture intégrale de $B$ dans $L \otimes_K K_1$. Observez que $A_1 \to B_1$ est
faiblement non ramifiée d'après le Lemme \ref{lemma-weakly-unramified-goes-up-along-totally-ramified}. Dans $B_1$ nous avons
$$
u \pi =
\left(\sum\nolimits_{0 \leq i < m} b_i (\pi')^{i + 1}\right)^p +
b (\pi')^{(m + 1)p} + (\pi')^{np} b_1
$$
pour un certain $b_1 \in B_1$ (calcul omis). Nous concluons que $M_1$ est obtenu à
partir de $L_1$ en adjoignant une $p$-ième racine de
$$
b + (\pi')^{n - m - 1} b_1
$$
Puisque le corps résiduel de $B_1$ est égal au corps résiduel de $B$, nous
voyons d'après le Lemme \ref{lemma-pre-purely-inseparable-case} que $M_1/L_1$ est de degré $p$ et que la clôture
intégrale $C_1$ de $B_1$ est faiblement non ramifiée sur $B_1$. Ainsi, nous
concluons dans ce cas.

\medskip\noindent Si aucun entier $m$ comme dans le paragraphe précédent
n'existe, alors $u$ est une $p$-ième puissance dans le complété
$\pi$-adique de $B_1$. Puisque $B$ est Nagata, ceci signifie que
$u$ est une $p$-ième puissance dans $B_1$ d'après l'Algèbre, Lemme
\ref{algebra-lemma-nagata-pth-roots}. D'où le deuxième cas de l'énoncé du lemme s'applique.
\end{proof}

\begin{lemma}
\label{lemma-cohen}
Soit $A$ un anneau local annihilé par un nombre premier $p$ dont l'idéal
maximal est nilpotent. Il existe une application d'anneau $\sigma : \kappa_A \to A$ qui est une
section de l'application résiduelle $A \to \kappa_A$. Si $A \to A'$ est un homomorphisme local
d'anneaux locaux, alors nous pouvons choisir une application d'anneau similaire
$\sigma' : \kappa_{A'} \to A'$ compatible avec $\sigma$ à condition que l'extension $\kappa_{A'}/\kappa_A$ soit
séparable.
\end{lemma}

\begin{proof}
Les extensions séparables sont formellement lisses selon Algèbre, Proposition
\ref{algebra-proposition-characterize-separable-field-extensions}. Ainsi l'existence de $\sigma$ découle du fait que $\mathbf{F}_p \to \kappa_A$ est séparable.
De même pour l'existence de $\sigma'$ compatible avec $\sigma$.
\end{proof}

\begin{lemma}
\label{lemma-pre-characteristic-p-case}
Soit $A$ un anneau de valuation discrète avec corps des fractions $K$ de
caractéristique $p > 0$. Soit $\xi \in K$. Soit $L$ une extension de $K$
obtenue en adjoignant une racine de $z^p - z = \xi$. Alors $L/K$ est galoisienne et l'une
des situations suivantes se produit
\begin{enumerate}
\item $L = K$,
\item $L/K$ est non ramifiée par rapport à $A$ de degré $p$,
\item $L/K$ est totalement ramifiée par rapport à $A$ avec un indice de
ramification $p$, et
\item la clôture intégrale $B$ de $A$ dans $L$ est un anneau de
valuation discrète, $A \subset B$ est faiblement non ramifiée, et $A \to B$ induit une
extension de corps résiduel purement inséparable de degré $p$.
\end{enumerate}
Soit $\pi$ un uniformisant de $A$. Nous avons les implications suivantes
:
\begin{enumerate}
\item[(A)] Si $\xi \in A$, alors nous sommes dans le cas (1) ou (2).
\item[(B)] Si $\xi = \pi^{-n}a$ où $n > 0$ n'est pas divisible par $p$ et $a$ est une
unité dans $A$, alors nous sommes dans le cas (3)
\item[(C)] Si $\xi = \pi^{-n} a$ où $n > 0$ est divisible par $p$ et l'image de $a$
dans $\kappa_A$ n'est pas une puissance $p$, alors nous sommes dans le cas (4).
\end{enumerate}
\end{lemma}

\begin{proof}
L'extension est galoisienne d'ordre divisant $p$ d'après la discussion de
Corps, Section \ref{fields-section-Artin-Schreier}. Il résulte immédiatement de la discussion de la Section
\ref{section-ramification} que nous sommes dans l'un des cas (1) -- (4) listés dans le lemme.

\medskip\noindent Cas (A). Ici nous voyons que $A \to A[x]/(x^p - x - \xi)$ est une extension d'anneaux
finie étale. Par conséquent, nous sommes dans les cas (1) ou (2).

\medskip\noindent Cas (B). Écrivez $\xi = \pi^{-n}a$ où $p$ ne divise pas $n$. Soit
$B \subset L$ la clôture intégrale de $A$ dans $L$. Si $C = B_\mathfrak m$ pour un certain
idéal maximal $\mathfrak m$, alors il est clair que $p \text{ord}_C(z) = -n \text{ord}_C(\pi)$. En particulier, $A \subset C$ a
un indice de ramification divisible par $p$. Il s'ensuit que c'est $p$
et que $B = C$.

\medskip\noindent Cas (C). Posons $k = n/p$. Alors nous pouvons réécrire l'équation
comme
$$
(\pi^kz)^p - \pi^{n - k} (\pi^kz) = a
$$
. Puisque $A[y]/(y^p - \pi^{n - k}y - a)$ est un anneau de valuation discrète faiblement non ramifié sur
$A$, le lemme s'ensuit.
\end{proof}

\begin{lemma}
\label{lemma-characteristic-p-case}
Soit $A \subset B \subset C$ des extensions d'anneaux de valuation discrète avec des corps des
fractions $K \subset L \subset M$. Supposons
\begin{enumerate}
\item $A \subset B$ soit faiblement non ramifié,
\item la caractéristique de $K$ est $p$,
\item $M$ est une extension galoisienne de degré $p$ de $L$, et
\item $\kappa_A = \bigcap_{n \geq 1} \kappa_B^{p^n}$.
\end{enumerate}
Alors, il existe une extension de Galois finie $K_1/K$ totalement ramifiée par
rapport à $A$ qui est une solution faible pour $A \to C$.
\end{lemma}

\begin{proof}
Puisque la caractéristique de $L$ est $p$, nous savons que $M$ est une
extension d'Artin-Schreier de $L$ (Corps, Lemme \ref{fields-lemma-Artin-Schreier}). Ainsi, nous pouvons
choisir $z \in M$, $z \not \in L$ de sorte que $\xi = z^p - z \in L$. Choisissez $n \geq 0$ de telle
sorte que $\pi^n\xi \in B$. Nous choisissons $z$ de sorte que $n$ soit minimal.
Si $n = 0$, alors $M/L$ est non ramifiée par rapport à $B$ (Lemme
\ref{lemma-pre-characteristic-p-case}) et nous avons terminé. Ainsi, nous avons $n > 0$.

\medskip\noindent L'hypothèse (4) implique que $\kappa_A$ est parfait. Ainsi, nous
pouvons choisir des morphismes d'anneaux compatibles $\overline{\sigma} : \kappa_A \to A/\pi^n A$ et $\overline{\sigma} : \kappa_B \to B/\pi^n B$ comme dans
le Lemme \ref{lemma-cohen}. Nous relevons la seconde de celles-ci à une application
d'ensembles $\sigma : \kappa_B \to B$\footnote{Si $B$ est complet, alors nous pouvons choisir
$\sigma$ comme morphisme d'anneaux. Si $A$ est également complet et $\sigma$ est
un morphisme d'anneaux, alors $\sigma$ envoie $\kappa_A$ dans $A$.}. Ensuite,
nous pouvons écrire
$$
\xi = \sum\nolimits_{i = n, \ldots, 1} \sigma(\lambda_i) \pi^{-i} + b
$$
pour certains $\lambda_i \in \kappa_B$ et $b \in B$. Soient
$$
I = \{i \in \{n, \ldots, 1\} \mid \lambda_i \in \kappa_A\}
$$
et
$$
J = \{j \in \{n, \ldots, 1\} \mid \lambda_i \not \in \kappa_A\}
$$
. Nous argumenterons par induction sur la taille de l'ensemble fini $J$.

\medskip\noindent Le cas $J = \emptyset$. Ici, pour tous $i \in \{n, \ldots, 1\}$, nous avons $\sigma(\lambda_i) = a_i + \pi^n b_i$ pour
certains $a_i \in A$ et $b_i \in B$ selon notre choix de $\sigma$. Ainsi $\xi = \pi^{-n} a + b$ pour
certains $a \in A$ et $b \in B$. Si $p | n$, alors nous écrivons $a = a_0^p + \pi a_1$ pour
certains $a_0, a_1 \in A$ (comme le corps résiduel de $A$ est parfait). Nous
calculons
$$
(z - \pi^{-n/p}a_0)^p - (z - \pi^{-n/p}a_0) =
\pi^{-(n - 1)}(a_1 + \pi^{n - 1 - n/p}a_0) + b'
$$
pour un certain $b' \in B$. Cela contredirait la minimalité de $n$. Ainsi $p$
ne divise pas $n$. Considérons l'extension de degré $p$ $K_1$ de $K$
donnée par $w^p - w = \pi^{-n}a$. D'après le Lemme \ref{lemma-pre-characteristic-p-case}, cette extension est galoisienne et
totalement ramifiée par rapport à $A$. Ainsi $L_1 = L \otimes_K K_1$ est un corps et $A_1 \subset B_1$
est faiblement non ramifiée (Lemme \ref{lemma-weakly-unramified-goes-up-along-totally-ramified}). D'après le Lemme \ref{lemma-pre-characteristic-p-case}, l'algèbre
$M_1 = M \otimes_K K_1$ est soit un produit de $p$ copies de $L_1$ (dans ce cas nous
avons terminé), soit une extension de corps de $L_1$
de degré $p$. De plus, dans le second cas, soit $C_1$ est faiblement non
ramifié sur $B_1$ (dans ce cas, nous avons terminé), soit $M_1/L_1$ est de degré
$p$, galoisien et totalement ramifié par rapport à $B_1$. Dans ce dernier
cas, l'extension $M_1/L_1$ est engendrée par l'élément $z - w$ et
$$
(z - w)^p - (z - w) = z^p - z - (w^p - w) = b
$$
avec $b \in B$ (voir ci-dessus). Ainsi, d'après le Lemme \ref{lemma-pre-characteristic-p-case}, une fois encore,
l'extension $M_1/L_1$ est non ramifiée par rapport à $B_1$ et nous concluons que
$K_1$ est une solution faible pour $A \to C$. Dorénavant, nous supposons
$J \not = \emptyset$.

\medskip\noindent Supposons qu'il existe $j', j \in J$ tels que $j' = p^r j$ pour un certain $r > 0$.
Alors nous changeons notre choix de $z$ en
$$
z' = z -
(\sigma(\lambda_j) \pi^{-j} + \sigma(\lambda_j^p) \pi^{-pj} + \ldots +
\sigma(\lambda_j^{p^{r - 1}}) \pi^{-p^{r - 1}j})
$$
. Ensuite, $\xi$ change en $\xi' = (z')^p - (z')$ comme suit
$$
\xi' =
\xi - \sigma(\lambda_j) \pi^{-j} + \sigma(\lambda_j^{p^r}) \pi^{-j'}
+ \text{un élément de }B
$$
En écrivant $\xi' = \sum\nolimits_{i = n, \ldots, 1} \sigma(\lambda'_i) \pi^{-i} + b'$ comme avant, nous constatons que $\lambda'_i = \lambda_i$ pour $i \not = j, j'$ et
$\lambda'_j = 0$. Ainsi, l'ensemble $J$ est devenu plus petit. Par induction sur la
taille de $J$, nous pouvons supposer qu'aucune telle paire $j, j'$ n'existe.
(Veuillez observer que, dans cette procédure, nous pouvons être ramenés au
cas $J = \emptyset$, traité ci-dessus.)

\medskip\noindent Pour chaque $j \in J$, écrivons $\lambda_j = \mu_j^{p^{r_j}}$ avec $r_j \geq 0$ et $\mu_j \in \kappa_B$, qui
n’est pas une $p$-ième puissance. Cela est possible par notre hypothèse (4). Soit
$j \in J$ l’indice unique tel que $j p^{-r_j}$ soit maximal. (L’indice est unique par le
résultat du paragraphe précédent.) Choisissons $r > \max(r_j + 1)$ de sorte que $j p^{r - r_j} > n$ pour
tout $j \in J$. Choisissons une extension séparable $K_1/K$ totalement ramifiée par
rapport à $A$ de degré $p^r$ telle que l’anneau de valuation discrète
correspondant $A_1 \subset K_1$ possède un uniformisant $\pi'$ avec $(\pi')^{p^r} = \pi + \pi^{n + 1}a$ pour un
certain $a \in A_1$ (Lemme \ref{lemma-make-degree-q-extension}).
Observez que $L_1 = L \otimes_K K_1$ est un corps et que $L_1/L$ est totalement ramifié par
rapport à $B$ (Lemme \ref{lemma-weakly-unramified-goes-up-along-totally-ramified}). En calculant dans la clôture intégrale $B_1$,
nous obtenons
$$
\xi = \sum\nolimits_{i \in I} \sigma(\lambda_i) (\pi')^{-i p^r} +
\sum\nolimits_{j \in J} \sigma(\mu_j)^{p^{r_j}} (\pi')^{-j p^r} + b_1
$$
pour un certain $b_1 \in B_1$. Notez que $\sigma(\lambda_i)$ pour $i \in I$ est une puissance $q$-ième
modulo $\pi^n$, c'est-à-dire modulo $(\pi')^{n p^r}$. Par conséquent, nous pouvons
réécrire ce qui précède comme
$$
\xi = \sum\nolimits_{i \in I} x_i^{p^r} (\pi')^{-i p^r} +
\sum\nolimits_{j \in J} \sigma(\mu_j)^{p^{r_j}} (\pi')^{- j p^r}
+ b_1
$$
. Comme dans le paragraphe précédent, nous changeons notre choix de $z$ en
\begin{align*}
z' & = z \\
& -
\sum\nolimits_{i \in I}
\left(x_i (\pi')^{-i} + \ldots + x_i^{p^{r - 1}} (\pi')^{-i p^{r - 1}}\right)
\\
& -
\sum\nolimits_{j \in J}
\left(
\sigma(\mu_j) (\pi')^{- j p^{r - r_j}}
+ \ldots +
\sigma(\mu_j)^{p^{r_j - 1}} (\pi')^{- j p^{r - 1}}
\right)
\end{align*}
pour obtenir
$$
(z')^p - z' =
\sum\nolimits_{i \in I} x_i (\pi')^{-i} +
\sum\nolimits_{j \in J} \sigma(\mu_j) (\pi')^{- j p^{r - r_j}} + b_1'
$$
pour un certain $b'_1 \in B_1$. Puisqu'il existe un unique $j$ tel que $j p^{r - r_j}$
soit maximal et que $j p^{r - r_j}$ est plus grand que $i \in I$ et divisible par
$p$, nous voyons que $M_1 / L_1$ relève du cas (C) du Lemme \ref{lemma-pre-characteristic-p-case}. Cela
termine la démonstration.
\end{proof}

\begin{lemma}
\label{lemma-prepare}
Soit $A$ un anneau qui contient une racine primitive $p$-ième de l'unité
$\zeta$. Posons $w = 1 - \zeta$. Alors
$$
P(z) = \frac{(1 + wz)^p - 1}{w^p} =
z^p - z + \sum\nolimits_{0 < i < p} a_i z^i
$$
est un élément de $A[z]$ et en fait $a_i \in (w)$. De plus, nous avons
$$
P(z_1 + z_2 + w z_1 z_2) = P(z_1) + P(z_2) + w^p P(z_1) P(z_2)
$$
dans l'anneau de polynômes $A[z_1, z_2]$.
\end{lemma}

\begin{proof}
Il suffit de le prouver lorsque
$$
A = \mathbf{Z}[\zeta] = \mathbf{Z}[x]/(x^{p - 1} + \ldots + x + 1)
$$
est l'anneau des entiers du corps cyclotomique. L'identité polynomiale $t^p - 1 = (t - 1)(t - \zeta) \ldots (t - \zeta^{p - 1})$
(qui est prouvée en regardant les racines de chaque côté) montre que $t^{p - 1} + \ldots + t + 1 = (t - \zeta) \ldots (t - \zeta^{p - 1})$. En
substituant $t = 1$, nous obtenons $p = (1 - \zeta)(1 - \zeta^2) \ldots (1 - \zeta^{p - 1})$. L'idéal maximal $(p, w) = (w)$ est l'unique
idéal premier de $A$ au-dessus de $p$ (car les corps de caractéristique
$p$ n'ont pas de racines $p$-ièmes non triviales de $1$). Il s'ensuit que
$p = u w^{p - 1}$ pour une certaine unité $u$. Cela implique que
$$
a_i = \frac{1}{p} {p \choose i} u w^{i - 1}
$$
pour $p > i > 1$ et $- 1 + a_1 = pw/w^p = u$. Puisque $P(-1) = 0$, nous voyons que $0 = (-1)^p - u$ modulo $(w)$.
Ainsi $a_1 \in (w)$ et la preuve de la première partie est terminée. La deuxième partie
découle d'un calcul direct que nous omettons.
\end{proof}

\begin{lemma}
\label{lemma-extension-defined-by-nice-polynial}
Soit $A$ un anneau de valuation discrète de caractéristique mixte $(0, p)$ qui
contient une racine primitive $p$-ième de $1$. Soit $P(t) \in A[t]$ le polynôme
du Lemme \ref{lemma-prepare}. Soit $\xi \in K$. Soit $L$ une extension de $K$ obtenue
en adjoignant une racine de $P(z) = \xi$. Alors $L/K$ est galoisienne et l'une des
situations suivantes se produit
\begin{enumerate}
\item $L = K$,
\item $L/K$ est non ramifiée par rapport à $A$ de degré $p$,
\item $L/K$ est totalement ramifiée par rapport à $A$ avec un indice de
ramification $p$, et
\item la clôture intégrale $B$ de $A$ dans $L$ est un anneau de
valuation discrète, $A \subset B$ est faiblement non ramifié, et $A \to B$ induit une
extension de corps résiduel purement inséparable de degré $p$.
\end{enumerate}
Soit $\pi$ un uniformisant de $A$. Nous avons les implications suivantes :
\begin{enumerate}
\item[(A)] Si $\xi \in A$, alors nous sommes dans le cas (1) ou (2).
\item[(B)] Si $\xi = \pi^{-n}a$ où $n > 0$ n'est pas divisible par $p$ et $a$ est
une unité dans $A$, alors nous sommes dans le cas (3)
\item[(C)] Si $\xi = \pi^{-n} a$ où $n > 0$ est divisible par $p$ et l'image de $a$
dans $\kappa_A$ n'est pas une $p$-ième puissance, alors nous sommes dans le cas
(4).
\end{enumerate}
\end{lemma}

\begin{proof}
Adjoindre une racine de $P(z) = \xi$ revient au même que d'adjoindre une racine de
$y^p = w^p(1 + \xi)$. Puisque $K$ contient une racine primitive $p$-ième de $1$,
l'extension est galoisienne d'ordre divisant $p$ d'après la discussion dans
Corps, Section \ref{fields-section-Kummer}. Il découle immédiatement de la discussion dans la section
\ref{section-ramification} que nous sommes dans l'un des cas (1) -- (4) énumérés dans le lemme.

\medskip\noindent Cas (A). Ici, nous voyons que $A \to A[x]/(P(x) - \xi)$ est une extension d'anneaux
finie étale. Par conséquent, nous sommes dans les cas (1) ou (2).

Cas \medskip\noindent (B). Écrivez $\xi = \pi^{-n}a$ où $p$ ne divise pas $n$. Soit
$B \subset L$ la clôture intégrale de $A$ dans $L$. Si $C = B_\mathfrak m$ pour un certain
idéal maximal $\mathfrak m$, alors il est clair que $p \text{ord}_C(z) = -n \text{ord}_C(\pi)$. En particulier, $A \subset C$ a
un indice de ramification divisible par $p$. Il s'ensuit que c'est $p$
et que $B = C$.

\medskip\noindent Cas (C). Posons $k = n/p$. Alors nous pouvons réécrire l'équation
comme
$$
(\pi^kz)^p - \pi^{n - k} (\pi^kz) + \sum a_i \pi^{n - ik} (\pi^kz)^i = a
$$
Comme $A[y]/(y^p - \pi^{n - k}y - \sum  a_i \pi^{n - ik} y^i - a)$ est un anneau de valuation discrète faiblement non ramifié sur
$A$, le lemme s'ensuit.
\end{proof}

\noindent
Soit $A$ un anneau de valuation discrète de caractéristique mixte $(0, p)$
contenant une racine primitive $p$-ième de $1$. Soient $w \in A$ et $P(t) \in A[t]$
comme dans le Lemme \ref{lemma-prepare}. Soit $L$ une extension finie de $K$. On dit que
$L/K$ est une extension de degré $p$ {\it de niveau fini} si $L$
est une extension de degré $p$ de $K$ obtenue en adjoignant une racine de
l’équation $P(z) = \xi$ où $\xi \in K$ est un élément avec $w^p \xi \in \mathfrak m_A$.

\medskip\noindent Cette définition est pertinente pour la discussion de cette
section en raison du lemme simple suivant.

\begin{lemma}
\label{lemma-make-finite-level}
Soient $A \subset B \subset C$ des extensions d’anneaux de valuation discrète avec des corps des
fractions $K \subset L \subset M$. Supposons que
\begin{enumerate}
\item $A$ soit de caractéristique mixte $(0, p)$,
\item $A \subset B$ soit faiblement non ramifié,
\item $B$ contient une racine primitive $p$-ième de $1$, et
\item $M/L$ est galoisienne de degré $p$.
\end{enumerate}
Il existe alors une extension galoisienne finie $K_1/K$ totalement ramifiée par
rapport à $A$ qui est soit une solution faible pour $A \to C$ soit telle que
$M_1/L_1$ est une extension de degré $p$ de niveau fini.
\end{lemma}

\begin{proof}
Soit $\pi \in A$ un uniformisant. D'après la théorie de Kummer (Corps, Lemme
\ref{fields-lemma-Kummer}) $M$ est obtenu à partir de $L$ en adjoignant une racine de
l'équation $y^p = b$ pour un certain $b \in L$.

\medskip\noindent Si $\text{ord}_B(b)$ est premier à $p$, alors nous choisissons une
extension séparable de degré $p$, $K_1/K$, totalement ramifiée par rapport à
$A$ (par exemple en utilisant le Lemme \ref{lemma-make-degree-q-extension}). Soit $A_1$ la clôture
intégrale de $A$ dans $K_1$. D'après le Lemme \ref{lemma-weakly-unramified-goes-up-along-totally-ramified}, la clôture intégrale
$B_1$ de $B$ dans $L_1 = L \otimes_K K_1$ est un anneau de valuation discrète faiblement
non ramifié sur $A_1$. Si $K_1/K$ n'est pas une solution faible pour $A \to C$,
alors la clôture intégrale $C_1$ de $C$ dans $M_1 = M \otimes_K K_1$ est un anneau de
valuation discrète et $B_1 \to C_1$ a un indice de ramification $p$.
Dans ce cas, le corps $M_1$ est obtenu à partir de $L_1$ en adjoignant la
$p$‑ième racine de $b$ avec $\text{ord}_{B_1}(b)$ divisible par $p$. En remplaçant
$A$ par $A_1$, etc., nous pouvons supposer que $b = \pi^n u$ où $u \in B$ est une
unité et $n$ est divisible par $p$. Bien sûr, dans ce cas, l'extension
$M$ est obtenue à partir de $L$ en adjoignant la $p$‑ième racine
d'une unité.

\medskip\noindent Supposons que $M$ soit obtenu à partir de $L$ en
adjoignant la racine de $y^p = u$ pour une certaine unité $u$ de $B$. Si
la classe résiduelle de $u$ dans $\kappa_B$ n'est pas une puissance
$p$‑ième, alors $B \subset C$ est faiblement non ramifiée (Lemme \ref{lemma-pre-purely-inseparable-case}) et nous
avons terminé. Sinon, nous pouvons remplacer notre choix de $y$ par $y/v$
où $v^p$ et $u$ ont la même image dans $\kappa_B$. Après un tel
remplacement, nous avons
$$
y^p = 1 + \pi b
$$
pour un certain $b \in B$. Ensuite, nous voyons que $P(z) = \pi b/ w^p$ où $z = (y - 1)/w$. Ainsi, nous
voyons que l'extension est une extension de degré $p$ de niveau fini avec
$\xi = \pi b / w^p$.
\end{proof}

\noindent
Soit $A$ un anneau de valuation discrète de caractéristique mixte $(0, p)$
contenant une racine primitive $p$-ième de $1$. Soient $w \in A$ et $P(t) \in A[t]$
comme dans le Lemme \ref{lemma-prepare}. Soit $L$ une extension de degré $p$ de
$K$ de niveau fini. Choisissez $z \in L$ générant $L$ sur $K$ avec
$\xi = P(z) \in K$. Choisissons un uniformisant $\pi$ de $A$ et écrivons $w = u \pi^{e_1}$
pour un certain entier $e_1 = \text{ord}_A(w)$ et une unité $u \in A$. Enfin, prenez $n \geq 0$ tel
que
$$
\pi^n \xi \in A
$$
Le {\it niveau} de $L/K$ est la plus petite valeur de la quantité $n/e_1$
prise sur tous les $z$ générant $L/K$ avec $\xi = P(z) \in K$.

\medskip\noindent Nous faisons quelques remarques. Puisque l'extension est de
niveau fini, nous savons que nous pouvons choisir $z$ de telle sorte que
$n < pe_1$. Ainsi, le niveau est un nombre rationnel contenu dans $[0, p)$. Si le
niveau est zéro, alors $L/K$ est non ramifiée par rapport à $A$ selon le
Lemme \ref{lemma-extension-defined-by-nice-polynial}. Notre objectif suivant est de réduire le niveau.

\begin{lemma}
\label{lemma-lowering-the-level}
Soient $A \subset B \subset C$ des extensions d'anneaux de valuation discrète avec des corps des
fractions $K \subset L \subset M$. Supposons
\begin{enumerate}
\item $A$ soit de caractéristique mixte $(0, p)$,
\item $A \subset B$ soit faiblement non ramifiée,
\item $B$ contient une racine primitive $p$-ième de $1$,
\item $M/L$ est une extension de degré $p$ de niveau fini $l > 0$,
\item $\kappa_A = \bigcap_{n \geq 1} \kappa_B^{p^n}$.
\end{enumerate}
Alors il existe une extension finie séparable $K_1$ de $K$ totalement
ramifiée par rapport à $A$ telle que $K_1$ soit une solution faible
pour $A \to C$, ou que l'extension $M_1/L_1$ soit une extension de degré $p$ de
niveau fini $\leq \max(0, l - 1, 2l - p)$.
\end{lemma}

\begin{proof}
Soit $\pi \in A$ un uniformisant. Soient $w \in B$ et $P \in B[t]$ comme dans le Lemme
\ref{lemma-prepare} (pour $B$). Posons $e_1 = \text{ord}_B(w)$, de sorte que $w$ et $\pi^{e_1}$ sont
associés dans $B$. Choisissons $z \in M$ générant $M$ sur $L$ avec
$\xi = P(z) \in K$, ainsi qu'un entier $n$ tel que $\pi^n\xi \in B$, comme dans la définition du niveau de
$M$ sur $L$, c'est-à-dire $l = n/e_1$.

\medskip\noindent La démonstration de ce lemme est complètement similaire à la
démonstration du Lemme \ref{lemma-characteristic-p-case}. Pour expliquer ce qui se passe, il faut observer
que
\begin{equation}
\label{equation-first-congruence}
P(z) \equiv z^p - z \bmod \pi^{-n + e_1}B
\end{equation}
pour tout $z \in L$ tel que $\pi^{-n} P(z) \in B$ (utilisez que $z$ a une valuation au pire
$-n/p$ et la forme du polynôme $P$). De plus, nous avons
\begin{equation}
\label{equation-second-congruence}
\xi_1 + \xi_2 + w^p \xi_1 \xi_2 \equiv \xi_1 + \xi_2 \bmod \pi^{-2n + pe_1}B 
\end{equation}
pour $\xi_1, \xi_2 \in \pi^{-n}B$. Enfin, observez que $n - e_1 = (l - 1)/e_1$ et $-2n + pe_1 = -(2l - p)e_1$. Écrivez $m = n - e_1 \max(0, l - 1, 2l - p)$. Ce qui
précède montre qu'en faisant des calculs dans $\pi^{-n}B / \pi^{-n + m}B$, le polynôme $P$ se
comporte exactement comme le polynôme $z^p -  z$. Cela explique pourquoi le lemme
est vrai mais nous donnons également les détails ci-dessous.

\medskip\noindent Hypothèse (4) implique que $\kappa_A$ est parfait. Observez que
$m \leq e_1$ et donc $A/\pi^m$ est anéanti par $w$ et donc $p$. Ainsi, nous
pouvons choisir des morphismes d'anneaux compatibles $\overline{\sigma} : \kappa_A \to A/\pi^mA$ et $\overline{\sigma} : \kappa_B \to B/\pi^mB$ comme
dans le Lemme \ref{lemma-cohen}. Nous relevons la seconde de ces applications en une
application d’ensembles $\sigma : \kappa_B \to B$. Nous pouvons alors écrire
$$
\xi =
\sum\nolimits_{i = n, \ldots, n - m + 1} \sigma(\lambda_i) \pi^{-i} +
\pi^{-n + m)} b
$$
pour certains $\lambda_i \in \kappa_B$ et $b \in B$. Soient
$$
I = \{i \in \{n, \ldots, n - m + 1\} \mid \lambda_i \in \kappa_A\}
$$
et
$$
J = \{j \in \{n, \ldots, n - m + 1\} \mid \lambda_i \not \in \kappa_A\}
$$
Nous argumenterons par induction sur la taille de l’ensemble fini $J$.

\medskip\noindent Le cas $J = \emptyset$. Ici, pour tout $i \in \{n, \ldots, n - m + 1\}$ nous avons $\sigma(\lambda_i) = a_i + \pi^{n - m}b_i$ pour
certains $a_i \in A$ et $b_i \in B$ selon notre choix de $\overline{\sigma}$. Ainsi $\xi = \pi^{-n} a + \pi^{-n + m} b$ pour
certains $a \in A$ et $b \in B$. Si $p | n$, alors nous écrivons $a = a_0^p + \pi a_1$ pour
certains $a_0, a_1 \in A$ (car le corps résiduel de $A$ est parfait). Posons
$z_1 = - \pi^{-n/p} a_0$. Notez que $P(z_1) \in \pi^{-n}B$ et que $z + z_1 + w z z_1$ est un élément générateur de $M$
sur $L$ (notez que $wz_1 \not = -1$ car $n < pe_1$). De plus, d'après le Lemme \ref{lemma-prepare}
nous avons
$$
P(z + z_1 + w z z_1) = P(z) + P(z_1) + w^p P(z) P(z_1) \in K
$$
et d'après les équations (\ref{equation-first-congruence}) et (\ref{equation-second-congruence}) nous avons
$$
P(z) + P(z_1) + w^p P(z) P(z_1)
\equiv
\xi + z_1^p - z_1 \bmod \pi^{-n + m}B
$$
pour un certain $b' \in B$. Cela contredit la minimalité de $n$ ! Ainsi $p$ ne
divise pas $n$. Considérons l'extension de degré $p$ $K_1$ de $K$
donnée par $P(y) = -\pi^{-n}a$. Selon le Lemme \ref{lemma-extension-defined-by-nice-polynial}, cette extension est séparable et
totalement ramifiée par rapport à $A$. Ainsi $L_1 = L \otimes_K K_1$ est un corps et $A_1 \subset B_1$
est faiblement non ramifié (Lemme \ref{lemma-weakly-unramified-goes-up-along-totally-ramified}). Selon le Lemme \ref{lemma-extension-defined-by-nice-polynial}, l'anneau
$M_1 = M \otimes_K K_1$ est soit un produit de $p$ copies de $L_1$ (dans ce cas, nous
avons terminé), soit une extension de corps de $L_1$
de degré $p$. De plus, dans le deuxième cas, soit $C_1$ est faiblement non
ramifiée sur $B_1$ (dans ce cas nous avons terminé), soit $M_1/L_1$ est de degré
$p$, galoisienne et totalement ramifiée par rapport à $B_1$. Dans ce dernier cas,
l'extension $M_1/L_1$ est engendrée par l'élément $z + y + wzy$ et nous voyons que $P(z + y + wzy) \in L_1$
et
\begin{align*}
P(z + y + wzy)
& = P(z) + P(y) + w^p P(z) P(y) \\
& \equiv
\xi - \pi^{-n}a \bmod \pi^{-n + m}B_1 \\
& \equiv
0 \bmod \pi^{-n + m}B_1
\end{align*}
exactement de la même manière que ci-dessus. Par notre choix de $m$, cela
signifie exactement que $M_1/L_1$ a un niveau au plus $\max(0, l - 1, 2l - p)$. À partir de
maintenant, nous supposons que $J \not = \emptyset$.

\medskip\noindent Supposons qu'il existe $j', j \in J$ tels que $j' = p^r j$ pour un certain $r > 0$.
Alors nous définissons
$$
z_1 = - \sigma(\lambda_j) \pi^{-j} - \sigma(\lambda_j^p) \pi^{-pj} -
\ldots - \sigma(\lambda_j^{p^{r - 1}}) \pi^{-p^{r - 1}j}
$$
et nous changeons $z$ en $z' = z + z_1 + wzz_1$. Observez que $z' \in M$ engendre $M$
sur $L$ et que nous avons $\xi' = P(z') = P(z) + P(z_1) + wP(z)P(z_1) \in L$ avec
$$
\xi' \equiv
\xi - \sigma(\lambda_j) \pi^{-j} + \sigma(\lambda_j^{p^r}) \pi^{-j'}
\bmod \pi^{-n + m}B
$$
en utilisant les équations (\ref{equation-first-congruence}) et (\ref{equation-second-congruence}) comme ci-dessus. En écrivant
$$
\xi' = \sum\nolimits_{i = n, \ldots, n - m + 1} \sigma(\lambda'_i) \pi^{-i}
+ \pi^{-n + m}b'
$$
comme auparavant, nous constatons que $\lambda'_i = \lambda_i$ pour $i \not = j, j'$ et $\lambda'_j = 0$. Ainsi,
l'ensemble $J$ est devenu plus petit. Par induction sur la taille de $J$,
nous pouvons supposer qu'il n'existe pas de paire $j, j'$ de $J$ telle que
$j'/j$ soit une puissance de $p$. (Veuillez noter que, dans cette procédure,
nous pouvons être ramenés au cas $J = \emptyset$, traité ci-dessus.)

\medskip\noindent Pour chaque $j \in J$, écrivons $\lambda_j = \mu_j^{p^{r_j}}$ avec $r_j \geq 0$ et $\mu_j \in \kappa_B$,
qui n'est pas une $p$-ième puissance. Cela est possible par notre hypothèse (4).
Soit $j \in J$ l'indice unique tel que $j p^{-r_j}$ soit maximal. (L'indice est unique
d'après le résultat du paragraphe précédent.) Choisissons $r > \max(r_j + 1)$ tel que
$j p^{r - r_j} > n$ pour tout $j \in J$. Soit $K_1/K$ l'extension de degré $p^r$, totalement
ramifiée par rapport à $A$, définie par $(\pi')^{p^r} = \pi$. Observez que $\pi'$ est
l'uniformisant de l'anneau de valuation discrète correspondant $A_1 \subset K_1$.
Observez que $L_1 = L \otimes_K K_1$ est un corps et $L_1/L$
est totalement ramifiée par rapport à $B$ (Lemme \ref{lemma-weakly-unramified-goes-up-along-totally-ramified}). En calculant dans
la clôture intégrale $B_1$, nous obtenons
$$
\xi = \sum\nolimits_{i \in I} \sigma(\lambda_i) (\pi')^{-i p^r} +
\sum\nolimits_{j \in J} \sigma(\mu_j)^{p^{r_j}} (\pi')^{-j p^r} +
\pi^{-n + m} b_1
$$
pour un certain $b_1 \in B_1$. Notez que $\sigma(\lambda_i)$ pour $i \in I$ est une puissance $q$-ième
modulo $\pi^m$, c’est-à-dire modulo $(\pi')^{m p^r}$. Par conséquent, nous pouvons réécrire
ce qui précède comme
$$
\xi = \sum\nolimits_{i \in I} x_i^{p^r} (\pi')^{-i p^r} +
\sum\nolimits_{j \in J} \sigma(\mu_j)^{p^{r_j}} (\pi')^{- j p^r}
+ \pi^{-n + m}b_1
$$
. De manière similaire à notre choix dans le paragraphe précédent nous définissons
\begin{align*}
z_1 & - \sum\nolimits_{i \in I}
\left(x_i (\pi')^{-i} + \ldots + x_i^{p^{r - 1}} (\pi')^{-i p^{r - 1}}\right)
\\
& - \sum\nolimits_{j \in J}
\left(
\sigma(\mu_j) (\pi')^{- j p^{r - r_j}}
+ \ldots +
\sigma(\mu_j)^{p^{r_j - 1}} (\pi')^{- j p^{r - 1}}
\right)
\end{align*}
et nous changeons notre choix de $z$ en $z' = z + z_1 + wzz_1$. Alors $z'$ engendre
$M_1$ sur $L_1$ et $\xi' = P(z') = P(z) + P(z_1) + w^p P(z) P(z_1) \in L_1$ et un calcul montre que
$$
\xi' \equiv
\sum\nolimits_{i \in I} x_i (\pi')^{-i} +
\sum\nolimits_{j \in J} \sigma(\mu_j) (\pi')^{- j p^{r - r_j}} +
(\pi')^{(-n + m)p^r}b'_1
$$
pour un certain $b'_1 \in B_1$. Il existe un unique $j$ tel que $j p^{r - r_j}$ soit maximal
et $j p^{r - r_j}$ soit plus grand que $i \in I$. Si $j p^{r - r_j} \leq (n - m)p^r$ alors le niveau de
l’extension $M_1/L_1$ est inférieur à $\max(0, l - 1, 2l - p)$. Sinon, comme $p$ divise
$j p^{r - r_j}$, nous voyons que $M_1 / L_1$ tombe dans le cas (C) du Lemme \ref{lemma-extension-defined-by-nice-polynial}. Cela
termine la démonstration.
\end{proof}

\begin{lemma}
\label{lemma-special-case}
Soit $A \subset B \subset C$ des extensions d'anneaux de valuation discrète avec des corps des
fractions $K \subset L \subset M$. Supposons que
\begin{enumerate}
\item le corps résiduel $k$ de $A$ est algébriquement clos de
caractéristique $p > 0$,
\item $A$ et $B$ sont complets,
\item $A \to B$ est faiblement non ramifié,
\item $M$ est une extension finie de $L$,
\item $k = \bigcap\nolimits_{n \geq 1} \kappa_B^{p^n}$
\end{enumerate}
Alors il existe une extension finie $K_1/K$ qui est une solution faible pour
$A \to C$.
\end{lemma}

\begin{proof}
Soit $M'$ une extension finie quelconque de $L$ et considérons la clôture intégrale
$C'$ de $B$ dans $M'$. Alors $C'$ est fini sur $B$ puisque
$B$ est un anneau de Nagata d'après Algèbre, Lemme \ref{algebra-lemma-Noetherian-complete-local-Nagata}. De plus, $C'$ est un
anneau de valuation discrète, voir la discussion dans la Remarque \ref{remark-construction}. En
outre, $C'$ est complet en tant que $B$-module, donc complet en tant
qu'anneau de valuation discrète, voir Algèbre, Section \ref{algebra-section-completion}. Il s'ensuit en
particulier que $C$ est la clôture intégrale de $B$ dans $M$ (par
définition des anneaux de valuation comme maximaux pour la relation de
domination).

\medskip\noindent Soit $M \subset M'$ une extension finie et soit $C' \subset M'$ la clôture
intégrale de $B$ comme ci-dessus. D'après le Lemme \ref{lemma-solutions-go-down}, il suffit de
prouver le résultat pour $A \to B \to C'$. Par conséquent, nous pouvons supposer que
$M/L$ est normal, voir Corps, Lemme \ref{fields-lemma-normal-closure}.

\medskip\noindent Si $M / L$ est normal, nous pouvons trouver une chaîne
d'extensions finies
$$
L = L^0 \subset L^1 \subset L^2 \subset \ldots \subset L^r = M
$$
telle que chaque extension $L^{j + 1}/L^j$ soit de l'un des types suivants :
\begin{enumerate}
\item[(a)] purement inséparable de degré $p$,
\item[(b)] totalement ramifiée par rapport à $B^j$ et galoisienne de degré
$p$,
\item[(c)] totalement ramifiée par rapport à $B^j$ et galoisienne cyclique
d'ordre premier à $p$,
\item[(d)] galoisienne et non ramifiée par rapport à $B^j$.
\end{enumerate}
Ici $B^j$ est la clôture intégrale de $B$ dans $L^j$. À savoir, puisque
$M/L$ est normal, nous pouvons l'écrire comme un compositum d'une extension
galoisienne et d'une extension purement inséparable (Corps, Lemme \ref{fields-lemma-normal-case}). Pour
l'extension purement inséparable, l'existence de la filtration est claire. Dans le
cas galoisien, notez que $G$ est « le » groupe de décomposition et soit
$I \subset G$ le groupe d'inertie. Alors, d'une part, $I$ est résoluble par le Lemme
\ref{lemma-galois-inertia} et d'autre part, l'extension $M^I/L$ est non ramifiée par rapport à
$B$ par le Lemme \ref{lemma-inertial-invariants-unramified}. Cela prouve que nous avons une filtration comme
indiqué.

\medskip\noindent Nous allons raisonner par induction sur l'entier $r$.
Supposons que nous puissions trouver une extension finie $K_1/K$ qui est une
solution faible pour $A \to B^1$ où $B^1$ est la clôture intégrale de $B$ dans
$L^1$. Soit $K'_1$ la clôture normale de $K_1/K$ (Corps, Lemme \ref{fields-lemma-normal-closure}).
Puisque $A$ est complet et le corps résiduel de $A$ est algébriquement
clos, nous voyons que $K'_1/K_1$ est séparable et totalement ramifiée par rapport à
$A_1$ (certains détails omis). Ainsi $K'_1/K$ est également une solution faible
pour $A \to B^1$ selon le Lemme \ref{lemma-weakly-unramified-goes-up-along-totally-ramified}. En d'autres termes, nous pouvons donc
supposer que $K_1$ est une extension normale de $K$.
Ayant fait cela, nous considérons la suite
$$
L^0_1 = (L^0 \otimes_K K_1)_{red} \subset
L^1_1 = (L^1 \otimes_K K_1)_{red} \subset \ldots \subset
L^r_1 = (L^r \otimes_K K_1)_{red}
$$
et les clôtures intégrales correspondantes $B^i_1$. Notez que $C_1 = B^r_1$ est un
produit d'anneaux de valuation discrète qui sont permutés de façon transitive par
$G = \text{Aut}(K_1/K)$ d'après le Lemme \ref{lemma-galois-relative}. En particulier, toutes les extensions d'anneaux
de valuation discrète $A_1 \to (C_1)_\mathfrak m$ sont isomorphes, et une solution faible pour l'une
sera une solution faible pour toutes. Nous pouvons appliquer l'hypothèse de
récurrence à la séquence
$$
A_1 \to (B^1_1)_{B^1_1 \cap \mathfrak m} \to
(B^2_1)_{B^2_1 \cap \mathfrak m} \to
\ldots \to
(B^r_1)_{B^r_1 \cap \mathfrak m} =
(C_1)_\mathfrak m
$$
pour obtenir une solution faible $K_2/K_1$ pour $A_1 \to (C_1)_\mathfrak m$. L'extension $K_2/K$ sera alors
une solution faible pour $A \to C$ d'après ce que nous avons dit précédemment. Notez
que l'hypothèse d'induction s'applique : l'homomorphisme d'anneaux $A_1 \to (B^1_1)_{B^1_1 \cap \mathfrak m}$ est
faiblement non ramifié selon notre choix de $K_1$ et la suite des extensions de
corps des fractions possède chacune encore l'une des propriétés (a), (b), (c) ou
(d) listées ci-dessus. De plus, observez que pour toute extension finie $\kappa_B \subset \kappa$
nous avons encore $k = \bigcap \kappa^{p^n}$.

\medskip\noindent Ainsi, tout se résume à trouver une solution faible pour $A \subset C$
lorsque l'extension de corps $M/L$ satisfait l'une des propriétés (a), (b), (c)
ou (d).

\medskip\noindent Cas (d). Ce cas est trivial car ici $B \to C$ est déjà non
ramifié.

\medskip\noindent Cas (c). Supposons que $M/L$ soit cyclique d’ordre $n$ premier
à $p$. Puisque $M/L$ est totalement ramifié par rapport à $B$, on voit
que l’indice de ramification de $B \subset C$ est $n$ et donc l’indice de
ramification de $A \subset C$ est également $n$. Choisissons un uniformisant $\pi \in A$
et posons $K_1 = K[\pi^{1/n}]$. Alors $K_1/K$ est une solution pour $A \subset C$ par le lemme
d’Abhyankar (Lemme \ref{lemma-abhyankar}).

Cas \medskip\noindent (b). Nous divisons ce cas en cas de caractéristique mixte et
cas équicaractéristique. Dans le cas équicaractéristique,
ceci est le Lemme \ref{lemma-characteristic-p-case}. Dans le cas de caractéristique mixte, nous remplaçons
d'abord $K$ par une extension finie pour arriver à la situation où $M/L$ est
une extension de degré $p$ de niveau fini en utilisant le Lemme \ref{lemma-make-finite-level}.
Ensuite, le niveau est un nombre rationnel $l \in [0, p)$, voir discussion précédant le
Lemme \ref{lemma-lowering-the-level}. Si le niveau est $0$, alors $B \to C$ est faiblement non ramifié
et nous avons terminé. Sinon, nous pouvons remplacer le corps $K$ par une
extension finie pour obtenir une nouvelle
situation avec le niveau $l' \leq \max(0, l - 1, 2l - p)$ par le Lemme \ref{lemma-lowering-the-level}. Si $l = p - \epsilon$ pour $\epsilon < 1$
alors nous voyons que $l' \leq p - 2\epsilon$. Ainsi, après un nombre fini de remplacements, nous
obtenons un cas avec le niveau $\leq p - 1$. Ensuite, après au plus $p - 1$
remplacements supplémentaires de ce type, nous atteignons la situation où le
niveau est zéro.

\medskip\noindent Le cas (a) est le Lemme \ref{lemma-purely-inseparable-case}. C'est le seul cas où nous avons
éventuellement besoin d'une extension purement inséparable de $K$, à savoir,
dans le cas (2) de l'énoncé du lemme, nous réussissons en adjoignant une racine
$p$-ième de l'élément $\pi$. Cela termine la démonstration du lemme.
\end{proof}

\noindent
À ce stade, nous avons rassemblé tous les lemmes dont nous avons besoin pour
prouver le résultat principal de cette section.

\begin{theorem}[Epp]
\label{theorem-epp}
Soit $A \subset B$ une extension d'anneaux à valuation discrète avec des corps des
fractions $K \subset L$. Si la caractéristique de $\kappa_A$ est $p > 0$, supposons que
chaque élément de
$$
\bigcap\nolimits_{n \geq 1} \kappa_B^{p^n}
$$
est algébrique et séparable sur $\kappa_A$. Alors il existe une extension finie
$K_1/K$ qui est une solution faible pour $A \to B$ au sens de la Définition
\ref{definition-solution}.
\end{theorem}

\begin{proof}
Si la caractéristique de $\kappa_A$ est nulle, ou si la caractéristique du corps
résiduel vaut $p$, que l'indice de ramification est premier à $p$ et que
l'extension du corps résiduel est séparable, alors cela découle du lemme
d'Abhyankar (Lemme \ref{lemma-abhyankar}). En effet, supposons que l'indice de ramification soit
$e$. Choisissons un uniformisant $\pi \in A$. Soit $K_1/K$ l'extension obtenue en
adjoignant une racine $e$-ième de $\pi$. D'après le Lemme \ref{lemma-pull-root-uniformizer}, nous voyons
que la clôture intégrale $A_1$ de $A$ dans $K_1$ est un anneau de
valuation discrète avec l'indice de ramification considéré sur $A$. Ainsi $A_1 \to (B_1)_\mathfrak m$ est
formellement lisse dans la topologie $\mathfrak m$-adique pour tous les idéaux maximaux
$\mathfrak m$ de $B_1$
d'après le Lemme \ref{lemma-abhyankar} et a fortiori ce sont des extensions faiblement non
ramifiées d'anneaux de valuation discrète.

\medskip\noindent Dorénavant, supposons que $p$ soit premier et que
$\kappa_A$ soit de caractéristique $p$. Nous appliquons d’abord le Lemme \ref{lemma-solution-after-strict-henselization}
pour réduire au cas où $A$ et $B$ ont des corps résiduels séparablement
fermés. Puisque $\kappa_A$ et $\kappa_B$ sont remplacés par leurs fermetures algébriques
séparables par cette procédure, on constate que l’on obtient
$$
\kappa_A \supset \bigcap\nolimits_{n \geq 1} \kappa_B^{p^n}
$$
à partir de la condition du théorème.

\medskip\noindent Soit $\pi \in A$ un uniformisant. Soient $A^\wedge$ et $B^\wedge$
les complétions de $A$ et $B$. Nous avons un diagramme commutatif
$$
\xymatrix{
B \ar[r] & B^\wedge \\
A \ar[u] \ar[r] & A^\wedge \ar[u]
}
$$
des extensions d'anneaux de valuation discrète. Soit $K^\wedge$ le corps des
fractions de $A^\wedge$. Supposons que nous puissions trouver une extension finie
$M/K^\wedge$ qui est (a) une solution faible pour $A^\wedge \to B^\wedge$ et (b) un compositum d'une
extension séparable et d'une extension obtenue en adjoignant une racine
$p$-ième de $\pi$. Alors, d'après le Lemme \ref{lemma-approximate-separable-extension}, nous pouvons trouver une
extension finie $K_1/K$ telle que $K^\wedge \otimes_K K_1 = M$. Soient $A_1$ et
$A_1^\wedge$ les clôtures intégrales respectives de $A$ et de $A^\wedge$ dans
$K_1$ et $M$. Puisque $A \to A^\wedge$ est formellement lisse dans la
topologie $\mathfrak m^\wedge$-adique (Lemme \ref{lemma-extension-dvrs-formally-smooth})
nous voyons que $A_1 \to A_1^\wedge$ est formellement lisse dans la topologie $\mathfrak m_1^\wedge$-adique
(par le Lemme \ref{lemma-formally-smooth-goes-up}, les anneaux $A_1$ et $A_1^\wedge$ étant des anneaux de valuation discrète
d'après la discussion de la Remarque \ref{remark-construction}). Nous concluons de la partie (2) du Lemme \ref{lemma-solutions-go-down}
que $K_1/K$ est une solution faible pour $A \to B^\wedge$. En appliquant la
partie (1) du Lemme \ref{lemma-solutions-go-down}, nous voyons que $K_1/K$ est une solution faible pour
$A \to B$.

\medskip\noindent Ainsi, nous pouvons supposer que $A$ et $B$ sont des
anneaux de valuation discrète complets avec des corps résiduels séparablement clos
de caractéristique $p$ et avec $\kappa_A \supset \bigcap\nolimits_{n \geq 1} \kappa_B^{p^n}$. Nous disposons également d’un
uniformisant $\pi \in A$ et nous devons trouver une solution faible pour $A \to B$ qui
est un compositum d’une extension séparable et d’un corps obtenu en adjoignant une
racine $p$-ième de $\pi$. Notez que la deuxième condition est automatique
si $A$ est de caractéristique mixte.

\medskip\noindent Posons $k = \bigcap\nolimits_{n \geq 1} \kappa_B^{p^n}$. Observez que $k$ est un corps algébriquement
clos de caractéristique $p$. Si $A$ est de caractéristique mixte, soit
$\Lambda$ un anneau de Cohen pour $k$ et, dans le cas équicaractéristique,
définissons $\Lambda = k[[t]]$. Nous pouvons choisir un morphisme d'anneaux $\Lambda \to A$ qui
envoie $t$ sur $\pi$ dans le cas équicaractéristique. Dans le cas
équicaractéristique, cela découle du théorème de structure de Cohen (Algèbre,
Théorème \ref{algebra-theorem-cohen-structure-theorem}) et dans le cas de caractéristique mixte, cela découle du fait
que $\mathbf{Z}_p \to \Lambda$ est formellement lisse pour la topologie adique (Lemmes \ref{lemma-extension-dvrs-formally-smooth} et
\ref{lemma-lift-continuous}). En appliquant le Lemme \ref{lemma-solutions-go-down}, nous voyons qu'il suffit de prouver
l'existence d'une solution faible pour $\Lambda \to B$, qui dans le cas
équicaractéristique $p$ est un compositum d'une extension séparable et d'un
corps obtenu en adjoignant une racine $p$-ième de $t$. Cependant, puisque
$\Lambda = k[[t]]$ dans le cas équicaractéristique et que toute extension de $k((t))$ est un tel
compositum, nous pouvons maintenant abandonner cette exigence !

\medskip\noindent Ainsi nous arrivons à la situation où $A$ et $B$ sont
complets, le corps résiduel $k$ de $A$ est algébriquement clos de
caractéristique $p > 0$, nous avons $k = \bigcap \kappa_B^{p^n}$, et dans le cas de caractéristique
mixte, $p$ est un uniformisant de $A$ (c'est-à-dire que $A$ est un
anneau de Cohen pour $k$). Si $A$ est de caractéristique mixte, choisissez
un anneau de Cohen $\Lambda$ pour $\kappa_B$ et dans le cas de caractéristique
équicaractéristique, définissez $\Lambda = \kappa_B[[t]]$. En argumentant comme ci-dessus, nous pouvons
choisir un morphisme d'anneaux $A \to \Lambda$ relevant $k \to \kappa_B$ et envoyant un
uniformisant sur un uniformisant. Puisque $k \subset \kappa_B$ est séparable, le morphisme
d'anneaux $A \to \Lambda$ est formellement lisse dans la topologie adique.
(Lemme \ref{lemma-extension-dvrs-formally-smooth}). Par conséquent, nous pouvons trouver un morphisme d'anneaux
$\Lambda \to B$ tel que la composition $A \to \Lambda \to B$ soit le morphisme d'anneaux donné $A \to B$
(voir le Lemme \ref{lemma-lift-continuous}). Comme $\Lambda$ et $B$ sont des anneaux de valuation
discrète complets avec le même corps résiduel, $B$ est fini sur $\Lambda$
(Algèbre, Lemme \ref{algebra-lemma-finite-over-complete-ring}). Cela nous ramène au cas particulier discuté dans le
Lemme \ref{lemma-special-case}.
\end{proof}







\section{Élimination de la ramification, II}
\label{section-eliminating-ramification-bis}

\noindent
Dans cette section, nous utilisons les résultats de la Section \ref{section-eliminating-ramification} pour
obtenir des solutions (séparables) dans certains cas.

\begin{lemma}
\label{lemma-solution-goes-up}
Soit $A \to B$ une extension d'anneaux de valuation discrète avec des corps des
fractions $K \subset L$. Si $K_1/K$ est une solution pour $A \subset B$, alors pour toute
extension finie $K_2/K_1$, l'extension $K_2/K$ est une solution pour $A \subset B$.
\end{lemma}

\begin{proof}
Ceci découle du Lemme \ref{lemma-formally-smooth-goes-up}. Détails omis.
\end{proof}

\begin{lemma}
\label{lemma-nagata-goes-down}
Soit $A \subset B$ une extension d'anneaux de valuation discrète. Si $B$ est Nagata
et que l'extension $L/K$ des corps des fractions est séparable, alors $A$
est Nagata.
\end{lemma}

\begin{proof}
Un anneau à valuation discrète est Nagata si et seulement s'il est N-2. Soit
$K_1/K$ une extension de corps finie et purement inséparable. Il nous faut montrer
que la clôture intégrale $A_1$ de $A$ dans $K_1$ est finie sur $A$,
voir Algèbre, Lemme \ref{algebra-lemma-domain-char-p-N-1-2}. Comme $L/K$ est séparable et $K_1/K$ est purement
inséparable, l'algèbre $L \otimes_K K_1$ est un corps (d'après Algèbre, Lemme \ref{algebra-lemma-separable-extension-preserves-reducedness} et
\ref{algebra-lemma-radicial-integral-bijective}). Soit $B_1$ la clôture intégrale de $B$ dans $L \otimes_K K_1$. Comme
$B$ est Nagata, $B_1$ est finie sur $B$. Comme $B \otimes_A A_1 \subset B_1$
et $B$ est noethérien, nous voyons que $B \otimes_A A_1$ est finie sur $B$. Comme
$A \to B$ est fidèlement plate, cela implique que $A_1$ est finie sur $A$,
voir Algèbre, Lemme \ref{algebra-lemma-descend-properties-modules}.
\end{proof}

\begin{lemma}
\label{lemma-construct-extension}
Soit $A' \subset A$ une extension d'anneaux. Soit $f \in A'$. Supposons que (a) $A$
soit fini sur $A'$, (b) $f$ soit un non-diviseur de zéro sur $A$, et
(c) $A'_f = A_f$. Il existe alors un entier $n_0 > 0$ tel que pour tout $n \geq n_0$ ce qui
suit est vrai : étant donnés un anneau $B'$ et un élément $g \in B'$ non-diviseur de zéro,
et un isomorphisme $\varphi' : A'/f^n A' \to B'/g^n B'$ avec $\varphi'(f) \equiv g$, il existe une extension finie $B' \subset B$
et un isomorphisme $\varphi : A/fA \to B/gB$ compatible avec $\varphi'$.
\end{lemma}

\begin{proof}
Puisque $A$ est fini sur $A'$ et puisque $A'_f = A_f$, nous pouvons choisir
$t > 0$ tel que $f^t A \subset A'$. Posons $n_0 = 2t$. Étant donné $n, B', g, \varphi'$ comme dans l'énoncé
du lemme, notons $N \subset B'$ l'ensemble des éléments $b \in B'$ tels que $b \bmod g^nB' \in \varphi'(f^tA)$. Posons
$B = g^{-t}N$. Comme $f^tA' \subset f^tA$ et que $\varphi'$ envoie $f$ à $g$, nous avons
$g^tB' \subset N$, donc $B' \subset B$. Puisque $f^tA \cdot f^tA \subset f^t \cdot f^tA$ et que $\varphi'$ envoie $f$ à
$g$, nous voyons que $N \cdot N \subset g^t N$. Ainsi nous obtenons une multiplication sur
$B$ étendant la multiplication de $B'$. Nous avons un isomorphisme de
$A'/f^nA'$-modules
$$
A/f^tA' \xrightarrow{f^t} f^tA/f^nA' \xrightarrow{\varphi'}
g^tB/g^nB' \xrightarrow{g^{-t}} B/g^tB'
$$
où les structures de module à droite sont définies en utilisant $\varphi'$. Puisque
$A/f^tA'$ est un $A'$-module fini, nous en concluons que $B/g^tB'$ est un
$B'$-module fini et donc nous voyons que $B' \to B$ est fini. Enfin, nous laissons
au lecteur le soin de voir que l'isomorphisme affiché de modules envoie $fA$
dans $gB$ et induit un isomorphisme d'anneaux $\varphi : A/fA \to B/gB$ compatible avec $\varphi'$
(il induit même un isomorphisme $A/f^tA \to B/g^tB$ mais nous n'en avons pas besoin).
\end{proof}

\begin{remark}
\label{remark-functoriality-construct-extension}
La construction dans le Lemme \ref{lemma-construct-extension} satisfait la ``fonctorialité'' suivante.
Supposons que nous ayons un diagramme commutatif
$$
\xymatrix{
A'_2 \ar[r] & A_2 \\
A'_1 \ar[r] \ar[u] & A_1 \ar[u]
}
$$
avec des flèches horizontales injectives. Supposons qu’un élément $f \in A'_1$ soit tel
que $(A'_1 \subset A_1, f)$ et $(A'_2 \subset A_2, f)$ satisfaient les propriétés (a), (b), (c) du lemme \ref{lemma-construct-extension}.
Soient $n_{0, 1}$ et $n_{0, 2}$ les entiers trouvés dans le lemme pour ces deux
situations. Enfin, soit $B'_1 \to B'_2$ une application d'anneaux, soit $g \in B'_1$ un
non-diviseur de zéro sur $B_1$ et $B_2$, soit $n \geq \max(n_{0, 1}, n_{0, 2})$, et soit donné un
diagramme commutatif
$$
\xymatrix{
A'_2/f^nA'_2 \ar[r]_{\varphi'_2} & B'_2/g^nB'_2 \\
A'_1/f^nA'_1 \ar[r]^{\varphi'_1} \ar[u] & B'_2/g^nB'_2 \ar[u]
}
$$
dont les flèches horizontales sont des isomorphismes et où $\varphi'_1(f) \equiv g$. On obtient
alors des diagrammes commutatifs
$$
\vcenter{
\xymatrix{
B'_2 \ar[r] & B_2 \\
B'_1 \ar[r] \ar[u] & B_1 \ar[u]
}
}
\quad\text{et}\quad
\vcenter{
\xymatrix{
A_2/fA_2 \ar[r]_{\varphi_2} & B_2/gB_2 \\
A_1/fA_1 \ar[r]^{\varphi_1} \ar[u] & B_2/gB_2 \ar[u]
}
}
$$
où $(B'_1 \subset B_1, \varphi_1)$ et $(B'_2 \subset B_2, \varphi_2)$ sont construits comme dans la preuve du Lemme \ref{lemma-construct-extension}.
Nous omettons la vérification détaillée.
\end{remark}


\begin{lemma}
\label{lemma-approximate-solution}
Soit $p$ un nombre premier. Soit $A \subset B$ une extension d'anneaux de valuation
discrète avec extension de corps des fractions $L/K$. Soit $K_2/K_1/K$ une tour
d'extensions de corps finis. Supposons que
\begin{enumerate}
\item $K$ ait pour caractéristique $p$,
\item $L/K$ soit séparable,
\item $B$ soit Nagata,
\item $K_2$ soit une solution pour $A \subset B$,
\item $K_2/K_1$ soit purement inséparable de degré $p$.
\end{enumerate}
Alors il existe une extension séparable $K_3/K_1$ qui est une solution pour
$A \subset B$.
\end{lemma}

\begin{proof}
Reprenons les notations de la Remarque \ref{remark-construction} ; nous utiliserons toutes
les observations qui y ont été faites. Comme $L/K$ est séparable, l'algèbre $L_1 = L \otimes_K K_1$
est réduite (Algèbre, Lemme \ref{algebra-lemma-separable-extension-preserves-reducedness}). Comme $B$ est un anneau de Nagata, l'extension
d'anneaux $B \subset B_1$ est finie, où $B_1$ est la clôture intégrale de $B$ dans
$L_1$ et $B_1$ est un anneau de Nagata. De même, l'anneau $A$ est un anneau de
Nagata d'après le Lemme \ref{lemma-nagata-goes-down}, donc $A \subset A_1$ est fini et $A_1$ est aussi un
anneau de Nagata. De plus, les mêmes affirmations sont vraies pour $K_2$,
c'est-à-dire que $L_2 = L \otimes_K K_2$ est réduite, les extensions d'anneaux $A_1 \subset A_2$ et
$B_1 \subset B_2$
sont finies, où $A_2$, respectivement $B_2$, est la clôture intégrale de
$A$, respectivement de $B$, dans $K_2$, respectivement $L_2$.

\medskip\noindent Soit $\pi \in A$ un uniformisant. Remarquons que $\pi$ est un
non-diviseur de zéro sur $K_1$, $K_2$, $A_1$, $A_2$, $L_1$, $L_2$,
$B_1$, et $B_2$ et nous avons $K_1 = (A_1)_\pi$, $K_2 = (A_2)_\pi$, $L_1 = (B_1)_\pi$, et $L_2 = (B_2)_\pi$. Nous
pouvons écrire $K_2 = K_1(\alpha)$ où $\alpha^p = a_1 \in K_1$, voir Corps, Lemme \ref{fields-lemma-finite-purely-inseparable}. Après avoir
multiplié $\alpha$ par une puissance de $\pi$, nous pouvons supposer, et nous supposerons, que
$a_1 \in A_1$. Pour le reste de la preuve, il est commode d'écrire $K_2 = K_1[x]/(x^p - a_1)$ et
$L_2 = L_1[x]/(x^p - a_1)$. Considérons les extensions d'anneaux
$$
A'_2 = A_1[x]/(x^p - a_1) \subset A_2
\quad\text{et}\quad
B'_2 = B_1[x]/(x^p - a_1) \subset B_2
$$
Nous pouvons appliquer le Lemme \ref{lemma-construct-extension} à $A'_2 \subset A_2$ et $f = \pi^2$ et à $B'_2 \subset B_2$ et
$f = \pi^2$. Choisissons un entier $n$ suffisamment grand qui convienne dans ces
deux cas.

\medskip\noindent Considérons les algèbres
$$
K_3 = K_1[x]/(x^p - \pi^{2n} x - a_1)
\quad\text{et}\quad
L_3 = L_1[x]/(x^p - \pi^{2n} x - a_1)
$$
Remarquez que $K_3/K_1$ et $L_3/L_1$ sont des extensions algébriques \'etales
finies de degré $p$. Considérons les sous-anneaux
$$
A'_3 = A_1[x]/(x^p - \pi^n x - a_1)
\quad\text{et}\quad
B'_3 = B_1[x]/(x^p - \pi^n x - a_1)
$$
de $K_3 = (A'_2)_\pi$ et $L_3 = (B'_3)_\pi$. Nous allons construire un diagramme commutatif
$$
\xymatrix{
B'_2/\pi^{2n} B'_2 \ar[r]_{\psi'} & B'_3/\pi^{2n} B'_3 \\
A'_2/\pi^{2n} A'_2 \ar[r]^{\varphi'} \ar[u] & A'_3/\pi^{2n} A'_3 \ar[u]
}
$$
. À savoir, $\varphi'$ est l'unique isomorphisme de $A_1$-algèbres envoyant la classe
de $x$ sur la classe de $x$. De même, $\psi'$ est l'unique isomorphisme
de $B_1$-algèbres envoyant la classe de $x$ sur la classe de $x$. Par
notre choix de $n$, nous obtenons, via le Lemme \ref{lemma-construct-extension} et la Remarque
\ref{remark-functoriality-construct-extension}, des extensions d'anneaux finies $A'_3 \subset A_3$ et $B'_3 \subset B_3$ telles que $A'_3 \to B'_3$
se prolonge en un homomorphisme d'anneaux $A_3 \to B_3$ et un diagramme commutatif
$$
\xymatrix{
B_2/\pi^2 B_2 \ar[r]_\psi & B_3/\pi^2B_3 \\
A_2/\pi^2 A_2 \ar[r]^\varphi \ar[u] & A_3/\pi^2A_3 \ar[u]
}
$$
avec toutes les propriétés affirmées dans les références mentionnées ci-dessus (en
particulier $\varphi$ et $\psi$ sont des isomorphismes).

\medskip\noindent Avec toutes ces données en main, nous pouvons terminer la
démonstration. À savoir, nous observons d'abord que $A_3$ et $B_3$ sont des
produits finis de domaines de Dedekind, et $\pi$ appartient à tous leurs idéaux
maximaux. En effet, si $\mathfrak p \subset A_3$ est un idéal maximal, alors $\pi \in \mathfrak p$, puisque $A \to A_3$
est fini. Ensuite, $\mathfrak p/\pi^2 A_3$ correspond via $\varphi$ à un idéal maximal dans $A_2 / \pi^2A_2$
qui est principal car $A_2$ est un produit fini de domaines de Dedekind. Nous
concluons que $\mathfrak p/\pi^2 A_3$ est principal et donc, par Nakayama, nous voyons que
$\mathfrak p (A_3)_\mathfrak p$ est principal. Le même argument fonctionne pour $B_3$. Nous concluons
que
$A_3$ est la clôture intégrale de $A$ dans $K_3$ et que $B_3$ est la
clôture intégrale de $B$ dans $L_3$. Soit $\mathfrak q \subset B_3$ un idéal maximal au-dessus
de $\mathfrak p \subset A_3$. Pour terminer la démonstration, nous devons montrer que $(A_3)_\mathfrak p \to (B_3)_\mathfrak q$ est
formellement lisse dans la topologie $\mathfrak q$-adique. Selon le critère du Lemme
\ref{lemma-extension-dvrs-formally-smooth}, il suffit de montrer $\mathfrak p (B_3)_\mathfrak q = \mathfrak q (B_3)_\mathfrak q$ et que l'extension de corps $\kappa(\mathfrak q)/\kappa(\mathfrak p)$ est
séparable. Cela est vrai car nous pouvons vérifier les deux assertions en
examinant l'homomorphisme d'anneaux $A_3/\pi^2 A_3 \to B_3/\pi^2 B_3$
et celui-ci est isomorphe à l'homomorphisme d'anneaux $A_2/\pi^2 A_2 \to B_2/\pi^2 B_2$, où l'énoncé
correspondant tient par notre hypothèse que $K_2$ est une solution pour
$A \subset B$. Quelques détails omis.
\end{proof}

\begin{lemma}
\label{lemma-separable-solution-separable-solution}
Soit $A \subset B$ une extension d'anneaux à valuation discrète. Supposons que
\begin{enumerate}
\item l'extension $L/K$ des corps des fractions soit séparable,
\item $B$ soit Nagata, et
\item qu'il existe une solution pour $A \subset B$.
\end{enumerate}
Alors il existe une solution séparable pour $A \subset B$.
\end{lemma}

\begin{proof}
Le lemme est trivial si la caractéristique de $K$ est nulle ; ainsi, nous
pouvons supposer, et nous supposerons, que la caractéristique de $K$ est $p > 0$.

\medskip\noindent Soit $K_2/K$ une solution pour $A \to B$. Nous allons utiliser
l'induction sur le degré inséparable $[K_2 : K]_i$ (Corps, Définition \ref{fields-definition-insep-degree}) de
$K_2/K$. Si $[K_2 : K]_i = 1$, alors $K_2$ est séparable sur $K$ et nous avons
terminé. Sinon, il existe un sous-corps $K_2/K_1/K$ tel que $K_2/K_1$ est purement
inséparable de degré $p$ (Corps, Lemme \ref{fields-lemma-separable-first} et \ref{fields-lemma-finite-purely-inseparable}). D'après le
lemme \ref{lemma-approximate-solution}, il existe une extension séparable $K_3/K_1$ qui est une solution
pour $A \subset B$. Alors $[K_3 : K]_i = [K_1 : K]_i = [K_2 : K]_i/p$
(Corps, Lemme \ref{fields-lemma-multiplicativity-all-degrees}) est plus petit et nous concluons par induction.
\end{proof}

\begin{lemma}
\label{lemma-big-extension-is-ok}
Soit $A \to B$ une extension d'anneaux de valuation discrète avec des corps des
fractions $K \subset L$. Supposons que $B$ soit essentiellement de type fini sur
$A$. Soit $K'/K$ une extension algébrique de corps telle que la clôture
intégrale $A'$ de $A$ dans $K'$ soit noethérienne. Alors la clôture
intégrale $B'$ de $B$ dans $L' = (L \otimes_K K')_{red}$ est aussi noethérienne. De plus,
l'application $\Spec(B') \to \Spec(A')$ est surjective et les extensions de corps résiduels
correspondantes sont des extensions de corps de type fini.
\end{lemma}

\begin{proof}
Soit $A \to C$ une application d’anneaux de type fini telle que $B$ soit une
localisation de $C$ à un $\mathfrak p$ premier. Alors $C' = C \otimes_A A'$ est une algèbre de
type fini sur $A'$, donc noethérienne. Puisque $A \to A'$ est intégrale,
$C \to C'$ l’est aussi. Ainsi, l'extension $B = C_\mathfrak p \subset C'_\mathfrak p$ est elle aussi intégrale. Il s’ensuit que la
dimension de $C'_\mathfrak p$ est $1$ (Algèbre, Lemme \ref{algebra-lemma-integral-sub-dim-equal}). Bien sûr, $C'_\mathfrak p$ est
noethérienne. Soient $\mathfrak q_1, \ldots, \mathfrak q_n$ les idéaux premiers minimaux de $C'_\mathfrak p$. Soit
$B'_i$ la clôture intégrale de $B = C_\mathfrak p$, ou, ce qui est équivalent d'après ce qui précède,
de $C'_\mathfrak p$
dans le corps des fractions de $C'_{\mathfrak p'}/\mathfrak q_i$. Il résulte du lemme de Krull-Akizuki
(Algèbre, Lemme \ref{algebra-lemma-krull-akizuki} appliqué au nombre fini de localisations de $C'_\mathfrak p$ en ses idéaux
maximaux) que chaque $B'_i$ est noethérien. De plus, les extensions de corps
résiduels dans $C'_\mathfrak p \to B'_i$ sont finies selon Algèbre, Lemme \ref{algebra-lemma-finite-extension-residue-fields-dimension-1}. Enfin, on observe
que $B' = \prod B'_i$ est la clôture intégrale de $B$ dans $L' = (L \otimes_K K')_{red}$.
\end{proof}

\begin{proposition}
\label{proposition-epp-essentially-finite-type}
\begin{reference}
Voir \cite[Lemme 2.13]{alterations} pour un cas particulier.
\end{reference}
Soit $A \to B$ une extension d'anneaux de valuation discrète avec des corps des
fractions $K \subset L$. Si $B$ est essentiellement de type fini sur $A$,
alors il existe une extension finie $K_1/K$ qui est une solution pour $A \to B$
comme défini dans la Définition \ref{definition-solution}.
\end{proposition}

\begin{proof}
Observez qu'une solution faible est une solution si le corps résiduel de $A$
est parfait, voir le Lemme \ref{lemma-extension-dvrs-formally-smooth}. Ainsi, la proposition découle immédiatement du
Théorème \ref{theorem-epp} si la caractéristique résiduelle de $A$ est $0$ (et en
fait, nous n'avons pas besoin de l'hypothèse que $A \to B$ est essentiellement de
type fini). Si la caractéristique résiduelle de $A$ est $p > 0$, nous la
déduirons également du théorème d'Epp.

\medskip\noindent Soit donnée une famille d'éléments $x_i \in A$, $i \in I$, dont les images forment
une $p$-base du corps résiduel $\kappa$ de $A$. Posons
$$
A' = \bigcup\nolimits_{n \geq 1} A[t_{i, n}]/(t_{i, n}^{p^n} - x_i)
$$
où les applications de transition envoient $t_{i, n + 1}$ vers $t_{i, n}^p$. Observez que
$A'$ est une colimite filtrante d'extensions finies faiblement non ramifiées de
$A$ par des anneaux de valuation discrète. Ainsi $A'$ est un anneau de
valuation discrète et $A \to A'$ est faiblement non ramifié. Par construction, le
corps résiduel $\kappa' = A'/\mathfrak m_A A'$ est la perfection de $\kappa$.

\medskip\noindent Soit $K'$ le corps des fractions de $A'$. Nous pouvons
appliquer le Lemme \ref{lemma-big-extension-is-ok} à l’extension $K'/K$. Ainsi, $B'$ est un produit fini
de domaines de Dedekind. Soient $\mathfrak m_1, \ldots, \mathfrak m_n$ les idéaux maximaux de $B'$. En
utilisant le théorème d’Epp (théorème \ref{theorem-epp}), on trouve une solution faible
$K'_i/K'$ pour chacune des extensions $A' \subset B'_{\mathfrak m_i}$. Puisque le corps résiduel de $A'$
est parfait, ce sont en réalité des solutions. Soit $K'_1/K'$ une extension finie
qui contient chaque $K'_i$. Alors $K'_1/K'$ reste une solution pour chaque
$A' \subset B'_{\mathfrak m_i}$ par le Lemme \ref{lemma-solution-goes-up}.

\medskip\noindent Soit $A'_1$ la clôture intégrale de $A$ dans $K'_1$. Notez
que $A'_1$ est un domaine de Dedekind d'après la discussion dans la Remarque
\ref{remark-construction} appliquée à $K' \subset K'_1$. Ainsi, le Lemme \ref{lemma-big-extension-is-ok} s'applique à $K'_1/K$. Par
conséquent, la clôture intégrale $B'_1$ de $B$ dans $L'_1 = (L \otimes_K K'_1)_{red}$ est un domaine
de Dedekind et, puisque $K'_1/K'$ est une solution pour chaque $A' \subset B'_{\mathfrak m_i}$, nous
voyons que $(A'_1)_{A'_1 \cap \mathfrak m} \to (B'_1)_{\mathfrak m}$ est formellement lisse dans la topologie $\mathfrak m$-adique pour
chaque idéal maximal $\mathfrak m \subset B'_1$.

\medskip\noindent Par construction, le corps $K'_1$ est une colimite filtrante
d'extensions finies de $K$. Disons $K'_1 = \colim_{i \in I} K_i$. Pour chaque $i$, soit $A_i$,
resp. $B_i$, la clôture intégrale de $A$, resp. $B$ dans $K_i$, resp.
$L_i = (L \otimes_K K_i)_{red}$. Alors il est clair que
$$
A'_1 =  \colim A_i\quad\text{et}\quad B'_1 = \colim B_i
$$
Puisque les morphismes d'anneaux $A_i \to A'_1$ et $B_i \to B'_1$ sont des homomorphismes d'anneaux
injectifs et intégraux et que $A'_1$ et $B'_1$ ont des spectres finis, nous
voyons que pour tout $i$ suffisamment grand, les morphismes d'anneaux $A_i \to A'_1$
et $B_i \to B'_1$ induisent des bijections sur les spectres. Une fois cela vrai, pour tout $i$
suffisamment grand, les applications $A_i \to A'_1$ et $B_i \to B'_1$ seront faiblement non
ramifiées (une fois que l'uniformisant appartient à l'image). Il résulte de la
multiplicativité des indices de ramification que $A_i \to B_i$ induit des applications
faiblement non ramifiées sur toutes les localisations aux idéaux maximaux de
$B_i$ pour de tels $i$. En augmentant un peu plus $i$, nous voyons
que
$$
B_i \otimes_{A_i} A'_1 \longrightarrow B'_1
$$
induit des applications surjectives sur les corps résiduels (parce que les corps
résiduels de $B'_1$ sont de type fini sur ceux de $A'_1$ selon le Lemme
\ref{lemma-big-extension-is-ok}). Voici le diagramme des corps résiduels aux idéaux maximaux situés sous un idéal
maximal choisi de $B'_1$ :
$$
\xymatrix{
\kappa_{B_i} \ar[r] &
\kappa_{B_{i'}} \ar[r] &
 & \ldots &
\kappa_{B'_1} \\
\kappa_{A_i} \ar[r] \ar[u] &
\kappa_{A_{i'}} \ar[r] \ar[u] &
 & \ldots &
\kappa_{A'_1} \ar[u]
}
$$
Ainsi $\kappa_{B_i}$ est une extension finiment engendrée de $\kappa_{A_i}$ telle que le
compositum de $\kappa_{B_i}$ et $\kappa_{A'_1}$ dans $\kappa_{B'_1}$ est séparable sur $\kappa_{A'_1}$. Cela se
produit déjà à un stade fini : par exemple, supposons que $\kappa_{B'_1}$ soit une extension finie et
séparable de $\kappa_{A'_1}(x_1, \ldots, x_n)$ ; il suffit alors d'augmenter $i$ de telle sorte que
$x_1, \ldots, x_n$ soient dans $\kappa_{B_i}$ et que tous les générateurs satisfassent des
équations polynomiales séparables sur $\kappa_{A_i}(x_1, \ldots, x_n)$. Cela signifie que $A_i \to (B_i)_\mathfrak m$ est
formellement lisse dans la topologie $\mathfrak m$-adique pour tous les idéaux maximaux
$\mathfrak m$ de $B_i$ et la preuve est complète.
\end{proof}

\begin{lemma}
\label{lemma-epp-essentially-finite-type-separable}
Soit $A \to B$ une extension d'anneaux de valuation discrète avec des corps des
fractions $K \subset L$. Supposons
\begin{enumerate}
\item $B$ est essentiellement de type fini sur $A$,
\item soit $A$, soit $B$, est un anneau de Nagata, et
\item $L/K$ est séparable.
\end{enumerate}
Alors il existe une solution séparable pour $A \to B$ (Définition \ref{definition-solution}).
\end{lemma}

\begin{proof}
Remarquez que si $A$ est un anneau de Nagata, alors $B$ l'est aussi (Algèbre, Lemme
\ref{algebra-lemma-nagata-localize} et Proposition \ref{algebra-proposition-nagata-universally-japanese}). Ainsi, le lemme s'ensuit en combinant la
Proposition \ref{proposition-epp-essentially-finite-type} et le Lemme \ref{lemma-separable-solution-separable-solution}.
\end{proof}







\section{Groupes de Picard des anneaux}
\label{section-picard}

\noindent
Nous définissons d’abord les modules inversibles comme suit.

\begin{definition}
\label{definition-invertible}
Soit $R$ un anneau. Un $R$-module $M$ est {\it inversible}
si le foncteur
$$
\text{Mod}_R \longrightarrow \text{Mod}_R,\quad
N \longmapsto M \otimes_R N
$$
est une équivalence de catégories. Un $R$-module inversible est dit {\it
trivial} s’il est isomorphe à $R$ en tant que $R$-module.
\end{definition}

\begin{lemma}
\label{lemma-invertible}
Soit $R$ un anneau. Soit $M$ un $R$-module. Les conditions suivantes sont
équivalentes :
\begin{enumerate}
\item $M$ est un module fini localement libre de rang $1$,
\item $M$ est inversible, et
\item il existe un $R$-module $N$ tel que $M \otimes_R N \cong R$.
\end{enumerate}
De plus, dans ce cas, le module $N$ dans (3) est isomorphe à $\Hom_R(M, R)$.
\end{lemma}

\begin{proof}
Supposons (1). Considérons le module $N = \Hom_R(M, R)$ et l'application d'évaluation
$M \otimes_R N = M \otimes_R \Hom_R(M, R) \to R$. Si $f \in R$ tel que $M_f \cong R_f$, alors l'application d'évaluation devient un
isomorphisme après localisation en $f$ (détails omis). Ainsi, nous voyons que
l'application d'évaluation est un isomorphisme d'après le Lemme \ref{algebra-lemma-cover} d'Algèbre. Ainsi
(1) $\Rightarrow$ (3).

\medskip\noindent Supposons (3). Alors le foncteur $K \mapsto K \otimes_R N$ est un quasi-inverse au
foncteur $K \mapsto K \otimes_R M$. Ainsi (3) $\Rightarrow$ (2). Inversement, si (2) est vrai, alors
$K \mapsto K \otimes_R M$ est essentiellement surjectif et nous voyons que (3) est vrai.

\medskip\noindent Supposons que les conditions équivalentes (2) et (3) soient
vérifiées. Notons $\psi : M \otimes_R N \to R$ l'isomorphisme provenant de (3). Choisissons un élément
$\xi = \sum_{i = 1, \ldots, n} x_i \otimes y_i$ tel que $\psi(\xi) = 1$. Considérons les isomorphismes
$$
M \to M \otimes_R M \otimes_R N \to M
$$
où la première flèche envoie $x$ sur $\sum x_i \otimes x \otimes y_i$ et la deuxième flèche envoie
$x \otimes x' \otimes y$ sur $\psi(x' \otimes y)x$. Nous en concluons que $x \mapsto \sum \psi(x \otimes y_i)x_i$ est un automorphisme de
$M$. Cet automorphisme se factorise comme
$$
M \to R^{\oplus n} \to M
$$
où la première flèche est donnée par $x \mapsto (\psi(x \otimes y_1), \ldots, \psi(x \otimes y_n))$ et la deuxième flèche par $(a_1, \ldots, a_n) \mapsto \sum a_i x_i$. De
cette façon, nous concluons que $M$ est un facteur direct d'un $R$-module
libre de type fini. Cela signifie que $M$ est localement libre de type fini
(Algèbre, Lemme \ref{algebra-lemma-finite-projective}). Comme il en est de même pour $N$ par symétrie et
puisque $M \otimes_R N \cong R$, nous voyons que $M$ et $N$ doivent tous deux avoir un
rang $1$.
\end{proof}

\noindent
L'ensemble des classes d'isomorphisme de ces modules est souvent appelé le {groupe
de classes \it} ou {groupe de Picard \it} de $R$. La structure de
groupe est déterminée en assignant aux classes d'isomorphisme des modules
inversibles $L$ et $L'$ la classe d'isomorphisme de $L \otimes_R L'$. L'inverse
d'un module inversible $L$ est le module
$$
L^{\otimes -1} = \Hom_R(L, R),
$$
car, comme on le voit dans la démonstration du Lemme \ref{lemma-invertible}, l'application
d'évaluation $L \otimes_R L^{\otimes -1} \to R$ est un isomorphisme. Le groupe de Picard de
$R$ sera noté $\Pic(R)$.

\begin{lemma}
\label{lemma-UFD-Pic-trivial}
Soit $R$ un anneau factoriel. Alors $\Pic(R)$ est trivial.
\end{lemma}

\begin{proof}
Soit $L$ un $R$-module inversible. D'après le Lemme \ref{lemma-invertible}, nous voyons
que $L$ est un $R$-module localement libre de type fini. En particulier,
$L$ est sans torsion et de type fini sur $R$. Choisissons un élément non
nul $\varphi \in \Hom_R(L, R)$ du module inversible dual. Alors $I = \varphi(L) \subset R$ est un idéal qui est un
module inversible. Choisissons un $f \in I$ non nul et écrivons
$$
f = u p_1^{e_1} \ldots p_r^{e_r}
$$
sa factorisation en éléments premiers, où les $p_i$ sont deux à deux distincts.
Puisque $L$ est un module localement libre de type fini, il existe $a_i \in R$, $a_i \not \in (p_i)$ tel que
$I_{a_i} = (g_i)$ pour un certain $g_i \in R_{a_i}$. Alors $p_i$ reste un élément premier de
l'anneau factoriel $R_{a_i}$ et nous pouvons écrire $g_i = p_i^{c_i} g'_i$ pour un certain $g'_i \in R_{a_i}$ non divisible par
$p_i$. Comme $f \in I_{a_i}$, on voit que $e_i \geq c_i$. Nous affirmons que $I$ est
engendré par $h = p_1^{c_1} \ldots p_r^{c_r}$, ce qui termine la démonstration.

\medskip\noindent Pour prouver l'énoncé, il suffit de montrer que $I_a$ est
engendré par $h$ pour tout $a \in R$ tel que $I_a$ soit un idéal principal
(Algèbre, Lemme \ref{algebra-lemma-cover}). Disons $I_a = (g)$. Soit $J \subset \{1, \ldots, r\}$ l'ensemble des $i$ tels
que $p_i$ soit une non-unité (et donc un élément premier) dans $R_a$. Comme
$f \in I_a = (g)$, nous obtenons la factorisation en éléments premiers $g = v \prod_{i \in J} p_j^{b_j}$ avec
$v$ une unité et $b_j \leq e_j$. Pour chaque $j \in J$, nous avons $I_{aa_j} = g R_{aa_j} = g_j R_{aa_j}$, en
d'autres termes, $g$ et $g_j$ sont associés dans
$R_{aa_j}$. Par unicité de la factorisation, cela implique que $b_j = c_j$ et la preuve
est complète.
\end{proof}







\section{Déterminants}
\label{section-determinants}

\noindent
Soit $R$ un anneau. Soit $M$ un $R$-module projectif fini. Il existe
une décomposition en produit $R = R_0 \times \ldots \times R_t$ telle que dans la décomposition
correspondante $M = M_0 \times \ldots \times M_t$ de $M$, nous avons que $M_i$ est localement libre de
rang fini $i$ sur $R_i$. Cela découle du Lemme \ref{algebra-lemma-finite-projective} d’Algèbre (pour voir
que le rang est localement constant) et des Lemmes \ref{algebra-lemma-disjoint-decomposition} et \ref{algebra-lemma-disjoint-implies-product}
d’Algèbre (pour décomposer $R$ en un produit). Dans cette situation, nous
définissons
$$
\det(M) = \wedge^0_{R_0}(M_0) \times \ldots \times \wedge^t_{R_t}(M_t)
$$
comme un $R$-module. Il s’agit d’un module localement libre de rang
$1$ puisque chaque terme est localement libre de rang $1$. Si
$\varphi : M \to N$ est un isomorphisme de $R$-modules projectifs finis, alors nous
obtenons un isomorphisme canonique :
$$
\det(\varphi) : \det(M) \longrightarrow \det(N)
$$
de modules localement libres de rang $1$. Plus généralement, si pour tous les
idéaux premiers $\mathfrak p$ de $R$ les rangs des modules libres $M_\mathfrak p$ et
$N_\mathfrak p$ sont les mêmes, alors tout homomorphisme de $R$-modules $\varphi : M \to N$ induit
une application de $R$-modules $\det(\varphi) : \det(M) \to \det(N)$. Enfin, si $M = N$ alors $\det(\varphi) : \det(M) \to \det(M)$ est
un endomorphisme d'un $R$-module inversible. Puisque $R = \Hom_R(L, L)$ pour un
$R$-module inversible, nous pouvons donc considérer $\det(\varphi)$ comme un élément
de $R$. De cette manière, nous obtenons le {\it déterminant}
$$
\det : \Hom_R(M, M) \longrightarrow R
$$
qui est une application multiplicative.

\begin{remark}
\label{remark-determinant-as-socle}
Soit $R$ un anneau. Soit $M$ un $R$-module projectif fini. Alors nous
pouvons considérer la $R$-algèbre graduée commutative qu'est l'algèbre extérieure
$\wedge^*_R(M)$ de $M$ sur $R$. Une formule pour $\det(M)$ est que
$\det(M) \subset \wedge^*_R(M)$ est
l'annulateur de $M \subset \wedge^*_R(M)$. Cela est parfois utile car cela ne fait pas référence à
la décomposition de $R$ en un produit. Bien sûr, pour prouver que cela
satisfait les propriétés souhaitées, il faut soit décomposer $R$ en un
produit (comme ci-dessus), soit examiner les localisations aux idéaux premiers de
$R$.
\end{remark}

\noindent
Ensuite, nous considérons le comportement du déterminant associé à une
suite exacte courte de modules projectifs finis.

\begin{lemma}
\label{lemma-det-ses}
Soit $R$ un anneau. Soit
$$
0 \to M' \to M \to M'' \to 0
$$
une suite exacte courte de $R$-modules projectifs finis. Alors il existe un
isomorphisme canonique
$$
\gamma : \det(M') \otimes \det(M'') \longrightarrow \det(M)
$$
\end{lemma}

\begin{proof}[Première preuve]
Première preuve. Décomposez $R$ en un produit d'anneaux $R_{ij}$ tel que
$M' = \prod M'_{ij}$ et $M'' = \prod M''_{ij}$ où $M'_{ij}$ a un rang $i$ et $M''_{ij}$ a un rang $j$. Bien
sûr, alors $M = \prod M_{ij}$ et $M_{ij}$ a un rang $i + j$. Cela nous réduit au cas où
$M'$ et $M''$ sont de rang constant, disons $i$ et $j$. Dans ce
cas, nous devons construire une application canonique
$$
\wedge^i(M') \otimes \wedge^j(M'') \longrightarrow \wedge^{i + j}(M)
$$
. Pour ce faire, choisissez $m'_1, \ldots, m'_i$ dans $M'$ et $m''_1, \ldots, m''_j$ dans $M''$.
Notons $m_1, \ldots, m_i \in M$ les images de $m'_1, \ldots, m'_i$ et choisissons $m_{i + 1}, \ldots , m_{i + j} \in M$, dont les images
sont $m''_1, \ldots, m''_j$ dans $M''$. Notre règle sera que
$$
m'_1 \wedge \ldots \wedge m'_i \otimes
m''_1 \wedge \ldots \wedge m''_j
\longmapsto
m_1 \wedge \ldots \wedge m_{i + j}
$$
. Nous omettons la preuve détaillée que ceci est bien défini et un isomorphisme.
\end{proof}

\begin{proof}[Deuxième preuve]
Nous utiliserons la description de $\det(M)$, $\det(M')$ et $\det(M'')$ donnée dans la
Remarque \ref{remark-determinant-as-socle}. Considérons les homomorphismes de $R$-algèbres $\wedge^*_R(M') \to \wedge^*_R(M)$ et
$\wedge^*_R(M) \to \wedge^*_R(M'')$. La première est injective et la seconde est surjective. Prenons un
élément $x' \in \det(M') \subset \wedge^*_R(M')$ et un élément $x'' \in \det(M'') \subset \wedge^*_R(M'')$. Choisissons un élément $y'' \in \wedge^*(M)$ ayant pour image
sur $x''$ et posons
$$
\gamma(x' \otimes x'') = x' \wedge y'' \in \det(M) \subset \wedge^*_R(M)
$$
. Le lecteur vérifie facilement en regardant les localisations aux idéaux
premiers que ceci est bien défini et un isomorphisme. De plus, cette construction
donne la même application que la construction donnée dans la première
démonstration.
\end{proof}

\begin{lemma}
\label{lemma-det-ses-functorial}
Soit $R$ un anneau. Soit
$$
\xymatrix{
0 \ar[r] &
M' \ar[r] \ar[d]^u &
M \ar[r] \ar[d]^v &
M'' \ar[r] \ar[d]^w &
0 \\
0 \ar[r] &
K' \ar[r] &
K \ar[r] &
K'' \ar[r] &
0
}
$$
un diagramme commutatif de modules projectifs finis sur $R$ dont les flèches
verticales sont des isomorphismes. Alors nous obtenons un diagramme commutatif
d'isomorphismes.
$$
\xymatrix{
\det(M') \otimes \det(M'') \ar[r]_-\gamma \ar[d]_{\det(u) \otimes \det(w)} &
\det(M) \ar[d]^{\det(v)} \\
\det(K') \otimes \det(K'') \ar[r]^-\gamma & \det(K)
}
$$
où les flèches horizontales sont celles construites dans le Lemme \ref{lemma-det-ses}.
\end{lemma}

\begin{proof}
Omis. Indication : utilisez la deuxième construction des applications $\gamma$ dans le
Lemme \ref{lemma-det-ses}.
\end{proof}

\begin{lemma}
\label{lemma-det-filtration}
Soit $R$ un anneau. Soient
$$
K \subset L \subset M
$$
des $R$-modules tels que $K$, $L/K$ et $M/L$ soient des
$R$-modules projectifs finis. Alors le diagramme
$$
\xymatrix{
\det(K) \otimes \det(L/K) \otimes \det(M/L) \ar[r] \ar[d] &
\det(L) \otimes \det(M/L) \ar[d] \\
\det(K) \otimes \det(M/K) \ar[r] &
\det(M)
}
$$
commute où les applications sont celles du Lemme \ref{lemma-det-ses}.
\end{lemma}

\begin{proof}
Omis. Indication : après localisation en un idéal premier de $R$, nous pouvons
supposer que $K \subset L \subset M$ est isomorphe à $R^{\oplus a} \subset R^{\oplus a + b}  \subset R^{\oplus a + b + c}$ et dans ce cas le résultat est un
calcul évident.
\end{proof}

\begin{lemma}
\label{lemma-det-direct-sum}
Soit $R$ un anneau. Soient $M'$ et $M''$ deux $R$-modules
projectifs finis. Alors le diagramme
$$
\xymatrix{
\det(M') \otimes \det(M'') \ar[r]
\ar[d]_{\epsilon \cdot (\text{permutation des tenseurs})} &
\det(M' \oplus M'') \ar[d]^{\det(\text{permutation des facteurs})} \\
\det(M'') \otimes \det(M') \ar[r] &
\det(M'' \oplus M')
}
$$
est commutatif, où $\epsilon = \det( -\text{id}_{M' \otimes M''}) \in R^*$ et les flèches horizontales sont celles du Lemme \ref{lemma-det-ses}.
\end{lemma}

\begin{proof}
Omis.
\end{proof}

\begin{lemma}
\label{lemma-det-switch}
Soit $R$ un anneau. Soient $M$, $N$ des $R$-modules projectifs
finis. Soient $a : M \to N$ et $b : N \to M$ des applications $R$-linéaires. Alors
$$
\det(\text{id} + a \circ b) = \det(\text{id} + b \circ a)
$$
en tant qu'éléments de $R$.
\end{lemma}

\begin{proof}
Il suffit de prouver l'assertion après avoir remplacé $R$ par une localisation
en un idéal premier. Ainsi nous pouvons supposer que $R$ est local et que
$M$ et $N$ sont libres finis. Dans ce cas, nous devons prouver l'égalité
$$
\det(I_n + AB) = \det(I_m + BA)
$$
des déterminants usuels de matrices où $A$ est de taille $n \times m$ et $B$
de taille $m \times n$. Cela se réduit au cas de l'anneau $R = \mathbf{Z}[a_{ij}, b_{ji}; 1 \leq i \leq n, 1 \leq j \leq m]$ où $a_{ij}$ et
$b_{ij}$ sont des variables et les entrées des matrices $A$ et $B$. En
prenant le corps des fractions, cela se réduit au cas d'un corps de
caractéristique zéro. En caractéristique zéro, il existe un polynôme universel
exprimant le déterminant d'une matrice de taille $\leq N$ en fonction des traces
des puissances de ladite matrice. Par conséquent, il suffit de prouver
$$
\text{Trace}((I_n + AB)^k) = \text{Trace}((I_m + BA)^k)
$$
pour tous les $k \geq 1$. En développant, nous voyons qu'il suffit de prouver
$\text{Trace}((AB)^k) = \text{Trace}((BA)^k)$ pour tous les $k \geq 0$. Pour $k = 1$, c'est le fait bien connu que
$\text{Trace}(AB) = \text{Trace}(BA)$. Pour $k > 1$, cela découle de cela en écrivant $(AB)^k = A(BA)^{k - 1}B$ et $(BA)^k = (BA)^{k - 1} A B$.
\end{proof}

\noindent
Rappelons que nous avons défini, dans Algèbre, Section \ref{algebra-section-K-groups}, un groupe $K_0(R)$
comme le groupe libre sur les classes d'isomorphisme de modules projectifs finis
$R$ modulo les relations $[M'] + [M''] = [M' \oplus M'']$.

\begin{lemma}
\label{lemma-det}
Soit $R$ un anneau. Il existe une application
$$
\det : K_0(R) \longrightarrow \Pic(R)
$$
qui envoie $[M]$ sur la classe du module inversible $\wedge^n(M)$ si $M$ est un
module localement libre fini de rang $n$.
\end{lemma}

\begin{proof}
Cela découle immédiatement des constructions ci-dessus ; en particulier, le Lemme
\ref{lemma-det-ses} montre que les relations sont envoyées sur $0$.
\end{proof}





\section{Complexes parfaits et groupes K}
\label{section-perfect-K-group}

\noindent
Nous montrons rapidement que le zéroième groupe K de la catégorie dérivée des
complexes parfaits d'un anneau $R$ est le même que $K_0(R)$ défini en Algèbre,
Section \ref{algebra-section-K-groups}.

\begin{lemma}
\label{lemma-perfect-to-K-group}
Soit $R$ un anneau. Il existe une application
$$
c : \text{complexes parfaits sur }R \longrightarrow K_0(R)
$$
avec les propriétés suivantes
\begin{enumerate}
\item $c(K[n]) = (-1)^nc(K)$ pour un complexe parfait $K$,
\item si $K \to L \to M \to K[1]$ est un triangle distingué de complexes parfaits, alors $c(L) = c(K) + c(M)$,
\item si $K$ est représenté par un complexe fini $M^\bullet$ consistant en
des modules finis projectifs, alors $c(K) = \sum (-1)^i[M_i]$.
\end{enumerate}
\end{lemma}

\begin{proof}
Soit $K$ un objet parfait de $D(R)$. Par définition, nous pouvons
représenter $K$ par un complexe fini $M^\bullet$ de $R$-modules
projectifs finis. Nous définissons $c$ en posant
$$
c(K) = \sum (-1)^n[M^n]
$$
en $K_0(R)$. Bien sûr, il faut montrer que c’est bien défini, mais une fois bien
défini, alors (1) et (3) sont immédiats. Pour l’instant, nous considérons
l’application $c$ telle que définie sur des complexes de $R$-modules
projectifs finis.

\medskip\noindent Supposons que $L^\bullet \to M^\bullet$ soit une application surjective de
complexes finis de $R$-modules projectifs finis. Soit $K^\bullet$ le noyau. On
obtient alors de courtes suites exactes de $R$-modules
$$
0 \to K^n \to L^n \to M^n \to 0
$$
qui sont scindées parce que $M^n$ est projectif. Ainsi, $K^\bullet$ est aussi un
complexe fini de $R$-modules projectifs finis et $c(L^\bullet) = c(K^\bullet) + c(M^\bullet)$ dans $K_0(R)$.

\medskip\noindent Supposons donné un complexe fini $M^\bullet$ de $R$-modules
projectifs finis, qui est acyclique. Supposons que $M^n = 0$ pour $n \not \in [a, b]$. On peut alors
décomposer $M^\bullet$ en courtes suites exactes
$$
\begin{matrix}
0 \to M^a \to M^{a + 1} \to N^{a + 1} \to 0, \\
0 \to N^{a + 1} \to M^{a + 2} \to N^{a + 3} \to 0, \\
\ldots \\
0 \to N^{b - 3} \to M^{b - 2} \to N^{b - 2} \to 0, \\
0 \to N^{b - 2} \to M^{b - 1} \to M^b \to 0
\end{matrix}
$$
En arguant par induction descendante, on voit que $N^{b - 2}, \ldots, N^{a + 1}$ sont des $R$-modules
projectifs finis, les suites sont exactes scindées, et
$$
c(M^\bullet) = \sum (-1)[M^n] = \sum (-1)^n([N^{n - 1}] + [N^n]) = 0
$$
Ainsi, notre construction donne zéro sur les complexes acycliques.

\medskip\noindent Il découle formellement des résultats des deux paragraphes
précédents que $c$ est bien défini et satisfait (2). À savoir, supposons que
les complexes finis $M^\bullet$ et $L^\bullet$ de $R$-modules projectifs
finis représentent le même objet de $D(R)$. Nous pouvons alors représenter
l’isomorphisme par une application $f : M^\bullet \to L^\bullet$ de complexes, voir Catégories dérivées,
Lemme \ref{derived-lemma-morphisms-from-projective-complex}. Nous obtenons une courte suite exacte de complexes
$$
0 \to L^\bullet \to C(f)^\bullet \to K^\bullet[1] \to 0
$$
voir Catégories dérivées, Définition \ref{derived-definition-cone}. Puisque $f$ est un
quasi-isomorphisme, le cône $C(f)^\bullet$ est acyclique (cela découle par exemple de la
discussion dans Catégories Dérivées, Section \ref{derived-section-canonical-delta-functor}). Ainsi
$$
0 = c(C(f)^\bullet) = c(L^\bullet) + c(K^\bullet[1]) =
c(L^\bullet) - c(K^\bullet)
$$
comme souhaité. Nous omettons la preuve de (2) qui est similaire.
\end{proof}

\noindent
Le lemme suivant montre que $K_0(R)$ est égal à $K_0(D_{perf}(R))$.

\begin{lemma}
\label{lemma-perfect-to-K-group-universal}
Soit $R$ un anneau. Soit $D_{perf}(R)$ la catégorie dérivée des objets parfaits,
voir le Lemme \ref{lemma-perfect-ring-classical-generator}. L'application $c$ du Lemme \ref{lemma-perfect-to-K-group} donne un isomorphisme
$K_0(D_{perf}(R)) = K_0(R)$.
\end{lemma}

\begin{proof}
Il résulte de la définition de $K_0(D_{perf}(R))$ (Catégories Dérivées, Définition \ref{derived-definition-K-zero})
que $c$ induit un homomorphisme $K_0(D_{perf}(R)) \to K_0(R)$.

\medskip\noindent Étant donné un module projectif fini $M$ sur $R$, notons
$M[0]$ le complexe parfait sur $R$ qui a $M$ en degré $0$ et zéro
dans les autres degrés. Étant donné une suite exacte courte $0 \to M \to M' \to M'' \to 0$ de modules
projectifs finis, nous obtenons un triangle distingué $M[0] \to M'[0] \to M''[0] \to M[1]$, voir Catégories
Dérivées, Section \ref{derived-section-canonical-delta-functor}. Cela montre que nous obtenons une application $K_0(R) \to K_0(D_{perf}(R))$
en envoyant $[M]$ sur $[M[0]]$, avec des excuses pour la notation épouvantable.

\medskip\noindent Il est clair que $K_0(R) \to K_0(D_{perf}(R)) \to K_0(R)$ est l'identité. D'autre part, si
$M^\bullet$ est un complexe borné de $R$-modules projectifs finis, alors
l'existence des triangles distingués des « troncatures stupides » (voir Homologie,
Section \ref{homology-section-truncations})
$$
\sigma_{\geq n}M^\bullet \to \sigma_{\geq n - 1}M^\bullet \to
M^{n - 1}[-n + 1] \to (\sigma_{\geq n}M^\bullet)[1]
$$
et l'induction montre que
$$
[M^\bullet] = \sum (-1)^i[M^i[0]]
$$
dans $K_0(D_{perf}(R))$ (avec encore des excuses pour la notation). Ainsi, l'application
$K_0(R) \to K_0(D_{perf}(R))$ est surjective, ce qui termine la démonstration.
\end{proof}
\section{Déterminants des endomorphismes de modules de longueur finie}
\label{section-determinants-finite-length}

\noindent
Soit $(R, \mathfrak m, \kappa)$ un anneau local. Considérons la catégorie des
paires $(M, \varphi)$ formées d'un $R$-module de longueur finie et d'un
endomorphisme $\varphi : M \to M$. Cette catégorie est abélienne, et chacun de
ses objets est à la fois artinien et noethérien. Pour les définitions, voir
Homologie, Section \ref{homology-section-jordan-holder}.

\medskip\noindent
Si $(M, \varphi)$ est un objet simple de cette catégorie, alors $M$ est annulé
par $\mathfrak m$, car sinon $(\mathfrak m M, \varphi|_{\mathfrak m M})$
serait un sous-objet non trivial. De plus,
$\dim_\kappa(M) = \text{length}_R(M)$ est fini. On peut donc définir le
déterminant et la trace
$$
\det\nolimits_\kappa(\varphi),\quad \text{Trace}_\kappa(\varphi)
$$
comme des éléments de $\kappa$ au moyen de l'algèbre linéaire. On définit de
même, dans ce cas, le polynôme caractéristique de $\varphi$.

\medskip\noindent
D'après Homologie, Lemme \ref{homology-lemma-finite-length}, tout objet
$(M, \varphi)$ de notre catégorie possède une filtration finie
$$
0 \subset M_1 \subset \ldots \subset M_n = M
$$
par des sous-modules stables par $\varphi$, telle que
$(M_i/M_{i - 1}, \varphi_i)$ soit un objet simple de la catégorie, où
$\varphi_i : M_i/M_{i - 1} \to M_i/M_{i - 1}$ est l'application induite.
On définit le {\it déterminant} de $(M, \varphi)$ sur $\kappa$ par
$$
\det\nolimits_\kappa(\varphi) = \prod \det\nolimits_\kappa(\varphi_i)
$$
où $\det_\kappa(\varphi_i)$ est défini comme au paragraphe précédent. On
définit la {\it trace} de $(M, \varphi)$ sur $\kappa$ par
$$
\text{Trace}_\kappa(\varphi) = \sum \text{Trace}_\kappa(\varphi_i)
$$
où $\text{Trace}_\kappa(\varphi_i)$ est défini comme au paragraphe précédent.
On peut définir de même le polynôme caractéristique de $\varphi$ sur $\kappa$
comme le produit des polynômes caractéristiques des $\varphi_i$ définis au
paragraphe précédent. Le théorème de Jordan--Hölder (Homologie, Lemme
\ref{homology-lemma-jordan-holder}) assure que ces définitions ne dépendent pas
des choix.

\begin{lemma}
\label{lemma-ses}
Soit $(R, \mathfrak m, \kappa)$ un anneau local. Soit
$0 \to (M, \varphi) \to (M', \varphi') \to (M'', \varphi'') \to 0$ une suite
exacte dans la catégorie considérée ci-dessus. Alors
$$
\det\nolimits_\kappa(\varphi') =
\det\nolimits_\kappa(\varphi)\det\nolimits_\kappa(\varphi''),\quad
\text{Trace}_\kappa(\varphi') = \text{Trace}_\kappa(\varphi) + 
\text{Trace}_\kappa(\varphi'')
$$
En outre, le polynôme caractéristique de $\varphi'$ sur $\kappa$ est le produit
des polynômes caractéristiques de $\varphi$ et de $\varphi''$.
\end{lemma}

\begin{proof}
La démonstration est laissée en exercice.
\end{proof}

\begin{lemma}
\label{lemma-multiplication}
Soit $(R, \mathfrak m, \kappa) \to (R', \mathfrak m', \kappa')$ un
homomorphisme local d'anneaux locaux. Supposons que $\kappa'/\kappa$ soit une
extension finie. Soit $u \in R'$. Alors, pour tout $R'$-module $M'$ de
longueur finie, on a
$$
\det\nolimits_\kappa(u : M' \to M') =
\text{Norm}_{\kappa'/\kappa}(u \bmod \mathfrak m')^m
$$
où $m = \text{length}_{R'}(M')$.
\end{lemma}

\begin{proof}
Remarquons que l'énoncé a un sens puisque
$\text{length}_R(M') = \text{length}_{R'}(M') [\kappa' : \kappa]$. Si
$M' = \kappa'$, l'égalité résulte de la définition de la norme comme
déterminant de l'opérateur linéaire donné par la multiplication par $u$. Dans
le cas général, on se ramène à ce cas en choisissant une filtration convenable
et en utilisant la multiplicativité du Lemme \ref{lemma-ses}. Nous omettons
certains détails.
\end{proof}

\begin{lemma}
\label{lemma-flat-base-change-det}
Soit $(R, \mathfrak m, \kappa) \to (R', \mathfrak m', \kappa')$ un
homomorphisme local plat d'anneaux locaux tel que
$m = \text{length}_{R'}(R'/\mathfrak mR') < \infty$. Pour tout
$(M, \varphi)$ comme ci-dessus, l'élément $\det_\kappa(\varphi)^m$ s'envoie
sur $\det_{\kappa'}(\varphi \otimes 1 : M \otimes_R R' \to M \otimes_R R')$
dans $\kappa'$.
\end{lemma}

\begin{proof}
La platitude de $R \to R'$ assure que les suites exactes comme celle du
Lemme \ref{lemma-ses} donnent, après changement de base, des suites exactes
sur $R'$. La multiplicativité du Lemme \ref{lemma-ses} permet donc de supposer
que $(M, \varphi)$ est un objet simple de notre catégorie (voir l'introduction
de cette section). Dans ce cas, $M$ est annulé par $\mathfrak m$. Choisissons
une filtration
$$
0 \subset I_1 \subset I_2 \subset \ldots \subset I_{m - 1} \subset
R'/\mathfrak mR'
$$
dont les quotients successifs sont isomorphes à $\kappa'$ comme $R'$-modules.
On obtient alors la filtration
$$
0 \subset
M \otimes_\kappa I_1 \subset
M \otimes_\kappa I_2 \subset
\ldots \subset
M \otimes_\kappa I_{m - 1} \subset
M \otimes_\kappa R'/\mathfrak mR' = M \otimes_R R'
$$
dont les quotients successifs sont isomorphes à $M \otimes_\kappa \kappa'$.
De plus, ces sous-modules sont stables par $\varphi \otimes 1$. Le
Lemme \ref{lemma-ses} donne
$$
\det\nolimits_{\kappa'}(\varphi \otimes 1 : M \otimes_R R' \to M \otimes_R R')
=
\det\nolimits_{\kappa'}(\varphi \otimes 1 :
M \otimes_\kappa \kappa' \to M \otimes_\kappa \kappa')^m =
\det\nolimits_\kappa(\varphi)^m
$$
La dernière égalité résulte de la compatibilité des déterminants des
applications linéaires aux extensions de corps. Ceci démontre le lemme.
\end{proof}







\section{Un anneau local régulier est factoriel}
\label{section-regular-local-UFD}

\noindent
Nous démontrons le résultat annoncé dans le titre de la section.

\begin{lemma}
\label{lemma-regular-local-Pic-zero}
Soit $R$ un anneau local régulier. Soit $f \in R$.
Alors $\Pic(R_f) = 0$.
\end{lemma}

\begin{proof}
Soit $L$ un $R_f$-module inversible. En particulier, $L$ est un
$R_f$-module fini. Il existe un $R$-module fini $M$ tel que
$M_f \cong L$ ; voir Algèbre, Lemme
\ref{algebra-lemma-construct-fp-module}.
La Proposition \ref{algebra-proposition-regular-finite-gl-dim} d'Algèbre
montre que $M$ possède une résolution libre finie $F_\bullet$ sur $R$.
Il s'ensuit que $L$ est quasi-isomorphe à un complexe fini de
$R_f$-modules {\it libres}. Le Lemme
\ref{lemma-perfect-to-K-group} donne donc
$[L_f] = n[R_f]$ dans $K_0(R_f)$ pour un certain $n \in \mathbf{Z}$.
En appliquant l'application du Lemme \ref{lemma-det},
on voit que $L$ est trivial.
\end{proof}

\begin{lemma}
\label{lemma-regular-local-UFD}
Un anneau local régulier est factoriel.
\end{lemma}

\begin{proof}
Rappelons qu'un anneau local régulier est intègre ; voir
Algèbre, Lemme \ref{algebra-lemma-regular-domain}.
Nous démontrons la propriété de factorisation unique par récurrence
sur la dimension de l'anneau local régulier $R$.
Si $\dim(R) = 0$, alors $R$ est un corps, donc en particulier un anneau
factoriel. Supposons $\dim(R) > 0$. Soit $x \in \mathfrak m$ avec
$x \not \in \mathfrak m^2$. Alors $R/(x)$ est régulier par Algèbre,
Lemme \ref{algebra-lemma-regular-ring-CM}, donc intègre par Algèbre,
Lemme \ref{algebra-lemma-regular-domain} ; ainsi $x$ est un élément premier.
Soit $\mathfrak p \subset R$ un idéal premier de hauteur $1$. Il faut
montrer que $\mathfrak p$ est principal ; voir Algèbre, Lemme
\ref{algebra-lemma-characterize-UFD-height-1}.
Nous pouvons supposer $x \not \in \mathfrak p$, car si
$x \in \mathfrak p$, alors $\mathfrak p = (x)$ et la preuve est achevée.
Pour tout idéal premier non maximal $\mathfrak q \subset R$, l'anneau local
$R_\mathfrak q$ est un anneau local régulier ; voir Algèbre, Lemme
\ref{algebra-lemma-localization-of-regular-local-is-regular}.
Par récurrence, $\mathfrak pR_\mathfrak q$ est principal. En particulier,
le $R_x$-module $\mathfrak p_x = \mathfrak pR_x \subset R_x$ est un
$R_x$-module de présentation finie dont la localisation en tout idéal premier est libre de
rang $1$. Le Lemme \ref{algebra-lemma-finite-projective} d'Algèbre montre
donc que $\mathfrak p_x$ est un $R_x$-module inversible. Par le Lemme
\ref{lemma-regular-local-Pic-zero}, on a $\mathfrak p_x = (y)$ pour un
certain $y \in R_x$. Nous pouvons écrire $y = x^e f$ avec
$f \in \mathfrak p$ et $e \in \mathbf{Z}$. Décomposons
$f = a_1 \ldots a_r$ en éléments irréductibles de $R$ (Algèbre,
Lemme \ref{algebra-lemma-factorization-exists}). Comme $\mathfrak p$ est
premier, $a_i \in \mathfrak p$ pour un certain $i$. Puisque
$\mathfrak p_x = (y)$ est premier et que $a_i | y$ dans $R_x$, il s'ensuit
que $\mathfrak p_x$ est engendré par $a_i$ dans $R_x$, c'est-à-dire que
l'image de $a_i$ dans $R_x$ est première. Comme $x$ est un élément premier,
le Lemme \ref{algebra-lemma-invert-prime-elements} d'Algèbre montre que
$a_i$ est premier dans $R$. Enfin, $(a_i) \subset \mathfrak p$ et
$\mathfrak p$ est de hauteur $1$ ; ainsi $(a_i) = \mathfrak p$, comme voulu.
\end{proof}

\begin{lemma}
\label{lemma-picard-group-generic-fibre-regular}
Soit $R$ un anneau de valuation, de corps des fractions $K$ et de corps
résiduel $\kappa$. Soit $R \to A$ un homomorphisme d'anneaux tel que
\begin{enumerate}
\item $A$ est local et $R \to A$ est local,
\item $A$ est plat et essentiellement de type fini sur $R$,
\item $A \otimes_R \kappa$ est régulier.
\end{enumerate}
Alors $\Pic(A \otimes_R K) = 0$.
\end{lemma}

\begin{proof}
Soit $L$ un $A \otimes_R K$-module inversible. En particulier, $L$ est
un module fini. Il existe un $A$-module fini $M$ tel que
$M \otimes_R K \cong L$ ; voir Algèbre, Lemme
\ref{algebra-lemma-construct-fp-module}. Nous pouvons supposer que $M$ est
sans torsion comme $R$-module. Alors $M$ est plat comme $R$-module (Lemme
\ref{lemma-valuation-ring-torsion-free-flat}). Le Lemme
\ref{lemma-flat-finite-type-valuation-ring-finite-presentation} montre que
$M$ est de présentation finie comme $A$-module et que $A$ est essentiellement
de présentation finie comme $R$-algèbre. Par le Lemme
\ref{lemma-structure-relatively-perfect}, $M$ est parfait relativement à
$R$ ; en particulier, $M$ est pseudo-cohérent comme $A$-module. Le Lemme
\ref{lemma-perfect-over-regular-local-ring} montre alors que $M$ est parfait ;
ainsi $M$ possède une résolution libre finie $F_\bullet$ sur $A$. Il s'ensuit que
$L$ est quasi-isomorphe à un complexe fini de $A \otimes_R K$-modules
{\it libres}. Le Lemme \ref{lemma-perfect-to-K-group} donne donc
$[L] = n[A \otimes_R K]$ dans $K_0(A \otimes_R K)$ pour un certain
$n \in \mathbf{Z}$. En appliquant l'application du Lemme \ref{lemma-det},
on voit que $L$ est trivial.
\end{proof}






\section{Déterminants des complexes}
\label{section-determinants-complexes}

\noindent
Dans la Section \ref{section-perfect-K-group}, nous avons vu qu'à un complexe
parfait $K$ sur un anneau $R$ est associée une classe d'isomorphisme de
$R$-modules inversibles, c'est-à-dire un élément de $\Pic(R)$. En fait,
de manière analogue à la Section \ref{section-determinants}, il existe
un foncteur
$$
\det :
\left\{
\begin{matrix}
\text{catégorie des complexes parfaits} \\
\text{les morphismes sont les isomorphismes}
\end{matrix}
\right\}
\longrightarrow
\left\{
\begin{matrix}
\text{catégorie des modules inversibles} \\
\text{les morphismes sont les isomorphismes}
\end{matrix}
\right\}
$$
De plus, si $(L, F)$ est un objet de la catégorie dérivée filtrée $DF(R)$
de $R$, dont la filtration est finie et dont les gradués sont des complexes
parfaits, il existe un isomorphisme canonique
$\det(\text{gr}L) \to \det(L)$. Voir \cite{determinant} pour l'exposé
original. Nous ajouterons ce matériau ultérieurement (insérer une référence
future).

\medskip\noindent
Pour le moment, nous présentons une construction ad hoc dans le cas des objets
parfaits $L$ de $D(R)$ d'amplitude de Tor contenue dans $[-1, 0]$. Un tel objet
peut être représenté par un complexe
$$
L^\bullet = \ldots \to 0 \to L^{-1} \to L^0 \to 0 \to \ldots
$$
où $L^{-1}$ et $L^0$ sont des $R$-modules projectifs finis ; voir le Lemme
\ref{lemma-perfect}. Dans ce cas, nous posons
$$
\det(L^\bullet) = \det(L^0) \otimes_R \det(L^{-1})^{\otimes -1} =
\Hom_R(\det(L^{-1}), \det(L^0))
$$
Nous dirons qu'un complexe de cette forme est de {\it rang $0$} si
$L^{-1}_\mathfrak p$ et $L^0_\mathfrak p$ ont le même rang pour tout idéal
premier de $R$. Si $L^\bullet$ est de rang $0$, la Section
\ref{section-determinants} fournit un élément canonique
$$
\delta(L^\bullet) \in \det(L^\bullet)
$$
qui est simplement le déterminant de $d : L^{-1} \to L^0$.
Remarquons que $\delta(L^\bullet)$ trivialise $\det(L^\bullet)$ si et
seulement si $L^\bullet$ est acyclique.

\medskip\noindent
Considérons un morphisme de complexes $a^\bullet : K^\bullet \to L^\bullet$
tel que
\begin{enumerate}
\item $a^\bullet$ est un quasi-isomorphisme,
\item $a^n : K^n \to L^n$ est surjectif pour tout $n$,
\item $K^n$ et $L^n$ sont des $R$-modules projectifs finis, non nuls seulement
pour
$n \in \{-1, 0\}$.
\end{enumerate}
Dans cette situation, nous construisons un isomorphisme
$$
\det(a^\bullet) : \det(K^\bullet) \longrightarrow \det(L^\bullet)
$$
Les suites exactes $0 \to \Ker(a^i) \to K^i \to L^i \to 0$ donnent des
isomorphismes
$$
\gamma^i : \det(\Ker(a^i)) \otimes \det(L^i) \to \det(K^i)
$$
pour $i = -1, 0$, par le Lemme \ref{lemma-det-ses}. Comme $a^\bullet$ est
un quasi-isomorphisme, le complexe $\Ker(a^\bullet)$ est acyclique et de
rang $0$. Par conséquent, l'élément canonique
$\delta(\Ker(a^\bullet))$ trivialise le $R$-module inversible
$\det(\Ker(a^\bullet))$ ; voir ci-dessus. Nous définissons
$\det(a^\bullet) : \det(K^\bullet) \to \det(L^\bullet)$ comme l'unique
isomorphisme rendant commutatif le diagramme
$$
\xymatrix{
\det(K^\bullet)
\ar[rr]_{\det(a^\bullet)}
& &
\det(L^\bullet)
\ar[ld]^{\delta(\Ker(a^\bullet))}
\\
& \det(L^\bullet) \otimes \det(\Ker(a^\bullet))
\ar[lu]^{\gamma^0 \otimes (\gamma^{-1})^{\otimes -1}}
}
$$
.

\begin{lemma}
\label{lemma-canonical-element-well-defined}
Soit $R$ un anneau. Soit $a^\bullet : K^\bullet \to L^\bullet$ un morphisme
de complexes de $R$-modules satisfaisant aux conditions (1), (2) et (3)
ci-dessus. Si $L^\bullet$ est de rang $0$, alors $\det(a^\bullet)$ envoie
l'élément canonique $\delta(K^\bullet)$ sur $\delta(L^\bullet)$.
\end{lemma}

\begin{proof}
Posons $M^i = \Ker(a^i)$. Nous avons donc un morphisme de suites exactes
courtes
$$
\xymatrix{
0 \ar[r] &
M^{-1} \ar[r] \ar[d]_{d_M} &
K^{-1} \ar[r] \ar[d]_{d_K} &
L^{-1} \ar[r] \ar[d]_{d_L} &
0 \\
0 \ar[r] &
M^0 \ar[r] &
K^0 \ar[r] &
L^0 \ar[r] &
0
}
$$
Le Lemme \ref{lemma-det-ses-functorial} montre que le morphisme
$\det(d_K)$ correspond à $\det(d_M) \otimes \det(d_L)$. En explicitant les
définitions, on obtient l'égalité voulue.
\end{proof}

\begin{lemma}
\label{lemma-homotopic-surjections}
Soit $R$ un anneau. Soit $a^\bullet : K^\bullet \to L^\bullet$ un morphisme
de complexes de $R$-modules satisfaisant aux conditions (1), (2) et (3)
ci-dessus. Soit $h : K^0 \to L^{-1}$ un morphisme tel que
$b^0 = a^0 + d \circ h$ et $b^{-1} = a^{-1} + h \circ d$ soient surjectifs.
Alors $\det(a^\bullet) = \det(b^\bullet)$ comme morphismes
$\det(K^\bullet) \to \det(L^\bullet)$.
\end{lemma}

\begin{proof}
Supposons qu'il existe un morphisme $\tilde h : K^0 \to K^{-1}$ tel que
$h = a^{-1} \circ \tilde h$ et tel que
$k^0 = \text{id} + d \circ \tilde h : K^0 \to K^0$ et
$k^1 = \text{id} + \tilde h \circ d : K^{-1} \to K^{-1}$ soient des
isomorphismes. Nous obtenons alors un diagramme commutatif
$$
\xymatrix{
0 \ar[r] &
\Ker(b^\bullet) \ar[r] \ar[d]_{c^\bullet} &
K^\bullet \ar[r]_{b^\bullet}
\ar[d]_{k^\bullet} &
L^\bullet \ar[r] \ar[d]^{\text{id}} &
0 \\
0 \ar[r] &
\Ker(a^\bullet) \ar[r] &
K^\bullet \ar[r]^{a^\bullet} &
L^\bullet \ar[r] &
0
}
$$
de complexes, où $c^\bullet$ est l'isomorphisme induit sur les noyaux.
Le Lemme \ref{lemma-det-ses-functorial} montre que
$$
\xymatrix{
\det(\Ker(b^i)) \otimes \det(L^i) \ar[r] \ar[d]_{\det(c^i) \otimes 1} &
\det(K^i) \ar[d]^{\det(k^i)} \\
\det(\Ker(a^i)) \otimes \det(L^i) \ar[r] &
\det(K^i)
}
$$
est commutatif. Comme $\det(c^\bullet)$ envoie la trivialisation canonique
de $\det(\Ker(a^\bullet))$ sur la trivialisation canonique de
$\Ker(b^\bullet)$ (Lemme \ref{lemma-canonical-element-well-defined}),
la conclusion équivaut à l'égalité
$$
\det(k^0) = \det(k^{-1})
$$
dans $R$, qui résulte du Lemme \ref{lemma-det-switch}.

\medskip\noindent
Supposons qu'il existe un facteur direct $U \subset K^{-1}$ tel que
$a^{-1}|_U : U \to L^{-1}$ et $b^{-1}|_U : U \to L^{-1}$ soient tous deux
des isomorphismes. Définissons $\tilde h$ comme la composée de $h$ avec
l'inverse de $a^{-1}|_U$. Nous affirmons que $\tilde h$ est un morphisme
comme dans le premier paragraphe de la preuve. En effet,
$h = a^{-1} \circ \tilde h$ par construction. Pour montrer que
$k^{-1} : K^{-1} \to K^{-1}$ est un isomorphisme, il suffit de montrer qu'il
est surjectif (Algèbre, Lemme \ref{algebra-lemma-fun}). Soit $u \in U$.
Nous pouvons choisir $u' \in U$ tel que $b^{-1}(u') = a^{-1}(u)$.
Alors $u = k^{-1}(u')$. En effet, $u$ et $k^{-1}(u')$ appartiennent à $U$,
et un calcul donne $a^{-1}(u) = a^{-1}(k^{-1}(u'))$\footnote{
$a^{-1}(k^{-1}(u')) =
a^{-1}(u') + a^{-1}(\tilde h(d(u'))) =
a^{-1}(u') + h(d(u')) = b^{-1}(u') = a^{-1}(u)$}
Comme $a^{-1}|_U$ est un isomorphisme, nous obtenons l'égalité. Ainsi
$U \subset \Im(k^{-1})$. D'autre part, si $x \in \Ker(a^{-1})$, alors
$x = k^{-1}(x) \bmod U$. Comme $K^{-1} = \Ker(a^{-1}) + U$, il s'ensuit
que $k^{-1}$ est surjectif. Montrons enfin que $k^0 : K^0 \to K^0$ est
surjectif. Puisque $a^0 \circ k^0 = b^0$, le morphisme
$a^0 \circ k^0$ est surjectif. Si $x \in \Ker(a^0)$, alors
$x = d(y)$ pour un certain $y \in \Ker(a^{-1})$. D'après ce qui précède,
nous pouvons écrire $y = k^{-1}(z)$ avec $z \in K^{-1}$. Alors
$x = k^0(d(z))$, ce qui conclut.

\medskip\noindent
Dernière étape de la preuve. Il suffit de trouver $U$ comme au paragraphe
précédent, mais cela n'est pas toujours possible. Toutefois, pour montrer
l'égalité de deux morphismes de $R$-modules, il suffit de la vérifier après
localisation aux idéaux premiers de $R$. Nous pouvons donc supposer $R$ local.
Le problème devient alors le suivant : soient
$$
\alpha, \beta : R^{\oplus n} \longrightarrow R^{\oplus m}
$$
deux applications $R$-linéaires surjectives. Il faut trouver un facteur
direct $U \subset R^{\oplus n}$ tel que $\alpha|_U$ et $\beta|_U$ soient
des isomorphismes. Si $R$ est un corps, cela est possible par algèbre linéaire.
En général, on choisit une solution sur le corps résiduel puis on la
relève en une solution sur l'anneau local $R$. Certains détails sont omis.
\end{proof}

\begin{lemma}
\label{lemma-compose-surjections}
Soit $R$ un anneau. Soient $a^\bullet : K^\bullet \to L^\bullet$ et
$b^\bullet : L^\bullet \to M^\bullet$ des morphismes de complexes de
$R$-modules satisfaisant aux conditions (1), (2) et (3) ci-dessus. Alors
$\det(b^\bullet) \circ \det(a^\bullet) =
\det(b^\bullet \circ a^\bullet)$ comme morphismes
$\det(M^\bullet) \to \det(K^\bullet)$.
\end{lemma}

\begin{proof}
Omis. Indications : cela résulte directement des Lemmes
\ref{lemma-det-ses}, \ref{lemma-det-ses-functorial} et
\ref{lemma-det-filtration}.
\end{proof}

\begin{lemma}
\label{lemma-determinant-two-term-complexes}
Soit $R$ un anneau. Les constructions précédentes déterminent un foncteur
$$
\det :
\left\{
\begin{matrix}
\text{catégorie des complexes parfaits} \\
\text{d'amplitude de Tor contenue dans }[-1, 0] \\
\text{les morphismes sont les isomorphismes}
\end{matrix}
\right\}
\longrightarrow
\left\{
\begin{matrix}
\text{catégorie des modules inversibles} \\
\text{les morphismes sont les isomorphismes}
\end{matrix}
\right\}
$$
De plus, si l'on se donne un objet parfait de rang $0$, noté $L$, de $D(R)$, d'amplitude de
Tor contenue dans $[-1, 0]$, il existe un élément canonique
$\delta(L) \in \det(L)$ tel que, pour tout isomorphisme $a : L \to K$
dans $D(R)$, on ait $\det(a)(\delta(L)) = \delta(K)$.
\end{lemma}

\begin{proof}
Par le Lemme \ref{lemma-perfect}, tout objet de la catégorie source peut être
représenté par un complexe
$$
L^\bullet = \ldots \to 0 \to L^{-1} \to L^0 \to 0 \to \ldots
$$
où $L^{-1}$ et $L^0$ sont des $R$-modules projectifs finis. Appelons
provisoirement {\it bon} un complexe de ce type. Par Catégories dérivées,
Lemme \ref{derived-lemma-morphisms-from-projective-complex}, les morphismes
entre bons complexes dans la catégorie dérivée sont les classes d'homotopie
de morphismes de complexes. Nous pouvons donc travailler avec les bons
complexes et employer le déterminant
$\det(L^\bullet) = \det(L^0) \otimes \det(L^{-1})^{\otimes -1}$
étudié ci-dessus.

\medskip\noindent
Soit $a^\bullet : L^\bullet \to K^\bullet$ un morphisme de bons complexes
qui est un isomorphisme dans $D(R)$, c'est-à-dire un quasi-isomorphisme.
Nous dirons que
$$
\xymatrix{
L^\bullet \ar[rr]_{a^\bullet} & & K^\bullet \\
& M^\bullet \ar[lu]^{b^\bullet} \ar[ru]_{c^\bullet}
}
$$
est un bon diagramme s'il commute à homotopie près et si $b^\bullet$ et
$c^\bullet$ satisfont aux conditions (1), (2) et (3) ci-dessus. Pour un tel
diagramme, on peut définir
$$
\det(a^\bullet) = \det(c^\bullet) \circ \det(b^\bullet)^{-1}
$$
où $\det(c^\bullet)$ et $\det(b^\bullet)$ sont les isomorphismes construits
ci-dessus. Nous montrerons que les bons diagrammes existent toujours et que
le morphisme $\det(a^\bullet)$ ainsi obtenu ne dépend pas du bon diagramme
choisi.

\medskip\noindent
Existence de bons diagrammes pour un quasi-isomorphisme
$a^\bullet : L^\bullet \to K^\bullet$ de bons complexes. Choisissons une
surjection $p : R^{\oplus n} \to K^{-1}$. Nous pouvons alors considérer le
nouveau bon complexe
$$
M^\bullet = \ldots \to 0 \to
L^{-1} \oplus R^{\oplus n} \xrightarrow{d \oplus 1}
L^0 \oplus R^{\oplus n} \to 0 \to \ldots
$$
muni du morphisme de projection $b^\bullet : M^\bullet \to L^\bullet$ et du
morphisme $c^\bullet : M^\bullet \to K^\bullet$ donné par
$a^{-1} \oplus p$ en degré $-1$ et par $a^0 \oplus d \circ p$ en degré $0$.
Les morphismes $b^\bullet : M^\bullet \to L^\bullet$ et
$c^\bullet : M^\bullet \to K^\bullet$ satisfont aux conditions (1), (2) et
(3) ci-dessus, ce qui fournit un bon diagramme.

\medskip\noindent
Supposons donné un bon diagramme
$$
\xymatrix{
L^\bullet \ar[rr]_{\text{id}^\bullet} & & L^\bullet \\
& M^\bullet \ar[lu]^{b^\bullet} \ar[ru]_{c^\bullet}
}
$$
Le Lemme \ref{lemma-homotopic-surjections} donne alors
$\det(c^\bullet) = \det(b^\bullet)$. Ainsi,
$\det(\text{id}^\bullet) = \text{id}$ indépendamment du bon diagramme choisi.

\medskip\noindent
Avant de démontrer l'indépendance en général, examinons la composition.
Supposons donnés des quasi-isomorphismes de bons complexes
$L_1^\bullet \to L_2^\bullet$ et $L_2^\bullet \to L_3^\bullet$, ainsi que de
bons diagrammes
$$
\vcenter{
\xymatrix{
L_1^\bullet \ar[rr] & &
L_2^\bullet \\
& M_{12}^\bullet \ar[lu] \ar[ru]
}
}
\quad\text{et}\quad
\vcenter{
\xymatrix{
L_2^\bullet \ar[rr] & &
L_3^\bullet \\
& M_{23}^\bullet \ar[lu] \ar[ru]
}
}
$$
Nous pouvons les compléter en un diagramme
$$
\xymatrix{
L_1^\bullet \ar[rr] & &
L_2^\bullet \ar[rr] & &
L_3^\bullet \\
& M_{12}^\bullet \ar[lu] \ar[ru]
& & M_{23}^\bullet \ar[lu] \ar[ru] \\
& &
M_{123}^\bullet \ar[lu] \ar[ru]
}
$$
où $M_{123}^\bullet \to M_{12}^\bullet$ et
$M_{123}^\bullet \to M_{23}^\bullet$ possèdent les propriétés (1), (2) et
(3), et où le carré est commutatif : il suffit de prendre
$M_{123}^n = M_{12}^n \times_{L_2^n} M_{23}^n$.
Le Lemme \ref{lemma-compose-surjections} montre alors que
$$
\xymatrix{
\det(L_2^\bullet) &
\det(M_{23}^\bullet) \ar[l] \\
\det(M_{12}^\bullet) \ar[u] &
\det(M_{123}^\bullet) \ar[l] \ar[u]
}
$$
est commutatif. Une chasse au diagramme montre que la composée
$\det(L_1^\bullet) \to \det(L_2^\bullet) \to \det(L_3^\bullet)$ des
morphismes associés aux deux bons diagrammes utilisant $M_{12}^\bullet$ et
$M_{23}^\bullet$ est égale au morphisme associé au bon diagramme
$$
\xymatrix{
L_1^\bullet \ar[rr] & & L_3^\bullet \\
& M_{123}^\bullet \ar[lu] \ar[ru]
}
$$
Par conséquent, si ces morphismes sont indépendants des choix, la loi de
composition est elle aussi satisfaite et nous obtenons le foncteur recherché.

\medskip\noindent
Indépendance. Soit $a^\bullet : L^\bullet \to K^\bullet$ un
quasi-isomorphisme de bons complexes. Choisissons un quasi-isomorphisme
inverse $b^\bullet : K^\bullet \to L^\bullet$. En posant
$L_1^\bullet = L$, $L_2^\bullet = K^\bullet$ et
$L_3^\bullet = L^\bullet$, nous pouvons fixer notre choix d'un bon diagramme
pour $b^\bullet$ et faire varier les bons diagrammes pour $a^\bullet$.
Les paragraphes précédents montrent alors que, quels que soient les choix,
la composée est toujours l'identité de $\det(L^\bullet)$. Cela établit
clairement l'indépendance à l'égard de ces choix.

\medskip\noindent
L'assertion concernant les éléments canoniques résulte immédiatement du
Lemme \ref{lemma-canonical-element-well-defined} et de notre construction.
\end{proof}







\section{Extensions d'anneaux de valuation}
\label{section-valuation-rings}

\noindent
Cette section est l'analogue de la Section
\ref{section-discrete-valuation-rings} pour les anneaux de valuation généraux.

\begin{definition}
\label{definition-extension-valuation-rings}
Nous dirons que $A \to B$, ou $A \subset B$, est une
{\it extension d'anneaux de valuation} si $A$ et $B$ sont des anneaux de
valuation et si $A \to B$ est injectif et local. Une telle extension induit
un diagramme commutatif
$$
\xymatrix{
A \setminus \{0\} \ar[r] \ar[d]_v & B \setminus \{0\} \ar[d]^v \\
\Gamma_A \ar[r] & \Gamma_B
}
$$
où $\Gamma_A$ et $\Gamma_B$ sont les groupes de valeurs. Nous dirons que $B$
est {\it faiblement non ramifié} sur $A$ si la flèche horizontale inférieure
est bijective. Si l'extension de corps résiduels
$\kappa_A = A/\mathfrak m_A \subset \kappa_B = B/\mathfrak m_B$ est finie,
nous posons $f = [\kappa_B : \kappa_A]$ et l'appelons le
{\it degré résiduel} de l'extension $A \subset B$.
\end{definition}

\noindent
Remarquons que $\Gamma_A \to \Gamma_B$ est injectif, car les unités de $A$
sont les images réciproques des unités de $B$ par l'application $A \to B$.
Remarquons aussi que nous n'exigeons pas que l'extension des corps des
fractions soit finie.

\begin{lemma}
\label{lemma-inequality-general}
Soit $A \subset B$ une extension d'anneaux de valuation de corps des
fractions respectifs $K \subset L$. Si l'extension $L/K$ est finie, alors
l'extension des corps résiduels est finie, l'indice de $\Gamma_A$ dans
$\Gamma_B$ est fini, et
$$
[\Gamma_B : \Gamma_A] [\kappa_B : \kappa_A] \leq [L : K].
$$
\end{lemma}

\begin{proof}
Soient $b_1, \ldots, b_n \in B$ des unités dont les images dans $\kappa_B$
sont linéairement indépendantes sur $\kappa_A$. Soient
$c_1, \ldots, c_m \in B$ des éléments non nuls dont les images dans
$\Gamma_B/\Gamma_A$ sont deux à deux distinctes. Nous affirmons que les
$b_i c_j$ sont linéairement indépendants sur $K$ dans $L$. En effet, une
somme
$$
\sum a_{ij} b_i c_j
$$
où les $a_{ij} \in K$ ne sont pas tous nuls ne peut être nulle. Choisissons
$(i_0, j_0)$ tel que $v(a_{i_0j_0}b_{i_0}c_{j_0})$ soit minimal. Remplaçons
$a_{ij}$ par $a_{ij}/a_{i_0j_0}$, de sorte que $a_{i_0 j_0} = 1$. Posons
$$
P = \{(i, j) \mid
v(a_{ij}b_ic_j) = v(a_{i_0j_0}b_{i_0}c_{j_0}) \}
$$
Par le choix des $c_1, \ldots, c_m$, la relation $(i, j) \in P$ implique
$j = j_0$. Ainsi, si $(i, j) \in P$, alors
$v(a_{ij}) = v(a_{i_0j_0}) = 0$, c'est-à-dire que $a_{ij}$ est une unité.
Par le choix des $b_1, \ldots, b_n$, la somme
$$
\sum\nolimits_{(i, j) \in P} a_{ij}b_i
$$
est une unité de $B$. La valuation de
$\sum\nolimits_{(i, j) \in P} a_{ij}b_ic_j$ est donc
$v(c_{j_0}) = v(a_{i_0j_0}b_{i_0}c_{j_0})$. Comme les termes d'indices
$(i, j) \not \in P$ de la première somme affichée ont une valuation
strictement plus grande, cette somme ne peut être nulle, ce qui démontre
le lemme.
\end{proof}

\begin{lemma}
\label{lemma-valuation-ring-purely-inseparable}
Soit $A$ un anneau de valuation de corps des fractions $K$ de caractéristique
$p > 0$. Soit $L/K$ une extension purement inséparable. Alors la clôture
intégrale $B$ de $A$ dans $L$ est un anneau de valuation de corps des
fractions $L$, et $A \subset B$ est une extension d'anneaux de valuation.
\end{lemma}

\begin{proof}
Omis. Indication : utiliser, par exemple, les Lemmes d'Algèbre
\ref{algebra-lemma-x-or-x-inverse-valuation-ring} and
\ref{algebra-lemma-integral-overring-surjective} for example.
\end{proof}

\begin{lemma}
\label{lemma-extension-normal-domains-and-roots}
Soit $A \to B$ un homomorphisme local plat d'anneaux locaux noethériens
normaux et intègres. Soient $f \in A$ et $h \in B$ tels que
$f = w h^n$ pour un entier $n > 1$ et une unité $w$ de $B$. Supposons que,
pour tout idéal premier de hauteur $1$, $\mathfrak p \subset A$, il existe
un idéal premier de hauteur $1$, $\mathfrak q \subset B$, au-dessus de
$\mathfrak p$ tel que l'extension
$A_\mathfrak p \subset B_\mathfrak q$ soit faiblement non ramifiée. Alors
$f = u g^n$ pour un certain $g \in A$ et une unité $u$ de $A$.
\end{lemma}

\begin{proof}
Les anneaux locaux de $A$ et $B$ en leurs idéaux premiers de hauteur $1$
sont des anneaux de valuation discrète (Algèbre, Lemme
\ref{algebra-lemma-characterize-dvr}). L'hypothèse a donc un sens (par la
Définition \ref{definition-extension-discrete-valuation-rings}). Soient
$\mathfrak p_1, \ldots, \mathfrak p_r$ les idéaux premiers de $A$ minimaux
parmi ceux contenant $f$. Ils sont de hauteur $1$ par Algèbre, Lemme
\ref{algebra-lemma-minimal-over-1}. Pour chaque $i$, soient
$\mathfrak q_{i, j} \subset B$, $j = 1, \ldots, r_i$, les idéaux premiers
de hauteur $1$ de $B$ au-dessus de $\mathfrak p_i$. Numérotons-les de sorte
que $A_{\mathfrak p_i} \to B_{\mathfrak q_{i, 1}}$ soit faiblement non
ramifiée. Puisque l'image de $f$ est le produit d'une puissance $n$-ième et
d'une unité dans $B_{\mathfrak q_{i, 1}}$, la valuation $v_i$ de $f$ dans
$A_{\mathfrak p_i}$ est divisible par $n$. Écrivons
$v_i = n w_i$ avec $w_i \geq 0$. Considérons la suite exacte
$$
0 \to I \to A \to
\prod\nolimits_{i = 1, \ldots, r}
A_{\mathfrak p_i}/\mathfrak p_i^{w_i}A_{\mathfrak p_i}
$$
qui définit l'idéal $I$. En appliquant le foncteur exact $- \otimes_A B$,
nous obtenons
une suite exacte
$$
0 \to I \otimes_A B \to B \to
\prod\nolimits_{i = 1, \ldots, r}
(A_{\mathfrak p_i}/\mathfrak p_i^{w_i}A_{\mathfrak p_i}) \otimes_A B
$$
Fixons $i$. Nous affirmons que l'application canonique
$$
(A_{\mathfrak p_i}/\mathfrak p_i^{w_i}A_{\mathfrak p_i}) \otimes_A B
\to
\prod\nolimits_{j = 1, \ldots, r_i}
B_{\mathfrak q_{i, j}}/\mathfrak q_{i, j}^{e_{i, j}w_i}B_{\mathfrak q_{i, j}}
$$
est injective. Ici, $e_{i, j}$ est l'indice de ramification de
$A_{\mathfrak p_i} \to B_{\mathfrak q_{i, j}}$. L'affirmation signifie que
$\mathfrak p_i^{w_i}B_{\mathfrak p_i}$ est l'ensemble des éléments $b$ de
$B_{\mathfrak p_i}$ dont la valuation en $\mathfrak q_{i, j}$ est
$\geq e_{i, j}w_i$. Choisissons un générateur $a \in A_{\mathfrak p_i}$
de l'idéal principal $\mathfrak p_i^{w_i}$. La valuation de $a$ en
$\mathfrak q_{i, j}$ est alors égale à $e_{i, j}w_i$. Or
$B_{\mathfrak p_i}$ est un anneau normal intègre dont les idéaux premiers de
hauteur un sont les $\mathfrak q_{i, j}$, $j = 1, \ldots, r_i$ ; pour $b$
comme ci-dessus, le Lemme d'Algèbre
\ref{algebra-lemma-normal-domain-intersection-localizations-height-1} donne
donc $b/a \in B_{\mathfrak p_i}$. D'où l'affirmation.

\medskip\noindent
L'affirmation, combinée avec la seconde suite exacte ci-dessus, fournit une
suite exacte
$$
0 \to I \otimes_A B \to B \to
\prod\nolimits_{i = 1, \ldots, r}
\prod\nolimits_{j = 1, \ldots, r_i}
B_{\mathfrak q_{i, j}}/\mathfrak q_{i, j}^{e_{i, j}w_i}B_{\mathfrak q_{i, j}}
$$
Il s'ensuit que $I \otimes_A B$ est l'ensemble des éléments $h'$ de $B$
dont la valuation est $\geq e_{i, j}w_i$ en $\mathfrak q_{i, j}$.
Puisque $f = wh^n$ dans $B$, la valuation de $h$ vaut
$e_{i, j}w_i$ en $\mathfrak q_{i, j}$. Par le Lemme d'Algèbre
\ref{algebra-lemma-normal-domain-intersection-localizations-height-1}, on a
donc $h'/h \in B$. Ainsi, $I \otimes_A B$ est un $B$-module libre de rang
$1$, engendré par $h$. Par conséquent, $I$ est un $A$-module libre de rang
$1$ ; voir Algèbre, Lemme \ref{algebra-lemma-finite-projective-descends}.
Soit $g \in I$ un générateur. Alors $g$ et $h$ diffèrent par une unité de
$B$. En remontant le raisonnement, on conclut que la valuation de $g$ dans
$A_{\mathfrak p_i}$ est $w_i = v_i/n$. Ainsi, $g^n$ et $f$ diffèrent par
une unité de $A$ (par Algèbre, Lemme
\ref{algebra-lemma-normal-domain-intersection-localizations-height-1}),
comme voulu.
\end{proof}

\begin{lemma}
\label{lemma-etale-extension-valuation-ring}
Soit $A$ un anneau de valuation. Soit $A \to B$ un homomorphisme d'anneaux
étale et soit $\mathfrak m \subset B$ un idéal premier au-dessus de l'idéal
maximal de $A$. Alors $A \subset B_\mathfrak m$ est une extension d'anneaux
de valuation faiblement non ramifiée.
\end{lemma}

\begin{proof}
Le Lemme \ref{lemma-weak-dimension-at-most-1} donne à $A$ une dimension
faible $\leq 1$. Les Lemmes \ref{lemma-weak-dimension-goes-up} et
\ref{lemma-when-weakly-etale} donnent alors à $B$ une dimension faible
$\leq 1$ ; l'anneau local $B_\mathfrak m$ est donc un anneau de valuation
par le Lemme \ref{lemma-weak-dimension-at-most-1}. Comme l'extension
$A \subset B_\mathfrak m$ induit une extension finie des corps des
fractions, $\Gamma_A$ est d'indice fini dans le groupe de valeurs de
$B_{\mathfrak m}$. Ainsi, pour tout $h \in B_\mathfrak m$, il existe
$n > 0$, un élément $f \in A$ et une unité $w \in B_\mathfrak m$ tels que
$f = w h^n$ dans $B_\mathfrak m$. Montrons que cela implique
$f = ug^n$ pour un certain $g \in A$ et une unité $u \in A$ ; cela montrera
que les groupes de valeurs de $A$ et $B_\mathfrak m$ coïncident, comme
l'affirme le lemme.

\medskip\noindent
Écrivons $A = \colim A_i$ comme colimite de ses sous-anneaux locaux qui sont
essentiellement de type fini sur $\mathbf{Z}$. Comme $A$ est un anneau normal
intègre (Algèbre, Lemme \ref{algebra-lemma-valuation-ring-normal}), nous
pouvons supposer chaque $A_i$ normal (nous utilisons ici le fait qu'après
normalisation les anneaux locaux restent essentiellement de type fini sur
$\mathbf{Z}$, par Algèbre, Proposition
\ref{algebra-proposition-ubiquity-nagata}). Pour un certain $i$, il existe une
extension étale $A_i \to B_i$ telle que $B = A \otimes_{A_i} B_i$ ; voir
Algèbre, Lemme \ref{algebra-lemma-etale}. Soit $\mathfrak m_i$ l'intersection
de $B_i$ avec $\mathfrak m$. Nous pouvons alors conclure en appliquant le
Lemme \ref{lemma-extension-normal-domains-and-roots} à l'homomorphisme
$A_i \to (B_i)_{\mathfrak m_i}$. Ses hypothèses sont satisfaites car :
\begin{enumerate}
\item $A_i$ et $(B_i)_{\mathfrak m_i}$ sont noethériens, puisqu'ils sont
essentiellement de type fini sur $\mathbf{Z}$,
\item $A_i \to (B_i)_{\mathfrak m_i}$ est plat, puisque $A_i \to B_i$ est
étale,
\item $B_i$ est normal, puisque $A_i \to B_i$ est étale ; voir Algèbre,
Lemme \ref{algebra-lemma-normal-goes-up},
\item pour tout idéal premier de hauteur $1$ de $A_i$, il existe un idéal
premier de hauteur $1$ de $(B_i)_{\mathfrak m_i}$ au-dessus de lui, par
Algèbre, Lemme \ref{algebra-lemma-finite-in-codim-1}, et parce que
$\Spec((B_i)_{\mathfrak m_i}) \to \Spec(A_i)$ est surjectif,
\item les extensions induites
$(A_i)_\mathfrak p \to (B_i)_\mathfrak q$ sont non ramifiées pour tout idéal
premier $\mathfrak q$ au-dessus d'un idéal premier $\mathfrak p$, puisque
$A_i \to B_i$ est étale.
\end{enumerate}
Cela achève la preuve du lemme.
\end{proof}

\begin{lemma}
\label{lemma-henselization-valuation-ring}
Soit $A$ un anneau de valuation. Soient $A^h$, resp.\ $A^{sh}$, son
hensélisé, resp.\ son hensélisé strict. Alors
$$
A \subset A^h \subset A^{sh}
$$
sont des extensions d'anneaux de valuation qui induisent des bijections sur
les groupes de valeurs, c'est-à-dire qui sont faiblement non ramifiées.
\end{lemma}

\begin{proof}
Écrivons $A^h = \colim (B_i)_{\mathfrak q_i}$, où $A \to B_i$ est étale et
$\mathfrak q_i \subset B_i$ est un idéal premier au-dessus de
$\mathfrak m_A$ ; voir Algèbre, Lemme
\ref{algebra-lemma-henselization-different}. Le Lemme
\ref{lemma-etale-extension-valuation-ring} montre alors que
$(B_i)_{\mathfrak q_i}$ est un anneau de valuation et que l'application
induite
$$
(A \setminus \{0\})/A^* \longrightarrow
((B_i)_{\mathfrak q_i} \setminus \{0\}) / (B_i)_{\mathfrak q_i}^*
$$
est bijective. Le Lemme \ref{algebra-lemma-colimit-valuation-rings} d'Algèbre
montre que $A^h$ est un anneau de valuation. Il s'ensuit également que
$(A \setminus \{0\})/A^* \to (A^h \setminus \{0\})/(A^h)^*$
est bijective. Cela démontre le lemme pour l'inclusion $A \subset A^h$.
Pour $A \subset A^{sh}$, on emploie exactement le même argument, en
remplaçant le Lemme \ref{algebra-lemma-henselization-different} d'Algèbre par
le Lemme \ref{algebra-lemma-strict-henselization-different} d'Algèbre.
Comme $A^{sh} = (A^h)^{sh}$, cela démontre aussi les assertions du lemme
pour l'inclusion $A^h \subset A^{sh}$.
\end{proof}










\section{Structure des modules sur un anneau principal}
\label{section-structure-modules}

\noindent
Nous travaillons dans un cadre un peu plus général (en suivant les articles
\cite{Warfield-Purity} et \cite{Warfield-Decomposition}
de Warfield), de sorte que les démonstrations fonctionnent sur les anneaux de valuation.

\begin{lemma}
\label{lemma-characterize-PD-modules}
\begin{reference}
\cite[Corollary 1]{Warfield-Purity}
\end{reference}
Soit $P$ un module sur un anneau $R$. Les assertions suivantes sont équivalentes
\begin{enumerate}
\item $P$ est un facteur direct d'une somme directe de modules de la
forme $R/fR$, où $f \in R$ varie.
\item pour toute suite exacte courte $0 \to A \to B \to C \to 0$
de $R$-modules telle que $fA = A \cap fB$ pour tout $f \in R$,
l'application $\Hom_R(P, B) \to \Hom_R(P, C)$ est surjective.
\end{enumerate}
\end{lemma}

\begin{proof}
Soit $0 \to A \to B \to C \to 0$ une suite exacte comme en (2).
Pour prouver que (1) implique (2), il suffit de prouver que
$\Hom_R(R/fR, B) \to \Hom_R(R/fR, C)$ est surjective pour tout $f \in R$.
Soit $\psi : R/fR \to C$ une application. Supposons que $\psi(1)$ soit l'image
de $b \in B$. Alors $fb \in A$. Il existe donc un $a \in A$
tel que $fa = fb$. Alors $f(b - a) = 0$ ; on obtient donc un morphisme
$\varphi : R/fR \to B$ envoyant $1$ sur $b - a$ qui relève $\psi$.

\medskip\noindent
Réciproquement, supposons que (2) soit vérifiée. Soit $I$ l'ensemble des couples
$(f, \varphi)$ où $f \in R$ et $\varphi : R/fR \to P$. Pour
$i \in I$, notons $(f_i, \varphi_i)$ le couple correspondant.
Considérons l'application
$$
B = \bigoplus\nolimits_{i \in I} R/f_iR \longrightarrow P
$$
qui envoie l'élément $r$ du facteur $R/f_iR$ sur $\varphi_i(r)$ dans $P$.
Soit $A = \Ker(B \to P)$. On voit alors que (1) est vraie si la suite
$$
0 \to A \to B \to P \to 0
$$
est une suite exacte comme en (2). Pour le voir, supposons que $f \in R$ et que
$a \in A$ ait pour image $f b$ dans $B$. Écrivons $b = (r_i)_{i \in I}$, avec
presque tous les $r_i = 0$. On voit alors que
$$
f\sum \varphi_i(r_i) = 0
$$
dans $P$. Il existe donc un $i_0 \in I$ tel que $f_{i_0} = f$ et
$\varphi_{i_0}(1) = \sum \varphi_i(r_i)$. Soit $x_{i_0} \in R/f_{i_0}R$
la classe de $1$. On voit alors que
$$
a' = (r_i)_{i \in I} - (0, \ldots, 0, x_{i_0}, 0, \ldots )
$$
est un élément de $A$ et que $fa' = a$, comme souhaité.
\end{proof}

\begin{lemma}[Anneaux de valuation généralisés]
\label{lemma-generalized-valuation-ring}
\begin{reference}
\cite{Warfield-Decomposition}
\end{reference}
Soit $R$ un anneau non nul. Les assertions suivantes sont équivalentes
\begin{enumerate}
\item Pour $a, b \in R$, soit $a$ divise $b$, soit $b$ divise $a$.
\item Tout idéal de type fini est principal et $R$ est local.
\item L'ensemble des idéaux de $R$ est totalement ordonné par inclusion.
\end{enumerate}
Ceci vaut en particulier si $R$ est un anneau de valuation.
\end{lemma}

\begin{proof}
Supposons (2) et soient $a, b \in R$. Alors $(a, b) = (c)$. Si $c = 0$,
alors $a = b = 0$ et $a$ divise $b$. Supposons $c \not = 0$. Écrivons
$c = ua + vb$, $a = wc$ et $b = zc$. Alors $c(1 - uw - vz) = 0$.
Puisque $R$ est local, cela implique que $1 - uw - vz \in \mathfrak m$.
Ainsi, $w$ ou $z$ est une unité, et donc soit $a$ divise $b$, soit
$b$ divise $a$. Ainsi, (2) implique (1).

\medskip\noindent
Supposons (1). Si $R$ possède deux idéaux maximaux $\mathfrak m_i$,
nous pouvons choisir $a \in \mathfrak m_1$ avec $a \not \in \mathfrak m_2$
et $b \in \mathfrak m_2$ avec $b \not \in \mathfrak m_1$.
Alors $a$ ne divise pas $b$ et $b$ ne divise pas $a$.
Ainsi, $R$ possède un unique idéal maximal et est local.
Il résulte aisément de la condition (1) et d'une récurrence que tout
idéal de type fini est principal. Ainsi, (1) implique (2).

\medskip\noindent
Il est immédiat de prouver que (1) et (3) sont équivalentes.
La dernière assertion est Algèbre, Lemme
\ref{algebra-lemma-valuation-ring-x-or-x-inverse}.
\end{proof}

\begin{lemma}
\label{lemma-generalized-valuation-ring-modules}
\begin{reference}
\cite[Theorem 1]{Warfield-Decomposition}
\end{reference}
Soit $R$ un anneau vérifiant les conditions équivalentes du
Lemme \ref{lemma-generalized-valuation-ring}.
Alors tout $R$-module de présentation finie
est isomorphe à une somme directe finie de modules de la forme $R/fR$.
\end{lemma}

\begin{proof}
Soit $M$ un $R$-module de présentation finie. Nous utiliserons toutes
les propriétés équivalentes de $R$ données dans le
Lemme \ref{lemma-generalized-valuation-ring}
sans plus de précisions. Notons $\mathfrak m \subset R$
l'idéal maximal et $\kappa = R/\mathfrak m$ le corps résiduel.
Soit $I \subset R$ l'annulateur de $M$.
Choisissons une base $y_1, \ldots, y_n$ du
$\kappa$-espace vectoriel de dimension finie $M/\mathfrak m M$. Nous allons raisonner
par récurrence sur $n$.

\medskip\noindent
Par le lemme de Nakayama, toute famille d'éléments $x_1, \ldots, x_n \in M$
relevant les éléments $y_1, \ldots, y_n$ de $M/\mathfrak m M$
engendre $M$ ; voir Algèbre, Lemme \ref{algebra-lemma-NAK}.
Ceci prouve immédiatement le cas initial $n = 0$ de la récurrence.

\medskip\noindent
Nous affirmons qu'il existe un indice $i$ tel que, pour tout choix
de $x_i \in M$ s'envoyant sur $y_i$, l'annulateur de $x_i$ soit $I$.
En effet, sinon, nous pourrions choisir $x_1, \ldots, x_n$ tels
que $I_i = \text{Ann}(x_i) \not = I$ pour tout $i$. Mais comme
$I \subset I_i$ pour tout $i$, le fait que les idéaux soient totalement ordonnés
implique que $I_i$ est strictement plus grand que $I$ pour $i = 1, \ldots, n$,
et, par un nouvel usage de l'ordre total, on verrait que
$\text{Ann}(M) = I_1 \cap \ldots \cap I_n$ est plus grand que $I$,
ce qui est une contradiction. Après renumérotation, nous pouvons supposer que
$y_1$ a la propriété suivante : pour tout $x_1 \in M$ relevant $y_1$,
l'annulateur de $x_1$ est $I$.

\medskip\noindent
Posons $A = Rx_1 \subset M$. Considérons la suite exacte
$0 \to A \to M \to M/A \to 0$. Puisque $A$ est de type fini, on voit que
$M/A$ est un $R$-module de présentation finie
(Algèbre, Lemme \ref{algebra-lemma-extension}) ayant moins de générateurs.
Ainsi, $M/A \cong \bigoplus_{j = 1, \ldots, m} R/f_jR$ par récurrence.
D'autre part, nous affirmons que $A \to M$ vérifie la propriété suivante :
si $f \in R$, alors $fA = A \cap fM$. L'inclusion $fA \subset A \cap fM$
est triviale. Réciproquement, si $x \in A \cap fM$, alors $x = gx_1 = f y$
pour certains $g \in R$ et $y \in M$. Si $f$ divise $g$, alors $x \in fA$,
comme souhaité. Sinon, nous pouvons écrire $f = hg$ pour un certain $h \in \mathfrak m$.
L'élément $x'_1 = x_1 - hy$ a pour annulateur $I$ d'après le paragraphe
précédent. Ainsi, $g \in I$ et l'on voit que $x = 0$, comme souhaité.
L'affirmation et le Lemme \ref{lemma-characterize-PD-modules}
impliquent que la suite $0 \to A \to M \to M/A \to 0$ est scindée,
et l'on obtient $M \cong A \oplus \bigoplus_{j = 1, \ldots, m} R/f_jR$.
Alors $A = R/I$ est de présentation finie (en tant que facteur direct de $M$),
et donc $I$ est de type fini, donc principal.
Ceci achève la démonstration.
\end{proof}

\begin{lemma}
\label{lemma-warfield}
\begin{reference}
\cite[Theorem 3]{Warfield-Decomposition}
\end{reference}
Soit $R$ un anneau tel que toute localisation de $R$ en un idéal
maximal vérifie les conditions équivalentes du
Lemme \ref{lemma-generalized-valuation-ring}.
Alors tout $R$-module de présentation finie est un facteur direct d'une 
somme directe finie de modules de la forme $R/fR$, où $f$ varie dans $R$.
\end{lemma}

\begin{proof}
Soit $M$ un $R$-module de présentation finie. Nous montrons d'abord que $M$ est un
facteur direct d'une somme directe de modules de la forme $R/fR$, puis, à la fin, nous
montrons que la somme directe peut être choisie finie. Soit
$$
0 \to A \to B \to C \to 0
$$
une suite exacte courte de $R$-modules telle que $fA = A \cap fB$
pour tout $f \in R$. D'après le Lemme \ref{lemma-characterize-PD-modules},
nous devons montrer que $\Hom_R(M, B) \to \Hom_R(M, C)$ est surjective.
Il suffit de le prouver après localisation aux idéaux maximaux
$\mathfrak m$ ; voir
Algèbre, Lemme \ref{algebra-lemma-characterize-zero-local}.
Notons que les suites localisées
$0 \to A_\mathfrak m \to B_\mathfrak m \to C_\mathfrak m \to 0$
vérifient la condition $fA_\mathfrak m = A_\mathfrak m \cap fB_\mathfrak m$
pour tout $f \in R_\mathfrak m$ (car nous pouvons écrire $f = uf'$ avec
$u \in R_\mathfrak m$ une unité et $f' \in R$, et parce que la localisation
est exacte). Puisque $M$ est de présentation finie, on a
$$
\Hom_R(M, B)_\mathfrak m = \Hom_{R_\mathfrak m}(M_\mathfrak m, B_\mathfrak m)
\quad\text{et}\quad
\Hom_R(M, C)_\mathfrak m = \Hom_{R_\mathfrak m}(M_\mathfrak m, C_\mathfrak m)
$$
d'après Algèbre, Lemme \ref{algebra-lemma-hom-from-finitely-presented}.
Le module $M_\mathfrak m$ est un $R_\mathfrak m$-module de présentation finie. D'après le
Lemme \ref{lemma-generalized-valuation-ring-modules},
on voit que $M_\mathfrak m$ est une somme directe de modules
de la forme $R_\mathfrak m/fR_\mathfrak m$. Nous concluons donc du
Lemme \ref{lemma-characterize-PD-modules} que l'application entre les
localisations est surjective.

\medskip\noindent
À ce stade, nous savons que $M$ est un facteur direct de
$\bigoplus_{i \in I} R/f_i R$. Considérons l'application
$M \to \bigoplus_{i \in I} R/f_i R$. Puisque $M$ est un $R$-module de type fini,
l'image est contenue dans $\bigoplus_{i \in I'} R/f_i R$ pour un certain sous-ensemble
fini $I' \subset I$. Ceci achève la démonstration.
\end{proof}

\begin{definition}
\label{definition-bezout}
Soit $R$ un anneau intègre.
\begin{enumerate}
\item Nous disons que $R$ est un {\it anneau de B\'ezout} si tout idéal de type fini
de $R$ est principal.
\item Nous disons que $R$ est un {\it anneau à diviseurs élémentaires} si, pour
tous $n , m \geq 1$ et toute matrice $n \times m$ $A$, il
existe des matrices inversibles $U, V$ de tailles $n \times n, m \times m$
telles que
$$
U A V =
\left(
\begin{matrix}
f_1 & 0 & 0 & \ldots \\
0 & f_2 & 0 & \ldots \\
0 & 0 & f_3 & \ldots \\
\ldots & \ldots & \ldots & \ldots
\end{matrix}
\right)
$$
avec $f_1, \ldots, f_{\min(n, m)} \in R$ et $f_1 | f_2 | \ldots$.
\end{enumerate}
\end{definition}

\noindent
Il semble que la question de savoir si tout anneau de B\'ezout
$R$ est un anneau à diviseurs élémentaires (ou non) soit encore ouverte. Ceci équivaut à la
question de savoir si tout module de présentation finie sur $R$ est une somme
directe de modules cycliques. L'implication réciproque est vraie.

\begin{lemma}
\label{lemma-elementary-divisor-is-bezout}
Un anneau à diviseurs élémentaires est de B\'ezout.
\end{lemma}

\begin{proof}
Soient $a, b \in R$ non nuls. Considérons la matrice $1 \times 2$ $A = (a\ b)$.
On voit alors que $u(a\ b)V = (f\ 0)$ avec $u \in R$ inversible
et $V = (g_{ij})$ une matrice inversible $2 \times 2$.
Alors $f = u a g_{11} + u b g_{2 1}$ et $(g_{11}, g_{2 1}) = R$.
Il en résulte que $(a, b) = (f)$. Un raisonnement par récurrence (omis)
montre alors que tout idéal de type fini de $R$ est engendré par un élément.
\end{proof}

\begin{lemma}
\label{lemma-localize-bezout}
Toute localisation d'un anneau de B\'ezout est de B\'ezout.
Tout anneau local d'un anneau de B\'ezout est un anneau de valuation.
Un anneau local intègre est de B\'ezout si et seulement s'il est un anneau de valuation.
\end{lemma}

\begin{proof}
Nous omettons la démonstration de l'assertion sur les localisations. La dernière
assertion est Algèbre, Lemme \ref{algebra-lemma-characterize-valuation-ring}.
La deuxième assertion résulte des deux autres.
\end{proof}

\begin{lemma}
\label{lemma-split-off-free-part}
Soit $R$ un anneau de B\'ezout.
\begin{enumerate}
\item Tout sous-module de type fini d'un module libre est libre de type fini.
\item Tout $R$-module de présentation finie $M$ est somme directe d'un
module libre de type fini et d'un module de torsion $M_{tors}$ qui est un
facteur direct d'un module de la forme $\bigoplus_{i = 1, \ldots, n} R/f_iR$,
où $f_1, \ldots, f_n \in R$ sont non nuls.
\end{enumerate}
\end{lemma}

\begin{proof}
Démonstration de (1). Soit $M \subset F$ un sous-module de type fini d'un module libre $F$.
Puisque $M$ est de type fini, nous pouvons supposer que $F$ est libre de type fini
(détails omis). Posons $F = R^{\oplus n}$. Nous raisonnons par récurrence
sur $n$. Si $n = 1$, alors $M$ est un idéal de type fini, donc
principal puisque, par hypothèse, $R$ est de B\'ezout. Si $n > 1$, alors
nous considérons l'image $I$ de $M$ par la projection
$R^{\oplus n} \to R$ sur le dernier facteur. Si $I = (0)$, alors
$M \subset R^{\oplus n - 1}$ et nous concluons par récurrence.
Si $I \not = 0$, alors $I = (f) \cong R$. Ainsi,
$M \cong R \oplus \Ker(M \to I)$ et nous concluons encore par récurrence.

\medskip\noindent
Soit $M$ un $R$-module de présentation finie. Puisque les localisations
de $R$ aux idéaux maximaux sont des anneaux de valuation
(Lemme \ref{lemma-localize-bezout}),
nous pouvons appliquer le Lemme \ref{lemma-warfield}.
Ainsi, $M$ est un facteur direct d'un module de la forme
$R^{\oplus r} \oplus \bigoplus_{i = 1, \ldots, n} R/f_iR$,
avec $f_i \not = 0$. Puisque la prise du sous-module de torsion est
un foncteur, on voit que $M_{tors}$ est
un facteur direct du module $\bigoplus_{i = 1, \ldots, n} R/f_iR$,
et que $M/M_{tors}$ est un facteur direct de $R^{\oplus r}$.
D'après la première partie de la démonstration, on voit que $M/M_{tors}$ est
libre de type fini. Ainsi, $M \cong M_{tors} \oplus M/M_{tors}$,
comme souhaité.
\end{proof}

\begin{lemma}
\label{lemma-modules-PID}
Soit $R$ un anneau principal. Tout $R$-module de type fini $M$ est isomorphe
à un module de la forme
$$
R^{\oplus r} \oplus \bigoplus\nolimits_{i = 1, \ldots, n} R/f_iR
$$
pour certains $r, n \geq 0$ et certains $f_1, \ldots, f_n \in R$ non nuls.
\end{lemma}

\begin{proof}
Un anneau principal est un anneau de B\'ezout noethérien. D'après le Lemme \ref{lemma-split-off-free-part},
il suffit de prouver le résultat lorsque $M$ est de torsion. Puisque $M$ est de type fini, cela
signifie que l'annulateur de $M$ est non nul. Supposons que $fM = 0$ pour un certain
$f \in R$ non nul. Nous pouvons alors considérer $M$ comme un module sur $R/fR$.
Puisque $R/fR$ est noethérien de dimension $0$ (petit détail omis),
on voit que $R/fR = \prod R_j$ est un produit fini d'anneaux artiniens
locaux $R_j$
(Algèbre, Proposition \ref{algebra-proposition-dimension-zero-ring}).
Écrivons $R_j = R/f_jR$ avec $f_j \in R$.
Chaque $R_j$, étant un anneau local et un quotient d'un anneau principal, est un anneau de
valuation généralisé au sens du
Lemme \ref{lemma-generalized-valuation-ring} (petit détail omis).
Écrivons $M = \prod M_j$, où $M_j = e_j M$ et où $e_j \in R/fR$ est
l'idempotent correspondant au facteur $R_j$.
D'après le Lemme \ref{lemma-generalized-valuation-ring-modules},
on voit que $M_j = \bigoplus_{i = 1, \ldots, n_j} R_j/\overline{f}_{ji}R_j$
pour certains $\overline{f}_{ji} \in R_j$. Choisissons des relèvements $f_{ji} \in R$
et choisissons $g_{ji} \in R$ tel que $(g_{ji}) = (f_j, f_{ji})$.
On en conclut que
$$
M \cong \bigoplus R/g_{ji}R
$$
comme $R$-module, ce qui achève la démonstration.
\end{proof}

\noindent
On peut également prouver qu'un anneau principal est un anneau à diviseurs élémentaires (insérer ici une
référence ultérieure), en démontrant des lemmes semblables au suivant.

\begin{lemma}
\label{lemma-unimodular-vector}
Soit $R$ un anneau de B\'ezout. Soient $n \geq 1$ et $f_1, \ldots, f_n \in R$
qui engendrent l'idéal unité. Il existe une matrice inversible $n \times n$ à
coefficients dans $R$ dont la première ligne est $f_1 \ldots f_n$.
\end{lemma}

\begin{proof}
Ceci résulte du Lemme \ref{lemma-split-off-free-part},
mais nous pouvons aussi le prouver directement comme suit.
Raisonnons par récurrence sur $n$. Le résultat vaut pour $n = 1$. Supposons $n > 1$.
Après renumérotation, nous pouvons supposer que $f_1 \not = 0$.
Choisissons $f \in R$ tel que $(f) = (f_1, \ldots, f_{n - 1})$.
Soit $A$ une matrice $(n - 1) \times (n - 1)$ dont la première ligne
est $f_1/f, \ldots, f_{n - 1}/f$. Choisissons $a, b \in R$ tels que
$af - bf_n = 1$, ce qui est possible puisque
$1 \in (f_1, \ldots, f_n) = (f, f_n)$. Une solution est alors
la matrice
$$
\left(
\begin{matrix}
f & 0 & \ldots & 0 & f_n \\
0 & 1 & \ldots & 0 & 0 \\
  &   & \ldots \\
0 & 0 & \ldots & 1 & 0 \\
b & 0 & \ldots & 0 & a
\end{matrix}
\right)
\left(
\begin{matrix}
  & & & 0 \\
  & A \\
  & & & 0 \\
0 & \ldots & 0 & 1
\end{matrix}
\right)
$$
Observons que la matrice de gauche est inversible car son déterminant vaut $1$.
\end{proof}
 
  
   
    
     
      
\section{Idéaux radicaux principaux}
\label{section-principal-radical-ideals}

\noindent
Dans cette section, nous démontrons que dans tout domaine intègre local normal
noethérien caténaire, il existe un idéal radical principal non trivial. Ce résultat se
trouve dans \cite{Artin-Lipman}.

\begin{lemma}
\label{lemma-polypoly}
Soit $(R,\mathfrak m)$ un anneau local noethérien de dimension un, et
soit $x\in\mathfrak m$ un élément n'appartenant à aucun idéal premier minimal
de $R$. Alors
\begin{enumerate}
\item la fonction $P : n \mapsto \text{longueur}_R(R/x^n R)$
satisfait $P(n) \leq n P(1)$ pour $n \geq 0$,
\item si $x$ est un non-diviseur de zéro, alors $P(n) = nP(1)$ pour $n \geq 0$.
\end{enumerate}
\end{lemma}

\begin{proof}
Comme $\dim(R) = 1$, on a $\dim(R/x^n R) = 0$
et donc $\text{longueur}_R(R/x^n R)$ est finie pour tout $n$
(Algèbre, Lemme \ref{algebra-lemma-support-point}).
Nous allons démontrer le lemme par récurrence sur $n$. Comme $x^0 R = R$, on a
$P(0) = \text{longueur}_R(R/x^0R) = \text{longueur}_R 0 = 0$.
L'assertion est aussi vraie pour $n = 1$.
Soit maintenant $n \geq 2$, et supposons l'assertion vraie pour $n - 1$.
La suite suivante est exacte
$$
R/x^{n-1}R \xrightarrow{x} R/x^nR \to R/xR \to 0
$$
où $x$ désigne l'application de multiplication par $x$.
Comme la longueur est additive
(Algèbre, Lemme \ref{algebra-lemma-length-additive}),
on a $P(n) \leq P(n - 1) + P(1)$. Par récurrence,
$P(n - 1) \leq (n - 1)P(1)$, d'où $P(n) \leq nP(1)$.
Ceci établit l'étape de récurrence.

\medskip\noindent
Si $x$ est un non-diviseur de zéro, alors la suite exacte affichée
ci-dessus est aussi exacte à gauche. Par conséquent, on obtient
$P(n) = P(n - 1) + P(1)$ pour tout $n \geq 1$.
\end{proof}

\begin{lemma}
\label{lemma-minprimespoly}
Soit $(R, \mathfrak m)$ un anneau local noethérien de dimension $1$.
Soit $x \in \mathfrak m$ un élément n'appartenant à aucun idéal premier
minimal de $R$. Soit $t$ le nombre d'idéaux premiers minimaux de $R$.
Alors $t \leq \text{longueur}_R(R/xR)$.
\end{lemma}

\begin{proof}
Soient $\mathfrak p_1, \ldots, \mathfrak p_t$ les idéaux premiers minimaux
de $R$. Posons $R' = R/\sqrt{0} = R/(\bigcap_{i = 1}^t \mathfrak p_i)$.
Nous affirmons qu'il suffit de démontrer le lemme pour $R'$. En effet, il est clair
que $R'$ possède lui aussi $t$ idéaux premiers minimaux
et que $\text{longueur}_{R'}(R'/xR') = \text{longueur}_R(R'/xR')$
est inférieure à $\text{longueur}_R(R/xR)$ puisqu'il existe une surjection
$R/xR \to R'/xR'$. On peut donc supposer que $R$ est réduit.

\medskip\noindent
Supposons que $R$ est réduit, d'idéaux premiers minimaux
$\mathfrak p_1, \ldots, \mathfrak p_t$.
Cela signifie qu'il existe une suite exacte
$$
0 \to R \to
\prod\nolimits_{i = 1}^t R/\mathfrak p_i \to Q \to 0
$$
Ici, $Q$ est le conoyau de la première application.
Posons $M = \prod_{i = 1}^t R/\mathfrak p_i$.
En localisant en $\mathfrak p_j$, on voit que
$$
R_{\mathfrak p_j} \to M_{\mathfrak p_j} =
\left(\prod\nolimits_{i=1}^t R/\mathfrak p_i\right)_{\mathfrak p_j} =
(R/\mathfrak p_j)_{\mathfrak p_j}
$$
est surjective. Ainsi, $Q_{\mathfrak p_j} = 0$ pour tout $j$.
On en déduit que $\text{Supp}(Q) = \{\mathfrak m\}$, puisque $\mathfrak m$
est le seul idéal premier de $R$ distinct des $\mathfrak p_i$.
Il s'ensuit que $Q$ est de longueur finie
(Algèbre, Lemme \ref{algebra-lemma-support-point}).
Comme $\text{Supp}(Q) = \{\mathfrak m\}$, on
peut choisir $n \gg 0$ tel que $x^n$
agisse comme $0$ sur $Q$
(Algèbre, Lemme \ref{algebra-lemma-Noetherian-power-ideal-kills-module}).
Considérons maintenant le diagramme
$$
\xymatrix{
0 \ar[r] & R \ar[r] \ar[d]^-{x^n} & M
\ar[r] \ar[d]^-{x^n} & Q \ar[r] \ar[d]^-{x^n} & 0 \\
0 \ar[r] & R \ar[r] & M \ar[r] & Q \ar[r] & 0
}
$$
où les applications verticales sont la multiplication par $x^n$. Celle-ci est injective sur
$R$ et sur $M$, car $x$ n'appartient à aucun des $\mathfrak p_i$.
Par le lemme du serpent (Algèbre, Lemme \ref{algebra-lemma-snake}),
la suite suivante est exacte :
$$
0 \to Q \to R/x^nR \to M/x^nM \to Q \to 0
$$
On obtient donc $\text{longueur}_R(R/x^nR) = \text{longueur}_R(M/x^nM)$
pour $n$ assez grand. En posant $R_i = R/\mathfrak p_i$, on voit
que $\text{longueur}(M/x^nM) =
\sum_{i = 1}^t \text{longueur}_R(R_i/x^nR_i)$.
En appliquant le Lemme \ref{lemma-polypoly} et en utilisant que $x$ est un non-diviseur de zéro
sur $R$ et sur $R_i$, on conclut que
$$
n \text{longueur}_R(R/xR) =
\sum\nolimits_{i = 1}^t n \text{longueur}_{R_i}(R_i/x R_i)
$$
Comme $\text{longueur}_{R_i}(R_i/x R_i) \geq 1$, le lemme est démontré.
\end{proof}

\begin{lemma}
\label{lemma-sopexists}
Soit $(R,\mathfrak m)$ un anneau local noethérien de dimension $d > 1$,
soit $f \in \mathfrak m$ un élément n'appartenant à aucun idéal premier minimal
de $R$, et soit $k\in\mathbf{N}$. Il existe alors des éléments
$g_1, \ldots, g_{d - 1} \in \mathfrak m^k$ tels que
$f, g_1, \ldots, g_{d - 1}$ soit un système de paramètres.
\end{lemma}

\begin{proof}
On a $\dim(R/fR) = d - 1$ d'après
Algèbre, Lemme \ref{algebra-lemma-one-equation}.
Choisissons un système de paramètres
$\overline{g}_1, \ldots, \overline{g}_{d - 1}$
dans $R/fR$ (Algèbre, Proposition \ref{algebra-proposition-dimension})
et choisissons des relèvements $g_1, \ldots, g_{d - 1}$ dans $R$.
On vérifie immédiatement que
$f, g_1, \ldots, g_{d - 1}$ est un système de paramètres de $R$.
Alors $f, g_1^k, \ldots, g_{d - 1}^k$ est aussi un système de
paramètres, ce qui achève la démonstration.
\end{proof}

\begin{lemma}
\label{lemma-syspar}
Soit $(R,\mathfrak m)$ un anneau local noethérien de dimension
deux, et soit $f \in \mathfrak m$ un élément n'appartenant à
aucun idéal premier minimal de $R$. Il existe alors
$g \in \mathfrak m$ et $N \in \mathbf{N}$ tels que
\begin{enumerate}
\item[(a)] $f,g$ forment un système de paramètres de $R$.
\item[(b)] Si $h \in \mathfrak m^N$, alors $f + h, g$ est un
système de paramètres et
$\text{longueur}_R (R/(f, g)) = \text{longueur}_R(R/(f + h, g))$.
\end{enumerate}
\end{lemma}

\begin{proof}
D'après le Lemme \ref{lemma-sopexists}, il existe $g \in \mathfrak m$
tel que $f, g$ soit un système de paramètres de $R$.
Alors $\mathfrak m = \sqrt{(f, g)}$. Il existe donc $n$
tel que $\mathfrak m^n \subset (f, g)$, voir
Algèbre, Lemme \ref{algebra-lemma-Noetherian-power}. Nous affirmons que
$N = n + 1$ convient.
En effet, soit $h \in \mathfrak m^N$. Par le choix de $N$, on peut écrire
$h = af + bg$ avec $a, b \in \mathfrak m$. Ainsi
$$
(f + h, g) = (f + af + bg, g) = ((1 + a)f, g) = (f, g)
$$
car $1 + a$ est inversible dans $R$. Cela démontre l'égalité
des longueurs et le fait que $f + h, g$ est un système de paramètres.
\end{proof}

\begin{lemma}
\label{lemma-radical-element}
Soit $R$ un domaine intègre local normal noethérien de dimension $2$.
Soient $\mathfrak p_1, \ldots, \mathfrak p_r$ des idéaux premiers distincts deux à deux
de hauteur $1$. Il existe un élément non nul
$f \in \mathfrak p_1 \cap \ldots \cap \mathfrak p_r$ tel
que $R/fR$ soit réduit.
\end{lemma}

\begin{proof}
Soit $f \in \mathfrak p_1 \cap \ldots \cap \mathfrak p_r$ un élément non nul.
Nous allons modifier légèrement $f$ afin d'obtenir un élément qui engendre un idéal radical.
Le localisé $R_\mathfrak p$ de $R$ en chaque idéal premier de hauteur un
$\mathfrak p$ est un anneau de valuation discrète, voir
Algèbre, Lemme \ref{algebra-lemma-characterize-dvr} ou
Algèbre, Lemme \ref{algebra-lemma-criterion-normal}.
Nous notons $\text{ord}_\mathfrak p(f)$ la
valuation correspondante de $f$ dans $R_{\mathfrak p}$. Soient
$\mathfrak q_1, \ldots, \mathfrak q_s$
les idéaux premiers distincts de hauteur un contenant $f$.
Écrivons $\text{ord}_{\mathfrak q_j}(f) = m_j \geq 1$ pour tout $j$.
Nous définissons alors $\text{div}(f) = \sum_{j = 1}^s m_j\mathfrak q_j$
comme une combinaison linéaire formelle d'idéaux premiers
de hauteur un à coefficients entiers.
Remarquons pour un usage ultérieur que chacun des idéaux premiers $\mathfrak p_i$
figure parmi les idéaux premiers $\mathfrak q_j$.
L'anneau $R/fR$ est réduit si et seulement si
$m_j = 1$ pour $j = 1, \ldots, s$. En effet, si $m_j$ vaut $1$, alors
$(R/fR)_{\mathfrak q_j}$ est réduit et
$R/fR \subset \prod (R/fR)_{\mathfrak q_j}$ puisque
$\mathfrak q_1, \ldots, \mathfrak q_s$ sont les idéaux premiers associés
de $R/fR$, voir les lemmes d'Algèbre
\ref{algebra-lemma-zero-at-ass-zero} et
\ref{algebra-lemma-normal-domain-intersection-localizations-height-1}.

\medskip\noindent
Choisissons et fixons $g$ et $N$ comme dans le Lemme \ref{lemma-syspar}.
Pour $y \in R$ non nul, notons $t(y)$ le nombre d'idéaux premiers minimaux au-dessus de $y$.
Comme $R$ est un domaine intègre normal, ces idéaux premiers
sont de hauteur un et correspondent $1$-à-$1$ aux idéaux premiers minimaux de
$R/yR$ (Algèbre, Lemmes \ref{algebra-lemma-minimal-over-1} et
\ref{algebra-lemma-normal-domain-intersection-localizations-height-1}).
Par exemple, $t(f) = s$ est le nombre
d'idéaux premiers $\mathfrak q_j$ apparaissant dans $\text{div}(f)$.
Soit $h \in \mathfrak m^N$. D'après le Lemme \ref{lemma-minprimespoly}, on a
\begin{align*}
t(f + h) & \leq \text{longueur}_{R/(f + h)}(R/(f + h, g)) \\
& = \text{longueur}_R(R/(f + h, g)) \\
& = \text{longueur}_R(R/(f, g))
\end{align*}
voir Algèbre, Lemme \ref{algebra-lemma-length-independent}
pour la première égalité.
On voit donc que $t(f + h)$ est borné indépendamment de
$h \in \mathfrak m^N$.

\medskip\noindent
Grâce à la bornitude démontrée ci-dessus, on peut choisir
$h \in \mathfrak m^N \cap \mathfrak p_1 \cap \ldots \cap \mathfrak p_r$
tel que $t(f + h)$ soit maximal parmi de tels $h$. Posons $f' = f + h$.
Étant donné $h' \in \mathfrak m^N \cap \mathfrak p_1 \cap \ldots \cap \mathfrak p_r$,
on voit que le nombre $t(f' + h') \leq t(f + h)$.
Ainsi, après avoir remplacé $f$ par $f'$, on peut supposer que pour tout
$h \in \mathfrak m^N \cap \mathfrak p_1 \cap \ldots \cap \mathfrak p_r$,
on a $t(f + h) \leq s$.

\medskip\noindent
Supposons ensuite que nous puissions trouver un élément $h \in \mathfrak m^N$ tel que
pour tout $j$, on ait $\text{ord}_{\mathfrak q_j}(h) \geq 1$ et
$\text{ord}_{\mathfrak q_j}(h) = 1 \Leftrightarrow m_j > 1$.
Remarquons que
$h \in \mathfrak m^N \cap \mathfrak p_1 \cap \ldots \cap \mathfrak p_r$.
Alors $\text{ord}_{\mathfrak q_j}(f + h) = 1$
pour tout $j$, d'après les propriétés élémentaires des valuations.
Ainsi
$$
\text{div}(f + h) = \sum\nolimits_{j = 1}^s \mathfrak q_j +
\sum\nolimits_{k = 1}^v e_k \mathfrak r_k
$$
pour certains idéaux premiers de hauteur un distincts deux à deux
$\mathfrak r_1, \ldots, \mathfrak r_v$ et $e_k \geq 1$.
Cependant, comme $s = t(f) \geq t(f + h)$, on voit que $v = 0$
et nous avons trouvé l'élément voulu.

\medskip\noindent
Choisissons maintenant $h$ satisfaisant les critères ci-dessus.
Par évitement des idéaux premiers (Algèbre, Lemme \ref{algebra-lemma-silly}),
pour tout $1 \leq j \leq s$, on peut trouver un élément $a_j \in \mathfrak q_j$
tel que $a_j \not \in \mathfrak q_{j'}$ pour $j' \not = j$
et $a_j \not \in \mathfrak q_j^{(2)}$. Ici,
$\mathfrak q_j^{(2)} = \{x \in R \mid \text{ord}_{\mathfrak q_j}(x) \geq 2\}$
est la deuxième puissance symbolique de $\mathfrak q_j$.
Prenons alors
$$
h = \prod\nolimits_{m_j = 1} a_j^2 \times
\prod\nolimits_{m_j > 1} a_j
$$
Alors $h$ satisfait manifestement les conditions imposées ci-dessus aux valuations.
Si $h \not \in \mathfrak m^N$, multiplions-le par un élément de
$\mathfrak m^N$ qui n'appartient à $\mathfrak q_j$ pour tout $j$.
\end{proof}

\begin{lemma}
\label{lemma-divides-radical}
Soit $(A, \mathfrak m, \kappa)$ un domaine intègre local normal noethérien
de dimension $2$. Si $a \in \mathfrak m$ est non nul, alors il existe un
élément $c \in A$ tel que $A/cA$ soit réduit et que $a$ divise
$c^n$ pour un certain $n$.
\end{lemma}

\begin{proof}
Soit $\text{div}(a) = \sum_{i = 1}^r n_i \mathfrak p_i$
avec les notations de la démonstration du Lemme \ref{lemma-radical-element}.
Choisissons $c \in \mathfrak p_1 \cap \ldots \cap \mathfrak p_r$ tel que $A/cA$
soit réduit, voir le Lemme \ref{lemma-radical-element}. Pour $n \geq \max(n_i)$,
on voit que $-\text{div}(a) + \text{div}(c^n)$
est un diviseur effectif (tous ses coefficients sont non négatifs).
Ainsi, $c^n/a \in A$ d'après Algèbre, Lemme
\ref{algebra-lemma-normal-domain-intersection-localizations-height-1}.
\end{proof}

\noindent
Dans la suite de cette section, nous démontrons le résultat en dimension $> 2$.

\begin{lemma}
\label{lemma-multiplicity}
Soit $(R, \mathfrak m)$ un anneau local noethérien de dimension $d$, soit
$g_1, \ldots, g_d$ un système de paramètres, et soit
$I = (g_1, \ldots, g_d)$. Si $e_I/d!$ est le coefficient dominant du
polynôme numérique
$n \mapsto \text{longueur}_R(R/I^{n+1})$, alors $e_I \leq \text{longueur}_R(R/I)$.
\end{lemma}

\begin{proof}
La fonction est un polynôme numérique d'après
Algèbre, Proposition \ref{algebra-proposition-hilbert-function-polynomial}.
Elle est de degré $d$ d'après
Algèbre, Proposition \ref{algebra-proposition-dimension}.
Si $d = 0$, le résultat est trivial.
Si $d = 1$, le résultat est le Lemme \ref{lemma-polypoly}.
Pour le démontrer en général, remarquons qu'il existe une surjection
$$
\bigoplus\nolimits_{i_1, \ldots, i_d \geq 0,\ \sum i_j = n} R/I
\longrightarrow
I^n/I^{n + 1}
$$
qui envoie l'élément de base correspondant à $i_1, \ldots, i_d$
sur la classe de $g_1^{i_1} \ldots g_d^{i_d}$ dans $I^n/I^{n + 1}$.
On voit ainsi que
$$
\text{longueur}_R(R/I^{n + 1}) - \text{longueur}_R(R/I^n)
\leq  \text{longueur}_R(R/I) {n + d - 1 \choose d - 1}
$$
Comme $d \geq 2$, le polynôme numérique de gauche est de
degré $d - 1$ et son coefficient dominant est $e_I / (d - 1)!$.
Le polynôme de droite est de degré $d - 1$ et son
coefficient dominant est $\text{longueur}_R(R/I)/ (d - 1)!$.
Ceci démontre le lemme.
\end{proof}

\begin{lemma}
\label{lemma-minprimespolyhigher}
Soit $(R, \mathfrak m)$ un anneau local noethérien de dimension $d$, soit $t$
le nombre d'idéaux premiers minimaux de $R$ de dimension $d$, et soit
$(g_1,\ldots,g_d)$ un système de paramètres. Alors
$t \leq \text{longueur}_R(R/(g_1,\ldots,g_d))$.
\end{lemma}

\begin{proof}
Si $d = 0$, le lemme est trivial. Si $d = 1$, le lemme est le
Lemme \ref{lemma-minprimespoly}. On peut donc supposer que $d > 1$.
Soient $\mathfrak p_1, \ldots, \mathfrak p_s$ les idéaux premiers minimaux de
$R$, les $t$ premiers étant de dimension $d$, et posons
$I = (g_1, \ldots, g_d)$. En raisonnant exactement comme dans
la démonstration du Lemme \ref{lemma-minprimespoly}, on peut supposer que $R$ est réduit.

\medskip\noindent
Supposons que $R$ est réduit, d'idéaux premiers minimaux
$\mathfrak p_1, \ldots, \mathfrak p_s$.
Cela signifie qu'il existe une suite exacte
$$
0 \to R \to
\prod\nolimits_{i = 1}^s R/\mathfrak p_i \to Q \to 0
$$
Ici, $Q$ est le conoyau de la première application.
Posons $M = \prod_{i = 1}^s R/\mathfrak p_i$.
En localisant en $\mathfrak p_j$, on voit que
$$
R_{\mathfrak p_j} \to M_{\mathfrak p_j} =
\left(\prod\nolimits_{i=1}^s R/\mathfrak p_i\right)_{\mathfrak p_j} =
(R/\mathfrak p_j)_{\mathfrak p_j}
$$
est surjective. Ainsi, $Q_{\mathfrak p_j} = 0$ pour tout $j$. Par conséquent, aucun
idéal premier de hauteur $0$ de $R$ n'appartient au support de $Q$. Il s'ensuit que
le degré du polynôme numérique
$n \mapsto \text{longueur}_R(Q/I^nQ)$ est $\dim(\text{Supp}(Q)) < d$, voir
Algèbre, Lemme \ref{algebra-lemma-support-dimension-d}.
D'après Algèbre, Lemme \ref{algebra-lemma-hilbert-ses-chi}
(qui s'applique car $R$ n'est pas de longueur finie), le polynôme
$$
n \longmapsto
\text{longueur}_R(M/I^nM) - \text{longueur}_R(R/I^n) - \text{longueur}_R(Q/I^nQ)
$$
est de degré $< d$. Comme $M = \prod R/\mathfrak p_i$ et comme
$n \to \text{longueur}_R(R/\mathfrak p_i + I^n)$ est un polynôme numérique
de degré exactement(!) $d$ pour $i = 1, \ldots, t$ (d'après
Algèbre, Lemme \ref{algebra-lemma-support-dimension-d}),
on voit que le coefficient dominant de $n \mapsto \text{longueur}_R(M/I^nM)$
est au moins $t/d!$. On conclut donc par le Lemme \ref{lemma-multiplicity}.
\end{proof}

\begin{lemma}
\label{lemma-sysparhigher}
Soit $(R, \mathfrak m)$ un anneau local noethérien de dimension $d$, et soit
$f \in \mathfrak m$ un élément n'appartenant à aucun idéal premier
minimal de $R$. Il existe alors des éléments
$g_1, \ldots, g_{d - 1} \in \mathfrak m$ et $N \in \mathbf{N}$ tels que
\begin{enumerate}
\item $f, g_1, \ldots, g_{d - 1}$ forment un système de paramètres de $R$
\item Si $h \in \mathfrak m^N$, alors $f + h, g_1, \ldots, g_{d - 1}$ est un
système de paramètres et on a
$\text{longueur}_R R/(f, g_1, \ldots, g_{d-1}) =
\text{longueur}_R R/(f + h, g_1, \ldots, g_{d-1})$.
\end{enumerate}
\end{lemma}

\begin{proof}
D'après le Lemme \ref{lemma-sopexists}, il existe
$g_1, \ldots, g_{d - 1} \in \mathfrak m$
tels que $f, g_1, \ldots, g_{d - 1}$ soit un système de paramètres de $R$.
Alors $\mathfrak m = \sqrt{(f, g_1, \ldots, g_{d - 1})}$.
Il existe donc $n$ tel que
$\mathfrak m^n \subset (f, g_1, \ldots, g_{d - 1})$, voir
Algèbre, Lemme \ref{algebra-lemma-Noetherian-power}. Nous affirmons que
$N = n + 1$ convient.
En effet, soit $h \in \mathfrak m^N$. Par le choix de $N$, on peut écrire
$h = af + \sum b_ig_i$ avec $a, b_i \in \mathfrak m$. Ainsi
\begin{align*}
(f + h, g_1, \ldots, g_{d - 1})
& =
(f + af + \sum b_ig_i, g_1, \ldots, g_{d - 1}) \\
& =
((1 + a)f, g_1, \ldots, g_{d - 1}) \\
& =
(f, g_1, \ldots, g_{d - 1})
\end{align*}
car $1 + a$ est inversible dans $R$. Cela démontre l'égalité
des longueurs et le fait que $f + h, g_1, \ldots, g_{d - 1}$
est un système de paramètres.
\end{proof}

\begin{proposition}
\label{proposition-propdimd}
\begin{reference}
Ce résultat est établi dans \cite[Lemma 3.14]{Artin-Lipman} sans
l'hypothèse que l'anneau soit caténaire
\end{reference}
Soit $R$ un domaine intègre local normal noethérien caténaire.
Soit $J \subset R$ un idéal radical.
Alors il existe un élément non nul $f \in J$
tel que $R/fR$ soit réduit.
\end{proposition}

\begin{proof}
La démonstration est la même que celle du
Lemme \ref{lemma-radical-element},
en utilisant le Lemme \ref{lemma-minprimespolyhigher} à la place du
Lemme \ref{lemma-minprimespoly} et le
Lemme \ref{lemma-sysparhigher} à la place du
Lemme \ref{lemma-syspar}.
On peut utiliser le Lemme \ref{lemma-minprimespolyhigher} puisque $R$
est un domaine intègre caténaire, de sorte que tout idéal premier de hauteur un de $R$
est de dimension $d - 1$, et que le spectre de $R/(f + h)$ est donc
équidimensionnel. Pour la commodité du lecteur, nous donnons
les détails.

\medskip\noindent
Soit $f \in J$ un élément non nul.
Nous allons modifier légèrement $f$ afin d'obtenir un élément qui engendre un idéal radical.
Le localisé $R_\mathfrak p$ de $R$ en chaque idéal premier de hauteur un
$\mathfrak p$ est un anneau de valuation discrète, voir
Algèbre, Lemme \ref{algebra-lemma-characterize-dvr} ou
Algèbre, Lemme \ref{algebra-lemma-criterion-normal}.
Nous notons $\text{ord}_\mathfrak p(f)$ la
valuation correspondante de $f$ dans $R_{\mathfrak p}$. Soient
$\mathfrak q_1, \ldots, \mathfrak q_s$
les idéaux premiers distincts de hauteur un contenant $f$.
Écrivons $\text{ord}_{\mathfrak q_j}(f) = m_j \geq 1$
pour tout $j$. Nous définissons alors
$\text{div}(f) = \sum_{j = 1}^s m_j\mathfrak q_j$
comme une combinaison linéaire formelle d'idéaux premiers
de hauteur un à coefficients entiers.
L'anneau $R/fR$ est réduit si et seulement si
$m_j = 1$ pour $j = 1, \ldots, s$. En effet, si $m_j$ vaut $1$, alors
$(R/fR)_{\mathfrak q_j}$ est réduit et
$R/fR \subset \prod (R/fR)_{\mathfrak q_j}$ puisque
$\mathfrak q_1, \ldots, \mathfrak q_s$ sont les idéaux premiers associés
de $R/fR$, voir les lemmes d'Algèbre
\ref{algebra-lemma-zero-at-ass-zero} et
\ref{algebra-lemma-normal-domain-intersection-localizations-height-1}.

\medskip\noindent
Choisissons et fixons $g_1, \ldots, g_{d - 1}$ et $N$ comme dans le
Lemme \ref{lemma-sysparhigher}.
Pour $y \in R$ non nul, notons $t(y)$ le nombre d'idéaux premiers minimaux au-dessus de $y$.
Comme $R$ est un domaine intègre normal, ces idéaux premiers
sont de hauteur un et correspondent $1$-à-$1$ aux idéaux premiers minimaux de
$R/yR$ (Algèbre, Lemmes \ref{algebra-lemma-minimal-over-1} et
\ref{algebra-lemma-normal-domain-intersection-localizations-height-1}).
Par exemple, $t(f) = s$ est le nombre
d'idéaux premiers $\mathfrak q_j$ apparaissant dans $\text{div}(f)$.
Soit $h \in \mathfrak m^N$. Puisque $R$ est caténaire, pour tout
idéal premier de hauteur un $\mathfrak p$ de $R$, on a
$\dim(R/\mathfrak p) = d - 1$. Par conséquent, d'après le
Lemme \ref{lemma-minprimespolyhigher},
on a
\begin{align*}
t(f + h) & \leq \text{longueur}_{R/(f + h)}(R/(f + h, g_1, \ldots, g_{d - 1})) \\
& = \text{longueur}_R(R/(f + h, g_1, \ldots, g_{d - 1})) \\
& = \text{longueur}_R(R/(f, g_1, \ldots, g_{d - 1}))
\end{align*}
voir Algèbre, Lemme \ref{algebra-lemma-length-independent}
pour la première égalité.
On voit donc que $t(f + h)$ est borné indépendamment de
$h \in \mathfrak m^N$.

\medskip\noindent
Grâce à la bornitude démontrée ci-dessus, on peut choisir $h \in \mathfrak m^N \cap J$
tel que $t(f + h)$ soit maximal parmi de tels $h$. Posons $f' = f + h$.
Étant donné $h' \in \mathfrak m^N \cap J$,
on voit que le nombre $t(f' + h') \leq t(f + h)$.
Ainsi, après avoir remplacé $f$ par $f'$, on peut supposer que pour tout
$h \in \mathfrak m^N \cap J$, on a $t(f + h) \leq s$.

\medskip\noindent
Supposons ensuite que nous puissions trouver un élément $h \in \mathfrak m^N \cap J$
tel que, pour tout $j$, on ait $\text{ord}_{\mathfrak q_j}(h) \geq 1$ et
$\text{ord}_{\mathfrak q_j}(h) = 1 \Leftrightarrow m_j > 1$.
Alors $\text{ord}_{\mathfrak q_j}(f + h) = 1$
pour tout $j$, d'après les propriétés élémentaires des valuations.
Ainsi
$$
\text{div}(f + h) = \sum\nolimits_{j = 1}^s \mathfrak q_j +
\sum\nolimits_{k = 1}^v e_k \mathfrak r_k
$$
pour certains idéaux premiers de hauteur un distincts deux à deux
$\mathfrak r_1, \ldots, \mathfrak r_v$ et $e_k \geq 1$.
Cependant, comme $s = t(f) \geq t(f + h)$, on voit que $v = 0$
et nous avons trouvé l'élément voulu.

\medskip\noindent
Choisissons maintenant $h$ satisfaisant les critères ci-dessus.
Par évitement des idéaux premiers (Algèbre, Lemme \ref{algebra-lemma-silly}),
pour tout $1 \leq j \leq s$, on peut trouver un élément
$a_j \in \mathfrak q_j \cap J$
tel que $a_j \not \in \mathfrak q_{j'}$ pour $j' \not = j$.
Ensuite, on peut choisir
$b_j \in J \cap \mathfrak q_1 \cap \ldots \cap \mathfrak q_s$
tel que $b_j \not \in \mathfrak q_j^{(2)}$. Ici,
$\mathfrak q_j^{(2)} = \{x \in R \mid \text{ord}_{\mathfrak q_j}(x) \geq 2\}$
est la deuxième puissance symbolique de $\mathfrak q_j$.
L'évitement des idéaux premiers s'applique car l'idéal
$J' = J \cap \mathfrak q_1 \cap \ldots \cap \mathfrak q_s$
est radical, donc $R/J'$ est réduit, donc $(R/J')_{\mathfrak q_j}$
est réduit, donc $J'$ contient un élément $x$ tel que
$\text{ord}_{\mathfrak q_j}(x) = 1$, donc
$J' \not \subset \mathfrak q_j^{(2)}$. Alors l'élément
$$
c = \sum\nolimits_{j = 1, \ldots, s}
b_j \times \prod\nolimits_{j' \not = j} a_{j'}
$$
est un élément de $J$
tel que $\text{ord}_{\mathfrak q_j}(c) = 1$ pour tout $j = 1, \ldots, s$,
d'après les propriétés élémentaires des valuations. Enfin, posons
$$
h = c \times \prod\nolimits_{m_j = 1} a_j \times y
$$
où $y \in \mathfrak m^N$ est un élément
qui n'appartient à $\mathfrak q_j$ pour tout $j$.
\end{proof}
\section{Objets inversibles dans la catégorie dérivée}
\label{section-invertible-D-or-R}

\noindent
Nous caractérisons les objets inversibles dans la catégorie dérivée
d'un anneau.

\begin{lemma}
\label{lemma-symmetric-monoidal-derived}
Soit $R$ un anneau. La catégorie dérivée $D(R)$ de $R$
est une catégorie monoïdale symétrique dont le produit tensoriel
est donné par le produit tensoriel dérivé et dont les contraintes d'associativité et
de commutativité sont celles de la section \ref{section-sign-rules}.
\end{lemma}

\begin{proof}
Omis. Indications : la contrainte d'associativité est l'isomorphisme du
lemme \ref{lemma-triple-tensor-product}
et la contrainte de commutativité est l'isomorphisme du
lemme \ref{lemma-flip-douoble-tensor-product}.
Cela étant dit, la commutativité des divers diagrammes
résulte du résultat correspondant pour la catégorie
des complexes de $R$-modules ; voir la section \ref{section-symmetric-monoidal}.
\end{proof}

\noindent
Ainsi, nous savons ce que signifie, pour un objet de $D(R)$, posséder un dual (à gauche)
ou être inversible. Avant de déterminer ce que cela implique, nous
avons besoin d'un lemme simple.

\begin{lemma}
\label{lemma-complex-bounded-above-free-colim-bounded-finite-free}
Soit $R$ un anneau. Soit $F^\bullet$ un complexe borné supérieurement de
$R$-modules libres. Étant donnés des couples $(n_i, f_i)$, $i = 1, \ldots, N$
avec $n_i \in \mathbf{Z}$ et $f_i \in F^{n_i}$, il existe
un sous-complexe $G^\bullet \subset F^\bullet$ contenant
tous les $f_i$, qui est borné et constitué de $R$-modules libres de type fini.
\end{lemma}

\begin{proof}
Par récurrence descendante sur $a = \min(n_i; i = 1, \ldots, N)$.
Si $F^n = 0$ pour $n \geq a$, le résultat est vrai en prenant
pour $G^\bullet$ le complexe nul. En général, après renumérotation,
nous pouvons supposer qu'il existe $1 \leq r \leq N$ tel que
$n_1 = \ldots = n_r = a$ et $n_i > a$ pour $i > r$.
Choisissons une base $b_j, j \in J$ de $F^a$. Nous pouvons choisir une partie finie
$J' \subset J$ telle que $f_i \in \bigoplus_{j \in J'} Rb_j$
pour $i = 1, \ldots, r$. Choisissons une base $c_k, k \in K$ de $F^{a + 1}$.
Nous pouvons choisir une partie finie $K' \subset K$ telle que
$\text{d}_F^a(b_j) \in \bigoplus_{k \in K'} Rc_k$ pour
$j \in J'$. Nous pouvons alors appliquer l'hypothèse de
récurrence pour trouver un sous-complexe $H^\bullet \subset F^\bullet$
contenant $c_k \in F^{a + 1}$ pour $k \in K'$ et
$f_i \in F^{n_i}$ pour $i > r$. Prenons $G^\bullet$ égal à
$H^\bullet$ aux degrés $> a$ et à $\bigoplus_{j \in J'} Rb_j$
au degré $a$.
\end{proof}

\begin{lemma}
\label{lemma-have-dual-derived}
Soit $R$ un anneau. Soit $M$ un objet de $D(R)$. Les conditions suivantes
sont équivalentes :
\begin{enumerate}
\item $M$ possède un dual à gauche dans $D(R)$ au sens de
Catégories, définition \ref{categories-definition-dual},
\item $M$ est un objet parfait de $D(R)$.
\end{enumerate}
De plus, dans ce cas, le dual à gauche de $M$ est l'objet
$M^\vee$ du lemme \ref{lemma-dual-perfect-complex}.
\end{lemma}

\begin{proof}
Si $M$ est parfait, nous pouvons représenter $M$ par un complexe borné
$M^\bullet$ de $R$-modules projectifs de type fini. Dans ce cas, $M^\bullet$
possède un dual à gauche dans la catégorie des complexes d'après le
lemme \ref{lemma-left-dual-complex},
qui est a fortiori un dual à gauche dans $D(R)$.

\medskip\noindent
Supposons (1). Soient $N$, $\eta : R \to M \otimes_R^\mathbf{L} N$ et
$\epsilon : M \otimes_R^\mathbf{L} N \to R$ les données d'un dual à gauche
au sens de Catégories, définition \ref{categories-definition-dual}.
Choisissons un complexe $M^\bullet$ représentant $M$. Choisissons un
complexe K-plat $N^\bullet$ à termes plats
représentant $N$ ; voir le lemme \ref{lemma-K-flat-resolution}.
Alors $\eta$ est représenté par un morphisme de complexes
$$
\eta : R \longrightarrow \text{Tot}(M^\bullet \otimes_R N^\bullet)
$$
Nous pouvons écrire l'image de $1$ comme une somme finie
$$
\eta(1) = \sum\nolimits_n \sum\nolimits_i m_{n, i} \otimes n_{-n, i}
$$
avec $m_{n, i} \in M^n$ et $n_{-n, i} \in N^{-n}$. Soit
$K^\bullet \subset M^\bullet$ le sous-complexe
engendré par tous les éléments $m_{n, i}$ et $\text{d}(m_{n, i})$.
Par notre choix de $N^\bullet$, nous avons
$\text{Tot}(K^\bullet \otimes_R N^\bullet) \subset
\text{Tot}(M^\bullet \otimes_R N^\bullet)$ et
$\eta(1)$ appartient au sous-complexe par le choix ci-dessus.
Notons $K$ l'objet de $D(R)$ représenté par $K^\bullet$.
Nous voyons alors que $\eta$ se factorise par un morphisme
$\tilde \eta : R \longrightarrow K \otimes_R^\mathbf{L} N$.
Puisque
$(1 \otimes \epsilon) \circ (\eta \otimes 1) = \text{id}_M$
nous en concluons que l'identité de $M$ se factorise par $K$
dans le diagramme commutatif
$$
\xymatrix{
M \ar[rr]_-{\eta \otimes 1} \ar[rrd]_{\tilde \eta \otimes 1} & &
M \otimes_R^\mathbf{L} N \otimes_R^\mathbf{L} M
\ar[r]_-{1 \otimes \epsilon} &
M \\
& &
K \otimes_R^\mathbf{L} N \otimes_R^\mathbf{L} M
\ar[u] \ar[r]^-{1 \otimes \epsilon} &
K \ar[u]
}
$$
Puisque $K$ est borné supérieurement, il s'ensuit que $M \in D^-(R)$.
Nous pouvons donc représenter $M$ par un complexe borné supérieurement
$M^\bullet$ de $R$-modules libres ; voir par exemple
Catégories dérivées, lemme \ref{derived-lemma-subcategory-left-resolution}.
Écrivons $\eta(1) = \sum\nolimits_n \sum\nolimits_i m_{n, i} \otimes n_{-n, i}$
comme précédemment.
D'après le lemme \ref{lemma-complex-bounded-above-free-colim-bounded-finite-free},
nous pouvons trouver un sous-complexe $K^\bullet \subset M^\bullet$
contenant tous les éléments $m_{n, i}$,
qui est borné et constitué de $R$-modules libres de type fini.
Comme ci-dessus, nous voyons que l'identité de $M$ se factorise par $K$.
Puisque $K$ est parfait, nous en concluons que $M$ est également parfait ; voir le
lemme \ref{lemma-summands-perfect}.
\end{proof}

\begin{lemma}
\label{lemma-invertible-derived}
Soit $R$ un anneau. Soit $M$ un objet de $D(R)$. Les conditions suivantes
sont équivalentes :
\begin{enumerate}
\item $M$ est inversible dans $D(R)$ ; voir
Catégories, définition \ref{categories-definition-invertible}, et
\item pour tout idéal premier $\mathfrak p \subset R$, il
existe $f \in R$, $f \not \in \mathfrak p$ tel que
$M_f \cong R_f[-n]$ pour un certain $n \in \mathbf{Z}$.
\end{enumerate}
De plus, dans ce cas,
\begin{enumerate}
\item[(a)] $M$ est un objet parfait de $D(R)$,
\item[(b)] $M = \bigoplus H^n(M)[-n]$ dans $D(R)$,
\item[(c)] chaque $H^n(M)$ est un $R$-module projectif de type fini,
\item[(d)] nous pouvons écrire $R = \prod_{a \leq n \leq b} R_n$
de sorte que $H^n(M)$ corresponde à un $R_n$-module inversible.
\end{enumerate}
\end{lemma}

\begin{proof}
Supposons (2). Considérons l'objet $R\Hom_R(M, R)$ et le morphisme de composition
$$
R\Hom(M, R) \otimes_R^\mathbf{L} M \to R
$$
Une vérification locale montre qu'il s'agit d'un isomorphisme ; nous omettons les détails.
Comme $D(R)$ est monoïdale symétrique, nous voyons que $M$ est inversible.

\medskip\noindent
Supposons (1). Observons qu'un objet inversible d'une catégorie monoïdale
possède un dual à gauche, à savoir son inverse. Ainsi, $M$ est parfait d'après le
lemme \ref{lemma-have-dual-derived}. Considérons un idéal premier
$\mathfrak p \subset R$ de corps résiduel $\kappa$.
Nous voyons alors que $M \otimes_R^\mathbf{L} \kappa$ est un
objet inversible de $D(\kappa)$. Cela implique manifestement que
$\dim H^i(M \otimes_R^\mathbf{L} \kappa)$ est non nulle pour un unique
$i$, et égale à $1$ dans ce cas. D'après le
lemme \ref{lemma-lift-perfect-from-residue-field},
nous obtenons (2).

\medskip\noindent
Dans la preuve ci-dessus, nous avons vu que (a) est vérifiée. Soit $U_n \subset \Spec(R)$
la réunion des ouverts de la forme $D(f)$ tels que
$M_f \cong R_f[-n]$. Manifestement, $U_n \cap U_{n'} = \emptyset$ si $n \not = n'$.
Si $M$ a une amplitude de Tor contenue dans $[a, b]$, alors $U_n = \emptyset$
si $n \not \in [a, b]$. Nous obtenons donc une décomposition en produit
$R = \prod_{a \leq n \leq b} R_n$ comme en (d), telle que
$U_n$ corresponde à $\Spec(R_n)$ ; voir
Algèbre, lemme \ref{algebra-lemma-disjoint-implies-product}.
Puisque $D(R) = \prod_{a \leq n \leq b} D(R_n)$, et de même
pour la catégorie des modules, les assertions (b), (c) et (d) en résultent immédiatement.
\end{proof}





\section{Extraction d'un facteur direct libre}
\label{section-serre-splitting}

\noindent
Les arguments de cette section sont dus à Serre ; voir \cite{Serre-projective}.

\begin{situation}
\label{situation-splitting}
Ici, $R$ est un anneau et $M$ est un $R$-module de présentation finie.
Notons $\Omega \subset \Spec(R)$ l'ensemble des points fermés
muni de la topologie induite. Pour $x \in \Omega$, notons
$M(x) = M/xM$ la fibre de $M$ en $x$. C'est un espace vectoriel
de dimension finie sur le corps résiduel $\kappa(x)$ en $x$.
Étant donné $s \in M$, nous notons $s(x)$ l'image de $s$ dans $M(x)$.
\end{situation}

\begin{lemma}
\label{lemma-which-elements-split}
Dans la situation \ref{situation-splitting}, soit $x \in \Omega$. Il existe une
suite exacte courte canonique
$$
0 \to B(x) \to M(x) \to V(x) \to 0
$$
d'espaces vectoriels sur $\kappa(x)$ qui possède la
propriété suivante : pour $s_1, \ldots, s_r \in M$, les conditions suivantes sont équivalentes :
\begin{enumerate}
\item il existe $f \in R$, $f \not \in x$ tel que
le morphisme $s_1, \ldots, s_r : R^{\oplus r} \to M$ devienne
l'inclusion d'un facteur direct après inversion de $f$, et
\item les éléments $s_1(x), \ldots, s_r(x)$ s'envoient sur des éléments linéairement indépendants
de $V(x)$.
\end{enumerate}
\end{lemma}

\begin{proof}
Définissons $B(x) \subset M(x)$ comme l'orthogonal de l'image du morphisme
$$
\Hom_R(M, R) \to \Hom_{\kappa(x)}(M(x), \kappa(x))
$$
et posons $V(x) = M(x)/B(x)$. Alors toute application $R$-linéaire $\varphi : M \to R$
induit une application $\overline{\varphi} : V(x) \to \kappa(x)$ et, réciproquement,
toute application $\kappa(x)$-linéaire $\lambda : V(x) \to \kappa(x)$ est égale
à $\overline{\varphi}$ pour un certain $\varphi$. Soient $s_1, \ldots, s_r \in M$.

\medskip\noindent
Supposons que $s_1, \ldots, s_r$ s'envoient sur des éléments linéairement indépendants de
$V(x)$. Nous pouvons alors trouver $\varphi_1, \ldots, \varphi_r \in \Hom_R(M, R)$
tels que $\varphi_i(s_j)$ s'envoie sur $\delta_{ij}$\footnote{Delta de Kronecker.}
dans $\kappa(x)$. Par conséquent, la matrice de la composée
$$
R^{\oplus r}
\xrightarrow{s_1, \ldots, s_r}
M
\xrightarrow{\varphi_1, \ldots, \varphi_r}
R^{\oplus r}
$$
a un déterminant $f \in R$ qui s'envoie sur $1$ dans $\kappa(x)$.
Manifestement, cela implique que $s_1, \ldots, s_r : R^{\oplus r} \to M$
est l'inclusion d'un facteur direct après inversion de $f$.

\medskip\noindent
Réciproquement, supposons qu'il existe $f \in R$, $f \not \in x$
tel que $s_1, \ldots, s_r : R^{\oplus r} \to M$
soit l'inclusion d'un facteur direct après inversion de $f$.
Nous pouvons donc trouver des applications $R_f$-linéaires $\varphi_i : M_f \to R_f$
telles que $\varphi_i(s_j) = \delta_{ij} \in R_f$.
Puisque $\Hom_R(M, R)_f = \Hom_{R_f}(M_f, R_f)$ d'après
Algèbre, lemme \ref{algebra-lemma-hom-from-finitely-presented},
nous en concluons qu'il existe $n \geq 0$ et $\varphi'_i \in \Hom_R(M, R)$
tels que $\varphi'_i(s_j) = f^n\delta_{ij} \in R$.
Il s'ensuit que $s_1, \ldots, s_r$ s'envoient sur des éléments linéairement indépendants
de $V(x)$, puisque $\overline{\varphi}'_i(s_j) = f^n\delta_{ij}$.
\end{proof}

\noindent
Dans la situation \ref{situation-splitting}, étant donnés $s_1, \ldots, s_r \in M$,
nous notons $Z(s_1, \ldots, s_r) \subset \Omega$ l'ensemble des $x \in \Omega$
tels que $s_1(x), \ldots, s_r(x)$ s'envoient sur des éléments linéairement
dépendants de $V(x)$. D'après le lemme, c'est un sous-ensemble fermé de $\Omega$.

\begin{lemma}
\label{lemma-choose-values}
Dans la situation \ref{situation-splitting}, soient $x_1, \ldots, x_n \in \Omega$
deux à deux distincts. Soient $v_i \in V(x_i)$. Il existe alors
$s \in M$ tel que $s(x_i)$ s'envoie sur $v_i$ pour $i = 1, \ldots, n$.
\end{lemma}

\begin{proof}
Puisque $x_i$ est un idéal maximal de $R$, nous pouvons utiliser
Algèbre, lemme \ref{algebra-lemma-chinese-remainder},
pour voir que $M(x_1) \oplus \ldots \oplus M(x_n)$
est un quotient de $M$.
\end{proof}

\begin{proposition}
\label{proposition-splitting}
\begin{reference}
\cite[Theorem 2]{Serre-projective}
\end{reference}
Dans la situation \ref{situation-splitting}, supposons que $\Omega$ soit
un espace topologique noethérien. Soient $s_1, \ldots, s_h \in M$.
Soient $Z(s_1, \ldots, s_h) \subset F \subset \Omega$ des fermés.
Soient $x_1, \ldots, x_n \in F$ deux à deux distincts.
Soit $v_i  \in V(x_i)$.
Soit $k \geq 0$ un entier tel que
$$
(*)\quad h + k \leq \dim_{\kappa(x)} V(x)\text{ pour tout }x \in \Omega
$$
Alors il existe $s \in M$ et un fermé $F' \subset \Omega$ tels que
\begin{enumerate}
\item[(a)] $s(x_i)$ s'envoie sur $v_i$,
\item[(b)] $Z(s_1, \ldots, s_h, s) \subset F \cup F'$, et
\item[(c)] toute composante irréductible de $F'$ est de
codimension $\geq k$ dans $\Omega$.
\end{enumerate}
\end{proposition}

\begin{proof}
Rappelons que la codimension a été définie dans
Topologie, section \ref{topology-section-catenary-spaces},
et que nous utiliserons certains résultats sur les espaces topologiques
noethériens contenus dans
Topologie, section \ref{topology-section-noetherian}.

\medskip\noindent
La preuve procède par récurrence sur $k$. Si $k = 0$, nous choisissons
$s \in M$ comme dans le lemme \ref{lemma-choose-values} et nous choisissons
$F' = \Omega$.

\medskip\noindent
Supposons $k > 0$. Par l'hypothèse de récurrence, nous pouvons
choisir $u \in M$ et un fermé $G \subset \Omega$ satisfaisant
(a), (b), (c) pour $s_1, \ldots, s_h$, $F$, $x_1, \ldots, x_n$,
$v_1, \ldots, v_n$ et $k - 1$.

\medskip\noindent
Soit $G = G_1 \cup \ldots \cup G_m$ la décomposition de $G$ en
ses composantes irréductibles. Si $G_j \subset F$, nous pouvons la retirer
de la liste. Nous pouvons donc supposer que $G_j$ n'est pas contenu dans $F$
pour $j = 1, \ldots, m$.
Pour $j = 1, \ldots, m$, choisissons $y_j \in G_j$ avec $y_j \not \in F$
et $y_j \not \in G_{j'}$ pour $j' \not = j$. Cela est possible puisqu'il
n'existe aucune inclusion entre les composantes irréductibles de $G$.
Choisissons $w_j \in V(y_j)$ n'appartenant pas au sous-espace engendré par les images de
$s_1(y_j), \ldots, s_h(y_j)$ ; cela est possible parce que
$h + k \leq \dim V(y_j)$ et $k > 0$.

\medskip\noindent
Appliquons l'hypothèse de récurrence aux $h + 1$ sections
$s_1, \ldots, s_h, u$, au fermé $F \cup G$,
aux points $x_1, \ldots, x_n, y_1, \ldots, y_m \in F \cup G$,
aux éléments $0 \in V(x_i)$ et $w_j \in V(y_j)$, et
à l'entier $k - 1$. Remarquons que nous avons augmenté $h$ de $1$ et
diminué $k$ de $1$ ; l'hypothèse $(*)$ de la proposition
reste donc vérifiée. Nous obtenons ainsi $t \in M$ et un fermé $H \subset \Omega$
satisfaisant (a), (b), (c) pour $s_1, \ldots, s_h, u$, $F \cup G$,
$x_1, \ldots, x_n, y_1, \ldots, y_m$, $0, \ldots, 0, w_1, \ldots, w_m$
et $k - 1$.

\medskip\noindent
Soient $H_1, \ldots, H_p \subset H$ les composantes irréductibles de $H$
qui ne sont pas contenues dans $F \cup G$. Comme précédemment, choisissons $z_l \in H_l$,
$z_l \not \in F \cup G$ et $z_l \not \in H_{l'}$ pour $l' \not = l$.
En utilisant Algèbre, lemme \ref{algebra-lemma-chinese-remainder},
nous pouvons choisir $f \in R$ tel que
$f(y_j) = 1$, $j = 1, \ldots, m$ et $f(z_l) = 0$, $l = 1, \ldots, p$.
Affirmation : l'élément $s = u + f t$ convient.

\medskip\noindent
Tout d'abord, la valeur $s(x_i)$ coïncide avec $u(x_i)$ parce que
$t(x_i) = 0$ ; nous voyons donc que $s(x_i)$ s'envoie sur $v_i$.
Cela prouve (a).
Pour achever la preuve, il suffit de montrer que toute composante irréductible
$Z$ de $Z(s_1, \ldots, s_h, s)$ non contenue dans $F$
est de codimension $\geq k$ dans $\Omega$. En effet, nous pouvons alors prendre $F'$
égal à leur réunion et nous obtenons (b) et (c).
Nous pouvons voir qu'il n'existe aucune composante irréductible $Z$ de $Z(s_1, \ldots, s_h, s)$
de codimension $\leq k - 1$ comme suit :
\begin{enumerate}
\item Observons que
$Z(s_1, \ldots, s_h, s) \subset Z(s_1, \ldots, s_h, u, t) = F \cup H$
puisque $s = u + ft$. Ainsi, $Z \subset H$.
\item Les composantes irréductibles de $H$ sont de codimension $\geq k - 1$.
Par conséquent, $Z$ est égale à une composante irréductible de $H$
puisque $Z$ est de codimension $\leq k - 1$.
Ainsi, $Z = H_l$ pour un certain $l \in \{1, \ldots, p\}$ ou
$Z = G_j$ pour un certain $j \in \{1, \ldots, m\}$.
\item Mais $Z = G_j$ est impossible, car $s_1(y_j), \ldots, s_h(y_j)$
s'envoient sur des éléments linéairement indépendants de $V(y_j)$ et
$s(y_j) = u(y_j) + f(y_j) t(y_j) = u(y_j) + t(y_j)$
s'envoie sur un élément de la forme
$$
\text{combinaison linéaire des images des }s_i(y_j) + w_j
$$
qui est linéairement indépendant des images de $s_1(y_j), \ldots, s_h(y_j)$
dans $V(y_j)$ par notre choix de $w_j$.
\item De même, $Z = H_l$ est impossible. En effet, à nouveau,
$s_1(z_l), \ldots, s_h(z_l)$ s'envoient sur des éléments linéairement indépendants
de $V(z_l)$ et $s(z_l) = u(z_l) + f(z_l) t(z_l) = u(z_l)$
s'envoie sur un élément de $V(z_l)$ linéairement indépendant de celles-ci,
puisque $z_l \not \in F \cup G$.
\end{enumerate}
Ceci achève la preuve.
\end{proof}

\begin{theorem}
\label{theorem-splitting}
\begin{reference}
\cite[Theorem 1]{Serre-projective}
\end{reference}
Soit $R$ un anneau dont le spectre maximal $\Omega \subset \Spec(R)$
est un espace topologique noethérien de dimension $d < \infty$.
Soit $M$ un $R$-module de présentation finie tel que,
pour tout $\mathfrak m \in \Omega$, le $R_\mathfrak m$-module
$M_\mathfrak m$ possède un facteur direct libre de rang $> d$.
Alors $M \cong R \oplus M'$.
\end{theorem}

\begin{proof}
Pour $\mathfrak m \in \Omega$, supposons que $R_\mathfrak m^{\oplus r}$
soit un facteur direct de $M_\mathfrak m$. Alors, d'après les lemmes d'Algèbre
\ref{algebra-lemma-localization-colimit} et
\ref{algebra-lemma-colimit-category-fp-modules},
nous voyons que $R_f^{\oplus r}$ est un facteur direct de $M_f$
pour un certain $f \in R$, $f \not \in \mathfrak m$. Par conséquent,
l'hypothèse signifie que $\dim V(x) > d$ pour tout $x \in \Omega$,
où $V(x)$ est défini comme dans le lemme \ref{lemma-which-elements-split}.
En appliquant la proposition \ref{proposition-splitting} avec $F = \emptyset$,
$h = 0$ et aucun $s_i$, $n = 0$ et aucun $x_i, v_i$, et $k = d + 1$,
nous trouvons $s \in M$ et $F' \subset \Omega$
tels que toute composante irréductible de $F'$ soit de codimension $\geq d + 1$
et que $Z(s) \subset F'$. Puisque $d = \dim(\Omega)$, cela impose
$F' = \emptyset$. Ainsi, $s : R \to M$ est l'inclusion d'un
facteur direct en tout idéal maximal. Il s'ensuit que $s$ est universellement
injectif ; voir Algèbre, lemme
\ref{algebra-lemma-universally-injective-check-stalks}. Alors
$s$ est un monomorphisme scindé d'après
Algèbre, lemme \ref{algebra-lemma-universally-exact-split}.
\end{proof}








\section{Les grands modules projectifs sont libres}
\label{section-big-projective-free}

\noindent
Dans cette section, nous examinons l'un des résultats de \cite{Bass} ;
nous suggérons au lecteur de consulter l'article original. Notre argument
utilisera la preuve légèrement simplifiée donnée dans les articles
\cite{Akasaki} et \cite{Hinohara}.

\begin{lemma}[Lemme d'Eilenberg]
\label{lemma-eilenberg-swindle}
\begin{slogan}
Astuce d'Eilenberg
\end{slogan}
\begin{history}
Dans \cite{Bass}, on lit : ``...est une élégante petite astuce,
observée plusieurs années auparavant par Eilenberg, et qui
aurait bien pu jaillir du front de Barry Mazur.''
\end{history}
\begin{reference}
\cite[Eilenberg's lemma]{Bass}
\end{reference}
Si $P \oplus Q \cong F$, où $F$ est un module libre qui n'est pas de type fini,
alors $P \oplus F \cong F$.
\end{lemma}

\begin{proof}
$$
F \cong F \oplus F \oplus \ldots \cong
P \oplus Q \oplus P \oplus Q \oplus \ldots \cong
P \oplus F \oplus F \oplus \ldots \cong P \oplus F
$$
\end{proof}

\begin{lemma}
\label{lemma-projective-plus-free-is-free}
Soit $R$ un anneau. Soit $P$ un module projectif.
Il existe un module libre $F$ tel que $P \oplus F$ soit libre.
\end{lemma}

\begin{proof}
Puisque $P$ est projectif, nous voyons que $F_0 = P \oplus Q$
est un module libre pour un certain module $Q$. Posons
$F = \bigoplus_{n \geq 1} F_0$. Alors
$P \oplus F \cong F$ d'après le lemme \ref{lemma-eilenberg-swindle}.
\end{proof}

\begin{lemma}
\label{lemma-element-projective}
Soit $R$ un anneau. Soit $P$ un module projectif.
Soit $s \in P$. Il existe un module libre de type fini $F$
et un facteur direct libre de type fini $K \subset F \oplus P$
tel que $(0, s) \in K$.
\end{lemma}

\begin{proof}
D'après le lemme \ref{lemma-projective-plus-free-is-free}, nous pouvons
trouver un module libre $F$ (éventuellement de rang infini) tel que
$F \oplus P$ soit libre. Alors, bien sûr, $(0, s)$ est contenu
dans un facteur direct libre de type fini $K \subset F \oplus P$. À son tour, $K$
est contenu dans $F' \oplus P$, où $F' \subset F$ est
un facteur direct libre de type fini.
\end{proof}

\begin{lemma}
\label{lemma-trick-to-find-good-element}
Soit $R$ un anneau de radical de Jacobson $J$ tel que $R/J$
soit noethérien. Soit $P$ un $R$-module projectif
tel que $P_\mathfrak m$ soit de rang infini pour tout idéal maximal
$\mathfrak m$ de $R$. Soient $s \in P$ et $M \subset P$ tels que
$Rs + M = P$. Alors nous pouvons trouver $m \in M$ tel que $R(s + m)$
soit un facteur direct libre de $P$.
\end{lemma}

\begin{proof}
L'énoncé a un sens, car $P_\mathfrak m$ est libre d'après
Algèbre, théorème \ref{algebra-theorem-projective-free-over-local-ring}.

\medskip\noindent
Notons $M' \subset P/JP$ l'image de $M$ et $s' \in P/JP$
l'image de $s$. Observons que $R/J s' + M' = P/JP$.
Supposons que nous puissions trouver $m' \in M'$ tel que
$R/J(s' + m')$ soit un facteur direct libre de $P/JP$.
Choisissons $\varphi' : P/JP \to R/J$ donnant une scission, c'est-à-dire que
nous avons $\varphi'(s' + m') = 1$ dans $R/J$.
Alors, puisque $P$ est un $R$-module projectif, nous pouvons trouver
un relèvement $\varphi : P \to R$ de $\varphi'$.
Choisissons $m \in M$ s'envoyant sur $m'$.
Alors $\varphi(s + m) \in R$ est congru à $1$ modulo $J$ et est donc
une unité de $R$ (Algèbre, lemme \ref{algebra-lemma-contained-in-radical}).
Par conséquent, $R(s + m)$ est un facteur direct libre de $P$.
Cela nous ramène au cas examiné dans le paragraphe suivant.

\medskip\noindent
Supposons $R$ noethérien. Soit $m \in M$ un élément et soient
$\varphi_1, \ldots, \varphi_n : P \to R$ des applications $R$-linéaires. Notons
$$
Z(s + m, \varphi_1, \ldots, \varphi_n) \subset \Spec(R)
$$
le lieu d'annulation de $\varphi_1(s + m), \ldots, \varphi_n(s + m) \in R$.

\medskip\noindent
Supposons que $\mathfrak m$ soit un idéal maximal de $R$ et que
$\mathfrak m \in Z(s + m, \varphi_1, \ldots, \varphi_n)$.
Posons $K = M \cap \bigcap \Ker(\varphi_i)$. Nous affirmons que l'image de
$$
K/\mathfrak mK \to P/\mathfrak m P
$$
est de dimension infinie. En effet, le quotient $P/K$ est un $R$-module de type fini,
car il est isomorphe à un sous-module de $P/M \oplus R^{\oplus n}$.
Nous voyons ainsi que
le noyau de la flèche affichée est un quotient de
$\text{Tor}_1^R(P/K, \kappa(\mathfrak m))$, qui est de type fini
d'après Algèbre, lemme \ref{algebra-lemma-tor-noetherian}.
En combinant ceci avec le fait que $P/\mathfrak mP$ est de dimension
infinie, nous obtenons notre affirmation. Nous pouvons donc trouver $t \in K$
qui s'envoie sur un élément $\overline{t}$
de l'espace vectoriel $P/\mathfrak mP$ linéairement indépendant
de l'image $\overline{s + m}$ de $s + m$.
Par algèbre linéaire, nous trouvons une application $R$-linéaire
$\overline{\varphi} : P \to \kappa(\mathfrak m)$
telle que $\overline{\varphi}(\overline{t}) = 1$ et
$\overline{\varphi}(\overline{s + m}) = 0$.
Puisque $P$ est projectif, nous pouvons trouver une application $R$-linéaire
$\varphi : P \to R$ relevant $\overline{\varphi}$.
Nous voyons alors que le lieu d'annulation
$Z(s + m + t, \varphi_1, \ldots, \varphi_n, \varphi)$ est
contenu dans $Z(s + m, \varphi_1, \ldots, \varphi_n)$ mais
ne contient pas $\mathfrak m$, c'est-à-dire qu'il
est strictement plus petit que $Z(s + m, \varphi_1, \ldots, \varphi_n)$.

\medskip\noindent
Puisque $\Spec(R)$ est un espace topologique noethérien, les arguments
ci-dessus montrent que nous pouvons trouver
$m \in M$ et $\varphi_1, \ldots, \varphi_n : P \to R$
tels que le fermé $Z(s + m, \varphi_1, \ldots, \varphi_n)$
ne contienne aucun point fermé de $\Spec(R)$.
Ainsi, $Z(s + m, \varphi_1, \ldots, \varphi_n) = \emptyset$.
Nous pouvons donc trouver $r_1, \ldots, r_n \in R$ tels que
$\sum r_i\varphi_i(s + m) = 1$. Par conséquent,
$$
R \xrightarrow{s + m} P \xrightarrow{\sum r_i \varphi_i} R
$$
est la scission recherchée.
\end{proof}

\begin{lemma}
\label{lemma-element-in-free-summand}
Soit $R$ un anneau de radical de Jacobson $J$ tel que $R/J$
soit noethérien. Soit $P$ un $R$-module projectif
tel que $P_\mathfrak m$ soit de rang infini pour tout idéal maximal
$\mathfrak m$ de $R$. Soit $s \in P$. Nous pouvons alors trouver
un facteur direct stablement libre de type fini $M \subset P$ tel que $s \in M$.
\end{lemma}

\begin{proof}
D'après le lemme \ref{lemma-element-projective}, nous pouvons trouver un module libre
de type fini $F$ et un facteur direct libre de type fini $K \subset F \oplus P$ tels
que $(0, s) \in K$. Par récurrence sur le rang de $F$, nous nous ramenons
au cas examiné dans le paragraphe suivant.

\medskip\noindent
Supposons qu'il existe un facteur direct stablement libre de type fini $K \subset R \oplus P$
tel que $(0, s) \in K$. Choisissons un complémentaire $K'$ de $K$, c'est-à-dire
tel que $R \oplus P = K \oplus K'$. La projection
$\pi : R \oplus P \to K'$ est surjective ; ainsi, d'après le
lemme \ref{lemma-trick-to-find-good-element},
nous trouvons $p \in P$ tel que $\pi(1, p) \in K'$ engendre
un facteur direct libre. En conséquence, écrivons
$K' = R\pi(1, p) \oplus K''$. Nous voyons que
$$
R \oplus P = K \oplus K' = K \oplus R\pi(1, p) \oplus K''
$$
La projection $\pi' : P \to K''$ est surjective\footnote{En effet,
si $k'' \in K''$, alors $k''$, considéré comme un élément de $K'$,
peut s'écrire $k'' = \lambda \pi(1, 0) + \pi(0, q)$
pour certains $\lambda \in R$ et $q \in P$.
Cela signifie que $k'' = \lambda \pi(1, p) + \pi(0, q - \lambda p)$.
Cela signifie à son tour que $q - \lambda p$ s'envoie sur $k''$ par
la composée $P \to R \oplus P \xrightarrow{\pi} K' \to K''$,
puisque $K' \to K''$ annule $\pi(1, p)$.}
et est donc scindée
(car $K''$ est projectif). Ainsi, $\Ker(\pi') \subset P$ est un
facteur direct contenant $s$. Enfin, par construction, nous
avons un isomorphisme
$$
R \oplus \Ker(\pi') \cong K \oplus R\pi(1, p)
$$
et donc, puisque $K$ est de type fini et stablement libre, $\Ker(\pi')$ l'est aussi.
\end{proof}

\begin{theorem}
\label{theorem-countable-free}
\begin{reference}
Cas commutatif de \cite[Theorem 3.1]{Bass}
\end{reference}
Soit $R$ un anneau de radical de Jacobson $J$ tel que $R/J$ soit noethérien.
Soit $P$ un $R$-module projectif dénombrablement engendré tel que
$P_\mathfrak m$ soit de rang infini pour tout idéal maximal
$\mathfrak m$ de $R$. Alors $P$ est libre.
\end{theorem}

\begin{proof}
Nous montrons d'abord que $P$ est une somme directe dénombrable de modules
stablement libres de type fini. Soit $x_1, x_2, \ldots$ un ensemble dénombrable de générateurs de $P$.
Nous construisons par récurrence des facteurs directs stablement libres de type fini
$F_1, F_2, \ldots$ de $P$ tels que, pour tout $n$,
$F_1 \oplus \ldots \oplus F_n$ soit un facteur direct de $P$
contenant $x_1, \ldots, x_n$. En effet, étant donnés $F_1, \ldots, F_n$
ayant les propriétés voulues, écrivons
$$
P = F_1 \oplus \ldots \oplus F_n \oplus P'
$$
et soit $s \in P'$ l'image de $x_{n + 1}$.
D'après le lemme \ref{lemma-element-in-free-summand},
nous pouvons trouver un facteur direct stablement libre de type fini $F_{n + 1} \subset P'$
contenant $s$. Alors $P = \bigoplus_{i = 1}^{\infty} F_i$.

\medskip\noindent
Supposons que $P$ soit une somme directe infinie
$P = \bigoplus_{i = 1}^{\infty} F_i$ de modules non nuls
stablement libres de type fini. La liberté stable
des modules $F_i$ sera utilisée de la manière suivante :
le rang de chaque $F_i$ est constant (et positif). Nous voyons donc que
$P_\mathfrak m$ est libre
de rang infini dénombrable pour tout idéal maximal
$\mathfrak m$ de $R$. En appliquant le lemme \ref{lemma-trick-to-find-good-element}
avec $s = 0$ et $M = P$, nous pouvons trouver
$t_1 \in P$ tel que $Rt_1$ soit un facteur direct libre de $P$.
Alors $t_1$ est contenu dans $F_1 \oplus \ldots \oplus F_{n_1}$
pour un certain $n_1 > n_0 = 0$.
Le même raisonnement appliqué à $\bigoplus_{n > n_1} F_n$
fournit $n_1 < n_2$ et $t_2 \in F_{n_1 + 1} \oplus \ldots \oplus F_{n_2}$
qui engendre un facteur direct libre.
En poursuivant de cette manière, nous obtenons un facteur direct libre
$$
\bigoplus\nolimits_{i \geq 1} t_i :
\bigoplus\nolimits_{i \geq 1} R
\longrightarrow
\bigoplus\nolimits_{i \geq 1}
\bigoplus\nolimits_{n_i \geq n > n_{i - 1}} F_n = P
$$
de rang infini. Nous voyons donc que $P \cong Q \oplus F$ pour un certain
$R$-module libre $F$ de rang dénombrable. Puisque $Q$ est dénombrablement
engendré, il s'ensuit que $Q \oplus Q' \cong F$ pour
un certain module $Q'$. Alors l'astuce d'Eilenberg
(lemme \ref{lemma-eilenberg-swindle}) implique que
$Q \oplus F \cong F$ et que $P$ est libre.
\end{proof}







\begin{multicols}{2}[\section{Autres chapitres}]
\noindent
Préliminaires
\begin{enumerate}
\item \hyperref[introduction-section-phantom]{Introduction}
\item \hyperref[conventions-section-phantom]{Conventions}
\item \hyperref[sets-section-phantom]{Théorie des ensembles}
\item \hyperref[categories-section-phantom]{Catégories}
\item \hyperref[topology-section-phantom]{Topologie}
\item \hyperref[sheaves-section-phantom]{Faisceaux sur les espaces}
\item \hyperref[sites-section-phantom]{Sites et faisceaux}
\item \hyperref[stacks-section-phantom]{Champs}
\item \hyperref[fields-section-phantom]{Corps}
\item \hyperref[algebra-section-phantom]{Algèbre commutative}
\item \hyperref[brauer-section-phantom]{Groupes de Brauer}
\item \hyperref[homology-section-phantom]{Algèbre homologique}
\item \hyperref[derived-section-phantom]{Catégories dérivées}
\item \hyperref[simplicial-section-phantom]{Méthodes simpliciales}
\item \hyperref[more-algebra-section-phantom]{Compléments d'algèbre}
\item \hyperref[smoothing-section-phantom]{Lissage des morphismes d'anneaux}
\item \hyperref[modules-section-phantom]{Faisceaux de modules}
\item \hyperref[sites-modules-section-phantom]{Modules sur les sites}
\item \hyperref[injectives-section-phantom]{Objets injectifs}
\item \hyperref[cohomology-section-phantom]{Cohomologie des faisceaux}
\item \hyperref[sites-cohomology-section-phantom]{Cohomologie sur les sites}
\item \hyperref[dga-section-phantom]{Algèbre différentielle graduée}
\item \hyperref[dpa-section-phantom]{Algèbre à puissances divisées}
\item \hyperref[sdga-section-phantom]{Faisceaux différentiels gradués}
\item \hyperref[hypercovering-section-phantom]{Hyperrecouvrements}
\end{enumerate}
Schémas
\begin{enumerate}
\setcounter{enumi}{25}
\item \hyperref[schemes-section-phantom]{Schémas}
\item \hyperref[constructions-section-phantom]{Constructions de schémas}
\item \hyperref[properties-section-phantom]{Propriétés des schémas}
\item \hyperref[morphisms-section-phantom]{Morphismes de schémas}
\item \hyperref[coherent-section-phantom]{Cohomologie des schémas}
\item \hyperref[divisors-section-phantom]{Diviseurs}
\item \hyperref[limits-section-phantom]{Limites de schémas}
\item \hyperref[varieties-section-phantom]{Variétés}
\item \hyperref[topologies-section-phantom]{Topologies sur les schémas}
\item \hyperref[descent-section-phantom]{Descente}
\item \hyperref[perfect-section-phantom]{Catégories dérivées de schémas}
\item \hyperref[more-morphisms-section-phantom]{Compléments sur les morphismes}
\item \hyperref[flat-section-phantom]{Compléments sur la platitude}
\item \hyperref[groupoids-section-phantom]{Schémas en groupoïdes}
\item \hyperref[more-groupoids-section-phantom]{Compléments sur les schémas en groupoïdes}
\item \hyperref[etale-section-phantom]{Morphismes étales de schémas}
\end{enumerate}
Thèmes de théorie des schémas
\begin{enumerate}
\setcounter{enumi}{41}
\item \hyperref[chow-section-phantom]{Homologie de Chow}
\item \hyperref[intersection-section-phantom]{Théorie de l'intersection}
\item \hyperref[pic-section-phantom]{Schémas de Picard des courbes}
\item \hyperref[weil-section-phantom]{Théories cohomologiques de Weil}
\item \hyperref[adequate-section-phantom]{Modules adéquats}
\item \hyperref[dualizing-section-phantom]{Complexes dualisants}
\item \hyperref[duality-section-phantom]{Dualité pour les schémas}
\item \hyperref[discriminant-section-phantom]{Discriminants et différentes}
\item \hyperref[derham-section-phantom]{Cohomologie de de Rham}
\item \hyperref[local-cohomology-section-phantom]{Cohomologie locale}
\item \hyperref[algebraization-section-phantom]{Géométrie algébrique et formelle}
\item \hyperref[curves-section-phantom]{Courbes algébriques}
\item \hyperref[resolve-section-phantom]{Résolution des surfaces}
\item \hyperref[models-section-phantom]{Réduction semi-stable}
\item \hyperref[functors-section-phantom]{Foncteurs et morphismes}
\item \hyperref[equiv-section-phantom]{Catégories dérivées de variétés}
\item \hyperref[pione-section-phantom]{Groupes fondamentaux de schémas}
\item \hyperref[etale-cohomology-section-phantom]{Cohomologie étale}
\item \hyperref[crystalline-section-phantom]{Cohomologie cristalline}
\item \hyperref[proetale-section-phantom]{Cohomologie pro-étale}
\item \hyperref[relative-cycles-section-phantom]{Cycles relatifs}
\item \hyperref[more-etale-section-phantom]{Compléments de cohomologie étale}
\item \hyperref[trace-section-phantom]{La formule des traces}
\end{enumerate}
Espaces algébriques
\begin{enumerate}
\setcounter{enumi}{64}
\item \hyperref[spaces-section-phantom]{Espaces algébriques}
\item \hyperref[spaces-properties-section-phantom]{Propriétés des espaces algébriques}
\item \hyperref[spaces-morphisms-section-phantom]{Morphismes d'espaces algébriques}
\item \hyperref[decent-spaces-section-phantom]{Espaces algébriques décents}
\item \hyperref[spaces-cohomology-section-phantom]{Cohomologie des espaces algébriques}
\item \hyperref[spaces-limits-section-phantom]{Limites d'espaces algébriques}
\item \hyperref[spaces-divisors-section-phantom]{Diviseurs sur les espaces algébriques}
\item \hyperref[spaces-over-fields-section-phantom]{Espaces algébriques sur des corps}
\item \hyperref[spaces-topologies-section-phantom]{Topologies sur les espaces algébriques}
\item \hyperref[spaces-descent-section-phantom]{Descente et espaces algébriques}
\item \hyperref[spaces-perfect-section-phantom]{Catégories dérivées d'espaces}
\item \hyperref[spaces-more-morphisms-section-phantom]{Compléments sur les morphismes d'espaces}
\item \hyperref[spaces-flat-section-phantom]{Platitude sur les espaces algébriques}
\item \hyperref[spaces-groupoids-section-phantom]{Groupoïdes dans les espaces algébriques}
\item \hyperref[spaces-more-groupoids-section-phantom]{Compléments sur les groupoïdes dans les espaces}
\item \hyperref[bootstrap-section-phantom]{Amorçage}
\item \hyperref[spaces-pushouts-section-phantom]{Sommes amalgamées d'espaces algébriques}
\end{enumerate}
Thèmes de géométrie
\begin{enumerate}
\setcounter{enumi}{81}
\item \hyperref[spaces-chow-section-phantom]{Groupes de Chow des espaces}
\item \hyperref[groupoids-quotients-section-phantom]{Quotients de groupoïdes}
\item \hyperref[spaces-more-cohomology-section-phantom]{Compléments sur la cohomologie des espaces}
\item \hyperref[spaces-simplicial-section-phantom]{Espaces simpliciaux}
\item \hyperref[spaces-duality-section-phantom]{Dualité pour les espaces}
\item \hyperref[formal-spaces-section-phantom]{Espaces algébriques formels}
\item \hyperref[restricted-section-phantom]{Algébrisation des espaces formels}
\item \hyperref[spaces-resolve-section-phantom]{Résolution des surfaces revisitée}
\end{enumerate}
Théorie des déformations
\begin{enumerate}
\setcounter{enumi}{89}
\item \hyperref[formal-defos-section-phantom]{Théorie formelle des déformations}
\item \hyperref[defos-section-phantom]{Théorie des déformations}
\item \hyperref[cotangent-section-phantom]{Le complexe cotangent}
\item \hyperref[examples-defos-section-phantom]{Problèmes de déformation}
\end{enumerate}
Champs algébriques
\begin{enumerate}
\setcounter{enumi}{93}
\item \hyperref[algebraic-section-phantom]{Champs algébriques}
\item \hyperref[examples-stacks-section-phantom]{Exemples de champs}
\item \hyperref[stacks-sheaves-section-phantom]{Faisceaux sur les champs algébriques}
\item \hyperref[criteria-section-phantom]{Critères de représentabilité}
\item \hyperref[artin-section-phantom]{Axiomes d'Artin}
\item \hyperref[quot-section-phantom]{Espaces de Quot et de Hilbert}
\item \hyperref[stacks-properties-section-phantom]{Propriétés des champs algébriques}
\item \hyperref[stacks-morphisms-section-phantom]{Morphismes de champs algébriques}
\item \hyperref[stacks-limits-section-phantom]{Limites de champs algébriques}
\item \hyperref[stacks-cohomology-section-phantom]{Cohomologie des champs algébriques}
\item \hyperref[stacks-perfect-section-phantom]{Catégories dérivées de champs}
\item \hyperref[stacks-introduction-section-phantom]{Introduction aux champs algébriques}
\item \hyperref[stacks-more-morphisms-section-phantom]{Compléments sur les morphismes de champs}
\item \hyperref[stacks-geometry-section-phantom]{La géométrie des champs}
\end{enumerate}
Thèmes de théorie des espaces de modules
\begin{enumerate}
\setcounter{enumi}{107}
\item \hyperref[moduli-section-phantom]{Champs de modules}
\item \hyperref[moduli-curves-section-phantom]{Espaces de modules de courbes}
\end{enumerate}
Divers
\begin{enumerate}
\setcounter{enumi}{109}
\item \hyperref[examples-section-phantom]{Exemples}
\item \hyperref[exercises-section-phantom]{Exercices}
\item \hyperref[guide-section-phantom]{Guide bibliographique}
\item \hyperref[desirables-section-phantom]{Desiderata}
\item \hyperref[coding-section-phantom]{Style de programmation}
\item \hyperref[obsolete-section-phantom]{Obsolète}
\item \hyperref[fdl-section-phantom]{Licence de documentation libre GNU}
\item \hyperref[index-section-phantom]{Index généré automatiquement}
\end{enumerate}
\end{multicols}


\bibliography{my}
\bibliographystyle{amsalpha}

\end{document}
